% Évaluation pratique — Réseaux locaux (BTS CIEL 1re année) — sujets A et B
% Compilation (4 PDF) :
%   pdflatex -jobname=Eval_Reseau_SujetA   "\def\variante{A}\def\modecorrige{0}% Évaluation pratique — Réseaux locaux (BTS CIEL 1re année) — sujets A et B
% Compilation (4 PDF) :
%   pdflatex -jobname=Eval_Reseau_SujetA   "\def\variante{A}\def\modecorrige{0}% Évaluation pratique — Réseaux locaux (BTS CIEL 1re année) — sujets A et B
% Compilation (4 PDF) :
%   pdflatex -jobname=Eval_Reseau_SujetA   "\def\variante{A}\def\modecorrige{0}% Évaluation pratique — Réseaux locaux (BTS CIEL 1re année) — sujets A et B
% Compilation (4 PDF) :
%   pdflatex -jobname=Eval_Reseau_SujetA   "\def\variante{A}\def\modecorrige{0}\input{Eval_TP_Reseau_LAN.tex}"
%   pdflatex -jobname=Eval_Reseau_SujetB   "\def\variante{B}\def\modecorrige{0}\input{Eval_TP_Reseau_LAN.tex}"
%   pdflatex -jobname=Eval_Reseau_CorrigeA "\def\variante{A}\def\modecorrige{1}\input{Eval_TP_Reseau_LAN.tex}"
%   pdflatex -jobname=Eval_Reseau_CorrigeB "\def\variante{B}\def\modecorrige{1}\input{Eval_TP_Reseau_LAN.tex}"
\documentclass[10pt,a4paper]{article}

\usepackage[T1]{fontenc}
\usepackage[utf8]{inputenc}
\usepackage{lmodern}
\usepackage[french]{babel}
\usepackage[a4paper,hmargin=1.5cm,top=1.3cm,bottom=1.6cm,footskip=0.7cm]{geometry}
\usepackage[table]{xcolor}
\usepackage{booktabs,tabularx,array}
\usepackage{enumitem}
\usepackage{fancyhdr}
\usepackage[most]{tcolorbox}
\usepackage{tikz}
\usetikzlibrary{positioning}
\usepackage{titlesec}

\providecommand{\modecorrige}{0}
\providecommand{\variante}{A}
\newif\ifcorrige
\ifnum\modecorrige=1 \corrigetrue \else \corrigefalse \fi

\definecolor{ipred}{HTML}{C0392B}
\definecolor{cdcbar}{HTML}{A67C3D}
\definecolor{cdcbg}{HTML}{F4F1EA}
\definecolor{qbar}{HTML}{4A6B5D}
\definecolor{qbg}{HTML}{EDF1EF}
\definecolor{corr}{HTML}{1F5FA8}
\definecolor{corrbg}{HTML}{E8F0FA}
\definecolor{grisfoot}{HTML}{828282}

\newcommand{\ip}[1]{\textcolor{ipred}{\texttt{#1}}}
\newcommand{\pts}[1]{\hfill\mbox{\textbf{(#1)}}}
\newcommand{\sol}[1]{\ifcorrige\textcolor{corr}{#1}\fi}
\newcommand{\bareme}[1]{\ifcorrige\par{\footnotesize\color{corr}#1\par}\fi}
\newcommand{\hl}{\ifcorrige\rule[-1.6mm]{0pt}{5.4mm}\else\rule[-2.2mm]{0pt}{6.6mm}\fi}
\newcolumntype{C}{>{\centering\arraybackslash}X}
\newcolumntype{L}{>{\raggedright\arraybackslash}X}

\titleformat{\section}{\large\bfseries}{}{0pt}{}
\titlespacing*{\section}{0pt}{6pt}{3pt}

\tcbset{qstyle/.style={enhanced,boxrule=0pt,leftrule=2.5pt,arc=0pt,outer arc=0pt,
  left=5pt,right=5pt,top=2pt,bottom=2pt,before skip=4pt,after skip=3pt}}
\newtcolorbox{cdc}{qstyle,colback=cdcbg,colframe=cdcbar}
\newtcolorbox{quest}{qstyle,colback=qbg,colframe=qbar}

% Zone de réponse : vierge dans le sujet, remplie dans le corrigé
\newcommand{\zone}[2]{%
  \ifcorrige
    \begin{tcolorbox}[colback=corrbg,colframe=corr,boxrule=0.5pt,arc=1pt,left=4pt,right=4pt,top=1pt,bottom=1pt,
      before skip=2pt,after skip=2pt,fontupper=\footnotesize\color{corr}]#2\end{tcolorbox}%
  \else
    \begin{tcolorbox}[colback=white,colframe=black!45,boxrule=0.4pt,arc=1pt,height=#1,
      before skip=2pt,after skip=2pt]\end{tcolorbox}%
  \fi}

\newcommand{\question}[3]{\begin{quest}\textbf{#1.} #3\pts{#2}\end{quest}}

%==========================================================================
% Données propres à chaque variante
%==========================================================================
\if\variante A
  \def\entreprise{Le cabinet d'architectes \textbf{ArchiNova}}
  \def\Rnom{R-ArchiNova}
  \def\resUn{Bureau d'études}\def\pUn{BE}\def\netUn{192.168.40}
  \def\resDeux{Accueil}\def\pDeux{ACC}\def\netDeux{172.16}
  \def\srvNom{SRV-Intranet}\def\netSrv{192.168.200}\def\srvH{10}
  \def\ssid{ArchiNova-WiFi}
  \def\AIrows{%
    \hl\ip{192.168.40.25} & \sol{C} & \sol{255.255.255.0} & \sol{192.168.40.0} & \sol{192.168.40.255} & \sol{254} \\ \hline
    \hl\ip{172.16.5.10}   & \sol{B} & \sol{255.255.0.0}   & \sol{172.16.0.0}   & \sol{172.16.255.255} & \sol{65\,534} \\ \hline
    \hl\ip{10.3.2.1}      & \sol{A} & \sol{255.0.0.0}     & \sol{10.0.0.0}     & \sol{10.255.255.255} & \sol{16\,777\,214} \\ \hline}
  \def\apipa{169.254.12.7}
  \def\Cnet{192.168.60}
  \def\Crows{%
    PC-A & \ip{192.168.60.10} & \ip{255.255.255.0} & \ip{192.168.60.254} & — \\
    PC-B & \ip{192.168.60.20} & \ip{255.255.255.0} & \ip{192.168.60.1}   & — \\
    PC-C & \ip{192.168.61.30} & \ip{255.255.255.0} & \ip{192.168.60.254} & — \\
    PC-D & \ip{192.168.60.10} & \ip{255.255.255.0} & \ip{192.168.60.254} & — \\
    SRV-C & \ip{10.0.0.5}     & \ip{255.0.0.0}     & \ip{10.0.0.254}     & — \\
    R1 — G0/0 & \ip{192.168.60.254} & \ip{255.255.255.0} & — & On \\
    R1 — G0/1 & \ip{10.0.0.254}     & \ip{255.0.0.0}     & — & Off \\}
  \def\Csol{%
    \hl 1 & \sol{PC-B} & \sol{Passerelle .60.1 : pas l'adresse du routeur.} & \sol{Joint son réseau, mais aucun autre (serveur).} & \sol{Passerelle 192.168.60.254.} \\ \hline
    \hl 2 & \sol{PC-C} & \sol{IP hors du réseau 192.168.60.0/24.} & \sol{Ne joint ni les postes ni le routeur.} & \sol{Ex. 192.168.60.30.} \\ \hline
    \hl 3 & \sol{PC-D} & \sol{Même IP que PC-A (doublon).} & \sol{Conflit d'adresse : PC-A et PC-D perturbés.} & \sol{Adresse unique, ex. .60.40.} \\ \hline
    \hl 4 & \sol{R1 G0/1} & \sol{Interface désactivée (Off).} & \sol{Réseau du serveur injoignable pour tous.} & \sol{Port Status : On.} \\ \hline}
  \def\Cbareme{PC-A et SRV-C sont correctement configurés ; accepter PC-A cité à la place de PC-D pour le doublon.}
\else
  \def\entreprise{La clinique vétérinaire \textbf{VetoPlus}}
  \def\Rnom{R-VetoPlus}
  \def\resUn{Soins}\def\pUn{SOI}\def\netUn{192.168.70}
  \def\resDeux{Secrétariat}\def\pDeux{SEC}\def\netDeux{172.20}
  \def\srvNom{SRV-Dossiers}\def\netSrv{192.168.250}\def\srvH{20}
  \def\ssid{VetoPlus-WiFi}
  \def\AIrows{%
    \hl\ip{192.168.15.130} & \sol{C} & \sol{255.255.255.0} & \sol{192.168.15.0} & \sol{192.168.15.255} & \sol{254} \\ \hline
    \hl\ip{172.31.200.7}   & \sol{B} & \sol{255.255.0.0}   & \sol{172.31.0.0}   & \sol{172.31.255.255} & \sol{65\,534} \\ \hline
    \hl\ip{10.45.6.9}      & \sol{A} & \sol{255.0.0.0}     & \sol{10.0.0.0}     & \sol{10.255.255.255} & \sol{16\,777\,214} \\ \hline}
  \def\apipa{169.254.88.3}
  \def\Cnet{192.168.80}
  \def\Crows{%
    PC-A & \ip{192.168.8.11}   & \ip{255.255.255.0} & \ip{192.168.80.254} & — \\
    PC-B & \ip{192.168.80.12}  & \ip{255.255.255.0} & \ip{192.168.80.254} & — \\
    PC-C & \ip{192.168.80.13}  & \ip{255.255.255.0} & \ip{192.168.80.100} & — \\
    PC-D & \ip{192.168.80.12}  & \ip{255.255.255.0} & \ip{192.168.80.254} & — \\
    SRV-C & \ip{10.0.0.8}      & \ip{255.0.0.0}     & \ip{10.0.0.254}     & — \\
    R1 — G0/0 & \ip{192.168.80.254} & \ip{255.255.255.0} & — & Off \\
    R1 — G0/1 & \ip{10.0.0.254}     & \ip{255.0.0.0}     & — & On \\}
  \def\Csol{%
    \hl 1 & \sol{PC-A} & \sol{IP hors du réseau 192.168.80.0/24.} & \sol{Ne joint ni les postes ni le routeur.} & \sol{Ex. 192.168.80.11.} \\ \hline
    \hl 2 & \sol{PC-C} & \sol{Passerelle .80.100 : pas l'adresse du routeur.} & \sol{Joint son réseau, mais aucun autre (serveur).} & \sol{Passerelle 192.168.80.254.} \\ \hline
    \hl 3 & \sol{PC-D} & \sol{Même IP que PC-B (doublon).} & \sol{Conflit d'adresse : PC-B et PC-D perturbés.} & \sol{Adresse unique, ex. .80.14.} \\ \hline
    \hl 4 & \sol{R1 G0/0} & \sol{Interface côté LAN désactivée (Off).} & \sol{Aucun poste ne joint le routeur ni le serveur.} & \sol{Port Status : On.} \\ \hline}
  \def\Cbareme{PC-B et SRV-C sont correctement configurés ; accepter PC-B cité à la place de PC-D pour le doublon.}
\fi

\pagestyle{fancy}
\fancyhf{}
\renewcommand{\headrulewidth}{0pt}
\fancyfoot[L]{\footnotesize\itshape\color{grisfoot}Évaluation — Réseaux locaux — Sujet \variante\ifcorrige\ — CORRIGÉ\fi}
\fancyfoot[R]{\footnotesize\color{grisfoot}\thepage/3}

\setlist{itemsep=0pt,topsep=2pt,parsep=0pt}
\setlength{\parindent}{0pt}
\setlength{\parskip}{2pt}

\begin{document}

\begin{center}
  {\Large\bfseries Évaluation pratique — Réseaux locaux \quad
   \fcolorbox{black}{white}{\ SUJET \variante\ }\par}
  \vspace{1pt}
  {\bfseries Adressage IP, commutation et routage — BTS CIEL 1re année — durée : 2 heures\par}
  \ifcorrige\vspace{2pt}{\bfseries\color{corr}CORRIGÉ ET BARÈME — document enseignant}\par\fi
\end{center}
\vspace{-2pt}
{\renewcommand{\arraystretch}{1.35}
\begin{tabularx}{\linewidth}{|X|X|p{2.2cm}|p{2.6cm}|}
  \hline
  \textbf{Nom :} & \textbf{Prénom :} & \textbf{Poste n°} & \textbf{Note :} \hfill /20 \\ \hline
\end{tabularx}}

\textbf{Consignes :} travail \textbf{individuel}, documents et Internet \textbf{non autorisés}. Ce sujet est le
\textbf{document-réponse}. Fichier Packet Tracer : \ip{Nom\_Prenom\_EvalLAN\_\variante.pkt}, à enregistrer
régulièrement dans le dossier de rendu. Les « visas » sont apposés par le professeur, que vous appelez à chaque test réussi.
Durées conseillées : A 30 min — B 1 h 10 — C 20 min.

%==========================================================================
\section{Partie A — Questions de cours et calculs \hfill /6}

\question{A1}{2 pts}{Complétez le tableau en considérant le \textbf{masque par défaut} de la classe de chaque adresse.}
{\small\renewcommand{\arraystretch}{1.15}
\begin{tabularx}{\linewidth}{|l|c|C|C|C|c|}
  \hline
  \textbf{Adresse IP} & \textbf{Classe} & \textbf{Masque par défaut} & \textbf{Adresse réseau} &
  \textbf{Adresse de diffusion} & \textbf{Nb d'hôtes} \\ \hline
  \AIrows
\end{tabularx}}
\bareme{0,5 pt par ligne juste + 0,5 pt si les trois nombres d'hôtes sont justes ($2^{n}-2$, $n$ = nombre de bits d'hôte).}

\question{A2}{1 pt}{Indiquez le type de câble (\textbf{droit} ou \textbf{croisé}) pour chaque liaison.}
{\small\renewcommand{\arraystretch}{1.15}
\begin{tabularx}{\linewidth}{|X|c||X|c|}
  \hline
  \hl PC $\leftrightarrow$ commutateur & \makebox[1.6cm]{\sol{droit}} & Commutateur $\leftrightarrow$ routeur & \makebox[1.6cm]{\sol{droit}} \\ \hline
  \hl PC $\leftrightarrow$ PC (réseau P2P) & \sol{croisé} & Concentrateur $\leftrightarrow$ concentrateur & \sol{croisé} \\ \hline
\end{tabularx}}
\bareme{0,25 pt par case. Équipements de types différents : câble droit ; de même type : câble croisé.}

\question{A3}{1 pt}{Comment un \textbf{concentrateur (hub)} et un \textbf{commutateur (switch)} traitent-ils une trame
reçue ? Lequel limite les collisions, et pourquoi ?}
\zone{2.9cm}{Le \textbf{hub} (couche 1) répète la trame sur \textbf{tous ses autres ports} : un seul domaine de collision
partagé. Le \textbf{switch} (couche 2) lit l'adresse MAC de destination et, grâce à sa \textbf{table MAC}, ne l'envoie
\textbf{que sur le port du destinataire} : un domaine de collision par port, c'est lui qui limite les collisions.
\emph{0,5 pt hub, 0,5 pt switch.}}

\question{A4}{1 pt}{Qu'est-ce que la \textbf{passerelle par défaut} d'un poste ? Quelle condition son adresse doit-elle respecter ?}
\zone{2.4cm}{Adresse IP de l'interface du \textbf{routeur} à laquelle le poste envoie les paquets destinés à un
\textbf{autre réseau}. Elle doit \textbf{appartenir au même réseau} que le poste. \emph{0,5 + 0,5 pt.}}

\question{A5}{0,5 pt}{Un portable configuré en DHCP affiche \ip{\apipa} / \ip{255.255.0.0}. Que s'est-il passé ?}
\zone{1.5cm}{Aucun \textbf{serveur DHCP} n'a répondu : le poste s'est attribué une adresse \textbf{APIPA} (169.254.0.0/16).}

\question{A6}{0,5 pt}{Quel protocole utilise \texttt{ping} ? Nommez les deux messages échangés et la couche OSI des adresses IP.}
\zone{1.5cm}{\textbf{ICMP} ; \textbf{Echo Request} / \textbf{Echo Reply} ; adresses IP en \textbf{couche 3 (Réseau)}.}

%==========================================================================
\newpage
\section{Partie B — Réalisation sous Packet Tracer \hfill /10}

\begin{cdc}
\entreprise{} vous confie son réseau : \textbf{trois réseaux} reliés par un \textbf{routeur unique Cisco 2911}
nommé \ip{\Rnom} (interfaces G0/0, G0/1, G0/2).
\begin{itemize}[leftmargin=12pt]
  \item \textbf{Réseau \resUn} : \textbf{4 postes} PC-\pUn1 à PC-\pUn4, commutateur \ip{SW-\pUn}, adressage statique,
        réseau \ip{\netUn.0/24}.
  \item \textbf{Réseau \resDeux} : \textbf{2 postes} PC-\pDeux1, PC-\pDeux2 et \textbf{1 imprimante} IMP-\pDeux,
        commutateur \ip{SW-\pDeux}, adressage statique, réseau de classe B \ip{\netDeux.0.0/16}.
  \item \textbf{Réseau Serveur} : serveur \ip{\srvNom} d'adresse \ip{\netSrv.\srvH/24}, relié \textbf{directement} au routeur.
  \item La passerelle de chaque réseau est sa \textbf{dernière adresse utilisable}. Tous les équipements communiquent
        entre eux et accèdent à la page web du serveur. Nommage imposé (« Display Name »).
\end{itemize}
\end{cdc}

\question{B1}{2 pts}{Complétez le plan d'adressage \textbf{avant} la réalisation.}
{\footnotesize\setlength{\tabcolsep}{3pt}\renewcommand{\arraystretch}{1.15}
\begin{tabularx}{\linewidth}{|l|C|C|C|C|C|C|}
  \hline
  \textbf{Réseau} & \textbf{Adresse réseau} & \textbf{Masque} & \textbf{1re @ utilisable} &
  \textbf{Dernière @ utilisable} & \textbf{Diffusion} & \textbf{Passerelle} \\ \hline
  \hl \resUn & \sol{\netUn.0} & \sol{255.255.255.0} & \sol{\netUn.1} & \sol{\netUn.254} & \sol{\netUn.255} & \sol{\netUn.254} \\ \hline
  \hl \resDeux & \sol{\netDeux.0.0} & \sol{255.255.0.0} & \sol{\netDeux.0.1} & \sol{\netDeux.255.254} & \sol{\netDeux.255.255} & \sol{\netDeux.255.254} \\ \hline
  \hl Serveur & \sol{\netSrv.0} & \sol{255.255.255.0} & \sol{\netSrv.1} & \sol{\netSrv.254} & \sol{\netSrv.255} & \sol{\netSrv.254} \\ \hline
\end{tabularx}}
\bareme{0,5 pt par ligne juste + 0,5 pt si la convention de passerelle est respectée partout.}

\question{B2}{2 pts}{Construisez la topologie (équipements, câblage, nommage) et relevez l'adressage configuré.}
{\footnotesize\setlength{\tabcolsep}{3pt}\renewcommand{\arraystretch}{1.1}
\begin{tabularx}{\linewidth}{|l|C|C|C||l|C|C|}
  \hline
  \textbf{Hôte} & \textbf{Adresse IP} & \textbf{Masque} & \textbf{Passerelle} & \textbf{Interface} & \textbf{Adresse IP} & \textbf{Masque} \\ \hline
  \hl PC-\pUn1 & \sol{\netUn.1} & \sol{255.255.255.0} & \sol{\netUn.254} & \Rnom{} G0/0 & \sol{\netUn.254} & \sol{255.255.255.0} \\ \hline
  \hl PC-\pDeux1 & \sol{\netDeux.0.1} & \sol{255.255.0.0} & \sol{\netDeux.255.254} & \Rnom{} G0/1 & \sol{\netDeux.255.254} & \sol{255.255.0.0} \\ \hline
  \hl IMP-\pDeux & \sol{\netDeux.0.10} & \sol{255.255.0.0} & \sol{\netDeux.255.254} & \Rnom{} G0/2 & \sol{\netSrv.254} & \sol{255.255.255.0} \\ \hline
  \hl \srvNom & \ip{\netSrv.\srvH} & \sol{255.255.255.0} & \sol{\netSrv.254} & \multicolumn{3}{c|}{\cellcolor{black!8}} \\ \hline
\end{tabularx}}
\bareme{Adresses d'hôtes libres (uniques, dans le bon réseau) ; affectation G0/x libre. Sur le .pkt : équipements 0,5 ;
câblage 1 (câbles droits vers les switchs et switch–routeur, \textbf{croisé} serveur–routeur, liens verts) ; nommage 0,5.}

\question{B3}{3 pts}{Configurez tous les équipements : IP, masques, passerelles des hôtes ; interfaces du routeur
(« Port Status : On »).}
\bareme{Sur le .pkt : interfaces routeur 1 pt ; IP/masques des hôtes 1 pt ; passerelles (imprimante et serveur compris)
1 pt. $-0,25$ par erreur, sans descendre sous 0 par item.}

\question{B4}{3 pts}{Réalisez les tests, notez la commande et le résultat, puis faites viser chaque test réussi.}
{\small\renewcommand{\arraystretch}{1.1}
\begin{tabularx}{\linewidth}{|c|l|l|C|c|c|}
  \hline
  \textbf{N°} & \textbf{Source} & \textbf{Destination} & \textbf{Commande / action} & \textbf{Résultat} & \textbf{Visa prof.} \\ \hline
  \hl 1 & PC-\pUn1 & PC-\pUn4 & \sol{ping \netUn.4} & \sol{réussi} & \hspace{1.4cm} \\ \hline
  \hl 2 & PC-\pDeux1 & IMP-\pDeux & \sol{ping \netDeux.0.10} & \sol{réussi} & \\ \hline
  \hl 3 & PC-\pUn2 & PC-\pDeux2 & \sol{ping \netDeux.0.2} & \sol{réussi} & \\ \hline
  \hl 4 & PC-\pDeux2 & \srvNom & \sol{ping \netSrv.\srvH} & \sol{réussi} & \\ \hline
  \hl 5 & PC-\pUn3 & page web du serveur & \sol{navigateur : {\NoAutoSpacing http://\netSrv.\srvH}} & \sol{page affichée} & \\ \hline
\end{tabularx}}
\bareme{0,5 pt par test visé + 0,5 pt pour B5. Le premier ping inter-réseaux peut échouer (résolution ARP) ; les suivants doivent réussir.}

\question{B5}{incl. B4}{Test n°3 : par quels équipements passe le paquet ? Pourquoi PC-\pUn2 l'envoie-t-il au routeur ?}
\zone{2.3cm}{PC-\pUn2 $\to$ SW-\pUn{} $\to$ \Rnom{} $\to$ SW-\pDeux{} $\to$ PC-\pDeux2. Grâce à son masque, PC-\pUn2 constate que
la destination (\netDeux.0.0/16) n'est \textbf{pas dans son réseau} (\netUn.0/24) : il envoie le paquet à sa
\textbf{passerelle par défaut}, qui le route.}

%==========================================================================
\newpage
\section{Partie C — Diagnostic d'un réseau en panne \hfill /4}

Un technicien constate qu'\textbf{aucun poste ne parvient à joindre le serveur SRV-C} et que certains postes
communiquent mal entre eux. Les câbles sont du bon type.

\begin{center}
\begin{minipage}{0.37\linewidth}\centering
\begin{tikzpicture}[font=\footnotesize,
  eq/.style={draw,fill=qbg,rounded corners=2pt,minimum width=9mm,minimum height=5mm,inner sep=2pt},
  every path/.style={thick}]
  \node[eq] (SW) {SW-1};
  \node[eq,right=9mm of SW] (R) {R1};
  \node[eq,right=7mm of R] (SRV) {SRV-C};
  \node[eq,above left=5mm and -3mm of SW] (A) {PC-A};
  \node[eq,above right=5mm and -3mm of SW] (B) {PC-B};
  \node[eq,below left=5mm and -3mm of SW] (C) {PC-C};
  \node[eq,below right=5mm and -3mm of SW] (D) {PC-D};
  \foreach \n in {A,B,C,D} \draw (\n) -- (SW);
  \draw (SW) -- node[above,font=\tiny]{G0/0} (R);
  \draw (R) -- node[above,font=\tiny]{G0/1} (SRV);
\end{tikzpicture}
\end{minipage}\hfill
\begin{minipage}{0.61\linewidth}\centering
{\small\renewcommand{\arraystretch}{1.05}\setlength{\tabcolsep}{4pt}
\begin{tabular}{lllll}
  \toprule
  \textbf{Équipement} & \textbf{Adresse IP} & \textbf{Masque} & \textbf{Passerelle} & \textbf{Port} \\
  \midrule
  \Crows
  \bottomrule
\end{tabular}}
\end{minipage}
\end{center}

\question{C1}{4 pts}{Relevez les \textbf{quatre anomalies}. Pour chacune : conséquence et correction.}
{\small\renewcommand{\arraystretch}{1.1}
\ifcorrige\def\hc{}\else\def\hc{\rule[-10mm]{0pt}{14mm}}\fi
\begin{tabularx}{\linewidth}{|c|>{\raggedright\arraybackslash}p{2cm}|L|L|L|}
  \hline
  \textbf{N°} & \textbf{Équipement} & \textbf{Anomalie} & \textbf{Conséquence} & \textbf{Correction} \\ \hline
  \ifcorrige\Csol\else
  \hc 1 & & & & \\ \hline
  \hc 2 & & & & \\ \hline
  \hc 3 & & & & \\ \hline
  \hc 4 & & & & \\ \hline
  \fi
\end{tabularx}}
\bareme{1 pt par anomalie (0,25 équipement, anomalie, conséquence, correction ; ordre indifférent). \Cbareme}

%==========================================================================
\section{Bonus — Extension sans fil \hfill +1}

\begin{cdc}Ajoutez au réseau \resUn{} un \textbf{AccessPoint-PT} relié à SW-\pUn{} (SSID \ip{\ssid}, \textbf{WPA2-PSK})
et un portable \ip{Laptop-\pUn} équipé d'un module \texttt{WPC300N}, en adressage \textbf{statique}.
Vérifiez qu'il joint \srvNom.\end{cdc}
{\small\renewcommand{\arraystretch}{1.1}
\begin{tabularx}{\linewidth}{|l|C|C|C|c|}
  \hline
  \textbf{Équipement} & \textbf{Adresse IP} & \textbf{Masque} & \textbf{Passerelle} & \textbf{Visa prof.} \\ \hline
  \hl Laptop-\pUn & \sol{\netUn.50} & \sol{255.255.255.0} & \sol{\netUn.254} & \hspace{1.4cm} \\ \hline
\end{tabularx}}
\bareme{Le point d'accès fait office de pont : le portable est dans \netUn.0/24. Bonus si associé en WPA2-PSK et ping
vers \netSrv.\srvH{} réussi. Note plafonnée à 20.}

\end{document}
"
%   pdflatex -jobname=Eval_Reseau_SujetB   "\def\variante{B}\def\modecorrige{0}% Évaluation pratique — Réseaux locaux (BTS CIEL 1re année) — sujets A et B
% Compilation (4 PDF) :
%   pdflatex -jobname=Eval_Reseau_SujetA   "\def\variante{A}\def\modecorrige{0}\input{Eval_TP_Reseau_LAN.tex}"
%   pdflatex -jobname=Eval_Reseau_SujetB   "\def\variante{B}\def\modecorrige{0}\input{Eval_TP_Reseau_LAN.tex}"
%   pdflatex -jobname=Eval_Reseau_CorrigeA "\def\variante{A}\def\modecorrige{1}\input{Eval_TP_Reseau_LAN.tex}"
%   pdflatex -jobname=Eval_Reseau_CorrigeB "\def\variante{B}\def\modecorrige{1}\input{Eval_TP_Reseau_LAN.tex}"
\documentclass[10pt,a4paper]{article}

\usepackage[T1]{fontenc}
\usepackage[utf8]{inputenc}
\usepackage{lmodern}
\usepackage[french]{babel}
\usepackage[a4paper,hmargin=1.5cm,top=1.3cm,bottom=1.6cm,footskip=0.7cm]{geometry}
\usepackage[table]{xcolor}
\usepackage{booktabs,tabularx,array}
\usepackage{enumitem}
\usepackage{fancyhdr}
\usepackage[most]{tcolorbox}
\usepackage{tikz}
\usetikzlibrary{positioning}
\usepackage{titlesec}

\providecommand{\modecorrige}{0}
\providecommand{\variante}{A}
\newif\ifcorrige
\ifnum\modecorrige=1 \corrigetrue \else \corrigefalse \fi

\definecolor{ipred}{HTML}{C0392B}
\definecolor{cdcbar}{HTML}{A67C3D}
\definecolor{cdcbg}{HTML}{F4F1EA}
\definecolor{qbar}{HTML}{4A6B5D}
\definecolor{qbg}{HTML}{EDF1EF}
\definecolor{corr}{HTML}{1F5FA8}
\definecolor{corrbg}{HTML}{E8F0FA}
\definecolor{grisfoot}{HTML}{828282}

\newcommand{\ip}[1]{\textcolor{ipred}{\texttt{#1}}}
\newcommand{\pts}[1]{\hfill\mbox{\textbf{(#1)}}}
\newcommand{\sol}[1]{\ifcorrige\textcolor{corr}{#1}\fi}
\newcommand{\bareme}[1]{\ifcorrige\par{\footnotesize\color{corr}#1\par}\fi}
\newcommand{\hl}{\ifcorrige\rule[-1.6mm]{0pt}{5.4mm}\else\rule[-2.2mm]{0pt}{6.6mm}\fi}
\newcolumntype{C}{>{\centering\arraybackslash}X}
\newcolumntype{L}{>{\raggedright\arraybackslash}X}

\titleformat{\section}{\large\bfseries}{}{0pt}{}
\titlespacing*{\section}{0pt}{6pt}{3pt}

\tcbset{qstyle/.style={enhanced,boxrule=0pt,leftrule=2.5pt,arc=0pt,outer arc=0pt,
  left=5pt,right=5pt,top=2pt,bottom=2pt,before skip=4pt,after skip=3pt}}
\newtcolorbox{cdc}{qstyle,colback=cdcbg,colframe=cdcbar}
\newtcolorbox{quest}{qstyle,colback=qbg,colframe=qbar}

% Zone de réponse : vierge dans le sujet, remplie dans le corrigé
\newcommand{\zone}[2]{%
  \ifcorrige
    \begin{tcolorbox}[colback=corrbg,colframe=corr,boxrule=0.5pt,arc=1pt,left=4pt,right=4pt,top=1pt,bottom=1pt,
      before skip=2pt,after skip=2pt,fontupper=\footnotesize\color{corr}]#2\end{tcolorbox}%
  \else
    \begin{tcolorbox}[colback=white,colframe=black!45,boxrule=0.4pt,arc=1pt,height=#1,
      before skip=2pt,after skip=2pt]\end{tcolorbox}%
  \fi}

\newcommand{\question}[3]{\begin{quest}\textbf{#1.} #3\pts{#2}\end{quest}}

%==========================================================================
% Données propres à chaque variante
%==========================================================================
\if\variante A
  \def\entreprise{Le cabinet d'architectes \textbf{ArchiNova}}
  \def\Rnom{R-ArchiNova}
  \def\resUn{Bureau d'études}\def\pUn{BE}\def\netUn{192.168.40}
  \def\resDeux{Accueil}\def\pDeux{ACC}\def\netDeux{172.16}
  \def\srvNom{SRV-Intranet}\def\netSrv{192.168.200}\def\srvH{10}
  \def\ssid{ArchiNova-WiFi}
  \def\AIrows{%
    \hl\ip{192.168.40.25} & \sol{C} & \sol{255.255.255.0} & \sol{192.168.40.0} & \sol{192.168.40.255} & \sol{254} \\ \hline
    \hl\ip{172.16.5.10}   & \sol{B} & \sol{255.255.0.0}   & \sol{172.16.0.0}   & \sol{172.16.255.255} & \sol{65\,534} \\ \hline
    \hl\ip{10.3.2.1}      & \sol{A} & \sol{255.0.0.0}     & \sol{10.0.0.0}     & \sol{10.255.255.255} & \sol{16\,777\,214} \\ \hline}
  \def\apipa{169.254.12.7}
  \def\Cnet{192.168.60}
  \def\Crows{%
    PC-A & \ip{192.168.60.10} & \ip{255.255.255.0} & \ip{192.168.60.254} & — \\
    PC-B & \ip{192.168.60.20} & \ip{255.255.255.0} & \ip{192.168.60.1}   & — \\
    PC-C & \ip{192.168.61.30} & \ip{255.255.255.0} & \ip{192.168.60.254} & — \\
    PC-D & \ip{192.168.60.10} & \ip{255.255.255.0} & \ip{192.168.60.254} & — \\
    SRV-C & \ip{10.0.0.5}     & \ip{255.0.0.0}     & \ip{10.0.0.254}     & — \\
    R1 — G0/0 & \ip{192.168.60.254} & \ip{255.255.255.0} & — & On \\
    R1 — G0/1 & \ip{10.0.0.254}     & \ip{255.0.0.0}     & — & Off \\}
  \def\Csol{%
    \hl 1 & \sol{PC-B} & \sol{Passerelle .60.1 : pas l'adresse du routeur.} & \sol{Joint son réseau, mais aucun autre (serveur).} & \sol{Passerelle 192.168.60.254.} \\ \hline
    \hl 2 & \sol{PC-C} & \sol{IP hors du réseau 192.168.60.0/24.} & \sol{Ne joint ni les postes ni le routeur.} & \sol{Ex. 192.168.60.30.} \\ \hline
    \hl 3 & \sol{PC-D} & \sol{Même IP que PC-A (doublon).} & \sol{Conflit d'adresse : PC-A et PC-D perturbés.} & \sol{Adresse unique, ex. .60.40.} \\ \hline
    \hl 4 & \sol{R1 G0/1} & \sol{Interface désactivée (Off).} & \sol{Réseau du serveur injoignable pour tous.} & \sol{Port Status : On.} \\ \hline}
  \def\Cbareme{PC-A et SRV-C sont correctement configurés ; accepter PC-A cité à la place de PC-D pour le doublon.}
\else
  \def\entreprise{La clinique vétérinaire \textbf{VetoPlus}}
  \def\Rnom{R-VetoPlus}
  \def\resUn{Soins}\def\pUn{SOI}\def\netUn{192.168.70}
  \def\resDeux{Secrétariat}\def\pDeux{SEC}\def\netDeux{172.20}
  \def\srvNom{SRV-Dossiers}\def\netSrv{192.168.250}\def\srvH{20}
  \def\ssid{VetoPlus-WiFi}
  \def\AIrows{%
    \hl\ip{192.168.15.130} & \sol{C} & \sol{255.255.255.0} & \sol{192.168.15.0} & \sol{192.168.15.255} & \sol{254} \\ \hline
    \hl\ip{172.31.200.7}   & \sol{B} & \sol{255.255.0.0}   & \sol{172.31.0.0}   & \sol{172.31.255.255} & \sol{65\,534} \\ \hline
    \hl\ip{10.45.6.9}      & \sol{A} & \sol{255.0.0.0}     & \sol{10.0.0.0}     & \sol{10.255.255.255} & \sol{16\,777\,214} \\ \hline}
  \def\apipa{169.254.88.3}
  \def\Cnet{192.168.80}
  \def\Crows{%
    PC-A & \ip{192.168.8.11}   & \ip{255.255.255.0} & \ip{192.168.80.254} & — \\
    PC-B & \ip{192.168.80.12}  & \ip{255.255.255.0} & \ip{192.168.80.254} & — \\
    PC-C & \ip{192.168.80.13}  & \ip{255.255.255.0} & \ip{192.168.80.100} & — \\
    PC-D & \ip{192.168.80.12}  & \ip{255.255.255.0} & \ip{192.168.80.254} & — \\
    SRV-C & \ip{10.0.0.8}      & \ip{255.0.0.0}     & \ip{10.0.0.254}     & — \\
    R1 — G0/0 & \ip{192.168.80.254} & \ip{255.255.255.0} & — & Off \\
    R1 — G0/1 & \ip{10.0.0.254}     & \ip{255.0.0.0}     & — & On \\}
  \def\Csol{%
    \hl 1 & \sol{PC-A} & \sol{IP hors du réseau 192.168.80.0/24.} & \sol{Ne joint ni les postes ni le routeur.} & \sol{Ex. 192.168.80.11.} \\ \hline
    \hl 2 & \sol{PC-C} & \sol{Passerelle .80.100 : pas l'adresse du routeur.} & \sol{Joint son réseau, mais aucun autre (serveur).} & \sol{Passerelle 192.168.80.254.} \\ \hline
    \hl 3 & \sol{PC-D} & \sol{Même IP que PC-B (doublon).} & \sol{Conflit d'adresse : PC-B et PC-D perturbés.} & \sol{Adresse unique, ex. .80.14.} \\ \hline
    \hl 4 & \sol{R1 G0/0} & \sol{Interface côté LAN désactivée (Off).} & \sol{Aucun poste ne joint le routeur ni le serveur.} & \sol{Port Status : On.} \\ \hline}
  \def\Cbareme{PC-B et SRV-C sont correctement configurés ; accepter PC-B cité à la place de PC-D pour le doublon.}
\fi

\pagestyle{fancy}
\fancyhf{}
\renewcommand{\headrulewidth}{0pt}
\fancyfoot[L]{\footnotesize\itshape\color{grisfoot}Évaluation — Réseaux locaux — Sujet \variante\ifcorrige\ — CORRIGÉ\fi}
\fancyfoot[R]{\footnotesize\color{grisfoot}\thepage/3}

\setlist{itemsep=0pt,topsep=2pt,parsep=0pt}
\setlength{\parindent}{0pt}
\setlength{\parskip}{2pt}

\begin{document}

\begin{center}
  {\Large\bfseries Évaluation pratique — Réseaux locaux \quad
   \fcolorbox{black}{white}{\ SUJET \variante\ }\par}
  \vspace{1pt}
  {\bfseries Adressage IP, commutation et routage — BTS CIEL 1re année — durée : 2 heures\par}
  \ifcorrige\vspace{2pt}{\bfseries\color{corr}CORRIGÉ ET BARÈME — document enseignant}\par\fi
\end{center}
\vspace{-2pt}
{\renewcommand{\arraystretch}{1.35}
\begin{tabularx}{\linewidth}{|X|X|p{2.2cm}|p{2.6cm}|}
  \hline
  \textbf{Nom :} & \textbf{Prénom :} & \textbf{Poste n°} & \textbf{Note :} \hfill /20 \\ \hline
\end{tabularx}}

\textbf{Consignes :} travail \textbf{individuel}, documents et Internet \textbf{non autorisés}. Ce sujet est le
\textbf{document-réponse}. Fichier Packet Tracer : \ip{Nom\_Prenom\_EvalLAN\_\variante.pkt}, à enregistrer
régulièrement dans le dossier de rendu. Les « visas » sont apposés par le professeur, que vous appelez à chaque test réussi.
Durées conseillées : A 30 min — B 1 h 10 — C 20 min.

%==========================================================================
\section{Partie A — Questions de cours et calculs \hfill /6}

\question{A1}{2 pts}{Complétez le tableau en considérant le \textbf{masque par défaut} de la classe de chaque adresse.}
{\small\renewcommand{\arraystretch}{1.15}
\begin{tabularx}{\linewidth}{|l|c|C|C|C|c|}
  \hline
  \textbf{Adresse IP} & \textbf{Classe} & \textbf{Masque par défaut} & \textbf{Adresse réseau} &
  \textbf{Adresse de diffusion} & \textbf{Nb d'hôtes} \\ \hline
  \AIrows
\end{tabularx}}
\bareme{0,5 pt par ligne juste + 0,5 pt si les trois nombres d'hôtes sont justes ($2^{n}-2$, $n$ = nombre de bits d'hôte).}

\question{A2}{1 pt}{Indiquez le type de câble (\textbf{droit} ou \textbf{croisé}) pour chaque liaison.}
{\small\renewcommand{\arraystretch}{1.15}
\begin{tabularx}{\linewidth}{|X|c||X|c|}
  \hline
  \hl PC $\leftrightarrow$ commutateur & \makebox[1.6cm]{\sol{droit}} & Commutateur $\leftrightarrow$ routeur & \makebox[1.6cm]{\sol{droit}} \\ \hline
  \hl PC $\leftrightarrow$ PC (réseau P2P) & \sol{croisé} & Concentrateur $\leftrightarrow$ concentrateur & \sol{croisé} \\ \hline
\end{tabularx}}
\bareme{0,25 pt par case. Équipements de types différents : câble droit ; de même type : câble croisé.}

\question{A3}{1 pt}{Comment un \textbf{concentrateur (hub)} et un \textbf{commutateur (switch)} traitent-ils une trame
reçue ? Lequel limite les collisions, et pourquoi ?}
\zone{2.9cm}{Le \textbf{hub} (couche 1) répète la trame sur \textbf{tous ses autres ports} : un seul domaine de collision
partagé. Le \textbf{switch} (couche 2) lit l'adresse MAC de destination et, grâce à sa \textbf{table MAC}, ne l'envoie
\textbf{que sur le port du destinataire} : un domaine de collision par port, c'est lui qui limite les collisions.
\emph{0,5 pt hub, 0,5 pt switch.}}

\question{A4}{1 pt}{Qu'est-ce que la \textbf{passerelle par défaut} d'un poste ? Quelle condition son adresse doit-elle respecter ?}
\zone{2.4cm}{Adresse IP de l'interface du \textbf{routeur} à laquelle le poste envoie les paquets destinés à un
\textbf{autre réseau}. Elle doit \textbf{appartenir au même réseau} que le poste. \emph{0,5 + 0,5 pt.}}

\question{A5}{0,5 pt}{Un portable configuré en DHCP affiche \ip{\apipa} / \ip{255.255.0.0}. Que s'est-il passé ?}
\zone{1.5cm}{Aucun \textbf{serveur DHCP} n'a répondu : le poste s'est attribué une adresse \textbf{APIPA} (169.254.0.0/16).}

\question{A6}{0,5 pt}{Quel protocole utilise \texttt{ping} ? Nommez les deux messages échangés et la couche OSI des adresses IP.}
\zone{1.5cm}{\textbf{ICMP} ; \textbf{Echo Request} / \textbf{Echo Reply} ; adresses IP en \textbf{couche 3 (Réseau)}.}

%==========================================================================
\newpage
\section{Partie B — Réalisation sous Packet Tracer \hfill /10}

\begin{cdc}
\entreprise{} vous confie son réseau : \textbf{trois réseaux} reliés par un \textbf{routeur unique Cisco 2911}
nommé \ip{\Rnom} (interfaces G0/0, G0/1, G0/2).
\begin{itemize}[leftmargin=12pt]
  \item \textbf{Réseau \resUn} : \textbf{4 postes} PC-\pUn1 à PC-\pUn4, commutateur \ip{SW-\pUn}, adressage statique,
        réseau \ip{\netUn.0/24}.
  \item \textbf{Réseau \resDeux} : \textbf{2 postes} PC-\pDeux1, PC-\pDeux2 et \textbf{1 imprimante} IMP-\pDeux,
        commutateur \ip{SW-\pDeux}, adressage statique, réseau de classe B \ip{\netDeux.0.0/16}.
  \item \textbf{Réseau Serveur} : serveur \ip{\srvNom} d'adresse \ip{\netSrv.\srvH/24}, relié \textbf{directement} au routeur.
  \item La passerelle de chaque réseau est sa \textbf{dernière adresse utilisable}. Tous les équipements communiquent
        entre eux et accèdent à la page web du serveur. Nommage imposé (« Display Name »).
\end{itemize}
\end{cdc}

\question{B1}{2 pts}{Complétez le plan d'adressage \textbf{avant} la réalisation.}
{\footnotesize\setlength{\tabcolsep}{3pt}\renewcommand{\arraystretch}{1.15}
\begin{tabularx}{\linewidth}{|l|C|C|C|C|C|C|}
  \hline
  \textbf{Réseau} & \textbf{Adresse réseau} & \textbf{Masque} & \textbf{1re @ utilisable} &
  \textbf{Dernière @ utilisable} & \textbf{Diffusion} & \textbf{Passerelle} \\ \hline
  \hl \resUn & \sol{\netUn.0} & \sol{255.255.255.0} & \sol{\netUn.1} & \sol{\netUn.254} & \sol{\netUn.255} & \sol{\netUn.254} \\ \hline
  \hl \resDeux & \sol{\netDeux.0.0} & \sol{255.255.0.0} & \sol{\netDeux.0.1} & \sol{\netDeux.255.254} & \sol{\netDeux.255.255} & \sol{\netDeux.255.254} \\ \hline
  \hl Serveur & \sol{\netSrv.0} & \sol{255.255.255.0} & \sol{\netSrv.1} & \sol{\netSrv.254} & \sol{\netSrv.255} & \sol{\netSrv.254} \\ \hline
\end{tabularx}}
\bareme{0,5 pt par ligne juste + 0,5 pt si la convention de passerelle est respectée partout.}

\question{B2}{2 pts}{Construisez la topologie (équipements, câblage, nommage) et relevez l'adressage configuré.}
{\footnotesize\setlength{\tabcolsep}{3pt}\renewcommand{\arraystretch}{1.1}
\begin{tabularx}{\linewidth}{|l|C|C|C||l|C|C|}
  \hline
  \textbf{Hôte} & \textbf{Adresse IP} & \textbf{Masque} & \textbf{Passerelle} & \textbf{Interface} & \textbf{Adresse IP} & \textbf{Masque} \\ \hline
  \hl PC-\pUn1 & \sol{\netUn.1} & \sol{255.255.255.0} & \sol{\netUn.254} & \Rnom{} G0/0 & \sol{\netUn.254} & \sol{255.255.255.0} \\ \hline
  \hl PC-\pDeux1 & \sol{\netDeux.0.1} & \sol{255.255.0.0} & \sol{\netDeux.255.254} & \Rnom{} G0/1 & \sol{\netDeux.255.254} & \sol{255.255.0.0} \\ \hline
  \hl IMP-\pDeux & \sol{\netDeux.0.10} & \sol{255.255.0.0} & \sol{\netDeux.255.254} & \Rnom{} G0/2 & \sol{\netSrv.254} & \sol{255.255.255.0} \\ \hline
  \hl \srvNom & \ip{\netSrv.\srvH} & \sol{255.255.255.0} & \sol{\netSrv.254} & \multicolumn{3}{c|}{\cellcolor{black!8}} \\ \hline
\end{tabularx}}
\bareme{Adresses d'hôtes libres (uniques, dans le bon réseau) ; affectation G0/x libre. Sur le .pkt : équipements 0,5 ;
câblage 1 (câbles droits vers les switchs et switch–routeur, \textbf{croisé} serveur–routeur, liens verts) ; nommage 0,5.}

\question{B3}{3 pts}{Configurez tous les équipements : IP, masques, passerelles des hôtes ; interfaces du routeur
(« Port Status : On »).}
\bareme{Sur le .pkt : interfaces routeur 1 pt ; IP/masques des hôtes 1 pt ; passerelles (imprimante et serveur compris)
1 pt. $-0,25$ par erreur, sans descendre sous 0 par item.}

\question{B4}{3 pts}{Réalisez les tests, notez la commande et le résultat, puis faites viser chaque test réussi.}
{\small\renewcommand{\arraystretch}{1.1}
\begin{tabularx}{\linewidth}{|c|l|l|C|c|c|}
  \hline
  \textbf{N°} & \textbf{Source} & \textbf{Destination} & \textbf{Commande / action} & \textbf{Résultat} & \textbf{Visa prof.} \\ \hline
  \hl 1 & PC-\pUn1 & PC-\pUn4 & \sol{ping \netUn.4} & \sol{réussi} & \hspace{1.4cm} \\ \hline
  \hl 2 & PC-\pDeux1 & IMP-\pDeux & \sol{ping \netDeux.0.10} & \sol{réussi} & \\ \hline
  \hl 3 & PC-\pUn2 & PC-\pDeux2 & \sol{ping \netDeux.0.2} & \sol{réussi} & \\ \hline
  \hl 4 & PC-\pDeux2 & \srvNom & \sol{ping \netSrv.\srvH} & \sol{réussi} & \\ \hline
  \hl 5 & PC-\pUn3 & page web du serveur & \sol{navigateur : {\NoAutoSpacing http://\netSrv.\srvH}} & \sol{page affichée} & \\ \hline
\end{tabularx}}
\bareme{0,5 pt par test visé + 0,5 pt pour B5. Le premier ping inter-réseaux peut échouer (résolution ARP) ; les suivants doivent réussir.}

\question{B5}{incl. B4}{Test n°3 : par quels équipements passe le paquet ? Pourquoi PC-\pUn2 l'envoie-t-il au routeur ?}
\zone{2.3cm}{PC-\pUn2 $\to$ SW-\pUn{} $\to$ \Rnom{} $\to$ SW-\pDeux{} $\to$ PC-\pDeux2. Grâce à son masque, PC-\pUn2 constate que
la destination (\netDeux.0.0/16) n'est \textbf{pas dans son réseau} (\netUn.0/24) : il envoie le paquet à sa
\textbf{passerelle par défaut}, qui le route.}

%==========================================================================
\newpage
\section{Partie C — Diagnostic d'un réseau en panne \hfill /4}

Un technicien constate qu'\textbf{aucun poste ne parvient à joindre le serveur SRV-C} et que certains postes
communiquent mal entre eux. Les câbles sont du bon type.

\begin{center}
\begin{minipage}{0.37\linewidth}\centering
\begin{tikzpicture}[font=\footnotesize,
  eq/.style={draw,fill=qbg,rounded corners=2pt,minimum width=9mm,minimum height=5mm,inner sep=2pt},
  every path/.style={thick}]
  \node[eq] (SW) {SW-1};
  \node[eq,right=9mm of SW] (R) {R1};
  \node[eq,right=7mm of R] (SRV) {SRV-C};
  \node[eq,above left=5mm and -3mm of SW] (A) {PC-A};
  \node[eq,above right=5mm and -3mm of SW] (B) {PC-B};
  \node[eq,below left=5mm and -3mm of SW] (C) {PC-C};
  \node[eq,below right=5mm and -3mm of SW] (D) {PC-D};
  \foreach \n in {A,B,C,D} \draw (\n) -- (SW);
  \draw (SW) -- node[above,font=\tiny]{G0/0} (R);
  \draw (R) -- node[above,font=\tiny]{G0/1} (SRV);
\end{tikzpicture}
\end{minipage}\hfill
\begin{minipage}{0.61\linewidth}\centering
{\small\renewcommand{\arraystretch}{1.05}\setlength{\tabcolsep}{4pt}
\begin{tabular}{lllll}
  \toprule
  \textbf{Équipement} & \textbf{Adresse IP} & \textbf{Masque} & \textbf{Passerelle} & \textbf{Port} \\
  \midrule
  \Crows
  \bottomrule
\end{tabular}}
\end{minipage}
\end{center}

\question{C1}{4 pts}{Relevez les \textbf{quatre anomalies}. Pour chacune : conséquence et correction.}
{\small\renewcommand{\arraystretch}{1.1}
\ifcorrige\def\hc{}\else\def\hc{\rule[-10mm]{0pt}{14mm}}\fi
\begin{tabularx}{\linewidth}{|c|>{\raggedright\arraybackslash}p{2cm}|L|L|L|}
  \hline
  \textbf{N°} & \textbf{Équipement} & \textbf{Anomalie} & \textbf{Conséquence} & \textbf{Correction} \\ \hline
  \ifcorrige\Csol\else
  \hc 1 & & & & \\ \hline
  \hc 2 & & & & \\ \hline
  \hc 3 & & & & \\ \hline
  \hc 4 & & & & \\ \hline
  \fi
\end{tabularx}}
\bareme{1 pt par anomalie (0,25 équipement, anomalie, conséquence, correction ; ordre indifférent). \Cbareme}

%==========================================================================
\section{Bonus — Extension sans fil \hfill +1}

\begin{cdc}Ajoutez au réseau \resUn{} un \textbf{AccessPoint-PT} relié à SW-\pUn{} (SSID \ip{\ssid}, \textbf{WPA2-PSK})
et un portable \ip{Laptop-\pUn} équipé d'un module \texttt{WPC300N}, en adressage \textbf{statique}.
Vérifiez qu'il joint \srvNom.\end{cdc}
{\small\renewcommand{\arraystretch}{1.1}
\begin{tabularx}{\linewidth}{|l|C|C|C|c|}
  \hline
  \textbf{Équipement} & \textbf{Adresse IP} & \textbf{Masque} & \textbf{Passerelle} & \textbf{Visa prof.} \\ \hline
  \hl Laptop-\pUn & \sol{\netUn.50} & \sol{255.255.255.0} & \sol{\netUn.254} & \hspace{1.4cm} \\ \hline
\end{tabularx}}
\bareme{Le point d'accès fait office de pont : le portable est dans \netUn.0/24. Bonus si associé en WPA2-PSK et ping
vers \netSrv.\srvH{} réussi. Note plafonnée à 20.}

\end{document}
"
%   pdflatex -jobname=Eval_Reseau_CorrigeA "\def\variante{A}\def\modecorrige{1}% Évaluation pratique — Réseaux locaux (BTS CIEL 1re année) — sujets A et B
% Compilation (4 PDF) :
%   pdflatex -jobname=Eval_Reseau_SujetA   "\def\variante{A}\def\modecorrige{0}\input{Eval_TP_Reseau_LAN.tex}"
%   pdflatex -jobname=Eval_Reseau_SujetB   "\def\variante{B}\def\modecorrige{0}\input{Eval_TP_Reseau_LAN.tex}"
%   pdflatex -jobname=Eval_Reseau_CorrigeA "\def\variante{A}\def\modecorrige{1}\input{Eval_TP_Reseau_LAN.tex}"
%   pdflatex -jobname=Eval_Reseau_CorrigeB "\def\variante{B}\def\modecorrige{1}\input{Eval_TP_Reseau_LAN.tex}"
\documentclass[10pt,a4paper]{article}

\usepackage[T1]{fontenc}
\usepackage[utf8]{inputenc}
\usepackage{lmodern}
\usepackage[french]{babel}
\usepackage[a4paper,hmargin=1.5cm,top=1.3cm,bottom=1.6cm,footskip=0.7cm]{geometry}
\usepackage[table]{xcolor}
\usepackage{booktabs,tabularx,array}
\usepackage{enumitem}
\usepackage{fancyhdr}
\usepackage[most]{tcolorbox}
\usepackage{tikz}
\usetikzlibrary{positioning}
\usepackage{titlesec}

\providecommand{\modecorrige}{0}
\providecommand{\variante}{A}
\newif\ifcorrige
\ifnum\modecorrige=1 \corrigetrue \else \corrigefalse \fi

\definecolor{ipred}{HTML}{C0392B}
\definecolor{cdcbar}{HTML}{A67C3D}
\definecolor{cdcbg}{HTML}{F4F1EA}
\definecolor{qbar}{HTML}{4A6B5D}
\definecolor{qbg}{HTML}{EDF1EF}
\definecolor{corr}{HTML}{1F5FA8}
\definecolor{corrbg}{HTML}{E8F0FA}
\definecolor{grisfoot}{HTML}{828282}

\newcommand{\ip}[1]{\textcolor{ipred}{\texttt{#1}}}
\newcommand{\pts}[1]{\hfill\mbox{\textbf{(#1)}}}
\newcommand{\sol}[1]{\ifcorrige\textcolor{corr}{#1}\fi}
\newcommand{\bareme}[1]{\ifcorrige\par{\footnotesize\color{corr}#1\par}\fi}
\newcommand{\hl}{\ifcorrige\rule[-1.6mm]{0pt}{5.4mm}\else\rule[-2.2mm]{0pt}{6.6mm}\fi}
\newcolumntype{C}{>{\centering\arraybackslash}X}
\newcolumntype{L}{>{\raggedright\arraybackslash}X}

\titleformat{\section}{\large\bfseries}{}{0pt}{}
\titlespacing*{\section}{0pt}{6pt}{3pt}

\tcbset{qstyle/.style={enhanced,boxrule=0pt,leftrule=2.5pt,arc=0pt,outer arc=0pt,
  left=5pt,right=5pt,top=2pt,bottom=2pt,before skip=4pt,after skip=3pt}}
\newtcolorbox{cdc}{qstyle,colback=cdcbg,colframe=cdcbar}
\newtcolorbox{quest}{qstyle,colback=qbg,colframe=qbar}

% Zone de réponse : vierge dans le sujet, remplie dans le corrigé
\newcommand{\zone}[2]{%
  \ifcorrige
    \begin{tcolorbox}[colback=corrbg,colframe=corr,boxrule=0.5pt,arc=1pt,left=4pt,right=4pt,top=1pt,bottom=1pt,
      before skip=2pt,after skip=2pt,fontupper=\footnotesize\color{corr}]#2\end{tcolorbox}%
  \else
    \begin{tcolorbox}[colback=white,colframe=black!45,boxrule=0.4pt,arc=1pt,height=#1,
      before skip=2pt,after skip=2pt]\end{tcolorbox}%
  \fi}

\newcommand{\question}[3]{\begin{quest}\textbf{#1.} #3\pts{#2}\end{quest}}

%==========================================================================
% Données propres à chaque variante
%==========================================================================
\if\variante A
  \def\entreprise{Le cabinet d'architectes \textbf{ArchiNova}}
  \def\Rnom{R-ArchiNova}
  \def\resUn{Bureau d'études}\def\pUn{BE}\def\netUn{192.168.40}
  \def\resDeux{Accueil}\def\pDeux{ACC}\def\netDeux{172.16}
  \def\srvNom{SRV-Intranet}\def\netSrv{192.168.200}\def\srvH{10}
  \def\ssid{ArchiNova-WiFi}
  \def\AIrows{%
    \hl\ip{192.168.40.25} & \sol{C} & \sol{255.255.255.0} & \sol{192.168.40.0} & \sol{192.168.40.255} & \sol{254} \\ \hline
    \hl\ip{172.16.5.10}   & \sol{B} & \sol{255.255.0.0}   & \sol{172.16.0.0}   & \sol{172.16.255.255} & \sol{65\,534} \\ \hline
    \hl\ip{10.3.2.1}      & \sol{A} & \sol{255.0.0.0}     & \sol{10.0.0.0}     & \sol{10.255.255.255} & \sol{16\,777\,214} \\ \hline}
  \def\apipa{169.254.12.7}
  \def\Cnet{192.168.60}
  \def\Crows{%
    PC-A & \ip{192.168.60.10} & \ip{255.255.255.0} & \ip{192.168.60.254} & — \\
    PC-B & \ip{192.168.60.20} & \ip{255.255.255.0} & \ip{192.168.60.1}   & — \\
    PC-C & \ip{192.168.61.30} & \ip{255.255.255.0} & \ip{192.168.60.254} & — \\
    PC-D & \ip{192.168.60.10} & \ip{255.255.255.0} & \ip{192.168.60.254} & — \\
    SRV-C & \ip{10.0.0.5}     & \ip{255.0.0.0}     & \ip{10.0.0.254}     & — \\
    R1 — G0/0 & \ip{192.168.60.254} & \ip{255.255.255.0} & — & On \\
    R1 — G0/1 & \ip{10.0.0.254}     & \ip{255.0.0.0}     & — & Off \\}
  \def\Csol{%
    \hl 1 & \sol{PC-B} & \sol{Passerelle .60.1 : pas l'adresse du routeur.} & \sol{Joint son réseau, mais aucun autre (serveur).} & \sol{Passerelle 192.168.60.254.} \\ \hline
    \hl 2 & \sol{PC-C} & \sol{IP hors du réseau 192.168.60.0/24.} & \sol{Ne joint ni les postes ni le routeur.} & \sol{Ex. 192.168.60.30.} \\ \hline
    \hl 3 & \sol{PC-D} & \sol{Même IP que PC-A (doublon).} & \sol{Conflit d'adresse : PC-A et PC-D perturbés.} & \sol{Adresse unique, ex. .60.40.} \\ \hline
    \hl 4 & \sol{R1 G0/1} & \sol{Interface désactivée (Off).} & \sol{Réseau du serveur injoignable pour tous.} & \sol{Port Status : On.} \\ \hline}
  \def\Cbareme{PC-A et SRV-C sont correctement configurés ; accepter PC-A cité à la place de PC-D pour le doublon.}
\else
  \def\entreprise{La clinique vétérinaire \textbf{VetoPlus}}
  \def\Rnom{R-VetoPlus}
  \def\resUn{Soins}\def\pUn{SOI}\def\netUn{192.168.70}
  \def\resDeux{Secrétariat}\def\pDeux{SEC}\def\netDeux{172.20}
  \def\srvNom{SRV-Dossiers}\def\netSrv{192.168.250}\def\srvH{20}
  \def\ssid{VetoPlus-WiFi}
  \def\AIrows{%
    \hl\ip{192.168.15.130} & \sol{C} & \sol{255.255.255.0} & \sol{192.168.15.0} & \sol{192.168.15.255} & \sol{254} \\ \hline
    \hl\ip{172.31.200.7}   & \sol{B} & \sol{255.255.0.0}   & \sol{172.31.0.0}   & \sol{172.31.255.255} & \sol{65\,534} \\ \hline
    \hl\ip{10.45.6.9}      & \sol{A} & \sol{255.0.0.0}     & \sol{10.0.0.0}     & \sol{10.255.255.255} & \sol{16\,777\,214} \\ \hline}
  \def\apipa{169.254.88.3}
  \def\Cnet{192.168.80}
  \def\Crows{%
    PC-A & \ip{192.168.8.11}   & \ip{255.255.255.0} & \ip{192.168.80.254} & — \\
    PC-B & \ip{192.168.80.12}  & \ip{255.255.255.0} & \ip{192.168.80.254} & — \\
    PC-C & \ip{192.168.80.13}  & \ip{255.255.255.0} & \ip{192.168.80.100} & — \\
    PC-D & \ip{192.168.80.12}  & \ip{255.255.255.0} & \ip{192.168.80.254} & — \\
    SRV-C & \ip{10.0.0.8}      & \ip{255.0.0.0}     & \ip{10.0.0.254}     & — \\
    R1 — G0/0 & \ip{192.168.80.254} & \ip{255.255.255.0} & — & Off \\
    R1 — G0/1 & \ip{10.0.0.254}     & \ip{255.0.0.0}     & — & On \\}
  \def\Csol{%
    \hl 1 & \sol{PC-A} & \sol{IP hors du réseau 192.168.80.0/24.} & \sol{Ne joint ni les postes ni le routeur.} & \sol{Ex. 192.168.80.11.} \\ \hline
    \hl 2 & \sol{PC-C} & \sol{Passerelle .80.100 : pas l'adresse du routeur.} & \sol{Joint son réseau, mais aucun autre (serveur).} & \sol{Passerelle 192.168.80.254.} \\ \hline
    \hl 3 & \sol{PC-D} & \sol{Même IP que PC-B (doublon).} & \sol{Conflit d'adresse : PC-B et PC-D perturbés.} & \sol{Adresse unique, ex. .80.14.} \\ \hline
    \hl 4 & \sol{R1 G0/0} & \sol{Interface côté LAN désactivée (Off).} & \sol{Aucun poste ne joint le routeur ni le serveur.} & \sol{Port Status : On.} \\ \hline}
  \def\Cbareme{PC-B et SRV-C sont correctement configurés ; accepter PC-B cité à la place de PC-D pour le doublon.}
\fi

\pagestyle{fancy}
\fancyhf{}
\renewcommand{\headrulewidth}{0pt}
\fancyfoot[L]{\footnotesize\itshape\color{grisfoot}Évaluation — Réseaux locaux — Sujet \variante\ifcorrige\ — CORRIGÉ\fi}
\fancyfoot[R]{\footnotesize\color{grisfoot}\thepage/3}

\setlist{itemsep=0pt,topsep=2pt,parsep=0pt}
\setlength{\parindent}{0pt}
\setlength{\parskip}{2pt}

\begin{document}

\begin{center}
  {\Large\bfseries Évaluation pratique — Réseaux locaux \quad
   \fcolorbox{black}{white}{\ SUJET \variante\ }\par}
  \vspace{1pt}
  {\bfseries Adressage IP, commutation et routage — BTS CIEL 1re année — durée : 2 heures\par}
  \ifcorrige\vspace{2pt}{\bfseries\color{corr}CORRIGÉ ET BARÈME — document enseignant}\par\fi
\end{center}
\vspace{-2pt}
{\renewcommand{\arraystretch}{1.35}
\begin{tabularx}{\linewidth}{|X|X|p{2.2cm}|p{2.6cm}|}
  \hline
  \textbf{Nom :} & \textbf{Prénom :} & \textbf{Poste n°} & \textbf{Note :} \hfill /20 \\ \hline
\end{tabularx}}

\textbf{Consignes :} travail \textbf{individuel}, documents et Internet \textbf{non autorisés}. Ce sujet est le
\textbf{document-réponse}. Fichier Packet Tracer : \ip{Nom\_Prenom\_EvalLAN\_\variante.pkt}, à enregistrer
régulièrement dans le dossier de rendu. Les « visas » sont apposés par le professeur, que vous appelez à chaque test réussi.
Durées conseillées : A 30 min — B 1 h 10 — C 20 min.

%==========================================================================
\section{Partie A — Questions de cours et calculs \hfill /6}

\question{A1}{2 pts}{Complétez le tableau en considérant le \textbf{masque par défaut} de la classe de chaque adresse.}
{\small\renewcommand{\arraystretch}{1.15}
\begin{tabularx}{\linewidth}{|l|c|C|C|C|c|}
  \hline
  \textbf{Adresse IP} & \textbf{Classe} & \textbf{Masque par défaut} & \textbf{Adresse réseau} &
  \textbf{Adresse de diffusion} & \textbf{Nb d'hôtes} \\ \hline
  \AIrows
\end{tabularx}}
\bareme{0,5 pt par ligne juste + 0,5 pt si les trois nombres d'hôtes sont justes ($2^{n}-2$, $n$ = nombre de bits d'hôte).}

\question{A2}{1 pt}{Indiquez le type de câble (\textbf{droit} ou \textbf{croisé}) pour chaque liaison.}
{\small\renewcommand{\arraystretch}{1.15}
\begin{tabularx}{\linewidth}{|X|c||X|c|}
  \hline
  \hl PC $\leftrightarrow$ commutateur & \makebox[1.6cm]{\sol{droit}} & Commutateur $\leftrightarrow$ routeur & \makebox[1.6cm]{\sol{droit}} \\ \hline
  \hl PC $\leftrightarrow$ PC (réseau P2P) & \sol{croisé} & Concentrateur $\leftrightarrow$ concentrateur & \sol{croisé} \\ \hline
\end{tabularx}}
\bareme{0,25 pt par case. Équipements de types différents : câble droit ; de même type : câble croisé.}

\question{A3}{1 pt}{Comment un \textbf{concentrateur (hub)} et un \textbf{commutateur (switch)} traitent-ils une trame
reçue ? Lequel limite les collisions, et pourquoi ?}
\zone{2.9cm}{Le \textbf{hub} (couche 1) répète la trame sur \textbf{tous ses autres ports} : un seul domaine de collision
partagé. Le \textbf{switch} (couche 2) lit l'adresse MAC de destination et, grâce à sa \textbf{table MAC}, ne l'envoie
\textbf{que sur le port du destinataire} : un domaine de collision par port, c'est lui qui limite les collisions.
\emph{0,5 pt hub, 0,5 pt switch.}}

\question{A4}{1 pt}{Qu'est-ce que la \textbf{passerelle par défaut} d'un poste ? Quelle condition son adresse doit-elle respecter ?}
\zone{2.4cm}{Adresse IP de l'interface du \textbf{routeur} à laquelle le poste envoie les paquets destinés à un
\textbf{autre réseau}. Elle doit \textbf{appartenir au même réseau} que le poste. \emph{0,5 + 0,5 pt.}}

\question{A5}{0,5 pt}{Un portable configuré en DHCP affiche \ip{\apipa} / \ip{255.255.0.0}. Que s'est-il passé ?}
\zone{1.5cm}{Aucun \textbf{serveur DHCP} n'a répondu : le poste s'est attribué une adresse \textbf{APIPA} (169.254.0.0/16).}

\question{A6}{0,5 pt}{Quel protocole utilise \texttt{ping} ? Nommez les deux messages échangés et la couche OSI des adresses IP.}
\zone{1.5cm}{\textbf{ICMP} ; \textbf{Echo Request} / \textbf{Echo Reply} ; adresses IP en \textbf{couche 3 (Réseau)}.}

%==========================================================================
\newpage
\section{Partie B — Réalisation sous Packet Tracer \hfill /10}

\begin{cdc}
\entreprise{} vous confie son réseau : \textbf{trois réseaux} reliés par un \textbf{routeur unique Cisco 2911}
nommé \ip{\Rnom} (interfaces G0/0, G0/1, G0/2).
\begin{itemize}[leftmargin=12pt]
  \item \textbf{Réseau \resUn} : \textbf{4 postes} PC-\pUn1 à PC-\pUn4, commutateur \ip{SW-\pUn}, adressage statique,
        réseau \ip{\netUn.0/24}.
  \item \textbf{Réseau \resDeux} : \textbf{2 postes} PC-\pDeux1, PC-\pDeux2 et \textbf{1 imprimante} IMP-\pDeux,
        commutateur \ip{SW-\pDeux}, adressage statique, réseau de classe B \ip{\netDeux.0.0/16}.
  \item \textbf{Réseau Serveur} : serveur \ip{\srvNom} d'adresse \ip{\netSrv.\srvH/24}, relié \textbf{directement} au routeur.
  \item La passerelle de chaque réseau est sa \textbf{dernière adresse utilisable}. Tous les équipements communiquent
        entre eux et accèdent à la page web du serveur. Nommage imposé (« Display Name »).
\end{itemize}
\end{cdc}

\question{B1}{2 pts}{Complétez le plan d'adressage \textbf{avant} la réalisation.}
{\footnotesize\setlength{\tabcolsep}{3pt}\renewcommand{\arraystretch}{1.15}
\begin{tabularx}{\linewidth}{|l|C|C|C|C|C|C|}
  \hline
  \textbf{Réseau} & \textbf{Adresse réseau} & \textbf{Masque} & \textbf{1re @ utilisable} &
  \textbf{Dernière @ utilisable} & \textbf{Diffusion} & \textbf{Passerelle} \\ \hline
  \hl \resUn & \sol{\netUn.0} & \sol{255.255.255.0} & \sol{\netUn.1} & \sol{\netUn.254} & \sol{\netUn.255} & \sol{\netUn.254} \\ \hline
  \hl \resDeux & \sol{\netDeux.0.0} & \sol{255.255.0.0} & \sol{\netDeux.0.1} & \sol{\netDeux.255.254} & \sol{\netDeux.255.255} & \sol{\netDeux.255.254} \\ \hline
  \hl Serveur & \sol{\netSrv.0} & \sol{255.255.255.0} & \sol{\netSrv.1} & \sol{\netSrv.254} & \sol{\netSrv.255} & \sol{\netSrv.254} \\ \hline
\end{tabularx}}
\bareme{0,5 pt par ligne juste + 0,5 pt si la convention de passerelle est respectée partout.}

\question{B2}{2 pts}{Construisez la topologie (équipements, câblage, nommage) et relevez l'adressage configuré.}
{\footnotesize\setlength{\tabcolsep}{3pt}\renewcommand{\arraystretch}{1.1}
\begin{tabularx}{\linewidth}{|l|C|C|C||l|C|C|}
  \hline
  \textbf{Hôte} & \textbf{Adresse IP} & \textbf{Masque} & \textbf{Passerelle} & \textbf{Interface} & \textbf{Adresse IP} & \textbf{Masque} \\ \hline
  \hl PC-\pUn1 & \sol{\netUn.1} & \sol{255.255.255.0} & \sol{\netUn.254} & \Rnom{} G0/0 & \sol{\netUn.254} & \sol{255.255.255.0} \\ \hline
  \hl PC-\pDeux1 & \sol{\netDeux.0.1} & \sol{255.255.0.0} & \sol{\netDeux.255.254} & \Rnom{} G0/1 & \sol{\netDeux.255.254} & \sol{255.255.0.0} \\ \hline
  \hl IMP-\pDeux & \sol{\netDeux.0.10} & \sol{255.255.0.0} & \sol{\netDeux.255.254} & \Rnom{} G0/2 & \sol{\netSrv.254} & \sol{255.255.255.0} \\ \hline
  \hl \srvNom & \ip{\netSrv.\srvH} & \sol{255.255.255.0} & \sol{\netSrv.254} & \multicolumn{3}{c|}{\cellcolor{black!8}} \\ \hline
\end{tabularx}}
\bareme{Adresses d'hôtes libres (uniques, dans le bon réseau) ; affectation G0/x libre. Sur le .pkt : équipements 0,5 ;
câblage 1 (câbles droits vers les switchs et switch–routeur, \textbf{croisé} serveur–routeur, liens verts) ; nommage 0,5.}

\question{B3}{3 pts}{Configurez tous les équipements : IP, masques, passerelles des hôtes ; interfaces du routeur
(« Port Status : On »).}
\bareme{Sur le .pkt : interfaces routeur 1 pt ; IP/masques des hôtes 1 pt ; passerelles (imprimante et serveur compris)
1 pt. $-0,25$ par erreur, sans descendre sous 0 par item.}

\question{B4}{3 pts}{Réalisez les tests, notez la commande et le résultat, puis faites viser chaque test réussi.}
{\small\renewcommand{\arraystretch}{1.1}
\begin{tabularx}{\linewidth}{|c|l|l|C|c|c|}
  \hline
  \textbf{N°} & \textbf{Source} & \textbf{Destination} & \textbf{Commande / action} & \textbf{Résultat} & \textbf{Visa prof.} \\ \hline
  \hl 1 & PC-\pUn1 & PC-\pUn4 & \sol{ping \netUn.4} & \sol{réussi} & \hspace{1.4cm} \\ \hline
  \hl 2 & PC-\pDeux1 & IMP-\pDeux & \sol{ping \netDeux.0.10} & \sol{réussi} & \\ \hline
  \hl 3 & PC-\pUn2 & PC-\pDeux2 & \sol{ping \netDeux.0.2} & \sol{réussi} & \\ \hline
  \hl 4 & PC-\pDeux2 & \srvNom & \sol{ping \netSrv.\srvH} & \sol{réussi} & \\ \hline
  \hl 5 & PC-\pUn3 & page web du serveur & \sol{navigateur : {\NoAutoSpacing http://\netSrv.\srvH}} & \sol{page affichée} & \\ \hline
\end{tabularx}}
\bareme{0,5 pt par test visé + 0,5 pt pour B5. Le premier ping inter-réseaux peut échouer (résolution ARP) ; les suivants doivent réussir.}

\question{B5}{incl. B4}{Test n°3 : par quels équipements passe le paquet ? Pourquoi PC-\pUn2 l'envoie-t-il au routeur ?}
\zone{2.3cm}{PC-\pUn2 $\to$ SW-\pUn{} $\to$ \Rnom{} $\to$ SW-\pDeux{} $\to$ PC-\pDeux2. Grâce à son masque, PC-\pUn2 constate que
la destination (\netDeux.0.0/16) n'est \textbf{pas dans son réseau} (\netUn.0/24) : il envoie le paquet à sa
\textbf{passerelle par défaut}, qui le route.}

%==========================================================================
\newpage
\section{Partie C — Diagnostic d'un réseau en panne \hfill /4}

Un technicien constate qu'\textbf{aucun poste ne parvient à joindre le serveur SRV-C} et que certains postes
communiquent mal entre eux. Les câbles sont du bon type.

\begin{center}
\begin{minipage}{0.37\linewidth}\centering
\begin{tikzpicture}[font=\footnotesize,
  eq/.style={draw,fill=qbg,rounded corners=2pt,minimum width=9mm,minimum height=5mm,inner sep=2pt},
  every path/.style={thick}]
  \node[eq] (SW) {SW-1};
  \node[eq,right=9mm of SW] (R) {R1};
  \node[eq,right=7mm of R] (SRV) {SRV-C};
  \node[eq,above left=5mm and -3mm of SW] (A) {PC-A};
  \node[eq,above right=5mm and -3mm of SW] (B) {PC-B};
  \node[eq,below left=5mm and -3mm of SW] (C) {PC-C};
  \node[eq,below right=5mm and -3mm of SW] (D) {PC-D};
  \foreach \n in {A,B,C,D} \draw (\n) -- (SW);
  \draw (SW) -- node[above,font=\tiny]{G0/0} (R);
  \draw (R) -- node[above,font=\tiny]{G0/1} (SRV);
\end{tikzpicture}
\end{minipage}\hfill
\begin{minipage}{0.61\linewidth}\centering
{\small\renewcommand{\arraystretch}{1.05}\setlength{\tabcolsep}{4pt}
\begin{tabular}{lllll}
  \toprule
  \textbf{Équipement} & \textbf{Adresse IP} & \textbf{Masque} & \textbf{Passerelle} & \textbf{Port} \\
  \midrule
  \Crows
  \bottomrule
\end{tabular}}
\end{minipage}
\end{center}

\question{C1}{4 pts}{Relevez les \textbf{quatre anomalies}. Pour chacune : conséquence et correction.}
{\small\renewcommand{\arraystretch}{1.1}
\ifcorrige\def\hc{}\else\def\hc{\rule[-10mm]{0pt}{14mm}}\fi
\begin{tabularx}{\linewidth}{|c|>{\raggedright\arraybackslash}p{2cm}|L|L|L|}
  \hline
  \textbf{N°} & \textbf{Équipement} & \textbf{Anomalie} & \textbf{Conséquence} & \textbf{Correction} \\ \hline
  \ifcorrige\Csol\else
  \hc 1 & & & & \\ \hline
  \hc 2 & & & & \\ \hline
  \hc 3 & & & & \\ \hline
  \hc 4 & & & & \\ \hline
  \fi
\end{tabularx}}
\bareme{1 pt par anomalie (0,25 équipement, anomalie, conséquence, correction ; ordre indifférent). \Cbareme}

%==========================================================================
\section{Bonus — Extension sans fil \hfill +1}

\begin{cdc}Ajoutez au réseau \resUn{} un \textbf{AccessPoint-PT} relié à SW-\pUn{} (SSID \ip{\ssid}, \textbf{WPA2-PSK})
et un portable \ip{Laptop-\pUn} équipé d'un module \texttt{WPC300N}, en adressage \textbf{statique}.
Vérifiez qu'il joint \srvNom.\end{cdc}
{\small\renewcommand{\arraystretch}{1.1}
\begin{tabularx}{\linewidth}{|l|C|C|C|c|}
  \hline
  \textbf{Équipement} & \textbf{Adresse IP} & \textbf{Masque} & \textbf{Passerelle} & \textbf{Visa prof.} \\ \hline
  \hl Laptop-\pUn & \sol{\netUn.50} & \sol{255.255.255.0} & \sol{\netUn.254} & \hspace{1.4cm} \\ \hline
\end{tabularx}}
\bareme{Le point d'accès fait office de pont : le portable est dans \netUn.0/24. Bonus si associé en WPA2-PSK et ping
vers \netSrv.\srvH{} réussi. Note plafonnée à 20.}

\end{document}
"
%   pdflatex -jobname=Eval_Reseau_CorrigeB "\def\variante{B}\def\modecorrige{1}% Évaluation pratique — Réseaux locaux (BTS CIEL 1re année) — sujets A et B
% Compilation (4 PDF) :
%   pdflatex -jobname=Eval_Reseau_SujetA   "\def\variante{A}\def\modecorrige{0}\input{Eval_TP_Reseau_LAN.tex}"
%   pdflatex -jobname=Eval_Reseau_SujetB   "\def\variante{B}\def\modecorrige{0}\input{Eval_TP_Reseau_LAN.tex}"
%   pdflatex -jobname=Eval_Reseau_CorrigeA "\def\variante{A}\def\modecorrige{1}\input{Eval_TP_Reseau_LAN.tex}"
%   pdflatex -jobname=Eval_Reseau_CorrigeB "\def\variante{B}\def\modecorrige{1}\input{Eval_TP_Reseau_LAN.tex}"
\documentclass[10pt,a4paper]{article}

\usepackage[T1]{fontenc}
\usepackage[utf8]{inputenc}
\usepackage{lmodern}
\usepackage[french]{babel}
\usepackage[a4paper,hmargin=1.5cm,top=1.3cm,bottom=1.6cm,footskip=0.7cm]{geometry}
\usepackage[table]{xcolor}
\usepackage{booktabs,tabularx,array}
\usepackage{enumitem}
\usepackage{fancyhdr}
\usepackage[most]{tcolorbox}
\usepackage{tikz}
\usetikzlibrary{positioning}
\usepackage{titlesec}

\providecommand{\modecorrige}{0}
\providecommand{\variante}{A}
\newif\ifcorrige
\ifnum\modecorrige=1 \corrigetrue \else \corrigefalse \fi

\definecolor{ipred}{HTML}{C0392B}
\definecolor{cdcbar}{HTML}{A67C3D}
\definecolor{cdcbg}{HTML}{F4F1EA}
\definecolor{qbar}{HTML}{4A6B5D}
\definecolor{qbg}{HTML}{EDF1EF}
\definecolor{corr}{HTML}{1F5FA8}
\definecolor{corrbg}{HTML}{E8F0FA}
\definecolor{grisfoot}{HTML}{828282}

\newcommand{\ip}[1]{\textcolor{ipred}{\texttt{#1}}}
\newcommand{\pts}[1]{\hfill\mbox{\textbf{(#1)}}}
\newcommand{\sol}[1]{\ifcorrige\textcolor{corr}{#1}\fi}
\newcommand{\bareme}[1]{\ifcorrige\par{\footnotesize\color{corr}#1\par}\fi}
\newcommand{\hl}{\ifcorrige\rule[-1.6mm]{0pt}{5.4mm}\else\rule[-2.2mm]{0pt}{6.6mm}\fi}
\newcolumntype{C}{>{\centering\arraybackslash}X}
\newcolumntype{L}{>{\raggedright\arraybackslash}X}

\titleformat{\section}{\large\bfseries}{}{0pt}{}
\titlespacing*{\section}{0pt}{6pt}{3pt}

\tcbset{qstyle/.style={enhanced,boxrule=0pt,leftrule=2.5pt,arc=0pt,outer arc=0pt,
  left=5pt,right=5pt,top=2pt,bottom=2pt,before skip=4pt,after skip=3pt}}
\newtcolorbox{cdc}{qstyle,colback=cdcbg,colframe=cdcbar}
\newtcolorbox{quest}{qstyle,colback=qbg,colframe=qbar}

% Zone de réponse : vierge dans le sujet, remplie dans le corrigé
\newcommand{\zone}[2]{%
  \ifcorrige
    \begin{tcolorbox}[colback=corrbg,colframe=corr,boxrule=0.5pt,arc=1pt,left=4pt,right=4pt,top=1pt,bottom=1pt,
      before skip=2pt,after skip=2pt,fontupper=\footnotesize\color{corr}]#2\end{tcolorbox}%
  \else
    \begin{tcolorbox}[colback=white,colframe=black!45,boxrule=0.4pt,arc=1pt,height=#1,
      before skip=2pt,after skip=2pt]\end{tcolorbox}%
  \fi}

\newcommand{\question}[3]{\begin{quest}\textbf{#1.} #3\pts{#2}\end{quest}}

%==========================================================================
% Données propres à chaque variante
%==========================================================================
\if\variante A
  \def\entreprise{Le cabinet d'architectes \textbf{ArchiNova}}
  \def\Rnom{R-ArchiNova}
  \def\resUn{Bureau d'études}\def\pUn{BE}\def\netUn{192.168.40}
  \def\resDeux{Accueil}\def\pDeux{ACC}\def\netDeux{172.16}
  \def\srvNom{SRV-Intranet}\def\netSrv{192.168.200}\def\srvH{10}
  \def\ssid{ArchiNova-WiFi}
  \def\AIrows{%
    \hl\ip{192.168.40.25} & \sol{C} & \sol{255.255.255.0} & \sol{192.168.40.0} & \sol{192.168.40.255} & \sol{254} \\ \hline
    \hl\ip{172.16.5.10}   & \sol{B} & \sol{255.255.0.0}   & \sol{172.16.0.0}   & \sol{172.16.255.255} & \sol{65\,534} \\ \hline
    \hl\ip{10.3.2.1}      & \sol{A} & \sol{255.0.0.0}     & \sol{10.0.0.0}     & \sol{10.255.255.255} & \sol{16\,777\,214} \\ \hline}
  \def\apipa{169.254.12.7}
  \def\Cnet{192.168.60}
  \def\Crows{%
    PC-A & \ip{192.168.60.10} & \ip{255.255.255.0} & \ip{192.168.60.254} & — \\
    PC-B & \ip{192.168.60.20} & \ip{255.255.255.0} & \ip{192.168.60.1}   & — \\
    PC-C & \ip{192.168.61.30} & \ip{255.255.255.0} & \ip{192.168.60.254} & — \\
    PC-D & \ip{192.168.60.10} & \ip{255.255.255.0} & \ip{192.168.60.254} & — \\
    SRV-C & \ip{10.0.0.5}     & \ip{255.0.0.0}     & \ip{10.0.0.254}     & — \\
    R1 — G0/0 & \ip{192.168.60.254} & \ip{255.255.255.0} & — & On \\
    R1 — G0/1 & \ip{10.0.0.254}     & \ip{255.0.0.0}     & — & Off \\}
  \def\Csol{%
    \hl 1 & \sol{PC-B} & \sol{Passerelle .60.1 : pas l'adresse du routeur.} & \sol{Joint son réseau, mais aucun autre (serveur).} & \sol{Passerelle 192.168.60.254.} \\ \hline
    \hl 2 & \sol{PC-C} & \sol{IP hors du réseau 192.168.60.0/24.} & \sol{Ne joint ni les postes ni le routeur.} & \sol{Ex. 192.168.60.30.} \\ \hline
    \hl 3 & \sol{PC-D} & \sol{Même IP que PC-A (doublon).} & \sol{Conflit d'adresse : PC-A et PC-D perturbés.} & \sol{Adresse unique, ex. .60.40.} \\ \hline
    \hl 4 & \sol{R1 G0/1} & \sol{Interface désactivée (Off).} & \sol{Réseau du serveur injoignable pour tous.} & \sol{Port Status : On.} \\ \hline}
  \def\Cbareme{PC-A et SRV-C sont correctement configurés ; accepter PC-A cité à la place de PC-D pour le doublon.}
\else
  \def\entreprise{La clinique vétérinaire \textbf{VetoPlus}}
  \def\Rnom{R-VetoPlus}
  \def\resUn{Soins}\def\pUn{SOI}\def\netUn{192.168.70}
  \def\resDeux{Secrétariat}\def\pDeux{SEC}\def\netDeux{172.20}
  \def\srvNom{SRV-Dossiers}\def\netSrv{192.168.250}\def\srvH{20}
  \def\ssid{VetoPlus-WiFi}
  \def\AIrows{%
    \hl\ip{192.168.15.130} & \sol{C} & \sol{255.255.255.0} & \sol{192.168.15.0} & \sol{192.168.15.255} & \sol{254} \\ \hline
    \hl\ip{172.31.200.7}   & \sol{B} & \sol{255.255.0.0}   & \sol{172.31.0.0}   & \sol{172.31.255.255} & \sol{65\,534} \\ \hline
    \hl\ip{10.45.6.9}      & \sol{A} & \sol{255.0.0.0}     & \sol{10.0.0.0}     & \sol{10.255.255.255} & \sol{16\,777\,214} \\ \hline}
  \def\apipa{169.254.88.3}
  \def\Cnet{192.168.80}
  \def\Crows{%
    PC-A & \ip{192.168.8.11}   & \ip{255.255.255.0} & \ip{192.168.80.254} & — \\
    PC-B & \ip{192.168.80.12}  & \ip{255.255.255.0} & \ip{192.168.80.254} & — \\
    PC-C & \ip{192.168.80.13}  & \ip{255.255.255.0} & \ip{192.168.80.100} & — \\
    PC-D & \ip{192.168.80.12}  & \ip{255.255.255.0} & \ip{192.168.80.254} & — \\
    SRV-C & \ip{10.0.0.8}      & \ip{255.0.0.0}     & \ip{10.0.0.254}     & — \\
    R1 — G0/0 & \ip{192.168.80.254} & \ip{255.255.255.0} & — & Off \\
    R1 — G0/1 & \ip{10.0.0.254}     & \ip{255.0.0.0}     & — & On \\}
  \def\Csol{%
    \hl 1 & \sol{PC-A} & \sol{IP hors du réseau 192.168.80.0/24.} & \sol{Ne joint ni les postes ni le routeur.} & \sol{Ex. 192.168.80.11.} \\ \hline
    \hl 2 & \sol{PC-C} & \sol{Passerelle .80.100 : pas l'adresse du routeur.} & \sol{Joint son réseau, mais aucun autre (serveur).} & \sol{Passerelle 192.168.80.254.} \\ \hline
    \hl 3 & \sol{PC-D} & \sol{Même IP que PC-B (doublon).} & \sol{Conflit d'adresse : PC-B et PC-D perturbés.} & \sol{Adresse unique, ex. .80.14.} \\ \hline
    \hl 4 & \sol{R1 G0/0} & \sol{Interface côté LAN désactivée (Off).} & \sol{Aucun poste ne joint le routeur ni le serveur.} & \sol{Port Status : On.} \\ \hline}
  \def\Cbareme{PC-B et SRV-C sont correctement configurés ; accepter PC-B cité à la place de PC-D pour le doublon.}
\fi

\pagestyle{fancy}
\fancyhf{}
\renewcommand{\headrulewidth}{0pt}
\fancyfoot[L]{\footnotesize\itshape\color{grisfoot}Évaluation — Réseaux locaux — Sujet \variante\ifcorrige\ — CORRIGÉ\fi}
\fancyfoot[R]{\footnotesize\color{grisfoot}\thepage/3}

\setlist{itemsep=0pt,topsep=2pt,parsep=0pt}
\setlength{\parindent}{0pt}
\setlength{\parskip}{2pt}

\begin{document}

\begin{center}
  {\Large\bfseries Évaluation pratique — Réseaux locaux \quad
   \fcolorbox{black}{white}{\ SUJET \variante\ }\par}
  \vspace{1pt}
  {\bfseries Adressage IP, commutation et routage — BTS CIEL 1re année — durée : 2 heures\par}
  \ifcorrige\vspace{2pt}{\bfseries\color{corr}CORRIGÉ ET BARÈME — document enseignant}\par\fi
\end{center}
\vspace{-2pt}
{\renewcommand{\arraystretch}{1.35}
\begin{tabularx}{\linewidth}{|X|X|p{2.2cm}|p{2.6cm}|}
  \hline
  \textbf{Nom :} & \textbf{Prénom :} & \textbf{Poste n°} & \textbf{Note :} \hfill /20 \\ \hline
\end{tabularx}}

\textbf{Consignes :} travail \textbf{individuel}, documents et Internet \textbf{non autorisés}. Ce sujet est le
\textbf{document-réponse}. Fichier Packet Tracer : \ip{Nom\_Prenom\_EvalLAN\_\variante.pkt}, à enregistrer
régulièrement dans le dossier de rendu. Les « visas » sont apposés par le professeur, que vous appelez à chaque test réussi.
Durées conseillées : A 30 min — B 1 h 10 — C 20 min.

%==========================================================================
\section{Partie A — Questions de cours et calculs \hfill /6}

\question{A1}{2 pts}{Complétez le tableau en considérant le \textbf{masque par défaut} de la classe de chaque adresse.}
{\small\renewcommand{\arraystretch}{1.15}
\begin{tabularx}{\linewidth}{|l|c|C|C|C|c|}
  \hline
  \textbf{Adresse IP} & \textbf{Classe} & \textbf{Masque par défaut} & \textbf{Adresse réseau} &
  \textbf{Adresse de diffusion} & \textbf{Nb d'hôtes} \\ \hline
  \AIrows
\end{tabularx}}
\bareme{0,5 pt par ligne juste + 0,5 pt si les trois nombres d'hôtes sont justes ($2^{n}-2$, $n$ = nombre de bits d'hôte).}

\question{A2}{1 pt}{Indiquez le type de câble (\textbf{droit} ou \textbf{croisé}) pour chaque liaison.}
{\small\renewcommand{\arraystretch}{1.15}
\begin{tabularx}{\linewidth}{|X|c||X|c|}
  \hline
  \hl PC $\leftrightarrow$ commutateur & \makebox[1.6cm]{\sol{droit}} & Commutateur $\leftrightarrow$ routeur & \makebox[1.6cm]{\sol{droit}} \\ \hline
  \hl PC $\leftrightarrow$ PC (réseau P2P) & \sol{croisé} & Concentrateur $\leftrightarrow$ concentrateur & \sol{croisé} \\ \hline
\end{tabularx}}
\bareme{0,25 pt par case. Équipements de types différents : câble droit ; de même type : câble croisé.}

\question{A3}{1 pt}{Comment un \textbf{concentrateur (hub)} et un \textbf{commutateur (switch)} traitent-ils une trame
reçue ? Lequel limite les collisions, et pourquoi ?}
\zone{2.9cm}{Le \textbf{hub} (couche 1) répète la trame sur \textbf{tous ses autres ports} : un seul domaine de collision
partagé. Le \textbf{switch} (couche 2) lit l'adresse MAC de destination et, grâce à sa \textbf{table MAC}, ne l'envoie
\textbf{que sur le port du destinataire} : un domaine de collision par port, c'est lui qui limite les collisions.
\emph{0,5 pt hub, 0,5 pt switch.}}

\question{A4}{1 pt}{Qu'est-ce que la \textbf{passerelle par défaut} d'un poste ? Quelle condition son adresse doit-elle respecter ?}
\zone{2.4cm}{Adresse IP de l'interface du \textbf{routeur} à laquelle le poste envoie les paquets destinés à un
\textbf{autre réseau}. Elle doit \textbf{appartenir au même réseau} que le poste. \emph{0,5 + 0,5 pt.}}

\question{A5}{0,5 pt}{Un portable configuré en DHCP affiche \ip{\apipa} / \ip{255.255.0.0}. Que s'est-il passé ?}
\zone{1.5cm}{Aucun \textbf{serveur DHCP} n'a répondu : le poste s'est attribué une adresse \textbf{APIPA} (169.254.0.0/16).}

\question{A6}{0,5 pt}{Quel protocole utilise \texttt{ping} ? Nommez les deux messages échangés et la couche OSI des adresses IP.}
\zone{1.5cm}{\textbf{ICMP} ; \textbf{Echo Request} / \textbf{Echo Reply} ; adresses IP en \textbf{couche 3 (Réseau)}.}

%==========================================================================
\newpage
\section{Partie B — Réalisation sous Packet Tracer \hfill /10}

\begin{cdc}
\entreprise{} vous confie son réseau : \textbf{trois réseaux} reliés par un \textbf{routeur unique Cisco 2911}
nommé \ip{\Rnom} (interfaces G0/0, G0/1, G0/2).
\begin{itemize}[leftmargin=12pt]
  \item \textbf{Réseau \resUn} : \textbf{4 postes} PC-\pUn1 à PC-\pUn4, commutateur \ip{SW-\pUn}, adressage statique,
        réseau \ip{\netUn.0/24}.
  \item \textbf{Réseau \resDeux} : \textbf{2 postes} PC-\pDeux1, PC-\pDeux2 et \textbf{1 imprimante} IMP-\pDeux,
        commutateur \ip{SW-\pDeux}, adressage statique, réseau de classe B \ip{\netDeux.0.0/16}.
  \item \textbf{Réseau Serveur} : serveur \ip{\srvNom} d'adresse \ip{\netSrv.\srvH/24}, relié \textbf{directement} au routeur.
  \item La passerelle de chaque réseau est sa \textbf{dernière adresse utilisable}. Tous les équipements communiquent
        entre eux et accèdent à la page web du serveur. Nommage imposé (« Display Name »).
\end{itemize}
\end{cdc}

\question{B1}{2 pts}{Complétez le plan d'adressage \textbf{avant} la réalisation.}
{\footnotesize\setlength{\tabcolsep}{3pt}\renewcommand{\arraystretch}{1.15}
\begin{tabularx}{\linewidth}{|l|C|C|C|C|C|C|}
  \hline
  \textbf{Réseau} & \textbf{Adresse réseau} & \textbf{Masque} & \textbf{1re @ utilisable} &
  \textbf{Dernière @ utilisable} & \textbf{Diffusion} & \textbf{Passerelle} \\ \hline
  \hl \resUn & \sol{\netUn.0} & \sol{255.255.255.0} & \sol{\netUn.1} & \sol{\netUn.254} & \sol{\netUn.255} & \sol{\netUn.254} \\ \hline
  \hl \resDeux & \sol{\netDeux.0.0} & \sol{255.255.0.0} & \sol{\netDeux.0.1} & \sol{\netDeux.255.254} & \sol{\netDeux.255.255} & \sol{\netDeux.255.254} \\ \hline
  \hl Serveur & \sol{\netSrv.0} & \sol{255.255.255.0} & \sol{\netSrv.1} & \sol{\netSrv.254} & \sol{\netSrv.255} & \sol{\netSrv.254} \\ \hline
\end{tabularx}}
\bareme{0,5 pt par ligne juste + 0,5 pt si la convention de passerelle est respectée partout.}

\question{B2}{2 pts}{Construisez la topologie (équipements, câblage, nommage) et relevez l'adressage configuré.}
{\footnotesize\setlength{\tabcolsep}{3pt}\renewcommand{\arraystretch}{1.1}
\begin{tabularx}{\linewidth}{|l|C|C|C||l|C|C|}
  \hline
  \textbf{Hôte} & \textbf{Adresse IP} & \textbf{Masque} & \textbf{Passerelle} & \textbf{Interface} & \textbf{Adresse IP} & \textbf{Masque} \\ \hline
  \hl PC-\pUn1 & \sol{\netUn.1} & \sol{255.255.255.0} & \sol{\netUn.254} & \Rnom{} G0/0 & \sol{\netUn.254} & \sol{255.255.255.0} \\ \hline
  \hl PC-\pDeux1 & \sol{\netDeux.0.1} & \sol{255.255.0.0} & \sol{\netDeux.255.254} & \Rnom{} G0/1 & \sol{\netDeux.255.254} & \sol{255.255.0.0} \\ \hline
  \hl IMP-\pDeux & \sol{\netDeux.0.10} & \sol{255.255.0.0} & \sol{\netDeux.255.254} & \Rnom{} G0/2 & \sol{\netSrv.254} & \sol{255.255.255.0} \\ \hline
  \hl \srvNom & \ip{\netSrv.\srvH} & \sol{255.255.255.0} & \sol{\netSrv.254} & \multicolumn{3}{c|}{\cellcolor{black!8}} \\ \hline
\end{tabularx}}
\bareme{Adresses d'hôtes libres (uniques, dans le bon réseau) ; affectation G0/x libre. Sur le .pkt : équipements 0,5 ;
câblage 1 (câbles droits vers les switchs et switch–routeur, \textbf{croisé} serveur–routeur, liens verts) ; nommage 0,5.}

\question{B3}{3 pts}{Configurez tous les équipements : IP, masques, passerelles des hôtes ; interfaces du routeur
(« Port Status : On »).}
\bareme{Sur le .pkt : interfaces routeur 1 pt ; IP/masques des hôtes 1 pt ; passerelles (imprimante et serveur compris)
1 pt. $-0,25$ par erreur, sans descendre sous 0 par item.}

\question{B4}{3 pts}{Réalisez les tests, notez la commande et le résultat, puis faites viser chaque test réussi.}
{\small\renewcommand{\arraystretch}{1.1}
\begin{tabularx}{\linewidth}{|c|l|l|C|c|c|}
  \hline
  \textbf{N°} & \textbf{Source} & \textbf{Destination} & \textbf{Commande / action} & \textbf{Résultat} & \textbf{Visa prof.} \\ \hline
  \hl 1 & PC-\pUn1 & PC-\pUn4 & \sol{ping \netUn.4} & \sol{réussi} & \hspace{1.4cm} \\ \hline
  \hl 2 & PC-\pDeux1 & IMP-\pDeux & \sol{ping \netDeux.0.10} & \sol{réussi} & \\ \hline
  \hl 3 & PC-\pUn2 & PC-\pDeux2 & \sol{ping \netDeux.0.2} & \sol{réussi} & \\ \hline
  \hl 4 & PC-\pDeux2 & \srvNom & \sol{ping \netSrv.\srvH} & \sol{réussi} & \\ \hline
  \hl 5 & PC-\pUn3 & page web du serveur & \sol{navigateur : {\NoAutoSpacing http://\netSrv.\srvH}} & \sol{page affichée} & \\ \hline
\end{tabularx}}
\bareme{0,5 pt par test visé + 0,5 pt pour B5. Le premier ping inter-réseaux peut échouer (résolution ARP) ; les suivants doivent réussir.}

\question{B5}{incl. B4}{Test n°3 : par quels équipements passe le paquet ? Pourquoi PC-\pUn2 l'envoie-t-il au routeur ?}
\zone{2.3cm}{PC-\pUn2 $\to$ SW-\pUn{} $\to$ \Rnom{} $\to$ SW-\pDeux{} $\to$ PC-\pDeux2. Grâce à son masque, PC-\pUn2 constate que
la destination (\netDeux.0.0/16) n'est \textbf{pas dans son réseau} (\netUn.0/24) : il envoie le paquet à sa
\textbf{passerelle par défaut}, qui le route.}

%==========================================================================
\newpage
\section{Partie C — Diagnostic d'un réseau en panne \hfill /4}

Un technicien constate qu'\textbf{aucun poste ne parvient à joindre le serveur SRV-C} et que certains postes
communiquent mal entre eux. Les câbles sont du bon type.

\begin{center}
\begin{minipage}{0.37\linewidth}\centering
\begin{tikzpicture}[font=\footnotesize,
  eq/.style={draw,fill=qbg,rounded corners=2pt,minimum width=9mm,minimum height=5mm,inner sep=2pt},
  every path/.style={thick}]
  \node[eq] (SW) {SW-1};
  \node[eq,right=9mm of SW] (R) {R1};
  \node[eq,right=7mm of R] (SRV) {SRV-C};
  \node[eq,above left=5mm and -3mm of SW] (A) {PC-A};
  \node[eq,above right=5mm and -3mm of SW] (B) {PC-B};
  \node[eq,below left=5mm and -3mm of SW] (C) {PC-C};
  \node[eq,below right=5mm and -3mm of SW] (D) {PC-D};
  \foreach \n in {A,B,C,D} \draw (\n) -- (SW);
  \draw (SW) -- node[above,font=\tiny]{G0/0} (R);
  \draw (R) -- node[above,font=\tiny]{G0/1} (SRV);
\end{tikzpicture}
\end{minipage}\hfill
\begin{minipage}{0.61\linewidth}\centering
{\small\renewcommand{\arraystretch}{1.05}\setlength{\tabcolsep}{4pt}
\begin{tabular}{lllll}
  \toprule
  \textbf{Équipement} & \textbf{Adresse IP} & \textbf{Masque} & \textbf{Passerelle} & \textbf{Port} \\
  \midrule
  \Crows
  \bottomrule
\end{tabular}}
\end{minipage}
\end{center}

\question{C1}{4 pts}{Relevez les \textbf{quatre anomalies}. Pour chacune : conséquence et correction.}
{\small\renewcommand{\arraystretch}{1.1}
\ifcorrige\def\hc{}\else\def\hc{\rule[-10mm]{0pt}{14mm}}\fi
\begin{tabularx}{\linewidth}{|c|>{\raggedright\arraybackslash}p{2cm}|L|L|L|}
  \hline
  \textbf{N°} & \textbf{Équipement} & \textbf{Anomalie} & \textbf{Conséquence} & \textbf{Correction} \\ \hline
  \ifcorrige\Csol\else
  \hc 1 & & & & \\ \hline
  \hc 2 & & & & \\ \hline
  \hc 3 & & & & \\ \hline
  \hc 4 & & & & \\ \hline
  \fi
\end{tabularx}}
\bareme{1 pt par anomalie (0,25 équipement, anomalie, conséquence, correction ; ordre indifférent). \Cbareme}

%==========================================================================
\section{Bonus — Extension sans fil \hfill +1}

\begin{cdc}Ajoutez au réseau \resUn{} un \textbf{AccessPoint-PT} relié à SW-\pUn{} (SSID \ip{\ssid}, \textbf{WPA2-PSK})
et un portable \ip{Laptop-\pUn} équipé d'un module \texttt{WPC300N}, en adressage \textbf{statique}.
Vérifiez qu'il joint \srvNom.\end{cdc}
{\small\renewcommand{\arraystretch}{1.1}
\begin{tabularx}{\linewidth}{|l|C|C|C|c|}
  \hline
  \textbf{Équipement} & \textbf{Adresse IP} & \textbf{Masque} & \textbf{Passerelle} & \textbf{Visa prof.} \\ \hline
  \hl Laptop-\pUn & \sol{\netUn.50} & \sol{255.255.255.0} & \sol{\netUn.254} & \hspace{1.4cm} \\ \hline
\end{tabularx}}
\bareme{Le point d'accès fait office de pont : le portable est dans \netUn.0/24. Bonus si associé en WPA2-PSK et ping
vers \netSrv.\srvH{} réussi. Note plafonnée à 20.}

\end{document}
"
\documentclass[10pt,a4paper]{article}

\usepackage[T1]{fontenc}
\usepackage[utf8]{inputenc}
\usepackage{lmodern}
\usepackage[french]{babel}
\usepackage[a4paper,hmargin=1.5cm,top=1.3cm,bottom=1.6cm,footskip=0.7cm]{geometry}
\usepackage[table]{xcolor}
\usepackage{booktabs,tabularx,array}
\usepackage{enumitem}
\usepackage{fancyhdr}
\usepackage[most]{tcolorbox}
\usepackage{tikz}
\usetikzlibrary{positioning}
\usepackage{titlesec}

\providecommand{\modecorrige}{0}
\providecommand{\variante}{A}
\newif\ifcorrige
\ifnum\modecorrige=1 \corrigetrue \else \corrigefalse \fi

\definecolor{ipred}{HTML}{C0392B}
\definecolor{cdcbar}{HTML}{A67C3D}
\definecolor{cdcbg}{HTML}{F4F1EA}
\definecolor{qbar}{HTML}{4A6B5D}
\definecolor{qbg}{HTML}{EDF1EF}
\definecolor{corr}{HTML}{1F5FA8}
\definecolor{corrbg}{HTML}{E8F0FA}
\definecolor{grisfoot}{HTML}{828282}

\newcommand{\ip}[1]{\textcolor{ipred}{\texttt{#1}}}
\newcommand{\pts}[1]{\hfill\mbox{\textbf{(#1)}}}
\newcommand{\sol}[1]{\ifcorrige\textcolor{corr}{#1}\fi}
\newcommand{\bareme}[1]{\ifcorrige\par{\footnotesize\color{corr}#1\par}\fi}
\newcommand{\hl}{\ifcorrige\rule[-1.6mm]{0pt}{5.4mm}\else\rule[-2.2mm]{0pt}{6.6mm}\fi}
\newcolumntype{C}{>{\centering\arraybackslash}X}
\newcolumntype{L}{>{\raggedright\arraybackslash}X}

\titleformat{\section}{\large\bfseries}{}{0pt}{}
\titlespacing*{\section}{0pt}{6pt}{3pt}

\tcbset{qstyle/.style={enhanced,boxrule=0pt,leftrule=2.5pt,arc=0pt,outer arc=0pt,
  left=5pt,right=5pt,top=2pt,bottom=2pt,before skip=4pt,after skip=3pt}}
\newtcolorbox{cdc}{qstyle,colback=cdcbg,colframe=cdcbar}
\newtcolorbox{quest}{qstyle,colback=qbg,colframe=qbar}

% Zone de réponse : vierge dans le sujet, remplie dans le corrigé
\newcommand{\zone}[2]{%
  \ifcorrige
    \begin{tcolorbox}[colback=corrbg,colframe=corr,boxrule=0.5pt,arc=1pt,left=4pt,right=4pt,top=1pt,bottom=1pt,
      before skip=2pt,after skip=2pt,fontupper=\footnotesize\color{corr}]#2\end{tcolorbox}%
  \else
    \begin{tcolorbox}[colback=white,colframe=black!45,boxrule=0.4pt,arc=1pt,height=#1,
      before skip=2pt,after skip=2pt]\end{tcolorbox}%
  \fi}

\newcommand{\question}[3]{\begin{quest}\textbf{#1.} #3\pts{#2}\end{quest}}

%==========================================================================
% Données propres à chaque variante
%==========================================================================
\if\variante A
  \def\entreprise{Le cabinet d'architectes \textbf{ArchiNova}}
  \def\Rnom{R-ArchiNova}
  \def\resUn{Bureau d'études}\def\pUn{BE}\def\netUn{192.168.40}
  \def\resDeux{Accueil}\def\pDeux{ACC}\def\netDeux{172.16}
  \def\srvNom{SRV-Intranet}\def\netSrv{192.168.200}\def\srvH{10}
  \def\ssid{ArchiNova-WiFi}
  \def\AIrows{%
    \hl\ip{192.168.40.25} & \sol{C} & \sol{255.255.255.0} & \sol{192.168.40.0} & \sol{192.168.40.255} & \sol{254} \\ \hline
    \hl\ip{172.16.5.10}   & \sol{B} & \sol{255.255.0.0}   & \sol{172.16.0.0}   & \sol{172.16.255.255} & \sol{65\,534} \\ \hline
    \hl\ip{10.3.2.1}      & \sol{A} & \sol{255.0.0.0}     & \sol{10.0.0.0}     & \sol{10.255.255.255} & \sol{16\,777\,214} \\ \hline}
  \def\apipa{169.254.12.7}
  \def\Cnet{192.168.60}
  \def\Crows{%
    PC-A & \ip{192.168.60.10} & \ip{255.255.255.0} & \ip{192.168.60.254} & — \\
    PC-B & \ip{192.168.60.20} & \ip{255.255.255.0} & \ip{192.168.60.1}   & — \\
    PC-C & \ip{192.168.61.30} & \ip{255.255.255.0} & \ip{192.168.60.254} & — \\
    PC-D & \ip{192.168.60.10} & \ip{255.255.255.0} & \ip{192.168.60.254} & — \\
    SRV-C & \ip{10.0.0.5}     & \ip{255.0.0.0}     & \ip{10.0.0.254}     & — \\
    R1 — G0/0 & \ip{192.168.60.254} & \ip{255.255.255.0} & — & On \\
    R1 — G0/1 & \ip{10.0.0.254}     & \ip{255.0.0.0}     & — & Off \\}
  \def\Csol{%
    \hl 1 & \sol{PC-B} & \sol{Passerelle .60.1 : pas l'adresse du routeur.} & \sol{Joint son réseau, mais aucun autre (serveur).} & \sol{Passerelle 192.168.60.254.} \\ \hline
    \hl 2 & \sol{PC-C} & \sol{IP hors du réseau 192.168.60.0/24.} & \sol{Ne joint ni les postes ni le routeur.} & \sol{Ex. 192.168.60.30.} \\ \hline
    \hl 3 & \sol{PC-D} & \sol{Même IP que PC-A (doublon).} & \sol{Conflit d'adresse : PC-A et PC-D perturbés.} & \sol{Adresse unique, ex. .60.40.} \\ \hline
    \hl 4 & \sol{R1 G0/1} & \sol{Interface désactivée (Off).} & \sol{Réseau du serveur injoignable pour tous.} & \sol{Port Status : On.} \\ \hline}
  \def\Cbareme{PC-A et SRV-C sont correctement configurés ; accepter PC-A cité à la place de PC-D pour le doublon.}
\else
  \def\entreprise{La clinique vétérinaire \textbf{VetoPlus}}
  \def\Rnom{R-VetoPlus}
  \def\resUn{Soins}\def\pUn{SOI}\def\netUn{192.168.70}
  \def\resDeux{Secrétariat}\def\pDeux{SEC}\def\netDeux{172.20}
  \def\srvNom{SRV-Dossiers}\def\netSrv{192.168.250}\def\srvH{20}
  \def\ssid{VetoPlus-WiFi}
  \def\AIrows{%
    \hl\ip{192.168.15.130} & \sol{C} & \sol{255.255.255.0} & \sol{192.168.15.0} & \sol{192.168.15.255} & \sol{254} \\ \hline
    \hl\ip{172.31.200.7}   & \sol{B} & \sol{255.255.0.0}   & \sol{172.31.0.0}   & \sol{172.31.255.255} & \sol{65\,534} \\ \hline
    \hl\ip{10.45.6.9}      & \sol{A} & \sol{255.0.0.0}     & \sol{10.0.0.0}     & \sol{10.255.255.255} & \sol{16\,777\,214} \\ \hline}
  \def\apipa{169.254.88.3}
  \def\Cnet{192.168.80}
  \def\Crows{%
    PC-A & \ip{192.168.8.11}   & \ip{255.255.255.0} & \ip{192.168.80.254} & — \\
    PC-B & \ip{192.168.80.12}  & \ip{255.255.255.0} & \ip{192.168.80.254} & — \\
    PC-C & \ip{192.168.80.13}  & \ip{255.255.255.0} & \ip{192.168.80.100} & — \\
    PC-D & \ip{192.168.80.12}  & \ip{255.255.255.0} & \ip{192.168.80.254} & — \\
    SRV-C & \ip{10.0.0.8}      & \ip{255.0.0.0}     & \ip{10.0.0.254}     & — \\
    R1 — G0/0 & \ip{192.168.80.254} & \ip{255.255.255.0} & — & Off \\
    R1 — G0/1 & \ip{10.0.0.254}     & \ip{255.0.0.0}     & — & On \\}
  \def\Csol{%
    \hl 1 & \sol{PC-A} & \sol{IP hors du réseau 192.168.80.0/24.} & \sol{Ne joint ni les postes ni le routeur.} & \sol{Ex. 192.168.80.11.} \\ \hline
    \hl 2 & \sol{PC-C} & \sol{Passerelle .80.100 : pas l'adresse du routeur.} & \sol{Joint son réseau, mais aucun autre (serveur).} & \sol{Passerelle 192.168.80.254.} \\ \hline
    \hl 3 & \sol{PC-D} & \sol{Même IP que PC-B (doublon).} & \sol{Conflit d'adresse : PC-B et PC-D perturbés.} & \sol{Adresse unique, ex. .80.14.} \\ \hline
    \hl 4 & \sol{R1 G0/0} & \sol{Interface côté LAN désactivée (Off).} & \sol{Aucun poste ne joint le routeur ni le serveur.} & \sol{Port Status : On.} \\ \hline}
  \def\Cbareme{PC-B et SRV-C sont correctement configurés ; accepter PC-B cité à la place de PC-D pour le doublon.}
\fi

\pagestyle{fancy}
\fancyhf{}
\renewcommand{\headrulewidth}{0pt}
\fancyfoot[L]{\footnotesize\itshape\color{grisfoot}Évaluation — Réseaux locaux — Sujet \variante\ifcorrige\ — CORRIGÉ\fi}
\fancyfoot[R]{\footnotesize\color{grisfoot}\thepage/3}

\setlist{itemsep=0pt,topsep=2pt,parsep=0pt}
\setlength{\parindent}{0pt}
\setlength{\parskip}{2pt}

\begin{document}

\begin{center}
  {\Large\bfseries Évaluation pratique — Réseaux locaux \quad
   \fcolorbox{black}{white}{\ SUJET \variante\ }\par}
  \vspace{1pt}
  {\bfseries Adressage IP, commutation et routage — BTS CIEL 1re année — durée : 2 heures\par}
  \ifcorrige\vspace{2pt}{\bfseries\color{corr}CORRIGÉ ET BARÈME — document enseignant}\par\fi
\end{center}
\vspace{-2pt}
{\renewcommand{\arraystretch}{1.35}
\begin{tabularx}{\linewidth}{|X|X|p{2.2cm}|p{2.6cm}|}
  \hline
  \textbf{Nom :} & \textbf{Prénom :} & \textbf{Poste n°} & \textbf{Note :} \hfill /20 \\ \hline
\end{tabularx}}

\textbf{Consignes :} travail \textbf{individuel}, documents et Internet \textbf{non autorisés}. Ce sujet est le
\textbf{document-réponse}. Fichier Packet Tracer : \ip{Nom\_Prenom\_EvalLAN\_\variante.pkt}, à enregistrer
régulièrement dans le dossier de rendu. Les « visas » sont apposés par le professeur, que vous appelez à chaque test réussi.
Durées conseillées : A 30 min — B 1 h 10 — C 20 min.

%==========================================================================
\section{Partie A — Questions de cours et calculs \hfill /6}

\question{A1}{2 pts}{Complétez le tableau en considérant le \textbf{masque par défaut} de la classe de chaque adresse.}
{\small\renewcommand{\arraystretch}{1.15}
\begin{tabularx}{\linewidth}{|l|c|C|C|C|c|}
  \hline
  \textbf{Adresse IP} & \textbf{Classe} & \textbf{Masque par défaut} & \textbf{Adresse réseau} &
  \textbf{Adresse de diffusion} & \textbf{Nb d'hôtes} \\ \hline
  \AIrows
\end{tabularx}}
\bareme{0,5 pt par ligne juste + 0,5 pt si les trois nombres d'hôtes sont justes ($2^{n}-2$, $n$ = nombre de bits d'hôte).}

\question{A2}{1 pt}{Indiquez le type de câble (\textbf{droit} ou \textbf{croisé}) pour chaque liaison.}
{\small\renewcommand{\arraystretch}{1.15}
\begin{tabularx}{\linewidth}{|X|c||X|c|}
  \hline
  \hl PC $\leftrightarrow$ commutateur & \makebox[1.6cm]{\sol{droit}} & Commutateur $\leftrightarrow$ routeur & \makebox[1.6cm]{\sol{droit}} \\ \hline
  \hl PC $\leftrightarrow$ PC (réseau P2P) & \sol{croisé} & Concentrateur $\leftrightarrow$ concentrateur & \sol{croisé} \\ \hline
\end{tabularx}}
\bareme{0,25 pt par case. Équipements de types différents : câble droit ; de même type : câble croisé.}

\question{A3}{1 pt}{Comment un \textbf{concentrateur (hub)} et un \textbf{commutateur (switch)} traitent-ils une trame
reçue ? Lequel limite les collisions, et pourquoi ?}
\zone{2.9cm}{Le \textbf{hub} (couche 1) répète la trame sur \textbf{tous ses autres ports} : un seul domaine de collision
partagé. Le \textbf{switch} (couche 2) lit l'adresse MAC de destination et, grâce à sa \textbf{table MAC}, ne l'envoie
\textbf{que sur le port du destinataire} : un domaine de collision par port, c'est lui qui limite les collisions.
\emph{0,5 pt hub, 0,5 pt switch.}}

\question{A4}{1 pt}{Qu'est-ce que la \textbf{passerelle par défaut} d'un poste ? Quelle condition son adresse doit-elle respecter ?}
\zone{2.4cm}{Adresse IP de l'interface du \textbf{routeur} à laquelle le poste envoie les paquets destinés à un
\textbf{autre réseau}. Elle doit \textbf{appartenir au même réseau} que le poste. \emph{0,5 + 0,5 pt.}}

\question{A5}{0,5 pt}{Un portable configuré en DHCP affiche \ip{\apipa} / \ip{255.255.0.0}. Que s'est-il passé ?}
\zone{1.5cm}{Aucun \textbf{serveur DHCP} n'a répondu : le poste s'est attribué une adresse \textbf{APIPA} (169.254.0.0/16).}

\question{A6}{0,5 pt}{Quel protocole utilise \texttt{ping} ? Nommez les deux messages échangés et la couche OSI des adresses IP.}
\zone{1.5cm}{\textbf{ICMP} ; \textbf{Echo Request} / \textbf{Echo Reply} ; adresses IP en \textbf{couche 3 (Réseau)}.}

%==========================================================================
\newpage
\section{Partie B — Réalisation sous Packet Tracer \hfill /10}

\begin{cdc}
\entreprise{} vous confie son réseau : \textbf{trois réseaux} reliés par un \textbf{routeur unique Cisco 2911}
nommé \ip{\Rnom} (interfaces G0/0, G0/1, G0/2).
\begin{itemize}[leftmargin=12pt]
  \item \textbf{Réseau \resUn} : \textbf{4 postes} PC-\pUn1 à PC-\pUn4, commutateur \ip{SW-\pUn}, adressage statique,
        réseau \ip{\netUn.0/24}.
  \item \textbf{Réseau \resDeux} : \textbf{2 postes} PC-\pDeux1, PC-\pDeux2 et \textbf{1 imprimante} IMP-\pDeux,
        commutateur \ip{SW-\pDeux}, adressage statique, réseau de classe B \ip{\netDeux.0.0/16}.
  \item \textbf{Réseau Serveur} : serveur \ip{\srvNom} d'adresse \ip{\netSrv.\srvH/24}, relié \textbf{directement} au routeur.
  \item La passerelle de chaque réseau est sa \textbf{dernière adresse utilisable}. Tous les équipements communiquent
        entre eux et accèdent à la page web du serveur. Nommage imposé (« Display Name »).
\end{itemize}
\end{cdc}

\question{B1}{2 pts}{Complétez le plan d'adressage \textbf{avant} la réalisation.}
{\footnotesize\setlength{\tabcolsep}{3pt}\renewcommand{\arraystretch}{1.15}
\begin{tabularx}{\linewidth}{|l|C|C|C|C|C|C|}
  \hline
  \textbf{Réseau} & \textbf{Adresse réseau} & \textbf{Masque} & \textbf{1re @ utilisable} &
  \textbf{Dernière @ utilisable} & \textbf{Diffusion} & \textbf{Passerelle} \\ \hline
  \hl \resUn & \sol{\netUn.0} & \sol{255.255.255.0} & \sol{\netUn.1} & \sol{\netUn.254} & \sol{\netUn.255} & \sol{\netUn.254} \\ \hline
  \hl \resDeux & \sol{\netDeux.0.0} & \sol{255.255.0.0} & \sol{\netDeux.0.1} & \sol{\netDeux.255.254} & \sol{\netDeux.255.255} & \sol{\netDeux.255.254} \\ \hline
  \hl Serveur & \sol{\netSrv.0} & \sol{255.255.255.0} & \sol{\netSrv.1} & \sol{\netSrv.254} & \sol{\netSrv.255} & \sol{\netSrv.254} \\ \hline
\end{tabularx}}
\bareme{0,5 pt par ligne juste + 0,5 pt si la convention de passerelle est respectée partout.}

\question{B2}{2 pts}{Construisez la topologie (équipements, câblage, nommage) et relevez l'adressage configuré.}
{\footnotesize\setlength{\tabcolsep}{3pt}\renewcommand{\arraystretch}{1.1}
\begin{tabularx}{\linewidth}{|l|C|C|C||l|C|C|}
  \hline
  \textbf{Hôte} & \textbf{Adresse IP} & \textbf{Masque} & \textbf{Passerelle} & \textbf{Interface} & \textbf{Adresse IP} & \textbf{Masque} \\ \hline
  \hl PC-\pUn1 & \sol{\netUn.1} & \sol{255.255.255.0} & \sol{\netUn.254} & \Rnom{} G0/0 & \sol{\netUn.254} & \sol{255.255.255.0} \\ \hline
  \hl PC-\pDeux1 & \sol{\netDeux.0.1} & \sol{255.255.0.0} & \sol{\netDeux.255.254} & \Rnom{} G0/1 & \sol{\netDeux.255.254} & \sol{255.255.0.0} \\ \hline
  \hl IMP-\pDeux & \sol{\netDeux.0.10} & \sol{255.255.0.0} & \sol{\netDeux.255.254} & \Rnom{} G0/2 & \sol{\netSrv.254} & \sol{255.255.255.0} \\ \hline
  \hl \srvNom & \ip{\netSrv.\srvH} & \sol{255.255.255.0} & \sol{\netSrv.254} & \multicolumn{3}{c|}{\cellcolor{black!8}} \\ \hline
\end{tabularx}}
\bareme{Adresses d'hôtes libres (uniques, dans le bon réseau) ; affectation G0/x libre. Sur le .pkt : équipements 0,5 ;
câblage 1 (câbles droits vers les switchs et switch–routeur, \textbf{croisé} serveur–routeur, liens verts) ; nommage 0,5.}

\question{B3}{3 pts}{Configurez tous les équipements : IP, masques, passerelles des hôtes ; interfaces du routeur
(« Port Status : On »).}
\bareme{Sur le .pkt : interfaces routeur 1 pt ; IP/masques des hôtes 1 pt ; passerelles (imprimante et serveur compris)
1 pt. $-0,25$ par erreur, sans descendre sous 0 par item.}

\question{B4}{3 pts}{Réalisez les tests, notez la commande et le résultat, puis faites viser chaque test réussi.}
{\small\renewcommand{\arraystretch}{1.1}
\begin{tabularx}{\linewidth}{|c|l|l|C|c|c|}
  \hline
  \textbf{N°} & \textbf{Source} & \textbf{Destination} & \textbf{Commande / action} & \textbf{Résultat} & \textbf{Visa prof.} \\ \hline
  \hl 1 & PC-\pUn1 & PC-\pUn4 & \sol{ping \netUn.4} & \sol{réussi} & \hspace{1.4cm} \\ \hline
  \hl 2 & PC-\pDeux1 & IMP-\pDeux & \sol{ping \netDeux.0.10} & \sol{réussi} & \\ \hline
  \hl 3 & PC-\pUn2 & PC-\pDeux2 & \sol{ping \netDeux.0.2} & \sol{réussi} & \\ \hline
  \hl 4 & PC-\pDeux2 & \srvNom & \sol{ping \netSrv.\srvH} & \sol{réussi} & \\ \hline
  \hl 5 & PC-\pUn3 & page web du serveur & \sol{navigateur : {\NoAutoSpacing http://\netSrv.\srvH}} & \sol{page affichée} & \\ \hline
\end{tabularx}}
\bareme{0,5 pt par test visé + 0,5 pt pour B5. Le premier ping inter-réseaux peut échouer (résolution ARP) ; les suivants doivent réussir.}

\question{B5}{incl. B4}{Test n°3 : par quels équipements passe le paquet ? Pourquoi PC-\pUn2 l'envoie-t-il au routeur ?}
\zone{2.3cm}{PC-\pUn2 $\to$ SW-\pUn{} $\to$ \Rnom{} $\to$ SW-\pDeux{} $\to$ PC-\pDeux2. Grâce à son masque, PC-\pUn2 constate que
la destination (\netDeux.0.0/16) n'est \textbf{pas dans son réseau} (\netUn.0/24) : il envoie le paquet à sa
\textbf{passerelle par défaut}, qui le route.}

%==========================================================================
\newpage
\section{Partie C — Diagnostic d'un réseau en panne \hfill /4}

Un technicien constate qu'\textbf{aucun poste ne parvient à joindre le serveur SRV-C} et que certains postes
communiquent mal entre eux. Les câbles sont du bon type.

\begin{center}
\begin{minipage}{0.37\linewidth}\centering
\begin{tikzpicture}[font=\footnotesize,
  eq/.style={draw,fill=qbg,rounded corners=2pt,minimum width=9mm,minimum height=5mm,inner sep=2pt},
  every path/.style={thick}]
  \node[eq] (SW) {SW-1};
  \node[eq,right=9mm of SW] (R) {R1};
  \node[eq,right=7mm of R] (SRV) {SRV-C};
  \node[eq,above left=5mm and -3mm of SW] (A) {PC-A};
  \node[eq,above right=5mm and -3mm of SW] (B) {PC-B};
  \node[eq,below left=5mm and -3mm of SW] (C) {PC-C};
  \node[eq,below right=5mm and -3mm of SW] (D) {PC-D};
  \foreach \n in {A,B,C,D} \draw (\n) -- (SW);
  \draw (SW) -- node[above,font=\tiny]{G0/0} (R);
  \draw (R) -- node[above,font=\tiny]{G0/1} (SRV);
\end{tikzpicture}
\end{minipage}\hfill
\begin{minipage}{0.61\linewidth}\centering
{\small\renewcommand{\arraystretch}{1.05}\setlength{\tabcolsep}{4pt}
\begin{tabular}{lllll}
  \toprule
  \textbf{Équipement} & \textbf{Adresse IP} & \textbf{Masque} & \textbf{Passerelle} & \textbf{Port} \\
  \midrule
  \Crows
  \bottomrule
\end{tabular}}
\end{minipage}
\end{center}

\question{C1}{4 pts}{Relevez les \textbf{quatre anomalies}. Pour chacune : conséquence et correction.}
{\small\renewcommand{\arraystretch}{1.1}
\ifcorrige\def\hc{}\else\def\hc{\rule[-10mm]{0pt}{14mm}}\fi
\begin{tabularx}{\linewidth}{|c|>{\raggedright\arraybackslash}p{2cm}|L|L|L|}
  \hline
  \textbf{N°} & \textbf{Équipement} & \textbf{Anomalie} & \textbf{Conséquence} & \textbf{Correction} \\ \hline
  \ifcorrige\Csol\else
  \hc 1 & & & & \\ \hline
  \hc 2 & & & & \\ \hline
  \hc 3 & & & & \\ \hline
  \hc 4 & & & & \\ \hline
  \fi
\end{tabularx}}
\bareme{1 pt par anomalie (0,25 équipement, anomalie, conséquence, correction ; ordre indifférent). \Cbareme}

%==========================================================================
\section{Bonus — Extension sans fil \hfill +1}

\begin{cdc}Ajoutez au réseau \resUn{} un \textbf{AccessPoint-PT} relié à SW-\pUn{} (SSID \ip{\ssid}, \textbf{WPA2-PSK})
et un portable \ip{Laptop-\pUn} équipé d'un module \texttt{WPC300N}, en adressage \textbf{statique}.
Vérifiez qu'il joint \srvNom.\end{cdc}
{\small\renewcommand{\arraystretch}{1.1}
\begin{tabularx}{\linewidth}{|l|C|C|C|c|}
  \hline
  \textbf{Équipement} & \textbf{Adresse IP} & \textbf{Masque} & \textbf{Passerelle} & \textbf{Visa prof.} \\ \hline
  \hl Laptop-\pUn & \sol{\netUn.50} & \sol{255.255.255.0} & \sol{\netUn.254} & \hspace{1.4cm} \\ \hline
\end{tabularx}}
\bareme{Le point d'accès fait office de pont : le portable est dans \netUn.0/24. Bonus si associé en WPA2-PSK et ping
vers \netSrv.\srvH{} réussi. Note plafonnée à 20.}

\end{document}
"
%   pdflatex -jobname=Eval_Reseau_SujetB   "\def\variante{B}\def\modecorrige{0}% Évaluation pratique — Réseaux locaux (BTS CIEL 1re année) — sujets A et B
% Compilation (4 PDF) :
%   pdflatex -jobname=Eval_Reseau_SujetA   "\def\variante{A}\def\modecorrige{0}% Évaluation pratique — Réseaux locaux (BTS CIEL 1re année) — sujets A et B
% Compilation (4 PDF) :
%   pdflatex -jobname=Eval_Reseau_SujetA   "\def\variante{A}\def\modecorrige{0}\input{Eval_TP_Reseau_LAN.tex}"
%   pdflatex -jobname=Eval_Reseau_SujetB   "\def\variante{B}\def\modecorrige{0}\input{Eval_TP_Reseau_LAN.tex}"
%   pdflatex -jobname=Eval_Reseau_CorrigeA "\def\variante{A}\def\modecorrige{1}\input{Eval_TP_Reseau_LAN.tex}"
%   pdflatex -jobname=Eval_Reseau_CorrigeB "\def\variante{B}\def\modecorrige{1}\input{Eval_TP_Reseau_LAN.tex}"
\documentclass[10pt,a4paper]{article}

\usepackage[T1]{fontenc}
\usepackage[utf8]{inputenc}
\usepackage{lmodern}
\usepackage[french]{babel}
\usepackage[a4paper,hmargin=1.5cm,top=1.3cm,bottom=1.6cm,footskip=0.7cm]{geometry}
\usepackage[table]{xcolor}
\usepackage{booktabs,tabularx,array}
\usepackage{enumitem}
\usepackage{fancyhdr}
\usepackage[most]{tcolorbox}
\usepackage{tikz}
\usetikzlibrary{positioning}
\usepackage{titlesec}

\providecommand{\modecorrige}{0}
\providecommand{\variante}{A}
\newif\ifcorrige
\ifnum\modecorrige=1 \corrigetrue \else \corrigefalse \fi

\definecolor{ipred}{HTML}{C0392B}
\definecolor{cdcbar}{HTML}{A67C3D}
\definecolor{cdcbg}{HTML}{F4F1EA}
\definecolor{qbar}{HTML}{4A6B5D}
\definecolor{qbg}{HTML}{EDF1EF}
\definecolor{corr}{HTML}{1F5FA8}
\definecolor{corrbg}{HTML}{E8F0FA}
\definecolor{grisfoot}{HTML}{828282}

\newcommand{\ip}[1]{\textcolor{ipred}{\texttt{#1}}}
\newcommand{\pts}[1]{\hfill\mbox{\textbf{(#1)}}}
\newcommand{\sol}[1]{\ifcorrige\textcolor{corr}{#1}\fi}
\newcommand{\bareme}[1]{\ifcorrige\par{\footnotesize\color{corr}#1\par}\fi}
\newcommand{\hl}{\ifcorrige\rule[-1.6mm]{0pt}{5.4mm}\else\rule[-2.2mm]{0pt}{6.6mm}\fi}
\newcolumntype{C}{>{\centering\arraybackslash}X}
\newcolumntype{L}{>{\raggedright\arraybackslash}X}

\titleformat{\section}{\large\bfseries}{}{0pt}{}
\titlespacing*{\section}{0pt}{6pt}{3pt}

\tcbset{qstyle/.style={enhanced,boxrule=0pt,leftrule=2.5pt,arc=0pt,outer arc=0pt,
  left=5pt,right=5pt,top=2pt,bottom=2pt,before skip=4pt,after skip=3pt}}
\newtcolorbox{cdc}{qstyle,colback=cdcbg,colframe=cdcbar}
\newtcolorbox{quest}{qstyle,colback=qbg,colframe=qbar}

% Zone de réponse : vierge dans le sujet, remplie dans le corrigé
\newcommand{\zone}[2]{%
  \ifcorrige
    \begin{tcolorbox}[colback=corrbg,colframe=corr,boxrule=0.5pt,arc=1pt,left=4pt,right=4pt,top=1pt,bottom=1pt,
      before skip=2pt,after skip=2pt,fontupper=\footnotesize\color{corr}]#2\end{tcolorbox}%
  \else
    \begin{tcolorbox}[colback=white,colframe=black!45,boxrule=0.4pt,arc=1pt,height=#1,
      before skip=2pt,after skip=2pt]\end{tcolorbox}%
  \fi}

\newcommand{\question}[3]{\begin{quest}\textbf{#1.} #3\pts{#2}\end{quest}}

%==========================================================================
% Données propres à chaque variante
%==========================================================================
\if\variante A
  \def\entreprise{Le cabinet d'architectes \textbf{ArchiNova}}
  \def\Rnom{R-ArchiNova}
  \def\resUn{Bureau d'études}\def\pUn{BE}\def\netUn{192.168.40}
  \def\resDeux{Accueil}\def\pDeux{ACC}\def\netDeux{172.16}
  \def\srvNom{SRV-Intranet}\def\netSrv{192.168.200}\def\srvH{10}
  \def\ssid{ArchiNova-WiFi}
  \def\AIrows{%
    \hl\ip{192.168.40.25} & \sol{C} & \sol{255.255.255.0} & \sol{192.168.40.0} & \sol{192.168.40.255} & \sol{254} \\ \hline
    \hl\ip{172.16.5.10}   & \sol{B} & \sol{255.255.0.0}   & \sol{172.16.0.0}   & \sol{172.16.255.255} & \sol{65\,534} \\ \hline
    \hl\ip{10.3.2.1}      & \sol{A} & \sol{255.0.0.0}     & \sol{10.0.0.0}     & \sol{10.255.255.255} & \sol{16\,777\,214} \\ \hline}
  \def\apipa{169.254.12.7}
  \def\Cnet{192.168.60}
  \def\Crows{%
    PC-A & \ip{192.168.60.10} & \ip{255.255.255.0} & \ip{192.168.60.254} & — \\
    PC-B & \ip{192.168.60.20} & \ip{255.255.255.0} & \ip{192.168.60.1}   & — \\
    PC-C & \ip{192.168.61.30} & \ip{255.255.255.0} & \ip{192.168.60.254} & — \\
    PC-D & \ip{192.168.60.10} & \ip{255.255.255.0} & \ip{192.168.60.254} & — \\
    SRV-C & \ip{10.0.0.5}     & \ip{255.0.0.0}     & \ip{10.0.0.254}     & — \\
    R1 — G0/0 & \ip{192.168.60.254} & \ip{255.255.255.0} & — & On \\
    R1 — G0/1 & \ip{10.0.0.254}     & \ip{255.0.0.0}     & — & Off \\}
  \def\Csol{%
    \hl 1 & \sol{PC-B} & \sol{Passerelle .60.1 : pas l'adresse du routeur.} & \sol{Joint son réseau, mais aucun autre (serveur).} & \sol{Passerelle 192.168.60.254.} \\ \hline
    \hl 2 & \sol{PC-C} & \sol{IP hors du réseau 192.168.60.0/24.} & \sol{Ne joint ni les postes ni le routeur.} & \sol{Ex. 192.168.60.30.} \\ \hline
    \hl 3 & \sol{PC-D} & \sol{Même IP que PC-A (doublon).} & \sol{Conflit d'adresse : PC-A et PC-D perturbés.} & \sol{Adresse unique, ex. .60.40.} \\ \hline
    \hl 4 & \sol{R1 G0/1} & \sol{Interface désactivée (Off).} & \sol{Réseau du serveur injoignable pour tous.} & \sol{Port Status : On.} \\ \hline}
  \def\Cbareme{PC-A et SRV-C sont correctement configurés ; accepter PC-A cité à la place de PC-D pour le doublon.}
\else
  \def\entreprise{La clinique vétérinaire \textbf{VetoPlus}}
  \def\Rnom{R-VetoPlus}
  \def\resUn{Soins}\def\pUn{SOI}\def\netUn{192.168.70}
  \def\resDeux{Secrétariat}\def\pDeux{SEC}\def\netDeux{172.20}
  \def\srvNom{SRV-Dossiers}\def\netSrv{192.168.250}\def\srvH{20}
  \def\ssid{VetoPlus-WiFi}
  \def\AIrows{%
    \hl\ip{192.168.15.130} & \sol{C} & \sol{255.255.255.0} & \sol{192.168.15.0} & \sol{192.168.15.255} & \sol{254} \\ \hline
    \hl\ip{172.31.200.7}   & \sol{B} & \sol{255.255.0.0}   & \sol{172.31.0.0}   & \sol{172.31.255.255} & \sol{65\,534} \\ \hline
    \hl\ip{10.45.6.9}      & \sol{A} & \sol{255.0.0.0}     & \sol{10.0.0.0}     & \sol{10.255.255.255} & \sol{16\,777\,214} \\ \hline}
  \def\apipa{169.254.88.3}
  \def\Cnet{192.168.80}
  \def\Crows{%
    PC-A & \ip{192.168.8.11}   & \ip{255.255.255.0} & \ip{192.168.80.254} & — \\
    PC-B & \ip{192.168.80.12}  & \ip{255.255.255.0} & \ip{192.168.80.254} & — \\
    PC-C & \ip{192.168.80.13}  & \ip{255.255.255.0} & \ip{192.168.80.100} & — \\
    PC-D & \ip{192.168.80.12}  & \ip{255.255.255.0} & \ip{192.168.80.254} & — \\
    SRV-C & \ip{10.0.0.8}      & \ip{255.0.0.0}     & \ip{10.0.0.254}     & — \\
    R1 — G0/0 & \ip{192.168.80.254} & \ip{255.255.255.0} & — & Off \\
    R1 — G0/1 & \ip{10.0.0.254}     & \ip{255.0.0.0}     & — & On \\}
  \def\Csol{%
    \hl 1 & \sol{PC-A} & \sol{IP hors du réseau 192.168.80.0/24.} & \sol{Ne joint ni les postes ni le routeur.} & \sol{Ex. 192.168.80.11.} \\ \hline
    \hl 2 & \sol{PC-C} & \sol{Passerelle .80.100 : pas l'adresse du routeur.} & \sol{Joint son réseau, mais aucun autre (serveur).} & \sol{Passerelle 192.168.80.254.} \\ \hline
    \hl 3 & \sol{PC-D} & \sol{Même IP que PC-B (doublon).} & \sol{Conflit d'adresse : PC-B et PC-D perturbés.} & \sol{Adresse unique, ex. .80.14.} \\ \hline
    \hl 4 & \sol{R1 G0/0} & \sol{Interface côté LAN désactivée (Off).} & \sol{Aucun poste ne joint le routeur ni le serveur.} & \sol{Port Status : On.} \\ \hline}
  \def\Cbareme{PC-B et SRV-C sont correctement configurés ; accepter PC-B cité à la place de PC-D pour le doublon.}
\fi

\pagestyle{fancy}
\fancyhf{}
\renewcommand{\headrulewidth}{0pt}
\fancyfoot[L]{\footnotesize\itshape\color{grisfoot}Évaluation — Réseaux locaux — Sujet \variante\ifcorrige\ — CORRIGÉ\fi}
\fancyfoot[R]{\footnotesize\color{grisfoot}\thepage/3}

\setlist{itemsep=0pt,topsep=2pt,parsep=0pt}
\setlength{\parindent}{0pt}
\setlength{\parskip}{2pt}

\begin{document}

\begin{center}
  {\Large\bfseries Évaluation pratique — Réseaux locaux \quad
   \fcolorbox{black}{white}{\ SUJET \variante\ }\par}
  \vspace{1pt}
  {\bfseries Adressage IP, commutation et routage — BTS CIEL 1re année — durée : 2 heures\par}
  \ifcorrige\vspace{2pt}{\bfseries\color{corr}CORRIGÉ ET BARÈME — document enseignant}\par\fi
\end{center}
\vspace{-2pt}
{\renewcommand{\arraystretch}{1.35}
\begin{tabularx}{\linewidth}{|X|X|p{2.2cm}|p{2.6cm}|}
  \hline
  \textbf{Nom :} & \textbf{Prénom :} & \textbf{Poste n°} & \textbf{Note :} \hfill /20 \\ \hline
\end{tabularx}}

\textbf{Consignes :} travail \textbf{individuel}, documents et Internet \textbf{non autorisés}. Ce sujet est le
\textbf{document-réponse}. Fichier Packet Tracer : \ip{Nom\_Prenom\_EvalLAN\_\variante.pkt}, à enregistrer
régulièrement dans le dossier de rendu. Les « visas » sont apposés par le professeur, que vous appelez à chaque test réussi.
Durées conseillées : A 30 min — B 1 h 10 — C 20 min.

%==========================================================================
\section{Partie A — Questions de cours et calculs \hfill /6}

\question{A1}{2 pts}{Complétez le tableau en considérant le \textbf{masque par défaut} de la classe de chaque adresse.}
{\small\renewcommand{\arraystretch}{1.15}
\begin{tabularx}{\linewidth}{|l|c|C|C|C|c|}
  \hline
  \textbf{Adresse IP} & \textbf{Classe} & \textbf{Masque par défaut} & \textbf{Adresse réseau} &
  \textbf{Adresse de diffusion} & \textbf{Nb d'hôtes} \\ \hline
  \AIrows
\end{tabularx}}
\bareme{0,5 pt par ligne juste + 0,5 pt si les trois nombres d'hôtes sont justes ($2^{n}-2$, $n$ = nombre de bits d'hôte).}

\question{A2}{1 pt}{Indiquez le type de câble (\textbf{droit} ou \textbf{croisé}) pour chaque liaison.}
{\small\renewcommand{\arraystretch}{1.15}
\begin{tabularx}{\linewidth}{|X|c||X|c|}
  \hline
  \hl PC $\leftrightarrow$ commutateur & \makebox[1.6cm]{\sol{droit}} & Commutateur $\leftrightarrow$ routeur & \makebox[1.6cm]{\sol{droit}} \\ \hline
  \hl PC $\leftrightarrow$ PC (réseau P2P) & \sol{croisé} & Concentrateur $\leftrightarrow$ concentrateur & \sol{croisé} \\ \hline
\end{tabularx}}
\bareme{0,25 pt par case. Équipements de types différents : câble droit ; de même type : câble croisé.}

\question{A3}{1 pt}{Comment un \textbf{concentrateur (hub)} et un \textbf{commutateur (switch)} traitent-ils une trame
reçue ? Lequel limite les collisions, et pourquoi ?}
\zone{2.9cm}{Le \textbf{hub} (couche 1) répète la trame sur \textbf{tous ses autres ports} : un seul domaine de collision
partagé. Le \textbf{switch} (couche 2) lit l'adresse MAC de destination et, grâce à sa \textbf{table MAC}, ne l'envoie
\textbf{que sur le port du destinataire} : un domaine de collision par port, c'est lui qui limite les collisions.
\emph{0,5 pt hub, 0,5 pt switch.}}

\question{A4}{1 pt}{Qu'est-ce que la \textbf{passerelle par défaut} d'un poste ? Quelle condition son adresse doit-elle respecter ?}
\zone{2.4cm}{Adresse IP de l'interface du \textbf{routeur} à laquelle le poste envoie les paquets destinés à un
\textbf{autre réseau}. Elle doit \textbf{appartenir au même réseau} que le poste. \emph{0,5 + 0,5 pt.}}

\question{A5}{0,5 pt}{Un portable configuré en DHCP affiche \ip{\apipa} / \ip{255.255.0.0}. Que s'est-il passé ?}
\zone{1.5cm}{Aucun \textbf{serveur DHCP} n'a répondu : le poste s'est attribué une adresse \textbf{APIPA} (169.254.0.0/16).}

\question{A6}{0,5 pt}{Quel protocole utilise \texttt{ping} ? Nommez les deux messages échangés et la couche OSI des adresses IP.}
\zone{1.5cm}{\textbf{ICMP} ; \textbf{Echo Request} / \textbf{Echo Reply} ; adresses IP en \textbf{couche 3 (Réseau)}.}

%==========================================================================
\newpage
\section{Partie B — Réalisation sous Packet Tracer \hfill /10}

\begin{cdc}
\entreprise{} vous confie son réseau : \textbf{trois réseaux} reliés par un \textbf{routeur unique Cisco 2911}
nommé \ip{\Rnom} (interfaces G0/0, G0/1, G0/2).
\begin{itemize}[leftmargin=12pt]
  \item \textbf{Réseau \resUn} : \textbf{4 postes} PC-\pUn1 à PC-\pUn4, commutateur \ip{SW-\pUn}, adressage statique,
        réseau \ip{\netUn.0/24}.
  \item \textbf{Réseau \resDeux} : \textbf{2 postes} PC-\pDeux1, PC-\pDeux2 et \textbf{1 imprimante} IMP-\pDeux,
        commutateur \ip{SW-\pDeux}, adressage statique, réseau de classe B \ip{\netDeux.0.0/16}.
  \item \textbf{Réseau Serveur} : serveur \ip{\srvNom} d'adresse \ip{\netSrv.\srvH/24}, relié \textbf{directement} au routeur.
  \item La passerelle de chaque réseau est sa \textbf{dernière adresse utilisable}. Tous les équipements communiquent
        entre eux et accèdent à la page web du serveur. Nommage imposé (« Display Name »).
\end{itemize}
\end{cdc}

\question{B1}{2 pts}{Complétez le plan d'adressage \textbf{avant} la réalisation.}
{\footnotesize\setlength{\tabcolsep}{3pt}\renewcommand{\arraystretch}{1.15}
\begin{tabularx}{\linewidth}{|l|C|C|C|C|C|C|}
  \hline
  \textbf{Réseau} & \textbf{Adresse réseau} & \textbf{Masque} & \textbf{1re @ utilisable} &
  \textbf{Dernière @ utilisable} & \textbf{Diffusion} & \textbf{Passerelle} \\ \hline
  \hl \resUn & \sol{\netUn.0} & \sol{255.255.255.0} & \sol{\netUn.1} & \sol{\netUn.254} & \sol{\netUn.255} & \sol{\netUn.254} \\ \hline
  \hl \resDeux & \sol{\netDeux.0.0} & \sol{255.255.0.0} & \sol{\netDeux.0.1} & \sol{\netDeux.255.254} & \sol{\netDeux.255.255} & \sol{\netDeux.255.254} \\ \hline
  \hl Serveur & \sol{\netSrv.0} & \sol{255.255.255.0} & \sol{\netSrv.1} & \sol{\netSrv.254} & \sol{\netSrv.255} & \sol{\netSrv.254} \\ \hline
\end{tabularx}}
\bareme{0,5 pt par ligne juste + 0,5 pt si la convention de passerelle est respectée partout.}

\question{B2}{2 pts}{Construisez la topologie (équipements, câblage, nommage) et relevez l'adressage configuré.}
{\footnotesize\setlength{\tabcolsep}{3pt}\renewcommand{\arraystretch}{1.1}
\begin{tabularx}{\linewidth}{|l|C|C|C||l|C|C|}
  \hline
  \textbf{Hôte} & \textbf{Adresse IP} & \textbf{Masque} & \textbf{Passerelle} & \textbf{Interface} & \textbf{Adresse IP} & \textbf{Masque} \\ \hline
  \hl PC-\pUn1 & \sol{\netUn.1} & \sol{255.255.255.0} & \sol{\netUn.254} & \Rnom{} G0/0 & \sol{\netUn.254} & \sol{255.255.255.0} \\ \hline
  \hl PC-\pDeux1 & \sol{\netDeux.0.1} & \sol{255.255.0.0} & \sol{\netDeux.255.254} & \Rnom{} G0/1 & \sol{\netDeux.255.254} & \sol{255.255.0.0} \\ \hline
  \hl IMP-\pDeux & \sol{\netDeux.0.10} & \sol{255.255.0.0} & \sol{\netDeux.255.254} & \Rnom{} G0/2 & \sol{\netSrv.254} & \sol{255.255.255.0} \\ \hline
  \hl \srvNom & \ip{\netSrv.\srvH} & \sol{255.255.255.0} & \sol{\netSrv.254} & \multicolumn{3}{c|}{\cellcolor{black!8}} \\ \hline
\end{tabularx}}
\bareme{Adresses d'hôtes libres (uniques, dans le bon réseau) ; affectation G0/x libre. Sur le .pkt : équipements 0,5 ;
câblage 1 (câbles droits vers les switchs et switch–routeur, \textbf{croisé} serveur–routeur, liens verts) ; nommage 0,5.}

\question{B3}{3 pts}{Configurez tous les équipements : IP, masques, passerelles des hôtes ; interfaces du routeur
(« Port Status : On »).}
\bareme{Sur le .pkt : interfaces routeur 1 pt ; IP/masques des hôtes 1 pt ; passerelles (imprimante et serveur compris)
1 pt. $-0,25$ par erreur, sans descendre sous 0 par item.}

\question{B4}{3 pts}{Réalisez les tests, notez la commande et le résultat, puis faites viser chaque test réussi.}
{\small\renewcommand{\arraystretch}{1.1}
\begin{tabularx}{\linewidth}{|c|l|l|C|c|c|}
  \hline
  \textbf{N°} & \textbf{Source} & \textbf{Destination} & \textbf{Commande / action} & \textbf{Résultat} & \textbf{Visa prof.} \\ \hline
  \hl 1 & PC-\pUn1 & PC-\pUn4 & \sol{ping \netUn.4} & \sol{réussi} & \hspace{1.4cm} \\ \hline
  \hl 2 & PC-\pDeux1 & IMP-\pDeux & \sol{ping \netDeux.0.10} & \sol{réussi} & \\ \hline
  \hl 3 & PC-\pUn2 & PC-\pDeux2 & \sol{ping \netDeux.0.2} & \sol{réussi} & \\ \hline
  \hl 4 & PC-\pDeux2 & \srvNom & \sol{ping \netSrv.\srvH} & \sol{réussi} & \\ \hline
  \hl 5 & PC-\pUn3 & page web du serveur & \sol{navigateur : {\NoAutoSpacing http://\netSrv.\srvH}} & \sol{page affichée} & \\ \hline
\end{tabularx}}
\bareme{0,5 pt par test visé + 0,5 pt pour B5. Le premier ping inter-réseaux peut échouer (résolution ARP) ; les suivants doivent réussir.}

\question{B5}{incl. B4}{Test n°3 : par quels équipements passe le paquet ? Pourquoi PC-\pUn2 l'envoie-t-il au routeur ?}
\zone{2.3cm}{PC-\pUn2 $\to$ SW-\pUn{} $\to$ \Rnom{} $\to$ SW-\pDeux{} $\to$ PC-\pDeux2. Grâce à son masque, PC-\pUn2 constate que
la destination (\netDeux.0.0/16) n'est \textbf{pas dans son réseau} (\netUn.0/24) : il envoie le paquet à sa
\textbf{passerelle par défaut}, qui le route.}

%==========================================================================
\newpage
\section{Partie C — Diagnostic d'un réseau en panne \hfill /4}

Un technicien constate qu'\textbf{aucun poste ne parvient à joindre le serveur SRV-C} et que certains postes
communiquent mal entre eux. Les câbles sont du bon type.

\begin{center}
\begin{minipage}{0.37\linewidth}\centering
\begin{tikzpicture}[font=\footnotesize,
  eq/.style={draw,fill=qbg,rounded corners=2pt,minimum width=9mm,minimum height=5mm,inner sep=2pt},
  every path/.style={thick}]
  \node[eq] (SW) {SW-1};
  \node[eq,right=9mm of SW] (R) {R1};
  \node[eq,right=7mm of R] (SRV) {SRV-C};
  \node[eq,above left=5mm and -3mm of SW] (A) {PC-A};
  \node[eq,above right=5mm and -3mm of SW] (B) {PC-B};
  \node[eq,below left=5mm and -3mm of SW] (C) {PC-C};
  \node[eq,below right=5mm and -3mm of SW] (D) {PC-D};
  \foreach \n in {A,B,C,D} \draw (\n) -- (SW);
  \draw (SW) -- node[above,font=\tiny]{G0/0} (R);
  \draw (R) -- node[above,font=\tiny]{G0/1} (SRV);
\end{tikzpicture}
\end{minipage}\hfill
\begin{minipage}{0.61\linewidth}\centering
{\small\renewcommand{\arraystretch}{1.05}\setlength{\tabcolsep}{4pt}
\begin{tabular}{lllll}
  \toprule
  \textbf{Équipement} & \textbf{Adresse IP} & \textbf{Masque} & \textbf{Passerelle} & \textbf{Port} \\
  \midrule
  \Crows
  \bottomrule
\end{tabular}}
\end{minipage}
\end{center}

\question{C1}{4 pts}{Relevez les \textbf{quatre anomalies}. Pour chacune : conséquence et correction.}
{\small\renewcommand{\arraystretch}{1.1}
\ifcorrige\def\hc{}\else\def\hc{\rule[-10mm]{0pt}{14mm}}\fi
\begin{tabularx}{\linewidth}{|c|>{\raggedright\arraybackslash}p{2cm}|L|L|L|}
  \hline
  \textbf{N°} & \textbf{Équipement} & \textbf{Anomalie} & \textbf{Conséquence} & \textbf{Correction} \\ \hline
  \ifcorrige\Csol\else
  \hc 1 & & & & \\ \hline
  \hc 2 & & & & \\ \hline
  \hc 3 & & & & \\ \hline
  \hc 4 & & & & \\ \hline
  \fi
\end{tabularx}}
\bareme{1 pt par anomalie (0,25 équipement, anomalie, conséquence, correction ; ordre indifférent). \Cbareme}

%==========================================================================
\section{Bonus — Extension sans fil \hfill +1}

\begin{cdc}Ajoutez au réseau \resUn{} un \textbf{AccessPoint-PT} relié à SW-\pUn{} (SSID \ip{\ssid}, \textbf{WPA2-PSK})
et un portable \ip{Laptop-\pUn} équipé d'un module \texttt{WPC300N}, en adressage \textbf{statique}.
Vérifiez qu'il joint \srvNom.\end{cdc}
{\small\renewcommand{\arraystretch}{1.1}
\begin{tabularx}{\linewidth}{|l|C|C|C|c|}
  \hline
  \textbf{Équipement} & \textbf{Adresse IP} & \textbf{Masque} & \textbf{Passerelle} & \textbf{Visa prof.} \\ \hline
  \hl Laptop-\pUn & \sol{\netUn.50} & \sol{255.255.255.0} & \sol{\netUn.254} & \hspace{1.4cm} \\ \hline
\end{tabularx}}
\bareme{Le point d'accès fait office de pont : le portable est dans \netUn.0/24. Bonus si associé en WPA2-PSK et ping
vers \netSrv.\srvH{} réussi. Note plafonnée à 20.}

\end{document}
"
%   pdflatex -jobname=Eval_Reseau_SujetB   "\def\variante{B}\def\modecorrige{0}% Évaluation pratique — Réseaux locaux (BTS CIEL 1re année) — sujets A et B
% Compilation (4 PDF) :
%   pdflatex -jobname=Eval_Reseau_SujetA   "\def\variante{A}\def\modecorrige{0}\input{Eval_TP_Reseau_LAN.tex}"
%   pdflatex -jobname=Eval_Reseau_SujetB   "\def\variante{B}\def\modecorrige{0}\input{Eval_TP_Reseau_LAN.tex}"
%   pdflatex -jobname=Eval_Reseau_CorrigeA "\def\variante{A}\def\modecorrige{1}\input{Eval_TP_Reseau_LAN.tex}"
%   pdflatex -jobname=Eval_Reseau_CorrigeB "\def\variante{B}\def\modecorrige{1}\input{Eval_TP_Reseau_LAN.tex}"
\documentclass[10pt,a4paper]{article}

\usepackage[T1]{fontenc}
\usepackage[utf8]{inputenc}
\usepackage{lmodern}
\usepackage[french]{babel}
\usepackage[a4paper,hmargin=1.5cm,top=1.3cm,bottom=1.6cm,footskip=0.7cm]{geometry}
\usepackage[table]{xcolor}
\usepackage{booktabs,tabularx,array}
\usepackage{enumitem}
\usepackage{fancyhdr}
\usepackage[most]{tcolorbox}
\usepackage{tikz}
\usetikzlibrary{positioning}
\usepackage{titlesec}

\providecommand{\modecorrige}{0}
\providecommand{\variante}{A}
\newif\ifcorrige
\ifnum\modecorrige=1 \corrigetrue \else \corrigefalse \fi

\definecolor{ipred}{HTML}{C0392B}
\definecolor{cdcbar}{HTML}{A67C3D}
\definecolor{cdcbg}{HTML}{F4F1EA}
\definecolor{qbar}{HTML}{4A6B5D}
\definecolor{qbg}{HTML}{EDF1EF}
\definecolor{corr}{HTML}{1F5FA8}
\definecolor{corrbg}{HTML}{E8F0FA}
\definecolor{grisfoot}{HTML}{828282}

\newcommand{\ip}[1]{\textcolor{ipred}{\texttt{#1}}}
\newcommand{\pts}[1]{\hfill\mbox{\textbf{(#1)}}}
\newcommand{\sol}[1]{\ifcorrige\textcolor{corr}{#1}\fi}
\newcommand{\bareme}[1]{\ifcorrige\par{\footnotesize\color{corr}#1\par}\fi}
\newcommand{\hl}{\ifcorrige\rule[-1.6mm]{0pt}{5.4mm}\else\rule[-2.2mm]{0pt}{6.6mm}\fi}
\newcolumntype{C}{>{\centering\arraybackslash}X}
\newcolumntype{L}{>{\raggedright\arraybackslash}X}

\titleformat{\section}{\large\bfseries}{}{0pt}{}
\titlespacing*{\section}{0pt}{6pt}{3pt}

\tcbset{qstyle/.style={enhanced,boxrule=0pt,leftrule=2.5pt,arc=0pt,outer arc=0pt,
  left=5pt,right=5pt,top=2pt,bottom=2pt,before skip=4pt,after skip=3pt}}
\newtcolorbox{cdc}{qstyle,colback=cdcbg,colframe=cdcbar}
\newtcolorbox{quest}{qstyle,colback=qbg,colframe=qbar}

% Zone de réponse : vierge dans le sujet, remplie dans le corrigé
\newcommand{\zone}[2]{%
  \ifcorrige
    \begin{tcolorbox}[colback=corrbg,colframe=corr,boxrule=0.5pt,arc=1pt,left=4pt,right=4pt,top=1pt,bottom=1pt,
      before skip=2pt,after skip=2pt,fontupper=\footnotesize\color{corr}]#2\end{tcolorbox}%
  \else
    \begin{tcolorbox}[colback=white,colframe=black!45,boxrule=0.4pt,arc=1pt,height=#1,
      before skip=2pt,after skip=2pt]\end{tcolorbox}%
  \fi}

\newcommand{\question}[3]{\begin{quest}\textbf{#1.} #3\pts{#2}\end{quest}}

%==========================================================================
% Données propres à chaque variante
%==========================================================================
\if\variante A
  \def\entreprise{Le cabinet d'architectes \textbf{ArchiNova}}
  \def\Rnom{R-ArchiNova}
  \def\resUn{Bureau d'études}\def\pUn{BE}\def\netUn{192.168.40}
  \def\resDeux{Accueil}\def\pDeux{ACC}\def\netDeux{172.16}
  \def\srvNom{SRV-Intranet}\def\netSrv{192.168.200}\def\srvH{10}
  \def\ssid{ArchiNova-WiFi}
  \def\AIrows{%
    \hl\ip{192.168.40.25} & \sol{C} & \sol{255.255.255.0} & \sol{192.168.40.0} & \sol{192.168.40.255} & \sol{254} \\ \hline
    \hl\ip{172.16.5.10}   & \sol{B} & \sol{255.255.0.0}   & \sol{172.16.0.0}   & \sol{172.16.255.255} & \sol{65\,534} \\ \hline
    \hl\ip{10.3.2.1}      & \sol{A} & \sol{255.0.0.0}     & \sol{10.0.0.0}     & \sol{10.255.255.255} & \sol{16\,777\,214} \\ \hline}
  \def\apipa{169.254.12.7}
  \def\Cnet{192.168.60}
  \def\Crows{%
    PC-A & \ip{192.168.60.10} & \ip{255.255.255.0} & \ip{192.168.60.254} & — \\
    PC-B & \ip{192.168.60.20} & \ip{255.255.255.0} & \ip{192.168.60.1}   & — \\
    PC-C & \ip{192.168.61.30} & \ip{255.255.255.0} & \ip{192.168.60.254} & — \\
    PC-D & \ip{192.168.60.10} & \ip{255.255.255.0} & \ip{192.168.60.254} & — \\
    SRV-C & \ip{10.0.0.5}     & \ip{255.0.0.0}     & \ip{10.0.0.254}     & — \\
    R1 — G0/0 & \ip{192.168.60.254} & \ip{255.255.255.0} & — & On \\
    R1 — G0/1 & \ip{10.0.0.254}     & \ip{255.0.0.0}     & — & Off \\}
  \def\Csol{%
    \hl 1 & \sol{PC-B} & \sol{Passerelle .60.1 : pas l'adresse du routeur.} & \sol{Joint son réseau, mais aucun autre (serveur).} & \sol{Passerelle 192.168.60.254.} \\ \hline
    \hl 2 & \sol{PC-C} & \sol{IP hors du réseau 192.168.60.0/24.} & \sol{Ne joint ni les postes ni le routeur.} & \sol{Ex. 192.168.60.30.} \\ \hline
    \hl 3 & \sol{PC-D} & \sol{Même IP que PC-A (doublon).} & \sol{Conflit d'adresse : PC-A et PC-D perturbés.} & \sol{Adresse unique, ex. .60.40.} \\ \hline
    \hl 4 & \sol{R1 G0/1} & \sol{Interface désactivée (Off).} & \sol{Réseau du serveur injoignable pour tous.} & \sol{Port Status : On.} \\ \hline}
  \def\Cbareme{PC-A et SRV-C sont correctement configurés ; accepter PC-A cité à la place de PC-D pour le doublon.}
\else
  \def\entreprise{La clinique vétérinaire \textbf{VetoPlus}}
  \def\Rnom{R-VetoPlus}
  \def\resUn{Soins}\def\pUn{SOI}\def\netUn{192.168.70}
  \def\resDeux{Secrétariat}\def\pDeux{SEC}\def\netDeux{172.20}
  \def\srvNom{SRV-Dossiers}\def\netSrv{192.168.250}\def\srvH{20}
  \def\ssid{VetoPlus-WiFi}
  \def\AIrows{%
    \hl\ip{192.168.15.130} & \sol{C} & \sol{255.255.255.0} & \sol{192.168.15.0} & \sol{192.168.15.255} & \sol{254} \\ \hline
    \hl\ip{172.31.200.7}   & \sol{B} & \sol{255.255.0.0}   & \sol{172.31.0.0}   & \sol{172.31.255.255} & \sol{65\,534} \\ \hline
    \hl\ip{10.45.6.9}      & \sol{A} & \sol{255.0.0.0}     & \sol{10.0.0.0}     & \sol{10.255.255.255} & \sol{16\,777\,214} \\ \hline}
  \def\apipa{169.254.88.3}
  \def\Cnet{192.168.80}
  \def\Crows{%
    PC-A & \ip{192.168.8.11}   & \ip{255.255.255.0} & \ip{192.168.80.254} & — \\
    PC-B & \ip{192.168.80.12}  & \ip{255.255.255.0} & \ip{192.168.80.254} & — \\
    PC-C & \ip{192.168.80.13}  & \ip{255.255.255.0} & \ip{192.168.80.100} & — \\
    PC-D & \ip{192.168.80.12}  & \ip{255.255.255.0} & \ip{192.168.80.254} & — \\
    SRV-C & \ip{10.0.0.8}      & \ip{255.0.0.0}     & \ip{10.0.0.254}     & — \\
    R1 — G0/0 & \ip{192.168.80.254} & \ip{255.255.255.0} & — & Off \\
    R1 — G0/1 & \ip{10.0.0.254}     & \ip{255.0.0.0}     & — & On \\}
  \def\Csol{%
    \hl 1 & \sol{PC-A} & \sol{IP hors du réseau 192.168.80.0/24.} & \sol{Ne joint ni les postes ni le routeur.} & \sol{Ex. 192.168.80.11.} \\ \hline
    \hl 2 & \sol{PC-C} & \sol{Passerelle .80.100 : pas l'adresse du routeur.} & \sol{Joint son réseau, mais aucun autre (serveur).} & \sol{Passerelle 192.168.80.254.} \\ \hline
    \hl 3 & \sol{PC-D} & \sol{Même IP que PC-B (doublon).} & \sol{Conflit d'adresse : PC-B et PC-D perturbés.} & \sol{Adresse unique, ex. .80.14.} \\ \hline
    \hl 4 & \sol{R1 G0/0} & \sol{Interface côté LAN désactivée (Off).} & \sol{Aucun poste ne joint le routeur ni le serveur.} & \sol{Port Status : On.} \\ \hline}
  \def\Cbareme{PC-B et SRV-C sont correctement configurés ; accepter PC-B cité à la place de PC-D pour le doublon.}
\fi

\pagestyle{fancy}
\fancyhf{}
\renewcommand{\headrulewidth}{0pt}
\fancyfoot[L]{\footnotesize\itshape\color{grisfoot}Évaluation — Réseaux locaux — Sujet \variante\ifcorrige\ — CORRIGÉ\fi}
\fancyfoot[R]{\footnotesize\color{grisfoot}\thepage/3}

\setlist{itemsep=0pt,topsep=2pt,parsep=0pt}
\setlength{\parindent}{0pt}
\setlength{\parskip}{2pt}

\begin{document}

\begin{center}
  {\Large\bfseries Évaluation pratique — Réseaux locaux \quad
   \fcolorbox{black}{white}{\ SUJET \variante\ }\par}
  \vspace{1pt}
  {\bfseries Adressage IP, commutation et routage — BTS CIEL 1re année — durée : 2 heures\par}
  \ifcorrige\vspace{2pt}{\bfseries\color{corr}CORRIGÉ ET BARÈME — document enseignant}\par\fi
\end{center}
\vspace{-2pt}
{\renewcommand{\arraystretch}{1.35}
\begin{tabularx}{\linewidth}{|X|X|p{2.2cm}|p{2.6cm}|}
  \hline
  \textbf{Nom :} & \textbf{Prénom :} & \textbf{Poste n°} & \textbf{Note :} \hfill /20 \\ \hline
\end{tabularx}}

\textbf{Consignes :} travail \textbf{individuel}, documents et Internet \textbf{non autorisés}. Ce sujet est le
\textbf{document-réponse}. Fichier Packet Tracer : \ip{Nom\_Prenom\_EvalLAN\_\variante.pkt}, à enregistrer
régulièrement dans le dossier de rendu. Les « visas » sont apposés par le professeur, que vous appelez à chaque test réussi.
Durées conseillées : A 30 min — B 1 h 10 — C 20 min.

%==========================================================================
\section{Partie A — Questions de cours et calculs \hfill /6}

\question{A1}{2 pts}{Complétez le tableau en considérant le \textbf{masque par défaut} de la classe de chaque adresse.}
{\small\renewcommand{\arraystretch}{1.15}
\begin{tabularx}{\linewidth}{|l|c|C|C|C|c|}
  \hline
  \textbf{Adresse IP} & \textbf{Classe} & \textbf{Masque par défaut} & \textbf{Adresse réseau} &
  \textbf{Adresse de diffusion} & \textbf{Nb d'hôtes} \\ \hline
  \AIrows
\end{tabularx}}
\bareme{0,5 pt par ligne juste + 0,5 pt si les trois nombres d'hôtes sont justes ($2^{n}-2$, $n$ = nombre de bits d'hôte).}

\question{A2}{1 pt}{Indiquez le type de câble (\textbf{droit} ou \textbf{croisé}) pour chaque liaison.}
{\small\renewcommand{\arraystretch}{1.15}
\begin{tabularx}{\linewidth}{|X|c||X|c|}
  \hline
  \hl PC $\leftrightarrow$ commutateur & \makebox[1.6cm]{\sol{droit}} & Commutateur $\leftrightarrow$ routeur & \makebox[1.6cm]{\sol{droit}} \\ \hline
  \hl PC $\leftrightarrow$ PC (réseau P2P) & \sol{croisé} & Concentrateur $\leftrightarrow$ concentrateur & \sol{croisé} \\ \hline
\end{tabularx}}
\bareme{0,25 pt par case. Équipements de types différents : câble droit ; de même type : câble croisé.}

\question{A3}{1 pt}{Comment un \textbf{concentrateur (hub)} et un \textbf{commutateur (switch)} traitent-ils une trame
reçue ? Lequel limite les collisions, et pourquoi ?}
\zone{2.9cm}{Le \textbf{hub} (couche 1) répète la trame sur \textbf{tous ses autres ports} : un seul domaine de collision
partagé. Le \textbf{switch} (couche 2) lit l'adresse MAC de destination et, grâce à sa \textbf{table MAC}, ne l'envoie
\textbf{que sur le port du destinataire} : un domaine de collision par port, c'est lui qui limite les collisions.
\emph{0,5 pt hub, 0,5 pt switch.}}

\question{A4}{1 pt}{Qu'est-ce que la \textbf{passerelle par défaut} d'un poste ? Quelle condition son adresse doit-elle respecter ?}
\zone{2.4cm}{Adresse IP de l'interface du \textbf{routeur} à laquelle le poste envoie les paquets destinés à un
\textbf{autre réseau}. Elle doit \textbf{appartenir au même réseau} que le poste. \emph{0,5 + 0,5 pt.}}

\question{A5}{0,5 pt}{Un portable configuré en DHCP affiche \ip{\apipa} / \ip{255.255.0.0}. Que s'est-il passé ?}
\zone{1.5cm}{Aucun \textbf{serveur DHCP} n'a répondu : le poste s'est attribué une adresse \textbf{APIPA} (169.254.0.0/16).}

\question{A6}{0,5 pt}{Quel protocole utilise \texttt{ping} ? Nommez les deux messages échangés et la couche OSI des adresses IP.}
\zone{1.5cm}{\textbf{ICMP} ; \textbf{Echo Request} / \textbf{Echo Reply} ; adresses IP en \textbf{couche 3 (Réseau)}.}

%==========================================================================
\newpage
\section{Partie B — Réalisation sous Packet Tracer \hfill /10}

\begin{cdc}
\entreprise{} vous confie son réseau : \textbf{trois réseaux} reliés par un \textbf{routeur unique Cisco 2911}
nommé \ip{\Rnom} (interfaces G0/0, G0/1, G0/2).
\begin{itemize}[leftmargin=12pt]
  \item \textbf{Réseau \resUn} : \textbf{4 postes} PC-\pUn1 à PC-\pUn4, commutateur \ip{SW-\pUn}, adressage statique,
        réseau \ip{\netUn.0/24}.
  \item \textbf{Réseau \resDeux} : \textbf{2 postes} PC-\pDeux1, PC-\pDeux2 et \textbf{1 imprimante} IMP-\pDeux,
        commutateur \ip{SW-\pDeux}, adressage statique, réseau de classe B \ip{\netDeux.0.0/16}.
  \item \textbf{Réseau Serveur} : serveur \ip{\srvNom} d'adresse \ip{\netSrv.\srvH/24}, relié \textbf{directement} au routeur.
  \item La passerelle de chaque réseau est sa \textbf{dernière adresse utilisable}. Tous les équipements communiquent
        entre eux et accèdent à la page web du serveur. Nommage imposé (« Display Name »).
\end{itemize}
\end{cdc}

\question{B1}{2 pts}{Complétez le plan d'adressage \textbf{avant} la réalisation.}
{\footnotesize\setlength{\tabcolsep}{3pt}\renewcommand{\arraystretch}{1.15}
\begin{tabularx}{\linewidth}{|l|C|C|C|C|C|C|}
  \hline
  \textbf{Réseau} & \textbf{Adresse réseau} & \textbf{Masque} & \textbf{1re @ utilisable} &
  \textbf{Dernière @ utilisable} & \textbf{Diffusion} & \textbf{Passerelle} \\ \hline
  \hl \resUn & \sol{\netUn.0} & \sol{255.255.255.0} & \sol{\netUn.1} & \sol{\netUn.254} & \sol{\netUn.255} & \sol{\netUn.254} \\ \hline
  \hl \resDeux & \sol{\netDeux.0.0} & \sol{255.255.0.0} & \sol{\netDeux.0.1} & \sol{\netDeux.255.254} & \sol{\netDeux.255.255} & \sol{\netDeux.255.254} \\ \hline
  \hl Serveur & \sol{\netSrv.0} & \sol{255.255.255.0} & \sol{\netSrv.1} & \sol{\netSrv.254} & \sol{\netSrv.255} & \sol{\netSrv.254} \\ \hline
\end{tabularx}}
\bareme{0,5 pt par ligne juste + 0,5 pt si la convention de passerelle est respectée partout.}

\question{B2}{2 pts}{Construisez la topologie (équipements, câblage, nommage) et relevez l'adressage configuré.}
{\footnotesize\setlength{\tabcolsep}{3pt}\renewcommand{\arraystretch}{1.1}
\begin{tabularx}{\linewidth}{|l|C|C|C||l|C|C|}
  \hline
  \textbf{Hôte} & \textbf{Adresse IP} & \textbf{Masque} & \textbf{Passerelle} & \textbf{Interface} & \textbf{Adresse IP} & \textbf{Masque} \\ \hline
  \hl PC-\pUn1 & \sol{\netUn.1} & \sol{255.255.255.0} & \sol{\netUn.254} & \Rnom{} G0/0 & \sol{\netUn.254} & \sol{255.255.255.0} \\ \hline
  \hl PC-\pDeux1 & \sol{\netDeux.0.1} & \sol{255.255.0.0} & \sol{\netDeux.255.254} & \Rnom{} G0/1 & \sol{\netDeux.255.254} & \sol{255.255.0.0} \\ \hline
  \hl IMP-\pDeux & \sol{\netDeux.0.10} & \sol{255.255.0.0} & \sol{\netDeux.255.254} & \Rnom{} G0/2 & \sol{\netSrv.254} & \sol{255.255.255.0} \\ \hline
  \hl \srvNom & \ip{\netSrv.\srvH} & \sol{255.255.255.0} & \sol{\netSrv.254} & \multicolumn{3}{c|}{\cellcolor{black!8}} \\ \hline
\end{tabularx}}
\bareme{Adresses d'hôtes libres (uniques, dans le bon réseau) ; affectation G0/x libre. Sur le .pkt : équipements 0,5 ;
câblage 1 (câbles droits vers les switchs et switch–routeur, \textbf{croisé} serveur–routeur, liens verts) ; nommage 0,5.}

\question{B3}{3 pts}{Configurez tous les équipements : IP, masques, passerelles des hôtes ; interfaces du routeur
(« Port Status : On »).}
\bareme{Sur le .pkt : interfaces routeur 1 pt ; IP/masques des hôtes 1 pt ; passerelles (imprimante et serveur compris)
1 pt. $-0,25$ par erreur, sans descendre sous 0 par item.}

\question{B4}{3 pts}{Réalisez les tests, notez la commande et le résultat, puis faites viser chaque test réussi.}
{\small\renewcommand{\arraystretch}{1.1}
\begin{tabularx}{\linewidth}{|c|l|l|C|c|c|}
  \hline
  \textbf{N°} & \textbf{Source} & \textbf{Destination} & \textbf{Commande / action} & \textbf{Résultat} & \textbf{Visa prof.} \\ \hline
  \hl 1 & PC-\pUn1 & PC-\pUn4 & \sol{ping \netUn.4} & \sol{réussi} & \hspace{1.4cm} \\ \hline
  \hl 2 & PC-\pDeux1 & IMP-\pDeux & \sol{ping \netDeux.0.10} & \sol{réussi} & \\ \hline
  \hl 3 & PC-\pUn2 & PC-\pDeux2 & \sol{ping \netDeux.0.2} & \sol{réussi} & \\ \hline
  \hl 4 & PC-\pDeux2 & \srvNom & \sol{ping \netSrv.\srvH} & \sol{réussi} & \\ \hline
  \hl 5 & PC-\pUn3 & page web du serveur & \sol{navigateur : {\NoAutoSpacing http://\netSrv.\srvH}} & \sol{page affichée} & \\ \hline
\end{tabularx}}
\bareme{0,5 pt par test visé + 0,5 pt pour B5. Le premier ping inter-réseaux peut échouer (résolution ARP) ; les suivants doivent réussir.}

\question{B5}{incl. B4}{Test n°3 : par quels équipements passe le paquet ? Pourquoi PC-\pUn2 l'envoie-t-il au routeur ?}
\zone{2.3cm}{PC-\pUn2 $\to$ SW-\pUn{} $\to$ \Rnom{} $\to$ SW-\pDeux{} $\to$ PC-\pDeux2. Grâce à son masque, PC-\pUn2 constate que
la destination (\netDeux.0.0/16) n'est \textbf{pas dans son réseau} (\netUn.0/24) : il envoie le paquet à sa
\textbf{passerelle par défaut}, qui le route.}

%==========================================================================
\newpage
\section{Partie C — Diagnostic d'un réseau en panne \hfill /4}

Un technicien constate qu'\textbf{aucun poste ne parvient à joindre le serveur SRV-C} et que certains postes
communiquent mal entre eux. Les câbles sont du bon type.

\begin{center}
\begin{minipage}{0.37\linewidth}\centering
\begin{tikzpicture}[font=\footnotesize,
  eq/.style={draw,fill=qbg,rounded corners=2pt,minimum width=9mm,minimum height=5mm,inner sep=2pt},
  every path/.style={thick}]
  \node[eq] (SW) {SW-1};
  \node[eq,right=9mm of SW] (R) {R1};
  \node[eq,right=7mm of R] (SRV) {SRV-C};
  \node[eq,above left=5mm and -3mm of SW] (A) {PC-A};
  \node[eq,above right=5mm and -3mm of SW] (B) {PC-B};
  \node[eq,below left=5mm and -3mm of SW] (C) {PC-C};
  \node[eq,below right=5mm and -3mm of SW] (D) {PC-D};
  \foreach \n in {A,B,C,D} \draw (\n) -- (SW);
  \draw (SW) -- node[above,font=\tiny]{G0/0} (R);
  \draw (R) -- node[above,font=\tiny]{G0/1} (SRV);
\end{tikzpicture}
\end{minipage}\hfill
\begin{minipage}{0.61\linewidth}\centering
{\small\renewcommand{\arraystretch}{1.05}\setlength{\tabcolsep}{4pt}
\begin{tabular}{lllll}
  \toprule
  \textbf{Équipement} & \textbf{Adresse IP} & \textbf{Masque} & \textbf{Passerelle} & \textbf{Port} \\
  \midrule
  \Crows
  \bottomrule
\end{tabular}}
\end{minipage}
\end{center}

\question{C1}{4 pts}{Relevez les \textbf{quatre anomalies}. Pour chacune : conséquence et correction.}
{\small\renewcommand{\arraystretch}{1.1}
\ifcorrige\def\hc{}\else\def\hc{\rule[-10mm]{0pt}{14mm}}\fi
\begin{tabularx}{\linewidth}{|c|>{\raggedright\arraybackslash}p{2cm}|L|L|L|}
  \hline
  \textbf{N°} & \textbf{Équipement} & \textbf{Anomalie} & \textbf{Conséquence} & \textbf{Correction} \\ \hline
  \ifcorrige\Csol\else
  \hc 1 & & & & \\ \hline
  \hc 2 & & & & \\ \hline
  \hc 3 & & & & \\ \hline
  \hc 4 & & & & \\ \hline
  \fi
\end{tabularx}}
\bareme{1 pt par anomalie (0,25 équipement, anomalie, conséquence, correction ; ordre indifférent). \Cbareme}

%==========================================================================
\section{Bonus — Extension sans fil \hfill +1}

\begin{cdc}Ajoutez au réseau \resUn{} un \textbf{AccessPoint-PT} relié à SW-\pUn{} (SSID \ip{\ssid}, \textbf{WPA2-PSK})
et un portable \ip{Laptop-\pUn} équipé d'un module \texttt{WPC300N}, en adressage \textbf{statique}.
Vérifiez qu'il joint \srvNom.\end{cdc}
{\small\renewcommand{\arraystretch}{1.1}
\begin{tabularx}{\linewidth}{|l|C|C|C|c|}
  \hline
  \textbf{Équipement} & \textbf{Adresse IP} & \textbf{Masque} & \textbf{Passerelle} & \textbf{Visa prof.} \\ \hline
  \hl Laptop-\pUn & \sol{\netUn.50} & \sol{255.255.255.0} & \sol{\netUn.254} & \hspace{1.4cm} \\ \hline
\end{tabularx}}
\bareme{Le point d'accès fait office de pont : le portable est dans \netUn.0/24. Bonus si associé en WPA2-PSK et ping
vers \netSrv.\srvH{} réussi. Note plafonnée à 20.}

\end{document}
"
%   pdflatex -jobname=Eval_Reseau_CorrigeA "\def\variante{A}\def\modecorrige{1}% Évaluation pratique — Réseaux locaux (BTS CIEL 1re année) — sujets A et B
% Compilation (4 PDF) :
%   pdflatex -jobname=Eval_Reseau_SujetA   "\def\variante{A}\def\modecorrige{0}\input{Eval_TP_Reseau_LAN.tex}"
%   pdflatex -jobname=Eval_Reseau_SujetB   "\def\variante{B}\def\modecorrige{0}\input{Eval_TP_Reseau_LAN.tex}"
%   pdflatex -jobname=Eval_Reseau_CorrigeA "\def\variante{A}\def\modecorrige{1}\input{Eval_TP_Reseau_LAN.tex}"
%   pdflatex -jobname=Eval_Reseau_CorrigeB "\def\variante{B}\def\modecorrige{1}\input{Eval_TP_Reseau_LAN.tex}"
\documentclass[10pt,a4paper]{article}

\usepackage[T1]{fontenc}
\usepackage[utf8]{inputenc}
\usepackage{lmodern}
\usepackage[french]{babel}
\usepackage[a4paper,hmargin=1.5cm,top=1.3cm,bottom=1.6cm,footskip=0.7cm]{geometry}
\usepackage[table]{xcolor}
\usepackage{booktabs,tabularx,array}
\usepackage{enumitem}
\usepackage{fancyhdr}
\usepackage[most]{tcolorbox}
\usepackage{tikz}
\usetikzlibrary{positioning}
\usepackage{titlesec}

\providecommand{\modecorrige}{0}
\providecommand{\variante}{A}
\newif\ifcorrige
\ifnum\modecorrige=1 \corrigetrue \else \corrigefalse \fi

\definecolor{ipred}{HTML}{C0392B}
\definecolor{cdcbar}{HTML}{A67C3D}
\definecolor{cdcbg}{HTML}{F4F1EA}
\definecolor{qbar}{HTML}{4A6B5D}
\definecolor{qbg}{HTML}{EDF1EF}
\definecolor{corr}{HTML}{1F5FA8}
\definecolor{corrbg}{HTML}{E8F0FA}
\definecolor{grisfoot}{HTML}{828282}

\newcommand{\ip}[1]{\textcolor{ipred}{\texttt{#1}}}
\newcommand{\pts}[1]{\hfill\mbox{\textbf{(#1)}}}
\newcommand{\sol}[1]{\ifcorrige\textcolor{corr}{#1}\fi}
\newcommand{\bareme}[1]{\ifcorrige\par{\footnotesize\color{corr}#1\par}\fi}
\newcommand{\hl}{\ifcorrige\rule[-1.6mm]{0pt}{5.4mm}\else\rule[-2.2mm]{0pt}{6.6mm}\fi}
\newcolumntype{C}{>{\centering\arraybackslash}X}
\newcolumntype{L}{>{\raggedright\arraybackslash}X}

\titleformat{\section}{\large\bfseries}{}{0pt}{}
\titlespacing*{\section}{0pt}{6pt}{3pt}

\tcbset{qstyle/.style={enhanced,boxrule=0pt,leftrule=2.5pt,arc=0pt,outer arc=0pt,
  left=5pt,right=5pt,top=2pt,bottom=2pt,before skip=4pt,after skip=3pt}}
\newtcolorbox{cdc}{qstyle,colback=cdcbg,colframe=cdcbar}
\newtcolorbox{quest}{qstyle,colback=qbg,colframe=qbar}

% Zone de réponse : vierge dans le sujet, remplie dans le corrigé
\newcommand{\zone}[2]{%
  \ifcorrige
    \begin{tcolorbox}[colback=corrbg,colframe=corr,boxrule=0.5pt,arc=1pt,left=4pt,right=4pt,top=1pt,bottom=1pt,
      before skip=2pt,after skip=2pt,fontupper=\footnotesize\color{corr}]#2\end{tcolorbox}%
  \else
    \begin{tcolorbox}[colback=white,colframe=black!45,boxrule=0.4pt,arc=1pt,height=#1,
      before skip=2pt,after skip=2pt]\end{tcolorbox}%
  \fi}

\newcommand{\question}[3]{\begin{quest}\textbf{#1.} #3\pts{#2}\end{quest}}

%==========================================================================
% Données propres à chaque variante
%==========================================================================
\if\variante A
  \def\entreprise{Le cabinet d'architectes \textbf{ArchiNova}}
  \def\Rnom{R-ArchiNova}
  \def\resUn{Bureau d'études}\def\pUn{BE}\def\netUn{192.168.40}
  \def\resDeux{Accueil}\def\pDeux{ACC}\def\netDeux{172.16}
  \def\srvNom{SRV-Intranet}\def\netSrv{192.168.200}\def\srvH{10}
  \def\ssid{ArchiNova-WiFi}
  \def\AIrows{%
    \hl\ip{192.168.40.25} & \sol{C} & \sol{255.255.255.0} & \sol{192.168.40.0} & \sol{192.168.40.255} & \sol{254} \\ \hline
    \hl\ip{172.16.5.10}   & \sol{B} & \sol{255.255.0.0}   & \sol{172.16.0.0}   & \sol{172.16.255.255} & \sol{65\,534} \\ \hline
    \hl\ip{10.3.2.1}      & \sol{A} & \sol{255.0.0.0}     & \sol{10.0.0.0}     & \sol{10.255.255.255} & \sol{16\,777\,214} \\ \hline}
  \def\apipa{169.254.12.7}
  \def\Cnet{192.168.60}
  \def\Crows{%
    PC-A & \ip{192.168.60.10} & \ip{255.255.255.0} & \ip{192.168.60.254} & — \\
    PC-B & \ip{192.168.60.20} & \ip{255.255.255.0} & \ip{192.168.60.1}   & — \\
    PC-C & \ip{192.168.61.30} & \ip{255.255.255.0} & \ip{192.168.60.254} & — \\
    PC-D & \ip{192.168.60.10} & \ip{255.255.255.0} & \ip{192.168.60.254} & — \\
    SRV-C & \ip{10.0.0.5}     & \ip{255.0.0.0}     & \ip{10.0.0.254}     & — \\
    R1 — G0/0 & \ip{192.168.60.254} & \ip{255.255.255.0} & — & On \\
    R1 — G0/1 & \ip{10.0.0.254}     & \ip{255.0.0.0}     & — & Off \\}
  \def\Csol{%
    \hl 1 & \sol{PC-B} & \sol{Passerelle .60.1 : pas l'adresse du routeur.} & \sol{Joint son réseau, mais aucun autre (serveur).} & \sol{Passerelle 192.168.60.254.} \\ \hline
    \hl 2 & \sol{PC-C} & \sol{IP hors du réseau 192.168.60.0/24.} & \sol{Ne joint ni les postes ni le routeur.} & \sol{Ex. 192.168.60.30.} \\ \hline
    \hl 3 & \sol{PC-D} & \sol{Même IP que PC-A (doublon).} & \sol{Conflit d'adresse : PC-A et PC-D perturbés.} & \sol{Adresse unique, ex. .60.40.} \\ \hline
    \hl 4 & \sol{R1 G0/1} & \sol{Interface désactivée (Off).} & \sol{Réseau du serveur injoignable pour tous.} & \sol{Port Status : On.} \\ \hline}
  \def\Cbareme{PC-A et SRV-C sont correctement configurés ; accepter PC-A cité à la place de PC-D pour le doublon.}
\else
  \def\entreprise{La clinique vétérinaire \textbf{VetoPlus}}
  \def\Rnom{R-VetoPlus}
  \def\resUn{Soins}\def\pUn{SOI}\def\netUn{192.168.70}
  \def\resDeux{Secrétariat}\def\pDeux{SEC}\def\netDeux{172.20}
  \def\srvNom{SRV-Dossiers}\def\netSrv{192.168.250}\def\srvH{20}
  \def\ssid{VetoPlus-WiFi}
  \def\AIrows{%
    \hl\ip{192.168.15.130} & \sol{C} & \sol{255.255.255.0} & \sol{192.168.15.0} & \sol{192.168.15.255} & \sol{254} \\ \hline
    \hl\ip{172.31.200.7}   & \sol{B} & \sol{255.255.0.0}   & \sol{172.31.0.0}   & \sol{172.31.255.255} & \sol{65\,534} \\ \hline
    \hl\ip{10.45.6.9}      & \sol{A} & \sol{255.0.0.0}     & \sol{10.0.0.0}     & \sol{10.255.255.255} & \sol{16\,777\,214} \\ \hline}
  \def\apipa{169.254.88.3}
  \def\Cnet{192.168.80}
  \def\Crows{%
    PC-A & \ip{192.168.8.11}   & \ip{255.255.255.0} & \ip{192.168.80.254} & — \\
    PC-B & \ip{192.168.80.12}  & \ip{255.255.255.0} & \ip{192.168.80.254} & — \\
    PC-C & \ip{192.168.80.13}  & \ip{255.255.255.0} & \ip{192.168.80.100} & — \\
    PC-D & \ip{192.168.80.12}  & \ip{255.255.255.0} & \ip{192.168.80.254} & — \\
    SRV-C & \ip{10.0.0.8}      & \ip{255.0.0.0}     & \ip{10.0.0.254}     & — \\
    R1 — G0/0 & \ip{192.168.80.254} & \ip{255.255.255.0} & — & Off \\
    R1 — G0/1 & \ip{10.0.0.254}     & \ip{255.0.0.0}     & — & On \\}
  \def\Csol{%
    \hl 1 & \sol{PC-A} & \sol{IP hors du réseau 192.168.80.0/24.} & \sol{Ne joint ni les postes ni le routeur.} & \sol{Ex. 192.168.80.11.} \\ \hline
    \hl 2 & \sol{PC-C} & \sol{Passerelle .80.100 : pas l'adresse du routeur.} & \sol{Joint son réseau, mais aucun autre (serveur).} & \sol{Passerelle 192.168.80.254.} \\ \hline
    \hl 3 & \sol{PC-D} & \sol{Même IP que PC-B (doublon).} & \sol{Conflit d'adresse : PC-B et PC-D perturbés.} & \sol{Adresse unique, ex. .80.14.} \\ \hline
    \hl 4 & \sol{R1 G0/0} & \sol{Interface côté LAN désactivée (Off).} & \sol{Aucun poste ne joint le routeur ni le serveur.} & \sol{Port Status : On.} \\ \hline}
  \def\Cbareme{PC-B et SRV-C sont correctement configurés ; accepter PC-B cité à la place de PC-D pour le doublon.}
\fi

\pagestyle{fancy}
\fancyhf{}
\renewcommand{\headrulewidth}{0pt}
\fancyfoot[L]{\footnotesize\itshape\color{grisfoot}Évaluation — Réseaux locaux — Sujet \variante\ifcorrige\ — CORRIGÉ\fi}
\fancyfoot[R]{\footnotesize\color{grisfoot}\thepage/3}

\setlist{itemsep=0pt,topsep=2pt,parsep=0pt}
\setlength{\parindent}{0pt}
\setlength{\parskip}{2pt}

\begin{document}

\begin{center}
  {\Large\bfseries Évaluation pratique — Réseaux locaux \quad
   \fcolorbox{black}{white}{\ SUJET \variante\ }\par}
  \vspace{1pt}
  {\bfseries Adressage IP, commutation et routage — BTS CIEL 1re année — durée : 2 heures\par}
  \ifcorrige\vspace{2pt}{\bfseries\color{corr}CORRIGÉ ET BARÈME — document enseignant}\par\fi
\end{center}
\vspace{-2pt}
{\renewcommand{\arraystretch}{1.35}
\begin{tabularx}{\linewidth}{|X|X|p{2.2cm}|p{2.6cm}|}
  \hline
  \textbf{Nom :} & \textbf{Prénom :} & \textbf{Poste n°} & \textbf{Note :} \hfill /20 \\ \hline
\end{tabularx}}

\textbf{Consignes :} travail \textbf{individuel}, documents et Internet \textbf{non autorisés}. Ce sujet est le
\textbf{document-réponse}. Fichier Packet Tracer : \ip{Nom\_Prenom\_EvalLAN\_\variante.pkt}, à enregistrer
régulièrement dans le dossier de rendu. Les « visas » sont apposés par le professeur, que vous appelez à chaque test réussi.
Durées conseillées : A 30 min — B 1 h 10 — C 20 min.

%==========================================================================
\section{Partie A — Questions de cours et calculs \hfill /6}

\question{A1}{2 pts}{Complétez le tableau en considérant le \textbf{masque par défaut} de la classe de chaque adresse.}
{\small\renewcommand{\arraystretch}{1.15}
\begin{tabularx}{\linewidth}{|l|c|C|C|C|c|}
  \hline
  \textbf{Adresse IP} & \textbf{Classe} & \textbf{Masque par défaut} & \textbf{Adresse réseau} &
  \textbf{Adresse de diffusion} & \textbf{Nb d'hôtes} \\ \hline
  \AIrows
\end{tabularx}}
\bareme{0,5 pt par ligne juste + 0,5 pt si les trois nombres d'hôtes sont justes ($2^{n}-2$, $n$ = nombre de bits d'hôte).}

\question{A2}{1 pt}{Indiquez le type de câble (\textbf{droit} ou \textbf{croisé}) pour chaque liaison.}
{\small\renewcommand{\arraystretch}{1.15}
\begin{tabularx}{\linewidth}{|X|c||X|c|}
  \hline
  \hl PC $\leftrightarrow$ commutateur & \makebox[1.6cm]{\sol{droit}} & Commutateur $\leftrightarrow$ routeur & \makebox[1.6cm]{\sol{droit}} \\ \hline
  \hl PC $\leftrightarrow$ PC (réseau P2P) & \sol{croisé} & Concentrateur $\leftrightarrow$ concentrateur & \sol{croisé} \\ \hline
\end{tabularx}}
\bareme{0,25 pt par case. Équipements de types différents : câble droit ; de même type : câble croisé.}

\question{A3}{1 pt}{Comment un \textbf{concentrateur (hub)} et un \textbf{commutateur (switch)} traitent-ils une trame
reçue ? Lequel limite les collisions, et pourquoi ?}
\zone{2.9cm}{Le \textbf{hub} (couche 1) répète la trame sur \textbf{tous ses autres ports} : un seul domaine de collision
partagé. Le \textbf{switch} (couche 2) lit l'adresse MAC de destination et, grâce à sa \textbf{table MAC}, ne l'envoie
\textbf{que sur le port du destinataire} : un domaine de collision par port, c'est lui qui limite les collisions.
\emph{0,5 pt hub, 0,5 pt switch.}}

\question{A4}{1 pt}{Qu'est-ce que la \textbf{passerelle par défaut} d'un poste ? Quelle condition son adresse doit-elle respecter ?}
\zone{2.4cm}{Adresse IP de l'interface du \textbf{routeur} à laquelle le poste envoie les paquets destinés à un
\textbf{autre réseau}. Elle doit \textbf{appartenir au même réseau} que le poste. \emph{0,5 + 0,5 pt.}}

\question{A5}{0,5 pt}{Un portable configuré en DHCP affiche \ip{\apipa} / \ip{255.255.0.0}. Que s'est-il passé ?}
\zone{1.5cm}{Aucun \textbf{serveur DHCP} n'a répondu : le poste s'est attribué une adresse \textbf{APIPA} (169.254.0.0/16).}

\question{A6}{0,5 pt}{Quel protocole utilise \texttt{ping} ? Nommez les deux messages échangés et la couche OSI des adresses IP.}
\zone{1.5cm}{\textbf{ICMP} ; \textbf{Echo Request} / \textbf{Echo Reply} ; adresses IP en \textbf{couche 3 (Réseau)}.}

%==========================================================================
\newpage
\section{Partie B — Réalisation sous Packet Tracer \hfill /10}

\begin{cdc}
\entreprise{} vous confie son réseau : \textbf{trois réseaux} reliés par un \textbf{routeur unique Cisco 2911}
nommé \ip{\Rnom} (interfaces G0/0, G0/1, G0/2).
\begin{itemize}[leftmargin=12pt]
  \item \textbf{Réseau \resUn} : \textbf{4 postes} PC-\pUn1 à PC-\pUn4, commutateur \ip{SW-\pUn}, adressage statique,
        réseau \ip{\netUn.0/24}.
  \item \textbf{Réseau \resDeux} : \textbf{2 postes} PC-\pDeux1, PC-\pDeux2 et \textbf{1 imprimante} IMP-\pDeux,
        commutateur \ip{SW-\pDeux}, adressage statique, réseau de classe B \ip{\netDeux.0.0/16}.
  \item \textbf{Réseau Serveur} : serveur \ip{\srvNom} d'adresse \ip{\netSrv.\srvH/24}, relié \textbf{directement} au routeur.
  \item La passerelle de chaque réseau est sa \textbf{dernière adresse utilisable}. Tous les équipements communiquent
        entre eux et accèdent à la page web du serveur. Nommage imposé (« Display Name »).
\end{itemize}
\end{cdc}

\question{B1}{2 pts}{Complétez le plan d'adressage \textbf{avant} la réalisation.}
{\footnotesize\setlength{\tabcolsep}{3pt}\renewcommand{\arraystretch}{1.15}
\begin{tabularx}{\linewidth}{|l|C|C|C|C|C|C|}
  \hline
  \textbf{Réseau} & \textbf{Adresse réseau} & \textbf{Masque} & \textbf{1re @ utilisable} &
  \textbf{Dernière @ utilisable} & \textbf{Diffusion} & \textbf{Passerelle} \\ \hline
  \hl \resUn & \sol{\netUn.0} & \sol{255.255.255.0} & \sol{\netUn.1} & \sol{\netUn.254} & \sol{\netUn.255} & \sol{\netUn.254} \\ \hline
  \hl \resDeux & \sol{\netDeux.0.0} & \sol{255.255.0.0} & \sol{\netDeux.0.1} & \sol{\netDeux.255.254} & \sol{\netDeux.255.255} & \sol{\netDeux.255.254} \\ \hline
  \hl Serveur & \sol{\netSrv.0} & \sol{255.255.255.0} & \sol{\netSrv.1} & \sol{\netSrv.254} & \sol{\netSrv.255} & \sol{\netSrv.254} \\ \hline
\end{tabularx}}
\bareme{0,5 pt par ligne juste + 0,5 pt si la convention de passerelle est respectée partout.}

\question{B2}{2 pts}{Construisez la topologie (équipements, câblage, nommage) et relevez l'adressage configuré.}
{\footnotesize\setlength{\tabcolsep}{3pt}\renewcommand{\arraystretch}{1.1}
\begin{tabularx}{\linewidth}{|l|C|C|C||l|C|C|}
  \hline
  \textbf{Hôte} & \textbf{Adresse IP} & \textbf{Masque} & \textbf{Passerelle} & \textbf{Interface} & \textbf{Adresse IP} & \textbf{Masque} \\ \hline
  \hl PC-\pUn1 & \sol{\netUn.1} & \sol{255.255.255.0} & \sol{\netUn.254} & \Rnom{} G0/0 & \sol{\netUn.254} & \sol{255.255.255.0} \\ \hline
  \hl PC-\pDeux1 & \sol{\netDeux.0.1} & \sol{255.255.0.0} & \sol{\netDeux.255.254} & \Rnom{} G0/1 & \sol{\netDeux.255.254} & \sol{255.255.0.0} \\ \hline
  \hl IMP-\pDeux & \sol{\netDeux.0.10} & \sol{255.255.0.0} & \sol{\netDeux.255.254} & \Rnom{} G0/2 & \sol{\netSrv.254} & \sol{255.255.255.0} \\ \hline
  \hl \srvNom & \ip{\netSrv.\srvH} & \sol{255.255.255.0} & \sol{\netSrv.254} & \multicolumn{3}{c|}{\cellcolor{black!8}} \\ \hline
\end{tabularx}}
\bareme{Adresses d'hôtes libres (uniques, dans le bon réseau) ; affectation G0/x libre. Sur le .pkt : équipements 0,5 ;
câblage 1 (câbles droits vers les switchs et switch–routeur, \textbf{croisé} serveur–routeur, liens verts) ; nommage 0,5.}

\question{B3}{3 pts}{Configurez tous les équipements : IP, masques, passerelles des hôtes ; interfaces du routeur
(« Port Status : On »).}
\bareme{Sur le .pkt : interfaces routeur 1 pt ; IP/masques des hôtes 1 pt ; passerelles (imprimante et serveur compris)
1 pt. $-0,25$ par erreur, sans descendre sous 0 par item.}

\question{B4}{3 pts}{Réalisez les tests, notez la commande et le résultat, puis faites viser chaque test réussi.}
{\small\renewcommand{\arraystretch}{1.1}
\begin{tabularx}{\linewidth}{|c|l|l|C|c|c|}
  \hline
  \textbf{N°} & \textbf{Source} & \textbf{Destination} & \textbf{Commande / action} & \textbf{Résultat} & \textbf{Visa prof.} \\ \hline
  \hl 1 & PC-\pUn1 & PC-\pUn4 & \sol{ping \netUn.4} & \sol{réussi} & \hspace{1.4cm} \\ \hline
  \hl 2 & PC-\pDeux1 & IMP-\pDeux & \sol{ping \netDeux.0.10} & \sol{réussi} & \\ \hline
  \hl 3 & PC-\pUn2 & PC-\pDeux2 & \sol{ping \netDeux.0.2} & \sol{réussi} & \\ \hline
  \hl 4 & PC-\pDeux2 & \srvNom & \sol{ping \netSrv.\srvH} & \sol{réussi} & \\ \hline
  \hl 5 & PC-\pUn3 & page web du serveur & \sol{navigateur : {\NoAutoSpacing http://\netSrv.\srvH}} & \sol{page affichée} & \\ \hline
\end{tabularx}}
\bareme{0,5 pt par test visé + 0,5 pt pour B5. Le premier ping inter-réseaux peut échouer (résolution ARP) ; les suivants doivent réussir.}

\question{B5}{incl. B4}{Test n°3 : par quels équipements passe le paquet ? Pourquoi PC-\pUn2 l'envoie-t-il au routeur ?}
\zone{2.3cm}{PC-\pUn2 $\to$ SW-\pUn{} $\to$ \Rnom{} $\to$ SW-\pDeux{} $\to$ PC-\pDeux2. Grâce à son masque, PC-\pUn2 constate que
la destination (\netDeux.0.0/16) n'est \textbf{pas dans son réseau} (\netUn.0/24) : il envoie le paquet à sa
\textbf{passerelle par défaut}, qui le route.}

%==========================================================================
\newpage
\section{Partie C — Diagnostic d'un réseau en panne \hfill /4}

Un technicien constate qu'\textbf{aucun poste ne parvient à joindre le serveur SRV-C} et que certains postes
communiquent mal entre eux. Les câbles sont du bon type.

\begin{center}
\begin{minipage}{0.37\linewidth}\centering
\begin{tikzpicture}[font=\footnotesize,
  eq/.style={draw,fill=qbg,rounded corners=2pt,minimum width=9mm,minimum height=5mm,inner sep=2pt},
  every path/.style={thick}]
  \node[eq] (SW) {SW-1};
  \node[eq,right=9mm of SW] (R) {R1};
  \node[eq,right=7mm of R] (SRV) {SRV-C};
  \node[eq,above left=5mm and -3mm of SW] (A) {PC-A};
  \node[eq,above right=5mm and -3mm of SW] (B) {PC-B};
  \node[eq,below left=5mm and -3mm of SW] (C) {PC-C};
  \node[eq,below right=5mm and -3mm of SW] (D) {PC-D};
  \foreach \n in {A,B,C,D} \draw (\n) -- (SW);
  \draw (SW) -- node[above,font=\tiny]{G0/0} (R);
  \draw (R) -- node[above,font=\tiny]{G0/1} (SRV);
\end{tikzpicture}
\end{minipage}\hfill
\begin{minipage}{0.61\linewidth}\centering
{\small\renewcommand{\arraystretch}{1.05}\setlength{\tabcolsep}{4pt}
\begin{tabular}{lllll}
  \toprule
  \textbf{Équipement} & \textbf{Adresse IP} & \textbf{Masque} & \textbf{Passerelle} & \textbf{Port} \\
  \midrule
  \Crows
  \bottomrule
\end{tabular}}
\end{minipage}
\end{center}

\question{C1}{4 pts}{Relevez les \textbf{quatre anomalies}. Pour chacune : conséquence et correction.}
{\small\renewcommand{\arraystretch}{1.1}
\ifcorrige\def\hc{}\else\def\hc{\rule[-10mm]{0pt}{14mm}}\fi
\begin{tabularx}{\linewidth}{|c|>{\raggedright\arraybackslash}p{2cm}|L|L|L|}
  \hline
  \textbf{N°} & \textbf{Équipement} & \textbf{Anomalie} & \textbf{Conséquence} & \textbf{Correction} \\ \hline
  \ifcorrige\Csol\else
  \hc 1 & & & & \\ \hline
  \hc 2 & & & & \\ \hline
  \hc 3 & & & & \\ \hline
  \hc 4 & & & & \\ \hline
  \fi
\end{tabularx}}
\bareme{1 pt par anomalie (0,25 équipement, anomalie, conséquence, correction ; ordre indifférent). \Cbareme}

%==========================================================================
\section{Bonus — Extension sans fil \hfill +1}

\begin{cdc}Ajoutez au réseau \resUn{} un \textbf{AccessPoint-PT} relié à SW-\pUn{} (SSID \ip{\ssid}, \textbf{WPA2-PSK})
et un portable \ip{Laptop-\pUn} équipé d'un module \texttt{WPC300N}, en adressage \textbf{statique}.
Vérifiez qu'il joint \srvNom.\end{cdc}
{\small\renewcommand{\arraystretch}{1.1}
\begin{tabularx}{\linewidth}{|l|C|C|C|c|}
  \hline
  \textbf{Équipement} & \textbf{Adresse IP} & \textbf{Masque} & \textbf{Passerelle} & \textbf{Visa prof.} \\ \hline
  \hl Laptop-\pUn & \sol{\netUn.50} & \sol{255.255.255.0} & \sol{\netUn.254} & \hspace{1.4cm} \\ \hline
\end{tabularx}}
\bareme{Le point d'accès fait office de pont : le portable est dans \netUn.0/24. Bonus si associé en WPA2-PSK et ping
vers \netSrv.\srvH{} réussi. Note plafonnée à 20.}

\end{document}
"
%   pdflatex -jobname=Eval_Reseau_CorrigeB "\def\variante{B}\def\modecorrige{1}% Évaluation pratique — Réseaux locaux (BTS CIEL 1re année) — sujets A et B
% Compilation (4 PDF) :
%   pdflatex -jobname=Eval_Reseau_SujetA   "\def\variante{A}\def\modecorrige{0}\input{Eval_TP_Reseau_LAN.tex}"
%   pdflatex -jobname=Eval_Reseau_SujetB   "\def\variante{B}\def\modecorrige{0}\input{Eval_TP_Reseau_LAN.tex}"
%   pdflatex -jobname=Eval_Reseau_CorrigeA "\def\variante{A}\def\modecorrige{1}\input{Eval_TP_Reseau_LAN.tex}"
%   pdflatex -jobname=Eval_Reseau_CorrigeB "\def\variante{B}\def\modecorrige{1}\input{Eval_TP_Reseau_LAN.tex}"
\documentclass[10pt,a4paper]{article}

\usepackage[T1]{fontenc}
\usepackage[utf8]{inputenc}
\usepackage{lmodern}
\usepackage[french]{babel}
\usepackage[a4paper,hmargin=1.5cm,top=1.3cm,bottom=1.6cm,footskip=0.7cm]{geometry}
\usepackage[table]{xcolor}
\usepackage{booktabs,tabularx,array}
\usepackage{enumitem}
\usepackage{fancyhdr}
\usepackage[most]{tcolorbox}
\usepackage{tikz}
\usetikzlibrary{positioning}
\usepackage{titlesec}

\providecommand{\modecorrige}{0}
\providecommand{\variante}{A}
\newif\ifcorrige
\ifnum\modecorrige=1 \corrigetrue \else \corrigefalse \fi

\definecolor{ipred}{HTML}{C0392B}
\definecolor{cdcbar}{HTML}{A67C3D}
\definecolor{cdcbg}{HTML}{F4F1EA}
\definecolor{qbar}{HTML}{4A6B5D}
\definecolor{qbg}{HTML}{EDF1EF}
\definecolor{corr}{HTML}{1F5FA8}
\definecolor{corrbg}{HTML}{E8F0FA}
\definecolor{grisfoot}{HTML}{828282}

\newcommand{\ip}[1]{\textcolor{ipred}{\texttt{#1}}}
\newcommand{\pts}[1]{\hfill\mbox{\textbf{(#1)}}}
\newcommand{\sol}[1]{\ifcorrige\textcolor{corr}{#1}\fi}
\newcommand{\bareme}[1]{\ifcorrige\par{\footnotesize\color{corr}#1\par}\fi}
\newcommand{\hl}{\ifcorrige\rule[-1.6mm]{0pt}{5.4mm}\else\rule[-2.2mm]{0pt}{6.6mm}\fi}
\newcolumntype{C}{>{\centering\arraybackslash}X}
\newcolumntype{L}{>{\raggedright\arraybackslash}X}

\titleformat{\section}{\large\bfseries}{}{0pt}{}
\titlespacing*{\section}{0pt}{6pt}{3pt}

\tcbset{qstyle/.style={enhanced,boxrule=0pt,leftrule=2.5pt,arc=0pt,outer arc=0pt,
  left=5pt,right=5pt,top=2pt,bottom=2pt,before skip=4pt,after skip=3pt}}
\newtcolorbox{cdc}{qstyle,colback=cdcbg,colframe=cdcbar}
\newtcolorbox{quest}{qstyle,colback=qbg,colframe=qbar}

% Zone de réponse : vierge dans le sujet, remplie dans le corrigé
\newcommand{\zone}[2]{%
  \ifcorrige
    \begin{tcolorbox}[colback=corrbg,colframe=corr,boxrule=0.5pt,arc=1pt,left=4pt,right=4pt,top=1pt,bottom=1pt,
      before skip=2pt,after skip=2pt,fontupper=\footnotesize\color{corr}]#2\end{tcolorbox}%
  \else
    \begin{tcolorbox}[colback=white,colframe=black!45,boxrule=0.4pt,arc=1pt,height=#1,
      before skip=2pt,after skip=2pt]\end{tcolorbox}%
  \fi}

\newcommand{\question}[3]{\begin{quest}\textbf{#1.} #3\pts{#2}\end{quest}}

%==========================================================================
% Données propres à chaque variante
%==========================================================================
\if\variante A
  \def\entreprise{Le cabinet d'architectes \textbf{ArchiNova}}
  \def\Rnom{R-ArchiNova}
  \def\resUn{Bureau d'études}\def\pUn{BE}\def\netUn{192.168.40}
  \def\resDeux{Accueil}\def\pDeux{ACC}\def\netDeux{172.16}
  \def\srvNom{SRV-Intranet}\def\netSrv{192.168.200}\def\srvH{10}
  \def\ssid{ArchiNova-WiFi}
  \def\AIrows{%
    \hl\ip{192.168.40.25} & \sol{C} & \sol{255.255.255.0} & \sol{192.168.40.0} & \sol{192.168.40.255} & \sol{254} \\ \hline
    \hl\ip{172.16.5.10}   & \sol{B} & \sol{255.255.0.0}   & \sol{172.16.0.0}   & \sol{172.16.255.255} & \sol{65\,534} \\ \hline
    \hl\ip{10.3.2.1}      & \sol{A} & \sol{255.0.0.0}     & \sol{10.0.0.0}     & \sol{10.255.255.255} & \sol{16\,777\,214} \\ \hline}
  \def\apipa{169.254.12.7}
  \def\Cnet{192.168.60}
  \def\Crows{%
    PC-A & \ip{192.168.60.10} & \ip{255.255.255.0} & \ip{192.168.60.254} & — \\
    PC-B & \ip{192.168.60.20} & \ip{255.255.255.0} & \ip{192.168.60.1}   & — \\
    PC-C & \ip{192.168.61.30} & \ip{255.255.255.0} & \ip{192.168.60.254} & — \\
    PC-D & \ip{192.168.60.10} & \ip{255.255.255.0} & \ip{192.168.60.254} & — \\
    SRV-C & \ip{10.0.0.5}     & \ip{255.0.0.0}     & \ip{10.0.0.254}     & — \\
    R1 — G0/0 & \ip{192.168.60.254} & \ip{255.255.255.0} & — & On \\
    R1 — G0/1 & \ip{10.0.0.254}     & \ip{255.0.0.0}     & — & Off \\}
  \def\Csol{%
    \hl 1 & \sol{PC-B} & \sol{Passerelle .60.1 : pas l'adresse du routeur.} & \sol{Joint son réseau, mais aucun autre (serveur).} & \sol{Passerelle 192.168.60.254.} \\ \hline
    \hl 2 & \sol{PC-C} & \sol{IP hors du réseau 192.168.60.0/24.} & \sol{Ne joint ni les postes ni le routeur.} & \sol{Ex. 192.168.60.30.} \\ \hline
    \hl 3 & \sol{PC-D} & \sol{Même IP que PC-A (doublon).} & \sol{Conflit d'adresse : PC-A et PC-D perturbés.} & \sol{Adresse unique, ex. .60.40.} \\ \hline
    \hl 4 & \sol{R1 G0/1} & \sol{Interface désactivée (Off).} & \sol{Réseau du serveur injoignable pour tous.} & \sol{Port Status : On.} \\ \hline}
  \def\Cbareme{PC-A et SRV-C sont correctement configurés ; accepter PC-A cité à la place de PC-D pour le doublon.}
\else
  \def\entreprise{La clinique vétérinaire \textbf{VetoPlus}}
  \def\Rnom{R-VetoPlus}
  \def\resUn{Soins}\def\pUn{SOI}\def\netUn{192.168.70}
  \def\resDeux{Secrétariat}\def\pDeux{SEC}\def\netDeux{172.20}
  \def\srvNom{SRV-Dossiers}\def\netSrv{192.168.250}\def\srvH{20}
  \def\ssid{VetoPlus-WiFi}
  \def\AIrows{%
    \hl\ip{192.168.15.130} & \sol{C} & \sol{255.255.255.0} & \sol{192.168.15.0} & \sol{192.168.15.255} & \sol{254} \\ \hline
    \hl\ip{172.31.200.7}   & \sol{B} & \sol{255.255.0.0}   & \sol{172.31.0.0}   & \sol{172.31.255.255} & \sol{65\,534} \\ \hline
    \hl\ip{10.45.6.9}      & \sol{A} & \sol{255.0.0.0}     & \sol{10.0.0.0}     & \sol{10.255.255.255} & \sol{16\,777\,214} \\ \hline}
  \def\apipa{169.254.88.3}
  \def\Cnet{192.168.80}
  \def\Crows{%
    PC-A & \ip{192.168.8.11}   & \ip{255.255.255.0} & \ip{192.168.80.254} & — \\
    PC-B & \ip{192.168.80.12}  & \ip{255.255.255.0} & \ip{192.168.80.254} & — \\
    PC-C & \ip{192.168.80.13}  & \ip{255.255.255.0} & \ip{192.168.80.100} & — \\
    PC-D & \ip{192.168.80.12}  & \ip{255.255.255.0} & \ip{192.168.80.254} & — \\
    SRV-C & \ip{10.0.0.8}      & \ip{255.0.0.0}     & \ip{10.0.0.254}     & — \\
    R1 — G0/0 & \ip{192.168.80.254} & \ip{255.255.255.0} & — & Off \\
    R1 — G0/1 & \ip{10.0.0.254}     & \ip{255.0.0.0}     & — & On \\}
  \def\Csol{%
    \hl 1 & \sol{PC-A} & \sol{IP hors du réseau 192.168.80.0/24.} & \sol{Ne joint ni les postes ni le routeur.} & \sol{Ex. 192.168.80.11.} \\ \hline
    \hl 2 & \sol{PC-C} & \sol{Passerelle .80.100 : pas l'adresse du routeur.} & \sol{Joint son réseau, mais aucun autre (serveur).} & \sol{Passerelle 192.168.80.254.} \\ \hline
    \hl 3 & \sol{PC-D} & \sol{Même IP que PC-B (doublon).} & \sol{Conflit d'adresse : PC-B et PC-D perturbés.} & \sol{Adresse unique, ex. .80.14.} \\ \hline
    \hl 4 & \sol{R1 G0/0} & \sol{Interface côté LAN désactivée (Off).} & \sol{Aucun poste ne joint le routeur ni le serveur.} & \sol{Port Status : On.} \\ \hline}
  \def\Cbareme{PC-B et SRV-C sont correctement configurés ; accepter PC-B cité à la place de PC-D pour le doublon.}
\fi

\pagestyle{fancy}
\fancyhf{}
\renewcommand{\headrulewidth}{0pt}
\fancyfoot[L]{\footnotesize\itshape\color{grisfoot}Évaluation — Réseaux locaux — Sujet \variante\ifcorrige\ — CORRIGÉ\fi}
\fancyfoot[R]{\footnotesize\color{grisfoot}\thepage/3}

\setlist{itemsep=0pt,topsep=2pt,parsep=0pt}
\setlength{\parindent}{0pt}
\setlength{\parskip}{2pt}

\begin{document}

\begin{center}
  {\Large\bfseries Évaluation pratique — Réseaux locaux \quad
   \fcolorbox{black}{white}{\ SUJET \variante\ }\par}
  \vspace{1pt}
  {\bfseries Adressage IP, commutation et routage — BTS CIEL 1re année — durée : 2 heures\par}
  \ifcorrige\vspace{2pt}{\bfseries\color{corr}CORRIGÉ ET BARÈME — document enseignant}\par\fi
\end{center}
\vspace{-2pt}
{\renewcommand{\arraystretch}{1.35}
\begin{tabularx}{\linewidth}{|X|X|p{2.2cm}|p{2.6cm}|}
  \hline
  \textbf{Nom :} & \textbf{Prénom :} & \textbf{Poste n°} & \textbf{Note :} \hfill /20 \\ \hline
\end{tabularx}}

\textbf{Consignes :} travail \textbf{individuel}, documents et Internet \textbf{non autorisés}. Ce sujet est le
\textbf{document-réponse}. Fichier Packet Tracer : \ip{Nom\_Prenom\_EvalLAN\_\variante.pkt}, à enregistrer
régulièrement dans le dossier de rendu. Les « visas » sont apposés par le professeur, que vous appelez à chaque test réussi.
Durées conseillées : A 30 min — B 1 h 10 — C 20 min.

%==========================================================================
\section{Partie A — Questions de cours et calculs \hfill /6}

\question{A1}{2 pts}{Complétez le tableau en considérant le \textbf{masque par défaut} de la classe de chaque adresse.}
{\small\renewcommand{\arraystretch}{1.15}
\begin{tabularx}{\linewidth}{|l|c|C|C|C|c|}
  \hline
  \textbf{Adresse IP} & \textbf{Classe} & \textbf{Masque par défaut} & \textbf{Adresse réseau} &
  \textbf{Adresse de diffusion} & \textbf{Nb d'hôtes} \\ \hline
  \AIrows
\end{tabularx}}
\bareme{0,5 pt par ligne juste + 0,5 pt si les trois nombres d'hôtes sont justes ($2^{n}-2$, $n$ = nombre de bits d'hôte).}

\question{A2}{1 pt}{Indiquez le type de câble (\textbf{droit} ou \textbf{croisé}) pour chaque liaison.}
{\small\renewcommand{\arraystretch}{1.15}
\begin{tabularx}{\linewidth}{|X|c||X|c|}
  \hline
  \hl PC $\leftrightarrow$ commutateur & \makebox[1.6cm]{\sol{droit}} & Commutateur $\leftrightarrow$ routeur & \makebox[1.6cm]{\sol{droit}} \\ \hline
  \hl PC $\leftrightarrow$ PC (réseau P2P) & \sol{croisé} & Concentrateur $\leftrightarrow$ concentrateur & \sol{croisé} \\ \hline
\end{tabularx}}
\bareme{0,25 pt par case. Équipements de types différents : câble droit ; de même type : câble croisé.}

\question{A3}{1 pt}{Comment un \textbf{concentrateur (hub)} et un \textbf{commutateur (switch)} traitent-ils une trame
reçue ? Lequel limite les collisions, et pourquoi ?}
\zone{2.9cm}{Le \textbf{hub} (couche 1) répète la trame sur \textbf{tous ses autres ports} : un seul domaine de collision
partagé. Le \textbf{switch} (couche 2) lit l'adresse MAC de destination et, grâce à sa \textbf{table MAC}, ne l'envoie
\textbf{que sur le port du destinataire} : un domaine de collision par port, c'est lui qui limite les collisions.
\emph{0,5 pt hub, 0,5 pt switch.}}

\question{A4}{1 pt}{Qu'est-ce que la \textbf{passerelle par défaut} d'un poste ? Quelle condition son adresse doit-elle respecter ?}
\zone{2.4cm}{Adresse IP de l'interface du \textbf{routeur} à laquelle le poste envoie les paquets destinés à un
\textbf{autre réseau}. Elle doit \textbf{appartenir au même réseau} que le poste. \emph{0,5 + 0,5 pt.}}

\question{A5}{0,5 pt}{Un portable configuré en DHCP affiche \ip{\apipa} / \ip{255.255.0.0}. Que s'est-il passé ?}
\zone{1.5cm}{Aucun \textbf{serveur DHCP} n'a répondu : le poste s'est attribué une adresse \textbf{APIPA} (169.254.0.0/16).}

\question{A6}{0,5 pt}{Quel protocole utilise \texttt{ping} ? Nommez les deux messages échangés et la couche OSI des adresses IP.}
\zone{1.5cm}{\textbf{ICMP} ; \textbf{Echo Request} / \textbf{Echo Reply} ; adresses IP en \textbf{couche 3 (Réseau)}.}

%==========================================================================
\newpage
\section{Partie B — Réalisation sous Packet Tracer \hfill /10}

\begin{cdc}
\entreprise{} vous confie son réseau : \textbf{trois réseaux} reliés par un \textbf{routeur unique Cisco 2911}
nommé \ip{\Rnom} (interfaces G0/0, G0/1, G0/2).
\begin{itemize}[leftmargin=12pt]
  \item \textbf{Réseau \resUn} : \textbf{4 postes} PC-\pUn1 à PC-\pUn4, commutateur \ip{SW-\pUn}, adressage statique,
        réseau \ip{\netUn.0/24}.
  \item \textbf{Réseau \resDeux} : \textbf{2 postes} PC-\pDeux1, PC-\pDeux2 et \textbf{1 imprimante} IMP-\pDeux,
        commutateur \ip{SW-\pDeux}, adressage statique, réseau de classe B \ip{\netDeux.0.0/16}.
  \item \textbf{Réseau Serveur} : serveur \ip{\srvNom} d'adresse \ip{\netSrv.\srvH/24}, relié \textbf{directement} au routeur.
  \item La passerelle de chaque réseau est sa \textbf{dernière adresse utilisable}. Tous les équipements communiquent
        entre eux et accèdent à la page web du serveur. Nommage imposé (« Display Name »).
\end{itemize}
\end{cdc}

\question{B1}{2 pts}{Complétez le plan d'adressage \textbf{avant} la réalisation.}
{\footnotesize\setlength{\tabcolsep}{3pt}\renewcommand{\arraystretch}{1.15}
\begin{tabularx}{\linewidth}{|l|C|C|C|C|C|C|}
  \hline
  \textbf{Réseau} & \textbf{Adresse réseau} & \textbf{Masque} & \textbf{1re @ utilisable} &
  \textbf{Dernière @ utilisable} & \textbf{Diffusion} & \textbf{Passerelle} \\ \hline
  \hl \resUn & \sol{\netUn.0} & \sol{255.255.255.0} & \sol{\netUn.1} & \sol{\netUn.254} & \sol{\netUn.255} & \sol{\netUn.254} \\ \hline
  \hl \resDeux & \sol{\netDeux.0.0} & \sol{255.255.0.0} & \sol{\netDeux.0.1} & \sol{\netDeux.255.254} & \sol{\netDeux.255.255} & \sol{\netDeux.255.254} \\ \hline
  \hl Serveur & \sol{\netSrv.0} & \sol{255.255.255.0} & \sol{\netSrv.1} & \sol{\netSrv.254} & \sol{\netSrv.255} & \sol{\netSrv.254} \\ \hline
\end{tabularx}}
\bareme{0,5 pt par ligne juste + 0,5 pt si la convention de passerelle est respectée partout.}

\question{B2}{2 pts}{Construisez la topologie (équipements, câblage, nommage) et relevez l'adressage configuré.}
{\footnotesize\setlength{\tabcolsep}{3pt}\renewcommand{\arraystretch}{1.1}
\begin{tabularx}{\linewidth}{|l|C|C|C||l|C|C|}
  \hline
  \textbf{Hôte} & \textbf{Adresse IP} & \textbf{Masque} & \textbf{Passerelle} & \textbf{Interface} & \textbf{Adresse IP} & \textbf{Masque} \\ \hline
  \hl PC-\pUn1 & \sol{\netUn.1} & \sol{255.255.255.0} & \sol{\netUn.254} & \Rnom{} G0/0 & \sol{\netUn.254} & \sol{255.255.255.0} \\ \hline
  \hl PC-\pDeux1 & \sol{\netDeux.0.1} & \sol{255.255.0.0} & \sol{\netDeux.255.254} & \Rnom{} G0/1 & \sol{\netDeux.255.254} & \sol{255.255.0.0} \\ \hline
  \hl IMP-\pDeux & \sol{\netDeux.0.10} & \sol{255.255.0.0} & \sol{\netDeux.255.254} & \Rnom{} G0/2 & \sol{\netSrv.254} & \sol{255.255.255.0} \\ \hline
  \hl \srvNom & \ip{\netSrv.\srvH} & \sol{255.255.255.0} & \sol{\netSrv.254} & \multicolumn{3}{c|}{\cellcolor{black!8}} \\ \hline
\end{tabularx}}
\bareme{Adresses d'hôtes libres (uniques, dans le bon réseau) ; affectation G0/x libre. Sur le .pkt : équipements 0,5 ;
câblage 1 (câbles droits vers les switchs et switch–routeur, \textbf{croisé} serveur–routeur, liens verts) ; nommage 0,5.}

\question{B3}{3 pts}{Configurez tous les équipements : IP, masques, passerelles des hôtes ; interfaces du routeur
(« Port Status : On »).}
\bareme{Sur le .pkt : interfaces routeur 1 pt ; IP/masques des hôtes 1 pt ; passerelles (imprimante et serveur compris)
1 pt. $-0,25$ par erreur, sans descendre sous 0 par item.}

\question{B4}{3 pts}{Réalisez les tests, notez la commande et le résultat, puis faites viser chaque test réussi.}
{\small\renewcommand{\arraystretch}{1.1}
\begin{tabularx}{\linewidth}{|c|l|l|C|c|c|}
  \hline
  \textbf{N°} & \textbf{Source} & \textbf{Destination} & \textbf{Commande / action} & \textbf{Résultat} & \textbf{Visa prof.} \\ \hline
  \hl 1 & PC-\pUn1 & PC-\pUn4 & \sol{ping \netUn.4} & \sol{réussi} & \hspace{1.4cm} \\ \hline
  \hl 2 & PC-\pDeux1 & IMP-\pDeux & \sol{ping \netDeux.0.10} & \sol{réussi} & \\ \hline
  \hl 3 & PC-\pUn2 & PC-\pDeux2 & \sol{ping \netDeux.0.2} & \sol{réussi} & \\ \hline
  \hl 4 & PC-\pDeux2 & \srvNom & \sol{ping \netSrv.\srvH} & \sol{réussi} & \\ \hline
  \hl 5 & PC-\pUn3 & page web du serveur & \sol{navigateur : {\NoAutoSpacing http://\netSrv.\srvH}} & \sol{page affichée} & \\ \hline
\end{tabularx}}
\bareme{0,5 pt par test visé + 0,5 pt pour B5. Le premier ping inter-réseaux peut échouer (résolution ARP) ; les suivants doivent réussir.}

\question{B5}{incl. B4}{Test n°3 : par quels équipements passe le paquet ? Pourquoi PC-\pUn2 l'envoie-t-il au routeur ?}
\zone{2.3cm}{PC-\pUn2 $\to$ SW-\pUn{} $\to$ \Rnom{} $\to$ SW-\pDeux{} $\to$ PC-\pDeux2. Grâce à son masque, PC-\pUn2 constate que
la destination (\netDeux.0.0/16) n'est \textbf{pas dans son réseau} (\netUn.0/24) : il envoie le paquet à sa
\textbf{passerelle par défaut}, qui le route.}

%==========================================================================
\newpage
\section{Partie C — Diagnostic d'un réseau en panne \hfill /4}

Un technicien constate qu'\textbf{aucun poste ne parvient à joindre le serveur SRV-C} et que certains postes
communiquent mal entre eux. Les câbles sont du bon type.

\begin{center}
\begin{minipage}{0.37\linewidth}\centering
\begin{tikzpicture}[font=\footnotesize,
  eq/.style={draw,fill=qbg,rounded corners=2pt,minimum width=9mm,minimum height=5mm,inner sep=2pt},
  every path/.style={thick}]
  \node[eq] (SW) {SW-1};
  \node[eq,right=9mm of SW] (R) {R1};
  \node[eq,right=7mm of R] (SRV) {SRV-C};
  \node[eq,above left=5mm and -3mm of SW] (A) {PC-A};
  \node[eq,above right=5mm and -3mm of SW] (B) {PC-B};
  \node[eq,below left=5mm and -3mm of SW] (C) {PC-C};
  \node[eq,below right=5mm and -3mm of SW] (D) {PC-D};
  \foreach \n in {A,B,C,D} \draw (\n) -- (SW);
  \draw (SW) -- node[above,font=\tiny]{G0/0} (R);
  \draw (R) -- node[above,font=\tiny]{G0/1} (SRV);
\end{tikzpicture}
\end{minipage}\hfill
\begin{minipage}{0.61\linewidth}\centering
{\small\renewcommand{\arraystretch}{1.05}\setlength{\tabcolsep}{4pt}
\begin{tabular}{lllll}
  \toprule
  \textbf{Équipement} & \textbf{Adresse IP} & \textbf{Masque} & \textbf{Passerelle} & \textbf{Port} \\
  \midrule
  \Crows
  \bottomrule
\end{tabular}}
\end{minipage}
\end{center}

\question{C1}{4 pts}{Relevez les \textbf{quatre anomalies}. Pour chacune : conséquence et correction.}
{\small\renewcommand{\arraystretch}{1.1}
\ifcorrige\def\hc{}\else\def\hc{\rule[-10mm]{0pt}{14mm}}\fi
\begin{tabularx}{\linewidth}{|c|>{\raggedright\arraybackslash}p{2cm}|L|L|L|}
  \hline
  \textbf{N°} & \textbf{Équipement} & \textbf{Anomalie} & \textbf{Conséquence} & \textbf{Correction} \\ \hline
  \ifcorrige\Csol\else
  \hc 1 & & & & \\ \hline
  \hc 2 & & & & \\ \hline
  \hc 3 & & & & \\ \hline
  \hc 4 & & & & \\ \hline
  \fi
\end{tabularx}}
\bareme{1 pt par anomalie (0,25 équipement, anomalie, conséquence, correction ; ordre indifférent). \Cbareme}

%==========================================================================
\section{Bonus — Extension sans fil \hfill +1}

\begin{cdc}Ajoutez au réseau \resUn{} un \textbf{AccessPoint-PT} relié à SW-\pUn{} (SSID \ip{\ssid}, \textbf{WPA2-PSK})
et un portable \ip{Laptop-\pUn} équipé d'un module \texttt{WPC300N}, en adressage \textbf{statique}.
Vérifiez qu'il joint \srvNom.\end{cdc}
{\small\renewcommand{\arraystretch}{1.1}
\begin{tabularx}{\linewidth}{|l|C|C|C|c|}
  \hline
  \textbf{Équipement} & \textbf{Adresse IP} & \textbf{Masque} & \textbf{Passerelle} & \textbf{Visa prof.} \\ \hline
  \hl Laptop-\pUn & \sol{\netUn.50} & \sol{255.255.255.0} & \sol{\netUn.254} & \hspace{1.4cm} \\ \hline
\end{tabularx}}
\bareme{Le point d'accès fait office de pont : le portable est dans \netUn.0/24. Bonus si associé en WPA2-PSK et ping
vers \netSrv.\srvH{} réussi. Note plafonnée à 20.}

\end{document}
"
\documentclass[10pt,a4paper]{article}

\usepackage[T1]{fontenc}
\usepackage[utf8]{inputenc}
\usepackage{lmodern}
\usepackage[french]{babel}
\usepackage[a4paper,hmargin=1.5cm,top=1.3cm,bottom=1.6cm,footskip=0.7cm]{geometry}
\usepackage[table]{xcolor}
\usepackage{booktabs,tabularx,array}
\usepackage{enumitem}
\usepackage{fancyhdr}
\usepackage[most]{tcolorbox}
\usepackage{tikz}
\usetikzlibrary{positioning}
\usepackage{titlesec}

\providecommand{\modecorrige}{0}
\providecommand{\variante}{A}
\newif\ifcorrige
\ifnum\modecorrige=1 \corrigetrue \else \corrigefalse \fi

\definecolor{ipred}{HTML}{C0392B}
\definecolor{cdcbar}{HTML}{A67C3D}
\definecolor{cdcbg}{HTML}{F4F1EA}
\definecolor{qbar}{HTML}{4A6B5D}
\definecolor{qbg}{HTML}{EDF1EF}
\definecolor{corr}{HTML}{1F5FA8}
\definecolor{corrbg}{HTML}{E8F0FA}
\definecolor{grisfoot}{HTML}{828282}

\newcommand{\ip}[1]{\textcolor{ipred}{\texttt{#1}}}
\newcommand{\pts}[1]{\hfill\mbox{\textbf{(#1)}}}
\newcommand{\sol}[1]{\ifcorrige\textcolor{corr}{#1}\fi}
\newcommand{\bareme}[1]{\ifcorrige\par{\footnotesize\color{corr}#1\par}\fi}
\newcommand{\hl}{\ifcorrige\rule[-1.6mm]{0pt}{5.4mm}\else\rule[-2.2mm]{0pt}{6.6mm}\fi}
\newcolumntype{C}{>{\centering\arraybackslash}X}
\newcolumntype{L}{>{\raggedright\arraybackslash}X}

\titleformat{\section}{\large\bfseries}{}{0pt}{}
\titlespacing*{\section}{0pt}{6pt}{3pt}

\tcbset{qstyle/.style={enhanced,boxrule=0pt,leftrule=2.5pt,arc=0pt,outer arc=0pt,
  left=5pt,right=5pt,top=2pt,bottom=2pt,before skip=4pt,after skip=3pt}}
\newtcolorbox{cdc}{qstyle,colback=cdcbg,colframe=cdcbar}
\newtcolorbox{quest}{qstyle,colback=qbg,colframe=qbar}

% Zone de réponse : vierge dans le sujet, remplie dans le corrigé
\newcommand{\zone}[2]{%
  \ifcorrige
    \begin{tcolorbox}[colback=corrbg,colframe=corr,boxrule=0.5pt,arc=1pt,left=4pt,right=4pt,top=1pt,bottom=1pt,
      before skip=2pt,after skip=2pt,fontupper=\footnotesize\color{corr}]#2\end{tcolorbox}%
  \else
    \begin{tcolorbox}[colback=white,colframe=black!45,boxrule=0.4pt,arc=1pt,height=#1,
      before skip=2pt,after skip=2pt]\end{tcolorbox}%
  \fi}

\newcommand{\question}[3]{\begin{quest}\textbf{#1.} #3\pts{#2}\end{quest}}

%==========================================================================
% Données propres à chaque variante
%==========================================================================
\if\variante A
  \def\entreprise{Le cabinet d'architectes \textbf{ArchiNova}}
  \def\Rnom{R-ArchiNova}
  \def\resUn{Bureau d'études}\def\pUn{BE}\def\netUn{192.168.40}
  \def\resDeux{Accueil}\def\pDeux{ACC}\def\netDeux{172.16}
  \def\srvNom{SRV-Intranet}\def\netSrv{192.168.200}\def\srvH{10}
  \def\ssid{ArchiNova-WiFi}
  \def\AIrows{%
    \hl\ip{192.168.40.25} & \sol{C} & \sol{255.255.255.0} & \sol{192.168.40.0} & \sol{192.168.40.255} & \sol{254} \\ \hline
    \hl\ip{172.16.5.10}   & \sol{B} & \sol{255.255.0.0}   & \sol{172.16.0.0}   & \sol{172.16.255.255} & \sol{65\,534} \\ \hline
    \hl\ip{10.3.2.1}      & \sol{A} & \sol{255.0.0.0}     & \sol{10.0.0.0}     & \sol{10.255.255.255} & \sol{16\,777\,214} \\ \hline}
  \def\apipa{169.254.12.7}
  \def\Cnet{192.168.60}
  \def\Crows{%
    PC-A & \ip{192.168.60.10} & \ip{255.255.255.0} & \ip{192.168.60.254} & — \\
    PC-B & \ip{192.168.60.20} & \ip{255.255.255.0} & \ip{192.168.60.1}   & — \\
    PC-C & \ip{192.168.61.30} & \ip{255.255.255.0} & \ip{192.168.60.254} & — \\
    PC-D & \ip{192.168.60.10} & \ip{255.255.255.0} & \ip{192.168.60.254} & — \\
    SRV-C & \ip{10.0.0.5}     & \ip{255.0.0.0}     & \ip{10.0.0.254}     & — \\
    R1 — G0/0 & \ip{192.168.60.254} & \ip{255.255.255.0} & — & On \\
    R1 — G0/1 & \ip{10.0.0.254}     & \ip{255.0.0.0}     & — & Off \\}
  \def\Csol{%
    \hl 1 & \sol{PC-B} & \sol{Passerelle .60.1 : pas l'adresse du routeur.} & \sol{Joint son réseau, mais aucun autre (serveur).} & \sol{Passerelle 192.168.60.254.} \\ \hline
    \hl 2 & \sol{PC-C} & \sol{IP hors du réseau 192.168.60.0/24.} & \sol{Ne joint ni les postes ni le routeur.} & \sol{Ex. 192.168.60.30.} \\ \hline
    \hl 3 & \sol{PC-D} & \sol{Même IP que PC-A (doublon).} & \sol{Conflit d'adresse : PC-A et PC-D perturbés.} & \sol{Adresse unique, ex. .60.40.} \\ \hline
    \hl 4 & \sol{R1 G0/1} & \sol{Interface désactivée (Off).} & \sol{Réseau du serveur injoignable pour tous.} & \sol{Port Status : On.} \\ \hline}
  \def\Cbareme{PC-A et SRV-C sont correctement configurés ; accepter PC-A cité à la place de PC-D pour le doublon.}
\else
  \def\entreprise{La clinique vétérinaire \textbf{VetoPlus}}
  \def\Rnom{R-VetoPlus}
  \def\resUn{Soins}\def\pUn{SOI}\def\netUn{192.168.70}
  \def\resDeux{Secrétariat}\def\pDeux{SEC}\def\netDeux{172.20}
  \def\srvNom{SRV-Dossiers}\def\netSrv{192.168.250}\def\srvH{20}
  \def\ssid{VetoPlus-WiFi}
  \def\AIrows{%
    \hl\ip{192.168.15.130} & \sol{C} & \sol{255.255.255.0} & \sol{192.168.15.0} & \sol{192.168.15.255} & \sol{254} \\ \hline
    \hl\ip{172.31.200.7}   & \sol{B} & \sol{255.255.0.0}   & \sol{172.31.0.0}   & \sol{172.31.255.255} & \sol{65\,534} \\ \hline
    \hl\ip{10.45.6.9}      & \sol{A} & \sol{255.0.0.0}     & \sol{10.0.0.0}     & \sol{10.255.255.255} & \sol{16\,777\,214} \\ \hline}
  \def\apipa{169.254.88.3}
  \def\Cnet{192.168.80}
  \def\Crows{%
    PC-A & \ip{192.168.8.11}   & \ip{255.255.255.0} & \ip{192.168.80.254} & — \\
    PC-B & \ip{192.168.80.12}  & \ip{255.255.255.0} & \ip{192.168.80.254} & — \\
    PC-C & \ip{192.168.80.13}  & \ip{255.255.255.0} & \ip{192.168.80.100} & — \\
    PC-D & \ip{192.168.80.12}  & \ip{255.255.255.0} & \ip{192.168.80.254} & — \\
    SRV-C & \ip{10.0.0.8}      & \ip{255.0.0.0}     & \ip{10.0.0.254}     & — \\
    R1 — G0/0 & \ip{192.168.80.254} & \ip{255.255.255.0} & — & Off \\
    R1 — G0/1 & \ip{10.0.0.254}     & \ip{255.0.0.0}     & — & On \\}
  \def\Csol{%
    \hl 1 & \sol{PC-A} & \sol{IP hors du réseau 192.168.80.0/24.} & \sol{Ne joint ni les postes ni le routeur.} & \sol{Ex. 192.168.80.11.} \\ \hline
    \hl 2 & \sol{PC-C} & \sol{Passerelle .80.100 : pas l'adresse du routeur.} & \sol{Joint son réseau, mais aucun autre (serveur).} & \sol{Passerelle 192.168.80.254.} \\ \hline
    \hl 3 & \sol{PC-D} & \sol{Même IP que PC-B (doublon).} & \sol{Conflit d'adresse : PC-B et PC-D perturbés.} & \sol{Adresse unique, ex. .80.14.} \\ \hline
    \hl 4 & \sol{R1 G0/0} & \sol{Interface côté LAN désactivée (Off).} & \sol{Aucun poste ne joint le routeur ni le serveur.} & \sol{Port Status : On.} \\ \hline}
  \def\Cbareme{PC-B et SRV-C sont correctement configurés ; accepter PC-B cité à la place de PC-D pour le doublon.}
\fi

\pagestyle{fancy}
\fancyhf{}
\renewcommand{\headrulewidth}{0pt}
\fancyfoot[L]{\footnotesize\itshape\color{grisfoot}Évaluation — Réseaux locaux — Sujet \variante\ifcorrige\ — CORRIGÉ\fi}
\fancyfoot[R]{\footnotesize\color{grisfoot}\thepage/3}

\setlist{itemsep=0pt,topsep=2pt,parsep=0pt}
\setlength{\parindent}{0pt}
\setlength{\parskip}{2pt}

\begin{document}

\begin{center}
  {\Large\bfseries Évaluation pratique — Réseaux locaux \quad
   \fcolorbox{black}{white}{\ SUJET \variante\ }\par}
  \vspace{1pt}
  {\bfseries Adressage IP, commutation et routage — BTS CIEL 1re année — durée : 2 heures\par}
  \ifcorrige\vspace{2pt}{\bfseries\color{corr}CORRIGÉ ET BARÈME — document enseignant}\par\fi
\end{center}
\vspace{-2pt}
{\renewcommand{\arraystretch}{1.35}
\begin{tabularx}{\linewidth}{|X|X|p{2.2cm}|p{2.6cm}|}
  \hline
  \textbf{Nom :} & \textbf{Prénom :} & \textbf{Poste n°} & \textbf{Note :} \hfill /20 \\ \hline
\end{tabularx}}

\textbf{Consignes :} travail \textbf{individuel}, documents et Internet \textbf{non autorisés}. Ce sujet est le
\textbf{document-réponse}. Fichier Packet Tracer : \ip{Nom\_Prenom\_EvalLAN\_\variante.pkt}, à enregistrer
régulièrement dans le dossier de rendu. Les « visas » sont apposés par le professeur, que vous appelez à chaque test réussi.
Durées conseillées : A 30 min — B 1 h 10 — C 20 min.

%==========================================================================
\section{Partie A — Questions de cours et calculs \hfill /6}

\question{A1}{2 pts}{Complétez le tableau en considérant le \textbf{masque par défaut} de la classe de chaque adresse.}
{\small\renewcommand{\arraystretch}{1.15}
\begin{tabularx}{\linewidth}{|l|c|C|C|C|c|}
  \hline
  \textbf{Adresse IP} & \textbf{Classe} & \textbf{Masque par défaut} & \textbf{Adresse réseau} &
  \textbf{Adresse de diffusion} & \textbf{Nb d'hôtes} \\ \hline
  \AIrows
\end{tabularx}}
\bareme{0,5 pt par ligne juste + 0,5 pt si les trois nombres d'hôtes sont justes ($2^{n}-2$, $n$ = nombre de bits d'hôte).}

\question{A2}{1 pt}{Indiquez le type de câble (\textbf{droit} ou \textbf{croisé}) pour chaque liaison.}
{\small\renewcommand{\arraystretch}{1.15}
\begin{tabularx}{\linewidth}{|X|c||X|c|}
  \hline
  \hl PC $\leftrightarrow$ commutateur & \makebox[1.6cm]{\sol{droit}} & Commutateur $\leftrightarrow$ routeur & \makebox[1.6cm]{\sol{droit}} \\ \hline
  \hl PC $\leftrightarrow$ PC (réseau P2P) & \sol{croisé} & Concentrateur $\leftrightarrow$ concentrateur & \sol{croisé} \\ \hline
\end{tabularx}}
\bareme{0,25 pt par case. Équipements de types différents : câble droit ; de même type : câble croisé.}

\question{A3}{1 pt}{Comment un \textbf{concentrateur (hub)} et un \textbf{commutateur (switch)} traitent-ils une trame
reçue ? Lequel limite les collisions, et pourquoi ?}
\zone{2.9cm}{Le \textbf{hub} (couche 1) répète la trame sur \textbf{tous ses autres ports} : un seul domaine de collision
partagé. Le \textbf{switch} (couche 2) lit l'adresse MAC de destination et, grâce à sa \textbf{table MAC}, ne l'envoie
\textbf{que sur le port du destinataire} : un domaine de collision par port, c'est lui qui limite les collisions.
\emph{0,5 pt hub, 0,5 pt switch.}}

\question{A4}{1 pt}{Qu'est-ce que la \textbf{passerelle par défaut} d'un poste ? Quelle condition son adresse doit-elle respecter ?}
\zone{2.4cm}{Adresse IP de l'interface du \textbf{routeur} à laquelle le poste envoie les paquets destinés à un
\textbf{autre réseau}. Elle doit \textbf{appartenir au même réseau} que le poste. \emph{0,5 + 0,5 pt.}}

\question{A5}{0,5 pt}{Un portable configuré en DHCP affiche \ip{\apipa} / \ip{255.255.0.0}. Que s'est-il passé ?}
\zone{1.5cm}{Aucun \textbf{serveur DHCP} n'a répondu : le poste s'est attribué une adresse \textbf{APIPA} (169.254.0.0/16).}

\question{A6}{0,5 pt}{Quel protocole utilise \texttt{ping} ? Nommez les deux messages échangés et la couche OSI des adresses IP.}
\zone{1.5cm}{\textbf{ICMP} ; \textbf{Echo Request} / \textbf{Echo Reply} ; adresses IP en \textbf{couche 3 (Réseau)}.}

%==========================================================================
\newpage
\section{Partie B — Réalisation sous Packet Tracer \hfill /10}

\begin{cdc}
\entreprise{} vous confie son réseau : \textbf{trois réseaux} reliés par un \textbf{routeur unique Cisco 2911}
nommé \ip{\Rnom} (interfaces G0/0, G0/1, G0/2).
\begin{itemize}[leftmargin=12pt]
  \item \textbf{Réseau \resUn} : \textbf{4 postes} PC-\pUn1 à PC-\pUn4, commutateur \ip{SW-\pUn}, adressage statique,
        réseau \ip{\netUn.0/24}.
  \item \textbf{Réseau \resDeux} : \textbf{2 postes} PC-\pDeux1, PC-\pDeux2 et \textbf{1 imprimante} IMP-\pDeux,
        commutateur \ip{SW-\pDeux}, adressage statique, réseau de classe B \ip{\netDeux.0.0/16}.
  \item \textbf{Réseau Serveur} : serveur \ip{\srvNom} d'adresse \ip{\netSrv.\srvH/24}, relié \textbf{directement} au routeur.
  \item La passerelle de chaque réseau est sa \textbf{dernière adresse utilisable}. Tous les équipements communiquent
        entre eux et accèdent à la page web du serveur. Nommage imposé (« Display Name »).
\end{itemize}
\end{cdc}

\question{B1}{2 pts}{Complétez le plan d'adressage \textbf{avant} la réalisation.}
{\footnotesize\setlength{\tabcolsep}{3pt}\renewcommand{\arraystretch}{1.15}
\begin{tabularx}{\linewidth}{|l|C|C|C|C|C|C|}
  \hline
  \textbf{Réseau} & \textbf{Adresse réseau} & \textbf{Masque} & \textbf{1re @ utilisable} &
  \textbf{Dernière @ utilisable} & \textbf{Diffusion} & \textbf{Passerelle} \\ \hline
  \hl \resUn & \sol{\netUn.0} & \sol{255.255.255.0} & \sol{\netUn.1} & \sol{\netUn.254} & \sol{\netUn.255} & \sol{\netUn.254} \\ \hline
  \hl \resDeux & \sol{\netDeux.0.0} & \sol{255.255.0.0} & \sol{\netDeux.0.1} & \sol{\netDeux.255.254} & \sol{\netDeux.255.255} & \sol{\netDeux.255.254} \\ \hline
  \hl Serveur & \sol{\netSrv.0} & \sol{255.255.255.0} & \sol{\netSrv.1} & \sol{\netSrv.254} & \sol{\netSrv.255} & \sol{\netSrv.254} \\ \hline
\end{tabularx}}
\bareme{0,5 pt par ligne juste + 0,5 pt si la convention de passerelle est respectée partout.}

\question{B2}{2 pts}{Construisez la topologie (équipements, câblage, nommage) et relevez l'adressage configuré.}
{\footnotesize\setlength{\tabcolsep}{3pt}\renewcommand{\arraystretch}{1.1}
\begin{tabularx}{\linewidth}{|l|C|C|C||l|C|C|}
  \hline
  \textbf{Hôte} & \textbf{Adresse IP} & \textbf{Masque} & \textbf{Passerelle} & \textbf{Interface} & \textbf{Adresse IP} & \textbf{Masque} \\ \hline
  \hl PC-\pUn1 & \sol{\netUn.1} & \sol{255.255.255.0} & \sol{\netUn.254} & \Rnom{} G0/0 & \sol{\netUn.254} & \sol{255.255.255.0} \\ \hline
  \hl PC-\pDeux1 & \sol{\netDeux.0.1} & \sol{255.255.0.0} & \sol{\netDeux.255.254} & \Rnom{} G0/1 & \sol{\netDeux.255.254} & \sol{255.255.0.0} \\ \hline
  \hl IMP-\pDeux & \sol{\netDeux.0.10} & \sol{255.255.0.0} & \sol{\netDeux.255.254} & \Rnom{} G0/2 & \sol{\netSrv.254} & \sol{255.255.255.0} \\ \hline
  \hl \srvNom & \ip{\netSrv.\srvH} & \sol{255.255.255.0} & \sol{\netSrv.254} & \multicolumn{3}{c|}{\cellcolor{black!8}} \\ \hline
\end{tabularx}}
\bareme{Adresses d'hôtes libres (uniques, dans le bon réseau) ; affectation G0/x libre. Sur le .pkt : équipements 0,5 ;
câblage 1 (câbles droits vers les switchs et switch–routeur, \textbf{croisé} serveur–routeur, liens verts) ; nommage 0,5.}

\question{B3}{3 pts}{Configurez tous les équipements : IP, masques, passerelles des hôtes ; interfaces du routeur
(« Port Status : On »).}
\bareme{Sur le .pkt : interfaces routeur 1 pt ; IP/masques des hôtes 1 pt ; passerelles (imprimante et serveur compris)
1 pt. $-0,25$ par erreur, sans descendre sous 0 par item.}

\question{B4}{3 pts}{Réalisez les tests, notez la commande et le résultat, puis faites viser chaque test réussi.}
{\small\renewcommand{\arraystretch}{1.1}
\begin{tabularx}{\linewidth}{|c|l|l|C|c|c|}
  \hline
  \textbf{N°} & \textbf{Source} & \textbf{Destination} & \textbf{Commande / action} & \textbf{Résultat} & \textbf{Visa prof.} \\ \hline
  \hl 1 & PC-\pUn1 & PC-\pUn4 & \sol{ping \netUn.4} & \sol{réussi} & \hspace{1.4cm} \\ \hline
  \hl 2 & PC-\pDeux1 & IMP-\pDeux & \sol{ping \netDeux.0.10} & \sol{réussi} & \\ \hline
  \hl 3 & PC-\pUn2 & PC-\pDeux2 & \sol{ping \netDeux.0.2} & \sol{réussi} & \\ \hline
  \hl 4 & PC-\pDeux2 & \srvNom & \sol{ping \netSrv.\srvH} & \sol{réussi} & \\ \hline
  \hl 5 & PC-\pUn3 & page web du serveur & \sol{navigateur : {\NoAutoSpacing http://\netSrv.\srvH}} & \sol{page affichée} & \\ \hline
\end{tabularx}}
\bareme{0,5 pt par test visé + 0,5 pt pour B5. Le premier ping inter-réseaux peut échouer (résolution ARP) ; les suivants doivent réussir.}

\question{B5}{incl. B4}{Test n°3 : par quels équipements passe le paquet ? Pourquoi PC-\pUn2 l'envoie-t-il au routeur ?}
\zone{2.3cm}{PC-\pUn2 $\to$ SW-\pUn{} $\to$ \Rnom{} $\to$ SW-\pDeux{} $\to$ PC-\pDeux2. Grâce à son masque, PC-\pUn2 constate que
la destination (\netDeux.0.0/16) n'est \textbf{pas dans son réseau} (\netUn.0/24) : il envoie le paquet à sa
\textbf{passerelle par défaut}, qui le route.}

%==========================================================================
\newpage
\section{Partie C — Diagnostic d'un réseau en panne \hfill /4}

Un technicien constate qu'\textbf{aucun poste ne parvient à joindre le serveur SRV-C} et que certains postes
communiquent mal entre eux. Les câbles sont du bon type.

\begin{center}
\begin{minipage}{0.37\linewidth}\centering
\begin{tikzpicture}[font=\footnotesize,
  eq/.style={draw,fill=qbg,rounded corners=2pt,minimum width=9mm,minimum height=5mm,inner sep=2pt},
  every path/.style={thick}]
  \node[eq] (SW) {SW-1};
  \node[eq,right=9mm of SW] (R) {R1};
  \node[eq,right=7mm of R] (SRV) {SRV-C};
  \node[eq,above left=5mm and -3mm of SW] (A) {PC-A};
  \node[eq,above right=5mm and -3mm of SW] (B) {PC-B};
  \node[eq,below left=5mm and -3mm of SW] (C) {PC-C};
  \node[eq,below right=5mm and -3mm of SW] (D) {PC-D};
  \foreach \n in {A,B,C,D} \draw (\n) -- (SW);
  \draw (SW) -- node[above,font=\tiny]{G0/0} (R);
  \draw (R) -- node[above,font=\tiny]{G0/1} (SRV);
\end{tikzpicture}
\end{minipage}\hfill
\begin{minipage}{0.61\linewidth}\centering
{\small\renewcommand{\arraystretch}{1.05}\setlength{\tabcolsep}{4pt}
\begin{tabular}{lllll}
  \toprule
  \textbf{Équipement} & \textbf{Adresse IP} & \textbf{Masque} & \textbf{Passerelle} & \textbf{Port} \\
  \midrule
  \Crows
  \bottomrule
\end{tabular}}
\end{minipage}
\end{center}

\question{C1}{4 pts}{Relevez les \textbf{quatre anomalies}. Pour chacune : conséquence et correction.}
{\small\renewcommand{\arraystretch}{1.1}
\ifcorrige\def\hc{}\else\def\hc{\rule[-10mm]{0pt}{14mm}}\fi
\begin{tabularx}{\linewidth}{|c|>{\raggedright\arraybackslash}p{2cm}|L|L|L|}
  \hline
  \textbf{N°} & \textbf{Équipement} & \textbf{Anomalie} & \textbf{Conséquence} & \textbf{Correction} \\ \hline
  \ifcorrige\Csol\else
  \hc 1 & & & & \\ \hline
  \hc 2 & & & & \\ \hline
  \hc 3 & & & & \\ \hline
  \hc 4 & & & & \\ \hline
  \fi
\end{tabularx}}
\bareme{1 pt par anomalie (0,25 équipement, anomalie, conséquence, correction ; ordre indifférent). \Cbareme}

%==========================================================================
\section{Bonus — Extension sans fil \hfill +1}

\begin{cdc}Ajoutez au réseau \resUn{} un \textbf{AccessPoint-PT} relié à SW-\pUn{} (SSID \ip{\ssid}, \textbf{WPA2-PSK})
et un portable \ip{Laptop-\pUn} équipé d'un module \texttt{WPC300N}, en adressage \textbf{statique}.
Vérifiez qu'il joint \srvNom.\end{cdc}
{\small\renewcommand{\arraystretch}{1.1}
\begin{tabularx}{\linewidth}{|l|C|C|C|c|}
  \hline
  \textbf{Équipement} & \textbf{Adresse IP} & \textbf{Masque} & \textbf{Passerelle} & \textbf{Visa prof.} \\ \hline
  \hl Laptop-\pUn & \sol{\netUn.50} & \sol{255.255.255.0} & \sol{\netUn.254} & \hspace{1.4cm} \\ \hline
\end{tabularx}}
\bareme{Le point d'accès fait office de pont : le portable est dans \netUn.0/24. Bonus si associé en WPA2-PSK et ping
vers \netSrv.\srvH{} réussi. Note plafonnée à 20.}

\end{document}
"
%   pdflatex -jobname=Eval_Reseau_CorrigeA "\def\variante{A}\def\modecorrige{1}% Évaluation pratique — Réseaux locaux (BTS CIEL 1re année) — sujets A et B
% Compilation (4 PDF) :
%   pdflatex -jobname=Eval_Reseau_SujetA   "\def\variante{A}\def\modecorrige{0}% Évaluation pratique — Réseaux locaux (BTS CIEL 1re année) — sujets A et B
% Compilation (4 PDF) :
%   pdflatex -jobname=Eval_Reseau_SujetA   "\def\variante{A}\def\modecorrige{0}\input{Eval_TP_Reseau_LAN.tex}"
%   pdflatex -jobname=Eval_Reseau_SujetB   "\def\variante{B}\def\modecorrige{0}\input{Eval_TP_Reseau_LAN.tex}"
%   pdflatex -jobname=Eval_Reseau_CorrigeA "\def\variante{A}\def\modecorrige{1}\input{Eval_TP_Reseau_LAN.tex}"
%   pdflatex -jobname=Eval_Reseau_CorrigeB "\def\variante{B}\def\modecorrige{1}\input{Eval_TP_Reseau_LAN.tex}"
\documentclass[10pt,a4paper]{article}

\usepackage[T1]{fontenc}
\usepackage[utf8]{inputenc}
\usepackage{lmodern}
\usepackage[french]{babel}
\usepackage[a4paper,hmargin=1.5cm,top=1.3cm,bottom=1.6cm,footskip=0.7cm]{geometry}
\usepackage[table]{xcolor}
\usepackage{booktabs,tabularx,array}
\usepackage{enumitem}
\usepackage{fancyhdr}
\usepackage[most]{tcolorbox}
\usepackage{tikz}
\usetikzlibrary{positioning}
\usepackage{titlesec}

\providecommand{\modecorrige}{0}
\providecommand{\variante}{A}
\newif\ifcorrige
\ifnum\modecorrige=1 \corrigetrue \else \corrigefalse \fi

\definecolor{ipred}{HTML}{C0392B}
\definecolor{cdcbar}{HTML}{A67C3D}
\definecolor{cdcbg}{HTML}{F4F1EA}
\definecolor{qbar}{HTML}{4A6B5D}
\definecolor{qbg}{HTML}{EDF1EF}
\definecolor{corr}{HTML}{1F5FA8}
\definecolor{corrbg}{HTML}{E8F0FA}
\definecolor{grisfoot}{HTML}{828282}

\newcommand{\ip}[1]{\textcolor{ipred}{\texttt{#1}}}
\newcommand{\pts}[1]{\hfill\mbox{\textbf{(#1)}}}
\newcommand{\sol}[1]{\ifcorrige\textcolor{corr}{#1}\fi}
\newcommand{\bareme}[1]{\ifcorrige\par{\footnotesize\color{corr}#1\par}\fi}
\newcommand{\hl}{\ifcorrige\rule[-1.6mm]{0pt}{5.4mm}\else\rule[-2.2mm]{0pt}{6.6mm}\fi}
\newcolumntype{C}{>{\centering\arraybackslash}X}
\newcolumntype{L}{>{\raggedright\arraybackslash}X}

\titleformat{\section}{\large\bfseries}{}{0pt}{}
\titlespacing*{\section}{0pt}{6pt}{3pt}

\tcbset{qstyle/.style={enhanced,boxrule=0pt,leftrule=2.5pt,arc=0pt,outer arc=0pt,
  left=5pt,right=5pt,top=2pt,bottom=2pt,before skip=4pt,after skip=3pt}}
\newtcolorbox{cdc}{qstyle,colback=cdcbg,colframe=cdcbar}
\newtcolorbox{quest}{qstyle,colback=qbg,colframe=qbar}

% Zone de réponse : vierge dans le sujet, remplie dans le corrigé
\newcommand{\zone}[2]{%
  \ifcorrige
    \begin{tcolorbox}[colback=corrbg,colframe=corr,boxrule=0.5pt,arc=1pt,left=4pt,right=4pt,top=1pt,bottom=1pt,
      before skip=2pt,after skip=2pt,fontupper=\footnotesize\color{corr}]#2\end{tcolorbox}%
  \else
    \begin{tcolorbox}[colback=white,colframe=black!45,boxrule=0.4pt,arc=1pt,height=#1,
      before skip=2pt,after skip=2pt]\end{tcolorbox}%
  \fi}

\newcommand{\question}[3]{\begin{quest}\textbf{#1.} #3\pts{#2}\end{quest}}

%==========================================================================
% Données propres à chaque variante
%==========================================================================
\if\variante A
  \def\entreprise{Le cabinet d'architectes \textbf{ArchiNova}}
  \def\Rnom{R-ArchiNova}
  \def\resUn{Bureau d'études}\def\pUn{BE}\def\netUn{192.168.40}
  \def\resDeux{Accueil}\def\pDeux{ACC}\def\netDeux{172.16}
  \def\srvNom{SRV-Intranet}\def\netSrv{192.168.200}\def\srvH{10}
  \def\ssid{ArchiNova-WiFi}
  \def\AIrows{%
    \hl\ip{192.168.40.25} & \sol{C} & \sol{255.255.255.0} & \sol{192.168.40.0} & \sol{192.168.40.255} & \sol{254} \\ \hline
    \hl\ip{172.16.5.10}   & \sol{B} & \sol{255.255.0.0}   & \sol{172.16.0.0}   & \sol{172.16.255.255} & \sol{65\,534} \\ \hline
    \hl\ip{10.3.2.1}      & \sol{A} & \sol{255.0.0.0}     & \sol{10.0.0.0}     & \sol{10.255.255.255} & \sol{16\,777\,214} \\ \hline}
  \def\apipa{169.254.12.7}
  \def\Cnet{192.168.60}
  \def\Crows{%
    PC-A & \ip{192.168.60.10} & \ip{255.255.255.0} & \ip{192.168.60.254} & — \\
    PC-B & \ip{192.168.60.20} & \ip{255.255.255.0} & \ip{192.168.60.1}   & — \\
    PC-C & \ip{192.168.61.30} & \ip{255.255.255.0} & \ip{192.168.60.254} & — \\
    PC-D & \ip{192.168.60.10} & \ip{255.255.255.0} & \ip{192.168.60.254} & — \\
    SRV-C & \ip{10.0.0.5}     & \ip{255.0.0.0}     & \ip{10.0.0.254}     & — \\
    R1 — G0/0 & \ip{192.168.60.254} & \ip{255.255.255.0} & — & On \\
    R1 — G0/1 & \ip{10.0.0.254}     & \ip{255.0.0.0}     & — & Off \\}
  \def\Csol{%
    \hl 1 & \sol{PC-B} & \sol{Passerelle .60.1 : pas l'adresse du routeur.} & \sol{Joint son réseau, mais aucun autre (serveur).} & \sol{Passerelle 192.168.60.254.} \\ \hline
    \hl 2 & \sol{PC-C} & \sol{IP hors du réseau 192.168.60.0/24.} & \sol{Ne joint ni les postes ni le routeur.} & \sol{Ex. 192.168.60.30.} \\ \hline
    \hl 3 & \sol{PC-D} & \sol{Même IP que PC-A (doublon).} & \sol{Conflit d'adresse : PC-A et PC-D perturbés.} & \sol{Adresse unique, ex. .60.40.} \\ \hline
    \hl 4 & \sol{R1 G0/1} & \sol{Interface désactivée (Off).} & \sol{Réseau du serveur injoignable pour tous.} & \sol{Port Status : On.} \\ \hline}
  \def\Cbareme{PC-A et SRV-C sont correctement configurés ; accepter PC-A cité à la place de PC-D pour le doublon.}
\else
  \def\entreprise{La clinique vétérinaire \textbf{VetoPlus}}
  \def\Rnom{R-VetoPlus}
  \def\resUn{Soins}\def\pUn{SOI}\def\netUn{192.168.70}
  \def\resDeux{Secrétariat}\def\pDeux{SEC}\def\netDeux{172.20}
  \def\srvNom{SRV-Dossiers}\def\netSrv{192.168.250}\def\srvH{20}
  \def\ssid{VetoPlus-WiFi}
  \def\AIrows{%
    \hl\ip{192.168.15.130} & \sol{C} & \sol{255.255.255.0} & \sol{192.168.15.0} & \sol{192.168.15.255} & \sol{254} \\ \hline
    \hl\ip{172.31.200.7}   & \sol{B} & \sol{255.255.0.0}   & \sol{172.31.0.0}   & \sol{172.31.255.255} & \sol{65\,534} \\ \hline
    \hl\ip{10.45.6.9}      & \sol{A} & \sol{255.0.0.0}     & \sol{10.0.0.0}     & \sol{10.255.255.255} & \sol{16\,777\,214} \\ \hline}
  \def\apipa{169.254.88.3}
  \def\Cnet{192.168.80}
  \def\Crows{%
    PC-A & \ip{192.168.8.11}   & \ip{255.255.255.0} & \ip{192.168.80.254} & — \\
    PC-B & \ip{192.168.80.12}  & \ip{255.255.255.0} & \ip{192.168.80.254} & — \\
    PC-C & \ip{192.168.80.13}  & \ip{255.255.255.0} & \ip{192.168.80.100} & — \\
    PC-D & \ip{192.168.80.12}  & \ip{255.255.255.0} & \ip{192.168.80.254} & — \\
    SRV-C & \ip{10.0.0.8}      & \ip{255.0.0.0}     & \ip{10.0.0.254}     & — \\
    R1 — G0/0 & \ip{192.168.80.254} & \ip{255.255.255.0} & — & Off \\
    R1 — G0/1 & \ip{10.0.0.254}     & \ip{255.0.0.0}     & — & On \\}
  \def\Csol{%
    \hl 1 & \sol{PC-A} & \sol{IP hors du réseau 192.168.80.0/24.} & \sol{Ne joint ni les postes ni le routeur.} & \sol{Ex. 192.168.80.11.} \\ \hline
    \hl 2 & \sol{PC-C} & \sol{Passerelle .80.100 : pas l'adresse du routeur.} & \sol{Joint son réseau, mais aucun autre (serveur).} & \sol{Passerelle 192.168.80.254.} \\ \hline
    \hl 3 & \sol{PC-D} & \sol{Même IP que PC-B (doublon).} & \sol{Conflit d'adresse : PC-B et PC-D perturbés.} & \sol{Adresse unique, ex. .80.14.} \\ \hline
    \hl 4 & \sol{R1 G0/0} & \sol{Interface côté LAN désactivée (Off).} & \sol{Aucun poste ne joint le routeur ni le serveur.} & \sol{Port Status : On.} \\ \hline}
  \def\Cbareme{PC-B et SRV-C sont correctement configurés ; accepter PC-B cité à la place de PC-D pour le doublon.}
\fi

\pagestyle{fancy}
\fancyhf{}
\renewcommand{\headrulewidth}{0pt}
\fancyfoot[L]{\footnotesize\itshape\color{grisfoot}Évaluation — Réseaux locaux — Sujet \variante\ifcorrige\ — CORRIGÉ\fi}
\fancyfoot[R]{\footnotesize\color{grisfoot}\thepage/3}

\setlist{itemsep=0pt,topsep=2pt,parsep=0pt}
\setlength{\parindent}{0pt}
\setlength{\parskip}{2pt}

\begin{document}

\begin{center}
  {\Large\bfseries Évaluation pratique — Réseaux locaux \quad
   \fcolorbox{black}{white}{\ SUJET \variante\ }\par}
  \vspace{1pt}
  {\bfseries Adressage IP, commutation et routage — BTS CIEL 1re année — durée : 2 heures\par}
  \ifcorrige\vspace{2pt}{\bfseries\color{corr}CORRIGÉ ET BARÈME — document enseignant}\par\fi
\end{center}
\vspace{-2pt}
{\renewcommand{\arraystretch}{1.35}
\begin{tabularx}{\linewidth}{|X|X|p{2.2cm}|p{2.6cm}|}
  \hline
  \textbf{Nom :} & \textbf{Prénom :} & \textbf{Poste n°} & \textbf{Note :} \hfill /20 \\ \hline
\end{tabularx}}

\textbf{Consignes :} travail \textbf{individuel}, documents et Internet \textbf{non autorisés}. Ce sujet est le
\textbf{document-réponse}. Fichier Packet Tracer : \ip{Nom\_Prenom\_EvalLAN\_\variante.pkt}, à enregistrer
régulièrement dans le dossier de rendu. Les « visas » sont apposés par le professeur, que vous appelez à chaque test réussi.
Durées conseillées : A 30 min — B 1 h 10 — C 20 min.

%==========================================================================
\section{Partie A — Questions de cours et calculs \hfill /6}

\question{A1}{2 pts}{Complétez le tableau en considérant le \textbf{masque par défaut} de la classe de chaque adresse.}
{\small\renewcommand{\arraystretch}{1.15}
\begin{tabularx}{\linewidth}{|l|c|C|C|C|c|}
  \hline
  \textbf{Adresse IP} & \textbf{Classe} & \textbf{Masque par défaut} & \textbf{Adresse réseau} &
  \textbf{Adresse de diffusion} & \textbf{Nb d'hôtes} \\ \hline
  \AIrows
\end{tabularx}}
\bareme{0,5 pt par ligne juste + 0,5 pt si les trois nombres d'hôtes sont justes ($2^{n}-2$, $n$ = nombre de bits d'hôte).}

\question{A2}{1 pt}{Indiquez le type de câble (\textbf{droit} ou \textbf{croisé}) pour chaque liaison.}
{\small\renewcommand{\arraystretch}{1.15}
\begin{tabularx}{\linewidth}{|X|c||X|c|}
  \hline
  \hl PC $\leftrightarrow$ commutateur & \makebox[1.6cm]{\sol{droit}} & Commutateur $\leftrightarrow$ routeur & \makebox[1.6cm]{\sol{droit}} \\ \hline
  \hl PC $\leftrightarrow$ PC (réseau P2P) & \sol{croisé} & Concentrateur $\leftrightarrow$ concentrateur & \sol{croisé} \\ \hline
\end{tabularx}}
\bareme{0,25 pt par case. Équipements de types différents : câble droit ; de même type : câble croisé.}

\question{A3}{1 pt}{Comment un \textbf{concentrateur (hub)} et un \textbf{commutateur (switch)} traitent-ils une trame
reçue ? Lequel limite les collisions, et pourquoi ?}
\zone{2.9cm}{Le \textbf{hub} (couche 1) répète la trame sur \textbf{tous ses autres ports} : un seul domaine de collision
partagé. Le \textbf{switch} (couche 2) lit l'adresse MAC de destination et, grâce à sa \textbf{table MAC}, ne l'envoie
\textbf{que sur le port du destinataire} : un domaine de collision par port, c'est lui qui limite les collisions.
\emph{0,5 pt hub, 0,5 pt switch.}}

\question{A4}{1 pt}{Qu'est-ce que la \textbf{passerelle par défaut} d'un poste ? Quelle condition son adresse doit-elle respecter ?}
\zone{2.4cm}{Adresse IP de l'interface du \textbf{routeur} à laquelle le poste envoie les paquets destinés à un
\textbf{autre réseau}. Elle doit \textbf{appartenir au même réseau} que le poste. \emph{0,5 + 0,5 pt.}}

\question{A5}{0,5 pt}{Un portable configuré en DHCP affiche \ip{\apipa} / \ip{255.255.0.0}. Que s'est-il passé ?}
\zone{1.5cm}{Aucun \textbf{serveur DHCP} n'a répondu : le poste s'est attribué une adresse \textbf{APIPA} (169.254.0.0/16).}

\question{A6}{0,5 pt}{Quel protocole utilise \texttt{ping} ? Nommez les deux messages échangés et la couche OSI des adresses IP.}
\zone{1.5cm}{\textbf{ICMP} ; \textbf{Echo Request} / \textbf{Echo Reply} ; adresses IP en \textbf{couche 3 (Réseau)}.}

%==========================================================================
\newpage
\section{Partie B — Réalisation sous Packet Tracer \hfill /10}

\begin{cdc}
\entreprise{} vous confie son réseau : \textbf{trois réseaux} reliés par un \textbf{routeur unique Cisco 2911}
nommé \ip{\Rnom} (interfaces G0/0, G0/1, G0/2).
\begin{itemize}[leftmargin=12pt]
  \item \textbf{Réseau \resUn} : \textbf{4 postes} PC-\pUn1 à PC-\pUn4, commutateur \ip{SW-\pUn}, adressage statique,
        réseau \ip{\netUn.0/24}.
  \item \textbf{Réseau \resDeux} : \textbf{2 postes} PC-\pDeux1, PC-\pDeux2 et \textbf{1 imprimante} IMP-\pDeux,
        commutateur \ip{SW-\pDeux}, adressage statique, réseau de classe B \ip{\netDeux.0.0/16}.
  \item \textbf{Réseau Serveur} : serveur \ip{\srvNom} d'adresse \ip{\netSrv.\srvH/24}, relié \textbf{directement} au routeur.
  \item La passerelle de chaque réseau est sa \textbf{dernière adresse utilisable}. Tous les équipements communiquent
        entre eux et accèdent à la page web du serveur. Nommage imposé (« Display Name »).
\end{itemize}
\end{cdc}

\question{B1}{2 pts}{Complétez le plan d'adressage \textbf{avant} la réalisation.}
{\footnotesize\setlength{\tabcolsep}{3pt}\renewcommand{\arraystretch}{1.15}
\begin{tabularx}{\linewidth}{|l|C|C|C|C|C|C|}
  \hline
  \textbf{Réseau} & \textbf{Adresse réseau} & \textbf{Masque} & \textbf{1re @ utilisable} &
  \textbf{Dernière @ utilisable} & \textbf{Diffusion} & \textbf{Passerelle} \\ \hline
  \hl \resUn & \sol{\netUn.0} & \sol{255.255.255.0} & \sol{\netUn.1} & \sol{\netUn.254} & \sol{\netUn.255} & \sol{\netUn.254} \\ \hline
  \hl \resDeux & \sol{\netDeux.0.0} & \sol{255.255.0.0} & \sol{\netDeux.0.1} & \sol{\netDeux.255.254} & \sol{\netDeux.255.255} & \sol{\netDeux.255.254} \\ \hline
  \hl Serveur & \sol{\netSrv.0} & \sol{255.255.255.0} & \sol{\netSrv.1} & \sol{\netSrv.254} & \sol{\netSrv.255} & \sol{\netSrv.254} \\ \hline
\end{tabularx}}
\bareme{0,5 pt par ligne juste + 0,5 pt si la convention de passerelle est respectée partout.}

\question{B2}{2 pts}{Construisez la topologie (équipements, câblage, nommage) et relevez l'adressage configuré.}
{\footnotesize\setlength{\tabcolsep}{3pt}\renewcommand{\arraystretch}{1.1}
\begin{tabularx}{\linewidth}{|l|C|C|C||l|C|C|}
  \hline
  \textbf{Hôte} & \textbf{Adresse IP} & \textbf{Masque} & \textbf{Passerelle} & \textbf{Interface} & \textbf{Adresse IP} & \textbf{Masque} \\ \hline
  \hl PC-\pUn1 & \sol{\netUn.1} & \sol{255.255.255.0} & \sol{\netUn.254} & \Rnom{} G0/0 & \sol{\netUn.254} & \sol{255.255.255.0} \\ \hline
  \hl PC-\pDeux1 & \sol{\netDeux.0.1} & \sol{255.255.0.0} & \sol{\netDeux.255.254} & \Rnom{} G0/1 & \sol{\netDeux.255.254} & \sol{255.255.0.0} \\ \hline
  \hl IMP-\pDeux & \sol{\netDeux.0.10} & \sol{255.255.0.0} & \sol{\netDeux.255.254} & \Rnom{} G0/2 & \sol{\netSrv.254} & \sol{255.255.255.0} \\ \hline
  \hl \srvNom & \ip{\netSrv.\srvH} & \sol{255.255.255.0} & \sol{\netSrv.254} & \multicolumn{3}{c|}{\cellcolor{black!8}} \\ \hline
\end{tabularx}}
\bareme{Adresses d'hôtes libres (uniques, dans le bon réseau) ; affectation G0/x libre. Sur le .pkt : équipements 0,5 ;
câblage 1 (câbles droits vers les switchs et switch–routeur, \textbf{croisé} serveur–routeur, liens verts) ; nommage 0,5.}

\question{B3}{3 pts}{Configurez tous les équipements : IP, masques, passerelles des hôtes ; interfaces du routeur
(« Port Status : On »).}
\bareme{Sur le .pkt : interfaces routeur 1 pt ; IP/masques des hôtes 1 pt ; passerelles (imprimante et serveur compris)
1 pt. $-0,25$ par erreur, sans descendre sous 0 par item.}

\question{B4}{3 pts}{Réalisez les tests, notez la commande et le résultat, puis faites viser chaque test réussi.}
{\small\renewcommand{\arraystretch}{1.1}
\begin{tabularx}{\linewidth}{|c|l|l|C|c|c|}
  \hline
  \textbf{N°} & \textbf{Source} & \textbf{Destination} & \textbf{Commande / action} & \textbf{Résultat} & \textbf{Visa prof.} \\ \hline
  \hl 1 & PC-\pUn1 & PC-\pUn4 & \sol{ping \netUn.4} & \sol{réussi} & \hspace{1.4cm} \\ \hline
  \hl 2 & PC-\pDeux1 & IMP-\pDeux & \sol{ping \netDeux.0.10} & \sol{réussi} & \\ \hline
  \hl 3 & PC-\pUn2 & PC-\pDeux2 & \sol{ping \netDeux.0.2} & \sol{réussi} & \\ \hline
  \hl 4 & PC-\pDeux2 & \srvNom & \sol{ping \netSrv.\srvH} & \sol{réussi} & \\ \hline
  \hl 5 & PC-\pUn3 & page web du serveur & \sol{navigateur : {\NoAutoSpacing http://\netSrv.\srvH}} & \sol{page affichée} & \\ \hline
\end{tabularx}}
\bareme{0,5 pt par test visé + 0,5 pt pour B5. Le premier ping inter-réseaux peut échouer (résolution ARP) ; les suivants doivent réussir.}

\question{B5}{incl. B4}{Test n°3 : par quels équipements passe le paquet ? Pourquoi PC-\pUn2 l'envoie-t-il au routeur ?}
\zone{2.3cm}{PC-\pUn2 $\to$ SW-\pUn{} $\to$ \Rnom{} $\to$ SW-\pDeux{} $\to$ PC-\pDeux2. Grâce à son masque, PC-\pUn2 constate que
la destination (\netDeux.0.0/16) n'est \textbf{pas dans son réseau} (\netUn.0/24) : il envoie le paquet à sa
\textbf{passerelle par défaut}, qui le route.}

%==========================================================================
\newpage
\section{Partie C — Diagnostic d'un réseau en panne \hfill /4}

Un technicien constate qu'\textbf{aucun poste ne parvient à joindre le serveur SRV-C} et que certains postes
communiquent mal entre eux. Les câbles sont du bon type.

\begin{center}
\begin{minipage}{0.37\linewidth}\centering
\begin{tikzpicture}[font=\footnotesize,
  eq/.style={draw,fill=qbg,rounded corners=2pt,minimum width=9mm,minimum height=5mm,inner sep=2pt},
  every path/.style={thick}]
  \node[eq] (SW) {SW-1};
  \node[eq,right=9mm of SW] (R) {R1};
  \node[eq,right=7mm of R] (SRV) {SRV-C};
  \node[eq,above left=5mm and -3mm of SW] (A) {PC-A};
  \node[eq,above right=5mm and -3mm of SW] (B) {PC-B};
  \node[eq,below left=5mm and -3mm of SW] (C) {PC-C};
  \node[eq,below right=5mm and -3mm of SW] (D) {PC-D};
  \foreach \n in {A,B,C,D} \draw (\n) -- (SW);
  \draw (SW) -- node[above,font=\tiny]{G0/0} (R);
  \draw (R) -- node[above,font=\tiny]{G0/1} (SRV);
\end{tikzpicture}
\end{minipage}\hfill
\begin{minipage}{0.61\linewidth}\centering
{\small\renewcommand{\arraystretch}{1.05}\setlength{\tabcolsep}{4pt}
\begin{tabular}{lllll}
  \toprule
  \textbf{Équipement} & \textbf{Adresse IP} & \textbf{Masque} & \textbf{Passerelle} & \textbf{Port} \\
  \midrule
  \Crows
  \bottomrule
\end{tabular}}
\end{minipage}
\end{center}

\question{C1}{4 pts}{Relevez les \textbf{quatre anomalies}. Pour chacune : conséquence et correction.}
{\small\renewcommand{\arraystretch}{1.1}
\ifcorrige\def\hc{}\else\def\hc{\rule[-10mm]{0pt}{14mm}}\fi
\begin{tabularx}{\linewidth}{|c|>{\raggedright\arraybackslash}p{2cm}|L|L|L|}
  \hline
  \textbf{N°} & \textbf{Équipement} & \textbf{Anomalie} & \textbf{Conséquence} & \textbf{Correction} \\ \hline
  \ifcorrige\Csol\else
  \hc 1 & & & & \\ \hline
  \hc 2 & & & & \\ \hline
  \hc 3 & & & & \\ \hline
  \hc 4 & & & & \\ \hline
  \fi
\end{tabularx}}
\bareme{1 pt par anomalie (0,25 équipement, anomalie, conséquence, correction ; ordre indifférent). \Cbareme}

%==========================================================================
\section{Bonus — Extension sans fil \hfill +1}

\begin{cdc}Ajoutez au réseau \resUn{} un \textbf{AccessPoint-PT} relié à SW-\pUn{} (SSID \ip{\ssid}, \textbf{WPA2-PSK})
et un portable \ip{Laptop-\pUn} équipé d'un module \texttt{WPC300N}, en adressage \textbf{statique}.
Vérifiez qu'il joint \srvNom.\end{cdc}
{\small\renewcommand{\arraystretch}{1.1}
\begin{tabularx}{\linewidth}{|l|C|C|C|c|}
  \hline
  \textbf{Équipement} & \textbf{Adresse IP} & \textbf{Masque} & \textbf{Passerelle} & \textbf{Visa prof.} \\ \hline
  \hl Laptop-\pUn & \sol{\netUn.50} & \sol{255.255.255.0} & \sol{\netUn.254} & \hspace{1.4cm} \\ \hline
\end{tabularx}}
\bareme{Le point d'accès fait office de pont : le portable est dans \netUn.0/24. Bonus si associé en WPA2-PSK et ping
vers \netSrv.\srvH{} réussi. Note plafonnée à 20.}

\end{document}
"
%   pdflatex -jobname=Eval_Reseau_SujetB   "\def\variante{B}\def\modecorrige{0}% Évaluation pratique — Réseaux locaux (BTS CIEL 1re année) — sujets A et B
% Compilation (4 PDF) :
%   pdflatex -jobname=Eval_Reseau_SujetA   "\def\variante{A}\def\modecorrige{0}\input{Eval_TP_Reseau_LAN.tex}"
%   pdflatex -jobname=Eval_Reseau_SujetB   "\def\variante{B}\def\modecorrige{0}\input{Eval_TP_Reseau_LAN.tex}"
%   pdflatex -jobname=Eval_Reseau_CorrigeA "\def\variante{A}\def\modecorrige{1}\input{Eval_TP_Reseau_LAN.tex}"
%   pdflatex -jobname=Eval_Reseau_CorrigeB "\def\variante{B}\def\modecorrige{1}\input{Eval_TP_Reseau_LAN.tex}"
\documentclass[10pt,a4paper]{article}

\usepackage[T1]{fontenc}
\usepackage[utf8]{inputenc}
\usepackage{lmodern}
\usepackage[french]{babel}
\usepackage[a4paper,hmargin=1.5cm,top=1.3cm,bottom=1.6cm,footskip=0.7cm]{geometry}
\usepackage[table]{xcolor}
\usepackage{booktabs,tabularx,array}
\usepackage{enumitem}
\usepackage{fancyhdr}
\usepackage[most]{tcolorbox}
\usepackage{tikz}
\usetikzlibrary{positioning}
\usepackage{titlesec}

\providecommand{\modecorrige}{0}
\providecommand{\variante}{A}
\newif\ifcorrige
\ifnum\modecorrige=1 \corrigetrue \else \corrigefalse \fi

\definecolor{ipred}{HTML}{C0392B}
\definecolor{cdcbar}{HTML}{A67C3D}
\definecolor{cdcbg}{HTML}{F4F1EA}
\definecolor{qbar}{HTML}{4A6B5D}
\definecolor{qbg}{HTML}{EDF1EF}
\definecolor{corr}{HTML}{1F5FA8}
\definecolor{corrbg}{HTML}{E8F0FA}
\definecolor{grisfoot}{HTML}{828282}

\newcommand{\ip}[1]{\textcolor{ipred}{\texttt{#1}}}
\newcommand{\pts}[1]{\hfill\mbox{\textbf{(#1)}}}
\newcommand{\sol}[1]{\ifcorrige\textcolor{corr}{#1}\fi}
\newcommand{\bareme}[1]{\ifcorrige\par{\footnotesize\color{corr}#1\par}\fi}
\newcommand{\hl}{\ifcorrige\rule[-1.6mm]{0pt}{5.4mm}\else\rule[-2.2mm]{0pt}{6.6mm}\fi}
\newcolumntype{C}{>{\centering\arraybackslash}X}
\newcolumntype{L}{>{\raggedright\arraybackslash}X}

\titleformat{\section}{\large\bfseries}{}{0pt}{}
\titlespacing*{\section}{0pt}{6pt}{3pt}

\tcbset{qstyle/.style={enhanced,boxrule=0pt,leftrule=2.5pt,arc=0pt,outer arc=0pt,
  left=5pt,right=5pt,top=2pt,bottom=2pt,before skip=4pt,after skip=3pt}}
\newtcolorbox{cdc}{qstyle,colback=cdcbg,colframe=cdcbar}
\newtcolorbox{quest}{qstyle,colback=qbg,colframe=qbar}

% Zone de réponse : vierge dans le sujet, remplie dans le corrigé
\newcommand{\zone}[2]{%
  \ifcorrige
    \begin{tcolorbox}[colback=corrbg,colframe=corr,boxrule=0.5pt,arc=1pt,left=4pt,right=4pt,top=1pt,bottom=1pt,
      before skip=2pt,after skip=2pt,fontupper=\footnotesize\color{corr}]#2\end{tcolorbox}%
  \else
    \begin{tcolorbox}[colback=white,colframe=black!45,boxrule=0.4pt,arc=1pt,height=#1,
      before skip=2pt,after skip=2pt]\end{tcolorbox}%
  \fi}

\newcommand{\question}[3]{\begin{quest}\textbf{#1.} #3\pts{#2}\end{quest}}

%==========================================================================
% Données propres à chaque variante
%==========================================================================
\if\variante A
  \def\entreprise{Le cabinet d'architectes \textbf{ArchiNova}}
  \def\Rnom{R-ArchiNova}
  \def\resUn{Bureau d'études}\def\pUn{BE}\def\netUn{192.168.40}
  \def\resDeux{Accueil}\def\pDeux{ACC}\def\netDeux{172.16}
  \def\srvNom{SRV-Intranet}\def\netSrv{192.168.200}\def\srvH{10}
  \def\ssid{ArchiNova-WiFi}
  \def\AIrows{%
    \hl\ip{192.168.40.25} & \sol{C} & \sol{255.255.255.0} & \sol{192.168.40.0} & \sol{192.168.40.255} & \sol{254} \\ \hline
    \hl\ip{172.16.5.10}   & \sol{B} & \sol{255.255.0.0}   & \sol{172.16.0.0}   & \sol{172.16.255.255} & \sol{65\,534} \\ \hline
    \hl\ip{10.3.2.1}      & \sol{A} & \sol{255.0.0.0}     & \sol{10.0.0.0}     & \sol{10.255.255.255} & \sol{16\,777\,214} \\ \hline}
  \def\apipa{169.254.12.7}
  \def\Cnet{192.168.60}
  \def\Crows{%
    PC-A & \ip{192.168.60.10} & \ip{255.255.255.0} & \ip{192.168.60.254} & — \\
    PC-B & \ip{192.168.60.20} & \ip{255.255.255.0} & \ip{192.168.60.1}   & — \\
    PC-C & \ip{192.168.61.30} & \ip{255.255.255.0} & \ip{192.168.60.254} & — \\
    PC-D & \ip{192.168.60.10} & \ip{255.255.255.0} & \ip{192.168.60.254} & — \\
    SRV-C & \ip{10.0.0.5}     & \ip{255.0.0.0}     & \ip{10.0.0.254}     & — \\
    R1 — G0/0 & \ip{192.168.60.254} & \ip{255.255.255.0} & — & On \\
    R1 — G0/1 & \ip{10.0.0.254}     & \ip{255.0.0.0}     & — & Off \\}
  \def\Csol{%
    \hl 1 & \sol{PC-B} & \sol{Passerelle .60.1 : pas l'adresse du routeur.} & \sol{Joint son réseau, mais aucun autre (serveur).} & \sol{Passerelle 192.168.60.254.} \\ \hline
    \hl 2 & \sol{PC-C} & \sol{IP hors du réseau 192.168.60.0/24.} & \sol{Ne joint ni les postes ni le routeur.} & \sol{Ex. 192.168.60.30.} \\ \hline
    \hl 3 & \sol{PC-D} & \sol{Même IP que PC-A (doublon).} & \sol{Conflit d'adresse : PC-A et PC-D perturbés.} & \sol{Adresse unique, ex. .60.40.} \\ \hline
    \hl 4 & \sol{R1 G0/1} & \sol{Interface désactivée (Off).} & \sol{Réseau du serveur injoignable pour tous.} & \sol{Port Status : On.} \\ \hline}
  \def\Cbareme{PC-A et SRV-C sont correctement configurés ; accepter PC-A cité à la place de PC-D pour le doublon.}
\else
  \def\entreprise{La clinique vétérinaire \textbf{VetoPlus}}
  \def\Rnom{R-VetoPlus}
  \def\resUn{Soins}\def\pUn{SOI}\def\netUn{192.168.70}
  \def\resDeux{Secrétariat}\def\pDeux{SEC}\def\netDeux{172.20}
  \def\srvNom{SRV-Dossiers}\def\netSrv{192.168.250}\def\srvH{20}
  \def\ssid{VetoPlus-WiFi}
  \def\AIrows{%
    \hl\ip{192.168.15.130} & \sol{C} & \sol{255.255.255.0} & \sol{192.168.15.0} & \sol{192.168.15.255} & \sol{254} \\ \hline
    \hl\ip{172.31.200.7}   & \sol{B} & \sol{255.255.0.0}   & \sol{172.31.0.0}   & \sol{172.31.255.255} & \sol{65\,534} \\ \hline
    \hl\ip{10.45.6.9}      & \sol{A} & \sol{255.0.0.0}     & \sol{10.0.0.0}     & \sol{10.255.255.255} & \sol{16\,777\,214} \\ \hline}
  \def\apipa{169.254.88.3}
  \def\Cnet{192.168.80}
  \def\Crows{%
    PC-A & \ip{192.168.8.11}   & \ip{255.255.255.0} & \ip{192.168.80.254} & — \\
    PC-B & \ip{192.168.80.12}  & \ip{255.255.255.0} & \ip{192.168.80.254} & — \\
    PC-C & \ip{192.168.80.13}  & \ip{255.255.255.0} & \ip{192.168.80.100} & — \\
    PC-D & \ip{192.168.80.12}  & \ip{255.255.255.0} & \ip{192.168.80.254} & — \\
    SRV-C & \ip{10.0.0.8}      & \ip{255.0.0.0}     & \ip{10.0.0.254}     & — \\
    R1 — G0/0 & \ip{192.168.80.254} & \ip{255.255.255.0} & — & Off \\
    R1 — G0/1 & \ip{10.0.0.254}     & \ip{255.0.0.0}     & — & On \\}
  \def\Csol{%
    \hl 1 & \sol{PC-A} & \sol{IP hors du réseau 192.168.80.0/24.} & \sol{Ne joint ni les postes ni le routeur.} & \sol{Ex. 192.168.80.11.} \\ \hline
    \hl 2 & \sol{PC-C} & \sol{Passerelle .80.100 : pas l'adresse du routeur.} & \sol{Joint son réseau, mais aucun autre (serveur).} & \sol{Passerelle 192.168.80.254.} \\ \hline
    \hl 3 & \sol{PC-D} & \sol{Même IP que PC-B (doublon).} & \sol{Conflit d'adresse : PC-B et PC-D perturbés.} & \sol{Adresse unique, ex. .80.14.} \\ \hline
    \hl 4 & \sol{R1 G0/0} & \sol{Interface côté LAN désactivée (Off).} & \sol{Aucun poste ne joint le routeur ni le serveur.} & \sol{Port Status : On.} \\ \hline}
  \def\Cbareme{PC-B et SRV-C sont correctement configurés ; accepter PC-B cité à la place de PC-D pour le doublon.}
\fi

\pagestyle{fancy}
\fancyhf{}
\renewcommand{\headrulewidth}{0pt}
\fancyfoot[L]{\footnotesize\itshape\color{grisfoot}Évaluation — Réseaux locaux — Sujet \variante\ifcorrige\ — CORRIGÉ\fi}
\fancyfoot[R]{\footnotesize\color{grisfoot}\thepage/3}

\setlist{itemsep=0pt,topsep=2pt,parsep=0pt}
\setlength{\parindent}{0pt}
\setlength{\parskip}{2pt}

\begin{document}

\begin{center}
  {\Large\bfseries Évaluation pratique — Réseaux locaux \quad
   \fcolorbox{black}{white}{\ SUJET \variante\ }\par}
  \vspace{1pt}
  {\bfseries Adressage IP, commutation et routage — BTS CIEL 1re année — durée : 2 heures\par}
  \ifcorrige\vspace{2pt}{\bfseries\color{corr}CORRIGÉ ET BARÈME — document enseignant}\par\fi
\end{center}
\vspace{-2pt}
{\renewcommand{\arraystretch}{1.35}
\begin{tabularx}{\linewidth}{|X|X|p{2.2cm}|p{2.6cm}|}
  \hline
  \textbf{Nom :} & \textbf{Prénom :} & \textbf{Poste n°} & \textbf{Note :} \hfill /20 \\ \hline
\end{tabularx}}

\textbf{Consignes :} travail \textbf{individuel}, documents et Internet \textbf{non autorisés}. Ce sujet est le
\textbf{document-réponse}. Fichier Packet Tracer : \ip{Nom\_Prenom\_EvalLAN\_\variante.pkt}, à enregistrer
régulièrement dans le dossier de rendu. Les « visas » sont apposés par le professeur, que vous appelez à chaque test réussi.
Durées conseillées : A 30 min — B 1 h 10 — C 20 min.

%==========================================================================
\section{Partie A — Questions de cours et calculs \hfill /6}

\question{A1}{2 pts}{Complétez le tableau en considérant le \textbf{masque par défaut} de la classe de chaque adresse.}
{\small\renewcommand{\arraystretch}{1.15}
\begin{tabularx}{\linewidth}{|l|c|C|C|C|c|}
  \hline
  \textbf{Adresse IP} & \textbf{Classe} & \textbf{Masque par défaut} & \textbf{Adresse réseau} &
  \textbf{Adresse de diffusion} & \textbf{Nb d'hôtes} \\ \hline
  \AIrows
\end{tabularx}}
\bareme{0,5 pt par ligne juste + 0,5 pt si les trois nombres d'hôtes sont justes ($2^{n}-2$, $n$ = nombre de bits d'hôte).}

\question{A2}{1 pt}{Indiquez le type de câble (\textbf{droit} ou \textbf{croisé}) pour chaque liaison.}
{\small\renewcommand{\arraystretch}{1.15}
\begin{tabularx}{\linewidth}{|X|c||X|c|}
  \hline
  \hl PC $\leftrightarrow$ commutateur & \makebox[1.6cm]{\sol{droit}} & Commutateur $\leftrightarrow$ routeur & \makebox[1.6cm]{\sol{droit}} \\ \hline
  \hl PC $\leftrightarrow$ PC (réseau P2P) & \sol{croisé} & Concentrateur $\leftrightarrow$ concentrateur & \sol{croisé} \\ \hline
\end{tabularx}}
\bareme{0,25 pt par case. Équipements de types différents : câble droit ; de même type : câble croisé.}

\question{A3}{1 pt}{Comment un \textbf{concentrateur (hub)} et un \textbf{commutateur (switch)} traitent-ils une trame
reçue ? Lequel limite les collisions, et pourquoi ?}
\zone{2.9cm}{Le \textbf{hub} (couche 1) répète la trame sur \textbf{tous ses autres ports} : un seul domaine de collision
partagé. Le \textbf{switch} (couche 2) lit l'adresse MAC de destination et, grâce à sa \textbf{table MAC}, ne l'envoie
\textbf{que sur le port du destinataire} : un domaine de collision par port, c'est lui qui limite les collisions.
\emph{0,5 pt hub, 0,5 pt switch.}}

\question{A4}{1 pt}{Qu'est-ce que la \textbf{passerelle par défaut} d'un poste ? Quelle condition son adresse doit-elle respecter ?}
\zone{2.4cm}{Adresse IP de l'interface du \textbf{routeur} à laquelle le poste envoie les paquets destinés à un
\textbf{autre réseau}. Elle doit \textbf{appartenir au même réseau} que le poste. \emph{0,5 + 0,5 pt.}}

\question{A5}{0,5 pt}{Un portable configuré en DHCP affiche \ip{\apipa} / \ip{255.255.0.0}. Que s'est-il passé ?}
\zone{1.5cm}{Aucun \textbf{serveur DHCP} n'a répondu : le poste s'est attribué une adresse \textbf{APIPA} (169.254.0.0/16).}

\question{A6}{0,5 pt}{Quel protocole utilise \texttt{ping} ? Nommez les deux messages échangés et la couche OSI des adresses IP.}
\zone{1.5cm}{\textbf{ICMP} ; \textbf{Echo Request} / \textbf{Echo Reply} ; adresses IP en \textbf{couche 3 (Réseau)}.}

%==========================================================================
\newpage
\section{Partie B — Réalisation sous Packet Tracer \hfill /10}

\begin{cdc}
\entreprise{} vous confie son réseau : \textbf{trois réseaux} reliés par un \textbf{routeur unique Cisco 2911}
nommé \ip{\Rnom} (interfaces G0/0, G0/1, G0/2).
\begin{itemize}[leftmargin=12pt]
  \item \textbf{Réseau \resUn} : \textbf{4 postes} PC-\pUn1 à PC-\pUn4, commutateur \ip{SW-\pUn}, adressage statique,
        réseau \ip{\netUn.0/24}.
  \item \textbf{Réseau \resDeux} : \textbf{2 postes} PC-\pDeux1, PC-\pDeux2 et \textbf{1 imprimante} IMP-\pDeux,
        commutateur \ip{SW-\pDeux}, adressage statique, réseau de classe B \ip{\netDeux.0.0/16}.
  \item \textbf{Réseau Serveur} : serveur \ip{\srvNom} d'adresse \ip{\netSrv.\srvH/24}, relié \textbf{directement} au routeur.
  \item La passerelle de chaque réseau est sa \textbf{dernière adresse utilisable}. Tous les équipements communiquent
        entre eux et accèdent à la page web du serveur. Nommage imposé (« Display Name »).
\end{itemize}
\end{cdc}

\question{B1}{2 pts}{Complétez le plan d'adressage \textbf{avant} la réalisation.}
{\footnotesize\setlength{\tabcolsep}{3pt}\renewcommand{\arraystretch}{1.15}
\begin{tabularx}{\linewidth}{|l|C|C|C|C|C|C|}
  \hline
  \textbf{Réseau} & \textbf{Adresse réseau} & \textbf{Masque} & \textbf{1re @ utilisable} &
  \textbf{Dernière @ utilisable} & \textbf{Diffusion} & \textbf{Passerelle} \\ \hline
  \hl \resUn & \sol{\netUn.0} & \sol{255.255.255.0} & \sol{\netUn.1} & \sol{\netUn.254} & \sol{\netUn.255} & \sol{\netUn.254} \\ \hline
  \hl \resDeux & \sol{\netDeux.0.0} & \sol{255.255.0.0} & \sol{\netDeux.0.1} & \sol{\netDeux.255.254} & \sol{\netDeux.255.255} & \sol{\netDeux.255.254} \\ \hline
  \hl Serveur & \sol{\netSrv.0} & \sol{255.255.255.0} & \sol{\netSrv.1} & \sol{\netSrv.254} & \sol{\netSrv.255} & \sol{\netSrv.254} \\ \hline
\end{tabularx}}
\bareme{0,5 pt par ligne juste + 0,5 pt si la convention de passerelle est respectée partout.}

\question{B2}{2 pts}{Construisez la topologie (équipements, câblage, nommage) et relevez l'adressage configuré.}
{\footnotesize\setlength{\tabcolsep}{3pt}\renewcommand{\arraystretch}{1.1}
\begin{tabularx}{\linewidth}{|l|C|C|C||l|C|C|}
  \hline
  \textbf{Hôte} & \textbf{Adresse IP} & \textbf{Masque} & \textbf{Passerelle} & \textbf{Interface} & \textbf{Adresse IP} & \textbf{Masque} \\ \hline
  \hl PC-\pUn1 & \sol{\netUn.1} & \sol{255.255.255.0} & \sol{\netUn.254} & \Rnom{} G0/0 & \sol{\netUn.254} & \sol{255.255.255.0} \\ \hline
  \hl PC-\pDeux1 & \sol{\netDeux.0.1} & \sol{255.255.0.0} & \sol{\netDeux.255.254} & \Rnom{} G0/1 & \sol{\netDeux.255.254} & \sol{255.255.0.0} \\ \hline
  \hl IMP-\pDeux & \sol{\netDeux.0.10} & \sol{255.255.0.0} & \sol{\netDeux.255.254} & \Rnom{} G0/2 & \sol{\netSrv.254} & \sol{255.255.255.0} \\ \hline
  \hl \srvNom & \ip{\netSrv.\srvH} & \sol{255.255.255.0} & \sol{\netSrv.254} & \multicolumn{3}{c|}{\cellcolor{black!8}} \\ \hline
\end{tabularx}}
\bareme{Adresses d'hôtes libres (uniques, dans le bon réseau) ; affectation G0/x libre. Sur le .pkt : équipements 0,5 ;
câblage 1 (câbles droits vers les switchs et switch–routeur, \textbf{croisé} serveur–routeur, liens verts) ; nommage 0,5.}

\question{B3}{3 pts}{Configurez tous les équipements : IP, masques, passerelles des hôtes ; interfaces du routeur
(« Port Status : On »).}
\bareme{Sur le .pkt : interfaces routeur 1 pt ; IP/masques des hôtes 1 pt ; passerelles (imprimante et serveur compris)
1 pt. $-0,25$ par erreur, sans descendre sous 0 par item.}

\question{B4}{3 pts}{Réalisez les tests, notez la commande et le résultat, puis faites viser chaque test réussi.}
{\small\renewcommand{\arraystretch}{1.1}
\begin{tabularx}{\linewidth}{|c|l|l|C|c|c|}
  \hline
  \textbf{N°} & \textbf{Source} & \textbf{Destination} & \textbf{Commande / action} & \textbf{Résultat} & \textbf{Visa prof.} \\ \hline
  \hl 1 & PC-\pUn1 & PC-\pUn4 & \sol{ping \netUn.4} & \sol{réussi} & \hspace{1.4cm} \\ \hline
  \hl 2 & PC-\pDeux1 & IMP-\pDeux & \sol{ping \netDeux.0.10} & \sol{réussi} & \\ \hline
  \hl 3 & PC-\pUn2 & PC-\pDeux2 & \sol{ping \netDeux.0.2} & \sol{réussi} & \\ \hline
  \hl 4 & PC-\pDeux2 & \srvNom & \sol{ping \netSrv.\srvH} & \sol{réussi} & \\ \hline
  \hl 5 & PC-\pUn3 & page web du serveur & \sol{navigateur : {\NoAutoSpacing http://\netSrv.\srvH}} & \sol{page affichée} & \\ \hline
\end{tabularx}}
\bareme{0,5 pt par test visé + 0,5 pt pour B5. Le premier ping inter-réseaux peut échouer (résolution ARP) ; les suivants doivent réussir.}

\question{B5}{incl. B4}{Test n°3 : par quels équipements passe le paquet ? Pourquoi PC-\pUn2 l'envoie-t-il au routeur ?}
\zone{2.3cm}{PC-\pUn2 $\to$ SW-\pUn{} $\to$ \Rnom{} $\to$ SW-\pDeux{} $\to$ PC-\pDeux2. Grâce à son masque, PC-\pUn2 constate que
la destination (\netDeux.0.0/16) n'est \textbf{pas dans son réseau} (\netUn.0/24) : il envoie le paquet à sa
\textbf{passerelle par défaut}, qui le route.}

%==========================================================================
\newpage
\section{Partie C — Diagnostic d'un réseau en panne \hfill /4}

Un technicien constate qu'\textbf{aucun poste ne parvient à joindre le serveur SRV-C} et que certains postes
communiquent mal entre eux. Les câbles sont du bon type.

\begin{center}
\begin{minipage}{0.37\linewidth}\centering
\begin{tikzpicture}[font=\footnotesize,
  eq/.style={draw,fill=qbg,rounded corners=2pt,minimum width=9mm,minimum height=5mm,inner sep=2pt},
  every path/.style={thick}]
  \node[eq] (SW) {SW-1};
  \node[eq,right=9mm of SW] (R) {R1};
  \node[eq,right=7mm of R] (SRV) {SRV-C};
  \node[eq,above left=5mm and -3mm of SW] (A) {PC-A};
  \node[eq,above right=5mm and -3mm of SW] (B) {PC-B};
  \node[eq,below left=5mm and -3mm of SW] (C) {PC-C};
  \node[eq,below right=5mm and -3mm of SW] (D) {PC-D};
  \foreach \n in {A,B,C,D} \draw (\n) -- (SW);
  \draw (SW) -- node[above,font=\tiny]{G0/0} (R);
  \draw (R) -- node[above,font=\tiny]{G0/1} (SRV);
\end{tikzpicture}
\end{minipage}\hfill
\begin{minipage}{0.61\linewidth}\centering
{\small\renewcommand{\arraystretch}{1.05}\setlength{\tabcolsep}{4pt}
\begin{tabular}{lllll}
  \toprule
  \textbf{Équipement} & \textbf{Adresse IP} & \textbf{Masque} & \textbf{Passerelle} & \textbf{Port} \\
  \midrule
  \Crows
  \bottomrule
\end{tabular}}
\end{minipage}
\end{center}

\question{C1}{4 pts}{Relevez les \textbf{quatre anomalies}. Pour chacune : conséquence et correction.}
{\small\renewcommand{\arraystretch}{1.1}
\ifcorrige\def\hc{}\else\def\hc{\rule[-10mm]{0pt}{14mm}}\fi
\begin{tabularx}{\linewidth}{|c|>{\raggedright\arraybackslash}p{2cm}|L|L|L|}
  \hline
  \textbf{N°} & \textbf{Équipement} & \textbf{Anomalie} & \textbf{Conséquence} & \textbf{Correction} \\ \hline
  \ifcorrige\Csol\else
  \hc 1 & & & & \\ \hline
  \hc 2 & & & & \\ \hline
  \hc 3 & & & & \\ \hline
  \hc 4 & & & & \\ \hline
  \fi
\end{tabularx}}
\bareme{1 pt par anomalie (0,25 équipement, anomalie, conséquence, correction ; ordre indifférent). \Cbareme}

%==========================================================================
\section{Bonus — Extension sans fil \hfill +1}

\begin{cdc}Ajoutez au réseau \resUn{} un \textbf{AccessPoint-PT} relié à SW-\pUn{} (SSID \ip{\ssid}, \textbf{WPA2-PSK})
et un portable \ip{Laptop-\pUn} équipé d'un module \texttt{WPC300N}, en adressage \textbf{statique}.
Vérifiez qu'il joint \srvNom.\end{cdc}
{\small\renewcommand{\arraystretch}{1.1}
\begin{tabularx}{\linewidth}{|l|C|C|C|c|}
  \hline
  \textbf{Équipement} & \textbf{Adresse IP} & \textbf{Masque} & \textbf{Passerelle} & \textbf{Visa prof.} \\ \hline
  \hl Laptop-\pUn & \sol{\netUn.50} & \sol{255.255.255.0} & \sol{\netUn.254} & \hspace{1.4cm} \\ \hline
\end{tabularx}}
\bareme{Le point d'accès fait office de pont : le portable est dans \netUn.0/24. Bonus si associé en WPA2-PSK et ping
vers \netSrv.\srvH{} réussi. Note plafonnée à 20.}

\end{document}
"
%   pdflatex -jobname=Eval_Reseau_CorrigeA "\def\variante{A}\def\modecorrige{1}% Évaluation pratique — Réseaux locaux (BTS CIEL 1re année) — sujets A et B
% Compilation (4 PDF) :
%   pdflatex -jobname=Eval_Reseau_SujetA   "\def\variante{A}\def\modecorrige{0}\input{Eval_TP_Reseau_LAN.tex}"
%   pdflatex -jobname=Eval_Reseau_SujetB   "\def\variante{B}\def\modecorrige{0}\input{Eval_TP_Reseau_LAN.tex}"
%   pdflatex -jobname=Eval_Reseau_CorrigeA "\def\variante{A}\def\modecorrige{1}\input{Eval_TP_Reseau_LAN.tex}"
%   pdflatex -jobname=Eval_Reseau_CorrigeB "\def\variante{B}\def\modecorrige{1}\input{Eval_TP_Reseau_LAN.tex}"
\documentclass[10pt,a4paper]{article}

\usepackage[T1]{fontenc}
\usepackage[utf8]{inputenc}
\usepackage{lmodern}
\usepackage[french]{babel}
\usepackage[a4paper,hmargin=1.5cm,top=1.3cm,bottom=1.6cm,footskip=0.7cm]{geometry}
\usepackage[table]{xcolor}
\usepackage{booktabs,tabularx,array}
\usepackage{enumitem}
\usepackage{fancyhdr}
\usepackage[most]{tcolorbox}
\usepackage{tikz}
\usetikzlibrary{positioning}
\usepackage{titlesec}

\providecommand{\modecorrige}{0}
\providecommand{\variante}{A}
\newif\ifcorrige
\ifnum\modecorrige=1 \corrigetrue \else \corrigefalse \fi

\definecolor{ipred}{HTML}{C0392B}
\definecolor{cdcbar}{HTML}{A67C3D}
\definecolor{cdcbg}{HTML}{F4F1EA}
\definecolor{qbar}{HTML}{4A6B5D}
\definecolor{qbg}{HTML}{EDF1EF}
\definecolor{corr}{HTML}{1F5FA8}
\definecolor{corrbg}{HTML}{E8F0FA}
\definecolor{grisfoot}{HTML}{828282}

\newcommand{\ip}[1]{\textcolor{ipred}{\texttt{#1}}}
\newcommand{\pts}[1]{\hfill\mbox{\textbf{(#1)}}}
\newcommand{\sol}[1]{\ifcorrige\textcolor{corr}{#1}\fi}
\newcommand{\bareme}[1]{\ifcorrige\par{\footnotesize\color{corr}#1\par}\fi}
\newcommand{\hl}{\ifcorrige\rule[-1.6mm]{0pt}{5.4mm}\else\rule[-2.2mm]{0pt}{6.6mm}\fi}
\newcolumntype{C}{>{\centering\arraybackslash}X}
\newcolumntype{L}{>{\raggedright\arraybackslash}X}

\titleformat{\section}{\large\bfseries}{}{0pt}{}
\titlespacing*{\section}{0pt}{6pt}{3pt}

\tcbset{qstyle/.style={enhanced,boxrule=0pt,leftrule=2.5pt,arc=0pt,outer arc=0pt,
  left=5pt,right=5pt,top=2pt,bottom=2pt,before skip=4pt,after skip=3pt}}
\newtcolorbox{cdc}{qstyle,colback=cdcbg,colframe=cdcbar}
\newtcolorbox{quest}{qstyle,colback=qbg,colframe=qbar}

% Zone de réponse : vierge dans le sujet, remplie dans le corrigé
\newcommand{\zone}[2]{%
  \ifcorrige
    \begin{tcolorbox}[colback=corrbg,colframe=corr,boxrule=0.5pt,arc=1pt,left=4pt,right=4pt,top=1pt,bottom=1pt,
      before skip=2pt,after skip=2pt,fontupper=\footnotesize\color{corr}]#2\end{tcolorbox}%
  \else
    \begin{tcolorbox}[colback=white,colframe=black!45,boxrule=0.4pt,arc=1pt,height=#1,
      before skip=2pt,after skip=2pt]\end{tcolorbox}%
  \fi}

\newcommand{\question}[3]{\begin{quest}\textbf{#1.} #3\pts{#2}\end{quest}}

%==========================================================================
% Données propres à chaque variante
%==========================================================================
\if\variante A
  \def\entreprise{Le cabinet d'architectes \textbf{ArchiNova}}
  \def\Rnom{R-ArchiNova}
  \def\resUn{Bureau d'études}\def\pUn{BE}\def\netUn{192.168.40}
  \def\resDeux{Accueil}\def\pDeux{ACC}\def\netDeux{172.16}
  \def\srvNom{SRV-Intranet}\def\netSrv{192.168.200}\def\srvH{10}
  \def\ssid{ArchiNova-WiFi}
  \def\AIrows{%
    \hl\ip{192.168.40.25} & \sol{C} & \sol{255.255.255.0} & \sol{192.168.40.0} & \sol{192.168.40.255} & \sol{254} \\ \hline
    \hl\ip{172.16.5.10}   & \sol{B} & \sol{255.255.0.0}   & \sol{172.16.0.0}   & \sol{172.16.255.255} & \sol{65\,534} \\ \hline
    \hl\ip{10.3.2.1}      & \sol{A} & \sol{255.0.0.0}     & \sol{10.0.0.0}     & \sol{10.255.255.255} & \sol{16\,777\,214} \\ \hline}
  \def\apipa{169.254.12.7}
  \def\Cnet{192.168.60}
  \def\Crows{%
    PC-A & \ip{192.168.60.10} & \ip{255.255.255.0} & \ip{192.168.60.254} & — \\
    PC-B & \ip{192.168.60.20} & \ip{255.255.255.0} & \ip{192.168.60.1}   & — \\
    PC-C & \ip{192.168.61.30} & \ip{255.255.255.0} & \ip{192.168.60.254} & — \\
    PC-D & \ip{192.168.60.10} & \ip{255.255.255.0} & \ip{192.168.60.254} & — \\
    SRV-C & \ip{10.0.0.5}     & \ip{255.0.0.0}     & \ip{10.0.0.254}     & — \\
    R1 — G0/0 & \ip{192.168.60.254} & \ip{255.255.255.0} & — & On \\
    R1 — G0/1 & \ip{10.0.0.254}     & \ip{255.0.0.0}     & — & Off \\}
  \def\Csol{%
    \hl 1 & \sol{PC-B} & \sol{Passerelle .60.1 : pas l'adresse du routeur.} & \sol{Joint son réseau, mais aucun autre (serveur).} & \sol{Passerelle 192.168.60.254.} \\ \hline
    \hl 2 & \sol{PC-C} & \sol{IP hors du réseau 192.168.60.0/24.} & \sol{Ne joint ni les postes ni le routeur.} & \sol{Ex. 192.168.60.30.} \\ \hline
    \hl 3 & \sol{PC-D} & \sol{Même IP que PC-A (doublon).} & \sol{Conflit d'adresse : PC-A et PC-D perturbés.} & \sol{Adresse unique, ex. .60.40.} \\ \hline
    \hl 4 & \sol{R1 G0/1} & \sol{Interface désactivée (Off).} & \sol{Réseau du serveur injoignable pour tous.} & \sol{Port Status : On.} \\ \hline}
  \def\Cbareme{PC-A et SRV-C sont correctement configurés ; accepter PC-A cité à la place de PC-D pour le doublon.}
\else
  \def\entreprise{La clinique vétérinaire \textbf{VetoPlus}}
  \def\Rnom{R-VetoPlus}
  \def\resUn{Soins}\def\pUn{SOI}\def\netUn{192.168.70}
  \def\resDeux{Secrétariat}\def\pDeux{SEC}\def\netDeux{172.20}
  \def\srvNom{SRV-Dossiers}\def\netSrv{192.168.250}\def\srvH{20}
  \def\ssid{VetoPlus-WiFi}
  \def\AIrows{%
    \hl\ip{192.168.15.130} & \sol{C} & \sol{255.255.255.0} & \sol{192.168.15.0} & \sol{192.168.15.255} & \sol{254} \\ \hline
    \hl\ip{172.31.200.7}   & \sol{B} & \sol{255.255.0.0}   & \sol{172.31.0.0}   & \sol{172.31.255.255} & \sol{65\,534} \\ \hline
    \hl\ip{10.45.6.9}      & \sol{A} & \sol{255.0.0.0}     & \sol{10.0.0.0}     & \sol{10.255.255.255} & \sol{16\,777\,214} \\ \hline}
  \def\apipa{169.254.88.3}
  \def\Cnet{192.168.80}
  \def\Crows{%
    PC-A & \ip{192.168.8.11}   & \ip{255.255.255.0} & \ip{192.168.80.254} & — \\
    PC-B & \ip{192.168.80.12}  & \ip{255.255.255.0} & \ip{192.168.80.254} & — \\
    PC-C & \ip{192.168.80.13}  & \ip{255.255.255.0} & \ip{192.168.80.100} & — \\
    PC-D & \ip{192.168.80.12}  & \ip{255.255.255.0} & \ip{192.168.80.254} & — \\
    SRV-C & \ip{10.0.0.8}      & \ip{255.0.0.0}     & \ip{10.0.0.254}     & — \\
    R1 — G0/0 & \ip{192.168.80.254} & \ip{255.255.255.0} & — & Off \\
    R1 — G0/1 & \ip{10.0.0.254}     & \ip{255.0.0.0}     & — & On \\}
  \def\Csol{%
    \hl 1 & \sol{PC-A} & \sol{IP hors du réseau 192.168.80.0/24.} & \sol{Ne joint ni les postes ni le routeur.} & \sol{Ex. 192.168.80.11.} \\ \hline
    \hl 2 & \sol{PC-C} & \sol{Passerelle .80.100 : pas l'adresse du routeur.} & \sol{Joint son réseau, mais aucun autre (serveur).} & \sol{Passerelle 192.168.80.254.} \\ \hline
    \hl 3 & \sol{PC-D} & \sol{Même IP que PC-B (doublon).} & \sol{Conflit d'adresse : PC-B et PC-D perturbés.} & \sol{Adresse unique, ex. .80.14.} \\ \hline
    \hl 4 & \sol{R1 G0/0} & \sol{Interface côté LAN désactivée (Off).} & \sol{Aucun poste ne joint le routeur ni le serveur.} & \sol{Port Status : On.} \\ \hline}
  \def\Cbareme{PC-B et SRV-C sont correctement configurés ; accepter PC-B cité à la place de PC-D pour le doublon.}
\fi

\pagestyle{fancy}
\fancyhf{}
\renewcommand{\headrulewidth}{0pt}
\fancyfoot[L]{\footnotesize\itshape\color{grisfoot}Évaluation — Réseaux locaux — Sujet \variante\ifcorrige\ — CORRIGÉ\fi}
\fancyfoot[R]{\footnotesize\color{grisfoot}\thepage/3}

\setlist{itemsep=0pt,topsep=2pt,parsep=0pt}
\setlength{\parindent}{0pt}
\setlength{\parskip}{2pt}

\begin{document}

\begin{center}
  {\Large\bfseries Évaluation pratique — Réseaux locaux \quad
   \fcolorbox{black}{white}{\ SUJET \variante\ }\par}
  \vspace{1pt}
  {\bfseries Adressage IP, commutation et routage — BTS CIEL 1re année — durée : 2 heures\par}
  \ifcorrige\vspace{2pt}{\bfseries\color{corr}CORRIGÉ ET BARÈME — document enseignant}\par\fi
\end{center}
\vspace{-2pt}
{\renewcommand{\arraystretch}{1.35}
\begin{tabularx}{\linewidth}{|X|X|p{2.2cm}|p{2.6cm}|}
  \hline
  \textbf{Nom :} & \textbf{Prénom :} & \textbf{Poste n°} & \textbf{Note :} \hfill /20 \\ \hline
\end{tabularx}}

\textbf{Consignes :} travail \textbf{individuel}, documents et Internet \textbf{non autorisés}. Ce sujet est le
\textbf{document-réponse}. Fichier Packet Tracer : \ip{Nom\_Prenom\_EvalLAN\_\variante.pkt}, à enregistrer
régulièrement dans le dossier de rendu. Les « visas » sont apposés par le professeur, que vous appelez à chaque test réussi.
Durées conseillées : A 30 min — B 1 h 10 — C 20 min.

%==========================================================================
\section{Partie A — Questions de cours et calculs \hfill /6}

\question{A1}{2 pts}{Complétez le tableau en considérant le \textbf{masque par défaut} de la classe de chaque adresse.}
{\small\renewcommand{\arraystretch}{1.15}
\begin{tabularx}{\linewidth}{|l|c|C|C|C|c|}
  \hline
  \textbf{Adresse IP} & \textbf{Classe} & \textbf{Masque par défaut} & \textbf{Adresse réseau} &
  \textbf{Adresse de diffusion} & \textbf{Nb d'hôtes} \\ \hline
  \AIrows
\end{tabularx}}
\bareme{0,5 pt par ligne juste + 0,5 pt si les trois nombres d'hôtes sont justes ($2^{n}-2$, $n$ = nombre de bits d'hôte).}

\question{A2}{1 pt}{Indiquez le type de câble (\textbf{droit} ou \textbf{croisé}) pour chaque liaison.}
{\small\renewcommand{\arraystretch}{1.15}
\begin{tabularx}{\linewidth}{|X|c||X|c|}
  \hline
  \hl PC $\leftrightarrow$ commutateur & \makebox[1.6cm]{\sol{droit}} & Commutateur $\leftrightarrow$ routeur & \makebox[1.6cm]{\sol{droit}} \\ \hline
  \hl PC $\leftrightarrow$ PC (réseau P2P) & \sol{croisé} & Concentrateur $\leftrightarrow$ concentrateur & \sol{croisé} \\ \hline
\end{tabularx}}
\bareme{0,25 pt par case. Équipements de types différents : câble droit ; de même type : câble croisé.}

\question{A3}{1 pt}{Comment un \textbf{concentrateur (hub)} et un \textbf{commutateur (switch)} traitent-ils une trame
reçue ? Lequel limite les collisions, et pourquoi ?}
\zone{2.9cm}{Le \textbf{hub} (couche 1) répète la trame sur \textbf{tous ses autres ports} : un seul domaine de collision
partagé. Le \textbf{switch} (couche 2) lit l'adresse MAC de destination et, grâce à sa \textbf{table MAC}, ne l'envoie
\textbf{que sur le port du destinataire} : un domaine de collision par port, c'est lui qui limite les collisions.
\emph{0,5 pt hub, 0,5 pt switch.}}

\question{A4}{1 pt}{Qu'est-ce que la \textbf{passerelle par défaut} d'un poste ? Quelle condition son adresse doit-elle respecter ?}
\zone{2.4cm}{Adresse IP de l'interface du \textbf{routeur} à laquelle le poste envoie les paquets destinés à un
\textbf{autre réseau}. Elle doit \textbf{appartenir au même réseau} que le poste. \emph{0,5 + 0,5 pt.}}

\question{A5}{0,5 pt}{Un portable configuré en DHCP affiche \ip{\apipa} / \ip{255.255.0.0}. Que s'est-il passé ?}
\zone{1.5cm}{Aucun \textbf{serveur DHCP} n'a répondu : le poste s'est attribué une adresse \textbf{APIPA} (169.254.0.0/16).}

\question{A6}{0,5 pt}{Quel protocole utilise \texttt{ping} ? Nommez les deux messages échangés et la couche OSI des adresses IP.}
\zone{1.5cm}{\textbf{ICMP} ; \textbf{Echo Request} / \textbf{Echo Reply} ; adresses IP en \textbf{couche 3 (Réseau)}.}

%==========================================================================
\newpage
\section{Partie B — Réalisation sous Packet Tracer \hfill /10}

\begin{cdc}
\entreprise{} vous confie son réseau : \textbf{trois réseaux} reliés par un \textbf{routeur unique Cisco 2911}
nommé \ip{\Rnom} (interfaces G0/0, G0/1, G0/2).
\begin{itemize}[leftmargin=12pt]
  \item \textbf{Réseau \resUn} : \textbf{4 postes} PC-\pUn1 à PC-\pUn4, commutateur \ip{SW-\pUn}, adressage statique,
        réseau \ip{\netUn.0/24}.
  \item \textbf{Réseau \resDeux} : \textbf{2 postes} PC-\pDeux1, PC-\pDeux2 et \textbf{1 imprimante} IMP-\pDeux,
        commutateur \ip{SW-\pDeux}, adressage statique, réseau de classe B \ip{\netDeux.0.0/16}.
  \item \textbf{Réseau Serveur} : serveur \ip{\srvNom} d'adresse \ip{\netSrv.\srvH/24}, relié \textbf{directement} au routeur.
  \item La passerelle de chaque réseau est sa \textbf{dernière adresse utilisable}. Tous les équipements communiquent
        entre eux et accèdent à la page web du serveur. Nommage imposé (« Display Name »).
\end{itemize}
\end{cdc}

\question{B1}{2 pts}{Complétez le plan d'adressage \textbf{avant} la réalisation.}
{\footnotesize\setlength{\tabcolsep}{3pt}\renewcommand{\arraystretch}{1.15}
\begin{tabularx}{\linewidth}{|l|C|C|C|C|C|C|}
  \hline
  \textbf{Réseau} & \textbf{Adresse réseau} & \textbf{Masque} & \textbf{1re @ utilisable} &
  \textbf{Dernière @ utilisable} & \textbf{Diffusion} & \textbf{Passerelle} \\ \hline
  \hl \resUn & \sol{\netUn.0} & \sol{255.255.255.0} & \sol{\netUn.1} & \sol{\netUn.254} & \sol{\netUn.255} & \sol{\netUn.254} \\ \hline
  \hl \resDeux & \sol{\netDeux.0.0} & \sol{255.255.0.0} & \sol{\netDeux.0.1} & \sol{\netDeux.255.254} & \sol{\netDeux.255.255} & \sol{\netDeux.255.254} \\ \hline
  \hl Serveur & \sol{\netSrv.0} & \sol{255.255.255.0} & \sol{\netSrv.1} & \sol{\netSrv.254} & \sol{\netSrv.255} & \sol{\netSrv.254} \\ \hline
\end{tabularx}}
\bareme{0,5 pt par ligne juste + 0,5 pt si la convention de passerelle est respectée partout.}

\question{B2}{2 pts}{Construisez la topologie (équipements, câblage, nommage) et relevez l'adressage configuré.}
{\footnotesize\setlength{\tabcolsep}{3pt}\renewcommand{\arraystretch}{1.1}
\begin{tabularx}{\linewidth}{|l|C|C|C||l|C|C|}
  \hline
  \textbf{Hôte} & \textbf{Adresse IP} & \textbf{Masque} & \textbf{Passerelle} & \textbf{Interface} & \textbf{Adresse IP} & \textbf{Masque} \\ \hline
  \hl PC-\pUn1 & \sol{\netUn.1} & \sol{255.255.255.0} & \sol{\netUn.254} & \Rnom{} G0/0 & \sol{\netUn.254} & \sol{255.255.255.0} \\ \hline
  \hl PC-\pDeux1 & \sol{\netDeux.0.1} & \sol{255.255.0.0} & \sol{\netDeux.255.254} & \Rnom{} G0/1 & \sol{\netDeux.255.254} & \sol{255.255.0.0} \\ \hline
  \hl IMP-\pDeux & \sol{\netDeux.0.10} & \sol{255.255.0.0} & \sol{\netDeux.255.254} & \Rnom{} G0/2 & \sol{\netSrv.254} & \sol{255.255.255.0} \\ \hline
  \hl \srvNom & \ip{\netSrv.\srvH} & \sol{255.255.255.0} & \sol{\netSrv.254} & \multicolumn{3}{c|}{\cellcolor{black!8}} \\ \hline
\end{tabularx}}
\bareme{Adresses d'hôtes libres (uniques, dans le bon réseau) ; affectation G0/x libre. Sur le .pkt : équipements 0,5 ;
câblage 1 (câbles droits vers les switchs et switch–routeur, \textbf{croisé} serveur–routeur, liens verts) ; nommage 0,5.}

\question{B3}{3 pts}{Configurez tous les équipements : IP, masques, passerelles des hôtes ; interfaces du routeur
(« Port Status : On »).}
\bareme{Sur le .pkt : interfaces routeur 1 pt ; IP/masques des hôtes 1 pt ; passerelles (imprimante et serveur compris)
1 pt. $-0,25$ par erreur, sans descendre sous 0 par item.}

\question{B4}{3 pts}{Réalisez les tests, notez la commande et le résultat, puis faites viser chaque test réussi.}
{\small\renewcommand{\arraystretch}{1.1}
\begin{tabularx}{\linewidth}{|c|l|l|C|c|c|}
  \hline
  \textbf{N°} & \textbf{Source} & \textbf{Destination} & \textbf{Commande / action} & \textbf{Résultat} & \textbf{Visa prof.} \\ \hline
  \hl 1 & PC-\pUn1 & PC-\pUn4 & \sol{ping \netUn.4} & \sol{réussi} & \hspace{1.4cm} \\ \hline
  \hl 2 & PC-\pDeux1 & IMP-\pDeux & \sol{ping \netDeux.0.10} & \sol{réussi} & \\ \hline
  \hl 3 & PC-\pUn2 & PC-\pDeux2 & \sol{ping \netDeux.0.2} & \sol{réussi} & \\ \hline
  \hl 4 & PC-\pDeux2 & \srvNom & \sol{ping \netSrv.\srvH} & \sol{réussi} & \\ \hline
  \hl 5 & PC-\pUn3 & page web du serveur & \sol{navigateur : {\NoAutoSpacing http://\netSrv.\srvH}} & \sol{page affichée} & \\ \hline
\end{tabularx}}
\bareme{0,5 pt par test visé + 0,5 pt pour B5. Le premier ping inter-réseaux peut échouer (résolution ARP) ; les suivants doivent réussir.}

\question{B5}{incl. B4}{Test n°3 : par quels équipements passe le paquet ? Pourquoi PC-\pUn2 l'envoie-t-il au routeur ?}
\zone{2.3cm}{PC-\pUn2 $\to$ SW-\pUn{} $\to$ \Rnom{} $\to$ SW-\pDeux{} $\to$ PC-\pDeux2. Grâce à son masque, PC-\pUn2 constate que
la destination (\netDeux.0.0/16) n'est \textbf{pas dans son réseau} (\netUn.0/24) : il envoie le paquet à sa
\textbf{passerelle par défaut}, qui le route.}

%==========================================================================
\newpage
\section{Partie C — Diagnostic d'un réseau en panne \hfill /4}

Un technicien constate qu'\textbf{aucun poste ne parvient à joindre le serveur SRV-C} et que certains postes
communiquent mal entre eux. Les câbles sont du bon type.

\begin{center}
\begin{minipage}{0.37\linewidth}\centering
\begin{tikzpicture}[font=\footnotesize,
  eq/.style={draw,fill=qbg,rounded corners=2pt,minimum width=9mm,minimum height=5mm,inner sep=2pt},
  every path/.style={thick}]
  \node[eq] (SW) {SW-1};
  \node[eq,right=9mm of SW] (R) {R1};
  \node[eq,right=7mm of R] (SRV) {SRV-C};
  \node[eq,above left=5mm and -3mm of SW] (A) {PC-A};
  \node[eq,above right=5mm and -3mm of SW] (B) {PC-B};
  \node[eq,below left=5mm and -3mm of SW] (C) {PC-C};
  \node[eq,below right=5mm and -3mm of SW] (D) {PC-D};
  \foreach \n in {A,B,C,D} \draw (\n) -- (SW);
  \draw (SW) -- node[above,font=\tiny]{G0/0} (R);
  \draw (R) -- node[above,font=\tiny]{G0/1} (SRV);
\end{tikzpicture}
\end{minipage}\hfill
\begin{minipage}{0.61\linewidth}\centering
{\small\renewcommand{\arraystretch}{1.05}\setlength{\tabcolsep}{4pt}
\begin{tabular}{lllll}
  \toprule
  \textbf{Équipement} & \textbf{Adresse IP} & \textbf{Masque} & \textbf{Passerelle} & \textbf{Port} \\
  \midrule
  \Crows
  \bottomrule
\end{tabular}}
\end{minipage}
\end{center}

\question{C1}{4 pts}{Relevez les \textbf{quatre anomalies}. Pour chacune : conséquence et correction.}
{\small\renewcommand{\arraystretch}{1.1}
\ifcorrige\def\hc{}\else\def\hc{\rule[-10mm]{0pt}{14mm}}\fi
\begin{tabularx}{\linewidth}{|c|>{\raggedright\arraybackslash}p{2cm}|L|L|L|}
  \hline
  \textbf{N°} & \textbf{Équipement} & \textbf{Anomalie} & \textbf{Conséquence} & \textbf{Correction} \\ \hline
  \ifcorrige\Csol\else
  \hc 1 & & & & \\ \hline
  \hc 2 & & & & \\ \hline
  \hc 3 & & & & \\ \hline
  \hc 4 & & & & \\ \hline
  \fi
\end{tabularx}}
\bareme{1 pt par anomalie (0,25 équipement, anomalie, conséquence, correction ; ordre indifférent). \Cbareme}

%==========================================================================
\section{Bonus — Extension sans fil \hfill +1}

\begin{cdc}Ajoutez au réseau \resUn{} un \textbf{AccessPoint-PT} relié à SW-\pUn{} (SSID \ip{\ssid}, \textbf{WPA2-PSK})
et un portable \ip{Laptop-\pUn} équipé d'un module \texttt{WPC300N}, en adressage \textbf{statique}.
Vérifiez qu'il joint \srvNom.\end{cdc}
{\small\renewcommand{\arraystretch}{1.1}
\begin{tabularx}{\linewidth}{|l|C|C|C|c|}
  \hline
  \textbf{Équipement} & \textbf{Adresse IP} & \textbf{Masque} & \textbf{Passerelle} & \textbf{Visa prof.} \\ \hline
  \hl Laptop-\pUn & \sol{\netUn.50} & \sol{255.255.255.0} & \sol{\netUn.254} & \hspace{1.4cm} \\ \hline
\end{tabularx}}
\bareme{Le point d'accès fait office de pont : le portable est dans \netUn.0/24. Bonus si associé en WPA2-PSK et ping
vers \netSrv.\srvH{} réussi. Note plafonnée à 20.}

\end{document}
"
%   pdflatex -jobname=Eval_Reseau_CorrigeB "\def\variante{B}\def\modecorrige{1}% Évaluation pratique — Réseaux locaux (BTS CIEL 1re année) — sujets A et B
% Compilation (4 PDF) :
%   pdflatex -jobname=Eval_Reseau_SujetA   "\def\variante{A}\def\modecorrige{0}\input{Eval_TP_Reseau_LAN.tex}"
%   pdflatex -jobname=Eval_Reseau_SujetB   "\def\variante{B}\def\modecorrige{0}\input{Eval_TP_Reseau_LAN.tex}"
%   pdflatex -jobname=Eval_Reseau_CorrigeA "\def\variante{A}\def\modecorrige{1}\input{Eval_TP_Reseau_LAN.tex}"
%   pdflatex -jobname=Eval_Reseau_CorrigeB "\def\variante{B}\def\modecorrige{1}\input{Eval_TP_Reseau_LAN.tex}"
\documentclass[10pt,a4paper]{article}

\usepackage[T1]{fontenc}
\usepackage[utf8]{inputenc}
\usepackage{lmodern}
\usepackage[french]{babel}
\usepackage[a4paper,hmargin=1.5cm,top=1.3cm,bottom=1.6cm,footskip=0.7cm]{geometry}
\usepackage[table]{xcolor}
\usepackage{booktabs,tabularx,array}
\usepackage{enumitem}
\usepackage{fancyhdr}
\usepackage[most]{tcolorbox}
\usepackage{tikz}
\usetikzlibrary{positioning}
\usepackage{titlesec}

\providecommand{\modecorrige}{0}
\providecommand{\variante}{A}
\newif\ifcorrige
\ifnum\modecorrige=1 \corrigetrue \else \corrigefalse \fi

\definecolor{ipred}{HTML}{C0392B}
\definecolor{cdcbar}{HTML}{A67C3D}
\definecolor{cdcbg}{HTML}{F4F1EA}
\definecolor{qbar}{HTML}{4A6B5D}
\definecolor{qbg}{HTML}{EDF1EF}
\definecolor{corr}{HTML}{1F5FA8}
\definecolor{corrbg}{HTML}{E8F0FA}
\definecolor{grisfoot}{HTML}{828282}

\newcommand{\ip}[1]{\textcolor{ipred}{\texttt{#1}}}
\newcommand{\pts}[1]{\hfill\mbox{\textbf{(#1)}}}
\newcommand{\sol}[1]{\ifcorrige\textcolor{corr}{#1}\fi}
\newcommand{\bareme}[1]{\ifcorrige\par{\footnotesize\color{corr}#1\par}\fi}
\newcommand{\hl}{\ifcorrige\rule[-1.6mm]{0pt}{5.4mm}\else\rule[-2.2mm]{0pt}{6.6mm}\fi}
\newcolumntype{C}{>{\centering\arraybackslash}X}
\newcolumntype{L}{>{\raggedright\arraybackslash}X}

\titleformat{\section}{\large\bfseries}{}{0pt}{}
\titlespacing*{\section}{0pt}{6pt}{3pt}

\tcbset{qstyle/.style={enhanced,boxrule=0pt,leftrule=2.5pt,arc=0pt,outer arc=0pt,
  left=5pt,right=5pt,top=2pt,bottom=2pt,before skip=4pt,after skip=3pt}}
\newtcolorbox{cdc}{qstyle,colback=cdcbg,colframe=cdcbar}
\newtcolorbox{quest}{qstyle,colback=qbg,colframe=qbar}

% Zone de réponse : vierge dans le sujet, remplie dans le corrigé
\newcommand{\zone}[2]{%
  \ifcorrige
    \begin{tcolorbox}[colback=corrbg,colframe=corr,boxrule=0.5pt,arc=1pt,left=4pt,right=4pt,top=1pt,bottom=1pt,
      before skip=2pt,after skip=2pt,fontupper=\footnotesize\color{corr}]#2\end{tcolorbox}%
  \else
    \begin{tcolorbox}[colback=white,colframe=black!45,boxrule=0.4pt,arc=1pt,height=#1,
      before skip=2pt,after skip=2pt]\end{tcolorbox}%
  \fi}

\newcommand{\question}[3]{\begin{quest}\textbf{#1.} #3\pts{#2}\end{quest}}

%==========================================================================
% Données propres à chaque variante
%==========================================================================
\if\variante A
  \def\entreprise{Le cabinet d'architectes \textbf{ArchiNova}}
  \def\Rnom{R-ArchiNova}
  \def\resUn{Bureau d'études}\def\pUn{BE}\def\netUn{192.168.40}
  \def\resDeux{Accueil}\def\pDeux{ACC}\def\netDeux{172.16}
  \def\srvNom{SRV-Intranet}\def\netSrv{192.168.200}\def\srvH{10}
  \def\ssid{ArchiNova-WiFi}
  \def\AIrows{%
    \hl\ip{192.168.40.25} & \sol{C} & \sol{255.255.255.0} & \sol{192.168.40.0} & \sol{192.168.40.255} & \sol{254} \\ \hline
    \hl\ip{172.16.5.10}   & \sol{B} & \sol{255.255.0.0}   & \sol{172.16.0.0}   & \sol{172.16.255.255} & \sol{65\,534} \\ \hline
    \hl\ip{10.3.2.1}      & \sol{A} & \sol{255.0.0.0}     & \sol{10.0.0.0}     & \sol{10.255.255.255} & \sol{16\,777\,214} \\ \hline}
  \def\apipa{169.254.12.7}
  \def\Cnet{192.168.60}
  \def\Crows{%
    PC-A & \ip{192.168.60.10} & \ip{255.255.255.0} & \ip{192.168.60.254} & — \\
    PC-B & \ip{192.168.60.20} & \ip{255.255.255.0} & \ip{192.168.60.1}   & — \\
    PC-C & \ip{192.168.61.30} & \ip{255.255.255.0} & \ip{192.168.60.254} & — \\
    PC-D & \ip{192.168.60.10} & \ip{255.255.255.0} & \ip{192.168.60.254} & — \\
    SRV-C & \ip{10.0.0.5}     & \ip{255.0.0.0}     & \ip{10.0.0.254}     & — \\
    R1 — G0/0 & \ip{192.168.60.254} & \ip{255.255.255.0} & — & On \\
    R1 — G0/1 & \ip{10.0.0.254}     & \ip{255.0.0.0}     & — & Off \\}
  \def\Csol{%
    \hl 1 & \sol{PC-B} & \sol{Passerelle .60.1 : pas l'adresse du routeur.} & \sol{Joint son réseau, mais aucun autre (serveur).} & \sol{Passerelle 192.168.60.254.} \\ \hline
    \hl 2 & \sol{PC-C} & \sol{IP hors du réseau 192.168.60.0/24.} & \sol{Ne joint ni les postes ni le routeur.} & \sol{Ex. 192.168.60.30.} \\ \hline
    \hl 3 & \sol{PC-D} & \sol{Même IP que PC-A (doublon).} & \sol{Conflit d'adresse : PC-A et PC-D perturbés.} & \sol{Adresse unique, ex. .60.40.} \\ \hline
    \hl 4 & \sol{R1 G0/1} & \sol{Interface désactivée (Off).} & \sol{Réseau du serveur injoignable pour tous.} & \sol{Port Status : On.} \\ \hline}
  \def\Cbareme{PC-A et SRV-C sont correctement configurés ; accepter PC-A cité à la place de PC-D pour le doublon.}
\else
  \def\entreprise{La clinique vétérinaire \textbf{VetoPlus}}
  \def\Rnom{R-VetoPlus}
  \def\resUn{Soins}\def\pUn{SOI}\def\netUn{192.168.70}
  \def\resDeux{Secrétariat}\def\pDeux{SEC}\def\netDeux{172.20}
  \def\srvNom{SRV-Dossiers}\def\netSrv{192.168.250}\def\srvH{20}
  \def\ssid{VetoPlus-WiFi}
  \def\AIrows{%
    \hl\ip{192.168.15.130} & \sol{C} & \sol{255.255.255.0} & \sol{192.168.15.0} & \sol{192.168.15.255} & \sol{254} \\ \hline
    \hl\ip{172.31.200.7}   & \sol{B} & \sol{255.255.0.0}   & \sol{172.31.0.0}   & \sol{172.31.255.255} & \sol{65\,534} \\ \hline
    \hl\ip{10.45.6.9}      & \sol{A} & \sol{255.0.0.0}     & \sol{10.0.0.0}     & \sol{10.255.255.255} & \sol{16\,777\,214} \\ \hline}
  \def\apipa{169.254.88.3}
  \def\Cnet{192.168.80}
  \def\Crows{%
    PC-A & \ip{192.168.8.11}   & \ip{255.255.255.0} & \ip{192.168.80.254} & — \\
    PC-B & \ip{192.168.80.12}  & \ip{255.255.255.0} & \ip{192.168.80.254} & — \\
    PC-C & \ip{192.168.80.13}  & \ip{255.255.255.0} & \ip{192.168.80.100} & — \\
    PC-D & \ip{192.168.80.12}  & \ip{255.255.255.0} & \ip{192.168.80.254} & — \\
    SRV-C & \ip{10.0.0.8}      & \ip{255.0.0.0}     & \ip{10.0.0.254}     & — \\
    R1 — G0/0 & \ip{192.168.80.254} & \ip{255.255.255.0} & — & Off \\
    R1 — G0/1 & \ip{10.0.0.254}     & \ip{255.0.0.0}     & — & On \\}
  \def\Csol{%
    \hl 1 & \sol{PC-A} & \sol{IP hors du réseau 192.168.80.0/24.} & \sol{Ne joint ni les postes ni le routeur.} & \sol{Ex. 192.168.80.11.} \\ \hline
    \hl 2 & \sol{PC-C} & \sol{Passerelle .80.100 : pas l'adresse du routeur.} & \sol{Joint son réseau, mais aucun autre (serveur).} & \sol{Passerelle 192.168.80.254.} \\ \hline
    \hl 3 & \sol{PC-D} & \sol{Même IP que PC-B (doublon).} & \sol{Conflit d'adresse : PC-B et PC-D perturbés.} & \sol{Adresse unique, ex. .80.14.} \\ \hline
    \hl 4 & \sol{R1 G0/0} & \sol{Interface côté LAN désactivée (Off).} & \sol{Aucun poste ne joint le routeur ni le serveur.} & \sol{Port Status : On.} \\ \hline}
  \def\Cbareme{PC-B et SRV-C sont correctement configurés ; accepter PC-B cité à la place de PC-D pour le doublon.}
\fi

\pagestyle{fancy}
\fancyhf{}
\renewcommand{\headrulewidth}{0pt}
\fancyfoot[L]{\footnotesize\itshape\color{grisfoot}Évaluation — Réseaux locaux — Sujet \variante\ifcorrige\ — CORRIGÉ\fi}
\fancyfoot[R]{\footnotesize\color{grisfoot}\thepage/3}

\setlist{itemsep=0pt,topsep=2pt,parsep=0pt}
\setlength{\parindent}{0pt}
\setlength{\parskip}{2pt}

\begin{document}

\begin{center}
  {\Large\bfseries Évaluation pratique — Réseaux locaux \quad
   \fcolorbox{black}{white}{\ SUJET \variante\ }\par}
  \vspace{1pt}
  {\bfseries Adressage IP, commutation et routage — BTS CIEL 1re année — durée : 2 heures\par}
  \ifcorrige\vspace{2pt}{\bfseries\color{corr}CORRIGÉ ET BARÈME — document enseignant}\par\fi
\end{center}
\vspace{-2pt}
{\renewcommand{\arraystretch}{1.35}
\begin{tabularx}{\linewidth}{|X|X|p{2.2cm}|p{2.6cm}|}
  \hline
  \textbf{Nom :} & \textbf{Prénom :} & \textbf{Poste n°} & \textbf{Note :} \hfill /20 \\ \hline
\end{tabularx}}

\textbf{Consignes :} travail \textbf{individuel}, documents et Internet \textbf{non autorisés}. Ce sujet est le
\textbf{document-réponse}. Fichier Packet Tracer : \ip{Nom\_Prenom\_EvalLAN\_\variante.pkt}, à enregistrer
régulièrement dans le dossier de rendu. Les « visas » sont apposés par le professeur, que vous appelez à chaque test réussi.
Durées conseillées : A 30 min — B 1 h 10 — C 20 min.

%==========================================================================
\section{Partie A — Questions de cours et calculs \hfill /6}

\question{A1}{2 pts}{Complétez le tableau en considérant le \textbf{masque par défaut} de la classe de chaque adresse.}
{\small\renewcommand{\arraystretch}{1.15}
\begin{tabularx}{\linewidth}{|l|c|C|C|C|c|}
  \hline
  \textbf{Adresse IP} & \textbf{Classe} & \textbf{Masque par défaut} & \textbf{Adresse réseau} &
  \textbf{Adresse de diffusion} & \textbf{Nb d'hôtes} \\ \hline
  \AIrows
\end{tabularx}}
\bareme{0,5 pt par ligne juste + 0,5 pt si les trois nombres d'hôtes sont justes ($2^{n}-2$, $n$ = nombre de bits d'hôte).}

\question{A2}{1 pt}{Indiquez le type de câble (\textbf{droit} ou \textbf{croisé}) pour chaque liaison.}
{\small\renewcommand{\arraystretch}{1.15}
\begin{tabularx}{\linewidth}{|X|c||X|c|}
  \hline
  \hl PC $\leftrightarrow$ commutateur & \makebox[1.6cm]{\sol{droit}} & Commutateur $\leftrightarrow$ routeur & \makebox[1.6cm]{\sol{droit}} \\ \hline
  \hl PC $\leftrightarrow$ PC (réseau P2P) & \sol{croisé} & Concentrateur $\leftrightarrow$ concentrateur & \sol{croisé} \\ \hline
\end{tabularx}}
\bareme{0,25 pt par case. Équipements de types différents : câble droit ; de même type : câble croisé.}

\question{A3}{1 pt}{Comment un \textbf{concentrateur (hub)} et un \textbf{commutateur (switch)} traitent-ils une trame
reçue ? Lequel limite les collisions, et pourquoi ?}
\zone{2.9cm}{Le \textbf{hub} (couche 1) répète la trame sur \textbf{tous ses autres ports} : un seul domaine de collision
partagé. Le \textbf{switch} (couche 2) lit l'adresse MAC de destination et, grâce à sa \textbf{table MAC}, ne l'envoie
\textbf{que sur le port du destinataire} : un domaine de collision par port, c'est lui qui limite les collisions.
\emph{0,5 pt hub, 0,5 pt switch.}}

\question{A4}{1 pt}{Qu'est-ce que la \textbf{passerelle par défaut} d'un poste ? Quelle condition son adresse doit-elle respecter ?}
\zone{2.4cm}{Adresse IP de l'interface du \textbf{routeur} à laquelle le poste envoie les paquets destinés à un
\textbf{autre réseau}. Elle doit \textbf{appartenir au même réseau} que le poste. \emph{0,5 + 0,5 pt.}}

\question{A5}{0,5 pt}{Un portable configuré en DHCP affiche \ip{\apipa} / \ip{255.255.0.0}. Que s'est-il passé ?}
\zone{1.5cm}{Aucun \textbf{serveur DHCP} n'a répondu : le poste s'est attribué une adresse \textbf{APIPA} (169.254.0.0/16).}

\question{A6}{0,5 pt}{Quel protocole utilise \texttt{ping} ? Nommez les deux messages échangés et la couche OSI des adresses IP.}
\zone{1.5cm}{\textbf{ICMP} ; \textbf{Echo Request} / \textbf{Echo Reply} ; adresses IP en \textbf{couche 3 (Réseau)}.}

%==========================================================================
\newpage
\section{Partie B — Réalisation sous Packet Tracer \hfill /10}

\begin{cdc}
\entreprise{} vous confie son réseau : \textbf{trois réseaux} reliés par un \textbf{routeur unique Cisco 2911}
nommé \ip{\Rnom} (interfaces G0/0, G0/1, G0/2).
\begin{itemize}[leftmargin=12pt]
  \item \textbf{Réseau \resUn} : \textbf{4 postes} PC-\pUn1 à PC-\pUn4, commutateur \ip{SW-\pUn}, adressage statique,
        réseau \ip{\netUn.0/24}.
  \item \textbf{Réseau \resDeux} : \textbf{2 postes} PC-\pDeux1, PC-\pDeux2 et \textbf{1 imprimante} IMP-\pDeux,
        commutateur \ip{SW-\pDeux}, adressage statique, réseau de classe B \ip{\netDeux.0.0/16}.
  \item \textbf{Réseau Serveur} : serveur \ip{\srvNom} d'adresse \ip{\netSrv.\srvH/24}, relié \textbf{directement} au routeur.
  \item La passerelle de chaque réseau est sa \textbf{dernière adresse utilisable}. Tous les équipements communiquent
        entre eux et accèdent à la page web du serveur. Nommage imposé (« Display Name »).
\end{itemize}
\end{cdc}

\question{B1}{2 pts}{Complétez le plan d'adressage \textbf{avant} la réalisation.}
{\footnotesize\setlength{\tabcolsep}{3pt}\renewcommand{\arraystretch}{1.15}
\begin{tabularx}{\linewidth}{|l|C|C|C|C|C|C|}
  \hline
  \textbf{Réseau} & \textbf{Adresse réseau} & \textbf{Masque} & \textbf{1re @ utilisable} &
  \textbf{Dernière @ utilisable} & \textbf{Diffusion} & \textbf{Passerelle} \\ \hline
  \hl \resUn & \sol{\netUn.0} & \sol{255.255.255.0} & \sol{\netUn.1} & \sol{\netUn.254} & \sol{\netUn.255} & \sol{\netUn.254} \\ \hline
  \hl \resDeux & \sol{\netDeux.0.0} & \sol{255.255.0.0} & \sol{\netDeux.0.1} & \sol{\netDeux.255.254} & \sol{\netDeux.255.255} & \sol{\netDeux.255.254} \\ \hline
  \hl Serveur & \sol{\netSrv.0} & \sol{255.255.255.0} & \sol{\netSrv.1} & \sol{\netSrv.254} & \sol{\netSrv.255} & \sol{\netSrv.254} \\ \hline
\end{tabularx}}
\bareme{0,5 pt par ligne juste + 0,5 pt si la convention de passerelle est respectée partout.}

\question{B2}{2 pts}{Construisez la topologie (équipements, câblage, nommage) et relevez l'adressage configuré.}
{\footnotesize\setlength{\tabcolsep}{3pt}\renewcommand{\arraystretch}{1.1}
\begin{tabularx}{\linewidth}{|l|C|C|C||l|C|C|}
  \hline
  \textbf{Hôte} & \textbf{Adresse IP} & \textbf{Masque} & \textbf{Passerelle} & \textbf{Interface} & \textbf{Adresse IP} & \textbf{Masque} \\ \hline
  \hl PC-\pUn1 & \sol{\netUn.1} & \sol{255.255.255.0} & \sol{\netUn.254} & \Rnom{} G0/0 & \sol{\netUn.254} & \sol{255.255.255.0} \\ \hline
  \hl PC-\pDeux1 & \sol{\netDeux.0.1} & \sol{255.255.0.0} & \sol{\netDeux.255.254} & \Rnom{} G0/1 & \sol{\netDeux.255.254} & \sol{255.255.0.0} \\ \hline
  \hl IMP-\pDeux & \sol{\netDeux.0.10} & \sol{255.255.0.0} & \sol{\netDeux.255.254} & \Rnom{} G0/2 & \sol{\netSrv.254} & \sol{255.255.255.0} \\ \hline
  \hl \srvNom & \ip{\netSrv.\srvH} & \sol{255.255.255.0} & \sol{\netSrv.254} & \multicolumn{3}{c|}{\cellcolor{black!8}} \\ \hline
\end{tabularx}}
\bareme{Adresses d'hôtes libres (uniques, dans le bon réseau) ; affectation G0/x libre. Sur le .pkt : équipements 0,5 ;
câblage 1 (câbles droits vers les switchs et switch–routeur, \textbf{croisé} serveur–routeur, liens verts) ; nommage 0,5.}

\question{B3}{3 pts}{Configurez tous les équipements : IP, masques, passerelles des hôtes ; interfaces du routeur
(« Port Status : On »).}
\bareme{Sur le .pkt : interfaces routeur 1 pt ; IP/masques des hôtes 1 pt ; passerelles (imprimante et serveur compris)
1 pt. $-0,25$ par erreur, sans descendre sous 0 par item.}

\question{B4}{3 pts}{Réalisez les tests, notez la commande et le résultat, puis faites viser chaque test réussi.}
{\small\renewcommand{\arraystretch}{1.1}
\begin{tabularx}{\linewidth}{|c|l|l|C|c|c|}
  \hline
  \textbf{N°} & \textbf{Source} & \textbf{Destination} & \textbf{Commande / action} & \textbf{Résultat} & \textbf{Visa prof.} \\ \hline
  \hl 1 & PC-\pUn1 & PC-\pUn4 & \sol{ping \netUn.4} & \sol{réussi} & \hspace{1.4cm} \\ \hline
  \hl 2 & PC-\pDeux1 & IMP-\pDeux & \sol{ping \netDeux.0.10} & \sol{réussi} & \\ \hline
  \hl 3 & PC-\pUn2 & PC-\pDeux2 & \sol{ping \netDeux.0.2} & \sol{réussi} & \\ \hline
  \hl 4 & PC-\pDeux2 & \srvNom & \sol{ping \netSrv.\srvH} & \sol{réussi} & \\ \hline
  \hl 5 & PC-\pUn3 & page web du serveur & \sol{navigateur : {\NoAutoSpacing http://\netSrv.\srvH}} & \sol{page affichée} & \\ \hline
\end{tabularx}}
\bareme{0,5 pt par test visé + 0,5 pt pour B5. Le premier ping inter-réseaux peut échouer (résolution ARP) ; les suivants doivent réussir.}

\question{B5}{incl. B4}{Test n°3 : par quels équipements passe le paquet ? Pourquoi PC-\pUn2 l'envoie-t-il au routeur ?}
\zone{2.3cm}{PC-\pUn2 $\to$ SW-\pUn{} $\to$ \Rnom{} $\to$ SW-\pDeux{} $\to$ PC-\pDeux2. Grâce à son masque, PC-\pUn2 constate que
la destination (\netDeux.0.0/16) n'est \textbf{pas dans son réseau} (\netUn.0/24) : il envoie le paquet à sa
\textbf{passerelle par défaut}, qui le route.}

%==========================================================================
\newpage
\section{Partie C — Diagnostic d'un réseau en panne \hfill /4}

Un technicien constate qu'\textbf{aucun poste ne parvient à joindre le serveur SRV-C} et que certains postes
communiquent mal entre eux. Les câbles sont du bon type.

\begin{center}
\begin{minipage}{0.37\linewidth}\centering
\begin{tikzpicture}[font=\footnotesize,
  eq/.style={draw,fill=qbg,rounded corners=2pt,minimum width=9mm,minimum height=5mm,inner sep=2pt},
  every path/.style={thick}]
  \node[eq] (SW) {SW-1};
  \node[eq,right=9mm of SW] (R) {R1};
  \node[eq,right=7mm of R] (SRV) {SRV-C};
  \node[eq,above left=5mm and -3mm of SW] (A) {PC-A};
  \node[eq,above right=5mm and -3mm of SW] (B) {PC-B};
  \node[eq,below left=5mm and -3mm of SW] (C) {PC-C};
  \node[eq,below right=5mm and -3mm of SW] (D) {PC-D};
  \foreach \n in {A,B,C,D} \draw (\n) -- (SW);
  \draw (SW) -- node[above,font=\tiny]{G0/0} (R);
  \draw (R) -- node[above,font=\tiny]{G0/1} (SRV);
\end{tikzpicture}
\end{minipage}\hfill
\begin{minipage}{0.61\linewidth}\centering
{\small\renewcommand{\arraystretch}{1.05}\setlength{\tabcolsep}{4pt}
\begin{tabular}{lllll}
  \toprule
  \textbf{Équipement} & \textbf{Adresse IP} & \textbf{Masque} & \textbf{Passerelle} & \textbf{Port} \\
  \midrule
  \Crows
  \bottomrule
\end{tabular}}
\end{minipage}
\end{center}

\question{C1}{4 pts}{Relevez les \textbf{quatre anomalies}. Pour chacune : conséquence et correction.}
{\small\renewcommand{\arraystretch}{1.1}
\ifcorrige\def\hc{}\else\def\hc{\rule[-10mm]{0pt}{14mm}}\fi
\begin{tabularx}{\linewidth}{|c|>{\raggedright\arraybackslash}p{2cm}|L|L|L|}
  \hline
  \textbf{N°} & \textbf{Équipement} & \textbf{Anomalie} & \textbf{Conséquence} & \textbf{Correction} \\ \hline
  \ifcorrige\Csol\else
  \hc 1 & & & & \\ \hline
  \hc 2 & & & & \\ \hline
  \hc 3 & & & & \\ \hline
  \hc 4 & & & & \\ \hline
  \fi
\end{tabularx}}
\bareme{1 pt par anomalie (0,25 équipement, anomalie, conséquence, correction ; ordre indifférent). \Cbareme}

%==========================================================================
\section{Bonus — Extension sans fil \hfill +1}

\begin{cdc}Ajoutez au réseau \resUn{} un \textbf{AccessPoint-PT} relié à SW-\pUn{} (SSID \ip{\ssid}, \textbf{WPA2-PSK})
et un portable \ip{Laptop-\pUn} équipé d'un module \texttt{WPC300N}, en adressage \textbf{statique}.
Vérifiez qu'il joint \srvNom.\end{cdc}
{\small\renewcommand{\arraystretch}{1.1}
\begin{tabularx}{\linewidth}{|l|C|C|C|c|}
  \hline
  \textbf{Équipement} & \textbf{Adresse IP} & \textbf{Masque} & \textbf{Passerelle} & \textbf{Visa prof.} \\ \hline
  \hl Laptop-\pUn & \sol{\netUn.50} & \sol{255.255.255.0} & \sol{\netUn.254} & \hspace{1.4cm} \\ \hline
\end{tabularx}}
\bareme{Le point d'accès fait office de pont : le portable est dans \netUn.0/24. Bonus si associé en WPA2-PSK et ping
vers \netSrv.\srvH{} réussi. Note plafonnée à 20.}

\end{document}
"
\documentclass[10pt,a4paper]{article}

\usepackage[T1]{fontenc}
\usepackage[utf8]{inputenc}
\usepackage{lmodern}
\usepackage[french]{babel}
\usepackage[a4paper,hmargin=1.5cm,top=1.3cm,bottom=1.6cm,footskip=0.7cm]{geometry}
\usepackage[table]{xcolor}
\usepackage{booktabs,tabularx,array}
\usepackage{enumitem}
\usepackage{fancyhdr}
\usepackage[most]{tcolorbox}
\usepackage{tikz}
\usetikzlibrary{positioning}
\usepackage{titlesec}

\providecommand{\modecorrige}{0}
\providecommand{\variante}{A}
\newif\ifcorrige
\ifnum\modecorrige=1 \corrigetrue \else \corrigefalse \fi

\definecolor{ipred}{HTML}{C0392B}
\definecolor{cdcbar}{HTML}{A67C3D}
\definecolor{cdcbg}{HTML}{F4F1EA}
\definecolor{qbar}{HTML}{4A6B5D}
\definecolor{qbg}{HTML}{EDF1EF}
\definecolor{corr}{HTML}{1F5FA8}
\definecolor{corrbg}{HTML}{E8F0FA}
\definecolor{grisfoot}{HTML}{828282}

\newcommand{\ip}[1]{\textcolor{ipred}{\texttt{#1}}}
\newcommand{\pts}[1]{\hfill\mbox{\textbf{(#1)}}}
\newcommand{\sol}[1]{\ifcorrige\textcolor{corr}{#1}\fi}
\newcommand{\bareme}[1]{\ifcorrige\par{\footnotesize\color{corr}#1\par}\fi}
\newcommand{\hl}{\ifcorrige\rule[-1.6mm]{0pt}{5.4mm}\else\rule[-2.2mm]{0pt}{6.6mm}\fi}
\newcolumntype{C}{>{\centering\arraybackslash}X}
\newcolumntype{L}{>{\raggedright\arraybackslash}X}

\titleformat{\section}{\large\bfseries}{}{0pt}{}
\titlespacing*{\section}{0pt}{6pt}{3pt}

\tcbset{qstyle/.style={enhanced,boxrule=0pt,leftrule=2.5pt,arc=0pt,outer arc=0pt,
  left=5pt,right=5pt,top=2pt,bottom=2pt,before skip=4pt,after skip=3pt}}
\newtcolorbox{cdc}{qstyle,colback=cdcbg,colframe=cdcbar}
\newtcolorbox{quest}{qstyle,colback=qbg,colframe=qbar}

% Zone de réponse : vierge dans le sujet, remplie dans le corrigé
\newcommand{\zone}[2]{%
  \ifcorrige
    \begin{tcolorbox}[colback=corrbg,colframe=corr,boxrule=0.5pt,arc=1pt,left=4pt,right=4pt,top=1pt,bottom=1pt,
      before skip=2pt,after skip=2pt,fontupper=\footnotesize\color{corr}]#2\end{tcolorbox}%
  \else
    \begin{tcolorbox}[colback=white,colframe=black!45,boxrule=0.4pt,arc=1pt,height=#1,
      before skip=2pt,after skip=2pt]\end{tcolorbox}%
  \fi}

\newcommand{\question}[3]{\begin{quest}\textbf{#1.} #3\pts{#2}\end{quest}}

%==========================================================================
% Données propres à chaque variante
%==========================================================================
\if\variante A
  \def\entreprise{Le cabinet d'architectes \textbf{ArchiNova}}
  \def\Rnom{R-ArchiNova}
  \def\resUn{Bureau d'études}\def\pUn{BE}\def\netUn{192.168.40}
  \def\resDeux{Accueil}\def\pDeux{ACC}\def\netDeux{172.16}
  \def\srvNom{SRV-Intranet}\def\netSrv{192.168.200}\def\srvH{10}
  \def\ssid{ArchiNova-WiFi}
  \def\AIrows{%
    \hl\ip{192.168.40.25} & \sol{C} & \sol{255.255.255.0} & \sol{192.168.40.0} & \sol{192.168.40.255} & \sol{254} \\ \hline
    \hl\ip{172.16.5.10}   & \sol{B} & \sol{255.255.0.0}   & \sol{172.16.0.0}   & \sol{172.16.255.255} & \sol{65\,534} \\ \hline
    \hl\ip{10.3.2.1}      & \sol{A} & \sol{255.0.0.0}     & \sol{10.0.0.0}     & \sol{10.255.255.255} & \sol{16\,777\,214} \\ \hline}
  \def\apipa{169.254.12.7}
  \def\Cnet{192.168.60}
  \def\Crows{%
    PC-A & \ip{192.168.60.10} & \ip{255.255.255.0} & \ip{192.168.60.254} & — \\
    PC-B & \ip{192.168.60.20} & \ip{255.255.255.0} & \ip{192.168.60.1}   & — \\
    PC-C & \ip{192.168.61.30} & \ip{255.255.255.0} & \ip{192.168.60.254} & — \\
    PC-D & \ip{192.168.60.10} & \ip{255.255.255.0} & \ip{192.168.60.254} & — \\
    SRV-C & \ip{10.0.0.5}     & \ip{255.0.0.0}     & \ip{10.0.0.254}     & — \\
    R1 — G0/0 & \ip{192.168.60.254} & \ip{255.255.255.0} & — & On \\
    R1 — G0/1 & \ip{10.0.0.254}     & \ip{255.0.0.0}     & — & Off \\}
  \def\Csol{%
    \hl 1 & \sol{PC-B} & \sol{Passerelle .60.1 : pas l'adresse du routeur.} & \sol{Joint son réseau, mais aucun autre (serveur).} & \sol{Passerelle 192.168.60.254.} \\ \hline
    \hl 2 & \sol{PC-C} & \sol{IP hors du réseau 192.168.60.0/24.} & \sol{Ne joint ni les postes ni le routeur.} & \sol{Ex. 192.168.60.30.} \\ \hline
    \hl 3 & \sol{PC-D} & \sol{Même IP que PC-A (doublon).} & \sol{Conflit d'adresse : PC-A et PC-D perturbés.} & \sol{Adresse unique, ex. .60.40.} \\ \hline
    \hl 4 & \sol{R1 G0/1} & \sol{Interface désactivée (Off).} & \sol{Réseau du serveur injoignable pour tous.} & \sol{Port Status : On.} \\ \hline}
  \def\Cbareme{PC-A et SRV-C sont correctement configurés ; accepter PC-A cité à la place de PC-D pour le doublon.}
\else
  \def\entreprise{La clinique vétérinaire \textbf{VetoPlus}}
  \def\Rnom{R-VetoPlus}
  \def\resUn{Soins}\def\pUn{SOI}\def\netUn{192.168.70}
  \def\resDeux{Secrétariat}\def\pDeux{SEC}\def\netDeux{172.20}
  \def\srvNom{SRV-Dossiers}\def\netSrv{192.168.250}\def\srvH{20}
  \def\ssid{VetoPlus-WiFi}
  \def\AIrows{%
    \hl\ip{192.168.15.130} & \sol{C} & \sol{255.255.255.0} & \sol{192.168.15.0} & \sol{192.168.15.255} & \sol{254} \\ \hline
    \hl\ip{172.31.200.7}   & \sol{B} & \sol{255.255.0.0}   & \sol{172.31.0.0}   & \sol{172.31.255.255} & \sol{65\,534} \\ \hline
    \hl\ip{10.45.6.9}      & \sol{A} & \sol{255.0.0.0}     & \sol{10.0.0.0}     & \sol{10.255.255.255} & \sol{16\,777\,214} \\ \hline}
  \def\apipa{169.254.88.3}
  \def\Cnet{192.168.80}
  \def\Crows{%
    PC-A & \ip{192.168.8.11}   & \ip{255.255.255.0} & \ip{192.168.80.254} & — \\
    PC-B & \ip{192.168.80.12}  & \ip{255.255.255.0} & \ip{192.168.80.254} & — \\
    PC-C & \ip{192.168.80.13}  & \ip{255.255.255.0} & \ip{192.168.80.100} & — \\
    PC-D & \ip{192.168.80.12}  & \ip{255.255.255.0} & \ip{192.168.80.254} & — \\
    SRV-C & \ip{10.0.0.8}      & \ip{255.0.0.0}     & \ip{10.0.0.254}     & — \\
    R1 — G0/0 & \ip{192.168.80.254} & \ip{255.255.255.0} & — & Off \\
    R1 — G0/1 & \ip{10.0.0.254}     & \ip{255.0.0.0}     & — & On \\}
  \def\Csol{%
    \hl 1 & \sol{PC-A} & \sol{IP hors du réseau 192.168.80.0/24.} & \sol{Ne joint ni les postes ni le routeur.} & \sol{Ex. 192.168.80.11.} \\ \hline
    \hl 2 & \sol{PC-C} & \sol{Passerelle .80.100 : pas l'adresse du routeur.} & \sol{Joint son réseau, mais aucun autre (serveur).} & \sol{Passerelle 192.168.80.254.} \\ \hline
    \hl 3 & \sol{PC-D} & \sol{Même IP que PC-B (doublon).} & \sol{Conflit d'adresse : PC-B et PC-D perturbés.} & \sol{Adresse unique, ex. .80.14.} \\ \hline
    \hl 4 & \sol{R1 G0/0} & \sol{Interface côté LAN désactivée (Off).} & \sol{Aucun poste ne joint le routeur ni le serveur.} & \sol{Port Status : On.} \\ \hline}
  \def\Cbareme{PC-B et SRV-C sont correctement configurés ; accepter PC-B cité à la place de PC-D pour le doublon.}
\fi

\pagestyle{fancy}
\fancyhf{}
\renewcommand{\headrulewidth}{0pt}
\fancyfoot[L]{\footnotesize\itshape\color{grisfoot}Évaluation — Réseaux locaux — Sujet \variante\ifcorrige\ — CORRIGÉ\fi}
\fancyfoot[R]{\footnotesize\color{grisfoot}\thepage/3}

\setlist{itemsep=0pt,topsep=2pt,parsep=0pt}
\setlength{\parindent}{0pt}
\setlength{\parskip}{2pt}

\begin{document}

\begin{center}
  {\Large\bfseries Évaluation pratique — Réseaux locaux \quad
   \fcolorbox{black}{white}{\ SUJET \variante\ }\par}
  \vspace{1pt}
  {\bfseries Adressage IP, commutation et routage — BTS CIEL 1re année — durée : 2 heures\par}
  \ifcorrige\vspace{2pt}{\bfseries\color{corr}CORRIGÉ ET BARÈME — document enseignant}\par\fi
\end{center}
\vspace{-2pt}
{\renewcommand{\arraystretch}{1.35}
\begin{tabularx}{\linewidth}{|X|X|p{2.2cm}|p{2.6cm}|}
  \hline
  \textbf{Nom :} & \textbf{Prénom :} & \textbf{Poste n°} & \textbf{Note :} \hfill /20 \\ \hline
\end{tabularx}}

\textbf{Consignes :} travail \textbf{individuel}, documents et Internet \textbf{non autorisés}. Ce sujet est le
\textbf{document-réponse}. Fichier Packet Tracer : \ip{Nom\_Prenom\_EvalLAN\_\variante.pkt}, à enregistrer
régulièrement dans le dossier de rendu. Les « visas » sont apposés par le professeur, que vous appelez à chaque test réussi.
Durées conseillées : A 30 min — B 1 h 10 — C 20 min.

%==========================================================================
\section{Partie A — Questions de cours et calculs \hfill /6}

\question{A1}{2 pts}{Complétez le tableau en considérant le \textbf{masque par défaut} de la classe de chaque adresse.}
{\small\renewcommand{\arraystretch}{1.15}
\begin{tabularx}{\linewidth}{|l|c|C|C|C|c|}
  \hline
  \textbf{Adresse IP} & \textbf{Classe} & \textbf{Masque par défaut} & \textbf{Adresse réseau} &
  \textbf{Adresse de diffusion} & \textbf{Nb d'hôtes} \\ \hline
  \AIrows
\end{tabularx}}
\bareme{0,5 pt par ligne juste + 0,5 pt si les trois nombres d'hôtes sont justes ($2^{n}-2$, $n$ = nombre de bits d'hôte).}

\question{A2}{1 pt}{Indiquez le type de câble (\textbf{droit} ou \textbf{croisé}) pour chaque liaison.}
{\small\renewcommand{\arraystretch}{1.15}
\begin{tabularx}{\linewidth}{|X|c||X|c|}
  \hline
  \hl PC $\leftrightarrow$ commutateur & \makebox[1.6cm]{\sol{droit}} & Commutateur $\leftrightarrow$ routeur & \makebox[1.6cm]{\sol{droit}} \\ \hline
  \hl PC $\leftrightarrow$ PC (réseau P2P) & \sol{croisé} & Concentrateur $\leftrightarrow$ concentrateur & \sol{croisé} \\ \hline
\end{tabularx}}
\bareme{0,25 pt par case. Équipements de types différents : câble droit ; de même type : câble croisé.}

\question{A3}{1 pt}{Comment un \textbf{concentrateur (hub)} et un \textbf{commutateur (switch)} traitent-ils une trame
reçue ? Lequel limite les collisions, et pourquoi ?}
\zone{2.9cm}{Le \textbf{hub} (couche 1) répète la trame sur \textbf{tous ses autres ports} : un seul domaine de collision
partagé. Le \textbf{switch} (couche 2) lit l'adresse MAC de destination et, grâce à sa \textbf{table MAC}, ne l'envoie
\textbf{que sur le port du destinataire} : un domaine de collision par port, c'est lui qui limite les collisions.
\emph{0,5 pt hub, 0,5 pt switch.}}

\question{A4}{1 pt}{Qu'est-ce que la \textbf{passerelle par défaut} d'un poste ? Quelle condition son adresse doit-elle respecter ?}
\zone{2.4cm}{Adresse IP de l'interface du \textbf{routeur} à laquelle le poste envoie les paquets destinés à un
\textbf{autre réseau}. Elle doit \textbf{appartenir au même réseau} que le poste. \emph{0,5 + 0,5 pt.}}

\question{A5}{0,5 pt}{Un portable configuré en DHCP affiche \ip{\apipa} / \ip{255.255.0.0}. Que s'est-il passé ?}
\zone{1.5cm}{Aucun \textbf{serveur DHCP} n'a répondu : le poste s'est attribué une adresse \textbf{APIPA} (169.254.0.0/16).}

\question{A6}{0,5 pt}{Quel protocole utilise \texttt{ping} ? Nommez les deux messages échangés et la couche OSI des adresses IP.}
\zone{1.5cm}{\textbf{ICMP} ; \textbf{Echo Request} / \textbf{Echo Reply} ; adresses IP en \textbf{couche 3 (Réseau)}.}

%==========================================================================
\newpage
\section{Partie B — Réalisation sous Packet Tracer \hfill /10}

\begin{cdc}
\entreprise{} vous confie son réseau : \textbf{trois réseaux} reliés par un \textbf{routeur unique Cisco 2911}
nommé \ip{\Rnom} (interfaces G0/0, G0/1, G0/2).
\begin{itemize}[leftmargin=12pt]
  \item \textbf{Réseau \resUn} : \textbf{4 postes} PC-\pUn1 à PC-\pUn4, commutateur \ip{SW-\pUn}, adressage statique,
        réseau \ip{\netUn.0/24}.
  \item \textbf{Réseau \resDeux} : \textbf{2 postes} PC-\pDeux1, PC-\pDeux2 et \textbf{1 imprimante} IMP-\pDeux,
        commutateur \ip{SW-\pDeux}, adressage statique, réseau de classe B \ip{\netDeux.0.0/16}.
  \item \textbf{Réseau Serveur} : serveur \ip{\srvNom} d'adresse \ip{\netSrv.\srvH/24}, relié \textbf{directement} au routeur.
  \item La passerelle de chaque réseau est sa \textbf{dernière adresse utilisable}. Tous les équipements communiquent
        entre eux et accèdent à la page web du serveur. Nommage imposé (« Display Name »).
\end{itemize}
\end{cdc}

\question{B1}{2 pts}{Complétez le plan d'adressage \textbf{avant} la réalisation.}
{\footnotesize\setlength{\tabcolsep}{3pt}\renewcommand{\arraystretch}{1.15}
\begin{tabularx}{\linewidth}{|l|C|C|C|C|C|C|}
  \hline
  \textbf{Réseau} & \textbf{Adresse réseau} & \textbf{Masque} & \textbf{1re @ utilisable} &
  \textbf{Dernière @ utilisable} & \textbf{Diffusion} & \textbf{Passerelle} \\ \hline
  \hl \resUn & \sol{\netUn.0} & \sol{255.255.255.0} & \sol{\netUn.1} & \sol{\netUn.254} & \sol{\netUn.255} & \sol{\netUn.254} \\ \hline
  \hl \resDeux & \sol{\netDeux.0.0} & \sol{255.255.0.0} & \sol{\netDeux.0.1} & \sol{\netDeux.255.254} & \sol{\netDeux.255.255} & \sol{\netDeux.255.254} \\ \hline
  \hl Serveur & \sol{\netSrv.0} & \sol{255.255.255.0} & \sol{\netSrv.1} & \sol{\netSrv.254} & \sol{\netSrv.255} & \sol{\netSrv.254} \\ \hline
\end{tabularx}}
\bareme{0,5 pt par ligne juste + 0,5 pt si la convention de passerelle est respectée partout.}

\question{B2}{2 pts}{Construisez la topologie (équipements, câblage, nommage) et relevez l'adressage configuré.}
{\footnotesize\setlength{\tabcolsep}{3pt}\renewcommand{\arraystretch}{1.1}
\begin{tabularx}{\linewidth}{|l|C|C|C||l|C|C|}
  \hline
  \textbf{Hôte} & \textbf{Adresse IP} & \textbf{Masque} & \textbf{Passerelle} & \textbf{Interface} & \textbf{Adresse IP} & \textbf{Masque} \\ \hline
  \hl PC-\pUn1 & \sol{\netUn.1} & \sol{255.255.255.0} & \sol{\netUn.254} & \Rnom{} G0/0 & \sol{\netUn.254} & \sol{255.255.255.0} \\ \hline
  \hl PC-\pDeux1 & \sol{\netDeux.0.1} & \sol{255.255.0.0} & \sol{\netDeux.255.254} & \Rnom{} G0/1 & \sol{\netDeux.255.254} & \sol{255.255.0.0} \\ \hline
  \hl IMP-\pDeux & \sol{\netDeux.0.10} & \sol{255.255.0.0} & \sol{\netDeux.255.254} & \Rnom{} G0/2 & \sol{\netSrv.254} & \sol{255.255.255.0} \\ \hline
  \hl \srvNom & \ip{\netSrv.\srvH} & \sol{255.255.255.0} & \sol{\netSrv.254} & \multicolumn{3}{c|}{\cellcolor{black!8}} \\ \hline
\end{tabularx}}
\bareme{Adresses d'hôtes libres (uniques, dans le bon réseau) ; affectation G0/x libre. Sur le .pkt : équipements 0,5 ;
câblage 1 (câbles droits vers les switchs et switch–routeur, \textbf{croisé} serveur–routeur, liens verts) ; nommage 0,5.}

\question{B3}{3 pts}{Configurez tous les équipements : IP, masques, passerelles des hôtes ; interfaces du routeur
(« Port Status : On »).}
\bareme{Sur le .pkt : interfaces routeur 1 pt ; IP/masques des hôtes 1 pt ; passerelles (imprimante et serveur compris)
1 pt. $-0,25$ par erreur, sans descendre sous 0 par item.}

\question{B4}{3 pts}{Réalisez les tests, notez la commande et le résultat, puis faites viser chaque test réussi.}
{\small\renewcommand{\arraystretch}{1.1}
\begin{tabularx}{\linewidth}{|c|l|l|C|c|c|}
  \hline
  \textbf{N°} & \textbf{Source} & \textbf{Destination} & \textbf{Commande / action} & \textbf{Résultat} & \textbf{Visa prof.} \\ \hline
  \hl 1 & PC-\pUn1 & PC-\pUn4 & \sol{ping \netUn.4} & \sol{réussi} & \hspace{1.4cm} \\ \hline
  \hl 2 & PC-\pDeux1 & IMP-\pDeux & \sol{ping \netDeux.0.10} & \sol{réussi} & \\ \hline
  \hl 3 & PC-\pUn2 & PC-\pDeux2 & \sol{ping \netDeux.0.2} & \sol{réussi} & \\ \hline
  \hl 4 & PC-\pDeux2 & \srvNom & \sol{ping \netSrv.\srvH} & \sol{réussi} & \\ \hline
  \hl 5 & PC-\pUn3 & page web du serveur & \sol{navigateur : {\NoAutoSpacing http://\netSrv.\srvH}} & \sol{page affichée} & \\ \hline
\end{tabularx}}
\bareme{0,5 pt par test visé + 0,5 pt pour B5. Le premier ping inter-réseaux peut échouer (résolution ARP) ; les suivants doivent réussir.}

\question{B5}{incl. B4}{Test n°3 : par quels équipements passe le paquet ? Pourquoi PC-\pUn2 l'envoie-t-il au routeur ?}
\zone{2.3cm}{PC-\pUn2 $\to$ SW-\pUn{} $\to$ \Rnom{} $\to$ SW-\pDeux{} $\to$ PC-\pDeux2. Grâce à son masque, PC-\pUn2 constate que
la destination (\netDeux.0.0/16) n'est \textbf{pas dans son réseau} (\netUn.0/24) : il envoie le paquet à sa
\textbf{passerelle par défaut}, qui le route.}

%==========================================================================
\newpage
\section{Partie C — Diagnostic d'un réseau en panne \hfill /4}

Un technicien constate qu'\textbf{aucun poste ne parvient à joindre le serveur SRV-C} et que certains postes
communiquent mal entre eux. Les câbles sont du bon type.

\begin{center}
\begin{minipage}{0.37\linewidth}\centering
\begin{tikzpicture}[font=\footnotesize,
  eq/.style={draw,fill=qbg,rounded corners=2pt,minimum width=9mm,minimum height=5mm,inner sep=2pt},
  every path/.style={thick}]
  \node[eq] (SW) {SW-1};
  \node[eq,right=9mm of SW] (R) {R1};
  \node[eq,right=7mm of R] (SRV) {SRV-C};
  \node[eq,above left=5mm and -3mm of SW] (A) {PC-A};
  \node[eq,above right=5mm and -3mm of SW] (B) {PC-B};
  \node[eq,below left=5mm and -3mm of SW] (C) {PC-C};
  \node[eq,below right=5mm and -3mm of SW] (D) {PC-D};
  \foreach \n in {A,B,C,D} \draw (\n) -- (SW);
  \draw (SW) -- node[above,font=\tiny]{G0/0} (R);
  \draw (R) -- node[above,font=\tiny]{G0/1} (SRV);
\end{tikzpicture}
\end{minipage}\hfill
\begin{minipage}{0.61\linewidth}\centering
{\small\renewcommand{\arraystretch}{1.05}\setlength{\tabcolsep}{4pt}
\begin{tabular}{lllll}
  \toprule
  \textbf{Équipement} & \textbf{Adresse IP} & \textbf{Masque} & \textbf{Passerelle} & \textbf{Port} \\
  \midrule
  \Crows
  \bottomrule
\end{tabular}}
\end{minipage}
\end{center}

\question{C1}{4 pts}{Relevez les \textbf{quatre anomalies}. Pour chacune : conséquence et correction.}
{\small\renewcommand{\arraystretch}{1.1}
\ifcorrige\def\hc{}\else\def\hc{\rule[-10mm]{0pt}{14mm}}\fi
\begin{tabularx}{\linewidth}{|c|>{\raggedright\arraybackslash}p{2cm}|L|L|L|}
  \hline
  \textbf{N°} & \textbf{Équipement} & \textbf{Anomalie} & \textbf{Conséquence} & \textbf{Correction} \\ \hline
  \ifcorrige\Csol\else
  \hc 1 & & & & \\ \hline
  \hc 2 & & & & \\ \hline
  \hc 3 & & & & \\ \hline
  \hc 4 & & & & \\ \hline
  \fi
\end{tabularx}}
\bareme{1 pt par anomalie (0,25 équipement, anomalie, conséquence, correction ; ordre indifférent). \Cbareme}

%==========================================================================
\section{Bonus — Extension sans fil \hfill +1}

\begin{cdc}Ajoutez au réseau \resUn{} un \textbf{AccessPoint-PT} relié à SW-\pUn{} (SSID \ip{\ssid}, \textbf{WPA2-PSK})
et un portable \ip{Laptop-\pUn} équipé d'un module \texttt{WPC300N}, en adressage \textbf{statique}.
Vérifiez qu'il joint \srvNom.\end{cdc}
{\small\renewcommand{\arraystretch}{1.1}
\begin{tabularx}{\linewidth}{|l|C|C|C|c|}
  \hline
  \textbf{Équipement} & \textbf{Adresse IP} & \textbf{Masque} & \textbf{Passerelle} & \textbf{Visa prof.} \\ \hline
  \hl Laptop-\pUn & \sol{\netUn.50} & \sol{255.255.255.0} & \sol{\netUn.254} & \hspace{1.4cm} \\ \hline
\end{tabularx}}
\bareme{Le point d'accès fait office de pont : le portable est dans \netUn.0/24. Bonus si associé en WPA2-PSK et ping
vers \netSrv.\srvH{} réussi. Note plafonnée à 20.}

\end{document}
"
%   pdflatex -jobname=Eval_Reseau_CorrigeB "\def\variante{B}\def\modecorrige{1}% Évaluation pratique — Réseaux locaux (BTS CIEL 1re année) — sujets A et B
% Compilation (4 PDF) :
%   pdflatex -jobname=Eval_Reseau_SujetA   "\def\variante{A}\def\modecorrige{0}% Évaluation pratique — Réseaux locaux (BTS CIEL 1re année) — sujets A et B
% Compilation (4 PDF) :
%   pdflatex -jobname=Eval_Reseau_SujetA   "\def\variante{A}\def\modecorrige{0}\input{Eval_TP_Reseau_LAN.tex}"
%   pdflatex -jobname=Eval_Reseau_SujetB   "\def\variante{B}\def\modecorrige{0}\input{Eval_TP_Reseau_LAN.tex}"
%   pdflatex -jobname=Eval_Reseau_CorrigeA "\def\variante{A}\def\modecorrige{1}\input{Eval_TP_Reseau_LAN.tex}"
%   pdflatex -jobname=Eval_Reseau_CorrigeB "\def\variante{B}\def\modecorrige{1}\input{Eval_TP_Reseau_LAN.tex}"
\documentclass[10pt,a4paper]{article}

\usepackage[T1]{fontenc}
\usepackage[utf8]{inputenc}
\usepackage{lmodern}
\usepackage[french]{babel}
\usepackage[a4paper,hmargin=1.5cm,top=1.3cm,bottom=1.6cm,footskip=0.7cm]{geometry}
\usepackage[table]{xcolor}
\usepackage{booktabs,tabularx,array}
\usepackage{enumitem}
\usepackage{fancyhdr}
\usepackage[most]{tcolorbox}
\usepackage{tikz}
\usetikzlibrary{positioning}
\usepackage{titlesec}

\providecommand{\modecorrige}{0}
\providecommand{\variante}{A}
\newif\ifcorrige
\ifnum\modecorrige=1 \corrigetrue \else \corrigefalse \fi

\definecolor{ipred}{HTML}{C0392B}
\definecolor{cdcbar}{HTML}{A67C3D}
\definecolor{cdcbg}{HTML}{F4F1EA}
\definecolor{qbar}{HTML}{4A6B5D}
\definecolor{qbg}{HTML}{EDF1EF}
\definecolor{corr}{HTML}{1F5FA8}
\definecolor{corrbg}{HTML}{E8F0FA}
\definecolor{grisfoot}{HTML}{828282}

\newcommand{\ip}[1]{\textcolor{ipred}{\texttt{#1}}}
\newcommand{\pts}[1]{\hfill\mbox{\textbf{(#1)}}}
\newcommand{\sol}[1]{\ifcorrige\textcolor{corr}{#1}\fi}
\newcommand{\bareme}[1]{\ifcorrige\par{\footnotesize\color{corr}#1\par}\fi}
\newcommand{\hl}{\ifcorrige\rule[-1.6mm]{0pt}{5.4mm}\else\rule[-2.2mm]{0pt}{6.6mm}\fi}
\newcolumntype{C}{>{\centering\arraybackslash}X}
\newcolumntype{L}{>{\raggedright\arraybackslash}X}

\titleformat{\section}{\large\bfseries}{}{0pt}{}
\titlespacing*{\section}{0pt}{6pt}{3pt}

\tcbset{qstyle/.style={enhanced,boxrule=0pt,leftrule=2.5pt,arc=0pt,outer arc=0pt,
  left=5pt,right=5pt,top=2pt,bottom=2pt,before skip=4pt,after skip=3pt}}
\newtcolorbox{cdc}{qstyle,colback=cdcbg,colframe=cdcbar}
\newtcolorbox{quest}{qstyle,colback=qbg,colframe=qbar}

% Zone de réponse : vierge dans le sujet, remplie dans le corrigé
\newcommand{\zone}[2]{%
  \ifcorrige
    \begin{tcolorbox}[colback=corrbg,colframe=corr,boxrule=0.5pt,arc=1pt,left=4pt,right=4pt,top=1pt,bottom=1pt,
      before skip=2pt,after skip=2pt,fontupper=\footnotesize\color{corr}]#2\end{tcolorbox}%
  \else
    \begin{tcolorbox}[colback=white,colframe=black!45,boxrule=0.4pt,arc=1pt,height=#1,
      before skip=2pt,after skip=2pt]\end{tcolorbox}%
  \fi}

\newcommand{\question}[3]{\begin{quest}\textbf{#1.} #3\pts{#2}\end{quest}}

%==========================================================================
% Données propres à chaque variante
%==========================================================================
\if\variante A
  \def\entreprise{Le cabinet d'architectes \textbf{ArchiNova}}
  \def\Rnom{R-ArchiNova}
  \def\resUn{Bureau d'études}\def\pUn{BE}\def\netUn{192.168.40}
  \def\resDeux{Accueil}\def\pDeux{ACC}\def\netDeux{172.16}
  \def\srvNom{SRV-Intranet}\def\netSrv{192.168.200}\def\srvH{10}
  \def\ssid{ArchiNova-WiFi}
  \def\AIrows{%
    \hl\ip{192.168.40.25} & \sol{C} & \sol{255.255.255.0} & \sol{192.168.40.0} & \sol{192.168.40.255} & \sol{254} \\ \hline
    \hl\ip{172.16.5.10}   & \sol{B} & \sol{255.255.0.0}   & \sol{172.16.0.0}   & \sol{172.16.255.255} & \sol{65\,534} \\ \hline
    \hl\ip{10.3.2.1}      & \sol{A} & \sol{255.0.0.0}     & \sol{10.0.0.0}     & \sol{10.255.255.255} & \sol{16\,777\,214} \\ \hline}
  \def\apipa{169.254.12.7}
  \def\Cnet{192.168.60}
  \def\Crows{%
    PC-A & \ip{192.168.60.10} & \ip{255.255.255.0} & \ip{192.168.60.254} & — \\
    PC-B & \ip{192.168.60.20} & \ip{255.255.255.0} & \ip{192.168.60.1}   & — \\
    PC-C & \ip{192.168.61.30} & \ip{255.255.255.0} & \ip{192.168.60.254} & — \\
    PC-D & \ip{192.168.60.10} & \ip{255.255.255.0} & \ip{192.168.60.254} & — \\
    SRV-C & \ip{10.0.0.5}     & \ip{255.0.0.0}     & \ip{10.0.0.254}     & — \\
    R1 — G0/0 & \ip{192.168.60.254} & \ip{255.255.255.0} & — & On \\
    R1 — G0/1 & \ip{10.0.0.254}     & \ip{255.0.0.0}     & — & Off \\}
  \def\Csol{%
    \hl 1 & \sol{PC-B} & \sol{Passerelle .60.1 : pas l'adresse du routeur.} & \sol{Joint son réseau, mais aucun autre (serveur).} & \sol{Passerelle 192.168.60.254.} \\ \hline
    \hl 2 & \sol{PC-C} & \sol{IP hors du réseau 192.168.60.0/24.} & \sol{Ne joint ni les postes ni le routeur.} & \sol{Ex. 192.168.60.30.} \\ \hline
    \hl 3 & \sol{PC-D} & \sol{Même IP que PC-A (doublon).} & \sol{Conflit d'adresse : PC-A et PC-D perturbés.} & \sol{Adresse unique, ex. .60.40.} \\ \hline
    \hl 4 & \sol{R1 G0/1} & \sol{Interface désactivée (Off).} & \sol{Réseau du serveur injoignable pour tous.} & \sol{Port Status : On.} \\ \hline}
  \def\Cbareme{PC-A et SRV-C sont correctement configurés ; accepter PC-A cité à la place de PC-D pour le doublon.}
\else
  \def\entreprise{La clinique vétérinaire \textbf{VetoPlus}}
  \def\Rnom{R-VetoPlus}
  \def\resUn{Soins}\def\pUn{SOI}\def\netUn{192.168.70}
  \def\resDeux{Secrétariat}\def\pDeux{SEC}\def\netDeux{172.20}
  \def\srvNom{SRV-Dossiers}\def\netSrv{192.168.250}\def\srvH{20}
  \def\ssid{VetoPlus-WiFi}
  \def\AIrows{%
    \hl\ip{192.168.15.130} & \sol{C} & \sol{255.255.255.0} & \sol{192.168.15.0} & \sol{192.168.15.255} & \sol{254} \\ \hline
    \hl\ip{172.31.200.7}   & \sol{B} & \sol{255.255.0.0}   & \sol{172.31.0.0}   & \sol{172.31.255.255} & \sol{65\,534} \\ \hline
    \hl\ip{10.45.6.9}      & \sol{A} & \sol{255.0.0.0}     & \sol{10.0.0.0}     & \sol{10.255.255.255} & \sol{16\,777\,214} \\ \hline}
  \def\apipa{169.254.88.3}
  \def\Cnet{192.168.80}
  \def\Crows{%
    PC-A & \ip{192.168.8.11}   & \ip{255.255.255.0} & \ip{192.168.80.254} & — \\
    PC-B & \ip{192.168.80.12}  & \ip{255.255.255.0} & \ip{192.168.80.254} & — \\
    PC-C & \ip{192.168.80.13}  & \ip{255.255.255.0} & \ip{192.168.80.100} & — \\
    PC-D & \ip{192.168.80.12}  & \ip{255.255.255.0} & \ip{192.168.80.254} & — \\
    SRV-C & \ip{10.0.0.8}      & \ip{255.0.0.0}     & \ip{10.0.0.254}     & — \\
    R1 — G0/0 & \ip{192.168.80.254} & \ip{255.255.255.0} & — & Off \\
    R1 — G0/1 & \ip{10.0.0.254}     & \ip{255.0.0.0}     & — & On \\}
  \def\Csol{%
    \hl 1 & \sol{PC-A} & \sol{IP hors du réseau 192.168.80.0/24.} & \sol{Ne joint ni les postes ni le routeur.} & \sol{Ex. 192.168.80.11.} \\ \hline
    \hl 2 & \sol{PC-C} & \sol{Passerelle .80.100 : pas l'adresse du routeur.} & \sol{Joint son réseau, mais aucun autre (serveur).} & \sol{Passerelle 192.168.80.254.} \\ \hline
    \hl 3 & \sol{PC-D} & \sol{Même IP que PC-B (doublon).} & \sol{Conflit d'adresse : PC-B et PC-D perturbés.} & \sol{Adresse unique, ex. .80.14.} \\ \hline
    \hl 4 & \sol{R1 G0/0} & \sol{Interface côté LAN désactivée (Off).} & \sol{Aucun poste ne joint le routeur ni le serveur.} & \sol{Port Status : On.} \\ \hline}
  \def\Cbareme{PC-B et SRV-C sont correctement configurés ; accepter PC-B cité à la place de PC-D pour le doublon.}
\fi

\pagestyle{fancy}
\fancyhf{}
\renewcommand{\headrulewidth}{0pt}
\fancyfoot[L]{\footnotesize\itshape\color{grisfoot}Évaluation — Réseaux locaux — Sujet \variante\ifcorrige\ — CORRIGÉ\fi}
\fancyfoot[R]{\footnotesize\color{grisfoot}\thepage/3}

\setlist{itemsep=0pt,topsep=2pt,parsep=0pt}
\setlength{\parindent}{0pt}
\setlength{\parskip}{2pt}

\begin{document}

\begin{center}
  {\Large\bfseries Évaluation pratique — Réseaux locaux \quad
   \fcolorbox{black}{white}{\ SUJET \variante\ }\par}
  \vspace{1pt}
  {\bfseries Adressage IP, commutation et routage — BTS CIEL 1re année — durée : 2 heures\par}
  \ifcorrige\vspace{2pt}{\bfseries\color{corr}CORRIGÉ ET BARÈME — document enseignant}\par\fi
\end{center}
\vspace{-2pt}
{\renewcommand{\arraystretch}{1.35}
\begin{tabularx}{\linewidth}{|X|X|p{2.2cm}|p{2.6cm}|}
  \hline
  \textbf{Nom :} & \textbf{Prénom :} & \textbf{Poste n°} & \textbf{Note :} \hfill /20 \\ \hline
\end{tabularx}}

\textbf{Consignes :} travail \textbf{individuel}, documents et Internet \textbf{non autorisés}. Ce sujet est le
\textbf{document-réponse}. Fichier Packet Tracer : \ip{Nom\_Prenom\_EvalLAN\_\variante.pkt}, à enregistrer
régulièrement dans le dossier de rendu. Les « visas » sont apposés par le professeur, que vous appelez à chaque test réussi.
Durées conseillées : A 30 min — B 1 h 10 — C 20 min.

%==========================================================================
\section{Partie A — Questions de cours et calculs \hfill /6}

\question{A1}{2 pts}{Complétez le tableau en considérant le \textbf{masque par défaut} de la classe de chaque adresse.}
{\small\renewcommand{\arraystretch}{1.15}
\begin{tabularx}{\linewidth}{|l|c|C|C|C|c|}
  \hline
  \textbf{Adresse IP} & \textbf{Classe} & \textbf{Masque par défaut} & \textbf{Adresse réseau} &
  \textbf{Adresse de diffusion} & \textbf{Nb d'hôtes} \\ \hline
  \AIrows
\end{tabularx}}
\bareme{0,5 pt par ligne juste + 0,5 pt si les trois nombres d'hôtes sont justes ($2^{n}-2$, $n$ = nombre de bits d'hôte).}

\question{A2}{1 pt}{Indiquez le type de câble (\textbf{droit} ou \textbf{croisé}) pour chaque liaison.}
{\small\renewcommand{\arraystretch}{1.15}
\begin{tabularx}{\linewidth}{|X|c||X|c|}
  \hline
  \hl PC $\leftrightarrow$ commutateur & \makebox[1.6cm]{\sol{droit}} & Commutateur $\leftrightarrow$ routeur & \makebox[1.6cm]{\sol{droit}} \\ \hline
  \hl PC $\leftrightarrow$ PC (réseau P2P) & \sol{croisé} & Concentrateur $\leftrightarrow$ concentrateur & \sol{croisé} \\ \hline
\end{tabularx}}
\bareme{0,25 pt par case. Équipements de types différents : câble droit ; de même type : câble croisé.}

\question{A3}{1 pt}{Comment un \textbf{concentrateur (hub)} et un \textbf{commutateur (switch)} traitent-ils une trame
reçue ? Lequel limite les collisions, et pourquoi ?}
\zone{2.9cm}{Le \textbf{hub} (couche 1) répète la trame sur \textbf{tous ses autres ports} : un seul domaine de collision
partagé. Le \textbf{switch} (couche 2) lit l'adresse MAC de destination et, grâce à sa \textbf{table MAC}, ne l'envoie
\textbf{que sur le port du destinataire} : un domaine de collision par port, c'est lui qui limite les collisions.
\emph{0,5 pt hub, 0,5 pt switch.}}

\question{A4}{1 pt}{Qu'est-ce que la \textbf{passerelle par défaut} d'un poste ? Quelle condition son adresse doit-elle respecter ?}
\zone{2.4cm}{Adresse IP de l'interface du \textbf{routeur} à laquelle le poste envoie les paquets destinés à un
\textbf{autre réseau}. Elle doit \textbf{appartenir au même réseau} que le poste. \emph{0,5 + 0,5 pt.}}

\question{A5}{0,5 pt}{Un portable configuré en DHCP affiche \ip{\apipa} / \ip{255.255.0.0}. Que s'est-il passé ?}
\zone{1.5cm}{Aucun \textbf{serveur DHCP} n'a répondu : le poste s'est attribué une adresse \textbf{APIPA} (169.254.0.0/16).}

\question{A6}{0,5 pt}{Quel protocole utilise \texttt{ping} ? Nommez les deux messages échangés et la couche OSI des adresses IP.}
\zone{1.5cm}{\textbf{ICMP} ; \textbf{Echo Request} / \textbf{Echo Reply} ; adresses IP en \textbf{couche 3 (Réseau)}.}

%==========================================================================
\newpage
\section{Partie B — Réalisation sous Packet Tracer \hfill /10}

\begin{cdc}
\entreprise{} vous confie son réseau : \textbf{trois réseaux} reliés par un \textbf{routeur unique Cisco 2911}
nommé \ip{\Rnom} (interfaces G0/0, G0/1, G0/2).
\begin{itemize}[leftmargin=12pt]
  \item \textbf{Réseau \resUn} : \textbf{4 postes} PC-\pUn1 à PC-\pUn4, commutateur \ip{SW-\pUn}, adressage statique,
        réseau \ip{\netUn.0/24}.
  \item \textbf{Réseau \resDeux} : \textbf{2 postes} PC-\pDeux1, PC-\pDeux2 et \textbf{1 imprimante} IMP-\pDeux,
        commutateur \ip{SW-\pDeux}, adressage statique, réseau de classe B \ip{\netDeux.0.0/16}.
  \item \textbf{Réseau Serveur} : serveur \ip{\srvNom} d'adresse \ip{\netSrv.\srvH/24}, relié \textbf{directement} au routeur.
  \item La passerelle de chaque réseau est sa \textbf{dernière adresse utilisable}. Tous les équipements communiquent
        entre eux et accèdent à la page web du serveur. Nommage imposé (« Display Name »).
\end{itemize}
\end{cdc}

\question{B1}{2 pts}{Complétez le plan d'adressage \textbf{avant} la réalisation.}
{\footnotesize\setlength{\tabcolsep}{3pt}\renewcommand{\arraystretch}{1.15}
\begin{tabularx}{\linewidth}{|l|C|C|C|C|C|C|}
  \hline
  \textbf{Réseau} & \textbf{Adresse réseau} & \textbf{Masque} & \textbf{1re @ utilisable} &
  \textbf{Dernière @ utilisable} & \textbf{Diffusion} & \textbf{Passerelle} \\ \hline
  \hl \resUn & \sol{\netUn.0} & \sol{255.255.255.0} & \sol{\netUn.1} & \sol{\netUn.254} & \sol{\netUn.255} & \sol{\netUn.254} \\ \hline
  \hl \resDeux & \sol{\netDeux.0.0} & \sol{255.255.0.0} & \sol{\netDeux.0.1} & \sol{\netDeux.255.254} & \sol{\netDeux.255.255} & \sol{\netDeux.255.254} \\ \hline
  \hl Serveur & \sol{\netSrv.0} & \sol{255.255.255.0} & \sol{\netSrv.1} & \sol{\netSrv.254} & \sol{\netSrv.255} & \sol{\netSrv.254} \\ \hline
\end{tabularx}}
\bareme{0,5 pt par ligne juste + 0,5 pt si la convention de passerelle est respectée partout.}

\question{B2}{2 pts}{Construisez la topologie (équipements, câblage, nommage) et relevez l'adressage configuré.}
{\footnotesize\setlength{\tabcolsep}{3pt}\renewcommand{\arraystretch}{1.1}
\begin{tabularx}{\linewidth}{|l|C|C|C||l|C|C|}
  \hline
  \textbf{Hôte} & \textbf{Adresse IP} & \textbf{Masque} & \textbf{Passerelle} & \textbf{Interface} & \textbf{Adresse IP} & \textbf{Masque} \\ \hline
  \hl PC-\pUn1 & \sol{\netUn.1} & \sol{255.255.255.0} & \sol{\netUn.254} & \Rnom{} G0/0 & \sol{\netUn.254} & \sol{255.255.255.0} \\ \hline
  \hl PC-\pDeux1 & \sol{\netDeux.0.1} & \sol{255.255.0.0} & \sol{\netDeux.255.254} & \Rnom{} G0/1 & \sol{\netDeux.255.254} & \sol{255.255.0.0} \\ \hline
  \hl IMP-\pDeux & \sol{\netDeux.0.10} & \sol{255.255.0.0} & \sol{\netDeux.255.254} & \Rnom{} G0/2 & \sol{\netSrv.254} & \sol{255.255.255.0} \\ \hline
  \hl \srvNom & \ip{\netSrv.\srvH} & \sol{255.255.255.0} & \sol{\netSrv.254} & \multicolumn{3}{c|}{\cellcolor{black!8}} \\ \hline
\end{tabularx}}
\bareme{Adresses d'hôtes libres (uniques, dans le bon réseau) ; affectation G0/x libre. Sur le .pkt : équipements 0,5 ;
câblage 1 (câbles droits vers les switchs et switch–routeur, \textbf{croisé} serveur–routeur, liens verts) ; nommage 0,5.}

\question{B3}{3 pts}{Configurez tous les équipements : IP, masques, passerelles des hôtes ; interfaces du routeur
(« Port Status : On »).}
\bareme{Sur le .pkt : interfaces routeur 1 pt ; IP/masques des hôtes 1 pt ; passerelles (imprimante et serveur compris)
1 pt. $-0,25$ par erreur, sans descendre sous 0 par item.}

\question{B4}{3 pts}{Réalisez les tests, notez la commande et le résultat, puis faites viser chaque test réussi.}
{\small\renewcommand{\arraystretch}{1.1}
\begin{tabularx}{\linewidth}{|c|l|l|C|c|c|}
  \hline
  \textbf{N°} & \textbf{Source} & \textbf{Destination} & \textbf{Commande / action} & \textbf{Résultat} & \textbf{Visa prof.} \\ \hline
  \hl 1 & PC-\pUn1 & PC-\pUn4 & \sol{ping \netUn.4} & \sol{réussi} & \hspace{1.4cm} \\ \hline
  \hl 2 & PC-\pDeux1 & IMP-\pDeux & \sol{ping \netDeux.0.10} & \sol{réussi} & \\ \hline
  \hl 3 & PC-\pUn2 & PC-\pDeux2 & \sol{ping \netDeux.0.2} & \sol{réussi} & \\ \hline
  \hl 4 & PC-\pDeux2 & \srvNom & \sol{ping \netSrv.\srvH} & \sol{réussi} & \\ \hline
  \hl 5 & PC-\pUn3 & page web du serveur & \sol{navigateur : {\NoAutoSpacing http://\netSrv.\srvH}} & \sol{page affichée} & \\ \hline
\end{tabularx}}
\bareme{0,5 pt par test visé + 0,5 pt pour B5. Le premier ping inter-réseaux peut échouer (résolution ARP) ; les suivants doivent réussir.}

\question{B5}{incl. B4}{Test n°3 : par quels équipements passe le paquet ? Pourquoi PC-\pUn2 l'envoie-t-il au routeur ?}
\zone{2.3cm}{PC-\pUn2 $\to$ SW-\pUn{} $\to$ \Rnom{} $\to$ SW-\pDeux{} $\to$ PC-\pDeux2. Grâce à son masque, PC-\pUn2 constate que
la destination (\netDeux.0.0/16) n'est \textbf{pas dans son réseau} (\netUn.0/24) : il envoie le paquet à sa
\textbf{passerelle par défaut}, qui le route.}

%==========================================================================
\newpage
\section{Partie C — Diagnostic d'un réseau en panne \hfill /4}

Un technicien constate qu'\textbf{aucun poste ne parvient à joindre le serveur SRV-C} et que certains postes
communiquent mal entre eux. Les câbles sont du bon type.

\begin{center}
\begin{minipage}{0.37\linewidth}\centering
\begin{tikzpicture}[font=\footnotesize,
  eq/.style={draw,fill=qbg,rounded corners=2pt,minimum width=9mm,minimum height=5mm,inner sep=2pt},
  every path/.style={thick}]
  \node[eq] (SW) {SW-1};
  \node[eq,right=9mm of SW] (R) {R1};
  \node[eq,right=7mm of R] (SRV) {SRV-C};
  \node[eq,above left=5mm and -3mm of SW] (A) {PC-A};
  \node[eq,above right=5mm and -3mm of SW] (B) {PC-B};
  \node[eq,below left=5mm and -3mm of SW] (C) {PC-C};
  \node[eq,below right=5mm and -3mm of SW] (D) {PC-D};
  \foreach \n in {A,B,C,D} \draw (\n) -- (SW);
  \draw (SW) -- node[above,font=\tiny]{G0/0} (R);
  \draw (R) -- node[above,font=\tiny]{G0/1} (SRV);
\end{tikzpicture}
\end{minipage}\hfill
\begin{minipage}{0.61\linewidth}\centering
{\small\renewcommand{\arraystretch}{1.05}\setlength{\tabcolsep}{4pt}
\begin{tabular}{lllll}
  \toprule
  \textbf{Équipement} & \textbf{Adresse IP} & \textbf{Masque} & \textbf{Passerelle} & \textbf{Port} \\
  \midrule
  \Crows
  \bottomrule
\end{tabular}}
\end{minipage}
\end{center}

\question{C1}{4 pts}{Relevez les \textbf{quatre anomalies}. Pour chacune : conséquence et correction.}
{\small\renewcommand{\arraystretch}{1.1}
\ifcorrige\def\hc{}\else\def\hc{\rule[-10mm]{0pt}{14mm}}\fi
\begin{tabularx}{\linewidth}{|c|>{\raggedright\arraybackslash}p{2cm}|L|L|L|}
  \hline
  \textbf{N°} & \textbf{Équipement} & \textbf{Anomalie} & \textbf{Conséquence} & \textbf{Correction} \\ \hline
  \ifcorrige\Csol\else
  \hc 1 & & & & \\ \hline
  \hc 2 & & & & \\ \hline
  \hc 3 & & & & \\ \hline
  \hc 4 & & & & \\ \hline
  \fi
\end{tabularx}}
\bareme{1 pt par anomalie (0,25 équipement, anomalie, conséquence, correction ; ordre indifférent). \Cbareme}

%==========================================================================
\section{Bonus — Extension sans fil \hfill +1}

\begin{cdc}Ajoutez au réseau \resUn{} un \textbf{AccessPoint-PT} relié à SW-\pUn{} (SSID \ip{\ssid}, \textbf{WPA2-PSK})
et un portable \ip{Laptop-\pUn} équipé d'un module \texttt{WPC300N}, en adressage \textbf{statique}.
Vérifiez qu'il joint \srvNom.\end{cdc}
{\small\renewcommand{\arraystretch}{1.1}
\begin{tabularx}{\linewidth}{|l|C|C|C|c|}
  \hline
  \textbf{Équipement} & \textbf{Adresse IP} & \textbf{Masque} & \textbf{Passerelle} & \textbf{Visa prof.} \\ \hline
  \hl Laptop-\pUn & \sol{\netUn.50} & \sol{255.255.255.0} & \sol{\netUn.254} & \hspace{1.4cm} \\ \hline
\end{tabularx}}
\bareme{Le point d'accès fait office de pont : le portable est dans \netUn.0/24. Bonus si associé en WPA2-PSK et ping
vers \netSrv.\srvH{} réussi. Note plafonnée à 20.}

\end{document}
"
%   pdflatex -jobname=Eval_Reseau_SujetB   "\def\variante{B}\def\modecorrige{0}% Évaluation pratique — Réseaux locaux (BTS CIEL 1re année) — sujets A et B
% Compilation (4 PDF) :
%   pdflatex -jobname=Eval_Reseau_SujetA   "\def\variante{A}\def\modecorrige{0}\input{Eval_TP_Reseau_LAN.tex}"
%   pdflatex -jobname=Eval_Reseau_SujetB   "\def\variante{B}\def\modecorrige{0}\input{Eval_TP_Reseau_LAN.tex}"
%   pdflatex -jobname=Eval_Reseau_CorrigeA "\def\variante{A}\def\modecorrige{1}\input{Eval_TP_Reseau_LAN.tex}"
%   pdflatex -jobname=Eval_Reseau_CorrigeB "\def\variante{B}\def\modecorrige{1}\input{Eval_TP_Reseau_LAN.tex}"
\documentclass[10pt,a4paper]{article}

\usepackage[T1]{fontenc}
\usepackage[utf8]{inputenc}
\usepackage{lmodern}
\usepackage[french]{babel}
\usepackage[a4paper,hmargin=1.5cm,top=1.3cm,bottom=1.6cm,footskip=0.7cm]{geometry}
\usepackage[table]{xcolor}
\usepackage{booktabs,tabularx,array}
\usepackage{enumitem}
\usepackage{fancyhdr}
\usepackage[most]{tcolorbox}
\usepackage{tikz}
\usetikzlibrary{positioning}
\usepackage{titlesec}

\providecommand{\modecorrige}{0}
\providecommand{\variante}{A}
\newif\ifcorrige
\ifnum\modecorrige=1 \corrigetrue \else \corrigefalse \fi

\definecolor{ipred}{HTML}{C0392B}
\definecolor{cdcbar}{HTML}{A67C3D}
\definecolor{cdcbg}{HTML}{F4F1EA}
\definecolor{qbar}{HTML}{4A6B5D}
\definecolor{qbg}{HTML}{EDF1EF}
\definecolor{corr}{HTML}{1F5FA8}
\definecolor{corrbg}{HTML}{E8F0FA}
\definecolor{grisfoot}{HTML}{828282}

\newcommand{\ip}[1]{\textcolor{ipred}{\texttt{#1}}}
\newcommand{\pts}[1]{\hfill\mbox{\textbf{(#1)}}}
\newcommand{\sol}[1]{\ifcorrige\textcolor{corr}{#1}\fi}
\newcommand{\bareme}[1]{\ifcorrige\par{\footnotesize\color{corr}#1\par}\fi}
\newcommand{\hl}{\ifcorrige\rule[-1.6mm]{0pt}{5.4mm}\else\rule[-2.2mm]{0pt}{6.6mm}\fi}
\newcolumntype{C}{>{\centering\arraybackslash}X}
\newcolumntype{L}{>{\raggedright\arraybackslash}X}

\titleformat{\section}{\large\bfseries}{}{0pt}{}
\titlespacing*{\section}{0pt}{6pt}{3pt}

\tcbset{qstyle/.style={enhanced,boxrule=0pt,leftrule=2.5pt,arc=0pt,outer arc=0pt,
  left=5pt,right=5pt,top=2pt,bottom=2pt,before skip=4pt,after skip=3pt}}
\newtcolorbox{cdc}{qstyle,colback=cdcbg,colframe=cdcbar}
\newtcolorbox{quest}{qstyle,colback=qbg,colframe=qbar}

% Zone de réponse : vierge dans le sujet, remplie dans le corrigé
\newcommand{\zone}[2]{%
  \ifcorrige
    \begin{tcolorbox}[colback=corrbg,colframe=corr,boxrule=0.5pt,arc=1pt,left=4pt,right=4pt,top=1pt,bottom=1pt,
      before skip=2pt,after skip=2pt,fontupper=\footnotesize\color{corr}]#2\end{tcolorbox}%
  \else
    \begin{tcolorbox}[colback=white,colframe=black!45,boxrule=0.4pt,arc=1pt,height=#1,
      before skip=2pt,after skip=2pt]\end{tcolorbox}%
  \fi}

\newcommand{\question}[3]{\begin{quest}\textbf{#1.} #3\pts{#2}\end{quest}}

%==========================================================================
% Données propres à chaque variante
%==========================================================================
\if\variante A
  \def\entreprise{Le cabinet d'architectes \textbf{ArchiNova}}
  \def\Rnom{R-ArchiNova}
  \def\resUn{Bureau d'études}\def\pUn{BE}\def\netUn{192.168.40}
  \def\resDeux{Accueil}\def\pDeux{ACC}\def\netDeux{172.16}
  \def\srvNom{SRV-Intranet}\def\netSrv{192.168.200}\def\srvH{10}
  \def\ssid{ArchiNova-WiFi}
  \def\AIrows{%
    \hl\ip{192.168.40.25} & \sol{C} & \sol{255.255.255.0} & \sol{192.168.40.0} & \sol{192.168.40.255} & \sol{254} \\ \hline
    \hl\ip{172.16.5.10}   & \sol{B} & \sol{255.255.0.0}   & \sol{172.16.0.0}   & \sol{172.16.255.255} & \sol{65\,534} \\ \hline
    \hl\ip{10.3.2.1}      & \sol{A} & \sol{255.0.0.0}     & \sol{10.0.0.0}     & \sol{10.255.255.255} & \sol{16\,777\,214} \\ \hline}
  \def\apipa{169.254.12.7}
  \def\Cnet{192.168.60}
  \def\Crows{%
    PC-A & \ip{192.168.60.10} & \ip{255.255.255.0} & \ip{192.168.60.254} & — \\
    PC-B & \ip{192.168.60.20} & \ip{255.255.255.0} & \ip{192.168.60.1}   & — \\
    PC-C & \ip{192.168.61.30} & \ip{255.255.255.0} & \ip{192.168.60.254} & — \\
    PC-D & \ip{192.168.60.10} & \ip{255.255.255.0} & \ip{192.168.60.254} & — \\
    SRV-C & \ip{10.0.0.5}     & \ip{255.0.0.0}     & \ip{10.0.0.254}     & — \\
    R1 — G0/0 & \ip{192.168.60.254} & \ip{255.255.255.0} & — & On \\
    R1 — G0/1 & \ip{10.0.0.254}     & \ip{255.0.0.0}     & — & Off \\}
  \def\Csol{%
    \hl 1 & \sol{PC-B} & \sol{Passerelle .60.1 : pas l'adresse du routeur.} & \sol{Joint son réseau, mais aucun autre (serveur).} & \sol{Passerelle 192.168.60.254.} \\ \hline
    \hl 2 & \sol{PC-C} & \sol{IP hors du réseau 192.168.60.0/24.} & \sol{Ne joint ni les postes ni le routeur.} & \sol{Ex. 192.168.60.30.} \\ \hline
    \hl 3 & \sol{PC-D} & \sol{Même IP que PC-A (doublon).} & \sol{Conflit d'adresse : PC-A et PC-D perturbés.} & \sol{Adresse unique, ex. .60.40.} \\ \hline
    \hl 4 & \sol{R1 G0/1} & \sol{Interface désactivée (Off).} & \sol{Réseau du serveur injoignable pour tous.} & \sol{Port Status : On.} \\ \hline}
  \def\Cbareme{PC-A et SRV-C sont correctement configurés ; accepter PC-A cité à la place de PC-D pour le doublon.}
\else
  \def\entreprise{La clinique vétérinaire \textbf{VetoPlus}}
  \def\Rnom{R-VetoPlus}
  \def\resUn{Soins}\def\pUn{SOI}\def\netUn{192.168.70}
  \def\resDeux{Secrétariat}\def\pDeux{SEC}\def\netDeux{172.20}
  \def\srvNom{SRV-Dossiers}\def\netSrv{192.168.250}\def\srvH{20}
  \def\ssid{VetoPlus-WiFi}
  \def\AIrows{%
    \hl\ip{192.168.15.130} & \sol{C} & \sol{255.255.255.0} & \sol{192.168.15.0} & \sol{192.168.15.255} & \sol{254} \\ \hline
    \hl\ip{172.31.200.7}   & \sol{B} & \sol{255.255.0.0}   & \sol{172.31.0.0}   & \sol{172.31.255.255} & \sol{65\,534} \\ \hline
    \hl\ip{10.45.6.9}      & \sol{A} & \sol{255.0.0.0}     & \sol{10.0.0.0}     & \sol{10.255.255.255} & \sol{16\,777\,214} \\ \hline}
  \def\apipa{169.254.88.3}
  \def\Cnet{192.168.80}
  \def\Crows{%
    PC-A & \ip{192.168.8.11}   & \ip{255.255.255.0} & \ip{192.168.80.254} & — \\
    PC-B & \ip{192.168.80.12}  & \ip{255.255.255.0} & \ip{192.168.80.254} & — \\
    PC-C & \ip{192.168.80.13}  & \ip{255.255.255.0} & \ip{192.168.80.100} & — \\
    PC-D & \ip{192.168.80.12}  & \ip{255.255.255.0} & \ip{192.168.80.254} & — \\
    SRV-C & \ip{10.0.0.8}      & \ip{255.0.0.0}     & \ip{10.0.0.254}     & — \\
    R1 — G0/0 & \ip{192.168.80.254} & \ip{255.255.255.0} & — & Off \\
    R1 — G0/1 & \ip{10.0.0.254}     & \ip{255.0.0.0}     & — & On \\}
  \def\Csol{%
    \hl 1 & \sol{PC-A} & \sol{IP hors du réseau 192.168.80.0/24.} & \sol{Ne joint ni les postes ni le routeur.} & \sol{Ex. 192.168.80.11.} \\ \hline
    \hl 2 & \sol{PC-C} & \sol{Passerelle .80.100 : pas l'adresse du routeur.} & \sol{Joint son réseau, mais aucun autre (serveur).} & \sol{Passerelle 192.168.80.254.} \\ \hline
    \hl 3 & \sol{PC-D} & \sol{Même IP que PC-B (doublon).} & \sol{Conflit d'adresse : PC-B et PC-D perturbés.} & \sol{Adresse unique, ex. .80.14.} \\ \hline
    \hl 4 & \sol{R1 G0/0} & \sol{Interface côté LAN désactivée (Off).} & \sol{Aucun poste ne joint le routeur ni le serveur.} & \sol{Port Status : On.} \\ \hline}
  \def\Cbareme{PC-B et SRV-C sont correctement configurés ; accepter PC-B cité à la place de PC-D pour le doublon.}
\fi

\pagestyle{fancy}
\fancyhf{}
\renewcommand{\headrulewidth}{0pt}
\fancyfoot[L]{\footnotesize\itshape\color{grisfoot}Évaluation — Réseaux locaux — Sujet \variante\ifcorrige\ — CORRIGÉ\fi}
\fancyfoot[R]{\footnotesize\color{grisfoot}\thepage/3}

\setlist{itemsep=0pt,topsep=2pt,parsep=0pt}
\setlength{\parindent}{0pt}
\setlength{\parskip}{2pt}

\begin{document}

\begin{center}
  {\Large\bfseries Évaluation pratique — Réseaux locaux \quad
   \fcolorbox{black}{white}{\ SUJET \variante\ }\par}
  \vspace{1pt}
  {\bfseries Adressage IP, commutation et routage — BTS CIEL 1re année — durée : 2 heures\par}
  \ifcorrige\vspace{2pt}{\bfseries\color{corr}CORRIGÉ ET BARÈME — document enseignant}\par\fi
\end{center}
\vspace{-2pt}
{\renewcommand{\arraystretch}{1.35}
\begin{tabularx}{\linewidth}{|X|X|p{2.2cm}|p{2.6cm}|}
  \hline
  \textbf{Nom :} & \textbf{Prénom :} & \textbf{Poste n°} & \textbf{Note :} \hfill /20 \\ \hline
\end{tabularx}}

\textbf{Consignes :} travail \textbf{individuel}, documents et Internet \textbf{non autorisés}. Ce sujet est le
\textbf{document-réponse}. Fichier Packet Tracer : \ip{Nom\_Prenom\_EvalLAN\_\variante.pkt}, à enregistrer
régulièrement dans le dossier de rendu. Les « visas » sont apposés par le professeur, que vous appelez à chaque test réussi.
Durées conseillées : A 30 min — B 1 h 10 — C 20 min.

%==========================================================================
\section{Partie A — Questions de cours et calculs \hfill /6}

\question{A1}{2 pts}{Complétez le tableau en considérant le \textbf{masque par défaut} de la classe de chaque adresse.}
{\small\renewcommand{\arraystretch}{1.15}
\begin{tabularx}{\linewidth}{|l|c|C|C|C|c|}
  \hline
  \textbf{Adresse IP} & \textbf{Classe} & \textbf{Masque par défaut} & \textbf{Adresse réseau} &
  \textbf{Adresse de diffusion} & \textbf{Nb d'hôtes} \\ \hline
  \AIrows
\end{tabularx}}
\bareme{0,5 pt par ligne juste + 0,5 pt si les trois nombres d'hôtes sont justes ($2^{n}-2$, $n$ = nombre de bits d'hôte).}

\question{A2}{1 pt}{Indiquez le type de câble (\textbf{droit} ou \textbf{croisé}) pour chaque liaison.}
{\small\renewcommand{\arraystretch}{1.15}
\begin{tabularx}{\linewidth}{|X|c||X|c|}
  \hline
  \hl PC $\leftrightarrow$ commutateur & \makebox[1.6cm]{\sol{droit}} & Commutateur $\leftrightarrow$ routeur & \makebox[1.6cm]{\sol{droit}} \\ \hline
  \hl PC $\leftrightarrow$ PC (réseau P2P) & \sol{croisé} & Concentrateur $\leftrightarrow$ concentrateur & \sol{croisé} \\ \hline
\end{tabularx}}
\bareme{0,25 pt par case. Équipements de types différents : câble droit ; de même type : câble croisé.}

\question{A3}{1 pt}{Comment un \textbf{concentrateur (hub)} et un \textbf{commutateur (switch)} traitent-ils une trame
reçue ? Lequel limite les collisions, et pourquoi ?}
\zone{2.9cm}{Le \textbf{hub} (couche 1) répète la trame sur \textbf{tous ses autres ports} : un seul domaine de collision
partagé. Le \textbf{switch} (couche 2) lit l'adresse MAC de destination et, grâce à sa \textbf{table MAC}, ne l'envoie
\textbf{que sur le port du destinataire} : un domaine de collision par port, c'est lui qui limite les collisions.
\emph{0,5 pt hub, 0,5 pt switch.}}

\question{A4}{1 pt}{Qu'est-ce que la \textbf{passerelle par défaut} d'un poste ? Quelle condition son adresse doit-elle respecter ?}
\zone{2.4cm}{Adresse IP de l'interface du \textbf{routeur} à laquelle le poste envoie les paquets destinés à un
\textbf{autre réseau}. Elle doit \textbf{appartenir au même réseau} que le poste. \emph{0,5 + 0,5 pt.}}

\question{A5}{0,5 pt}{Un portable configuré en DHCP affiche \ip{\apipa} / \ip{255.255.0.0}. Que s'est-il passé ?}
\zone{1.5cm}{Aucun \textbf{serveur DHCP} n'a répondu : le poste s'est attribué une adresse \textbf{APIPA} (169.254.0.0/16).}

\question{A6}{0,5 pt}{Quel protocole utilise \texttt{ping} ? Nommez les deux messages échangés et la couche OSI des adresses IP.}
\zone{1.5cm}{\textbf{ICMP} ; \textbf{Echo Request} / \textbf{Echo Reply} ; adresses IP en \textbf{couche 3 (Réseau)}.}

%==========================================================================
\newpage
\section{Partie B — Réalisation sous Packet Tracer \hfill /10}

\begin{cdc}
\entreprise{} vous confie son réseau : \textbf{trois réseaux} reliés par un \textbf{routeur unique Cisco 2911}
nommé \ip{\Rnom} (interfaces G0/0, G0/1, G0/2).
\begin{itemize}[leftmargin=12pt]
  \item \textbf{Réseau \resUn} : \textbf{4 postes} PC-\pUn1 à PC-\pUn4, commutateur \ip{SW-\pUn}, adressage statique,
        réseau \ip{\netUn.0/24}.
  \item \textbf{Réseau \resDeux} : \textbf{2 postes} PC-\pDeux1, PC-\pDeux2 et \textbf{1 imprimante} IMP-\pDeux,
        commutateur \ip{SW-\pDeux}, adressage statique, réseau de classe B \ip{\netDeux.0.0/16}.
  \item \textbf{Réseau Serveur} : serveur \ip{\srvNom} d'adresse \ip{\netSrv.\srvH/24}, relié \textbf{directement} au routeur.
  \item La passerelle de chaque réseau est sa \textbf{dernière adresse utilisable}. Tous les équipements communiquent
        entre eux et accèdent à la page web du serveur. Nommage imposé (« Display Name »).
\end{itemize}
\end{cdc}

\question{B1}{2 pts}{Complétez le plan d'adressage \textbf{avant} la réalisation.}
{\footnotesize\setlength{\tabcolsep}{3pt}\renewcommand{\arraystretch}{1.15}
\begin{tabularx}{\linewidth}{|l|C|C|C|C|C|C|}
  \hline
  \textbf{Réseau} & \textbf{Adresse réseau} & \textbf{Masque} & \textbf{1re @ utilisable} &
  \textbf{Dernière @ utilisable} & \textbf{Diffusion} & \textbf{Passerelle} \\ \hline
  \hl \resUn & \sol{\netUn.0} & \sol{255.255.255.0} & \sol{\netUn.1} & \sol{\netUn.254} & \sol{\netUn.255} & \sol{\netUn.254} \\ \hline
  \hl \resDeux & \sol{\netDeux.0.0} & \sol{255.255.0.0} & \sol{\netDeux.0.1} & \sol{\netDeux.255.254} & \sol{\netDeux.255.255} & \sol{\netDeux.255.254} \\ \hline
  \hl Serveur & \sol{\netSrv.0} & \sol{255.255.255.0} & \sol{\netSrv.1} & \sol{\netSrv.254} & \sol{\netSrv.255} & \sol{\netSrv.254} \\ \hline
\end{tabularx}}
\bareme{0,5 pt par ligne juste + 0,5 pt si la convention de passerelle est respectée partout.}

\question{B2}{2 pts}{Construisez la topologie (équipements, câblage, nommage) et relevez l'adressage configuré.}
{\footnotesize\setlength{\tabcolsep}{3pt}\renewcommand{\arraystretch}{1.1}
\begin{tabularx}{\linewidth}{|l|C|C|C||l|C|C|}
  \hline
  \textbf{Hôte} & \textbf{Adresse IP} & \textbf{Masque} & \textbf{Passerelle} & \textbf{Interface} & \textbf{Adresse IP} & \textbf{Masque} \\ \hline
  \hl PC-\pUn1 & \sol{\netUn.1} & \sol{255.255.255.0} & \sol{\netUn.254} & \Rnom{} G0/0 & \sol{\netUn.254} & \sol{255.255.255.0} \\ \hline
  \hl PC-\pDeux1 & \sol{\netDeux.0.1} & \sol{255.255.0.0} & \sol{\netDeux.255.254} & \Rnom{} G0/1 & \sol{\netDeux.255.254} & \sol{255.255.0.0} \\ \hline
  \hl IMP-\pDeux & \sol{\netDeux.0.10} & \sol{255.255.0.0} & \sol{\netDeux.255.254} & \Rnom{} G0/2 & \sol{\netSrv.254} & \sol{255.255.255.0} \\ \hline
  \hl \srvNom & \ip{\netSrv.\srvH} & \sol{255.255.255.0} & \sol{\netSrv.254} & \multicolumn{3}{c|}{\cellcolor{black!8}} \\ \hline
\end{tabularx}}
\bareme{Adresses d'hôtes libres (uniques, dans le bon réseau) ; affectation G0/x libre. Sur le .pkt : équipements 0,5 ;
câblage 1 (câbles droits vers les switchs et switch–routeur, \textbf{croisé} serveur–routeur, liens verts) ; nommage 0,5.}

\question{B3}{3 pts}{Configurez tous les équipements : IP, masques, passerelles des hôtes ; interfaces du routeur
(« Port Status : On »).}
\bareme{Sur le .pkt : interfaces routeur 1 pt ; IP/masques des hôtes 1 pt ; passerelles (imprimante et serveur compris)
1 pt. $-0,25$ par erreur, sans descendre sous 0 par item.}

\question{B4}{3 pts}{Réalisez les tests, notez la commande et le résultat, puis faites viser chaque test réussi.}
{\small\renewcommand{\arraystretch}{1.1}
\begin{tabularx}{\linewidth}{|c|l|l|C|c|c|}
  \hline
  \textbf{N°} & \textbf{Source} & \textbf{Destination} & \textbf{Commande / action} & \textbf{Résultat} & \textbf{Visa prof.} \\ \hline
  \hl 1 & PC-\pUn1 & PC-\pUn4 & \sol{ping \netUn.4} & \sol{réussi} & \hspace{1.4cm} \\ \hline
  \hl 2 & PC-\pDeux1 & IMP-\pDeux & \sol{ping \netDeux.0.10} & \sol{réussi} & \\ \hline
  \hl 3 & PC-\pUn2 & PC-\pDeux2 & \sol{ping \netDeux.0.2} & \sol{réussi} & \\ \hline
  \hl 4 & PC-\pDeux2 & \srvNom & \sol{ping \netSrv.\srvH} & \sol{réussi} & \\ \hline
  \hl 5 & PC-\pUn3 & page web du serveur & \sol{navigateur : {\NoAutoSpacing http://\netSrv.\srvH}} & \sol{page affichée} & \\ \hline
\end{tabularx}}
\bareme{0,5 pt par test visé + 0,5 pt pour B5. Le premier ping inter-réseaux peut échouer (résolution ARP) ; les suivants doivent réussir.}

\question{B5}{incl. B4}{Test n°3 : par quels équipements passe le paquet ? Pourquoi PC-\pUn2 l'envoie-t-il au routeur ?}
\zone{2.3cm}{PC-\pUn2 $\to$ SW-\pUn{} $\to$ \Rnom{} $\to$ SW-\pDeux{} $\to$ PC-\pDeux2. Grâce à son masque, PC-\pUn2 constate que
la destination (\netDeux.0.0/16) n'est \textbf{pas dans son réseau} (\netUn.0/24) : il envoie le paquet à sa
\textbf{passerelle par défaut}, qui le route.}

%==========================================================================
\newpage
\section{Partie C — Diagnostic d'un réseau en panne \hfill /4}

Un technicien constate qu'\textbf{aucun poste ne parvient à joindre le serveur SRV-C} et que certains postes
communiquent mal entre eux. Les câbles sont du bon type.

\begin{center}
\begin{minipage}{0.37\linewidth}\centering
\begin{tikzpicture}[font=\footnotesize,
  eq/.style={draw,fill=qbg,rounded corners=2pt,minimum width=9mm,minimum height=5mm,inner sep=2pt},
  every path/.style={thick}]
  \node[eq] (SW) {SW-1};
  \node[eq,right=9mm of SW] (R) {R1};
  \node[eq,right=7mm of R] (SRV) {SRV-C};
  \node[eq,above left=5mm and -3mm of SW] (A) {PC-A};
  \node[eq,above right=5mm and -3mm of SW] (B) {PC-B};
  \node[eq,below left=5mm and -3mm of SW] (C) {PC-C};
  \node[eq,below right=5mm and -3mm of SW] (D) {PC-D};
  \foreach \n in {A,B,C,D} \draw (\n) -- (SW);
  \draw (SW) -- node[above,font=\tiny]{G0/0} (R);
  \draw (R) -- node[above,font=\tiny]{G0/1} (SRV);
\end{tikzpicture}
\end{minipage}\hfill
\begin{minipage}{0.61\linewidth}\centering
{\small\renewcommand{\arraystretch}{1.05}\setlength{\tabcolsep}{4pt}
\begin{tabular}{lllll}
  \toprule
  \textbf{Équipement} & \textbf{Adresse IP} & \textbf{Masque} & \textbf{Passerelle} & \textbf{Port} \\
  \midrule
  \Crows
  \bottomrule
\end{tabular}}
\end{minipage}
\end{center}

\question{C1}{4 pts}{Relevez les \textbf{quatre anomalies}. Pour chacune : conséquence et correction.}
{\small\renewcommand{\arraystretch}{1.1}
\ifcorrige\def\hc{}\else\def\hc{\rule[-10mm]{0pt}{14mm}}\fi
\begin{tabularx}{\linewidth}{|c|>{\raggedright\arraybackslash}p{2cm}|L|L|L|}
  \hline
  \textbf{N°} & \textbf{Équipement} & \textbf{Anomalie} & \textbf{Conséquence} & \textbf{Correction} \\ \hline
  \ifcorrige\Csol\else
  \hc 1 & & & & \\ \hline
  \hc 2 & & & & \\ \hline
  \hc 3 & & & & \\ \hline
  \hc 4 & & & & \\ \hline
  \fi
\end{tabularx}}
\bareme{1 pt par anomalie (0,25 équipement, anomalie, conséquence, correction ; ordre indifférent). \Cbareme}

%==========================================================================
\section{Bonus — Extension sans fil \hfill +1}

\begin{cdc}Ajoutez au réseau \resUn{} un \textbf{AccessPoint-PT} relié à SW-\pUn{} (SSID \ip{\ssid}, \textbf{WPA2-PSK})
et un portable \ip{Laptop-\pUn} équipé d'un module \texttt{WPC300N}, en adressage \textbf{statique}.
Vérifiez qu'il joint \srvNom.\end{cdc}
{\small\renewcommand{\arraystretch}{1.1}
\begin{tabularx}{\linewidth}{|l|C|C|C|c|}
  \hline
  \textbf{Équipement} & \textbf{Adresse IP} & \textbf{Masque} & \textbf{Passerelle} & \textbf{Visa prof.} \\ \hline
  \hl Laptop-\pUn & \sol{\netUn.50} & \sol{255.255.255.0} & \sol{\netUn.254} & \hspace{1.4cm} \\ \hline
\end{tabularx}}
\bareme{Le point d'accès fait office de pont : le portable est dans \netUn.0/24. Bonus si associé en WPA2-PSK et ping
vers \netSrv.\srvH{} réussi. Note plafonnée à 20.}

\end{document}
"
%   pdflatex -jobname=Eval_Reseau_CorrigeA "\def\variante{A}\def\modecorrige{1}% Évaluation pratique — Réseaux locaux (BTS CIEL 1re année) — sujets A et B
% Compilation (4 PDF) :
%   pdflatex -jobname=Eval_Reseau_SujetA   "\def\variante{A}\def\modecorrige{0}\input{Eval_TP_Reseau_LAN.tex}"
%   pdflatex -jobname=Eval_Reseau_SujetB   "\def\variante{B}\def\modecorrige{0}\input{Eval_TP_Reseau_LAN.tex}"
%   pdflatex -jobname=Eval_Reseau_CorrigeA "\def\variante{A}\def\modecorrige{1}\input{Eval_TP_Reseau_LAN.tex}"
%   pdflatex -jobname=Eval_Reseau_CorrigeB "\def\variante{B}\def\modecorrige{1}\input{Eval_TP_Reseau_LAN.tex}"
\documentclass[10pt,a4paper]{article}

\usepackage[T1]{fontenc}
\usepackage[utf8]{inputenc}
\usepackage{lmodern}
\usepackage[french]{babel}
\usepackage[a4paper,hmargin=1.5cm,top=1.3cm,bottom=1.6cm,footskip=0.7cm]{geometry}
\usepackage[table]{xcolor}
\usepackage{booktabs,tabularx,array}
\usepackage{enumitem}
\usepackage{fancyhdr}
\usepackage[most]{tcolorbox}
\usepackage{tikz}
\usetikzlibrary{positioning}
\usepackage{titlesec}

\providecommand{\modecorrige}{0}
\providecommand{\variante}{A}
\newif\ifcorrige
\ifnum\modecorrige=1 \corrigetrue \else \corrigefalse \fi

\definecolor{ipred}{HTML}{C0392B}
\definecolor{cdcbar}{HTML}{A67C3D}
\definecolor{cdcbg}{HTML}{F4F1EA}
\definecolor{qbar}{HTML}{4A6B5D}
\definecolor{qbg}{HTML}{EDF1EF}
\definecolor{corr}{HTML}{1F5FA8}
\definecolor{corrbg}{HTML}{E8F0FA}
\definecolor{grisfoot}{HTML}{828282}

\newcommand{\ip}[1]{\textcolor{ipred}{\texttt{#1}}}
\newcommand{\pts}[1]{\hfill\mbox{\textbf{(#1)}}}
\newcommand{\sol}[1]{\ifcorrige\textcolor{corr}{#1}\fi}
\newcommand{\bareme}[1]{\ifcorrige\par{\footnotesize\color{corr}#1\par}\fi}
\newcommand{\hl}{\ifcorrige\rule[-1.6mm]{0pt}{5.4mm}\else\rule[-2.2mm]{0pt}{6.6mm}\fi}
\newcolumntype{C}{>{\centering\arraybackslash}X}
\newcolumntype{L}{>{\raggedright\arraybackslash}X}

\titleformat{\section}{\large\bfseries}{}{0pt}{}
\titlespacing*{\section}{0pt}{6pt}{3pt}

\tcbset{qstyle/.style={enhanced,boxrule=0pt,leftrule=2.5pt,arc=0pt,outer arc=0pt,
  left=5pt,right=5pt,top=2pt,bottom=2pt,before skip=4pt,after skip=3pt}}
\newtcolorbox{cdc}{qstyle,colback=cdcbg,colframe=cdcbar}
\newtcolorbox{quest}{qstyle,colback=qbg,colframe=qbar}

% Zone de réponse : vierge dans le sujet, remplie dans le corrigé
\newcommand{\zone}[2]{%
  \ifcorrige
    \begin{tcolorbox}[colback=corrbg,colframe=corr,boxrule=0.5pt,arc=1pt,left=4pt,right=4pt,top=1pt,bottom=1pt,
      before skip=2pt,after skip=2pt,fontupper=\footnotesize\color{corr}]#2\end{tcolorbox}%
  \else
    \begin{tcolorbox}[colback=white,colframe=black!45,boxrule=0.4pt,arc=1pt,height=#1,
      before skip=2pt,after skip=2pt]\end{tcolorbox}%
  \fi}

\newcommand{\question}[3]{\begin{quest}\textbf{#1.} #3\pts{#2}\end{quest}}

%==========================================================================
% Données propres à chaque variante
%==========================================================================
\if\variante A
  \def\entreprise{Le cabinet d'architectes \textbf{ArchiNova}}
  \def\Rnom{R-ArchiNova}
  \def\resUn{Bureau d'études}\def\pUn{BE}\def\netUn{192.168.40}
  \def\resDeux{Accueil}\def\pDeux{ACC}\def\netDeux{172.16}
  \def\srvNom{SRV-Intranet}\def\netSrv{192.168.200}\def\srvH{10}
  \def\ssid{ArchiNova-WiFi}
  \def\AIrows{%
    \hl\ip{192.168.40.25} & \sol{C} & \sol{255.255.255.0} & \sol{192.168.40.0} & \sol{192.168.40.255} & \sol{254} \\ \hline
    \hl\ip{172.16.5.10}   & \sol{B} & \sol{255.255.0.0}   & \sol{172.16.0.0}   & \sol{172.16.255.255} & \sol{65\,534} \\ \hline
    \hl\ip{10.3.2.1}      & \sol{A} & \sol{255.0.0.0}     & \sol{10.0.0.0}     & \sol{10.255.255.255} & \sol{16\,777\,214} \\ \hline}
  \def\apipa{169.254.12.7}
  \def\Cnet{192.168.60}
  \def\Crows{%
    PC-A & \ip{192.168.60.10} & \ip{255.255.255.0} & \ip{192.168.60.254} & — \\
    PC-B & \ip{192.168.60.20} & \ip{255.255.255.0} & \ip{192.168.60.1}   & — \\
    PC-C & \ip{192.168.61.30} & \ip{255.255.255.0} & \ip{192.168.60.254} & — \\
    PC-D & \ip{192.168.60.10} & \ip{255.255.255.0} & \ip{192.168.60.254} & — \\
    SRV-C & \ip{10.0.0.5}     & \ip{255.0.0.0}     & \ip{10.0.0.254}     & — \\
    R1 — G0/0 & \ip{192.168.60.254} & \ip{255.255.255.0} & — & On \\
    R1 — G0/1 & \ip{10.0.0.254}     & \ip{255.0.0.0}     & — & Off \\}
  \def\Csol{%
    \hl 1 & \sol{PC-B} & \sol{Passerelle .60.1 : pas l'adresse du routeur.} & \sol{Joint son réseau, mais aucun autre (serveur).} & \sol{Passerelle 192.168.60.254.} \\ \hline
    \hl 2 & \sol{PC-C} & \sol{IP hors du réseau 192.168.60.0/24.} & \sol{Ne joint ni les postes ni le routeur.} & \sol{Ex. 192.168.60.30.} \\ \hline
    \hl 3 & \sol{PC-D} & \sol{Même IP que PC-A (doublon).} & \sol{Conflit d'adresse : PC-A et PC-D perturbés.} & \sol{Adresse unique, ex. .60.40.} \\ \hline
    \hl 4 & \sol{R1 G0/1} & \sol{Interface désactivée (Off).} & \sol{Réseau du serveur injoignable pour tous.} & \sol{Port Status : On.} \\ \hline}
  \def\Cbareme{PC-A et SRV-C sont correctement configurés ; accepter PC-A cité à la place de PC-D pour le doublon.}
\else
  \def\entreprise{La clinique vétérinaire \textbf{VetoPlus}}
  \def\Rnom{R-VetoPlus}
  \def\resUn{Soins}\def\pUn{SOI}\def\netUn{192.168.70}
  \def\resDeux{Secrétariat}\def\pDeux{SEC}\def\netDeux{172.20}
  \def\srvNom{SRV-Dossiers}\def\netSrv{192.168.250}\def\srvH{20}
  \def\ssid{VetoPlus-WiFi}
  \def\AIrows{%
    \hl\ip{192.168.15.130} & \sol{C} & \sol{255.255.255.0} & \sol{192.168.15.0} & \sol{192.168.15.255} & \sol{254} \\ \hline
    \hl\ip{172.31.200.7}   & \sol{B} & \sol{255.255.0.0}   & \sol{172.31.0.0}   & \sol{172.31.255.255} & \sol{65\,534} \\ \hline
    \hl\ip{10.45.6.9}      & \sol{A} & \sol{255.0.0.0}     & \sol{10.0.0.0}     & \sol{10.255.255.255} & \sol{16\,777\,214} \\ \hline}
  \def\apipa{169.254.88.3}
  \def\Cnet{192.168.80}
  \def\Crows{%
    PC-A & \ip{192.168.8.11}   & \ip{255.255.255.0} & \ip{192.168.80.254} & — \\
    PC-B & \ip{192.168.80.12}  & \ip{255.255.255.0} & \ip{192.168.80.254} & — \\
    PC-C & \ip{192.168.80.13}  & \ip{255.255.255.0} & \ip{192.168.80.100} & — \\
    PC-D & \ip{192.168.80.12}  & \ip{255.255.255.0} & \ip{192.168.80.254} & — \\
    SRV-C & \ip{10.0.0.8}      & \ip{255.0.0.0}     & \ip{10.0.0.254}     & — \\
    R1 — G0/0 & \ip{192.168.80.254} & \ip{255.255.255.0} & — & Off \\
    R1 — G0/1 & \ip{10.0.0.254}     & \ip{255.0.0.0}     & — & On \\}
  \def\Csol{%
    \hl 1 & \sol{PC-A} & \sol{IP hors du réseau 192.168.80.0/24.} & \sol{Ne joint ni les postes ni le routeur.} & \sol{Ex. 192.168.80.11.} \\ \hline
    \hl 2 & \sol{PC-C} & \sol{Passerelle .80.100 : pas l'adresse du routeur.} & \sol{Joint son réseau, mais aucun autre (serveur).} & \sol{Passerelle 192.168.80.254.} \\ \hline
    \hl 3 & \sol{PC-D} & \sol{Même IP que PC-B (doublon).} & \sol{Conflit d'adresse : PC-B et PC-D perturbés.} & \sol{Adresse unique, ex. .80.14.} \\ \hline
    \hl 4 & \sol{R1 G0/0} & \sol{Interface côté LAN désactivée (Off).} & \sol{Aucun poste ne joint le routeur ni le serveur.} & \sol{Port Status : On.} \\ \hline}
  \def\Cbareme{PC-B et SRV-C sont correctement configurés ; accepter PC-B cité à la place de PC-D pour le doublon.}
\fi

\pagestyle{fancy}
\fancyhf{}
\renewcommand{\headrulewidth}{0pt}
\fancyfoot[L]{\footnotesize\itshape\color{grisfoot}Évaluation — Réseaux locaux — Sujet \variante\ifcorrige\ — CORRIGÉ\fi}
\fancyfoot[R]{\footnotesize\color{grisfoot}\thepage/3}

\setlist{itemsep=0pt,topsep=2pt,parsep=0pt}
\setlength{\parindent}{0pt}
\setlength{\parskip}{2pt}

\begin{document}

\begin{center}
  {\Large\bfseries Évaluation pratique — Réseaux locaux \quad
   \fcolorbox{black}{white}{\ SUJET \variante\ }\par}
  \vspace{1pt}
  {\bfseries Adressage IP, commutation et routage — BTS CIEL 1re année — durée : 2 heures\par}
  \ifcorrige\vspace{2pt}{\bfseries\color{corr}CORRIGÉ ET BARÈME — document enseignant}\par\fi
\end{center}
\vspace{-2pt}
{\renewcommand{\arraystretch}{1.35}
\begin{tabularx}{\linewidth}{|X|X|p{2.2cm}|p{2.6cm}|}
  \hline
  \textbf{Nom :} & \textbf{Prénom :} & \textbf{Poste n°} & \textbf{Note :} \hfill /20 \\ \hline
\end{tabularx}}

\textbf{Consignes :} travail \textbf{individuel}, documents et Internet \textbf{non autorisés}. Ce sujet est le
\textbf{document-réponse}. Fichier Packet Tracer : \ip{Nom\_Prenom\_EvalLAN\_\variante.pkt}, à enregistrer
régulièrement dans le dossier de rendu. Les « visas » sont apposés par le professeur, que vous appelez à chaque test réussi.
Durées conseillées : A 30 min — B 1 h 10 — C 20 min.

%==========================================================================
\section{Partie A — Questions de cours et calculs \hfill /6}

\question{A1}{2 pts}{Complétez le tableau en considérant le \textbf{masque par défaut} de la classe de chaque adresse.}
{\small\renewcommand{\arraystretch}{1.15}
\begin{tabularx}{\linewidth}{|l|c|C|C|C|c|}
  \hline
  \textbf{Adresse IP} & \textbf{Classe} & \textbf{Masque par défaut} & \textbf{Adresse réseau} &
  \textbf{Adresse de diffusion} & \textbf{Nb d'hôtes} \\ \hline
  \AIrows
\end{tabularx}}
\bareme{0,5 pt par ligne juste + 0,5 pt si les trois nombres d'hôtes sont justes ($2^{n}-2$, $n$ = nombre de bits d'hôte).}

\question{A2}{1 pt}{Indiquez le type de câble (\textbf{droit} ou \textbf{croisé}) pour chaque liaison.}
{\small\renewcommand{\arraystretch}{1.15}
\begin{tabularx}{\linewidth}{|X|c||X|c|}
  \hline
  \hl PC $\leftrightarrow$ commutateur & \makebox[1.6cm]{\sol{droit}} & Commutateur $\leftrightarrow$ routeur & \makebox[1.6cm]{\sol{droit}} \\ \hline
  \hl PC $\leftrightarrow$ PC (réseau P2P) & \sol{croisé} & Concentrateur $\leftrightarrow$ concentrateur & \sol{croisé} \\ \hline
\end{tabularx}}
\bareme{0,25 pt par case. Équipements de types différents : câble droit ; de même type : câble croisé.}

\question{A3}{1 pt}{Comment un \textbf{concentrateur (hub)} et un \textbf{commutateur (switch)} traitent-ils une trame
reçue ? Lequel limite les collisions, et pourquoi ?}
\zone{2.9cm}{Le \textbf{hub} (couche 1) répète la trame sur \textbf{tous ses autres ports} : un seul domaine de collision
partagé. Le \textbf{switch} (couche 2) lit l'adresse MAC de destination et, grâce à sa \textbf{table MAC}, ne l'envoie
\textbf{que sur le port du destinataire} : un domaine de collision par port, c'est lui qui limite les collisions.
\emph{0,5 pt hub, 0,5 pt switch.}}

\question{A4}{1 pt}{Qu'est-ce que la \textbf{passerelle par défaut} d'un poste ? Quelle condition son adresse doit-elle respecter ?}
\zone{2.4cm}{Adresse IP de l'interface du \textbf{routeur} à laquelle le poste envoie les paquets destinés à un
\textbf{autre réseau}. Elle doit \textbf{appartenir au même réseau} que le poste. \emph{0,5 + 0,5 pt.}}

\question{A5}{0,5 pt}{Un portable configuré en DHCP affiche \ip{\apipa} / \ip{255.255.0.0}. Que s'est-il passé ?}
\zone{1.5cm}{Aucun \textbf{serveur DHCP} n'a répondu : le poste s'est attribué une adresse \textbf{APIPA} (169.254.0.0/16).}

\question{A6}{0,5 pt}{Quel protocole utilise \texttt{ping} ? Nommez les deux messages échangés et la couche OSI des adresses IP.}
\zone{1.5cm}{\textbf{ICMP} ; \textbf{Echo Request} / \textbf{Echo Reply} ; adresses IP en \textbf{couche 3 (Réseau)}.}

%==========================================================================
\newpage
\section{Partie B — Réalisation sous Packet Tracer \hfill /10}

\begin{cdc}
\entreprise{} vous confie son réseau : \textbf{trois réseaux} reliés par un \textbf{routeur unique Cisco 2911}
nommé \ip{\Rnom} (interfaces G0/0, G0/1, G0/2).
\begin{itemize}[leftmargin=12pt]
  \item \textbf{Réseau \resUn} : \textbf{4 postes} PC-\pUn1 à PC-\pUn4, commutateur \ip{SW-\pUn}, adressage statique,
        réseau \ip{\netUn.0/24}.
  \item \textbf{Réseau \resDeux} : \textbf{2 postes} PC-\pDeux1, PC-\pDeux2 et \textbf{1 imprimante} IMP-\pDeux,
        commutateur \ip{SW-\pDeux}, adressage statique, réseau de classe B \ip{\netDeux.0.0/16}.
  \item \textbf{Réseau Serveur} : serveur \ip{\srvNom} d'adresse \ip{\netSrv.\srvH/24}, relié \textbf{directement} au routeur.
  \item La passerelle de chaque réseau est sa \textbf{dernière adresse utilisable}. Tous les équipements communiquent
        entre eux et accèdent à la page web du serveur. Nommage imposé (« Display Name »).
\end{itemize}
\end{cdc}

\question{B1}{2 pts}{Complétez le plan d'adressage \textbf{avant} la réalisation.}
{\footnotesize\setlength{\tabcolsep}{3pt}\renewcommand{\arraystretch}{1.15}
\begin{tabularx}{\linewidth}{|l|C|C|C|C|C|C|}
  \hline
  \textbf{Réseau} & \textbf{Adresse réseau} & \textbf{Masque} & \textbf{1re @ utilisable} &
  \textbf{Dernière @ utilisable} & \textbf{Diffusion} & \textbf{Passerelle} \\ \hline
  \hl \resUn & \sol{\netUn.0} & \sol{255.255.255.0} & \sol{\netUn.1} & \sol{\netUn.254} & \sol{\netUn.255} & \sol{\netUn.254} \\ \hline
  \hl \resDeux & \sol{\netDeux.0.0} & \sol{255.255.0.0} & \sol{\netDeux.0.1} & \sol{\netDeux.255.254} & \sol{\netDeux.255.255} & \sol{\netDeux.255.254} \\ \hline
  \hl Serveur & \sol{\netSrv.0} & \sol{255.255.255.0} & \sol{\netSrv.1} & \sol{\netSrv.254} & \sol{\netSrv.255} & \sol{\netSrv.254} \\ \hline
\end{tabularx}}
\bareme{0,5 pt par ligne juste + 0,5 pt si la convention de passerelle est respectée partout.}

\question{B2}{2 pts}{Construisez la topologie (équipements, câblage, nommage) et relevez l'adressage configuré.}
{\footnotesize\setlength{\tabcolsep}{3pt}\renewcommand{\arraystretch}{1.1}
\begin{tabularx}{\linewidth}{|l|C|C|C||l|C|C|}
  \hline
  \textbf{Hôte} & \textbf{Adresse IP} & \textbf{Masque} & \textbf{Passerelle} & \textbf{Interface} & \textbf{Adresse IP} & \textbf{Masque} \\ \hline
  \hl PC-\pUn1 & \sol{\netUn.1} & \sol{255.255.255.0} & \sol{\netUn.254} & \Rnom{} G0/0 & \sol{\netUn.254} & \sol{255.255.255.0} \\ \hline
  \hl PC-\pDeux1 & \sol{\netDeux.0.1} & \sol{255.255.0.0} & \sol{\netDeux.255.254} & \Rnom{} G0/1 & \sol{\netDeux.255.254} & \sol{255.255.0.0} \\ \hline
  \hl IMP-\pDeux & \sol{\netDeux.0.10} & \sol{255.255.0.0} & \sol{\netDeux.255.254} & \Rnom{} G0/2 & \sol{\netSrv.254} & \sol{255.255.255.0} \\ \hline
  \hl \srvNom & \ip{\netSrv.\srvH} & \sol{255.255.255.0} & \sol{\netSrv.254} & \multicolumn{3}{c|}{\cellcolor{black!8}} \\ \hline
\end{tabularx}}
\bareme{Adresses d'hôtes libres (uniques, dans le bon réseau) ; affectation G0/x libre. Sur le .pkt : équipements 0,5 ;
câblage 1 (câbles droits vers les switchs et switch–routeur, \textbf{croisé} serveur–routeur, liens verts) ; nommage 0,5.}

\question{B3}{3 pts}{Configurez tous les équipements : IP, masques, passerelles des hôtes ; interfaces du routeur
(« Port Status : On »).}
\bareme{Sur le .pkt : interfaces routeur 1 pt ; IP/masques des hôtes 1 pt ; passerelles (imprimante et serveur compris)
1 pt. $-0,25$ par erreur, sans descendre sous 0 par item.}

\question{B4}{3 pts}{Réalisez les tests, notez la commande et le résultat, puis faites viser chaque test réussi.}
{\small\renewcommand{\arraystretch}{1.1}
\begin{tabularx}{\linewidth}{|c|l|l|C|c|c|}
  \hline
  \textbf{N°} & \textbf{Source} & \textbf{Destination} & \textbf{Commande / action} & \textbf{Résultat} & \textbf{Visa prof.} \\ \hline
  \hl 1 & PC-\pUn1 & PC-\pUn4 & \sol{ping \netUn.4} & \sol{réussi} & \hspace{1.4cm} \\ \hline
  \hl 2 & PC-\pDeux1 & IMP-\pDeux & \sol{ping \netDeux.0.10} & \sol{réussi} & \\ \hline
  \hl 3 & PC-\pUn2 & PC-\pDeux2 & \sol{ping \netDeux.0.2} & \sol{réussi} & \\ \hline
  \hl 4 & PC-\pDeux2 & \srvNom & \sol{ping \netSrv.\srvH} & \sol{réussi} & \\ \hline
  \hl 5 & PC-\pUn3 & page web du serveur & \sol{navigateur : {\NoAutoSpacing http://\netSrv.\srvH}} & \sol{page affichée} & \\ \hline
\end{tabularx}}
\bareme{0,5 pt par test visé + 0,5 pt pour B5. Le premier ping inter-réseaux peut échouer (résolution ARP) ; les suivants doivent réussir.}

\question{B5}{incl. B4}{Test n°3 : par quels équipements passe le paquet ? Pourquoi PC-\pUn2 l'envoie-t-il au routeur ?}
\zone{2.3cm}{PC-\pUn2 $\to$ SW-\pUn{} $\to$ \Rnom{} $\to$ SW-\pDeux{} $\to$ PC-\pDeux2. Grâce à son masque, PC-\pUn2 constate que
la destination (\netDeux.0.0/16) n'est \textbf{pas dans son réseau} (\netUn.0/24) : il envoie le paquet à sa
\textbf{passerelle par défaut}, qui le route.}

%==========================================================================
\newpage
\section{Partie C — Diagnostic d'un réseau en panne \hfill /4}

Un technicien constate qu'\textbf{aucun poste ne parvient à joindre le serveur SRV-C} et que certains postes
communiquent mal entre eux. Les câbles sont du bon type.

\begin{center}
\begin{minipage}{0.37\linewidth}\centering
\begin{tikzpicture}[font=\footnotesize,
  eq/.style={draw,fill=qbg,rounded corners=2pt,minimum width=9mm,minimum height=5mm,inner sep=2pt},
  every path/.style={thick}]
  \node[eq] (SW) {SW-1};
  \node[eq,right=9mm of SW] (R) {R1};
  \node[eq,right=7mm of R] (SRV) {SRV-C};
  \node[eq,above left=5mm and -3mm of SW] (A) {PC-A};
  \node[eq,above right=5mm and -3mm of SW] (B) {PC-B};
  \node[eq,below left=5mm and -3mm of SW] (C) {PC-C};
  \node[eq,below right=5mm and -3mm of SW] (D) {PC-D};
  \foreach \n in {A,B,C,D} \draw (\n) -- (SW);
  \draw (SW) -- node[above,font=\tiny]{G0/0} (R);
  \draw (R) -- node[above,font=\tiny]{G0/1} (SRV);
\end{tikzpicture}
\end{minipage}\hfill
\begin{minipage}{0.61\linewidth}\centering
{\small\renewcommand{\arraystretch}{1.05}\setlength{\tabcolsep}{4pt}
\begin{tabular}{lllll}
  \toprule
  \textbf{Équipement} & \textbf{Adresse IP} & \textbf{Masque} & \textbf{Passerelle} & \textbf{Port} \\
  \midrule
  \Crows
  \bottomrule
\end{tabular}}
\end{minipage}
\end{center}

\question{C1}{4 pts}{Relevez les \textbf{quatre anomalies}. Pour chacune : conséquence et correction.}
{\small\renewcommand{\arraystretch}{1.1}
\ifcorrige\def\hc{}\else\def\hc{\rule[-10mm]{0pt}{14mm}}\fi
\begin{tabularx}{\linewidth}{|c|>{\raggedright\arraybackslash}p{2cm}|L|L|L|}
  \hline
  \textbf{N°} & \textbf{Équipement} & \textbf{Anomalie} & \textbf{Conséquence} & \textbf{Correction} \\ \hline
  \ifcorrige\Csol\else
  \hc 1 & & & & \\ \hline
  \hc 2 & & & & \\ \hline
  \hc 3 & & & & \\ \hline
  \hc 4 & & & & \\ \hline
  \fi
\end{tabularx}}
\bareme{1 pt par anomalie (0,25 équipement, anomalie, conséquence, correction ; ordre indifférent). \Cbareme}

%==========================================================================
\section{Bonus — Extension sans fil \hfill +1}

\begin{cdc}Ajoutez au réseau \resUn{} un \textbf{AccessPoint-PT} relié à SW-\pUn{} (SSID \ip{\ssid}, \textbf{WPA2-PSK})
et un portable \ip{Laptop-\pUn} équipé d'un module \texttt{WPC300N}, en adressage \textbf{statique}.
Vérifiez qu'il joint \srvNom.\end{cdc}
{\small\renewcommand{\arraystretch}{1.1}
\begin{tabularx}{\linewidth}{|l|C|C|C|c|}
  \hline
  \textbf{Équipement} & \textbf{Adresse IP} & \textbf{Masque} & \textbf{Passerelle} & \textbf{Visa prof.} \\ \hline
  \hl Laptop-\pUn & \sol{\netUn.50} & \sol{255.255.255.0} & \sol{\netUn.254} & \hspace{1.4cm} \\ \hline
\end{tabularx}}
\bareme{Le point d'accès fait office de pont : le portable est dans \netUn.0/24. Bonus si associé en WPA2-PSK et ping
vers \netSrv.\srvH{} réussi. Note plafonnée à 20.}

\end{document}
"
%   pdflatex -jobname=Eval_Reseau_CorrigeB "\def\variante{B}\def\modecorrige{1}% Évaluation pratique — Réseaux locaux (BTS CIEL 1re année) — sujets A et B
% Compilation (4 PDF) :
%   pdflatex -jobname=Eval_Reseau_SujetA   "\def\variante{A}\def\modecorrige{0}\input{Eval_TP_Reseau_LAN.tex}"
%   pdflatex -jobname=Eval_Reseau_SujetB   "\def\variante{B}\def\modecorrige{0}\input{Eval_TP_Reseau_LAN.tex}"
%   pdflatex -jobname=Eval_Reseau_CorrigeA "\def\variante{A}\def\modecorrige{1}\input{Eval_TP_Reseau_LAN.tex}"
%   pdflatex -jobname=Eval_Reseau_CorrigeB "\def\variante{B}\def\modecorrige{1}\input{Eval_TP_Reseau_LAN.tex}"
\documentclass[10pt,a4paper]{article}

\usepackage[T1]{fontenc}
\usepackage[utf8]{inputenc}
\usepackage{lmodern}
\usepackage[french]{babel}
\usepackage[a4paper,hmargin=1.5cm,top=1.3cm,bottom=1.6cm,footskip=0.7cm]{geometry}
\usepackage[table]{xcolor}
\usepackage{booktabs,tabularx,array}
\usepackage{enumitem}
\usepackage{fancyhdr}
\usepackage[most]{tcolorbox}
\usepackage{tikz}
\usetikzlibrary{positioning}
\usepackage{titlesec}

\providecommand{\modecorrige}{0}
\providecommand{\variante}{A}
\newif\ifcorrige
\ifnum\modecorrige=1 \corrigetrue \else \corrigefalse \fi

\definecolor{ipred}{HTML}{C0392B}
\definecolor{cdcbar}{HTML}{A67C3D}
\definecolor{cdcbg}{HTML}{F4F1EA}
\definecolor{qbar}{HTML}{4A6B5D}
\definecolor{qbg}{HTML}{EDF1EF}
\definecolor{corr}{HTML}{1F5FA8}
\definecolor{corrbg}{HTML}{E8F0FA}
\definecolor{grisfoot}{HTML}{828282}

\newcommand{\ip}[1]{\textcolor{ipred}{\texttt{#1}}}
\newcommand{\pts}[1]{\hfill\mbox{\textbf{(#1)}}}
\newcommand{\sol}[1]{\ifcorrige\textcolor{corr}{#1}\fi}
\newcommand{\bareme}[1]{\ifcorrige\par{\footnotesize\color{corr}#1\par}\fi}
\newcommand{\hl}{\ifcorrige\rule[-1.6mm]{0pt}{5.4mm}\else\rule[-2.2mm]{0pt}{6.6mm}\fi}
\newcolumntype{C}{>{\centering\arraybackslash}X}
\newcolumntype{L}{>{\raggedright\arraybackslash}X}

\titleformat{\section}{\large\bfseries}{}{0pt}{}
\titlespacing*{\section}{0pt}{6pt}{3pt}

\tcbset{qstyle/.style={enhanced,boxrule=0pt,leftrule=2.5pt,arc=0pt,outer arc=0pt,
  left=5pt,right=5pt,top=2pt,bottom=2pt,before skip=4pt,after skip=3pt}}
\newtcolorbox{cdc}{qstyle,colback=cdcbg,colframe=cdcbar}
\newtcolorbox{quest}{qstyle,colback=qbg,colframe=qbar}

% Zone de réponse : vierge dans le sujet, remplie dans le corrigé
\newcommand{\zone}[2]{%
  \ifcorrige
    \begin{tcolorbox}[colback=corrbg,colframe=corr,boxrule=0.5pt,arc=1pt,left=4pt,right=4pt,top=1pt,bottom=1pt,
      before skip=2pt,after skip=2pt,fontupper=\footnotesize\color{corr}]#2\end{tcolorbox}%
  \else
    \begin{tcolorbox}[colback=white,colframe=black!45,boxrule=0.4pt,arc=1pt,height=#1,
      before skip=2pt,after skip=2pt]\end{tcolorbox}%
  \fi}

\newcommand{\question}[3]{\begin{quest}\textbf{#1.} #3\pts{#2}\end{quest}}

%==========================================================================
% Données propres à chaque variante
%==========================================================================
\if\variante A
  \def\entreprise{Le cabinet d'architectes \textbf{ArchiNova}}
  \def\Rnom{R-ArchiNova}
  \def\resUn{Bureau d'études}\def\pUn{BE}\def\netUn{192.168.40}
  \def\resDeux{Accueil}\def\pDeux{ACC}\def\netDeux{172.16}
  \def\srvNom{SRV-Intranet}\def\netSrv{192.168.200}\def\srvH{10}
  \def\ssid{ArchiNova-WiFi}
  \def\AIrows{%
    \hl\ip{192.168.40.25} & \sol{C} & \sol{255.255.255.0} & \sol{192.168.40.0} & \sol{192.168.40.255} & \sol{254} \\ \hline
    \hl\ip{172.16.5.10}   & \sol{B} & \sol{255.255.0.0}   & \sol{172.16.0.0}   & \sol{172.16.255.255} & \sol{65\,534} \\ \hline
    \hl\ip{10.3.2.1}      & \sol{A} & \sol{255.0.0.0}     & \sol{10.0.0.0}     & \sol{10.255.255.255} & \sol{16\,777\,214} \\ \hline}
  \def\apipa{169.254.12.7}
  \def\Cnet{192.168.60}
  \def\Crows{%
    PC-A & \ip{192.168.60.10} & \ip{255.255.255.0} & \ip{192.168.60.254} & — \\
    PC-B & \ip{192.168.60.20} & \ip{255.255.255.0} & \ip{192.168.60.1}   & — \\
    PC-C & \ip{192.168.61.30} & \ip{255.255.255.0} & \ip{192.168.60.254} & — \\
    PC-D & \ip{192.168.60.10} & \ip{255.255.255.0} & \ip{192.168.60.254} & — \\
    SRV-C & \ip{10.0.0.5}     & \ip{255.0.0.0}     & \ip{10.0.0.254}     & — \\
    R1 — G0/0 & \ip{192.168.60.254} & \ip{255.255.255.0} & — & On \\
    R1 — G0/1 & \ip{10.0.0.254}     & \ip{255.0.0.0}     & — & Off \\}
  \def\Csol{%
    \hl 1 & \sol{PC-B} & \sol{Passerelle .60.1 : pas l'adresse du routeur.} & \sol{Joint son réseau, mais aucun autre (serveur).} & \sol{Passerelle 192.168.60.254.} \\ \hline
    \hl 2 & \sol{PC-C} & \sol{IP hors du réseau 192.168.60.0/24.} & \sol{Ne joint ni les postes ni le routeur.} & \sol{Ex. 192.168.60.30.} \\ \hline
    \hl 3 & \sol{PC-D} & \sol{Même IP que PC-A (doublon).} & \sol{Conflit d'adresse : PC-A et PC-D perturbés.} & \sol{Adresse unique, ex. .60.40.} \\ \hline
    \hl 4 & \sol{R1 G0/1} & \sol{Interface désactivée (Off).} & \sol{Réseau du serveur injoignable pour tous.} & \sol{Port Status : On.} \\ \hline}
  \def\Cbareme{PC-A et SRV-C sont correctement configurés ; accepter PC-A cité à la place de PC-D pour le doublon.}
\else
  \def\entreprise{La clinique vétérinaire \textbf{VetoPlus}}
  \def\Rnom{R-VetoPlus}
  \def\resUn{Soins}\def\pUn{SOI}\def\netUn{192.168.70}
  \def\resDeux{Secrétariat}\def\pDeux{SEC}\def\netDeux{172.20}
  \def\srvNom{SRV-Dossiers}\def\netSrv{192.168.250}\def\srvH{20}
  \def\ssid{VetoPlus-WiFi}
  \def\AIrows{%
    \hl\ip{192.168.15.130} & \sol{C} & \sol{255.255.255.0} & \sol{192.168.15.0} & \sol{192.168.15.255} & \sol{254} \\ \hline
    \hl\ip{172.31.200.7}   & \sol{B} & \sol{255.255.0.0}   & \sol{172.31.0.0}   & \sol{172.31.255.255} & \sol{65\,534} \\ \hline
    \hl\ip{10.45.6.9}      & \sol{A} & \sol{255.0.0.0}     & \sol{10.0.0.0}     & \sol{10.255.255.255} & \sol{16\,777\,214} \\ \hline}
  \def\apipa{169.254.88.3}
  \def\Cnet{192.168.80}
  \def\Crows{%
    PC-A & \ip{192.168.8.11}   & \ip{255.255.255.0} & \ip{192.168.80.254} & — \\
    PC-B & \ip{192.168.80.12}  & \ip{255.255.255.0} & \ip{192.168.80.254} & — \\
    PC-C & \ip{192.168.80.13}  & \ip{255.255.255.0} & \ip{192.168.80.100} & — \\
    PC-D & \ip{192.168.80.12}  & \ip{255.255.255.0} & \ip{192.168.80.254} & — \\
    SRV-C & \ip{10.0.0.8}      & \ip{255.0.0.0}     & \ip{10.0.0.254}     & — \\
    R1 — G0/0 & \ip{192.168.80.254} & \ip{255.255.255.0} & — & Off \\
    R1 — G0/1 & \ip{10.0.0.254}     & \ip{255.0.0.0}     & — & On \\}
  \def\Csol{%
    \hl 1 & \sol{PC-A} & \sol{IP hors du réseau 192.168.80.0/24.} & \sol{Ne joint ni les postes ni le routeur.} & \sol{Ex. 192.168.80.11.} \\ \hline
    \hl 2 & \sol{PC-C} & \sol{Passerelle .80.100 : pas l'adresse du routeur.} & \sol{Joint son réseau, mais aucun autre (serveur).} & \sol{Passerelle 192.168.80.254.} \\ \hline
    \hl 3 & \sol{PC-D} & \sol{Même IP que PC-B (doublon).} & \sol{Conflit d'adresse : PC-B et PC-D perturbés.} & \sol{Adresse unique, ex. .80.14.} \\ \hline
    \hl 4 & \sol{R1 G0/0} & \sol{Interface côté LAN désactivée (Off).} & \sol{Aucun poste ne joint le routeur ni le serveur.} & \sol{Port Status : On.} \\ \hline}
  \def\Cbareme{PC-B et SRV-C sont correctement configurés ; accepter PC-B cité à la place de PC-D pour le doublon.}
\fi

\pagestyle{fancy}
\fancyhf{}
\renewcommand{\headrulewidth}{0pt}
\fancyfoot[L]{\footnotesize\itshape\color{grisfoot}Évaluation — Réseaux locaux — Sujet \variante\ifcorrige\ — CORRIGÉ\fi}
\fancyfoot[R]{\footnotesize\color{grisfoot}\thepage/3}

\setlist{itemsep=0pt,topsep=2pt,parsep=0pt}
\setlength{\parindent}{0pt}
\setlength{\parskip}{2pt}

\begin{document}

\begin{center}
  {\Large\bfseries Évaluation pratique — Réseaux locaux \quad
   \fcolorbox{black}{white}{\ SUJET \variante\ }\par}
  \vspace{1pt}
  {\bfseries Adressage IP, commutation et routage — BTS CIEL 1re année — durée : 2 heures\par}
  \ifcorrige\vspace{2pt}{\bfseries\color{corr}CORRIGÉ ET BARÈME — document enseignant}\par\fi
\end{center}
\vspace{-2pt}
{\renewcommand{\arraystretch}{1.35}
\begin{tabularx}{\linewidth}{|X|X|p{2.2cm}|p{2.6cm}|}
  \hline
  \textbf{Nom :} & \textbf{Prénom :} & \textbf{Poste n°} & \textbf{Note :} \hfill /20 \\ \hline
\end{tabularx}}

\textbf{Consignes :} travail \textbf{individuel}, documents et Internet \textbf{non autorisés}. Ce sujet est le
\textbf{document-réponse}. Fichier Packet Tracer : \ip{Nom\_Prenom\_EvalLAN\_\variante.pkt}, à enregistrer
régulièrement dans le dossier de rendu. Les « visas » sont apposés par le professeur, que vous appelez à chaque test réussi.
Durées conseillées : A 30 min — B 1 h 10 — C 20 min.

%==========================================================================
\section{Partie A — Questions de cours et calculs \hfill /6}

\question{A1}{2 pts}{Complétez le tableau en considérant le \textbf{masque par défaut} de la classe de chaque adresse.}
{\small\renewcommand{\arraystretch}{1.15}
\begin{tabularx}{\linewidth}{|l|c|C|C|C|c|}
  \hline
  \textbf{Adresse IP} & \textbf{Classe} & \textbf{Masque par défaut} & \textbf{Adresse réseau} &
  \textbf{Adresse de diffusion} & \textbf{Nb d'hôtes} \\ \hline
  \AIrows
\end{tabularx}}
\bareme{0,5 pt par ligne juste + 0,5 pt si les trois nombres d'hôtes sont justes ($2^{n}-2$, $n$ = nombre de bits d'hôte).}

\question{A2}{1 pt}{Indiquez le type de câble (\textbf{droit} ou \textbf{croisé}) pour chaque liaison.}
{\small\renewcommand{\arraystretch}{1.15}
\begin{tabularx}{\linewidth}{|X|c||X|c|}
  \hline
  \hl PC $\leftrightarrow$ commutateur & \makebox[1.6cm]{\sol{droit}} & Commutateur $\leftrightarrow$ routeur & \makebox[1.6cm]{\sol{droit}} \\ \hline
  \hl PC $\leftrightarrow$ PC (réseau P2P) & \sol{croisé} & Concentrateur $\leftrightarrow$ concentrateur & \sol{croisé} \\ \hline
\end{tabularx}}
\bareme{0,25 pt par case. Équipements de types différents : câble droit ; de même type : câble croisé.}

\question{A3}{1 pt}{Comment un \textbf{concentrateur (hub)} et un \textbf{commutateur (switch)} traitent-ils une trame
reçue ? Lequel limite les collisions, et pourquoi ?}
\zone{2.9cm}{Le \textbf{hub} (couche 1) répète la trame sur \textbf{tous ses autres ports} : un seul domaine de collision
partagé. Le \textbf{switch} (couche 2) lit l'adresse MAC de destination et, grâce à sa \textbf{table MAC}, ne l'envoie
\textbf{que sur le port du destinataire} : un domaine de collision par port, c'est lui qui limite les collisions.
\emph{0,5 pt hub, 0,5 pt switch.}}

\question{A4}{1 pt}{Qu'est-ce que la \textbf{passerelle par défaut} d'un poste ? Quelle condition son adresse doit-elle respecter ?}
\zone{2.4cm}{Adresse IP de l'interface du \textbf{routeur} à laquelle le poste envoie les paquets destinés à un
\textbf{autre réseau}. Elle doit \textbf{appartenir au même réseau} que le poste. \emph{0,5 + 0,5 pt.}}

\question{A5}{0,5 pt}{Un portable configuré en DHCP affiche \ip{\apipa} / \ip{255.255.0.0}. Que s'est-il passé ?}
\zone{1.5cm}{Aucun \textbf{serveur DHCP} n'a répondu : le poste s'est attribué une adresse \textbf{APIPA} (169.254.0.0/16).}

\question{A6}{0,5 pt}{Quel protocole utilise \texttt{ping} ? Nommez les deux messages échangés et la couche OSI des adresses IP.}
\zone{1.5cm}{\textbf{ICMP} ; \textbf{Echo Request} / \textbf{Echo Reply} ; adresses IP en \textbf{couche 3 (Réseau)}.}

%==========================================================================
\newpage
\section{Partie B — Réalisation sous Packet Tracer \hfill /10}

\begin{cdc}
\entreprise{} vous confie son réseau : \textbf{trois réseaux} reliés par un \textbf{routeur unique Cisco 2911}
nommé \ip{\Rnom} (interfaces G0/0, G0/1, G0/2).
\begin{itemize}[leftmargin=12pt]
  \item \textbf{Réseau \resUn} : \textbf{4 postes} PC-\pUn1 à PC-\pUn4, commutateur \ip{SW-\pUn}, adressage statique,
        réseau \ip{\netUn.0/24}.
  \item \textbf{Réseau \resDeux} : \textbf{2 postes} PC-\pDeux1, PC-\pDeux2 et \textbf{1 imprimante} IMP-\pDeux,
        commutateur \ip{SW-\pDeux}, adressage statique, réseau de classe B \ip{\netDeux.0.0/16}.
  \item \textbf{Réseau Serveur} : serveur \ip{\srvNom} d'adresse \ip{\netSrv.\srvH/24}, relié \textbf{directement} au routeur.
  \item La passerelle de chaque réseau est sa \textbf{dernière adresse utilisable}. Tous les équipements communiquent
        entre eux et accèdent à la page web du serveur. Nommage imposé (« Display Name »).
\end{itemize}
\end{cdc}

\question{B1}{2 pts}{Complétez le plan d'adressage \textbf{avant} la réalisation.}
{\footnotesize\setlength{\tabcolsep}{3pt}\renewcommand{\arraystretch}{1.15}
\begin{tabularx}{\linewidth}{|l|C|C|C|C|C|C|}
  \hline
  \textbf{Réseau} & \textbf{Adresse réseau} & \textbf{Masque} & \textbf{1re @ utilisable} &
  \textbf{Dernière @ utilisable} & \textbf{Diffusion} & \textbf{Passerelle} \\ \hline
  \hl \resUn & \sol{\netUn.0} & \sol{255.255.255.0} & \sol{\netUn.1} & \sol{\netUn.254} & \sol{\netUn.255} & \sol{\netUn.254} \\ \hline
  \hl \resDeux & \sol{\netDeux.0.0} & \sol{255.255.0.0} & \sol{\netDeux.0.1} & \sol{\netDeux.255.254} & \sol{\netDeux.255.255} & \sol{\netDeux.255.254} \\ \hline
  \hl Serveur & \sol{\netSrv.0} & \sol{255.255.255.0} & \sol{\netSrv.1} & \sol{\netSrv.254} & \sol{\netSrv.255} & \sol{\netSrv.254} \\ \hline
\end{tabularx}}
\bareme{0,5 pt par ligne juste + 0,5 pt si la convention de passerelle est respectée partout.}

\question{B2}{2 pts}{Construisez la topologie (équipements, câblage, nommage) et relevez l'adressage configuré.}
{\footnotesize\setlength{\tabcolsep}{3pt}\renewcommand{\arraystretch}{1.1}
\begin{tabularx}{\linewidth}{|l|C|C|C||l|C|C|}
  \hline
  \textbf{Hôte} & \textbf{Adresse IP} & \textbf{Masque} & \textbf{Passerelle} & \textbf{Interface} & \textbf{Adresse IP} & \textbf{Masque} \\ \hline
  \hl PC-\pUn1 & \sol{\netUn.1} & \sol{255.255.255.0} & \sol{\netUn.254} & \Rnom{} G0/0 & \sol{\netUn.254} & \sol{255.255.255.0} \\ \hline
  \hl PC-\pDeux1 & \sol{\netDeux.0.1} & \sol{255.255.0.0} & \sol{\netDeux.255.254} & \Rnom{} G0/1 & \sol{\netDeux.255.254} & \sol{255.255.0.0} \\ \hline
  \hl IMP-\pDeux & \sol{\netDeux.0.10} & \sol{255.255.0.0} & \sol{\netDeux.255.254} & \Rnom{} G0/2 & \sol{\netSrv.254} & \sol{255.255.255.0} \\ \hline
  \hl \srvNom & \ip{\netSrv.\srvH} & \sol{255.255.255.0} & \sol{\netSrv.254} & \multicolumn{3}{c|}{\cellcolor{black!8}} \\ \hline
\end{tabularx}}
\bareme{Adresses d'hôtes libres (uniques, dans le bon réseau) ; affectation G0/x libre. Sur le .pkt : équipements 0,5 ;
câblage 1 (câbles droits vers les switchs et switch–routeur, \textbf{croisé} serveur–routeur, liens verts) ; nommage 0,5.}

\question{B3}{3 pts}{Configurez tous les équipements : IP, masques, passerelles des hôtes ; interfaces du routeur
(« Port Status : On »).}
\bareme{Sur le .pkt : interfaces routeur 1 pt ; IP/masques des hôtes 1 pt ; passerelles (imprimante et serveur compris)
1 pt. $-0,25$ par erreur, sans descendre sous 0 par item.}

\question{B4}{3 pts}{Réalisez les tests, notez la commande et le résultat, puis faites viser chaque test réussi.}
{\small\renewcommand{\arraystretch}{1.1}
\begin{tabularx}{\linewidth}{|c|l|l|C|c|c|}
  \hline
  \textbf{N°} & \textbf{Source} & \textbf{Destination} & \textbf{Commande / action} & \textbf{Résultat} & \textbf{Visa prof.} \\ \hline
  \hl 1 & PC-\pUn1 & PC-\pUn4 & \sol{ping \netUn.4} & \sol{réussi} & \hspace{1.4cm} \\ \hline
  \hl 2 & PC-\pDeux1 & IMP-\pDeux & \sol{ping \netDeux.0.10} & \sol{réussi} & \\ \hline
  \hl 3 & PC-\pUn2 & PC-\pDeux2 & \sol{ping \netDeux.0.2} & \sol{réussi} & \\ \hline
  \hl 4 & PC-\pDeux2 & \srvNom & \sol{ping \netSrv.\srvH} & \sol{réussi} & \\ \hline
  \hl 5 & PC-\pUn3 & page web du serveur & \sol{navigateur : {\NoAutoSpacing http://\netSrv.\srvH}} & \sol{page affichée} & \\ \hline
\end{tabularx}}
\bareme{0,5 pt par test visé + 0,5 pt pour B5. Le premier ping inter-réseaux peut échouer (résolution ARP) ; les suivants doivent réussir.}

\question{B5}{incl. B4}{Test n°3 : par quels équipements passe le paquet ? Pourquoi PC-\pUn2 l'envoie-t-il au routeur ?}
\zone{2.3cm}{PC-\pUn2 $\to$ SW-\pUn{} $\to$ \Rnom{} $\to$ SW-\pDeux{} $\to$ PC-\pDeux2. Grâce à son masque, PC-\pUn2 constate que
la destination (\netDeux.0.0/16) n'est \textbf{pas dans son réseau} (\netUn.0/24) : il envoie le paquet à sa
\textbf{passerelle par défaut}, qui le route.}

%==========================================================================
\newpage
\section{Partie C — Diagnostic d'un réseau en panne \hfill /4}

Un technicien constate qu'\textbf{aucun poste ne parvient à joindre le serveur SRV-C} et que certains postes
communiquent mal entre eux. Les câbles sont du bon type.

\begin{center}
\begin{minipage}{0.37\linewidth}\centering
\begin{tikzpicture}[font=\footnotesize,
  eq/.style={draw,fill=qbg,rounded corners=2pt,minimum width=9mm,minimum height=5mm,inner sep=2pt},
  every path/.style={thick}]
  \node[eq] (SW) {SW-1};
  \node[eq,right=9mm of SW] (R) {R1};
  \node[eq,right=7mm of R] (SRV) {SRV-C};
  \node[eq,above left=5mm and -3mm of SW] (A) {PC-A};
  \node[eq,above right=5mm and -3mm of SW] (B) {PC-B};
  \node[eq,below left=5mm and -3mm of SW] (C) {PC-C};
  \node[eq,below right=5mm and -3mm of SW] (D) {PC-D};
  \foreach \n in {A,B,C,D} \draw (\n) -- (SW);
  \draw (SW) -- node[above,font=\tiny]{G0/0} (R);
  \draw (R) -- node[above,font=\tiny]{G0/1} (SRV);
\end{tikzpicture}
\end{minipage}\hfill
\begin{minipage}{0.61\linewidth}\centering
{\small\renewcommand{\arraystretch}{1.05}\setlength{\tabcolsep}{4pt}
\begin{tabular}{lllll}
  \toprule
  \textbf{Équipement} & \textbf{Adresse IP} & \textbf{Masque} & \textbf{Passerelle} & \textbf{Port} \\
  \midrule
  \Crows
  \bottomrule
\end{tabular}}
\end{minipage}
\end{center}

\question{C1}{4 pts}{Relevez les \textbf{quatre anomalies}. Pour chacune : conséquence et correction.}
{\small\renewcommand{\arraystretch}{1.1}
\ifcorrige\def\hc{}\else\def\hc{\rule[-10mm]{0pt}{14mm}}\fi
\begin{tabularx}{\linewidth}{|c|>{\raggedright\arraybackslash}p{2cm}|L|L|L|}
  \hline
  \textbf{N°} & \textbf{Équipement} & \textbf{Anomalie} & \textbf{Conséquence} & \textbf{Correction} \\ \hline
  \ifcorrige\Csol\else
  \hc 1 & & & & \\ \hline
  \hc 2 & & & & \\ \hline
  \hc 3 & & & & \\ \hline
  \hc 4 & & & & \\ \hline
  \fi
\end{tabularx}}
\bareme{1 pt par anomalie (0,25 équipement, anomalie, conséquence, correction ; ordre indifférent). \Cbareme}

%==========================================================================
\section{Bonus — Extension sans fil \hfill +1}

\begin{cdc}Ajoutez au réseau \resUn{} un \textbf{AccessPoint-PT} relié à SW-\pUn{} (SSID \ip{\ssid}, \textbf{WPA2-PSK})
et un portable \ip{Laptop-\pUn} équipé d'un module \texttt{WPC300N}, en adressage \textbf{statique}.
Vérifiez qu'il joint \srvNom.\end{cdc}
{\small\renewcommand{\arraystretch}{1.1}
\begin{tabularx}{\linewidth}{|l|C|C|C|c|}
  \hline
  \textbf{Équipement} & \textbf{Adresse IP} & \textbf{Masque} & \textbf{Passerelle} & \textbf{Visa prof.} \\ \hline
  \hl Laptop-\pUn & \sol{\netUn.50} & \sol{255.255.255.0} & \sol{\netUn.254} & \hspace{1.4cm} \\ \hline
\end{tabularx}}
\bareme{Le point d'accès fait office de pont : le portable est dans \netUn.0/24. Bonus si associé en WPA2-PSK et ping
vers \netSrv.\srvH{} réussi. Note plafonnée à 20.}

\end{document}
"
\documentclass[10pt,a4paper]{article}

\usepackage[T1]{fontenc}
\usepackage[utf8]{inputenc}
\usepackage{lmodern}
\usepackage[french]{babel}
\usepackage[a4paper,hmargin=1.5cm,top=1.3cm,bottom=1.6cm,footskip=0.7cm]{geometry}
\usepackage[table]{xcolor}
\usepackage{booktabs,tabularx,array}
\usepackage{enumitem}
\usepackage{fancyhdr}
\usepackage[most]{tcolorbox}
\usepackage{tikz}
\usetikzlibrary{positioning}
\usepackage{titlesec}

\providecommand{\modecorrige}{0}
\providecommand{\variante}{A}
\newif\ifcorrige
\ifnum\modecorrige=1 \corrigetrue \else \corrigefalse \fi

\definecolor{ipred}{HTML}{C0392B}
\definecolor{cdcbar}{HTML}{A67C3D}
\definecolor{cdcbg}{HTML}{F4F1EA}
\definecolor{qbar}{HTML}{4A6B5D}
\definecolor{qbg}{HTML}{EDF1EF}
\definecolor{corr}{HTML}{1F5FA8}
\definecolor{corrbg}{HTML}{E8F0FA}
\definecolor{grisfoot}{HTML}{828282}

\newcommand{\ip}[1]{\textcolor{ipred}{\texttt{#1}}}
\newcommand{\pts}[1]{\hfill\mbox{\textbf{(#1)}}}
\newcommand{\sol}[1]{\ifcorrige\textcolor{corr}{#1}\fi}
\newcommand{\bareme}[1]{\ifcorrige\par{\footnotesize\color{corr}#1\par}\fi}
\newcommand{\hl}{\ifcorrige\rule[-1.6mm]{0pt}{5.4mm}\else\rule[-2.2mm]{0pt}{6.6mm}\fi}
\newcolumntype{C}{>{\centering\arraybackslash}X}
\newcolumntype{L}{>{\raggedright\arraybackslash}X}

\titleformat{\section}{\large\bfseries}{}{0pt}{}
\titlespacing*{\section}{0pt}{6pt}{3pt}

\tcbset{qstyle/.style={enhanced,boxrule=0pt,leftrule=2.5pt,arc=0pt,outer arc=0pt,
  left=5pt,right=5pt,top=2pt,bottom=2pt,before skip=4pt,after skip=3pt}}
\newtcolorbox{cdc}{qstyle,colback=cdcbg,colframe=cdcbar}
\newtcolorbox{quest}{qstyle,colback=qbg,colframe=qbar}

% Zone de réponse : vierge dans le sujet, remplie dans le corrigé
\newcommand{\zone}[2]{%
  \ifcorrige
    \begin{tcolorbox}[colback=corrbg,colframe=corr,boxrule=0.5pt,arc=1pt,left=4pt,right=4pt,top=1pt,bottom=1pt,
      before skip=2pt,after skip=2pt,fontupper=\footnotesize\color{corr}]#2\end{tcolorbox}%
  \else
    \begin{tcolorbox}[colback=white,colframe=black!45,boxrule=0.4pt,arc=1pt,height=#1,
      before skip=2pt,after skip=2pt]\end{tcolorbox}%
  \fi}

\newcommand{\question}[3]{\begin{quest}\textbf{#1.} #3\pts{#2}\end{quest}}

%==========================================================================
% Données propres à chaque variante
%==========================================================================
\if\variante A
  \def\entreprise{Le cabinet d'architectes \textbf{ArchiNova}}
  \def\Rnom{R-ArchiNova}
  \def\resUn{Bureau d'études}\def\pUn{BE}\def\netUn{192.168.40}
  \def\resDeux{Accueil}\def\pDeux{ACC}\def\netDeux{172.16}
  \def\srvNom{SRV-Intranet}\def\netSrv{192.168.200}\def\srvH{10}
  \def\ssid{ArchiNova-WiFi}
  \def\AIrows{%
    \hl\ip{192.168.40.25} & \sol{C} & \sol{255.255.255.0} & \sol{192.168.40.0} & \sol{192.168.40.255} & \sol{254} \\ \hline
    \hl\ip{172.16.5.10}   & \sol{B} & \sol{255.255.0.0}   & \sol{172.16.0.0}   & \sol{172.16.255.255} & \sol{65\,534} \\ \hline
    \hl\ip{10.3.2.1}      & \sol{A} & \sol{255.0.0.0}     & \sol{10.0.0.0}     & \sol{10.255.255.255} & \sol{16\,777\,214} \\ \hline}
  \def\apipa{169.254.12.7}
  \def\Cnet{192.168.60}
  \def\Crows{%
    PC-A & \ip{192.168.60.10} & \ip{255.255.255.0} & \ip{192.168.60.254} & — \\
    PC-B & \ip{192.168.60.20} & \ip{255.255.255.0} & \ip{192.168.60.1}   & — \\
    PC-C & \ip{192.168.61.30} & \ip{255.255.255.0} & \ip{192.168.60.254} & — \\
    PC-D & \ip{192.168.60.10} & \ip{255.255.255.0} & \ip{192.168.60.254} & — \\
    SRV-C & \ip{10.0.0.5}     & \ip{255.0.0.0}     & \ip{10.0.0.254}     & — \\
    R1 — G0/0 & \ip{192.168.60.254} & \ip{255.255.255.0} & — & On \\
    R1 — G0/1 & \ip{10.0.0.254}     & \ip{255.0.0.0}     & — & Off \\}
  \def\Csol{%
    \hl 1 & \sol{PC-B} & \sol{Passerelle .60.1 : pas l'adresse du routeur.} & \sol{Joint son réseau, mais aucun autre (serveur).} & \sol{Passerelle 192.168.60.254.} \\ \hline
    \hl 2 & \sol{PC-C} & \sol{IP hors du réseau 192.168.60.0/24.} & \sol{Ne joint ni les postes ni le routeur.} & \sol{Ex. 192.168.60.30.} \\ \hline
    \hl 3 & \sol{PC-D} & \sol{Même IP que PC-A (doublon).} & \sol{Conflit d'adresse : PC-A et PC-D perturbés.} & \sol{Adresse unique, ex. .60.40.} \\ \hline
    \hl 4 & \sol{R1 G0/1} & \sol{Interface désactivée (Off).} & \sol{Réseau du serveur injoignable pour tous.} & \sol{Port Status : On.} \\ \hline}
  \def\Cbareme{PC-A et SRV-C sont correctement configurés ; accepter PC-A cité à la place de PC-D pour le doublon.}
\else
  \def\entreprise{La clinique vétérinaire \textbf{VetoPlus}}
  \def\Rnom{R-VetoPlus}
  \def\resUn{Soins}\def\pUn{SOI}\def\netUn{192.168.70}
  \def\resDeux{Secrétariat}\def\pDeux{SEC}\def\netDeux{172.20}
  \def\srvNom{SRV-Dossiers}\def\netSrv{192.168.250}\def\srvH{20}
  \def\ssid{VetoPlus-WiFi}
  \def\AIrows{%
    \hl\ip{192.168.15.130} & \sol{C} & \sol{255.255.255.0} & \sol{192.168.15.0} & \sol{192.168.15.255} & \sol{254} \\ \hline
    \hl\ip{172.31.200.7}   & \sol{B} & \sol{255.255.0.0}   & \sol{172.31.0.0}   & \sol{172.31.255.255} & \sol{65\,534} \\ \hline
    \hl\ip{10.45.6.9}      & \sol{A} & \sol{255.0.0.0}     & \sol{10.0.0.0}     & \sol{10.255.255.255} & \sol{16\,777\,214} \\ \hline}
  \def\apipa{169.254.88.3}
  \def\Cnet{192.168.80}
  \def\Crows{%
    PC-A & \ip{192.168.8.11}   & \ip{255.255.255.0} & \ip{192.168.80.254} & — \\
    PC-B & \ip{192.168.80.12}  & \ip{255.255.255.0} & \ip{192.168.80.254} & — \\
    PC-C & \ip{192.168.80.13}  & \ip{255.255.255.0} & \ip{192.168.80.100} & — \\
    PC-D & \ip{192.168.80.12}  & \ip{255.255.255.0} & \ip{192.168.80.254} & — \\
    SRV-C & \ip{10.0.0.8}      & \ip{255.0.0.0}     & \ip{10.0.0.254}     & — \\
    R1 — G0/0 & \ip{192.168.80.254} & \ip{255.255.255.0} & — & Off \\
    R1 — G0/1 & \ip{10.0.0.254}     & \ip{255.0.0.0}     & — & On \\}
  \def\Csol{%
    \hl 1 & \sol{PC-A} & \sol{IP hors du réseau 192.168.80.0/24.} & \sol{Ne joint ni les postes ni le routeur.} & \sol{Ex. 192.168.80.11.} \\ \hline
    \hl 2 & \sol{PC-C} & \sol{Passerelle .80.100 : pas l'adresse du routeur.} & \sol{Joint son réseau, mais aucun autre (serveur).} & \sol{Passerelle 192.168.80.254.} \\ \hline
    \hl 3 & \sol{PC-D} & \sol{Même IP que PC-B (doublon).} & \sol{Conflit d'adresse : PC-B et PC-D perturbés.} & \sol{Adresse unique, ex. .80.14.} \\ \hline
    \hl 4 & \sol{R1 G0/0} & \sol{Interface côté LAN désactivée (Off).} & \sol{Aucun poste ne joint le routeur ni le serveur.} & \sol{Port Status : On.} \\ \hline}
  \def\Cbareme{PC-B et SRV-C sont correctement configurés ; accepter PC-B cité à la place de PC-D pour le doublon.}
\fi

\pagestyle{fancy}
\fancyhf{}
\renewcommand{\headrulewidth}{0pt}
\fancyfoot[L]{\footnotesize\itshape\color{grisfoot}Évaluation — Réseaux locaux — Sujet \variante\ifcorrige\ — CORRIGÉ\fi}
\fancyfoot[R]{\footnotesize\color{grisfoot}\thepage/3}

\setlist{itemsep=0pt,topsep=2pt,parsep=0pt}
\setlength{\parindent}{0pt}
\setlength{\parskip}{2pt}

\begin{document}

\begin{center}
  {\Large\bfseries Évaluation pratique — Réseaux locaux \quad
   \fcolorbox{black}{white}{\ SUJET \variante\ }\par}
  \vspace{1pt}
  {\bfseries Adressage IP, commutation et routage — BTS CIEL 1re année — durée : 2 heures\par}
  \ifcorrige\vspace{2pt}{\bfseries\color{corr}CORRIGÉ ET BARÈME — document enseignant}\par\fi
\end{center}
\vspace{-2pt}
{\renewcommand{\arraystretch}{1.35}
\begin{tabularx}{\linewidth}{|X|X|p{2.2cm}|p{2.6cm}|}
  \hline
  \textbf{Nom :} & \textbf{Prénom :} & \textbf{Poste n°} & \textbf{Note :} \hfill /20 \\ \hline
\end{tabularx}}

\textbf{Consignes :} travail \textbf{individuel}, documents et Internet \textbf{non autorisés}. Ce sujet est le
\textbf{document-réponse}. Fichier Packet Tracer : \ip{Nom\_Prenom\_EvalLAN\_\variante.pkt}, à enregistrer
régulièrement dans le dossier de rendu. Les « visas » sont apposés par le professeur, que vous appelez à chaque test réussi.
Durées conseillées : A 30 min — B 1 h 10 — C 20 min.

%==========================================================================
\section{Partie A — Questions de cours et calculs \hfill /6}

\question{A1}{2 pts}{Complétez le tableau en considérant le \textbf{masque par défaut} de la classe de chaque adresse.}
{\small\renewcommand{\arraystretch}{1.15}
\begin{tabularx}{\linewidth}{|l|c|C|C|C|c|}
  \hline
  \textbf{Adresse IP} & \textbf{Classe} & \textbf{Masque par défaut} & \textbf{Adresse réseau} &
  \textbf{Adresse de diffusion} & \textbf{Nb d'hôtes} \\ \hline
  \AIrows
\end{tabularx}}
\bareme{0,5 pt par ligne juste + 0,5 pt si les trois nombres d'hôtes sont justes ($2^{n}-2$, $n$ = nombre de bits d'hôte).}

\question{A2}{1 pt}{Indiquez le type de câble (\textbf{droit} ou \textbf{croisé}) pour chaque liaison.}
{\small\renewcommand{\arraystretch}{1.15}
\begin{tabularx}{\linewidth}{|X|c||X|c|}
  \hline
  \hl PC $\leftrightarrow$ commutateur & \makebox[1.6cm]{\sol{droit}} & Commutateur $\leftrightarrow$ routeur & \makebox[1.6cm]{\sol{droit}} \\ \hline
  \hl PC $\leftrightarrow$ PC (réseau P2P) & \sol{croisé} & Concentrateur $\leftrightarrow$ concentrateur & \sol{croisé} \\ \hline
\end{tabularx}}
\bareme{0,25 pt par case. Équipements de types différents : câble droit ; de même type : câble croisé.}

\question{A3}{1 pt}{Comment un \textbf{concentrateur (hub)} et un \textbf{commutateur (switch)} traitent-ils une trame
reçue ? Lequel limite les collisions, et pourquoi ?}
\zone{2.9cm}{Le \textbf{hub} (couche 1) répète la trame sur \textbf{tous ses autres ports} : un seul domaine de collision
partagé. Le \textbf{switch} (couche 2) lit l'adresse MAC de destination et, grâce à sa \textbf{table MAC}, ne l'envoie
\textbf{que sur le port du destinataire} : un domaine de collision par port, c'est lui qui limite les collisions.
\emph{0,5 pt hub, 0,5 pt switch.}}

\question{A4}{1 pt}{Qu'est-ce que la \textbf{passerelle par défaut} d'un poste ? Quelle condition son adresse doit-elle respecter ?}
\zone{2.4cm}{Adresse IP de l'interface du \textbf{routeur} à laquelle le poste envoie les paquets destinés à un
\textbf{autre réseau}. Elle doit \textbf{appartenir au même réseau} que le poste. \emph{0,5 + 0,5 pt.}}

\question{A5}{0,5 pt}{Un portable configuré en DHCP affiche \ip{\apipa} / \ip{255.255.0.0}. Que s'est-il passé ?}
\zone{1.5cm}{Aucun \textbf{serveur DHCP} n'a répondu : le poste s'est attribué une adresse \textbf{APIPA} (169.254.0.0/16).}

\question{A6}{0,5 pt}{Quel protocole utilise \texttt{ping} ? Nommez les deux messages échangés et la couche OSI des adresses IP.}
\zone{1.5cm}{\textbf{ICMP} ; \textbf{Echo Request} / \textbf{Echo Reply} ; adresses IP en \textbf{couche 3 (Réseau)}.}

%==========================================================================
\newpage
\section{Partie B — Réalisation sous Packet Tracer \hfill /10}

\begin{cdc}
\entreprise{} vous confie son réseau : \textbf{trois réseaux} reliés par un \textbf{routeur unique Cisco 2911}
nommé \ip{\Rnom} (interfaces G0/0, G0/1, G0/2).
\begin{itemize}[leftmargin=12pt]
  \item \textbf{Réseau \resUn} : \textbf{4 postes} PC-\pUn1 à PC-\pUn4, commutateur \ip{SW-\pUn}, adressage statique,
        réseau \ip{\netUn.0/24}.
  \item \textbf{Réseau \resDeux} : \textbf{2 postes} PC-\pDeux1, PC-\pDeux2 et \textbf{1 imprimante} IMP-\pDeux,
        commutateur \ip{SW-\pDeux}, adressage statique, réseau de classe B \ip{\netDeux.0.0/16}.
  \item \textbf{Réseau Serveur} : serveur \ip{\srvNom} d'adresse \ip{\netSrv.\srvH/24}, relié \textbf{directement} au routeur.
  \item La passerelle de chaque réseau est sa \textbf{dernière adresse utilisable}. Tous les équipements communiquent
        entre eux et accèdent à la page web du serveur. Nommage imposé (« Display Name »).
\end{itemize}
\end{cdc}

\question{B1}{2 pts}{Complétez le plan d'adressage \textbf{avant} la réalisation.}
{\footnotesize\setlength{\tabcolsep}{3pt}\renewcommand{\arraystretch}{1.15}
\begin{tabularx}{\linewidth}{|l|C|C|C|C|C|C|}
  \hline
  \textbf{Réseau} & \textbf{Adresse réseau} & \textbf{Masque} & \textbf{1re @ utilisable} &
  \textbf{Dernière @ utilisable} & \textbf{Diffusion} & \textbf{Passerelle} \\ \hline
  \hl \resUn & \sol{\netUn.0} & \sol{255.255.255.0} & \sol{\netUn.1} & \sol{\netUn.254} & \sol{\netUn.255} & \sol{\netUn.254} \\ \hline
  \hl \resDeux & \sol{\netDeux.0.0} & \sol{255.255.0.0} & \sol{\netDeux.0.1} & \sol{\netDeux.255.254} & \sol{\netDeux.255.255} & \sol{\netDeux.255.254} \\ \hline
  \hl Serveur & \sol{\netSrv.0} & \sol{255.255.255.0} & \sol{\netSrv.1} & \sol{\netSrv.254} & \sol{\netSrv.255} & \sol{\netSrv.254} \\ \hline
\end{tabularx}}
\bareme{0,5 pt par ligne juste + 0,5 pt si la convention de passerelle est respectée partout.}

\question{B2}{2 pts}{Construisez la topologie (équipements, câblage, nommage) et relevez l'adressage configuré.}
{\footnotesize\setlength{\tabcolsep}{3pt}\renewcommand{\arraystretch}{1.1}
\begin{tabularx}{\linewidth}{|l|C|C|C||l|C|C|}
  \hline
  \textbf{Hôte} & \textbf{Adresse IP} & \textbf{Masque} & \textbf{Passerelle} & \textbf{Interface} & \textbf{Adresse IP} & \textbf{Masque} \\ \hline
  \hl PC-\pUn1 & \sol{\netUn.1} & \sol{255.255.255.0} & \sol{\netUn.254} & \Rnom{} G0/0 & \sol{\netUn.254} & \sol{255.255.255.0} \\ \hline
  \hl PC-\pDeux1 & \sol{\netDeux.0.1} & \sol{255.255.0.0} & \sol{\netDeux.255.254} & \Rnom{} G0/1 & \sol{\netDeux.255.254} & \sol{255.255.0.0} \\ \hline
  \hl IMP-\pDeux & \sol{\netDeux.0.10} & \sol{255.255.0.0} & \sol{\netDeux.255.254} & \Rnom{} G0/2 & \sol{\netSrv.254} & \sol{255.255.255.0} \\ \hline
  \hl \srvNom & \ip{\netSrv.\srvH} & \sol{255.255.255.0} & \sol{\netSrv.254} & \multicolumn{3}{c|}{\cellcolor{black!8}} \\ \hline
\end{tabularx}}
\bareme{Adresses d'hôtes libres (uniques, dans le bon réseau) ; affectation G0/x libre. Sur le .pkt : équipements 0,5 ;
câblage 1 (câbles droits vers les switchs et switch–routeur, \textbf{croisé} serveur–routeur, liens verts) ; nommage 0,5.}

\question{B3}{3 pts}{Configurez tous les équipements : IP, masques, passerelles des hôtes ; interfaces du routeur
(« Port Status : On »).}
\bareme{Sur le .pkt : interfaces routeur 1 pt ; IP/masques des hôtes 1 pt ; passerelles (imprimante et serveur compris)
1 pt. $-0,25$ par erreur, sans descendre sous 0 par item.}

\question{B4}{3 pts}{Réalisez les tests, notez la commande et le résultat, puis faites viser chaque test réussi.}
{\small\renewcommand{\arraystretch}{1.1}
\begin{tabularx}{\linewidth}{|c|l|l|C|c|c|}
  \hline
  \textbf{N°} & \textbf{Source} & \textbf{Destination} & \textbf{Commande / action} & \textbf{Résultat} & \textbf{Visa prof.} \\ \hline
  \hl 1 & PC-\pUn1 & PC-\pUn4 & \sol{ping \netUn.4} & \sol{réussi} & \hspace{1.4cm} \\ \hline
  \hl 2 & PC-\pDeux1 & IMP-\pDeux & \sol{ping \netDeux.0.10} & \sol{réussi} & \\ \hline
  \hl 3 & PC-\pUn2 & PC-\pDeux2 & \sol{ping \netDeux.0.2} & \sol{réussi} & \\ \hline
  \hl 4 & PC-\pDeux2 & \srvNom & \sol{ping \netSrv.\srvH} & \sol{réussi} & \\ \hline
  \hl 5 & PC-\pUn3 & page web du serveur & \sol{navigateur : {\NoAutoSpacing http://\netSrv.\srvH}} & \sol{page affichée} & \\ \hline
\end{tabularx}}
\bareme{0,5 pt par test visé + 0,5 pt pour B5. Le premier ping inter-réseaux peut échouer (résolution ARP) ; les suivants doivent réussir.}

\question{B5}{incl. B4}{Test n°3 : par quels équipements passe le paquet ? Pourquoi PC-\pUn2 l'envoie-t-il au routeur ?}
\zone{2.3cm}{PC-\pUn2 $\to$ SW-\pUn{} $\to$ \Rnom{} $\to$ SW-\pDeux{} $\to$ PC-\pDeux2. Grâce à son masque, PC-\pUn2 constate que
la destination (\netDeux.0.0/16) n'est \textbf{pas dans son réseau} (\netUn.0/24) : il envoie le paquet à sa
\textbf{passerelle par défaut}, qui le route.}

%==========================================================================
\newpage
\section{Partie C — Diagnostic d'un réseau en panne \hfill /4}

Un technicien constate qu'\textbf{aucun poste ne parvient à joindre le serveur SRV-C} et que certains postes
communiquent mal entre eux. Les câbles sont du bon type.

\begin{center}
\begin{minipage}{0.37\linewidth}\centering
\begin{tikzpicture}[font=\footnotesize,
  eq/.style={draw,fill=qbg,rounded corners=2pt,minimum width=9mm,minimum height=5mm,inner sep=2pt},
  every path/.style={thick}]
  \node[eq] (SW) {SW-1};
  \node[eq,right=9mm of SW] (R) {R1};
  \node[eq,right=7mm of R] (SRV) {SRV-C};
  \node[eq,above left=5mm and -3mm of SW] (A) {PC-A};
  \node[eq,above right=5mm and -3mm of SW] (B) {PC-B};
  \node[eq,below left=5mm and -3mm of SW] (C) {PC-C};
  \node[eq,below right=5mm and -3mm of SW] (D) {PC-D};
  \foreach \n in {A,B,C,D} \draw (\n) -- (SW);
  \draw (SW) -- node[above,font=\tiny]{G0/0} (R);
  \draw (R) -- node[above,font=\tiny]{G0/1} (SRV);
\end{tikzpicture}
\end{minipage}\hfill
\begin{minipage}{0.61\linewidth}\centering
{\small\renewcommand{\arraystretch}{1.05}\setlength{\tabcolsep}{4pt}
\begin{tabular}{lllll}
  \toprule
  \textbf{Équipement} & \textbf{Adresse IP} & \textbf{Masque} & \textbf{Passerelle} & \textbf{Port} \\
  \midrule
  \Crows
  \bottomrule
\end{tabular}}
\end{minipage}
\end{center}

\question{C1}{4 pts}{Relevez les \textbf{quatre anomalies}. Pour chacune : conséquence et correction.}
{\small\renewcommand{\arraystretch}{1.1}
\ifcorrige\def\hc{}\else\def\hc{\rule[-10mm]{0pt}{14mm}}\fi
\begin{tabularx}{\linewidth}{|c|>{\raggedright\arraybackslash}p{2cm}|L|L|L|}
  \hline
  \textbf{N°} & \textbf{Équipement} & \textbf{Anomalie} & \textbf{Conséquence} & \textbf{Correction} \\ \hline
  \ifcorrige\Csol\else
  \hc 1 & & & & \\ \hline
  \hc 2 & & & & \\ \hline
  \hc 3 & & & & \\ \hline
  \hc 4 & & & & \\ \hline
  \fi
\end{tabularx}}
\bareme{1 pt par anomalie (0,25 équipement, anomalie, conséquence, correction ; ordre indifférent). \Cbareme}

%==========================================================================
\section{Bonus — Extension sans fil \hfill +1}

\begin{cdc}Ajoutez au réseau \resUn{} un \textbf{AccessPoint-PT} relié à SW-\pUn{} (SSID \ip{\ssid}, \textbf{WPA2-PSK})
et un portable \ip{Laptop-\pUn} équipé d'un module \texttt{WPC300N}, en adressage \textbf{statique}.
Vérifiez qu'il joint \srvNom.\end{cdc}
{\small\renewcommand{\arraystretch}{1.1}
\begin{tabularx}{\linewidth}{|l|C|C|C|c|}
  \hline
  \textbf{Équipement} & \textbf{Adresse IP} & \textbf{Masque} & \textbf{Passerelle} & \textbf{Visa prof.} \\ \hline
  \hl Laptop-\pUn & \sol{\netUn.50} & \sol{255.255.255.0} & \sol{\netUn.254} & \hspace{1.4cm} \\ \hline
\end{tabularx}}
\bareme{Le point d'accès fait office de pont : le portable est dans \netUn.0/24. Bonus si associé en WPA2-PSK et ping
vers \netSrv.\srvH{} réussi. Note plafonnée à 20.}

\end{document}
"
\documentclass[10pt,a4paper]{article}

\usepackage[T1]{fontenc}
\usepackage[utf8]{inputenc}
\usepackage{lmodern}
\usepackage[french]{babel}
\usepackage[a4paper,hmargin=1.5cm,top=1.3cm,bottom=1.6cm,footskip=0.7cm]{geometry}
\usepackage[table]{xcolor}
\usepackage{booktabs,tabularx,array}
\usepackage{enumitem}
\usepackage{fancyhdr}
\usepackage[most]{tcolorbox}
\usepackage{tikz}
\usetikzlibrary{positioning}
\usepackage{titlesec}

\providecommand{\modecorrige}{0}
\providecommand{\variante}{A}
\newif\ifcorrige
\ifnum\modecorrige=1 \corrigetrue \else \corrigefalse \fi

\definecolor{ipred}{HTML}{C0392B}
\definecolor{cdcbar}{HTML}{A67C3D}
\definecolor{cdcbg}{HTML}{F4F1EA}
\definecolor{qbar}{HTML}{4A6B5D}
\definecolor{qbg}{HTML}{EDF1EF}
\definecolor{corr}{HTML}{1F5FA8}
\definecolor{corrbg}{HTML}{E8F0FA}
\definecolor{grisfoot}{HTML}{828282}

\newcommand{\ip}[1]{\textcolor{ipred}{\texttt{#1}}}
\newcommand{\pts}[1]{\hfill\mbox{\textbf{(#1)}}}
\newcommand{\sol}[1]{\ifcorrige\textcolor{corr}{#1}\fi}
\newcommand{\bareme}[1]{\ifcorrige\par{\footnotesize\color{corr}#1\par}\fi}
\newcommand{\hl}{\ifcorrige\rule[-1.6mm]{0pt}{5.4mm}\else\rule[-2.2mm]{0pt}{6.6mm}\fi}
\newcolumntype{C}{>{\centering\arraybackslash}X}
\newcolumntype{L}{>{\raggedright\arraybackslash}X}

\titleformat{\section}{\large\bfseries}{}{0pt}{}
\titlespacing*{\section}{0pt}{6pt}{3pt}

\tcbset{qstyle/.style={enhanced,boxrule=0pt,leftrule=2.5pt,arc=0pt,outer arc=0pt,
  left=5pt,right=5pt,top=2pt,bottom=2pt,before skip=4pt,after skip=3pt}}
\newtcolorbox{cdc}{qstyle,colback=cdcbg,colframe=cdcbar}
\newtcolorbox{quest}{qstyle,colback=qbg,colframe=qbar}

% Zone de réponse : vierge dans le sujet, remplie dans le corrigé
\newcommand{\zone}[2]{%
  \ifcorrige
    \begin{tcolorbox}[colback=corrbg,colframe=corr,boxrule=0.5pt,arc=1pt,left=4pt,right=4pt,top=1pt,bottom=1pt,
      before skip=2pt,after skip=2pt,fontupper=\footnotesize\color{corr}]#2\end{tcolorbox}%
  \else
    \begin{tcolorbox}[colback=white,colframe=black!45,boxrule=0.4pt,arc=1pt,height=#1,
      before skip=2pt,after skip=2pt]\end{tcolorbox}%
  \fi}

\newcommand{\question}[3]{\begin{quest}\textbf{#1.} #3\pts{#2}\end{quest}}

%==========================================================================
% Données propres à chaque variante
%==========================================================================
\if\variante A
  \def\entreprise{Le cabinet d'architectes \textbf{ArchiNova}}
  \def\Rnom{R-ArchiNova}
  \def\resUn{Bureau d'études}\def\pUn{BE}\def\netUn{192.168.40}
  \def\resDeux{Accueil}\def\pDeux{ACC}\def\netDeux{172.16}
  \def\srvNom{SRV-Intranet}\def\netSrv{192.168.200}\def\srvH{10}
  \def\ssid{ArchiNova-WiFi}
  \def\AIrows{%
    \hl\ip{192.168.40.25} & \sol{C} & \sol{255.255.255.0} & \sol{192.168.40.0} & \sol{192.168.40.255} & \sol{254} \\ \hline
    \hl\ip{172.16.5.10}   & \sol{B} & \sol{255.255.0.0}   & \sol{172.16.0.0}   & \sol{172.16.255.255} & \sol{65\,534} \\ \hline
    \hl\ip{10.3.2.1}      & \sol{A} & \sol{255.0.0.0}     & \sol{10.0.0.0}     & \sol{10.255.255.255} & \sol{16\,777\,214} \\ \hline}
  \def\apipa{169.254.12.7}
  \def\Cnet{192.168.60}
  \def\Crows{%
    PC-A & \ip{192.168.60.10} & \ip{255.255.255.0} & \ip{192.168.60.254} & — \\
    PC-B & \ip{192.168.60.20} & \ip{255.255.255.0} & \ip{192.168.60.1}   & — \\
    PC-C & \ip{192.168.61.30} & \ip{255.255.255.0} & \ip{192.168.60.254} & — \\
    PC-D & \ip{192.168.60.10} & \ip{255.255.255.0} & \ip{192.168.60.254} & — \\
    SRV-C & \ip{10.0.0.5}     & \ip{255.0.0.0}     & \ip{10.0.0.254}     & — \\
    R1 — G0/0 & \ip{192.168.60.254} & \ip{255.255.255.0} & — & On \\
    R1 — G0/1 & \ip{10.0.0.254}     & \ip{255.0.0.0}     & — & Off \\}
  \def\Csol{%
    \hl 1 & \sol{PC-B} & \sol{Passerelle .60.1 : pas l'adresse du routeur.} & \sol{Joint son réseau, mais aucun autre (serveur).} & \sol{Passerelle 192.168.60.254.} \\ \hline
    \hl 2 & \sol{PC-C} & \sol{IP hors du réseau 192.168.60.0/24.} & \sol{Ne joint ni les postes ni le routeur.} & \sol{Ex. 192.168.60.30.} \\ \hline
    \hl 3 & \sol{PC-D} & \sol{Même IP que PC-A (doublon).} & \sol{Conflit d'adresse : PC-A et PC-D perturbés.} & \sol{Adresse unique, ex. .60.40.} \\ \hline
    \hl 4 & \sol{R1 G0/1} & \sol{Interface désactivée (Off).} & \sol{Réseau du serveur injoignable pour tous.} & \sol{Port Status : On.} \\ \hline}
  \def\Cbareme{PC-A et SRV-C sont correctement configurés ; accepter PC-A cité à la place de PC-D pour le doublon.}
\else
  \def\entreprise{La clinique vétérinaire \textbf{VetoPlus}}
  \def\Rnom{R-VetoPlus}
  \def\resUn{Soins}\def\pUn{SOI}\def\netUn{192.168.70}
  \def\resDeux{Secrétariat}\def\pDeux{SEC}\def\netDeux{172.20}
  \def\srvNom{SRV-Dossiers}\def\netSrv{192.168.250}\def\srvH{20}
  \def\ssid{VetoPlus-WiFi}
  \def\AIrows{%
    \hl\ip{192.168.15.130} & \sol{C} & \sol{255.255.255.0} & \sol{192.168.15.0} & \sol{192.168.15.255} & \sol{254} \\ \hline
    \hl\ip{172.31.200.7}   & \sol{B} & \sol{255.255.0.0}   & \sol{172.31.0.0}   & \sol{172.31.255.255} & \sol{65\,534} \\ \hline
    \hl\ip{10.45.6.9}      & \sol{A} & \sol{255.0.0.0}     & \sol{10.0.0.0}     & \sol{10.255.255.255} & \sol{16\,777\,214} \\ \hline}
  \def\apipa{169.254.88.3}
  \def\Cnet{192.168.80}
  \def\Crows{%
    PC-A & \ip{192.168.8.11}   & \ip{255.255.255.0} & \ip{192.168.80.254} & — \\
    PC-B & \ip{192.168.80.12}  & \ip{255.255.255.0} & \ip{192.168.80.254} & — \\
    PC-C & \ip{192.168.80.13}  & \ip{255.255.255.0} & \ip{192.168.80.100} & — \\
    PC-D & \ip{192.168.80.12}  & \ip{255.255.255.0} & \ip{192.168.80.254} & — \\
    SRV-C & \ip{10.0.0.8}      & \ip{255.0.0.0}     & \ip{10.0.0.254}     & — \\
    R1 — G0/0 & \ip{192.168.80.254} & \ip{255.255.255.0} & — & Off \\
    R1 — G0/1 & \ip{10.0.0.254}     & \ip{255.0.0.0}     & — & On \\}
  \def\Csol{%
    \hl 1 & \sol{PC-A} & \sol{IP hors du réseau 192.168.80.0/24.} & \sol{Ne joint ni les postes ni le routeur.} & \sol{Ex. 192.168.80.11.} \\ \hline
    \hl 2 & \sol{PC-C} & \sol{Passerelle .80.100 : pas l'adresse du routeur.} & \sol{Joint son réseau, mais aucun autre (serveur).} & \sol{Passerelle 192.168.80.254.} \\ \hline
    \hl 3 & \sol{PC-D} & \sol{Même IP que PC-B (doublon).} & \sol{Conflit d'adresse : PC-B et PC-D perturbés.} & \sol{Adresse unique, ex. .80.14.} \\ \hline
    \hl 4 & \sol{R1 G0/0} & \sol{Interface côté LAN désactivée (Off).} & \sol{Aucun poste ne joint le routeur ni le serveur.} & \sol{Port Status : On.} \\ \hline}
  \def\Cbareme{PC-B et SRV-C sont correctement configurés ; accepter PC-B cité à la place de PC-D pour le doublon.}
\fi

\pagestyle{fancy}
\fancyhf{}
\renewcommand{\headrulewidth}{0pt}
\fancyfoot[L]{\footnotesize\itshape\color{grisfoot}Évaluation — Réseaux locaux — Sujet \variante\ifcorrige\ — CORRIGÉ\fi}
\fancyfoot[R]{\footnotesize\color{grisfoot}\thepage/3}

\setlist{itemsep=0pt,topsep=2pt,parsep=0pt}
\setlength{\parindent}{0pt}
\setlength{\parskip}{2pt}

\begin{document}

\begin{center}
  {\Large\bfseries Évaluation pratique — Réseaux locaux \quad
   \fcolorbox{black}{white}{\ SUJET \variante\ }\par}
  \vspace{1pt}
  {\bfseries Adressage IP, commutation et routage — BTS CIEL 1re année — durée : 2 heures\par}
  \ifcorrige\vspace{2pt}{\bfseries\color{corr}CORRIGÉ ET BARÈME — document enseignant}\par\fi
\end{center}
\vspace{-2pt}
{\renewcommand{\arraystretch}{1.35}
\begin{tabularx}{\linewidth}{|X|X|p{2.2cm}|p{2.6cm}|}
  \hline
  \textbf{Nom :} & \textbf{Prénom :} & \textbf{Poste n°} & \textbf{Note :} \hfill /20 \\ \hline
\end{tabularx}}

\textbf{Consignes :} travail \textbf{individuel}, documents et Internet \textbf{non autorisés}. Ce sujet est le
\textbf{document-réponse}. Fichier Packet Tracer : \ip{Nom\_Prenom\_EvalLAN\_\variante.pkt}, à enregistrer
régulièrement dans le dossier de rendu. Les « visas » sont apposés par le professeur, que vous appelez à chaque test réussi.
Durées conseillées : A 30 min — B 1 h 10 — C 20 min.

%==========================================================================
\section{Partie A — Questions de cours et calculs \hfill /6}

\question{A1}{2 pts}{Complétez le tableau en considérant le \textbf{masque par défaut} de la classe de chaque adresse.}
{\small\renewcommand{\arraystretch}{1.15}
\begin{tabularx}{\linewidth}{|l|c|C|C|C|c|}
  \hline
  \textbf{Adresse IP} & \textbf{Classe} & \textbf{Masque par défaut} & \textbf{Adresse réseau} &
  \textbf{Adresse de diffusion} & \textbf{Nb d'hôtes} \\ \hline
  \AIrows
\end{tabularx}}
\bareme{0,5 pt par ligne juste + 0,5 pt si les trois nombres d'hôtes sont justes ($2^{n}-2$, $n$ = nombre de bits d'hôte).}

\question{A2}{1 pt}{Indiquez le type de câble (\textbf{droit} ou \textbf{croisé}) pour chaque liaison.}
{\small\renewcommand{\arraystretch}{1.15}
\begin{tabularx}{\linewidth}{|X|c||X|c|}
  \hline
  \hl PC $\leftrightarrow$ commutateur & \makebox[1.6cm]{\sol{droit}} & Commutateur $\leftrightarrow$ routeur & \makebox[1.6cm]{\sol{droit}} \\ \hline
  \hl PC $\leftrightarrow$ PC (réseau P2P) & \sol{croisé} & Concentrateur $\leftrightarrow$ concentrateur & \sol{croisé} \\ \hline
\end{tabularx}}
\bareme{0,25 pt par case. Équipements de types différents : câble droit ; de même type : câble croisé.}

\question{A3}{1 pt}{Comment un \textbf{concentrateur (hub)} et un \textbf{commutateur (switch)} traitent-ils une trame
reçue ? Lequel limite les collisions, et pourquoi ?}
\zone{2.9cm}{Le \textbf{hub} (couche 1) répète la trame sur \textbf{tous ses autres ports} : un seul domaine de collision
partagé. Le \textbf{switch} (couche 2) lit l'adresse MAC de destination et, grâce à sa \textbf{table MAC}, ne l'envoie
\textbf{que sur le port du destinataire} : un domaine de collision par port, c'est lui qui limite les collisions.
\emph{0,5 pt hub, 0,5 pt switch.}}

\question{A4}{1 pt}{Qu'est-ce que la \textbf{passerelle par défaut} d'un poste ? Quelle condition son adresse doit-elle respecter ?}
\zone{2.4cm}{Adresse IP de l'interface du \textbf{routeur} à laquelle le poste envoie les paquets destinés à un
\textbf{autre réseau}. Elle doit \textbf{appartenir au même réseau} que le poste. \emph{0,5 + 0,5 pt.}}

\question{A5}{0,5 pt}{Un portable configuré en DHCP affiche \ip{\apipa} / \ip{255.255.0.0}. Que s'est-il passé ?}
\zone{1.5cm}{Aucun \textbf{serveur DHCP} n'a répondu : le poste s'est attribué une adresse \textbf{APIPA} (169.254.0.0/16).}

\question{A6}{0,5 pt}{Quel protocole utilise \texttt{ping} ? Nommez les deux messages échangés et la couche OSI des adresses IP.}
\zone{1.5cm}{\textbf{ICMP} ; \textbf{Echo Request} / \textbf{Echo Reply} ; adresses IP en \textbf{couche 3 (Réseau)}.}

%==========================================================================
\newpage
\section{Partie B — Réalisation sous Packet Tracer \hfill /10}

\begin{cdc}
\entreprise{} vous confie son réseau : \textbf{trois réseaux} reliés par un \textbf{routeur unique Cisco 2911}
nommé \ip{\Rnom} (interfaces G0/0, G0/1, G0/2).
\begin{itemize}[leftmargin=12pt]
  \item \textbf{Réseau \resUn} : \textbf{4 postes} PC-\pUn1 à PC-\pUn4, commutateur \ip{SW-\pUn}, adressage statique,
        réseau \ip{\netUn.0/24}.
  \item \textbf{Réseau \resDeux} : \textbf{2 postes} PC-\pDeux1, PC-\pDeux2 et \textbf{1 imprimante} IMP-\pDeux,
        commutateur \ip{SW-\pDeux}, adressage statique, réseau de classe B \ip{\netDeux.0.0/16}.
  \item \textbf{Réseau Serveur} : serveur \ip{\srvNom} d'adresse \ip{\netSrv.\srvH/24}, relié \textbf{directement} au routeur.
  \item La passerelle de chaque réseau est sa \textbf{dernière adresse utilisable}. Tous les équipements communiquent
        entre eux et accèdent à la page web du serveur. Nommage imposé (« Display Name »).
\end{itemize}
\end{cdc}

\question{B1}{2 pts}{Complétez le plan d'adressage \textbf{avant} la réalisation.}
{\footnotesize\setlength{\tabcolsep}{3pt}\renewcommand{\arraystretch}{1.15}
\begin{tabularx}{\linewidth}{|l|C|C|C|C|C|C|}
  \hline
  \textbf{Réseau} & \textbf{Adresse réseau} & \textbf{Masque} & \textbf{1re @ utilisable} &
  \textbf{Dernière @ utilisable} & \textbf{Diffusion} & \textbf{Passerelle} \\ \hline
  \hl \resUn & \sol{\netUn.0} & \sol{255.255.255.0} & \sol{\netUn.1} & \sol{\netUn.254} & \sol{\netUn.255} & \sol{\netUn.254} \\ \hline
  \hl \resDeux & \sol{\netDeux.0.0} & \sol{255.255.0.0} & \sol{\netDeux.0.1} & \sol{\netDeux.255.254} & \sol{\netDeux.255.255} & \sol{\netDeux.255.254} \\ \hline
  \hl Serveur & \sol{\netSrv.0} & \sol{255.255.255.0} & \sol{\netSrv.1} & \sol{\netSrv.254} & \sol{\netSrv.255} & \sol{\netSrv.254} \\ \hline
\end{tabularx}}
\bareme{0,5 pt par ligne juste + 0,5 pt si la convention de passerelle est respectée partout.}

\question{B2}{2 pts}{Construisez la topologie (équipements, câblage, nommage) et relevez l'adressage configuré.}
{\footnotesize\setlength{\tabcolsep}{3pt}\renewcommand{\arraystretch}{1.1}
\begin{tabularx}{\linewidth}{|l|C|C|C||l|C|C|}
  \hline
  \textbf{Hôte} & \textbf{Adresse IP} & \textbf{Masque} & \textbf{Passerelle} & \textbf{Interface} & \textbf{Adresse IP} & \textbf{Masque} \\ \hline
  \hl PC-\pUn1 & \sol{\netUn.1} & \sol{255.255.255.0} & \sol{\netUn.254} & \Rnom{} G0/0 & \sol{\netUn.254} & \sol{255.255.255.0} \\ \hline
  \hl PC-\pDeux1 & \sol{\netDeux.0.1} & \sol{255.255.0.0} & \sol{\netDeux.255.254} & \Rnom{} G0/1 & \sol{\netDeux.255.254} & \sol{255.255.0.0} \\ \hline
  \hl IMP-\pDeux & \sol{\netDeux.0.10} & \sol{255.255.0.0} & \sol{\netDeux.255.254} & \Rnom{} G0/2 & \sol{\netSrv.254} & \sol{255.255.255.0} \\ \hline
  \hl \srvNom & \ip{\netSrv.\srvH} & \sol{255.255.255.0} & \sol{\netSrv.254} & \multicolumn{3}{c|}{\cellcolor{black!8}} \\ \hline
\end{tabularx}}
\bareme{Adresses d'hôtes libres (uniques, dans le bon réseau) ; affectation G0/x libre. Sur le .pkt : équipements 0,5 ;
câblage 1 (câbles droits vers les switchs et switch–routeur, \textbf{croisé} serveur–routeur, liens verts) ; nommage 0,5.}

\question{B3}{3 pts}{Configurez tous les équipements : IP, masques, passerelles des hôtes ; interfaces du routeur
(« Port Status : On »).}
\bareme{Sur le .pkt : interfaces routeur 1 pt ; IP/masques des hôtes 1 pt ; passerelles (imprimante et serveur compris)
1 pt. $-0,25$ par erreur, sans descendre sous 0 par item.}

\question{B4}{3 pts}{Réalisez les tests, notez la commande et le résultat, puis faites viser chaque test réussi.}
{\small\renewcommand{\arraystretch}{1.1}
\begin{tabularx}{\linewidth}{|c|l|l|C|c|c|}
  \hline
  \textbf{N°} & \textbf{Source} & \textbf{Destination} & \textbf{Commande / action} & \textbf{Résultat} & \textbf{Visa prof.} \\ \hline
  \hl 1 & PC-\pUn1 & PC-\pUn4 & \sol{ping \netUn.4} & \sol{réussi} & \hspace{1.4cm} \\ \hline
  \hl 2 & PC-\pDeux1 & IMP-\pDeux & \sol{ping \netDeux.0.10} & \sol{réussi} & \\ \hline
  \hl 3 & PC-\pUn2 & PC-\pDeux2 & \sol{ping \netDeux.0.2} & \sol{réussi} & \\ \hline
  \hl 4 & PC-\pDeux2 & \srvNom & \sol{ping \netSrv.\srvH} & \sol{réussi} & \\ \hline
  \hl 5 & PC-\pUn3 & page web du serveur & \sol{navigateur : {\NoAutoSpacing http://\netSrv.\srvH}} & \sol{page affichée} & \\ \hline
\end{tabularx}}
\bareme{0,5 pt par test visé + 0,5 pt pour B5. Le premier ping inter-réseaux peut échouer (résolution ARP) ; les suivants doivent réussir.}

\question{B5}{incl. B4}{Test n°3 : par quels équipements passe le paquet ? Pourquoi PC-\pUn2 l'envoie-t-il au routeur ?}
\zone{2.3cm}{PC-\pUn2 $\to$ SW-\pUn{} $\to$ \Rnom{} $\to$ SW-\pDeux{} $\to$ PC-\pDeux2. Grâce à son masque, PC-\pUn2 constate que
la destination (\netDeux.0.0/16) n'est \textbf{pas dans son réseau} (\netUn.0/24) : il envoie le paquet à sa
\textbf{passerelle par défaut}, qui le route.}

%==========================================================================
\newpage
\section{Partie C — Diagnostic d'un réseau en panne \hfill /4}

Un technicien constate qu'\textbf{aucun poste ne parvient à joindre le serveur SRV-C} et que certains postes
communiquent mal entre eux. Les câbles sont du bon type.

\begin{center}
\begin{minipage}{0.37\linewidth}\centering
\begin{tikzpicture}[font=\footnotesize,
  eq/.style={draw,fill=qbg,rounded corners=2pt,minimum width=9mm,minimum height=5mm,inner sep=2pt},
  every path/.style={thick}]
  \node[eq] (SW) {SW-1};
  \node[eq,right=9mm of SW] (R) {R1};
  \node[eq,right=7mm of R] (SRV) {SRV-C};
  \node[eq,above left=5mm and -3mm of SW] (A) {PC-A};
  \node[eq,above right=5mm and -3mm of SW] (B) {PC-B};
  \node[eq,below left=5mm and -3mm of SW] (C) {PC-C};
  \node[eq,below right=5mm and -3mm of SW] (D) {PC-D};
  \foreach \n in {A,B,C,D} \draw (\n) -- (SW);
  \draw (SW) -- node[above,font=\tiny]{G0/0} (R);
  \draw (R) -- node[above,font=\tiny]{G0/1} (SRV);
\end{tikzpicture}
\end{minipage}\hfill
\begin{minipage}{0.61\linewidth}\centering
{\small\renewcommand{\arraystretch}{1.05}\setlength{\tabcolsep}{4pt}
\begin{tabular}{lllll}
  \toprule
  \textbf{Équipement} & \textbf{Adresse IP} & \textbf{Masque} & \textbf{Passerelle} & \textbf{Port} \\
  \midrule
  \Crows
  \bottomrule
\end{tabular}}
\end{minipage}
\end{center}

\question{C1}{4 pts}{Relevez les \textbf{quatre anomalies}. Pour chacune : conséquence et correction.}
{\small\renewcommand{\arraystretch}{1.1}
\ifcorrige\def\hc{}\else\def\hc{\rule[-10mm]{0pt}{14mm}}\fi
\begin{tabularx}{\linewidth}{|c|>{\raggedright\arraybackslash}p{2cm}|L|L|L|}
  \hline
  \textbf{N°} & \textbf{Équipement} & \textbf{Anomalie} & \textbf{Conséquence} & \textbf{Correction} \\ \hline
  \ifcorrige\Csol\else
  \hc 1 & & & & \\ \hline
  \hc 2 & & & & \\ \hline
  \hc 3 & & & & \\ \hline
  \hc 4 & & & & \\ \hline
  \fi
\end{tabularx}}
\bareme{1 pt par anomalie (0,25 équipement, anomalie, conséquence, correction ; ordre indifférent). \Cbareme}

%==========================================================================
\section{Bonus — Extension sans fil \hfill +1}

\begin{cdc}Ajoutez au réseau \resUn{} un \textbf{AccessPoint-PT} relié à SW-\pUn{} (SSID \ip{\ssid}, \textbf{WPA2-PSK})
et un portable \ip{Laptop-\pUn} équipé d'un module \texttt{WPC300N}, en adressage \textbf{statique}.
Vérifiez qu'il joint \srvNom.\end{cdc}
{\small\renewcommand{\arraystretch}{1.1}
\begin{tabularx}{\linewidth}{|l|C|C|C|c|}
  \hline
  \textbf{Équipement} & \textbf{Adresse IP} & \textbf{Masque} & \textbf{Passerelle} & \textbf{Visa prof.} \\ \hline
  \hl Laptop-\pUn & \sol{\netUn.50} & \sol{255.255.255.0} & \sol{\netUn.254} & \hspace{1.4cm} \\ \hline
\end{tabularx}}
\bareme{Le point d'accès fait office de pont : le portable est dans \netUn.0/24. Bonus si associé en WPA2-PSK et ping
vers \netSrv.\srvH{} réussi. Note plafonnée à 20.}

\end{document}
"
%   pdflatex -jobname=Eval_Reseau_SujetB   "\def\variante{B}\def\modecorrige{0}% Évaluation pratique — Réseaux locaux (BTS CIEL 1re année) — sujets A et B
% Compilation (4 PDF) :
%   pdflatex -jobname=Eval_Reseau_SujetA   "\def\variante{A}\def\modecorrige{0}% Évaluation pratique — Réseaux locaux (BTS CIEL 1re année) — sujets A et B
% Compilation (4 PDF) :
%   pdflatex -jobname=Eval_Reseau_SujetA   "\def\variante{A}\def\modecorrige{0}% Évaluation pratique — Réseaux locaux (BTS CIEL 1re année) — sujets A et B
% Compilation (4 PDF) :
%   pdflatex -jobname=Eval_Reseau_SujetA   "\def\variante{A}\def\modecorrige{0}\input{Eval_TP_Reseau_LAN.tex}"
%   pdflatex -jobname=Eval_Reseau_SujetB   "\def\variante{B}\def\modecorrige{0}\input{Eval_TP_Reseau_LAN.tex}"
%   pdflatex -jobname=Eval_Reseau_CorrigeA "\def\variante{A}\def\modecorrige{1}\input{Eval_TP_Reseau_LAN.tex}"
%   pdflatex -jobname=Eval_Reseau_CorrigeB "\def\variante{B}\def\modecorrige{1}\input{Eval_TP_Reseau_LAN.tex}"
\documentclass[10pt,a4paper]{article}

\usepackage[T1]{fontenc}
\usepackage[utf8]{inputenc}
\usepackage{lmodern}
\usepackage[french]{babel}
\usepackage[a4paper,hmargin=1.5cm,top=1.3cm,bottom=1.6cm,footskip=0.7cm]{geometry}
\usepackage[table]{xcolor}
\usepackage{booktabs,tabularx,array}
\usepackage{enumitem}
\usepackage{fancyhdr}
\usepackage[most]{tcolorbox}
\usepackage{tikz}
\usetikzlibrary{positioning}
\usepackage{titlesec}

\providecommand{\modecorrige}{0}
\providecommand{\variante}{A}
\newif\ifcorrige
\ifnum\modecorrige=1 \corrigetrue \else \corrigefalse \fi

\definecolor{ipred}{HTML}{C0392B}
\definecolor{cdcbar}{HTML}{A67C3D}
\definecolor{cdcbg}{HTML}{F4F1EA}
\definecolor{qbar}{HTML}{4A6B5D}
\definecolor{qbg}{HTML}{EDF1EF}
\definecolor{corr}{HTML}{1F5FA8}
\definecolor{corrbg}{HTML}{E8F0FA}
\definecolor{grisfoot}{HTML}{828282}

\newcommand{\ip}[1]{\textcolor{ipred}{\texttt{#1}}}
\newcommand{\pts}[1]{\hfill\mbox{\textbf{(#1)}}}
\newcommand{\sol}[1]{\ifcorrige\textcolor{corr}{#1}\fi}
\newcommand{\bareme}[1]{\ifcorrige\par{\footnotesize\color{corr}#1\par}\fi}
\newcommand{\hl}{\ifcorrige\rule[-1.6mm]{0pt}{5.4mm}\else\rule[-2.2mm]{0pt}{6.6mm}\fi}
\newcolumntype{C}{>{\centering\arraybackslash}X}
\newcolumntype{L}{>{\raggedright\arraybackslash}X}

\titleformat{\section}{\large\bfseries}{}{0pt}{}
\titlespacing*{\section}{0pt}{6pt}{3pt}

\tcbset{qstyle/.style={enhanced,boxrule=0pt,leftrule=2.5pt,arc=0pt,outer arc=0pt,
  left=5pt,right=5pt,top=2pt,bottom=2pt,before skip=4pt,after skip=3pt}}
\newtcolorbox{cdc}{qstyle,colback=cdcbg,colframe=cdcbar}
\newtcolorbox{quest}{qstyle,colback=qbg,colframe=qbar}

% Zone de réponse : vierge dans le sujet, remplie dans le corrigé
\newcommand{\zone}[2]{%
  \ifcorrige
    \begin{tcolorbox}[colback=corrbg,colframe=corr,boxrule=0.5pt,arc=1pt,left=4pt,right=4pt,top=1pt,bottom=1pt,
      before skip=2pt,after skip=2pt,fontupper=\footnotesize\color{corr}]#2\end{tcolorbox}%
  \else
    \begin{tcolorbox}[colback=white,colframe=black!45,boxrule=0.4pt,arc=1pt,height=#1,
      before skip=2pt,after skip=2pt]\end{tcolorbox}%
  \fi}

\newcommand{\question}[3]{\begin{quest}\textbf{#1.} #3\pts{#2}\end{quest}}

%==========================================================================
% Données propres à chaque variante
%==========================================================================
\if\variante A
  \def\entreprise{Le cabinet d'architectes \textbf{ArchiNova}}
  \def\Rnom{R-ArchiNova}
  \def\resUn{Bureau d'études}\def\pUn{BE}\def\netUn{192.168.40}
  \def\resDeux{Accueil}\def\pDeux{ACC}\def\netDeux{172.16}
  \def\srvNom{SRV-Intranet}\def\netSrv{192.168.200}\def\srvH{10}
  \def\ssid{ArchiNova-WiFi}
  \def\AIrows{%
    \hl\ip{192.168.40.25} & \sol{C} & \sol{255.255.255.0} & \sol{192.168.40.0} & \sol{192.168.40.255} & \sol{254} \\ \hline
    \hl\ip{172.16.5.10}   & \sol{B} & \sol{255.255.0.0}   & \sol{172.16.0.0}   & \sol{172.16.255.255} & \sol{65\,534} \\ \hline
    \hl\ip{10.3.2.1}      & \sol{A} & \sol{255.0.0.0}     & \sol{10.0.0.0}     & \sol{10.255.255.255} & \sol{16\,777\,214} \\ \hline}
  \def\apipa{169.254.12.7}
  \def\Cnet{192.168.60}
  \def\Crows{%
    PC-A & \ip{192.168.60.10} & \ip{255.255.255.0} & \ip{192.168.60.254} & — \\
    PC-B & \ip{192.168.60.20} & \ip{255.255.255.0} & \ip{192.168.60.1}   & — \\
    PC-C & \ip{192.168.61.30} & \ip{255.255.255.0} & \ip{192.168.60.254} & — \\
    PC-D & \ip{192.168.60.10} & \ip{255.255.255.0} & \ip{192.168.60.254} & — \\
    SRV-C & \ip{10.0.0.5}     & \ip{255.0.0.0}     & \ip{10.0.0.254}     & — \\
    R1 — G0/0 & \ip{192.168.60.254} & \ip{255.255.255.0} & — & On \\
    R1 — G0/1 & \ip{10.0.0.254}     & \ip{255.0.0.0}     & — & Off \\}
  \def\Csol{%
    \hl 1 & \sol{PC-B} & \sol{Passerelle .60.1 : pas l'adresse du routeur.} & \sol{Joint son réseau, mais aucun autre (serveur).} & \sol{Passerelle 192.168.60.254.} \\ \hline
    \hl 2 & \sol{PC-C} & \sol{IP hors du réseau 192.168.60.0/24.} & \sol{Ne joint ni les postes ni le routeur.} & \sol{Ex. 192.168.60.30.} \\ \hline
    \hl 3 & \sol{PC-D} & \sol{Même IP que PC-A (doublon).} & \sol{Conflit d'adresse : PC-A et PC-D perturbés.} & \sol{Adresse unique, ex. .60.40.} \\ \hline
    \hl 4 & \sol{R1 G0/1} & \sol{Interface désactivée (Off).} & \sol{Réseau du serveur injoignable pour tous.} & \sol{Port Status : On.} \\ \hline}
  \def\Cbareme{PC-A et SRV-C sont correctement configurés ; accepter PC-A cité à la place de PC-D pour le doublon.}
\else
  \def\entreprise{La clinique vétérinaire \textbf{VetoPlus}}
  \def\Rnom{R-VetoPlus}
  \def\resUn{Soins}\def\pUn{SOI}\def\netUn{192.168.70}
  \def\resDeux{Secrétariat}\def\pDeux{SEC}\def\netDeux{172.20}
  \def\srvNom{SRV-Dossiers}\def\netSrv{192.168.250}\def\srvH{20}
  \def\ssid{VetoPlus-WiFi}
  \def\AIrows{%
    \hl\ip{192.168.15.130} & \sol{C} & \sol{255.255.255.0} & \sol{192.168.15.0} & \sol{192.168.15.255} & \sol{254} \\ \hline
    \hl\ip{172.31.200.7}   & \sol{B} & \sol{255.255.0.0}   & \sol{172.31.0.0}   & \sol{172.31.255.255} & \sol{65\,534} \\ \hline
    \hl\ip{10.45.6.9}      & \sol{A} & \sol{255.0.0.0}     & \sol{10.0.0.0}     & \sol{10.255.255.255} & \sol{16\,777\,214} \\ \hline}
  \def\apipa{169.254.88.3}
  \def\Cnet{192.168.80}
  \def\Crows{%
    PC-A & \ip{192.168.8.11}   & \ip{255.255.255.0} & \ip{192.168.80.254} & — \\
    PC-B & \ip{192.168.80.12}  & \ip{255.255.255.0} & \ip{192.168.80.254} & — \\
    PC-C & \ip{192.168.80.13}  & \ip{255.255.255.0} & \ip{192.168.80.100} & — \\
    PC-D & \ip{192.168.80.12}  & \ip{255.255.255.0} & \ip{192.168.80.254} & — \\
    SRV-C & \ip{10.0.0.8}      & \ip{255.0.0.0}     & \ip{10.0.0.254}     & — \\
    R1 — G0/0 & \ip{192.168.80.254} & \ip{255.255.255.0} & — & Off \\
    R1 — G0/1 & \ip{10.0.0.254}     & \ip{255.0.0.0}     & — & On \\}
  \def\Csol{%
    \hl 1 & \sol{PC-A} & \sol{IP hors du réseau 192.168.80.0/24.} & \sol{Ne joint ni les postes ni le routeur.} & \sol{Ex. 192.168.80.11.} \\ \hline
    \hl 2 & \sol{PC-C} & \sol{Passerelle .80.100 : pas l'adresse du routeur.} & \sol{Joint son réseau, mais aucun autre (serveur).} & \sol{Passerelle 192.168.80.254.} \\ \hline
    \hl 3 & \sol{PC-D} & \sol{Même IP que PC-B (doublon).} & \sol{Conflit d'adresse : PC-B et PC-D perturbés.} & \sol{Adresse unique, ex. .80.14.} \\ \hline
    \hl 4 & \sol{R1 G0/0} & \sol{Interface côté LAN désactivée (Off).} & \sol{Aucun poste ne joint le routeur ni le serveur.} & \sol{Port Status : On.} \\ \hline}
  \def\Cbareme{PC-B et SRV-C sont correctement configurés ; accepter PC-B cité à la place de PC-D pour le doublon.}
\fi

\pagestyle{fancy}
\fancyhf{}
\renewcommand{\headrulewidth}{0pt}
\fancyfoot[L]{\footnotesize\itshape\color{grisfoot}Évaluation — Réseaux locaux — Sujet \variante\ifcorrige\ — CORRIGÉ\fi}
\fancyfoot[R]{\footnotesize\color{grisfoot}\thepage/3}

\setlist{itemsep=0pt,topsep=2pt,parsep=0pt}
\setlength{\parindent}{0pt}
\setlength{\parskip}{2pt}

\begin{document}

\begin{center}
  {\Large\bfseries Évaluation pratique — Réseaux locaux \quad
   \fcolorbox{black}{white}{\ SUJET \variante\ }\par}
  \vspace{1pt}
  {\bfseries Adressage IP, commutation et routage — BTS CIEL 1re année — durée : 2 heures\par}
  \ifcorrige\vspace{2pt}{\bfseries\color{corr}CORRIGÉ ET BARÈME — document enseignant}\par\fi
\end{center}
\vspace{-2pt}
{\renewcommand{\arraystretch}{1.35}
\begin{tabularx}{\linewidth}{|X|X|p{2.2cm}|p{2.6cm}|}
  \hline
  \textbf{Nom :} & \textbf{Prénom :} & \textbf{Poste n°} & \textbf{Note :} \hfill /20 \\ \hline
\end{tabularx}}

\textbf{Consignes :} travail \textbf{individuel}, documents et Internet \textbf{non autorisés}. Ce sujet est le
\textbf{document-réponse}. Fichier Packet Tracer : \ip{Nom\_Prenom\_EvalLAN\_\variante.pkt}, à enregistrer
régulièrement dans le dossier de rendu. Les « visas » sont apposés par le professeur, que vous appelez à chaque test réussi.
Durées conseillées : A 30 min — B 1 h 10 — C 20 min.

%==========================================================================
\section{Partie A — Questions de cours et calculs \hfill /6}

\question{A1}{2 pts}{Complétez le tableau en considérant le \textbf{masque par défaut} de la classe de chaque adresse.}
{\small\renewcommand{\arraystretch}{1.15}
\begin{tabularx}{\linewidth}{|l|c|C|C|C|c|}
  \hline
  \textbf{Adresse IP} & \textbf{Classe} & \textbf{Masque par défaut} & \textbf{Adresse réseau} &
  \textbf{Adresse de diffusion} & \textbf{Nb d'hôtes} \\ \hline
  \AIrows
\end{tabularx}}
\bareme{0,5 pt par ligne juste + 0,5 pt si les trois nombres d'hôtes sont justes ($2^{n}-2$, $n$ = nombre de bits d'hôte).}

\question{A2}{1 pt}{Indiquez le type de câble (\textbf{droit} ou \textbf{croisé}) pour chaque liaison.}
{\small\renewcommand{\arraystretch}{1.15}
\begin{tabularx}{\linewidth}{|X|c||X|c|}
  \hline
  \hl PC $\leftrightarrow$ commutateur & \makebox[1.6cm]{\sol{droit}} & Commutateur $\leftrightarrow$ routeur & \makebox[1.6cm]{\sol{droit}} \\ \hline
  \hl PC $\leftrightarrow$ PC (réseau P2P) & \sol{croisé} & Concentrateur $\leftrightarrow$ concentrateur & \sol{croisé} \\ \hline
\end{tabularx}}
\bareme{0,25 pt par case. Équipements de types différents : câble droit ; de même type : câble croisé.}

\question{A3}{1 pt}{Comment un \textbf{concentrateur (hub)} et un \textbf{commutateur (switch)} traitent-ils une trame
reçue ? Lequel limite les collisions, et pourquoi ?}
\zone{2.9cm}{Le \textbf{hub} (couche 1) répète la trame sur \textbf{tous ses autres ports} : un seul domaine de collision
partagé. Le \textbf{switch} (couche 2) lit l'adresse MAC de destination et, grâce à sa \textbf{table MAC}, ne l'envoie
\textbf{que sur le port du destinataire} : un domaine de collision par port, c'est lui qui limite les collisions.
\emph{0,5 pt hub, 0,5 pt switch.}}

\question{A4}{1 pt}{Qu'est-ce que la \textbf{passerelle par défaut} d'un poste ? Quelle condition son adresse doit-elle respecter ?}
\zone{2.4cm}{Adresse IP de l'interface du \textbf{routeur} à laquelle le poste envoie les paquets destinés à un
\textbf{autre réseau}. Elle doit \textbf{appartenir au même réseau} que le poste. \emph{0,5 + 0,5 pt.}}

\question{A5}{0,5 pt}{Un portable configuré en DHCP affiche \ip{\apipa} / \ip{255.255.0.0}. Que s'est-il passé ?}
\zone{1.5cm}{Aucun \textbf{serveur DHCP} n'a répondu : le poste s'est attribué une adresse \textbf{APIPA} (169.254.0.0/16).}

\question{A6}{0,5 pt}{Quel protocole utilise \texttt{ping} ? Nommez les deux messages échangés et la couche OSI des adresses IP.}
\zone{1.5cm}{\textbf{ICMP} ; \textbf{Echo Request} / \textbf{Echo Reply} ; adresses IP en \textbf{couche 3 (Réseau)}.}

%==========================================================================
\newpage
\section{Partie B — Réalisation sous Packet Tracer \hfill /10}

\begin{cdc}
\entreprise{} vous confie son réseau : \textbf{trois réseaux} reliés par un \textbf{routeur unique Cisco 2911}
nommé \ip{\Rnom} (interfaces G0/0, G0/1, G0/2).
\begin{itemize}[leftmargin=12pt]
  \item \textbf{Réseau \resUn} : \textbf{4 postes} PC-\pUn1 à PC-\pUn4, commutateur \ip{SW-\pUn}, adressage statique,
        réseau \ip{\netUn.0/24}.
  \item \textbf{Réseau \resDeux} : \textbf{2 postes} PC-\pDeux1, PC-\pDeux2 et \textbf{1 imprimante} IMP-\pDeux,
        commutateur \ip{SW-\pDeux}, adressage statique, réseau de classe B \ip{\netDeux.0.0/16}.
  \item \textbf{Réseau Serveur} : serveur \ip{\srvNom} d'adresse \ip{\netSrv.\srvH/24}, relié \textbf{directement} au routeur.
  \item La passerelle de chaque réseau est sa \textbf{dernière adresse utilisable}. Tous les équipements communiquent
        entre eux et accèdent à la page web du serveur. Nommage imposé (« Display Name »).
\end{itemize}
\end{cdc}

\question{B1}{2 pts}{Complétez le plan d'adressage \textbf{avant} la réalisation.}
{\footnotesize\setlength{\tabcolsep}{3pt}\renewcommand{\arraystretch}{1.15}
\begin{tabularx}{\linewidth}{|l|C|C|C|C|C|C|}
  \hline
  \textbf{Réseau} & \textbf{Adresse réseau} & \textbf{Masque} & \textbf{1re @ utilisable} &
  \textbf{Dernière @ utilisable} & \textbf{Diffusion} & \textbf{Passerelle} \\ \hline
  \hl \resUn & \sol{\netUn.0} & \sol{255.255.255.0} & \sol{\netUn.1} & \sol{\netUn.254} & \sol{\netUn.255} & \sol{\netUn.254} \\ \hline
  \hl \resDeux & \sol{\netDeux.0.0} & \sol{255.255.0.0} & \sol{\netDeux.0.1} & \sol{\netDeux.255.254} & \sol{\netDeux.255.255} & \sol{\netDeux.255.254} \\ \hline
  \hl Serveur & \sol{\netSrv.0} & \sol{255.255.255.0} & \sol{\netSrv.1} & \sol{\netSrv.254} & \sol{\netSrv.255} & \sol{\netSrv.254} \\ \hline
\end{tabularx}}
\bareme{0,5 pt par ligne juste + 0,5 pt si la convention de passerelle est respectée partout.}

\question{B2}{2 pts}{Construisez la topologie (équipements, câblage, nommage) et relevez l'adressage configuré.}
{\footnotesize\setlength{\tabcolsep}{3pt}\renewcommand{\arraystretch}{1.1}
\begin{tabularx}{\linewidth}{|l|C|C|C||l|C|C|}
  \hline
  \textbf{Hôte} & \textbf{Adresse IP} & \textbf{Masque} & \textbf{Passerelle} & \textbf{Interface} & \textbf{Adresse IP} & \textbf{Masque} \\ \hline
  \hl PC-\pUn1 & \sol{\netUn.1} & \sol{255.255.255.0} & \sol{\netUn.254} & \Rnom{} G0/0 & \sol{\netUn.254} & \sol{255.255.255.0} \\ \hline
  \hl PC-\pDeux1 & \sol{\netDeux.0.1} & \sol{255.255.0.0} & \sol{\netDeux.255.254} & \Rnom{} G0/1 & \sol{\netDeux.255.254} & \sol{255.255.0.0} \\ \hline
  \hl IMP-\pDeux & \sol{\netDeux.0.10} & \sol{255.255.0.0} & \sol{\netDeux.255.254} & \Rnom{} G0/2 & \sol{\netSrv.254} & \sol{255.255.255.0} \\ \hline
  \hl \srvNom & \ip{\netSrv.\srvH} & \sol{255.255.255.0} & \sol{\netSrv.254} & \multicolumn{3}{c|}{\cellcolor{black!8}} \\ \hline
\end{tabularx}}
\bareme{Adresses d'hôtes libres (uniques, dans le bon réseau) ; affectation G0/x libre. Sur le .pkt : équipements 0,5 ;
câblage 1 (câbles droits vers les switchs et switch–routeur, \textbf{croisé} serveur–routeur, liens verts) ; nommage 0,5.}

\question{B3}{3 pts}{Configurez tous les équipements : IP, masques, passerelles des hôtes ; interfaces du routeur
(« Port Status : On »).}
\bareme{Sur le .pkt : interfaces routeur 1 pt ; IP/masques des hôtes 1 pt ; passerelles (imprimante et serveur compris)
1 pt. $-0,25$ par erreur, sans descendre sous 0 par item.}

\question{B4}{3 pts}{Réalisez les tests, notez la commande et le résultat, puis faites viser chaque test réussi.}
{\small\renewcommand{\arraystretch}{1.1}
\begin{tabularx}{\linewidth}{|c|l|l|C|c|c|}
  \hline
  \textbf{N°} & \textbf{Source} & \textbf{Destination} & \textbf{Commande / action} & \textbf{Résultat} & \textbf{Visa prof.} \\ \hline
  \hl 1 & PC-\pUn1 & PC-\pUn4 & \sol{ping \netUn.4} & \sol{réussi} & \hspace{1.4cm} \\ \hline
  \hl 2 & PC-\pDeux1 & IMP-\pDeux & \sol{ping \netDeux.0.10} & \sol{réussi} & \\ \hline
  \hl 3 & PC-\pUn2 & PC-\pDeux2 & \sol{ping \netDeux.0.2} & \sol{réussi} & \\ \hline
  \hl 4 & PC-\pDeux2 & \srvNom & \sol{ping \netSrv.\srvH} & \sol{réussi} & \\ \hline
  \hl 5 & PC-\pUn3 & page web du serveur & \sol{navigateur : {\NoAutoSpacing http://\netSrv.\srvH}} & \sol{page affichée} & \\ \hline
\end{tabularx}}
\bareme{0,5 pt par test visé + 0,5 pt pour B5. Le premier ping inter-réseaux peut échouer (résolution ARP) ; les suivants doivent réussir.}

\question{B5}{incl. B4}{Test n°3 : par quels équipements passe le paquet ? Pourquoi PC-\pUn2 l'envoie-t-il au routeur ?}
\zone{2.3cm}{PC-\pUn2 $\to$ SW-\pUn{} $\to$ \Rnom{} $\to$ SW-\pDeux{} $\to$ PC-\pDeux2. Grâce à son masque, PC-\pUn2 constate que
la destination (\netDeux.0.0/16) n'est \textbf{pas dans son réseau} (\netUn.0/24) : il envoie le paquet à sa
\textbf{passerelle par défaut}, qui le route.}

%==========================================================================
\newpage
\section{Partie C — Diagnostic d'un réseau en panne \hfill /4}

Un technicien constate qu'\textbf{aucun poste ne parvient à joindre le serveur SRV-C} et que certains postes
communiquent mal entre eux. Les câbles sont du bon type.

\begin{center}
\begin{minipage}{0.37\linewidth}\centering
\begin{tikzpicture}[font=\footnotesize,
  eq/.style={draw,fill=qbg,rounded corners=2pt,minimum width=9mm,minimum height=5mm,inner sep=2pt},
  every path/.style={thick}]
  \node[eq] (SW) {SW-1};
  \node[eq,right=9mm of SW] (R) {R1};
  \node[eq,right=7mm of R] (SRV) {SRV-C};
  \node[eq,above left=5mm and -3mm of SW] (A) {PC-A};
  \node[eq,above right=5mm and -3mm of SW] (B) {PC-B};
  \node[eq,below left=5mm and -3mm of SW] (C) {PC-C};
  \node[eq,below right=5mm and -3mm of SW] (D) {PC-D};
  \foreach \n in {A,B,C,D} \draw (\n) -- (SW);
  \draw (SW) -- node[above,font=\tiny]{G0/0} (R);
  \draw (R) -- node[above,font=\tiny]{G0/1} (SRV);
\end{tikzpicture}
\end{minipage}\hfill
\begin{minipage}{0.61\linewidth}\centering
{\small\renewcommand{\arraystretch}{1.05}\setlength{\tabcolsep}{4pt}
\begin{tabular}{lllll}
  \toprule
  \textbf{Équipement} & \textbf{Adresse IP} & \textbf{Masque} & \textbf{Passerelle} & \textbf{Port} \\
  \midrule
  \Crows
  \bottomrule
\end{tabular}}
\end{minipage}
\end{center}

\question{C1}{4 pts}{Relevez les \textbf{quatre anomalies}. Pour chacune : conséquence et correction.}
{\small\renewcommand{\arraystretch}{1.1}
\ifcorrige\def\hc{}\else\def\hc{\rule[-10mm]{0pt}{14mm}}\fi
\begin{tabularx}{\linewidth}{|c|>{\raggedright\arraybackslash}p{2cm}|L|L|L|}
  \hline
  \textbf{N°} & \textbf{Équipement} & \textbf{Anomalie} & \textbf{Conséquence} & \textbf{Correction} \\ \hline
  \ifcorrige\Csol\else
  \hc 1 & & & & \\ \hline
  \hc 2 & & & & \\ \hline
  \hc 3 & & & & \\ \hline
  \hc 4 & & & & \\ \hline
  \fi
\end{tabularx}}
\bareme{1 pt par anomalie (0,25 équipement, anomalie, conséquence, correction ; ordre indifférent). \Cbareme}

%==========================================================================
\section{Bonus — Extension sans fil \hfill +1}

\begin{cdc}Ajoutez au réseau \resUn{} un \textbf{AccessPoint-PT} relié à SW-\pUn{} (SSID \ip{\ssid}, \textbf{WPA2-PSK})
et un portable \ip{Laptop-\pUn} équipé d'un module \texttt{WPC300N}, en adressage \textbf{statique}.
Vérifiez qu'il joint \srvNom.\end{cdc}
{\small\renewcommand{\arraystretch}{1.1}
\begin{tabularx}{\linewidth}{|l|C|C|C|c|}
  \hline
  \textbf{Équipement} & \textbf{Adresse IP} & \textbf{Masque} & \textbf{Passerelle} & \textbf{Visa prof.} \\ \hline
  \hl Laptop-\pUn & \sol{\netUn.50} & \sol{255.255.255.0} & \sol{\netUn.254} & \hspace{1.4cm} \\ \hline
\end{tabularx}}
\bareme{Le point d'accès fait office de pont : le portable est dans \netUn.0/24. Bonus si associé en WPA2-PSK et ping
vers \netSrv.\srvH{} réussi. Note plafonnée à 20.}

\end{document}
"
%   pdflatex -jobname=Eval_Reseau_SujetB   "\def\variante{B}\def\modecorrige{0}% Évaluation pratique — Réseaux locaux (BTS CIEL 1re année) — sujets A et B
% Compilation (4 PDF) :
%   pdflatex -jobname=Eval_Reseau_SujetA   "\def\variante{A}\def\modecorrige{0}\input{Eval_TP_Reseau_LAN.tex}"
%   pdflatex -jobname=Eval_Reseau_SujetB   "\def\variante{B}\def\modecorrige{0}\input{Eval_TP_Reseau_LAN.tex}"
%   pdflatex -jobname=Eval_Reseau_CorrigeA "\def\variante{A}\def\modecorrige{1}\input{Eval_TP_Reseau_LAN.tex}"
%   pdflatex -jobname=Eval_Reseau_CorrigeB "\def\variante{B}\def\modecorrige{1}\input{Eval_TP_Reseau_LAN.tex}"
\documentclass[10pt,a4paper]{article}

\usepackage[T1]{fontenc}
\usepackage[utf8]{inputenc}
\usepackage{lmodern}
\usepackage[french]{babel}
\usepackage[a4paper,hmargin=1.5cm,top=1.3cm,bottom=1.6cm,footskip=0.7cm]{geometry}
\usepackage[table]{xcolor}
\usepackage{booktabs,tabularx,array}
\usepackage{enumitem}
\usepackage{fancyhdr}
\usepackage[most]{tcolorbox}
\usepackage{tikz}
\usetikzlibrary{positioning}
\usepackage{titlesec}

\providecommand{\modecorrige}{0}
\providecommand{\variante}{A}
\newif\ifcorrige
\ifnum\modecorrige=1 \corrigetrue \else \corrigefalse \fi

\definecolor{ipred}{HTML}{C0392B}
\definecolor{cdcbar}{HTML}{A67C3D}
\definecolor{cdcbg}{HTML}{F4F1EA}
\definecolor{qbar}{HTML}{4A6B5D}
\definecolor{qbg}{HTML}{EDF1EF}
\definecolor{corr}{HTML}{1F5FA8}
\definecolor{corrbg}{HTML}{E8F0FA}
\definecolor{grisfoot}{HTML}{828282}

\newcommand{\ip}[1]{\textcolor{ipred}{\texttt{#1}}}
\newcommand{\pts}[1]{\hfill\mbox{\textbf{(#1)}}}
\newcommand{\sol}[1]{\ifcorrige\textcolor{corr}{#1}\fi}
\newcommand{\bareme}[1]{\ifcorrige\par{\footnotesize\color{corr}#1\par}\fi}
\newcommand{\hl}{\ifcorrige\rule[-1.6mm]{0pt}{5.4mm}\else\rule[-2.2mm]{0pt}{6.6mm}\fi}
\newcolumntype{C}{>{\centering\arraybackslash}X}
\newcolumntype{L}{>{\raggedright\arraybackslash}X}

\titleformat{\section}{\large\bfseries}{}{0pt}{}
\titlespacing*{\section}{0pt}{6pt}{3pt}

\tcbset{qstyle/.style={enhanced,boxrule=0pt,leftrule=2.5pt,arc=0pt,outer arc=0pt,
  left=5pt,right=5pt,top=2pt,bottom=2pt,before skip=4pt,after skip=3pt}}
\newtcolorbox{cdc}{qstyle,colback=cdcbg,colframe=cdcbar}
\newtcolorbox{quest}{qstyle,colback=qbg,colframe=qbar}

% Zone de réponse : vierge dans le sujet, remplie dans le corrigé
\newcommand{\zone}[2]{%
  \ifcorrige
    \begin{tcolorbox}[colback=corrbg,colframe=corr,boxrule=0.5pt,arc=1pt,left=4pt,right=4pt,top=1pt,bottom=1pt,
      before skip=2pt,after skip=2pt,fontupper=\footnotesize\color{corr}]#2\end{tcolorbox}%
  \else
    \begin{tcolorbox}[colback=white,colframe=black!45,boxrule=0.4pt,arc=1pt,height=#1,
      before skip=2pt,after skip=2pt]\end{tcolorbox}%
  \fi}

\newcommand{\question}[3]{\begin{quest}\textbf{#1.} #3\pts{#2}\end{quest}}

%==========================================================================
% Données propres à chaque variante
%==========================================================================
\if\variante A
  \def\entreprise{Le cabinet d'architectes \textbf{ArchiNova}}
  \def\Rnom{R-ArchiNova}
  \def\resUn{Bureau d'études}\def\pUn{BE}\def\netUn{192.168.40}
  \def\resDeux{Accueil}\def\pDeux{ACC}\def\netDeux{172.16}
  \def\srvNom{SRV-Intranet}\def\netSrv{192.168.200}\def\srvH{10}
  \def\ssid{ArchiNova-WiFi}
  \def\AIrows{%
    \hl\ip{192.168.40.25} & \sol{C} & \sol{255.255.255.0} & \sol{192.168.40.0} & \sol{192.168.40.255} & \sol{254} \\ \hline
    \hl\ip{172.16.5.10}   & \sol{B} & \sol{255.255.0.0}   & \sol{172.16.0.0}   & \sol{172.16.255.255} & \sol{65\,534} \\ \hline
    \hl\ip{10.3.2.1}      & \sol{A} & \sol{255.0.0.0}     & \sol{10.0.0.0}     & \sol{10.255.255.255} & \sol{16\,777\,214} \\ \hline}
  \def\apipa{169.254.12.7}
  \def\Cnet{192.168.60}
  \def\Crows{%
    PC-A & \ip{192.168.60.10} & \ip{255.255.255.0} & \ip{192.168.60.254} & — \\
    PC-B & \ip{192.168.60.20} & \ip{255.255.255.0} & \ip{192.168.60.1}   & — \\
    PC-C & \ip{192.168.61.30} & \ip{255.255.255.0} & \ip{192.168.60.254} & — \\
    PC-D & \ip{192.168.60.10} & \ip{255.255.255.0} & \ip{192.168.60.254} & — \\
    SRV-C & \ip{10.0.0.5}     & \ip{255.0.0.0}     & \ip{10.0.0.254}     & — \\
    R1 — G0/0 & \ip{192.168.60.254} & \ip{255.255.255.0} & — & On \\
    R1 — G0/1 & \ip{10.0.0.254}     & \ip{255.0.0.0}     & — & Off \\}
  \def\Csol{%
    \hl 1 & \sol{PC-B} & \sol{Passerelle .60.1 : pas l'adresse du routeur.} & \sol{Joint son réseau, mais aucun autre (serveur).} & \sol{Passerelle 192.168.60.254.} \\ \hline
    \hl 2 & \sol{PC-C} & \sol{IP hors du réseau 192.168.60.0/24.} & \sol{Ne joint ni les postes ni le routeur.} & \sol{Ex. 192.168.60.30.} \\ \hline
    \hl 3 & \sol{PC-D} & \sol{Même IP que PC-A (doublon).} & \sol{Conflit d'adresse : PC-A et PC-D perturbés.} & \sol{Adresse unique, ex. .60.40.} \\ \hline
    \hl 4 & \sol{R1 G0/1} & \sol{Interface désactivée (Off).} & \sol{Réseau du serveur injoignable pour tous.} & \sol{Port Status : On.} \\ \hline}
  \def\Cbareme{PC-A et SRV-C sont correctement configurés ; accepter PC-A cité à la place de PC-D pour le doublon.}
\else
  \def\entreprise{La clinique vétérinaire \textbf{VetoPlus}}
  \def\Rnom{R-VetoPlus}
  \def\resUn{Soins}\def\pUn{SOI}\def\netUn{192.168.70}
  \def\resDeux{Secrétariat}\def\pDeux{SEC}\def\netDeux{172.20}
  \def\srvNom{SRV-Dossiers}\def\netSrv{192.168.250}\def\srvH{20}
  \def\ssid{VetoPlus-WiFi}
  \def\AIrows{%
    \hl\ip{192.168.15.130} & \sol{C} & \sol{255.255.255.0} & \sol{192.168.15.0} & \sol{192.168.15.255} & \sol{254} \\ \hline
    \hl\ip{172.31.200.7}   & \sol{B} & \sol{255.255.0.0}   & \sol{172.31.0.0}   & \sol{172.31.255.255} & \sol{65\,534} \\ \hline
    \hl\ip{10.45.6.9}      & \sol{A} & \sol{255.0.0.0}     & \sol{10.0.0.0}     & \sol{10.255.255.255} & \sol{16\,777\,214} \\ \hline}
  \def\apipa{169.254.88.3}
  \def\Cnet{192.168.80}
  \def\Crows{%
    PC-A & \ip{192.168.8.11}   & \ip{255.255.255.0} & \ip{192.168.80.254} & — \\
    PC-B & \ip{192.168.80.12}  & \ip{255.255.255.0} & \ip{192.168.80.254} & — \\
    PC-C & \ip{192.168.80.13}  & \ip{255.255.255.0} & \ip{192.168.80.100} & — \\
    PC-D & \ip{192.168.80.12}  & \ip{255.255.255.0} & \ip{192.168.80.254} & — \\
    SRV-C & \ip{10.0.0.8}      & \ip{255.0.0.0}     & \ip{10.0.0.254}     & — \\
    R1 — G0/0 & \ip{192.168.80.254} & \ip{255.255.255.0} & — & Off \\
    R1 — G0/1 & \ip{10.0.0.254}     & \ip{255.0.0.0}     & — & On \\}
  \def\Csol{%
    \hl 1 & \sol{PC-A} & \sol{IP hors du réseau 192.168.80.0/24.} & \sol{Ne joint ni les postes ni le routeur.} & \sol{Ex. 192.168.80.11.} \\ \hline
    \hl 2 & \sol{PC-C} & \sol{Passerelle .80.100 : pas l'adresse du routeur.} & \sol{Joint son réseau, mais aucun autre (serveur).} & \sol{Passerelle 192.168.80.254.} \\ \hline
    \hl 3 & \sol{PC-D} & \sol{Même IP que PC-B (doublon).} & \sol{Conflit d'adresse : PC-B et PC-D perturbés.} & \sol{Adresse unique, ex. .80.14.} \\ \hline
    \hl 4 & \sol{R1 G0/0} & \sol{Interface côté LAN désactivée (Off).} & \sol{Aucun poste ne joint le routeur ni le serveur.} & \sol{Port Status : On.} \\ \hline}
  \def\Cbareme{PC-B et SRV-C sont correctement configurés ; accepter PC-B cité à la place de PC-D pour le doublon.}
\fi

\pagestyle{fancy}
\fancyhf{}
\renewcommand{\headrulewidth}{0pt}
\fancyfoot[L]{\footnotesize\itshape\color{grisfoot}Évaluation — Réseaux locaux — Sujet \variante\ifcorrige\ — CORRIGÉ\fi}
\fancyfoot[R]{\footnotesize\color{grisfoot}\thepage/3}

\setlist{itemsep=0pt,topsep=2pt,parsep=0pt}
\setlength{\parindent}{0pt}
\setlength{\parskip}{2pt}

\begin{document}

\begin{center}
  {\Large\bfseries Évaluation pratique — Réseaux locaux \quad
   \fcolorbox{black}{white}{\ SUJET \variante\ }\par}
  \vspace{1pt}
  {\bfseries Adressage IP, commutation et routage — BTS CIEL 1re année — durée : 2 heures\par}
  \ifcorrige\vspace{2pt}{\bfseries\color{corr}CORRIGÉ ET BARÈME — document enseignant}\par\fi
\end{center}
\vspace{-2pt}
{\renewcommand{\arraystretch}{1.35}
\begin{tabularx}{\linewidth}{|X|X|p{2.2cm}|p{2.6cm}|}
  \hline
  \textbf{Nom :} & \textbf{Prénom :} & \textbf{Poste n°} & \textbf{Note :} \hfill /20 \\ \hline
\end{tabularx}}

\textbf{Consignes :} travail \textbf{individuel}, documents et Internet \textbf{non autorisés}. Ce sujet est le
\textbf{document-réponse}. Fichier Packet Tracer : \ip{Nom\_Prenom\_EvalLAN\_\variante.pkt}, à enregistrer
régulièrement dans le dossier de rendu. Les « visas » sont apposés par le professeur, que vous appelez à chaque test réussi.
Durées conseillées : A 30 min — B 1 h 10 — C 20 min.

%==========================================================================
\section{Partie A — Questions de cours et calculs \hfill /6}

\question{A1}{2 pts}{Complétez le tableau en considérant le \textbf{masque par défaut} de la classe de chaque adresse.}
{\small\renewcommand{\arraystretch}{1.15}
\begin{tabularx}{\linewidth}{|l|c|C|C|C|c|}
  \hline
  \textbf{Adresse IP} & \textbf{Classe} & \textbf{Masque par défaut} & \textbf{Adresse réseau} &
  \textbf{Adresse de diffusion} & \textbf{Nb d'hôtes} \\ \hline
  \AIrows
\end{tabularx}}
\bareme{0,5 pt par ligne juste + 0,5 pt si les trois nombres d'hôtes sont justes ($2^{n}-2$, $n$ = nombre de bits d'hôte).}

\question{A2}{1 pt}{Indiquez le type de câble (\textbf{droit} ou \textbf{croisé}) pour chaque liaison.}
{\small\renewcommand{\arraystretch}{1.15}
\begin{tabularx}{\linewidth}{|X|c||X|c|}
  \hline
  \hl PC $\leftrightarrow$ commutateur & \makebox[1.6cm]{\sol{droit}} & Commutateur $\leftrightarrow$ routeur & \makebox[1.6cm]{\sol{droit}} \\ \hline
  \hl PC $\leftrightarrow$ PC (réseau P2P) & \sol{croisé} & Concentrateur $\leftrightarrow$ concentrateur & \sol{croisé} \\ \hline
\end{tabularx}}
\bareme{0,25 pt par case. Équipements de types différents : câble droit ; de même type : câble croisé.}

\question{A3}{1 pt}{Comment un \textbf{concentrateur (hub)} et un \textbf{commutateur (switch)} traitent-ils une trame
reçue ? Lequel limite les collisions, et pourquoi ?}
\zone{2.9cm}{Le \textbf{hub} (couche 1) répète la trame sur \textbf{tous ses autres ports} : un seul domaine de collision
partagé. Le \textbf{switch} (couche 2) lit l'adresse MAC de destination et, grâce à sa \textbf{table MAC}, ne l'envoie
\textbf{que sur le port du destinataire} : un domaine de collision par port, c'est lui qui limite les collisions.
\emph{0,5 pt hub, 0,5 pt switch.}}

\question{A4}{1 pt}{Qu'est-ce que la \textbf{passerelle par défaut} d'un poste ? Quelle condition son adresse doit-elle respecter ?}
\zone{2.4cm}{Adresse IP de l'interface du \textbf{routeur} à laquelle le poste envoie les paquets destinés à un
\textbf{autre réseau}. Elle doit \textbf{appartenir au même réseau} que le poste. \emph{0,5 + 0,5 pt.}}

\question{A5}{0,5 pt}{Un portable configuré en DHCP affiche \ip{\apipa} / \ip{255.255.0.0}. Que s'est-il passé ?}
\zone{1.5cm}{Aucun \textbf{serveur DHCP} n'a répondu : le poste s'est attribué une adresse \textbf{APIPA} (169.254.0.0/16).}

\question{A6}{0,5 pt}{Quel protocole utilise \texttt{ping} ? Nommez les deux messages échangés et la couche OSI des adresses IP.}
\zone{1.5cm}{\textbf{ICMP} ; \textbf{Echo Request} / \textbf{Echo Reply} ; adresses IP en \textbf{couche 3 (Réseau)}.}

%==========================================================================
\newpage
\section{Partie B — Réalisation sous Packet Tracer \hfill /10}

\begin{cdc}
\entreprise{} vous confie son réseau : \textbf{trois réseaux} reliés par un \textbf{routeur unique Cisco 2911}
nommé \ip{\Rnom} (interfaces G0/0, G0/1, G0/2).
\begin{itemize}[leftmargin=12pt]
  \item \textbf{Réseau \resUn} : \textbf{4 postes} PC-\pUn1 à PC-\pUn4, commutateur \ip{SW-\pUn}, adressage statique,
        réseau \ip{\netUn.0/24}.
  \item \textbf{Réseau \resDeux} : \textbf{2 postes} PC-\pDeux1, PC-\pDeux2 et \textbf{1 imprimante} IMP-\pDeux,
        commutateur \ip{SW-\pDeux}, adressage statique, réseau de classe B \ip{\netDeux.0.0/16}.
  \item \textbf{Réseau Serveur} : serveur \ip{\srvNom} d'adresse \ip{\netSrv.\srvH/24}, relié \textbf{directement} au routeur.
  \item La passerelle de chaque réseau est sa \textbf{dernière adresse utilisable}. Tous les équipements communiquent
        entre eux et accèdent à la page web du serveur. Nommage imposé (« Display Name »).
\end{itemize}
\end{cdc}

\question{B1}{2 pts}{Complétez le plan d'adressage \textbf{avant} la réalisation.}
{\footnotesize\setlength{\tabcolsep}{3pt}\renewcommand{\arraystretch}{1.15}
\begin{tabularx}{\linewidth}{|l|C|C|C|C|C|C|}
  \hline
  \textbf{Réseau} & \textbf{Adresse réseau} & \textbf{Masque} & \textbf{1re @ utilisable} &
  \textbf{Dernière @ utilisable} & \textbf{Diffusion} & \textbf{Passerelle} \\ \hline
  \hl \resUn & \sol{\netUn.0} & \sol{255.255.255.0} & \sol{\netUn.1} & \sol{\netUn.254} & \sol{\netUn.255} & \sol{\netUn.254} \\ \hline
  \hl \resDeux & \sol{\netDeux.0.0} & \sol{255.255.0.0} & \sol{\netDeux.0.1} & \sol{\netDeux.255.254} & \sol{\netDeux.255.255} & \sol{\netDeux.255.254} \\ \hline
  \hl Serveur & \sol{\netSrv.0} & \sol{255.255.255.0} & \sol{\netSrv.1} & \sol{\netSrv.254} & \sol{\netSrv.255} & \sol{\netSrv.254} \\ \hline
\end{tabularx}}
\bareme{0,5 pt par ligne juste + 0,5 pt si la convention de passerelle est respectée partout.}

\question{B2}{2 pts}{Construisez la topologie (équipements, câblage, nommage) et relevez l'adressage configuré.}
{\footnotesize\setlength{\tabcolsep}{3pt}\renewcommand{\arraystretch}{1.1}
\begin{tabularx}{\linewidth}{|l|C|C|C||l|C|C|}
  \hline
  \textbf{Hôte} & \textbf{Adresse IP} & \textbf{Masque} & \textbf{Passerelle} & \textbf{Interface} & \textbf{Adresse IP} & \textbf{Masque} \\ \hline
  \hl PC-\pUn1 & \sol{\netUn.1} & \sol{255.255.255.0} & \sol{\netUn.254} & \Rnom{} G0/0 & \sol{\netUn.254} & \sol{255.255.255.0} \\ \hline
  \hl PC-\pDeux1 & \sol{\netDeux.0.1} & \sol{255.255.0.0} & \sol{\netDeux.255.254} & \Rnom{} G0/1 & \sol{\netDeux.255.254} & \sol{255.255.0.0} \\ \hline
  \hl IMP-\pDeux & \sol{\netDeux.0.10} & \sol{255.255.0.0} & \sol{\netDeux.255.254} & \Rnom{} G0/2 & \sol{\netSrv.254} & \sol{255.255.255.0} \\ \hline
  \hl \srvNom & \ip{\netSrv.\srvH} & \sol{255.255.255.0} & \sol{\netSrv.254} & \multicolumn{3}{c|}{\cellcolor{black!8}} \\ \hline
\end{tabularx}}
\bareme{Adresses d'hôtes libres (uniques, dans le bon réseau) ; affectation G0/x libre. Sur le .pkt : équipements 0,5 ;
câblage 1 (câbles droits vers les switchs et switch–routeur, \textbf{croisé} serveur–routeur, liens verts) ; nommage 0,5.}

\question{B3}{3 pts}{Configurez tous les équipements : IP, masques, passerelles des hôtes ; interfaces du routeur
(« Port Status : On »).}
\bareme{Sur le .pkt : interfaces routeur 1 pt ; IP/masques des hôtes 1 pt ; passerelles (imprimante et serveur compris)
1 pt. $-0,25$ par erreur, sans descendre sous 0 par item.}

\question{B4}{3 pts}{Réalisez les tests, notez la commande et le résultat, puis faites viser chaque test réussi.}
{\small\renewcommand{\arraystretch}{1.1}
\begin{tabularx}{\linewidth}{|c|l|l|C|c|c|}
  \hline
  \textbf{N°} & \textbf{Source} & \textbf{Destination} & \textbf{Commande / action} & \textbf{Résultat} & \textbf{Visa prof.} \\ \hline
  \hl 1 & PC-\pUn1 & PC-\pUn4 & \sol{ping \netUn.4} & \sol{réussi} & \hspace{1.4cm} \\ \hline
  \hl 2 & PC-\pDeux1 & IMP-\pDeux & \sol{ping \netDeux.0.10} & \sol{réussi} & \\ \hline
  \hl 3 & PC-\pUn2 & PC-\pDeux2 & \sol{ping \netDeux.0.2} & \sol{réussi} & \\ \hline
  \hl 4 & PC-\pDeux2 & \srvNom & \sol{ping \netSrv.\srvH} & \sol{réussi} & \\ \hline
  \hl 5 & PC-\pUn3 & page web du serveur & \sol{navigateur : {\NoAutoSpacing http://\netSrv.\srvH}} & \sol{page affichée} & \\ \hline
\end{tabularx}}
\bareme{0,5 pt par test visé + 0,5 pt pour B5. Le premier ping inter-réseaux peut échouer (résolution ARP) ; les suivants doivent réussir.}

\question{B5}{incl. B4}{Test n°3 : par quels équipements passe le paquet ? Pourquoi PC-\pUn2 l'envoie-t-il au routeur ?}
\zone{2.3cm}{PC-\pUn2 $\to$ SW-\pUn{} $\to$ \Rnom{} $\to$ SW-\pDeux{} $\to$ PC-\pDeux2. Grâce à son masque, PC-\pUn2 constate que
la destination (\netDeux.0.0/16) n'est \textbf{pas dans son réseau} (\netUn.0/24) : il envoie le paquet à sa
\textbf{passerelle par défaut}, qui le route.}

%==========================================================================
\newpage
\section{Partie C — Diagnostic d'un réseau en panne \hfill /4}

Un technicien constate qu'\textbf{aucun poste ne parvient à joindre le serveur SRV-C} et que certains postes
communiquent mal entre eux. Les câbles sont du bon type.

\begin{center}
\begin{minipage}{0.37\linewidth}\centering
\begin{tikzpicture}[font=\footnotesize,
  eq/.style={draw,fill=qbg,rounded corners=2pt,minimum width=9mm,minimum height=5mm,inner sep=2pt},
  every path/.style={thick}]
  \node[eq] (SW) {SW-1};
  \node[eq,right=9mm of SW] (R) {R1};
  \node[eq,right=7mm of R] (SRV) {SRV-C};
  \node[eq,above left=5mm and -3mm of SW] (A) {PC-A};
  \node[eq,above right=5mm and -3mm of SW] (B) {PC-B};
  \node[eq,below left=5mm and -3mm of SW] (C) {PC-C};
  \node[eq,below right=5mm and -3mm of SW] (D) {PC-D};
  \foreach \n in {A,B,C,D} \draw (\n) -- (SW);
  \draw (SW) -- node[above,font=\tiny]{G0/0} (R);
  \draw (R) -- node[above,font=\tiny]{G0/1} (SRV);
\end{tikzpicture}
\end{minipage}\hfill
\begin{minipage}{0.61\linewidth}\centering
{\small\renewcommand{\arraystretch}{1.05}\setlength{\tabcolsep}{4pt}
\begin{tabular}{lllll}
  \toprule
  \textbf{Équipement} & \textbf{Adresse IP} & \textbf{Masque} & \textbf{Passerelle} & \textbf{Port} \\
  \midrule
  \Crows
  \bottomrule
\end{tabular}}
\end{minipage}
\end{center}

\question{C1}{4 pts}{Relevez les \textbf{quatre anomalies}. Pour chacune : conséquence et correction.}
{\small\renewcommand{\arraystretch}{1.1}
\ifcorrige\def\hc{}\else\def\hc{\rule[-10mm]{0pt}{14mm}}\fi
\begin{tabularx}{\linewidth}{|c|>{\raggedright\arraybackslash}p{2cm}|L|L|L|}
  \hline
  \textbf{N°} & \textbf{Équipement} & \textbf{Anomalie} & \textbf{Conséquence} & \textbf{Correction} \\ \hline
  \ifcorrige\Csol\else
  \hc 1 & & & & \\ \hline
  \hc 2 & & & & \\ \hline
  \hc 3 & & & & \\ \hline
  \hc 4 & & & & \\ \hline
  \fi
\end{tabularx}}
\bareme{1 pt par anomalie (0,25 équipement, anomalie, conséquence, correction ; ordre indifférent). \Cbareme}

%==========================================================================
\section{Bonus — Extension sans fil \hfill +1}

\begin{cdc}Ajoutez au réseau \resUn{} un \textbf{AccessPoint-PT} relié à SW-\pUn{} (SSID \ip{\ssid}, \textbf{WPA2-PSK})
et un portable \ip{Laptop-\pUn} équipé d'un module \texttt{WPC300N}, en adressage \textbf{statique}.
Vérifiez qu'il joint \srvNom.\end{cdc}
{\small\renewcommand{\arraystretch}{1.1}
\begin{tabularx}{\linewidth}{|l|C|C|C|c|}
  \hline
  \textbf{Équipement} & \textbf{Adresse IP} & \textbf{Masque} & \textbf{Passerelle} & \textbf{Visa prof.} \\ \hline
  \hl Laptop-\pUn & \sol{\netUn.50} & \sol{255.255.255.0} & \sol{\netUn.254} & \hspace{1.4cm} \\ \hline
\end{tabularx}}
\bareme{Le point d'accès fait office de pont : le portable est dans \netUn.0/24. Bonus si associé en WPA2-PSK et ping
vers \netSrv.\srvH{} réussi. Note plafonnée à 20.}

\end{document}
"
%   pdflatex -jobname=Eval_Reseau_CorrigeA "\def\variante{A}\def\modecorrige{1}% Évaluation pratique — Réseaux locaux (BTS CIEL 1re année) — sujets A et B
% Compilation (4 PDF) :
%   pdflatex -jobname=Eval_Reseau_SujetA   "\def\variante{A}\def\modecorrige{0}\input{Eval_TP_Reseau_LAN.tex}"
%   pdflatex -jobname=Eval_Reseau_SujetB   "\def\variante{B}\def\modecorrige{0}\input{Eval_TP_Reseau_LAN.tex}"
%   pdflatex -jobname=Eval_Reseau_CorrigeA "\def\variante{A}\def\modecorrige{1}\input{Eval_TP_Reseau_LAN.tex}"
%   pdflatex -jobname=Eval_Reseau_CorrigeB "\def\variante{B}\def\modecorrige{1}\input{Eval_TP_Reseau_LAN.tex}"
\documentclass[10pt,a4paper]{article}

\usepackage[T1]{fontenc}
\usepackage[utf8]{inputenc}
\usepackage{lmodern}
\usepackage[french]{babel}
\usepackage[a4paper,hmargin=1.5cm,top=1.3cm,bottom=1.6cm,footskip=0.7cm]{geometry}
\usepackage[table]{xcolor}
\usepackage{booktabs,tabularx,array}
\usepackage{enumitem}
\usepackage{fancyhdr}
\usepackage[most]{tcolorbox}
\usepackage{tikz}
\usetikzlibrary{positioning}
\usepackage{titlesec}

\providecommand{\modecorrige}{0}
\providecommand{\variante}{A}
\newif\ifcorrige
\ifnum\modecorrige=1 \corrigetrue \else \corrigefalse \fi

\definecolor{ipred}{HTML}{C0392B}
\definecolor{cdcbar}{HTML}{A67C3D}
\definecolor{cdcbg}{HTML}{F4F1EA}
\definecolor{qbar}{HTML}{4A6B5D}
\definecolor{qbg}{HTML}{EDF1EF}
\definecolor{corr}{HTML}{1F5FA8}
\definecolor{corrbg}{HTML}{E8F0FA}
\definecolor{grisfoot}{HTML}{828282}

\newcommand{\ip}[1]{\textcolor{ipred}{\texttt{#1}}}
\newcommand{\pts}[1]{\hfill\mbox{\textbf{(#1)}}}
\newcommand{\sol}[1]{\ifcorrige\textcolor{corr}{#1}\fi}
\newcommand{\bareme}[1]{\ifcorrige\par{\footnotesize\color{corr}#1\par}\fi}
\newcommand{\hl}{\ifcorrige\rule[-1.6mm]{0pt}{5.4mm}\else\rule[-2.2mm]{0pt}{6.6mm}\fi}
\newcolumntype{C}{>{\centering\arraybackslash}X}
\newcolumntype{L}{>{\raggedright\arraybackslash}X}

\titleformat{\section}{\large\bfseries}{}{0pt}{}
\titlespacing*{\section}{0pt}{6pt}{3pt}

\tcbset{qstyle/.style={enhanced,boxrule=0pt,leftrule=2.5pt,arc=0pt,outer arc=0pt,
  left=5pt,right=5pt,top=2pt,bottom=2pt,before skip=4pt,after skip=3pt}}
\newtcolorbox{cdc}{qstyle,colback=cdcbg,colframe=cdcbar}
\newtcolorbox{quest}{qstyle,colback=qbg,colframe=qbar}

% Zone de réponse : vierge dans le sujet, remplie dans le corrigé
\newcommand{\zone}[2]{%
  \ifcorrige
    \begin{tcolorbox}[colback=corrbg,colframe=corr,boxrule=0.5pt,arc=1pt,left=4pt,right=4pt,top=1pt,bottom=1pt,
      before skip=2pt,after skip=2pt,fontupper=\footnotesize\color{corr}]#2\end{tcolorbox}%
  \else
    \begin{tcolorbox}[colback=white,colframe=black!45,boxrule=0.4pt,arc=1pt,height=#1,
      before skip=2pt,after skip=2pt]\end{tcolorbox}%
  \fi}

\newcommand{\question}[3]{\begin{quest}\textbf{#1.} #3\pts{#2}\end{quest}}

%==========================================================================
% Données propres à chaque variante
%==========================================================================
\if\variante A
  \def\entreprise{Le cabinet d'architectes \textbf{ArchiNova}}
  \def\Rnom{R-ArchiNova}
  \def\resUn{Bureau d'études}\def\pUn{BE}\def\netUn{192.168.40}
  \def\resDeux{Accueil}\def\pDeux{ACC}\def\netDeux{172.16}
  \def\srvNom{SRV-Intranet}\def\netSrv{192.168.200}\def\srvH{10}
  \def\ssid{ArchiNova-WiFi}
  \def\AIrows{%
    \hl\ip{192.168.40.25} & \sol{C} & \sol{255.255.255.0} & \sol{192.168.40.0} & \sol{192.168.40.255} & \sol{254} \\ \hline
    \hl\ip{172.16.5.10}   & \sol{B} & \sol{255.255.0.0}   & \sol{172.16.0.0}   & \sol{172.16.255.255} & \sol{65\,534} \\ \hline
    \hl\ip{10.3.2.1}      & \sol{A} & \sol{255.0.0.0}     & \sol{10.0.0.0}     & \sol{10.255.255.255} & \sol{16\,777\,214} \\ \hline}
  \def\apipa{169.254.12.7}
  \def\Cnet{192.168.60}
  \def\Crows{%
    PC-A & \ip{192.168.60.10} & \ip{255.255.255.0} & \ip{192.168.60.254} & — \\
    PC-B & \ip{192.168.60.20} & \ip{255.255.255.0} & \ip{192.168.60.1}   & — \\
    PC-C & \ip{192.168.61.30} & \ip{255.255.255.0} & \ip{192.168.60.254} & — \\
    PC-D & \ip{192.168.60.10} & \ip{255.255.255.0} & \ip{192.168.60.254} & — \\
    SRV-C & \ip{10.0.0.5}     & \ip{255.0.0.0}     & \ip{10.0.0.254}     & — \\
    R1 — G0/0 & \ip{192.168.60.254} & \ip{255.255.255.0} & — & On \\
    R1 — G0/1 & \ip{10.0.0.254}     & \ip{255.0.0.0}     & — & Off \\}
  \def\Csol{%
    \hl 1 & \sol{PC-B} & \sol{Passerelle .60.1 : pas l'adresse du routeur.} & \sol{Joint son réseau, mais aucun autre (serveur).} & \sol{Passerelle 192.168.60.254.} \\ \hline
    \hl 2 & \sol{PC-C} & \sol{IP hors du réseau 192.168.60.0/24.} & \sol{Ne joint ni les postes ni le routeur.} & \sol{Ex. 192.168.60.30.} \\ \hline
    \hl 3 & \sol{PC-D} & \sol{Même IP que PC-A (doublon).} & \sol{Conflit d'adresse : PC-A et PC-D perturbés.} & \sol{Adresse unique, ex. .60.40.} \\ \hline
    \hl 4 & \sol{R1 G0/1} & \sol{Interface désactivée (Off).} & \sol{Réseau du serveur injoignable pour tous.} & \sol{Port Status : On.} \\ \hline}
  \def\Cbareme{PC-A et SRV-C sont correctement configurés ; accepter PC-A cité à la place de PC-D pour le doublon.}
\else
  \def\entreprise{La clinique vétérinaire \textbf{VetoPlus}}
  \def\Rnom{R-VetoPlus}
  \def\resUn{Soins}\def\pUn{SOI}\def\netUn{192.168.70}
  \def\resDeux{Secrétariat}\def\pDeux{SEC}\def\netDeux{172.20}
  \def\srvNom{SRV-Dossiers}\def\netSrv{192.168.250}\def\srvH{20}
  \def\ssid{VetoPlus-WiFi}
  \def\AIrows{%
    \hl\ip{192.168.15.130} & \sol{C} & \sol{255.255.255.0} & \sol{192.168.15.0} & \sol{192.168.15.255} & \sol{254} \\ \hline
    \hl\ip{172.31.200.7}   & \sol{B} & \sol{255.255.0.0}   & \sol{172.31.0.0}   & \sol{172.31.255.255} & \sol{65\,534} \\ \hline
    \hl\ip{10.45.6.9}      & \sol{A} & \sol{255.0.0.0}     & \sol{10.0.0.0}     & \sol{10.255.255.255} & \sol{16\,777\,214} \\ \hline}
  \def\apipa{169.254.88.3}
  \def\Cnet{192.168.80}
  \def\Crows{%
    PC-A & \ip{192.168.8.11}   & \ip{255.255.255.0} & \ip{192.168.80.254} & — \\
    PC-B & \ip{192.168.80.12}  & \ip{255.255.255.0} & \ip{192.168.80.254} & — \\
    PC-C & \ip{192.168.80.13}  & \ip{255.255.255.0} & \ip{192.168.80.100} & — \\
    PC-D & \ip{192.168.80.12}  & \ip{255.255.255.0} & \ip{192.168.80.254} & — \\
    SRV-C & \ip{10.0.0.8}      & \ip{255.0.0.0}     & \ip{10.0.0.254}     & — \\
    R1 — G0/0 & \ip{192.168.80.254} & \ip{255.255.255.0} & — & Off \\
    R1 — G0/1 & \ip{10.0.0.254}     & \ip{255.0.0.0}     & — & On \\}
  \def\Csol{%
    \hl 1 & \sol{PC-A} & \sol{IP hors du réseau 192.168.80.0/24.} & \sol{Ne joint ni les postes ni le routeur.} & \sol{Ex. 192.168.80.11.} \\ \hline
    \hl 2 & \sol{PC-C} & \sol{Passerelle .80.100 : pas l'adresse du routeur.} & \sol{Joint son réseau, mais aucun autre (serveur).} & \sol{Passerelle 192.168.80.254.} \\ \hline
    \hl 3 & \sol{PC-D} & \sol{Même IP que PC-B (doublon).} & \sol{Conflit d'adresse : PC-B et PC-D perturbés.} & \sol{Adresse unique, ex. .80.14.} \\ \hline
    \hl 4 & \sol{R1 G0/0} & \sol{Interface côté LAN désactivée (Off).} & \sol{Aucun poste ne joint le routeur ni le serveur.} & \sol{Port Status : On.} \\ \hline}
  \def\Cbareme{PC-B et SRV-C sont correctement configurés ; accepter PC-B cité à la place de PC-D pour le doublon.}
\fi

\pagestyle{fancy}
\fancyhf{}
\renewcommand{\headrulewidth}{0pt}
\fancyfoot[L]{\footnotesize\itshape\color{grisfoot}Évaluation — Réseaux locaux — Sujet \variante\ifcorrige\ — CORRIGÉ\fi}
\fancyfoot[R]{\footnotesize\color{grisfoot}\thepage/3}

\setlist{itemsep=0pt,topsep=2pt,parsep=0pt}
\setlength{\parindent}{0pt}
\setlength{\parskip}{2pt}

\begin{document}

\begin{center}
  {\Large\bfseries Évaluation pratique — Réseaux locaux \quad
   \fcolorbox{black}{white}{\ SUJET \variante\ }\par}
  \vspace{1pt}
  {\bfseries Adressage IP, commutation et routage — BTS CIEL 1re année — durée : 2 heures\par}
  \ifcorrige\vspace{2pt}{\bfseries\color{corr}CORRIGÉ ET BARÈME — document enseignant}\par\fi
\end{center}
\vspace{-2pt}
{\renewcommand{\arraystretch}{1.35}
\begin{tabularx}{\linewidth}{|X|X|p{2.2cm}|p{2.6cm}|}
  \hline
  \textbf{Nom :} & \textbf{Prénom :} & \textbf{Poste n°} & \textbf{Note :} \hfill /20 \\ \hline
\end{tabularx}}

\textbf{Consignes :} travail \textbf{individuel}, documents et Internet \textbf{non autorisés}. Ce sujet est le
\textbf{document-réponse}. Fichier Packet Tracer : \ip{Nom\_Prenom\_EvalLAN\_\variante.pkt}, à enregistrer
régulièrement dans le dossier de rendu. Les « visas » sont apposés par le professeur, que vous appelez à chaque test réussi.
Durées conseillées : A 30 min — B 1 h 10 — C 20 min.

%==========================================================================
\section{Partie A — Questions de cours et calculs \hfill /6}

\question{A1}{2 pts}{Complétez le tableau en considérant le \textbf{masque par défaut} de la classe de chaque adresse.}
{\small\renewcommand{\arraystretch}{1.15}
\begin{tabularx}{\linewidth}{|l|c|C|C|C|c|}
  \hline
  \textbf{Adresse IP} & \textbf{Classe} & \textbf{Masque par défaut} & \textbf{Adresse réseau} &
  \textbf{Adresse de diffusion} & \textbf{Nb d'hôtes} \\ \hline
  \AIrows
\end{tabularx}}
\bareme{0,5 pt par ligne juste + 0,5 pt si les trois nombres d'hôtes sont justes ($2^{n}-2$, $n$ = nombre de bits d'hôte).}

\question{A2}{1 pt}{Indiquez le type de câble (\textbf{droit} ou \textbf{croisé}) pour chaque liaison.}
{\small\renewcommand{\arraystretch}{1.15}
\begin{tabularx}{\linewidth}{|X|c||X|c|}
  \hline
  \hl PC $\leftrightarrow$ commutateur & \makebox[1.6cm]{\sol{droit}} & Commutateur $\leftrightarrow$ routeur & \makebox[1.6cm]{\sol{droit}} \\ \hline
  \hl PC $\leftrightarrow$ PC (réseau P2P) & \sol{croisé} & Concentrateur $\leftrightarrow$ concentrateur & \sol{croisé} \\ \hline
\end{tabularx}}
\bareme{0,25 pt par case. Équipements de types différents : câble droit ; de même type : câble croisé.}

\question{A3}{1 pt}{Comment un \textbf{concentrateur (hub)} et un \textbf{commutateur (switch)} traitent-ils une trame
reçue ? Lequel limite les collisions, et pourquoi ?}
\zone{2.9cm}{Le \textbf{hub} (couche 1) répète la trame sur \textbf{tous ses autres ports} : un seul domaine de collision
partagé. Le \textbf{switch} (couche 2) lit l'adresse MAC de destination et, grâce à sa \textbf{table MAC}, ne l'envoie
\textbf{que sur le port du destinataire} : un domaine de collision par port, c'est lui qui limite les collisions.
\emph{0,5 pt hub, 0,5 pt switch.}}

\question{A4}{1 pt}{Qu'est-ce que la \textbf{passerelle par défaut} d'un poste ? Quelle condition son adresse doit-elle respecter ?}
\zone{2.4cm}{Adresse IP de l'interface du \textbf{routeur} à laquelle le poste envoie les paquets destinés à un
\textbf{autre réseau}. Elle doit \textbf{appartenir au même réseau} que le poste. \emph{0,5 + 0,5 pt.}}

\question{A5}{0,5 pt}{Un portable configuré en DHCP affiche \ip{\apipa} / \ip{255.255.0.0}. Que s'est-il passé ?}
\zone{1.5cm}{Aucun \textbf{serveur DHCP} n'a répondu : le poste s'est attribué une adresse \textbf{APIPA} (169.254.0.0/16).}

\question{A6}{0,5 pt}{Quel protocole utilise \texttt{ping} ? Nommez les deux messages échangés et la couche OSI des adresses IP.}
\zone{1.5cm}{\textbf{ICMP} ; \textbf{Echo Request} / \textbf{Echo Reply} ; adresses IP en \textbf{couche 3 (Réseau)}.}

%==========================================================================
\newpage
\section{Partie B — Réalisation sous Packet Tracer \hfill /10}

\begin{cdc}
\entreprise{} vous confie son réseau : \textbf{trois réseaux} reliés par un \textbf{routeur unique Cisco 2911}
nommé \ip{\Rnom} (interfaces G0/0, G0/1, G0/2).
\begin{itemize}[leftmargin=12pt]
  \item \textbf{Réseau \resUn} : \textbf{4 postes} PC-\pUn1 à PC-\pUn4, commutateur \ip{SW-\pUn}, adressage statique,
        réseau \ip{\netUn.0/24}.
  \item \textbf{Réseau \resDeux} : \textbf{2 postes} PC-\pDeux1, PC-\pDeux2 et \textbf{1 imprimante} IMP-\pDeux,
        commutateur \ip{SW-\pDeux}, adressage statique, réseau de classe B \ip{\netDeux.0.0/16}.
  \item \textbf{Réseau Serveur} : serveur \ip{\srvNom} d'adresse \ip{\netSrv.\srvH/24}, relié \textbf{directement} au routeur.
  \item La passerelle de chaque réseau est sa \textbf{dernière adresse utilisable}. Tous les équipements communiquent
        entre eux et accèdent à la page web du serveur. Nommage imposé (« Display Name »).
\end{itemize}
\end{cdc}

\question{B1}{2 pts}{Complétez le plan d'adressage \textbf{avant} la réalisation.}
{\footnotesize\setlength{\tabcolsep}{3pt}\renewcommand{\arraystretch}{1.15}
\begin{tabularx}{\linewidth}{|l|C|C|C|C|C|C|}
  \hline
  \textbf{Réseau} & \textbf{Adresse réseau} & \textbf{Masque} & \textbf{1re @ utilisable} &
  \textbf{Dernière @ utilisable} & \textbf{Diffusion} & \textbf{Passerelle} \\ \hline
  \hl \resUn & \sol{\netUn.0} & \sol{255.255.255.0} & \sol{\netUn.1} & \sol{\netUn.254} & \sol{\netUn.255} & \sol{\netUn.254} \\ \hline
  \hl \resDeux & \sol{\netDeux.0.0} & \sol{255.255.0.0} & \sol{\netDeux.0.1} & \sol{\netDeux.255.254} & \sol{\netDeux.255.255} & \sol{\netDeux.255.254} \\ \hline
  \hl Serveur & \sol{\netSrv.0} & \sol{255.255.255.0} & \sol{\netSrv.1} & \sol{\netSrv.254} & \sol{\netSrv.255} & \sol{\netSrv.254} \\ \hline
\end{tabularx}}
\bareme{0,5 pt par ligne juste + 0,5 pt si la convention de passerelle est respectée partout.}

\question{B2}{2 pts}{Construisez la topologie (équipements, câblage, nommage) et relevez l'adressage configuré.}
{\footnotesize\setlength{\tabcolsep}{3pt}\renewcommand{\arraystretch}{1.1}
\begin{tabularx}{\linewidth}{|l|C|C|C||l|C|C|}
  \hline
  \textbf{Hôte} & \textbf{Adresse IP} & \textbf{Masque} & \textbf{Passerelle} & \textbf{Interface} & \textbf{Adresse IP} & \textbf{Masque} \\ \hline
  \hl PC-\pUn1 & \sol{\netUn.1} & \sol{255.255.255.0} & \sol{\netUn.254} & \Rnom{} G0/0 & \sol{\netUn.254} & \sol{255.255.255.0} \\ \hline
  \hl PC-\pDeux1 & \sol{\netDeux.0.1} & \sol{255.255.0.0} & \sol{\netDeux.255.254} & \Rnom{} G0/1 & \sol{\netDeux.255.254} & \sol{255.255.0.0} \\ \hline
  \hl IMP-\pDeux & \sol{\netDeux.0.10} & \sol{255.255.0.0} & \sol{\netDeux.255.254} & \Rnom{} G0/2 & \sol{\netSrv.254} & \sol{255.255.255.0} \\ \hline
  \hl \srvNom & \ip{\netSrv.\srvH} & \sol{255.255.255.0} & \sol{\netSrv.254} & \multicolumn{3}{c|}{\cellcolor{black!8}} \\ \hline
\end{tabularx}}
\bareme{Adresses d'hôtes libres (uniques, dans le bon réseau) ; affectation G0/x libre. Sur le .pkt : équipements 0,5 ;
câblage 1 (câbles droits vers les switchs et switch–routeur, \textbf{croisé} serveur–routeur, liens verts) ; nommage 0,5.}

\question{B3}{3 pts}{Configurez tous les équipements : IP, masques, passerelles des hôtes ; interfaces du routeur
(« Port Status : On »).}
\bareme{Sur le .pkt : interfaces routeur 1 pt ; IP/masques des hôtes 1 pt ; passerelles (imprimante et serveur compris)
1 pt. $-0,25$ par erreur, sans descendre sous 0 par item.}

\question{B4}{3 pts}{Réalisez les tests, notez la commande et le résultat, puis faites viser chaque test réussi.}
{\small\renewcommand{\arraystretch}{1.1}
\begin{tabularx}{\linewidth}{|c|l|l|C|c|c|}
  \hline
  \textbf{N°} & \textbf{Source} & \textbf{Destination} & \textbf{Commande / action} & \textbf{Résultat} & \textbf{Visa prof.} \\ \hline
  \hl 1 & PC-\pUn1 & PC-\pUn4 & \sol{ping \netUn.4} & \sol{réussi} & \hspace{1.4cm} \\ \hline
  \hl 2 & PC-\pDeux1 & IMP-\pDeux & \sol{ping \netDeux.0.10} & \sol{réussi} & \\ \hline
  \hl 3 & PC-\pUn2 & PC-\pDeux2 & \sol{ping \netDeux.0.2} & \sol{réussi} & \\ \hline
  \hl 4 & PC-\pDeux2 & \srvNom & \sol{ping \netSrv.\srvH} & \sol{réussi} & \\ \hline
  \hl 5 & PC-\pUn3 & page web du serveur & \sol{navigateur : {\NoAutoSpacing http://\netSrv.\srvH}} & \sol{page affichée} & \\ \hline
\end{tabularx}}
\bareme{0,5 pt par test visé + 0,5 pt pour B5. Le premier ping inter-réseaux peut échouer (résolution ARP) ; les suivants doivent réussir.}

\question{B5}{incl. B4}{Test n°3 : par quels équipements passe le paquet ? Pourquoi PC-\pUn2 l'envoie-t-il au routeur ?}
\zone{2.3cm}{PC-\pUn2 $\to$ SW-\pUn{} $\to$ \Rnom{} $\to$ SW-\pDeux{} $\to$ PC-\pDeux2. Grâce à son masque, PC-\pUn2 constate que
la destination (\netDeux.0.0/16) n'est \textbf{pas dans son réseau} (\netUn.0/24) : il envoie le paquet à sa
\textbf{passerelle par défaut}, qui le route.}

%==========================================================================
\newpage
\section{Partie C — Diagnostic d'un réseau en panne \hfill /4}

Un technicien constate qu'\textbf{aucun poste ne parvient à joindre le serveur SRV-C} et que certains postes
communiquent mal entre eux. Les câbles sont du bon type.

\begin{center}
\begin{minipage}{0.37\linewidth}\centering
\begin{tikzpicture}[font=\footnotesize,
  eq/.style={draw,fill=qbg,rounded corners=2pt,minimum width=9mm,minimum height=5mm,inner sep=2pt},
  every path/.style={thick}]
  \node[eq] (SW) {SW-1};
  \node[eq,right=9mm of SW] (R) {R1};
  \node[eq,right=7mm of R] (SRV) {SRV-C};
  \node[eq,above left=5mm and -3mm of SW] (A) {PC-A};
  \node[eq,above right=5mm and -3mm of SW] (B) {PC-B};
  \node[eq,below left=5mm and -3mm of SW] (C) {PC-C};
  \node[eq,below right=5mm and -3mm of SW] (D) {PC-D};
  \foreach \n in {A,B,C,D} \draw (\n) -- (SW);
  \draw (SW) -- node[above,font=\tiny]{G0/0} (R);
  \draw (R) -- node[above,font=\tiny]{G0/1} (SRV);
\end{tikzpicture}
\end{minipage}\hfill
\begin{minipage}{0.61\linewidth}\centering
{\small\renewcommand{\arraystretch}{1.05}\setlength{\tabcolsep}{4pt}
\begin{tabular}{lllll}
  \toprule
  \textbf{Équipement} & \textbf{Adresse IP} & \textbf{Masque} & \textbf{Passerelle} & \textbf{Port} \\
  \midrule
  \Crows
  \bottomrule
\end{tabular}}
\end{minipage}
\end{center}

\question{C1}{4 pts}{Relevez les \textbf{quatre anomalies}. Pour chacune : conséquence et correction.}
{\small\renewcommand{\arraystretch}{1.1}
\ifcorrige\def\hc{}\else\def\hc{\rule[-10mm]{0pt}{14mm}}\fi
\begin{tabularx}{\linewidth}{|c|>{\raggedright\arraybackslash}p{2cm}|L|L|L|}
  \hline
  \textbf{N°} & \textbf{Équipement} & \textbf{Anomalie} & \textbf{Conséquence} & \textbf{Correction} \\ \hline
  \ifcorrige\Csol\else
  \hc 1 & & & & \\ \hline
  \hc 2 & & & & \\ \hline
  \hc 3 & & & & \\ \hline
  \hc 4 & & & & \\ \hline
  \fi
\end{tabularx}}
\bareme{1 pt par anomalie (0,25 équipement, anomalie, conséquence, correction ; ordre indifférent). \Cbareme}

%==========================================================================
\section{Bonus — Extension sans fil \hfill +1}

\begin{cdc}Ajoutez au réseau \resUn{} un \textbf{AccessPoint-PT} relié à SW-\pUn{} (SSID \ip{\ssid}, \textbf{WPA2-PSK})
et un portable \ip{Laptop-\pUn} équipé d'un module \texttt{WPC300N}, en adressage \textbf{statique}.
Vérifiez qu'il joint \srvNom.\end{cdc}
{\small\renewcommand{\arraystretch}{1.1}
\begin{tabularx}{\linewidth}{|l|C|C|C|c|}
  \hline
  \textbf{Équipement} & \textbf{Adresse IP} & \textbf{Masque} & \textbf{Passerelle} & \textbf{Visa prof.} \\ \hline
  \hl Laptop-\pUn & \sol{\netUn.50} & \sol{255.255.255.0} & \sol{\netUn.254} & \hspace{1.4cm} \\ \hline
\end{tabularx}}
\bareme{Le point d'accès fait office de pont : le portable est dans \netUn.0/24. Bonus si associé en WPA2-PSK et ping
vers \netSrv.\srvH{} réussi. Note plafonnée à 20.}

\end{document}
"
%   pdflatex -jobname=Eval_Reseau_CorrigeB "\def\variante{B}\def\modecorrige{1}% Évaluation pratique — Réseaux locaux (BTS CIEL 1re année) — sujets A et B
% Compilation (4 PDF) :
%   pdflatex -jobname=Eval_Reseau_SujetA   "\def\variante{A}\def\modecorrige{0}\input{Eval_TP_Reseau_LAN.tex}"
%   pdflatex -jobname=Eval_Reseau_SujetB   "\def\variante{B}\def\modecorrige{0}\input{Eval_TP_Reseau_LAN.tex}"
%   pdflatex -jobname=Eval_Reseau_CorrigeA "\def\variante{A}\def\modecorrige{1}\input{Eval_TP_Reseau_LAN.tex}"
%   pdflatex -jobname=Eval_Reseau_CorrigeB "\def\variante{B}\def\modecorrige{1}\input{Eval_TP_Reseau_LAN.tex}"
\documentclass[10pt,a4paper]{article}

\usepackage[T1]{fontenc}
\usepackage[utf8]{inputenc}
\usepackage{lmodern}
\usepackage[french]{babel}
\usepackage[a4paper,hmargin=1.5cm,top=1.3cm,bottom=1.6cm,footskip=0.7cm]{geometry}
\usepackage[table]{xcolor}
\usepackage{booktabs,tabularx,array}
\usepackage{enumitem}
\usepackage{fancyhdr}
\usepackage[most]{tcolorbox}
\usepackage{tikz}
\usetikzlibrary{positioning}
\usepackage{titlesec}

\providecommand{\modecorrige}{0}
\providecommand{\variante}{A}
\newif\ifcorrige
\ifnum\modecorrige=1 \corrigetrue \else \corrigefalse \fi

\definecolor{ipred}{HTML}{C0392B}
\definecolor{cdcbar}{HTML}{A67C3D}
\definecolor{cdcbg}{HTML}{F4F1EA}
\definecolor{qbar}{HTML}{4A6B5D}
\definecolor{qbg}{HTML}{EDF1EF}
\definecolor{corr}{HTML}{1F5FA8}
\definecolor{corrbg}{HTML}{E8F0FA}
\definecolor{grisfoot}{HTML}{828282}

\newcommand{\ip}[1]{\textcolor{ipred}{\texttt{#1}}}
\newcommand{\pts}[1]{\hfill\mbox{\textbf{(#1)}}}
\newcommand{\sol}[1]{\ifcorrige\textcolor{corr}{#1}\fi}
\newcommand{\bareme}[1]{\ifcorrige\par{\footnotesize\color{corr}#1\par}\fi}
\newcommand{\hl}{\ifcorrige\rule[-1.6mm]{0pt}{5.4mm}\else\rule[-2.2mm]{0pt}{6.6mm}\fi}
\newcolumntype{C}{>{\centering\arraybackslash}X}
\newcolumntype{L}{>{\raggedright\arraybackslash}X}

\titleformat{\section}{\large\bfseries}{}{0pt}{}
\titlespacing*{\section}{0pt}{6pt}{3pt}

\tcbset{qstyle/.style={enhanced,boxrule=0pt,leftrule=2.5pt,arc=0pt,outer arc=0pt,
  left=5pt,right=5pt,top=2pt,bottom=2pt,before skip=4pt,after skip=3pt}}
\newtcolorbox{cdc}{qstyle,colback=cdcbg,colframe=cdcbar}
\newtcolorbox{quest}{qstyle,colback=qbg,colframe=qbar}

% Zone de réponse : vierge dans le sujet, remplie dans le corrigé
\newcommand{\zone}[2]{%
  \ifcorrige
    \begin{tcolorbox}[colback=corrbg,colframe=corr,boxrule=0.5pt,arc=1pt,left=4pt,right=4pt,top=1pt,bottom=1pt,
      before skip=2pt,after skip=2pt,fontupper=\footnotesize\color{corr}]#2\end{tcolorbox}%
  \else
    \begin{tcolorbox}[colback=white,colframe=black!45,boxrule=0.4pt,arc=1pt,height=#1,
      before skip=2pt,after skip=2pt]\end{tcolorbox}%
  \fi}

\newcommand{\question}[3]{\begin{quest}\textbf{#1.} #3\pts{#2}\end{quest}}

%==========================================================================
% Données propres à chaque variante
%==========================================================================
\if\variante A
  \def\entreprise{Le cabinet d'architectes \textbf{ArchiNova}}
  \def\Rnom{R-ArchiNova}
  \def\resUn{Bureau d'études}\def\pUn{BE}\def\netUn{192.168.40}
  \def\resDeux{Accueil}\def\pDeux{ACC}\def\netDeux{172.16}
  \def\srvNom{SRV-Intranet}\def\netSrv{192.168.200}\def\srvH{10}
  \def\ssid{ArchiNova-WiFi}
  \def\AIrows{%
    \hl\ip{192.168.40.25} & \sol{C} & \sol{255.255.255.0} & \sol{192.168.40.0} & \sol{192.168.40.255} & \sol{254} \\ \hline
    \hl\ip{172.16.5.10}   & \sol{B} & \sol{255.255.0.0}   & \sol{172.16.0.0}   & \sol{172.16.255.255} & \sol{65\,534} \\ \hline
    \hl\ip{10.3.2.1}      & \sol{A} & \sol{255.0.0.0}     & \sol{10.0.0.0}     & \sol{10.255.255.255} & \sol{16\,777\,214} \\ \hline}
  \def\apipa{169.254.12.7}
  \def\Cnet{192.168.60}
  \def\Crows{%
    PC-A & \ip{192.168.60.10} & \ip{255.255.255.0} & \ip{192.168.60.254} & — \\
    PC-B & \ip{192.168.60.20} & \ip{255.255.255.0} & \ip{192.168.60.1}   & — \\
    PC-C & \ip{192.168.61.30} & \ip{255.255.255.0} & \ip{192.168.60.254} & — \\
    PC-D & \ip{192.168.60.10} & \ip{255.255.255.0} & \ip{192.168.60.254} & — \\
    SRV-C & \ip{10.0.0.5}     & \ip{255.0.0.0}     & \ip{10.0.0.254}     & — \\
    R1 — G0/0 & \ip{192.168.60.254} & \ip{255.255.255.0} & — & On \\
    R1 — G0/1 & \ip{10.0.0.254}     & \ip{255.0.0.0}     & — & Off \\}
  \def\Csol{%
    \hl 1 & \sol{PC-B} & \sol{Passerelle .60.1 : pas l'adresse du routeur.} & \sol{Joint son réseau, mais aucun autre (serveur).} & \sol{Passerelle 192.168.60.254.} \\ \hline
    \hl 2 & \sol{PC-C} & \sol{IP hors du réseau 192.168.60.0/24.} & \sol{Ne joint ni les postes ni le routeur.} & \sol{Ex. 192.168.60.30.} \\ \hline
    \hl 3 & \sol{PC-D} & \sol{Même IP que PC-A (doublon).} & \sol{Conflit d'adresse : PC-A et PC-D perturbés.} & \sol{Adresse unique, ex. .60.40.} \\ \hline
    \hl 4 & \sol{R1 G0/1} & \sol{Interface désactivée (Off).} & \sol{Réseau du serveur injoignable pour tous.} & \sol{Port Status : On.} \\ \hline}
  \def\Cbareme{PC-A et SRV-C sont correctement configurés ; accepter PC-A cité à la place de PC-D pour le doublon.}
\else
  \def\entreprise{La clinique vétérinaire \textbf{VetoPlus}}
  \def\Rnom{R-VetoPlus}
  \def\resUn{Soins}\def\pUn{SOI}\def\netUn{192.168.70}
  \def\resDeux{Secrétariat}\def\pDeux{SEC}\def\netDeux{172.20}
  \def\srvNom{SRV-Dossiers}\def\netSrv{192.168.250}\def\srvH{20}
  \def\ssid{VetoPlus-WiFi}
  \def\AIrows{%
    \hl\ip{192.168.15.130} & \sol{C} & \sol{255.255.255.0} & \sol{192.168.15.0} & \sol{192.168.15.255} & \sol{254} \\ \hline
    \hl\ip{172.31.200.7}   & \sol{B} & \sol{255.255.0.0}   & \sol{172.31.0.0}   & \sol{172.31.255.255} & \sol{65\,534} \\ \hline
    \hl\ip{10.45.6.9}      & \sol{A} & \sol{255.0.0.0}     & \sol{10.0.0.0}     & \sol{10.255.255.255} & \sol{16\,777\,214} \\ \hline}
  \def\apipa{169.254.88.3}
  \def\Cnet{192.168.80}
  \def\Crows{%
    PC-A & \ip{192.168.8.11}   & \ip{255.255.255.0} & \ip{192.168.80.254} & — \\
    PC-B & \ip{192.168.80.12}  & \ip{255.255.255.0} & \ip{192.168.80.254} & — \\
    PC-C & \ip{192.168.80.13}  & \ip{255.255.255.0} & \ip{192.168.80.100} & — \\
    PC-D & \ip{192.168.80.12}  & \ip{255.255.255.0} & \ip{192.168.80.254} & — \\
    SRV-C & \ip{10.0.0.8}      & \ip{255.0.0.0}     & \ip{10.0.0.254}     & — \\
    R1 — G0/0 & \ip{192.168.80.254} & \ip{255.255.255.0} & — & Off \\
    R1 — G0/1 & \ip{10.0.0.254}     & \ip{255.0.0.0}     & — & On \\}
  \def\Csol{%
    \hl 1 & \sol{PC-A} & \sol{IP hors du réseau 192.168.80.0/24.} & \sol{Ne joint ni les postes ni le routeur.} & \sol{Ex. 192.168.80.11.} \\ \hline
    \hl 2 & \sol{PC-C} & \sol{Passerelle .80.100 : pas l'adresse du routeur.} & \sol{Joint son réseau, mais aucun autre (serveur).} & \sol{Passerelle 192.168.80.254.} \\ \hline
    \hl 3 & \sol{PC-D} & \sol{Même IP que PC-B (doublon).} & \sol{Conflit d'adresse : PC-B et PC-D perturbés.} & \sol{Adresse unique, ex. .80.14.} \\ \hline
    \hl 4 & \sol{R1 G0/0} & \sol{Interface côté LAN désactivée (Off).} & \sol{Aucun poste ne joint le routeur ni le serveur.} & \sol{Port Status : On.} \\ \hline}
  \def\Cbareme{PC-B et SRV-C sont correctement configurés ; accepter PC-B cité à la place de PC-D pour le doublon.}
\fi

\pagestyle{fancy}
\fancyhf{}
\renewcommand{\headrulewidth}{0pt}
\fancyfoot[L]{\footnotesize\itshape\color{grisfoot}Évaluation — Réseaux locaux — Sujet \variante\ifcorrige\ — CORRIGÉ\fi}
\fancyfoot[R]{\footnotesize\color{grisfoot}\thepage/3}

\setlist{itemsep=0pt,topsep=2pt,parsep=0pt}
\setlength{\parindent}{0pt}
\setlength{\parskip}{2pt}

\begin{document}

\begin{center}
  {\Large\bfseries Évaluation pratique — Réseaux locaux \quad
   \fcolorbox{black}{white}{\ SUJET \variante\ }\par}
  \vspace{1pt}
  {\bfseries Adressage IP, commutation et routage — BTS CIEL 1re année — durée : 2 heures\par}
  \ifcorrige\vspace{2pt}{\bfseries\color{corr}CORRIGÉ ET BARÈME — document enseignant}\par\fi
\end{center}
\vspace{-2pt}
{\renewcommand{\arraystretch}{1.35}
\begin{tabularx}{\linewidth}{|X|X|p{2.2cm}|p{2.6cm}|}
  \hline
  \textbf{Nom :} & \textbf{Prénom :} & \textbf{Poste n°} & \textbf{Note :} \hfill /20 \\ \hline
\end{tabularx}}

\textbf{Consignes :} travail \textbf{individuel}, documents et Internet \textbf{non autorisés}. Ce sujet est le
\textbf{document-réponse}. Fichier Packet Tracer : \ip{Nom\_Prenom\_EvalLAN\_\variante.pkt}, à enregistrer
régulièrement dans le dossier de rendu. Les « visas » sont apposés par le professeur, que vous appelez à chaque test réussi.
Durées conseillées : A 30 min — B 1 h 10 — C 20 min.

%==========================================================================
\section{Partie A — Questions de cours et calculs \hfill /6}

\question{A1}{2 pts}{Complétez le tableau en considérant le \textbf{masque par défaut} de la classe de chaque adresse.}
{\small\renewcommand{\arraystretch}{1.15}
\begin{tabularx}{\linewidth}{|l|c|C|C|C|c|}
  \hline
  \textbf{Adresse IP} & \textbf{Classe} & \textbf{Masque par défaut} & \textbf{Adresse réseau} &
  \textbf{Adresse de diffusion} & \textbf{Nb d'hôtes} \\ \hline
  \AIrows
\end{tabularx}}
\bareme{0,5 pt par ligne juste + 0,5 pt si les trois nombres d'hôtes sont justes ($2^{n}-2$, $n$ = nombre de bits d'hôte).}

\question{A2}{1 pt}{Indiquez le type de câble (\textbf{droit} ou \textbf{croisé}) pour chaque liaison.}
{\small\renewcommand{\arraystretch}{1.15}
\begin{tabularx}{\linewidth}{|X|c||X|c|}
  \hline
  \hl PC $\leftrightarrow$ commutateur & \makebox[1.6cm]{\sol{droit}} & Commutateur $\leftrightarrow$ routeur & \makebox[1.6cm]{\sol{droit}} \\ \hline
  \hl PC $\leftrightarrow$ PC (réseau P2P) & \sol{croisé} & Concentrateur $\leftrightarrow$ concentrateur & \sol{croisé} \\ \hline
\end{tabularx}}
\bareme{0,25 pt par case. Équipements de types différents : câble droit ; de même type : câble croisé.}

\question{A3}{1 pt}{Comment un \textbf{concentrateur (hub)} et un \textbf{commutateur (switch)} traitent-ils une trame
reçue ? Lequel limite les collisions, et pourquoi ?}
\zone{2.9cm}{Le \textbf{hub} (couche 1) répète la trame sur \textbf{tous ses autres ports} : un seul domaine de collision
partagé. Le \textbf{switch} (couche 2) lit l'adresse MAC de destination et, grâce à sa \textbf{table MAC}, ne l'envoie
\textbf{que sur le port du destinataire} : un domaine de collision par port, c'est lui qui limite les collisions.
\emph{0,5 pt hub, 0,5 pt switch.}}

\question{A4}{1 pt}{Qu'est-ce que la \textbf{passerelle par défaut} d'un poste ? Quelle condition son adresse doit-elle respecter ?}
\zone{2.4cm}{Adresse IP de l'interface du \textbf{routeur} à laquelle le poste envoie les paquets destinés à un
\textbf{autre réseau}. Elle doit \textbf{appartenir au même réseau} que le poste. \emph{0,5 + 0,5 pt.}}

\question{A5}{0,5 pt}{Un portable configuré en DHCP affiche \ip{\apipa} / \ip{255.255.0.0}. Que s'est-il passé ?}
\zone{1.5cm}{Aucun \textbf{serveur DHCP} n'a répondu : le poste s'est attribué une adresse \textbf{APIPA} (169.254.0.0/16).}

\question{A6}{0,5 pt}{Quel protocole utilise \texttt{ping} ? Nommez les deux messages échangés et la couche OSI des adresses IP.}
\zone{1.5cm}{\textbf{ICMP} ; \textbf{Echo Request} / \textbf{Echo Reply} ; adresses IP en \textbf{couche 3 (Réseau)}.}

%==========================================================================
\newpage
\section{Partie B — Réalisation sous Packet Tracer \hfill /10}

\begin{cdc}
\entreprise{} vous confie son réseau : \textbf{trois réseaux} reliés par un \textbf{routeur unique Cisco 2911}
nommé \ip{\Rnom} (interfaces G0/0, G0/1, G0/2).
\begin{itemize}[leftmargin=12pt]
  \item \textbf{Réseau \resUn} : \textbf{4 postes} PC-\pUn1 à PC-\pUn4, commutateur \ip{SW-\pUn}, adressage statique,
        réseau \ip{\netUn.0/24}.
  \item \textbf{Réseau \resDeux} : \textbf{2 postes} PC-\pDeux1, PC-\pDeux2 et \textbf{1 imprimante} IMP-\pDeux,
        commutateur \ip{SW-\pDeux}, adressage statique, réseau de classe B \ip{\netDeux.0.0/16}.
  \item \textbf{Réseau Serveur} : serveur \ip{\srvNom} d'adresse \ip{\netSrv.\srvH/24}, relié \textbf{directement} au routeur.
  \item La passerelle de chaque réseau est sa \textbf{dernière adresse utilisable}. Tous les équipements communiquent
        entre eux et accèdent à la page web du serveur. Nommage imposé (« Display Name »).
\end{itemize}
\end{cdc}

\question{B1}{2 pts}{Complétez le plan d'adressage \textbf{avant} la réalisation.}
{\footnotesize\setlength{\tabcolsep}{3pt}\renewcommand{\arraystretch}{1.15}
\begin{tabularx}{\linewidth}{|l|C|C|C|C|C|C|}
  \hline
  \textbf{Réseau} & \textbf{Adresse réseau} & \textbf{Masque} & \textbf{1re @ utilisable} &
  \textbf{Dernière @ utilisable} & \textbf{Diffusion} & \textbf{Passerelle} \\ \hline
  \hl \resUn & \sol{\netUn.0} & \sol{255.255.255.0} & \sol{\netUn.1} & \sol{\netUn.254} & \sol{\netUn.255} & \sol{\netUn.254} \\ \hline
  \hl \resDeux & \sol{\netDeux.0.0} & \sol{255.255.0.0} & \sol{\netDeux.0.1} & \sol{\netDeux.255.254} & \sol{\netDeux.255.255} & \sol{\netDeux.255.254} \\ \hline
  \hl Serveur & \sol{\netSrv.0} & \sol{255.255.255.0} & \sol{\netSrv.1} & \sol{\netSrv.254} & \sol{\netSrv.255} & \sol{\netSrv.254} \\ \hline
\end{tabularx}}
\bareme{0,5 pt par ligne juste + 0,5 pt si la convention de passerelle est respectée partout.}

\question{B2}{2 pts}{Construisez la topologie (équipements, câblage, nommage) et relevez l'adressage configuré.}
{\footnotesize\setlength{\tabcolsep}{3pt}\renewcommand{\arraystretch}{1.1}
\begin{tabularx}{\linewidth}{|l|C|C|C||l|C|C|}
  \hline
  \textbf{Hôte} & \textbf{Adresse IP} & \textbf{Masque} & \textbf{Passerelle} & \textbf{Interface} & \textbf{Adresse IP} & \textbf{Masque} \\ \hline
  \hl PC-\pUn1 & \sol{\netUn.1} & \sol{255.255.255.0} & \sol{\netUn.254} & \Rnom{} G0/0 & \sol{\netUn.254} & \sol{255.255.255.0} \\ \hline
  \hl PC-\pDeux1 & \sol{\netDeux.0.1} & \sol{255.255.0.0} & \sol{\netDeux.255.254} & \Rnom{} G0/1 & \sol{\netDeux.255.254} & \sol{255.255.0.0} \\ \hline
  \hl IMP-\pDeux & \sol{\netDeux.0.10} & \sol{255.255.0.0} & \sol{\netDeux.255.254} & \Rnom{} G0/2 & \sol{\netSrv.254} & \sol{255.255.255.0} \\ \hline
  \hl \srvNom & \ip{\netSrv.\srvH} & \sol{255.255.255.0} & \sol{\netSrv.254} & \multicolumn{3}{c|}{\cellcolor{black!8}} \\ \hline
\end{tabularx}}
\bareme{Adresses d'hôtes libres (uniques, dans le bon réseau) ; affectation G0/x libre. Sur le .pkt : équipements 0,5 ;
câblage 1 (câbles droits vers les switchs et switch–routeur, \textbf{croisé} serveur–routeur, liens verts) ; nommage 0,5.}

\question{B3}{3 pts}{Configurez tous les équipements : IP, masques, passerelles des hôtes ; interfaces du routeur
(« Port Status : On »).}
\bareme{Sur le .pkt : interfaces routeur 1 pt ; IP/masques des hôtes 1 pt ; passerelles (imprimante et serveur compris)
1 pt. $-0,25$ par erreur, sans descendre sous 0 par item.}

\question{B4}{3 pts}{Réalisez les tests, notez la commande et le résultat, puis faites viser chaque test réussi.}
{\small\renewcommand{\arraystretch}{1.1}
\begin{tabularx}{\linewidth}{|c|l|l|C|c|c|}
  \hline
  \textbf{N°} & \textbf{Source} & \textbf{Destination} & \textbf{Commande / action} & \textbf{Résultat} & \textbf{Visa prof.} \\ \hline
  \hl 1 & PC-\pUn1 & PC-\pUn4 & \sol{ping \netUn.4} & \sol{réussi} & \hspace{1.4cm} \\ \hline
  \hl 2 & PC-\pDeux1 & IMP-\pDeux & \sol{ping \netDeux.0.10} & \sol{réussi} & \\ \hline
  \hl 3 & PC-\pUn2 & PC-\pDeux2 & \sol{ping \netDeux.0.2} & \sol{réussi} & \\ \hline
  \hl 4 & PC-\pDeux2 & \srvNom & \sol{ping \netSrv.\srvH} & \sol{réussi} & \\ \hline
  \hl 5 & PC-\pUn3 & page web du serveur & \sol{navigateur : {\NoAutoSpacing http://\netSrv.\srvH}} & \sol{page affichée} & \\ \hline
\end{tabularx}}
\bareme{0,5 pt par test visé + 0,5 pt pour B5. Le premier ping inter-réseaux peut échouer (résolution ARP) ; les suivants doivent réussir.}

\question{B5}{incl. B4}{Test n°3 : par quels équipements passe le paquet ? Pourquoi PC-\pUn2 l'envoie-t-il au routeur ?}
\zone{2.3cm}{PC-\pUn2 $\to$ SW-\pUn{} $\to$ \Rnom{} $\to$ SW-\pDeux{} $\to$ PC-\pDeux2. Grâce à son masque, PC-\pUn2 constate que
la destination (\netDeux.0.0/16) n'est \textbf{pas dans son réseau} (\netUn.0/24) : il envoie le paquet à sa
\textbf{passerelle par défaut}, qui le route.}

%==========================================================================
\newpage
\section{Partie C — Diagnostic d'un réseau en panne \hfill /4}

Un technicien constate qu'\textbf{aucun poste ne parvient à joindre le serveur SRV-C} et que certains postes
communiquent mal entre eux. Les câbles sont du bon type.

\begin{center}
\begin{minipage}{0.37\linewidth}\centering
\begin{tikzpicture}[font=\footnotesize,
  eq/.style={draw,fill=qbg,rounded corners=2pt,minimum width=9mm,minimum height=5mm,inner sep=2pt},
  every path/.style={thick}]
  \node[eq] (SW) {SW-1};
  \node[eq,right=9mm of SW] (R) {R1};
  \node[eq,right=7mm of R] (SRV) {SRV-C};
  \node[eq,above left=5mm and -3mm of SW] (A) {PC-A};
  \node[eq,above right=5mm and -3mm of SW] (B) {PC-B};
  \node[eq,below left=5mm and -3mm of SW] (C) {PC-C};
  \node[eq,below right=5mm and -3mm of SW] (D) {PC-D};
  \foreach \n in {A,B,C,D} \draw (\n) -- (SW);
  \draw (SW) -- node[above,font=\tiny]{G0/0} (R);
  \draw (R) -- node[above,font=\tiny]{G0/1} (SRV);
\end{tikzpicture}
\end{minipage}\hfill
\begin{minipage}{0.61\linewidth}\centering
{\small\renewcommand{\arraystretch}{1.05}\setlength{\tabcolsep}{4pt}
\begin{tabular}{lllll}
  \toprule
  \textbf{Équipement} & \textbf{Adresse IP} & \textbf{Masque} & \textbf{Passerelle} & \textbf{Port} \\
  \midrule
  \Crows
  \bottomrule
\end{tabular}}
\end{minipage}
\end{center}

\question{C1}{4 pts}{Relevez les \textbf{quatre anomalies}. Pour chacune : conséquence et correction.}
{\small\renewcommand{\arraystretch}{1.1}
\ifcorrige\def\hc{}\else\def\hc{\rule[-10mm]{0pt}{14mm}}\fi
\begin{tabularx}{\linewidth}{|c|>{\raggedright\arraybackslash}p{2cm}|L|L|L|}
  \hline
  \textbf{N°} & \textbf{Équipement} & \textbf{Anomalie} & \textbf{Conséquence} & \textbf{Correction} \\ \hline
  \ifcorrige\Csol\else
  \hc 1 & & & & \\ \hline
  \hc 2 & & & & \\ \hline
  \hc 3 & & & & \\ \hline
  \hc 4 & & & & \\ \hline
  \fi
\end{tabularx}}
\bareme{1 pt par anomalie (0,25 équipement, anomalie, conséquence, correction ; ordre indifférent). \Cbareme}

%==========================================================================
\section{Bonus — Extension sans fil \hfill +1}

\begin{cdc}Ajoutez au réseau \resUn{} un \textbf{AccessPoint-PT} relié à SW-\pUn{} (SSID \ip{\ssid}, \textbf{WPA2-PSK})
et un portable \ip{Laptop-\pUn} équipé d'un module \texttt{WPC300N}, en adressage \textbf{statique}.
Vérifiez qu'il joint \srvNom.\end{cdc}
{\small\renewcommand{\arraystretch}{1.1}
\begin{tabularx}{\linewidth}{|l|C|C|C|c|}
  \hline
  \textbf{Équipement} & \textbf{Adresse IP} & \textbf{Masque} & \textbf{Passerelle} & \textbf{Visa prof.} \\ \hline
  \hl Laptop-\pUn & \sol{\netUn.50} & \sol{255.255.255.0} & \sol{\netUn.254} & \hspace{1.4cm} \\ \hline
\end{tabularx}}
\bareme{Le point d'accès fait office de pont : le portable est dans \netUn.0/24. Bonus si associé en WPA2-PSK et ping
vers \netSrv.\srvH{} réussi. Note plafonnée à 20.}

\end{document}
"
\documentclass[10pt,a4paper]{article}

\usepackage[T1]{fontenc}
\usepackage[utf8]{inputenc}
\usepackage{lmodern}
\usepackage[french]{babel}
\usepackage[a4paper,hmargin=1.5cm,top=1.3cm,bottom=1.6cm,footskip=0.7cm]{geometry}
\usepackage[table]{xcolor}
\usepackage{booktabs,tabularx,array}
\usepackage{enumitem}
\usepackage{fancyhdr}
\usepackage[most]{tcolorbox}
\usepackage{tikz}
\usetikzlibrary{positioning}
\usepackage{titlesec}

\providecommand{\modecorrige}{0}
\providecommand{\variante}{A}
\newif\ifcorrige
\ifnum\modecorrige=1 \corrigetrue \else \corrigefalse \fi

\definecolor{ipred}{HTML}{C0392B}
\definecolor{cdcbar}{HTML}{A67C3D}
\definecolor{cdcbg}{HTML}{F4F1EA}
\definecolor{qbar}{HTML}{4A6B5D}
\definecolor{qbg}{HTML}{EDF1EF}
\definecolor{corr}{HTML}{1F5FA8}
\definecolor{corrbg}{HTML}{E8F0FA}
\definecolor{grisfoot}{HTML}{828282}

\newcommand{\ip}[1]{\textcolor{ipred}{\texttt{#1}}}
\newcommand{\pts}[1]{\hfill\mbox{\textbf{(#1)}}}
\newcommand{\sol}[1]{\ifcorrige\textcolor{corr}{#1}\fi}
\newcommand{\bareme}[1]{\ifcorrige\par{\footnotesize\color{corr}#1\par}\fi}
\newcommand{\hl}{\ifcorrige\rule[-1.6mm]{0pt}{5.4mm}\else\rule[-2.2mm]{0pt}{6.6mm}\fi}
\newcolumntype{C}{>{\centering\arraybackslash}X}
\newcolumntype{L}{>{\raggedright\arraybackslash}X}

\titleformat{\section}{\large\bfseries}{}{0pt}{}
\titlespacing*{\section}{0pt}{6pt}{3pt}

\tcbset{qstyle/.style={enhanced,boxrule=0pt,leftrule=2.5pt,arc=0pt,outer arc=0pt,
  left=5pt,right=5pt,top=2pt,bottom=2pt,before skip=4pt,after skip=3pt}}
\newtcolorbox{cdc}{qstyle,colback=cdcbg,colframe=cdcbar}
\newtcolorbox{quest}{qstyle,colback=qbg,colframe=qbar}

% Zone de réponse : vierge dans le sujet, remplie dans le corrigé
\newcommand{\zone}[2]{%
  \ifcorrige
    \begin{tcolorbox}[colback=corrbg,colframe=corr,boxrule=0.5pt,arc=1pt,left=4pt,right=4pt,top=1pt,bottom=1pt,
      before skip=2pt,after skip=2pt,fontupper=\footnotesize\color{corr}]#2\end{tcolorbox}%
  \else
    \begin{tcolorbox}[colback=white,colframe=black!45,boxrule=0.4pt,arc=1pt,height=#1,
      before skip=2pt,after skip=2pt]\end{tcolorbox}%
  \fi}

\newcommand{\question}[3]{\begin{quest}\textbf{#1.} #3\pts{#2}\end{quest}}

%==========================================================================
% Données propres à chaque variante
%==========================================================================
\if\variante A
  \def\entreprise{Le cabinet d'architectes \textbf{ArchiNova}}
  \def\Rnom{R-ArchiNova}
  \def\resUn{Bureau d'études}\def\pUn{BE}\def\netUn{192.168.40}
  \def\resDeux{Accueil}\def\pDeux{ACC}\def\netDeux{172.16}
  \def\srvNom{SRV-Intranet}\def\netSrv{192.168.200}\def\srvH{10}
  \def\ssid{ArchiNova-WiFi}
  \def\AIrows{%
    \hl\ip{192.168.40.25} & \sol{C} & \sol{255.255.255.0} & \sol{192.168.40.0} & \sol{192.168.40.255} & \sol{254} \\ \hline
    \hl\ip{172.16.5.10}   & \sol{B} & \sol{255.255.0.0}   & \sol{172.16.0.0}   & \sol{172.16.255.255} & \sol{65\,534} \\ \hline
    \hl\ip{10.3.2.1}      & \sol{A} & \sol{255.0.0.0}     & \sol{10.0.0.0}     & \sol{10.255.255.255} & \sol{16\,777\,214} \\ \hline}
  \def\apipa{169.254.12.7}
  \def\Cnet{192.168.60}
  \def\Crows{%
    PC-A & \ip{192.168.60.10} & \ip{255.255.255.0} & \ip{192.168.60.254} & — \\
    PC-B & \ip{192.168.60.20} & \ip{255.255.255.0} & \ip{192.168.60.1}   & — \\
    PC-C & \ip{192.168.61.30} & \ip{255.255.255.0} & \ip{192.168.60.254} & — \\
    PC-D & \ip{192.168.60.10} & \ip{255.255.255.0} & \ip{192.168.60.254} & — \\
    SRV-C & \ip{10.0.0.5}     & \ip{255.0.0.0}     & \ip{10.0.0.254}     & — \\
    R1 — G0/0 & \ip{192.168.60.254} & \ip{255.255.255.0} & — & On \\
    R1 — G0/1 & \ip{10.0.0.254}     & \ip{255.0.0.0}     & — & Off \\}
  \def\Csol{%
    \hl 1 & \sol{PC-B} & \sol{Passerelle .60.1 : pas l'adresse du routeur.} & \sol{Joint son réseau, mais aucun autre (serveur).} & \sol{Passerelle 192.168.60.254.} \\ \hline
    \hl 2 & \sol{PC-C} & \sol{IP hors du réseau 192.168.60.0/24.} & \sol{Ne joint ni les postes ni le routeur.} & \sol{Ex. 192.168.60.30.} \\ \hline
    \hl 3 & \sol{PC-D} & \sol{Même IP que PC-A (doublon).} & \sol{Conflit d'adresse : PC-A et PC-D perturbés.} & \sol{Adresse unique, ex. .60.40.} \\ \hline
    \hl 4 & \sol{R1 G0/1} & \sol{Interface désactivée (Off).} & \sol{Réseau du serveur injoignable pour tous.} & \sol{Port Status : On.} \\ \hline}
  \def\Cbareme{PC-A et SRV-C sont correctement configurés ; accepter PC-A cité à la place de PC-D pour le doublon.}
\else
  \def\entreprise{La clinique vétérinaire \textbf{VetoPlus}}
  \def\Rnom{R-VetoPlus}
  \def\resUn{Soins}\def\pUn{SOI}\def\netUn{192.168.70}
  \def\resDeux{Secrétariat}\def\pDeux{SEC}\def\netDeux{172.20}
  \def\srvNom{SRV-Dossiers}\def\netSrv{192.168.250}\def\srvH{20}
  \def\ssid{VetoPlus-WiFi}
  \def\AIrows{%
    \hl\ip{192.168.15.130} & \sol{C} & \sol{255.255.255.0} & \sol{192.168.15.0} & \sol{192.168.15.255} & \sol{254} \\ \hline
    \hl\ip{172.31.200.7}   & \sol{B} & \sol{255.255.0.0}   & \sol{172.31.0.0}   & \sol{172.31.255.255} & \sol{65\,534} \\ \hline
    \hl\ip{10.45.6.9}      & \sol{A} & \sol{255.0.0.0}     & \sol{10.0.0.0}     & \sol{10.255.255.255} & \sol{16\,777\,214} \\ \hline}
  \def\apipa{169.254.88.3}
  \def\Cnet{192.168.80}
  \def\Crows{%
    PC-A & \ip{192.168.8.11}   & \ip{255.255.255.0} & \ip{192.168.80.254} & — \\
    PC-B & \ip{192.168.80.12}  & \ip{255.255.255.0} & \ip{192.168.80.254} & — \\
    PC-C & \ip{192.168.80.13}  & \ip{255.255.255.0} & \ip{192.168.80.100} & — \\
    PC-D & \ip{192.168.80.12}  & \ip{255.255.255.0} & \ip{192.168.80.254} & — \\
    SRV-C & \ip{10.0.0.8}      & \ip{255.0.0.0}     & \ip{10.0.0.254}     & — \\
    R1 — G0/0 & \ip{192.168.80.254} & \ip{255.255.255.0} & — & Off \\
    R1 — G0/1 & \ip{10.0.0.254}     & \ip{255.0.0.0}     & — & On \\}
  \def\Csol{%
    \hl 1 & \sol{PC-A} & \sol{IP hors du réseau 192.168.80.0/24.} & \sol{Ne joint ni les postes ni le routeur.} & \sol{Ex. 192.168.80.11.} \\ \hline
    \hl 2 & \sol{PC-C} & \sol{Passerelle .80.100 : pas l'adresse du routeur.} & \sol{Joint son réseau, mais aucun autre (serveur).} & \sol{Passerelle 192.168.80.254.} \\ \hline
    \hl 3 & \sol{PC-D} & \sol{Même IP que PC-B (doublon).} & \sol{Conflit d'adresse : PC-B et PC-D perturbés.} & \sol{Adresse unique, ex. .80.14.} \\ \hline
    \hl 4 & \sol{R1 G0/0} & \sol{Interface côté LAN désactivée (Off).} & \sol{Aucun poste ne joint le routeur ni le serveur.} & \sol{Port Status : On.} \\ \hline}
  \def\Cbareme{PC-B et SRV-C sont correctement configurés ; accepter PC-B cité à la place de PC-D pour le doublon.}
\fi

\pagestyle{fancy}
\fancyhf{}
\renewcommand{\headrulewidth}{0pt}
\fancyfoot[L]{\footnotesize\itshape\color{grisfoot}Évaluation — Réseaux locaux — Sujet \variante\ifcorrige\ — CORRIGÉ\fi}
\fancyfoot[R]{\footnotesize\color{grisfoot}\thepage/3}

\setlist{itemsep=0pt,topsep=2pt,parsep=0pt}
\setlength{\parindent}{0pt}
\setlength{\parskip}{2pt}

\begin{document}

\begin{center}
  {\Large\bfseries Évaluation pratique — Réseaux locaux \quad
   \fcolorbox{black}{white}{\ SUJET \variante\ }\par}
  \vspace{1pt}
  {\bfseries Adressage IP, commutation et routage — BTS CIEL 1re année — durée : 2 heures\par}
  \ifcorrige\vspace{2pt}{\bfseries\color{corr}CORRIGÉ ET BARÈME — document enseignant}\par\fi
\end{center}
\vspace{-2pt}
{\renewcommand{\arraystretch}{1.35}
\begin{tabularx}{\linewidth}{|X|X|p{2.2cm}|p{2.6cm}|}
  \hline
  \textbf{Nom :} & \textbf{Prénom :} & \textbf{Poste n°} & \textbf{Note :} \hfill /20 \\ \hline
\end{tabularx}}

\textbf{Consignes :} travail \textbf{individuel}, documents et Internet \textbf{non autorisés}. Ce sujet est le
\textbf{document-réponse}. Fichier Packet Tracer : \ip{Nom\_Prenom\_EvalLAN\_\variante.pkt}, à enregistrer
régulièrement dans le dossier de rendu. Les « visas » sont apposés par le professeur, que vous appelez à chaque test réussi.
Durées conseillées : A 30 min — B 1 h 10 — C 20 min.

%==========================================================================
\section{Partie A — Questions de cours et calculs \hfill /6}

\question{A1}{2 pts}{Complétez le tableau en considérant le \textbf{masque par défaut} de la classe de chaque adresse.}
{\small\renewcommand{\arraystretch}{1.15}
\begin{tabularx}{\linewidth}{|l|c|C|C|C|c|}
  \hline
  \textbf{Adresse IP} & \textbf{Classe} & \textbf{Masque par défaut} & \textbf{Adresse réseau} &
  \textbf{Adresse de diffusion} & \textbf{Nb d'hôtes} \\ \hline
  \AIrows
\end{tabularx}}
\bareme{0,5 pt par ligne juste + 0,5 pt si les trois nombres d'hôtes sont justes ($2^{n}-2$, $n$ = nombre de bits d'hôte).}

\question{A2}{1 pt}{Indiquez le type de câble (\textbf{droit} ou \textbf{croisé}) pour chaque liaison.}
{\small\renewcommand{\arraystretch}{1.15}
\begin{tabularx}{\linewidth}{|X|c||X|c|}
  \hline
  \hl PC $\leftrightarrow$ commutateur & \makebox[1.6cm]{\sol{droit}} & Commutateur $\leftrightarrow$ routeur & \makebox[1.6cm]{\sol{droit}} \\ \hline
  \hl PC $\leftrightarrow$ PC (réseau P2P) & \sol{croisé} & Concentrateur $\leftrightarrow$ concentrateur & \sol{croisé} \\ \hline
\end{tabularx}}
\bareme{0,25 pt par case. Équipements de types différents : câble droit ; de même type : câble croisé.}

\question{A3}{1 pt}{Comment un \textbf{concentrateur (hub)} et un \textbf{commutateur (switch)} traitent-ils une trame
reçue ? Lequel limite les collisions, et pourquoi ?}
\zone{2.9cm}{Le \textbf{hub} (couche 1) répète la trame sur \textbf{tous ses autres ports} : un seul domaine de collision
partagé. Le \textbf{switch} (couche 2) lit l'adresse MAC de destination et, grâce à sa \textbf{table MAC}, ne l'envoie
\textbf{que sur le port du destinataire} : un domaine de collision par port, c'est lui qui limite les collisions.
\emph{0,5 pt hub, 0,5 pt switch.}}

\question{A4}{1 pt}{Qu'est-ce que la \textbf{passerelle par défaut} d'un poste ? Quelle condition son adresse doit-elle respecter ?}
\zone{2.4cm}{Adresse IP de l'interface du \textbf{routeur} à laquelle le poste envoie les paquets destinés à un
\textbf{autre réseau}. Elle doit \textbf{appartenir au même réseau} que le poste. \emph{0,5 + 0,5 pt.}}

\question{A5}{0,5 pt}{Un portable configuré en DHCP affiche \ip{\apipa} / \ip{255.255.0.0}. Que s'est-il passé ?}
\zone{1.5cm}{Aucun \textbf{serveur DHCP} n'a répondu : le poste s'est attribué une adresse \textbf{APIPA} (169.254.0.0/16).}

\question{A6}{0,5 pt}{Quel protocole utilise \texttt{ping} ? Nommez les deux messages échangés et la couche OSI des adresses IP.}
\zone{1.5cm}{\textbf{ICMP} ; \textbf{Echo Request} / \textbf{Echo Reply} ; adresses IP en \textbf{couche 3 (Réseau)}.}

%==========================================================================
\newpage
\section{Partie B — Réalisation sous Packet Tracer \hfill /10}

\begin{cdc}
\entreprise{} vous confie son réseau : \textbf{trois réseaux} reliés par un \textbf{routeur unique Cisco 2911}
nommé \ip{\Rnom} (interfaces G0/0, G0/1, G0/2).
\begin{itemize}[leftmargin=12pt]
  \item \textbf{Réseau \resUn} : \textbf{4 postes} PC-\pUn1 à PC-\pUn4, commutateur \ip{SW-\pUn}, adressage statique,
        réseau \ip{\netUn.0/24}.
  \item \textbf{Réseau \resDeux} : \textbf{2 postes} PC-\pDeux1, PC-\pDeux2 et \textbf{1 imprimante} IMP-\pDeux,
        commutateur \ip{SW-\pDeux}, adressage statique, réseau de classe B \ip{\netDeux.0.0/16}.
  \item \textbf{Réseau Serveur} : serveur \ip{\srvNom} d'adresse \ip{\netSrv.\srvH/24}, relié \textbf{directement} au routeur.
  \item La passerelle de chaque réseau est sa \textbf{dernière adresse utilisable}. Tous les équipements communiquent
        entre eux et accèdent à la page web du serveur. Nommage imposé (« Display Name »).
\end{itemize}
\end{cdc}

\question{B1}{2 pts}{Complétez le plan d'adressage \textbf{avant} la réalisation.}
{\footnotesize\setlength{\tabcolsep}{3pt}\renewcommand{\arraystretch}{1.15}
\begin{tabularx}{\linewidth}{|l|C|C|C|C|C|C|}
  \hline
  \textbf{Réseau} & \textbf{Adresse réseau} & \textbf{Masque} & \textbf{1re @ utilisable} &
  \textbf{Dernière @ utilisable} & \textbf{Diffusion} & \textbf{Passerelle} \\ \hline
  \hl \resUn & \sol{\netUn.0} & \sol{255.255.255.0} & \sol{\netUn.1} & \sol{\netUn.254} & \sol{\netUn.255} & \sol{\netUn.254} \\ \hline
  \hl \resDeux & \sol{\netDeux.0.0} & \sol{255.255.0.0} & \sol{\netDeux.0.1} & \sol{\netDeux.255.254} & \sol{\netDeux.255.255} & \sol{\netDeux.255.254} \\ \hline
  \hl Serveur & \sol{\netSrv.0} & \sol{255.255.255.0} & \sol{\netSrv.1} & \sol{\netSrv.254} & \sol{\netSrv.255} & \sol{\netSrv.254} \\ \hline
\end{tabularx}}
\bareme{0,5 pt par ligne juste + 0,5 pt si la convention de passerelle est respectée partout.}

\question{B2}{2 pts}{Construisez la topologie (équipements, câblage, nommage) et relevez l'adressage configuré.}
{\footnotesize\setlength{\tabcolsep}{3pt}\renewcommand{\arraystretch}{1.1}
\begin{tabularx}{\linewidth}{|l|C|C|C||l|C|C|}
  \hline
  \textbf{Hôte} & \textbf{Adresse IP} & \textbf{Masque} & \textbf{Passerelle} & \textbf{Interface} & \textbf{Adresse IP} & \textbf{Masque} \\ \hline
  \hl PC-\pUn1 & \sol{\netUn.1} & \sol{255.255.255.0} & \sol{\netUn.254} & \Rnom{} G0/0 & \sol{\netUn.254} & \sol{255.255.255.0} \\ \hline
  \hl PC-\pDeux1 & \sol{\netDeux.0.1} & \sol{255.255.0.0} & \sol{\netDeux.255.254} & \Rnom{} G0/1 & \sol{\netDeux.255.254} & \sol{255.255.0.0} \\ \hline
  \hl IMP-\pDeux & \sol{\netDeux.0.10} & \sol{255.255.0.0} & \sol{\netDeux.255.254} & \Rnom{} G0/2 & \sol{\netSrv.254} & \sol{255.255.255.0} \\ \hline
  \hl \srvNom & \ip{\netSrv.\srvH} & \sol{255.255.255.0} & \sol{\netSrv.254} & \multicolumn{3}{c|}{\cellcolor{black!8}} \\ \hline
\end{tabularx}}
\bareme{Adresses d'hôtes libres (uniques, dans le bon réseau) ; affectation G0/x libre. Sur le .pkt : équipements 0,5 ;
câblage 1 (câbles droits vers les switchs et switch–routeur, \textbf{croisé} serveur–routeur, liens verts) ; nommage 0,5.}

\question{B3}{3 pts}{Configurez tous les équipements : IP, masques, passerelles des hôtes ; interfaces du routeur
(« Port Status : On »).}
\bareme{Sur le .pkt : interfaces routeur 1 pt ; IP/masques des hôtes 1 pt ; passerelles (imprimante et serveur compris)
1 pt. $-0,25$ par erreur, sans descendre sous 0 par item.}

\question{B4}{3 pts}{Réalisez les tests, notez la commande et le résultat, puis faites viser chaque test réussi.}
{\small\renewcommand{\arraystretch}{1.1}
\begin{tabularx}{\linewidth}{|c|l|l|C|c|c|}
  \hline
  \textbf{N°} & \textbf{Source} & \textbf{Destination} & \textbf{Commande / action} & \textbf{Résultat} & \textbf{Visa prof.} \\ \hline
  \hl 1 & PC-\pUn1 & PC-\pUn4 & \sol{ping \netUn.4} & \sol{réussi} & \hspace{1.4cm} \\ \hline
  \hl 2 & PC-\pDeux1 & IMP-\pDeux & \sol{ping \netDeux.0.10} & \sol{réussi} & \\ \hline
  \hl 3 & PC-\pUn2 & PC-\pDeux2 & \sol{ping \netDeux.0.2} & \sol{réussi} & \\ \hline
  \hl 4 & PC-\pDeux2 & \srvNom & \sol{ping \netSrv.\srvH} & \sol{réussi} & \\ \hline
  \hl 5 & PC-\pUn3 & page web du serveur & \sol{navigateur : {\NoAutoSpacing http://\netSrv.\srvH}} & \sol{page affichée} & \\ \hline
\end{tabularx}}
\bareme{0,5 pt par test visé + 0,5 pt pour B5. Le premier ping inter-réseaux peut échouer (résolution ARP) ; les suivants doivent réussir.}

\question{B5}{incl. B4}{Test n°3 : par quels équipements passe le paquet ? Pourquoi PC-\pUn2 l'envoie-t-il au routeur ?}
\zone{2.3cm}{PC-\pUn2 $\to$ SW-\pUn{} $\to$ \Rnom{} $\to$ SW-\pDeux{} $\to$ PC-\pDeux2. Grâce à son masque, PC-\pUn2 constate que
la destination (\netDeux.0.0/16) n'est \textbf{pas dans son réseau} (\netUn.0/24) : il envoie le paquet à sa
\textbf{passerelle par défaut}, qui le route.}

%==========================================================================
\newpage
\section{Partie C — Diagnostic d'un réseau en panne \hfill /4}

Un technicien constate qu'\textbf{aucun poste ne parvient à joindre le serveur SRV-C} et que certains postes
communiquent mal entre eux. Les câbles sont du bon type.

\begin{center}
\begin{minipage}{0.37\linewidth}\centering
\begin{tikzpicture}[font=\footnotesize,
  eq/.style={draw,fill=qbg,rounded corners=2pt,minimum width=9mm,minimum height=5mm,inner sep=2pt},
  every path/.style={thick}]
  \node[eq] (SW) {SW-1};
  \node[eq,right=9mm of SW] (R) {R1};
  \node[eq,right=7mm of R] (SRV) {SRV-C};
  \node[eq,above left=5mm and -3mm of SW] (A) {PC-A};
  \node[eq,above right=5mm and -3mm of SW] (B) {PC-B};
  \node[eq,below left=5mm and -3mm of SW] (C) {PC-C};
  \node[eq,below right=5mm and -3mm of SW] (D) {PC-D};
  \foreach \n in {A,B,C,D} \draw (\n) -- (SW);
  \draw (SW) -- node[above,font=\tiny]{G0/0} (R);
  \draw (R) -- node[above,font=\tiny]{G0/1} (SRV);
\end{tikzpicture}
\end{minipage}\hfill
\begin{minipage}{0.61\linewidth}\centering
{\small\renewcommand{\arraystretch}{1.05}\setlength{\tabcolsep}{4pt}
\begin{tabular}{lllll}
  \toprule
  \textbf{Équipement} & \textbf{Adresse IP} & \textbf{Masque} & \textbf{Passerelle} & \textbf{Port} \\
  \midrule
  \Crows
  \bottomrule
\end{tabular}}
\end{minipage}
\end{center}

\question{C1}{4 pts}{Relevez les \textbf{quatre anomalies}. Pour chacune : conséquence et correction.}
{\small\renewcommand{\arraystretch}{1.1}
\ifcorrige\def\hc{}\else\def\hc{\rule[-10mm]{0pt}{14mm}}\fi
\begin{tabularx}{\linewidth}{|c|>{\raggedright\arraybackslash}p{2cm}|L|L|L|}
  \hline
  \textbf{N°} & \textbf{Équipement} & \textbf{Anomalie} & \textbf{Conséquence} & \textbf{Correction} \\ \hline
  \ifcorrige\Csol\else
  \hc 1 & & & & \\ \hline
  \hc 2 & & & & \\ \hline
  \hc 3 & & & & \\ \hline
  \hc 4 & & & & \\ \hline
  \fi
\end{tabularx}}
\bareme{1 pt par anomalie (0,25 équipement, anomalie, conséquence, correction ; ordre indifférent). \Cbareme}

%==========================================================================
\section{Bonus — Extension sans fil \hfill +1}

\begin{cdc}Ajoutez au réseau \resUn{} un \textbf{AccessPoint-PT} relié à SW-\pUn{} (SSID \ip{\ssid}, \textbf{WPA2-PSK})
et un portable \ip{Laptop-\pUn} équipé d'un module \texttt{WPC300N}, en adressage \textbf{statique}.
Vérifiez qu'il joint \srvNom.\end{cdc}
{\small\renewcommand{\arraystretch}{1.1}
\begin{tabularx}{\linewidth}{|l|C|C|C|c|}
  \hline
  \textbf{Équipement} & \textbf{Adresse IP} & \textbf{Masque} & \textbf{Passerelle} & \textbf{Visa prof.} \\ \hline
  \hl Laptop-\pUn & \sol{\netUn.50} & \sol{255.255.255.0} & \sol{\netUn.254} & \hspace{1.4cm} \\ \hline
\end{tabularx}}
\bareme{Le point d'accès fait office de pont : le portable est dans \netUn.0/24. Bonus si associé en WPA2-PSK et ping
vers \netSrv.\srvH{} réussi. Note plafonnée à 20.}

\end{document}
"
%   pdflatex -jobname=Eval_Reseau_SujetB   "\def\variante{B}\def\modecorrige{0}% Évaluation pratique — Réseaux locaux (BTS CIEL 1re année) — sujets A et B
% Compilation (4 PDF) :
%   pdflatex -jobname=Eval_Reseau_SujetA   "\def\variante{A}\def\modecorrige{0}% Évaluation pratique — Réseaux locaux (BTS CIEL 1re année) — sujets A et B
% Compilation (4 PDF) :
%   pdflatex -jobname=Eval_Reseau_SujetA   "\def\variante{A}\def\modecorrige{0}\input{Eval_TP_Reseau_LAN.tex}"
%   pdflatex -jobname=Eval_Reseau_SujetB   "\def\variante{B}\def\modecorrige{0}\input{Eval_TP_Reseau_LAN.tex}"
%   pdflatex -jobname=Eval_Reseau_CorrigeA "\def\variante{A}\def\modecorrige{1}\input{Eval_TP_Reseau_LAN.tex}"
%   pdflatex -jobname=Eval_Reseau_CorrigeB "\def\variante{B}\def\modecorrige{1}\input{Eval_TP_Reseau_LAN.tex}"
\documentclass[10pt,a4paper]{article}

\usepackage[T1]{fontenc}
\usepackage[utf8]{inputenc}
\usepackage{lmodern}
\usepackage[french]{babel}
\usepackage[a4paper,hmargin=1.5cm,top=1.3cm,bottom=1.6cm,footskip=0.7cm]{geometry}
\usepackage[table]{xcolor}
\usepackage{booktabs,tabularx,array}
\usepackage{enumitem}
\usepackage{fancyhdr}
\usepackage[most]{tcolorbox}
\usepackage{tikz}
\usetikzlibrary{positioning}
\usepackage{titlesec}

\providecommand{\modecorrige}{0}
\providecommand{\variante}{A}
\newif\ifcorrige
\ifnum\modecorrige=1 \corrigetrue \else \corrigefalse \fi

\definecolor{ipred}{HTML}{C0392B}
\definecolor{cdcbar}{HTML}{A67C3D}
\definecolor{cdcbg}{HTML}{F4F1EA}
\definecolor{qbar}{HTML}{4A6B5D}
\definecolor{qbg}{HTML}{EDF1EF}
\definecolor{corr}{HTML}{1F5FA8}
\definecolor{corrbg}{HTML}{E8F0FA}
\definecolor{grisfoot}{HTML}{828282}

\newcommand{\ip}[1]{\textcolor{ipred}{\texttt{#1}}}
\newcommand{\pts}[1]{\hfill\mbox{\textbf{(#1)}}}
\newcommand{\sol}[1]{\ifcorrige\textcolor{corr}{#1}\fi}
\newcommand{\bareme}[1]{\ifcorrige\par{\footnotesize\color{corr}#1\par}\fi}
\newcommand{\hl}{\ifcorrige\rule[-1.6mm]{0pt}{5.4mm}\else\rule[-2.2mm]{0pt}{6.6mm}\fi}
\newcolumntype{C}{>{\centering\arraybackslash}X}
\newcolumntype{L}{>{\raggedright\arraybackslash}X}

\titleformat{\section}{\large\bfseries}{}{0pt}{}
\titlespacing*{\section}{0pt}{6pt}{3pt}

\tcbset{qstyle/.style={enhanced,boxrule=0pt,leftrule=2.5pt,arc=0pt,outer arc=0pt,
  left=5pt,right=5pt,top=2pt,bottom=2pt,before skip=4pt,after skip=3pt}}
\newtcolorbox{cdc}{qstyle,colback=cdcbg,colframe=cdcbar}
\newtcolorbox{quest}{qstyle,colback=qbg,colframe=qbar}

% Zone de réponse : vierge dans le sujet, remplie dans le corrigé
\newcommand{\zone}[2]{%
  \ifcorrige
    \begin{tcolorbox}[colback=corrbg,colframe=corr,boxrule=0.5pt,arc=1pt,left=4pt,right=4pt,top=1pt,bottom=1pt,
      before skip=2pt,after skip=2pt,fontupper=\footnotesize\color{corr}]#2\end{tcolorbox}%
  \else
    \begin{tcolorbox}[colback=white,colframe=black!45,boxrule=0.4pt,arc=1pt,height=#1,
      before skip=2pt,after skip=2pt]\end{tcolorbox}%
  \fi}

\newcommand{\question}[3]{\begin{quest}\textbf{#1.} #3\pts{#2}\end{quest}}

%==========================================================================
% Données propres à chaque variante
%==========================================================================
\if\variante A
  \def\entreprise{Le cabinet d'architectes \textbf{ArchiNova}}
  \def\Rnom{R-ArchiNova}
  \def\resUn{Bureau d'études}\def\pUn{BE}\def\netUn{192.168.40}
  \def\resDeux{Accueil}\def\pDeux{ACC}\def\netDeux{172.16}
  \def\srvNom{SRV-Intranet}\def\netSrv{192.168.200}\def\srvH{10}
  \def\ssid{ArchiNova-WiFi}
  \def\AIrows{%
    \hl\ip{192.168.40.25} & \sol{C} & \sol{255.255.255.0} & \sol{192.168.40.0} & \sol{192.168.40.255} & \sol{254} \\ \hline
    \hl\ip{172.16.5.10}   & \sol{B} & \sol{255.255.0.0}   & \sol{172.16.0.0}   & \sol{172.16.255.255} & \sol{65\,534} \\ \hline
    \hl\ip{10.3.2.1}      & \sol{A} & \sol{255.0.0.0}     & \sol{10.0.0.0}     & \sol{10.255.255.255} & \sol{16\,777\,214} \\ \hline}
  \def\apipa{169.254.12.7}
  \def\Cnet{192.168.60}
  \def\Crows{%
    PC-A & \ip{192.168.60.10} & \ip{255.255.255.0} & \ip{192.168.60.254} & — \\
    PC-B & \ip{192.168.60.20} & \ip{255.255.255.0} & \ip{192.168.60.1}   & — \\
    PC-C & \ip{192.168.61.30} & \ip{255.255.255.0} & \ip{192.168.60.254} & — \\
    PC-D & \ip{192.168.60.10} & \ip{255.255.255.0} & \ip{192.168.60.254} & — \\
    SRV-C & \ip{10.0.0.5}     & \ip{255.0.0.0}     & \ip{10.0.0.254}     & — \\
    R1 — G0/0 & \ip{192.168.60.254} & \ip{255.255.255.0} & — & On \\
    R1 — G0/1 & \ip{10.0.0.254}     & \ip{255.0.0.0}     & — & Off \\}
  \def\Csol{%
    \hl 1 & \sol{PC-B} & \sol{Passerelle .60.1 : pas l'adresse du routeur.} & \sol{Joint son réseau, mais aucun autre (serveur).} & \sol{Passerelle 192.168.60.254.} \\ \hline
    \hl 2 & \sol{PC-C} & \sol{IP hors du réseau 192.168.60.0/24.} & \sol{Ne joint ni les postes ni le routeur.} & \sol{Ex. 192.168.60.30.} \\ \hline
    \hl 3 & \sol{PC-D} & \sol{Même IP que PC-A (doublon).} & \sol{Conflit d'adresse : PC-A et PC-D perturbés.} & \sol{Adresse unique, ex. .60.40.} \\ \hline
    \hl 4 & \sol{R1 G0/1} & \sol{Interface désactivée (Off).} & \sol{Réseau du serveur injoignable pour tous.} & \sol{Port Status : On.} \\ \hline}
  \def\Cbareme{PC-A et SRV-C sont correctement configurés ; accepter PC-A cité à la place de PC-D pour le doublon.}
\else
  \def\entreprise{La clinique vétérinaire \textbf{VetoPlus}}
  \def\Rnom{R-VetoPlus}
  \def\resUn{Soins}\def\pUn{SOI}\def\netUn{192.168.70}
  \def\resDeux{Secrétariat}\def\pDeux{SEC}\def\netDeux{172.20}
  \def\srvNom{SRV-Dossiers}\def\netSrv{192.168.250}\def\srvH{20}
  \def\ssid{VetoPlus-WiFi}
  \def\AIrows{%
    \hl\ip{192.168.15.130} & \sol{C} & \sol{255.255.255.0} & \sol{192.168.15.0} & \sol{192.168.15.255} & \sol{254} \\ \hline
    \hl\ip{172.31.200.7}   & \sol{B} & \sol{255.255.0.0}   & \sol{172.31.0.0}   & \sol{172.31.255.255} & \sol{65\,534} \\ \hline
    \hl\ip{10.45.6.9}      & \sol{A} & \sol{255.0.0.0}     & \sol{10.0.0.0}     & \sol{10.255.255.255} & \sol{16\,777\,214} \\ \hline}
  \def\apipa{169.254.88.3}
  \def\Cnet{192.168.80}
  \def\Crows{%
    PC-A & \ip{192.168.8.11}   & \ip{255.255.255.0} & \ip{192.168.80.254} & — \\
    PC-B & \ip{192.168.80.12}  & \ip{255.255.255.0} & \ip{192.168.80.254} & — \\
    PC-C & \ip{192.168.80.13}  & \ip{255.255.255.0} & \ip{192.168.80.100} & — \\
    PC-D & \ip{192.168.80.12}  & \ip{255.255.255.0} & \ip{192.168.80.254} & — \\
    SRV-C & \ip{10.0.0.8}      & \ip{255.0.0.0}     & \ip{10.0.0.254}     & — \\
    R1 — G0/0 & \ip{192.168.80.254} & \ip{255.255.255.0} & — & Off \\
    R1 — G0/1 & \ip{10.0.0.254}     & \ip{255.0.0.0}     & — & On \\}
  \def\Csol{%
    \hl 1 & \sol{PC-A} & \sol{IP hors du réseau 192.168.80.0/24.} & \sol{Ne joint ni les postes ni le routeur.} & \sol{Ex. 192.168.80.11.} \\ \hline
    \hl 2 & \sol{PC-C} & \sol{Passerelle .80.100 : pas l'adresse du routeur.} & \sol{Joint son réseau, mais aucun autre (serveur).} & \sol{Passerelle 192.168.80.254.} \\ \hline
    \hl 3 & \sol{PC-D} & \sol{Même IP que PC-B (doublon).} & \sol{Conflit d'adresse : PC-B et PC-D perturbés.} & \sol{Adresse unique, ex. .80.14.} \\ \hline
    \hl 4 & \sol{R1 G0/0} & \sol{Interface côté LAN désactivée (Off).} & \sol{Aucun poste ne joint le routeur ni le serveur.} & \sol{Port Status : On.} \\ \hline}
  \def\Cbareme{PC-B et SRV-C sont correctement configurés ; accepter PC-B cité à la place de PC-D pour le doublon.}
\fi

\pagestyle{fancy}
\fancyhf{}
\renewcommand{\headrulewidth}{0pt}
\fancyfoot[L]{\footnotesize\itshape\color{grisfoot}Évaluation — Réseaux locaux — Sujet \variante\ifcorrige\ — CORRIGÉ\fi}
\fancyfoot[R]{\footnotesize\color{grisfoot}\thepage/3}

\setlist{itemsep=0pt,topsep=2pt,parsep=0pt}
\setlength{\parindent}{0pt}
\setlength{\parskip}{2pt}

\begin{document}

\begin{center}
  {\Large\bfseries Évaluation pratique — Réseaux locaux \quad
   \fcolorbox{black}{white}{\ SUJET \variante\ }\par}
  \vspace{1pt}
  {\bfseries Adressage IP, commutation et routage — BTS CIEL 1re année — durée : 2 heures\par}
  \ifcorrige\vspace{2pt}{\bfseries\color{corr}CORRIGÉ ET BARÈME — document enseignant}\par\fi
\end{center}
\vspace{-2pt}
{\renewcommand{\arraystretch}{1.35}
\begin{tabularx}{\linewidth}{|X|X|p{2.2cm}|p{2.6cm}|}
  \hline
  \textbf{Nom :} & \textbf{Prénom :} & \textbf{Poste n°} & \textbf{Note :} \hfill /20 \\ \hline
\end{tabularx}}

\textbf{Consignes :} travail \textbf{individuel}, documents et Internet \textbf{non autorisés}. Ce sujet est le
\textbf{document-réponse}. Fichier Packet Tracer : \ip{Nom\_Prenom\_EvalLAN\_\variante.pkt}, à enregistrer
régulièrement dans le dossier de rendu. Les « visas » sont apposés par le professeur, que vous appelez à chaque test réussi.
Durées conseillées : A 30 min — B 1 h 10 — C 20 min.

%==========================================================================
\section{Partie A — Questions de cours et calculs \hfill /6}

\question{A1}{2 pts}{Complétez le tableau en considérant le \textbf{masque par défaut} de la classe de chaque adresse.}
{\small\renewcommand{\arraystretch}{1.15}
\begin{tabularx}{\linewidth}{|l|c|C|C|C|c|}
  \hline
  \textbf{Adresse IP} & \textbf{Classe} & \textbf{Masque par défaut} & \textbf{Adresse réseau} &
  \textbf{Adresse de diffusion} & \textbf{Nb d'hôtes} \\ \hline
  \AIrows
\end{tabularx}}
\bareme{0,5 pt par ligne juste + 0,5 pt si les trois nombres d'hôtes sont justes ($2^{n}-2$, $n$ = nombre de bits d'hôte).}

\question{A2}{1 pt}{Indiquez le type de câble (\textbf{droit} ou \textbf{croisé}) pour chaque liaison.}
{\small\renewcommand{\arraystretch}{1.15}
\begin{tabularx}{\linewidth}{|X|c||X|c|}
  \hline
  \hl PC $\leftrightarrow$ commutateur & \makebox[1.6cm]{\sol{droit}} & Commutateur $\leftrightarrow$ routeur & \makebox[1.6cm]{\sol{droit}} \\ \hline
  \hl PC $\leftrightarrow$ PC (réseau P2P) & \sol{croisé} & Concentrateur $\leftrightarrow$ concentrateur & \sol{croisé} \\ \hline
\end{tabularx}}
\bareme{0,25 pt par case. Équipements de types différents : câble droit ; de même type : câble croisé.}

\question{A3}{1 pt}{Comment un \textbf{concentrateur (hub)} et un \textbf{commutateur (switch)} traitent-ils une trame
reçue ? Lequel limite les collisions, et pourquoi ?}
\zone{2.9cm}{Le \textbf{hub} (couche 1) répète la trame sur \textbf{tous ses autres ports} : un seul domaine de collision
partagé. Le \textbf{switch} (couche 2) lit l'adresse MAC de destination et, grâce à sa \textbf{table MAC}, ne l'envoie
\textbf{que sur le port du destinataire} : un domaine de collision par port, c'est lui qui limite les collisions.
\emph{0,5 pt hub, 0,5 pt switch.}}

\question{A4}{1 pt}{Qu'est-ce que la \textbf{passerelle par défaut} d'un poste ? Quelle condition son adresse doit-elle respecter ?}
\zone{2.4cm}{Adresse IP de l'interface du \textbf{routeur} à laquelle le poste envoie les paquets destinés à un
\textbf{autre réseau}. Elle doit \textbf{appartenir au même réseau} que le poste. \emph{0,5 + 0,5 pt.}}

\question{A5}{0,5 pt}{Un portable configuré en DHCP affiche \ip{\apipa} / \ip{255.255.0.0}. Que s'est-il passé ?}
\zone{1.5cm}{Aucun \textbf{serveur DHCP} n'a répondu : le poste s'est attribué une adresse \textbf{APIPA} (169.254.0.0/16).}

\question{A6}{0,5 pt}{Quel protocole utilise \texttt{ping} ? Nommez les deux messages échangés et la couche OSI des adresses IP.}
\zone{1.5cm}{\textbf{ICMP} ; \textbf{Echo Request} / \textbf{Echo Reply} ; adresses IP en \textbf{couche 3 (Réseau)}.}

%==========================================================================
\newpage
\section{Partie B — Réalisation sous Packet Tracer \hfill /10}

\begin{cdc}
\entreprise{} vous confie son réseau : \textbf{trois réseaux} reliés par un \textbf{routeur unique Cisco 2911}
nommé \ip{\Rnom} (interfaces G0/0, G0/1, G0/2).
\begin{itemize}[leftmargin=12pt]
  \item \textbf{Réseau \resUn} : \textbf{4 postes} PC-\pUn1 à PC-\pUn4, commutateur \ip{SW-\pUn}, adressage statique,
        réseau \ip{\netUn.0/24}.
  \item \textbf{Réseau \resDeux} : \textbf{2 postes} PC-\pDeux1, PC-\pDeux2 et \textbf{1 imprimante} IMP-\pDeux,
        commutateur \ip{SW-\pDeux}, adressage statique, réseau de classe B \ip{\netDeux.0.0/16}.
  \item \textbf{Réseau Serveur} : serveur \ip{\srvNom} d'adresse \ip{\netSrv.\srvH/24}, relié \textbf{directement} au routeur.
  \item La passerelle de chaque réseau est sa \textbf{dernière adresse utilisable}. Tous les équipements communiquent
        entre eux et accèdent à la page web du serveur. Nommage imposé (« Display Name »).
\end{itemize}
\end{cdc}

\question{B1}{2 pts}{Complétez le plan d'adressage \textbf{avant} la réalisation.}
{\footnotesize\setlength{\tabcolsep}{3pt}\renewcommand{\arraystretch}{1.15}
\begin{tabularx}{\linewidth}{|l|C|C|C|C|C|C|}
  \hline
  \textbf{Réseau} & \textbf{Adresse réseau} & \textbf{Masque} & \textbf{1re @ utilisable} &
  \textbf{Dernière @ utilisable} & \textbf{Diffusion} & \textbf{Passerelle} \\ \hline
  \hl \resUn & \sol{\netUn.0} & \sol{255.255.255.0} & \sol{\netUn.1} & \sol{\netUn.254} & \sol{\netUn.255} & \sol{\netUn.254} \\ \hline
  \hl \resDeux & \sol{\netDeux.0.0} & \sol{255.255.0.0} & \sol{\netDeux.0.1} & \sol{\netDeux.255.254} & \sol{\netDeux.255.255} & \sol{\netDeux.255.254} \\ \hline
  \hl Serveur & \sol{\netSrv.0} & \sol{255.255.255.0} & \sol{\netSrv.1} & \sol{\netSrv.254} & \sol{\netSrv.255} & \sol{\netSrv.254} \\ \hline
\end{tabularx}}
\bareme{0,5 pt par ligne juste + 0,5 pt si la convention de passerelle est respectée partout.}

\question{B2}{2 pts}{Construisez la topologie (équipements, câblage, nommage) et relevez l'adressage configuré.}
{\footnotesize\setlength{\tabcolsep}{3pt}\renewcommand{\arraystretch}{1.1}
\begin{tabularx}{\linewidth}{|l|C|C|C||l|C|C|}
  \hline
  \textbf{Hôte} & \textbf{Adresse IP} & \textbf{Masque} & \textbf{Passerelle} & \textbf{Interface} & \textbf{Adresse IP} & \textbf{Masque} \\ \hline
  \hl PC-\pUn1 & \sol{\netUn.1} & \sol{255.255.255.0} & \sol{\netUn.254} & \Rnom{} G0/0 & \sol{\netUn.254} & \sol{255.255.255.0} \\ \hline
  \hl PC-\pDeux1 & \sol{\netDeux.0.1} & \sol{255.255.0.0} & \sol{\netDeux.255.254} & \Rnom{} G0/1 & \sol{\netDeux.255.254} & \sol{255.255.0.0} \\ \hline
  \hl IMP-\pDeux & \sol{\netDeux.0.10} & \sol{255.255.0.0} & \sol{\netDeux.255.254} & \Rnom{} G0/2 & \sol{\netSrv.254} & \sol{255.255.255.0} \\ \hline
  \hl \srvNom & \ip{\netSrv.\srvH} & \sol{255.255.255.0} & \sol{\netSrv.254} & \multicolumn{3}{c|}{\cellcolor{black!8}} \\ \hline
\end{tabularx}}
\bareme{Adresses d'hôtes libres (uniques, dans le bon réseau) ; affectation G0/x libre. Sur le .pkt : équipements 0,5 ;
câblage 1 (câbles droits vers les switchs et switch–routeur, \textbf{croisé} serveur–routeur, liens verts) ; nommage 0,5.}

\question{B3}{3 pts}{Configurez tous les équipements : IP, masques, passerelles des hôtes ; interfaces du routeur
(« Port Status : On »).}
\bareme{Sur le .pkt : interfaces routeur 1 pt ; IP/masques des hôtes 1 pt ; passerelles (imprimante et serveur compris)
1 pt. $-0,25$ par erreur, sans descendre sous 0 par item.}

\question{B4}{3 pts}{Réalisez les tests, notez la commande et le résultat, puis faites viser chaque test réussi.}
{\small\renewcommand{\arraystretch}{1.1}
\begin{tabularx}{\linewidth}{|c|l|l|C|c|c|}
  \hline
  \textbf{N°} & \textbf{Source} & \textbf{Destination} & \textbf{Commande / action} & \textbf{Résultat} & \textbf{Visa prof.} \\ \hline
  \hl 1 & PC-\pUn1 & PC-\pUn4 & \sol{ping \netUn.4} & \sol{réussi} & \hspace{1.4cm} \\ \hline
  \hl 2 & PC-\pDeux1 & IMP-\pDeux & \sol{ping \netDeux.0.10} & \sol{réussi} & \\ \hline
  \hl 3 & PC-\pUn2 & PC-\pDeux2 & \sol{ping \netDeux.0.2} & \sol{réussi} & \\ \hline
  \hl 4 & PC-\pDeux2 & \srvNom & \sol{ping \netSrv.\srvH} & \sol{réussi} & \\ \hline
  \hl 5 & PC-\pUn3 & page web du serveur & \sol{navigateur : {\NoAutoSpacing http://\netSrv.\srvH}} & \sol{page affichée} & \\ \hline
\end{tabularx}}
\bareme{0,5 pt par test visé + 0,5 pt pour B5. Le premier ping inter-réseaux peut échouer (résolution ARP) ; les suivants doivent réussir.}

\question{B5}{incl. B4}{Test n°3 : par quels équipements passe le paquet ? Pourquoi PC-\pUn2 l'envoie-t-il au routeur ?}
\zone{2.3cm}{PC-\pUn2 $\to$ SW-\pUn{} $\to$ \Rnom{} $\to$ SW-\pDeux{} $\to$ PC-\pDeux2. Grâce à son masque, PC-\pUn2 constate que
la destination (\netDeux.0.0/16) n'est \textbf{pas dans son réseau} (\netUn.0/24) : il envoie le paquet à sa
\textbf{passerelle par défaut}, qui le route.}

%==========================================================================
\newpage
\section{Partie C — Diagnostic d'un réseau en panne \hfill /4}

Un technicien constate qu'\textbf{aucun poste ne parvient à joindre le serveur SRV-C} et que certains postes
communiquent mal entre eux. Les câbles sont du bon type.

\begin{center}
\begin{minipage}{0.37\linewidth}\centering
\begin{tikzpicture}[font=\footnotesize,
  eq/.style={draw,fill=qbg,rounded corners=2pt,minimum width=9mm,minimum height=5mm,inner sep=2pt},
  every path/.style={thick}]
  \node[eq] (SW) {SW-1};
  \node[eq,right=9mm of SW] (R) {R1};
  \node[eq,right=7mm of R] (SRV) {SRV-C};
  \node[eq,above left=5mm and -3mm of SW] (A) {PC-A};
  \node[eq,above right=5mm and -3mm of SW] (B) {PC-B};
  \node[eq,below left=5mm and -3mm of SW] (C) {PC-C};
  \node[eq,below right=5mm and -3mm of SW] (D) {PC-D};
  \foreach \n in {A,B,C,D} \draw (\n) -- (SW);
  \draw (SW) -- node[above,font=\tiny]{G0/0} (R);
  \draw (R) -- node[above,font=\tiny]{G0/1} (SRV);
\end{tikzpicture}
\end{minipage}\hfill
\begin{minipage}{0.61\linewidth}\centering
{\small\renewcommand{\arraystretch}{1.05}\setlength{\tabcolsep}{4pt}
\begin{tabular}{lllll}
  \toprule
  \textbf{Équipement} & \textbf{Adresse IP} & \textbf{Masque} & \textbf{Passerelle} & \textbf{Port} \\
  \midrule
  \Crows
  \bottomrule
\end{tabular}}
\end{minipage}
\end{center}

\question{C1}{4 pts}{Relevez les \textbf{quatre anomalies}. Pour chacune : conséquence et correction.}
{\small\renewcommand{\arraystretch}{1.1}
\ifcorrige\def\hc{}\else\def\hc{\rule[-10mm]{0pt}{14mm}}\fi
\begin{tabularx}{\linewidth}{|c|>{\raggedright\arraybackslash}p{2cm}|L|L|L|}
  \hline
  \textbf{N°} & \textbf{Équipement} & \textbf{Anomalie} & \textbf{Conséquence} & \textbf{Correction} \\ \hline
  \ifcorrige\Csol\else
  \hc 1 & & & & \\ \hline
  \hc 2 & & & & \\ \hline
  \hc 3 & & & & \\ \hline
  \hc 4 & & & & \\ \hline
  \fi
\end{tabularx}}
\bareme{1 pt par anomalie (0,25 équipement, anomalie, conséquence, correction ; ordre indifférent). \Cbareme}

%==========================================================================
\section{Bonus — Extension sans fil \hfill +1}

\begin{cdc}Ajoutez au réseau \resUn{} un \textbf{AccessPoint-PT} relié à SW-\pUn{} (SSID \ip{\ssid}, \textbf{WPA2-PSK})
et un portable \ip{Laptop-\pUn} équipé d'un module \texttt{WPC300N}, en adressage \textbf{statique}.
Vérifiez qu'il joint \srvNom.\end{cdc}
{\small\renewcommand{\arraystretch}{1.1}
\begin{tabularx}{\linewidth}{|l|C|C|C|c|}
  \hline
  \textbf{Équipement} & \textbf{Adresse IP} & \textbf{Masque} & \textbf{Passerelle} & \textbf{Visa prof.} \\ \hline
  \hl Laptop-\pUn & \sol{\netUn.50} & \sol{255.255.255.0} & \sol{\netUn.254} & \hspace{1.4cm} \\ \hline
\end{tabularx}}
\bareme{Le point d'accès fait office de pont : le portable est dans \netUn.0/24. Bonus si associé en WPA2-PSK et ping
vers \netSrv.\srvH{} réussi. Note plafonnée à 20.}

\end{document}
"
%   pdflatex -jobname=Eval_Reseau_SujetB   "\def\variante{B}\def\modecorrige{0}% Évaluation pratique — Réseaux locaux (BTS CIEL 1re année) — sujets A et B
% Compilation (4 PDF) :
%   pdflatex -jobname=Eval_Reseau_SujetA   "\def\variante{A}\def\modecorrige{0}\input{Eval_TP_Reseau_LAN.tex}"
%   pdflatex -jobname=Eval_Reseau_SujetB   "\def\variante{B}\def\modecorrige{0}\input{Eval_TP_Reseau_LAN.tex}"
%   pdflatex -jobname=Eval_Reseau_CorrigeA "\def\variante{A}\def\modecorrige{1}\input{Eval_TP_Reseau_LAN.tex}"
%   pdflatex -jobname=Eval_Reseau_CorrigeB "\def\variante{B}\def\modecorrige{1}\input{Eval_TP_Reseau_LAN.tex}"
\documentclass[10pt,a4paper]{article}

\usepackage[T1]{fontenc}
\usepackage[utf8]{inputenc}
\usepackage{lmodern}
\usepackage[french]{babel}
\usepackage[a4paper,hmargin=1.5cm,top=1.3cm,bottom=1.6cm,footskip=0.7cm]{geometry}
\usepackage[table]{xcolor}
\usepackage{booktabs,tabularx,array}
\usepackage{enumitem}
\usepackage{fancyhdr}
\usepackage[most]{tcolorbox}
\usepackage{tikz}
\usetikzlibrary{positioning}
\usepackage{titlesec}

\providecommand{\modecorrige}{0}
\providecommand{\variante}{A}
\newif\ifcorrige
\ifnum\modecorrige=1 \corrigetrue \else \corrigefalse \fi

\definecolor{ipred}{HTML}{C0392B}
\definecolor{cdcbar}{HTML}{A67C3D}
\definecolor{cdcbg}{HTML}{F4F1EA}
\definecolor{qbar}{HTML}{4A6B5D}
\definecolor{qbg}{HTML}{EDF1EF}
\definecolor{corr}{HTML}{1F5FA8}
\definecolor{corrbg}{HTML}{E8F0FA}
\definecolor{grisfoot}{HTML}{828282}

\newcommand{\ip}[1]{\textcolor{ipred}{\texttt{#1}}}
\newcommand{\pts}[1]{\hfill\mbox{\textbf{(#1)}}}
\newcommand{\sol}[1]{\ifcorrige\textcolor{corr}{#1}\fi}
\newcommand{\bareme}[1]{\ifcorrige\par{\footnotesize\color{corr}#1\par}\fi}
\newcommand{\hl}{\ifcorrige\rule[-1.6mm]{0pt}{5.4mm}\else\rule[-2.2mm]{0pt}{6.6mm}\fi}
\newcolumntype{C}{>{\centering\arraybackslash}X}
\newcolumntype{L}{>{\raggedright\arraybackslash}X}

\titleformat{\section}{\large\bfseries}{}{0pt}{}
\titlespacing*{\section}{0pt}{6pt}{3pt}

\tcbset{qstyle/.style={enhanced,boxrule=0pt,leftrule=2.5pt,arc=0pt,outer arc=0pt,
  left=5pt,right=5pt,top=2pt,bottom=2pt,before skip=4pt,after skip=3pt}}
\newtcolorbox{cdc}{qstyle,colback=cdcbg,colframe=cdcbar}
\newtcolorbox{quest}{qstyle,colback=qbg,colframe=qbar}

% Zone de réponse : vierge dans le sujet, remplie dans le corrigé
\newcommand{\zone}[2]{%
  \ifcorrige
    \begin{tcolorbox}[colback=corrbg,colframe=corr,boxrule=0.5pt,arc=1pt,left=4pt,right=4pt,top=1pt,bottom=1pt,
      before skip=2pt,after skip=2pt,fontupper=\footnotesize\color{corr}]#2\end{tcolorbox}%
  \else
    \begin{tcolorbox}[colback=white,colframe=black!45,boxrule=0.4pt,arc=1pt,height=#1,
      before skip=2pt,after skip=2pt]\end{tcolorbox}%
  \fi}

\newcommand{\question}[3]{\begin{quest}\textbf{#1.} #3\pts{#2}\end{quest}}

%==========================================================================
% Données propres à chaque variante
%==========================================================================
\if\variante A
  \def\entreprise{Le cabinet d'architectes \textbf{ArchiNova}}
  \def\Rnom{R-ArchiNova}
  \def\resUn{Bureau d'études}\def\pUn{BE}\def\netUn{192.168.40}
  \def\resDeux{Accueil}\def\pDeux{ACC}\def\netDeux{172.16}
  \def\srvNom{SRV-Intranet}\def\netSrv{192.168.200}\def\srvH{10}
  \def\ssid{ArchiNova-WiFi}
  \def\AIrows{%
    \hl\ip{192.168.40.25} & \sol{C} & \sol{255.255.255.0} & \sol{192.168.40.0} & \sol{192.168.40.255} & \sol{254} \\ \hline
    \hl\ip{172.16.5.10}   & \sol{B} & \sol{255.255.0.0}   & \sol{172.16.0.0}   & \sol{172.16.255.255} & \sol{65\,534} \\ \hline
    \hl\ip{10.3.2.1}      & \sol{A} & \sol{255.0.0.0}     & \sol{10.0.0.0}     & \sol{10.255.255.255} & \sol{16\,777\,214} \\ \hline}
  \def\apipa{169.254.12.7}
  \def\Cnet{192.168.60}
  \def\Crows{%
    PC-A & \ip{192.168.60.10} & \ip{255.255.255.0} & \ip{192.168.60.254} & — \\
    PC-B & \ip{192.168.60.20} & \ip{255.255.255.0} & \ip{192.168.60.1}   & — \\
    PC-C & \ip{192.168.61.30} & \ip{255.255.255.0} & \ip{192.168.60.254} & — \\
    PC-D & \ip{192.168.60.10} & \ip{255.255.255.0} & \ip{192.168.60.254} & — \\
    SRV-C & \ip{10.0.0.5}     & \ip{255.0.0.0}     & \ip{10.0.0.254}     & — \\
    R1 — G0/0 & \ip{192.168.60.254} & \ip{255.255.255.0} & — & On \\
    R1 — G0/1 & \ip{10.0.0.254}     & \ip{255.0.0.0}     & — & Off \\}
  \def\Csol{%
    \hl 1 & \sol{PC-B} & \sol{Passerelle .60.1 : pas l'adresse du routeur.} & \sol{Joint son réseau, mais aucun autre (serveur).} & \sol{Passerelle 192.168.60.254.} \\ \hline
    \hl 2 & \sol{PC-C} & \sol{IP hors du réseau 192.168.60.0/24.} & \sol{Ne joint ni les postes ni le routeur.} & \sol{Ex. 192.168.60.30.} \\ \hline
    \hl 3 & \sol{PC-D} & \sol{Même IP que PC-A (doublon).} & \sol{Conflit d'adresse : PC-A et PC-D perturbés.} & \sol{Adresse unique, ex. .60.40.} \\ \hline
    \hl 4 & \sol{R1 G0/1} & \sol{Interface désactivée (Off).} & \sol{Réseau du serveur injoignable pour tous.} & \sol{Port Status : On.} \\ \hline}
  \def\Cbareme{PC-A et SRV-C sont correctement configurés ; accepter PC-A cité à la place de PC-D pour le doublon.}
\else
  \def\entreprise{La clinique vétérinaire \textbf{VetoPlus}}
  \def\Rnom{R-VetoPlus}
  \def\resUn{Soins}\def\pUn{SOI}\def\netUn{192.168.70}
  \def\resDeux{Secrétariat}\def\pDeux{SEC}\def\netDeux{172.20}
  \def\srvNom{SRV-Dossiers}\def\netSrv{192.168.250}\def\srvH{20}
  \def\ssid{VetoPlus-WiFi}
  \def\AIrows{%
    \hl\ip{192.168.15.130} & \sol{C} & \sol{255.255.255.0} & \sol{192.168.15.0} & \sol{192.168.15.255} & \sol{254} \\ \hline
    \hl\ip{172.31.200.7}   & \sol{B} & \sol{255.255.0.0}   & \sol{172.31.0.0}   & \sol{172.31.255.255} & \sol{65\,534} \\ \hline
    \hl\ip{10.45.6.9}      & \sol{A} & \sol{255.0.0.0}     & \sol{10.0.0.0}     & \sol{10.255.255.255} & \sol{16\,777\,214} \\ \hline}
  \def\apipa{169.254.88.3}
  \def\Cnet{192.168.80}
  \def\Crows{%
    PC-A & \ip{192.168.8.11}   & \ip{255.255.255.0} & \ip{192.168.80.254} & — \\
    PC-B & \ip{192.168.80.12}  & \ip{255.255.255.0} & \ip{192.168.80.254} & — \\
    PC-C & \ip{192.168.80.13}  & \ip{255.255.255.0} & \ip{192.168.80.100} & — \\
    PC-D & \ip{192.168.80.12}  & \ip{255.255.255.0} & \ip{192.168.80.254} & — \\
    SRV-C & \ip{10.0.0.8}      & \ip{255.0.0.0}     & \ip{10.0.0.254}     & — \\
    R1 — G0/0 & \ip{192.168.80.254} & \ip{255.255.255.0} & — & Off \\
    R1 — G0/1 & \ip{10.0.0.254}     & \ip{255.0.0.0}     & — & On \\}
  \def\Csol{%
    \hl 1 & \sol{PC-A} & \sol{IP hors du réseau 192.168.80.0/24.} & \sol{Ne joint ni les postes ni le routeur.} & \sol{Ex. 192.168.80.11.} \\ \hline
    \hl 2 & \sol{PC-C} & \sol{Passerelle .80.100 : pas l'adresse du routeur.} & \sol{Joint son réseau, mais aucun autre (serveur).} & \sol{Passerelle 192.168.80.254.} \\ \hline
    \hl 3 & \sol{PC-D} & \sol{Même IP que PC-B (doublon).} & \sol{Conflit d'adresse : PC-B et PC-D perturbés.} & \sol{Adresse unique, ex. .80.14.} \\ \hline
    \hl 4 & \sol{R1 G0/0} & \sol{Interface côté LAN désactivée (Off).} & \sol{Aucun poste ne joint le routeur ni le serveur.} & \sol{Port Status : On.} \\ \hline}
  \def\Cbareme{PC-B et SRV-C sont correctement configurés ; accepter PC-B cité à la place de PC-D pour le doublon.}
\fi

\pagestyle{fancy}
\fancyhf{}
\renewcommand{\headrulewidth}{0pt}
\fancyfoot[L]{\footnotesize\itshape\color{grisfoot}Évaluation — Réseaux locaux — Sujet \variante\ifcorrige\ — CORRIGÉ\fi}
\fancyfoot[R]{\footnotesize\color{grisfoot}\thepage/3}

\setlist{itemsep=0pt,topsep=2pt,parsep=0pt}
\setlength{\parindent}{0pt}
\setlength{\parskip}{2pt}

\begin{document}

\begin{center}
  {\Large\bfseries Évaluation pratique — Réseaux locaux \quad
   \fcolorbox{black}{white}{\ SUJET \variante\ }\par}
  \vspace{1pt}
  {\bfseries Adressage IP, commutation et routage — BTS CIEL 1re année — durée : 2 heures\par}
  \ifcorrige\vspace{2pt}{\bfseries\color{corr}CORRIGÉ ET BARÈME — document enseignant}\par\fi
\end{center}
\vspace{-2pt}
{\renewcommand{\arraystretch}{1.35}
\begin{tabularx}{\linewidth}{|X|X|p{2.2cm}|p{2.6cm}|}
  \hline
  \textbf{Nom :} & \textbf{Prénom :} & \textbf{Poste n°} & \textbf{Note :} \hfill /20 \\ \hline
\end{tabularx}}

\textbf{Consignes :} travail \textbf{individuel}, documents et Internet \textbf{non autorisés}. Ce sujet est le
\textbf{document-réponse}. Fichier Packet Tracer : \ip{Nom\_Prenom\_EvalLAN\_\variante.pkt}, à enregistrer
régulièrement dans le dossier de rendu. Les « visas » sont apposés par le professeur, que vous appelez à chaque test réussi.
Durées conseillées : A 30 min — B 1 h 10 — C 20 min.

%==========================================================================
\section{Partie A — Questions de cours et calculs \hfill /6}

\question{A1}{2 pts}{Complétez le tableau en considérant le \textbf{masque par défaut} de la classe de chaque adresse.}
{\small\renewcommand{\arraystretch}{1.15}
\begin{tabularx}{\linewidth}{|l|c|C|C|C|c|}
  \hline
  \textbf{Adresse IP} & \textbf{Classe} & \textbf{Masque par défaut} & \textbf{Adresse réseau} &
  \textbf{Adresse de diffusion} & \textbf{Nb d'hôtes} \\ \hline
  \AIrows
\end{tabularx}}
\bareme{0,5 pt par ligne juste + 0,5 pt si les trois nombres d'hôtes sont justes ($2^{n}-2$, $n$ = nombre de bits d'hôte).}

\question{A2}{1 pt}{Indiquez le type de câble (\textbf{droit} ou \textbf{croisé}) pour chaque liaison.}
{\small\renewcommand{\arraystretch}{1.15}
\begin{tabularx}{\linewidth}{|X|c||X|c|}
  \hline
  \hl PC $\leftrightarrow$ commutateur & \makebox[1.6cm]{\sol{droit}} & Commutateur $\leftrightarrow$ routeur & \makebox[1.6cm]{\sol{droit}} \\ \hline
  \hl PC $\leftrightarrow$ PC (réseau P2P) & \sol{croisé} & Concentrateur $\leftrightarrow$ concentrateur & \sol{croisé} \\ \hline
\end{tabularx}}
\bareme{0,25 pt par case. Équipements de types différents : câble droit ; de même type : câble croisé.}

\question{A3}{1 pt}{Comment un \textbf{concentrateur (hub)} et un \textbf{commutateur (switch)} traitent-ils une trame
reçue ? Lequel limite les collisions, et pourquoi ?}
\zone{2.9cm}{Le \textbf{hub} (couche 1) répète la trame sur \textbf{tous ses autres ports} : un seul domaine de collision
partagé. Le \textbf{switch} (couche 2) lit l'adresse MAC de destination et, grâce à sa \textbf{table MAC}, ne l'envoie
\textbf{que sur le port du destinataire} : un domaine de collision par port, c'est lui qui limite les collisions.
\emph{0,5 pt hub, 0,5 pt switch.}}

\question{A4}{1 pt}{Qu'est-ce que la \textbf{passerelle par défaut} d'un poste ? Quelle condition son adresse doit-elle respecter ?}
\zone{2.4cm}{Adresse IP de l'interface du \textbf{routeur} à laquelle le poste envoie les paquets destinés à un
\textbf{autre réseau}. Elle doit \textbf{appartenir au même réseau} que le poste. \emph{0,5 + 0,5 pt.}}

\question{A5}{0,5 pt}{Un portable configuré en DHCP affiche \ip{\apipa} / \ip{255.255.0.0}. Que s'est-il passé ?}
\zone{1.5cm}{Aucun \textbf{serveur DHCP} n'a répondu : le poste s'est attribué une adresse \textbf{APIPA} (169.254.0.0/16).}

\question{A6}{0,5 pt}{Quel protocole utilise \texttt{ping} ? Nommez les deux messages échangés et la couche OSI des adresses IP.}
\zone{1.5cm}{\textbf{ICMP} ; \textbf{Echo Request} / \textbf{Echo Reply} ; adresses IP en \textbf{couche 3 (Réseau)}.}

%==========================================================================
\newpage
\section{Partie B — Réalisation sous Packet Tracer \hfill /10}

\begin{cdc}
\entreprise{} vous confie son réseau : \textbf{trois réseaux} reliés par un \textbf{routeur unique Cisco 2911}
nommé \ip{\Rnom} (interfaces G0/0, G0/1, G0/2).
\begin{itemize}[leftmargin=12pt]
  \item \textbf{Réseau \resUn} : \textbf{4 postes} PC-\pUn1 à PC-\pUn4, commutateur \ip{SW-\pUn}, adressage statique,
        réseau \ip{\netUn.0/24}.
  \item \textbf{Réseau \resDeux} : \textbf{2 postes} PC-\pDeux1, PC-\pDeux2 et \textbf{1 imprimante} IMP-\pDeux,
        commutateur \ip{SW-\pDeux}, adressage statique, réseau de classe B \ip{\netDeux.0.0/16}.
  \item \textbf{Réseau Serveur} : serveur \ip{\srvNom} d'adresse \ip{\netSrv.\srvH/24}, relié \textbf{directement} au routeur.
  \item La passerelle de chaque réseau est sa \textbf{dernière adresse utilisable}. Tous les équipements communiquent
        entre eux et accèdent à la page web du serveur. Nommage imposé (« Display Name »).
\end{itemize}
\end{cdc}

\question{B1}{2 pts}{Complétez le plan d'adressage \textbf{avant} la réalisation.}
{\footnotesize\setlength{\tabcolsep}{3pt}\renewcommand{\arraystretch}{1.15}
\begin{tabularx}{\linewidth}{|l|C|C|C|C|C|C|}
  \hline
  \textbf{Réseau} & \textbf{Adresse réseau} & \textbf{Masque} & \textbf{1re @ utilisable} &
  \textbf{Dernière @ utilisable} & \textbf{Diffusion} & \textbf{Passerelle} \\ \hline
  \hl \resUn & \sol{\netUn.0} & \sol{255.255.255.0} & \sol{\netUn.1} & \sol{\netUn.254} & \sol{\netUn.255} & \sol{\netUn.254} \\ \hline
  \hl \resDeux & \sol{\netDeux.0.0} & \sol{255.255.0.0} & \sol{\netDeux.0.1} & \sol{\netDeux.255.254} & \sol{\netDeux.255.255} & \sol{\netDeux.255.254} \\ \hline
  \hl Serveur & \sol{\netSrv.0} & \sol{255.255.255.0} & \sol{\netSrv.1} & \sol{\netSrv.254} & \sol{\netSrv.255} & \sol{\netSrv.254} \\ \hline
\end{tabularx}}
\bareme{0,5 pt par ligne juste + 0,5 pt si la convention de passerelle est respectée partout.}

\question{B2}{2 pts}{Construisez la topologie (équipements, câblage, nommage) et relevez l'adressage configuré.}
{\footnotesize\setlength{\tabcolsep}{3pt}\renewcommand{\arraystretch}{1.1}
\begin{tabularx}{\linewidth}{|l|C|C|C||l|C|C|}
  \hline
  \textbf{Hôte} & \textbf{Adresse IP} & \textbf{Masque} & \textbf{Passerelle} & \textbf{Interface} & \textbf{Adresse IP} & \textbf{Masque} \\ \hline
  \hl PC-\pUn1 & \sol{\netUn.1} & \sol{255.255.255.0} & \sol{\netUn.254} & \Rnom{} G0/0 & \sol{\netUn.254} & \sol{255.255.255.0} \\ \hline
  \hl PC-\pDeux1 & \sol{\netDeux.0.1} & \sol{255.255.0.0} & \sol{\netDeux.255.254} & \Rnom{} G0/1 & \sol{\netDeux.255.254} & \sol{255.255.0.0} \\ \hline
  \hl IMP-\pDeux & \sol{\netDeux.0.10} & \sol{255.255.0.0} & \sol{\netDeux.255.254} & \Rnom{} G0/2 & \sol{\netSrv.254} & \sol{255.255.255.0} \\ \hline
  \hl \srvNom & \ip{\netSrv.\srvH} & \sol{255.255.255.0} & \sol{\netSrv.254} & \multicolumn{3}{c|}{\cellcolor{black!8}} \\ \hline
\end{tabularx}}
\bareme{Adresses d'hôtes libres (uniques, dans le bon réseau) ; affectation G0/x libre. Sur le .pkt : équipements 0,5 ;
câblage 1 (câbles droits vers les switchs et switch–routeur, \textbf{croisé} serveur–routeur, liens verts) ; nommage 0,5.}

\question{B3}{3 pts}{Configurez tous les équipements : IP, masques, passerelles des hôtes ; interfaces du routeur
(« Port Status : On »).}
\bareme{Sur le .pkt : interfaces routeur 1 pt ; IP/masques des hôtes 1 pt ; passerelles (imprimante et serveur compris)
1 pt. $-0,25$ par erreur, sans descendre sous 0 par item.}

\question{B4}{3 pts}{Réalisez les tests, notez la commande et le résultat, puis faites viser chaque test réussi.}
{\small\renewcommand{\arraystretch}{1.1}
\begin{tabularx}{\linewidth}{|c|l|l|C|c|c|}
  \hline
  \textbf{N°} & \textbf{Source} & \textbf{Destination} & \textbf{Commande / action} & \textbf{Résultat} & \textbf{Visa prof.} \\ \hline
  \hl 1 & PC-\pUn1 & PC-\pUn4 & \sol{ping \netUn.4} & \sol{réussi} & \hspace{1.4cm} \\ \hline
  \hl 2 & PC-\pDeux1 & IMP-\pDeux & \sol{ping \netDeux.0.10} & \sol{réussi} & \\ \hline
  \hl 3 & PC-\pUn2 & PC-\pDeux2 & \sol{ping \netDeux.0.2} & \sol{réussi} & \\ \hline
  \hl 4 & PC-\pDeux2 & \srvNom & \sol{ping \netSrv.\srvH} & \sol{réussi} & \\ \hline
  \hl 5 & PC-\pUn3 & page web du serveur & \sol{navigateur : {\NoAutoSpacing http://\netSrv.\srvH}} & \sol{page affichée} & \\ \hline
\end{tabularx}}
\bareme{0,5 pt par test visé + 0,5 pt pour B5. Le premier ping inter-réseaux peut échouer (résolution ARP) ; les suivants doivent réussir.}

\question{B5}{incl. B4}{Test n°3 : par quels équipements passe le paquet ? Pourquoi PC-\pUn2 l'envoie-t-il au routeur ?}
\zone{2.3cm}{PC-\pUn2 $\to$ SW-\pUn{} $\to$ \Rnom{} $\to$ SW-\pDeux{} $\to$ PC-\pDeux2. Grâce à son masque, PC-\pUn2 constate que
la destination (\netDeux.0.0/16) n'est \textbf{pas dans son réseau} (\netUn.0/24) : il envoie le paquet à sa
\textbf{passerelle par défaut}, qui le route.}

%==========================================================================
\newpage
\section{Partie C — Diagnostic d'un réseau en panne \hfill /4}

Un technicien constate qu'\textbf{aucun poste ne parvient à joindre le serveur SRV-C} et que certains postes
communiquent mal entre eux. Les câbles sont du bon type.

\begin{center}
\begin{minipage}{0.37\linewidth}\centering
\begin{tikzpicture}[font=\footnotesize,
  eq/.style={draw,fill=qbg,rounded corners=2pt,minimum width=9mm,minimum height=5mm,inner sep=2pt},
  every path/.style={thick}]
  \node[eq] (SW) {SW-1};
  \node[eq,right=9mm of SW] (R) {R1};
  \node[eq,right=7mm of R] (SRV) {SRV-C};
  \node[eq,above left=5mm and -3mm of SW] (A) {PC-A};
  \node[eq,above right=5mm and -3mm of SW] (B) {PC-B};
  \node[eq,below left=5mm and -3mm of SW] (C) {PC-C};
  \node[eq,below right=5mm and -3mm of SW] (D) {PC-D};
  \foreach \n in {A,B,C,D} \draw (\n) -- (SW);
  \draw (SW) -- node[above,font=\tiny]{G0/0} (R);
  \draw (R) -- node[above,font=\tiny]{G0/1} (SRV);
\end{tikzpicture}
\end{minipage}\hfill
\begin{minipage}{0.61\linewidth}\centering
{\small\renewcommand{\arraystretch}{1.05}\setlength{\tabcolsep}{4pt}
\begin{tabular}{lllll}
  \toprule
  \textbf{Équipement} & \textbf{Adresse IP} & \textbf{Masque} & \textbf{Passerelle} & \textbf{Port} \\
  \midrule
  \Crows
  \bottomrule
\end{tabular}}
\end{minipage}
\end{center}

\question{C1}{4 pts}{Relevez les \textbf{quatre anomalies}. Pour chacune : conséquence et correction.}
{\small\renewcommand{\arraystretch}{1.1}
\ifcorrige\def\hc{}\else\def\hc{\rule[-10mm]{0pt}{14mm}}\fi
\begin{tabularx}{\linewidth}{|c|>{\raggedright\arraybackslash}p{2cm}|L|L|L|}
  \hline
  \textbf{N°} & \textbf{Équipement} & \textbf{Anomalie} & \textbf{Conséquence} & \textbf{Correction} \\ \hline
  \ifcorrige\Csol\else
  \hc 1 & & & & \\ \hline
  \hc 2 & & & & \\ \hline
  \hc 3 & & & & \\ \hline
  \hc 4 & & & & \\ \hline
  \fi
\end{tabularx}}
\bareme{1 pt par anomalie (0,25 équipement, anomalie, conséquence, correction ; ordre indifférent). \Cbareme}

%==========================================================================
\section{Bonus — Extension sans fil \hfill +1}

\begin{cdc}Ajoutez au réseau \resUn{} un \textbf{AccessPoint-PT} relié à SW-\pUn{} (SSID \ip{\ssid}, \textbf{WPA2-PSK})
et un portable \ip{Laptop-\pUn} équipé d'un module \texttt{WPC300N}, en adressage \textbf{statique}.
Vérifiez qu'il joint \srvNom.\end{cdc}
{\small\renewcommand{\arraystretch}{1.1}
\begin{tabularx}{\linewidth}{|l|C|C|C|c|}
  \hline
  \textbf{Équipement} & \textbf{Adresse IP} & \textbf{Masque} & \textbf{Passerelle} & \textbf{Visa prof.} \\ \hline
  \hl Laptop-\pUn & \sol{\netUn.50} & \sol{255.255.255.0} & \sol{\netUn.254} & \hspace{1.4cm} \\ \hline
\end{tabularx}}
\bareme{Le point d'accès fait office de pont : le portable est dans \netUn.0/24. Bonus si associé en WPA2-PSK et ping
vers \netSrv.\srvH{} réussi. Note plafonnée à 20.}

\end{document}
"
%   pdflatex -jobname=Eval_Reseau_CorrigeA "\def\variante{A}\def\modecorrige{1}% Évaluation pratique — Réseaux locaux (BTS CIEL 1re année) — sujets A et B
% Compilation (4 PDF) :
%   pdflatex -jobname=Eval_Reseau_SujetA   "\def\variante{A}\def\modecorrige{0}\input{Eval_TP_Reseau_LAN.tex}"
%   pdflatex -jobname=Eval_Reseau_SujetB   "\def\variante{B}\def\modecorrige{0}\input{Eval_TP_Reseau_LAN.tex}"
%   pdflatex -jobname=Eval_Reseau_CorrigeA "\def\variante{A}\def\modecorrige{1}\input{Eval_TP_Reseau_LAN.tex}"
%   pdflatex -jobname=Eval_Reseau_CorrigeB "\def\variante{B}\def\modecorrige{1}\input{Eval_TP_Reseau_LAN.tex}"
\documentclass[10pt,a4paper]{article}

\usepackage[T1]{fontenc}
\usepackage[utf8]{inputenc}
\usepackage{lmodern}
\usepackage[french]{babel}
\usepackage[a4paper,hmargin=1.5cm,top=1.3cm,bottom=1.6cm,footskip=0.7cm]{geometry}
\usepackage[table]{xcolor}
\usepackage{booktabs,tabularx,array}
\usepackage{enumitem}
\usepackage{fancyhdr}
\usepackage[most]{tcolorbox}
\usepackage{tikz}
\usetikzlibrary{positioning}
\usepackage{titlesec}

\providecommand{\modecorrige}{0}
\providecommand{\variante}{A}
\newif\ifcorrige
\ifnum\modecorrige=1 \corrigetrue \else \corrigefalse \fi

\definecolor{ipred}{HTML}{C0392B}
\definecolor{cdcbar}{HTML}{A67C3D}
\definecolor{cdcbg}{HTML}{F4F1EA}
\definecolor{qbar}{HTML}{4A6B5D}
\definecolor{qbg}{HTML}{EDF1EF}
\definecolor{corr}{HTML}{1F5FA8}
\definecolor{corrbg}{HTML}{E8F0FA}
\definecolor{grisfoot}{HTML}{828282}

\newcommand{\ip}[1]{\textcolor{ipred}{\texttt{#1}}}
\newcommand{\pts}[1]{\hfill\mbox{\textbf{(#1)}}}
\newcommand{\sol}[1]{\ifcorrige\textcolor{corr}{#1}\fi}
\newcommand{\bareme}[1]{\ifcorrige\par{\footnotesize\color{corr}#1\par}\fi}
\newcommand{\hl}{\ifcorrige\rule[-1.6mm]{0pt}{5.4mm}\else\rule[-2.2mm]{0pt}{6.6mm}\fi}
\newcolumntype{C}{>{\centering\arraybackslash}X}
\newcolumntype{L}{>{\raggedright\arraybackslash}X}

\titleformat{\section}{\large\bfseries}{}{0pt}{}
\titlespacing*{\section}{0pt}{6pt}{3pt}

\tcbset{qstyle/.style={enhanced,boxrule=0pt,leftrule=2.5pt,arc=0pt,outer arc=0pt,
  left=5pt,right=5pt,top=2pt,bottom=2pt,before skip=4pt,after skip=3pt}}
\newtcolorbox{cdc}{qstyle,colback=cdcbg,colframe=cdcbar}
\newtcolorbox{quest}{qstyle,colback=qbg,colframe=qbar}

% Zone de réponse : vierge dans le sujet, remplie dans le corrigé
\newcommand{\zone}[2]{%
  \ifcorrige
    \begin{tcolorbox}[colback=corrbg,colframe=corr,boxrule=0.5pt,arc=1pt,left=4pt,right=4pt,top=1pt,bottom=1pt,
      before skip=2pt,after skip=2pt,fontupper=\footnotesize\color{corr}]#2\end{tcolorbox}%
  \else
    \begin{tcolorbox}[colback=white,colframe=black!45,boxrule=0.4pt,arc=1pt,height=#1,
      before skip=2pt,after skip=2pt]\end{tcolorbox}%
  \fi}

\newcommand{\question}[3]{\begin{quest}\textbf{#1.} #3\pts{#2}\end{quest}}

%==========================================================================
% Données propres à chaque variante
%==========================================================================
\if\variante A
  \def\entreprise{Le cabinet d'architectes \textbf{ArchiNova}}
  \def\Rnom{R-ArchiNova}
  \def\resUn{Bureau d'études}\def\pUn{BE}\def\netUn{192.168.40}
  \def\resDeux{Accueil}\def\pDeux{ACC}\def\netDeux{172.16}
  \def\srvNom{SRV-Intranet}\def\netSrv{192.168.200}\def\srvH{10}
  \def\ssid{ArchiNova-WiFi}
  \def\AIrows{%
    \hl\ip{192.168.40.25} & \sol{C} & \sol{255.255.255.0} & \sol{192.168.40.0} & \sol{192.168.40.255} & \sol{254} \\ \hline
    \hl\ip{172.16.5.10}   & \sol{B} & \sol{255.255.0.0}   & \sol{172.16.0.0}   & \sol{172.16.255.255} & \sol{65\,534} \\ \hline
    \hl\ip{10.3.2.1}      & \sol{A} & \sol{255.0.0.0}     & \sol{10.0.0.0}     & \sol{10.255.255.255} & \sol{16\,777\,214} \\ \hline}
  \def\apipa{169.254.12.7}
  \def\Cnet{192.168.60}
  \def\Crows{%
    PC-A & \ip{192.168.60.10} & \ip{255.255.255.0} & \ip{192.168.60.254} & — \\
    PC-B & \ip{192.168.60.20} & \ip{255.255.255.0} & \ip{192.168.60.1}   & — \\
    PC-C & \ip{192.168.61.30} & \ip{255.255.255.0} & \ip{192.168.60.254} & — \\
    PC-D & \ip{192.168.60.10} & \ip{255.255.255.0} & \ip{192.168.60.254} & — \\
    SRV-C & \ip{10.0.0.5}     & \ip{255.0.0.0}     & \ip{10.0.0.254}     & — \\
    R1 — G0/0 & \ip{192.168.60.254} & \ip{255.255.255.0} & — & On \\
    R1 — G0/1 & \ip{10.0.0.254}     & \ip{255.0.0.0}     & — & Off \\}
  \def\Csol{%
    \hl 1 & \sol{PC-B} & \sol{Passerelle .60.1 : pas l'adresse du routeur.} & \sol{Joint son réseau, mais aucun autre (serveur).} & \sol{Passerelle 192.168.60.254.} \\ \hline
    \hl 2 & \sol{PC-C} & \sol{IP hors du réseau 192.168.60.0/24.} & \sol{Ne joint ni les postes ni le routeur.} & \sol{Ex. 192.168.60.30.} \\ \hline
    \hl 3 & \sol{PC-D} & \sol{Même IP que PC-A (doublon).} & \sol{Conflit d'adresse : PC-A et PC-D perturbés.} & \sol{Adresse unique, ex. .60.40.} \\ \hline
    \hl 4 & \sol{R1 G0/1} & \sol{Interface désactivée (Off).} & \sol{Réseau du serveur injoignable pour tous.} & \sol{Port Status : On.} \\ \hline}
  \def\Cbareme{PC-A et SRV-C sont correctement configurés ; accepter PC-A cité à la place de PC-D pour le doublon.}
\else
  \def\entreprise{La clinique vétérinaire \textbf{VetoPlus}}
  \def\Rnom{R-VetoPlus}
  \def\resUn{Soins}\def\pUn{SOI}\def\netUn{192.168.70}
  \def\resDeux{Secrétariat}\def\pDeux{SEC}\def\netDeux{172.20}
  \def\srvNom{SRV-Dossiers}\def\netSrv{192.168.250}\def\srvH{20}
  \def\ssid{VetoPlus-WiFi}
  \def\AIrows{%
    \hl\ip{192.168.15.130} & \sol{C} & \sol{255.255.255.0} & \sol{192.168.15.0} & \sol{192.168.15.255} & \sol{254} \\ \hline
    \hl\ip{172.31.200.7}   & \sol{B} & \sol{255.255.0.0}   & \sol{172.31.0.0}   & \sol{172.31.255.255} & \sol{65\,534} \\ \hline
    \hl\ip{10.45.6.9}      & \sol{A} & \sol{255.0.0.0}     & \sol{10.0.0.0}     & \sol{10.255.255.255} & \sol{16\,777\,214} \\ \hline}
  \def\apipa{169.254.88.3}
  \def\Cnet{192.168.80}
  \def\Crows{%
    PC-A & \ip{192.168.8.11}   & \ip{255.255.255.0} & \ip{192.168.80.254} & — \\
    PC-B & \ip{192.168.80.12}  & \ip{255.255.255.0} & \ip{192.168.80.254} & — \\
    PC-C & \ip{192.168.80.13}  & \ip{255.255.255.0} & \ip{192.168.80.100} & — \\
    PC-D & \ip{192.168.80.12}  & \ip{255.255.255.0} & \ip{192.168.80.254} & — \\
    SRV-C & \ip{10.0.0.8}      & \ip{255.0.0.0}     & \ip{10.0.0.254}     & — \\
    R1 — G0/0 & \ip{192.168.80.254} & \ip{255.255.255.0} & — & Off \\
    R1 — G0/1 & \ip{10.0.0.254}     & \ip{255.0.0.0}     & — & On \\}
  \def\Csol{%
    \hl 1 & \sol{PC-A} & \sol{IP hors du réseau 192.168.80.0/24.} & \sol{Ne joint ni les postes ni le routeur.} & \sol{Ex. 192.168.80.11.} \\ \hline
    \hl 2 & \sol{PC-C} & \sol{Passerelle .80.100 : pas l'adresse du routeur.} & \sol{Joint son réseau, mais aucun autre (serveur).} & \sol{Passerelle 192.168.80.254.} \\ \hline
    \hl 3 & \sol{PC-D} & \sol{Même IP que PC-B (doublon).} & \sol{Conflit d'adresse : PC-B et PC-D perturbés.} & \sol{Adresse unique, ex. .80.14.} \\ \hline
    \hl 4 & \sol{R1 G0/0} & \sol{Interface côté LAN désactivée (Off).} & \sol{Aucun poste ne joint le routeur ni le serveur.} & \sol{Port Status : On.} \\ \hline}
  \def\Cbareme{PC-B et SRV-C sont correctement configurés ; accepter PC-B cité à la place de PC-D pour le doublon.}
\fi

\pagestyle{fancy}
\fancyhf{}
\renewcommand{\headrulewidth}{0pt}
\fancyfoot[L]{\footnotesize\itshape\color{grisfoot}Évaluation — Réseaux locaux — Sujet \variante\ifcorrige\ — CORRIGÉ\fi}
\fancyfoot[R]{\footnotesize\color{grisfoot}\thepage/3}

\setlist{itemsep=0pt,topsep=2pt,parsep=0pt}
\setlength{\parindent}{0pt}
\setlength{\parskip}{2pt}

\begin{document}

\begin{center}
  {\Large\bfseries Évaluation pratique — Réseaux locaux \quad
   \fcolorbox{black}{white}{\ SUJET \variante\ }\par}
  \vspace{1pt}
  {\bfseries Adressage IP, commutation et routage — BTS CIEL 1re année — durée : 2 heures\par}
  \ifcorrige\vspace{2pt}{\bfseries\color{corr}CORRIGÉ ET BARÈME — document enseignant}\par\fi
\end{center}
\vspace{-2pt}
{\renewcommand{\arraystretch}{1.35}
\begin{tabularx}{\linewidth}{|X|X|p{2.2cm}|p{2.6cm}|}
  \hline
  \textbf{Nom :} & \textbf{Prénom :} & \textbf{Poste n°} & \textbf{Note :} \hfill /20 \\ \hline
\end{tabularx}}

\textbf{Consignes :} travail \textbf{individuel}, documents et Internet \textbf{non autorisés}. Ce sujet est le
\textbf{document-réponse}. Fichier Packet Tracer : \ip{Nom\_Prenom\_EvalLAN\_\variante.pkt}, à enregistrer
régulièrement dans le dossier de rendu. Les « visas » sont apposés par le professeur, que vous appelez à chaque test réussi.
Durées conseillées : A 30 min — B 1 h 10 — C 20 min.

%==========================================================================
\section{Partie A — Questions de cours et calculs \hfill /6}

\question{A1}{2 pts}{Complétez le tableau en considérant le \textbf{masque par défaut} de la classe de chaque adresse.}
{\small\renewcommand{\arraystretch}{1.15}
\begin{tabularx}{\linewidth}{|l|c|C|C|C|c|}
  \hline
  \textbf{Adresse IP} & \textbf{Classe} & \textbf{Masque par défaut} & \textbf{Adresse réseau} &
  \textbf{Adresse de diffusion} & \textbf{Nb d'hôtes} \\ \hline
  \AIrows
\end{tabularx}}
\bareme{0,5 pt par ligne juste + 0,5 pt si les trois nombres d'hôtes sont justes ($2^{n}-2$, $n$ = nombre de bits d'hôte).}

\question{A2}{1 pt}{Indiquez le type de câble (\textbf{droit} ou \textbf{croisé}) pour chaque liaison.}
{\small\renewcommand{\arraystretch}{1.15}
\begin{tabularx}{\linewidth}{|X|c||X|c|}
  \hline
  \hl PC $\leftrightarrow$ commutateur & \makebox[1.6cm]{\sol{droit}} & Commutateur $\leftrightarrow$ routeur & \makebox[1.6cm]{\sol{droit}} \\ \hline
  \hl PC $\leftrightarrow$ PC (réseau P2P) & \sol{croisé} & Concentrateur $\leftrightarrow$ concentrateur & \sol{croisé} \\ \hline
\end{tabularx}}
\bareme{0,25 pt par case. Équipements de types différents : câble droit ; de même type : câble croisé.}

\question{A3}{1 pt}{Comment un \textbf{concentrateur (hub)} et un \textbf{commutateur (switch)} traitent-ils une trame
reçue ? Lequel limite les collisions, et pourquoi ?}
\zone{2.9cm}{Le \textbf{hub} (couche 1) répète la trame sur \textbf{tous ses autres ports} : un seul domaine de collision
partagé. Le \textbf{switch} (couche 2) lit l'adresse MAC de destination et, grâce à sa \textbf{table MAC}, ne l'envoie
\textbf{que sur le port du destinataire} : un domaine de collision par port, c'est lui qui limite les collisions.
\emph{0,5 pt hub, 0,5 pt switch.}}

\question{A4}{1 pt}{Qu'est-ce que la \textbf{passerelle par défaut} d'un poste ? Quelle condition son adresse doit-elle respecter ?}
\zone{2.4cm}{Adresse IP de l'interface du \textbf{routeur} à laquelle le poste envoie les paquets destinés à un
\textbf{autre réseau}. Elle doit \textbf{appartenir au même réseau} que le poste. \emph{0,5 + 0,5 pt.}}

\question{A5}{0,5 pt}{Un portable configuré en DHCP affiche \ip{\apipa} / \ip{255.255.0.0}. Que s'est-il passé ?}
\zone{1.5cm}{Aucun \textbf{serveur DHCP} n'a répondu : le poste s'est attribué une adresse \textbf{APIPA} (169.254.0.0/16).}

\question{A6}{0,5 pt}{Quel protocole utilise \texttt{ping} ? Nommez les deux messages échangés et la couche OSI des adresses IP.}
\zone{1.5cm}{\textbf{ICMP} ; \textbf{Echo Request} / \textbf{Echo Reply} ; adresses IP en \textbf{couche 3 (Réseau)}.}

%==========================================================================
\newpage
\section{Partie B — Réalisation sous Packet Tracer \hfill /10}

\begin{cdc}
\entreprise{} vous confie son réseau : \textbf{trois réseaux} reliés par un \textbf{routeur unique Cisco 2911}
nommé \ip{\Rnom} (interfaces G0/0, G0/1, G0/2).
\begin{itemize}[leftmargin=12pt]
  \item \textbf{Réseau \resUn} : \textbf{4 postes} PC-\pUn1 à PC-\pUn4, commutateur \ip{SW-\pUn}, adressage statique,
        réseau \ip{\netUn.0/24}.
  \item \textbf{Réseau \resDeux} : \textbf{2 postes} PC-\pDeux1, PC-\pDeux2 et \textbf{1 imprimante} IMP-\pDeux,
        commutateur \ip{SW-\pDeux}, adressage statique, réseau de classe B \ip{\netDeux.0.0/16}.
  \item \textbf{Réseau Serveur} : serveur \ip{\srvNom} d'adresse \ip{\netSrv.\srvH/24}, relié \textbf{directement} au routeur.
  \item La passerelle de chaque réseau est sa \textbf{dernière adresse utilisable}. Tous les équipements communiquent
        entre eux et accèdent à la page web du serveur. Nommage imposé (« Display Name »).
\end{itemize}
\end{cdc}

\question{B1}{2 pts}{Complétez le plan d'adressage \textbf{avant} la réalisation.}
{\footnotesize\setlength{\tabcolsep}{3pt}\renewcommand{\arraystretch}{1.15}
\begin{tabularx}{\linewidth}{|l|C|C|C|C|C|C|}
  \hline
  \textbf{Réseau} & \textbf{Adresse réseau} & \textbf{Masque} & \textbf{1re @ utilisable} &
  \textbf{Dernière @ utilisable} & \textbf{Diffusion} & \textbf{Passerelle} \\ \hline
  \hl \resUn & \sol{\netUn.0} & \sol{255.255.255.0} & \sol{\netUn.1} & \sol{\netUn.254} & \sol{\netUn.255} & \sol{\netUn.254} \\ \hline
  \hl \resDeux & \sol{\netDeux.0.0} & \sol{255.255.0.0} & \sol{\netDeux.0.1} & \sol{\netDeux.255.254} & \sol{\netDeux.255.255} & \sol{\netDeux.255.254} \\ \hline
  \hl Serveur & \sol{\netSrv.0} & \sol{255.255.255.0} & \sol{\netSrv.1} & \sol{\netSrv.254} & \sol{\netSrv.255} & \sol{\netSrv.254} \\ \hline
\end{tabularx}}
\bareme{0,5 pt par ligne juste + 0,5 pt si la convention de passerelle est respectée partout.}

\question{B2}{2 pts}{Construisez la topologie (équipements, câblage, nommage) et relevez l'adressage configuré.}
{\footnotesize\setlength{\tabcolsep}{3pt}\renewcommand{\arraystretch}{1.1}
\begin{tabularx}{\linewidth}{|l|C|C|C||l|C|C|}
  \hline
  \textbf{Hôte} & \textbf{Adresse IP} & \textbf{Masque} & \textbf{Passerelle} & \textbf{Interface} & \textbf{Adresse IP} & \textbf{Masque} \\ \hline
  \hl PC-\pUn1 & \sol{\netUn.1} & \sol{255.255.255.0} & \sol{\netUn.254} & \Rnom{} G0/0 & \sol{\netUn.254} & \sol{255.255.255.0} \\ \hline
  \hl PC-\pDeux1 & \sol{\netDeux.0.1} & \sol{255.255.0.0} & \sol{\netDeux.255.254} & \Rnom{} G0/1 & \sol{\netDeux.255.254} & \sol{255.255.0.0} \\ \hline
  \hl IMP-\pDeux & \sol{\netDeux.0.10} & \sol{255.255.0.0} & \sol{\netDeux.255.254} & \Rnom{} G0/2 & \sol{\netSrv.254} & \sol{255.255.255.0} \\ \hline
  \hl \srvNom & \ip{\netSrv.\srvH} & \sol{255.255.255.0} & \sol{\netSrv.254} & \multicolumn{3}{c|}{\cellcolor{black!8}} \\ \hline
\end{tabularx}}
\bareme{Adresses d'hôtes libres (uniques, dans le bon réseau) ; affectation G0/x libre. Sur le .pkt : équipements 0,5 ;
câblage 1 (câbles droits vers les switchs et switch–routeur, \textbf{croisé} serveur–routeur, liens verts) ; nommage 0,5.}

\question{B3}{3 pts}{Configurez tous les équipements : IP, masques, passerelles des hôtes ; interfaces du routeur
(« Port Status : On »).}
\bareme{Sur le .pkt : interfaces routeur 1 pt ; IP/masques des hôtes 1 pt ; passerelles (imprimante et serveur compris)
1 pt. $-0,25$ par erreur, sans descendre sous 0 par item.}

\question{B4}{3 pts}{Réalisez les tests, notez la commande et le résultat, puis faites viser chaque test réussi.}
{\small\renewcommand{\arraystretch}{1.1}
\begin{tabularx}{\linewidth}{|c|l|l|C|c|c|}
  \hline
  \textbf{N°} & \textbf{Source} & \textbf{Destination} & \textbf{Commande / action} & \textbf{Résultat} & \textbf{Visa prof.} \\ \hline
  \hl 1 & PC-\pUn1 & PC-\pUn4 & \sol{ping \netUn.4} & \sol{réussi} & \hspace{1.4cm} \\ \hline
  \hl 2 & PC-\pDeux1 & IMP-\pDeux & \sol{ping \netDeux.0.10} & \sol{réussi} & \\ \hline
  \hl 3 & PC-\pUn2 & PC-\pDeux2 & \sol{ping \netDeux.0.2} & \sol{réussi} & \\ \hline
  \hl 4 & PC-\pDeux2 & \srvNom & \sol{ping \netSrv.\srvH} & \sol{réussi} & \\ \hline
  \hl 5 & PC-\pUn3 & page web du serveur & \sol{navigateur : {\NoAutoSpacing http://\netSrv.\srvH}} & \sol{page affichée} & \\ \hline
\end{tabularx}}
\bareme{0,5 pt par test visé + 0,5 pt pour B5. Le premier ping inter-réseaux peut échouer (résolution ARP) ; les suivants doivent réussir.}

\question{B5}{incl. B4}{Test n°3 : par quels équipements passe le paquet ? Pourquoi PC-\pUn2 l'envoie-t-il au routeur ?}
\zone{2.3cm}{PC-\pUn2 $\to$ SW-\pUn{} $\to$ \Rnom{} $\to$ SW-\pDeux{} $\to$ PC-\pDeux2. Grâce à son masque, PC-\pUn2 constate que
la destination (\netDeux.0.0/16) n'est \textbf{pas dans son réseau} (\netUn.0/24) : il envoie le paquet à sa
\textbf{passerelle par défaut}, qui le route.}

%==========================================================================
\newpage
\section{Partie C — Diagnostic d'un réseau en panne \hfill /4}

Un technicien constate qu'\textbf{aucun poste ne parvient à joindre le serveur SRV-C} et que certains postes
communiquent mal entre eux. Les câbles sont du bon type.

\begin{center}
\begin{minipage}{0.37\linewidth}\centering
\begin{tikzpicture}[font=\footnotesize,
  eq/.style={draw,fill=qbg,rounded corners=2pt,minimum width=9mm,minimum height=5mm,inner sep=2pt},
  every path/.style={thick}]
  \node[eq] (SW) {SW-1};
  \node[eq,right=9mm of SW] (R) {R1};
  \node[eq,right=7mm of R] (SRV) {SRV-C};
  \node[eq,above left=5mm and -3mm of SW] (A) {PC-A};
  \node[eq,above right=5mm and -3mm of SW] (B) {PC-B};
  \node[eq,below left=5mm and -3mm of SW] (C) {PC-C};
  \node[eq,below right=5mm and -3mm of SW] (D) {PC-D};
  \foreach \n in {A,B,C,D} \draw (\n) -- (SW);
  \draw (SW) -- node[above,font=\tiny]{G0/0} (R);
  \draw (R) -- node[above,font=\tiny]{G0/1} (SRV);
\end{tikzpicture}
\end{minipage}\hfill
\begin{minipage}{0.61\linewidth}\centering
{\small\renewcommand{\arraystretch}{1.05}\setlength{\tabcolsep}{4pt}
\begin{tabular}{lllll}
  \toprule
  \textbf{Équipement} & \textbf{Adresse IP} & \textbf{Masque} & \textbf{Passerelle} & \textbf{Port} \\
  \midrule
  \Crows
  \bottomrule
\end{tabular}}
\end{minipage}
\end{center}

\question{C1}{4 pts}{Relevez les \textbf{quatre anomalies}. Pour chacune : conséquence et correction.}
{\small\renewcommand{\arraystretch}{1.1}
\ifcorrige\def\hc{}\else\def\hc{\rule[-10mm]{0pt}{14mm}}\fi
\begin{tabularx}{\linewidth}{|c|>{\raggedright\arraybackslash}p{2cm}|L|L|L|}
  \hline
  \textbf{N°} & \textbf{Équipement} & \textbf{Anomalie} & \textbf{Conséquence} & \textbf{Correction} \\ \hline
  \ifcorrige\Csol\else
  \hc 1 & & & & \\ \hline
  \hc 2 & & & & \\ \hline
  \hc 3 & & & & \\ \hline
  \hc 4 & & & & \\ \hline
  \fi
\end{tabularx}}
\bareme{1 pt par anomalie (0,25 équipement, anomalie, conséquence, correction ; ordre indifférent). \Cbareme}

%==========================================================================
\section{Bonus — Extension sans fil \hfill +1}

\begin{cdc}Ajoutez au réseau \resUn{} un \textbf{AccessPoint-PT} relié à SW-\pUn{} (SSID \ip{\ssid}, \textbf{WPA2-PSK})
et un portable \ip{Laptop-\pUn} équipé d'un module \texttt{WPC300N}, en adressage \textbf{statique}.
Vérifiez qu'il joint \srvNom.\end{cdc}
{\small\renewcommand{\arraystretch}{1.1}
\begin{tabularx}{\linewidth}{|l|C|C|C|c|}
  \hline
  \textbf{Équipement} & \textbf{Adresse IP} & \textbf{Masque} & \textbf{Passerelle} & \textbf{Visa prof.} \\ \hline
  \hl Laptop-\pUn & \sol{\netUn.50} & \sol{255.255.255.0} & \sol{\netUn.254} & \hspace{1.4cm} \\ \hline
\end{tabularx}}
\bareme{Le point d'accès fait office de pont : le portable est dans \netUn.0/24. Bonus si associé en WPA2-PSK et ping
vers \netSrv.\srvH{} réussi. Note plafonnée à 20.}

\end{document}
"
%   pdflatex -jobname=Eval_Reseau_CorrigeB "\def\variante{B}\def\modecorrige{1}% Évaluation pratique — Réseaux locaux (BTS CIEL 1re année) — sujets A et B
% Compilation (4 PDF) :
%   pdflatex -jobname=Eval_Reseau_SujetA   "\def\variante{A}\def\modecorrige{0}\input{Eval_TP_Reseau_LAN.tex}"
%   pdflatex -jobname=Eval_Reseau_SujetB   "\def\variante{B}\def\modecorrige{0}\input{Eval_TP_Reseau_LAN.tex}"
%   pdflatex -jobname=Eval_Reseau_CorrigeA "\def\variante{A}\def\modecorrige{1}\input{Eval_TP_Reseau_LAN.tex}"
%   pdflatex -jobname=Eval_Reseau_CorrigeB "\def\variante{B}\def\modecorrige{1}\input{Eval_TP_Reseau_LAN.tex}"
\documentclass[10pt,a4paper]{article}

\usepackage[T1]{fontenc}
\usepackage[utf8]{inputenc}
\usepackage{lmodern}
\usepackage[french]{babel}
\usepackage[a4paper,hmargin=1.5cm,top=1.3cm,bottom=1.6cm,footskip=0.7cm]{geometry}
\usepackage[table]{xcolor}
\usepackage{booktabs,tabularx,array}
\usepackage{enumitem}
\usepackage{fancyhdr}
\usepackage[most]{tcolorbox}
\usepackage{tikz}
\usetikzlibrary{positioning}
\usepackage{titlesec}

\providecommand{\modecorrige}{0}
\providecommand{\variante}{A}
\newif\ifcorrige
\ifnum\modecorrige=1 \corrigetrue \else \corrigefalse \fi

\definecolor{ipred}{HTML}{C0392B}
\definecolor{cdcbar}{HTML}{A67C3D}
\definecolor{cdcbg}{HTML}{F4F1EA}
\definecolor{qbar}{HTML}{4A6B5D}
\definecolor{qbg}{HTML}{EDF1EF}
\definecolor{corr}{HTML}{1F5FA8}
\definecolor{corrbg}{HTML}{E8F0FA}
\definecolor{grisfoot}{HTML}{828282}

\newcommand{\ip}[1]{\textcolor{ipred}{\texttt{#1}}}
\newcommand{\pts}[1]{\hfill\mbox{\textbf{(#1)}}}
\newcommand{\sol}[1]{\ifcorrige\textcolor{corr}{#1}\fi}
\newcommand{\bareme}[1]{\ifcorrige\par{\footnotesize\color{corr}#1\par}\fi}
\newcommand{\hl}{\ifcorrige\rule[-1.6mm]{0pt}{5.4mm}\else\rule[-2.2mm]{0pt}{6.6mm}\fi}
\newcolumntype{C}{>{\centering\arraybackslash}X}
\newcolumntype{L}{>{\raggedright\arraybackslash}X}

\titleformat{\section}{\large\bfseries}{}{0pt}{}
\titlespacing*{\section}{0pt}{6pt}{3pt}

\tcbset{qstyle/.style={enhanced,boxrule=0pt,leftrule=2.5pt,arc=0pt,outer arc=0pt,
  left=5pt,right=5pt,top=2pt,bottom=2pt,before skip=4pt,after skip=3pt}}
\newtcolorbox{cdc}{qstyle,colback=cdcbg,colframe=cdcbar}
\newtcolorbox{quest}{qstyle,colback=qbg,colframe=qbar}

% Zone de réponse : vierge dans le sujet, remplie dans le corrigé
\newcommand{\zone}[2]{%
  \ifcorrige
    \begin{tcolorbox}[colback=corrbg,colframe=corr,boxrule=0.5pt,arc=1pt,left=4pt,right=4pt,top=1pt,bottom=1pt,
      before skip=2pt,after skip=2pt,fontupper=\footnotesize\color{corr}]#2\end{tcolorbox}%
  \else
    \begin{tcolorbox}[colback=white,colframe=black!45,boxrule=0.4pt,arc=1pt,height=#1,
      before skip=2pt,after skip=2pt]\end{tcolorbox}%
  \fi}

\newcommand{\question}[3]{\begin{quest}\textbf{#1.} #3\pts{#2}\end{quest}}

%==========================================================================
% Données propres à chaque variante
%==========================================================================
\if\variante A
  \def\entreprise{Le cabinet d'architectes \textbf{ArchiNova}}
  \def\Rnom{R-ArchiNova}
  \def\resUn{Bureau d'études}\def\pUn{BE}\def\netUn{192.168.40}
  \def\resDeux{Accueil}\def\pDeux{ACC}\def\netDeux{172.16}
  \def\srvNom{SRV-Intranet}\def\netSrv{192.168.200}\def\srvH{10}
  \def\ssid{ArchiNova-WiFi}
  \def\AIrows{%
    \hl\ip{192.168.40.25} & \sol{C} & \sol{255.255.255.0} & \sol{192.168.40.0} & \sol{192.168.40.255} & \sol{254} \\ \hline
    \hl\ip{172.16.5.10}   & \sol{B} & \sol{255.255.0.0}   & \sol{172.16.0.0}   & \sol{172.16.255.255} & \sol{65\,534} \\ \hline
    \hl\ip{10.3.2.1}      & \sol{A} & \sol{255.0.0.0}     & \sol{10.0.0.0}     & \sol{10.255.255.255} & \sol{16\,777\,214} \\ \hline}
  \def\apipa{169.254.12.7}
  \def\Cnet{192.168.60}
  \def\Crows{%
    PC-A & \ip{192.168.60.10} & \ip{255.255.255.0} & \ip{192.168.60.254} & — \\
    PC-B & \ip{192.168.60.20} & \ip{255.255.255.0} & \ip{192.168.60.1}   & — \\
    PC-C & \ip{192.168.61.30} & \ip{255.255.255.0} & \ip{192.168.60.254} & — \\
    PC-D & \ip{192.168.60.10} & \ip{255.255.255.0} & \ip{192.168.60.254} & — \\
    SRV-C & \ip{10.0.0.5}     & \ip{255.0.0.0}     & \ip{10.0.0.254}     & — \\
    R1 — G0/0 & \ip{192.168.60.254} & \ip{255.255.255.0} & — & On \\
    R1 — G0/1 & \ip{10.0.0.254}     & \ip{255.0.0.0}     & — & Off \\}
  \def\Csol{%
    \hl 1 & \sol{PC-B} & \sol{Passerelle .60.1 : pas l'adresse du routeur.} & \sol{Joint son réseau, mais aucun autre (serveur).} & \sol{Passerelle 192.168.60.254.} \\ \hline
    \hl 2 & \sol{PC-C} & \sol{IP hors du réseau 192.168.60.0/24.} & \sol{Ne joint ni les postes ni le routeur.} & \sol{Ex. 192.168.60.30.} \\ \hline
    \hl 3 & \sol{PC-D} & \sol{Même IP que PC-A (doublon).} & \sol{Conflit d'adresse : PC-A et PC-D perturbés.} & \sol{Adresse unique, ex. .60.40.} \\ \hline
    \hl 4 & \sol{R1 G0/1} & \sol{Interface désactivée (Off).} & \sol{Réseau du serveur injoignable pour tous.} & \sol{Port Status : On.} \\ \hline}
  \def\Cbareme{PC-A et SRV-C sont correctement configurés ; accepter PC-A cité à la place de PC-D pour le doublon.}
\else
  \def\entreprise{La clinique vétérinaire \textbf{VetoPlus}}
  \def\Rnom{R-VetoPlus}
  \def\resUn{Soins}\def\pUn{SOI}\def\netUn{192.168.70}
  \def\resDeux{Secrétariat}\def\pDeux{SEC}\def\netDeux{172.20}
  \def\srvNom{SRV-Dossiers}\def\netSrv{192.168.250}\def\srvH{20}
  \def\ssid{VetoPlus-WiFi}
  \def\AIrows{%
    \hl\ip{192.168.15.130} & \sol{C} & \sol{255.255.255.0} & \sol{192.168.15.0} & \sol{192.168.15.255} & \sol{254} \\ \hline
    \hl\ip{172.31.200.7}   & \sol{B} & \sol{255.255.0.0}   & \sol{172.31.0.0}   & \sol{172.31.255.255} & \sol{65\,534} \\ \hline
    \hl\ip{10.45.6.9}      & \sol{A} & \sol{255.0.0.0}     & \sol{10.0.0.0}     & \sol{10.255.255.255} & \sol{16\,777\,214} \\ \hline}
  \def\apipa{169.254.88.3}
  \def\Cnet{192.168.80}
  \def\Crows{%
    PC-A & \ip{192.168.8.11}   & \ip{255.255.255.0} & \ip{192.168.80.254} & — \\
    PC-B & \ip{192.168.80.12}  & \ip{255.255.255.0} & \ip{192.168.80.254} & — \\
    PC-C & \ip{192.168.80.13}  & \ip{255.255.255.0} & \ip{192.168.80.100} & — \\
    PC-D & \ip{192.168.80.12}  & \ip{255.255.255.0} & \ip{192.168.80.254} & — \\
    SRV-C & \ip{10.0.0.8}      & \ip{255.0.0.0}     & \ip{10.0.0.254}     & — \\
    R1 — G0/0 & \ip{192.168.80.254} & \ip{255.255.255.0} & — & Off \\
    R1 — G0/1 & \ip{10.0.0.254}     & \ip{255.0.0.0}     & — & On \\}
  \def\Csol{%
    \hl 1 & \sol{PC-A} & \sol{IP hors du réseau 192.168.80.0/24.} & \sol{Ne joint ni les postes ni le routeur.} & \sol{Ex. 192.168.80.11.} \\ \hline
    \hl 2 & \sol{PC-C} & \sol{Passerelle .80.100 : pas l'adresse du routeur.} & \sol{Joint son réseau, mais aucun autre (serveur).} & \sol{Passerelle 192.168.80.254.} \\ \hline
    \hl 3 & \sol{PC-D} & \sol{Même IP que PC-B (doublon).} & \sol{Conflit d'adresse : PC-B et PC-D perturbés.} & \sol{Adresse unique, ex. .80.14.} \\ \hline
    \hl 4 & \sol{R1 G0/0} & \sol{Interface côté LAN désactivée (Off).} & \sol{Aucun poste ne joint le routeur ni le serveur.} & \sol{Port Status : On.} \\ \hline}
  \def\Cbareme{PC-B et SRV-C sont correctement configurés ; accepter PC-B cité à la place de PC-D pour le doublon.}
\fi

\pagestyle{fancy}
\fancyhf{}
\renewcommand{\headrulewidth}{0pt}
\fancyfoot[L]{\footnotesize\itshape\color{grisfoot}Évaluation — Réseaux locaux — Sujet \variante\ifcorrige\ — CORRIGÉ\fi}
\fancyfoot[R]{\footnotesize\color{grisfoot}\thepage/3}

\setlist{itemsep=0pt,topsep=2pt,parsep=0pt}
\setlength{\parindent}{0pt}
\setlength{\parskip}{2pt}

\begin{document}

\begin{center}
  {\Large\bfseries Évaluation pratique — Réseaux locaux \quad
   \fcolorbox{black}{white}{\ SUJET \variante\ }\par}
  \vspace{1pt}
  {\bfseries Adressage IP, commutation et routage — BTS CIEL 1re année — durée : 2 heures\par}
  \ifcorrige\vspace{2pt}{\bfseries\color{corr}CORRIGÉ ET BARÈME — document enseignant}\par\fi
\end{center}
\vspace{-2pt}
{\renewcommand{\arraystretch}{1.35}
\begin{tabularx}{\linewidth}{|X|X|p{2.2cm}|p{2.6cm}|}
  \hline
  \textbf{Nom :} & \textbf{Prénom :} & \textbf{Poste n°} & \textbf{Note :} \hfill /20 \\ \hline
\end{tabularx}}

\textbf{Consignes :} travail \textbf{individuel}, documents et Internet \textbf{non autorisés}. Ce sujet est le
\textbf{document-réponse}. Fichier Packet Tracer : \ip{Nom\_Prenom\_EvalLAN\_\variante.pkt}, à enregistrer
régulièrement dans le dossier de rendu. Les « visas » sont apposés par le professeur, que vous appelez à chaque test réussi.
Durées conseillées : A 30 min — B 1 h 10 — C 20 min.

%==========================================================================
\section{Partie A — Questions de cours et calculs \hfill /6}

\question{A1}{2 pts}{Complétez le tableau en considérant le \textbf{masque par défaut} de la classe de chaque adresse.}
{\small\renewcommand{\arraystretch}{1.15}
\begin{tabularx}{\linewidth}{|l|c|C|C|C|c|}
  \hline
  \textbf{Adresse IP} & \textbf{Classe} & \textbf{Masque par défaut} & \textbf{Adresse réseau} &
  \textbf{Adresse de diffusion} & \textbf{Nb d'hôtes} \\ \hline
  \AIrows
\end{tabularx}}
\bareme{0,5 pt par ligne juste + 0,5 pt si les trois nombres d'hôtes sont justes ($2^{n}-2$, $n$ = nombre de bits d'hôte).}

\question{A2}{1 pt}{Indiquez le type de câble (\textbf{droit} ou \textbf{croisé}) pour chaque liaison.}
{\small\renewcommand{\arraystretch}{1.15}
\begin{tabularx}{\linewidth}{|X|c||X|c|}
  \hline
  \hl PC $\leftrightarrow$ commutateur & \makebox[1.6cm]{\sol{droit}} & Commutateur $\leftrightarrow$ routeur & \makebox[1.6cm]{\sol{droit}} \\ \hline
  \hl PC $\leftrightarrow$ PC (réseau P2P) & \sol{croisé} & Concentrateur $\leftrightarrow$ concentrateur & \sol{croisé} \\ \hline
\end{tabularx}}
\bareme{0,25 pt par case. Équipements de types différents : câble droit ; de même type : câble croisé.}

\question{A3}{1 pt}{Comment un \textbf{concentrateur (hub)} et un \textbf{commutateur (switch)} traitent-ils une trame
reçue ? Lequel limite les collisions, et pourquoi ?}
\zone{2.9cm}{Le \textbf{hub} (couche 1) répète la trame sur \textbf{tous ses autres ports} : un seul domaine de collision
partagé. Le \textbf{switch} (couche 2) lit l'adresse MAC de destination et, grâce à sa \textbf{table MAC}, ne l'envoie
\textbf{que sur le port du destinataire} : un domaine de collision par port, c'est lui qui limite les collisions.
\emph{0,5 pt hub, 0,5 pt switch.}}

\question{A4}{1 pt}{Qu'est-ce que la \textbf{passerelle par défaut} d'un poste ? Quelle condition son adresse doit-elle respecter ?}
\zone{2.4cm}{Adresse IP de l'interface du \textbf{routeur} à laquelle le poste envoie les paquets destinés à un
\textbf{autre réseau}. Elle doit \textbf{appartenir au même réseau} que le poste. \emph{0,5 + 0,5 pt.}}

\question{A5}{0,5 pt}{Un portable configuré en DHCP affiche \ip{\apipa} / \ip{255.255.0.0}. Que s'est-il passé ?}
\zone{1.5cm}{Aucun \textbf{serveur DHCP} n'a répondu : le poste s'est attribué une adresse \textbf{APIPA} (169.254.0.0/16).}

\question{A6}{0,5 pt}{Quel protocole utilise \texttt{ping} ? Nommez les deux messages échangés et la couche OSI des adresses IP.}
\zone{1.5cm}{\textbf{ICMP} ; \textbf{Echo Request} / \textbf{Echo Reply} ; adresses IP en \textbf{couche 3 (Réseau)}.}

%==========================================================================
\newpage
\section{Partie B — Réalisation sous Packet Tracer \hfill /10}

\begin{cdc}
\entreprise{} vous confie son réseau : \textbf{trois réseaux} reliés par un \textbf{routeur unique Cisco 2911}
nommé \ip{\Rnom} (interfaces G0/0, G0/1, G0/2).
\begin{itemize}[leftmargin=12pt]
  \item \textbf{Réseau \resUn} : \textbf{4 postes} PC-\pUn1 à PC-\pUn4, commutateur \ip{SW-\pUn}, adressage statique,
        réseau \ip{\netUn.0/24}.
  \item \textbf{Réseau \resDeux} : \textbf{2 postes} PC-\pDeux1, PC-\pDeux2 et \textbf{1 imprimante} IMP-\pDeux,
        commutateur \ip{SW-\pDeux}, adressage statique, réseau de classe B \ip{\netDeux.0.0/16}.
  \item \textbf{Réseau Serveur} : serveur \ip{\srvNom} d'adresse \ip{\netSrv.\srvH/24}, relié \textbf{directement} au routeur.
  \item La passerelle de chaque réseau est sa \textbf{dernière adresse utilisable}. Tous les équipements communiquent
        entre eux et accèdent à la page web du serveur. Nommage imposé (« Display Name »).
\end{itemize}
\end{cdc}

\question{B1}{2 pts}{Complétez le plan d'adressage \textbf{avant} la réalisation.}
{\footnotesize\setlength{\tabcolsep}{3pt}\renewcommand{\arraystretch}{1.15}
\begin{tabularx}{\linewidth}{|l|C|C|C|C|C|C|}
  \hline
  \textbf{Réseau} & \textbf{Adresse réseau} & \textbf{Masque} & \textbf{1re @ utilisable} &
  \textbf{Dernière @ utilisable} & \textbf{Diffusion} & \textbf{Passerelle} \\ \hline
  \hl \resUn & \sol{\netUn.0} & \sol{255.255.255.0} & \sol{\netUn.1} & \sol{\netUn.254} & \sol{\netUn.255} & \sol{\netUn.254} \\ \hline
  \hl \resDeux & \sol{\netDeux.0.0} & \sol{255.255.0.0} & \sol{\netDeux.0.1} & \sol{\netDeux.255.254} & \sol{\netDeux.255.255} & \sol{\netDeux.255.254} \\ \hline
  \hl Serveur & \sol{\netSrv.0} & \sol{255.255.255.0} & \sol{\netSrv.1} & \sol{\netSrv.254} & \sol{\netSrv.255} & \sol{\netSrv.254} \\ \hline
\end{tabularx}}
\bareme{0,5 pt par ligne juste + 0,5 pt si la convention de passerelle est respectée partout.}

\question{B2}{2 pts}{Construisez la topologie (équipements, câblage, nommage) et relevez l'adressage configuré.}
{\footnotesize\setlength{\tabcolsep}{3pt}\renewcommand{\arraystretch}{1.1}
\begin{tabularx}{\linewidth}{|l|C|C|C||l|C|C|}
  \hline
  \textbf{Hôte} & \textbf{Adresse IP} & \textbf{Masque} & \textbf{Passerelle} & \textbf{Interface} & \textbf{Adresse IP} & \textbf{Masque} \\ \hline
  \hl PC-\pUn1 & \sol{\netUn.1} & \sol{255.255.255.0} & \sol{\netUn.254} & \Rnom{} G0/0 & \sol{\netUn.254} & \sol{255.255.255.0} \\ \hline
  \hl PC-\pDeux1 & \sol{\netDeux.0.1} & \sol{255.255.0.0} & \sol{\netDeux.255.254} & \Rnom{} G0/1 & \sol{\netDeux.255.254} & \sol{255.255.0.0} \\ \hline
  \hl IMP-\pDeux & \sol{\netDeux.0.10} & \sol{255.255.0.0} & \sol{\netDeux.255.254} & \Rnom{} G0/2 & \sol{\netSrv.254} & \sol{255.255.255.0} \\ \hline
  \hl \srvNom & \ip{\netSrv.\srvH} & \sol{255.255.255.0} & \sol{\netSrv.254} & \multicolumn{3}{c|}{\cellcolor{black!8}} \\ \hline
\end{tabularx}}
\bareme{Adresses d'hôtes libres (uniques, dans le bon réseau) ; affectation G0/x libre. Sur le .pkt : équipements 0,5 ;
câblage 1 (câbles droits vers les switchs et switch–routeur, \textbf{croisé} serveur–routeur, liens verts) ; nommage 0,5.}

\question{B3}{3 pts}{Configurez tous les équipements : IP, masques, passerelles des hôtes ; interfaces du routeur
(« Port Status : On »).}
\bareme{Sur le .pkt : interfaces routeur 1 pt ; IP/masques des hôtes 1 pt ; passerelles (imprimante et serveur compris)
1 pt. $-0,25$ par erreur, sans descendre sous 0 par item.}

\question{B4}{3 pts}{Réalisez les tests, notez la commande et le résultat, puis faites viser chaque test réussi.}
{\small\renewcommand{\arraystretch}{1.1}
\begin{tabularx}{\linewidth}{|c|l|l|C|c|c|}
  \hline
  \textbf{N°} & \textbf{Source} & \textbf{Destination} & \textbf{Commande / action} & \textbf{Résultat} & \textbf{Visa prof.} \\ \hline
  \hl 1 & PC-\pUn1 & PC-\pUn4 & \sol{ping \netUn.4} & \sol{réussi} & \hspace{1.4cm} \\ \hline
  \hl 2 & PC-\pDeux1 & IMP-\pDeux & \sol{ping \netDeux.0.10} & \sol{réussi} & \\ \hline
  \hl 3 & PC-\pUn2 & PC-\pDeux2 & \sol{ping \netDeux.0.2} & \sol{réussi} & \\ \hline
  \hl 4 & PC-\pDeux2 & \srvNom & \sol{ping \netSrv.\srvH} & \sol{réussi} & \\ \hline
  \hl 5 & PC-\pUn3 & page web du serveur & \sol{navigateur : {\NoAutoSpacing http://\netSrv.\srvH}} & \sol{page affichée} & \\ \hline
\end{tabularx}}
\bareme{0,5 pt par test visé + 0,5 pt pour B5. Le premier ping inter-réseaux peut échouer (résolution ARP) ; les suivants doivent réussir.}

\question{B5}{incl. B4}{Test n°3 : par quels équipements passe le paquet ? Pourquoi PC-\pUn2 l'envoie-t-il au routeur ?}
\zone{2.3cm}{PC-\pUn2 $\to$ SW-\pUn{} $\to$ \Rnom{} $\to$ SW-\pDeux{} $\to$ PC-\pDeux2. Grâce à son masque, PC-\pUn2 constate que
la destination (\netDeux.0.0/16) n'est \textbf{pas dans son réseau} (\netUn.0/24) : il envoie le paquet à sa
\textbf{passerelle par défaut}, qui le route.}

%==========================================================================
\newpage
\section{Partie C — Diagnostic d'un réseau en panne \hfill /4}

Un technicien constate qu'\textbf{aucun poste ne parvient à joindre le serveur SRV-C} et que certains postes
communiquent mal entre eux. Les câbles sont du bon type.

\begin{center}
\begin{minipage}{0.37\linewidth}\centering
\begin{tikzpicture}[font=\footnotesize,
  eq/.style={draw,fill=qbg,rounded corners=2pt,minimum width=9mm,minimum height=5mm,inner sep=2pt},
  every path/.style={thick}]
  \node[eq] (SW) {SW-1};
  \node[eq,right=9mm of SW] (R) {R1};
  \node[eq,right=7mm of R] (SRV) {SRV-C};
  \node[eq,above left=5mm and -3mm of SW] (A) {PC-A};
  \node[eq,above right=5mm and -3mm of SW] (B) {PC-B};
  \node[eq,below left=5mm and -3mm of SW] (C) {PC-C};
  \node[eq,below right=5mm and -3mm of SW] (D) {PC-D};
  \foreach \n in {A,B,C,D} \draw (\n) -- (SW);
  \draw (SW) -- node[above,font=\tiny]{G0/0} (R);
  \draw (R) -- node[above,font=\tiny]{G0/1} (SRV);
\end{tikzpicture}
\end{minipage}\hfill
\begin{minipage}{0.61\linewidth}\centering
{\small\renewcommand{\arraystretch}{1.05}\setlength{\tabcolsep}{4pt}
\begin{tabular}{lllll}
  \toprule
  \textbf{Équipement} & \textbf{Adresse IP} & \textbf{Masque} & \textbf{Passerelle} & \textbf{Port} \\
  \midrule
  \Crows
  \bottomrule
\end{tabular}}
\end{minipage}
\end{center}

\question{C1}{4 pts}{Relevez les \textbf{quatre anomalies}. Pour chacune : conséquence et correction.}
{\small\renewcommand{\arraystretch}{1.1}
\ifcorrige\def\hc{}\else\def\hc{\rule[-10mm]{0pt}{14mm}}\fi
\begin{tabularx}{\linewidth}{|c|>{\raggedright\arraybackslash}p{2cm}|L|L|L|}
  \hline
  \textbf{N°} & \textbf{Équipement} & \textbf{Anomalie} & \textbf{Conséquence} & \textbf{Correction} \\ \hline
  \ifcorrige\Csol\else
  \hc 1 & & & & \\ \hline
  \hc 2 & & & & \\ \hline
  \hc 3 & & & & \\ \hline
  \hc 4 & & & & \\ \hline
  \fi
\end{tabularx}}
\bareme{1 pt par anomalie (0,25 équipement, anomalie, conséquence, correction ; ordre indifférent). \Cbareme}

%==========================================================================
\section{Bonus — Extension sans fil \hfill +1}

\begin{cdc}Ajoutez au réseau \resUn{} un \textbf{AccessPoint-PT} relié à SW-\pUn{} (SSID \ip{\ssid}, \textbf{WPA2-PSK})
et un portable \ip{Laptop-\pUn} équipé d'un module \texttt{WPC300N}, en adressage \textbf{statique}.
Vérifiez qu'il joint \srvNom.\end{cdc}
{\small\renewcommand{\arraystretch}{1.1}
\begin{tabularx}{\linewidth}{|l|C|C|C|c|}
  \hline
  \textbf{Équipement} & \textbf{Adresse IP} & \textbf{Masque} & \textbf{Passerelle} & \textbf{Visa prof.} \\ \hline
  \hl Laptop-\pUn & \sol{\netUn.50} & \sol{255.255.255.0} & \sol{\netUn.254} & \hspace{1.4cm} \\ \hline
\end{tabularx}}
\bareme{Le point d'accès fait office de pont : le portable est dans \netUn.0/24. Bonus si associé en WPA2-PSK et ping
vers \netSrv.\srvH{} réussi. Note plafonnée à 20.}

\end{document}
"
\documentclass[10pt,a4paper]{article}

\usepackage[T1]{fontenc}
\usepackage[utf8]{inputenc}
\usepackage{lmodern}
\usepackage[french]{babel}
\usepackage[a4paper,hmargin=1.5cm,top=1.3cm,bottom=1.6cm,footskip=0.7cm]{geometry}
\usepackage[table]{xcolor}
\usepackage{booktabs,tabularx,array}
\usepackage{enumitem}
\usepackage{fancyhdr}
\usepackage[most]{tcolorbox}
\usepackage{tikz}
\usetikzlibrary{positioning}
\usepackage{titlesec}

\providecommand{\modecorrige}{0}
\providecommand{\variante}{A}
\newif\ifcorrige
\ifnum\modecorrige=1 \corrigetrue \else \corrigefalse \fi

\definecolor{ipred}{HTML}{C0392B}
\definecolor{cdcbar}{HTML}{A67C3D}
\definecolor{cdcbg}{HTML}{F4F1EA}
\definecolor{qbar}{HTML}{4A6B5D}
\definecolor{qbg}{HTML}{EDF1EF}
\definecolor{corr}{HTML}{1F5FA8}
\definecolor{corrbg}{HTML}{E8F0FA}
\definecolor{grisfoot}{HTML}{828282}

\newcommand{\ip}[1]{\textcolor{ipred}{\texttt{#1}}}
\newcommand{\pts}[1]{\hfill\mbox{\textbf{(#1)}}}
\newcommand{\sol}[1]{\ifcorrige\textcolor{corr}{#1}\fi}
\newcommand{\bareme}[1]{\ifcorrige\par{\footnotesize\color{corr}#1\par}\fi}
\newcommand{\hl}{\ifcorrige\rule[-1.6mm]{0pt}{5.4mm}\else\rule[-2.2mm]{0pt}{6.6mm}\fi}
\newcolumntype{C}{>{\centering\arraybackslash}X}
\newcolumntype{L}{>{\raggedright\arraybackslash}X}

\titleformat{\section}{\large\bfseries}{}{0pt}{}
\titlespacing*{\section}{0pt}{6pt}{3pt}

\tcbset{qstyle/.style={enhanced,boxrule=0pt,leftrule=2.5pt,arc=0pt,outer arc=0pt,
  left=5pt,right=5pt,top=2pt,bottom=2pt,before skip=4pt,after skip=3pt}}
\newtcolorbox{cdc}{qstyle,colback=cdcbg,colframe=cdcbar}
\newtcolorbox{quest}{qstyle,colback=qbg,colframe=qbar}

% Zone de réponse : vierge dans le sujet, remplie dans le corrigé
\newcommand{\zone}[2]{%
  \ifcorrige
    \begin{tcolorbox}[colback=corrbg,colframe=corr,boxrule=0.5pt,arc=1pt,left=4pt,right=4pt,top=1pt,bottom=1pt,
      before skip=2pt,after skip=2pt,fontupper=\footnotesize\color{corr}]#2\end{tcolorbox}%
  \else
    \begin{tcolorbox}[colback=white,colframe=black!45,boxrule=0.4pt,arc=1pt,height=#1,
      before skip=2pt,after skip=2pt]\end{tcolorbox}%
  \fi}

\newcommand{\question}[3]{\begin{quest}\textbf{#1.} #3\pts{#2}\end{quest}}

%==========================================================================
% Données propres à chaque variante
%==========================================================================
\if\variante A
  \def\entreprise{Le cabinet d'architectes \textbf{ArchiNova}}
  \def\Rnom{R-ArchiNova}
  \def\resUn{Bureau d'études}\def\pUn{BE}\def\netUn{192.168.40}
  \def\resDeux{Accueil}\def\pDeux{ACC}\def\netDeux{172.16}
  \def\srvNom{SRV-Intranet}\def\netSrv{192.168.200}\def\srvH{10}
  \def\ssid{ArchiNova-WiFi}
  \def\AIrows{%
    \hl\ip{192.168.40.25} & \sol{C} & \sol{255.255.255.0} & \sol{192.168.40.0} & \sol{192.168.40.255} & \sol{254} \\ \hline
    \hl\ip{172.16.5.10}   & \sol{B} & \sol{255.255.0.0}   & \sol{172.16.0.0}   & \sol{172.16.255.255} & \sol{65\,534} \\ \hline
    \hl\ip{10.3.2.1}      & \sol{A} & \sol{255.0.0.0}     & \sol{10.0.0.0}     & \sol{10.255.255.255} & \sol{16\,777\,214} \\ \hline}
  \def\apipa{169.254.12.7}
  \def\Cnet{192.168.60}
  \def\Crows{%
    PC-A & \ip{192.168.60.10} & \ip{255.255.255.0} & \ip{192.168.60.254} & — \\
    PC-B & \ip{192.168.60.20} & \ip{255.255.255.0} & \ip{192.168.60.1}   & — \\
    PC-C & \ip{192.168.61.30} & \ip{255.255.255.0} & \ip{192.168.60.254} & — \\
    PC-D & \ip{192.168.60.10} & \ip{255.255.255.0} & \ip{192.168.60.254} & — \\
    SRV-C & \ip{10.0.0.5}     & \ip{255.0.0.0}     & \ip{10.0.0.254}     & — \\
    R1 — G0/0 & \ip{192.168.60.254} & \ip{255.255.255.0} & — & On \\
    R1 — G0/1 & \ip{10.0.0.254}     & \ip{255.0.0.0}     & — & Off \\}
  \def\Csol{%
    \hl 1 & \sol{PC-B} & \sol{Passerelle .60.1 : pas l'adresse du routeur.} & \sol{Joint son réseau, mais aucun autre (serveur).} & \sol{Passerelle 192.168.60.254.} \\ \hline
    \hl 2 & \sol{PC-C} & \sol{IP hors du réseau 192.168.60.0/24.} & \sol{Ne joint ni les postes ni le routeur.} & \sol{Ex. 192.168.60.30.} \\ \hline
    \hl 3 & \sol{PC-D} & \sol{Même IP que PC-A (doublon).} & \sol{Conflit d'adresse : PC-A et PC-D perturbés.} & \sol{Adresse unique, ex. .60.40.} \\ \hline
    \hl 4 & \sol{R1 G0/1} & \sol{Interface désactivée (Off).} & \sol{Réseau du serveur injoignable pour tous.} & \sol{Port Status : On.} \\ \hline}
  \def\Cbareme{PC-A et SRV-C sont correctement configurés ; accepter PC-A cité à la place de PC-D pour le doublon.}
\else
  \def\entreprise{La clinique vétérinaire \textbf{VetoPlus}}
  \def\Rnom{R-VetoPlus}
  \def\resUn{Soins}\def\pUn{SOI}\def\netUn{192.168.70}
  \def\resDeux{Secrétariat}\def\pDeux{SEC}\def\netDeux{172.20}
  \def\srvNom{SRV-Dossiers}\def\netSrv{192.168.250}\def\srvH{20}
  \def\ssid{VetoPlus-WiFi}
  \def\AIrows{%
    \hl\ip{192.168.15.130} & \sol{C} & \sol{255.255.255.0} & \sol{192.168.15.0} & \sol{192.168.15.255} & \sol{254} \\ \hline
    \hl\ip{172.31.200.7}   & \sol{B} & \sol{255.255.0.0}   & \sol{172.31.0.0}   & \sol{172.31.255.255} & \sol{65\,534} \\ \hline
    \hl\ip{10.45.6.9}      & \sol{A} & \sol{255.0.0.0}     & \sol{10.0.0.0}     & \sol{10.255.255.255} & \sol{16\,777\,214} \\ \hline}
  \def\apipa{169.254.88.3}
  \def\Cnet{192.168.80}
  \def\Crows{%
    PC-A & \ip{192.168.8.11}   & \ip{255.255.255.0} & \ip{192.168.80.254} & — \\
    PC-B & \ip{192.168.80.12}  & \ip{255.255.255.0} & \ip{192.168.80.254} & — \\
    PC-C & \ip{192.168.80.13}  & \ip{255.255.255.0} & \ip{192.168.80.100} & — \\
    PC-D & \ip{192.168.80.12}  & \ip{255.255.255.0} & \ip{192.168.80.254} & — \\
    SRV-C & \ip{10.0.0.8}      & \ip{255.0.0.0}     & \ip{10.0.0.254}     & — \\
    R1 — G0/0 & \ip{192.168.80.254} & \ip{255.255.255.0} & — & Off \\
    R1 — G0/1 & \ip{10.0.0.254}     & \ip{255.0.0.0}     & — & On \\}
  \def\Csol{%
    \hl 1 & \sol{PC-A} & \sol{IP hors du réseau 192.168.80.0/24.} & \sol{Ne joint ni les postes ni le routeur.} & \sol{Ex. 192.168.80.11.} \\ \hline
    \hl 2 & \sol{PC-C} & \sol{Passerelle .80.100 : pas l'adresse du routeur.} & \sol{Joint son réseau, mais aucun autre (serveur).} & \sol{Passerelle 192.168.80.254.} \\ \hline
    \hl 3 & \sol{PC-D} & \sol{Même IP que PC-B (doublon).} & \sol{Conflit d'adresse : PC-B et PC-D perturbés.} & \sol{Adresse unique, ex. .80.14.} \\ \hline
    \hl 4 & \sol{R1 G0/0} & \sol{Interface côté LAN désactivée (Off).} & \sol{Aucun poste ne joint le routeur ni le serveur.} & \sol{Port Status : On.} \\ \hline}
  \def\Cbareme{PC-B et SRV-C sont correctement configurés ; accepter PC-B cité à la place de PC-D pour le doublon.}
\fi

\pagestyle{fancy}
\fancyhf{}
\renewcommand{\headrulewidth}{0pt}
\fancyfoot[L]{\footnotesize\itshape\color{grisfoot}Évaluation — Réseaux locaux — Sujet \variante\ifcorrige\ — CORRIGÉ\fi}
\fancyfoot[R]{\footnotesize\color{grisfoot}\thepage/3}

\setlist{itemsep=0pt,topsep=2pt,parsep=0pt}
\setlength{\parindent}{0pt}
\setlength{\parskip}{2pt}

\begin{document}

\begin{center}
  {\Large\bfseries Évaluation pratique — Réseaux locaux \quad
   \fcolorbox{black}{white}{\ SUJET \variante\ }\par}
  \vspace{1pt}
  {\bfseries Adressage IP, commutation et routage — BTS CIEL 1re année — durée : 2 heures\par}
  \ifcorrige\vspace{2pt}{\bfseries\color{corr}CORRIGÉ ET BARÈME — document enseignant}\par\fi
\end{center}
\vspace{-2pt}
{\renewcommand{\arraystretch}{1.35}
\begin{tabularx}{\linewidth}{|X|X|p{2.2cm}|p{2.6cm}|}
  \hline
  \textbf{Nom :} & \textbf{Prénom :} & \textbf{Poste n°} & \textbf{Note :} \hfill /20 \\ \hline
\end{tabularx}}

\textbf{Consignes :} travail \textbf{individuel}, documents et Internet \textbf{non autorisés}. Ce sujet est le
\textbf{document-réponse}. Fichier Packet Tracer : \ip{Nom\_Prenom\_EvalLAN\_\variante.pkt}, à enregistrer
régulièrement dans le dossier de rendu. Les « visas » sont apposés par le professeur, que vous appelez à chaque test réussi.
Durées conseillées : A 30 min — B 1 h 10 — C 20 min.

%==========================================================================
\section{Partie A — Questions de cours et calculs \hfill /6}

\question{A1}{2 pts}{Complétez le tableau en considérant le \textbf{masque par défaut} de la classe de chaque adresse.}
{\small\renewcommand{\arraystretch}{1.15}
\begin{tabularx}{\linewidth}{|l|c|C|C|C|c|}
  \hline
  \textbf{Adresse IP} & \textbf{Classe} & \textbf{Masque par défaut} & \textbf{Adresse réseau} &
  \textbf{Adresse de diffusion} & \textbf{Nb d'hôtes} \\ \hline
  \AIrows
\end{tabularx}}
\bareme{0,5 pt par ligne juste + 0,5 pt si les trois nombres d'hôtes sont justes ($2^{n}-2$, $n$ = nombre de bits d'hôte).}

\question{A2}{1 pt}{Indiquez le type de câble (\textbf{droit} ou \textbf{croisé}) pour chaque liaison.}
{\small\renewcommand{\arraystretch}{1.15}
\begin{tabularx}{\linewidth}{|X|c||X|c|}
  \hline
  \hl PC $\leftrightarrow$ commutateur & \makebox[1.6cm]{\sol{droit}} & Commutateur $\leftrightarrow$ routeur & \makebox[1.6cm]{\sol{droit}} \\ \hline
  \hl PC $\leftrightarrow$ PC (réseau P2P) & \sol{croisé} & Concentrateur $\leftrightarrow$ concentrateur & \sol{croisé} \\ \hline
\end{tabularx}}
\bareme{0,25 pt par case. Équipements de types différents : câble droit ; de même type : câble croisé.}

\question{A3}{1 pt}{Comment un \textbf{concentrateur (hub)} et un \textbf{commutateur (switch)} traitent-ils une trame
reçue ? Lequel limite les collisions, et pourquoi ?}
\zone{2.9cm}{Le \textbf{hub} (couche 1) répète la trame sur \textbf{tous ses autres ports} : un seul domaine de collision
partagé. Le \textbf{switch} (couche 2) lit l'adresse MAC de destination et, grâce à sa \textbf{table MAC}, ne l'envoie
\textbf{que sur le port du destinataire} : un domaine de collision par port, c'est lui qui limite les collisions.
\emph{0,5 pt hub, 0,5 pt switch.}}

\question{A4}{1 pt}{Qu'est-ce que la \textbf{passerelle par défaut} d'un poste ? Quelle condition son adresse doit-elle respecter ?}
\zone{2.4cm}{Adresse IP de l'interface du \textbf{routeur} à laquelle le poste envoie les paquets destinés à un
\textbf{autre réseau}. Elle doit \textbf{appartenir au même réseau} que le poste. \emph{0,5 + 0,5 pt.}}

\question{A5}{0,5 pt}{Un portable configuré en DHCP affiche \ip{\apipa} / \ip{255.255.0.0}. Que s'est-il passé ?}
\zone{1.5cm}{Aucun \textbf{serveur DHCP} n'a répondu : le poste s'est attribué une adresse \textbf{APIPA} (169.254.0.0/16).}

\question{A6}{0,5 pt}{Quel protocole utilise \texttt{ping} ? Nommez les deux messages échangés et la couche OSI des adresses IP.}
\zone{1.5cm}{\textbf{ICMP} ; \textbf{Echo Request} / \textbf{Echo Reply} ; adresses IP en \textbf{couche 3 (Réseau)}.}

%==========================================================================
\newpage
\section{Partie B — Réalisation sous Packet Tracer \hfill /10}

\begin{cdc}
\entreprise{} vous confie son réseau : \textbf{trois réseaux} reliés par un \textbf{routeur unique Cisco 2911}
nommé \ip{\Rnom} (interfaces G0/0, G0/1, G0/2).
\begin{itemize}[leftmargin=12pt]
  \item \textbf{Réseau \resUn} : \textbf{4 postes} PC-\pUn1 à PC-\pUn4, commutateur \ip{SW-\pUn}, adressage statique,
        réseau \ip{\netUn.0/24}.
  \item \textbf{Réseau \resDeux} : \textbf{2 postes} PC-\pDeux1, PC-\pDeux2 et \textbf{1 imprimante} IMP-\pDeux,
        commutateur \ip{SW-\pDeux}, adressage statique, réseau de classe B \ip{\netDeux.0.0/16}.
  \item \textbf{Réseau Serveur} : serveur \ip{\srvNom} d'adresse \ip{\netSrv.\srvH/24}, relié \textbf{directement} au routeur.
  \item La passerelle de chaque réseau est sa \textbf{dernière adresse utilisable}. Tous les équipements communiquent
        entre eux et accèdent à la page web du serveur. Nommage imposé (« Display Name »).
\end{itemize}
\end{cdc}

\question{B1}{2 pts}{Complétez le plan d'adressage \textbf{avant} la réalisation.}
{\footnotesize\setlength{\tabcolsep}{3pt}\renewcommand{\arraystretch}{1.15}
\begin{tabularx}{\linewidth}{|l|C|C|C|C|C|C|}
  \hline
  \textbf{Réseau} & \textbf{Adresse réseau} & \textbf{Masque} & \textbf{1re @ utilisable} &
  \textbf{Dernière @ utilisable} & \textbf{Diffusion} & \textbf{Passerelle} \\ \hline
  \hl \resUn & \sol{\netUn.0} & \sol{255.255.255.0} & \sol{\netUn.1} & \sol{\netUn.254} & \sol{\netUn.255} & \sol{\netUn.254} \\ \hline
  \hl \resDeux & \sol{\netDeux.0.0} & \sol{255.255.0.0} & \sol{\netDeux.0.1} & \sol{\netDeux.255.254} & \sol{\netDeux.255.255} & \sol{\netDeux.255.254} \\ \hline
  \hl Serveur & \sol{\netSrv.0} & \sol{255.255.255.0} & \sol{\netSrv.1} & \sol{\netSrv.254} & \sol{\netSrv.255} & \sol{\netSrv.254} \\ \hline
\end{tabularx}}
\bareme{0,5 pt par ligne juste + 0,5 pt si la convention de passerelle est respectée partout.}

\question{B2}{2 pts}{Construisez la topologie (équipements, câblage, nommage) et relevez l'adressage configuré.}
{\footnotesize\setlength{\tabcolsep}{3pt}\renewcommand{\arraystretch}{1.1}
\begin{tabularx}{\linewidth}{|l|C|C|C||l|C|C|}
  \hline
  \textbf{Hôte} & \textbf{Adresse IP} & \textbf{Masque} & \textbf{Passerelle} & \textbf{Interface} & \textbf{Adresse IP} & \textbf{Masque} \\ \hline
  \hl PC-\pUn1 & \sol{\netUn.1} & \sol{255.255.255.0} & \sol{\netUn.254} & \Rnom{} G0/0 & \sol{\netUn.254} & \sol{255.255.255.0} \\ \hline
  \hl PC-\pDeux1 & \sol{\netDeux.0.1} & \sol{255.255.0.0} & \sol{\netDeux.255.254} & \Rnom{} G0/1 & \sol{\netDeux.255.254} & \sol{255.255.0.0} \\ \hline
  \hl IMP-\pDeux & \sol{\netDeux.0.10} & \sol{255.255.0.0} & \sol{\netDeux.255.254} & \Rnom{} G0/2 & \sol{\netSrv.254} & \sol{255.255.255.0} \\ \hline
  \hl \srvNom & \ip{\netSrv.\srvH} & \sol{255.255.255.0} & \sol{\netSrv.254} & \multicolumn{3}{c|}{\cellcolor{black!8}} \\ \hline
\end{tabularx}}
\bareme{Adresses d'hôtes libres (uniques, dans le bon réseau) ; affectation G0/x libre. Sur le .pkt : équipements 0,5 ;
câblage 1 (câbles droits vers les switchs et switch–routeur, \textbf{croisé} serveur–routeur, liens verts) ; nommage 0,5.}

\question{B3}{3 pts}{Configurez tous les équipements : IP, masques, passerelles des hôtes ; interfaces du routeur
(« Port Status : On »).}
\bareme{Sur le .pkt : interfaces routeur 1 pt ; IP/masques des hôtes 1 pt ; passerelles (imprimante et serveur compris)
1 pt. $-0,25$ par erreur, sans descendre sous 0 par item.}

\question{B4}{3 pts}{Réalisez les tests, notez la commande et le résultat, puis faites viser chaque test réussi.}
{\small\renewcommand{\arraystretch}{1.1}
\begin{tabularx}{\linewidth}{|c|l|l|C|c|c|}
  \hline
  \textbf{N°} & \textbf{Source} & \textbf{Destination} & \textbf{Commande / action} & \textbf{Résultat} & \textbf{Visa prof.} \\ \hline
  \hl 1 & PC-\pUn1 & PC-\pUn4 & \sol{ping \netUn.4} & \sol{réussi} & \hspace{1.4cm} \\ \hline
  \hl 2 & PC-\pDeux1 & IMP-\pDeux & \sol{ping \netDeux.0.10} & \sol{réussi} & \\ \hline
  \hl 3 & PC-\pUn2 & PC-\pDeux2 & \sol{ping \netDeux.0.2} & \sol{réussi} & \\ \hline
  \hl 4 & PC-\pDeux2 & \srvNom & \sol{ping \netSrv.\srvH} & \sol{réussi} & \\ \hline
  \hl 5 & PC-\pUn3 & page web du serveur & \sol{navigateur : {\NoAutoSpacing http://\netSrv.\srvH}} & \sol{page affichée} & \\ \hline
\end{tabularx}}
\bareme{0,5 pt par test visé + 0,5 pt pour B5. Le premier ping inter-réseaux peut échouer (résolution ARP) ; les suivants doivent réussir.}

\question{B5}{incl. B4}{Test n°3 : par quels équipements passe le paquet ? Pourquoi PC-\pUn2 l'envoie-t-il au routeur ?}
\zone{2.3cm}{PC-\pUn2 $\to$ SW-\pUn{} $\to$ \Rnom{} $\to$ SW-\pDeux{} $\to$ PC-\pDeux2. Grâce à son masque, PC-\pUn2 constate que
la destination (\netDeux.0.0/16) n'est \textbf{pas dans son réseau} (\netUn.0/24) : il envoie le paquet à sa
\textbf{passerelle par défaut}, qui le route.}

%==========================================================================
\newpage
\section{Partie C — Diagnostic d'un réseau en panne \hfill /4}

Un technicien constate qu'\textbf{aucun poste ne parvient à joindre le serveur SRV-C} et que certains postes
communiquent mal entre eux. Les câbles sont du bon type.

\begin{center}
\begin{minipage}{0.37\linewidth}\centering
\begin{tikzpicture}[font=\footnotesize,
  eq/.style={draw,fill=qbg,rounded corners=2pt,minimum width=9mm,minimum height=5mm,inner sep=2pt},
  every path/.style={thick}]
  \node[eq] (SW) {SW-1};
  \node[eq,right=9mm of SW] (R) {R1};
  \node[eq,right=7mm of R] (SRV) {SRV-C};
  \node[eq,above left=5mm and -3mm of SW] (A) {PC-A};
  \node[eq,above right=5mm and -3mm of SW] (B) {PC-B};
  \node[eq,below left=5mm and -3mm of SW] (C) {PC-C};
  \node[eq,below right=5mm and -3mm of SW] (D) {PC-D};
  \foreach \n in {A,B,C,D} \draw (\n) -- (SW);
  \draw (SW) -- node[above,font=\tiny]{G0/0} (R);
  \draw (R) -- node[above,font=\tiny]{G0/1} (SRV);
\end{tikzpicture}
\end{minipage}\hfill
\begin{minipage}{0.61\linewidth}\centering
{\small\renewcommand{\arraystretch}{1.05}\setlength{\tabcolsep}{4pt}
\begin{tabular}{lllll}
  \toprule
  \textbf{Équipement} & \textbf{Adresse IP} & \textbf{Masque} & \textbf{Passerelle} & \textbf{Port} \\
  \midrule
  \Crows
  \bottomrule
\end{tabular}}
\end{minipage}
\end{center}

\question{C1}{4 pts}{Relevez les \textbf{quatre anomalies}. Pour chacune : conséquence et correction.}
{\small\renewcommand{\arraystretch}{1.1}
\ifcorrige\def\hc{}\else\def\hc{\rule[-10mm]{0pt}{14mm}}\fi
\begin{tabularx}{\linewidth}{|c|>{\raggedright\arraybackslash}p{2cm}|L|L|L|}
  \hline
  \textbf{N°} & \textbf{Équipement} & \textbf{Anomalie} & \textbf{Conséquence} & \textbf{Correction} \\ \hline
  \ifcorrige\Csol\else
  \hc 1 & & & & \\ \hline
  \hc 2 & & & & \\ \hline
  \hc 3 & & & & \\ \hline
  \hc 4 & & & & \\ \hline
  \fi
\end{tabularx}}
\bareme{1 pt par anomalie (0,25 équipement, anomalie, conséquence, correction ; ordre indifférent). \Cbareme}

%==========================================================================
\section{Bonus — Extension sans fil \hfill +1}

\begin{cdc}Ajoutez au réseau \resUn{} un \textbf{AccessPoint-PT} relié à SW-\pUn{} (SSID \ip{\ssid}, \textbf{WPA2-PSK})
et un portable \ip{Laptop-\pUn} équipé d'un module \texttt{WPC300N}, en adressage \textbf{statique}.
Vérifiez qu'il joint \srvNom.\end{cdc}
{\small\renewcommand{\arraystretch}{1.1}
\begin{tabularx}{\linewidth}{|l|C|C|C|c|}
  \hline
  \textbf{Équipement} & \textbf{Adresse IP} & \textbf{Masque} & \textbf{Passerelle} & \textbf{Visa prof.} \\ \hline
  \hl Laptop-\pUn & \sol{\netUn.50} & \sol{255.255.255.0} & \sol{\netUn.254} & \hspace{1.4cm} \\ \hline
\end{tabularx}}
\bareme{Le point d'accès fait office de pont : le portable est dans \netUn.0/24. Bonus si associé en WPA2-PSK et ping
vers \netSrv.\srvH{} réussi. Note plafonnée à 20.}

\end{document}
"
%   pdflatex -jobname=Eval_Reseau_CorrigeA "\def\variante{A}\def\modecorrige{1}% Évaluation pratique — Réseaux locaux (BTS CIEL 1re année) — sujets A et B
% Compilation (4 PDF) :
%   pdflatex -jobname=Eval_Reseau_SujetA   "\def\variante{A}\def\modecorrige{0}% Évaluation pratique — Réseaux locaux (BTS CIEL 1re année) — sujets A et B
% Compilation (4 PDF) :
%   pdflatex -jobname=Eval_Reseau_SujetA   "\def\variante{A}\def\modecorrige{0}\input{Eval_TP_Reseau_LAN.tex}"
%   pdflatex -jobname=Eval_Reseau_SujetB   "\def\variante{B}\def\modecorrige{0}\input{Eval_TP_Reseau_LAN.tex}"
%   pdflatex -jobname=Eval_Reseau_CorrigeA "\def\variante{A}\def\modecorrige{1}\input{Eval_TP_Reseau_LAN.tex}"
%   pdflatex -jobname=Eval_Reseau_CorrigeB "\def\variante{B}\def\modecorrige{1}\input{Eval_TP_Reseau_LAN.tex}"
\documentclass[10pt,a4paper]{article}

\usepackage[T1]{fontenc}
\usepackage[utf8]{inputenc}
\usepackage{lmodern}
\usepackage[french]{babel}
\usepackage[a4paper,hmargin=1.5cm,top=1.3cm,bottom=1.6cm,footskip=0.7cm]{geometry}
\usepackage[table]{xcolor}
\usepackage{booktabs,tabularx,array}
\usepackage{enumitem}
\usepackage{fancyhdr}
\usepackage[most]{tcolorbox}
\usepackage{tikz}
\usetikzlibrary{positioning}
\usepackage{titlesec}

\providecommand{\modecorrige}{0}
\providecommand{\variante}{A}
\newif\ifcorrige
\ifnum\modecorrige=1 \corrigetrue \else \corrigefalse \fi

\definecolor{ipred}{HTML}{C0392B}
\definecolor{cdcbar}{HTML}{A67C3D}
\definecolor{cdcbg}{HTML}{F4F1EA}
\definecolor{qbar}{HTML}{4A6B5D}
\definecolor{qbg}{HTML}{EDF1EF}
\definecolor{corr}{HTML}{1F5FA8}
\definecolor{corrbg}{HTML}{E8F0FA}
\definecolor{grisfoot}{HTML}{828282}

\newcommand{\ip}[1]{\textcolor{ipred}{\texttt{#1}}}
\newcommand{\pts}[1]{\hfill\mbox{\textbf{(#1)}}}
\newcommand{\sol}[1]{\ifcorrige\textcolor{corr}{#1}\fi}
\newcommand{\bareme}[1]{\ifcorrige\par{\footnotesize\color{corr}#1\par}\fi}
\newcommand{\hl}{\ifcorrige\rule[-1.6mm]{0pt}{5.4mm}\else\rule[-2.2mm]{0pt}{6.6mm}\fi}
\newcolumntype{C}{>{\centering\arraybackslash}X}
\newcolumntype{L}{>{\raggedright\arraybackslash}X}

\titleformat{\section}{\large\bfseries}{}{0pt}{}
\titlespacing*{\section}{0pt}{6pt}{3pt}

\tcbset{qstyle/.style={enhanced,boxrule=0pt,leftrule=2.5pt,arc=0pt,outer arc=0pt,
  left=5pt,right=5pt,top=2pt,bottom=2pt,before skip=4pt,after skip=3pt}}
\newtcolorbox{cdc}{qstyle,colback=cdcbg,colframe=cdcbar}
\newtcolorbox{quest}{qstyle,colback=qbg,colframe=qbar}

% Zone de réponse : vierge dans le sujet, remplie dans le corrigé
\newcommand{\zone}[2]{%
  \ifcorrige
    \begin{tcolorbox}[colback=corrbg,colframe=corr,boxrule=0.5pt,arc=1pt,left=4pt,right=4pt,top=1pt,bottom=1pt,
      before skip=2pt,after skip=2pt,fontupper=\footnotesize\color{corr}]#2\end{tcolorbox}%
  \else
    \begin{tcolorbox}[colback=white,colframe=black!45,boxrule=0.4pt,arc=1pt,height=#1,
      before skip=2pt,after skip=2pt]\end{tcolorbox}%
  \fi}

\newcommand{\question}[3]{\begin{quest}\textbf{#1.} #3\pts{#2}\end{quest}}

%==========================================================================
% Données propres à chaque variante
%==========================================================================
\if\variante A
  \def\entreprise{Le cabinet d'architectes \textbf{ArchiNova}}
  \def\Rnom{R-ArchiNova}
  \def\resUn{Bureau d'études}\def\pUn{BE}\def\netUn{192.168.40}
  \def\resDeux{Accueil}\def\pDeux{ACC}\def\netDeux{172.16}
  \def\srvNom{SRV-Intranet}\def\netSrv{192.168.200}\def\srvH{10}
  \def\ssid{ArchiNova-WiFi}
  \def\AIrows{%
    \hl\ip{192.168.40.25} & \sol{C} & \sol{255.255.255.0} & \sol{192.168.40.0} & \sol{192.168.40.255} & \sol{254} \\ \hline
    \hl\ip{172.16.5.10}   & \sol{B} & \sol{255.255.0.0}   & \sol{172.16.0.0}   & \sol{172.16.255.255} & \sol{65\,534} \\ \hline
    \hl\ip{10.3.2.1}      & \sol{A} & \sol{255.0.0.0}     & \sol{10.0.0.0}     & \sol{10.255.255.255} & \sol{16\,777\,214} \\ \hline}
  \def\apipa{169.254.12.7}
  \def\Cnet{192.168.60}
  \def\Crows{%
    PC-A & \ip{192.168.60.10} & \ip{255.255.255.0} & \ip{192.168.60.254} & — \\
    PC-B & \ip{192.168.60.20} & \ip{255.255.255.0} & \ip{192.168.60.1}   & — \\
    PC-C & \ip{192.168.61.30} & \ip{255.255.255.0} & \ip{192.168.60.254} & — \\
    PC-D & \ip{192.168.60.10} & \ip{255.255.255.0} & \ip{192.168.60.254} & — \\
    SRV-C & \ip{10.0.0.5}     & \ip{255.0.0.0}     & \ip{10.0.0.254}     & — \\
    R1 — G0/0 & \ip{192.168.60.254} & \ip{255.255.255.0} & — & On \\
    R1 — G0/1 & \ip{10.0.0.254}     & \ip{255.0.0.0}     & — & Off \\}
  \def\Csol{%
    \hl 1 & \sol{PC-B} & \sol{Passerelle .60.1 : pas l'adresse du routeur.} & \sol{Joint son réseau, mais aucun autre (serveur).} & \sol{Passerelle 192.168.60.254.} \\ \hline
    \hl 2 & \sol{PC-C} & \sol{IP hors du réseau 192.168.60.0/24.} & \sol{Ne joint ni les postes ni le routeur.} & \sol{Ex. 192.168.60.30.} \\ \hline
    \hl 3 & \sol{PC-D} & \sol{Même IP que PC-A (doublon).} & \sol{Conflit d'adresse : PC-A et PC-D perturbés.} & \sol{Adresse unique, ex. .60.40.} \\ \hline
    \hl 4 & \sol{R1 G0/1} & \sol{Interface désactivée (Off).} & \sol{Réseau du serveur injoignable pour tous.} & \sol{Port Status : On.} \\ \hline}
  \def\Cbareme{PC-A et SRV-C sont correctement configurés ; accepter PC-A cité à la place de PC-D pour le doublon.}
\else
  \def\entreprise{La clinique vétérinaire \textbf{VetoPlus}}
  \def\Rnom{R-VetoPlus}
  \def\resUn{Soins}\def\pUn{SOI}\def\netUn{192.168.70}
  \def\resDeux{Secrétariat}\def\pDeux{SEC}\def\netDeux{172.20}
  \def\srvNom{SRV-Dossiers}\def\netSrv{192.168.250}\def\srvH{20}
  \def\ssid{VetoPlus-WiFi}
  \def\AIrows{%
    \hl\ip{192.168.15.130} & \sol{C} & \sol{255.255.255.0} & \sol{192.168.15.0} & \sol{192.168.15.255} & \sol{254} \\ \hline
    \hl\ip{172.31.200.7}   & \sol{B} & \sol{255.255.0.0}   & \sol{172.31.0.0}   & \sol{172.31.255.255} & \sol{65\,534} \\ \hline
    \hl\ip{10.45.6.9}      & \sol{A} & \sol{255.0.0.0}     & \sol{10.0.0.0}     & \sol{10.255.255.255} & \sol{16\,777\,214} \\ \hline}
  \def\apipa{169.254.88.3}
  \def\Cnet{192.168.80}
  \def\Crows{%
    PC-A & \ip{192.168.8.11}   & \ip{255.255.255.0} & \ip{192.168.80.254} & — \\
    PC-B & \ip{192.168.80.12}  & \ip{255.255.255.0} & \ip{192.168.80.254} & — \\
    PC-C & \ip{192.168.80.13}  & \ip{255.255.255.0} & \ip{192.168.80.100} & — \\
    PC-D & \ip{192.168.80.12}  & \ip{255.255.255.0} & \ip{192.168.80.254} & — \\
    SRV-C & \ip{10.0.0.8}      & \ip{255.0.0.0}     & \ip{10.0.0.254}     & — \\
    R1 — G0/0 & \ip{192.168.80.254} & \ip{255.255.255.0} & — & Off \\
    R1 — G0/1 & \ip{10.0.0.254}     & \ip{255.0.0.0}     & — & On \\}
  \def\Csol{%
    \hl 1 & \sol{PC-A} & \sol{IP hors du réseau 192.168.80.0/24.} & \sol{Ne joint ni les postes ni le routeur.} & \sol{Ex. 192.168.80.11.} \\ \hline
    \hl 2 & \sol{PC-C} & \sol{Passerelle .80.100 : pas l'adresse du routeur.} & \sol{Joint son réseau, mais aucun autre (serveur).} & \sol{Passerelle 192.168.80.254.} \\ \hline
    \hl 3 & \sol{PC-D} & \sol{Même IP que PC-B (doublon).} & \sol{Conflit d'adresse : PC-B et PC-D perturbés.} & \sol{Adresse unique, ex. .80.14.} \\ \hline
    \hl 4 & \sol{R1 G0/0} & \sol{Interface côté LAN désactivée (Off).} & \sol{Aucun poste ne joint le routeur ni le serveur.} & \sol{Port Status : On.} \\ \hline}
  \def\Cbareme{PC-B et SRV-C sont correctement configurés ; accepter PC-B cité à la place de PC-D pour le doublon.}
\fi

\pagestyle{fancy}
\fancyhf{}
\renewcommand{\headrulewidth}{0pt}
\fancyfoot[L]{\footnotesize\itshape\color{grisfoot}Évaluation — Réseaux locaux — Sujet \variante\ifcorrige\ — CORRIGÉ\fi}
\fancyfoot[R]{\footnotesize\color{grisfoot}\thepage/3}

\setlist{itemsep=0pt,topsep=2pt,parsep=0pt}
\setlength{\parindent}{0pt}
\setlength{\parskip}{2pt}

\begin{document}

\begin{center}
  {\Large\bfseries Évaluation pratique — Réseaux locaux \quad
   \fcolorbox{black}{white}{\ SUJET \variante\ }\par}
  \vspace{1pt}
  {\bfseries Adressage IP, commutation et routage — BTS CIEL 1re année — durée : 2 heures\par}
  \ifcorrige\vspace{2pt}{\bfseries\color{corr}CORRIGÉ ET BARÈME — document enseignant}\par\fi
\end{center}
\vspace{-2pt}
{\renewcommand{\arraystretch}{1.35}
\begin{tabularx}{\linewidth}{|X|X|p{2.2cm}|p{2.6cm}|}
  \hline
  \textbf{Nom :} & \textbf{Prénom :} & \textbf{Poste n°} & \textbf{Note :} \hfill /20 \\ \hline
\end{tabularx}}

\textbf{Consignes :} travail \textbf{individuel}, documents et Internet \textbf{non autorisés}. Ce sujet est le
\textbf{document-réponse}. Fichier Packet Tracer : \ip{Nom\_Prenom\_EvalLAN\_\variante.pkt}, à enregistrer
régulièrement dans le dossier de rendu. Les « visas » sont apposés par le professeur, que vous appelez à chaque test réussi.
Durées conseillées : A 30 min — B 1 h 10 — C 20 min.

%==========================================================================
\section{Partie A — Questions de cours et calculs \hfill /6}

\question{A1}{2 pts}{Complétez le tableau en considérant le \textbf{masque par défaut} de la classe de chaque adresse.}
{\small\renewcommand{\arraystretch}{1.15}
\begin{tabularx}{\linewidth}{|l|c|C|C|C|c|}
  \hline
  \textbf{Adresse IP} & \textbf{Classe} & \textbf{Masque par défaut} & \textbf{Adresse réseau} &
  \textbf{Adresse de diffusion} & \textbf{Nb d'hôtes} \\ \hline
  \AIrows
\end{tabularx}}
\bareme{0,5 pt par ligne juste + 0,5 pt si les trois nombres d'hôtes sont justes ($2^{n}-2$, $n$ = nombre de bits d'hôte).}

\question{A2}{1 pt}{Indiquez le type de câble (\textbf{droit} ou \textbf{croisé}) pour chaque liaison.}
{\small\renewcommand{\arraystretch}{1.15}
\begin{tabularx}{\linewidth}{|X|c||X|c|}
  \hline
  \hl PC $\leftrightarrow$ commutateur & \makebox[1.6cm]{\sol{droit}} & Commutateur $\leftrightarrow$ routeur & \makebox[1.6cm]{\sol{droit}} \\ \hline
  \hl PC $\leftrightarrow$ PC (réseau P2P) & \sol{croisé} & Concentrateur $\leftrightarrow$ concentrateur & \sol{croisé} \\ \hline
\end{tabularx}}
\bareme{0,25 pt par case. Équipements de types différents : câble droit ; de même type : câble croisé.}

\question{A3}{1 pt}{Comment un \textbf{concentrateur (hub)} et un \textbf{commutateur (switch)} traitent-ils une trame
reçue ? Lequel limite les collisions, et pourquoi ?}
\zone{2.9cm}{Le \textbf{hub} (couche 1) répète la trame sur \textbf{tous ses autres ports} : un seul domaine de collision
partagé. Le \textbf{switch} (couche 2) lit l'adresse MAC de destination et, grâce à sa \textbf{table MAC}, ne l'envoie
\textbf{que sur le port du destinataire} : un domaine de collision par port, c'est lui qui limite les collisions.
\emph{0,5 pt hub, 0,5 pt switch.}}

\question{A4}{1 pt}{Qu'est-ce que la \textbf{passerelle par défaut} d'un poste ? Quelle condition son adresse doit-elle respecter ?}
\zone{2.4cm}{Adresse IP de l'interface du \textbf{routeur} à laquelle le poste envoie les paquets destinés à un
\textbf{autre réseau}. Elle doit \textbf{appartenir au même réseau} que le poste. \emph{0,5 + 0,5 pt.}}

\question{A5}{0,5 pt}{Un portable configuré en DHCP affiche \ip{\apipa} / \ip{255.255.0.0}. Que s'est-il passé ?}
\zone{1.5cm}{Aucun \textbf{serveur DHCP} n'a répondu : le poste s'est attribué une adresse \textbf{APIPA} (169.254.0.0/16).}

\question{A6}{0,5 pt}{Quel protocole utilise \texttt{ping} ? Nommez les deux messages échangés et la couche OSI des adresses IP.}
\zone{1.5cm}{\textbf{ICMP} ; \textbf{Echo Request} / \textbf{Echo Reply} ; adresses IP en \textbf{couche 3 (Réseau)}.}

%==========================================================================
\newpage
\section{Partie B — Réalisation sous Packet Tracer \hfill /10}

\begin{cdc}
\entreprise{} vous confie son réseau : \textbf{trois réseaux} reliés par un \textbf{routeur unique Cisco 2911}
nommé \ip{\Rnom} (interfaces G0/0, G0/1, G0/2).
\begin{itemize}[leftmargin=12pt]
  \item \textbf{Réseau \resUn} : \textbf{4 postes} PC-\pUn1 à PC-\pUn4, commutateur \ip{SW-\pUn}, adressage statique,
        réseau \ip{\netUn.0/24}.
  \item \textbf{Réseau \resDeux} : \textbf{2 postes} PC-\pDeux1, PC-\pDeux2 et \textbf{1 imprimante} IMP-\pDeux,
        commutateur \ip{SW-\pDeux}, adressage statique, réseau de classe B \ip{\netDeux.0.0/16}.
  \item \textbf{Réseau Serveur} : serveur \ip{\srvNom} d'adresse \ip{\netSrv.\srvH/24}, relié \textbf{directement} au routeur.
  \item La passerelle de chaque réseau est sa \textbf{dernière adresse utilisable}. Tous les équipements communiquent
        entre eux et accèdent à la page web du serveur. Nommage imposé (« Display Name »).
\end{itemize}
\end{cdc}

\question{B1}{2 pts}{Complétez le plan d'adressage \textbf{avant} la réalisation.}
{\footnotesize\setlength{\tabcolsep}{3pt}\renewcommand{\arraystretch}{1.15}
\begin{tabularx}{\linewidth}{|l|C|C|C|C|C|C|}
  \hline
  \textbf{Réseau} & \textbf{Adresse réseau} & \textbf{Masque} & \textbf{1re @ utilisable} &
  \textbf{Dernière @ utilisable} & \textbf{Diffusion} & \textbf{Passerelle} \\ \hline
  \hl \resUn & \sol{\netUn.0} & \sol{255.255.255.0} & \sol{\netUn.1} & \sol{\netUn.254} & \sol{\netUn.255} & \sol{\netUn.254} \\ \hline
  \hl \resDeux & \sol{\netDeux.0.0} & \sol{255.255.0.0} & \sol{\netDeux.0.1} & \sol{\netDeux.255.254} & \sol{\netDeux.255.255} & \sol{\netDeux.255.254} \\ \hline
  \hl Serveur & \sol{\netSrv.0} & \sol{255.255.255.0} & \sol{\netSrv.1} & \sol{\netSrv.254} & \sol{\netSrv.255} & \sol{\netSrv.254} \\ \hline
\end{tabularx}}
\bareme{0,5 pt par ligne juste + 0,5 pt si la convention de passerelle est respectée partout.}

\question{B2}{2 pts}{Construisez la topologie (équipements, câblage, nommage) et relevez l'adressage configuré.}
{\footnotesize\setlength{\tabcolsep}{3pt}\renewcommand{\arraystretch}{1.1}
\begin{tabularx}{\linewidth}{|l|C|C|C||l|C|C|}
  \hline
  \textbf{Hôte} & \textbf{Adresse IP} & \textbf{Masque} & \textbf{Passerelle} & \textbf{Interface} & \textbf{Adresse IP} & \textbf{Masque} \\ \hline
  \hl PC-\pUn1 & \sol{\netUn.1} & \sol{255.255.255.0} & \sol{\netUn.254} & \Rnom{} G0/0 & \sol{\netUn.254} & \sol{255.255.255.0} \\ \hline
  \hl PC-\pDeux1 & \sol{\netDeux.0.1} & \sol{255.255.0.0} & \sol{\netDeux.255.254} & \Rnom{} G0/1 & \sol{\netDeux.255.254} & \sol{255.255.0.0} \\ \hline
  \hl IMP-\pDeux & \sol{\netDeux.0.10} & \sol{255.255.0.0} & \sol{\netDeux.255.254} & \Rnom{} G0/2 & \sol{\netSrv.254} & \sol{255.255.255.0} \\ \hline
  \hl \srvNom & \ip{\netSrv.\srvH} & \sol{255.255.255.0} & \sol{\netSrv.254} & \multicolumn{3}{c|}{\cellcolor{black!8}} \\ \hline
\end{tabularx}}
\bareme{Adresses d'hôtes libres (uniques, dans le bon réseau) ; affectation G0/x libre. Sur le .pkt : équipements 0,5 ;
câblage 1 (câbles droits vers les switchs et switch–routeur, \textbf{croisé} serveur–routeur, liens verts) ; nommage 0,5.}

\question{B3}{3 pts}{Configurez tous les équipements : IP, masques, passerelles des hôtes ; interfaces du routeur
(« Port Status : On »).}
\bareme{Sur le .pkt : interfaces routeur 1 pt ; IP/masques des hôtes 1 pt ; passerelles (imprimante et serveur compris)
1 pt. $-0,25$ par erreur, sans descendre sous 0 par item.}

\question{B4}{3 pts}{Réalisez les tests, notez la commande et le résultat, puis faites viser chaque test réussi.}
{\small\renewcommand{\arraystretch}{1.1}
\begin{tabularx}{\linewidth}{|c|l|l|C|c|c|}
  \hline
  \textbf{N°} & \textbf{Source} & \textbf{Destination} & \textbf{Commande / action} & \textbf{Résultat} & \textbf{Visa prof.} \\ \hline
  \hl 1 & PC-\pUn1 & PC-\pUn4 & \sol{ping \netUn.4} & \sol{réussi} & \hspace{1.4cm} \\ \hline
  \hl 2 & PC-\pDeux1 & IMP-\pDeux & \sol{ping \netDeux.0.10} & \sol{réussi} & \\ \hline
  \hl 3 & PC-\pUn2 & PC-\pDeux2 & \sol{ping \netDeux.0.2} & \sol{réussi} & \\ \hline
  \hl 4 & PC-\pDeux2 & \srvNom & \sol{ping \netSrv.\srvH} & \sol{réussi} & \\ \hline
  \hl 5 & PC-\pUn3 & page web du serveur & \sol{navigateur : {\NoAutoSpacing http://\netSrv.\srvH}} & \sol{page affichée} & \\ \hline
\end{tabularx}}
\bareme{0,5 pt par test visé + 0,5 pt pour B5. Le premier ping inter-réseaux peut échouer (résolution ARP) ; les suivants doivent réussir.}

\question{B5}{incl. B4}{Test n°3 : par quels équipements passe le paquet ? Pourquoi PC-\pUn2 l'envoie-t-il au routeur ?}
\zone{2.3cm}{PC-\pUn2 $\to$ SW-\pUn{} $\to$ \Rnom{} $\to$ SW-\pDeux{} $\to$ PC-\pDeux2. Grâce à son masque, PC-\pUn2 constate que
la destination (\netDeux.0.0/16) n'est \textbf{pas dans son réseau} (\netUn.0/24) : il envoie le paquet à sa
\textbf{passerelle par défaut}, qui le route.}

%==========================================================================
\newpage
\section{Partie C — Diagnostic d'un réseau en panne \hfill /4}

Un technicien constate qu'\textbf{aucun poste ne parvient à joindre le serveur SRV-C} et que certains postes
communiquent mal entre eux. Les câbles sont du bon type.

\begin{center}
\begin{minipage}{0.37\linewidth}\centering
\begin{tikzpicture}[font=\footnotesize,
  eq/.style={draw,fill=qbg,rounded corners=2pt,minimum width=9mm,minimum height=5mm,inner sep=2pt},
  every path/.style={thick}]
  \node[eq] (SW) {SW-1};
  \node[eq,right=9mm of SW] (R) {R1};
  \node[eq,right=7mm of R] (SRV) {SRV-C};
  \node[eq,above left=5mm and -3mm of SW] (A) {PC-A};
  \node[eq,above right=5mm and -3mm of SW] (B) {PC-B};
  \node[eq,below left=5mm and -3mm of SW] (C) {PC-C};
  \node[eq,below right=5mm and -3mm of SW] (D) {PC-D};
  \foreach \n in {A,B,C,D} \draw (\n) -- (SW);
  \draw (SW) -- node[above,font=\tiny]{G0/0} (R);
  \draw (R) -- node[above,font=\tiny]{G0/1} (SRV);
\end{tikzpicture}
\end{minipage}\hfill
\begin{minipage}{0.61\linewidth}\centering
{\small\renewcommand{\arraystretch}{1.05}\setlength{\tabcolsep}{4pt}
\begin{tabular}{lllll}
  \toprule
  \textbf{Équipement} & \textbf{Adresse IP} & \textbf{Masque} & \textbf{Passerelle} & \textbf{Port} \\
  \midrule
  \Crows
  \bottomrule
\end{tabular}}
\end{minipage}
\end{center}

\question{C1}{4 pts}{Relevez les \textbf{quatre anomalies}. Pour chacune : conséquence et correction.}
{\small\renewcommand{\arraystretch}{1.1}
\ifcorrige\def\hc{}\else\def\hc{\rule[-10mm]{0pt}{14mm}}\fi
\begin{tabularx}{\linewidth}{|c|>{\raggedright\arraybackslash}p{2cm}|L|L|L|}
  \hline
  \textbf{N°} & \textbf{Équipement} & \textbf{Anomalie} & \textbf{Conséquence} & \textbf{Correction} \\ \hline
  \ifcorrige\Csol\else
  \hc 1 & & & & \\ \hline
  \hc 2 & & & & \\ \hline
  \hc 3 & & & & \\ \hline
  \hc 4 & & & & \\ \hline
  \fi
\end{tabularx}}
\bareme{1 pt par anomalie (0,25 équipement, anomalie, conséquence, correction ; ordre indifférent). \Cbareme}

%==========================================================================
\section{Bonus — Extension sans fil \hfill +1}

\begin{cdc}Ajoutez au réseau \resUn{} un \textbf{AccessPoint-PT} relié à SW-\pUn{} (SSID \ip{\ssid}, \textbf{WPA2-PSK})
et un portable \ip{Laptop-\pUn} équipé d'un module \texttt{WPC300N}, en adressage \textbf{statique}.
Vérifiez qu'il joint \srvNom.\end{cdc}
{\small\renewcommand{\arraystretch}{1.1}
\begin{tabularx}{\linewidth}{|l|C|C|C|c|}
  \hline
  \textbf{Équipement} & \textbf{Adresse IP} & \textbf{Masque} & \textbf{Passerelle} & \textbf{Visa prof.} \\ \hline
  \hl Laptop-\pUn & \sol{\netUn.50} & \sol{255.255.255.0} & \sol{\netUn.254} & \hspace{1.4cm} \\ \hline
\end{tabularx}}
\bareme{Le point d'accès fait office de pont : le portable est dans \netUn.0/24. Bonus si associé en WPA2-PSK et ping
vers \netSrv.\srvH{} réussi. Note plafonnée à 20.}

\end{document}
"
%   pdflatex -jobname=Eval_Reseau_SujetB   "\def\variante{B}\def\modecorrige{0}% Évaluation pratique — Réseaux locaux (BTS CIEL 1re année) — sujets A et B
% Compilation (4 PDF) :
%   pdflatex -jobname=Eval_Reseau_SujetA   "\def\variante{A}\def\modecorrige{0}\input{Eval_TP_Reseau_LAN.tex}"
%   pdflatex -jobname=Eval_Reseau_SujetB   "\def\variante{B}\def\modecorrige{0}\input{Eval_TP_Reseau_LAN.tex}"
%   pdflatex -jobname=Eval_Reseau_CorrigeA "\def\variante{A}\def\modecorrige{1}\input{Eval_TP_Reseau_LAN.tex}"
%   pdflatex -jobname=Eval_Reseau_CorrigeB "\def\variante{B}\def\modecorrige{1}\input{Eval_TP_Reseau_LAN.tex}"
\documentclass[10pt,a4paper]{article}

\usepackage[T1]{fontenc}
\usepackage[utf8]{inputenc}
\usepackage{lmodern}
\usepackage[french]{babel}
\usepackage[a4paper,hmargin=1.5cm,top=1.3cm,bottom=1.6cm,footskip=0.7cm]{geometry}
\usepackage[table]{xcolor}
\usepackage{booktabs,tabularx,array}
\usepackage{enumitem}
\usepackage{fancyhdr}
\usepackage[most]{tcolorbox}
\usepackage{tikz}
\usetikzlibrary{positioning}
\usepackage{titlesec}

\providecommand{\modecorrige}{0}
\providecommand{\variante}{A}
\newif\ifcorrige
\ifnum\modecorrige=1 \corrigetrue \else \corrigefalse \fi

\definecolor{ipred}{HTML}{C0392B}
\definecolor{cdcbar}{HTML}{A67C3D}
\definecolor{cdcbg}{HTML}{F4F1EA}
\definecolor{qbar}{HTML}{4A6B5D}
\definecolor{qbg}{HTML}{EDF1EF}
\definecolor{corr}{HTML}{1F5FA8}
\definecolor{corrbg}{HTML}{E8F0FA}
\definecolor{grisfoot}{HTML}{828282}

\newcommand{\ip}[1]{\textcolor{ipred}{\texttt{#1}}}
\newcommand{\pts}[1]{\hfill\mbox{\textbf{(#1)}}}
\newcommand{\sol}[1]{\ifcorrige\textcolor{corr}{#1}\fi}
\newcommand{\bareme}[1]{\ifcorrige\par{\footnotesize\color{corr}#1\par}\fi}
\newcommand{\hl}{\ifcorrige\rule[-1.6mm]{0pt}{5.4mm}\else\rule[-2.2mm]{0pt}{6.6mm}\fi}
\newcolumntype{C}{>{\centering\arraybackslash}X}
\newcolumntype{L}{>{\raggedright\arraybackslash}X}

\titleformat{\section}{\large\bfseries}{}{0pt}{}
\titlespacing*{\section}{0pt}{6pt}{3pt}

\tcbset{qstyle/.style={enhanced,boxrule=0pt,leftrule=2.5pt,arc=0pt,outer arc=0pt,
  left=5pt,right=5pt,top=2pt,bottom=2pt,before skip=4pt,after skip=3pt}}
\newtcolorbox{cdc}{qstyle,colback=cdcbg,colframe=cdcbar}
\newtcolorbox{quest}{qstyle,colback=qbg,colframe=qbar}

% Zone de réponse : vierge dans le sujet, remplie dans le corrigé
\newcommand{\zone}[2]{%
  \ifcorrige
    \begin{tcolorbox}[colback=corrbg,colframe=corr,boxrule=0.5pt,arc=1pt,left=4pt,right=4pt,top=1pt,bottom=1pt,
      before skip=2pt,after skip=2pt,fontupper=\footnotesize\color{corr}]#2\end{tcolorbox}%
  \else
    \begin{tcolorbox}[colback=white,colframe=black!45,boxrule=0.4pt,arc=1pt,height=#1,
      before skip=2pt,after skip=2pt]\end{tcolorbox}%
  \fi}

\newcommand{\question}[3]{\begin{quest}\textbf{#1.} #3\pts{#2}\end{quest}}

%==========================================================================
% Données propres à chaque variante
%==========================================================================
\if\variante A
  \def\entreprise{Le cabinet d'architectes \textbf{ArchiNova}}
  \def\Rnom{R-ArchiNova}
  \def\resUn{Bureau d'études}\def\pUn{BE}\def\netUn{192.168.40}
  \def\resDeux{Accueil}\def\pDeux{ACC}\def\netDeux{172.16}
  \def\srvNom{SRV-Intranet}\def\netSrv{192.168.200}\def\srvH{10}
  \def\ssid{ArchiNova-WiFi}
  \def\AIrows{%
    \hl\ip{192.168.40.25} & \sol{C} & \sol{255.255.255.0} & \sol{192.168.40.0} & \sol{192.168.40.255} & \sol{254} \\ \hline
    \hl\ip{172.16.5.10}   & \sol{B} & \sol{255.255.0.0}   & \sol{172.16.0.0}   & \sol{172.16.255.255} & \sol{65\,534} \\ \hline
    \hl\ip{10.3.2.1}      & \sol{A} & \sol{255.0.0.0}     & \sol{10.0.0.0}     & \sol{10.255.255.255} & \sol{16\,777\,214} \\ \hline}
  \def\apipa{169.254.12.7}
  \def\Cnet{192.168.60}
  \def\Crows{%
    PC-A & \ip{192.168.60.10} & \ip{255.255.255.0} & \ip{192.168.60.254} & — \\
    PC-B & \ip{192.168.60.20} & \ip{255.255.255.0} & \ip{192.168.60.1}   & — \\
    PC-C & \ip{192.168.61.30} & \ip{255.255.255.0} & \ip{192.168.60.254} & — \\
    PC-D & \ip{192.168.60.10} & \ip{255.255.255.0} & \ip{192.168.60.254} & — \\
    SRV-C & \ip{10.0.0.5}     & \ip{255.0.0.0}     & \ip{10.0.0.254}     & — \\
    R1 — G0/0 & \ip{192.168.60.254} & \ip{255.255.255.0} & — & On \\
    R1 — G0/1 & \ip{10.0.0.254}     & \ip{255.0.0.0}     & — & Off \\}
  \def\Csol{%
    \hl 1 & \sol{PC-B} & \sol{Passerelle .60.1 : pas l'adresse du routeur.} & \sol{Joint son réseau, mais aucun autre (serveur).} & \sol{Passerelle 192.168.60.254.} \\ \hline
    \hl 2 & \sol{PC-C} & \sol{IP hors du réseau 192.168.60.0/24.} & \sol{Ne joint ni les postes ni le routeur.} & \sol{Ex. 192.168.60.30.} \\ \hline
    \hl 3 & \sol{PC-D} & \sol{Même IP que PC-A (doublon).} & \sol{Conflit d'adresse : PC-A et PC-D perturbés.} & \sol{Adresse unique, ex. .60.40.} \\ \hline
    \hl 4 & \sol{R1 G0/1} & \sol{Interface désactivée (Off).} & \sol{Réseau du serveur injoignable pour tous.} & \sol{Port Status : On.} \\ \hline}
  \def\Cbareme{PC-A et SRV-C sont correctement configurés ; accepter PC-A cité à la place de PC-D pour le doublon.}
\else
  \def\entreprise{La clinique vétérinaire \textbf{VetoPlus}}
  \def\Rnom{R-VetoPlus}
  \def\resUn{Soins}\def\pUn{SOI}\def\netUn{192.168.70}
  \def\resDeux{Secrétariat}\def\pDeux{SEC}\def\netDeux{172.20}
  \def\srvNom{SRV-Dossiers}\def\netSrv{192.168.250}\def\srvH{20}
  \def\ssid{VetoPlus-WiFi}
  \def\AIrows{%
    \hl\ip{192.168.15.130} & \sol{C} & \sol{255.255.255.0} & \sol{192.168.15.0} & \sol{192.168.15.255} & \sol{254} \\ \hline
    \hl\ip{172.31.200.7}   & \sol{B} & \sol{255.255.0.0}   & \sol{172.31.0.0}   & \sol{172.31.255.255} & \sol{65\,534} \\ \hline
    \hl\ip{10.45.6.9}      & \sol{A} & \sol{255.0.0.0}     & \sol{10.0.0.0}     & \sol{10.255.255.255} & \sol{16\,777\,214} \\ \hline}
  \def\apipa{169.254.88.3}
  \def\Cnet{192.168.80}
  \def\Crows{%
    PC-A & \ip{192.168.8.11}   & \ip{255.255.255.0} & \ip{192.168.80.254} & — \\
    PC-B & \ip{192.168.80.12}  & \ip{255.255.255.0} & \ip{192.168.80.254} & — \\
    PC-C & \ip{192.168.80.13}  & \ip{255.255.255.0} & \ip{192.168.80.100} & — \\
    PC-D & \ip{192.168.80.12}  & \ip{255.255.255.0} & \ip{192.168.80.254} & — \\
    SRV-C & \ip{10.0.0.8}      & \ip{255.0.0.0}     & \ip{10.0.0.254}     & — \\
    R1 — G0/0 & \ip{192.168.80.254} & \ip{255.255.255.0} & — & Off \\
    R1 — G0/1 & \ip{10.0.0.254}     & \ip{255.0.0.0}     & — & On \\}
  \def\Csol{%
    \hl 1 & \sol{PC-A} & \sol{IP hors du réseau 192.168.80.0/24.} & \sol{Ne joint ni les postes ni le routeur.} & \sol{Ex. 192.168.80.11.} \\ \hline
    \hl 2 & \sol{PC-C} & \sol{Passerelle .80.100 : pas l'adresse du routeur.} & \sol{Joint son réseau, mais aucun autre (serveur).} & \sol{Passerelle 192.168.80.254.} \\ \hline
    \hl 3 & \sol{PC-D} & \sol{Même IP que PC-B (doublon).} & \sol{Conflit d'adresse : PC-B et PC-D perturbés.} & \sol{Adresse unique, ex. .80.14.} \\ \hline
    \hl 4 & \sol{R1 G0/0} & \sol{Interface côté LAN désactivée (Off).} & \sol{Aucun poste ne joint le routeur ni le serveur.} & \sol{Port Status : On.} \\ \hline}
  \def\Cbareme{PC-B et SRV-C sont correctement configurés ; accepter PC-B cité à la place de PC-D pour le doublon.}
\fi

\pagestyle{fancy}
\fancyhf{}
\renewcommand{\headrulewidth}{0pt}
\fancyfoot[L]{\footnotesize\itshape\color{grisfoot}Évaluation — Réseaux locaux — Sujet \variante\ifcorrige\ — CORRIGÉ\fi}
\fancyfoot[R]{\footnotesize\color{grisfoot}\thepage/3}

\setlist{itemsep=0pt,topsep=2pt,parsep=0pt}
\setlength{\parindent}{0pt}
\setlength{\parskip}{2pt}

\begin{document}

\begin{center}
  {\Large\bfseries Évaluation pratique — Réseaux locaux \quad
   \fcolorbox{black}{white}{\ SUJET \variante\ }\par}
  \vspace{1pt}
  {\bfseries Adressage IP, commutation et routage — BTS CIEL 1re année — durée : 2 heures\par}
  \ifcorrige\vspace{2pt}{\bfseries\color{corr}CORRIGÉ ET BARÈME — document enseignant}\par\fi
\end{center}
\vspace{-2pt}
{\renewcommand{\arraystretch}{1.35}
\begin{tabularx}{\linewidth}{|X|X|p{2.2cm}|p{2.6cm}|}
  \hline
  \textbf{Nom :} & \textbf{Prénom :} & \textbf{Poste n°} & \textbf{Note :} \hfill /20 \\ \hline
\end{tabularx}}

\textbf{Consignes :} travail \textbf{individuel}, documents et Internet \textbf{non autorisés}. Ce sujet est le
\textbf{document-réponse}. Fichier Packet Tracer : \ip{Nom\_Prenom\_EvalLAN\_\variante.pkt}, à enregistrer
régulièrement dans le dossier de rendu. Les « visas » sont apposés par le professeur, que vous appelez à chaque test réussi.
Durées conseillées : A 30 min — B 1 h 10 — C 20 min.

%==========================================================================
\section{Partie A — Questions de cours et calculs \hfill /6}

\question{A1}{2 pts}{Complétez le tableau en considérant le \textbf{masque par défaut} de la classe de chaque adresse.}
{\small\renewcommand{\arraystretch}{1.15}
\begin{tabularx}{\linewidth}{|l|c|C|C|C|c|}
  \hline
  \textbf{Adresse IP} & \textbf{Classe} & \textbf{Masque par défaut} & \textbf{Adresse réseau} &
  \textbf{Adresse de diffusion} & \textbf{Nb d'hôtes} \\ \hline
  \AIrows
\end{tabularx}}
\bareme{0,5 pt par ligne juste + 0,5 pt si les trois nombres d'hôtes sont justes ($2^{n}-2$, $n$ = nombre de bits d'hôte).}

\question{A2}{1 pt}{Indiquez le type de câble (\textbf{droit} ou \textbf{croisé}) pour chaque liaison.}
{\small\renewcommand{\arraystretch}{1.15}
\begin{tabularx}{\linewidth}{|X|c||X|c|}
  \hline
  \hl PC $\leftrightarrow$ commutateur & \makebox[1.6cm]{\sol{droit}} & Commutateur $\leftrightarrow$ routeur & \makebox[1.6cm]{\sol{droit}} \\ \hline
  \hl PC $\leftrightarrow$ PC (réseau P2P) & \sol{croisé} & Concentrateur $\leftrightarrow$ concentrateur & \sol{croisé} \\ \hline
\end{tabularx}}
\bareme{0,25 pt par case. Équipements de types différents : câble droit ; de même type : câble croisé.}

\question{A3}{1 pt}{Comment un \textbf{concentrateur (hub)} et un \textbf{commutateur (switch)} traitent-ils une trame
reçue ? Lequel limite les collisions, et pourquoi ?}
\zone{2.9cm}{Le \textbf{hub} (couche 1) répète la trame sur \textbf{tous ses autres ports} : un seul domaine de collision
partagé. Le \textbf{switch} (couche 2) lit l'adresse MAC de destination et, grâce à sa \textbf{table MAC}, ne l'envoie
\textbf{que sur le port du destinataire} : un domaine de collision par port, c'est lui qui limite les collisions.
\emph{0,5 pt hub, 0,5 pt switch.}}

\question{A4}{1 pt}{Qu'est-ce que la \textbf{passerelle par défaut} d'un poste ? Quelle condition son adresse doit-elle respecter ?}
\zone{2.4cm}{Adresse IP de l'interface du \textbf{routeur} à laquelle le poste envoie les paquets destinés à un
\textbf{autre réseau}. Elle doit \textbf{appartenir au même réseau} que le poste. \emph{0,5 + 0,5 pt.}}

\question{A5}{0,5 pt}{Un portable configuré en DHCP affiche \ip{\apipa} / \ip{255.255.0.0}. Que s'est-il passé ?}
\zone{1.5cm}{Aucun \textbf{serveur DHCP} n'a répondu : le poste s'est attribué une adresse \textbf{APIPA} (169.254.0.0/16).}

\question{A6}{0,5 pt}{Quel protocole utilise \texttt{ping} ? Nommez les deux messages échangés et la couche OSI des adresses IP.}
\zone{1.5cm}{\textbf{ICMP} ; \textbf{Echo Request} / \textbf{Echo Reply} ; adresses IP en \textbf{couche 3 (Réseau)}.}

%==========================================================================
\newpage
\section{Partie B — Réalisation sous Packet Tracer \hfill /10}

\begin{cdc}
\entreprise{} vous confie son réseau : \textbf{trois réseaux} reliés par un \textbf{routeur unique Cisco 2911}
nommé \ip{\Rnom} (interfaces G0/0, G0/1, G0/2).
\begin{itemize}[leftmargin=12pt]
  \item \textbf{Réseau \resUn} : \textbf{4 postes} PC-\pUn1 à PC-\pUn4, commutateur \ip{SW-\pUn}, adressage statique,
        réseau \ip{\netUn.0/24}.
  \item \textbf{Réseau \resDeux} : \textbf{2 postes} PC-\pDeux1, PC-\pDeux2 et \textbf{1 imprimante} IMP-\pDeux,
        commutateur \ip{SW-\pDeux}, adressage statique, réseau de classe B \ip{\netDeux.0.0/16}.
  \item \textbf{Réseau Serveur} : serveur \ip{\srvNom} d'adresse \ip{\netSrv.\srvH/24}, relié \textbf{directement} au routeur.
  \item La passerelle de chaque réseau est sa \textbf{dernière adresse utilisable}. Tous les équipements communiquent
        entre eux et accèdent à la page web du serveur. Nommage imposé (« Display Name »).
\end{itemize}
\end{cdc}

\question{B1}{2 pts}{Complétez le plan d'adressage \textbf{avant} la réalisation.}
{\footnotesize\setlength{\tabcolsep}{3pt}\renewcommand{\arraystretch}{1.15}
\begin{tabularx}{\linewidth}{|l|C|C|C|C|C|C|}
  \hline
  \textbf{Réseau} & \textbf{Adresse réseau} & \textbf{Masque} & \textbf{1re @ utilisable} &
  \textbf{Dernière @ utilisable} & \textbf{Diffusion} & \textbf{Passerelle} \\ \hline
  \hl \resUn & \sol{\netUn.0} & \sol{255.255.255.0} & \sol{\netUn.1} & \sol{\netUn.254} & \sol{\netUn.255} & \sol{\netUn.254} \\ \hline
  \hl \resDeux & \sol{\netDeux.0.0} & \sol{255.255.0.0} & \sol{\netDeux.0.1} & \sol{\netDeux.255.254} & \sol{\netDeux.255.255} & \sol{\netDeux.255.254} \\ \hline
  \hl Serveur & \sol{\netSrv.0} & \sol{255.255.255.0} & \sol{\netSrv.1} & \sol{\netSrv.254} & \sol{\netSrv.255} & \sol{\netSrv.254} \\ \hline
\end{tabularx}}
\bareme{0,5 pt par ligne juste + 0,5 pt si la convention de passerelle est respectée partout.}

\question{B2}{2 pts}{Construisez la topologie (équipements, câblage, nommage) et relevez l'adressage configuré.}
{\footnotesize\setlength{\tabcolsep}{3pt}\renewcommand{\arraystretch}{1.1}
\begin{tabularx}{\linewidth}{|l|C|C|C||l|C|C|}
  \hline
  \textbf{Hôte} & \textbf{Adresse IP} & \textbf{Masque} & \textbf{Passerelle} & \textbf{Interface} & \textbf{Adresse IP} & \textbf{Masque} \\ \hline
  \hl PC-\pUn1 & \sol{\netUn.1} & \sol{255.255.255.0} & \sol{\netUn.254} & \Rnom{} G0/0 & \sol{\netUn.254} & \sol{255.255.255.0} \\ \hline
  \hl PC-\pDeux1 & \sol{\netDeux.0.1} & \sol{255.255.0.0} & \sol{\netDeux.255.254} & \Rnom{} G0/1 & \sol{\netDeux.255.254} & \sol{255.255.0.0} \\ \hline
  \hl IMP-\pDeux & \sol{\netDeux.0.10} & \sol{255.255.0.0} & \sol{\netDeux.255.254} & \Rnom{} G0/2 & \sol{\netSrv.254} & \sol{255.255.255.0} \\ \hline
  \hl \srvNom & \ip{\netSrv.\srvH} & \sol{255.255.255.0} & \sol{\netSrv.254} & \multicolumn{3}{c|}{\cellcolor{black!8}} \\ \hline
\end{tabularx}}
\bareme{Adresses d'hôtes libres (uniques, dans le bon réseau) ; affectation G0/x libre. Sur le .pkt : équipements 0,5 ;
câblage 1 (câbles droits vers les switchs et switch–routeur, \textbf{croisé} serveur–routeur, liens verts) ; nommage 0,5.}

\question{B3}{3 pts}{Configurez tous les équipements : IP, masques, passerelles des hôtes ; interfaces du routeur
(« Port Status : On »).}
\bareme{Sur le .pkt : interfaces routeur 1 pt ; IP/masques des hôtes 1 pt ; passerelles (imprimante et serveur compris)
1 pt. $-0,25$ par erreur, sans descendre sous 0 par item.}

\question{B4}{3 pts}{Réalisez les tests, notez la commande et le résultat, puis faites viser chaque test réussi.}
{\small\renewcommand{\arraystretch}{1.1}
\begin{tabularx}{\linewidth}{|c|l|l|C|c|c|}
  \hline
  \textbf{N°} & \textbf{Source} & \textbf{Destination} & \textbf{Commande / action} & \textbf{Résultat} & \textbf{Visa prof.} \\ \hline
  \hl 1 & PC-\pUn1 & PC-\pUn4 & \sol{ping \netUn.4} & \sol{réussi} & \hspace{1.4cm} \\ \hline
  \hl 2 & PC-\pDeux1 & IMP-\pDeux & \sol{ping \netDeux.0.10} & \sol{réussi} & \\ \hline
  \hl 3 & PC-\pUn2 & PC-\pDeux2 & \sol{ping \netDeux.0.2} & \sol{réussi} & \\ \hline
  \hl 4 & PC-\pDeux2 & \srvNom & \sol{ping \netSrv.\srvH} & \sol{réussi} & \\ \hline
  \hl 5 & PC-\pUn3 & page web du serveur & \sol{navigateur : {\NoAutoSpacing http://\netSrv.\srvH}} & \sol{page affichée} & \\ \hline
\end{tabularx}}
\bareme{0,5 pt par test visé + 0,5 pt pour B5. Le premier ping inter-réseaux peut échouer (résolution ARP) ; les suivants doivent réussir.}

\question{B5}{incl. B4}{Test n°3 : par quels équipements passe le paquet ? Pourquoi PC-\pUn2 l'envoie-t-il au routeur ?}
\zone{2.3cm}{PC-\pUn2 $\to$ SW-\pUn{} $\to$ \Rnom{} $\to$ SW-\pDeux{} $\to$ PC-\pDeux2. Grâce à son masque, PC-\pUn2 constate que
la destination (\netDeux.0.0/16) n'est \textbf{pas dans son réseau} (\netUn.0/24) : il envoie le paquet à sa
\textbf{passerelle par défaut}, qui le route.}

%==========================================================================
\newpage
\section{Partie C — Diagnostic d'un réseau en panne \hfill /4}

Un technicien constate qu'\textbf{aucun poste ne parvient à joindre le serveur SRV-C} et que certains postes
communiquent mal entre eux. Les câbles sont du bon type.

\begin{center}
\begin{minipage}{0.37\linewidth}\centering
\begin{tikzpicture}[font=\footnotesize,
  eq/.style={draw,fill=qbg,rounded corners=2pt,minimum width=9mm,minimum height=5mm,inner sep=2pt},
  every path/.style={thick}]
  \node[eq] (SW) {SW-1};
  \node[eq,right=9mm of SW] (R) {R1};
  \node[eq,right=7mm of R] (SRV) {SRV-C};
  \node[eq,above left=5mm and -3mm of SW] (A) {PC-A};
  \node[eq,above right=5mm and -3mm of SW] (B) {PC-B};
  \node[eq,below left=5mm and -3mm of SW] (C) {PC-C};
  \node[eq,below right=5mm and -3mm of SW] (D) {PC-D};
  \foreach \n in {A,B,C,D} \draw (\n) -- (SW);
  \draw (SW) -- node[above,font=\tiny]{G0/0} (R);
  \draw (R) -- node[above,font=\tiny]{G0/1} (SRV);
\end{tikzpicture}
\end{minipage}\hfill
\begin{minipage}{0.61\linewidth}\centering
{\small\renewcommand{\arraystretch}{1.05}\setlength{\tabcolsep}{4pt}
\begin{tabular}{lllll}
  \toprule
  \textbf{Équipement} & \textbf{Adresse IP} & \textbf{Masque} & \textbf{Passerelle} & \textbf{Port} \\
  \midrule
  \Crows
  \bottomrule
\end{tabular}}
\end{minipage}
\end{center}

\question{C1}{4 pts}{Relevez les \textbf{quatre anomalies}. Pour chacune : conséquence et correction.}
{\small\renewcommand{\arraystretch}{1.1}
\ifcorrige\def\hc{}\else\def\hc{\rule[-10mm]{0pt}{14mm}}\fi
\begin{tabularx}{\linewidth}{|c|>{\raggedright\arraybackslash}p{2cm}|L|L|L|}
  \hline
  \textbf{N°} & \textbf{Équipement} & \textbf{Anomalie} & \textbf{Conséquence} & \textbf{Correction} \\ \hline
  \ifcorrige\Csol\else
  \hc 1 & & & & \\ \hline
  \hc 2 & & & & \\ \hline
  \hc 3 & & & & \\ \hline
  \hc 4 & & & & \\ \hline
  \fi
\end{tabularx}}
\bareme{1 pt par anomalie (0,25 équipement, anomalie, conséquence, correction ; ordre indifférent). \Cbareme}

%==========================================================================
\section{Bonus — Extension sans fil \hfill +1}

\begin{cdc}Ajoutez au réseau \resUn{} un \textbf{AccessPoint-PT} relié à SW-\pUn{} (SSID \ip{\ssid}, \textbf{WPA2-PSK})
et un portable \ip{Laptop-\pUn} équipé d'un module \texttt{WPC300N}, en adressage \textbf{statique}.
Vérifiez qu'il joint \srvNom.\end{cdc}
{\small\renewcommand{\arraystretch}{1.1}
\begin{tabularx}{\linewidth}{|l|C|C|C|c|}
  \hline
  \textbf{Équipement} & \textbf{Adresse IP} & \textbf{Masque} & \textbf{Passerelle} & \textbf{Visa prof.} \\ \hline
  \hl Laptop-\pUn & \sol{\netUn.50} & \sol{255.255.255.0} & \sol{\netUn.254} & \hspace{1.4cm} \\ \hline
\end{tabularx}}
\bareme{Le point d'accès fait office de pont : le portable est dans \netUn.0/24. Bonus si associé en WPA2-PSK et ping
vers \netSrv.\srvH{} réussi. Note plafonnée à 20.}

\end{document}
"
%   pdflatex -jobname=Eval_Reseau_CorrigeA "\def\variante{A}\def\modecorrige{1}% Évaluation pratique — Réseaux locaux (BTS CIEL 1re année) — sujets A et B
% Compilation (4 PDF) :
%   pdflatex -jobname=Eval_Reseau_SujetA   "\def\variante{A}\def\modecorrige{0}\input{Eval_TP_Reseau_LAN.tex}"
%   pdflatex -jobname=Eval_Reseau_SujetB   "\def\variante{B}\def\modecorrige{0}\input{Eval_TP_Reseau_LAN.tex}"
%   pdflatex -jobname=Eval_Reseau_CorrigeA "\def\variante{A}\def\modecorrige{1}\input{Eval_TP_Reseau_LAN.tex}"
%   pdflatex -jobname=Eval_Reseau_CorrigeB "\def\variante{B}\def\modecorrige{1}\input{Eval_TP_Reseau_LAN.tex}"
\documentclass[10pt,a4paper]{article}

\usepackage[T1]{fontenc}
\usepackage[utf8]{inputenc}
\usepackage{lmodern}
\usepackage[french]{babel}
\usepackage[a4paper,hmargin=1.5cm,top=1.3cm,bottom=1.6cm,footskip=0.7cm]{geometry}
\usepackage[table]{xcolor}
\usepackage{booktabs,tabularx,array}
\usepackage{enumitem}
\usepackage{fancyhdr}
\usepackage[most]{tcolorbox}
\usepackage{tikz}
\usetikzlibrary{positioning}
\usepackage{titlesec}

\providecommand{\modecorrige}{0}
\providecommand{\variante}{A}
\newif\ifcorrige
\ifnum\modecorrige=1 \corrigetrue \else \corrigefalse \fi

\definecolor{ipred}{HTML}{C0392B}
\definecolor{cdcbar}{HTML}{A67C3D}
\definecolor{cdcbg}{HTML}{F4F1EA}
\definecolor{qbar}{HTML}{4A6B5D}
\definecolor{qbg}{HTML}{EDF1EF}
\definecolor{corr}{HTML}{1F5FA8}
\definecolor{corrbg}{HTML}{E8F0FA}
\definecolor{grisfoot}{HTML}{828282}

\newcommand{\ip}[1]{\textcolor{ipred}{\texttt{#1}}}
\newcommand{\pts}[1]{\hfill\mbox{\textbf{(#1)}}}
\newcommand{\sol}[1]{\ifcorrige\textcolor{corr}{#1}\fi}
\newcommand{\bareme}[1]{\ifcorrige\par{\footnotesize\color{corr}#1\par}\fi}
\newcommand{\hl}{\ifcorrige\rule[-1.6mm]{0pt}{5.4mm}\else\rule[-2.2mm]{0pt}{6.6mm}\fi}
\newcolumntype{C}{>{\centering\arraybackslash}X}
\newcolumntype{L}{>{\raggedright\arraybackslash}X}

\titleformat{\section}{\large\bfseries}{}{0pt}{}
\titlespacing*{\section}{0pt}{6pt}{3pt}

\tcbset{qstyle/.style={enhanced,boxrule=0pt,leftrule=2.5pt,arc=0pt,outer arc=0pt,
  left=5pt,right=5pt,top=2pt,bottom=2pt,before skip=4pt,after skip=3pt}}
\newtcolorbox{cdc}{qstyle,colback=cdcbg,colframe=cdcbar}
\newtcolorbox{quest}{qstyle,colback=qbg,colframe=qbar}

% Zone de réponse : vierge dans le sujet, remplie dans le corrigé
\newcommand{\zone}[2]{%
  \ifcorrige
    \begin{tcolorbox}[colback=corrbg,colframe=corr,boxrule=0.5pt,arc=1pt,left=4pt,right=4pt,top=1pt,bottom=1pt,
      before skip=2pt,after skip=2pt,fontupper=\footnotesize\color{corr}]#2\end{tcolorbox}%
  \else
    \begin{tcolorbox}[colback=white,colframe=black!45,boxrule=0.4pt,arc=1pt,height=#1,
      before skip=2pt,after skip=2pt]\end{tcolorbox}%
  \fi}

\newcommand{\question}[3]{\begin{quest}\textbf{#1.} #3\pts{#2}\end{quest}}

%==========================================================================
% Données propres à chaque variante
%==========================================================================
\if\variante A
  \def\entreprise{Le cabinet d'architectes \textbf{ArchiNova}}
  \def\Rnom{R-ArchiNova}
  \def\resUn{Bureau d'études}\def\pUn{BE}\def\netUn{192.168.40}
  \def\resDeux{Accueil}\def\pDeux{ACC}\def\netDeux{172.16}
  \def\srvNom{SRV-Intranet}\def\netSrv{192.168.200}\def\srvH{10}
  \def\ssid{ArchiNova-WiFi}
  \def\AIrows{%
    \hl\ip{192.168.40.25} & \sol{C} & \sol{255.255.255.0} & \sol{192.168.40.0} & \sol{192.168.40.255} & \sol{254} \\ \hline
    \hl\ip{172.16.5.10}   & \sol{B} & \sol{255.255.0.0}   & \sol{172.16.0.0}   & \sol{172.16.255.255} & \sol{65\,534} \\ \hline
    \hl\ip{10.3.2.1}      & \sol{A} & \sol{255.0.0.0}     & \sol{10.0.0.0}     & \sol{10.255.255.255} & \sol{16\,777\,214} \\ \hline}
  \def\apipa{169.254.12.7}
  \def\Cnet{192.168.60}
  \def\Crows{%
    PC-A & \ip{192.168.60.10} & \ip{255.255.255.0} & \ip{192.168.60.254} & — \\
    PC-B & \ip{192.168.60.20} & \ip{255.255.255.0} & \ip{192.168.60.1}   & — \\
    PC-C & \ip{192.168.61.30} & \ip{255.255.255.0} & \ip{192.168.60.254} & — \\
    PC-D & \ip{192.168.60.10} & \ip{255.255.255.0} & \ip{192.168.60.254} & — \\
    SRV-C & \ip{10.0.0.5}     & \ip{255.0.0.0}     & \ip{10.0.0.254}     & — \\
    R1 — G0/0 & \ip{192.168.60.254} & \ip{255.255.255.0} & — & On \\
    R1 — G0/1 & \ip{10.0.0.254}     & \ip{255.0.0.0}     & — & Off \\}
  \def\Csol{%
    \hl 1 & \sol{PC-B} & \sol{Passerelle .60.1 : pas l'adresse du routeur.} & \sol{Joint son réseau, mais aucun autre (serveur).} & \sol{Passerelle 192.168.60.254.} \\ \hline
    \hl 2 & \sol{PC-C} & \sol{IP hors du réseau 192.168.60.0/24.} & \sol{Ne joint ni les postes ni le routeur.} & \sol{Ex. 192.168.60.30.} \\ \hline
    \hl 3 & \sol{PC-D} & \sol{Même IP que PC-A (doublon).} & \sol{Conflit d'adresse : PC-A et PC-D perturbés.} & \sol{Adresse unique, ex. .60.40.} \\ \hline
    \hl 4 & \sol{R1 G0/1} & \sol{Interface désactivée (Off).} & \sol{Réseau du serveur injoignable pour tous.} & \sol{Port Status : On.} \\ \hline}
  \def\Cbareme{PC-A et SRV-C sont correctement configurés ; accepter PC-A cité à la place de PC-D pour le doublon.}
\else
  \def\entreprise{La clinique vétérinaire \textbf{VetoPlus}}
  \def\Rnom{R-VetoPlus}
  \def\resUn{Soins}\def\pUn{SOI}\def\netUn{192.168.70}
  \def\resDeux{Secrétariat}\def\pDeux{SEC}\def\netDeux{172.20}
  \def\srvNom{SRV-Dossiers}\def\netSrv{192.168.250}\def\srvH{20}
  \def\ssid{VetoPlus-WiFi}
  \def\AIrows{%
    \hl\ip{192.168.15.130} & \sol{C} & \sol{255.255.255.0} & \sol{192.168.15.0} & \sol{192.168.15.255} & \sol{254} \\ \hline
    \hl\ip{172.31.200.7}   & \sol{B} & \sol{255.255.0.0}   & \sol{172.31.0.0}   & \sol{172.31.255.255} & \sol{65\,534} \\ \hline
    \hl\ip{10.45.6.9}      & \sol{A} & \sol{255.0.0.0}     & \sol{10.0.0.0}     & \sol{10.255.255.255} & \sol{16\,777\,214} \\ \hline}
  \def\apipa{169.254.88.3}
  \def\Cnet{192.168.80}
  \def\Crows{%
    PC-A & \ip{192.168.8.11}   & \ip{255.255.255.0} & \ip{192.168.80.254} & — \\
    PC-B & \ip{192.168.80.12}  & \ip{255.255.255.0} & \ip{192.168.80.254} & — \\
    PC-C & \ip{192.168.80.13}  & \ip{255.255.255.0} & \ip{192.168.80.100} & — \\
    PC-D & \ip{192.168.80.12}  & \ip{255.255.255.0} & \ip{192.168.80.254} & — \\
    SRV-C & \ip{10.0.0.8}      & \ip{255.0.0.0}     & \ip{10.0.0.254}     & — \\
    R1 — G0/0 & \ip{192.168.80.254} & \ip{255.255.255.0} & — & Off \\
    R1 — G0/1 & \ip{10.0.0.254}     & \ip{255.0.0.0}     & — & On \\}
  \def\Csol{%
    \hl 1 & \sol{PC-A} & \sol{IP hors du réseau 192.168.80.0/24.} & \sol{Ne joint ni les postes ni le routeur.} & \sol{Ex. 192.168.80.11.} \\ \hline
    \hl 2 & \sol{PC-C} & \sol{Passerelle .80.100 : pas l'adresse du routeur.} & \sol{Joint son réseau, mais aucun autre (serveur).} & \sol{Passerelle 192.168.80.254.} \\ \hline
    \hl 3 & \sol{PC-D} & \sol{Même IP que PC-B (doublon).} & \sol{Conflit d'adresse : PC-B et PC-D perturbés.} & \sol{Adresse unique, ex. .80.14.} \\ \hline
    \hl 4 & \sol{R1 G0/0} & \sol{Interface côté LAN désactivée (Off).} & \sol{Aucun poste ne joint le routeur ni le serveur.} & \sol{Port Status : On.} \\ \hline}
  \def\Cbareme{PC-B et SRV-C sont correctement configurés ; accepter PC-B cité à la place de PC-D pour le doublon.}
\fi

\pagestyle{fancy}
\fancyhf{}
\renewcommand{\headrulewidth}{0pt}
\fancyfoot[L]{\footnotesize\itshape\color{grisfoot}Évaluation — Réseaux locaux — Sujet \variante\ifcorrige\ — CORRIGÉ\fi}
\fancyfoot[R]{\footnotesize\color{grisfoot}\thepage/3}

\setlist{itemsep=0pt,topsep=2pt,parsep=0pt}
\setlength{\parindent}{0pt}
\setlength{\parskip}{2pt}

\begin{document}

\begin{center}
  {\Large\bfseries Évaluation pratique — Réseaux locaux \quad
   \fcolorbox{black}{white}{\ SUJET \variante\ }\par}
  \vspace{1pt}
  {\bfseries Adressage IP, commutation et routage — BTS CIEL 1re année — durée : 2 heures\par}
  \ifcorrige\vspace{2pt}{\bfseries\color{corr}CORRIGÉ ET BARÈME — document enseignant}\par\fi
\end{center}
\vspace{-2pt}
{\renewcommand{\arraystretch}{1.35}
\begin{tabularx}{\linewidth}{|X|X|p{2.2cm}|p{2.6cm}|}
  \hline
  \textbf{Nom :} & \textbf{Prénom :} & \textbf{Poste n°} & \textbf{Note :} \hfill /20 \\ \hline
\end{tabularx}}

\textbf{Consignes :} travail \textbf{individuel}, documents et Internet \textbf{non autorisés}. Ce sujet est le
\textbf{document-réponse}. Fichier Packet Tracer : \ip{Nom\_Prenom\_EvalLAN\_\variante.pkt}, à enregistrer
régulièrement dans le dossier de rendu. Les « visas » sont apposés par le professeur, que vous appelez à chaque test réussi.
Durées conseillées : A 30 min — B 1 h 10 — C 20 min.

%==========================================================================
\section{Partie A — Questions de cours et calculs \hfill /6}

\question{A1}{2 pts}{Complétez le tableau en considérant le \textbf{masque par défaut} de la classe de chaque adresse.}
{\small\renewcommand{\arraystretch}{1.15}
\begin{tabularx}{\linewidth}{|l|c|C|C|C|c|}
  \hline
  \textbf{Adresse IP} & \textbf{Classe} & \textbf{Masque par défaut} & \textbf{Adresse réseau} &
  \textbf{Adresse de diffusion} & \textbf{Nb d'hôtes} \\ \hline
  \AIrows
\end{tabularx}}
\bareme{0,5 pt par ligne juste + 0,5 pt si les trois nombres d'hôtes sont justes ($2^{n}-2$, $n$ = nombre de bits d'hôte).}

\question{A2}{1 pt}{Indiquez le type de câble (\textbf{droit} ou \textbf{croisé}) pour chaque liaison.}
{\small\renewcommand{\arraystretch}{1.15}
\begin{tabularx}{\linewidth}{|X|c||X|c|}
  \hline
  \hl PC $\leftrightarrow$ commutateur & \makebox[1.6cm]{\sol{droit}} & Commutateur $\leftrightarrow$ routeur & \makebox[1.6cm]{\sol{droit}} \\ \hline
  \hl PC $\leftrightarrow$ PC (réseau P2P) & \sol{croisé} & Concentrateur $\leftrightarrow$ concentrateur & \sol{croisé} \\ \hline
\end{tabularx}}
\bareme{0,25 pt par case. Équipements de types différents : câble droit ; de même type : câble croisé.}

\question{A3}{1 pt}{Comment un \textbf{concentrateur (hub)} et un \textbf{commutateur (switch)} traitent-ils une trame
reçue ? Lequel limite les collisions, et pourquoi ?}
\zone{2.9cm}{Le \textbf{hub} (couche 1) répète la trame sur \textbf{tous ses autres ports} : un seul domaine de collision
partagé. Le \textbf{switch} (couche 2) lit l'adresse MAC de destination et, grâce à sa \textbf{table MAC}, ne l'envoie
\textbf{que sur le port du destinataire} : un domaine de collision par port, c'est lui qui limite les collisions.
\emph{0,5 pt hub, 0,5 pt switch.}}

\question{A4}{1 pt}{Qu'est-ce que la \textbf{passerelle par défaut} d'un poste ? Quelle condition son adresse doit-elle respecter ?}
\zone{2.4cm}{Adresse IP de l'interface du \textbf{routeur} à laquelle le poste envoie les paquets destinés à un
\textbf{autre réseau}. Elle doit \textbf{appartenir au même réseau} que le poste. \emph{0,5 + 0,5 pt.}}

\question{A5}{0,5 pt}{Un portable configuré en DHCP affiche \ip{\apipa} / \ip{255.255.0.0}. Que s'est-il passé ?}
\zone{1.5cm}{Aucun \textbf{serveur DHCP} n'a répondu : le poste s'est attribué une adresse \textbf{APIPA} (169.254.0.0/16).}

\question{A6}{0,5 pt}{Quel protocole utilise \texttt{ping} ? Nommez les deux messages échangés et la couche OSI des adresses IP.}
\zone{1.5cm}{\textbf{ICMP} ; \textbf{Echo Request} / \textbf{Echo Reply} ; adresses IP en \textbf{couche 3 (Réseau)}.}

%==========================================================================
\newpage
\section{Partie B — Réalisation sous Packet Tracer \hfill /10}

\begin{cdc}
\entreprise{} vous confie son réseau : \textbf{trois réseaux} reliés par un \textbf{routeur unique Cisco 2911}
nommé \ip{\Rnom} (interfaces G0/0, G0/1, G0/2).
\begin{itemize}[leftmargin=12pt]
  \item \textbf{Réseau \resUn} : \textbf{4 postes} PC-\pUn1 à PC-\pUn4, commutateur \ip{SW-\pUn}, adressage statique,
        réseau \ip{\netUn.0/24}.
  \item \textbf{Réseau \resDeux} : \textbf{2 postes} PC-\pDeux1, PC-\pDeux2 et \textbf{1 imprimante} IMP-\pDeux,
        commutateur \ip{SW-\pDeux}, adressage statique, réseau de classe B \ip{\netDeux.0.0/16}.
  \item \textbf{Réseau Serveur} : serveur \ip{\srvNom} d'adresse \ip{\netSrv.\srvH/24}, relié \textbf{directement} au routeur.
  \item La passerelle de chaque réseau est sa \textbf{dernière adresse utilisable}. Tous les équipements communiquent
        entre eux et accèdent à la page web du serveur. Nommage imposé (« Display Name »).
\end{itemize}
\end{cdc}

\question{B1}{2 pts}{Complétez le plan d'adressage \textbf{avant} la réalisation.}
{\footnotesize\setlength{\tabcolsep}{3pt}\renewcommand{\arraystretch}{1.15}
\begin{tabularx}{\linewidth}{|l|C|C|C|C|C|C|}
  \hline
  \textbf{Réseau} & \textbf{Adresse réseau} & \textbf{Masque} & \textbf{1re @ utilisable} &
  \textbf{Dernière @ utilisable} & \textbf{Diffusion} & \textbf{Passerelle} \\ \hline
  \hl \resUn & \sol{\netUn.0} & \sol{255.255.255.0} & \sol{\netUn.1} & \sol{\netUn.254} & \sol{\netUn.255} & \sol{\netUn.254} \\ \hline
  \hl \resDeux & \sol{\netDeux.0.0} & \sol{255.255.0.0} & \sol{\netDeux.0.1} & \sol{\netDeux.255.254} & \sol{\netDeux.255.255} & \sol{\netDeux.255.254} \\ \hline
  \hl Serveur & \sol{\netSrv.0} & \sol{255.255.255.0} & \sol{\netSrv.1} & \sol{\netSrv.254} & \sol{\netSrv.255} & \sol{\netSrv.254} \\ \hline
\end{tabularx}}
\bareme{0,5 pt par ligne juste + 0,5 pt si la convention de passerelle est respectée partout.}

\question{B2}{2 pts}{Construisez la topologie (équipements, câblage, nommage) et relevez l'adressage configuré.}
{\footnotesize\setlength{\tabcolsep}{3pt}\renewcommand{\arraystretch}{1.1}
\begin{tabularx}{\linewidth}{|l|C|C|C||l|C|C|}
  \hline
  \textbf{Hôte} & \textbf{Adresse IP} & \textbf{Masque} & \textbf{Passerelle} & \textbf{Interface} & \textbf{Adresse IP} & \textbf{Masque} \\ \hline
  \hl PC-\pUn1 & \sol{\netUn.1} & \sol{255.255.255.0} & \sol{\netUn.254} & \Rnom{} G0/0 & \sol{\netUn.254} & \sol{255.255.255.0} \\ \hline
  \hl PC-\pDeux1 & \sol{\netDeux.0.1} & \sol{255.255.0.0} & \sol{\netDeux.255.254} & \Rnom{} G0/1 & \sol{\netDeux.255.254} & \sol{255.255.0.0} \\ \hline
  \hl IMP-\pDeux & \sol{\netDeux.0.10} & \sol{255.255.0.0} & \sol{\netDeux.255.254} & \Rnom{} G0/2 & \sol{\netSrv.254} & \sol{255.255.255.0} \\ \hline
  \hl \srvNom & \ip{\netSrv.\srvH} & \sol{255.255.255.0} & \sol{\netSrv.254} & \multicolumn{3}{c|}{\cellcolor{black!8}} \\ \hline
\end{tabularx}}
\bareme{Adresses d'hôtes libres (uniques, dans le bon réseau) ; affectation G0/x libre. Sur le .pkt : équipements 0,5 ;
câblage 1 (câbles droits vers les switchs et switch–routeur, \textbf{croisé} serveur–routeur, liens verts) ; nommage 0,5.}

\question{B3}{3 pts}{Configurez tous les équipements : IP, masques, passerelles des hôtes ; interfaces du routeur
(« Port Status : On »).}
\bareme{Sur le .pkt : interfaces routeur 1 pt ; IP/masques des hôtes 1 pt ; passerelles (imprimante et serveur compris)
1 pt. $-0,25$ par erreur, sans descendre sous 0 par item.}

\question{B4}{3 pts}{Réalisez les tests, notez la commande et le résultat, puis faites viser chaque test réussi.}
{\small\renewcommand{\arraystretch}{1.1}
\begin{tabularx}{\linewidth}{|c|l|l|C|c|c|}
  \hline
  \textbf{N°} & \textbf{Source} & \textbf{Destination} & \textbf{Commande / action} & \textbf{Résultat} & \textbf{Visa prof.} \\ \hline
  \hl 1 & PC-\pUn1 & PC-\pUn4 & \sol{ping \netUn.4} & \sol{réussi} & \hspace{1.4cm} \\ \hline
  \hl 2 & PC-\pDeux1 & IMP-\pDeux & \sol{ping \netDeux.0.10} & \sol{réussi} & \\ \hline
  \hl 3 & PC-\pUn2 & PC-\pDeux2 & \sol{ping \netDeux.0.2} & \sol{réussi} & \\ \hline
  \hl 4 & PC-\pDeux2 & \srvNom & \sol{ping \netSrv.\srvH} & \sol{réussi} & \\ \hline
  \hl 5 & PC-\pUn3 & page web du serveur & \sol{navigateur : {\NoAutoSpacing http://\netSrv.\srvH}} & \sol{page affichée} & \\ \hline
\end{tabularx}}
\bareme{0,5 pt par test visé + 0,5 pt pour B5. Le premier ping inter-réseaux peut échouer (résolution ARP) ; les suivants doivent réussir.}

\question{B5}{incl. B4}{Test n°3 : par quels équipements passe le paquet ? Pourquoi PC-\pUn2 l'envoie-t-il au routeur ?}
\zone{2.3cm}{PC-\pUn2 $\to$ SW-\pUn{} $\to$ \Rnom{} $\to$ SW-\pDeux{} $\to$ PC-\pDeux2. Grâce à son masque, PC-\pUn2 constate que
la destination (\netDeux.0.0/16) n'est \textbf{pas dans son réseau} (\netUn.0/24) : il envoie le paquet à sa
\textbf{passerelle par défaut}, qui le route.}

%==========================================================================
\newpage
\section{Partie C — Diagnostic d'un réseau en panne \hfill /4}

Un technicien constate qu'\textbf{aucun poste ne parvient à joindre le serveur SRV-C} et que certains postes
communiquent mal entre eux. Les câbles sont du bon type.

\begin{center}
\begin{minipage}{0.37\linewidth}\centering
\begin{tikzpicture}[font=\footnotesize,
  eq/.style={draw,fill=qbg,rounded corners=2pt,minimum width=9mm,minimum height=5mm,inner sep=2pt},
  every path/.style={thick}]
  \node[eq] (SW) {SW-1};
  \node[eq,right=9mm of SW] (R) {R1};
  \node[eq,right=7mm of R] (SRV) {SRV-C};
  \node[eq,above left=5mm and -3mm of SW] (A) {PC-A};
  \node[eq,above right=5mm and -3mm of SW] (B) {PC-B};
  \node[eq,below left=5mm and -3mm of SW] (C) {PC-C};
  \node[eq,below right=5mm and -3mm of SW] (D) {PC-D};
  \foreach \n in {A,B,C,D} \draw (\n) -- (SW);
  \draw (SW) -- node[above,font=\tiny]{G0/0} (R);
  \draw (R) -- node[above,font=\tiny]{G0/1} (SRV);
\end{tikzpicture}
\end{minipage}\hfill
\begin{minipage}{0.61\linewidth}\centering
{\small\renewcommand{\arraystretch}{1.05}\setlength{\tabcolsep}{4pt}
\begin{tabular}{lllll}
  \toprule
  \textbf{Équipement} & \textbf{Adresse IP} & \textbf{Masque} & \textbf{Passerelle} & \textbf{Port} \\
  \midrule
  \Crows
  \bottomrule
\end{tabular}}
\end{minipage}
\end{center}

\question{C1}{4 pts}{Relevez les \textbf{quatre anomalies}. Pour chacune : conséquence et correction.}
{\small\renewcommand{\arraystretch}{1.1}
\ifcorrige\def\hc{}\else\def\hc{\rule[-10mm]{0pt}{14mm}}\fi
\begin{tabularx}{\linewidth}{|c|>{\raggedright\arraybackslash}p{2cm}|L|L|L|}
  \hline
  \textbf{N°} & \textbf{Équipement} & \textbf{Anomalie} & \textbf{Conséquence} & \textbf{Correction} \\ \hline
  \ifcorrige\Csol\else
  \hc 1 & & & & \\ \hline
  \hc 2 & & & & \\ \hline
  \hc 3 & & & & \\ \hline
  \hc 4 & & & & \\ \hline
  \fi
\end{tabularx}}
\bareme{1 pt par anomalie (0,25 équipement, anomalie, conséquence, correction ; ordre indifférent). \Cbareme}

%==========================================================================
\section{Bonus — Extension sans fil \hfill +1}

\begin{cdc}Ajoutez au réseau \resUn{} un \textbf{AccessPoint-PT} relié à SW-\pUn{} (SSID \ip{\ssid}, \textbf{WPA2-PSK})
et un portable \ip{Laptop-\pUn} équipé d'un module \texttt{WPC300N}, en adressage \textbf{statique}.
Vérifiez qu'il joint \srvNom.\end{cdc}
{\small\renewcommand{\arraystretch}{1.1}
\begin{tabularx}{\linewidth}{|l|C|C|C|c|}
  \hline
  \textbf{Équipement} & \textbf{Adresse IP} & \textbf{Masque} & \textbf{Passerelle} & \textbf{Visa prof.} \\ \hline
  \hl Laptop-\pUn & \sol{\netUn.50} & \sol{255.255.255.0} & \sol{\netUn.254} & \hspace{1.4cm} \\ \hline
\end{tabularx}}
\bareme{Le point d'accès fait office de pont : le portable est dans \netUn.0/24. Bonus si associé en WPA2-PSK et ping
vers \netSrv.\srvH{} réussi. Note plafonnée à 20.}

\end{document}
"
%   pdflatex -jobname=Eval_Reseau_CorrigeB "\def\variante{B}\def\modecorrige{1}% Évaluation pratique — Réseaux locaux (BTS CIEL 1re année) — sujets A et B
% Compilation (4 PDF) :
%   pdflatex -jobname=Eval_Reseau_SujetA   "\def\variante{A}\def\modecorrige{0}\input{Eval_TP_Reseau_LAN.tex}"
%   pdflatex -jobname=Eval_Reseau_SujetB   "\def\variante{B}\def\modecorrige{0}\input{Eval_TP_Reseau_LAN.tex}"
%   pdflatex -jobname=Eval_Reseau_CorrigeA "\def\variante{A}\def\modecorrige{1}\input{Eval_TP_Reseau_LAN.tex}"
%   pdflatex -jobname=Eval_Reseau_CorrigeB "\def\variante{B}\def\modecorrige{1}\input{Eval_TP_Reseau_LAN.tex}"
\documentclass[10pt,a4paper]{article}

\usepackage[T1]{fontenc}
\usepackage[utf8]{inputenc}
\usepackage{lmodern}
\usepackage[french]{babel}
\usepackage[a4paper,hmargin=1.5cm,top=1.3cm,bottom=1.6cm,footskip=0.7cm]{geometry}
\usepackage[table]{xcolor}
\usepackage{booktabs,tabularx,array}
\usepackage{enumitem}
\usepackage{fancyhdr}
\usepackage[most]{tcolorbox}
\usepackage{tikz}
\usetikzlibrary{positioning}
\usepackage{titlesec}

\providecommand{\modecorrige}{0}
\providecommand{\variante}{A}
\newif\ifcorrige
\ifnum\modecorrige=1 \corrigetrue \else \corrigefalse \fi

\definecolor{ipred}{HTML}{C0392B}
\definecolor{cdcbar}{HTML}{A67C3D}
\definecolor{cdcbg}{HTML}{F4F1EA}
\definecolor{qbar}{HTML}{4A6B5D}
\definecolor{qbg}{HTML}{EDF1EF}
\definecolor{corr}{HTML}{1F5FA8}
\definecolor{corrbg}{HTML}{E8F0FA}
\definecolor{grisfoot}{HTML}{828282}

\newcommand{\ip}[1]{\textcolor{ipred}{\texttt{#1}}}
\newcommand{\pts}[1]{\hfill\mbox{\textbf{(#1)}}}
\newcommand{\sol}[1]{\ifcorrige\textcolor{corr}{#1}\fi}
\newcommand{\bareme}[1]{\ifcorrige\par{\footnotesize\color{corr}#1\par}\fi}
\newcommand{\hl}{\ifcorrige\rule[-1.6mm]{0pt}{5.4mm}\else\rule[-2.2mm]{0pt}{6.6mm}\fi}
\newcolumntype{C}{>{\centering\arraybackslash}X}
\newcolumntype{L}{>{\raggedright\arraybackslash}X}

\titleformat{\section}{\large\bfseries}{}{0pt}{}
\titlespacing*{\section}{0pt}{6pt}{3pt}

\tcbset{qstyle/.style={enhanced,boxrule=0pt,leftrule=2.5pt,arc=0pt,outer arc=0pt,
  left=5pt,right=5pt,top=2pt,bottom=2pt,before skip=4pt,after skip=3pt}}
\newtcolorbox{cdc}{qstyle,colback=cdcbg,colframe=cdcbar}
\newtcolorbox{quest}{qstyle,colback=qbg,colframe=qbar}

% Zone de réponse : vierge dans le sujet, remplie dans le corrigé
\newcommand{\zone}[2]{%
  \ifcorrige
    \begin{tcolorbox}[colback=corrbg,colframe=corr,boxrule=0.5pt,arc=1pt,left=4pt,right=4pt,top=1pt,bottom=1pt,
      before skip=2pt,after skip=2pt,fontupper=\footnotesize\color{corr}]#2\end{tcolorbox}%
  \else
    \begin{tcolorbox}[colback=white,colframe=black!45,boxrule=0.4pt,arc=1pt,height=#1,
      before skip=2pt,after skip=2pt]\end{tcolorbox}%
  \fi}

\newcommand{\question}[3]{\begin{quest}\textbf{#1.} #3\pts{#2}\end{quest}}

%==========================================================================
% Données propres à chaque variante
%==========================================================================
\if\variante A
  \def\entreprise{Le cabinet d'architectes \textbf{ArchiNova}}
  \def\Rnom{R-ArchiNova}
  \def\resUn{Bureau d'études}\def\pUn{BE}\def\netUn{192.168.40}
  \def\resDeux{Accueil}\def\pDeux{ACC}\def\netDeux{172.16}
  \def\srvNom{SRV-Intranet}\def\netSrv{192.168.200}\def\srvH{10}
  \def\ssid{ArchiNova-WiFi}
  \def\AIrows{%
    \hl\ip{192.168.40.25} & \sol{C} & \sol{255.255.255.0} & \sol{192.168.40.0} & \sol{192.168.40.255} & \sol{254} \\ \hline
    \hl\ip{172.16.5.10}   & \sol{B} & \sol{255.255.0.0}   & \sol{172.16.0.0}   & \sol{172.16.255.255} & \sol{65\,534} \\ \hline
    \hl\ip{10.3.2.1}      & \sol{A} & \sol{255.0.0.0}     & \sol{10.0.0.0}     & \sol{10.255.255.255} & \sol{16\,777\,214} \\ \hline}
  \def\apipa{169.254.12.7}
  \def\Cnet{192.168.60}
  \def\Crows{%
    PC-A & \ip{192.168.60.10} & \ip{255.255.255.0} & \ip{192.168.60.254} & — \\
    PC-B & \ip{192.168.60.20} & \ip{255.255.255.0} & \ip{192.168.60.1}   & — \\
    PC-C & \ip{192.168.61.30} & \ip{255.255.255.0} & \ip{192.168.60.254} & — \\
    PC-D & \ip{192.168.60.10} & \ip{255.255.255.0} & \ip{192.168.60.254} & — \\
    SRV-C & \ip{10.0.0.5}     & \ip{255.0.0.0}     & \ip{10.0.0.254}     & — \\
    R1 — G0/0 & \ip{192.168.60.254} & \ip{255.255.255.0} & — & On \\
    R1 — G0/1 & \ip{10.0.0.254}     & \ip{255.0.0.0}     & — & Off \\}
  \def\Csol{%
    \hl 1 & \sol{PC-B} & \sol{Passerelle .60.1 : pas l'adresse du routeur.} & \sol{Joint son réseau, mais aucun autre (serveur).} & \sol{Passerelle 192.168.60.254.} \\ \hline
    \hl 2 & \sol{PC-C} & \sol{IP hors du réseau 192.168.60.0/24.} & \sol{Ne joint ni les postes ni le routeur.} & \sol{Ex. 192.168.60.30.} \\ \hline
    \hl 3 & \sol{PC-D} & \sol{Même IP que PC-A (doublon).} & \sol{Conflit d'adresse : PC-A et PC-D perturbés.} & \sol{Adresse unique, ex. .60.40.} \\ \hline
    \hl 4 & \sol{R1 G0/1} & \sol{Interface désactivée (Off).} & \sol{Réseau du serveur injoignable pour tous.} & \sol{Port Status : On.} \\ \hline}
  \def\Cbareme{PC-A et SRV-C sont correctement configurés ; accepter PC-A cité à la place de PC-D pour le doublon.}
\else
  \def\entreprise{La clinique vétérinaire \textbf{VetoPlus}}
  \def\Rnom{R-VetoPlus}
  \def\resUn{Soins}\def\pUn{SOI}\def\netUn{192.168.70}
  \def\resDeux{Secrétariat}\def\pDeux{SEC}\def\netDeux{172.20}
  \def\srvNom{SRV-Dossiers}\def\netSrv{192.168.250}\def\srvH{20}
  \def\ssid{VetoPlus-WiFi}
  \def\AIrows{%
    \hl\ip{192.168.15.130} & \sol{C} & \sol{255.255.255.0} & \sol{192.168.15.0} & \sol{192.168.15.255} & \sol{254} \\ \hline
    \hl\ip{172.31.200.7}   & \sol{B} & \sol{255.255.0.0}   & \sol{172.31.0.0}   & \sol{172.31.255.255} & \sol{65\,534} \\ \hline
    \hl\ip{10.45.6.9}      & \sol{A} & \sol{255.0.0.0}     & \sol{10.0.0.0}     & \sol{10.255.255.255} & \sol{16\,777\,214} \\ \hline}
  \def\apipa{169.254.88.3}
  \def\Cnet{192.168.80}
  \def\Crows{%
    PC-A & \ip{192.168.8.11}   & \ip{255.255.255.0} & \ip{192.168.80.254} & — \\
    PC-B & \ip{192.168.80.12}  & \ip{255.255.255.0} & \ip{192.168.80.254} & — \\
    PC-C & \ip{192.168.80.13}  & \ip{255.255.255.0} & \ip{192.168.80.100} & — \\
    PC-D & \ip{192.168.80.12}  & \ip{255.255.255.0} & \ip{192.168.80.254} & — \\
    SRV-C & \ip{10.0.0.8}      & \ip{255.0.0.0}     & \ip{10.0.0.254}     & — \\
    R1 — G0/0 & \ip{192.168.80.254} & \ip{255.255.255.0} & — & Off \\
    R1 — G0/1 & \ip{10.0.0.254}     & \ip{255.0.0.0}     & — & On \\}
  \def\Csol{%
    \hl 1 & \sol{PC-A} & \sol{IP hors du réseau 192.168.80.0/24.} & \sol{Ne joint ni les postes ni le routeur.} & \sol{Ex. 192.168.80.11.} \\ \hline
    \hl 2 & \sol{PC-C} & \sol{Passerelle .80.100 : pas l'adresse du routeur.} & \sol{Joint son réseau, mais aucun autre (serveur).} & \sol{Passerelle 192.168.80.254.} \\ \hline
    \hl 3 & \sol{PC-D} & \sol{Même IP que PC-B (doublon).} & \sol{Conflit d'adresse : PC-B et PC-D perturbés.} & \sol{Adresse unique, ex. .80.14.} \\ \hline
    \hl 4 & \sol{R1 G0/0} & \sol{Interface côté LAN désactivée (Off).} & \sol{Aucun poste ne joint le routeur ni le serveur.} & \sol{Port Status : On.} \\ \hline}
  \def\Cbareme{PC-B et SRV-C sont correctement configurés ; accepter PC-B cité à la place de PC-D pour le doublon.}
\fi

\pagestyle{fancy}
\fancyhf{}
\renewcommand{\headrulewidth}{0pt}
\fancyfoot[L]{\footnotesize\itshape\color{grisfoot}Évaluation — Réseaux locaux — Sujet \variante\ifcorrige\ — CORRIGÉ\fi}
\fancyfoot[R]{\footnotesize\color{grisfoot}\thepage/3}

\setlist{itemsep=0pt,topsep=2pt,parsep=0pt}
\setlength{\parindent}{0pt}
\setlength{\parskip}{2pt}

\begin{document}

\begin{center}
  {\Large\bfseries Évaluation pratique — Réseaux locaux \quad
   \fcolorbox{black}{white}{\ SUJET \variante\ }\par}
  \vspace{1pt}
  {\bfseries Adressage IP, commutation et routage — BTS CIEL 1re année — durée : 2 heures\par}
  \ifcorrige\vspace{2pt}{\bfseries\color{corr}CORRIGÉ ET BARÈME — document enseignant}\par\fi
\end{center}
\vspace{-2pt}
{\renewcommand{\arraystretch}{1.35}
\begin{tabularx}{\linewidth}{|X|X|p{2.2cm}|p{2.6cm}|}
  \hline
  \textbf{Nom :} & \textbf{Prénom :} & \textbf{Poste n°} & \textbf{Note :} \hfill /20 \\ \hline
\end{tabularx}}

\textbf{Consignes :} travail \textbf{individuel}, documents et Internet \textbf{non autorisés}. Ce sujet est le
\textbf{document-réponse}. Fichier Packet Tracer : \ip{Nom\_Prenom\_EvalLAN\_\variante.pkt}, à enregistrer
régulièrement dans le dossier de rendu. Les « visas » sont apposés par le professeur, que vous appelez à chaque test réussi.
Durées conseillées : A 30 min — B 1 h 10 — C 20 min.

%==========================================================================
\section{Partie A — Questions de cours et calculs \hfill /6}

\question{A1}{2 pts}{Complétez le tableau en considérant le \textbf{masque par défaut} de la classe de chaque adresse.}
{\small\renewcommand{\arraystretch}{1.15}
\begin{tabularx}{\linewidth}{|l|c|C|C|C|c|}
  \hline
  \textbf{Adresse IP} & \textbf{Classe} & \textbf{Masque par défaut} & \textbf{Adresse réseau} &
  \textbf{Adresse de diffusion} & \textbf{Nb d'hôtes} \\ \hline
  \AIrows
\end{tabularx}}
\bareme{0,5 pt par ligne juste + 0,5 pt si les trois nombres d'hôtes sont justes ($2^{n}-2$, $n$ = nombre de bits d'hôte).}

\question{A2}{1 pt}{Indiquez le type de câble (\textbf{droit} ou \textbf{croisé}) pour chaque liaison.}
{\small\renewcommand{\arraystretch}{1.15}
\begin{tabularx}{\linewidth}{|X|c||X|c|}
  \hline
  \hl PC $\leftrightarrow$ commutateur & \makebox[1.6cm]{\sol{droit}} & Commutateur $\leftrightarrow$ routeur & \makebox[1.6cm]{\sol{droit}} \\ \hline
  \hl PC $\leftrightarrow$ PC (réseau P2P) & \sol{croisé} & Concentrateur $\leftrightarrow$ concentrateur & \sol{croisé} \\ \hline
\end{tabularx}}
\bareme{0,25 pt par case. Équipements de types différents : câble droit ; de même type : câble croisé.}

\question{A3}{1 pt}{Comment un \textbf{concentrateur (hub)} et un \textbf{commutateur (switch)} traitent-ils une trame
reçue ? Lequel limite les collisions, et pourquoi ?}
\zone{2.9cm}{Le \textbf{hub} (couche 1) répète la trame sur \textbf{tous ses autres ports} : un seul domaine de collision
partagé. Le \textbf{switch} (couche 2) lit l'adresse MAC de destination et, grâce à sa \textbf{table MAC}, ne l'envoie
\textbf{que sur le port du destinataire} : un domaine de collision par port, c'est lui qui limite les collisions.
\emph{0,5 pt hub, 0,5 pt switch.}}

\question{A4}{1 pt}{Qu'est-ce que la \textbf{passerelle par défaut} d'un poste ? Quelle condition son adresse doit-elle respecter ?}
\zone{2.4cm}{Adresse IP de l'interface du \textbf{routeur} à laquelle le poste envoie les paquets destinés à un
\textbf{autre réseau}. Elle doit \textbf{appartenir au même réseau} que le poste. \emph{0,5 + 0,5 pt.}}

\question{A5}{0,5 pt}{Un portable configuré en DHCP affiche \ip{\apipa} / \ip{255.255.0.0}. Que s'est-il passé ?}
\zone{1.5cm}{Aucun \textbf{serveur DHCP} n'a répondu : le poste s'est attribué une adresse \textbf{APIPA} (169.254.0.0/16).}

\question{A6}{0,5 pt}{Quel protocole utilise \texttt{ping} ? Nommez les deux messages échangés et la couche OSI des adresses IP.}
\zone{1.5cm}{\textbf{ICMP} ; \textbf{Echo Request} / \textbf{Echo Reply} ; adresses IP en \textbf{couche 3 (Réseau)}.}

%==========================================================================
\newpage
\section{Partie B — Réalisation sous Packet Tracer \hfill /10}

\begin{cdc}
\entreprise{} vous confie son réseau : \textbf{trois réseaux} reliés par un \textbf{routeur unique Cisco 2911}
nommé \ip{\Rnom} (interfaces G0/0, G0/1, G0/2).
\begin{itemize}[leftmargin=12pt]
  \item \textbf{Réseau \resUn} : \textbf{4 postes} PC-\pUn1 à PC-\pUn4, commutateur \ip{SW-\pUn}, adressage statique,
        réseau \ip{\netUn.0/24}.
  \item \textbf{Réseau \resDeux} : \textbf{2 postes} PC-\pDeux1, PC-\pDeux2 et \textbf{1 imprimante} IMP-\pDeux,
        commutateur \ip{SW-\pDeux}, adressage statique, réseau de classe B \ip{\netDeux.0.0/16}.
  \item \textbf{Réseau Serveur} : serveur \ip{\srvNom} d'adresse \ip{\netSrv.\srvH/24}, relié \textbf{directement} au routeur.
  \item La passerelle de chaque réseau est sa \textbf{dernière adresse utilisable}. Tous les équipements communiquent
        entre eux et accèdent à la page web du serveur. Nommage imposé (« Display Name »).
\end{itemize}
\end{cdc}

\question{B1}{2 pts}{Complétez le plan d'adressage \textbf{avant} la réalisation.}
{\footnotesize\setlength{\tabcolsep}{3pt}\renewcommand{\arraystretch}{1.15}
\begin{tabularx}{\linewidth}{|l|C|C|C|C|C|C|}
  \hline
  \textbf{Réseau} & \textbf{Adresse réseau} & \textbf{Masque} & \textbf{1re @ utilisable} &
  \textbf{Dernière @ utilisable} & \textbf{Diffusion} & \textbf{Passerelle} \\ \hline
  \hl \resUn & \sol{\netUn.0} & \sol{255.255.255.0} & \sol{\netUn.1} & \sol{\netUn.254} & \sol{\netUn.255} & \sol{\netUn.254} \\ \hline
  \hl \resDeux & \sol{\netDeux.0.0} & \sol{255.255.0.0} & \sol{\netDeux.0.1} & \sol{\netDeux.255.254} & \sol{\netDeux.255.255} & \sol{\netDeux.255.254} \\ \hline
  \hl Serveur & \sol{\netSrv.0} & \sol{255.255.255.0} & \sol{\netSrv.1} & \sol{\netSrv.254} & \sol{\netSrv.255} & \sol{\netSrv.254} \\ \hline
\end{tabularx}}
\bareme{0,5 pt par ligne juste + 0,5 pt si la convention de passerelle est respectée partout.}

\question{B2}{2 pts}{Construisez la topologie (équipements, câblage, nommage) et relevez l'adressage configuré.}
{\footnotesize\setlength{\tabcolsep}{3pt}\renewcommand{\arraystretch}{1.1}
\begin{tabularx}{\linewidth}{|l|C|C|C||l|C|C|}
  \hline
  \textbf{Hôte} & \textbf{Adresse IP} & \textbf{Masque} & \textbf{Passerelle} & \textbf{Interface} & \textbf{Adresse IP} & \textbf{Masque} \\ \hline
  \hl PC-\pUn1 & \sol{\netUn.1} & \sol{255.255.255.0} & \sol{\netUn.254} & \Rnom{} G0/0 & \sol{\netUn.254} & \sol{255.255.255.0} \\ \hline
  \hl PC-\pDeux1 & \sol{\netDeux.0.1} & \sol{255.255.0.0} & \sol{\netDeux.255.254} & \Rnom{} G0/1 & \sol{\netDeux.255.254} & \sol{255.255.0.0} \\ \hline
  \hl IMP-\pDeux & \sol{\netDeux.0.10} & \sol{255.255.0.0} & \sol{\netDeux.255.254} & \Rnom{} G0/2 & \sol{\netSrv.254} & \sol{255.255.255.0} \\ \hline
  \hl \srvNom & \ip{\netSrv.\srvH} & \sol{255.255.255.0} & \sol{\netSrv.254} & \multicolumn{3}{c|}{\cellcolor{black!8}} \\ \hline
\end{tabularx}}
\bareme{Adresses d'hôtes libres (uniques, dans le bon réseau) ; affectation G0/x libre. Sur le .pkt : équipements 0,5 ;
câblage 1 (câbles droits vers les switchs et switch–routeur, \textbf{croisé} serveur–routeur, liens verts) ; nommage 0,5.}

\question{B3}{3 pts}{Configurez tous les équipements : IP, masques, passerelles des hôtes ; interfaces du routeur
(« Port Status : On »).}
\bareme{Sur le .pkt : interfaces routeur 1 pt ; IP/masques des hôtes 1 pt ; passerelles (imprimante et serveur compris)
1 pt. $-0,25$ par erreur, sans descendre sous 0 par item.}

\question{B4}{3 pts}{Réalisez les tests, notez la commande et le résultat, puis faites viser chaque test réussi.}
{\small\renewcommand{\arraystretch}{1.1}
\begin{tabularx}{\linewidth}{|c|l|l|C|c|c|}
  \hline
  \textbf{N°} & \textbf{Source} & \textbf{Destination} & \textbf{Commande / action} & \textbf{Résultat} & \textbf{Visa prof.} \\ \hline
  \hl 1 & PC-\pUn1 & PC-\pUn4 & \sol{ping \netUn.4} & \sol{réussi} & \hspace{1.4cm} \\ \hline
  \hl 2 & PC-\pDeux1 & IMP-\pDeux & \sol{ping \netDeux.0.10} & \sol{réussi} & \\ \hline
  \hl 3 & PC-\pUn2 & PC-\pDeux2 & \sol{ping \netDeux.0.2} & \sol{réussi} & \\ \hline
  \hl 4 & PC-\pDeux2 & \srvNom & \sol{ping \netSrv.\srvH} & \sol{réussi} & \\ \hline
  \hl 5 & PC-\pUn3 & page web du serveur & \sol{navigateur : {\NoAutoSpacing http://\netSrv.\srvH}} & \sol{page affichée} & \\ \hline
\end{tabularx}}
\bareme{0,5 pt par test visé + 0,5 pt pour B5. Le premier ping inter-réseaux peut échouer (résolution ARP) ; les suivants doivent réussir.}

\question{B5}{incl. B4}{Test n°3 : par quels équipements passe le paquet ? Pourquoi PC-\pUn2 l'envoie-t-il au routeur ?}
\zone{2.3cm}{PC-\pUn2 $\to$ SW-\pUn{} $\to$ \Rnom{} $\to$ SW-\pDeux{} $\to$ PC-\pDeux2. Grâce à son masque, PC-\pUn2 constate que
la destination (\netDeux.0.0/16) n'est \textbf{pas dans son réseau} (\netUn.0/24) : il envoie le paquet à sa
\textbf{passerelle par défaut}, qui le route.}

%==========================================================================
\newpage
\section{Partie C — Diagnostic d'un réseau en panne \hfill /4}

Un technicien constate qu'\textbf{aucun poste ne parvient à joindre le serveur SRV-C} et que certains postes
communiquent mal entre eux. Les câbles sont du bon type.

\begin{center}
\begin{minipage}{0.37\linewidth}\centering
\begin{tikzpicture}[font=\footnotesize,
  eq/.style={draw,fill=qbg,rounded corners=2pt,minimum width=9mm,minimum height=5mm,inner sep=2pt},
  every path/.style={thick}]
  \node[eq] (SW) {SW-1};
  \node[eq,right=9mm of SW] (R) {R1};
  \node[eq,right=7mm of R] (SRV) {SRV-C};
  \node[eq,above left=5mm and -3mm of SW] (A) {PC-A};
  \node[eq,above right=5mm and -3mm of SW] (B) {PC-B};
  \node[eq,below left=5mm and -3mm of SW] (C) {PC-C};
  \node[eq,below right=5mm and -3mm of SW] (D) {PC-D};
  \foreach \n in {A,B,C,D} \draw (\n) -- (SW);
  \draw (SW) -- node[above,font=\tiny]{G0/0} (R);
  \draw (R) -- node[above,font=\tiny]{G0/1} (SRV);
\end{tikzpicture}
\end{minipage}\hfill
\begin{minipage}{0.61\linewidth}\centering
{\small\renewcommand{\arraystretch}{1.05}\setlength{\tabcolsep}{4pt}
\begin{tabular}{lllll}
  \toprule
  \textbf{Équipement} & \textbf{Adresse IP} & \textbf{Masque} & \textbf{Passerelle} & \textbf{Port} \\
  \midrule
  \Crows
  \bottomrule
\end{tabular}}
\end{minipage}
\end{center}

\question{C1}{4 pts}{Relevez les \textbf{quatre anomalies}. Pour chacune : conséquence et correction.}
{\small\renewcommand{\arraystretch}{1.1}
\ifcorrige\def\hc{}\else\def\hc{\rule[-10mm]{0pt}{14mm}}\fi
\begin{tabularx}{\linewidth}{|c|>{\raggedright\arraybackslash}p{2cm}|L|L|L|}
  \hline
  \textbf{N°} & \textbf{Équipement} & \textbf{Anomalie} & \textbf{Conséquence} & \textbf{Correction} \\ \hline
  \ifcorrige\Csol\else
  \hc 1 & & & & \\ \hline
  \hc 2 & & & & \\ \hline
  \hc 3 & & & & \\ \hline
  \hc 4 & & & & \\ \hline
  \fi
\end{tabularx}}
\bareme{1 pt par anomalie (0,25 équipement, anomalie, conséquence, correction ; ordre indifférent). \Cbareme}

%==========================================================================
\section{Bonus — Extension sans fil \hfill +1}

\begin{cdc}Ajoutez au réseau \resUn{} un \textbf{AccessPoint-PT} relié à SW-\pUn{} (SSID \ip{\ssid}, \textbf{WPA2-PSK})
et un portable \ip{Laptop-\pUn} équipé d'un module \texttt{WPC300N}, en adressage \textbf{statique}.
Vérifiez qu'il joint \srvNom.\end{cdc}
{\small\renewcommand{\arraystretch}{1.1}
\begin{tabularx}{\linewidth}{|l|C|C|C|c|}
  \hline
  \textbf{Équipement} & \textbf{Adresse IP} & \textbf{Masque} & \textbf{Passerelle} & \textbf{Visa prof.} \\ \hline
  \hl Laptop-\pUn & \sol{\netUn.50} & \sol{255.255.255.0} & \sol{\netUn.254} & \hspace{1.4cm} \\ \hline
\end{tabularx}}
\bareme{Le point d'accès fait office de pont : le portable est dans \netUn.0/24. Bonus si associé en WPA2-PSK et ping
vers \netSrv.\srvH{} réussi. Note plafonnée à 20.}

\end{document}
"
\documentclass[10pt,a4paper]{article}

\usepackage[T1]{fontenc}
\usepackage[utf8]{inputenc}
\usepackage{lmodern}
\usepackage[french]{babel}
\usepackage[a4paper,hmargin=1.5cm,top=1.3cm,bottom=1.6cm,footskip=0.7cm]{geometry}
\usepackage[table]{xcolor}
\usepackage{booktabs,tabularx,array}
\usepackage{enumitem}
\usepackage{fancyhdr}
\usepackage[most]{tcolorbox}
\usepackage{tikz}
\usetikzlibrary{positioning}
\usepackage{titlesec}

\providecommand{\modecorrige}{0}
\providecommand{\variante}{A}
\newif\ifcorrige
\ifnum\modecorrige=1 \corrigetrue \else \corrigefalse \fi

\definecolor{ipred}{HTML}{C0392B}
\definecolor{cdcbar}{HTML}{A67C3D}
\definecolor{cdcbg}{HTML}{F4F1EA}
\definecolor{qbar}{HTML}{4A6B5D}
\definecolor{qbg}{HTML}{EDF1EF}
\definecolor{corr}{HTML}{1F5FA8}
\definecolor{corrbg}{HTML}{E8F0FA}
\definecolor{grisfoot}{HTML}{828282}

\newcommand{\ip}[1]{\textcolor{ipred}{\texttt{#1}}}
\newcommand{\pts}[1]{\hfill\mbox{\textbf{(#1)}}}
\newcommand{\sol}[1]{\ifcorrige\textcolor{corr}{#1}\fi}
\newcommand{\bareme}[1]{\ifcorrige\par{\footnotesize\color{corr}#1\par}\fi}
\newcommand{\hl}{\ifcorrige\rule[-1.6mm]{0pt}{5.4mm}\else\rule[-2.2mm]{0pt}{6.6mm}\fi}
\newcolumntype{C}{>{\centering\arraybackslash}X}
\newcolumntype{L}{>{\raggedright\arraybackslash}X}

\titleformat{\section}{\large\bfseries}{}{0pt}{}
\titlespacing*{\section}{0pt}{6pt}{3pt}

\tcbset{qstyle/.style={enhanced,boxrule=0pt,leftrule=2.5pt,arc=0pt,outer arc=0pt,
  left=5pt,right=5pt,top=2pt,bottom=2pt,before skip=4pt,after skip=3pt}}
\newtcolorbox{cdc}{qstyle,colback=cdcbg,colframe=cdcbar}
\newtcolorbox{quest}{qstyle,colback=qbg,colframe=qbar}

% Zone de réponse : vierge dans le sujet, remplie dans le corrigé
\newcommand{\zone}[2]{%
  \ifcorrige
    \begin{tcolorbox}[colback=corrbg,colframe=corr,boxrule=0.5pt,arc=1pt,left=4pt,right=4pt,top=1pt,bottom=1pt,
      before skip=2pt,after skip=2pt,fontupper=\footnotesize\color{corr}]#2\end{tcolorbox}%
  \else
    \begin{tcolorbox}[colback=white,colframe=black!45,boxrule=0.4pt,arc=1pt,height=#1,
      before skip=2pt,after skip=2pt]\end{tcolorbox}%
  \fi}

\newcommand{\question}[3]{\begin{quest}\textbf{#1.} #3\pts{#2}\end{quest}}

%==========================================================================
% Données propres à chaque variante
%==========================================================================
\if\variante A
  \def\entreprise{Le cabinet d'architectes \textbf{ArchiNova}}
  \def\Rnom{R-ArchiNova}
  \def\resUn{Bureau d'études}\def\pUn{BE}\def\netUn{192.168.40}
  \def\resDeux{Accueil}\def\pDeux{ACC}\def\netDeux{172.16}
  \def\srvNom{SRV-Intranet}\def\netSrv{192.168.200}\def\srvH{10}
  \def\ssid{ArchiNova-WiFi}
  \def\AIrows{%
    \hl\ip{192.168.40.25} & \sol{C} & \sol{255.255.255.0} & \sol{192.168.40.0} & \sol{192.168.40.255} & \sol{254} \\ \hline
    \hl\ip{172.16.5.10}   & \sol{B} & \sol{255.255.0.0}   & \sol{172.16.0.0}   & \sol{172.16.255.255} & \sol{65\,534} \\ \hline
    \hl\ip{10.3.2.1}      & \sol{A} & \sol{255.0.0.0}     & \sol{10.0.0.0}     & \sol{10.255.255.255} & \sol{16\,777\,214} \\ \hline}
  \def\apipa{169.254.12.7}
  \def\Cnet{192.168.60}
  \def\Crows{%
    PC-A & \ip{192.168.60.10} & \ip{255.255.255.0} & \ip{192.168.60.254} & — \\
    PC-B & \ip{192.168.60.20} & \ip{255.255.255.0} & \ip{192.168.60.1}   & — \\
    PC-C & \ip{192.168.61.30} & \ip{255.255.255.0} & \ip{192.168.60.254} & — \\
    PC-D & \ip{192.168.60.10} & \ip{255.255.255.0} & \ip{192.168.60.254} & — \\
    SRV-C & \ip{10.0.0.5}     & \ip{255.0.0.0}     & \ip{10.0.0.254}     & — \\
    R1 — G0/0 & \ip{192.168.60.254} & \ip{255.255.255.0} & — & On \\
    R1 — G0/1 & \ip{10.0.0.254}     & \ip{255.0.0.0}     & — & Off \\}
  \def\Csol{%
    \hl 1 & \sol{PC-B} & \sol{Passerelle .60.1 : pas l'adresse du routeur.} & \sol{Joint son réseau, mais aucun autre (serveur).} & \sol{Passerelle 192.168.60.254.} \\ \hline
    \hl 2 & \sol{PC-C} & \sol{IP hors du réseau 192.168.60.0/24.} & \sol{Ne joint ni les postes ni le routeur.} & \sol{Ex. 192.168.60.30.} \\ \hline
    \hl 3 & \sol{PC-D} & \sol{Même IP que PC-A (doublon).} & \sol{Conflit d'adresse : PC-A et PC-D perturbés.} & \sol{Adresse unique, ex. .60.40.} \\ \hline
    \hl 4 & \sol{R1 G0/1} & \sol{Interface désactivée (Off).} & \sol{Réseau du serveur injoignable pour tous.} & \sol{Port Status : On.} \\ \hline}
  \def\Cbareme{PC-A et SRV-C sont correctement configurés ; accepter PC-A cité à la place de PC-D pour le doublon.}
\else
  \def\entreprise{La clinique vétérinaire \textbf{VetoPlus}}
  \def\Rnom{R-VetoPlus}
  \def\resUn{Soins}\def\pUn{SOI}\def\netUn{192.168.70}
  \def\resDeux{Secrétariat}\def\pDeux{SEC}\def\netDeux{172.20}
  \def\srvNom{SRV-Dossiers}\def\netSrv{192.168.250}\def\srvH{20}
  \def\ssid{VetoPlus-WiFi}
  \def\AIrows{%
    \hl\ip{192.168.15.130} & \sol{C} & \sol{255.255.255.0} & \sol{192.168.15.0} & \sol{192.168.15.255} & \sol{254} \\ \hline
    \hl\ip{172.31.200.7}   & \sol{B} & \sol{255.255.0.0}   & \sol{172.31.0.0}   & \sol{172.31.255.255} & \sol{65\,534} \\ \hline
    \hl\ip{10.45.6.9}      & \sol{A} & \sol{255.0.0.0}     & \sol{10.0.0.0}     & \sol{10.255.255.255} & \sol{16\,777\,214} \\ \hline}
  \def\apipa{169.254.88.3}
  \def\Cnet{192.168.80}
  \def\Crows{%
    PC-A & \ip{192.168.8.11}   & \ip{255.255.255.0} & \ip{192.168.80.254} & — \\
    PC-B & \ip{192.168.80.12}  & \ip{255.255.255.0} & \ip{192.168.80.254} & — \\
    PC-C & \ip{192.168.80.13}  & \ip{255.255.255.0} & \ip{192.168.80.100} & — \\
    PC-D & \ip{192.168.80.12}  & \ip{255.255.255.0} & \ip{192.168.80.254} & — \\
    SRV-C & \ip{10.0.0.8}      & \ip{255.0.0.0}     & \ip{10.0.0.254}     & — \\
    R1 — G0/0 & \ip{192.168.80.254} & \ip{255.255.255.0} & — & Off \\
    R1 — G0/1 & \ip{10.0.0.254}     & \ip{255.0.0.0}     & — & On \\}
  \def\Csol{%
    \hl 1 & \sol{PC-A} & \sol{IP hors du réseau 192.168.80.0/24.} & \sol{Ne joint ni les postes ni le routeur.} & \sol{Ex. 192.168.80.11.} \\ \hline
    \hl 2 & \sol{PC-C} & \sol{Passerelle .80.100 : pas l'adresse du routeur.} & \sol{Joint son réseau, mais aucun autre (serveur).} & \sol{Passerelle 192.168.80.254.} \\ \hline
    \hl 3 & \sol{PC-D} & \sol{Même IP que PC-B (doublon).} & \sol{Conflit d'adresse : PC-B et PC-D perturbés.} & \sol{Adresse unique, ex. .80.14.} \\ \hline
    \hl 4 & \sol{R1 G0/0} & \sol{Interface côté LAN désactivée (Off).} & \sol{Aucun poste ne joint le routeur ni le serveur.} & \sol{Port Status : On.} \\ \hline}
  \def\Cbareme{PC-B et SRV-C sont correctement configurés ; accepter PC-B cité à la place de PC-D pour le doublon.}
\fi

\pagestyle{fancy}
\fancyhf{}
\renewcommand{\headrulewidth}{0pt}
\fancyfoot[L]{\footnotesize\itshape\color{grisfoot}Évaluation — Réseaux locaux — Sujet \variante\ifcorrige\ — CORRIGÉ\fi}
\fancyfoot[R]{\footnotesize\color{grisfoot}\thepage/3}

\setlist{itemsep=0pt,topsep=2pt,parsep=0pt}
\setlength{\parindent}{0pt}
\setlength{\parskip}{2pt}

\begin{document}

\begin{center}
  {\Large\bfseries Évaluation pratique — Réseaux locaux \quad
   \fcolorbox{black}{white}{\ SUJET \variante\ }\par}
  \vspace{1pt}
  {\bfseries Adressage IP, commutation et routage — BTS CIEL 1re année — durée : 2 heures\par}
  \ifcorrige\vspace{2pt}{\bfseries\color{corr}CORRIGÉ ET BARÈME — document enseignant}\par\fi
\end{center}
\vspace{-2pt}
{\renewcommand{\arraystretch}{1.35}
\begin{tabularx}{\linewidth}{|X|X|p{2.2cm}|p{2.6cm}|}
  \hline
  \textbf{Nom :} & \textbf{Prénom :} & \textbf{Poste n°} & \textbf{Note :} \hfill /20 \\ \hline
\end{tabularx}}

\textbf{Consignes :} travail \textbf{individuel}, documents et Internet \textbf{non autorisés}. Ce sujet est le
\textbf{document-réponse}. Fichier Packet Tracer : \ip{Nom\_Prenom\_EvalLAN\_\variante.pkt}, à enregistrer
régulièrement dans le dossier de rendu. Les « visas » sont apposés par le professeur, que vous appelez à chaque test réussi.
Durées conseillées : A 30 min — B 1 h 10 — C 20 min.

%==========================================================================
\section{Partie A — Questions de cours et calculs \hfill /6}

\question{A1}{2 pts}{Complétez le tableau en considérant le \textbf{masque par défaut} de la classe de chaque adresse.}
{\small\renewcommand{\arraystretch}{1.15}
\begin{tabularx}{\linewidth}{|l|c|C|C|C|c|}
  \hline
  \textbf{Adresse IP} & \textbf{Classe} & \textbf{Masque par défaut} & \textbf{Adresse réseau} &
  \textbf{Adresse de diffusion} & \textbf{Nb d'hôtes} \\ \hline
  \AIrows
\end{tabularx}}
\bareme{0,5 pt par ligne juste + 0,5 pt si les trois nombres d'hôtes sont justes ($2^{n}-2$, $n$ = nombre de bits d'hôte).}

\question{A2}{1 pt}{Indiquez le type de câble (\textbf{droit} ou \textbf{croisé}) pour chaque liaison.}
{\small\renewcommand{\arraystretch}{1.15}
\begin{tabularx}{\linewidth}{|X|c||X|c|}
  \hline
  \hl PC $\leftrightarrow$ commutateur & \makebox[1.6cm]{\sol{droit}} & Commutateur $\leftrightarrow$ routeur & \makebox[1.6cm]{\sol{droit}} \\ \hline
  \hl PC $\leftrightarrow$ PC (réseau P2P) & \sol{croisé} & Concentrateur $\leftrightarrow$ concentrateur & \sol{croisé} \\ \hline
\end{tabularx}}
\bareme{0,25 pt par case. Équipements de types différents : câble droit ; de même type : câble croisé.}

\question{A3}{1 pt}{Comment un \textbf{concentrateur (hub)} et un \textbf{commutateur (switch)} traitent-ils une trame
reçue ? Lequel limite les collisions, et pourquoi ?}
\zone{2.9cm}{Le \textbf{hub} (couche 1) répète la trame sur \textbf{tous ses autres ports} : un seul domaine de collision
partagé. Le \textbf{switch} (couche 2) lit l'adresse MAC de destination et, grâce à sa \textbf{table MAC}, ne l'envoie
\textbf{que sur le port du destinataire} : un domaine de collision par port, c'est lui qui limite les collisions.
\emph{0,5 pt hub, 0,5 pt switch.}}

\question{A4}{1 pt}{Qu'est-ce que la \textbf{passerelle par défaut} d'un poste ? Quelle condition son adresse doit-elle respecter ?}
\zone{2.4cm}{Adresse IP de l'interface du \textbf{routeur} à laquelle le poste envoie les paquets destinés à un
\textbf{autre réseau}. Elle doit \textbf{appartenir au même réseau} que le poste. \emph{0,5 + 0,5 pt.}}

\question{A5}{0,5 pt}{Un portable configuré en DHCP affiche \ip{\apipa} / \ip{255.255.0.0}. Que s'est-il passé ?}
\zone{1.5cm}{Aucun \textbf{serveur DHCP} n'a répondu : le poste s'est attribué une adresse \textbf{APIPA} (169.254.0.0/16).}

\question{A6}{0,5 pt}{Quel protocole utilise \texttt{ping} ? Nommez les deux messages échangés et la couche OSI des adresses IP.}
\zone{1.5cm}{\textbf{ICMP} ; \textbf{Echo Request} / \textbf{Echo Reply} ; adresses IP en \textbf{couche 3 (Réseau)}.}

%==========================================================================
\newpage
\section{Partie B — Réalisation sous Packet Tracer \hfill /10}

\begin{cdc}
\entreprise{} vous confie son réseau : \textbf{trois réseaux} reliés par un \textbf{routeur unique Cisco 2911}
nommé \ip{\Rnom} (interfaces G0/0, G0/1, G0/2).
\begin{itemize}[leftmargin=12pt]
  \item \textbf{Réseau \resUn} : \textbf{4 postes} PC-\pUn1 à PC-\pUn4, commutateur \ip{SW-\pUn}, adressage statique,
        réseau \ip{\netUn.0/24}.
  \item \textbf{Réseau \resDeux} : \textbf{2 postes} PC-\pDeux1, PC-\pDeux2 et \textbf{1 imprimante} IMP-\pDeux,
        commutateur \ip{SW-\pDeux}, adressage statique, réseau de classe B \ip{\netDeux.0.0/16}.
  \item \textbf{Réseau Serveur} : serveur \ip{\srvNom} d'adresse \ip{\netSrv.\srvH/24}, relié \textbf{directement} au routeur.
  \item La passerelle de chaque réseau est sa \textbf{dernière adresse utilisable}. Tous les équipements communiquent
        entre eux et accèdent à la page web du serveur. Nommage imposé (« Display Name »).
\end{itemize}
\end{cdc}

\question{B1}{2 pts}{Complétez le plan d'adressage \textbf{avant} la réalisation.}
{\footnotesize\setlength{\tabcolsep}{3pt}\renewcommand{\arraystretch}{1.15}
\begin{tabularx}{\linewidth}{|l|C|C|C|C|C|C|}
  \hline
  \textbf{Réseau} & \textbf{Adresse réseau} & \textbf{Masque} & \textbf{1re @ utilisable} &
  \textbf{Dernière @ utilisable} & \textbf{Diffusion} & \textbf{Passerelle} \\ \hline
  \hl \resUn & \sol{\netUn.0} & \sol{255.255.255.0} & \sol{\netUn.1} & \sol{\netUn.254} & \sol{\netUn.255} & \sol{\netUn.254} \\ \hline
  \hl \resDeux & \sol{\netDeux.0.0} & \sol{255.255.0.0} & \sol{\netDeux.0.1} & \sol{\netDeux.255.254} & \sol{\netDeux.255.255} & \sol{\netDeux.255.254} \\ \hline
  \hl Serveur & \sol{\netSrv.0} & \sol{255.255.255.0} & \sol{\netSrv.1} & \sol{\netSrv.254} & \sol{\netSrv.255} & \sol{\netSrv.254} \\ \hline
\end{tabularx}}
\bareme{0,5 pt par ligne juste + 0,5 pt si la convention de passerelle est respectée partout.}

\question{B2}{2 pts}{Construisez la topologie (équipements, câblage, nommage) et relevez l'adressage configuré.}
{\footnotesize\setlength{\tabcolsep}{3pt}\renewcommand{\arraystretch}{1.1}
\begin{tabularx}{\linewidth}{|l|C|C|C||l|C|C|}
  \hline
  \textbf{Hôte} & \textbf{Adresse IP} & \textbf{Masque} & \textbf{Passerelle} & \textbf{Interface} & \textbf{Adresse IP} & \textbf{Masque} \\ \hline
  \hl PC-\pUn1 & \sol{\netUn.1} & \sol{255.255.255.0} & \sol{\netUn.254} & \Rnom{} G0/0 & \sol{\netUn.254} & \sol{255.255.255.0} \\ \hline
  \hl PC-\pDeux1 & \sol{\netDeux.0.1} & \sol{255.255.0.0} & \sol{\netDeux.255.254} & \Rnom{} G0/1 & \sol{\netDeux.255.254} & \sol{255.255.0.0} \\ \hline
  \hl IMP-\pDeux & \sol{\netDeux.0.10} & \sol{255.255.0.0} & \sol{\netDeux.255.254} & \Rnom{} G0/2 & \sol{\netSrv.254} & \sol{255.255.255.0} \\ \hline
  \hl \srvNom & \ip{\netSrv.\srvH} & \sol{255.255.255.0} & \sol{\netSrv.254} & \multicolumn{3}{c|}{\cellcolor{black!8}} \\ \hline
\end{tabularx}}
\bareme{Adresses d'hôtes libres (uniques, dans le bon réseau) ; affectation G0/x libre. Sur le .pkt : équipements 0,5 ;
câblage 1 (câbles droits vers les switchs et switch–routeur, \textbf{croisé} serveur–routeur, liens verts) ; nommage 0,5.}

\question{B3}{3 pts}{Configurez tous les équipements : IP, masques, passerelles des hôtes ; interfaces du routeur
(« Port Status : On »).}
\bareme{Sur le .pkt : interfaces routeur 1 pt ; IP/masques des hôtes 1 pt ; passerelles (imprimante et serveur compris)
1 pt. $-0,25$ par erreur, sans descendre sous 0 par item.}

\question{B4}{3 pts}{Réalisez les tests, notez la commande et le résultat, puis faites viser chaque test réussi.}
{\small\renewcommand{\arraystretch}{1.1}
\begin{tabularx}{\linewidth}{|c|l|l|C|c|c|}
  \hline
  \textbf{N°} & \textbf{Source} & \textbf{Destination} & \textbf{Commande / action} & \textbf{Résultat} & \textbf{Visa prof.} \\ \hline
  \hl 1 & PC-\pUn1 & PC-\pUn4 & \sol{ping \netUn.4} & \sol{réussi} & \hspace{1.4cm} \\ \hline
  \hl 2 & PC-\pDeux1 & IMP-\pDeux & \sol{ping \netDeux.0.10} & \sol{réussi} & \\ \hline
  \hl 3 & PC-\pUn2 & PC-\pDeux2 & \sol{ping \netDeux.0.2} & \sol{réussi} & \\ \hline
  \hl 4 & PC-\pDeux2 & \srvNom & \sol{ping \netSrv.\srvH} & \sol{réussi} & \\ \hline
  \hl 5 & PC-\pUn3 & page web du serveur & \sol{navigateur : {\NoAutoSpacing http://\netSrv.\srvH}} & \sol{page affichée} & \\ \hline
\end{tabularx}}
\bareme{0,5 pt par test visé + 0,5 pt pour B5. Le premier ping inter-réseaux peut échouer (résolution ARP) ; les suivants doivent réussir.}

\question{B5}{incl. B4}{Test n°3 : par quels équipements passe le paquet ? Pourquoi PC-\pUn2 l'envoie-t-il au routeur ?}
\zone{2.3cm}{PC-\pUn2 $\to$ SW-\pUn{} $\to$ \Rnom{} $\to$ SW-\pDeux{} $\to$ PC-\pDeux2. Grâce à son masque, PC-\pUn2 constate que
la destination (\netDeux.0.0/16) n'est \textbf{pas dans son réseau} (\netUn.0/24) : il envoie le paquet à sa
\textbf{passerelle par défaut}, qui le route.}

%==========================================================================
\newpage
\section{Partie C — Diagnostic d'un réseau en panne \hfill /4}

Un technicien constate qu'\textbf{aucun poste ne parvient à joindre le serveur SRV-C} et que certains postes
communiquent mal entre eux. Les câbles sont du bon type.

\begin{center}
\begin{minipage}{0.37\linewidth}\centering
\begin{tikzpicture}[font=\footnotesize,
  eq/.style={draw,fill=qbg,rounded corners=2pt,minimum width=9mm,minimum height=5mm,inner sep=2pt},
  every path/.style={thick}]
  \node[eq] (SW) {SW-1};
  \node[eq,right=9mm of SW] (R) {R1};
  \node[eq,right=7mm of R] (SRV) {SRV-C};
  \node[eq,above left=5mm and -3mm of SW] (A) {PC-A};
  \node[eq,above right=5mm and -3mm of SW] (B) {PC-B};
  \node[eq,below left=5mm and -3mm of SW] (C) {PC-C};
  \node[eq,below right=5mm and -3mm of SW] (D) {PC-D};
  \foreach \n in {A,B,C,D} \draw (\n) -- (SW);
  \draw (SW) -- node[above,font=\tiny]{G0/0} (R);
  \draw (R) -- node[above,font=\tiny]{G0/1} (SRV);
\end{tikzpicture}
\end{minipage}\hfill
\begin{minipage}{0.61\linewidth}\centering
{\small\renewcommand{\arraystretch}{1.05}\setlength{\tabcolsep}{4pt}
\begin{tabular}{lllll}
  \toprule
  \textbf{Équipement} & \textbf{Adresse IP} & \textbf{Masque} & \textbf{Passerelle} & \textbf{Port} \\
  \midrule
  \Crows
  \bottomrule
\end{tabular}}
\end{minipage}
\end{center}

\question{C1}{4 pts}{Relevez les \textbf{quatre anomalies}. Pour chacune : conséquence et correction.}
{\small\renewcommand{\arraystretch}{1.1}
\ifcorrige\def\hc{}\else\def\hc{\rule[-10mm]{0pt}{14mm}}\fi
\begin{tabularx}{\linewidth}{|c|>{\raggedright\arraybackslash}p{2cm}|L|L|L|}
  \hline
  \textbf{N°} & \textbf{Équipement} & \textbf{Anomalie} & \textbf{Conséquence} & \textbf{Correction} \\ \hline
  \ifcorrige\Csol\else
  \hc 1 & & & & \\ \hline
  \hc 2 & & & & \\ \hline
  \hc 3 & & & & \\ \hline
  \hc 4 & & & & \\ \hline
  \fi
\end{tabularx}}
\bareme{1 pt par anomalie (0,25 équipement, anomalie, conséquence, correction ; ordre indifférent). \Cbareme}

%==========================================================================
\section{Bonus — Extension sans fil \hfill +1}

\begin{cdc}Ajoutez au réseau \resUn{} un \textbf{AccessPoint-PT} relié à SW-\pUn{} (SSID \ip{\ssid}, \textbf{WPA2-PSK})
et un portable \ip{Laptop-\pUn} équipé d'un module \texttt{WPC300N}, en adressage \textbf{statique}.
Vérifiez qu'il joint \srvNom.\end{cdc}
{\small\renewcommand{\arraystretch}{1.1}
\begin{tabularx}{\linewidth}{|l|C|C|C|c|}
  \hline
  \textbf{Équipement} & \textbf{Adresse IP} & \textbf{Masque} & \textbf{Passerelle} & \textbf{Visa prof.} \\ \hline
  \hl Laptop-\pUn & \sol{\netUn.50} & \sol{255.255.255.0} & \sol{\netUn.254} & \hspace{1.4cm} \\ \hline
\end{tabularx}}
\bareme{Le point d'accès fait office de pont : le portable est dans \netUn.0/24. Bonus si associé en WPA2-PSK et ping
vers \netSrv.\srvH{} réussi. Note plafonnée à 20.}

\end{document}
"
%   pdflatex -jobname=Eval_Reseau_CorrigeB "\def\variante{B}\def\modecorrige{1}% Évaluation pratique — Réseaux locaux (BTS CIEL 1re année) — sujets A et B
% Compilation (4 PDF) :
%   pdflatex -jobname=Eval_Reseau_SujetA   "\def\variante{A}\def\modecorrige{0}% Évaluation pratique — Réseaux locaux (BTS CIEL 1re année) — sujets A et B
% Compilation (4 PDF) :
%   pdflatex -jobname=Eval_Reseau_SujetA   "\def\variante{A}\def\modecorrige{0}\input{Eval_TP_Reseau_LAN.tex}"
%   pdflatex -jobname=Eval_Reseau_SujetB   "\def\variante{B}\def\modecorrige{0}\input{Eval_TP_Reseau_LAN.tex}"
%   pdflatex -jobname=Eval_Reseau_CorrigeA "\def\variante{A}\def\modecorrige{1}\input{Eval_TP_Reseau_LAN.tex}"
%   pdflatex -jobname=Eval_Reseau_CorrigeB "\def\variante{B}\def\modecorrige{1}\input{Eval_TP_Reseau_LAN.tex}"
\documentclass[10pt,a4paper]{article}

\usepackage[T1]{fontenc}
\usepackage[utf8]{inputenc}
\usepackage{lmodern}
\usepackage[french]{babel}
\usepackage[a4paper,hmargin=1.5cm,top=1.3cm,bottom=1.6cm,footskip=0.7cm]{geometry}
\usepackage[table]{xcolor}
\usepackage{booktabs,tabularx,array}
\usepackage{enumitem}
\usepackage{fancyhdr}
\usepackage[most]{tcolorbox}
\usepackage{tikz}
\usetikzlibrary{positioning}
\usepackage{titlesec}

\providecommand{\modecorrige}{0}
\providecommand{\variante}{A}
\newif\ifcorrige
\ifnum\modecorrige=1 \corrigetrue \else \corrigefalse \fi

\definecolor{ipred}{HTML}{C0392B}
\definecolor{cdcbar}{HTML}{A67C3D}
\definecolor{cdcbg}{HTML}{F4F1EA}
\definecolor{qbar}{HTML}{4A6B5D}
\definecolor{qbg}{HTML}{EDF1EF}
\definecolor{corr}{HTML}{1F5FA8}
\definecolor{corrbg}{HTML}{E8F0FA}
\definecolor{grisfoot}{HTML}{828282}

\newcommand{\ip}[1]{\textcolor{ipred}{\texttt{#1}}}
\newcommand{\pts}[1]{\hfill\mbox{\textbf{(#1)}}}
\newcommand{\sol}[1]{\ifcorrige\textcolor{corr}{#1}\fi}
\newcommand{\bareme}[1]{\ifcorrige\par{\footnotesize\color{corr}#1\par}\fi}
\newcommand{\hl}{\ifcorrige\rule[-1.6mm]{0pt}{5.4mm}\else\rule[-2.2mm]{0pt}{6.6mm}\fi}
\newcolumntype{C}{>{\centering\arraybackslash}X}
\newcolumntype{L}{>{\raggedright\arraybackslash}X}

\titleformat{\section}{\large\bfseries}{}{0pt}{}
\titlespacing*{\section}{0pt}{6pt}{3pt}

\tcbset{qstyle/.style={enhanced,boxrule=0pt,leftrule=2.5pt,arc=0pt,outer arc=0pt,
  left=5pt,right=5pt,top=2pt,bottom=2pt,before skip=4pt,after skip=3pt}}
\newtcolorbox{cdc}{qstyle,colback=cdcbg,colframe=cdcbar}
\newtcolorbox{quest}{qstyle,colback=qbg,colframe=qbar}

% Zone de réponse : vierge dans le sujet, remplie dans le corrigé
\newcommand{\zone}[2]{%
  \ifcorrige
    \begin{tcolorbox}[colback=corrbg,colframe=corr,boxrule=0.5pt,arc=1pt,left=4pt,right=4pt,top=1pt,bottom=1pt,
      before skip=2pt,after skip=2pt,fontupper=\footnotesize\color{corr}]#2\end{tcolorbox}%
  \else
    \begin{tcolorbox}[colback=white,colframe=black!45,boxrule=0.4pt,arc=1pt,height=#1,
      before skip=2pt,after skip=2pt]\end{tcolorbox}%
  \fi}

\newcommand{\question}[3]{\begin{quest}\textbf{#1.} #3\pts{#2}\end{quest}}

%==========================================================================
% Données propres à chaque variante
%==========================================================================
\if\variante A
  \def\entreprise{Le cabinet d'architectes \textbf{ArchiNova}}
  \def\Rnom{R-ArchiNova}
  \def\resUn{Bureau d'études}\def\pUn{BE}\def\netUn{192.168.40}
  \def\resDeux{Accueil}\def\pDeux{ACC}\def\netDeux{172.16}
  \def\srvNom{SRV-Intranet}\def\netSrv{192.168.200}\def\srvH{10}
  \def\ssid{ArchiNova-WiFi}
  \def\AIrows{%
    \hl\ip{192.168.40.25} & \sol{C} & \sol{255.255.255.0} & \sol{192.168.40.0} & \sol{192.168.40.255} & \sol{254} \\ \hline
    \hl\ip{172.16.5.10}   & \sol{B} & \sol{255.255.0.0}   & \sol{172.16.0.0}   & \sol{172.16.255.255} & \sol{65\,534} \\ \hline
    \hl\ip{10.3.2.1}      & \sol{A} & \sol{255.0.0.0}     & \sol{10.0.0.0}     & \sol{10.255.255.255} & \sol{16\,777\,214} \\ \hline}
  \def\apipa{169.254.12.7}
  \def\Cnet{192.168.60}
  \def\Crows{%
    PC-A & \ip{192.168.60.10} & \ip{255.255.255.0} & \ip{192.168.60.254} & — \\
    PC-B & \ip{192.168.60.20} & \ip{255.255.255.0} & \ip{192.168.60.1}   & — \\
    PC-C & \ip{192.168.61.30} & \ip{255.255.255.0} & \ip{192.168.60.254} & — \\
    PC-D & \ip{192.168.60.10} & \ip{255.255.255.0} & \ip{192.168.60.254} & — \\
    SRV-C & \ip{10.0.0.5}     & \ip{255.0.0.0}     & \ip{10.0.0.254}     & — \\
    R1 — G0/0 & \ip{192.168.60.254} & \ip{255.255.255.0} & — & On \\
    R1 — G0/1 & \ip{10.0.0.254}     & \ip{255.0.0.0}     & — & Off \\}
  \def\Csol{%
    \hl 1 & \sol{PC-B} & \sol{Passerelle .60.1 : pas l'adresse du routeur.} & \sol{Joint son réseau, mais aucun autre (serveur).} & \sol{Passerelle 192.168.60.254.} \\ \hline
    \hl 2 & \sol{PC-C} & \sol{IP hors du réseau 192.168.60.0/24.} & \sol{Ne joint ni les postes ni le routeur.} & \sol{Ex. 192.168.60.30.} \\ \hline
    \hl 3 & \sol{PC-D} & \sol{Même IP que PC-A (doublon).} & \sol{Conflit d'adresse : PC-A et PC-D perturbés.} & \sol{Adresse unique, ex. .60.40.} \\ \hline
    \hl 4 & \sol{R1 G0/1} & \sol{Interface désactivée (Off).} & \sol{Réseau du serveur injoignable pour tous.} & \sol{Port Status : On.} \\ \hline}
  \def\Cbareme{PC-A et SRV-C sont correctement configurés ; accepter PC-A cité à la place de PC-D pour le doublon.}
\else
  \def\entreprise{La clinique vétérinaire \textbf{VetoPlus}}
  \def\Rnom{R-VetoPlus}
  \def\resUn{Soins}\def\pUn{SOI}\def\netUn{192.168.70}
  \def\resDeux{Secrétariat}\def\pDeux{SEC}\def\netDeux{172.20}
  \def\srvNom{SRV-Dossiers}\def\netSrv{192.168.250}\def\srvH{20}
  \def\ssid{VetoPlus-WiFi}
  \def\AIrows{%
    \hl\ip{192.168.15.130} & \sol{C} & \sol{255.255.255.0} & \sol{192.168.15.0} & \sol{192.168.15.255} & \sol{254} \\ \hline
    \hl\ip{172.31.200.7}   & \sol{B} & \sol{255.255.0.0}   & \sol{172.31.0.0}   & \sol{172.31.255.255} & \sol{65\,534} \\ \hline
    \hl\ip{10.45.6.9}      & \sol{A} & \sol{255.0.0.0}     & \sol{10.0.0.0}     & \sol{10.255.255.255} & \sol{16\,777\,214} \\ \hline}
  \def\apipa{169.254.88.3}
  \def\Cnet{192.168.80}
  \def\Crows{%
    PC-A & \ip{192.168.8.11}   & \ip{255.255.255.0} & \ip{192.168.80.254} & — \\
    PC-B & \ip{192.168.80.12}  & \ip{255.255.255.0} & \ip{192.168.80.254} & — \\
    PC-C & \ip{192.168.80.13}  & \ip{255.255.255.0} & \ip{192.168.80.100} & — \\
    PC-D & \ip{192.168.80.12}  & \ip{255.255.255.0} & \ip{192.168.80.254} & — \\
    SRV-C & \ip{10.0.0.8}      & \ip{255.0.0.0}     & \ip{10.0.0.254}     & — \\
    R1 — G0/0 & \ip{192.168.80.254} & \ip{255.255.255.0} & — & Off \\
    R1 — G0/1 & \ip{10.0.0.254}     & \ip{255.0.0.0}     & — & On \\}
  \def\Csol{%
    \hl 1 & \sol{PC-A} & \sol{IP hors du réseau 192.168.80.0/24.} & \sol{Ne joint ni les postes ni le routeur.} & \sol{Ex. 192.168.80.11.} \\ \hline
    \hl 2 & \sol{PC-C} & \sol{Passerelle .80.100 : pas l'adresse du routeur.} & \sol{Joint son réseau, mais aucun autre (serveur).} & \sol{Passerelle 192.168.80.254.} \\ \hline
    \hl 3 & \sol{PC-D} & \sol{Même IP que PC-B (doublon).} & \sol{Conflit d'adresse : PC-B et PC-D perturbés.} & \sol{Adresse unique, ex. .80.14.} \\ \hline
    \hl 4 & \sol{R1 G0/0} & \sol{Interface côté LAN désactivée (Off).} & \sol{Aucun poste ne joint le routeur ni le serveur.} & \sol{Port Status : On.} \\ \hline}
  \def\Cbareme{PC-B et SRV-C sont correctement configurés ; accepter PC-B cité à la place de PC-D pour le doublon.}
\fi

\pagestyle{fancy}
\fancyhf{}
\renewcommand{\headrulewidth}{0pt}
\fancyfoot[L]{\footnotesize\itshape\color{grisfoot}Évaluation — Réseaux locaux — Sujet \variante\ifcorrige\ — CORRIGÉ\fi}
\fancyfoot[R]{\footnotesize\color{grisfoot}\thepage/3}

\setlist{itemsep=0pt,topsep=2pt,parsep=0pt}
\setlength{\parindent}{0pt}
\setlength{\parskip}{2pt}

\begin{document}

\begin{center}
  {\Large\bfseries Évaluation pratique — Réseaux locaux \quad
   \fcolorbox{black}{white}{\ SUJET \variante\ }\par}
  \vspace{1pt}
  {\bfseries Adressage IP, commutation et routage — BTS CIEL 1re année — durée : 2 heures\par}
  \ifcorrige\vspace{2pt}{\bfseries\color{corr}CORRIGÉ ET BARÈME — document enseignant}\par\fi
\end{center}
\vspace{-2pt}
{\renewcommand{\arraystretch}{1.35}
\begin{tabularx}{\linewidth}{|X|X|p{2.2cm}|p{2.6cm}|}
  \hline
  \textbf{Nom :} & \textbf{Prénom :} & \textbf{Poste n°} & \textbf{Note :} \hfill /20 \\ \hline
\end{tabularx}}

\textbf{Consignes :} travail \textbf{individuel}, documents et Internet \textbf{non autorisés}. Ce sujet est le
\textbf{document-réponse}. Fichier Packet Tracer : \ip{Nom\_Prenom\_EvalLAN\_\variante.pkt}, à enregistrer
régulièrement dans le dossier de rendu. Les « visas » sont apposés par le professeur, que vous appelez à chaque test réussi.
Durées conseillées : A 30 min — B 1 h 10 — C 20 min.

%==========================================================================
\section{Partie A — Questions de cours et calculs \hfill /6}

\question{A1}{2 pts}{Complétez le tableau en considérant le \textbf{masque par défaut} de la classe de chaque adresse.}
{\small\renewcommand{\arraystretch}{1.15}
\begin{tabularx}{\linewidth}{|l|c|C|C|C|c|}
  \hline
  \textbf{Adresse IP} & \textbf{Classe} & \textbf{Masque par défaut} & \textbf{Adresse réseau} &
  \textbf{Adresse de diffusion} & \textbf{Nb d'hôtes} \\ \hline
  \AIrows
\end{tabularx}}
\bareme{0,5 pt par ligne juste + 0,5 pt si les trois nombres d'hôtes sont justes ($2^{n}-2$, $n$ = nombre de bits d'hôte).}

\question{A2}{1 pt}{Indiquez le type de câble (\textbf{droit} ou \textbf{croisé}) pour chaque liaison.}
{\small\renewcommand{\arraystretch}{1.15}
\begin{tabularx}{\linewidth}{|X|c||X|c|}
  \hline
  \hl PC $\leftrightarrow$ commutateur & \makebox[1.6cm]{\sol{droit}} & Commutateur $\leftrightarrow$ routeur & \makebox[1.6cm]{\sol{droit}} \\ \hline
  \hl PC $\leftrightarrow$ PC (réseau P2P) & \sol{croisé} & Concentrateur $\leftrightarrow$ concentrateur & \sol{croisé} \\ \hline
\end{tabularx}}
\bareme{0,25 pt par case. Équipements de types différents : câble droit ; de même type : câble croisé.}

\question{A3}{1 pt}{Comment un \textbf{concentrateur (hub)} et un \textbf{commutateur (switch)} traitent-ils une trame
reçue ? Lequel limite les collisions, et pourquoi ?}
\zone{2.9cm}{Le \textbf{hub} (couche 1) répète la trame sur \textbf{tous ses autres ports} : un seul domaine de collision
partagé. Le \textbf{switch} (couche 2) lit l'adresse MAC de destination et, grâce à sa \textbf{table MAC}, ne l'envoie
\textbf{que sur le port du destinataire} : un domaine de collision par port, c'est lui qui limite les collisions.
\emph{0,5 pt hub, 0,5 pt switch.}}

\question{A4}{1 pt}{Qu'est-ce que la \textbf{passerelle par défaut} d'un poste ? Quelle condition son adresse doit-elle respecter ?}
\zone{2.4cm}{Adresse IP de l'interface du \textbf{routeur} à laquelle le poste envoie les paquets destinés à un
\textbf{autre réseau}. Elle doit \textbf{appartenir au même réseau} que le poste. \emph{0,5 + 0,5 pt.}}

\question{A5}{0,5 pt}{Un portable configuré en DHCP affiche \ip{\apipa} / \ip{255.255.0.0}. Que s'est-il passé ?}
\zone{1.5cm}{Aucun \textbf{serveur DHCP} n'a répondu : le poste s'est attribué une adresse \textbf{APIPA} (169.254.0.0/16).}

\question{A6}{0,5 pt}{Quel protocole utilise \texttt{ping} ? Nommez les deux messages échangés et la couche OSI des adresses IP.}
\zone{1.5cm}{\textbf{ICMP} ; \textbf{Echo Request} / \textbf{Echo Reply} ; adresses IP en \textbf{couche 3 (Réseau)}.}

%==========================================================================
\newpage
\section{Partie B — Réalisation sous Packet Tracer \hfill /10}

\begin{cdc}
\entreprise{} vous confie son réseau : \textbf{trois réseaux} reliés par un \textbf{routeur unique Cisco 2911}
nommé \ip{\Rnom} (interfaces G0/0, G0/1, G0/2).
\begin{itemize}[leftmargin=12pt]
  \item \textbf{Réseau \resUn} : \textbf{4 postes} PC-\pUn1 à PC-\pUn4, commutateur \ip{SW-\pUn}, adressage statique,
        réseau \ip{\netUn.0/24}.
  \item \textbf{Réseau \resDeux} : \textbf{2 postes} PC-\pDeux1, PC-\pDeux2 et \textbf{1 imprimante} IMP-\pDeux,
        commutateur \ip{SW-\pDeux}, adressage statique, réseau de classe B \ip{\netDeux.0.0/16}.
  \item \textbf{Réseau Serveur} : serveur \ip{\srvNom} d'adresse \ip{\netSrv.\srvH/24}, relié \textbf{directement} au routeur.
  \item La passerelle de chaque réseau est sa \textbf{dernière adresse utilisable}. Tous les équipements communiquent
        entre eux et accèdent à la page web du serveur. Nommage imposé (« Display Name »).
\end{itemize}
\end{cdc}

\question{B1}{2 pts}{Complétez le plan d'adressage \textbf{avant} la réalisation.}
{\footnotesize\setlength{\tabcolsep}{3pt}\renewcommand{\arraystretch}{1.15}
\begin{tabularx}{\linewidth}{|l|C|C|C|C|C|C|}
  \hline
  \textbf{Réseau} & \textbf{Adresse réseau} & \textbf{Masque} & \textbf{1re @ utilisable} &
  \textbf{Dernière @ utilisable} & \textbf{Diffusion} & \textbf{Passerelle} \\ \hline
  \hl \resUn & \sol{\netUn.0} & \sol{255.255.255.0} & \sol{\netUn.1} & \sol{\netUn.254} & \sol{\netUn.255} & \sol{\netUn.254} \\ \hline
  \hl \resDeux & \sol{\netDeux.0.0} & \sol{255.255.0.0} & \sol{\netDeux.0.1} & \sol{\netDeux.255.254} & \sol{\netDeux.255.255} & \sol{\netDeux.255.254} \\ \hline
  \hl Serveur & \sol{\netSrv.0} & \sol{255.255.255.0} & \sol{\netSrv.1} & \sol{\netSrv.254} & \sol{\netSrv.255} & \sol{\netSrv.254} \\ \hline
\end{tabularx}}
\bareme{0,5 pt par ligne juste + 0,5 pt si la convention de passerelle est respectée partout.}

\question{B2}{2 pts}{Construisez la topologie (équipements, câblage, nommage) et relevez l'adressage configuré.}
{\footnotesize\setlength{\tabcolsep}{3pt}\renewcommand{\arraystretch}{1.1}
\begin{tabularx}{\linewidth}{|l|C|C|C||l|C|C|}
  \hline
  \textbf{Hôte} & \textbf{Adresse IP} & \textbf{Masque} & \textbf{Passerelle} & \textbf{Interface} & \textbf{Adresse IP} & \textbf{Masque} \\ \hline
  \hl PC-\pUn1 & \sol{\netUn.1} & \sol{255.255.255.0} & \sol{\netUn.254} & \Rnom{} G0/0 & \sol{\netUn.254} & \sol{255.255.255.0} \\ \hline
  \hl PC-\pDeux1 & \sol{\netDeux.0.1} & \sol{255.255.0.0} & \sol{\netDeux.255.254} & \Rnom{} G0/1 & \sol{\netDeux.255.254} & \sol{255.255.0.0} \\ \hline
  \hl IMP-\pDeux & \sol{\netDeux.0.10} & \sol{255.255.0.0} & \sol{\netDeux.255.254} & \Rnom{} G0/2 & \sol{\netSrv.254} & \sol{255.255.255.0} \\ \hline
  \hl \srvNom & \ip{\netSrv.\srvH} & \sol{255.255.255.0} & \sol{\netSrv.254} & \multicolumn{3}{c|}{\cellcolor{black!8}} \\ \hline
\end{tabularx}}
\bareme{Adresses d'hôtes libres (uniques, dans le bon réseau) ; affectation G0/x libre. Sur le .pkt : équipements 0,5 ;
câblage 1 (câbles droits vers les switchs et switch–routeur, \textbf{croisé} serveur–routeur, liens verts) ; nommage 0,5.}

\question{B3}{3 pts}{Configurez tous les équipements : IP, masques, passerelles des hôtes ; interfaces du routeur
(« Port Status : On »).}
\bareme{Sur le .pkt : interfaces routeur 1 pt ; IP/masques des hôtes 1 pt ; passerelles (imprimante et serveur compris)
1 pt. $-0,25$ par erreur, sans descendre sous 0 par item.}

\question{B4}{3 pts}{Réalisez les tests, notez la commande et le résultat, puis faites viser chaque test réussi.}
{\small\renewcommand{\arraystretch}{1.1}
\begin{tabularx}{\linewidth}{|c|l|l|C|c|c|}
  \hline
  \textbf{N°} & \textbf{Source} & \textbf{Destination} & \textbf{Commande / action} & \textbf{Résultat} & \textbf{Visa prof.} \\ \hline
  \hl 1 & PC-\pUn1 & PC-\pUn4 & \sol{ping \netUn.4} & \sol{réussi} & \hspace{1.4cm} \\ \hline
  \hl 2 & PC-\pDeux1 & IMP-\pDeux & \sol{ping \netDeux.0.10} & \sol{réussi} & \\ \hline
  \hl 3 & PC-\pUn2 & PC-\pDeux2 & \sol{ping \netDeux.0.2} & \sol{réussi} & \\ \hline
  \hl 4 & PC-\pDeux2 & \srvNom & \sol{ping \netSrv.\srvH} & \sol{réussi} & \\ \hline
  \hl 5 & PC-\pUn3 & page web du serveur & \sol{navigateur : {\NoAutoSpacing http://\netSrv.\srvH}} & \sol{page affichée} & \\ \hline
\end{tabularx}}
\bareme{0,5 pt par test visé + 0,5 pt pour B5. Le premier ping inter-réseaux peut échouer (résolution ARP) ; les suivants doivent réussir.}

\question{B5}{incl. B4}{Test n°3 : par quels équipements passe le paquet ? Pourquoi PC-\pUn2 l'envoie-t-il au routeur ?}
\zone{2.3cm}{PC-\pUn2 $\to$ SW-\pUn{} $\to$ \Rnom{} $\to$ SW-\pDeux{} $\to$ PC-\pDeux2. Grâce à son masque, PC-\pUn2 constate que
la destination (\netDeux.0.0/16) n'est \textbf{pas dans son réseau} (\netUn.0/24) : il envoie le paquet à sa
\textbf{passerelle par défaut}, qui le route.}

%==========================================================================
\newpage
\section{Partie C — Diagnostic d'un réseau en panne \hfill /4}

Un technicien constate qu'\textbf{aucun poste ne parvient à joindre le serveur SRV-C} et que certains postes
communiquent mal entre eux. Les câbles sont du bon type.

\begin{center}
\begin{minipage}{0.37\linewidth}\centering
\begin{tikzpicture}[font=\footnotesize,
  eq/.style={draw,fill=qbg,rounded corners=2pt,minimum width=9mm,minimum height=5mm,inner sep=2pt},
  every path/.style={thick}]
  \node[eq] (SW) {SW-1};
  \node[eq,right=9mm of SW] (R) {R1};
  \node[eq,right=7mm of R] (SRV) {SRV-C};
  \node[eq,above left=5mm and -3mm of SW] (A) {PC-A};
  \node[eq,above right=5mm and -3mm of SW] (B) {PC-B};
  \node[eq,below left=5mm and -3mm of SW] (C) {PC-C};
  \node[eq,below right=5mm and -3mm of SW] (D) {PC-D};
  \foreach \n in {A,B,C,D} \draw (\n) -- (SW);
  \draw (SW) -- node[above,font=\tiny]{G0/0} (R);
  \draw (R) -- node[above,font=\tiny]{G0/1} (SRV);
\end{tikzpicture}
\end{minipage}\hfill
\begin{minipage}{0.61\linewidth}\centering
{\small\renewcommand{\arraystretch}{1.05}\setlength{\tabcolsep}{4pt}
\begin{tabular}{lllll}
  \toprule
  \textbf{Équipement} & \textbf{Adresse IP} & \textbf{Masque} & \textbf{Passerelle} & \textbf{Port} \\
  \midrule
  \Crows
  \bottomrule
\end{tabular}}
\end{minipage}
\end{center}

\question{C1}{4 pts}{Relevez les \textbf{quatre anomalies}. Pour chacune : conséquence et correction.}
{\small\renewcommand{\arraystretch}{1.1}
\ifcorrige\def\hc{}\else\def\hc{\rule[-10mm]{0pt}{14mm}}\fi
\begin{tabularx}{\linewidth}{|c|>{\raggedright\arraybackslash}p{2cm}|L|L|L|}
  \hline
  \textbf{N°} & \textbf{Équipement} & \textbf{Anomalie} & \textbf{Conséquence} & \textbf{Correction} \\ \hline
  \ifcorrige\Csol\else
  \hc 1 & & & & \\ \hline
  \hc 2 & & & & \\ \hline
  \hc 3 & & & & \\ \hline
  \hc 4 & & & & \\ \hline
  \fi
\end{tabularx}}
\bareme{1 pt par anomalie (0,25 équipement, anomalie, conséquence, correction ; ordre indifférent). \Cbareme}

%==========================================================================
\section{Bonus — Extension sans fil \hfill +1}

\begin{cdc}Ajoutez au réseau \resUn{} un \textbf{AccessPoint-PT} relié à SW-\pUn{} (SSID \ip{\ssid}, \textbf{WPA2-PSK})
et un portable \ip{Laptop-\pUn} équipé d'un module \texttt{WPC300N}, en adressage \textbf{statique}.
Vérifiez qu'il joint \srvNom.\end{cdc}
{\small\renewcommand{\arraystretch}{1.1}
\begin{tabularx}{\linewidth}{|l|C|C|C|c|}
  \hline
  \textbf{Équipement} & \textbf{Adresse IP} & \textbf{Masque} & \textbf{Passerelle} & \textbf{Visa prof.} \\ \hline
  \hl Laptop-\pUn & \sol{\netUn.50} & \sol{255.255.255.0} & \sol{\netUn.254} & \hspace{1.4cm} \\ \hline
\end{tabularx}}
\bareme{Le point d'accès fait office de pont : le portable est dans \netUn.0/24. Bonus si associé en WPA2-PSK et ping
vers \netSrv.\srvH{} réussi. Note plafonnée à 20.}

\end{document}
"
%   pdflatex -jobname=Eval_Reseau_SujetB   "\def\variante{B}\def\modecorrige{0}% Évaluation pratique — Réseaux locaux (BTS CIEL 1re année) — sujets A et B
% Compilation (4 PDF) :
%   pdflatex -jobname=Eval_Reseau_SujetA   "\def\variante{A}\def\modecorrige{0}\input{Eval_TP_Reseau_LAN.tex}"
%   pdflatex -jobname=Eval_Reseau_SujetB   "\def\variante{B}\def\modecorrige{0}\input{Eval_TP_Reseau_LAN.tex}"
%   pdflatex -jobname=Eval_Reseau_CorrigeA "\def\variante{A}\def\modecorrige{1}\input{Eval_TP_Reseau_LAN.tex}"
%   pdflatex -jobname=Eval_Reseau_CorrigeB "\def\variante{B}\def\modecorrige{1}\input{Eval_TP_Reseau_LAN.tex}"
\documentclass[10pt,a4paper]{article}

\usepackage[T1]{fontenc}
\usepackage[utf8]{inputenc}
\usepackage{lmodern}
\usepackage[french]{babel}
\usepackage[a4paper,hmargin=1.5cm,top=1.3cm,bottom=1.6cm,footskip=0.7cm]{geometry}
\usepackage[table]{xcolor}
\usepackage{booktabs,tabularx,array}
\usepackage{enumitem}
\usepackage{fancyhdr}
\usepackage[most]{tcolorbox}
\usepackage{tikz}
\usetikzlibrary{positioning}
\usepackage{titlesec}

\providecommand{\modecorrige}{0}
\providecommand{\variante}{A}
\newif\ifcorrige
\ifnum\modecorrige=1 \corrigetrue \else \corrigefalse \fi

\definecolor{ipred}{HTML}{C0392B}
\definecolor{cdcbar}{HTML}{A67C3D}
\definecolor{cdcbg}{HTML}{F4F1EA}
\definecolor{qbar}{HTML}{4A6B5D}
\definecolor{qbg}{HTML}{EDF1EF}
\definecolor{corr}{HTML}{1F5FA8}
\definecolor{corrbg}{HTML}{E8F0FA}
\definecolor{grisfoot}{HTML}{828282}

\newcommand{\ip}[1]{\textcolor{ipred}{\texttt{#1}}}
\newcommand{\pts}[1]{\hfill\mbox{\textbf{(#1)}}}
\newcommand{\sol}[1]{\ifcorrige\textcolor{corr}{#1}\fi}
\newcommand{\bareme}[1]{\ifcorrige\par{\footnotesize\color{corr}#1\par}\fi}
\newcommand{\hl}{\ifcorrige\rule[-1.6mm]{0pt}{5.4mm}\else\rule[-2.2mm]{0pt}{6.6mm}\fi}
\newcolumntype{C}{>{\centering\arraybackslash}X}
\newcolumntype{L}{>{\raggedright\arraybackslash}X}

\titleformat{\section}{\large\bfseries}{}{0pt}{}
\titlespacing*{\section}{0pt}{6pt}{3pt}

\tcbset{qstyle/.style={enhanced,boxrule=0pt,leftrule=2.5pt,arc=0pt,outer arc=0pt,
  left=5pt,right=5pt,top=2pt,bottom=2pt,before skip=4pt,after skip=3pt}}
\newtcolorbox{cdc}{qstyle,colback=cdcbg,colframe=cdcbar}
\newtcolorbox{quest}{qstyle,colback=qbg,colframe=qbar}

% Zone de réponse : vierge dans le sujet, remplie dans le corrigé
\newcommand{\zone}[2]{%
  \ifcorrige
    \begin{tcolorbox}[colback=corrbg,colframe=corr,boxrule=0.5pt,arc=1pt,left=4pt,right=4pt,top=1pt,bottom=1pt,
      before skip=2pt,after skip=2pt,fontupper=\footnotesize\color{corr}]#2\end{tcolorbox}%
  \else
    \begin{tcolorbox}[colback=white,colframe=black!45,boxrule=0.4pt,arc=1pt,height=#1,
      before skip=2pt,after skip=2pt]\end{tcolorbox}%
  \fi}

\newcommand{\question}[3]{\begin{quest}\textbf{#1.} #3\pts{#2}\end{quest}}

%==========================================================================
% Données propres à chaque variante
%==========================================================================
\if\variante A
  \def\entreprise{Le cabinet d'architectes \textbf{ArchiNova}}
  \def\Rnom{R-ArchiNova}
  \def\resUn{Bureau d'études}\def\pUn{BE}\def\netUn{192.168.40}
  \def\resDeux{Accueil}\def\pDeux{ACC}\def\netDeux{172.16}
  \def\srvNom{SRV-Intranet}\def\netSrv{192.168.200}\def\srvH{10}
  \def\ssid{ArchiNova-WiFi}
  \def\AIrows{%
    \hl\ip{192.168.40.25} & \sol{C} & \sol{255.255.255.0} & \sol{192.168.40.0} & \sol{192.168.40.255} & \sol{254} \\ \hline
    \hl\ip{172.16.5.10}   & \sol{B} & \sol{255.255.0.0}   & \sol{172.16.0.0}   & \sol{172.16.255.255} & \sol{65\,534} \\ \hline
    \hl\ip{10.3.2.1}      & \sol{A} & \sol{255.0.0.0}     & \sol{10.0.0.0}     & \sol{10.255.255.255} & \sol{16\,777\,214} \\ \hline}
  \def\apipa{169.254.12.7}
  \def\Cnet{192.168.60}
  \def\Crows{%
    PC-A & \ip{192.168.60.10} & \ip{255.255.255.0} & \ip{192.168.60.254} & — \\
    PC-B & \ip{192.168.60.20} & \ip{255.255.255.0} & \ip{192.168.60.1}   & — \\
    PC-C & \ip{192.168.61.30} & \ip{255.255.255.0} & \ip{192.168.60.254} & — \\
    PC-D & \ip{192.168.60.10} & \ip{255.255.255.0} & \ip{192.168.60.254} & — \\
    SRV-C & \ip{10.0.0.5}     & \ip{255.0.0.0}     & \ip{10.0.0.254}     & — \\
    R1 — G0/0 & \ip{192.168.60.254} & \ip{255.255.255.0} & — & On \\
    R1 — G0/1 & \ip{10.0.0.254}     & \ip{255.0.0.0}     & — & Off \\}
  \def\Csol{%
    \hl 1 & \sol{PC-B} & \sol{Passerelle .60.1 : pas l'adresse du routeur.} & \sol{Joint son réseau, mais aucun autre (serveur).} & \sol{Passerelle 192.168.60.254.} \\ \hline
    \hl 2 & \sol{PC-C} & \sol{IP hors du réseau 192.168.60.0/24.} & \sol{Ne joint ni les postes ni le routeur.} & \sol{Ex. 192.168.60.30.} \\ \hline
    \hl 3 & \sol{PC-D} & \sol{Même IP que PC-A (doublon).} & \sol{Conflit d'adresse : PC-A et PC-D perturbés.} & \sol{Adresse unique, ex. .60.40.} \\ \hline
    \hl 4 & \sol{R1 G0/1} & \sol{Interface désactivée (Off).} & \sol{Réseau du serveur injoignable pour tous.} & \sol{Port Status : On.} \\ \hline}
  \def\Cbareme{PC-A et SRV-C sont correctement configurés ; accepter PC-A cité à la place de PC-D pour le doublon.}
\else
  \def\entreprise{La clinique vétérinaire \textbf{VetoPlus}}
  \def\Rnom{R-VetoPlus}
  \def\resUn{Soins}\def\pUn{SOI}\def\netUn{192.168.70}
  \def\resDeux{Secrétariat}\def\pDeux{SEC}\def\netDeux{172.20}
  \def\srvNom{SRV-Dossiers}\def\netSrv{192.168.250}\def\srvH{20}
  \def\ssid{VetoPlus-WiFi}
  \def\AIrows{%
    \hl\ip{192.168.15.130} & \sol{C} & \sol{255.255.255.0} & \sol{192.168.15.0} & \sol{192.168.15.255} & \sol{254} \\ \hline
    \hl\ip{172.31.200.7}   & \sol{B} & \sol{255.255.0.0}   & \sol{172.31.0.0}   & \sol{172.31.255.255} & \sol{65\,534} \\ \hline
    \hl\ip{10.45.6.9}      & \sol{A} & \sol{255.0.0.0}     & \sol{10.0.0.0}     & \sol{10.255.255.255} & \sol{16\,777\,214} \\ \hline}
  \def\apipa{169.254.88.3}
  \def\Cnet{192.168.80}
  \def\Crows{%
    PC-A & \ip{192.168.8.11}   & \ip{255.255.255.0} & \ip{192.168.80.254} & — \\
    PC-B & \ip{192.168.80.12}  & \ip{255.255.255.0} & \ip{192.168.80.254} & — \\
    PC-C & \ip{192.168.80.13}  & \ip{255.255.255.0} & \ip{192.168.80.100} & — \\
    PC-D & \ip{192.168.80.12}  & \ip{255.255.255.0} & \ip{192.168.80.254} & — \\
    SRV-C & \ip{10.0.0.8}      & \ip{255.0.0.0}     & \ip{10.0.0.254}     & — \\
    R1 — G0/0 & \ip{192.168.80.254} & \ip{255.255.255.0} & — & Off \\
    R1 — G0/1 & \ip{10.0.0.254}     & \ip{255.0.0.0}     & — & On \\}
  \def\Csol{%
    \hl 1 & \sol{PC-A} & \sol{IP hors du réseau 192.168.80.0/24.} & \sol{Ne joint ni les postes ni le routeur.} & \sol{Ex. 192.168.80.11.} \\ \hline
    \hl 2 & \sol{PC-C} & \sol{Passerelle .80.100 : pas l'adresse du routeur.} & \sol{Joint son réseau, mais aucun autre (serveur).} & \sol{Passerelle 192.168.80.254.} \\ \hline
    \hl 3 & \sol{PC-D} & \sol{Même IP que PC-B (doublon).} & \sol{Conflit d'adresse : PC-B et PC-D perturbés.} & \sol{Adresse unique, ex. .80.14.} \\ \hline
    \hl 4 & \sol{R1 G0/0} & \sol{Interface côté LAN désactivée (Off).} & \sol{Aucun poste ne joint le routeur ni le serveur.} & \sol{Port Status : On.} \\ \hline}
  \def\Cbareme{PC-B et SRV-C sont correctement configurés ; accepter PC-B cité à la place de PC-D pour le doublon.}
\fi

\pagestyle{fancy}
\fancyhf{}
\renewcommand{\headrulewidth}{0pt}
\fancyfoot[L]{\footnotesize\itshape\color{grisfoot}Évaluation — Réseaux locaux — Sujet \variante\ifcorrige\ — CORRIGÉ\fi}
\fancyfoot[R]{\footnotesize\color{grisfoot}\thepage/3}

\setlist{itemsep=0pt,topsep=2pt,parsep=0pt}
\setlength{\parindent}{0pt}
\setlength{\parskip}{2pt}

\begin{document}

\begin{center}
  {\Large\bfseries Évaluation pratique — Réseaux locaux \quad
   \fcolorbox{black}{white}{\ SUJET \variante\ }\par}
  \vspace{1pt}
  {\bfseries Adressage IP, commutation et routage — BTS CIEL 1re année — durée : 2 heures\par}
  \ifcorrige\vspace{2pt}{\bfseries\color{corr}CORRIGÉ ET BARÈME — document enseignant}\par\fi
\end{center}
\vspace{-2pt}
{\renewcommand{\arraystretch}{1.35}
\begin{tabularx}{\linewidth}{|X|X|p{2.2cm}|p{2.6cm}|}
  \hline
  \textbf{Nom :} & \textbf{Prénom :} & \textbf{Poste n°} & \textbf{Note :} \hfill /20 \\ \hline
\end{tabularx}}

\textbf{Consignes :} travail \textbf{individuel}, documents et Internet \textbf{non autorisés}. Ce sujet est le
\textbf{document-réponse}. Fichier Packet Tracer : \ip{Nom\_Prenom\_EvalLAN\_\variante.pkt}, à enregistrer
régulièrement dans le dossier de rendu. Les « visas » sont apposés par le professeur, que vous appelez à chaque test réussi.
Durées conseillées : A 30 min — B 1 h 10 — C 20 min.

%==========================================================================
\section{Partie A — Questions de cours et calculs \hfill /6}

\question{A1}{2 pts}{Complétez le tableau en considérant le \textbf{masque par défaut} de la classe de chaque adresse.}
{\small\renewcommand{\arraystretch}{1.15}
\begin{tabularx}{\linewidth}{|l|c|C|C|C|c|}
  \hline
  \textbf{Adresse IP} & \textbf{Classe} & \textbf{Masque par défaut} & \textbf{Adresse réseau} &
  \textbf{Adresse de diffusion} & \textbf{Nb d'hôtes} \\ \hline
  \AIrows
\end{tabularx}}
\bareme{0,5 pt par ligne juste + 0,5 pt si les trois nombres d'hôtes sont justes ($2^{n}-2$, $n$ = nombre de bits d'hôte).}

\question{A2}{1 pt}{Indiquez le type de câble (\textbf{droit} ou \textbf{croisé}) pour chaque liaison.}
{\small\renewcommand{\arraystretch}{1.15}
\begin{tabularx}{\linewidth}{|X|c||X|c|}
  \hline
  \hl PC $\leftrightarrow$ commutateur & \makebox[1.6cm]{\sol{droit}} & Commutateur $\leftrightarrow$ routeur & \makebox[1.6cm]{\sol{droit}} \\ \hline
  \hl PC $\leftrightarrow$ PC (réseau P2P) & \sol{croisé} & Concentrateur $\leftrightarrow$ concentrateur & \sol{croisé} \\ \hline
\end{tabularx}}
\bareme{0,25 pt par case. Équipements de types différents : câble droit ; de même type : câble croisé.}

\question{A3}{1 pt}{Comment un \textbf{concentrateur (hub)} et un \textbf{commutateur (switch)} traitent-ils une trame
reçue ? Lequel limite les collisions, et pourquoi ?}
\zone{2.9cm}{Le \textbf{hub} (couche 1) répète la trame sur \textbf{tous ses autres ports} : un seul domaine de collision
partagé. Le \textbf{switch} (couche 2) lit l'adresse MAC de destination et, grâce à sa \textbf{table MAC}, ne l'envoie
\textbf{que sur le port du destinataire} : un domaine de collision par port, c'est lui qui limite les collisions.
\emph{0,5 pt hub, 0,5 pt switch.}}

\question{A4}{1 pt}{Qu'est-ce que la \textbf{passerelle par défaut} d'un poste ? Quelle condition son adresse doit-elle respecter ?}
\zone{2.4cm}{Adresse IP de l'interface du \textbf{routeur} à laquelle le poste envoie les paquets destinés à un
\textbf{autre réseau}. Elle doit \textbf{appartenir au même réseau} que le poste. \emph{0,5 + 0,5 pt.}}

\question{A5}{0,5 pt}{Un portable configuré en DHCP affiche \ip{\apipa} / \ip{255.255.0.0}. Que s'est-il passé ?}
\zone{1.5cm}{Aucun \textbf{serveur DHCP} n'a répondu : le poste s'est attribué une adresse \textbf{APIPA} (169.254.0.0/16).}

\question{A6}{0,5 pt}{Quel protocole utilise \texttt{ping} ? Nommez les deux messages échangés et la couche OSI des adresses IP.}
\zone{1.5cm}{\textbf{ICMP} ; \textbf{Echo Request} / \textbf{Echo Reply} ; adresses IP en \textbf{couche 3 (Réseau)}.}

%==========================================================================
\newpage
\section{Partie B — Réalisation sous Packet Tracer \hfill /10}

\begin{cdc}
\entreprise{} vous confie son réseau : \textbf{trois réseaux} reliés par un \textbf{routeur unique Cisco 2911}
nommé \ip{\Rnom} (interfaces G0/0, G0/1, G0/2).
\begin{itemize}[leftmargin=12pt]
  \item \textbf{Réseau \resUn} : \textbf{4 postes} PC-\pUn1 à PC-\pUn4, commutateur \ip{SW-\pUn}, adressage statique,
        réseau \ip{\netUn.0/24}.
  \item \textbf{Réseau \resDeux} : \textbf{2 postes} PC-\pDeux1, PC-\pDeux2 et \textbf{1 imprimante} IMP-\pDeux,
        commutateur \ip{SW-\pDeux}, adressage statique, réseau de classe B \ip{\netDeux.0.0/16}.
  \item \textbf{Réseau Serveur} : serveur \ip{\srvNom} d'adresse \ip{\netSrv.\srvH/24}, relié \textbf{directement} au routeur.
  \item La passerelle de chaque réseau est sa \textbf{dernière adresse utilisable}. Tous les équipements communiquent
        entre eux et accèdent à la page web du serveur. Nommage imposé (« Display Name »).
\end{itemize}
\end{cdc}

\question{B1}{2 pts}{Complétez le plan d'adressage \textbf{avant} la réalisation.}
{\footnotesize\setlength{\tabcolsep}{3pt}\renewcommand{\arraystretch}{1.15}
\begin{tabularx}{\linewidth}{|l|C|C|C|C|C|C|}
  \hline
  \textbf{Réseau} & \textbf{Adresse réseau} & \textbf{Masque} & \textbf{1re @ utilisable} &
  \textbf{Dernière @ utilisable} & \textbf{Diffusion} & \textbf{Passerelle} \\ \hline
  \hl \resUn & \sol{\netUn.0} & \sol{255.255.255.0} & \sol{\netUn.1} & \sol{\netUn.254} & \sol{\netUn.255} & \sol{\netUn.254} \\ \hline
  \hl \resDeux & \sol{\netDeux.0.0} & \sol{255.255.0.0} & \sol{\netDeux.0.1} & \sol{\netDeux.255.254} & \sol{\netDeux.255.255} & \sol{\netDeux.255.254} \\ \hline
  \hl Serveur & \sol{\netSrv.0} & \sol{255.255.255.0} & \sol{\netSrv.1} & \sol{\netSrv.254} & \sol{\netSrv.255} & \sol{\netSrv.254} \\ \hline
\end{tabularx}}
\bareme{0,5 pt par ligne juste + 0,5 pt si la convention de passerelle est respectée partout.}

\question{B2}{2 pts}{Construisez la topologie (équipements, câblage, nommage) et relevez l'adressage configuré.}
{\footnotesize\setlength{\tabcolsep}{3pt}\renewcommand{\arraystretch}{1.1}
\begin{tabularx}{\linewidth}{|l|C|C|C||l|C|C|}
  \hline
  \textbf{Hôte} & \textbf{Adresse IP} & \textbf{Masque} & \textbf{Passerelle} & \textbf{Interface} & \textbf{Adresse IP} & \textbf{Masque} \\ \hline
  \hl PC-\pUn1 & \sol{\netUn.1} & \sol{255.255.255.0} & \sol{\netUn.254} & \Rnom{} G0/0 & \sol{\netUn.254} & \sol{255.255.255.0} \\ \hline
  \hl PC-\pDeux1 & \sol{\netDeux.0.1} & \sol{255.255.0.0} & \sol{\netDeux.255.254} & \Rnom{} G0/1 & \sol{\netDeux.255.254} & \sol{255.255.0.0} \\ \hline
  \hl IMP-\pDeux & \sol{\netDeux.0.10} & \sol{255.255.0.0} & \sol{\netDeux.255.254} & \Rnom{} G0/2 & \sol{\netSrv.254} & \sol{255.255.255.0} \\ \hline
  \hl \srvNom & \ip{\netSrv.\srvH} & \sol{255.255.255.0} & \sol{\netSrv.254} & \multicolumn{3}{c|}{\cellcolor{black!8}} \\ \hline
\end{tabularx}}
\bareme{Adresses d'hôtes libres (uniques, dans le bon réseau) ; affectation G0/x libre. Sur le .pkt : équipements 0,5 ;
câblage 1 (câbles droits vers les switchs et switch–routeur, \textbf{croisé} serveur–routeur, liens verts) ; nommage 0,5.}

\question{B3}{3 pts}{Configurez tous les équipements : IP, masques, passerelles des hôtes ; interfaces du routeur
(« Port Status : On »).}
\bareme{Sur le .pkt : interfaces routeur 1 pt ; IP/masques des hôtes 1 pt ; passerelles (imprimante et serveur compris)
1 pt. $-0,25$ par erreur, sans descendre sous 0 par item.}

\question{B4}{3 pts}{Réalisez les tests, notez la commande et le résultat, puis faites viser chaque test réussi.}
{\small\renewcommand{\arraystretch}{1.1}
\begin{tabularx}{\linewidth}{|c|l|l|C|c|c|}
  \hline
  \textbf{N°} & \textbf{Source} & \textbf{Destination} & \textbf{Commande / action} & \textbf{Résultat} & \textbf{Visa prof.} \\ \hline
  \hl 1 & PC-\pUn1 & PC-\pUn4 & \sol{ping \netUn.4} & \sol{réussi} & \hspace{1.4cm} \\ \hline
  \hl 2 & PC-\pDeux1 & IMP-\pDeux & \sol{ping \netDeux.0.10} & \sol{réussi} & \\ \hline
  \hl 3 & PC-\pUn2 & PC-\pDeux2 & \sol{ping \netDeux.0.2} & \sol{réussi} & \\ \hline
  \hl 4 & PC-\pDeux2 & \srvNom & \sol{ping \netSrv.\srvH} & \sol{réussi} & \\ \hline
  \hl 5 & PC-\pUn3 & page web du serveur & \sol{navigateur : {\NoAutoSpacing http://\netSrv.\srvH}} & \sol{page affichée} & \\ \hline
\end{tabularx}}
\bareme{0,5 pt par test visé + 0,5 pt pour B5. Le premier ping inter-réseaux peut échouer (résolution ARP) ; les suivants doivent réussir.}

\question{B5}{incl. B4}{Test n°3 : par quels équipements passe le paquet ? Pourquoi PC-\pUn2 l'envoie-t-il au routeur ?}
\zone{2.3cm}{PC-\pUn2 $\to$ SW-\pUn{} $\to$ \Rnom{} $\to$ SW-\pDeux{} $\to$ PC-\pDeux2. Grâce à son masque, PC-\pUn2 constate que
la destination (\netDeux.0.0/16) n'est \textbf{pas dans son réseau} (\netUn.0/24) : il envoie le paquet à sa
\textbf{passerelle par défaut}, qui le route.}

%==========================================================================
\newpage
\section{Partie C — Diagnostic d'un réseau en panne \hfill /4}

Un technicien constate qu'\textbf{aucun poste ne parvient à joindre le serveur SRV-C} et que certains postes
communiquent mal entre eux. Les câbles sont du bon type.

\begin{center}
\begin{minipage}{0.37\linewidth}\centering
\begin{tikzpicture}[font=\footnotesize,
  eq/.style={draw,fill=qbg,rounded corners=2pt,minimum width=9mm,minimum height=5mm,inner sep=2pt},
  every path/.style={thick}]
  \node[eq] (SW) {SW-1};
  \node[eq,right=9mm of SW] (R) {R1};
  \node[eq,right=7mm of R] (SRV) {SRV-C};
  \node[eq,above left=5mm and -3mm of SW] (A) {PC-A};
  \node[eq,above right=5mm and -3mm of SW] (B) {PC-B};
  \node[eq,below left=5mm and -3mm of SW] (C) {PC-C};
  \node[eq,below right=5mm and -3mm of SW] (D) {PC-D};
  \foreach \n in {A,B,C,D} \draw (\n) -- (SW);
  \draw (SW) -- node[above,font=\tiny]{G0/0} (R);
  \draw (R) -- node[above,font=\tiny]{G0/1} (SRV);
\end{tikzpicture}
\end{minipage}\hfill
\begin{minipage}{0.61\linewidth}\centering
{\small\renewcommand{\arraystretch}{1.05}\setlength{\tabcolsep}{4pt}
\begin{tabular}{lllll}
  \toprule
  \textbf{Équipement} & \textbf{Adresse IP} & \textbf{Masque} & \textbf{Passerelle} & \textbf{Port} \\
  \midrule
  \Crows
  \bottomrule
\end{tabular}}
\end{minipage}
\end{center}

\question{C1}{4 pts}{Relevez les \textbf{quatre anomalies}. Pour chacune : conséquence et correction.}
{\small\renewcommand{\arraystretch}{1.1}
\ifcorrige\def\hc{}\else\def\hc{\rule[-10mm]{0pt}{14mm}}\fi
\begin{tabularx}{\linewidth}{|c|>{\raggedright\arraybackslash}p{2cm}|L|L|L|}
  \hline
  \textbf{N°} & \textbf{Équipement} & \textbf{Anomalie} & \textbf{Conséquence} & \textbf{Correction} \\ \hline
  \ifcorrige\Csol\else
  \hc 1 & & & & \\ \hline
  \hc 2 & & & & \\ \hline
  \hc 3 & & & & \\ \hline
  \hc 4 & & & & \\ \hline
  \fi
\end{tabularx}}
\bareme{1 pt par anomalie (0,25 équipement, anomalie, conséquence, correction ; ordre indifférent). \Cbareme}

%==========================================================================
\section{Bonus — Extension sans fil \hfill +1}

\begin{cdc}Ajoutez au réseau \resUn{} un \textbf{AccessPoint-PT} relié à SW-\pUn{} (SSID \ip{\ssid}, \textbf{WPA2-PSK})
et un portable \ip{Laptop-\pUn} équipé d'un module \texttt{WPC300N}, en adressage \textbf{statique}.
Vérifiez qu'il joint \srvNom.\end{cdc}
{\small\renewcommand{\arraystretch}{1.1}
\begin{tabularx}{\linewidth}{|l|C|C|C|c|}
  \hline
  \textbf{Équipement} & \textbf{Adresse IP} & \textbf{Masque} & \textbf{Passerelle} & \textbf{Visa prof.} \\ \hline
  \hl Laptop-\pUn & \sol{\netUn.50} & \sol{255.255.255.0} & \sol{\netUn.254} & \hspace{1.4cm} \\ \hline
\end{tabularx}}
\bareme{Le point d'accès fait office de pont : le portable est dans \netUn.0/24. Bonus si associé en WPA2-PSK et ping
vers \netSrv.\srvH{} réussi. Note plafonnée à 20.}

\end{document}
"
%   pdflatex -jobname=Eval_Reseau_CorrigeA "\def\variante{A}\def\modecorrige{1}% Évaluation pratique — Réseaux locaux (BTS CIEL 1re année) — sujets A et B
% Compilation (4 PDF) :
%   pdflatex -jobname=Eval_Reseau_SujetA   "\def\variante{A}\def\modecorrige{0}\input{Eval_TP_Reseau_LAN.tex}"
%   pdflatex -jobname=Eval_Reseau_SujetB   "\def\variante{B}\def\modecorrige{0}\input{Eval_TP_Reseau_LAN.tex}"
%   pdflatex -jobname=Eval_Reseau_CorrigeA "\def\variante{A}\def\modecorrige{1}\input{Eval_TP_Reseau_LAN.tex}"
%   pdflatex -jobname=Eval_Reseau_CorrigeB "\def\variante{B}\def\modecorrige{1}\input{Eval_TP_Reseau_LAN.tex}"
\documentclass[10pt,a4paper]{article}

\usepackage[T1]{fontenc}
\usepackage[utf8]{inputenc}
\usepackage{lmodern}
\usepackage[french]{babel}
\usepackage[a4paper,hmargin=1.5cm,top=1.3cm,bottom=1.6cm,footskip=0.7cm]{geometry}
\usepackage[table]{xcolor}
\usepackage{booktabs,tabularx,array}
\usepackage{enumitem}
\usepackage{fancyhdr}
\usepackage[most]{tcolorbox}
\usepackage{tikz}
\usetikzlibrary{positioning}
\usepackage{titlesec}

\providecommand{\modecorrige}{0}
\providecommand{\variante}{A}
\newif\ifcorrige
\ifnum\modecorrige=1 \corrigetrue \else \corrigefalse \fi

\definecolor{ipred}{HTML}{C0392B}
\definecolor{cdcbar}{HTML}{A67C3D}
\definecolor{cdcbg}{HTML}{F4F1EA}
\definecolor{qbar}{HTML}{4A6B5D}
\definecolor{qbg}{HTML}{EDF1EF}
\definecolor{corr}{HTML}{1F5FA8}
\definecolor{corrbg}{HTML}{E8F0FA}
\definecolor{grisfoot}{HTML}{828282}

\newcommand{\ip}[1]{\textcolor{ipred}{\texttt{#1}}}
\newcommand{\pts}[1]{\hfill\mbox{\textbf{(#1)}}}
\newcommand{\sol}[1]{\ifcorrige\textcolor{corr}{#1}\fi}
\newcommand{\bareme}[1]{\ifcorrige\par{\footnotesize\color{corr}#1\par}\fi}
\newcommand{\hl}{\ifcorrige\rule[-1.6mm]{0pt}{5.4mm}\else\rule[-2.2mm]{0pt}{6.6mm}\fi}
\newcolumntype{C}{>{\centering\arraybackslash}X}
\newcolumntype{L}{>{\raggedright\arraybackslash}X}

\titleformat{\section}{\large\bfseries}{}{0pt}{}
\titlespacing*{\section}{0pt}{6pt}{3pt}

\tcbset{qstyle/.style={enhanced,boxrule=0pt,leftrule=2.5pt,arc=0pt,outer arc=0pt,
  left=5pt,right=5pt,top=2pt,bottom=2pt,before skip=4pt,after skip=3pt}}
\newtcolorbox{cdc}{qstyle,colback=cdcbg,colframe=cdcbar}
\newtcolorbox{quest}{qstyle,colback=qbg,colframe=qbar}

% Zone de réponse : vierge dans le sujet, remplie dans le corrigé
\newcommand{\zone}[2]{%
  \ifcorrige
    \begin{tcolorbox}[colback=corrbg,colframe=corr,boxrule=0.5pt,arc=1pt,left=4pt,right=4pt,top=1pt,bottom=1pt,
      before skip=2pt,after skip=2pt,fontupper=\footnotesize\color{corr}]#2\end{tcolorbox}%
  \else
    \begin{tcolorbox}[colback=white,colframe=black!45,boxrule=0.4pt,arc=1pt,height=#1,
      before skip=2pt,after skip=2pt]\end{tcolorbox}%
  \fi}

\newcommand{\question}[3]{\begin{quest}\textbf{#1.} #3\pts{#2}\end{quest}}

%==========================================================================
% Données propres à chaque variante
%==========================================================================
\if\variante A
  \def\entreprise{Le cabinet d'architectes \textbf{ArchiNova}}
  \def\Rnom{R-ArchiNova}
  \def\resUn{Bureau d'études}\def\pUn{BE}\def\netUn{192.168.40}
  \def\resDeux{Accueil}\def\pDeux{ACC}\def\netDeux{172.16}
  \def\srvNom{SRV-Intranet}\def\netSrv{192.168.200}\def\srvH{10}
  \def\ssid{ArchiNova-WiFi}
  \def\AIrows{%
    \hl\ip{192.168.40.25} & \sol{C} & \sol{255.255.255.0} & \sol{192.168.40.0} & \sol{192.168.40.255} & \sol{254} \\ \hline
    \hl\ip{172.16.5.10}   & \sol{B} & \sol{255.255.0.0}   & \sol{172.16.0.0}   & \sol{172.16.255.255} & \sol{65\,534} \\ \hline
    \hl\ip{10.3.2.1}      & \sol{A} & \sol{255.0.0.0}     & \sol{10.0.0.0}     & \sol{10.255.255.255} & \sol{16\,777\,214} \\ \hline}
  \def\apipa{169.254.12.7}
  \def\Cnet{192.168.60}
  \def\Crows{%
    PC-A & \ip{192.168.60.10} & \ip{255.255.255.0} & \ip{192.168.60.254} & — \\
    PC-B & \ip{192.168.60.20} & \ip{255.255.255.0} & \ip{192.168.60.1}   & — \\
    PC-C & \ip{192.168.61.30} & \ip{255.255.255.0} & \ip{192.168.60.254} & — \\
    PC-D & \ip{192.168.60.10} & \ip{255.255.255.0} & \ip{192.168.60.254} & — \\
    SRV-C & \ip{10.0.0.5}     & \ip{255.0.0.0}     & \ip{10.0.0.254}     & — \\
    R1 — G0/0 & \ip{192.168.60.254} & \ip{255.255.255.0} & — & On \\
    R1 — G0/1 & \ip{10.0.0.254}     & \ip{255.0.0.0}     & — & Off \\}
  \def\Csol{%
    \hl 1 & \sol{PC-B} & \sol{Passerelle .60.1 : pas l'adresse du routeur.} & \sol{Joint son réseau, mais aucun autre (serveur).} & \sol{Passerelle 192.168.60.254.} \\ \hline
    \hl 2 & \sol{PC-C} & \sol{IP hors du réseau 192.168.60.0/24.} & \sol{Ne joint ni les postes ni le routeur.} & \sol{Ex. 192.168.60.30.} \\ \hline
    \hl 3 & \sol{PC-D} & \sol{Même IP que PC-A (doublon).} & \sol{Conflit d'adresse : PC-A et PC-D perturbés.} & \sol{Adresse unique, ex. .60.40.} \\ \hline
    \hl 4 & \sol{R1 G0/1} & \sol{Interface désactivée (Off).} & \sol{Réseau du serveur injoignable pour tous.} & \sol{Port Status : On.} \\ \hline}
  \def\Cbareme{PC-A et SRV-C sont correctement configurés ; accepter PC-A cité à la place de PC-D pour le doublon.}
\else
  \def\entreprise{La clinique vétérinaire \textbf{VetoPlus}}
  \def\Rnom{R-VetoPlus}
  \def\resUn{Soins}\def\pUn{SOI}\def\netUn{192.168.70}
  \def\resDeux{Secrétariat}\def\pDeux{SEC}\def\netDeux{172.20}
  \def\srvNom{SRV-Dossiers}\def\netSrv{192.168.250}\def\srvH{20}
  \def\ssid{VetoPlus-WiFi}
  \def\AIrows{%
    \hl\ip{192.168.15.130} & \sol{C} & \sol{255.255.255.0} & \sol{192.168.15.0} & \sol{192.168.15.255} & \sol{254} \\ \hline
    \hl\ip{172.31.200.7}   & \sol{B} & \sol{255.255.0.0}   & \sol{172.31.0.0}   & \sol{172.31.255.255} & \sol{65\,534} \\ \hline
    \hl\ip{10.45.6.9}      & \sol{A} & \sol{255.0.0.0}     & \sol{10.0.0.0}     & \sol{10.255.255.255} & \sol{16\,777\,214} \\ \hline}
  \def\apipa{169.254.88.3}
  \def\Cnet{192.168.80}
  \def\Crows{%
    PC-A & \ip{192.168.8.11}   & \ip{255.255.255.0} & \ip{192.168.80.254} & — \\
    PC-B & \ip{192.168.80.12}  & \ip{255.255.255.0} & \ip{192.168.80.254} & — \\
    PC-C & \ip{192.168.80.13}  & \ip{255.255.255.0} & \ip{192.168.80.100} & — \\
    PC-D & \ip{192.168.80.12}  & \ip{255.255.255.0} & \ip{192.168.80.254} & — \\
    SRV-C & \ip{10.0.0.8}      & \ip{255.0.0.0}     & \ip{10.0.0.254}     & — \\
    R1 — G0/0 & \ip{192.168.80.254} & \ip{255.255.255.0} & — & Off \\
    R1 — G0/1 & \ip{10.0.0.254}     & \ip{255.0.0.0}     & — & On \\}
  \def\Csol{%
    \hl 1 & \sol{PC-A} & \sol{IP hors du réseau 192.168.80.0/24.} & \sol{Ne joint ni les postes ni le routeur.} & \sol{Ex. 192.168.80.11.} \\ \hline
    \hl 2 & \sol{PC-C} & \sol{Passerelle .80.100 : pas l'adresse du routeur.} & \sol{Joint son réseau, mais aucun autre (serveur).} & \sol{Passerelle 192.168.80.254.} \\ \hline
    \hl 3 & \sol{PC-D} & \sol{Même IP que PC-B (doublon).} & \sol{Conflit d'adresse : PC-B et PC-D perturbés.} & \sol{Adresse unique, ex. .80.14.} \\ \hline
    \hl 4 & \sol{R1 G0/0} & \sol{Interface côté LAN désactivée (Off).} & \sol{Aucun poste ne joint le routeur ni le serveur.} & \sol{Port Status : On.} \\ \hline}
  \def\Cbareme{PC-B et SRV-C sont correctement configurés ; accepter PC-B cité à la place de PC-D pour le doublon.}
\fi

\pagestyle{fancy}
\fancyhf{}
\renewcommand{\headrulewidth}{0pt}
\fancyfoot[L]{\footnotesize\itshape\color{grisfoot}Évaluation — Réseaux locaux — Sujet \variante\ifcorrige\ — CORRIGÉ\fi}
\fancyfoot[R]{\footnotesize\color{grisfoot}\thepage/3}

\setlist{itemsep=0pt,topsep=2pt,parsep=0pt}
\setlength{\parindent}{0pt}
\setlength{\parskip}{2pt}

\begin{document}

\begin{center}
  {\Large\bfseries Évaluation pratique — Réseaux locaux \quad
   \fcolorbox{black}{white}{\ SUJET \variante\ }\par}
  \vspace{1pt}
  {\bfseries Adressage IP, commutation et routage — BTS CIEL 1re année — durée : 2 heures\par}
  \ifcorrige\vspace{2pt}{\bfseries\color{corr}CORRIGÉ ET BARÈME — document enseignant}\par\fi
\end{center}
\vspace{-2pt}
{\renewcommand{\arraystretch}{1.35}
\begin{tabularx}{\linewidth}{|X|X|p{2.2cm}|p{2.6cm}|}
  \hline
  \textbf{Nom :} & \textbf{Prénom :} & \textbf{Poste n°} & \textbf{Note :} \hfill /20 \\ \hline
\end{tabularx}}

\textbf{Consignes :} travail \textbf{individuel}, documents et Internet \textbf{non autorisés}. Ce sujet est le
\textbf{document-réponse}. Fichier Packet Tracer : \ip{Nom\_Prenom\_EvalLAN\_\variante.pkt}, à enregistrer
régulièrement dans le dossier de rendu. Les « visas » sont apposés par le professeur, que vous appelez à chaque test réussi.
Durées conseillées : A 30 min — B 1 h 10 — C 20 min.

%==========================================================================
\section{Partie A — Questions de cours et calculs \hfill /6}

\question{A1}{2 pts}{Complétez le tableau en considérant le \textbf{masque par défaut} de la classe de chaque adresse.}
{\small\renewcommand{\arraystretch}{1.15}
\begin{tabularx}{\linewidth}{|l|c|C|C|C|c|}
  \hline
  \textbf{Adresse IP} & \textbf{Classe} & \textbf{Masque par défaut} & \textbf{Adresse réseau} &
  \textbf{Adresse de diffusion} & \textbf{Nb d'hôtes} \\ \hline
  \AIrows
\end{tabularx}}
\bareme{0,5 pt par ligne juste + 0,5 pt si les trois nombres d'hôtes sont justes ($2^{n}-2$, $n$ = nombre de bits d'hôte).}

\question{A2}{1 pt}{Indiquez le type de câble (\textbf{droit} ou \textbf{croisé}) pour chaque liaison.}
{\small\renewcommand{\arraystretch}{1.15}
\begin{tabularx}{\linewidth}{|X|c||X|c|}
  \hline
  \hl PC $\leftrightarrow$ commutateur & \makebox[1.6cm]{\sol{droit}} & Commutateur $\leftrightarrow$ routeur & \makebox[1.6cm]{\sol{droit}} \\ \hline
  \hl PC $\leftrightarrow$ PC (réseau P2P) & \sol{croisé} & Concentrateur $\leftrightarrow$ concentrateur & \sol{croisé} \\ \hline
\end{tabularx}}
\bareme{0,25 pt par case. Équipements de types différents : câble droit ; de même type : câble croisé.}

\question{A3}{1 pt}{Comment un \textbf{concentrateur (hub)} et un \textbf{commutateur (switch)} traitent-ils une trame
reçue ? Lequel limite les collisions, et pourquoi ?}
\zone{2.9cm}{Le \textbf{hub} (couche 1) répète la trame sur \textbf{tous ses autres ports} : un seul domaine de collision
partagé. Le \textbf{switch} (couche 2) lit l'adresse MAC de destination et, grâce à sa \textbf{table MAC}, ne l'envoie
\textbf{que sur le port du destinataire} : un domaine de collision par port, c'est lui qui limite les collisions.
\emph{0,5 pt hub, 0,5 pt switch.}}

\question{A4}{1 pt}{Qu'est-ce que la \textbf{passerelle par défaut} d'un poste ? Quelle condition son adresse doit-elle respecter ?}
\zone{2.4cm}{Adresse IP de l'interface du \textbf{routeur} à laquelle le poste envoie les paquets destinés à un
\textbf{autre réseau}. Elle doit \textbf{appartenir au même réseau} que le poste. \emph{0,5 + 0,5 pt.}}

\question{A5}{0,5 pt}{Un portable configuré en DHCP affiche \ip{\apipa} / \ip{255.255.0.0}. Que s'est-il passé ?}
\zone{1.5cm}{Aucun \textbf{serveur DHCP} n'a répondu : le poste s'est attribué une adresse \textbf{APIPA} (169.254.0.0/16).}

\question{A6}{0,5 pt}{Quel protocole utilise \texttt{ping} ? Nommez les deux messages échangés et la couche OSI des adresses IP.}
\zone{1.5cm}{\textbf{ICMP} ; \textbf{Echo Request} / \textbf{Echo Reply} ; adresses IP en \textbf{couche 3 (Réseau)}.}

%==========================================================================
\newpage
\section{Partie B — Réalisation sous Packet Tracer \hfill /10}

\begin{cdc}
\entreprise{} vous confie son réseau : \textbf{trois réseaux} reliés par un \textbf{routeur unique Cisco 2911}
nommé \ip{\Rnom} (interfaces G0/0, G0/1, G0/2).
\begin{itemize}[leftmargin=12pt]
  \item \textbf{Réseau \resUn} : \textbf{4 postes} PC-\pUn1 à PC-\pUn4, commutateur \ip{SW-\pUn}, adressage statique,
        réseau \ip{\netUn.0/24}.
  \item \textbf{Réseau \resDeux} : \textbf{2 postes} PC-\pDeux1, PC-\pDeux2 et \textbf{1 imprimante} IMP-\pDeux,
        commutateur \ip{SW-\pDeux}, adressage statique, réseau de classe B \ip{\netDeux.0.0/16}.
  \item \textbf{Réseau Serveur} : serveur \ip{\srvNom} d'adresse \ip{\netSrv.\srvH/24}, relié \textbf{directement} au routeur.
  \item La passerelle de chaque réseau est sa \textbf{dernière adresse utilisable}. Tous les équipements communiquent
        entre eux et accèdent à la page web du serveur. Nommage imposé (« Display Name »).
\end{itemize}
\end{cdc}

\question{B1}{2 pts}{Complétez le plan d'adressage \textbf{avant} la réalisation.}
{\footnotesize\setlength{\tabcolsep}{3pt}\renewcommand{\arraystretch}{1.15}
\begin{tabularx}{\linewidth}{|l|C|C|C|C|C|C|}
  \hline
  \textbf{Réseau} & \textbf{Adresse réseau} & \textbf{Masque} & \textbf{1re @ utilisable} &
  \textbf{Dernière @ utilisable} & \textbf{Diffusion} & \textbf{Passerelle} \\ \hline
  \hl \resUn & \sol{\netUn.0} & \sol{255.255.255.0} & \sol{\netUn.1} & \sol{\netUn.254} & \sol{\netUn.255} & \sol{\netUn.254} \\ \hline
  \hl \resDeux & \sol{\netDeux.0.0} & \sol{255.255.0.0} & \sol{\netDeux.0.1} & \sol{\netDeux.255.254} & \sol{\netDeux.255.255} & \sol{\netDeux.255.254} \\ \hline
  \hl Serveur & \sol{\netSrv.0} & \sol{255.255.255.0} & \sol{\netSrv.1} & \sol{\netSrv.254} & \sol{\netSrv.255} & \sol{\netSrv.254} \\ \hline
\end{tabularx}}
\bareme{0,5 pt par ligne juste + 0,5 pt si la convention de passerelle est respectée partout.}

\question{B2}{2 pts}{Construisez la topologie (équipements, câblage, nommage) et relevez l'adressage configuré.}
{\footnotesize\setlength{\tabcolsep}{3pt}\renewcommand{\arraystretch}{1.1}
\begin{tabularx}{\linewidth}{|l|C|C|C||l|C|C|}
  \hline
  \textbf{Hôte} & \textbf{Adresse IP} & \textbf{Masque} & \textbf{Passerelle} & \textbf{Interface} & \textbf{Adresse IP} & \textbf{Masque} \\ \hline
  \hl PC-\pUn1 & \sol{\netUn.1} & \sol{255.255.255.0} & \sol{\netUn.254} & \Rnom{} G0/0 & \sol{\netUn.254} & \sol{255.255.255.0} \\ \hline
  \hl PC-\pDeux1 & \sol{\netDeux.0.1} & \sol{255.255.0.0} & \sol{\netDeux.255.254} & \Rnom{} G0/1 & \sol{\netDeux.255.254} & \sol{255.255.0.0} \\ \hline
  \hl IMP-\pDeux & \sol{\netDeux.0.10} & \sol{255.255.0.0} & \sol{\netDeux.255.254} & \Rnom{} G0/2 & \sol{\netSrv.254} & \sol{255.255.255.0} \\ \hline
  \hl \srvNom & \ip{\netSrv.\srvH} & \sol{255.255.255.0} & \sol{\netSrv.254} & \multicolumn{3}{c|}{\cellcolor{black!8}} \\ \hline
\end{tabularx}}
\bareme{Adresses d'hôtes libres (uniques, dans le bon réseau) ; affectation G0/x libre. Sur le .pkt : équipements 0,5 ;
câblage 1 (câbles droits vers les switchs et switch–routeur, \textbf{croisé} serveur–routeur, liens verts) ; nommage 0,5.}

\question{B3}{3 pts}{Configurez tous les équipements : IP, masques, passerelles des hôtes ; interfaces du routeur
(« Port Status : On »).}
\bareme{Sur le .pkt : interfaces routeur 1 pt ; IP/masques des hôtes 1 pt ; passerelles (imprimante et serveur compris)
1 pt. $-0,25$ par erreur, sans descendre sous 0 par item.}

\question{B4}{3 pts}{Réalisez les tests, notez la commande et le résultat, puis faites viser chaque test réussi.}
{\small\renewcommand{\arraystretch}{1.1}
\begin{tabularx}{\linewidth}{|c|l|l|C|c|c|}
  \hline
  \textbf{N°} & \textbf{Source} & \textbf{Destination} & \textbf{Commande / action} & \textbf{Résultat} & \textbf{Visa prof.} \\ \hline
  \hl 1 & PC-\pUn1 & PC-\pUn4 & \sol{ping \netUn.4} & \sol{réussi} & \hspace{1.4cm} \\ \hline
  \hl 2 & PC-\pDeux1 & IMP-\pDeux & \sol{ping \netDeux.0.10} & \sol{réussi} & \\ \hline
  \hl 3 & PC-\pUn2 & PC-\pDeux2 & \sol{ping \netDeux.0.2} & \sol{réussi} & \\ \hline
  \hl 4 & PC-\pDeux2 & \srvNom & \sol{ping \netSrv.\srvH} & \sol{réussi} & \\ \hline
  \hl 5 & PC-\pUn3 & page web du serveur & \sol{navigateur : {\NoAutoSpacing http://\netSrv.\srvH}} & \sol{page affichée} & \\ \hline
\end{tabularx}}
\bareme{0,5 pt par test visé + 0,5 pt pour B5. Le premier ping inter-réseaux peut échouer (résolution ARP) ; les suivants doivent réussir.}

\question{B5}{incl. B4}{Test n°3 : par quels équipements passe le paquet ? Pourquoi PC-\pUn2 l'envoie-t-il au routeur ?}
\zone{2.3cm}{PC-\pUn2 $\to$ SW-\pUn{} $\to$ \Rnom{} $\to$ SW-\pDeux{} $\to$ PC-\pDeux2. Grâce à son masque, PC-\pUn2 constate que
la destination (\netDeux.0.0/16) n'est \textbf{pas dans son réseau} (\netUn.0/24) : il envoie le paquet à sa
\textbf{passerelle par défaut}, qui le route.}

%==========================================================================
\newpage
\section{Partie C — Diagnostic d'un réseau en panne \hfill /4}

Un technicien constate qu'\textbf{aucun poste ne parvient à joindre le serveur SRV-C} et que certains postes
communiquent mal entre eux. Les câbles sont du bon type.

\begin{center}
\begin{minipage}{0.37\linewidth}\centering
\begin{tikzpicture}[font=\footnotesize,
  eq/.style={draw,fill=qbg,rounded corners=2pt,minimum width=9mm,minimum height=5mm,inner sep=2pt},
  every path/.style={thick}]
  \node[eq] (SW) {SW-1};
  \node[eq,right=9mm of SW] (R) {R1};
  \node[eq,right=7mm of R] (SRV) {SRV-C};
  \node[eq,above left=5mm and -3mm of SW] (A) {PC-A};
  \node[eq,above right=5mm and -3mm of SW] (B) {PC-B};
  \node[eq,below left=5mm and -3mm of SW] (C) {PC-C};
  \node[eq,below right=5mm and -3mm of SW] (D) {PC-D};
  \foreach \n in {A,B,C,D} \draw (\n) -- (SW);
  \draw (SW) -- node[above,font=\tiny]{G0/0} (R);
  \draw (R) -- node[above,font=\tiny]{G0/1} (SRV);
\end{tikzpicture}
\end{minipage}\hfill
\begin{minipage}{0.61\linewidth}\centering
{\small\renewcommand{\arraystretch}{1.05}\setlength{\tabcolsep}{4pt}
\begin{tabular}{lllll}
  \toprule
  \textbf{Équipement} & \textbf{Adresse IP} & \textbf{Masque} & \textbf{Passerelle} & \textbf{Port} \\
  \midrule
  \Crows
  \bottomrule
\end{tabular}}
\end{minipage}
\end{center}

\question{C1}{4 pts}{Relevez les \textbf{quatre anomalies}. Pour chacune : conséquence et correction.}
{\small\renewcommand{\arraystretch}{1.1}
\ifcorrige\def\hc{}\else\def\hc{\rule[-10mm]{0pt}{14mm}}\fi
\begin{tabularx}{\linewidth}{|c|>{\raggedright\arraybackslash}p{2cm}|L|L|L|}
  \hline
  \textbf{N°} & \textbf{Équipement} & \textbf{Anomalie} & \textbf{Conséquence} & \textbf{Correction} \\ \hline
  \ifcorrige\Csol\else
  \hc 1 & & & & \\ \hline
  \hc 2 & & & & \\ \hline
  \hc 3 & & & & \\ \hline
  \hc 4 & & & & \\ \hline
  \fi
\end{tabularx}}
\bareme{1 pt par anomalie (0,25 équipement, anomalie, conséquence, correction ; ordre indifférent). \Cbareme}

%==========================================================================
\section{Bonus — Extension sans fil \hfill +1}

\begin{cdc}Ajoutez au réseau \resUn{} un \textbf{AccessPoint-PT} relié à SW-\pUn{} (SSID \ip{\ssid}, \textbf{WPA2-PSK})
et un portable \ip{Laptop-\pUn} équipé d'un module \texttt{WPC300N}, en adressage \textbf{statique}.
Vérifiez qu'il joint \srvNom.\end{cdc}
{\small\renewcommand{\arraystretch}{1.1}
\begin{tabularx}{\linewidth}{|l|C|C|C|c|}
  \hline
  \textbf{Équipement} & \textbf{Adresse IP} & \textbf{Masque} & \textbf{Passerelle} & \textbf{Visa prof.} \\ \hline
  \hl Laptop-\pUn & \sol{\netUn.50} & \sol{255.255.255.0} & \sol{\netUn.254} & \hspace{1.4cm} \\ \hline
\end{tabularx}}
\bareme{Le point d'accès fait office de pont : le portable est dans \netUn.0/24. Bonus si associé en WPA2-PSK et ping
vers \netSrv.\srvH{} réussi. Note plafonnée à 20.}

\end{document}
"
%   pdflatex -jobname=Eval_Reseau_CorrigeB "\def\variante{B}\def\modecorrige{1}% Évaluation pratique — Réseaux locaux (BTS CIEL 1re année) — sujets A et B
% Compilation (4 PDF) :
%   pdflatex -jobname=Eval_Reseau_SujetA   "\def\variante{A}\def\modecorrige{0}\input{Eval_TP_Reseau_LAN.tex}"
%   pdflatex -jobname=Eval_Reseau_SujetB   "\def\variante{B}\def\modecorrige{0}\input{Eval_TP_Reseau_LAN.tex}"
%   pdflatex -jobname=Eval_Reseau_CorrigeA "\def\variante{A}\def\modecorrige{1}\input{Eval_TP_Reseau_LAN.tex}"
%   pdflatex -jobname=Eval_Reseau_CorrigeB "\def\variante{B}\def\modecorrige{1}\input{Eval_TP_Reseau_LAN.tex}"
\documentclass[10pt,a4paper]{article}

\usepackage[T1]{fontenc}
\usepackage[utf8]{inputenc}
\usepackage{lmodern}
\usepackage[french]{babel}
\usepackage[a4paper,hmargin=1.5cm,top=1.3cm,bottom=1.6cm,footskip=0.7cm]{geometry}
\usepackage[table]{xcolor}
\usepackage{booktabs,tabularx,array}
\usepackage{enumitem}
\usepackage{fancyhdr}
\usepackage[most]{tcolorbox}
\usepackage{tikz}
\usetikzlibrary{positioning}
\usepackage{titlesec}

\providecommand{\modecorrige}{0}
\providecommand{\variante}{A}
\newif\ifcorrige
\ifnum\modecorrige=1 \corrigetrue \else \corrigefalse \fi

\definecolor{ipred}{HTML}{C0392B}
\definecolor{cdcbar}{HTML}{A67C3D}
\definecolor{cdcbg}{HTML}{F4F1EA}
\definecolor{qbar}{HTML}{4A6B5D}
\definecolor{qbg}{HTML}{EDF1EF}
\definecolor{corr}{HTML}{1F5FA8}
\definecolor{corrbg}{HTML}{E8F0FA}
\definecolor{grisfoot}{HTML}{828282}

\newcommand{\ip}[1]{\textcolor{ipred}{\texttt{#1}}}
\newcommand{\pts}[1]{\hfill\mbox{\textbf{(#1)}}}
\newcommand{\sol}[1]{\ifcorrige\textcolor{corr}{#1}\fi}
\newcommand{\bareme}[1]{\ifcorrige\par{\footnotesize\color{corr}#1\par}\fi}
\newcommand{\hl}{\ifcorrige\rule[-1.6mm]{0pt}{5.4mm}\else\rule[-2.2mm]{0pt}{6.6mm}\fi}
\newcolumntype{C}{>{\centering\arraybackslash}X}
\newcolumntype{L}{>{\raggedright\arraybackslash}X}

\titleformat{\section}{\large\bfseries}{}{0pt}{}
\titlespacing*{\section}{0pt}{6pt}{3pt}

\tcbset{qstyle/.style={enhanced,boxrule=0pt,leftrule=2.5pt,arc=0pt,outer arc=0pt,
  left=5pt,right=5pt,top=2pt,bottom=2pt,before skip=4pt,after skip=3pt}}
\newtcolorbox{cdc}{qstyle,colback=cdcbg,colframe=cdcbar}
\newtcolorbox{quest}{qstyle,colback=qbg,colframe=qbar}

% Zone de réponse : vierge dans le sujet, remplie dans le corrigé
\newcommand{\zone}[2]{%
  \ifcorrige
    \begin{tcolorbox}[colback=corrbg,colframe=corr,boxrule=0.5pt,arc=1pt,left=4pt,right=4pt,top=1pt,bottom=1pt,
      before skip=2pt,after skip=2pt,fontupper=\footnotesize\color{corr}]#2\end{tcolorbox}%
  \else
    \begin{tcolorbox}[colback=white,colframe=black!45,boxrule=0.4pt,arc=1pt,height=#1,
      before skip=2pt,after skip=2pt]\end{tcolorbox}%
  \fi}

\newcommand{\question}[3]{\begin{quest}\textbf{#1.} #3\pts{#2}\end{quest}}

%==========================================================================
% Données propres à chaque variante
%==========================================================================
\if\variante A
  \def\entreprise{Le cabinet d'architectes \textbf{ArchiNova}}
  \def\Rnom{R-ArchiNova}
  \def\resUn{Bureau d'études}\def\pUn{BE}\def\netUn{192.168.40}
  \def\resDeux{Accueil}\def\pDeux{ACC}\def\netDeux{172.16}
  \def\srvNom{SRV-Intranet}\def\netSrv{192.168.200}\def\srvH{10}
  \def\ssid{ArchiNova-WiFi}
  \def\AIrows{%
    \hl\ip{192.168.40.25} & \sol{C} & \sol{255.255.255.0} & \sol{192.168.40.0} & \sol{192.168.40.255} & \sol{254} \\ \hline
    \hl\ip{172.16.5.10}   & \sol{B} & \sol{255.255.0.0}   & \sol{172.16.0.0}   & \sol{172.16.255.255} & \sol{65\,534} \\ \hline
    \hl\ip{10.3.2.1}      & \sol{A} & \sol{255.0.0.0}     & \sol{10.0.0.0}     & \sol{10.255.255.255} & \sol{16\,777\,214} \\ \hline}
  \def\apipa{169.254.12.7}
  \def\Cnet{192.168.60}
  \def\Crows{%
    PC-A & \ip{192.168.60.10} & \ip{255.255.255.0} & \ip{192.168.60.254} & — \\
    PC-B & \ip{192.168.60.20} & \ip{255.255.255.0} & \ip{192.168.60.1}   & — \\
    PC-C & \ip{192.168.61.30} & \ip{255.255.255.0} & \ip{192.168.60.254} & — \\
    PC-D & \ip{192.168.60.10} & \ip{255.255.255.0} & \ip{192.168.60.254} & — \\
    SRV-C & \ip{10.0.0.5}     & \ip{255.0.0.0}     & \ip{10.0.0.254}     & — \\
    R1 — G0/0 & \ip{192.168.60.254} & \ip{255.255.255.0} & — & On \\
    R1 — G0/1 & \ip{10.0.0.254}     & \ip{255.0.0.0}     & — & Off \\}
  \def\Csol{%
    \hl 1 & \sol{PC-B} & \sol{Passerelle .60.1 : pas l'adresse du routeur.} & \sol{Joint son réseau, mais aucun autre (serveur).} & \sol{Passerelle 192.168.60.254.} \\ \hline
    \hl 2 & \sol{PC-C} & \sol{IP hors du réseau 192.168.60.0/24.} & \sol{Ne joint ni les postes ni le routeur.} & \sol{Ex. 192.168.60.30.} \\ \hline
    \hl 3 & \sol{PC-D} & \sol{Même IP que PC-A (doublon).} & \sol{Conflit d'adresse : PC-A et PC-D perturbés.} & \sol{Adresse unique, ex. .60.40.} \\ \hline
    \hl 4 & \sol{R1 G0/1} & \sol{Interface désactivée (Off).} & \sol{Réseau du serveur injoignable pour tous.} & \sol{Port Status : On.} \\ \hline}
  \def\Cbareme{PC-A et SRV-C sont correctement configurés ; accepter PC-A cité à la place de PC-D pour le doublon.}
\else
  \def\entreprise{La clinique vétérinaire \textbf{VetoPlus}}
  \def\Rnom{R-VetoPlus}
  \def\resUn{Soins}\def\pUn{SOI}\def\netUn{192.168.70}
  \def\resDeux{Secrétariat}\def\pDeux{SEC}\def\netDeux{172.20}
  \def\srvNom{SRV-Dossiers}\def\netSrv{192.168.250}\def\srvH{20}
  \def\ssid{VetoPlus-WiFi}
  \def\AIrows{%
    \hl\ip{192.168.15.130} & \sol{C} & \sol{255.255.255.0} & \sol{192.168.15.0} & \sol{192.168.15.255} & \sol{254} \\ \hline
    \hl\ip{172.31.200.7}   & \sol{B} & \sol{255.255.0.0}   & \sol{172.31.0.0}   & \sol{172.31.255.255} & \sol{65\,534} \\ \hline
    \hl\ip{10.45.6.9}      & \sol{A} & \sol{255.0.0.0}     & \sol{10.0.0.0}     & \sol{10.255.255.255} & \sol{16\,777\,214} \\ \hline}
  \def\apipa{169.254.88.3}
  \def\Cnet{192.168.80}
  \def\Crows{%
    PC-A & \ip{192.168.8.11}   & \ip{255.255.255.0} & \ip{192.168.80.254} & — \\
    PC-B & \ip{192.168.80.12}  & \ip{255.255.255.0} & \ip{192.168.80.254} & — \\
    PC-C & \ip{192.168.80.13}  & \ip{255.255.255.0} & \ip{192.168.80.100} & — \\
    PC-D & \ip{192.168.80.12}  & \ip{255.255.255.0} & \ip{192.168.80.254} & — \\
    SRV-C & \ip{10.0.0.8}      & \ip{255.0.0.0}     & \ip{10.0.0.254}     & — \\
    R1 — G0/0 & \ip{192.168.80.254} & \ip{255.255.255.0} & — & Off \\
    R1 — G0/1 & \ip{10.0.0.254}     & \ip{255.0.0.0}     & — & On \\}
  \def\Csol{%
    \hl 1 & \sol{PC-A} & \sol{IP hors du réseau 192.168.80.0/24.} & \sol{Ne joint ni les postes ni le routeur.} & \sol{Ex. 192.168.80.11.} \\ \hline
    \hl 2 & \sol{PC-C} & \sol{Passerelle .80.100 : pas l'adresse du routeur.} & \sol{Joint son réseau, mais aucun autre (serveur).} & \sol{Passerelle 192.168.80.254.} \\ \hline
    \hl 3 & \sol{PC-D} & \sol{Même IP que PC-B (doublon).} & \sol{Conflit d'adresse : PC-B et PC-D perturbés.} & \sol{Adresse unique, ex. .80.14.} \\ \hline
    \hl 4 & \sol{R1 G0/0} & \sol{Interface côté LAN désactivée (Off).} & \sol{Aucun poste ne joint le routeur ni le serveur.} & \sol{Port Status : On.} \\ \hline}
  \def\Cbareme{PC-B et SRV-C sont correctement configurés ; accepter PC-B cité à la place de PC-D pour le doublon.}
\fi

\pagestyle{fancy}
\fancyhf{}
\renewcommand{\headrulewidth}{0pt}
\fancyfoot[L]{\footnotesize\itshape\color{grisfoot}Évaluation — Réseaux locaux — Sujet \variante\ifcorrige\ — CORRIGÉ\fi}
\fancyfoot[R]{\footnotesize\color{grisfoot}\thepage/3}

\setlist{itemsep=0pt,topsep=2pt,parsep=0pt}
\setlength{\parindent}{0pt}
\setlength{\parskip}{2pt}

\begin{document}

\begin{center}
  {\Large\bfseries Évaluation pratique — Réseaux locaux \quad
   \fcolorbox{black}{white}{\ SUJET \variante\ }\par}
  \vspace{1pt}
  {\bfseries Adressage IP, commutation et routage — BTS CIEL 1re année — durée : 2 heures\par}
  \ifcorrige\vspace{2pt}{\bfseries\color{corr}CORRIGÉ ET BARÈME — document enseignant}\par\fi
\end{center}
\vspace{-2pt}
{\renewcommand{\arraystretch}{1.35}
\begin{tabularx}{\linewidth}{|X|X|p{2.2cm}|p{2.6cm}|}
  \hline
  \textbf{Nom :} & \textbf{Prénom :} & \textbf{Poste n°} & \textbf{Note :} \hfill /20 \\ \hline
\end{tabularx}}

\textbf{Consignes :} travail \textbf{individuel}, documents et Internet \textbf{non autorisés}. Ce sujet est le
\textbf{document-réponse}. Fichier Packet Tracer : \ip{Nom\_Prenom\_EvalLAN\_\variante.pkt}, à enregistrer
régulièrement dans le dossier de rendu. Les « visas » sont apposés par le professeur, que vous appelez à chaque test réussi.
Durées conseillées : A 30 min — B 1 h 10 — C 20 min.

%==========================================================================
\section{Partie A — Questions de cours et calculs \hfill /6}

\question{A1}{2 pts}{Complétez le tableau en considérant le \textbf{masque par défaut} de la classe de chaque adresse.}
{\small\renewcommand{\arraystretch}{1.15}
\begin{tabularx}{\linewidth}{|l|c|C|C|C|c|}
  \hline
  \textbf{Adresse IP} & \textbf{Classe} & \textbf{Masque par défaut} & \textbf{Adresse réseau} &
  \textbf{Adresse de diffusion} & \textbf{Nb d'hôtes} \\ \hline
  \AIrows
\end{tabularx}}
\bareme{0,5 pt par ligne juste + 0,5 pt si les trois nombres d'hôtes sont justes ($2^{n}-2$, $n$ = nombre de bits d'hôte).}

\question{A2}{1 pt}{Indiquez le type de câble (\textbf{droit} ou \textbf{croisé}) pour chaque liaison.}
{\small\renewcommand{\arraystretch}{1.15}
\begin{tabularx}{\linewidth}{|X|c||X|c|}
  \hline
  \hl PC $\leftrightarrow$ commutateur & \makebox[1.6cm]{\sol{droit}} & Commutateur $\leftrightarrow$ routeur & \makebox[1.6cm]{\sol{droit}} \\ \hline
  \hl PC $\leftrightarrow$ PC (réseau P2P) & \sol{croisé} & Concentrateur $\leftrightarrow$ concentrateur & \sol{croisé} \\ \hline
\end{tabularx}}
\bareme{0,25 pt par case. Équipements de types différents : câble droit ; de même type : câble croisé.}

\question{A3}{1 pt}{Comment un \textbf{concentrateur (hub)} et un \textbf{commutateur (switch)} traitent-ils une trame
reçue ? Lequel limite les collisions, et pourquoi ?}
\zone{2.9cm}{Le \textbf{hub} (couche 1) répète la trame sur \textbf{tous ses autres ports} : un seul domaine de collision
partagé. Le \textbf{switch} (couche 2) lit l'adresse MAC de destination et, grâce à sa \textbf{table MAC}, ne l'envoie
\textbf{que sur le port du destinataire} : un domaine de collision par port, c'est lui qui limite les collisions.
\emph{0,5 pt hub, 0,5 pt switch.}}

\question{A4}{1 pt}{Qu'est-ce que la \textbf{passerelle par défaut} d'un poste ? Quelle condition son adresse doit-elle respecter ?}
\zone{2.4cm}{Adresse IP de l'interface du \textbf{routeur} à laquelle le poste envoie les paquets destinés à un
\textbf{autre réseau}. Elle doit \textbf{appartenir au même réseau} que le poste. \emph{0,5 + 0,5 pt.}}

\question{A5}{0,5 pt}{Un portable configuré en DHCP affiche \ip{\apipa} / \ip{255.255.0.0}. Que s'est-il passé ?}
\zone{1.5cm}{Aucun \textbf{serveur DHCP} n'a répondu : le poste s'est attribué une adresse \textbf{APIPA} (169.254.0.0/16).}

\question{A6}{0,5 pt}{Quel protocole utilise \texttt{ping} ? Nommez les deux messages échangés et la couche OSI des adresses IP.}
\zone{1.5cm}{\textbf{ICMP} ; \textbf{Echo Request} / \textbf{Echo Reply} ; adresses IP en \textbf{couche 3 (Réseau)}.}

%==========================================================================
\newpage
\section{Partie B — Réalisation sous Packet Tracer \hfill /10}

\begin{cdc}
\entreprise{} vous confie son réseau : \textbf{trois réseaux} reliés par un \textbf{routeur unique Cisco 2911}
nommé \ip{\Rnom} (interfaces G0/0, G0/1, G0/2).
\begin{itemize}[leftmargin=12pt]
  \item \textbf{Réseau \resUn} : \textbf{4 postes} PC-\pUn1 à PC-\pUn4, commutateur \ip{SW-\pUn}, adressage statique,
        réseau \ip{\netUn.0/24}.
  \item \textbf{Réseau \resDeux} : \textbf{2 postes} PC-\pDeux1, PC-\pDeux2 et \textbf{1 imprimante} IMP-\pDeux,
        commutateur \ip{SW-\pDeux}, adressage statique, réseau de classe B \ip{\netDeux.0.0/16}.
  \item \textbf{Réseau Serveur} : serveur \ip{\srvNom} d'adresse \ip{\netSrv.\srvH/24}, relié \textbf{directement} au routeur.
  \item La passerelle de chaque réseau est sa \textbf{dernière adresse utilisable}. Tous les équipements communiquent
        entre eux et accèdent à la page web du serveur. Nommage imposé (« Display Name »).
\end{itemize}
\end{cdc}

\question{B1}{2 pts}{Complétez le plan d'adressage \textbf{avant} la réalisation.}
{\footnotesize\setlength{\tabcolsep}{3pt}\renewcommand{\arraystretch}{1.15}
\begin{tabularx}{\linewidth}{|l|C|C|C|C|C|C|}
  \hline
  \textbf{Réseau} & \textbf{Adresse réseau} & \textbf{Masque} & \textbf{1re @ utilisable} &
  \textbf{Dernière @ utilisable} & \textbf{Diffusion} & \textbf{Passerelle} \\ \hline
  \hl \resUn & \sol{\netUn.0} & \sol{255.255.255.0} & \sol{\netUn.1} & \sol{\netUn.254} & \sol{\netUn.255} & \sol{\netUn.254} \\ \hline
  \hl \resDeux & \sol{\netDeux.0.0} & \sol{255.255.0.0} & \sol{\netDeux.0.1} & \sol{\netDeux.255.254} & \sol{\netDeux.255.255} & \sol{\netDeux.255.254} \\ \hline
  \hl Serveur & \sol{\netSrv.0} & \sol{255.255.255.0} & \sol{\netSrv.1} & \sol{\netSrv.254} & \sol{\netSrv.255} & \sol{\netSrv.254} \\ \hline
\end{tabularx}}
\bareme{0,5 pt par ligne juste + 0,5 pt si la convention de passerelle est respectée partout.}

\question{B2}{2 pts}{Construisez la topologie (équipements, câblage, nommage) et relevez l'adressage configuré.}
{\footnotesize\setlength{\tabcolsep}{3pt}\renewcommand{\arraystretch}{1.1}
\begin{tabularx}{\linewidth}{|l|C|C|C||l|C|C|}
  \hline
  \textbf{Hôte} & \textbf{Adresse IP} & \textbf{Masque} & \textbf{Passerelle} & \textbf{Interface} & \textbf{Adresse IP} & \textbf{Masque} \\ \hline
  \hl PC-\pUn1 & \sol{\netUn.1} & \sol{255.255.255.0} & \sol{\netUn.254} & \Rnom{} G0/0 & \sol{\netUn.254} & \sol{255.255.255.0} \\ \hline
  \hl PC-\pDeux1 & \sol{\netDeux.0.1} & \sol{255.255.0.0} & \sol{\netDeux.255.254} & \Rnom{} G0/1 & \sol{\netDeux.255.254} & \sol{255.255.0.0} \\ \hline
  \hl IMP-\pDeux & \sol{\netDeux.0.10} & \sol{255.255.0.0} & \sol{\netDeux.255.254} & \Rnom{} G0/2 & \sol{\netSrv.254} & \sol{255.255.255.0} \\ \hline
  \hl \srvNom & \ip{\netSrv.\srvH} & \sol{255.255.255.0} & \sol{\netSrv.254} & \multicolumn{3}{c|}{\cellcolor{black!8}} \\ \hline
\end{tabularx}}
\bareme{Adresses d'hôtes libres (uniques, dans le bon réseau) ; affectation G0/x libre. Sur le .pkt : équipements 0,5 ;
câblage 1 (câbles droits vers les switchs et switch–routeur, \textbf{croisé} serveur–routeur, liens verts) ; nommage 0,5.}

\question{B3}{3 pts}{Configurez tous les équipements : IP, masques, passerelles des hôtes ; interfaces du routeur
(« Port Status : On »).}
\bareme{Sur le .pkt : interfaces routeur 1 pt ; IP/masques des hôtes 1 pt ; passerelles (imprimante et serveur compris)
1 pt. $-0,25$ par erreur, sans descendre sous 0 par item.}

\question{B4}{3 pts}{Réalisez les tests, notez la commande et le résultat, puis faites viser chaque test réussi.}
{\small\renewcommand{\arraystretch}{1.1}
\begin{tabularx}{\linewidth}{|c|l|l|C|c|c|}
  \hline
  \textbf{N°} & \textbf{Source} & \textbf{Destination} & \textbf{Commande / action} & \textbf{Résultat} & \textbf{Visa prof.} \\ \hline
  \hl 1 & PC-\pUn1 & PC-\pUn4 & \sol{ping \netUn.4} & \sol{réussi} & \hspace{1.4cm} \\ \hline
  \hl 2 & PC-\pDeux1 & IMP-\pDeux & \sol{ping \netDeux.0.10} & \sol{réussi} & \\ \hline
  \hl 3 & PC-\pUn2 & PC-\pDeux2 & \sol{ping \netDeux.0.2} & \sol{réussi} & \\ \hline
  \hl 4 & PC-\pDeux2 & \srvNom & \sol{ping \netSrv.\srvH} & \sol{réussi} & \\ \hline
  \hl 5 & PC-\pUn3 & page web du serveur & \sol{navigateur : {\NoAutoSpacing http://\netSrv.\srvH}} & \sol{page affichée} & \\ \hline
\end{tabularx}}
\bareme{0,5 pt par test visé + 0,5 pt pour B5. Le premier ping inter-réseaux peut échouer (résolution ARP) ; les suivants doivent réussir.}

\question{B5}{incl. B4}{Test n°3 : par quels équipements passe le paquet ? Pourquoi PC-\pUn2 l'envoie-t-il au routeur ?}
\zone{2.3cm}{PC-\pUn2 $\to$ SW-\pUn{} $\to$ \Rnom{} $\to$ SW-\pDeux{} $\to$ PC-\pDeux2. Grâce à son masque, PC-\pUn2 constate que
la destination (\netDeux.0.0/16) n'est \textbf{pas dans son réseau} (\netUn.0/24) : il envoie le paquet à sa
\textbf{passerelle par défaut}, qui le route.}

%==========================================================================
\newpage
\section{Partie C — Diagnostic d'un réseau en panne \hfill /4}

Un technicien constate qu'\textbf{aucun poste ne parvient à joindre le serveur SRV-C} et que certains postes
communiquent mal entre eux. Les câbles sont du bon type.

\begin{center}
\begin{minipage}{0.37\linewidth}\centering
\begin{tikzpicture}[font=\footnotesize,
  eq/.style={draw,fill=qbg,rounded corners=2pt,minimum width=9mm,minimum height=5mm,inner sep=2pt},
  every path/.style={thick}]
  \node[eq] (SW) {SW-1};
  \node[eq,right=9mm of SW] (R) {R1};
  \node[eq,right=7mm of R] (SRV) {SRV-C};
  \node[eq,above left=5mm and -3mm of SW] (A) {PC-A};
  \node[eq,above right=5mm and -3mm of SW] (B) {PC-B};
  \node[eq,below left=5mm and -3mm of SW] (C) {PC-C};
  \node[eq,below right=5mm and -3mm of SW] (D) {PC-D};
  \foreach \n in {A,B,C,D} \draw (\n) -- (SW);
  \draw (SW) -- node[above,font=\tiny]{G0/0} (R);
  \draw (R) -- node[above,font=\tiny]{G0/1} (SRV);
\end{tikzpicture}
\end{minipage}\hfill
\begin{minipage}{0.61\linewidth}\centering
{\small\renewcommand{\arraystretch}{1.05}\setlength{\tabcolsep}{4pt}
\begin{tabular}{lllll}
  \toprule
  \textbf{Équipement} & \textbf{Adresse IP} & \textbf{Masque} & \textbf{Passerelle} & \textbf{Port} \\
  \midrule
  \Crows
  \bottomrule
\end{tabular}}
\end{minipage}
\end{center}

\question{C1}{4 pts}{Relevez les \textbf{quatre anomalies}. Pour chacune : conséquence et correction.}
{\small\renewcommand{\arraystretch}{1.1}
\ifcorrige\def\hc{}\else\def\hc{\rule[-10mm]{0pt}{14mm}}\fi
\begin{tabularx}{\linewidth}{|c|>{\raggedright\arraybackslash}p{2cm}|L|L|L|}
  \hline
  \textbf{N°} & \textbf{Équipement} & \textbf{Anomalie} & \textbf{Conséquence} & \textbf{Correction} \\ \hline
  \ifcorrige\Csol\else
  \hc 1 & & & & \\ \hline
  \hc 2 & & & & \\ \hline
  \hc 3 & & & & \\ \hline
  \hc 4 & & & & \\ \hline
  \fi
\end{tabularx}}
\bareme{1 pt par anomalie (0,25 équipement, anomalie, conséquence, correction ; ordre indifférent). \Cbareme}

%==========================================================================
\section{Bonus — Extension sans fil \hfill +1}

\begin{cdc}Ajoutez au réseau \resUn{} un \textbf{AccessPoint-PT} relié à SW-\pUn{} (SSID \ip{\ssid}, \textbf{WPA2-PSK})
et un portable \ip{Laptop-\pUn} équipé d'un module \texttt{WPC300N}, en adressage \textbf{statique}.
Vérifiez qu'il joint \srvNom.\end{cdc}
{\small\renewcommand{\arraystretch}{1.1}
\begin{tabularx}{\linewidth}{|l|C|C|C|c|}
  \hline
  \textbf{Équipement} & \textbf{Adresse IP} & \textbf{Masque} & \textbf{Passerelle} & \textbf{Visa prof.} \\ \hline
  \hl Laptop-\pUn & \sol{\netUn.50} & \sol{255.255.255.0} & \sol{\netUn.254} & \hspace{1.4cm} \\ \hline
\end{tabularx}}
\bareme{Le point d'accès fait office de pont : le portable est dans \netUn.0/24. Bonus si associé en WPA2-PSK et ping
vers \netSrv.\srvH{} réussi. Note plafonnée à 20.}

\end{document}
"
\documentclass[10pt,a4paper]{article}

\usepackage[T1]{fontenc}
\usepackage[utf8]{inputenc}
\usepackage{lmodern}
\usepackage[french]{babel}
\usepackage[a4paper,hmargin=1.5cm,top=1.3cm,bottom=1.6cm,footskip=0.7cm]{geometry}
\usepackage[table]{xcolor}
\usepackage{booktabs,tabularx,array}
\usepackage{enumitem}
\usepackage{fancyhdr}
\usepackage[most]{tcolorbox}
\usepackage{tikz}
\usetikzlibrary{positioning}
\usepackage{titlesec}

\providecommand{\modecorrige}{0}
\providecommand{\variante}{A}
\newif\ifcorrige
\ifnum\modecorrige=1 \corrigetrue \else \corrigefalse \fi

\definecolor{ipred}{HTML}{C0392B}
\definecolor{cdcbar}{HTML}{A67C3D}
\definecolor{cdcbg}{HTML}{F4F1EA}
\definecolor{qbar}{HTML}{4A6B5D}
\definecolor{qbg}{HTML}{EDF1EF}
\definecolor{corr}{HTML}{1F5FA8}
\definecolor{corrbg}{HTML}{E8F0FA}
\definecolor{grisfoot}{HTML}{828282}

\newcommand{\ip}[1]{\textcolor{ipred}{\texttt{#1}}}
\newcommand{\pts}[1]{\hfill\mbox{\textbf{(#1)}}}
\newcommand{\sol}[1]{\ifcorrige\textcolor{corr}{#1}\fi}
\newcommand{\bareme}[1]{\ifcorrige\par{\footnotesize\color{corr}#1\par}\fi}
\newcommand{\hl}{\ifcorrige\rule[-1.6mm]{0pt}{5.4mm}\else\rule[-2.2mm]{0pt}{6.6mm}\fi}
\newcolumntype{C}{>{\centering\arraybackslash}X}
\newcolumntype{L}{>{\raggedright\arraybackslash}X}

\titleformat{\section}{\large\bfseries}{}{0pt}{}
\titlespacing*{\section}{0pt}{6pt}{3pt}

\tcbset{qstyle/.style={enhanced,boxrule=0pt,leftrule=2.5pt,arc=0pt,outer arc=0pt,
  left=5pt,right=5pt,top=2pt,bottom=2pt,before skip=4pt,after skip=3pt}}
\newtcolorbox{cdc}{qstyle,colback=cdcbg,colframe=cdcbar}
\newtcolorbox{quest}{qstyle,colback=qbg,colframe=qbar}

% Zone de réponse : vierge dans le sujet, remplie dans le corrigé
\newcommand{\zone}[2]{%
  \ifcorrige
    \begin{tcolorbox}[colback=corrbg,colframe=corr,boxrule=0.5pt,arc=1pt,left=4pt,right=4pt,top=1pt,bottom=1pt,
      before skip=2pt,after skip=2pt,fontupper=\footnotesize\color{corr}]#2\end{tcolorbox}%
  \else
    \begin{tcolorbox}[colback=white,colframe=black!45,boxrule=0.4pt,arc=1pt,height=#1,
      before skip=2pt,after skip=2pt]\end{tcolorbox}%
  \fi}

\newcommand{\question}[3]{\begin{quest}\textbf{#1.} #3\pts{#2}\end{quest}}

%==========================================================================
% Données propres à chaque variante
%==========================================================================
\if\variante A
  \def\entreprise{Le cabinet d'architectes \textbf{ArchiNova}}
  \def\Rnom{R-ArchiNova}
  \def\resUn{Bureau d'études}\def\pUn{BE}\def\netUn{192.168.40}
  \def\resDeux{Accueil}\def\pDeux{ACC}\def\netDeux{172.16}
  \def\srvNom{SRV-Intranet}\def\netSrv{192.168.200}\def\srvH{10}
  \def\ssid{ArchiNova-WiFi}
  \def\AIrows{%
    \hl\ip{192.168.40.25} & \sol{C} & \sol{255.255.255.0} & \sol{192.168.40.0} & \sol{192.168.40.255} & \sol{254} \\ \hline
    \hl\ip{172.16.5.10}   & \sol{B} & \sol{255.255.0.0}   & \sol{172.16.0.0}   & \sol{172.16.255.255} & \sol{65\,534} \\ \hline
    \hl\ip{10.3.2.1}      & \sol{A} & \sol{255.0.0.0}     & \sol{10.0.0.0}     & \sol{10.255.255.255} & \sol{16\,777\,214} \\ \hline}
  \def\apipa{169.254.12.7}
  \def\Cnet{192.168.60}
  \def\Crows{%
    PC-A & \ip{192.168.60.10} & \ip{255.255.255.0} & \ip{192.168.60.254} & — \\
    PC-B & \ip{192.168.60.20} & \ip{255.255.255.0} & \ip{192.168.60.1}   & — \\
    PC-C & \ip{192.168.61.30} & \ip{255.255.255.0} & \ip{192.168.60.254} & — \\
    PC-D & \ip{192.168.60.10} & \ip{255.255.255.0} & \ip{192.168.60.254} & — \\
    SRV-C & \ip{10.0.0.5}     & \ip{255.0.0.0}     & \ip{10.0.0.254}     & — \\
    R1 — G0/0 & \ip{192.168.60.254} & \ip{255.255.255.0} & — & On \\
    R1 — G0/1 & \ip{10.0.0.254}     & \ip{255.0.0.0}     & — & Off \\}
  \def\Csol{%
    \hl 1 & \sol{PC-B} & \sol{Passerelle .60.1 : pas l'adresse du routeur.} & \sol{Joint son réseau, mais aucun autre (serveur).} & \sol{Passerelle 192.168.60.254.} \\ \hline
    \hl 2 & \sol{PC-C} & \sol{IP hors du réseau 192.168.60.0/24.} & \sol{Ne joint ni les postes ni le routeur.} & \sol{Ex. 192.168.60.30.} \\ \hline
    \hl 3 & \sol{PC-D} & \sol{Même IP que PC-A (doublon).} & \sol{Conflit d'adresse : PC-A et PC-D perturbés.} & \sol{Adresse unique, ex. .60.40.} \\ \hline
    \hl 4 & \sol{R1 G0/1} & \sol{Interface désactivée (Off).} & \sol{Réseau du serveur injoignable pour tous.} & \sol{Port Status : On.} \\ \hline}
  \def\Cbareme{PC-A et SRV-C sont correctement configurés ; accepter PC-A cité à la place de PC-D pour le doublon.}
\else
  \def\entreprise{La clinique vétérinaire \textbf{VetoPlus}}
  \def\Rnom{R-VetoPlus}
  \def\resUn{Soins}\def\pUn{SOI}\def\netUn{192.168.70}
  \def\resDeux{Secrétariat}\def\pDeux{SEC}\def\netDeux{172.20}
  \def\srvNom{SRV-Dossiers}\def\netSrv{192.168.250}\def\srvH{20}
  \def\ssid{VetoPlus-WiFi}
  \def\AIrows{%
    \hl\ip{192.168.15.130} & \sol{C} & \sol{255.255.255.0} & \sol{192.168.15.0} & \sol{192.168.15.255} & \sol{254} \\ \hline
    \hl\ip{172.31.200.7}   & \sol{B} & \sol{255.255.0.0}   & \sol{172.31.0.0}   & \sol{172.31.255.255} & \sol{65\,534} \\ \hline
    \hl\ip{10.45.6.9}      & \sol{A} & \sol{255.0.0.0}     & \sol{10.0.0.0}     & \sol{10.255.255.255} & \sol{16\,777\,214} \\ \hline}
  \def\apipa{169.254.88.3}
  \def\Cnet{192.168.80}
  \def\Crows{%
    PC-A & \ip{192.168.8.11}   & \ip{255.255.255.0} & \ip{192.168.80.254} & — \\
    PC-B & \ip{192.168.80.12}  & \ip{255.255.255.0} & \ip{192.168.80.254} & — \\
    PC-C & \ip{192.168.80.13}  & \ip{255.255.255.0} & \ip{192.168.80.100} & — \\
    PC-D & \ip{192.168.80.12}  & \ip{255.255.255.0} & \ip{192.168.80.254} & — \\
    SRV-C & \ip{10.0.0.8}      & \ip{255.0.0.0}     & \ip{10.0.0.254}     & — \\
    R1 — G0/0 & \ip{192.168.80.254} & \ip{255.255.255.0} & — & Off \\
    R1 — G0/1 & \ip{10.0.0.254}     & \ip{255.0.0.0}     & — & On \\}
  \def\Csol{%
    \hl 1 & \sol{PC-A} & \sol{IP hors du réseau 192.168.80.0/24.} & \sol{Ne joint ni les postes ni le routeur.} & \sol{Ex. 192.168.80.11.} \\ \hline
    \hl 2 & \sol{PC-C} & \sol{Passerelle .80.100 : pas l'adresse du routeur.} & \sol{Joint son réseau, mais aucun autre (serveur).} & \sol{Passerelle 192.168.80.254.} \\ \hline
    \hl 3 & \sol{PC-D} & \sol{Même IP que PC-B (doublon).} & \sol{Conflit d'adresse : PC-B et PC-D perturbés.} & \sol{Adresse unique, ex. .80.14.} \\ \hline
    \hl 4 & \sol{R1 G0/0} & \sol{Interface côté LAN désactivée (Off).} & \sol{Aucun poste ne joint le routeur ni le serveur.} & \sol{Port Status : On.} \\ \hline}
  \def\Cbareme{PC-B et SRV-C sont correctement configurés ; accepter PC-B cité à la place de PC-D pour le doublon.}
\fi

\pagestyle{fancy}
\fancyhf{}
\renewcommand{\headrulewidth}{0pt}
\fancyfoot[L]{\footnotesize\itshape\color{grisfoot}Évaluation — Réseaux locaux — Sujet \variante\ifcorrige\ — CORRIGÉ\fi}
\fancyfoot[R]{\footnotesize\color{grisfoot}\thepage/3}

\setlist{itemsep=0pt,topsep=2pt,parsep=0pt}
\setlength{\parindent}{0pt}
\setlength{\parskip}{2pt}

\begin{document}

\begin{center}
  {\Large\bfseries Évaluation pratique — Réseaux locaux \quad
   \fcolorbox{black}{white}{\ SUJET \variante\ }\par}
  \vspace{1pt}
  {\bfseries Adressage IP, commutation et routage — BTS CIEL 1re année — durée : 2 heures\par}
  \ifcorrige\vspace{2pt}{\bfseries\color{corr}CORRIGÉ ET BARÈME — document enseignant}\par\fi
\end{center}
\vspace{-2pt}
{\renewcommand{\arraystretch}{1.35}
\begin{tabularx}{\linewidth}{|X|X|p{2.2cm}|p{2.6cm}|}
  \hline
  \textbf{Nom :} & \textbf{Prénom :} & \textbf{Poste n°} & \textbf{Note :} \hfill /20 \\ \hline
\end{tabularx}}

\textbf{Consignes :} travail \textbf{individuel}, documents et Internet \textbf{non autorisés}. Ce sujet est le
\textbf{document-réponse}. Fichier Packet Tracer : \ip{Nom\_Prenom\_EvalLAN\_\variante.pkt}, à enregistrer
régulièrement dans le dossier de rendu. Les « visas » sont apposés par le professeur, que vous appelez à chaque test réussi.
Durées conseillées : A 30 min — B 1 h 10 — C 20 min.

%==========================================================================
\section{Partie A — Questions de cours et calculs \hfill /6}

\question{A1}{2 pts}{Complétez le tableau en considérant le \textbf{masque par défaut} de la classe de chaque adresse.}
{\small\renewcommand{\arraystretch}{1.15}
\begin{tabularx}{\linewidth}{|l|c|C|C|C|c|}
  \hline
  \textbf{Adresse IP} & \textbf{Classe} & \textbf{Masque par défaut} & \textbf{Adresse réseau} &
  \textbf{Adresse de diffusion} & \textbf{Nb d'hôtes} \\ \hline
  \AIrows
\end{tabularx}}
\bareme{0,5 pt par ligne juste + 0,5 pt si les trois nombres d'hôtes sont justes ($2^{n}-2$, $n$ = nombre de bits d'hôte).}

\question{A2}{1 pt}{Indiquez le type de câble (\textbf{droit} ou \textbf{croisé}) pour chaque liaison.}
{\small\renewcommand{\arraystretch}{1.15}
\begin{tabularx}{\linewidth}{|X|c||X|c|}
  \hline
  \hl PC $\leftrightarrow$ commutateur & \makebox[1.6cm]{\sol{droit}} & Commutateur $\leftrightarrow$ routeur & \makebox[1.6cm]{\sol{droit}} \\ \hline
  \hl PC $\leftrightarrow$ PC (réseau P2P) & \sol{croisé} & Concentrateur $\leftrightarrow$ concentrateur & \sol{croisé} \\ \hline
\end{tabularx}}
\bareme{0,25 pt par case. Équipements de types différents : câble droit ; de même type : câble croisé.}

\question{A3}{1 pt}{Comment un \textbf{concentrateur (hub)} et un \textbf{commutateur (switch)} traitent-ils une trame
reçue ? Lequel limite les collisions, et pourquoi ?}
\zone{2.9cm}{Le \textbf{hub} (couche 1) répète la trame sur \textbf{tous ses autres ports} : un seul domaine de collision
partagé. Le \textbf{switch} (couche 2) lit l'adresse MAC de destination et, grâce à sa \textbf{table MAC}, ne l'envoie
\textbf{que sur le port du destinataire} : un domaine de collision par port, c'est lui qui limite les collisions.
\emph{0,5 pt hub, 0,5 pt switch.}}

\question{A4}{1 pt}{Qu'est-ce que la \textbf{passerelle par défaut} d'un poste ? Quelle condition son adresse doit-elle respecter ?}
\zone{2.4cm}{Adresse IP de l'interface du \textbf{routeur} à laquelle le poste envoie les paquets destinés à un
\textbf{autre réseau}. Elle doit \textbf{appartenir au même réseau} que le poste. \emph{0,5 + 0,5 pt.}}

\question{A5}{0,5 pt}{Un portable configuré en DHCP affiche \ip{\apipa} / \ip{255.255.0.0}. Que s'est-il passé ?}
\zone{1.5cm}{Aucun \textbf{serveur DHCP} n'a répondu : le poste s'est attribué une adresse \textbf{APIPA} (169.254.0.0/16).}

\question{A6}{0,5 pt}{Quel protocole utilise \texttt{ping} ? Nommez les deux messages échangés et la couche OSI des adresses IP.}
\zone{1.5cm}{\textbf{ICMP} ; \textbf{Echo Request} / \textbf{Echo Reply} ; adresses IP en \textbf{couche 3 (Réseau)}.}

%==========================================================================
\newpage
\section{Partie B — Réalisation sous Packet Tracer \hfill /10}

\begin{cdc}
\entreprise{} vous confie son réseau : \textbf{trois réseaux} reliés par un \textbf{routeur unique Cisco 2911}
nommé \ip{\Rnom} (interfaces G0/0, G0/1, G0/2).
\begin{itemize}[leftmargin=12pt]
  \item \textbf{Réseau \resUn} : \textbf{4 postes} PC-\pUn1 à PC-\pUn4, commutateur \ip{SW-\pUn}, adressage statique,
        réseau \ip{\netUn.0/24}.
  \item \textbf{Réseau \resDeux} : \textbf{2 postes} PC-\pDeux1, PC-\pDeux2 et \textbf{1 imprimante} IMP-\pDeux,
        commutateur \ip{SW-\pDeux}, adressage statique, réseau de classe B \ip{\netDeux.0.0/16}.
  \item \textbf{Réseau Serveur} : serveur \ip{\srvNom} d'adresse \ip{\netSrv.\srvH/24}, relié \textbf{directement} au routeur.
  \item La passerelle de chaque réseau est sa \textbf{dernière adresse utilisable}. Tous les équipements communiquent
        entre eux et accèdent à la page web du serveur. Nommage imposé (« Display Name »).
\end{itemize}
\end{cdc}

\question{B1}{2 pts}{Complétez le plan d'adressage \textbf{avant} la réalisation.}
{\footnotesize\setlength{\tabcolsep}{3pt}\renewcommand{\arraystretch}{1.15}
\begin{tabularx}{\linewidth}{|l|C|C|C|C|C|C|}
  \hline
  \textbf{Réseau} & \textbf{Adresse réseau} & \textbf{Masque} & \textbf{1re @ utilisable} &
  \textbf{Dernière @ utilisable} & \textbf{Diffusion} & \textbf{Passerelle} \\ \hline
  \hl \resUn & \sol{\netUn.0} & \sol{255.255.255.0} & \sol{\netUn.1} & \sol{\netUn.254} & \sol{\netUn.255} & \sol{\netUn.254} \\ \hline
  \hl \resDeux & \sol{\netDeux.0.0} & \sol{255.255.0.0} & \sol{\netDeux.0.1} & \sol{\netDeux.255.254} & \sol{\netDeux.255.255} & \sol{\netDeux.255.254} \\ \hline
  \hl Serveur & \sol{\netSrv.0} & \sol{255.255.255.0} & \sol{\netSrv.1} & \sol{\netSrv.254} & \sol{\netSrv.255} & \sol{\netSrv.254} \\ \hline
\end{tabularx}}
\bareme{0,5 pt par ligne juste + 0,5 pt si la convention de passerelle est respectée partout.}

\question{B2}{2 pts}{Construisez la topologie (équipements, câblage, nommage) et relevez l'adressage configuré.}
{\footnotesize\setlength{\tabcolsep}{3pt}\renewcommand{\arraystretch}{1.1}
\begin{tabularx}{\linewidth}{|l|C|C|C||l|C|C|}
  \hline
  \textbf{Hôte} & \textbf{Adresse IP} & \textbf{Masque} & \textbf{Passerelle} & \textbf{Interface} & \textbf{Adresse IP} & \textbf{Masque} \\ \hline
  \hl PC-\pUn1 & \sol{\netUn.1} & \sol{255.255.255.0} & \sol{\netUn.254} & \Rnom{} G0/0 & \sol{\netUn.254} & \sol{255.255.255.0} \\ \hline
  \hl PC-\pDeux1 & \sol{\netDeux.0.1} & \sol{255.255.0.0} & \sol{\netDeux.255.254} & \Rnom{} G0/1 & \sol{\netDeux.255.254} & \sol{255.255.0.0} \\ \hline
  \hl IMP-\pDeux & \sol{\netDeux.0.10} & \sol{255.255.0.0} & \sol{\netDeux.255.254} & \Rnom{} G0/2 & \sol{\netSrv.254} & \sol{255.255.255.0} \\ \hline
  \hl \srvNom & \ip{\netSrv.\srvH} & \sol{255.255.255.0} & \sol{\netSrv.254} & \multicolumn{3}{c|}{\cellcolor{black!8}} \\ \hline
\end{tabularx}}
\bareme{Adresses d'hôtes libres (uniques, dans le bon réseau) ; affectation G0/x libre. Sur le .pkt : équipements 0,5 ;
câblage 1 (câbles droits vers les switchs et switch–routeur, \textbf{croisé} serveur–routeur, liens verts) ; nommage 0,5.}

\question{B3}{3 pts}{Configurez tous les équipements : IP, masques, passerelles des hôtes ; interfaces du routeur
(« Port Status : On »).}
\bareme{Sur le .pkt : interfaces routeur 1 pt ; IP/masques des hôtes 1 pt ; passerelles (imprimante et serveur compris)
1 pt. $-0,25$ par erreur, sans descendre sous 0 par item.}

\question{B4}{3 pts}{Réalisez les tests, notez la commande et le résultat, puis faites viser chaque test réussi.}
{\small\renewcommand{\arraystretch}{1.1}
\begin{tabularx}{\linewidth}{|c|l|l|C|c|c|}
  \hline
  \textbf{N°} & \textbf{Source} & \textbf{Destination} & \textbf{Commande / action} & \textbf{Résultat} & \textbf{Visa prof.} \\ \hline
  \hl 1 & PC-\pUn1 & PC-\pUn4 & \sol{ping \netUn.4} & \sol{réussi} & \hspace{1.4cm} \\ \hline
  \hl 2 & PC-\pDeux1 & IMP-\pDeux & \sol{ping \netDeux.0.10} & \sol{réussi} & \\ \hline
  \hl 3 & PC-\pUn2 & PC-\pDeux2 & \sol{ping \netDeux.0.2} & \sol{réussi} & \\ \hline
  \hl 4 & PC-\pDeux2 & \srvNom & \sol{ping \netSrv.\srvH} & \sol{réussi} & \\ \hline
  \hl 5 & PC-\pUn3 & page web du serveur & \sol{navigateur : {\NoAutoSpacing http://\netSrv.\srvH}} & \sol{page affichée} & \\ \hline
\end{tabularx}}
\bareme{0,5 pt par test visé + 0,5 pt pour B5. Le premier ping inter-réseaux peut échouer (résolution ARP) ; les suivants doivent réussir.}

\question{B5}{incl. B4}{Test n°3 : par quels équipements passe le paquet ? Pourquoi PC-\pUn2 l'envoie-t-il au routeur ?}
\zone{2.3cm}{PC-\pUn2 $\to$ SW-\pUn{} $\to$ \Rnom{} $\to$ SW-\pDeux{} $\to$ PC-\pDeux2. Grâce à son masque, PC-\pUn2 constate que
la destination (\netDeux.0.0/16) n'est \textbf{pas dans son réseau} (\netUn.0/24) : il envoie le paquet à sa
\textbf{passerelle par défaut}, qui le route.}

%==========================================================================
\newpage
\section{Partie C — Diagnostic d'un réseau en panne \hfill /4}

Un technicien constate qu'\textbf{aucun poste ne parvient à joindre le serveur SRV-C} et que certains postes
communiquent mal entre eux. Les câbles sont du bon type.

\begin{center}
\begin{minipage}{0.37\linewidth}\centering
\begin{tikzpicture}[font=\footnotesize,
  eq/.style={draw,fill=qbg,rounded corners=2pt,minimum width=9mm,minimum height=5mm,inner sep=2pt},
  every path/.style={thick}]
  \node[eq] (SW) {SW-1};
  \node[eq,right=9mm of SW] (R) {R1};
  \node[eq,right=7mm of R] (SRV) {SRV-C};
  \node[eq,above left=5mm and -3mm of SW] (A) {PC-A};
  \node[eq,above right=5mm and -3mm of SW] (B) {PC-B};
  \node[eq,below left=5mm and -3mm of SW] (C) {PC-C};
  \node[eq,below right=5mm and -3mm of SW] (D) {PC-D};
  \foreach \n in {A,B,C,D} \draw (\n) -- (SW);
  \draw (SW) -- node[above,font=\tiny]{G0/0} (R);
  \draw (R) -- node[above,font=\tiny]{G0/1} (SRV);
\end{tikzpicture}
\end{minipage}\hfill
\begin{minipage}{0.61\linewidth}\centering
{\small\renewcommand{\arraystretch}{1.05}\setlength{\tabcolsep}{4pt}
\begin{tabular}{lllll}
  \toprule
  \textbf{Équipement} & \textbf{Adresse IP} & \textbf{Masque} & \textbf{Passerelle} & \textbf{Port} \\
  \midrule
  \Crows
  \bottomrule
\end{tabular}}
\end{minipage}
\end{center}

\question{C1}{4 pts}{Relevez les \textbf{quatre anomalies}. Pour chacune : conséquence et correction.}
{\small\renewcommand{\arraystretch}{1.1}
\ifcorrige\def\hc{}\else\def\hc{\rule[-10mm]{0pt}{14mm}}\fi
\begin{tabularx}{\linewidth}{|c|>{\raggedright\arraybackslash}p{2cm}|L|L|L|}
  \hline
  \textbf{N°} & \textbf{Équipement} & \textbf{Anomalie} & \textbf{Conséquence} & \textbf{Correction} \\ \hline
  \ifcorrige\Csol\else
  \hc 1 & & & & \\ \hline
  \hc 2 & & & & \\ \hline
  \hc 3 & & & & \\ \hline
  \hc 4 & & & & \\ \hline
  \fi
\end{tabularx}}
\bareme{1 pt par anomalie (0,25 équipement, anomalie, conséquence, correction ; ordre indifférent). \Cbareme}

%==========================================================================
\section{Bonus — Extension sans fil \hfill +1}

\begin{cdc}Ajoutez au réseau \resUn{} un \textbf{AccessPoint-PT} relié à SW-\pUn{} (SSID \ip{\ssid}, \textbf{WPA2-PSK})
et un portable \ip{Laptop-\pUn} équipé d'un module \texttt{WPC300N}, en adressage \textbf{statique}.
Vérifiez qu'il joint \srvNom.\end{cdc}
{\small\renewcommand{\arraystretch}{1.1}
\begin{tabularx}{\linewidth}{|l|C|C|C|c|}
  \hline
  \textbf{Équipement} & \textbf{Adresse IP} & \textbf{Masque} & \textbf{Passerelle} & \textbf{Visa prof.} \\ \hline
  \hl Laptop-\pUn & \sol{\netUn.50} & \sol{255.255.255.0} & \sol{\netUn.254} & \hspace{1.4cm} \\ \hline
\end{tabularx}}
\bareme{Le point d'accès fait office de pont : le portable est dans \netUn.0/24. Bonus si associé en WPA2-PSK et ping
vers \netSrv.\srvH{} réussi. Note plafonnée à 20.}

\end{document}
"
\documentclass[10pt,a4paper]{article}

\usepackage[T1]{fontenc}
\usepackage[utf8]{inputenc}
\usepackage{lmodern}
\usepackage[french]{babel}
\usepackage[a4paper,hmargin=1.5cm,top=1.3cm,bottom=1.6cm,footskip=0.7cm]{geometry}
\usepackage[table]{xcolor}
\usepackage{booktabs,tabularx,array}
\usepackage{enumitem}
\usepackage{fancyhdr}
\usepackage[most]{tcolorbox}
\usepackage{tikz}
\usetikzlibrary{positioning}
\usepackage{titlesec}

\providecommand{\modecorrige}{0}
\providecommand{\variante}{A}
\newif\ifcorrige
\ifnum\modecorrige=1 \corrigetrue \else \corrigefalse \fi

\definecolor{ipred}{HTML}{C0392B}
\definecolor{cdcbar}{HTML}{A67C3D}
\definecolor{cdcbg}{HTML}{F4F1EA}
\definecolor{qbar}{HTML}{4A6B5D}
\definecolor{qbg}{HTML}{EDF1EF}
\definecolor{corr}{HTML}{1F5FA8}
\definecolor{corrbg}{HTML}{E8F0FA}
\definecolor{grisfoot}{HTML}{828282}

\newcommand{\ip}[1]{\textcolor{ipred}{\texttt{#1}}}
\newcommand{\pts}[1]{\hfill\mbox{\textbf{(#1)}}}
\newcommand{\sol}[1]{\ifcorrige\textcolor{corr}{#1}\fi}
\newcommand{\bareme}[1]{\ifcorrige\par{\footnotesize\color{corr}#1\par}\fi}
\newcommand{\hl}{\ifcorrige\rule[-1.6mm]{0pt}{5.4mm}\else\rule[-2.2mm]{0pt}{6.6mm}\fi}
\newcolumntype{C}{>{\centering\arraybackslash}X}
\newcolumntype{L}{>{\raggedright\arraybackslash}X}

\titleformat{\section}{\large\bfseries}{}{0pt}{}
\titlespacing*{\section}{0pt}{6pt}{3pt}

\tcbset{qstyle/.style={enhanced,boxrule=0pt,leftrule=2.5pt,arc=0pt,outer arc=0pt,
  left=5pt,right=5pt,top=2pt,bottom=2pt,before skip=4pt,after skip=3pt}}
\newtcolorbox{cdc}{qstyle,colback=cdcbg,colframe=cdcbar}
\newtcolorbox{quest}{qstyle,colback=qbg,colframe=qbar}

% Zone de réponse : vierge dans le sujet, remplie dans le corrigé
\newcommand{\zone}[2]{%
  \ifcorrige
    \begin{tcolorbox}[colback=corrbg,colframe=corr,boxrule=0.5pt,arc=1pt,left=4pt,right=4pt,top=1pt,bottom=1pt,
      before skip=2pt,after skip=2pt,fontupper=\footnotesize\color{corr}]#2\end{tcolorbox}%
  \else
    \begin{tcolorbox}[colback=white,colframe=black!45,boxrule=0.4pt,arc=1pt,height=#1,
      before skip=2pt,after skip=2pt]\end{tcolorbox}%
  \fi}

\newcommand{\question}[3]{\begin{quest}\textbf{#1.} #3\pts{#2}\end{quest}}

%==========================================================================
% Données propres à chaque variante
%==========================================================================
\if\variante A
  \def\entreprise{Le cabinet d'architectes \textbf{ArchiNova}}
  \def\Rnom{R-ArchiNova}
  \def\resUn{Bureau d'études}\def\pUn{BE}\def\netUn{192.168.40}
  \def\resDeux{Accueil}\def\pDeux{ACC}\def\netDeux{172.16}
  \def\srvNom{SRV-Intranet}\def\netSrv{192.168.200}\def\srvH{10}
  \def\ssid{ArchiNova-WiFi}
  \def\AIrows{%
    \hl\ip{192.168.40.25} & \sol{C} & \sol{255.255.255.0} & \sol{192.168.40.0} & \sol{192.168.40.255} & \sol{254} \\ \hline
    \hl\ip{172.16.5.10}   & \sol{B} & \sol{255.255.0.0}   & \sol{172.16.0.0}   & \sol{172.16.255.255} & \sol{65\,534} \\ \hline
    \hl\ip{10.3.2.1}      & \sol{A} & \sol{255.0.0.0}     & \sol{10.0.0.0}     & \sol{10.255.255.255} & \sol{16\,777\,214} \\ \hline}
  \def\apipa{169.254.12.7}
  \def\Cnet{192.168.60}
  \def\Crows{%
    PC-A & \ip{192.168.60.10} & \ip{255.255.255.0} & \ip{192.168.60.254} & — \\
    PC-B & \ip{192.168.60.20} & \ip{255.255.255.0} & \ip{192.168.60.1}   & — \\
    PC-C & \ip{192.168.61.30} & \ip{255.255.255.0} & \ip{192.168.60.254} & — \\
    PC-D & \ip{192.168.60.10} & \ip{255.255.255.0} & \ip{192.168.60.254} & — \\
    SRV-C & \ip{10.0.0.5}     & \ip{255.0.0.0}     & \ip{10.0.0.254}     & — \\
    R1 — G0/0 & \ip{192.168.60.254} & \ip{255.255.255.0} & — & On \\
    R1 — G0/1 & \ip{10.0.0.254}     & \ip{255.0.0.0}     & — & Off \\}
  \def\Csol{%
    \hl 1 & \sol{PC-B} & \sol{Passerelle .60.1 : pas l'adresse du routeur.} & \sol{Joint son réseau, mais aucun autre (serveur).} & \sol{Passerelle 192.168.60.254.} \\ \hline
    \hl 2 & \sol{PC-C} & \sol{IP hors du réseau 192.168.60.0/24.} & \sol{Ne joint ni les postes ni le routeur.} & \sol{Ex. 192.168.60.30.} \\ \hline
    \hl 3 & \sol{PC-D} & \sol{Même IP que PC-A (doublon).} & \sol{Conflit d'adresse : PC-A et PC-D perturbés.} & \sol{Adresse unique, ex. .60.40.} \\ \hline
    \hl 4 & \sol{R1 G0/1} & \sol{Interface désactivée (Off).} & \sol{Réseau du serveur injoignable pour tous.} & \sol{Port Status : On.} \\ \hline}
  \def\Cbareme{PC-A et SRV-C sont correctement configurés ; accepter PC-A cité à la place de PC-D pour le doublon.}
\else
  \def\entreprise{La clinique vétérinaire \textbf{VetoPlus}}
  \def\Rnom{R-VetoPlus}
  \def\resUn{Soins}\def\pUn{SOI}\def\netUn{192.168.70}
  \def\resDeux{Secrétariat}\def\pDeux{SEC}\def\netDeux{172.20}
  \def\srvNom{SRV-Dossiers}\def\netSrv{192.168.250}\def\srvH{20}
  \def\ssid{VetoPlus-WiFi}
  \def\AIrows{%
    \hl\ip{192.168.15.130} & \sol{C} & \sol{255.255.255.0} & \sol{192.168.15.0} & \sol{192.168.15.255} & \sol{254} \\ \hline
    \hl\ip{172.31.200.7}   & \sol{B} & \sol{255.255.0.0}   & \sol{172.31.0.0}   & \sol{172.31.255.255} & \sol{65\,534} \\ \hline
    \hl\ip{10.45.6.9}      & \sol{A} & \sol{255.0.0.0}     & \sol{10.0.0.0}     & \sol{10.255.255.255} & \sol{16\,777\,214} \\ \hline}
  \def\apipa{169.254.88.3}
  \def\Cnet{192.168.80}
  \def\Crows{%
    PC-A & \ip{192.168.8.11}   & \ip{255.255.255.0} & \ip{192.168.80.254} & — \\
    PC-B & \ip{192.168.80.12}  & \ip{255.255.255.0} & \ip{192.168.80.254} & — \\
    PC-C & \ip{192.168.80.13}  & \ip{255.255.255.0} & \ip{192.168.80.100} & — \\
    PC-D & \ip{192.168.80.12}  & \ip{255.255.255.0} & \ip{192.168.80.254} & — \\
    SRV-C & \ip{10.0.0.8}      & \ip{255.0.0.0}     & \ip{10.0.0.254}     & — \\
    R1 — G0/0 & \ip{192.168.80.254} & \ip{255.255.255.0} & — & Off \\
    R1 — G0/1 & \ip{10.0.0.254}     & \ip{255.0.0.0}     & — & On \\}
  \def\Csol{%
    \hl 1 & \sol{PC-A} & \sol{IP hors du réseau 192.168.80.0/24.} & \sol{Ne joint ni les postes ni le routeur.} & \sol{Ex. 192.168.80.11.} \\ \hline
    \hl 2 & \sol{PC-C} & \sol{Passerelle .80.100 : pas l'adresse du routeur.} & \sol{Joint son réseau, mais aucun autre (serveur).} & \sol{Passerelle 192.168.80.254.} \\ \hline
    \hl 3 & \sol{PC-D} & \sol{Même IP que PC-B (doublon).} & \sol{Conflit d'adresse : PC-B et PC-D perturbés.} & \sol{Adresse unique, ex. .80.14.} \\ \hline
    \hl 4 & \sol{R1 G0/0} & \sol{Interface côté LAN désactivée (Off).} & \sol{Aucun poste ne joint le routeur ni le serveur.} & \sol{Port Status : On.} \\ \hline}
  \def\Cbareme{PC-B et SRV-C sont correctement configurés ; accepter PC-B cité à la place de PC-D pour le doublon.}
\fi

\pagestyle{fancy}
\fancyhf{}
\renewcommand{\headrulewidth}{0pt}
\fancyfoot[L]{\footnotesize\itshape\color{grisfoot}Évaluation — Réseaux locaux — Sujet \variante\ifcorrige\ — CORRIGÉ\fi}
\fancyfoot[R]{\footnotesize\color{grisfoot}\thepage/3}

\setlist{itemsep=0pt,topsep=2pt,parsep=0pt}
\setlength{\parindent}{0pt}
\setlength{\parskip}{2pt}

\begin{document}

\begin{center}
  {\Large\bfseries Évaluation pratique — Réseaux locaux \quad
   \fcolorbox{black}{white}{\ SUJET \variante\ }\par}
  \vspace{1pt}
  {\bfseries Adressage IP, commutation et routage — BTS CIEL 1re année — durée : 2 heures\par}
  \ifcorrige\vspace{2pt}{\bfseries\color{corr}CORRIGÉ ET BARÈME — document enseignant}\par\fi
\end{center}
\vspace{-2pt}
{\renewcommand{\arraystretch}{1.35}
\begin{tabularx}{\linewidth}{|X|X|p{2.2cm}|p{2.6cm}|}
  \hline
  \textbf{Nom :} & \textbf{Prénom :} & \textbf{Poste n°} & \textbf{Note :} \hfill /20 \\ \hline
\end{tabularx}}

\textbf{Consignes :} travail \textbf{individuel}, documents et Internet \textbf{non autorisés}. Ce sujet est le
\textbf{document-réponse}. Fichier Packet Tracer : \ip{Nom\_Prenom\_EvalLAN\_\variante.pkt}, à enregistrer
régulièrement dans le dossier de rendu. Les « visas » sont apposés par le professeur, que vous appelez à chaque test réussi.
Durées conseillées : A 30 min — B 1 h 10 — C 20 min.

%==========================================================================
\section{Partie A — Questions de cours et calculs \hfill /6}

\question{A1}{2 pts}{Complétez le tableau en considérant le \textbf{masque par défaut} de la classe de chaque adresse.}
{\small\renewcommand{\arraystretch}{1.15}
\begin{tabularx}{\linewidth}{|l|c|C|C|C|c|}
  \hline
  \textbf{Adresse IP} & \textbf{Classe} & \textbf{Masque par défaut} & \textbf{Adresse réseau} &
  \textbf{Adresse de diffusion} & \textbf{Nb d'hôtes} \\ \hline
  \AIrows
\end{tabularx}}
\bareme{0,5 pt par ligne juste + 0,5 pt si les trois nombres d'hôtes sont justes ($2^{n}-2$, $n$ = nombre de bits d'hôte).}

\question{A2}{1 pt}{Indiquez le type de câble (\textbf{droit} ou \textbf{croisé}) pour chaque liaison.}
{\small\renewcommand{\arraystretch}{1.15}
\begin{tabularx}{\linewidth}{|X|c||X|c|}
  \hline
  \hl PC $\leftrightarrow$ commutateur & \makebox[1.6cm]{\sol{droit}} & Commutateur $\leftrightarrow$ routeur & \makebox[1.6cm]{\sol{droit}} \\ \hline
  \hl PC $\leftrightarrow$ PC (réseau P2P) & \sol{croisé} & Concentrateur $\leftrightarrow$ concentrateur & \sol{croisé} \\ \hline
\end{tabularx}}
\bareme{0,25 pt par case. Équipements de types différents : câble droit ; de même type : câble croisé.}

\question{A3}{1 pt}{Comment un \textbf{concentrateur (hub)} et un \textbf{commutateur (switch)} traitent-ils une trame
reçue ? Lequel limite les collisions, et pourquoi ?}
\zone{2.9cm}{Le \textbf{hub} (couche 1) répète la trame sur \textbf{tous ses autres ports} : un seul domaine de collision
partagé. Le \textbf{switch} (couche 2) lit l'adresse MAC de destination et, grâce à sa \textbf{table MAC}, ne l'envoie
\textbf{que sur le port du destinataire} : un domaine de collision par port, c'est lui qui limite les collisions.
\emph{0,5 pt hub, 0,5 pt switch.}}

\question{A4}{1 pt}{Qu'est-ce que la \textbf{passerelle par défaut} d'un poste ? Quelle condition son adresse doit-elle respecter ?}
\zone{2.4cm}{Adresse IP de l'interface du \textbf{routeur} à laquelle le poste envoie les paquets destinés à un
\textbf{autre réseau}. Elle doit \textbf{appartenir au même réseau} que le poste. \emph{0,5 + 0,5 pt.}}

\question{A5}{0,5 pt}{Un portable configuré en DHCP affiche \ip{\apipa} / \ip{255.255.0.0}. Que s'est-il passé ?}
\zone{1.5cm}{Aucun \textbf{serveur DHCP} n'a répondu : le poste s'est attribué une adresse \textbf{APIPA} (169.254.0.0/16).}

\question{A6}{0,5 pt}{Quel protocole utilise \texttt{ping} ? Nommez les deux messages échangés et la couche OSI des adresses IP.}
\zone{1.5cm}{\textbf{ICMP} ; \textbf{Echo Request} / \textbf{Echo Reply} ; adresses IP en \textbf{couche 3 (Réseau)}.}

%==========================================================================
\newpage
\section{Partie B — Réalisation sous Packet Tracer \hfill /10}

\begin{cdc}
\entreprise{} vous confie son réseau : \textbf{trois réseaux} reliés par un \textbf{routeur unique Cisco 2911}
nommé \ip{\Rnom} (interfaces G0/0, G0/1, G0/2).
\begin{itemize}[leftmargin=12pt]
  \item \textbf{Réseau \resUn} : \textbf{4 postes} PC-\pUn1 à PC-\pUn4, commutateur \ip{SW-\pUn}, adressage statique,
        réseau \ip{\netUn.0/24}.
  \item \textbf{Réseau \resDeux} : \textbf{2 postes} PC-\pDeux1, PC-\pDeux2 et \textbf{1 imprimante} IMP-\pDeux,
        commutateur \ip{SW-\pDeux}, adressage statique, réseau de classe B \ip{\netDeux.0.0/16}.
  \item \textbf{Réseau Serveur} : serveur \ip{\srvNom} d'adresse \ip{\netSrv.\srvH/24}, relié \textbf{directement} au routeur.
  \item La passerelle de chaque réseau est sa \textbf{dernière adresse utilisable}. Tous les équipements communiquent
        entre eux et accèdent à la page web du serveur. Nommage imposé (« Display Name »).
\end{itemize}
\end{cdc}

\question{B1}{2 pts}{Complétez le plan d'adressage \textbf{avant} la réalisation.}
{\footnotesize\setlength{\tabcolsep}{3pt}\renewcommand{\arraystretch}{1.15}
\begin{tabularx}{\linewidth}{|l|C|C|C|C|C|C|}
  \hline
  \textbf{Réseau} & \textbf{Adresse réseau} & \textbf{Masque} & \textbf{1re @ utilisable} &
  \textbf{Dernière @ utilisable} & \textbf{Diffusion} & \textbf{Passerelle} \\ \hline
  \hl \resUn & \sol{\netUn.0} & \sol{255.255.255.0} & \sol{\netUn.1} & \sol{\netUn.254} & \sol{\netUn.255} & \sol{\netUn.254} \\ \hline
  \hl \resDeux & \sol{\netDeux.0.0} & \sol{255.255.0.0} & \sol{\netDeux.0.1} & \sol{\netDeux.255.254} & \sol{\netDeux.255.255} & \sol{\netDeux.255.254} \\ \hline
  \hl Serveur & \sol{\netSrv.0} & \sol{255.255.255.0} & \sol{\netSrv.1} & \sol{\netSrv.254} & \sol{\netSrv.255} & \sol{\netSrv.254} \\ \hline
\end{tabularx}}
\bareme{0,5 pt par ligne juste + 0,5 pt si la convention de passerelle est respectée partout.}

\question{B2}{2 pts}{Construisez la topologie (équipements, câblage, nommage) et relevez l'adressage configuré.}
{\footnotesize\setlength{\tabcolsep}{3pt}\renewcommand{\arraystretch}{1.1}
\begin{tabularx}{\linewidth}{|l|C|C|C||l|C|C|}
  \hline
  \textbf{Hôte} & \textbf{Adresse IP} & \textbf{Masque} & \textbf{Passerelle} & \textbf{Interface} & \textbf{Adresse IP} & \textbf{Masque} \\ \hline
  \hl PC-\pUn1 & \sol{\netUn.1} & \sol{255.255.255.0} & \sol{\netUn.254} & \Rnom{} G0/0 & \sol{\netUn.254} & \sol{255.255.255.0} \\ \hline
  \hl PC-\pDeux1 & \sol{\netDeux.0.1} & \sol{255.255.0.0} & \sol{\netDeux.255.254} & \Rnom{} G0/1 & \sol{\netDeux.255.254} & \sol{255.255.0.0} \\ \hline
  \hl IMP-\pDeux & \sol{\netDeux.0.10} & \sol{255.255.0.0} & \sol{\netDeux.255.254} & \Rnom{} G0/2 & \sol{\netSrv.254} & \sol{255.255.255.0} \\ \hline
  \hl \srvNom & \ip{\netSrv.\srvH} & \sol{255.255.255.0} & \sol{\netSrv.254} & \multicolumn{3}{c|}{\cellcolor{black!8}} \\ \hline
\end{tabularx}}
\bareme{Adresses d'hôtes libres (uniques, dans le bon réseau) ; affectation G0/x libre. Sur le .pkt : équipements 0,5 ;
câblage 1 (câbles droits vers les switchs et switch–routeur, \textbf{croisé} serveur–routeur, liens verts) ; nommage 0,5.}

\question{B3}{3 pts}{Configurez tous les équipements : IP, masques, passerelles des hôtes ; interfaces du routeur
(« Port Status : On »).}
\bareme{Sur le .pkt : interfaces routeur 1 pt ; IP/masques des hôtes 1 pt ; passerelles (imprimante et serveur compris)
1 pt. $-0,25$ par erreur, sans descendre sous 0 par item.}

\question{B4}{3 pts}{Réalisez les tests, notez la commande et le résultat, puis faites viser chaque test réussi.}
{\small\renewcommand{\arraystretch}{1.1}
\begin{tabularx}{\linewidth}{|c|l|l|C|c|c|}
  \hline
  \textbf{N°} & \textbf{Source} & \textbf{Destination} & \textbf{Commande / action} & \textbf{Résultat} & \textbf{Visa prof.} \\ \hline
  \hl 1 & PC-\pUn1 & PC-\pUn4 & \sol{ping \netUn.4} & \sol{réussi} & \hspace{1.4cm} \\ \hline
  \hl 2 & PC-\pDeux1 & IMP-\pDeux & \sol{ping \netDeux.0.10} & \sol{réussi} & \\ \hline
  \hl 3 & PC-\pUn2 & PC-\pDeux2 & \sol{ping \netDeux.0.2} & \sol{réussi} & \\ \hline
  \hl 4 & PC-\pDeux2 & \srvNom & \sol{ping \netSrv.\srvH} & \sol{réussi} & \\ \hline
  \hl 5 & PC-\pUn3 & page web du serveur & \sol{navigateur : {\NoAutoSpacing http://\netSrv.\srvH}} & \sol{page affichée} & \\ \hline
\end{tabularx}}
\bareme{0,5 pt par test visé + 0,5 pt pour B5. Le premier ping inter-réseaux peut échouer (résolution ARP) ; les suivants doivent réussir.}

\question{B5}{incl. B4}{Test n°3 : par quels équipements passe le paquet ? Pourquoi PC-\pUn2 l'envoie-t-il au routeur ?}
\zone{2.3cm}{PC-\pUn2 $\to$ SW-\pUn{} $\to$ \Rnom{} $\to$ SW-\pDeux{} $\to$ PC-\pDeux2. Grâce à son masque, PC-\pUn2 constate que
la destination (\netDeux.0.0/16) n'est \textbf{pas dans son réseau} (\netUn.0/24) : il envoie le paquet à sa
\textbf{passerelle par défaut}, qui le route.}

%==========================================================================
\newpage
\section{Partie C — Diagnostic d'un réseau en panne \hfill /4}

Un technicien constate qu'\textbf{aucun poste ne parvient à joindre le serveur SRV-C} et que certains postes
communiquent mal entre eux. Les câbles sont du bon type.

\begin{center}
\begin{minipage}{0.37\linewidth}\centering
\begin{tikzpicture}[font=\footnotesize,
  eq/.style={draw,fill=qbg,rounded corners=2pt,minimum width=9mm,minimum height=5mm,inner sep=2pt},
  every path/.style={thick}]
  \node[eq] (SW) {SW-1};
  \node[eq,right=9mm of SW] (R) {R1};
  \node[eq,right=7mm of R] (SRV) {SRV-C};
  \node[eq,above left=5mm and -3mm of SW] (A) {PC-A};
  \node[eq,above right=5mm and -3mm of SW] (B) {PC-B};
  \node[eq,below left=5mm and -3mm of SW] (C) {PC-C};
  \node[eq,below right=5mm and -3mm of SW] (D) {PC-D};
  \foreach \n in {A,B,C,D} \draw (\n) -- (SW);
  \draw (SW) -- node[above,font=\tiny]{G0/0} (R);
  \draw (R) -- node[above,font=\tiny]{G0/1} (SRV);
\end{tikzpicture}
\end{minipage}\hfill
\begin{minipage}{0.61\linewidth}\centering
{\small\renewcommand{\arraystretch}{1.05}\setlength{\tabcolsep}{4pt}
\begin{tabular}{lllll}
  \toprule
  \textbf{Équipement} & \textbf{Adresse IP} & \textbf{Masque} & \textbf{Passerelle} & \textbf{Port} \\
  \midrule
  \Crows
  \bottomrule
\end{tabular}}
\end{minipage}
\end{center}

\question{C1}{4 pts}{Relevez les \textbf{quatre anomalies}. Pour chacune : conséquence et correction.}
{\small\renewcommand{\arraystretch}{1.1}
\ifcorrige\def\hc{}\else\def\hc{\rule[-10mm]{0pt}{14mm}}\fi
\begin{tabularx}{\linewidth}{|c|>{\raggedright\arraybackslash}p{2cm}|L|L|L|}
  \hline
  \textbf{N°} & \textbf{Équipement} & \textbf{Anomalie} & \textbf{Conséquence} & \textbf{Correction} \\ \hline
  \ifcorrige\Csol\else
  \hc 1 & & & & \\ \hline
  \hc 2 & & & & \\ \hline
  \hc 3 & & & & \\ \hline
  \hc 4 & & & & \\ \hline
  \fi
\end{tabularx}}
\bareme{1 pt par anomalie (0,25 équipement, anomalie, conséquence, correction ; ordre indifférent). \Cbareme}

%==========================================================================
\section{Bonus — Extension sans fil \hfill +1}

\begin{cdc}Ajoutez au réseau \resUn{} un \textbf{AccessPoint-PT} relié à SW-\pUn{} (SSID \ip{\ssid}, \textbf{WPA2-PSK})
et un portable \ip{Laptop-\pUn} équipé d'un module \texttt{WPC300N}, en adressage \textbf{statique}.
Vérifiez qu'il joint \srvNom.\end{cdc}
{\small\renewcommand{\arraystretch}{1.1}
\begin{tabularx}{\linewidth}{|l|C|C|C|c|}
  \hline
  \textbf{Équipement} & \textbf{Adresse IP} & \textbf{Masque} & \textbf{Passerelle} & \textbf{Visa prof.} \\ \hline
  \hl Laptop-\pUn & \sol{\netUn.50} & \sol{255.255.255.0} & \sol{\netUn.254} & \hspace{1.4cm} \\ \hline
\end{tabularx}}
\bareme{Le point d'accès fait office de pont : le portable est dans \netUn.0/24. Bonus si associé en WPA2-PSK et ping
vers \netSrv.\srvH{} réussi. Note plafonnée à 20.}

\end{document}
"
%   pdflatex -jobname=Eval_Reseau_CorrigeA "\def\variante{A}\def\modecorrige{1}% Évaluation pratique — Réseaux locaux (BTS CIEL 1re année) — sujets A et B
% Compilation (4 PDF) :
%   pdflatex -jobname=Eval_Reseau_SujetA   "\def\variante{A}\def\modecorrige{0}% Évaluation pratique — Réseaux locaux (BTS CIEL 1re année) — sujets A et B
% Compilation (4 PDF) :
%   pdflatex -jobname=Eval_Reseau_SujetA   "\def\variante{A}\def\modecorrige{0}% Évaluation pratique — Réseaux locaux (BTS CIEL 1re année) — sujets A et B
% Compilation (4 PDF) :
%   pdflatex -jobname=Eval_Reseau_SujetA   "\def\variante{A}\def\modecorrige{0}\input{Eval_TP_Reseau_LAN.tex}"
%   pdflatex -jobname=Eval_Reseau_SujetB   "\def\variante{B}\def\modecorrige{0}\input{Eval_TP_Reseau_LAN.tex}"
%   pdflatex -jobname=Eval_Reseau_CorrigeA "\def\variante{A}\def\modecorrige{1}\input{Eval_TP_Reseau_LAN.tex}"
%   pdflatex -jobname=Eval_Reseau_CorrigeB "\def\variante{B}\def\modecorrige{1}\input{Eval_TP_Reseau_LAN.tex}"
\documentclass[10pt,a4paper]{article}

\usepackage[T1]{fontenc}
\usepackage[utf8]{inputenc}
\usepackage{lmodern}
\usepackage[french]{babel}
\usepackage[a4paper,hmargin=1.5cm,top=1.3cm,bottom=1.6cm,footskip=0.7cm]{geometry}
\usepackage[table]{xcolor}
\usepackage{booktabs,tabularx,array}
\usepackage{enumitem}
\usepackage{fancyhdr}
\usepackage[most]{tcolorbox}
\usepackage{tikz}
\usetikzlibrary{positioning}
\usepackage{titlesec}

\providecommand{\modecorrige}{0}
\providecommand{\variante}{A}
\newif\ifcorrige
\ifnum\modecorrige=1 \corrigetrue \else \corrigefalse \fi

\definecolor{ipred}{HTML}{C0392B}
\definecolor{cdcbar}{HTML}{A67C3D}
\definecolor{cdcbg}{HTML}{F4F1EA}
\definecolor{qbar}{HTML}{4A6B5D}
\definecolor{qbg}{HTML}{EDF1EF}
\definecolor{corr}{HTML}{1F5FA8}
\definecolor{corrbg}{HTML}{E8F0FA}
\definecolor{grisfoot}{HTML}{828282}

\newcommand{\ip}[1]{\textcolor{ipred}{\texttt{#1}}}
\newcommand{\pts}[1]{\hfill\mbox{\textbf{(#1)}}}
\newcommand{\sol}[1]{\ifcorrige\textcolor{corr}{#1}\fi}
\newcommand{\bareme}[1]{\ifcorrige\par{\footnotesize\color{corr}#1\par}\fi}
\newcommand{\hl}{\ifcorrige\rule[-1.6mm]{0pt}{5.4mm}\else\rule[-2.2mm]{0pt}{6.6mm}\fi}
\newcolumntype{C}{>{\centering\arraybackslash}X}
\newcolumntype{L}{>{\raggedright\arraybackslash}X}

\titleformat{\section}{\large\bfseries}{}{0pt}{}
\titlespacing*{\section}{0pt}{6pt}{3pt}

\tcbset{qstyle/.style={enhanced,boxrule=0pt,leftrule=2.5pt,arc=0pt,outer arc=0pt,
  left=5pt,right=5pt,top=2pt,bottom=2pt,before skip=4pt,after skip=3pt}}
\newtcolorbox{cdc}{qstyle,colback=cdcbg,colframe=cdcbar}
\newtcolorbox{quest}{qstyle,colback=qbg,colframe=qbar}

% Zone de réponse : vierge dans le sujet, remplie dans le corrigé
\newcommand{\zone}[2]{%
  \ifcorrige
    \begin{tcolorbox}[colback=corrbg,colframe=corr,boxrule=0.5pt,arc=1pt,left=4pt,right=4pt,top=1pt,bottom=1pt,
      before skip=2pt,after skip=2pt,fontupper=\footnotesize\color{corr}]#2\end{tcolorbox}%
  \else
    \begin{tcolorbox}[colback=white,colframe=black!45,boxrule=0.4pt,arc=1pt,height=#1,
      before skip=2pt,after skip=2pt]\end{tcolorbox}%
  \fi}

\newcommand{\question}[3]{\begin{quest}\textbf{#1.} #3\pts{#2}\end{quest}}

%==========================================================================
% Données propres à chaque variante
%==========================================================================
\if\variante A
  \def\entreprise{Le cabinet d'architectes \textbf{ArchiNova}}
  \def\Rnom{R-ArchiNova}
  \def\resUn{Bureau d'études}\def\pUn{BE}\def\netUn{192.168.40}
  \def\resDeux{Accueil}\def\pDeux{ACC}\def\netDeux{172.16}
  \def\srvNom{SRV-Intranet}\def\netSrv{192.168.200}\def\srvH{10}
  \def\ssid{ArchiNova-WiFi}
  \def\AIrows{%
    \hl\ip{192.168.40.25} & \sol{C} & \sol{255.255.255.0} & \sol{192.168.40.0} & \sol{192.168.40.255} & \sol{254} \\ \hline
    \hl\ip{172.16.5.10}   & \sol{B} & \sol{255.255.0.0}   & \sol{172.16.0.0}   & \sol{172.16.255.255} & \sol{65\,534} \\ \hline
    \hl\ip{10.3.2.1}      & \sol{A} & \sol{255.0.0.0}     & \sol{10.0.0.0}     & \sol{10.255.255.255} & \sol{16\,777\,214} \\ \hline}
  \def\apipa{169.254.12.7}
  \def\Cnet{192.168.60}
  \def\Crows{%
    PC-A & \ip{192.168.60.10} & \ip{255.255.255.0} & \ip{192.168.60.254} & — \\
    PC-B & \ip{192.168.60.20} & \ip{255.255.255.0} & \ip{192.168.60.1}   & — \\
    PC-C & \ip{192.168.61.30} & \ip{255.255.255.0} & \ip{192.168.60.254} & — \\
    PC-D & \ip{192.168.60.10} & \ip{255.255.255.0} & \ip{192.168.60.254} & — \\
    SRV-C & \ip{10.0.0.5}     & \ip{255.0.0.0}     & \ip{10.0.0.254}     & — \\
    R1 — G0/0 & \ip{192.168.60.254} & \ip{255.255.255.0} & — & On \\
    R1 — G0/1 & \ip{10.0.0.254}     & \ip{255.0.0.0}     & — & Off \\}
  \def\Csol{%
    \hl 1 & \sol{PC-B} & \sol{Passerelle .60.1 : pas l'adresse du routeur.} & \sol{Joint son réseau, mais aucun autre (serveur).} & \sol{Passerelle 192.168.60.254.} \\ \hline
    \hl 2 & \sol{PC-C} & \sol{IP hors du réseau 192.168.60.0/24.} & \sol{Ne joint ni les postes ni le routeur.} & \sol{Ex. 192.168.60.30.} \\ \hline
    \hl 3 & \sol{PC-D} & \sol{Même IP que PC-A (doublon).} & \sol{Conflit d'adresse : PC-A et PC-D perturbés.} & \sol{Adresse unique, ex. .60.40.} \\ \hline
    \hl 4 & \sol{R1 G0/1} & \sol{Interface désactivée (Off).} & \sol{Réseau du serveur injoignable pour tous.} & \sol{Port Status : On.} \\ \hline}
  \def\Cbareme{PC-A et SRV-C sont correctement configurés ; accepter PC-A cité à la place de PC-D pour le doublon.}
\else
  \def\entreprise{La clinique vétérinaire \textbf{VetoPlus}}
  \def\Rnom{R-VetoPlus}
  \def\resUn{Soins}\def\pUn{SOI}\def\netUn{192.168.70}
  \def\resDeux{Secrétariat}\def\pDeux{SEC}\def\netDeux{172.20}
  \def\srvNom{SRV-Dossiers}\def\netSrv{192.168.250}\def\srvH{20}
  \def\ssid{VetoPlus-WiFi}
  \def\AIrows{%
    \hl\ip{192.168.15.130} & \sol{C} & \sol{255.255.255.0} & \sol{192.168.15.0} & \sol{192.168.15.255} & \sol{254} \\ \hline
    \hl\ip{172.31.200.7}   & \sol{B} & \sol{255.255.0.0}   & \sol{172.31.0.0}   & \sol{172.31.255.255} & \sol{65\,534} \\ \hline
    \hl\ip{10.45.6.9}      & \sol{A} & \sol{255.0.0.0}     & \sol{10.0.0.0}     & \sol{10.255.255.255} & \sol{16\,777\,214} \\ \hline}
  \def\apipa{169.254.88.3}
  \def\Cnet{192.168.80}
  \def\Crows{%
    PC-A & \ip{192.168.8.11}   & \ip{255.255.255.0} & \ip{192.168.80.254} & — \\
    PC-B & \ip{192.168.80.12}  & \ip{255.255.255.0} & \ip{192.168.80.254} & — \\
    PC-C & \ip{192.168.80.13}  & \ip{255.255.255.0} & \ip{192.168.80.100} & — \\
    PC-D & \ip{192.168.80.12}  & \ip{255.255.255.0} & \ip{192.168.80.254} & — \\
    SRV-C & \ip{10.0.0.8}      & \ip{255.0.0.0}     & \ip{10.0.0.254}     & — \\
    R1 — G0/0 & \ip{192.168.80.254} & \ip{255.255.255.0} & — & Off \\
    R1 — G0/1 & \ip{10.0.0.254}     & \ip{255.0.0.0}     & — & On \\}
  \def\Csol{%
    \hl 1 & \sol{PC-A} & \sol{IP hors du réseau 192.168.80.0/24.} & \sol{Ne joint ni les postes ni le routeur.} & \sol{Ex. 192.168.80.11.} \\ \hline
    \hl 2 & \sol{PC-C} & \sol{Passerelle .80.100 : pas l'adresse du routeur.} & \sol{Joint son réseau, mais aucun autre (serveur).} & \sol{Passerelle 192.168.80.254.} \\ \hline
    \hl 3 & \sol{PC-D} & \sol{Même IP que PC-B (doublon).} & \sol{Conflit d'adresse : PC-B et PC-D perturbés.} & \sol{Adresse unique, ex. .80.14.} \\ \hline
    \hl 4 & \sol{R1 G0/0} & \sol{Interface côté LAN désactivée (Off).} & \sol{Aucun poste ne joint le routeur ni le serveur.} & \sol{Port Status : On.} \\ \hline}
  \def\Cbareme{PC-B et SRV-C sont correctement configurés ; accepter PC-B cité à la place de PC-D pour le doublon.}
\fi

\pagestyle{fancy}
\fancyhf{}
\renewcommand{\headrulewidth}{0pt}
\fancyfoot[L]{\footnotesize\itshape\color{grisfoot}Évaluation — Réseaux locaux — Sujet \variante\ifcorrige\ — CORRIGÉ\fi}
\fancyfoot[R]{\footnotesize\color{grisfoot}\thepage/3}

\setlist{itemsep=0pt,topsep=2pt,parsep=0pt}
\setlength{\parindent}{0pt}
\setlength{\parskip}{2pt}

\begin{document}

\begin{center}
  {\Large\bfseries Évaluation pratique — Réseaux locaux \quad
   \fcolorbox{black}{white}{\ SUJET \variante\ }\par}
  \vspace{1pt}
  {\bfseries Adressage IP, commutation et routage — BTS CIEL 1re année — durée : 2 heures\par}
  \ifcorrige\vspace{2pt}{\bfseries\color{corr}CORRIGÉ ET BARÈME — document enseignant}\par\fi
\end{center}
\vspace{-2pt}
{\renewcommand{\arraystretch}{1.35}
\begin{tabularx}{\linewidth}{|X|X|p{2.2cm}|p{2.6cm}|}
  \hline
  \textbf{Nom :} & \textbf{Prénom :} & \textbf{Poste n°} & \textbf{Note :} \hfill /20 \\ \hline
\end{tabularx}}

\textbf{Consignes :} travail \textbf{individuel}, documents et Internet \textbf{non autorisés}. Ce sujet est le
\textbf{document-réponse}. Fichier Packet Tracer : \ip{Nom\_Prenom\_EvalLAN\_\variante.pkt}, à enregistrer
régulièrement dans le dossier de rendu. Les « visas » sont apposés par le professeur, que vous appelez à chaque test réussi.
Durées conseillées : A 30 min — B 1 h 10 — C 20 min.

%==========================================================================
\section{Partie A — Questions de cours et calculs \hfill /6}

\question{A1}{2 pts}{Complétez le tableau en considérant le \textbf{masque par défaut} de la classe de chaque adresse.}
{\small\renewcommand{\arraystretch}{1.15}
\begin{tabularx}{\linewidth}{|l|c|C|C|C|c|}
  \hline
  \textbf{Adresse IP} & \textbf{Classe} & \textbf{Masque par défaut} & \textbf{Adresse réseau} &
  \textbf{Adresse de diffusion} & \textbf{Nb d'hôtes} \\ \hline
  \AIrows
\end{tabularx}}
\bareme{0,5 pt par ligne juste + 0,5 pt si les trois nombres d'hôtes sont justes ($2^{n}-2$, $n$ = nombre de bits d'hôte).}

\question{A2}{1 pt}{Indiquez le type de câble (\textbf{droit} ou \textbf{croisé}) pour chaque liaison.}
{\small\renewcommand{\arraystretch}{1.15}
\begin{tabularx}{\linewidth}{|X|c||X|c|}
  \hline
  \hl PC $\leftrightarrow$ commutateur & \makebox[1.6cm]{\sol{droit}} & Commutateur $\leftrightarrow$ routeur & \makebox[1.6cm]{\sol{droit}} \\ \hline
  \hl PC $\leftrightarrow$ PC (réseau P2P) & \sol{croisé} & Concentrateur $\leftrightarrow$ concentrateur & \sol{croisé} \\ \hline
\end{tabularx}}
\bareme{0,25 pt par case. Équipements de types différents : câble droit ; de même type : câble croisé.}

\question{A3}{1 pt}{Comment un \textbf{concentrateur (hub)} et un \textbf{commutateur (switch)} traitent-ils une trame
reçue ? Lequel limite les collisions, et pourquoi ?}
\zone{2.9cm}{Le \textbf{hub} (couche 1) répète la trame sur \textbf{tous ses autres ports} : un seul domaine de collision
partagé. Le \textbf{switch} (couche 2) lit l'adresse MAC de destination et, grâce à sa \textbf{table MAC}, ne l'envoie
\textbf{que sur le port du destinataire} : un domaine de collision par port, c'est lui qui limite les collisions.
\emph{0,5 pt hub, 0,5 pt switch.}}

\question{A4}{1 pt}{Qu'est-ce que la \textbf{passerelle par défaut} d'un poste ? Quelle condition son adresse doit-elle respecter ?}
\zone{2.4cm}{Adresse IP de l'interface du \textbf{routeur} à laquelle le poste envoie les paquets destinés à un
\textbf{autre réseau}. Elle doit \textbf{appartenir au même réseau} que le poste. \emph{0,5 + 0,5 pt.}}

\question{A5}{0,5 pt}{Un portable configuré en DHCP affiche \ip{\apipa} / \ip{255.255.0.0}. Que s'est-il passé ?}
\zone{1.5cm}{Aucun \textbf{serveur DHCP} n'a répondu : le poste s'est attribué une adresse \textbf{APIPA} (169.254.0.0/16).}

\question{A6}{0,5 pt}{Quel protocole utilise \texttt{ping} ? Nommez les deux messages échangés et la couche OSI des adresses IP.}
\zone{1.5cm}{\textbf{ICMP} ; \textbf{Echo Request} / \textbf{Echo Reply} ; adresses IP en \textbf{couche 3 (Réseau)}.}

%==========================================================================
\newpage
\section{Partie B — Réalisation sous Packet Tracer \hfill /10}

\begin{cdc}
\entreprise{} vous confie son réseau : \textbf{trois réseaux} reliés par un \textbf{routeur unique Cisco 2911}
nommé \ip{\Rnom} (interfaces G0/0, G0/1, G0/2).
\begin{itemize}[leftmargin=12pt]
  \item \textbf{Réseau \resUn} : \textbf{4 postes} PC-\pUn1 à PC-\pUn4, commutateur \ip{SW-\pUn}, adressage statique,
        réseau \ip{\netUn.0/24}.
  \item \textbf{Réseau \resDeux} : \textbf{2 postes} PC-\pDeux1, PC-\pDeux2 et \textbf{1 imprimante} IMP-\pDeux,
        commutateur \ip{SW-\pDeux}, adressage statique, réseau de classe B \ip{\netDeux.0.0/16}.
  \item \textbf{Réseau Serveur} : serveur \ip{\srvNom} d'adresse \ip{\netSrv.\srvH/24}, relié \textbf{directement} au routeur.
  \item La passerelle de chaque réseau est sa \textbf{dernière adresse utilisable}. Tous les équipements communiquent
        entre eux et accèdent à la page web du serveur. Nommage imposé (« Display Name »).
\end{itemize}
\end{cdc}

\question{B1}{2 pts}{Complétez le plan d'adressage \textbf{avant} la réalisation.}
{\footnotesize\setlength{\tabcolsep}{3pt}\renewcommand{\arraystretch}{1.15}
\begin{tabularx}{\linewidth}{|l|C|C|C|C|C|C|}
  \hline
  \textbf{Réseau} & \textbf{Adresse réseau} & \textbf{Masque} & \textbf{1re @ utilisable} &
  \textbf{Dernière @ utilisable} & \textbf{Diffusion} & \textbf{Passerelle} \\ \hline
  \hl \resUn & \sol{\netUn.0} & \sol{255.255.255.0} & \sol{\netUn.1} & \sol{\netUn.254} & \sol{\netUn.255} & \sol{\netUn.254} \\ \hline
  \hl \resDeux & \sol{\netDeux.0.0} & \sol{255.255.0.0} & \sol{\netDeux.0.1} & \sol{\netDeux.255.254} & \sol{\netDeux.255.255} & \sol{\netDeux.255.254} \\ \hline
  \hl Serveur & \sol{\netSrv.0} & \sol{255.255.255.0} & \sol{\netSrv.1} & \sol{\netSrv.254} & \sol{\netSrv.255} & \sol{\netSrv.254} \\ \hline
\end{tabularx}}
\bareme{0,5 pt par ligne juste + 0,5 pt si la convention de passerelle est respectée partout.}

\question{B2}{2 pts}{Construisez la topologie (équipements, câblage, nommage) et relevez l'adressage configuré.}
{\footnotesize\setlength{\tabcolsep}{3pt}\renewcommand{\arraystretch}{1.1}
\begin{tabularx}{\linewidth}{|l|C|C|C||l|C|C|}
  \hline
  \textbf{Hôte} & \textbf{Adresse IP} & \textbf{Masque} & \textbf{Passerelle} & \textbf{Interface} & \textbf{Adresse IP} & \textbf{Masque} \\ \hline
  \hl PC-\pUn1 & \sol{\netUn.1} & \sol{255.255.255.0} & \sol{\netUn.254} & \Rnom{} G0/0 & \sol{\netUn.254} & \sol{255.255.255.0} \\ \hline
  \hl PC-\pDeux1 & \sol{\netDeux.0.1} & \sol{255.255.0.0} & \sol{\netDeux.255.254} & \Rnom{} G0/1 & \sol{\netDeux.255.254} & \sol{255.255.0.0} \\ \hline
  \hl IMP-\pDeux & \sol{\netDeux.0.10} & \sol{255.255.0.0} & \sol{\netDeux.255.254} & \Rnom{} G0/2 & \sol{\netSrv.254} & \sol{255.255.255.0} \\ \hline
  \hl \srvNom & \ip{\netSrv.\srvH} & \sol{255.255.255.0} & \sol{\netSrv.254} & \multicolumn{3}{c|}{\cellcolor{black!8}} \\ \hline
\end{tabularx}}
\bareme{Adresses d'hôtes libres (uniques, dans le bon réseau) ; affectation G0/x libre. Sur le .pkt : équipements 0,5 ;
câblage 1 (câbles droits vers les switchs et switch–routeur, \textbf{croisé} serveur–routeur, liens verts) ; nommage 0,5.}

\question{B3}{3 pts}{Configurez tous les équipements : IP, masques, passerelles des hôtes ; interfaces du routeur
(« Port Status : On »).}
\bareme{Sur le .pkt : interfaces routeur 1 pt ; IP/masques des hôtes 1 pt ; passerelles (imprimante et serveur compris)
1 pt. $-0,25$ par erreur, sans descendre sous 0 par item.}

\question{B4}{3 pts}{Réalisez les tests, notez la commande et le résultat, puis faites viser chaque test réussi.}
{\small\renewcommand{\arraystretch}{1.1}
\begin{tabularx}{\linewidth}{|c|l|l|C|c|c|}
  \hline
  \textbf{N°} & \textbf{Source} & \textbf{Destination} & \textbf{Commande / action} & \textbf{Résultat} & \textbf{Visa prof.} \\ \hline
  \hl 1 & PC-\pUn1 & PC-\pUn4 & \sol{ping \netUn.4} & \sol{réussi} & \hspace{1.4cm} \\ \hline
  \hl 2 & PC-\pDeux1 & IMP-\pDeux & \sol{ping \netDeux.0.10} & \sol{réussi} & \\ \hline
  \hl 3 & PC-\pUn2 & PC-\pDeux2 & \sol{ping \netDeux.0.2} & \sol{réussi} & \\ \hline
  \hl 4 & PC-\pDeux2 & \srvNom & \sol{ping \netSrv.\srvH} & \sol{réussi} & \\ \hline
  \hl 5 & PC-\pUn3 & page web du serveur & \sol{navigateur : {\NoAutoSpacing http://\netSrv.\srvH}} & \sol{page affichée} & \\ \hline
\end{tabularx}}
\bareme{0,5 pt par test visé + 0,5 pt pour B5. Le premier ping inter-réseaux peut échouer (résolution ARP) ; les suivants doivent réussir.}

\question{B5}{incl. B4}{Test n°3 : par quels équipements passe le paquet ? Pourquoi PC-\pUn2 l'envoie-t-il au routeur ?}
\zone{2.3cm}{PC-\pUn2 $\to$ SW-\pUn{} $\to$ \Rnom{} $\to$ SW-\pDeux{} $\to$ PC-\pDeux2. Grâce à son masque, PC-\pUn2 constate que
la destination (\netDeux.0.0/16) n'est \textbf{pas dans son réseau} (\netUn.0/24) : il envoie le paquet à sa
\textbf{passerelle par défaut}, qui le route.}

%==========================================================================
\newpage
\section{Partie C — Diagnostic d'un réseau en panne \hfill /4}

Un technicien constate qu'\textbf{aucun poste ne parvient à joindre le serveur SRV-C} et que certains postes
communiquent mal entre eux. Les câbles sont du bon type.

\begin{center}
\begin{minipage}{0.37\linewidth}\centering
\begin{tikzpicture}[font=\footnotesize,
  eq/.style={draw,fill=qbg,rounded corners=2pt,minimum width=9mm,minimum height=5mm,inner sep=2pt},
  every path/.style={thick}]
  \node[eq] (SW) {SW-1};
  \node[eq,right=9mm of SW] (R) {R1};
  \node[eq,right=7mm of R] (SRV) {SRV-C};
  \node[eq,above left=5mm and -3mm of SW] (A) {PC-A};
  \node[eq,above right=5mm and -3mm of SW] (B) {PC-B};
  \node[eq,below left=5mm and -3mm of SW] (C) {PC-C};
  \node[eq,below right=5mm and -3mm of SW] (D) {PC-D};
  \foreach \n in {A,B,C,D} \draw (\n) -- (SW);
  \draw (SW) -- node[above,font=\tiny]{G0/0} (R);
  \draw (R) -- node[above,font=\tiny]{G0/1} (SRV);
\end{tikzpicture}
\end{minipage}\hfill
\begin{minipage}{0.61\linewidth}\centering
{\small\renewcommand{\arraystretch}{1.05}\setlength{\tabcolsep}{4pt}
\begin{tabular}{lllll}
  \toprule
  \textbf{Équipement} & \textbf{Adresse IP} & \textbf{Masque} & \textbf{Passerelle} & \textbf{Port} \\
  \midrule
  \Crows
  \bottomrule
\end{tabular}}
\end{minipage}
\end{center}

\question{C1}{4 pts}{Relevez les \textbf{quatre anomalies}. Pour chacune : conséquence et correction.}
{\small\renewcommand{\arraystretch}{1.1}
\ifcorrige\def\hc{}\else\def\hc{\rule[-10mm]{0pt}{14mm}}\fi
\begin{tabularx}{\linewidth}{|c|>{\raggedright\arraybackslash}p{2cm}|L|L|L|}
  \hline
  \textbf{N°} & \textbf{Équipement} & \textbf{Anomalie} & \textbf{Conséquence} & \textbf{Correction} \\ \hline
  \ifcorrige\Csol\else
  \hc 1 & & & & \\ \hline
  \hc 2 & & & & \\ \hline
  \hc 3 & & & & \\ \hline
  \hc 4 & & & & \\ \hline
  \fi
\end{tabularx}}
\bareme{1 pt par anomalie (0,25 équipement, anomalie, conséquence, correction ; ordre indifférent). \Cbareme}

%==========================================================================
\section{Bonus — Extension sans fil \hfill +1}

\begin{cdc}Ajoutez au réseau \resUn{} un \textbf{AccessPoint-PT} relié à SW-\pUn{} (SSID \ip{\ssid}, \textbf{WPA2-PSK})
et un portable \ip{Laptop-\pUn} équipé d'un module \texttt{WPC300N}, en adressage \textbf{statique}.
Vérifiez qu'il joint \srvNom.\end{cdc}
{\small\renewcommand{\arraystretch}{1.1}
\begin{tabularx}{\linewidth}{|l|C|C|C|c|}
  \hline
  \textbf{Équipement} & \textbf{Adresse IP} & \textbf{Masque} & \textbf{Passerelle} & \textbf{Visa prof.} \\ \hline
  \hl Laptop-\pUn & \sol{\netUn.50} & \sol{255.255.255.0} & \sol{\netUn.254} & \hspace{1.4cm} \\ \hline
\end{tabularx}}
\bareme{Le point d'accès fait office de pont : le portable est dans \netUn.0/24. Bonus si associé en WPA2-PSK et ping
vers \netSrv.\srvH{} réussi. Note plafonnée à 20.}

\end{document}
"
%   pdflatex -jobname=Eval_Reseau_SujetB   "\def\variante{B}\def\modecorrige{0}% Évaluation pratique — Réseaux locaux (BTS CIEL 1re année) — sujets A et B
% Compilation (4 PDF) :
%   pdflatex -jobname=Eval_Reseau_SujetA   "\def\variante{A}\def\modecorrige{0}\input{Eval_TP_Reseau_LAN.tex}"
%   pdflatex -jobname=Eval_Reseau_SujetB   "\def\variante{B}\def\modecorrige{0}\input{Eval_TP_Reseau_LAN.tex}"
%   pdflatex -jobname=Eval_Reseau_CorrigeA "\def\variante{A}\def\modecorrige{1}\input{Eval_TP_Reseau_LAN.tex}"
%   pdflatex -jobname=Eval_Reseau_CorrigeB "\def\variante{B}\def\modecorrige{1}\input{Eval_TP_Reseau_LAN.tex}"
\documentclass[10pt,a4paper]{article}

\usepackage[T1]{fontenc}
\usepackage[utf8]{inputenc}
\usepackage{lmodern}
\usepackage[french]{babel}
\usepackage[a4paper,hmargin=1.5cm,top=1.3cm,bottom=1.6cm,footskip=0.7cm]{geometry}
\usepackage[table]{xcolor}
\usepackage{booktabs,tabularx,array}
\usepackage{enumitem}
\usepackage{fancyhdr}
\usepackage[most]{tcolorbox}
\usepackage{tikz}
\usetikzlibrary{positioning}
\usepackage{titlesec}

\providecommand{\modecorrige}{0}
\providecommand{\variante}{A}
\newif\ifcorrige
\ifnum\modecorrige=1 \corrigetrue \else \corrigefalse \fi

\definecolor{ipred}{HTML}{C0392B}
\definecolor{cdcbar}{HTML}{A67C3D}
\definecolor{cdcbg}{HTML}{F4F1EA}
\definecolor{qbar}{HTML}{4A6B5D}
\definecolor{qbg}{HTML}{EDF1EF}
\definecolor{corr}{HTML}{1F5FA8}
\definecolor{corrbg}{HTML}{E8F0FA}
\definecolor{grisfoot}{HTML}{828282}

\newcommand{\ip}[1]{\textcolor{ipred}{\texttt{#1}}}
\newcommand{\pts}[1]{\hfill\mbox{\textbf{(#1)}}}
\newcommand{\sol}[1]{\ifcorrige\textcolor{corr}{#1}\fi}
\newcommand{\bareme}[1]{\ifcorrige\par{\footnotesize\color{corr}#1\par}\fi}
\newcommand{\hl}{\ifcorrige\rule[-1.6mm]{0pt}{5.4mm}\else\rule[-2.2mm]{0pt}{6.6mm}\fi}
\newcolumntype{C}{>{\centering\arraybackslash}X}
\newcolumntype{L}{>{\raggedright\arraybackslash}X}

\titleformat{\section}{\large\bfseries}{}{0pt}{}
\titlespacing*{\section}{0pt}{6pt}{3pt}

\tcbset{qstyle/.style={enhanced,boxrule=0pt,leftrule=2.5pt,arc=0pt,outer arc=0pt,
  left=5pt,right=5pt,top=2pt,bottom=2pt,before skip=4pt,after skip=3pt}}
\newtcolorbox{cdc}{qstyle,colback=cdcbg,colframe=cdcbar}
\newtcolorbox{quest}{qstyle,colback=qbg,colframe=qbar}

% Zone de réponse : vierge dans le sujet, remplie dans le corrigé
\newcommand{\zone}[2]{%
  \ifcorrige
    \begin{tcolorbox}[colback=corrbg,colframe=corr,boxrule=0.5pt,arc=1pt,left=4pt,right=4pt,top=1pt,bottom=1pt,
      before skip=2pt,after skip=2pt,fontupper=\footnotesize\color{corr}]#2\end{tcolorbox}%
  \else
    \begin{tcolorbox}[colback=white,colframe=black!45,boxrule=0.4pt,arc=1pt,height=#1,
      before skip=2pt,after skip=2pt]\end{tcolorbox}%
  \fi}

\newcommand{\question}[3]{\begin{quest}\textbf{#1.} #3\pts{#2}\end{quest}}

%==========================================================================
% Données propres à chaque variante
%==========================================================================
\if\variante A
  \def\entreprise{Le cabinet d'architectes \textbf{ArchiNova}}
  \def\Rnom{R-ArchiNova}
  \def\resUn{Bureau d'études}\def\pUn{BE}\def\netUn{192.168.40}
  \def\resDeux{Accueil}\def\pDeux{ACC}\def\netDeux{172.16}
  \def\srvNom{SRV-Intranet}\def\netSrv{192.168.200}\def\srvH{10}
  \def\ssid{ArchiNova-WiFi}
  \def\AIrows{%
    \hl\ip{192.168.40.25} & \sol{C} & \sol{255.255.255.0} & \sol{192.168.40.0} & \sol{192.168.40.255} & \sol{254} \\ \hline
    \hl\ip{172.16.5.10}   & \sol{B} & \sol{255.255.0.0}   & \sol{172.16.0.0}   & \sol{172.16.255.255} & \sol{65\,534} \\ \hline
    \hl\ip{10.3.2.1}      & \sol{A} & \sol{255.0.0.0}     & \sol{10.0.0.0}     & \sol{10.255.255.255} & \sol{16\,777\,214} \\ \hline}
  \def\apipa{169.254.12.7}
  \def\Cnet{192.168.60}
  \def\Crows{%
    PC-A & \ip{192.168.60.10} & \ip{255.255.255.0} & \ip{192.168.60.254} & — \\
    PC-B & \ip{192.168.60.20} & \ip{255.255.255.0} & \ip{192.168.60.1}   & — \\
    PC-C & \ip{192.168.61.30} & \ip{255.255.255.0} & \ip{192.168.60.254} & — \\
    PC-D & \ip{192.168.60.10} & \ip{255.255.255.0} & \ip{192.168.60.254} & — \\
    SRV-C & \ip{10.0.0.5}     & \ip{255.0.0.0}     & \ip{10.0.0.254}     & — \\
    R1 — G0/0 & \ip{192.168.60.254} & \ip{255.255.255.0} & — & On \\
    R1 — G0/1 & \ip{10.0.0.254}     & \ip{255.0.0.0}     & — & Off \\}
  \def\Csol{%
    \hl 1 & \sol{PC-B} & \sol{Passerelle .60.1 : pas l'adresse du routeur.} & \sol{Joint son réseau, mais aucun autre (serveur).} & \sol{Passerelle 192.168.60.254.} \\ \hline
    \hl 2 & \sol{PC-C} & \sol{IP hors du réseau 192.168.60.0/24.} & \sol{Ne joint ni les postes ni le routeur.} & \sol{Ex. 192.168.60.30.} \\ \hline
    \hl 3 & \sol{PC-D} & \sol{Même IP que PC-A (doublon).} & \sol{Conflit d'adresse : PC-A et PC-D perturbés.} & \sol{Adresse unique, ex. .60.40.} \\ \hline
    \hl 4 & \sol{R1 G0/1} & \sol{Interface désactivée (Off).} & \sol{Réseau du serveur injoignable pour tous.} & \sol{Port Status : On.} \\ \hline}
  \def\Cbareme{PC-A et SRV-C sont correctement configurés ; accepter PC-A cité à la place de PC-D pour le doublon.}
\else
  \def\entreprise{La clinique vétérinaire \textbf{VetoPlus}}
  \def\Rnom{R-VetoPlus}
  \def\resUn{Soins}\def\pUn{SOI}\def\netUn{192.168.70}
  \def\resDeux{Secrétariat}\def\pDeux{SEC}\def\netDeux{172.20}
  \def\srvNom{SRV-Dossiers}\def\netSrv{192.168.250}\def\srvH{20}
  \def\ssid{VetoPlus-WiFi}
  \def\AIrows{%
    \hl\ip{192.168.15.130} & \sol{C} & \sol{255.255.255.0} & \sol{192.168.15.0} & \sol{192.168.15.255} & \sol{254} \\ \hline
    \hl\ip{172.31.200.7}   & \sol{B} & \sol{255.255.0.0}   & \sol{172.31.0.0}   & \sol{172.31.255.255} & \sol{65\,534} \\ \hline
    \hl\ip{10.45.6.9}      & \sol{A} & \sol{255.0.0.0}     & \sol{10.0.0.0}     & \sol{10.255.255.255} & \sol{16\,777\,214} \\ \hline}
  \def\apipa{169.254.88.3}
  \def\Cnet{192.168.80}
  \def\Crows{%
    PC-A & \ip{192.168.8.11}   & \ip{255.255.255.0} & \ip{192.168.80.254} & — \\
    PC-B & \ip{192.168.80.12}  & \ip{255.255.255.0} & \ip{192.168.80.254} & — \\
    PC-C & \ip{192.168.80.13}  & \ip{255.255.255.0} & \ip{192.168.80.100} & — \\
    PC-D & \ip{192.168.80.12}  & \ip{255.255.255.0} & \ip{192.168.80.254} & — \\
    SRV-C & \ip{10.0.0.8}      & \ip{255.0.0.0}     & \ip{10.0.0.254}     & — \\
    R1 — G0/0 & \ip{192.168.80.254} & \ip{255.255.255.0} & — & Off \\
    R1 — G0/1 & \ip{10.0.0.254}     & \ip{255.0.0.0}     & — & On \\}
  \def\Csol{%
    \hl 1 & \sol{PC-A} & \sol{IP hors du réseau 192.168.80.0/24.} & \sol{Ne joint ni les postes ni le routeur.} & \sol{Ex. 192.168.80.11.} \\ \hline
    \hl 2 & \sol{PC-C} & \sol{Passerelle .80.100 : pas l'adresse du routeur.} & \sol{Joint son réseau, mais aucun autre (serveur).} & \sol{Passerelle 192.168.80.254.} \\ \hline
    \hl 3 & \sol{PC-D} & \sol{Même IP que PC-B (doublon).} & \sol{Conflit d'adresse : PC-B et PC-D perturbés.} & \sol{Adresse unique, ex. .80.14.} \\ \hline
    \hl 4 & \sol{R1 G0/0} & \sol{Interface côté LAN désactivée (Off).} & \sol{Aucun poste ne joint le routeur ni le serveur.} & \sol{Port Status : On.} \\ \hline}
  \def\Cbareme{PC-B et SRV-C sont correctement configurés ; accepter PC-B cité à la place de PC-D pour le doublon.}
\fi

\pagestyle{fancy}
\fancyhf{}
\renewcommand{\headrulewidth}{0pt}
\fancyfoot[L]{\footnotesize\itshape\color{grisfoot}Évaluation — Réseaux locaux — Sujet \variante\ifcorrige\ — CORRIGÉ\fi}
\fancyfoot[R]{\footnotesize\color{grisfoot}\thepage/3}

\setlist{itemsep=0pt,topsep=2pt,parsep=0pt}
\setlength{\parindent}{0pt}
\setlength{\parskip}{2pt}

\begin{document}

\begin{center}
  {\Large\bfseries Évaluation pratique — Réseaux locaux \quad
   \fcolorbox{black}{white}{\ SUJET \variante\ }\par}
  \vspace{1pt}
  {\bfseries Adressage IP, commutation et routage — BTS CIEL 1re année — durée : 2 heures\par}
  \ifcorrige\vspace{2pt}{\bfseries\color{corr}CORRIGÉ ET BARÈME — document enseignant}\par\fi
\end{center}
\vspace{-2pt}
{\renewcommand{\arraystretch}{1.35}
\begin{tabularx}{\linewidth}{|X|X|p{2.2cm}|p{2.6cm}|}
  \hline
  \textbf{Nom :} & \textbf{Prénom :} & \textbf{Poste n°} & \textbf{Note :} \hfill /20 \\ \hline
\end{tabularx}}

\textbf{Consignes :} travail \textbf{individuel}, documents et Internet \textbf{non autorisés}. Ce sujet est le
\textbf{document-réponse}. Fichier Packet Tracer : \ip{Nom\_Prenom\_EvalLAN\_\variante.pkt}, à enregistrer
régulièrement dans le dossier de rendu. Les « visas » sont apposés par le professeur, que vous appelez à chaque test réussi.
Durées conseillées : A 30 min — B 1 h 10 — C 20 min.

%==========================================================================
\section{Partie A — Questions de cours et calculs \hfill /6}

\question{A1}{2 pts}{Complétez le tableau en considérant le \textbf{masque par défaut} de la classe de chaque adresse.}
{\small\renewcommand{\arraystretch}{1.15}
\begin{tabularx}{\linewidth}{|l|c|C|C|C|c|}
  \hline
  \textbf{Adresse IP} & \textbf{Classe} & \textbf{Masque par défaut} & \textbf{Adresse réseau} &
  \textbf{Adresse de diffusion} & \textbf{Nb d'hôtes} \\ \hline
  \AIrows
\end{tabularx}}
\bareme{0,5 pt par ligne juste + 0,5 pt si les trois nombres d'hôtes sont justes ($2^{n}-2$, $n$ = nombre de bits d'hôte).}

\question{A2}{1 pt}{Indiquez le type de câble (\textbf{droit} ou \textbf{croisé}) pour chaque liaison.}
{\small\renewcommand{\arraystretch}{1.15}
\begin{tabularx}{\linewidth}{|X|c||X|c|}
  \hline
  \hl PC $\leftrightarrow$ commutateur & \makebox[1.6cm]{\sol{droit}} & Commutateur $\leftrightarrow$ routeur & \makebox[1.6cm]{\sol{droit}} \\ \hline
  \hl PC $\leftrightarrow$ PC (réseau P2P) & \sol{croisé} & Concentrateur $\leftrightarrow$ concentrateur & \sol{croisé} \\ \hline
\end{tabularx}}
\bareme{0,25 pt par case. Équipements de types différents : câble droit ; de même type : câble croisé.}

\question{A3}{1 pt}{Comment un \textbf{concentrateur (hub)} et un \textbf{commutateur (switch)} traitent-ils une trame
reçue ? Lequel limite les collisions, et pourquoi ?}
\zone{2.9cm}{Le \textbf{hub} (couche 1) répète la trame sur \textbf{tous ses autres ports} : un seul domaine de collision
partagé. Le \textbf{switch} (couche 2) lit l'adresse MAC de destination et, grâce à sa \textbf{table MAC}, ne l'envoie
\textbf{que sur le port du destinataire} : un domaine de collision par port, c'est lui qui limite les collisions.
\emph{0,5 pt hub, 0,5 pt switch.}}

\question{A4}{1 pt}{Qu'est-ce que la \textbf{passerelle par défaut} d'un poste ? Quelle condition son adresse doit-elle respecter ?}
\zone{2.4cm}{Adresse IP de l'interface du \textbf{routeur} à laquelle le poste envoie les paquets destinés à un
\textbf{autre réseau}. Elle doit \textbf{appartenir au même réseau} que le poste. \emph{0,5 + 0,5 pt.}}

\question{A5}{0,5 pt}{Un portable configuré en DHCP affiche \ip{\apipa} / \ip{255.255.0.0}. Que s'est-il passé ?}
\zone{1.5cm}{Aucun \textbf{serveur DHCP} n'a répondu : le poste s'est attribué une adresse \textbf{APIPA} (169.254.0.0/16).}

\question{A6}{0,5 pt}{Quel protocole utilise \texttt{ping} ? Nommez les deux messages échangés et la couche OSI des adresses IP.}
\zone{1.5cm}{\textbf{ICMP} ; \textbf{Echo Request} / \textbf{Echo Reply} ; adresses IP en \textbf{couche 3 (Réseau)}.}

%==========================================================================
\newpage
\section{Partie B — Réalisation sous Packet Tracer \hfill /10}

\begin{cdc}
\entreprise{} vous confie son réseau : \textbf{trois réseaux} reliés par un \textbf{routeur unique Cisco 2911}
nommé \ip{\Rnom} (interfaces G0/0, G0/1, G0/2).
\begin{itemize}[leftmargin=12pt]
  \item \textbf{Réseau \resUn} : \textbf{4 postes} PC-\pUn1 à PC-\pUn4, commutateur \ip{SW-\pUn}, adressage statique,
        réseau \ip{\netUn.0/24}.
  \item \textbf{Réseau \resDeux} : \textbf{2 postes} PC-\pDeux1, PC-\pDeux2 et \textbf{1 imprimante} IMP-\pDeux,
        commutateur \ip{SW-\pDeux}, adressage statique, réseau de classe B \ip{\netDeux.0.0/16}.
  \item \textbf{Réseau Serveur} : serveur \ip{\srvNom} d'adresse \ip{\netSrv.\srvH/24}, relié \textbf{directement} au routeur.
  \item La passerelle de chaque réseau est sa \textbf{dernière adresse utilisable}. Tous les équipements communiquent
        entre eux et accèdent à la page web du serveur. Nommage imposé (« Display Name »).
\end{itemize}
\end{cdc}

\question{B1}{2 pts}{Complétez le plan d'adressage \textbf{avant} la réalisation.}
{\footnotesize\setlength{\tabcolsep}{3pt}\renewcommand{\arraystretch}{1.15}
\begin{tabularx}{\linewidth}{|l|C|C|C|C|C|C|}
  \hline
  \textbf{Réseau} & \textbf{Adresse réseau} & \textbf{Masque} & \textbf{1re @ utilisable} &
  \textbf{Dernière @ utilisable} & \textbf{Diffusion} & \textbf{Passerelle} \\ \hline
  \hl \resUn & \sol{\netUn.0} & \sol{255.255.255.0} & \sol{\netUn.1} & \sol{\netUn.254} & \sol{\netUn.255} & \sol{\netUn.254} \\ \hline
  \hl \resDeux & \sol{\netDeux.0.0} & \sol{255.255.0.0} & \sol{\netDeux.0.1} & \sol{\netDeux.255.254} & \sol{\netDeux.255.255} & \sol{\netDeux.255.254} \\ \hline
  \hl Serveur & \sol{\netSrv.0} & \sol{255.255.255.0} & \sol{\netSrv.1} & \sol{\netSrv.254} & \sol{\netSrv.255} & \sol{\netSrv.254} \\ \hline
\end{tabularx}}
\bareme{0,5 pt par ligne juste + 0,5 pt si la convention de passerelle est respectée partout.}

\question{B2}{2 pts}{Construisez la topologie (équipements, câblage, nommage) et relevez l'adressage configuré.}
{\footnotesize\setlength{\tabcolsep}{3pt}\renewcommand{\arraystretch}{1.1}
\begin{tabularx}{\linewidth}{|l|C|C|C||l|C|C|}
  \hline
  \textbf{Hôte} & \textbf{Adresse IP} & \textbf{Masque} & \textbf{Passerelle} & \textbf{Interface} & \textbf{Adresse IP} & \textbf{Masque} \\ \hline
  \hl PC-\pUn1 & \sol{\netUn.1} & \sol{255.255.255.0} & \sol{\netUn.254} & \Rnom{} G0/0 & \sol{\netUn.254} & \sol{255.255.255.0} \\ \hline
  \hl PC-\pDeux1 & \sol{\netDeux.0.1} & \sol{255.255.0.0} & \sol{\netDeux.255.254} & \Rnom{} G0/1 & \sol{\netDeux.255.254} & \sol{255.255.0.0} \\ \hline
  \hl IMP-\pDeux & \sol{\netDeux.0.10} & \sol{255.255.0.0} & \sol{\netDeux.255.254} & \Rnom{} G0/2 & \sol{\netSrv.254} & \sol{255.255.255.0} \\ \hline
  \hl \srvNom & \ip{\netSrv.\srvH} & \sol{255.255.255.0} & \sol{\netSrv.254} & \multicolumn{3}{c|}{\cellcolor{black!8}} \\ \hline
\end{tabularx}}
\bareme{Adresses d'hôtes libres (uniques, dans le bon réseau) ; affectation G0/x libre. Sur le .pkt : équipements 0,5 ;
câblage 1 (câbles droits vers les switchs et switch–routeur, \textbf{croisé} serveur–routeur, liens verts) ; nommage 0,5.}

\question{B3}{3 pts}{Configurez tous les équipements : IP, masques, passerelles des hôtes ; interfaces du routeur
(« Port Status : On »).}
\bareme{Sur le .pkt : interfaces routeur 1 pt ; IP/masques des hôtes 1 pt ; passerelles (imprimante et serveur compris)
1 pt. $-0,25$ par erreur, sans descendre sous 0 par item.}

\question{B4}{3 pts}{Réalisez les tests, notez la commande et le résultat, puis faites viser chaque test réussi.}
{\small\renewcommand{\arraystretch}{1.1}
\begin{tabularx}{\linewidth}{|c|l|l|C|c|c|}
  \hline
  \textbf{N°} & \textbf{Source} & \textbf{Destination} & \textbf{Commande / action} & \textbf{Résultat} & \textbf{Visa prof.} \\ \hline
  \hl 1 & PC-\pUn1 & PC-\pUn4 & \sol{ping \netUn.4} & \sol{réussi} & \hspace{1.4cm} \\ \hline
  \hl 2 & PC-\pDeux1 & IMP-\pDeux & \sol{ping \netDeux.0.10} & \sol{réussi} & \\ \hline
  \hl 3 & PC-\pUn2 & PC-\pDeux2 & \sol{ping \netDeux.0.2} & \sol{réussi} & \\ \hline
  \hl 4 & PC-\pDeux2 & \srvNom & \sol{ping \netSrv.\srvH} & \sol{réussi} & \\ \hline
  \hl 5 & PC-\pUn3 & page web du serveur & \sol{navigateur : {\NoAutoSpacing http://\netSrv.\srvH}} & \sol{page affichée} & \\ \hline
\end{tabularx}}
\bareme{0,5 pt par test visé + 0,5 pt pour B5. Le premier ping inter-réseaux peut échouer (résolution ARP) ; les suivants doivent réussir.}

\question{B5}{incl. B4}{Test n°3 : par quels équipements passe le paquet ? Pourquoi PC-\pUn2 l'envoie-t-il au routeur ?}
\zone{2.3cm}{PC-\pUn2 $\to$ SW-\pUn{} $\to$ \Rnom{} $\to$ SW-\pDeux{} $\to$ PC-\pDeux2. Grâce à son masque, PC-\pUn2 constate que
la destination (\netDeux.0.0/16) n'est \textbf{pas dans son réseau} (\netUn.0/24) : il envoie le paquet à sa
\textbf{passerelle par défaut}, qui le route.}

%==========================================================================
\newpage
\section{Partie C — Diagnostic d'un réseau en panne \hfill /4}

Un technicien constate qu'\textbf{aucun poste ne parvient à joindre le serveur SRV-C} et que certains postes
communiquent mal entre eux. Les câbles sont du bon type.

\begin{center}
\begin{minipage}{0.37\linewidth}\centering
\begin{tikzpicture}[font=\footnotesize,
  eq/.style={draw,fill=qbg,rounded corners=2pt,minimum width=9mm,minimum height=5mm,inner sep=2pt},
  every path/.style={thick}]
  \node[eq] (SW) {SW-1};
  \node[eq,right=9mm of SW] (R) {R1};
  \node[eq,right=7mm of R] (SRV) {SRV-C};
  \node[eq,above left=5mm and -3mm of SW] (A) {PC-A};
  \node[eq,above right=5mm and -3mm of SW] (B) {PC-B};
  \node[eq,below left=5mm and -3mm of SW] (C) {PC-C};
  \node[eq,below right=5mm and -3mm of SW] (D) {PC-D};
  \foreach \n in {A,B,C,D} \draw (\n) -- (SW);
  \draw (SW) -- node[above,font=\tiny]{G0/0} (R);
  \draw (R) -- node[above,font=\tiny]{G0/1} (SRV);
\end{tikzpicture}
\end{minipage}\hfill
\begin{minipage}{0.61\linewidth}\centering
{\small\renewcommand{\arraystretch}{1.05}\setlength{\tabcolsep}{4pt}
\begin{tabular}{lllll}
  \toprule
  \textbf{Équipement} & \textbf{Adresse IP} & \textbf{Masque} & \textbf{Passerelle} & \textbf{Port} \\
  \midrule
  \Crows
  \bottomrule
\end{tabular}}
\end{minipage}
\end{center}

\question{C1}{4 pts}{Relevez les \textbf{quatre anomalies}. Pour chacune : conséquence et correction.}
{\small\renewcommand{\arraystretch}{1.1}
\ifcorrige\def\hc{}\else\def\hc{\rule[-10mm]{0pt}{14mm}}\fi
\begin{tabularx}{\linewidth}{|c|>{\raggedright\arraybackslash}p{2cm}|L|L|L|}
  \hline
  \textbf{N°} & \textbf{Équipement} & \textbf{Anomalie} & \textbf{Conséquence} & \textbf{Correction} \\ \hline
  \ifcorrige\Csol\else
  \hc 1 & & & & \\ \hline
  \hc 2 & & & & \\ \hline
  \hc 3 & & & & \\ \hline
  \hc 4 & & & & \\ \hline
  \fi
\end{tabularx}}
\bareme{1 pt par anomalie (0,25 équipement, anomalie, conséquence, correction ; ordre indifférent). \Cbareme}

%==========================================================================
\section{Bonus — Extension sans fil \hfill +1}

\begin{cdc}Ajoutez au réseau \resUn{} un \textbf{AccessPoint-PT} relié à SW-\pUn{} (SSID \ip{\ssid}, \textbf{WPA2-PSK})
et un portable \ip{Laptop-\pUn} équipé d'un module \texttt{WPC300N}, en adressage \textbf{statique}.
Vérifiez qu'il joint \srvNom.\end{cdc}
{\small\renewcommand{\arraystretch}{1.1}
\begin{tabularx}{\linewidth}{|l|C|C|C|c|}
  \hline
  \textbf{Équipement} & \textbf{Adresse IP} & \textbf{Masque} & \textbf{Passerelle} & \textbf{Visa prof.} \\ \hline
  \hl Laptop-\pUn & \sol{\netUn.50} & \sol{255.255.255.0} & \sol{\netUn.254} & \hspace{1.4cm} \\ \hline
\end{tabularx}}
\bareme{Le point d'accès fait office de pont : le portable est dans \netUn.0/24. Bonus si associé en WPA2-PSK et ping
vers \netSrv.\srvH{} réussi. Note plafonnée à 20.}

\end{document}
"
%   pdflatex -jobname=Eval_Reseau_CorrigeA "\def\variante{A}\def\modecorrige{1}% Évaluation pratique — Réseaux locaux (BTS CIEL 1re année) — sujets A et B
% Compilation (4 PDF) :
%   pdflatex -jobname=Eval_Reseau_SujetA   "\def\variante{A}\def\modecorrige{0}\input{Eval_TP_Reseau_LAN.tex}"
%   pdflatex -jobname=Eval_Reseau_SujetB   "\def\variante{B}\def\modecorrige{0}\input{Eval_TP_Reseau_LAN.tex}"
%   pdflatex -jobname=Eval_Reseau_CorrigeA "\def\variante{A}\def\modecorrige{1}\input{Eval_TP_Reseau_LAN.tex}"
%   pdflatex -jobname=Eval_Reseau_CorrigeB "\def\variante{B}\def\modecorrige{1}\input{Eval_TP_Reseau_LAN.tex}"
\documentclass[10pt,a4paper]{article}

\usepackage[T1]{fontenc}
\usepackage[utf8]{inputenc}
\usepackage{lmodern}
\usepackage[french]{babel}
\usepackage[a4paper,hmargin=1.5cm,top=1.3cm,bottom=1.6cm,footskip=0.7cm]{geometry}
\usepackage[table]{xcolor}
\usepackage{booktabs,tabularx,array}
\usepackage{enumitem}
\usepackage{fancyhdr}
\usepackage[most]{tcolorbox}
\usepackage{tikz}
\usetikzlibrary{positioning}
\usepackage{titlesec}

\providecommand{\modecorrige}{0}
\providecommand{\variante}{A}
\newif\ifcorrige
\ifnum\modecorrige=1 \corrigetrue \else \corrigefalse \fi

\definecolor{ipred}{HTML}{C0392B}
\definecolor{cdcbar}{HTML}{A67C3D}
\definecolor{cdcbg}{HTML}{F4F1EA}
\definecolor{qbar}{HTML}{4A6B5D}
\definecolor{qbg}{HTML}{EDF1EF}
\definecolor{corr}{HTML}{1F5FA8}
\definecolor{corrbg}{HTML}{E8F0FA}
\definecolor{grisfoot}{HTML}{828282}

\newcommand{\ip}[1]{\textcolor{ipred}{\texttt{#1}}}
\newcommand{\pts}[1]{\hfill\mbox{\textbf{(#1)}}}
\newcommand{\sol}[1]{\ifcorrige\textcolor{corr}{#1}\fi}
\newcommand{\bareme}[1]{\ifcorrige\par{\footnotesize\color{corr}#1\par}\fi}
\newcommand{\hl}{\ifcorrige\rule[-1.6mm]{0pt}{5.4mm}\else\rule[-2.2mm]{0pt}{6.6mm}\fi}
\newcolumntype{C}{>{\centering\arraybackslash}X}
\newcolumntype{L}{>{\raggedright\arraybackslash}X}

\titleformat{\section}{\large\bfseries}{}{0pt}{}
\titlespacing*{\section}{0pt}{6pt}{3pt}

\tcbset{qstyle/.style={enhanced,boxrule=0pt,leftrule=2.5pt,arc=0pt,outer arc=0pt,
  left=5pt,right=5pt,top=2pt,bottom=2pt,before skip=4pt,after skip=3pt}}
\newtcolorbox{cdc}{qstyle,colback=cdcbg,colframe=cdcbar}
\newtcolorbox{quest}{qstyle,colback=qbg,colframe=qbar}

% Zone de réponse : vierge dans le sujet, remplie dans le corrigé
\newcommand{\zone}[2]{%
  \ifcorrige
    \begin{tcolorbox}[colback=corrbg,colframe=corr,boxrule=0.5pt,arc=1pt,left=4pt,right=4pt,top=1pt,bottom=1pt,
      before skip=2pt,after skip=2pt,fontupper=\footnotesize\color{corr}]#2\end{tcolorbox}%
  \else
    \begin{tcolorbox}[colback=white,colframe=black!45,boxrule=0.4pt,arc=1pt,height=#1,
      before skip=2pt,after skip=2pt]\end{tcolorbox}%
  \fi}

\newcommand{\question}[3]{\begin{quest}\textbf{#1.} #3\pts{#2}\end{quest}}

%==========================================================================
% Données propres à chaque variante
%==========================================================================
\if\variante A
  \def\entreprise{Le cabinet d'architectes \textbf{ArchiNova}}
  \def\Rnom{R-ArchiNova}
  \def\resUn{Bureau d'études}\def\pUn{BE}\def\netUn{192.168.40}
  \def\resDeux{Accueil}\def\pDeux{ACC}\def\netDeux{172.16}
  \def\srvNom{SRV-Intranet}\def\netSrv{192.168.200}\def\srvH{10}
  \def\ssid{ArchiNova-WiFi}
  \def\AIrows{%
    \hl\ip{192.168.40.25} & \sol{C} & \sol{255.255.255.0} & \sol{192.168.40.0} & \sol{192.168.40.255} & \sol{254} \\ \hline
    \hl\ip{172.16.5.10}   & \sol{B} & \sol{255.255.0.0}   & \sol{172.16.0.0}   & \sol{172.16.255.255} & \sol{65\,534} \\ \hline
    \hl\ip{10.3.2.1}      & \sol{A} & \sol{255.0.0.0}     & \sol{10.0.0.0}     & \sol{10.255.255.255} & \sol{16\,777\,214} \\ \hline}
  \def\apipa{169.254.12.7}
  \def\Cnet{192.168.60}
  \def\Crows{%
    PC-A & \ip{192.168.60.10} & \ip{255.255.255.0} & \ip{192.168.60.254} & — \\
    PC-B & \ip{192.168.60.20} & \ip{255.255.255.0} & \ip{192.168.60.1}   & — \\
    PC-C & \ip{192.168.61.30} & \ip{255.255.255.0} & \ip{192.168.60.254} & — \\
    PC-D & \ip{192.168.60.10} & \ip{255.255.255.0} & \ip{192.168.60.254} & — \\
    SRV-C & \ip{10.0.0.5}     & \ip{255.0.0.0}     & \ip{10.0.0.254}     & — \\
    R1 — G0/0 & \ip{192.168.60.254} & \ip{255.255.255.0} & — & On \\
    R1 — G0/1 & \ip{10.0.0.254}     & \ip{255.0.0.0}     & — & Off \\}
  \def\Csol{%
    \hl 1 & \sol{PC-B} & \sol{Passerelle .60.1 : pas l'adresse du routeur.} & \sol{Joint son réseau, mais aucun autre (serveur).} & \sol{Passerelle 192.168.60.254.} \\ \hline
    \hl 2 & \sol{PC-C} & \sol{IP hors du réseau 192.168.60.0/24.} & \sol{Ne joint ni les postes ni le routeur.} & \sol{Ex. 192.168.60.30.} \\ \hline
    \hl 3 & \sol{PC-D} & \sol{Même IP que PC-A (doublon).} & \sol{Conflit d'adresse : PC-A et PC-D perturbés.} & \sol{Adresse unique, ex. .60.40.} \\ \hline
    \hl 4 & \sol{R1 G0/1} & \sol{Interface désactivée (Off).} & \sol{Réseau du serveur injoignable pour tous.} & \sol{Port Status : On.} \\ \hline}
  \def\Cbareme{PC-A et SRV-C sont correctement configurés ; accepter PC-A cité à la place de PC-D pour le doublon.}
\else
  \def\entreprise{La clinique vétérinaire \textbf{VetoPlus}}
  \def\Rnom{R-VetoPlus}
  \def\resUn{Soins}\def\pUn{SOI}\def\netUn{192.168.70}
  \def\resDeux{Secrétariat}\def\pDeux{SEC}\def\netDeux{172.20}
  \def\srvNom{SRV-Dossiers}\def\netSrv{192.168.250}\def\srvH{20}
  \def\ssid{VetoPlus-WiFi}
  \def\AIrows{%
    \hl\ip{192.168.15.130} & \sol{C} & \sol{255.255.255.0} & \sol{192.168.15.0} & \sol{192.168.15.255} & \sol{254} \\ \hline
    \hl\ip{172.31.200.7}   & \sol{B} & \sol{255.255.0.0}   & \sol{172.31.0.0}   & \sol{172.31.255.255} & \sol{65\,534} \\ \hline
    \hl\ip{10.45.6.9}      & \sol{A} & \sol{255.0.0.0}     & \sol{10.0.0.0}     & \sol{10.255.255.255} & \sol{16\,777\,214} \\ \hline}
  \def\apipa{169.254.88.3}
  \def\Cnet{192.168.80}
  \def\Crows{%
    PC-A & \ip{192.168.8.11}   & \ip{255.255.255.0} & \ip{192.168.80.254} & — \\
    PC-B & \ip{192.168.80.12}  & \ip{255.255.255.0} & \ip{192.168.80.254} & — \\
    PC-C & \ip{192.168.80.13}  & \ip{255.255.255.0} & \ip{192.168.80.100} & — \\
    PC-D & \ip{192.168.80.12}  & \ip{255.255.255.0} & \ip{192.168.80.254} & — \\
    SRV-C & \ip{10.0.0.8}      & \ip{255.0.0.0}     & \ip{10.0.0.254}     & — \\
    R1 — G0/0 & \ip{192.168.80.254} & \ip{255.255.255.0} & — & Off \\
    R1 — G0/1 & \ip{10.0.0.254}     & \ip{255.0.0.0}     & — & On \\}
  \def\Csol{%
    \hl 1 & \sol{PC-A} & \sol{IP hors du réseau 192.168.80.0/24.} & \sol{Ne joint ni les postes ni le routeur.} & \sol{Ex. 192.168.80.11.} \\ \hline
    \hl 2 & \sol{PC-C} & \sol{Passerelle .80.100 : pas l'adresse du routeur.} & \sol{Joint son réseau, mais aucun autre (serveur).} & \sol{Passerelle 192.168.80.254.} \\ \hline
    \hl 3 & \sol{PC-D} & \sol{Même IP que PC-B (doublon).} & \sol{Conflit d'adresse : PC-B et PC-D perturbés.} & \sol{Adresse unique, ex. .80.14.} \\ \hline
    \hl 4 & \sol{R1 G0/0} & \sol{Interface côté LAN désactivée (Off).} & \sol{Aucun poste ne joint le routeur ni le serveur.} & \sol{Port Status : On.} \\ \hline}
  \def\Cbareme{PC-B et SRV-C sont correctement configurés ; accepter PC-B cité à la place de PC-D pour le doublon.}
\fi

\pagestyle{fancy}
\fancyhf{}
\renewcommand{\headrulewidth}{0pt}
\fancyfoot[L]{\footnotesize\itshape\color{grisfoot}Évaluation — Réseaux locaux — Sujet \variante\ifcorrige\ — CORRIGÉ\fi}
\fancyfoot[R]{\footnotesize\color{grisfoot}\thepage/3}

\setlist{itemsep=0pt,topsep=2pt,parsep=0pt}
\setlength{\parindent}{0pt}
\setlength{\parskip}{2pt}

\begin{document}

\begin{center}
  {\Large\bfseries Évaluation pratique — Réseaux locaux \quad
   \fcolorbox{black}{white}{\ SUJET \variante\ }\par}
  \vspace{1pt}
  {\bfseries Adressage IP, commutation et routage — BTS CIEL 1re année — durée : 2 heures\par}
  \ifcorrige\vspace{2pt}{\bfseries\color{corr}CORRIGÉ ET BARÈME — document enseignant}\par\fi
\end{center}
\vspace{-2pt}
{\renewcommand{\arraystretch}{1.35}
\begin{tabularx}{\linewidth}{|X|X|p{2.2cm}|p{2.6cm}|}
  \hline
  \textbf{Nom :} & \textbf{Prénom :} & \textbf{Poste n°} & \textbf{Note :} \hfill /20 \\ \hline
\end{tabularx}}

\textbf{Consignes :} travail \textbf{individuel}, documents et Internet \textbf{non autorisés}. Ce sujet est le
\textbf{document-réponse}. Fichier Packet Tracer : \ip{Nom\_Prenom\_EvalLAN\_\variante.pkt}, à enregistrer
régulièrement dans le dossier de rendu. Les « visas » sont apposés par le professeur, que vous appelez à chaque test réussi.
Durées conseillées : A 30 min — B 1 h 10 — C 20 min.

%==========================================================================
\section{Partie A — Questions de cours et calculs \hfill /6}

\question{A1}{2 pts}{Complétez le tableau en considérant le \textbf{masque par défaut} de la classe de chaque adresse.}
{\small\renewcommand{\arraystretch}{1.15}
\begin{tabularx}{\linewidth}{|l|c|C|C|C|c|}
  \hline
  \textbf{Adresse IP} & \textbf{Classe} & \textbf{Masque par défaut} & \textbf{Adresse réseau} &
  \textbf{Adresse de diffusion} & \textbf{Nb d'hôtes} \\ \hline
  \AIrows
\end{tabularx}}
\bareme{0,5 pt par ligne juste + 0,5 pt si les trois nombres d'hôtes sont justes ($2^{n}-2$, $n$ = nombre de bits d'hôte).}

\question{A2}{1 pt}{Indiquez le type de câble (\textbf{droit} ou \textbf{croisé}) pour chaque liaison.}
{\small\renewcommand{\arraystretch}{1.15}
\begin{tabularx}{\linewidth}{|X|c||X|c|}
  \hline
  \hl PC $\leftrightarrow$ commutateur & \makebox[1.6cm]{\sol{droit}} & Commutateur $\leftrightarrow$ routeur & \makebox[1.6cm]{\sol{droit}} \\ \hline
  \hl PC $\leftrightarrow$ PC (réseau P2P) & \sol{croisé} & Concentrateur $\leftrightarrow$ concentrateur & \sol{croisé} \\ \hline
\end{tabularx}}
\bareme{0,25 pt par case. Équipements de types différents : câble droit ; de même type : câble croisé.}

\question{A3}{1 pt}{Comment un \textbf{concentrateur (hub)} et un \textbf{commutateur (switch)} traitent-ils une trame
reçue ? Lequel limite les collisions, et pourquoi ?}
\zone{2.9cm}{Le \textbf{hub} (couche 1) répète la trame sur \textbf{tous ses autres ports} : un seul domaine de collision
partagé. Le \textbf{switch} (couche 2) lit l'adresse MAC de destination et, grâce à sa \textbf{table MAC}, ne l'envoie
\textbf{que sur le port du destinataire} : un domaine de collision par port, c'est lui qui limite les collisions.
\emph{0,5 pt hub, 0,5 pt switch.}}

\question{A4}{1 pt}{Qu'est-ce que la \textbf{passerelle par défaut} d'un poste ? Quelle condition son adresse doit-elle respecter ?}
\zone{2.4cm}{Adresse IP de l'interface du \textbf{routeur} à laquelle le poste envoie les paquets destinés à un
\textbf{autre réseau}. Elle doit \textbf{appartenir au même réseau} que le poste. \emph{0,5 + 0,5 pt.}}

\question{A5}{0,5 pt}{Un portable configuré en DHCP affiche \ip{\apipa} / \ip{255.255.0.0}. Que s'est-il passé ?}
\zone{1.5cm}{Aucun \textbf{serveur DHCP} n'a répondu : le poste s'est attribué une adresse \textbf{APIPA} (169.254.0.0/16).}

\question{A6}{0,5 pt}{Quel protocole utilise \texttt{ping} ? Nommez les deux messages échangés et la couche OSI des adresses IP.}
\zone{1.5cm}{\textbf{ICMP} ; \textbf{Echo Request} / \textbf{Echo Reply} ; adresses IP en \textbf{couche 3 (Réseau)}.}

%==========================================================================
\newpage
\section{Partie B — Réalisation sous Packet Tracer \hfill /10}

\begin{cdc}
\entreprise{} vous confie son réseau : \textbf{trois réseaux} reliés par un \textbf{routeur unique Cisco 2911}
nommé \ip{\Rnom} (interfaces G0/0, G0/1, G0/2).
\begin{itemize}[leftmargin=12pt]
  \item \textbf{Réseau \resUn} : \textbf{4 postes} PC-\pUn1 à PC-\pUn4, commutateur \ip{SW-\pUn}, adressage statique,
        réseau \ip{\netUn.0/24}.
  \item \textbf{Réseau \resDeux} : \textbf{2 postes} PC-\pDeux1, PC-\pDeux2 et \textbf{1 imprimante} IMP-\pDeux,
        commutateur \ip{SW-\pDeux}, adressage statique, réseau de classe B \ip{\netDeux.0.0/16}.
  \item \textbf{Réseau Serveur} : serveur \ip{\srvNom} d'adresse \ip{\netSrv.\srvH/24}, relié \textbf{directement} au routeur.
  \item La passerelle de chaque réseau est sa \textbf{dernière adresse utilisable}. Tous les équipements communiquent
        entre eux et accèdent à la page web du serveur. Nommage imposé (« Display Name »).
\end{itemize}
\end{cdc}

\question{B1}{2 pts}{Complétez le plan d'adressage \textbf{avant} la réalisation.}
{\footnotesize\setlength{\tabcolsep}{3pt}\renewcommand{\arraystretch}{1.15}
\begin{tabularx}{\linewidth}{|l|C|C|C|C|C|C|}
  \hline
  \textbf{Réseau} & \textbf{Adresse réseau} & \textbf{Masque} & \textbf{1re @ utilisable} &
  \textbf{Dernière @ utilisable} & \textbf{Diffusion} & \textbf{Passerelle} \\ \hline
  \hl \resUn & \sol{\netUn.0} & \sol{255.255.255.0} & \sol{\netUn.1} & \sol{\netUn.254} & \sol{\netUn.255} & \sol{\netUn.254} \\ \hline
  \hl \resDeux & \sol{\netDeux.0.0} & \sol{255.255.0.0} & \sol{\netDeux.0.1} & \sol{\netDeux.255.254} & \sol{\netDeux.255.255} & \sol{\netDeux.255.254} \\ \hline
  \hl Serveur & \sol{\netSrv.0} & \sol{255.255.255.0} & \sol{\netSrv.1} & \sol{\netSrv.254} & \sol{\netSrv.255} & \sol{\netSrv.254} \\ \hline
\end{tabularx}}
\bareme{0,5 pt par ligne juste + 0,5 pt si la convention de passerelle est respectée partout.}

\question{B2}{2 pts}{Construisez la topologie (équipements, câblage, nommage) et relevez l'adressage configuré.}
{\footnotesize\setlength{\tabcolsep}{3pt}\renewcommand{\arraystretch}{1.1}
\begin{tabularx}{\linewidth}{|l|C|C|C||l|C|C|}
  \hline
  \textbf{Hôte} & \textbf{Adresse IP} & \textbf{Masque} & \textbf{Passerelle} & \textbf{Interface} & \textbf{Adresse IP} & \textbf{Masque} \\ \hline
  \hl PC-\pUn1 & \sol{\netUn.1} & \sol{255.255.255.0} & \sol{\netUn.254} & \Rnom{} G0/0 & \sol{\netUn.254} & \sol{255.255.255.0} \\ \hline
  \hl PC-\pDeux1 & \sol{\netDeux.0.1} & \sol{255.255.0.0} & \sol{\netDeux.255.254} & \Rnom{} G0/1 & \sol{\netDeux.255.254} & \sol{255.255.0.0} \\ \hline
  \hl IMP-\pDeux & \sol{\netDeux.0.10} & \sol{255.255.0.0} & \sol{\netDeux.255.254} & \Rnom{} G0/2 & \sol{\netSrv.254} & \sol{255.255.255.0} \\ \hline
  \hl \srvNom & \ip{\netSrv.\srvH} & \sol{255.255.255.0} & \sol{\netSrv.254} & \multicolumn{3}{c|}{\cellcolor{black!8}} \\ \hline
\end{tabularx}}
\bareme{Adresses d'hôtes libres (uniques, dans le bon réseau) ; affectation G0/x libre. Sur le .pkt : équipements 0,5 ;
câblage 1 (câbles droits vers les switchs et switch–routeur, \textbf{croisé} serveur–routeur, liens verts) ; nommage 0,5.}

\question{B3}{3 pts}{Configurez tous les équipements : IP, masques, passerelles des hôtes ; interfaces du routeur
(« Port Status : On »).}
\bareme{Sur le .pkt : interfaces routeur 1 pt ; IP/masques des hôtes 1 pt ; passerelles (imprimante et serveur compris)
1 pt. $-0,25$ par erreur, sans descendre sous 0 par item.}

\question{B4}{3 pts}{Réalisez les tests, notez la commande et le résultat, puis faites viser chaque test réussi.}
{\small\renewcommand{\arraystretch}{1.1}
\begin{tabularx}{\linewidth}{|c|l|l|C|c|c|}
  \hline
  \textbf{N°} & \textbf{Source} & \textbf{Destination} & \textbf{Commande / action} & \textbf{Résultat} & \textbf{Visa prof.} \\ \hline
  \hl 1 & PC-\pUn1 & PC-\pUn4 & \sol{ping \netUn.4} & \sol{réussi} & \hspace{1.4cm} \\ \hline
  \hl 2 & PC-\pDeux1 & IMP-\pDeux & \sol{ping \netDeux.0.10} & \sol{réussi} & \\ \hline
  \hl 3 & PC-\pUn2 & PC-\pDeux2 & \sol{ping \netDeux.0.2} & \sol{réussi} & \\ \hline
  \hl 4 & PC-\pDeux2 & \srvNom & \sol{ping \netSrv.\srvH} & \sol{réussi} & \\ \hline
  \hl 5 & PC-\pUn3 & page web du serveur & \sol{navigateur : {\NoAutoSpacing http://\netSrv.\srvH}} & \sol{page affichée} & \\ \hline
\end{tabularx}}
\bareme{0,5 pt par test visé + 0,5 pt pour B5. Le premier ping inter-réseaux peut échouer (résolution ARP) ; les suivants doivent réussir.}

\question{B5}{incl. B4}{Test n°3 : par quels équipements passe le paquet ? Pourquoi PC-\pUn2 l'envoie-t-il au routeur ?}
\zone{2.3cm}{PC-\pUn2 $\to$ SW-\pUn{} $\to$ \Rnom{} $\to$ SW-\pDeux{} $\to$ PC-\pDeux2. Grâce à son masque, PC-\pUn2 constate que
la destination (\netDeux.0.0/16) n'est \textbf{pas dans son réseau} (\netUn.0/24) : il envoie le paquet à sa
\textbf{passerelle par défaut}, qui le route.}

%==========================================================================
\newpage
\section{Partie C — Diagnostic d'un réseau en panne \hfill /4}

Un technicien constate qu'\textbf{aucun poste ne parvient à joindre le serveur SRV-C} et que certains postes
communiquent mal entre eux. Les câbles sont du bon type.

\begin{center}
\begin{minipage}{0.37\linewidth}\centering
\begin{tikzpicture}[font=\footnotesize,
  eq/.style={draw,fill=qbg,rounded corners=2pt,minimum width=9mm,minimum height=5mm,inner sep=2pt},
  every path/.style={thick}]
  \node[eq] (SW) {SW-1};
  \node[eq,right=9mm of SW] (R) {R1};
  \node[eq,right=7mm of R] (SRV) {SRV-C};
  \node[eq,above left=5mm and -3mm of SW] (A) {PC-A};
  \node[eq,above right=5mm and -3mm of SW] (B) {PC-B};
  \node[eq,below left=5mm and -3mm of SW] (C) {PC-C};
  \node[eq,below right=5mm and -3mm of SW] (D) {PC-D};
  \foreach \n in {A,B,C,D} \draw (\n) -- (SW);
  \draw (SW) -- node[above,font=\tiny]{G0/0} (R);
  \draw (R) -- node[above,font=\tiny]{G0/1} (SRV);
\end{tikzpicture}
\end{minipage}\hfill
\begin{minipage}{0.61\linewidth}\centering
{\small\renewcommand{\arraystretch}{1.05}\setlength{\tabcolsep}{4pt}
\begin{tabular}{lllll}
  \toprule
  \textbf{Équipement} & \textbf{Adresse IP} & \textbf{Masque} & \textbf{Passerelle} & \textbf{Port} \\
  \midrule
  \Crows
  \bottomrule
\end{tabular}}
\end{minipage}
\end{center}

\question{C1}{4 pts}{Relevez les \textbf{quatre anomalies}. Pour chacune : conséquence et correction.}
{\small\renewcommand{\arraystretch}{1.1}
\ifcorrige\def\hc{}\else\def\hc{\rule[-10mm]{0pt}{14mm}}\fi
\begin{tabularx}{\linewidth}{|c|>{\raggedright\arraybackslash}p{2cm}|L|L|L|}
  \hline
  \textbf{N°} & \textbf{Équipement} & \textbf{Anomalie} & \textbf{Conséquence} & \textbf{Correction} \\ \hline
  \ifcorrige\Csol\else
  \hc 1 & & & & \\ \hline
  \hc 2 & & & & \\ \hline
  \hc 3 & & & & \\ \hline
  \hc 4 & & & & \\ \hline
  \fi
\end{tabularx}}
\bareme{1 pt par anomalie (0,25 équipement, anomalie, conséquence, correction ; ordre indifférent). \Cbareme}

%==========================================================================
\section{Bonus — Extension sans fil \hfill +1}

\begin{cdc}Ajoutez au réseau \resUn{} un \textbf{AccessPoint-PT} relié à SW-\pUn{} (SSID \ip{\ssid}, \textbf{WPA2-PSK})
et un portable \ip{Laptop-\pUn} équipé d'un module \texttt{WPC300N}, en adressage \textbf{statique}.
Vérifiez qu'il joint \srvNom.\end{cdc}
{\small\renewcommand{\arraystretch}{1.1}
\begin{tabularx}{\linewidth}{|l|C|C|C|c|}
  \hline
  \textbf{Équipement} & \textbf{Adresse IP} & \textbf{Masque} & \textbf{Passerelle} & \textbf{Visa prof.} \\ \hline
  \hl Laptop-\pUn & \sol{\netUn.50} & \sol{255.255.255.0} & \sol{\netUn.254} & \hspace{1.4cm} \\ \hline
\end{tabularx}}
\bareme{Le point d'accès fait office de pont : le portable est dans \netUn.0/24. Bonus si associé en WPA2-PSK et ping
vers \netSrv.\srvH{} réussi. Note plafonnée à 20.}

\end{document}
"
%   pdflatex -jobname=Eval_Reseau_CorrigeB "\def\variante{B}\def\modecorrige{1}% Évaluation pratique — Réseaux locaux (BTS CIEL 1re année) — sujets A et B
% Compilation (4 PDF) :
%   pdflatex -jobname=Eval_Reseau_SujetA   "\def\variante{A}\def\modecorrige{0}\input{Eval_TP_Reseau_LAN.tex}"
%   pdflatex -jobname=Eval_Reseau_SujetB   "\def\variante{B}\def\modecorrige{0}\input{Eval_TP_Reseau_LAN.tex}"
%   pdflatex -jobname=Eval_Reseau_CorrigeA "\def\variante{A}\def\modecorrige{1}\input{Eval_TP_Reseau_LAN.tex}"
%   pdflatex -jobname=Eval_Reseau_CorrigeB "\def\variante{B}\def\modecorrige{1}\input{Eval_TP_Reseau_LAN.tex}"
\documentclass[10pt,a4paper]{article}

\usepackage[T1]{fontenc}
\usepackage[utf8]{inputenc}
\usepackage{lmodern}
\usepackage[french]{babel}
\usepackage[a4paper,hmargin=1.5cm,top=1.3cm,bottom=1.6cm,footskip=0.7cm]{geometry}
\usepackage[table]{xcolor}
\usepackage{booktabs,tabularx,array}
\usepackage{enumitem}
\usepackage{fancyhdr}
\usepackage[most]{tcolorbox}
\usepackage{tikz}
\usetikzlibrary{positioning}
\usepackage{titlesec}

\providecommand{\modecorrige}{0}
\providecommand{\variante}{A}
\newif\ifcorrige
\ifnum\modecorrige=1 \corrigetrue \else \corrigefalse \fi

\definecolor{ipred}{HTML}{C0392B}
\definecolor{cdcbar}{HTML}{A67C3D}
\definecolor{cdcbg}{HTML}{F4F1EA}
\definecolor{qbar}{HTML}{4A6B5D}
\definecolor{qbg}{HTML}{EDF1EF}
\definecolor{corr}{HTML}{1F5FA8}
\definecolor{corrbg}{HTML}{E8F0FA}
\definecolor{grisfoot}{HTML}{828282}

\newcommand{\ip}[1]{\textcolor{ipred}{\texttt{#1}}}
\newcommand{\pts}[1]{\hfill\mbox{\textbf{(#1)}}}
\newcommand{\sol}[1]{\ifcorrige\textcolor{corr}{#1}\fi}
\newcommand{\bareme}[1]{\ifcorrige\par{\footnotesize\color{corr}#1\par}\fi}
\newcommand{\hl}{\ifcorrige\rule[-1.6mm]{0pt}{5.4mm}\else\rule[-2.2mm]{0pt}{6.6mm}\fi}
\newcolumntype{C}{>{\centering\arraybackslash}X}
\newcolumntype{L}{>{\raggedright\arraybackslash}X}

\titleformat{\section}{\large\bfseries}{}{0pt}{}
\titlespacing*{\section}{0pt}{6pt}{3pt}

\tcbset{qstyle/.style={enhanced,boxrule=0pt,leftrule=2.5pt,arc=0pt,outer arc=0pt,
  left=5pt,right=5pt,top=2pt,bottom=2pt,before skip=4pt,after skip=3pt}}
\newtcolorbox{cdc}{qstyle,colback=cdcbg,colframe=cdcbar}
\newtcolorbox{quest}{qstyle,colback=qbg,colframe=qbar}

% Zone de réponse : vierge dans le sujet, remplie dans le corrigé
\newcommand{\zone}[2]{%
  \ifcorrige
    \begin{tcolorbox}[colback=corrbg,colframe=corr,boxrule=0.5pt,arc=1pt,left=4pt,right=4pt,top=1pt,bottom=1pt,
      before skip=2pt,after skip=2pt,fontupper=\footnotesize\color{corr}]#2\end{tcolorbox}%
  \else
    \begin{tcolorbox}[colback=white,colframe=black!45,boxrule=0.4pt,arc=1pt,height=#1,
      before skip=2pt,after skip=2pt]\end{tcolorbox}%
  \fi}

\newcommand{\question}[3]{\begin{quest}\textbf{#1.} #3\pts{#2}\end{quest}}

%==========================================================================
% Données propres à chaque variante
%==========================================================================
\if\variante A
  \def\entreprise{Le cabinet d'architectes \textbf{ArchiNova}}
  \def\Rnom{R-ArchiNova}
  \def\resUn{Bureau d'études}\def\pUn{BE}\def\netUn{192.168.40}
  \def\resDeux{Accueil}\def\pDeux{ACC}\def\netDeux{172.16}
  \def\srvNom{SRV-Intranet}\def\netSrv{192.168.200}\def\srvH{10}
  \def\ssid{ArchiNova-WiFi}
  \def\AIrows{%
    \hl\ip{192.168.40.25} & \sol{C} & \sol{255.255.255.0} & \sol{192.168.40.0} & \sol{192.168.40.255} & \sol{254} \\ \hline
    \hl\ip{172.16.5.10}   & \sol{B} & \sol{255.255.0.0}   & \sol{172.16.0.0}   & \sol{172.16.255.255} & \sol{65\,534} \\ \hline
    \hl\ip{10.3.2.1}      & \sol{A} & \sol{255.0.0.0}     & \sol{10.0.0.0}     & \sol{10.255.255.255} & \sol{16\,777\,214} \\ \hline}
  \def\apipa{169.254.12.7}
  \def\Cnet{192.168.60}
  \def\Crows{%
    PC-A & \ip{192.168.60.10} & \ip{255.255.255.0} & \ip{192.168.60.254} & — \\
    PC-B & \ip{192.168.60.20} & \ip{255.255.255.0} & \ip{192.168.60.1}   & — \\
    PC-C & \ip{192.168.61.30} & \ip{255.255.255.0} & \ip{192.168.60.254} & — \\
    PC-D & \ip{192.168.60.10} & \ip{255.255.255.0} & \ip{192.168.60.254} & — \\
    SRV-C & \ip{10.0.0.5}     & \ip{255.0.0.0}     & \ip{10.0.0.254}     & — \\
    R1 — G0/0 & \ip{192.168.60.254} & \ip{255.255.255.0} & — & On \\
    R1 — G0/1 & \ip{10.0.0.254}     & \ip{255.0.0.0}     & — & Off \\}
  \def\Csol{%
    \hl 1 & \sol{PC-B} & \sol{Passerelle .60.1 : pas l'adresse du routeur.} & \sol{Joint son réseau, mais aucun autre (serveur).} & \sol{Passerelle 192.168.60.254.} \\ \hline
    \hl 2 & \sol{PC-C} & \sol{IP hors du réseau 192.168.60.0/24.} & \sol{Ne joint ni les postes ni le routeur.} & \sol{Ex. 192.168.60.30.} \\ \hline
    \hl 3 & \sol{PC-D} & \sol{Même IP que PC-A (doublon).} & \sol{Conflit d'adresse : PC-A et PC-D perturbés.} & \sol{Adresse unique, ex. .60.40.} \\ \hline
    \hl 4 & \sol{R1 G0/1} & \sol{Interface désactivée (Off).} & \sol{Réseau du serveur injoignable pour tous.} & \sol{Port Status : On.} \\ \hline}
  \def\Cbareme{PC-A et SRV-C sont correctement configurés ; accepter PC-A cité à la place de PC-D pour le doublon.}
\else
  \def\entreprise{La clinique vétérinaire \textbf{VetoPlus}}
  \def\Rnom{R-VetoPlus}
  \def\resUn{Soins}\def\pUn{SOI}\def\netUn{192.168.70}
  \def\resDeux{Secrétariat}\def\pDeux{SEC}\def\netDeux{172.20}
  \def\srvNom{SRV-Dossiers}\def\netSrv{192.168.250}\def\srvH{20}
  \def\ssid{VetoPlus-WiFi}
  \def\AIrows{%
    \hl\ip{192.168.15.130} & \sol{C} & \sol{255.255.255.0} & \sol{192.168.15.0} & \sol{192.168.15.255} & \sol{254} \\ \hline
    \hl\ip{172.31.200.7}   & \sol{B} & \sol{255.255.0.0}   & \sol{172.31.0.0}   & \sol{172.31.255.255} & \sol{65\,534} \\ \hline
    \hl\ip{10.45.6.9}      & \sol{A} & \sol{255.0.0.0}     & \sol{10.0.0.0}     & \sol{10.255.255.255} & \sol{16\,777\,214} \\ \hline}
  \def\apipa{169.254.88.3}
  \def\Cnet{192.168.80}
  \def\Crows{%
    PC-A & \ip{192.168.8.11}   & \ip{255.255.255.0} & \ip{192.168.80.254} & — \\
    PC-B & \ip{192.168.80.12}  & \ip{255.255.255.0} & \ip{192.168.80.254} & — \\
    PC-C & \ip{192.168.80.13}  & \ip{255.255.255.0} & \ip{192.168.80.100} & — \\
    PC-D & \ip{192.168.80.12}  & \ip{255.255.255.0} & \ip{192.168.80.254} & — \\
    SRV-C & \ip{10.0.0.8}      & \ip{255.0.0.0}     & \ip{10.0.0.254}     & — \\
    R1 — G0/0 & \ip{192.168.80.254} & \ip{255.255.255.0} & — & Off \\
    R1 — G0/1 & \ip{10.0.0.254}     & \ip{255.0.0.0}     & — & On \\}
  \def\Csol{%
    \hl 1 & \sol{PC-A} & \sol{IP hors du réseau 192.168.80.0/24.} & \sol{Ne joint ni les postes ni le routeur.} & \sol{Ex. 192.168.80.11.} \\ \hline
    \hl 2 & \sol{PC-C} & \sol{Passerelle .80.100 : pas l'adresse du routeur.} & \sol{Joint son réseau, mais aucun autre (serveur).} & \sol{Passerelle 192.168.80.254.} \\ \hline
    \hl 3 & \sol{PC-D} & \sol{Même IP que PC-B (doublon).} & \sol{Conflit d'adresse : PC-B et PC-D perturbés.} & \sol{Adresse unique, ex. .80.14.} \\ \hline
    \hl 4 & \sol{R1 G0/0} & \sol{Interface côté LAN désactivée (Off).} & \sol{Aucun poste ne joint le routeur ni le serveur.} & \sol{Port Status : On.} \\ \hline}
  \def\Cbareme{PC-B et SRV-C sont correctement configurés ; accepter PC-B cité à la place de PC-D pour le doublon.}
\fi

\pagestyle{fancy}
\fancyhf{}
\renewcommand{\headrulewidth}{0pt}
\fancyfoot[L]{\footnotesize\itshape\color{grisfoot}Évaluation — Réseaux locaux — Sujet \variante\ifcorrige\ — CORRIGÉ\fi}
\fancyfoot[R]{\footnotesize\color{grisfoot}\thepage/3}

\setlist{itemsep=0pt,topsep=2pt,parsep=0pt}
\setlength{\parindent}{0pt}
\setlength{\parskip}{2pt}

\begin{document}

\begin{center}
  {\Large\bfseries Évaluation pratique — Réseaux locaux \quad
   \fcolorbox{black}{white}{\ SUJET \variante\ }\par}
  \vspace{1pt}
  {\bfseries Adressage IP, commutation et routage — BTS CIEL 1re année — durée : 2 heures\par}
  \ifcorrige\vspace{2pt}{\bfseries\color{corr}CORRIGÉ ET BARÈME — document enseignant}\par\fi
\end{center}
\vspace{-2pt}
{\renewcommand{\arraystretch}{1.35}
\begin{tabularx}{\linewidth}{|X|X|p{2.2cm}|p{2.6cm}|}
  \hline
  \textbf{Nom :} & \textbf{Prénom :} & \textbf{Poste n°} & \textbf{Note :} \hfill /20 \\ \hline
\end{tabularx}}

\textbf{Consignes :} travail \textbf{individuel}, documents et Internet \textbf{non autorisés}. Ce sujet est le
\textbf{document-réponse}. Fichier Packet Tracer : \ip{Nom\_Prenom\_EvalLAN\_\variante.pkt}, à enregistrer
régulièrement dans le dossier de rendu. Les « visas » sont apposés par le professeur, que vous appelez à chaque test réussi.
Durées conseillées : A 30 min — B 1 h 10 — C 20 min.

%==========================================================================
\section{Partie A — Questions de cours et calculs \hfill /6}

\question{A1}{2 pts}{Complétez le tableau en considérant le \textbf{masque par défaut} de la classe de chaque adresse.}
{\small\renewcommand{\arraystretch}{1.15}
\begin{tabularx}{\linewidth}{|l|c|C|C|C|c|}
  \hline
  \textbf{Adresse IP} & \textbf{Classe} & \textbf{Masque par défaut} & \textbf{Adresse réseau} &
  \textbf{Adresse de diffusion} & \textbf{Nb d'hôtes} \\ \hline
  \AIrows
\end{tabularx}}
\bareme{0,5 pt par ligne juste + 0,5 pt si les trois nombres d'hôtes sont justes ($2^{n}-2$, $n$ = nombre de bits d'hôte).}

\question{A2}{1 pt}{Indiquez le type de câble (\textbf{droit} ou \textbf{croisé}) pour chaque liaison.}
{\small\renewcommand{\arraystretch}{1.15}
\begin{tabularx}{\linewidth}{|X|c||X|c|}
  \hline
  \hl PC $\leftrightarrow$ commutateur & \makebox[1.6cm]{\sol{droit}} & Commutateur $\leftrightarrow$ routeur & \makebox[1.6cm]{\sol{droit}} \\ \hline
  \hl PC $\leftrightarrow$ PC (réseau P2P) & \sol{croisé} & Concentrateur $\leftrightarrow$ concentrateur & \sol{croisé} \\ \hline
\end{tabularx}}
\bareme{0,25 pt par case. Équipements de types différents : câble droit ; de même type : câble croisé.}

\question{A3}{1 pt}{Comment un \textbf{concentrateur (hub)} et un \textbf{commutateur (switch)} traitent-ils une trame
reçue ? Lequel limite les collisions, et pourquoi ?}
\zone{2.9cm}{Le \textbf{hub} (couche 1) répète la trame sur \textbf{tous ses autres ports} : un seul domaine de collision
partagé. Le \textbf{switch} (couche 2) lit l'adresse MAC de destination et, grâce à sa \textbf{table MAC}, ne l'envoie
\textbf{que sur le port du destinataire} : un domaine de collision par port, c'est lui qui limite les collisions.
\emph{0,5 pt hub, 0,5 pt switch.}}

\question{A4}{1 pt}{Qu'est-ce que la \textbf{passerelle par défaut} d'un poste ? Quelle condition son adresse doit-elle respecter ?}
\zone{2.4cm}{Adresse IP de l'interface du \textbf{routeur} à laquelle le poste envoie les paquets destinés à un
\textbf{autre réseau}. Elle doit \textbf{appartenir au même réseau} que le poste. \emph{0,5 + 0,5 pt.}}

\question{A5}{0,5 pt}{Un portable configuré en DHCP affiche \ip{\apipa} / \ip{255.255.0.0}. Que s'est-il passé ?}
\zone{1.5cm}{Aucun \textbf{serveur DHCP} n'a répondu : le poste s'est attribué une adresse \textbf{APIPA} (169.254.0.0/16).}

\question{A6}{0,5 pt}{Quel protocole utilise \texttt{ping} ? Nommez les deux messages échangés et la couche OSI des adresses IP.}
\zone{1.5cm}{\textbf{ICMP} ; \textbf{Echo Request} / \textbf{Echo Reply} ; adresses IP en \textbf{couche 3 (Réseau)}.}

%==========================================================================
\newpage
\section{Partie B — Réalisation sous Packet Tracer \hfill /10}

\begin{cdc}
\entreprise{} vous confie son réseau : \textbf{trois réseaux} reliés par un \textbf{routeur unique Cisco 2911}
nommé \ip{\Rnom} (interfaces G0/0, G0/1, G0/2).
\begin{itemize}[leftmargin=12pt]
  \item \textbf{Réseau \resUn} : \textbf{4 postes} PC-\pUn1 à PC-\pUn4, commutateur \ip{SW-\pUn}, adressage statique,
        réseau \ip{\netUn.0/24}.
  \item \textbf{Réseau \resDeux} : \textbf{2 postes} PC-\pDeux1, PC-\pDeux2 et \textbf{1 imprimante} IMP-\pDeux,
        commutateur \ip{SW-\pDeux}, adressage statique, réseau de classe B \ip{\netDeux.0.0/16}.
  \item \textbf{Réseau Serveur} : serveur \ip{\srvNom} d'adresse \ip{\netSrv.\srvH/24}, relié \textbf{directement} au routeur.
  \item La passerelle de chaque réseau est sa \textbf{dernière adresse utilisable}. Tous les équipements communiquent
        entre eux et accèdent à la page web du serveur. Nommage imposé (« Display Name »).
\end{itemize}
\end{cdc}

\question{B1}{2 pts}{Complétez le plan d'adressage \textbf{avant} la réalisation.}
{\footnotesize\setlength{\tabcolsep}{3pt}\renewcommand{\arraystretch}{1.15}
\begin{tabularx}{\linewidth}{|l|C|C|C|C|C|C|}
  \hline
  \textbf{Réseau} & \textbf{Adresse réseau} & \textbf{Masque} & \textbf{1re @ utilisable} &
  \textbf{Dernière @ utilisable} & \textbf{Diffusion} & \textbf{Passerelle} \\ \hline
  \hl \resUn & \sol{\netUn.0} & \sol{255.255.255.0} & \sol{\netUn.1} & \sol{\netUn.254} & \sol{\netUn.255} & \sol{\netUn.254} \\ \hline
  \hl \resDeux & \sol{\netDeux.0.0} & \sol{255.255.0.0} & \sol{\netDeux.0.1} & \sol{\netDeux.255.254} & \sol{\netDeux.255.255} & \sol{\netDeux.255.254} \\ \hline
  \hl Serveur & \sol{\netSrv.0} & \sol{255.255.255.0} & \sol{\netSrv.1} & \sol{\netSrv.254} & \sol{\netSrv.255} & \sol{\netSrv.254} \\ \hline
\end{tabularx}}
\bareme{0,5 pt par ligne juste + 0,5 pt si la convention de passerelle est respectée partout.}

\question{B2}{2 pts}{Construisez la topologie (équipements, câblage, nommage) et relevez l'adressage configuré.}
{\footnotesize\setlength{\tabcolsep}{3pt}\renewcommand{\arraystretch}{1.1}
\begin{tabularx}{\linewidth}{|l|C|C|C||l|C|C|}
  \hline
  \textbf{Hôte} & \textbf{Adresse IP} & \textbf{Masque} & \textbf{Passerelle} & \textbf{Interface} & \textbf{Adresse IP} & \textbf{Masque} \\ \hline
  \hl PC-\pUn1 & \sol{\netUn.1} & \sol{255.255.255.0} & \sol{\netUn.254} & \Rnom{} G0/0 & \sol{\netUn.254} & \sol{255.255.255.0} \\ \hline
  \hl PC-\pDeux1 & \sol{\netDeux.0.1} & \sol{255.255.0.0} & \sol{\netDeux.255.254} & \Rnom{} G0/1 & \sol{\netDeux.255.254} & \sol{255.255.0.0} \\ \hline
  \hl IMP-\pDeux & \sol{\netDeux.0.10} & \sol{255.255.0.0} & \sol{\netDeux.255.254} & \Rnom{} G0/2 & \sol{\netSrv.254} & \sol{255.255.255.0} \\ \hline
  \hl \srvNom & \ip{\netSrv.\srvH} & \sol{255.255.255.0} & \sol{\netSrv.254} & \multicolumn{3}{c|}{\cellcolor{black!8}} \\ \hline
\end{tabularx}}
\bareme{Adresses d'hôtes libres (uniques, dans le bon réseau) ; affectation G0/x libre. Sur le .pkt : équipements 0,5 ;
câblage 1 (câbles droits vers les switchs et switch–routeur, \textbf{croisé} serveur–routeur, liens verts) ; nommage 0,5.}

\question{B3}{3 pts}{Configurez tous les équipements : IP, masques, passerelles des hôtes ; interfaces du routeur
(« Port Status : On »).}
\bareme{Sur le .pkt : interfaces routeur 1 pt ; IP/masques des hôtes 1 pt ; passerelles (imprimante et serveur compris)
1 pt. $-0,25$ par erreur, sans descendre sous 0 par item.}

\question{B4}{3 pts}{Réalisez les tests, notez la commande et le résultat, puis faites viser chaque test réussi.}
{\small\renewcommand{\arraystretch}{1.1}
\begin{tabularx}{\linewidth}{|c|l|l|C|c|c|}
  \hline
  \textbf{N°} & \textbf{Source} & \textbf{Destination} & \textbf{Commande / action} & \textbf{Résultat} & \textbf{Visa prof.} \\ \hline
  \hl 1 & PC-\pUn1 & PC-\pUn4 & \sol{ping \netUn.4} & \sol{réussi} & \hspace{1.4cm} \\ \hline
  \hl 2 & PC-\pDeux1 & IMP-\pDeux & \sol{ping \netDeux.0.10} & \sol{réussi} & \\ \hline
  \hl 3 & PC-\pUn2 & PC-\pDeux2 & \sol{ping \netDeux.0.2} & \sol{réussi} & \\ \hline
  \hl 4 & PC-\pDeux2 & \srvNom & \sol{ping \netSrv.\srvH} & \sol{réussi} & \\ \hline
  \hl 5 & PC-\pUn3 & page web du serveur & \sol{navigateur : {\NoAutoSpacing http://\netSrv.\srvH}} & \sol{page affichée} & \\ \hline
\end{tabularx}}
\bareme{0,5 pt par test visé + 0,5 pt pour B5. Le premier ping inter-réseaux peut échouer (résolution ARP) ; les suivants doivent réussir.}

\question{B5}{incl. B4}{Test n°3 : par quels équipements passe le paquet ? Pourquoi PC-\pUn2 l'envoie-t-il au routeur ?}
\zone{2.3cm}{PC-\pUn2 $\to$ SW-\pUn{} $\to$ \Rnom{} $\to$ SW-\pDeux{} $\to$ PC-\pDeux2. Grâce à son masque, PC-\pUn2 constate que
la destination (\netDeux.0.0/16) n'est \textbf{pas dans son réseau} (\netUn.0/24) : il envoie le paquet à sa
\textbf{passerelle par défaut}, qui le route.}

%==========================================================================
\newpage
\section{Partie C — Diagnostic d'un réseau en panne \hfill /4}

Un technicien constate qu'\textbf{aucun poste ne parvient à joindre le serveur SRV-C} et que certains postes
communiquent mal entre eux. Les câbles sont du bon type.

\begin{center}
\begin{minipage}{0.37\linewidth}\centering
\begin{tikzpicture}[font=\footnotesize,
  eq/.style={draw,fill=qbg,rounded corners=2pt,minimum width=9mm,minimum height=5mm,inner sep=2pt},
  every path/.style={thick}]
  \node[eq] (SW) {SW-1};
  \node[eq,right=9mm of SW] (R) {R1};
  \node[eq,right=7mm of R] (SRV) {SRV-C};
  \node[eq,above left=5mm and -3mm of SW] (A) {PC-A};
  \node[eq,above right=5mm and -3mm of SW] (B) {PC-B};
  \node[eq,below left=5mm and -3mm of SW] (C) {PC-C};
  \node[eq,below right=5mm and -3mm of SW] (D) {PC-D};
  \foreach \n in {A,B,C,D} \draw (\n) -- (SW);
  \draw (SW) -- node[above,font=\tiny]{G0/0} (R);
  \draw (R) -- node[above,font=\tiny]{G0/1} (SRV);
\end{tikzpicture}
\end{minipage}\hfill
\begin{minipage}{0.61\linewidth}\centering
{\small\renewcommand{\arraystretch}{1.05}\setlength{\tabcolsep}{4pt}
\begin{tabular}{lllll}
  \toprule
  \textbf{Équipement} & \textbf{Adresse IP} & \textbf{Masque} & \textbf{Passerelle} & \textbf{Port} \\
  \midrule
  \Crows
  \bottomrule
\end{tabular}}
\end{minipage}
\end{center}

\question{C1}{4 pts}{Relevez les \textbf{quatre anomalies}. Pour chacune : conséquence et correction.}
{\small\renewcommand{\arraystretch}{1.1}
\ifcorrige\def\hc{}\else\def\hc{\rule[-10mm]{0pt}{14mm}}\fi
\begin{tabularx}{\linewidth}{|c|>{\raggedright\arraybackslash}p{2cm}|L|L|L|}
  \hline
  \textbf{N°} & \textbf{Équipement} & \textbf{Anomalie} & \textbf{Conséquence} & \textbf{Correction} \\ \hline
  \ifcorrige\Csol\else
  \hc 1 & & & & \\ \hline
  \hc 2 & & & & \\ \hline
  \hc 3 & & & & \\ \hline
  \hc 4 & & & & \\ \hline
  \fi
\end{tabularx}}
\bareme{1 pt par anomalie (0,25 équipement, anomalie, conséquence, correction ; ordre indifférent). \Cbareme}

%==========================================================================
\section{Bonus — Extension sans fil \hfill +1}

\begin{cdc}Ajoutez au réseau \resUn{} un \textbf{AccessPoint-PT} relié à SW-\pUn{} (SSID \ip{\ssid}, \textbf{WPA2-PSK})
et un portable \ip{Laptop-\pUn} équipé d'un module \texttt{WPC300N}, en adressage \textbf{statique}.
Vérifiez qu'il joint \srvNom.\end{cdc}
{\small\renewcommand{\arraystretch}{1.1}
\begin{tabularx}{\linewidth}{|l|C|C|C|c|}
  \hline
  \textbf{Équipement} & \textbf{Adresse IP} & \textbf{Masque} & \textbf{Passerelle} & \textbf{Visa prof.} \\ \hline
  \hl Laptop-\pUn & \sol{\netUn.50} & \sol{255.255.255.0} & \sol{\netUn.254} & \hspace{1.4cm} \\ \hline
\end{tabularx}}
\bareme{Le point d'accès fait office de pont : le portable est dans \netUn.0/24. Bonus si associé en WPA2-PSK et ping
vers \netSrv.\srvH{} réussi. Note plafonnée à 20.}

\end{document}
"
\documentclass[10pt,a4paper]{article}

\usepackage[T1]{fontenc}
\usepackage[utf8]{inputenc}
\usepackage{lmodern}
\usepackage[french]{babel}
\usepackage[a4paper,hmargin=1.5cm,top=1.3cm,bottom=1.6cm,footskip=0.7cm]{geometry}
\usepackage[table]{xcolor}
\usepackage{booktabs,tabularx,array}
\usepackage{enumitem}
\usepackage{fancyhdr}
\usepackage[most]{tcolorbox}
\usepackage{tikz}
\usetikzlibrary{positioning}
\usepackage{titlesec}

\providecommand{\modecorrige}{0}
\providecommand{\variante}{A}
\newif\ifcorrige
\ifnum\modecorrige=1 \corrigetrue \else \corrigefalse \fi

\definecolor{ipred}{HTML}{C0392B}
\definecolor{cdcbar}{HTML}{A67C3D}
\definecolor{cdcbg}{HTML}{F4F1EA}
\definecolor{qbar}{HTML}{4A6B5D}
\definecolor{qbg}{HTML}{EDF1EF}
\definecolor{corr}{HTML}{1F5FA8}
\definecolor{corrbg}{HTML}{E8F0FA}
\definecolor{grisfoot}{HTML}{828282}

\newcommand{\ip}[1]{\textcolor{ipred}{\texttt{#1}}}
\newcommand{\pts}[1]{\hfill\mbox{\textbf{(#1)}}}
\newcommand{\sol}[1]{\ifcorrige\textcolor{corr}{#1}\fi}
\newcommand{\bareme}[1]{\ifcorrige\par{\footnotesize\color{corr}#1\par}\fi}
\newcommand{\hl}{\ifcorrige\rule[-1.6mm]{0pt}{5.4mm}\else\rule[-2.2mm]{0pt}{6.6mm}\fi}
\newcolumntype{C}{>{\centering\arraybackslash}X}
\newcolumntype{L}{>{\raggedright\arraybackslash}X}

\titleformat{\section}{\large\bfseries}{}{0pt}{}
\titlespacing*{\section}{0pt}{6pt}{3pt}

\tcbset{qstyle/.style={enhanced,boxrule=0pt,leftrule=2.5pt,arc=0pt,outer arc=0pt,
  left=5pt,right=5pt,top=2pt,bottom=2pt,before skip=4pt,after skip=3pt}}
\newtcolorbox{cdc}{qstyle,colback=cdcbg,colframe=cdcbar}
\newtcolorbox{quest}{qstyle,colback=qbg,colframe=qbar}

% Zone de réponse : vierge dans le sujet, remplie dans le corrigé
\newcommand{\zone}[2]{%
  \ifcorrige
    \begin{tcolorbox}[colback=corrbg,colframe=corr,boxrule=0.5pt,arc=1pt,left=4pt,right=4pt,top=1pt,bottom=1pt,
      before skip=2pt,after skip=2pt,fontupper=\footnotesize\color{corr}]#2\end{tcolorbox}%
  \else
    \begin{tcolorbox}[colback=white,colframe=black!45,boxrule=0.4pt,arc=1pt,height=#1,
      before skip=2pt,after skip=2pt]\end{tcolorbox}%
  \fi}

\newcommand{\question}[3]{\begin{quest}\textbf{#1.} #3\pts{#2}\end{quest}}

%==========================================================================
% Données propres à chaque variante
%==========================================================================
\if\variante A
  \def\entreprise{Le cabinet d'architectes \textbf{ArchiNova}}
  \def\Rnom{R-ArchiNova}
  \def\resUn{Bureau d'études}\def\pUn{BE}\def\netUn{192.168.40}
  \def\resDeux{Accueil}\def\pDeux{ACC}\def\netDeux{172.16}
  \def\srvNom{SRV-Intranet}\def\netSrv{192.168.200}\def\srvH{10}
  \def\ssid{ArchiNova-WiFi}
  \def\AIrows{%
    \hl\ip{192.168.40.25} & \sol{C} & \sol{255.255.255.0} & \sol{192.168.40.0} & \sol{192.168.40.255} & \sol{254} \\ \hline
    \hl\ip{172.16.5.10}   & \sol{B} & \sol{255.255.0.0}   & \sol{172.16.0.0}   & \sol{172.16.255.255} & \sol{65\,534} \\ \hline
    \hl\ip{10.3.2.1}      & \sol{A} & \sol{255.0.0.0}     & \sol{10.0.0.0}     & \sol{10.255.255.255} & \sol{16\,777\,214} \\ \hline}
  \def\apipa{169.254.12.7}
  \def\Cnet{192.168.60}
  \def\Crows{%
    PC-A & \ip{192.168.60.10} & \ip{255.255.255.0} & \ip{192.168.60.254} & — \\
    PC-B & \ip{192.168.60.20} & \ip{255.255.255.0} & \ip{192.168.60.1}   & — \\
    PC-C & \ip{192.168.61.30} & \ip{255.255.255.0} & \ip{192.168.60.254} & — \\
    PC-D & \ip{192.168.60.10} & \ip{255.255.255.0} & \ip{192.168.60.254} & — \\
    SRV-C & \ip{10.0.0.5}     & \ip{255.0.0.0}     & \ip{10.0.0.254}     & — \\
    R1 — G0/0 & \ip{192.168.60.254} & \ip{255.255.255.0} & — & On \\
    R1 — G0/1 & \ip{10.0.0.254}     & \ip{255.0.0.0}     & — & Off \\}
  \def\Csol{%
    \hl 1 & \sol{PC-B} & \sol{Passerelle .60.1 : pas l'adresse du routeur.} & \sol{Joint son réseau, mais aucun autre (serveur).} & \sol{Passerelle 192.168.60.254.} \\ \hline
    \hl 2 & \sol{PC-C} & \sol{IP hors du réseau 192.168.60.0/24.} & \sol{Ne joint ni les postes ni le routeur.} & \sol{Ex. 192.168.60.30.} \\ \hline
    \hl 3 & \sol{PC-D} & \sol{Même IP que PC-A (doublon).} & \sol{Conflit d'adresse : PC-A et PC-D perturbés.} & \sol{Adresse unique, ex. .60.40.} \\ \hline
    \hl 4 & \sol{R1 G0/1} & \sol{Interface désactivée (Off).} & \sol{Réseau du serveur injoignable pour tous.} & \sol{Port Status : On.} \\ \hline}
  \def\Cbareme{PC-A et SRV-C sont correctement configurés ; accepter PC-A cité à la place de PC-D pour le doublon.}
\else
  \def\entreprise{La clinique vétérinaire \textbf{VetoPlus}}
  \def\Rnom{R-VetoPlus}
  \def\resUn{Soins}\def\pUn{SOI}\def\netUn{192.168.70}
  \def\resDeux{Secrétariat}\def\pDeux{SEC}\def\netDeux{172.20}
  \def\srvNom{SRV-Dossiers}\def\netSrv{192.168.250}\def\srvH{20}
  \def\ssid{VetoPlus-WiFi}
  \def\AIrows{%
    \hl\ip{192.168.15.130} & \sol{C} & \sol{255.255.255.0} & \sol{192.168.15.0} & \sol{192.168.15.255} & \sol{254} \\ \hline
    \hl\ip{172.31.200.7}   & \sol{B} & \sol{255.255.0.0}   & \sol{172.31.0.0}   & \sol{172.31.255.255} & \sol{65\,534} \\ \hline
    \hl\ip{10.45.6.9}      & \sol{A} & \sol{255.0.0.0}     & \sol{10.0.0.0}     & \sol{10.255.255.255} & \sol{16\,777\,214} \\ \hline}
  \def\apipa{169.254.88.3}
  \def\Cnet{192.168.80}
  \def\Crows{%
    PC-A & \ip{192.168.8.11}   & \ip{255.255.255.0} & \ip{192.168.80.254} & — \\
    PC-B & \ip{192.168.80.12}  & \ip{255.255.255.0} & \ip{192.168.80.254} & — \\
    PC-C & \ip{192.168.80.13}  & \ip{255.255.255.0} & \ip{192.168.80.100} & — \\
    PC-D & \ip{192.168.80.12}  & \ip{255.255.255.0} & \ip{192.168.80.254} & — \\
    SRV-C & \ip{10.0.0.8}      & \ip{255.0.0.0}     & \ip{10.0.0.254}     & — \\
    R1 — G0/0 & \ip{192.168.80.254} & \ip{255.255.255.0} & — & Off \\
    R1 — G0/1 & \ip{10.0.0.254}     & \ip{255.0.0.0}     & — & On \\}
  \def\Csol{%
    \hl 1 & \sol{PC-A} & \sol{IP hors du réseau 192.168.80.0/24.} & \sol{Ne joint ni les postes ni le routeur.} & \sol{Ex. 192.168.80.11.} \\ \hline
    \hl 2 & \sol{PC-C} & \sol{Passerelle .80.100 : pas l'adresse du routeur.} & \sol{Joint son réseau, mais aucun autre (serveur).} & \sol{Passerelle 192.168.80.254.} \\ \hline
    \hl 3 & \sol{PC-D} & \sol{Même IP que PC-B (doublon).} & \sol{Conflit d'adresse : PC-B et PC-D perturbés.} & \sol{Adresse unique, ex. .80.14.} \\ \hline
    \hl 4 & \sol{R1 G0/0} & \sol{Interface côté LAN désactivée (Off).} & \sol{Aucun poste ne joint le routeur ni le serveur.} & \sol{Port Status : On.} \\ \hline}
  \def\Cbareme{PC-B et SRV-C sont correctement configurés ; accepter PC-B cité à la place de PC-D pour le doublon.}
\fi

\pagestyle{fancy}
\fancyhf{}
\renewcommand{\headrulewidth}{0pt}
\fancyfoot[L]{\footnotesize\itshape\color{grisfoot}Évaluation — Réseaux locaux — Sujet \variante\ifcorrige\ — CORRIGÉ\fi}
\fancyfoot[R]{\footnotesize\color{grisfoot}\thepage/3}

\setlist{itemsep=0pt,topsep=2pt,parsep=0pt}
\setlength{\parindent}{0pt}
\setlength{\parskip}{2pt}

\begin{document}

\begin{center}
  {\Large\bfseries Évaluation pratique — Réseaux locaux \quad
   \fcolorbox{black}{white}{\ SUJET \variante\ }\par}
  \vspace{1pt}
  {\bfseries Adressage IP, commutation et routage — BTS CIEL 1re année — durée : 2 heures\par}
  \ifcorrige\vspace{2pt}{\bfseries\color{corr}CORRIGÉ ET BARÈME — document enseignant}\par\fi
\end{center}
\vspace{-2pt}
{\renewcommand{\arraystretch}{1.35}
\begin{tabularx}{\linewidth}{|X|X|p{2.2cm}|p{2.6cm}|}
  \hline
  \textbf{Nom :} & \textbf{Prénom :} & \textbf{Poste n°} & \textbf{Note :} \hfill /20 \\ \hline
\end{tabularx}}

\textbf{Consignes :} travail \textbf{individuel}, documents et Internet \textbf{non autorisés}. Ce sujet est le
\textbf{document-réponse}. Fichier Packet Tracer : \ip{Nom\_Prenom\_EvalLAN\_\variante.pkt}, à enregistrer
régulièrement dans le dossier de rendu. Les « visas » sont apposés par le professeur, que vous appelez à chaque test réussi.
Durées conseillées : A 30 min — B 1 h 10 — C 20 min.

%==========================================================================
\section{Partie A — Questions de cours et calculs \hfill /6}

\question{A1}{2 pts}{Complétez le tableau en considérant le \textbf{masque par défaut} de la classe de chaque adresse.}
{\small\renewcommand{\arraystretch}{1.15}
\begin{tabularx}{\linewidth}{|l|c|C|C|C|c|}
  \hline
  \textbf{Adresse IP} & \textbf{Classe} & \textbf{Masque par défaut} & \textbf{Adresse réseau} &
  \textbf{Adresse de diffusion} & \textbf{Nb d'hôtes} \\ \hline
  \AIrows
\end{tabularx}}
\bareme{0,5 pt par ligne juste + 0,5 pt si les trois nombres d'hôtes sont justes ($2^{n}-2$, $n$ = nombre de bits d'hôte).}

\question{A2}{1 pt}{Indiquez le type de câble (\textbf{droit} ou \textbf{croisé}) pour chaque liaison.}
{\small\renewcommand{\arraystretch}{1.15}
\begin{tabularx}{\linewidth}{|X|c||X|c|}
  \hline
  \hl PC $\leftrightarrow$ commutateur & \makebox[1.6cm]{\sol{droit}} & Commutateur $\leftrightarrow$ routeur & \makebox[1.6cm]{\sol{droit}} \\ \hline
  \hl PC $\leftrightarrow$ PC (réseau P2P) & \sol{croisé} & Concentrateur $\leftrightarrow$ concentrateur & \sol{croisé} \\ \hline
\end{tabularx}}
\bareme{0,25 pt par case. Équipements de types différents : câble droit ; de même type : câble croisé.}

\question{A3}{1 pt}{Comment un \textbf{concentrateur (hub)} et un \textbf{commutateur (switch)} traitent-ils une trame
reçue ? Lequel limite les collisions, et pourquoi ?}
\zone{2.9cm}{Le \textbf{hub} (couche 1) répète la trame sur \textbf{tous ses autres ports} : un seul domaine de collision
partagé. Le \textbf{switch} (couche 2) lit l'adresse MAC de destination et, grâce à sa \textbf{table MAC}, ne l'envoie
\textbf{que sur le port du destinataire} : un domaine de collision par port, c'est lui qui limite les collisions.
\emph{0,5 pt hub, 0,5 pt switch.}}

\question{A4}{1 pt}{Qu'est-ce que la \textbf{passerelle par défaut} d'un poste ? Quelle condition son adresse doit-elle respecter ?}
\zone{2.4cm}{Adresse IP de l'interface du \textbf{routeur} à laquelle le poste envoie les paquets destinés à un
\textbf{autre réseau}. Elle doit \textbf{appartenir au même réseau} que le poste. \emph{0,5 + 0,5 pt.}}

\question{A5}{0,5 pt}{Un portable configuré en DHCP affiche \ip{\apipa} / \ip{255.255.0.0}. Que s'est-il passé ?}
\zone{1.5cm}{Aucun \textbf{serveur DHCP} n'a répondu : le poste s'est attribué une adresse \textbf{APIPA} (169.254.0.0/16).}

\question{A6}{0,5 pt}{Quel protocole utilise \texttt{ping} ? Nommez les deux messages échangés et la couche OSI des adresses IP.}
\zone{1.5cm}{\textbf{ICMP} ; \textbf{Echo Request} / \textbf{Echo Reply} ; adresses IP en \textbf{couche 3 (Réseau)}.}

%==========================================================================
\newpage
\section{Partie B — Réalisation sous Packet Tracer \hfill /10}

\begin{cdc}
\entreprise{} vous confie son réseau : \textbf{trois réseaux} reliés par un \textbf{routeur unique Cisco 2911}
nommé \ip{\Rnom} (interfaces G0/0, G0/1, G0/2).
\begin{itemize}[leftmargin=12pt]
  \item \textbf{Réseau \resUn} : \textbf{4 postes} PC-\pUn1 à PC-\pUn4, commutateur \ip{SW-\pUn}, adressage statique,
        réseau \ip{\netUn.0/24}.
  \item \textbf{Réseau \resDeux} : \textbf{2 postes} PC-\pDeux1, PC-\pDeux2 et \textbf{1 imprimante} IMP-\pDeux,
        commutateur \ip{SW-\pDeux}, adressage statique, réseau de classe B \ip{\netDeux.0.0/16}.
  \item \textbf{Réseau Serveur} : serveur \ip{\srvNom} d'adresse \ip{\netSrv.\srvH/24}, relié \textbf{directement} au routeur.
  \item La passerelle de chaque réseau est sa \textbf{dernière adresse utilisable}. Tous les équipements communiquent
        entre eux et accèdent à la page web du serveur. Nommage imposé (« Display Name »).
\end{itemize}
\end{cdc}

\question{B1}{2 pts}{Complétez le plan d'adressage \textbf{avant} la réalisation.}
{\footnotesize\setlength{\tabcolsep}{3pt}\renewcommand{\arraystretch}{1.15}
\begin{tabularx}{\linewidth}{|l|C|C|C|C|C|C|}
  \hline
  \textbf{Réseau} & \textbf{Adresse réseau} & \textbf{Masque} & \textbf{1re @ utilisable} &
  \textbf{Dernière @ utilisable} & \textbf{Diffusion} & \textbf{Passerelle} \\ \hline
  \hl \resUn & \sol{\netUn.0} & \sol{255.255.255.0} & \sol{\netUn.1} & \sol{\netUn.254} & \sol{\netUn.255} & \sol{\netUn.254} \\ \hline
  \hl \resDeux & \sol{\netDeux.0.0} & \sol{255.255.0.0} & \sol{\netDeux.0.1} & \sol{\netDeux.255.254} & \sol{\netDeux.255.255} & \sol{\netDeux.255.254} \\ \hline
  \hl Serveur & \sol{\netSrv.0} & \sol{255.255.255.0} & \sol{\netSrv.1} & \sol{\netSrv.254} & \sol{\netSrv.255} & \sol{\netSrv.254} \\ \hline
\end{tabularx}}
\bareme{0,5 pt par ligne juste + 0,5 pt si la convention de passerelle est respectée partout.}

\question{B2}{2 pts}{Construisez la topologie (équipements, câblage, nommage) et relevez l'adressage configuré.}
{\footnotesize\setlength{\tabcolsep}{3pt}\renewcommand{\arraystretch}{1.1}
\begin{tabularx}{\linewidth}{|l|C|C|C||l|C|C|}
  \hline
  \textbf{Hôte} & \textbf{Adresse IP} & \textbf{Masque} & \textbf{Passerelle} & \textbf{Interface} & \textbf{Adresse IP} & \textbf{Masque} \\ \hline
  \hl PC-\pUn1 & \sol{\netUn.1} & \sol{255.255.255.0} & \sol{\netUn.254} & \Rnom{} G0/0 & \sol{\netUn.254} & \sol{255.255.255.0} \\ \hline
  \hl PC-\pDeux1 & \sol{\netDeux.0.1} & \sol{255.255.0.0} & \sol{\netDeux.255.254} & \Rnom{} G0/1 & \sol{\netDeux.255.254} & \sol{255.255.0.0} \\ \hline
  \hl IMP-\pDeux & \sol{\netDeux.0.10} & \sol{255.255.0.0} & \sol{\netDeux.255.254} & \Rnom{} G0/2 & \sol{\netSrv.254} & \sol{255.255.255.0} \\ \hline
  \hl \srvNom & \ip{\netSrv.\srvH} & \sol{255.255.255.0} & \sol{\netSrv.254} & \multicolumn{3}{c|}{\cellcolor{black!8}} \\ \hline
\end{tabularx}}
\bareme{Adresses d'hôtes libres (uniques, dans le bon réseau) ; affectation G0/x libre. Sur le .pkt : équipements 0,5 ;
câblage 1 (câbles droits vers les switchs et switch–routeur, \textbf{croisé} serveur–routeur, liens verts) ; nommage 0,5.}

\question{B3}{3 pts}{Configurez tous les équipements : IP, masques, passerelles des hôtes ; interfaces du routeur
(« Port Status : On »).}
\bareme{Sur le .pkt : interfaces routeur 1 pt ; IP/masques des hôtes 1 pt ; passerelles (imprimante et serveur compris)
1 pt. $-0,25$ par erreur, sans descendre sous 0 par item.}

\question{B4}{3 pts}{Réalisez les tests, notez la commande et le résultat, puis faites viser chaque test réussi.}
{\small\renewcommand{\arraystretch}{1.1}
\begin{tabularx}{\linewidth}{|c|l|l|C|c|c|}
  \hline
  \textbf{N°} & \textbf{Source} & \textbf{Destination} & \textbf{Commande / action} & \textbf{Résultat} & \textbf{Visa prof.} \\ \hline
  \hl 1 & PC-\pUn1 & PC-\pUn4 & \sol{ping \netUn.4} & \sol{réussi} & \hspace{1.4cm} \\ \hline
  \hl 2 & PC-\pDeux1 & IMP-\pDeux & \sol{ping \netDeux.0.10} & \sol{réussi} & \\ \hline
  \hl 3 & PC-\pUn2 & PC-\pDeux2 & \sol{ping \netDeux.0.2} & \sol{réussi} & \\ \hline
  \hl 4 & PC-\pDeux2 & \srvNom & \sol{ping \netSrv.\srvH} & \sol{réussi} & \\ \hline
  \hl 5 & PC-\pUn3 & page web du serveur & \sol{navigateur : {\NoAutoSpacing http://\netSrv.\srvH}} & \sol{page affichée} & \\ \hline
\end{tabularx}}
\bareme{0,5 pt par test visé + 0,5 pt pour B5. Le premier ping inter-réseaux peut échouer (résolution ARP) ; les suivants doivent réussir.}

\question{B5}{incl. B4}{Test n°3 : par quels équipements passe le paquet ? Pourquoi PC-\pUn2 l'envoie-t-il au routeur ?}
\zone{2.3cm}{PC-\pUn2 $\to$ SW-\pUn{} $\to$ \Rnom{} $\to$ SW-\pDeux{} $\to$ PC-\pDeux2. Grâce à son masque, PC-\pUn2 constate que
la destination (\netDeux.0.0/16) n'est \textbf{pas dans son réseau} (\netUn.0/24) : il envoie le paquet à sa
\textbf{passerelle par défaut}, qui le route.}

%==========================================================================
\newpage
\section{Partie C — Diagnostic d'un réseau en panne \hfill /4}

Un technicien constate qu'\textbf{aucun poste ne parvient à joindre le serveur SRV-C} et que certains postes
communiquent mal entre eux. Les câbles sont du bon type.

\begin{center}
\begin{minipage}{0.37\linewidth}\centering
\begin{tikzpicture}[font=\footnotesize,
  eq/.style={draw,fill=qbg,rounded corners=2pt,minimum width=9mm,minimum height=5mm,inner sep=2pt},
  every path/.style={thick}]
  \node[eq] (SW) {SW-1};
  \node[eq,right=9mm of SW] (R) {R1};
  \node[eq,right=7mm of R] (SRV) {SRV-C};
  \node[eq,above left=5mm and -3mm of SW] (A) {PC-A};
  \node[eq,above right=5mm and -3mm of SW] (B) {PC-B};
  \node[eq,below left=5mm and -3mm of SW] (C) {PC-C};
  \node[eq,below right=5mm and -3mm of SW] (D) {PC-D};
  \foreach \n in {A,B,C,D} \draw (\n) -- (SW);
  \draw (SW) -- node[above,font=\tiny]{G0/0} (R);
  \draw (R) -- node[above,font=\tiny]{G0/1} (SRV);
\end{tikzpicture}
\end{minipage}\hfill
\begin{minipage}{0.61\linewidth}\centering
{\small\renewcommand{\arraystretch}{1.05}\setlength{\tabcolsep}{4pt}
\begin{tabular}{lllll}
  \toprule
  \textbf{Équipement} & \textbf{Adresse IP} & \textbf{Masque} & \textbf{Passerelle} & \textbf{Port} \\
  \midrule
  \Crows
  \bottomrule
\end{tabular}}
\end{minipage}
\end{center}

\question{C1}{4 pts}{Relevez les \textbf{quatre anomalies}. Pour chacune : conséquence et correction.}
{\small\renewcommand{\arraystretch}{1.1}
\ifcorrige\def\hc{}\else\def\hc{\rule[-10mm]{0pt}{14mm}}\fi
\begin{tabularx}{\linewidth}{|c|>{\raggedright\arraybackslash}p{2cm}|L|L|L|}
  \hline
  \textbf{N°} & \textbf{Équipement} & \textbf{Anomalie} & \textbf{Conséquence} & \textbf{Correction} \\ \hline
  \ifcorrige\Csol\else
  \hc 1 & & & & \\ \hline
  \hc 2 & & & & \\ \hline
  \hc 3 & & & & \\ \hline
  \hc 4 & & & & \\ \hline
  \fi
\end{tabularx}}
\bareme{1 pt par anomalie (0,25 équipement, anomalie, conséquence, correction ; ordre indifférent). \Cbareme}

%==========================================================================
\section{Bonus — Extension sans fil \hfill +1}

\begin{cdc}Ajoutez au réseau \resUn{} un \textbf{AccessPoint-PT} relié à SW-\pUn{} (SSID \ip{\ssid}, \textbf{WPA2-PSK})
et un portable \ip{Laptop-\pUn} équipé d'un module \texttt{WPC300N}, en adressage \textbf{statique}.
Vérifiez qu'il joint \srvNom.\end{cdc}
{\small\renewcommand{\arraystretch}{1.1}
\begin{tabularx}{\linewidth}{|l|C|C|C|c|}
  \hline
  \textbf{Équipement} & \textbf{Adresse IP} & \textbf{Masque} & \textbf{Passerelle} & \textbf{Visa prof.} \\ \hline
  \hl Laptop-\pUn & \sol{\netUn.50} & \sol{255.255.255.0} & \sol{\netUn.254} & \hspace{1.4cm} \\ \hline
\end{tabularx}}
\bareme{Le point d'accès fait office de pont : le portable est dans \netUn.0/24. Bonus si associé en WPA2-PSK et ping
vers \netSrv.\srvH{} réussi. Note plafonnée à 20.}

\end{document}
"
%   pdflatex -jobname=Eval_Reseau_SujetB   "\def\variante{B}\def\modecorrige{0}% Évaluation pratique — Réseaux locaux (BTS CIEL 1re année) — sujets A et B
% Compilation (4 PDF) :
%   pdflatex -jobname=Eval_Reseau_SujetA   "\def\variante{A}\def\modecorrige{0}% Évaluation pratique — Réseaux locaux (BTS CIEL 1re année) — sujets A et B
% Compilation (4 PDF) :
%   pdflatex -jobname=Eval_Reseau_SujetA   "\def\variante{A}\def\modecorrige{0}\input{Eval_TP_Reseau_LAN.tex}"
%   pdflatex -jobname=Eval_Reseau_SujetB   "\def\variante{B}\def\modecorrige{0}\input{Eval_TP_Reseau_LAN.tex}"
%   pdflatex -jobname=Eval_Reseau_CorrigeA "\def\variante{A}\def\modecorrige{1}\input{Eval_TP_Reseau_LAN.tex}"
%   pdflatex -jobname=Eval_Reseau_CorrigeB "\def\variante{B}\def\modecorrige{1}\input{Eval_TP_Reseau_LAN.tex}"
\documentclass[10pt,a4paper]{article}

\usepackage[T1]{fontenc}
\usepackage[utf8]{inputenc}
\usepackage{lmodern}
\usepackage[french]{babel}
\usepackage[a4paper,hmargin=1.5cm,top=1.3cm,bottom=1.6cm,footskip=0.7cm]{geometry}
\usepackage[table]{xcolor}
\usepackage{booktabs,tabularx,array}
\usepackage{enumitem}
\usepackage{fancyhdr}
\usepackage[most]{tcolorbox}
\usepackage{tikz}
\usetikzlibrary{positioning}
\usepackage{titlesec}

\providecommand{\modecorrige}{0}
\providecommand{\variante}{A}
\newif\ifcorrige
\ifnum\modecorrige=1 \corrigetrue \else \corrigefalse \fi

\definecolor{ipred}{HTML}{C0392B}
\definecolor{cdcbar}{HTML}{A67C3D}
\definecolor{cdcbg}{HTML}{F4F1EA}
\definecolor{qbar}{HTML}{4A6B5D}
\definecolor{qbg}{HTML}{EDF1EF}
\definecolor{corr}{HTML}{1F5FA8}
\definecolor{corrbg}{HTML}{E8F0FA}
\definecolor{grisfoot}{HTML}{828282}

\newcommand{\ip}[1]{\textcolor{ipred}{\texttt{#1}}}
\newcommand{\pts}[1]{\hfill\mbox{\textbf{(#1)}}}
\newcommand{\sol}[1]{\ifcorrige\textcolor{corr}{#1}\fi}
\newcommand{\bareme}[1]{\ifcorrige\par{\footnotesize\color{corr}#1\par}\fi}
\newcommand{\hl}{\ifcorrige\rule[-1.6mm]{0pt}{5.4mm}\else\rule[-2.2mm]{0pt}{6.6mm}\fi}
\newcolumntype{C}{>{\centering\arraybackslash}X}
\newcolumntype{L}{>{\raggedright\arraybackslash}X}

\titleformat{\section}{\large\bfseries}{}{0pt}{}
\titlespacing*{\section}{0pt}{6pt}{3pt}

\tcbset{qstyle/.style={enhanced,boxrule=0pt,leftrule=2.5pt,arc=0pt,outer arc=0pt,
  left=5pt,right=5pt,top=2pt,bottom=2pt,before skip=4pt,after skip=3pt}}
\newtcolorbox{cdc}{qstyle,colback=cdcbg,colframe=cdcbar}
\newtcolorbox{quest}{qstyle,colback=qbg,colframe=qbar}

% Zone de réponse : vierge dans le sujet, remplie dans le corrigé
\newcommand{\zone}[2]{%
  \ifcorrige
    \begin{tcolorbox}[colback=corrbg,colframe=corr,boxrule=0.5pt,arc=1pt,left=4pt,right=4pt,top=1pt,bottom=1pt,
      before skip=2pt,after skip=2pt,fontupper=\footnotesize\color{corr}]#2\end{tcolorbox}%
  \else
    \begin{tcolorbox}[colback=white,colframe=black!45,boxrule=0.4pt,arc=1pt,height=#1,
      before skip=2pt,after skip=2pt]\end{tcolorbox}%
  \fi}

\newcommand{\question}[3]{\begin{quest}\textbf{#1.} #3\pts{#2}\end{quest}}

%==========================================================================
% Données propres à chaque variante
%==========================================================================
\if\variante A
  \def\entreprise{Le cabinet d'architectes \textbf{ArchiNova}}
  \def\Rnom{R-ArchiNova}
  \def\resUn{Bureau d'études}\def\pUn{BE}\def\netUn{192.168.40}
  \def\resDeux{Accueil}\def\pDeux{ACC}\def\netDeux{172.16}
  \def\srvNom{SRV-Intranet}\def\netSrv{192.168.200}\def\srvH{10}
  \def\ssid{ArchiNova-WiFi}
  \def\AIrows{%
    \hl\ip{192.168.40.25} & \sol{C} & \sol{255.255.255.0} & \sol{192.168.40.0} & \sol{192.168.40.255} & \sol{254} \\ \hline
    \hl\ip{172.16.5.10}   & \sol{B} & \sol{255.255.0.0}   & \sol{172.16.0.0}   & \sol{172.16.255.255} & \sol{65\,534} \\ \hline
    \hl\ip{10.3.2.1}      & \sol{A} & \sol{255.0.0.0}     & \sol{10.0.0.0}     & \sol{10.255.255.255} & \sol{16\,777\,214} \\ \hline}
  \def\apipa{169.254.12.7}
  \def\Cnet{192.168.60}
  \def\Crows{%
    PC-A & \ip{192.168.60.10} & \ip{255.255.255.0} & \ip{192.168.60.254} & — \\
    PC-B & \ip{192.168.60.20} & \ip{255.255.255.0} & \ip{192.168.60.1}   & — \\
    PC-C & \ip{192.168.61.30} & \ip{255.255.255.0} & \ip{192.168.60.254} & — \\
    PC-D & \ip{192.168.60.10} & \ip{255.255.255.0} & \ip{192.168.60.254} & — \\
    SRV-C & \ip{10.0.0.5}     & \ip{255.0.0.0}     & \ip{10.0.0.254}     & — \\
    R1 — G0/0 & \ip{192.168.60.254} & \ip{255.255.255.0} & — & On \\
    R1 — G0/1 & \ip{10.0.0.254}     & \ip{255.0.0.0}     & — & Off \\}
  \def\Csol{%
    \hl 1 & \sol{PC-B} & \sol{Passerelle .60.1 : pas l'adresse du routeur.} & \sol{Joint son réseau, mais aucun autre (serveur).} & \sol{Passerelle 192.168.60.254.} \\ \hline
    \hl 2 & \sol{PC-C} & \sol{IP hors du réseau 192.168.60.0/24.} & \sol{Ne joint ni les postes ni le routeur.} & \sol{Ex. 192.168.60.30.} \\ \hline
    \hl 3 & \sol{PC-D} & \sol{Même IP que PC-A (doublon).} & \sol{Conflit d'adresse : PC-A et PC-D perturbés.} & \sol{Adresse unique, ex. .60.40.} \\ \hline
    \hl 4 & \sol{R1 G0/1} & \sol{Interface désactivée (Off).} & \sol{Réseau du serveur injoignable pour tous.} & \sol{Port Status : On.} \\ \hline}
  \def\Cbareme{PC-A et SRV-C sont correctement configurés ; accepter PC-A cité à la place de PC-D pour le doublon.}
\else
  \def\entreprise{La clinique vétérinaire \textbf{VetoPlus}}
  \def\Rnom{R-VetoPlus}
  \def\resUn{Soins}\def\pUn{SOI}\def\netUn{192.168.70}
  \def\resDeux{Secrétariat}\def\pDeux{SEC}\def\netDeux{172.20}
  \def\srvNom{SRV-Dossiers}\def\netSrv{192.168.250}\def\srvH{20}
  \def\ssid{VetoPlus-WiFi}
  \def\AIrows{%
    \hl\ip{192.168.15.130} & \sol{C} & \sol{255.255.255.0} & \sol{192.168.15.0} & \sol{192.168.15.255} & \sol{254} \\ \hline
    \hl\ip{172.31.200.7}   & \sol{B} & \sol{255.255.0.0}   & \sol{172.31.0.0}   & \sol{172.31.255.255} & \sol{65\,534} \\ \hline
    \hl\ip{10.45.6.9}      & \sol{A} & \sol{255.0.0.0}     & \sol{10.0.0.0}     & \sol{10.255.255.255} & \sol{16\,777\,214} \\ \hline}
  \def\apipa{169.254.88.3}
  \def\Cnet{192.168.80}
  \def\Crows{%
    PC-A & \ip{192.168.8.11}   & \ip{255.255.255.0} & \ip{192.168.80.254} & — \\
    PC-B & \ip{192.168.80.12}  & \ip{255.255.255.0} & \ip{192.168.80.254} & — \\
    PC-C & \ip{192.168.80.13}  & \ip{255.255.255.0} & \ip{192.168.80.100} & — \\
    PC-D & \ip{192.168.80.12}  & \ip{255.255.255.0} & \ip{192.168.80.254} & — \\
    SRV-C & \ip{10.0.0.8}      & \ip{255.0.0.0}     & \ip{10.0.0.254}     & — \\
    R1 — G0/0 & \ip{192.168.80.254} & \ip{255.255.255.0} & — & Off \\
    R1 — G0/1 & \ip{10.0.0.254}     & \ip{255.0.0.0}     & — & On \\}
  \def\Csol{%
    \hl 1 & \sol{PC-A} & \sol{IP hors du réseau 192.168.80.0/24.} & \sol{Ne joint ni les postes ni le routeur.} & \sol{Ex. 192.168.80.11.} \\ \hline
    \hl 2 & \sol{PC-C} & \sol{Passerelle .80.100 : pas l'adresse du routeur.} & \sol{Joint son réseau, mais aucun autre (serveur).} & \sol{Passerelle 192.168.80.254.} \\ \hline
    \hl 3 & \sol{PC-D} & \sol{Même IP que PC-B (doublon).} & \sol{Conflit d'adresse : PC-B et PC-D perturbés.} & \sol{Adresse unique, ex. .80.14.} \\ \hline
    \hl 4 & \sol{R1 G0/0} & \sol{Interface côté LAN désactivée (Off).} & \sol{Aucun poste ne joint le routeur ni le serveur.} & \sol{Port Status : On.} \\ \hline}
  \def\Cbareme{PC-B et SRV-C sont correctement configurés ; accepter PC-B cité à la place de PC-D pour le doublon.}
\fi

\pagestyle{fancy}
\fancyhf{}
\renewcommand{\headrulewidth}{0pt}
\fancyfoot[L]{\footnotesize\itshape\color{grisfoot}Évaluation — Réseaux locaux — Sujet \variante\ifcorrige\ — CORRIGÉ\fi}
\fancyfoot[R]{\footnotesize\color{grisfoot}\thepage/3}

\setlist{itemsep=0pt,topsep=2pt,parsep=0pt}
\setlength{\parindent}{0pt}
\setlength{\parskip}{2pt}

\begin{document}

\begin{center}
  {\Large\bfseries Évaluation pratique — Réseaux locaux \quad
   \fcolorbox{black}{white}{\ SUJET \variante\ }\par}
  \vspace{1pt}
  {\bfseries Adressage IP, commutation et routage — BTS CIEL 1re année — durée : 2 heures\par}
  \ifcorrige\vspace{2pt}{\bfseries\color{corr}CORRIGÉ ET BARÈME — document enseignant}\par\fi
\end{center}
\vspace{-2pt}
{\renewcommand{\arraystretch}{1.35}
\begin{tabularx}{\linewidth}{|X|X|p{2.2cm}|p{2.6cm}|}
  \hline
  \textbf{Nom :} & \textbf{Prénom :} & \textbf{Poste n°} & \textbf{Note :} \hfill /20 \\ \hline
\end{tabularx}}

\textbf{Consignes :} travail \textbf{individuel}, documents et Internet \textbf{non autorisés}. Ce sujet est le
\textbf{document-réponse}. Fichier Packet Tracer : \ip{Nom\_Prenom\_EvalLAN\_\variante.pkt}, à enregistrer
régulièrement dans le dossier de rendu. Les « visas » sont apposés par le professeur, que vous appelez à chaque test réussi.
Durées conseillées : A 30 min — B 1 h 10 — C 20 min.

%==========================================================================
\section{Partie A — Questions de cours et calculs \hfill /6}

\question{A1}{2 pts}{Complétez le tableau en considérant le \textbf{masque par défaut} de la classe de chaque adresse.}
{\small\renewcommand{\arraystretch}{1.15}
\begin{tabularx}{\linewidth}{|l|c|C|C|C|c|}
  \hline
  \textbf{Adresse IP} & \textbf{Classe} & \textbf{Masque par défaut} & \textbf{Adresse réseau} &
  \textbf{Adresse de diffusion} & \textbf{Nb d'hôtes} \\ \hline
  \AIrows
\end{tabularx}}
\bareme{0,5 pt par ligne juste + 0,5 pt si les trois nombres d'hôtes sont justes ($2^{n}-2$, $n$ = nombre de bits d'hôte).}

\question{A2}{1 pt}{Indiquez le type de câble (\textbf{droit} ou \textbf{croisé}) pour chaque liaison.}
{\small\renewcommand{\arraystretch}{1.15}
\begin{tabularx}{\linewidth}{|X|c||X|c|}
  \hline
  \hl PC $\leftrightarrow$ commutateur & \makebox[1.6cm]{\sol{droit}} & Commutateur $\leftrightarrow$ routeur & \makebox[1.6cm]{\sol{droit}} \\ \hline
  \hl PC $\leftrightarrow$ PC (réseau P2P) & \sol{croisé} & Concentrateur $\leftrightarrow$ concentrateur & \sol{croisé} \\ \hline
\end{tabularx}}
\bareme{0,25 pt par case. Équipements de types différents : câble droit ; de même type : câble croisé.}

\question{A3}{1 pt}{Comment un \textbf{concentrateur (hub)} et un \textbf{commutateur (switch)} traitent-ils une trame
reçue ? Lequel limite les collisions, et pourquoi ?}
\zone{2.9cm}{Le \textbf{hub} (couche 1) répète la trame sur \textbf{tous ses autres ports} : un seul domaine de collision
partagé. Le \textbf{switch} (couche 2) lit l'adresse MAC de destination et, grâce à sa \textbf{table MAC}, ne l'envoie
\textbf{que sur le port du destinataire} : un domaine de collision par port, c'est lui qui limite les collisions.
\emph{0,5 pt hub, 0,5 pt switch.}}

\question{A4}{1 pt}{Qu'est-ce que la \textbf{passerelle par défaut} d'un poste ? Quelle condition son adresse doit-elle respecter ?}
\zone{2.4cm}{Adresse IP de l'interface du \textbf{routeur} à laquelle le poste envoie les paquets destinés à un
\textbf{autre réseau}. Elle doit \textbf{appartenir au même réseau} que le poste. \emph{0,5 + 0,5 pt.}}

\question{A5}{0,5 pt}{Un portable configuré en DHCP affiche \ip{\apipa} / \ip{255.255.0.0}. Que s'est-il passé ?}
\zone{1.5cm}{Aucun \textbf{serveur DHCP} n'a répondu : le poste s'est attribué une adresse \textbf{APIPA} (169.254.0.0/16).}

\question{A6}{0,5 pt}{Quel protocole utilise \texttt{ping} ? Nommez les deux messages échangés et la couche OSI des adresses IP.}
\zone{1.5cm}{\textbf{ICMP} ; \textbf{Echo Request} / \textbf{Echo Reply} ; adresses IP en \textbf{couche 3 (Réseau)}.}

%==========================================================================
\newpage
\section{Partie B — Réalisation sous Packet Tracer \hfill /10}

\begin{cdc}
\entreprise{} vous confie son réseau : \textbf{trois réseaux} reliés par un \textbf{routeur unique Cisco 2911}
nommé \ip{\Rnom} (interfaces G0/0, G0/1, G0/2).
\begin{itemize}[leftmargin=12pt]
  \item \textbf{Réseau \resUn} : \textbf{4 postes} PC-\pUn1 à PC-\pUn4, commutateur \ip{SW-\pUn}, adressage statique,
        réseau \ip{\netUn.0/24}.
  \item \textbf{Réseau \resDeux} : \textbf{2 postes} PC-\pDeux1, PC-\pDeux2 et \textbf{1 imprimante} IMP-\pDeux,
        commutateur \ip{SW-\pDeux}, adressage statique, réseau de classe B \ip{\netDeux.0.0/16}.
  \item \textbf{Réseau Serveur} : serveur \ip{\srvNom} d'adresse \ip{\netSrv.\srvH/24}, relié \textbf{directement} au routeur.
  \item La passerelle de chaque réseau est sa \textbf{dernière adresse utilisable}. Tous les équipements communiquent
        entre eux et accèdent à la page web du serveur. Nommage imposé (« Display Name »).
\end{itemize}
\end{cdc}

\question{B1}{2 pts}{Complétez le plan d'adressage \textbf{avant} la réalisation.}
{\footnotesize\setlength{\tabcolsep}{3pt}\renewcommand{\arraystretch}{1.15}
\begin{tabularx}{\linewidth}{|l|C|C|C|C|C|C|}
  \hline
  \textbf{Réseau} & \textbf{Adresse réseau} & \textbf{Masque} & \textbf{1re @ utilisable} &
  \textbf{Dernière @ utilisable} & \textbf{Diffusion} & \textbf{Passerelle} \\ \hline
  \hl \resUn & \sol{\netUn.0} & \sol{255.255.255.0} & \sol{\netUn.1} & \sol{\netUn.254} & \sol{\netUn.255} & \sol{\netUn.254} \\ \hline
  \hl \resDeux & \sol{\netDeux.0.0} & \sol{255.255.0.0} & \sol{\netDeux.0.1} & \sol{\netDeux.255.254} & \sol{\netDeux.255.255} & \sol{\netDeux.255.254} \\ \hline
  \hl Serveur & \sol{\netSrv.0} & \sol{255.255.255.0} & \sol{\netSrv.1} & \sol{\netSrv.254} & \sol{\netSrv.255} & \sol{\netSrv.254} \\ \hline
\end{tabularx}}
\bareme{0,5 pt par ligne juste + 0,5 pt si la convention de passerelle est respectée partout.}

\question{B2}{2 pts}{Construisez la topologie (équipements, câblage, nommage) et relevez l'adressage configuré.}
{\footnotesize\setlength{\tabcolsep}{3pt}\renewcommand{\arraystretch}{1.1}
\begin{tabularx}{\linewidth}{|l|C|C|C||l|C|C|}
  \hline
  \textbf{Hôte} & \textbf{Adresse IP} & \textbf{Masque} & \textbf{Passerelle} & \textbf{Interface} & \textbf{Adresse IP} & \textbf{Masque} \\ \hline
  \hl PC-\pUn1 & \sol{\netUn.1} & \sol{255.255.255.0} & \sol{\netUn.254} & \Rnom{} G0/0 & \sol{\netUn.254} & \sol{255.255.255.0} \\ \hline
  \hl PC-\pDeux1 & \sol{\netDeux.0.1} & \sol{255.255.0.0} & \sol{\netDeux.255.254} & \Rnom{} G0/1 & \sol{\netDeux.255.254} & \sol{255.255.0.0} \\ \hline
  \hl IMP-\pDeux & \sol{\netDeux.0.10} & \sol{255.255.0.0} & \sol{\netDeux.255.254} & \Rnom{} G0/2 & \sol{\netSrv.254} & \sol{255.255.255.0} \\ \hline
  \hl \srvNom & \ip{\netSrv.\srvH} & \sol{255.255.255.0} & \sol{\netSrv.254} & \multicolumn{3}{c|}{\cellcolor{black!8}} \\ \hline
\end{tabularx}}
\bareme{Adresses d'hôtes libres (uniques, dans le bon réseau) ; affectation G0/x libre. Sur le .pkt : équipements 0,5 ;
câblage 1 (câbles droits vers les switchs et switch–routeur, \textbf{croisé} serveur–routeur, liens verts) ; nommage 0,5.}

\question{B3}{3 pts}{Configurez tous les équipements : IP, masques, passerelles des hôtes ; interfaces du routeur
(« Port Status : On »).}
\bareme{Sur le .pkt : interfaces routeur 1 pt ; IP/masques des hôtes 1 pt ; passerelles (imprimante et serveur compris)
1 pt. $-0,25$ par erreur, sans descendre sous 0 par item.}

\question{B4}{3 pts}{Réalisez les tests, notez la commande et le résultat, puis faites viser chaque test réussi.}
{\small\renewcommand{\arraystretch}{1.1}
\begin{tabularx}{\linewidth}{|c|l|l|C|c|c|}
  \hline
  \textbf{N°} & \textbf{Source} & \textbf{Destination} & \textbf{Commande / action} & \textbf{Résultat} & \textbf{Visa prof.} \\ \hline
  \hl 1 & PC-\pUn1 & PC-\pUn4 & \sol{ping \netUn.4} & \sol{réussi} & \hspace{1.4cm} \\ \hline
  \hl 2 & PC-\pDeux1 & IMP-\pDeux & \sol{ping \netDeux.0.10} & \sol{réussi} & \\ \hline
  \hl 3 & PC-\pUn2 & PC-\pDeux2 & \sol{ping \netDeux.0.2} & \sol{réussi} & \\ \hline
  \hl 4 & PC-\pDeux2 & \srvNom & \sol{ping \netSrv.\srvH} & \sol{réussi} & \\ \hline
  \hl 5 & PC-\pUn3 & page web du serveur & \sol{navigateur : {\NoAutoSpacing http://\netSrv.\srvH}} & \sol{page affichée} & \\ \hline
\end{tabularx}}
\bareme{0,5 pt par test visé + 0,5 pt pour B5. Le premier ping inter-réseaux peut échouer (résolution ARP) ; les suivants doivent réussir.}

\question{B5}{incl. B4}{Test n°3 : par quels équipements passe le paquet ? Pourquoi PC-\pUn2 l'envoie-t-il au routeur ?}
\zone{2.3cm}{PC-\pUn2 $\to$ SW-\pUn{} $\to$ \Rnom{} $\to$ SW-\pDeux{} $\to$ PC-\pDeux2. Grâce à son masque, PC-\pUn2 constate que
la destination (\netDeux.0.0/16) n'est \textbf{pas dans son réseau} (\netUn.0/24) : il envoie le paquet à sa
\textbf{passerelle par défaut}, qui le route.}

%==========================================================================
\newpage
\section{Partie C — Diagnostic d'un réseau en panne \hfill /4}

Un technicien constate qu'\textbf{aucun poste ne parvient à joindre le serveur SRV-C} et que certains postes
communiquent mal entre eux. Les câbles sont du bon type.

\begin{center}
\begin{minipage}{0.37\linewidth}\centering
\begin{tikzpicture}[font=\footnotesize,
  eq/.style={draw,fill=qbg,rounded corners=2pt,minimum width=9mm,minimum height=5mm,inner sep=2pt},
  every path/.style={thick}]
  \node[eq] (SW) {SW-1};
  \node[eq,right=9mm of SW] (R) {R1};
  \node[eq,right=7mm of R] (SRV) {SRV-C};
  \node[eq,above left=5mm and -3mm of SW] (A) {PC-A};
  \node[eq,above right=5mm and -3mm of SW] (B) {PC-B};
  \node[eq,below left=5mm and -3mm of SW] (C) {PC-C};
  \node[eq,below right=5mm and -3mm of SW] (D) {PC-D};
  \foreach \n in {A,B,C,D} \draw (\n) -- (SW);
  \draw (SW) -- node[above,font=\tiny]{G0/0} (R);
  \draw (R) -- node[above,font=\tiny]{G0/1} (SRV);
\end{tikzpicture}
\end{minipage}\hfill
\begin{minipage}{0.61\linewidth}\centering
{\small\renewcommand{\arraystretch}{1.05}\setlength{\tabcolsep}{4pt}
\begin{tabular}{lllll}
  \toprule
  \textbf{Équipement} & \textbf{Adresse IP} & \textbf{Masque} & \textbf{Passerelle} & \textbf{Port} \\
  \midrule
  \Crows
  \bottomrule
\end{tabular}}
\end{minipage}
\end{center}

\question{C1}{4 pts}{Relevez les \textbf{quatre anomalies}. Pour chacune : conséquence et correction.}
{\small\renewcommand{\arraystretch}{1.1}
\ifcorrige\def\hc{}\else\def\hc{\rule[-10mm]{0pt}{14mm}}\fi
\begin{tabularx}{\linewidth}{|c|>{\raggedright\arraybackslash}p{2cm}|L|L|L|}
  \hline
  \textbf{N°} & \textbf{Équipement} & \textbf{Anomalie} & \textbf{Conséquence} & \textbf{Correction} \\ \hline
  \ifcorrige\Csol\else
  \hc 1 & & & & \\ \hline
  \hc 2 & & & & \\ \hline
  \hc 3 & & & & \\ \hline
  \hc 4 & & & & \\ \hline
  \fi
\end{tabularx}}
\bareme{1 pt par anomalie (0,25 équipement, anomalie, conséquence, correction ; ordre indifférent). \Cbareme}

%==========================================================================
\section{Bonus — Extension sans fil \hfill +1}

\begin{cdc}Ajoutez au réseau \resUn{} un \textbf{AccessPoint-PT} relié à SW-\pUn{} (SSID \ip{\ssid}, \textbf{WPA2-PSK})
et un portable \ip{Laptop-\pUn} équipé d'un module \texttt{WPC300N}, en adressage \textbf{statique}.
Vérifiez qu'il joint \srvNom.\end{cdc}
{\small\renewcommand{\arraystretch}{1.1}
\begin{tabularx}{\linewidth}{|l|C|C|C|c|}
  \hline
  \textbf{Équipement} & \textbf{Adresse IP} & \textbf{Masque} & \textbf{Passerelle} & \textbf{Visa prof.} \\ \hline
  \hl Laptop-\pUn & \sol{\netUn.50} & \sol{255.255.255.0} & \sol{\netUn.254} & \hspace{1.4cm} \\ \hline
\end{tabularx}}
\bareme{Le point d'accès fait office de pont : le portable est dans \netUn.0/24. Bonus si associé en WPA2-PSK et ping
vers \netSrv.\srvH{} réussi. Note plafonnée à 20.}

\end{document}
"
%   pdflatex -jobname=Eval_Reseau_SujetB   "\def\variante{B}\def\modecorrige{0}% Évaluation pratique — Réseaux locaux (BTS CIEL 1re année) — sujets A et B
% Compilation (4 PDF) :
%   pdflatex -jobname=Eval_Reseau_SujetA   "\def\variante{A}\def\modecorrige{0}\input{Eval_TP_Reseau_LAN.tex}"
%   pdflatex -jobname=Eval_Reseau_SujetB   "\def\variante{B}\def\modecorrige{0}\input{Eval_TP_Reseau_LAN.tex}"
%   pdflatex -jobname=Eval_Reseau_CorrigeA "\def\variante{A}\def\modecorrige{1}\input{Eval_TP_Reseau_LAN.tex}"
%   pdflatex -jobname=Eval_Reseau_CorrigeB "\def\variante{B}\def\modecorrige{1}\input{Eval_TP_Reseau_LAN.tex}"
\documentclass[10pt,a4paper]{article}

\usepackage[T1]{fontenc}
\usepackage[utf8]{inputenc}
\usepackage{lmodern}
\usepackage[french]{babel}
\usepackage[a4paper,hmargin=1.5cm,top=1.3cm,bottom=1.6cm,footskip=0.7cm]{geometry}
\usepackage[table]{xcolor}
\usepackage{booktabs,tabularx,array}
\usepackage{enumitem}
\usepackage{fancyhdr}
\usepackage[most]{tcolorbox}
\usepackage{tikz}
\usetikzlibrary{positioning}
\usepackage{titlesec}

\providecommand{\modecorrige}{0}
\providecommand{\variante}{A}
\newif\ifcorrige
\ifnum\modecorrige=1 \corrigetrue \else \corrigefalse \fi

\definecolor{ipred}{HTML}{C0392B}
\definecolor{cdcbar}{HTML}{A67C3D}
\definecolor{cdcbg}{HTML}{F4F1EA}
\definecolor{qbar}{HTML}{4A6B5D}
\definecolor{qbg}{HTML}{EDF1EF}
\definecolor{corr}{HTML}{1F5FA8}
\definecolor{corrbg}{HTML}{E8F0FA}
\definecolor{grisfoot}{HTML}{828282}

\newcommand{\ip}[1]{\textcolor{ipred}{\texttt{#1}}}
\newcommand{\pts}[1]{\hfill\mbox{\textbf{(#1)}}}
\newcommand{\sol}[1]{\ifcorrige\textcolor{corr}{#1}\fi}
\newcommand{\bareme}[1]{\ifcorrige\par{\footnotesize\color{corr}#1\par}\fi}
\newcommand{\hl}{\ifcorrige\rule[-1.6mm]{0pt}{5.4mm}\else\rule[-2.2mm]{0pt}{6.6mm}\fi}
\newcolumntype{C}{>{\centering\arraybackslash}X}
\newcolumntype{L}{>{\raggedright\arraybackslash}X}

\titleformat{\section}{\large\bfseries}{}{0pt}{}
\titlespacing*{\section}{0pt}{6pt}{3pt}

\tcbset{qstyle/.style={enhanced,boxrule=0pt,leftrule=2.5pt,arc=0pt,outer arc=0pt,
  left=5pt,right=5pt,top=2pt,bottom=2pt,before skip=4pt,after skip=3pt}}
\newtcolorbox{cdc}{qstyle,colback=cdcbg,colframe=cdcbar}
\newtcolorbox{quest}{qstyle,colback=qbg,colframe=qbar}

% Zone de réponse : vierge dans le sujet, remplie dans le corrigé
\newcommand{\zone}[2]{%
  \ifcorrige
    \begin{tcolorbox}[colback=corrbg,colframe=corr,boxrule=0.5pt,arc=1pt,left=4pt,right=4pt,top=1pt,bottom=1pt,
      before skip=2pt,after skip=2pt,fontupper=\footnotesize\color{corr}]#2\end{tcolorbox}%
  \else
    \begin{tcolorbox}[colback=white,colframe=black!45,boxrule=0.4pt,arc=1pt,height=#1,
      before skip=2pt,after skip=2pt]\end{tcolorbox}%
  \fi}

\newcommand{\question}[3]{\begin{quest}\textbf{#1.} #3\pts{#2}\end{quest}}

%==========================================================================
% Données propres à chaque variante
%==========================================================================
\if\variante A
  \def\entreprise{Le cabinet d'architectes \textbf{ArchiNova}}
  \def\Rnom{R-ArchiNova}
  \def\resUn{Bureau d'études}\def\pUn{BE}\def\netUn{192.168.40}
  \def\resDeux{Accueil}\def\pDeux{ACC}\def\netDeux{172.16}
  \def\srvNom{SRV-Intranet}\def\netSrv{192.168.200}\def\srvH{10}
  \def\ssid{ArchiNova-WiFi}
  \def\AIrows{%
    \hl\ip{192.168.40.25} & \sol{C} & \sol{255.255.255.0} & \sol{192.168.40.0} & \sol{192.168.40.255} & \sol{254} \\ \hline
    \hl\ip{172.16.5.10}   & \sol{B} & \sol{255.255.0.0}   & \sol{172.16.0.0}   & \sol{172.16.255.255} & \sol{65\,534} \\ \hline
    \hl\ip{10.3.2.1}      & \sol{A} & \sol{255.0.0.0}     & \sol{10.0.0.0}     & \sol{10.255.255.255} & \sol{16\,777\,214} \\ \hline}
  \def\apipa{169.254.12.7}
  \def\Cnet{192.168.60}
  \def\Crows{%
    PC-A & \ip{192.168.60.10} & \ip{255.255.255.0} & \ip{192.168.60.254} & — \\
    PC-B & \ip{192.168.60.20} & \ip{255.255.255.0} & \ip{192.168.60.1}   & — \\
    PC-C & \ip{192.168.61.30} & \ip{255.255.255.0} & \ip{192.168.60.254} & — \\
    PC-D & \ip{192.168.60.10} & \ip{255.255.255.0} & \ip{192.168.60.254} & — \\
    SRV-C & \ip{10.0.0.5}     & \ip{255.0.0.0}     & \ip{10.0.0.254}     & — \\
    R1 — G0/0 & \ip{192.168.60.254} & \ip{255.255.255.0} & — & On \\
    R1 — G0/1 & \ip{10.0.0.254}     & \ip{255.0.0.0}     & — & Off \\}
  \def\Csol{%
    \hl 1 & \sol{PC-B} & \sol{Passerelle .60.1 : pas l'adresse du routeur.} & \sol{Joint son réseau, mais aucun autre (serveur).} & \sol{Passerelle 192.168.60.254.} \\ \hline
    \hl 2 & \sol{PC-C} & \sol{IP hors du réseau 192.168.60.0/24.} & \sol{Ne joint ni les postes ni le routeur.} & \sol{Ex. 192.168.60.30.} \\ \hline
    \hl 3 & \sol{PC-D} & \sol{Même IP que PC-A (doublon).} & \sol{Conflit d'adresse : PC-A et PC-D perturbés.} & \sol{Adresse unique, ex. .60.40.} \\ \hline
    \hl 4 & \sol{R1 G0/1} & \sol{Interface désactivée (Off).} & \sol{Réseau du serveur injoignable pour tous.} & \sol{Port Status : On.} \\ \hline}
  \def\Cbareme{PC-A et SRV-C sont correctement configurés ; accepter PC-A cité à la place de PC-D pour le doublon.}
\else
  \def\entreprise{La clinique vétérinaire \textbf{VetoPlus}}
  \def\Rnom{R-VetoPlus}
  \def\resUn{Soins}\def\pUn{SOI}\def\netUn{192.168.70}
  \def\resDeux{Secrétariat}\def\pDeux{SEC}\def\netDeux{172.20}
  \def\srvNom{SRV-Dossiers}\def\netSrv{192.168.250}\def\srvH{20}
  \def\ssid{VetoPlus-WiFi}
  \def\AIrows{%
    \hl\ip{192.168.15.130} & \sol{C} & \sol{255.255.255.0} & \sol{192.168.15.0} & \sol{192.168.15.255} & \sol{254} \\ \hline
    \hl\ip{172.31.200.7}   & \sol{B} & \sol{255.255.0.0}   & \sol{172.31.0.0}   & \sol{172.31.255.255} & \sol{65\,534} \\ \hline
    \hl\ip{10.45.6.9}      & \sol{A} & \sol{255.0.0.0}     & \sol{10.0.0.0}     & \sol{10.255.255.255} & \sol{16\,777\,214} \\ \hline}
  \def\apipa{169.254.88.3}
  \def\Cnet{192.168.80}
  \def\Crows{%
    PC-A & \ip{192.168.8.11}   & \ip{255.255.255.0} & \ip{192.168.80.254} & — \\
    PC-B & \ip{192.168.80.12}  & \ip{255.255.255.0} & \ip{192.168.80.254} & — \\
    PC-C & \ip{192.168.80.13}  & \ip{255.255.255.0} & \ip{192.168.80.100} & — \\
    PC-D & \ip{192.168.80.12}  & \ip{255.255.255.0} & \ip{192.168.80.254} & — \\
    SRV-C & \ip{10.0.0.8}      & \ip{255.0.0.0}     & \ip{10.0.0.254}     & — \\
    R1 — G0/0 & \ip{192.168.80.254} & \ip{255.255.255.0} & — & Off \\
    R1 — G0/1 & \ip{10.0.0.254}     & \ip{255.0.0.0}     & — & On \\}
  \def\Csol{%
    \hl 1 & \sol{PC-A} & \sol{IP hors du réseau 192.168.80.0/24.} & \sol{Ne joint ni les postes ni le routeur.} & \sol{Ex. 192.168.80.11.} \\ \hline
    \hl 2 & \sol{PC-C} & \sol{Passerelle .80.100 : pas l'adresse du routeur.} & \sol{Joint son réseau, mais aucun autre (serveur).} & \sol{Passerelle 192.168.80.254.} \\ \hline
    \hl 3 & \sol{PC-D} & \sol{Même IP que PC-B (doublon).} & \sol{Conflit d'adresse : PC-B et PC-D perturbés.} & \sol{Adresse unique, ex. .80.14.} \\ \hline
    \hl 4 & \sol{R1 G0/0} & \sol{Interface côté LAN désactivée (Off).} & \sol{Aucun poste ne joint le routeur ni le serveur.} & \sol{Port Status : On.} \\ \hline}
  \def\Cbareme{PC-B et SRV-C sont correctement configurés ; accepter PC-B cité à la place de PC-D pour le doublon.}
\fi

\pagestyle{fancy}
\fancyhf{}
\renewcommand{\headrulewidth}{0pt}
\fancyfoot[L]{\footnotesize\itshape\color{grisfoot}Évaluation — Réseaux locaux — Sujet \variante\ifcorrige\ — CORRIGÉ\fi}
\fancyfoot[R]{\footnotesize\color{grisfoot}\thepage/3}

\setlist{itemsep=0pt,topsep=2pt,parsep=0pt}
\setlength{\parindent}{0pt}
\setlength{\parskip}{2pt}

\begin{document}

\begin{center}
  {\Large\bfseries Évaluation pratique — Réseaux locaux \quad
   \fcolorbox{black}{white}{\ SUJET \variante\ }\par}
  \vspace{1pt}
  {\bfseries Adressage IP, commutation et routage — BTS CIEL 1re année — durée : 2 heures\par}
  \ifcorrige\vspace{2pt}{\bfseries\color{corr}CORRIGÉ ET BARÈME — document enseignant}\par\fi
\end{center}
\vspace{-2pt}
{\renewcommand{\arraystretch}{1.35}
\begin{tabularx}{\linewidth}{|X|X|p{2.2cm}|p{2.6cm}|}
  \hline
  \textbf{Nom :} & \textbf{Prénom :} & \textbf{Poste n°} & \textbf{Note :} \hfill /20 \\ \hline
\end{tabularx}}

\textbf{Consignes :} travail \textbf{individuel}, documents et Internet \textbf{non autorisés}. Ce sujet est le
\textbf{document-réponse}. Fichier Packet Tracer : \ip{Nom\_Prenom\_EvalLAN\_\variante.pkt}, à enregistrer
régulièrement dans le dossier de rendu. Les « visas » sont apposés par le professeur, que vous appelez à chaque test réussi.
Durées conseillées : A 30 min — B 1 h 10 — C 20 min.

%==========================================================================
\section{Partie A — Questions de cours et calculs \hfill /6}

\question{A1}{2 pts}{Complétez le tableau en considérant le \textbf{masque par défaut} de la classe de chaque adresse.}
{\small\renewcommand{\arraystretch}{1.15}
\begin{tabularx}{\linewidth}{|l|c|C|C|C|c|}
  \hline
  \textbf{Adresse IP} & \textbf{Classe} & \textbf{Masque par défaut} & \textbf{Adresse réseau} &
  \textbf{Adresse de diffusion} & \textbf{Nb d'hôtes} \\ \hline
  \AIrows
\end{tabularx}}
\bareme{0,5 pt par ligne juste + 0,5 pt si les trois nombres d'hôtes sont justes ($2^{n}-2$, $n$ = nombre de bits d'hôte).}

\question{A2}{1 pt}{Indiquez le type de câble (\textbf{droit} ou \textbf{croisé}) pour chaque liaison.}
{\small\renewcommand{\arraystretch}{1.15}
\begin{tabularx}{\linewidth}{|X|c||X|c|}
  \hline
  \hl PC $\leftrightarrow$ commutateur & \makebox[1.6cm]{\sol{droit}} & Commutateur $\leftrightarrow$ routeur & \makebox[1.6cm]{\sol{droit}} \\ \hline
  \hl PC $\leftrightarrow$ PC (réseau P2P) & \sol{croisé} & Concentrateur $\leftrightarrow$ concentrateur & \sol{croisé} \\ \hline
\end{tabularx}}
\bareme{0,25 pt par case. Équipements de types différents : câble droit ; de même type : câble croisé.}

\question{A3}{1 pt}{Comment un \textbf{concentrateur (hub)} et un \textbf{commutateur (switch)} traitent-ils une trame
reçue ? Lequel limite les collisions, et pourquoi ?}
\zone{2.9cm}{Le \textbf{hub} (couche 1) répète la trame sur \textbf{tous ses autres ports} : un seul domaine de collision
partagé. Le \textbf{switch} (couche 2) lit l'adresse MAC de destination et, grâce à sa \textbf{table MAC}, ne l'envoie
\textbf{que sur le port du destinataire} : un domaine de collision par port, c'est lui qui limite les collisions.
\emph{0,5 pt hub, 0,5 pt switch.}}

\question{A4}{1 pt}{Qu'est-ce que la \textbf{passerelle par défaut} d'un poste ? Quelle condition son adresse doit-elle respecter ?}
\zone{2.4cm}{Adresse IP de l'interface du \textbf{routeur} à laquelle le poste envoie les paquets destinés à un
\textbf{autre réseau}. Elle doit \textbf{appartenir au même réseau} que le poste. \emph{0,5 + 0,5 pt.}}

\question{A5}{0,5 pt}{Un portable configuré en DHCP affiche \ip{\apipa} / \ip{255.255.0.0}. Que s'est-il passé ?}
\zone{1.5cm}{Aucun \textbf{serveur DHCP} n'a répondu : le poste s'est attribué une adresse \textbf{APIPA} (169.254.0.0/16).}

\question{A6}{0,5 pt}{Quel protocole utilise \texttt{ping} ? Nommez les deux messages échangés et la couche OSI des adresses IP.}
\zone{1.5cm}{\textbf{ICMP} ; \textbf{Echo Request} / \textbf{Echo Reply} ; adresses IP en \textbf{couche 3 (Réseau)}.}

%==========================================================================
\newpage
\section{Partie B — Réalisation sous Packet Tracer \hfill /10}

\begin{cdc}
\entreprise{} vous confie son réseau : \textbf{trois réseaux} reliés par un \textbf{routeur unique Cisco 2911}
nommé \ip{\Rnom} (interfaces G0/0, G0/1, G0/2).
\begin{itemize}[leftmargin=12pt]
  \item \textbf{Réseau \resUn} : \textbf{4 postes} PC-\pUn1 à PC-\pUn4, commutateur \ip{SW-\pUn}, adressage statique,
        réseau \ip{\netUn.0/24}.
  \item \textbf{Réseau \resDeux} : \textbf{2 postes} PC-\pDeux1, PC-\pDeux2 et \textbf{1 imprimante} IMP-\pDeux,
        commutateur \ip{SW-\pDeux}, adressage statique, réseau de classe B \ip{\netDeux.0.0/16}.
  \item \textbf{Réseau Serveur} : serveur \ip{\srvNom} d'adresse \ip{\netSrv.\srvH/24}, relié \textbf{directement} au routeur.
  \item La passerelle de chaque réseau est sa \textbf{dernière adresse utilisable}. Tous les équipements communiquent
        entre eux et accèdent à la page web du serveur. Nommage imposé (« Display Name »).
\end{itemize}
\end{cdc}

\question{B1}{2 pts}{Complétez le plan d'adressage \textbf{avant} la réalisation.}
{\footnotesize\setlength{\tabcolsep}{3pt}\renewcommand{\arraystretch}{1.15}
\begin{tabularx}{\linewidth}{|l|C|C|C|C|C|C|}
  \hline
  \textbf{Réseau} & \textbf{Adresse réseau} & \textbf{Masque} & \textbf{1re @ utilisable} &
  \textbf{Dernière @ utilisable} & \textbf{Diffusion} & \textbf{Passerelle} \\ \hline
  \hl \resUn & \sol{\netUn.0} & \sol{255.255.255.0} & \sol{\netUn.1} & \sol{\netUn.254} & \sol{\netUn.255} & \sol{\netUn.254} \\ \hline
  \hl \resDeux & \sol{\netDeux.0.0} & \sol{255.255.0.0} & \sol{\netDeux.0.1} & \sol{\netDeux.255.254} & \sol{\netDeux.255.255} & \sol{\netDeux.255.254} \\ \hline
  \hl Serveur & \sol{\netSrv.0} & \sol{255.255.255.0} & \sol{\netSrv.1} & \sol{\netSrv.254} & \sol{\netSrv.255} & \sol{\netSrv.254} \\ \hline
\end{tabularx}}
\bareme{0,5 pt par ligne juste + 0,5 pt si la convention de passerelle est respectée partout.}

\question{B2}{2 pts}{Construisez la topologie (équipements, câblage, nommage) et relevez l'adressage configuré.}
{\footnotesize\setlength{\tabcolsep}{3pt}\renewcommand{\arraystretch}{1.1}
\begin{tabularx}{\linewidth}{|l|C|C|C||l|C|C|}
  \hline
  \textbf{Hôte} & \textbf{Adresse IP} & \textbf{Masque} & \textbf{Passerelle} & \textbf{Interface} & \textbf{Adresse IP} & \textbf{Masque} \\ \hline
  \hl PC-\pUn1 & \sol{\netUn.1} & \sol{255.255.255.0} & \sol{\netUn.254} & \Rnom{} G0/0 & \sol{\netUn.254} & \sol{255.255.255.0} \\ \hline
  \hl PC-\pDeux1 & \sol{\netDeux.0.1} & \sol{255.255.0.0} & \sol{\netDeux.255.254} & \Rnom{} G0/1 & \sol{\netDeux.255.254} & \sol{255.255.0.0} \\ \hline
  \hl IMP-\pDeux & \sol{\netDeux.0.10} & \sol{255.255.0.0} & \sol{\netDeux.255.254} & \Rnom{} G0/2 & \sol{\netSrv.254} & \sol{255.255.255.0} \\ \hline
  \hl \srvNom & \ip{\netSrv.\srvH} & \sol{255.255.255.0} & \sol{\netSrv.254} & \multicolumn{3}{c|}{\cellcolor{black!8}} \\ \hline
\end{tabularx}}
\bareme{Adresses d'hôtes libres (uniques, dans le bon réseau) ; affectation G0/x libre. Sur le .pkt : équipements 0,5 ;
câblage 1 (câbles droits vers les switchs et switch–routeur, \textbf{croisé} serveur–routeur, liens verts) ; nommage 0,5.}

\question{B3}{3 pts}{Configurez tous les équipements : IP, masques, passerelles des hôtes ; interfaces du routeur
(« Port Status : On »).}
\bareme{Sur le .pkt : interfaces routeur 1 pt ; IP/masques des hôtes 1 pt ; passerelles (imprimante et serveur compris)
1 pt. $-0,25$ par erreur, sans descendre sous 0 par item.}

\question{B4}{3 pts}{Réalisez les tests, notez la commande et le résultat, puis faites viser chaque test réussi.}
{\small\renewcommand{\arraystretch}{1.1}
\begin{tabularx}{\linewidth}{|c|l|l|C|c|c|}
  \hline
  \textbf{N°} & \textbf{Source} & \textbf{Destination} & \textbf{Commande / action} & \textbf{Résultat} & \textbf{Visa prof.} \\ \hline
  \hl 1 & PC-\pUn1 & PC-\pUn4 & \sol{ping \netUn.4} & \sol{réussi} & \hspace{1.4cm} \\ \hline
  \hl 2 & PC-\pDeux1 & IMP-\pDeux & \sol{ping \netDeux.0.10} & \sol{réussi} & \\ \hline
  \hl 3 & PC-\pUn2 & PC-\pDeux2 & \sol{ping \netDeux.0.2} & \sol{réussi} & \\ \hline
  \hl 4 & PC-\pDeux2 & \srvNom & \sol{ping \netSrv.\srvH} & \sol{réussi} & \\ \hline
  \hl 5 & PC-\pUn3 & page web du serveur & \sol{navigateur : {\NoAutoSpacing http://\netSrv.\srvH}} & \sol{page affichée} & \\ \hline
\end{tabularx}}
\bareme{0,5 pt par test visé + 0,5 pt pour B5. Le premier ping inter-réseaux peut échouer (résolution ARP) ; les suivants doivent réussir.}

\question{B5}{incl. B4}{Test n°3 : par quels équipements passe le paquet ? Pourquoi PC-\pUn2 l'envoie-t-il au routeur ?}
\zone{2.3cm}{PC-\pUn2 $\to$ SW-\pUn{} $\to$ \Rnom{} $\to$ SW-\pDeux{} $\to$ PC-\pDeux2. Grâce à son masque, PC-\pUn2 constate que
la destination (\netDeux.0.0/16) n'est \textbf{pas dans son réseau} (\netUn.0/24) : il envoie le paquet à sa
\textbf{passerelle par défaut}, qui le route.}

%==========================================================================
\newpage
\section{Partie C — Diagnostic d'un réseau en panne \hfill /4}

Un technicien constate qu'\textbf{aucun poste ne parvient à joindre le serveur SRV-C} et que certains postes
communiquent mal entre eux. Les câbles sont du bon type.

\begin{center}
\begin{minipage}{0.37\linewidth}\centering
\begin{tikzpicture}[font=\footnotesize,
  eq/.style={draw,fill=qbg,rounded corners=2pt,minimum width=9mm,minimum height=5mm,inner sep=2pt},
  every path/.style={thick}]
  \node[eq] (SW) {SW-1};
  \node[eq,right=9mm of SW] (R) {R1};
  \node[eq,right=7mm of R] (SRV) {SRV-C};
  \node[eq,above left=5mm and -3mm of SW] (A) {PC-A};
  \node[eq,above right=5mm and -3mm of SW] (B) {PC-B};
  \node[eq,below left=5mm and -3mm of SW] (C) {PC-C};
  \node[eq,below right=5mm and -3mm of SW] (D) {PC-D};
  \foreach \n in {A,B,C,D} \draw (\n) -- (SW);
  \draw (SW) -- node[above,font=\tiny]{G0/0} (R);
  \draw (R) -- node[above,font=\tiny]{G0/1} (SRV);
\end{tikzpicture}
\end{minipage}\hfill
\begin{minipage}{0.61\linewidth}\centering
{\small\renewcommand{\arraystretch}{1.05}\setlength{\tabcolsep}{4pt}
\begin{tabular}{lllll}
  \toprule
  \textbf{Équipement} & \textbf{Adresse IP} & \textbf{Masque} & \textbf{Passerelle} & \textbf{Port} \\
  \midrule
  \Crows
  \bottomrule
\end{tabular}}
\end{minipage}
\end{center}

\question{C1}{4 pts}{Relevez les \textbf{quatre anomalies}. Pour chacune : conséquence et correction.}
{\small\renewcommand{\arraystretch}{1.1}
\ifcorrige\def\hc{}\else\def\hc{\rule[-10mm]{0pt}{14mm}}\fi
\begin{tabularx}{\linewidth}{|c|>{\raggedright\arraybackslash}p{2cm}|L|L|L|}
  \hline
  \textbf{N°} & \textbf{Équipement} & \textbf{Anomalie} & \textbf{Conséquence} & \textbf{Correction} \\ \hline
  \ifcorrige\Csol\else
  \hc 1 & & & & \\ \hline
  \hc 2 & & & & \\ \hline
  \hc 3 & & & & \\ \hline
  \hc 4 & & & & \\ \hline
  \fi
\end{tabularx}}
\bareme{1 pt par anomalie (0,25 équipement, anomalie, conséquence, correction ; ordre indifférent). \Cbareme}

%==========================================================================
\section{Bonus — Extension sans fil \hfill +1}

\begin{cdc}Ajoutez au réseau \resUn{} un \textbf{AccessPoint-PT} relié à SW-\pUn{} (SSID \ip{\ssid}, \textbf{WPA2-PSK})
et un portable \ip{Laptop-\pUn} équipé d'un module \texttt{WPC300N}, en adressage \textbf{statique}.
Vérifiez qu'il joint \srvNom.\end{cdc}
{\small\renewcommand{\arraystretch}{1.1}
\begin{tabularx}{\linewidth}{|l|C|C|C|c|}
  \hline
  \textbf{Équipement} & \textbf{Adresse IP} & \textbf{Masque} & \textbf{Passerelle} & \textbf{Visa prof.} \\ \hline
  \hl Laptop-\pUn & \sol{\netUn.50} & \sol{255.255.255.0} & \sol{\netUn.254} & \hspace{1.4cm} \\ \hline
\end{tabularx}}
\bareme{Le point d'accès fait office de pont : le portable est dans \netUn.0/24. Bonus si associé en WPA2-PSK et ping
vers \netSrv.\srvH{} réussi. Note plafonnée à 20.}

\end{document}
"
%   pdflatex -jobname=Eval_Reseau_CorrigeA "\def\variante{A}\def\modecorrige{1}% Évaluation pratique — Réseaux locaux (BTS CIEL 1re année) — sujets A et B
% Compilation (4 PDF) :
%   pdflatex -jobname=Eval_Reseau_SujetA   "\def\variante{A}\def\modecorrige{0}\input{Eval_TP_Reseau_LAN.tex}"
%   pdflatex -jobname=Eval_Reseau_SujetB   "\def\variante{B}\def\modecorrige{0}\input{Eval_TP_Reseau_LAN.tex}"
%   pdflatex -jobname=Eval_Reseau_CorrigeA "\def\variante{A}\def\modecorrige{1}\input{Eval_TP_Reseau_LAN.tex}"
%   pdflatex -jobname=Eval_Reseau_CorrigeB "\def\variante{B}\def\modecorrige{1}\input{Eval_TP_Reseau_LAN.tex}"
\documentclass[10pt,a4paper]{article}

\usepackage[T1]{fontenc}
\usepackage[utf8]{inputenc}
\usepackage{lmodern}
\usepackage[french]{babel}
\usepackage[a4paper,hmargin=1.5cm,top=1.3cm,bottom=1.6cm,footskip=0.7cm]{geometry}
\usepackage[table]{xcolor}
\usepackage{booktabs,tabularx,array}
\usepackage{enumitem}
\usepackage{fancyhdr}
\usepackage[most]{tcolorbox}
\usepackage{tikz}
\usetikzlibrary{positioning}
\usepackage{titlesec}

\providecommand{\modecorrige}{0}
\providecommand{\variante}{A}
\newif\ifcorrige
\ifnum\modecorrige=1 \corrigetrue \else \corrigefalse \fi

\definecolor{ipred}{HTML}{C0392B}
\definecolor{cdcbar}{HTML}{A67C3D}
\definecolor{cdcbg}{HTML}{F4F1EA}
\definecolor{qbar}{HTML}{4A6B5D}
\definecolor{qbg}{HTML}{EDF1EF}
\definecolor{corr}{HTML}{1F5FA8}
\definecolor{corrbg}{HTML}{E8F0FA}
\definecolor{grisfoot}{HTML}{828282}

\newcommand{\ip}[1]{\textcolor{ipred}{\texttt{#1}}}
\newcommand{\pts}[1]{\hfill\mbox{\textbf{(#1)}}}
\newcommand{\sol}[1]{\ifcorrige\textcolor{corr}{#1}\fi}
\newcommand{\bareme}[1]{\ifcorrige\par{\footnotesize\color{corr}#1\par}\fi}
\newcommand{\hl}{\ifcorrige\rule[-1.6mm]{0pt}{5.4mm}\else\rule[-2.2mm]{0pt}{6.6mm}\fi}
\newcolumntype{C}{>{\centering\arraybackslash}X}
\newcolumntype{L}{>{\raggedright\arraybackslash}X}

\titleformat{\section}{\large\bfseries}{}{0pt}{}
\titlespacing*{\section}{0pt}{6pt}{3pt}

\tcbset{qstyle/.style={enhanced,boxrule=0pt,leftrule=2.5pt,arc=0pt,outer arc=0pt,
  left=5pt,right=5pt,top=2pt,bottom=2pt,before skip=4pt,after skip=3pt}}
\newtcolorbox{cdc}{qstyle,colback=cdcbg,colframe=cdcbar}
\newtcolorbox{quest}{qstyle,colback=qbg,colframe=qbar}

% Zone de réponse : vierge dans le sujet, remplie dans le corrigé
\newcommand{\zone}[2]{%
  \ifcorrige
    \begin{tcolorbox}[colback=corrbg,colframe=corr,boxrule=0.5pt,arc=1pt,left=4pt,right=4pt,top=1pt,bottom=1pt,
      before skip=2pt,after skip=2pt,fontupper=\footnotesize\color{corr}]#2\end{tcolorbox}%
  \else
    \begin{tcolorbox}[colback=white,colframe=black!45,boxrule=0.4pt,arc=1pt,height=#1,
      before skip=2pt,after skip=2pt]\end{tcolorbox}%
  \fi}

\newcommand{\question}[3]{\begin{quest}\textbf{#1.} #3\pts{#2}\end{quest}}

%==========================================================================
% Données propres à chaque variante
%==========================================================================
\if\variante A
  \def\entreprise{Le cabinet d'architectes \textbf{ArchiNova}}
  \def\Rnom{R-ArchiNova}
  \def\resUn{Bureau d'études}\def\pUn{BE}\def\netUn{192.168.40}
  \def\resDeux{Accueil}\def\pDeux{ACC}\def\netDeux{172.16}
  \def\srvNom{SRV-Intranet}\def\netSrv{192.168.200}\def\srvH{10}
  \def\ssid{ArchiNova-WiFi}
  \def\AIrows{%
    \hl\ip{192.168.40.25} & \sol{C} & \sol{255.255.255.0} & \sol{192.168.40.0} & \sol{192.168.40.255} & \sol{254} \\ \hline
    \hl\ip{172.16.5.10}   & \sol{B} & \sol{255.255.0.0}   & \sol{172.16.0.0}   & \sol{172.16.255.255} & \sol{65\,534} \\ \hline
    \hl\ip{10.3.2.1}      & \sol{A} & \sol{255.0.0.0}     & \sol{10.0.0.0}     & \sol{10.255.255.255} & \sol{16\,777\,214} \\ \hline}
  \def\apipa{169.254.12.7}
  \def\Cnet{192.168.60}
  \def\Crows{%
    PC-A & \ip{192.168.60.10} & \ip{255.255.255.0} & \ip{192.168.60.254} & — \\
    PC-B & \ip{192.168.60.20} & \ip{255.255.255.0} & \ip{192.168.60.1}   & — \\
    PC-C & \ip{192.168.61.30} & \ip{255.255.255.0} & \ip{192.168.60.254} & — \\
    PC-D & \ip{192.168.60.10} & \ip{255.255.255.0} & \ip{192.168.60.254} & — \\
    SRV-C & \ip{10.0.0.5}     & \ip{255.0.0.0}     & \ip{10.0.0.254}     & — \\
    R1 — G0/0 & \ip{192.168.60.254} & \ip{255.255.255.0} & — & On \\
    R1 — G0/1 & \ip{10.0.0.254}     & \ip{255.0.0.0}     & — & Off \\}
  \def\Csol{%
    \hl 1 & \sol{PC-B} & \sol{Passerelle .60.1 : pas l'adresse du routeur.} & \sol{Joint son réseau, mais aucun autre (serveur).} & \sol{Passerelle 192.168.60.254.} \\ \hline
    \hl 2 & \sol{PC-C} & \sol{IP hors du réseau 192.168.60.0/24.} & \sol{Ne joint ni les postes ni le routeur.} & \sol{Ex. 192.168.60.30.} \\ \hline
    \hl 3 & \sol{PC-D} & \sol{Même IP que PC-A (doublon).} & \sol{Conflit d'adresse : PC-A et PC-D perturbés.} & \sol{Adresse unique, ex. .60.40.} \\ \hline
    \hl 4 & \sol{R1 G0/1} & \sol{Interface désactivée (Off).} & \sol{Réseau du serveur injoignable pour tous.} & \sol{Port Status : On.} \\ \hline}
  \def\Cbareme{PC-A et SRV-C sont correctement configurés ; accepter PC-A cité à la place de PC-D pour le doublon.}
\else
  \def\entreprise{La clinique vétérinaire \textbf{VetoPlus}}
  \def\Rnom{R-VetoPlus}
  \def\resUn{Soins}\def\pUn{SOI}\def\netUn{192.168.70}
  \def\resDeux{Secrétariat}\def\pDeux{SEC}\def\netDeux{172.20}
  \def\srvNom{SRV-Dossiers}\def\netSrv{192.168.250}\def\srvH{20}
  \def\ssid{VetoPlus-WiFi}
  \def\AIrows{%
    \hl\ip{192.168.15.130} & \sol{C} & \sol{255.255.255.0} & \sol{192.168.15.0} & \sol{192.168.15.255} & \sol{254} \\ \hline
    \hl\ip{172.31.200.7}   & \sol{B} & \sol{255.255.0.0}   & \sol{172.31.0.0}   & \sol{172.31.255.255} & \sol{65\,534} \\ \hline
    \hl\ip{10.45.6.9}      & \sol{A} & \sol{255.0.0.0}     & \sol{10.0.0.0}     & \sol{10.255.255.255} & \sol{16\,777\,214} \\ \hline}
  \def\apipa{169.254.88.3}
  \def\Cnet{192.168.80}
  \def\Crows{%
    PC-A & \ip{192.168.8.11}   & \ip{255.255.255.0} & \ip{192.168.80.254} & — \\
    PC-B & \ip{192.168.80.12}  & \ip{255.255.255.0} & \ip{192.168.80.254} & — \\
    PC-C & \ip{192.168.80.13}  & \ip{255.255.255.0} & \ip{192.168.80.100} & — \\
    PC-D & \ip{192.168.80.12}  & \ip{255.255.255.0} & \ip{192.168.80.254} & — \\
    SRV-C & \ip{10.0.0.8}      & \ip{255.0.0.0}     & \ip{10.0.0.254}     & — \\
    R1 — G0/0 & \ip{192.168.80.254} & \ip{255.255.255.0} & — & Off \\
    R1 — G0/1 & \ip{10.0.0.254}     & \ip{255.0.0.0}     & — & On \\}
  \def\Csol{%
    \hl 1 & \sol{PC-A} & \sol{IP hors du réseau 192.168.80.0/24.} & \sol{Ne joint ni les postes ni le routeur.} & \sol{Ex. 192.168.80.11.} \\ \hline
    \hl 2 & \sol{PC-C} & \sol{Passerelle .80.100 : pas l'adresse du routeur.} & \sol{Joint son réseau, mais aucun autre (serveur).} & \sol{Passerelle 192.168.80.254.} \\ \hline
    \hl 3 & \sol{PC-D} & \sol{Même IP que PC-B (doublon).} & \sol{Conflit d'adresse : PC-B et PC-D perturbés.} & \sol{Adresse unique, ex. .80.14.} \\ \hline
    \hl 4 & \sol{R1 G0/0} & \sol{Interface côté LAN désactivée (Off).} & \sol{Aucun poste ne joint le routeur ni le serveur.} & \sol{Port Status : On.} \\ \hline}
  \def\Cbareme{PC-B et SRV-C sont correctement configurés ; accepter PC-B cité à la place de PC-D pour le doublon.}
\fi

\pagestyle{fancy}
\fancyhf{}
\renewcommand{\headrulewidth}{0pt}
\fancyfoot[L]{\footnotesize\itshape\color{grisfoot}Évaluation — Réseaux locaux — Sujet \variante\ifcorrige\ — CORRIGÉ\fi}
\fancyfoot[R]{\footnotesize\color{grisfoot}\thepage/3}

\setlist{itemsep=0pt,topsep=2pt,parsep=0pt}
\setlength{\parindent}{0pt}
\setlength{\parskip}{2pt}

\begin{document}

\begin{center}
  {\Large\bfseries Évaluation pratique — Réseaux locaux \quad
   \fcolorbox{black}{white}{\ SUJET \variante\ }\par}
  \vspace{1pt}
  {\bfseries Adressage IP, commutation et routage — BTS CIEL 1re année — durée : 2 heures\par}
  \ifcorrige\vspace{2pt}{\bfseries\color{corr}CORRIGÉ ET BARÈME — document enseignant}\par\fi
\end{center}
\vspace{-2pt}
{\renewcommand{\arraystretch}{1.35}
\begin{tabularx}{\linewidth}{|X|X|p{2.2cm}|p{2.6cm}|}
  \hline
  \textbf{Nom :} & \textbf{Prénom :} & \textbf{Poste n°} & \textbf{Note :} \hfill /20 \\ \hline
\end{tabularx}}

\textbf{Consignes :} travail \textbf{individuel}, documents et Internet \textbf{non autorisés}. Ce sujet est le
\textbf{document-réponse}. Fichier Packet Tracer : \ip{Nom\_Prenom\_EvalLAN\_\variante.pkt}, à enregistrer
régulièrement dans le dossier de rendu. Les « visas » sont apposés par le professeur, que vous appelez à chaque test réussi.
Durées conseillées : A 30 min — B 1 h 10 — C 20 min.

%==========================================================================
\section{Partie A — Questions de cours et calculs \hfill /6}

\question{A1}{2 pts}{Complétez le tableau en considérant le \textbf{masque par défaut} de la classe de chaque adresse.}
{\small\renewcommand{\arraystretch}{1.15}
\begin{tabularx}{\linewidth}{|l|c|C|C|C|c|}
  \hline
  \textbf{Adresse IP} & \textbf{Classe} & \textbf{Masque par défaut} & \textbf{Adresse réseau} &
  \textbf{Adresse de diffusion} & \textbf{Nb d'hôtes} \\ \hline
  \AIrows
\end{tabularx}}
\bareme{0,5 pt par ligne juste + 0,5 pt si les trois nombres d'hôtes sont justes ($2^{n}-2$, $n$ = nombre de bits d'hôte).}

\question{A2}{1 pt}{Indiquez le type de câble (\textbf{droit} ou \textbf{croisé}) pour chaque liaison.}
{\small\renewcommand{\arraystretch}{1.15}
\begin{tabularx}{\linewidth}{|X|c||X|c|}
  \hline
  \hl PC $\leftrightarrow$ commutateur & \makebox[1.6cm]{\sol{droit}} & Commutateur $\leftrightarrow$ routeur & \makebox[1.6cm]{\sol{droit}} \\ \hline
  \hl PC $\leftrightarrow$ PC (réseau P2P) & \sol{croisé} & Concentrateur $\leftrightarrow$ concentrateur & \sol{croisé} \\ \hline
\end{tabularx}}
\bareme{0,25 pt par case. Équipements de types différents : câble droit ; de même type : câble croisé.}

\question{A3}{1 pt}{Comment un \textbf{concentrateur (hub)} et un \textbf{commutateur (switch)} traitent-ils une trame
reçue ? Lequel limite les collisions, et pourquoi ?}
\zone{2.9cm}{Le \textbf{hub} (couche 1) répète la trame sur \textbf{tous ses autres ports} : un seul domaine de collision
partagé. Le \textbf{switch} (couche 2) lit l'adresse MAC de destination et, grâce à sa \textbf{table MAC}, ne l'envoie
\textbf{que sur le port du destinataire} : un domaine de collision par port, c'est lui qui limite les collisions.
\emph{0,5 pt hub, 0,5 pt switch.}}

\question{A4}{1 pt}{Qu'est-ce que la \textbf{passerelle par défaut} d'un poste ? Quelle condition son adresse doit-elle respecter ?}
\zone{2.4cm}{Adresse IP de l'interface du \textbf{routeur} à laquelle le poste envoie les paquets destinés à un
\textbf{autre réseau}. Elle doit \textbf{appartenir au même réseau} que le poste. \emph{0,5 + 0,5 pt.}}

\question{A5}{0,5 pt}{Un portable configuré en DHCP affiche \ip{\apipa} / \ip{255.255.0.0}. Que s'est-il passé ?}
\zone{1.5cm}{Aucun \textbf{serveur DHCP} n'a répondu : le poste s'est attribué une adresse \textbf{APIPA} (169.254.0.0/16).}

\question{A6}{0,5 pt}{Quel protocole utilise \texttt{ping} ? Nommez les deux messages échangés et la couche OSI des adresses IP.}
\zone{1.5cm}{\textbf{ICMP} ; \textbf{Echo Request} / \textbf{Echo Reply} ; adresses IP en \textbf{couche 3 (Réseau)}.}

%==========================================================================
\newpage
\section{Partie B — Réalisation sous Packet Tracer \hfill /10}

\begin{cdc}
\entreprise{} vous confie son réseau : \textbf{trois réseaux} reliés par un \textbf{routeur unique Cisco 2911}
nommé \ip{\Rnom} (interfaces G0/0, G0/1, G0/2).
\begin{itemize}[leftmargin=12pt]
  \item \textbf{Réseau \resUn} : \textbf{4 postes} PC-\pUn1 à PC-\pUn4, commutateur \ip{SW-\pUn}, adressage statique,
        réseau \ip{\netUn.0/24}.
  \item \textbf{Réseau \resDeux} : \textbf{2 postes} PC-\pDeux1, PC-\pDeux2 et \textbf{1 imprimante} IMP-\pDeux,
        commutateur \ip{SW-\pDeux}, adressage statique, réseau de classe B \ip{\netDeux.0.0/16}.
  \item \textbf{Réseau Serveur} : serveur \ip{\srvNom} d'adresse \ip{\netSrv.\srvH/24}, relié \textbf{directement} au routeur.
  \item La passerelle de chaque réseau est sa \textbf{dernière adresse utilisable}. Tous les équipements communiquent
        entre eux et accèdent à la page web du serveur. Nommage imposé (« Display Name »).
\end{itemize}
\end{cdc}

\question{B1}{2 pts}{Complétez le plan d'adressage \textbf{avant} la réalisation.}
{\footnotesize\setlength{\tabcolsep}{3pt}\renewcommand{\arraystretch}{1.15}
\begin{tabularx}{\linewidth}{|l|C|C|C|C|C|C|}
  \hline
  \textbf{Réseau} & \textbf{Adresse réseau} & \textbf{Masque} & \textbf{1re @ utilisable} &
  \textbf{Dernière @ utilisable} & \textbf{Diffusion} & \textbf{Passerelle} \\ \hline
  \hl \resUn & \sol{\netUn.0} & \sol{255.255.255.0} & \sol{\netUn.1} & \sol{\netUn.254} & \sol{\netUn.255} & \sol{\netUn.254} \\ \hline
  \hl \resDeux & \sol{\netDeux.0.0} & \sol{255.255.0.0} & \sol{\netDeux.0.1} & \sol{\netDeux.255.254} & \sol{\netDeux.255.255} & \sol{\netDeux.255.254} \\ \hline
  \hl Serveur & \sol{\netSrv.0} & \sol{255.255.255.0} & \sol{\netSrv.1} & \sol{\netSrv.254} & \sol{\netSrv.255} & \sol{\netSrv.254} \\ \hline
\end{tabularx}}
\bareme{0,5 pt par ligne juste + 0,5 pt si la convention de passerelle est respectée partout.}

\question{B2}{2 pts}{Construisez la topologie (équipements, câblage, nommage) et relevez l'adressage configuré.}
{\footnotesize\setlength{\tabcolsep}{3pt}\renewcommand{\arraystretch}{1.1}
\begin{tabularx}{\linewidth}{|l|C|C|C||l|C|C|}
  \hline
  \textbf{Hôte} & \textbf{Adresse IP} & \textbf{Masque} & \textbf{Passerelle} & \textbf{Interface} & \textbf{Adresse IP} & \textbf{Masque} \\ \hline
  \hl PC-\pUn1 & \sol{\netUn.1} & \sol{255.255.255.0} & \sol{\netUn.254} & \Rnom{} G0/0 & \sol{\netUn.254} & \sol{255.255.255.0} \\ \hline
  \hl PC-\pDeux1 & \sol{\netDeux.0.1} & \sol{255.255.0.0} & \sol{\netDeux.255.254} & \Rnom{} G0/1 & \sol{\netDeux.255.254} & \sol{255.255.0.0} \\ \hline
  \hl IMP-\pDeux & \sol{\netDeux.0.10} & \sol{255.255.0.0} & \sol{\netDeux.255.254} & \Rnom{} G0/2 & \sol{\netSrv.254} & \sol{255.255.255.0} \\ \hline
  \hl \srvNom & \ip{\netSrv.\srvH} & \sol{255.255.255.0} & \sol{\netSrv.254} & \multicolumn{3}{c|}{\cellcolor{black!8}} \\ \hline
\end{tabularx}}
\bareme{Adresses d'hôtes libres (uniques, dans le bon réseau) ; affectation G0/x libre. Sur le .pkt : équipements 0,5 ;
câblage 1 (câbles droits vers les switchs et switch–routeur, \textbf{croisé} serveur–routeur, liens verts) ; nommage 0,5.}

\question{B3}{3 pts}{Configurez tous les équipements : IP, masques, passerelles des hôtes ; interfaces du routeur
(« Port Status : On »).}
\bareme{Sur le .pkt : interfaces routeur 1 pt ; IP/masques des hôtes 1 pt ; passerelles (imprimante et serveur compris)
1 pt. $-0,25$ par erreur, sans descendre sous 0 par item.}

\question{B4}{3 pts}{Réalisez les tests, notez la commande et le résultat, puis faites viser chaque test réussi.}
{\small\renewcommand{\arraystretch}{1.1}
\begin{tabularx}{\linewidth}{|c|l|l|C|c|c|}
  \hline
  \textbf{N°} & \textbf{Source} & \textbf{Destination} & \textbf{Commande / action} & \textbf{Résultat} & \textbf{Visa prof.} \\ \hline
  \hl 1 & PC-\pUn1 & PC-\pUn4 & \sol{ping \netUn.4} & \sol{réussi} & \hspace{1.4cm} \\ \hline
  \hl 2 & PC-\pDeux1 & IMP-\pDeux & \sol{ping \netDeux.0.10} & \sol{réussi} & \\ \hline
  \hl 3 & PC-\pUn2 & PC-\pDeux2 & \sol{ping \netDeux.0.2} & \sol{réussi} & \\ \hline
  \hl 4 & PC-\pDeux2 & \srvNom & \sol{ping \netSrv.\srvH} & \sol{réussi} & \\ \hline
  \hl 5 & PC-\pUn3 & page web du serveur & \sol{navigateur : {\NoAutoSpacing http://\netSrv.\srvH}} & \sol{page affichée} & \\ \hline
\end{tabularx}}
\bareme{0,5 pt par test visé + 0,5 pt pour B5. Le premier ping inter-réseaux peut échouer (résolution ARP) ; les suivants doivent réussir.}

\question{B5}{incl. B4}{Test n°3 : par quels équipements passe le paquet ? Pourquoi PC-\pUn2 l'envoie-t-il au routeur ?}
\zone{2.3cm}{PC-\pUn2 $\to$ SW-\pUn{} $\to$ \Rnom{} $\to$ SW-\pDeux{} $\to$ PC-\pDeux2. Grâce à son masque, PC-\pUn2 constate que
la destination (\netDeux.0.0/16) n'est \textbf{pas dans son réseau} (\netUn.0/24) : il envoie le paquet à sa
\textbf{passerelle par défaut}, qui le route.}

%==========================================================================
\newpage
\section{Partie C — Diagnostic d'un réseau en panne \hfill /4}

Un technicien constate qu'\textbf{aucun poste ne parvient à joindre le serveur SRV-C} et que certains postes
communiquent mal entre eux. Les câbles sont du bon type.

\begin{center}
\begin{minipage}{0.37\linewidth}\centering
\begin{tikzpicture}[font=\footnotesize,
  eq/.style={draw,fill=qbg,rounded corners=2pt,minimum width=9mm,minimum height=5mm,inner sep=2pt},
  every path/.style={thick}]
  \node[eq] (SW) {SW-1};
  \node[eq,right=9mm of SW] (R) {R1};
  \node[eq,right=7mm of R] (SRV) {SRV-C};
  \node[eq,above left=5mm and -3mm of SW] (A) {PC-A};
  \node[eq,above right=5mm and -3mm of SW] (B) {PC-B};
  \node[eq,below left=5mm and -3mm of SW] (C) {PC-C};
  \node[eq,below right=5mm and -3mm of SW] (D) {PC-D};
  \foreach \n in {A,B,C,D} \draw (\n) -- (SW);
  \draw (SW) -- node[above,font=\tiny]{G0/0} (R);
  \draw (R) -- node[above,font=\tiny]{G0/1} (SRV);
\end{tikzpicture}
\end{minipage}\hfill
\begin{minipage}{0.61\linewidth}\centering
{\small\renewcommand{\arraystretch}{1.05}\setlength{\tabcolsep}{4pt}
\begin{tabular}{lllll}
  \toprule
  \textbf{Équipement} & \textbf{Adresse IP} & \textbf{Masque} & \textbf{Passerelle} & \textbf{Port} \\
  \midrule
  \Crows
  \bottomrule
\end{tabular}}
\end{minipage}
\end{center}

\question{C1}{4 pts}{Relevez les \textbf{quatre anomalies}. Pour chacune : conséquence et correction.}
{\small\renewcommand{\arraystretch}{1.1}
\ifcorrige\def\hc{}\else\def\hc{\rule[-10mm]{0pt}{14mm}}\fi
\begin{tabularx}{\linewidth}{|c|>{\raggedright\arraybackslash}p{2cm}|L|L|L|}
  \hline
  \textbf{N°} & \textbf{Équipement} & \textbf{Anomalie} & \textbf{Conséquence} & \textbf{Correction} \\ \hline
  \ifcorrige\Csol\else
  \hc 1 & & & & \\ \hline
  \hc 2 & & & & \\ \hline
  \hc 3 & & & & \\ \hline
  \hc 4 & & & & \\ \hline
  \fi
\end{tabularx}}
\bareme{1 pt par anomalie (0,25 équipement, anomalie, conséquence, correction ; ordre indifférent). \Cbareme}

%==========================================================================
\section{Bonus — Extension sans fil \hfill +1}

\begin{cdc}Ajoutez au réseau \resUn{} un \textbf{AccessPoint-PT} relié à SW-\pUn{} (SSID \ip{\ssid}, \textbf{WPA2-PSK})
et un portable \ip{Laptop-\pUn} équipé d'un module \texttt{WPC300N}, en adressage \textbf{statique}.
Vérifiez qu'il joint \srvNom.\end{cdc}
{\small\renewcommand{\arraystretch}{1.1}
\begin{tabularx}{\linewidth}{|l|C|C|C|c|}
  \hline
  \textbf{Équipement} & \textbf{Adresse IP} & \textbf{Masque} & \textbf{Passerelle} & \textbf{Visa prof.} \\ \hline
  \hl Laptop-\pUn & \sol{\netUn.50} & \sol{255.255.255.0} & \sol{\netUn.254} & \hspace{1.4cm} \\ \hline
\end{tabularx}}
\bareme{Le point d'accès fait office de pont : le portable est dans \netUn.0/24. Bonus si associé en WPA2-PSK et ping
vers \netSrv.\srvH{} réussi. Note plafonnée à 20.}

\end{document}
"
%   pdflatex -jobname=Eval_Reseau_CorrigeB "\def\variante{B}\def\modecorrige{1}% Évaluation pratique — Réseaux locaux (BTS CIEL 1re année) — sujets A et B
% Compilation (4 PDF) :
%   pdflatex -jobname=Eval_Reseau_SujetA   "\def\variante{A}\def\modecorrige{0}\input{Eval_TP_Reseau_LAN.tex}"
%   pdflatex -jobname=Eval_Reseau_SujetB   "\def\variante{B}\def\modecorrige{0}\input{Eval_TP_Reseau_LAN.tex}"
%   pdflatex -jobname=Eval_Reseau_CorrigeA "\def\variante{A}\def\modecorrige{1}\input{Eval_TP_Reseau_LAN.tex}"
%   pdflatex -jobname=Eval_Reseau_CorrigeB "\def\variante{B}\def\modecorrige{1}\input{Eval_TP_Reseau_LAN.tex}"
\documentclass[10pt,a4paper]{article}

\usepackage[T1]{fontenc}
\usepackage[utf8]{inputenc}
\usepackage{lmodern}
\usepackage[french]{babel}
\usepackage[a4paper,hmargin=1.5cm,top=1.3cm,bottom=1.6cm,footskip=0.7cm]{geometry}
\usepackage[table]{xcolor}
\usepackage{booktabs,tabularx,array}
\usepackage{enumitem}
\usepackage{fancyhdr}
\usepackage[most]{tcolorbox}
\usepackage{tikz}
\usetikzlibrary{positioning}
\usepackage{titlesec}

\providecommand{\modecorrige}{0}
\providecommand{\variante}{A}
\newif\ifcorrige
\ifnum\modecorrige=1 \corrigetrue \else \corrigefalse \fi

\definecolor{ipred}{HTML}{C0392B}
\definecolor{cdcbar}{HTML}{A67C3D}
\definecolor{cdcbg}{HTML}{F4F1EA}
\definecolor{qbar}{HTML}{4A6B5D}
\definecolor{qbg}{HTML}{EDF1EF}
\definecolor{corr}{HTML}{1F5FA8}
\definecolor{corrbg}{HTML}{E8F0FA}
\definecolor{grisfoot}{HTML}{828282}

\newcommand{\ip}[1]{\textcolor{ipred}{\texttt{#1}}}
\newcommand{\pts}[1]{\hfill\mbox{\textbf{(#1)}}}
\newcommand{\sol}[1]{\ifcorrige\textcolor{corr}{#1}\fi}
\newcommand{\bareme}[1]{\ifcorrige\par{\footnotesize\color{corr}#1\par}\fi}
\newcommand{\hl}{\ifcorrige\rule[-1.6mm]{0pt}{5.4mm}\else\rule[-2.2mm]{0pt}{6.6mm}\fi}
\newcolumntype{C}{>{\centering\arraybackslash}X}
\newcolumntype{L}{>{\raggedright\arraybackslash}X}

\titleformat{\section}{\large\bfseries}{}{0pt}{}
\titlespacing*{\section}{0pt}{6pt}{3pt}

\tcbset{qstyle/.style={enhanced,boxrule=0pt,leftrule=2.5pt,arc=0pt,outer arc=0pt,
  left=5pt,right=5pt,top=2pt,bottom=2pt,before skip=4pt,after skip=3pt}}
\newtcolorbox{cdc}{qstyle,colback=cdcbg,colframe=cdcbar}
\newtcolorbox{quest}{qstyle,colback=qbg,colframe=qbar}

% Zone de réponse : vierge dans le sujet, remplie dans le corrigé
\newcommand{\zone}[2]{%
  \ifcorrige
    \begin{tcolorbox}[colback=corrbg,colframe=corr,boxrule=0.5pt,arc=1pt,left=4pt,right=4pt,top=1pt,bottom=1pt,
      before skip=2pt,after skip=2pt,fontupper=\footnotesize\color{corr}]#2\end{tcolorbox}%
  \else
    \begin{tcolorbox}[colback=white,colframe=black!45,boxrule=0.4pt,arc=1pt,height=#1,
      before skip=2pt,after skip=2pt]\end{tcolorbox}%
  \fi}

\newcommand{\question}[3]{\begin{quest}\textbf{#1.} #3\pts{#2}\end{quest}}

%==========================================================================
% Données propres à chaque variante
%==========================================================================
\if\variante A
  \def\entreprise{Le cabinet d'architectes \textbf{ArchiNova}}
  \def\Rnom{R-ArchiNova}
  \def\resUn{Bureau d'études}\def\pUn{BE}\def\netUn{192.168.40}
  \def\resDeux{Accueil}\def\pDeux{ACC}\def\netDeux{172.16}
  \def\srvNom{SRV-Intranet}\def\netSrv{192.168.200}\def\srvH{10}
  \def\ssid{ArchiNova-WiFi}
  \def\AIrows{%
    \hl\ip{192.168.40.25} & \sol{C} & \sol{255.255.255.0} & \sol{192.168.40.0} & \sol{192.168.40.255} & \sol{254} \\ \hline
    \hl\ip{172.16.5.10}   & \sol{B} & \sol{255.255.0.0}   & \sol{172.16.0.0}   & \sol{172.16.255.255} & \sol{65\,534} \\ \hline
    \hl\ip{10.3.2.1}      & \sol{A} & \sol{255.0.0.0}     & \sol{10.0.0.0}     & \sol{10.255.255.255} & \sol{16\,777\,214} \\ \hline}
  \def\apipa{169.254.12.7}
  \def\Cnet{192.168.60}
  \def\Crows{%
    PC-A & \ip{192.168.60.10} & \ip{255.255.255.0} & \ip{192.168.60.254} & — \\
    PC-B & \ip{192.168.60.20} & \ip{255.255.255.0} & \ip{192.168.60.1}   & — \\
    PC-C & \ip{192.168.61.30} & \ip{255.255.255.0} & \ip{192.168.60.254} & — \\
    PC-D & \ip{192.168.60.10} & \ip{255.255.255.0} & \ip{192.168.60.254} & — \\
    SRV-C & \ip{10.0.0.5}     & \ip{255.0.0.0}     & \ip{10.0.0.254}     & — \\
    R1 — G0/0 & \ip{192.168.60.254} & \ip{255.255.255.0} & — & On \\
    R1 — G0/1 & \ip{10.0.0.254}     & \ip{255.0.0.0}     & — & Off \\}
  \def\Csol{%
    \hl 1 & \sol{PC-B} & \sol{Passerelle .60.1 : pas l'adresse du routeur.} & \sol{Joint son réseau, mais aucun autre (serveur).} & \sol{Passerelle 192.168.60.254.} \\ \hline
    \hl 2 & \sol{PC-C} & \sol{IP hors du réseau 192.168.60.0/24.} & \sol{Ne joint ni les postes ni le routeur.} & \sol{Ex. 192.168.60.30.} \\ \hline
    \hl 3 & \sol{PC-D} & \sol{Même IP que PC-A (doublon).} & \sol{Conflit d'adresse : PC-A et PC-D perturbés.} & \sol{Adresse unique, ex. .60.40.} \\ \hline
    \hl 4 & \sol{R1 G0/1} & \sol{Interface désactivée (Off).} & \sol{Réseau du serveur injoignable pour tous.} & \sol{Port Status : On.} \\ \hline}
  \def\Cbareme{PC-A et SRV-C sont correctement configurés ; accepter PC-A cité à la place de PC-D pour le doublon.}
\else
  \def\entreprise{La clinique vétérinaire \textbf{VetoPlus}}
  \def\Rnom{R-VetoPlus}
  \def\resUn{Soins}\def\pUn{SOI}\def\netUn{192.168.70}
  \def\resDeux{Secrétariat}\def\pDeux{SEC}\def\netDeux{172.20}
  \def\srvNom{SRV-Dossiers}\def\netSrv{192.168.250}\def\srvH{20}
  \def\ssid{VetoPlus-WiFi}
  \def\AIrows{%
    \hl\ip{192.168.15.130} & \sol{C} & \sol{255.255.255.0} & \sol{192.168.15.0} & \sol{192.168.15.255} & \sol{254} \\ \hline
    \hl\ip{172.31.200.7}   & \sol{B} & \sol{255.255.0.0}   & \sol{172.31.0.0}   & \sol{172.31.255.255} & \sol{65\,534} \\ \hline
    \hl\ip{10.45.6.9}      & \sol{A} & \sol{255.0.0.0}     & \sol{10.0.0.0}     & \sol{10.255.255.255} & \sol{16\,777\,214} \\ \hline}
  \def\apipa{169.254.88.3}
  \def\Cnet{192.168.80}
  \def\Crows{%
    PC-A & \ip{192.168.8.11}   & \ip{255.255.255.0} & \ip{192.168.80.254} & — \\
    PC-B & \ip{192.168.80.12}  & \ip{255.255.255.0} & \ip{192.168.80.254} & — \\
    PC-C & \ip{192.168.80.13}  & \ip{255.255.255.0} & \ip{192.168.80.100} & — \\
    PC-D & \ip{192.168.80.12}  & \ip{255.255.255.0} & \ip{192.168.80.254} & — \\
    SRV-C & \ip{10.0.0.8}      & \ip{255.0.0.0}     & \ip{10.0.0.254}     & — \\
    R1 — G0/0 & \ip{192.168.80.254} & \ip{255.255.255.0} & — & Off \\
    R1 — G0/1 & \ip{10.0.0.254}     & \ip{255.0.0.0}     & — & On \\}
  \def\Csol{%
    \hl 1 & \sol{PC-A} & \sol{IP hors du réseau 192.168.80.0/24.} & \sol{Ne joint ni les postes ni le routeur.} & \sol{Ex. 192.168.80.11.} \\ \hline
    \hl 2 & \sol{PC-C} & \sol{Passerelle .80.100 : pas l'adresse du routeur.} & \sol{Joint son réseau, mais aucun autre (serveur).} & \sol{Passerelle 192.168.80.254.} \\ \hline
    \hl 3 & \sol{PC-D} & \sol{Même IP que PC-B (doublon).} & \sol{Conflit d'adresse : PC-B et PC-D perturbés.} & \sol{Adresse unique, ex. .80.14.} \\ \hline
    \hl 4 & \sol{R1 G0/0} & \sol{Interface côté LAN désactivée (Off).} & \sol{Aucun poste ne joint le routeur ni le serveur.} & \sol{Port Status : On.} \\ \hline}
  \def\Cbareme{PC-B et SRV-C sont correctement configurés ; accepter PC-B cité à la place de PC-D pour le doublon.}
\fi

\pagestyle{fancy}
\fancyhf{}
\renewcommand{\headrulewidth}{0pt}
\fancyfoot[L]{\footnotesize\itshape\color{grisfoot}Évaluation — Réseaux locaux — Sujet \variante\ifcorrige\ — CORRIGÉ\fi}
\fancyfoot[R]{\footnotesize\color{grisfoot}\thepage/3}

\setlist{itemsep=0pt,topsep=2pt,parsep=0pt}
\setlength{\parindent}{0pt}
\setlength{\parskip}{2pt}

\begin{document}

\begin{center}
  {\Large\bfseries Évaluation pratique — Réseaux locaux \quad
   \fcolorbox{black}{white}{\ SUJET \variante\ }\par}
  \vspace{1pt}
  {\bfseries Adressage IP, commutation et routage — BTS CIEL 1re année — durée : 2 heures\par}
  \ifcorrige\vspace{2pt}{\bfseries\color{corr}CORRIGÉ ET BARÈME — document enseignant}\par\fi
\end{center}
\vspace{-2pt}
{\renewcommand{\arraystretch}{1.35}
\begin{tabularx}{\linewidth}{|X|X|p{2.2cm}|p{2.6cm}|}
  \hline
  \textbf{Nom :} & \textbf{Prénom :} & \textbf{Poste n°} & \textbf{Note :} \hfill /20 \\ \hline
\end{tabularx}}

\textbf{Consignes :} travail \textbf{individuel}, documents et Internet \textbf{non autorisés}. Ce sujet est le
\textbf{document-réponse}. Fichier Packet Tracer : \ip{Nom\_Prenom\_EvalLAN\_\variante.pkt}, à enregistrer
régulièrement dans le dossier de rendu. Les « visas » sont apposés par le professeur, que vous appelez à chaque test réussi.
Durées conseillées : A 30 min — B 1 h 10 — C 20 min.

%==========================================================================
\section{Partie A — Questions de cours et calculs \hfill /6}

\question{A1}{2 pts}{Complétez le tableau en considérant le \textbf{masque par défaut} de la classe de chaque adresse.}
{\small\renewcommand{\arraystretch}{1.15}
\begin{tabularx}{\linewidth}{|l|c|C|C|C|c|}
  \hline
  \textbf{Adresse IP} & \textbf{Classe} & \textbf{Masque par défaut} & \textbf{Adresse réseau} &
  \textbf{Adresse de diffusion} & \textbf{Nb d'hôtes} \\ \hline
  \AIrows
\end{tabularx}}
\bareme{0,5 pt par ligne juste + 0,5 pt si les trois nombres d'hôtes sont justes ($2^{n}-2$, $n$ = nombre de bits d'hôte).}

\question{A2}{1 pt}{Indiquez le type de câble (\textbf{droit} ou \textbf{croisé}) pour chaque liaison.}
{\small\renewcommand{\arraystretch}{1.15}
\begin{tabularx}{\linewidth}{|X|c||X|c|}
  \hline
  \hl PC $\leftrightarrow$ commutateur & \makebox[1.6cm]{\sol{droit}} & Commutateur $\leftrightarrow$ routeur & \makebox[1.6cm]{\sol{droit}} \\ \hline
  \hl PC $\leftrightarrow$ PC (réseau P2P) & \sol{croisé} & Concentrateur $\leftrightarrow$ concentrateur & \sol{croisé} \\ \hline
\end{tabularx}}
\bareme{0,25 pt par case. Équipements de types différents : câble droit ; de même type : câble croisé.}

\question{A3}{1 pt}{Comment un \textbf{concentrateur (hub)} et un \textbf{commutateur (switch)} traitent-ils une trame
reçue ? Lequel limite les collisions, et pourquoi ?}
\zone{2.9cm}{Le \textbf{hub} (couche 1) répète la trame sur \textbf{tous ses autres ports} : un seul domaine de collision
partagé. Le \textbf{switch} (couche 2) lit l'adresse MAC de destination et, grâce à sa \textbf{table MAC}, ne l'envoie
\textbf{que sur le port du destinataire} : un domaine de collision par port, c'est lui qui limite les collisions.
\emph{0,5 pt hub, 0,5 pt switch.}}

\question{A4}{1 pt}{Qu'est-ce que la \textbf{passerelle par défaut} d'un poste ? Quelle condition son adresse doit-elle respecter ?}
\zone{2.4cm}{Adresse IP de l'interface du \textbf{routeur} à laquelle le poste envoie les paquets destinés à un
\textbf{autre réseau}. Elle doit \textbf{appartenir au même réseau} que le poste. \emph{0,5 + 0,5 pt.}}

\question{A5}{0,5 pt}{Un portable configuré en DHCP affiche \ip{\apipa} / \ip{255.255.0.0}. Que s'est-il passé ?}
\zone{1.5cm}{Aucun \textbf{serveur DHCP} n'a répondu : le poste s'est attribué une adresse \textbf{APIPA} (169.254.0.0/16).}

\question{A6}{0,5 pt}{Quel protocole utilise \texttt{ping} ? Nommez les deux messages échangés et la couche OSI des adresses IP.}
\zone{1.5cm}{\textbf{ICMP} ; \textbf{Echo Request} / \textbf{Echo Reply} ; adresses IP en \textbf{couche 3 (Réseau)}.}

%==========================================================================
\newpage
\section{Partie B — Réalisation sous Packet Tracer \hfill /10}

\begin{cdc}
\entreprise{} vous confie son réseau : \textbf{trois réseaux} reliés par un \textbf{routeur unique Cisco 2911}
nommé \ip{\Rnom} (interfaces G0/0, G0/1, G0/2).
\begin{itemize}[leftmargin=12pt]
  \item \textbf{Réseau \resUn} : \textbf{4 postes} PC-\pUn1 à PC-\pUn4, commutateur \ip{SW-\pUn}, adressage statique,
        réseau \ip{\netUn.0/24}.
  \item \textbf{Réseau \resDeux} : \textbf{2 postes} PC-\pDeux1, PC-\pDeux2 et \textbf{1 imprimante} IMP-\pDeux,
        commutateur \ip{SW-\pDeux}, adressage statique, réseau de classe B \ip{\netDeux.0.0/16}.
  \item \textbf{Réseau Serveur} : serveur \ip{\srvNom} d'adresse \ip{\netSrv.\srvH/24}, relié \textbf{directement} au routeur.
  \item La passerelle de chaque réseau est sa \textbf{dernière adresse utilisable}. Tous les équipements communiquent
        entre eux et accèdent à la page web du serveur. Nommage imposé (« Display Name »).
\end{itemize}
\end{cdc}

\question{B1}{2 pts}{Complétez le plan d'adressage \textbf{avant} la réalisation.}
{\footnotesize\setlength{\tabcolsep}{3pt}\renewcommand{\arraystretch}{1.15}
\begin{tabularx}{\linewidth}{|l|C|C|C|C|C|C|}
  \hline
  \textbf{Réseau} & \textbf{Adresse réseau} & \textbf{Masque} & \textbf{1re @ utilisable} &
  \textbf{Dernière @ utilisable} & \textbf{Diffusion} & \textbf{Passerelle} \\ \hline
  \hl \resUn & \sol{\netUn.0} & \sol{255.255.255.0} & \sol{\netUn.1} & \sol{\netUn.254} & \sol{\netUn.255} & \sol{\netUn.254} \\ \hline
  \hl \resDeux & \sol{\netDeux.0.0} & \sol{255.255.0.0} & \sol{\netDeux.0.1} & \sol{\netDeux.255.254} & \sol{\netDeux.255.255} & \sol{\netDeux.255.254} \\ \hline
  \hl Serveur & \sol{\netSrv.0} & \sol{255.255.255.0} & \sol{\netSrv.1} & \sol{\netSrv.254} & \sol{\netSrv.255} & \sol{\netSrv.254} \\ \hline
\end{tabularx}}
\bareme{0,5 pt par ligne juste + 0,5 pt si la convention de passerelle est respectée partout.}

\question{B2}{2 pts}{Construisez la topologie (équipements, câblage, nommage) et relevez l'adressage configuré.}
{\footnotesize\setlength{\tabcolsep}{3pt}\renewcommand{\arraystretch}{1.1}
\begin{tabularx}{\linewidth}{|l|C|C|C||l|C|C|}
  \hline
  \textbf{Hôte} & \textbf{Adresse IP} & \textbf{Masque} & \textbf{Passerelle} & \textbf{Interface} & \textbf{Adresse IP} & \textbf{Masque} \\ \hline
  \hl PC-\pUn1 & \sol{\netUn.1} & \sol{255.255.255.0} & \sol{\netUn.254} & \Rnom{} G0/0 & \sol{\netUn.254} & \sol{255.255.255.0} \\ \hline
  \hl PC-\pDeux1 & \sol{\netDeux.0.1} & \sol{255.255.0.0} & \sol{\netDeux.255.254} & \Rnom{} G0/1 & \sol{\netDeux.255.254} & \sol{255.255.0.0} \\ \hline
  \hl IMP-\pDeux & \sol{\netDeux.0.10} & \sol{255.255.0.0} & \sol{\netDeux.255.254} & \Rnom{} G0/2 & \sol{\netSrv.254} & \sol{255.255.255.0} \\ \hline
  \hl \srvNom & \ip{\netSrv.\srvH} & \sol{255.255.255.0} & \sol{\netSrv.254} & \multicolumn{3}{c|}{\cellcolor{black!8}} \\ \hline
\end{tabularx}}
\bareme{Adresses d'hôtes libres (uniques, dans le bon réseau) ; affectation G0/x libre. Sur le .pkt : équipements 0,5 ;
câblage 1 (câbles droits vers les switchs et switch–routeur, \textbf{croisé} serveur–routeur, liens verts) ; nommage 0,5.}

\question{B3}{3 pts}{Configurez tous les équipements : IP, masques, passerelles des hôtes ; interfaces du routeur
(« Port Status : On »).}
\bareme{Sur le .pkt : interfaces routeur 1 pt ; IP/masques des hôtes 1 pt ; passerelles (imprimante et serveur compris)
1 pt. $-0,25$ par erreur, sans descendre sous 0 par item.}

\question{B4}{3 pts}{Réalisez les tests, notez la commande et le résultat, puis faites viser chaque test réussi.}
{\small\renewcommand{\arraystretch}{1.1}
\begin{tabularx}{\linewidth}{|c|l|l|C|c|c|}
  \hline
  \textbf{N°} & \textbf{Source} & \textbf{Destination} & \textbf{Commande / action} & \textbf{Résultat} & \textbf{Visa prof.} \\ \hline
  \hl 1 & PC-\pUn1 & PC-\pUn4 & \sol{ping \netUn.4} & \sol{réussi} & \hspace{1.4cm} \\ \hline
  \hl 2 & PC-\pDeux1 & IMP-\pDeux & \sol{ping \netDeux.0.10} & \sol{réussi} & \\ \hline
  \hl 3 & PC-\pUn2 & PC-\pDeux2 & \sol{ping \netDeux.0.2} & \sol{réussi} & \\ \hline
  \hl 4 & PC-\pDeux2 & \srvNom & \sol{ping \netSrv.\srvH} & \sol{réussi} & \\ \hline
  \hl 5 & PC-\pUn3 & page web du serveur & \sol{navigateur : {\NoAutoSpacing http://\netSrv.\srvH}} & \sol{page affichée} & \\ \hline
\end{tabularx}}
\bareme{0,5 pt par test visé + 0,5 pt pour B5. Le premier ping inter-réseaux peut échouer (résolution ARP) ; les suivants doivent réussir.}

\question{B5}{incl. B4}{Test n°3 : par quels équipements passe le paquet ? Pourquoi PC-\pUn2 l'envoie-t-il au routeur ?}
\zone{2.3cm}{PC-\pUn2 $\to$ SW-\pUn{} $\to$ \Rnom{} $\to$ SW-\pDeux{} $\to$ PC-\pDeux2. Grâce à son masque, PC-\pUn2 constate que
la destination (\netDeux.0.0/16) n'est \textbf{pas dans son réseau} (\netUn.0/24) : il envoie le paquet à sa
\textbf{passerelle par défaut}, qui le route.}

%==========================================================================
\newpage
\section{Partie C — Diagnostic d'un réseau en panne \hfill /4}

Un technicien constate qu'\textbf{aucun poste ne parvient à joindre le serveur SRV-C} et que certains postes
communiquent mal entre eux. Les câbles sont du bon type.

\begin{center}
\begin{minipage}{0.37\linewidth}\centering
\begin{tikzpicture}[font=\footnotesize,
  eq/.style={draw,fill=qbg,rounded corners=2pt,minimum width=9mm,minimum height=5mm,inner sep=2pt},
  every path/.style={thick}]
  \node[eq] (SW) {SW-1};
  \node[eq,right=9mm of SW] (R) {R1};
  \node[eq,right=7mm of R] (SRV) {SRV-C};
  \node[eq,above left=5mm and -3mm of SW] (A) {PC-A};
  \node[eq,above right=5mm and -3mm of SW] (B) {PC-B};
  \node[eq,below left=5mm and -3mm of SW] (C) {PC-C};
  \node[eq,below right=5mm and -3mm of SW] (D) {PC-D};
  \foreach \n in {A,B,C,D} \draw (\n) -- (SW);
  \draw (SW) -- node[above,font=\tiny]{G0/0} (R);
  \draw (R) -- node[above,font=\tiny]{G0/1} (SRV);
\end{tikzpicture}
\end{minipage}\hfill
\begin{minipage}{0.61\linewidth}\centering
{\small\renewcommand{\arraystretch}{1.05}\setlength{\tabcolsep}{4pt}
\begin{tabular}{lllll}
  \toprule
  \textbf{Équipement} & \textbf{Adresse IP} & \textbf{Masque} & \textbf{Passerelle} & \textbf{Port} \\
  \midrule
  \Crows
  \bottomrule
\end{tabular}}
\end{minipage}
\end{center}

\question{C1}{4 pts}{Relevez les \textbf{quatre anomalies}. Pour chacune : conséquence et correction.}
{\small\renewcommand{\arraystretch}{1.1}
\ifcorrige\def\hc{}\else\def\hc{\rule[-10mm]{0pt}{14mm}}\fi
\begin{tabularx}{\linewidth}{|c|>{\raggedright\arraybackslash}p{2cm}|L|L|L|}
  \hline
  \textbf{N°} & \textbf{Équipement} & \textbf{Anomalie} & \textbf{Conséquence} & \textbf{Correction} \\ \hline
  \ifcorrige\Csol\else
  \hc 1 & & & & \\ \hline
  \hc 2 & & & & \\ \hline
  \hc 3 & & & & \\ \hline
  \hc 4 & & & & \\ \hline
  \fi
\end{tabularx}}
\bareme{1 pt par anomalie (0,25 équipement, anomalie, conséquence, correction ; ordre indifférent). \Cbareme}

%==========================================================================
\section{Bonus — Extension sans fil \hfill +1}

\begin{cdc}Ajoutez au réseau \resUn{} un \textbf{AccessPoint-PT} relié à SW-\pUn{} (SSID \ip{\ssid}, \textbf{WPA2-PSK})
et un portable \ip{Laptop-\pUn} équipé d'un module \texttt{WPC300N}, en adressage \textbf{statique}.
Vérifiez qu'il joint \srvNom.\end{cdc}
{\small\renewcommand{\arraystretch}{1.1}
\begin{tabularx}{\linewidth}{|l|C|C|C|c|}
  \hline
  \textbf{Équipement} & \textbf{Adresse IP} & \textbf{Masque} & \textbf{Passerelle} & \textbf{Visa prof.} \\ \hline
  \hl Laptop-\pUn & \sol{\netUn.50} & \sol{255.255.255.0} & \sol{\netUn.254} & \hspace{1.4cm} \\ \hline
\end{tabularx}}
\bareme{Le point d'accès fait office de pont : le portable est dans \netUn.0/24. Bonus si associé en WPA2-PSK et ping
vers \netSrv.\srvH{} réussi. Note plafonnée à 20.}

\end{document}
"
\documentclass[10pt,a4paper]{article}

\usepackage[T1]{fontenc}
\usepackage[utf8]{inputenc}
\usepackage{lmodern}
\usepackage[french]{babel}
\usepackage[a4paper,hmargin=1.5cm,top=1.3cm,bottom=1.6cm,footskip=0.7cm]{geometry}
\usepackage[table]{xcolor}
\usepackage{booktabs,tabularx,array}
\usepackage{enumitem}
\usepackage{fancyhdr}
\usepackage[most]{tcolorbox}
\usepackage{tikz}
\usetikzlibrary{positioning}
\usepackage{titlesec}

\providecommand{\modecorrige}{0}
\providecommand{\variante}{A}
\newif\ifcorrige
\ifnum\modecorrige=1 \corrigetrue \else \corrigefalse \fi

\definecolor{ipred}{HTML}{C0392B}
\definecolor{cdcbar}{HTML}{A67C3D}
\definecolor{cdcbg}{HTML}{F4F1EA}
\definecolor{qbar}{HTML}{4A6B5D}
\definecolor{qbg}{HTML}{EDF1EF}
\definecolor{corr}{HTML}{1F5FA8}
\definecolor{corrbg}{HTML}{E8F0FA}
\definecolor{grisfoot}{HTML}{828282}

\newcommand{\ip}[1]{\textcolor{ipred}{\texttt{#1}}}
\newcommand{\pts}[1]{\hfill\mbox{\textbf{(#1)}}}
\newcommand{\sol}[1]{\ifcorrige\textcolor{corr}{#1}\fi}
\newcommand{\bareme}[1]{\ifcorrige\par{\footnotesize\color{corr}#1\par}\fi}
\newcommand{\hl}{\ifcorrige\rule[-1.6mm]{0pt}{5.4mm}\else\rule[-2.2mm]{0pt}{6.6mm}\fi}
\newcolumntype{C}{>{\centering\arraybackslash}X}
\newcolumntype{L}{>{\raggedright\arraybackslash}X}

\titleformat{\section}{\large\bfseries}{}{0pt}{}
\titlespacing*{\section}{0pt}{6pt}{3pt}

\tcbset{qstyle/.style={enhanced,boxrule=0pt,leftrule=2.5pt,arc=0pt,outer arc=0pt,
  left=5pt,right=5pt,top=2pt,bottom=2pt,before skip=4pt,after skip=3pt}}
\newtcolorbox{cdc}{qstyle,colback=cdcbg,colframe=cdcbar}
\newtcolorbox{quest}{qstyle,colback=qbg,colframe=qbar}

% Zone de réponse : vierge dans le sujet, remplie dans le corrigé
\newcommand{\zone}[2]{%
  \ifcorrige
    \begin{tcolorbox}[colback=corrbg,colframe=corr,boxrule=0.5pt,arc=1pt,left=4pt,right=4pt,top=1pt,bottom=1pt,
      before skip=2pt,after skip=2pt,fontupper=\footnotesize\color{corr}]#2\end{tcolorbox}%
  \else
    \begin{tcolorbox}[colback=white,colframe=black!45,boxrule=0.4pt,arc=1pt,height=#1,
      before skip=2pt,after skip=2pt]\end{tcolorbox}%
  \fi}

\newcommand{\question}[3]{\begin{quest}\textbf{#1.} #3\pts{#2}\end{quest}}

%==========================================================================
% Données propres à chaque variante
%==========================================================================
\if\variante A
  \def\entreprise{Le cabinet d'architectes \textbf{ArchiNova}}
  \def\Rnom{R-ArchiNova}
  \def\resUn{Bureau d'études}\def\pUn{BE}\def\netUn{192.168.40}
  \def\resDeux{Accueil}\def\pDeux{ACC}\def\netDeux{172.16}
  \def\srvNom{SRV-Intranet}\def\netSrv{192.168.200}\def\srvH{10}
  \def\ssid{ArchiNova-WiFi}
  \def\AIrows{%
    \hl\ip{192.168.40.25} & \sol{C} & \sol{255.255.255.0} & \sol{192.168.40.0} & \sol{192.168.40.255} & \sol{254} \\ \hline
    \hl\ip{172.16.5.10}   & \sol{B} & \sol{255.255.0.0}   & \sol{172.16.0.0}   & \sol{172.16.255.255} & \sol{65\,534} \\ \hline
    \hl\ip{10.3.2.1}      & \sol{A} & \sol{255.0.0.0}     & \sol{10.0.0.0}     & \sol{10.255.255.255} & \sol{16\,777\,214} \\ \hline}
  \def\apipa{169.254.12.7}
  \def\Cnet{192.168.60}
  \def\Crows{%
    PC-A & \ip{192.168.60.10} & \ip{255.255.255.0} & \ip{192.168.60.254} & — \\
    PC-B & \ip{192.168.60.20} & \ip{255.255.255.0} & \ip{192.168.60.1}   & — \\
    PC-C & \ip{192.168.61.30} & \ip{255.255.255.0} & \ip{192.168.60.254} & — \\
    PC-D & \ip{192.168.60.10} & \ip{255.255.255.0} & \ip{192.168.60.254} & — \\
    SRV-C & \ip{10.0.0.5}     & \ip{255.0.0.0}     & \ip{10.0.0.254}     & — \\
    R1 — G0/0 & \ip{192.168.60.254} & \ip{255.255.255.0} & — & On \\
    R1 — G0/1 & \ip{10.0.0.254}     & \ip{255.0.0.0}     & — & Off \\}
  \def\Csol{%
    \hl 1 & \sol{PC-B} & \sol{Passerelle .60.1 : pas l'adresse du routeur.} & \sol{Joint son réseau, mais aucun autre (serveur).} & \sol{Passerelle 192.168.60.254.} \\ \hline
    \hl 2 & \sol{PC-C} & \sol{IP hors du réseau 192.168.60.0/24.} & \sol{Ne joint ni les postes ni le routeur.} & \sol{Ex. 192.168.60.30.} \\ \hline
    \hl 3 & \sol{PC-D} & \sol{Même IP que PC-A (doublon).} & \sol{Conflit d'adresse : PC-A et PC-D perturbés.} & \sol{Adresse unique, ex. .60.40.} \\ \hline
    \hl 4 & \sol{R1 G0/1} & \sol{Interface désactivée (Off).} & \sol{Réseau du serveur injoignable pour tous.} & \sol{Port Status : On.} \\ \hline}
  \def\Cbareme{PC-A et SRV-C sont correctement configurés ; accepter PC-A cité à la place de PC-D pour le doublon.}
\else
  \def\entreprise{La clinique vétérinaire \textbf{VetoPlus}}
  \def\Rnom{R-VetoPlus}
  \def\resUn{Soins}\def\pUn{SOI}\def\netUn{192.168.70}
  \def\resDeux{Secrétariat}\def\pDeux{SEC}\def\netDeux{172.20}
  \def\srvNom{SRV-Dossiers}\def\netSrv{192.168.250}\def\srvH{20}
  \def\ssid{VetoPlus-WiFi}
  \def\AIrows{%
    \hl\ip{192.168.15.130} & \sol{C} & \sol{255.255.255.0} & \sol{192.168.15.0} & \sol{192.168.15.255} & \sol{254} \\ \hline
    \hl\ip{172.31.200.7}   & \sol{B} & \sol{255.255.0.0}   & \sol{172.31.0.0}   & \sol{172.31.255.255} & \sol{65\,534} \\ \hline
    \hl\ip{10.45.6.9}      & \sol{A} & \sol{255.0.0.0}     & \sol{10.0.0.0}     & \sol{10.255.255.255} & \sol{16\,777\,214} \\ \hline}
  \def\apipa{169.254.88.3}
  \def\Cnet{192.168.80}
  \def\Crows{%
    PC-A & \ip{192.168.8.11}   & \ip{255.255.255.0} & \ip{192.168.80.254} & — \\
    PC-B & \ip{192.168.80.12}  & \ip{255.255.255.0} & \ip{192.168.80.254} & — \\
    PC-C & \ip{192.168.80.13}  & \ip{255.255.255.0} & \ip{192.168.80.100} & — \\
    PC-D & \ip{192.168.80.12}  & \ip{255.255.255.0} & \ip{192.168.80.254} & — \\
    SRV-C & \ip{10.0.0.8}      & \ip{255.0.0.0}     & \ip{10.0.0.254}     & — \\
    R1 — G0/0 & \ip{192.168.80.254} & \ip{255.255.255.0} & — & Off \\
    R1 — G0/1 & \ip{10.0.0.254}     & \ip{255.0.0.0}     & — & On \\}
  \def\Csol{%
    \hl 1 & \sol{PC-A} & \sol{IP hors du réseau 192.168.80.0/24.} & \sol{Ne joint ni les postes ni le routeur.} & \sol{Ex. 192.168.80.11.} \\ \hline
    \hl 2 & \sol{PC-C} & \sol{Passerelle .80.100 : pas l'adresse du routeur.} & \sol{Joint son réseau, mais aucun autre (serveur).} & \sol{Passerelle 192.168.80.254.} \\ \hline
    \hl 3 & \sol{PC-D} & \sol{Même IP que PC-B (doublon).} & \sol{Conflit d'adresse : PC-B et PC-D perturbés.} & \sol{Adresse unique, ex. .80.14.} \\ \hline
    \hl 4 & \sol{R1 G0/0} & \sol{Interface côté LAN désactivée (Off).} & \sol{Aucun poste ne joint le routeur ni le serveur.} & \sol{Port Status : On.} \\ \hline}
  \def\Cbareme{PC-B et SRV-C sont correctement configurés ; accepter PC-B cité à la place de PC-D pour le doublon.}
\fi

\pagestyle{fancy}
\fancyhf{}
\renewcommand{\headrulewidth}{0pt}
\fancyfoot[L]{\footnotesize\itshape\color{grisfoot}Évaluation — Réseaux locaux — Sujet \variante\ifcorrige\ — CORRIGÉ\fi}
\fancyfoot[R]{\footnotesize\color{grisfoot}\thepage/3}

\setlist{itemsep=0pt,topsep=2pt,parsep=0pt}
\setlength{\parindent}{0pt}
\setlength{\parskip}{2pt}

\begin{document}

\begin{center}
  {\Large\bfseries Évaluation pratique — Réseaux locaux \quad
   \fcolorbox{black}{white}{\ SUJET \variante\ }\par}
  \vspace{1pt}
  {\bfseries Adressage IP, commutation et routage — BTS CIEL 1re année — durée : 2 heures\par}
  \ifcorrige\vspace{2pt}{\bfseries\color{corr}CORRIGÉ ET BARÈME — document enseignant}\par\fi
\end{center}
\vspace{-2pt}
{\renewcommand{\arraystretch}{1.35}
\begin{tabularx}{\linewidth}{|X|X|p{2.2cm}|p{2.6cm}|}
  \hline
  \textbf{Nom :} & \textbf{Prénom :} & \textbf{Poste n°} & \textbf{Note :} \hfill /20 \\ \hline
\end{tabularx}}

\textbf{Consignes :} travail \textbf{individuel}, documents et Internet \textbf{non autorisés}. Ce sujet est le
\textbf{document-réponse}. Fichier Packet Tracer : \ip{Nom\_Prenom\_EvalLAN\_\variante.pkt}, à enregistrer
régulièrement dans le dossier de rendu. Les « visas » sont apposés par le professeur, que vous appelez à chaque test réussi.
Durées conseillées : A 30 min — B 1 h 10 — C 20 min.

%==========================================================================
\section{Partie A — Questions de cours et calculs \hfill /6}

\question{A1}{2 pts}{Complétez le tableau en considérant le \textbf{masque par défaut} de la classe de chaque adresse.}
{\small\renewcommand{\arraystretch}{1.15}
\begin{tabularx}{\linewidth}{|l|c|C|C|C|c|}
  \hline
  \textbf{Adresse IP} & \textbf{Classe} & \textbf{Masque par défaut} & \textbf{Adresse réseau} &
  \textbf{Adresse de diffusion} & \textbf{Nb d'hôtes} \\ \hline
  \AIrows
\end{tabularx}}
\bareme{0,5 pt par ligne juste + 0,5 pt si les trois nombres d'hôtes sont justes ($2^{n}-2$, $n$ = nombre de bits d'hôte).}

\question{A2}{1 pt}{Indiquez le type de câble (\textbf{droit} ou \textbf{croisé}) pour chaque liaison.}
{\small\renewcommand{\arraystretch}{1.15}
\begin{tabularx}{\linewidth}{|X|c||X|c|}
  \hline
  \hl PC $\leftrightarrow$ commutateur & \makebox[1.6cm]{\sol{droit}} & Commutateur $\leftrightarrow$ routeur & \makebox[1.6cm]{\sol{droit}} \\ \hline
  \hl PC $\leftrightarrow$ PC (réseau P2P) & \sol{croisé} & Concentrateur $\leftrightarrow$ concentrateur & \sol{croisé} \\ \hline
\end{tabularx}}
\bareme{0,25 pt par case. Équipements de types différents : câble droit ; de même type : câble croisé.}

\question{A3}{1 pt}{Comment un \textbf{concentrateur (hub)} et un \textbf{commutateur (switch)} traitent-ils une trame
reçue ? Lequel limite les collisions, et pourquoi ?}
\zone{2.9cm}{Le \textbf{hub} (couche 1) répète la trame sur \textbf{tous ses autres ports} : un seul domaine de collision
partagé. Le \textbf{switch} (couche 2) lit l'adresse MAC de destination et, grâce à sa \textbf{table MAC}, ne l'envoie
\textbf{que sur le port du destinataire} : un domaine de collision par port, c'est lui qui limite les collisions.
\emph{0,5 pt hub, 0,5 pt switch.}}

\question{A4}{1 pt}{Qu'est-ce que la \textbf{passerelle par défaut} d'un poste ? Quelle condition son adresse doit-elle respecter ?}
\zone{2.4cm}{Adresse IP de l'interface du \textbf{routeur} à laquelle le poste envoie les paquets destinés à un
\textbf{autre réseau}. Elle doit \textbf{appartenir au même réseau} que le poste. \emph{0,5 + 0,5 pt.}}

\question{A5}{0,5 pt}{Un portable configuré en DHCP affiche \ip{\apipa} / \ip{255.255.0.0}. Que s'est-il passé ?}
\zone{1.5cm}{Aucun \textbf{serveur DHCP} n'a répondu : le poste s'est attribué une adresse \textbf{APIPA} (169.254.0.0/16).}

\question{A6}{0,5 pt}{Quel protocole utilise \texttt{ping} ? Nommez les deux messages échangés et la couche OSI des adresses IP.}
\zone{1.5cm}{\textbf{ICMP} ; \textbf{Echo Request} / \textbf{Echo Reply} ; adresses IP en \textbf{couche 3 (Réseau)}.}

%==========================================================================
\newpage
\section{Partie B — Réalisation sous Packet Tracer \hfill /10}

\begin{cdc}
\entreprise{} vous confie son réseau : \textbf{trois réseaux} reliés par un \textbf{routeur unique Cisco 2911}
nommé \ip{\Rnom} (interfaces G0/0, G0/1, G0/2).
\begin{itemize}[leftmargin=12pt]
  \item \textbf{Réseau \resUn} : \textbf{4 postes} PC-\pUn1 à PC-\pUn4, commutateur \ip{SW-\pUn}, adressage statique,
        réseau \ip{\netUn.0/24}.
  \item \textbf{Réseau \resDeux} : \textbf{2 postes} PC-\pDeux1, PC-\pDeux2 et \textbf{1 imprimante} IMP-\pDeux,
        commutateur \ip{SW-\pDeux}, adressage statique, réseau de classe B \ip{\netDeux.0.0/16}.
  \item \textbf{Réseau Serveur} : serveur \ip{\srvNom} d'adresse \ip{\netSrv.\srvH/24}, relié \textbf{directement} au routeur.
  \item La passerelle de chaque réseau est sa \textbf{dernière adresse utilisable}. Tous les équipements communiquent
        entre eux et accèdent à la page web du serveur. Nommage imposé (« Display Name »).
\end{itemize}
\end{cdc}

\question{B1}{2 pts}{Complétez le plan d'adressage \textbf{avant} la réalisation.}
{\footnotesize\setlength{\tabcolsep}{3pt}\renewcommand{\arraystretch}{1.15}
\begin{tabularx}{\linewidth}{|l|C|C|C|C|C|C|}
  \hline
  \textbf{Réseau} & \textbf{Adresse réseau} & \textbf{Masque} & \textbf{1re @ utilisable} &
  \textbf{Dernière @ utilisable} & \textbf{Diffusion} & \textbf{Passerelle} \\ \hline
  \hl \resUn & \sol{\netUn.0} & \sol{255.255.255.0} & \sol{\netUn.1} & \sol{\netUn.254} & \sol{\netUn.255} & \sol{\netUn.254} \\ \hline
  \hl \resDeux & \sol{\netDeux.0.0} & \sol{255.255.0.0} & \sol{\netDeux.0.1} & \sol{\netDeux.255.254} & \sol{\netDeux.255.255} & \sol{\netDeux.255.254} \\ \hline
  \hl Serveur & \sol{\netSrv.0} & \sol{255.255.255.0} & \sol{\netSrv.1} & \sol{\netSrv.254} & \sol{\netSrv.255} & \sol{\netSrv.254} \\ \hline
\end{tabularx}}
\bareme{0,5 pt par ligne juste + 0,5 pt si la convention de passerelle est respectée partout.}

\question{B2}{2 pts}{Construisez la topologie (équipements, câblage, nommage) et relevez l'adressage configuré.}
{\footnotesize\setlength{\tabcolsep}{3pt}\renewcommand{\arraystretch}{1.1}
\begin{tabularx}{\linewidth}{|l|C|C|C||l|C|C|}
  \hline
  \textbf{Hôte} & \textbf{Adresse IP} & \textbf{Masque} & \textbf{Passerelle} & \textbf{Interface} & \textbf{Adresse IP} & \textbf{Masque} \\ \hline
  \hl PC-\pUn1 & \sol{\netUn.1} & \sol{255.255.255.0} & \sol{\netUn.254} & \Rnom{} G0/0 & \sol{\netUn.254} & \sol{255.255.255.0} \\ \hline
  \hl PC-\pDeux1 & \sol{\netDeux.0.1} & \sol{255.255.0.0} & \sol{\netDeux.255.254} & \Rnom{} G0/1 & \sol{\netDeux.255.254} & \sol{255.255.0.0} \\ \hline
  \hl IMP-\pDeux & \sol{\netDeux.0.10} & \sol{255.255.0.0} & \sol{\netDeux.255.254} & \Rnom{} G0/2 & \sol{\netSrv.254} & \sol{255.255.255.0} \\ \hline
  \hl \srvNom & \ip{\netSrv.\srvH} & \sol{255.255.255.0} & \sol{\netSrv.254} & \multicolumn{3}{c|}{\cellcolor{black!8}} \\ \hline
\end{tabularx}}
\bareme{Adresses d'hôtes libres (uniques, dans le bon réseau) ; affectation G0/x libre. Sur le .pkt : équipements 0,5 ;
câblage 1 (câbles droits vers les switchs et switch–routeur, \textbf{croisé} serveur–routeur, liens verts) ; nommage 0,5.}

\question{B3}{3 pts}{Configurez tous les équipements : IP, masques, passerelles des hôtes ; interfaces du routeur
(« Port Status : On »).}
\bareme{Sur le .pkt : interfaces routeur 1 pt ; IP/masques des hôtes 1 pt ; passerelles (imprimante et serveur compris)
1 pt. $-0,25$ par erreur, sans descendre sous 0 par item.}

\question{B4}{3 pts}{Réalisez les tests, notez la commande et le résultat, puis faites viser chaque test réussi.}
{\small\renewcommand{\arraystretch}{1.1}
\begin{tabularx}{\linewidth}{|c|l|l|C|c|c|}
  \hline
  \textbf{N°} & \textbf{Source} & \textbf{Destination} & \textbf{Commande / action} & \textbf{Résultat} & \textbf{Visa prof.} \\ \hline
  \hl 1 & PC-\pUn1 & PC-\pUn4 & \sol{ping \netUn.4} & \sol{réussi} & \hspace{1.4cm} \\ \hline
  \hl 2 & PC-\pDeux1 & IMP-\pDeux & \sol{ping \netDeux.0.10} & \sol{réussi} & \\ \hline
  \hl 3 & PC-\pUn2 & PC-\pDeux2 & \sol{ping \netDeux.0.2} & \sol{réussi} & \\ \hline
  \hl 4 & PC-\pDeux2 & \srvNom & \sol{ping \netSrv.\srvH} & \sol{réussi} & \\ \hline
  \hl 5 & PC-\pUn3 & page web du serveur & \sol{navigateur : {\NoAutoSpacing http://\netSrv.\srvH}} & \sol{page affichée} & \\ \hline
\end{tabularx}}
\bareme{0,5 pt par test visé + 0,5 pt pour B5. Le premier ping inter-réseaux peut échouer (résolution ARP) ; les suivants doivent réussir.}

\question{B5}{incl. B4}{Test n°3 : par quels équipements passe le paquet ? Pourquoi PC-\pUn2 l'envoie-t-il au routeur ?}
\zone{2.3cm}{PC-\pUn2 $\to$ SW-\pUn{} $\to$ \Rnom{} $\to$ SW-\pDeux{} $\to$ PC-\pDeux2. Grâce à son masque, PC-\pUn2 constate que
la destination (\netDeux.0.0/16) n'est \textbf{pas dans son réseau} (\netUn.0/24) : il envoie le paquet à sa
\textbf{passerelle par défaut}, qui le route.}

%==========================================================================
\newpage
\section{Partie C — Diagnostic d'un réseau en panne \hfill /4}

Un technicien constate qu'\textbf{aucun poste ne parvient à joindre le serveur SRV-C} et que certains postes
communiquent mal entre eux. Les câbles sont du bon type.

\begin{center}
\begin{minipage}{0.37\linewidth}\centering
\begin{tikzpicture}[font=\footnotesize,
  eq/.style={draw,fill=qbg,rounded corners=2pt,minimum width=9mm,minimum height=5mm,inner sep=2pt},
  every path/.style={thick}]
  \node[eq] (SW) {SW-1};
  \node[eq,right=9mm of SW] (R) {R1};
  \node[eq,right=7mm of R] (SRV) {SRV-C};
  \node[eq,above left=5mm and -3mm of SW] (A) {PC-A};
  \node[eq,above right=5mm and -3mm of SW] (B) {PC-B};
  \node[eq,below left=5mm and -3mm of SW] (C) {PC-C};
  \node[eq,below right=5mm and -3mm of SW] (D) {PC-D};
  \foreach \n in {A,B,C,D} \draw (\n) -- (SW);
  \draw (SW) -- node[above,font=\tiny]{G0/0} (R);
  \draw (R) -- node[above,font=\tiny]{G0/1} (SRV);
\end{tikzpicture}
\end{minipage}\hfill
\begin{minipage}{0.61\linewidth}\centering
{\small\renewcommand{\arraystretch}{1.05}\setlength{\tabcolsep}{4pt}
\begin{tabular}{lllll}
  \toprule
  \textbf{Équipement} & \textbf{Adresse IP} & \textbf{Masque} & \textbf{Passerelle} & \textbf{Port} \\
  \midrule
  \Crows
  \bottomrule
\end{tabular}}
\end{minipage}
\end{center}

\question{C1}{4 pts}{Relevez les \textbf{quatre anomalies}. Pour chacune : conséquence et correction.}
{\small\renewcommand{\arraystretch}{1.1}
\ifcorrige\def\hc{}\else\def\hc{\rule[-10mm]{0pt}{14mm}}\fi
\begin{tabularx}{\linewidth}{|c|>{\raggedright\arraybackslash}p{2cm}|L|L|L|}
  \hline
  \textbf{N°} & \textbf{Équipement} & \textbf{Anomalie} & \textbf{Conséquence} & \textbf{Correction} \\ \hline
  \ifcorrige\Csol\else
  \hc 1 & & & & \\ \hline
  \hc 2 & & & & \\ \hline
  \hc 3 & & & & \\ \hline
  \hc 4 & & & & \\ \hline
  \fi
\end{tabularx}}
\bareme{1 pt par anomalie (0,25 équipement, anomalie, conséquence, correction ; ordre indifférent). \Cbareme}

%==========================================================================
\section{Bonus — Extension sans fil \hfill +1}

\begin{cdc}Ajoutez au réseau \resUn{} un \textbf{AccessPoint-PT} relié à SW-\pUn{} (SSID \ip{\ssid}, \textbf{WPA2-PSK})
et un portable \ip{Laptop-\pUn} équipé d'un module \texttt{WPC300N}, en adressage \textbf{statique}.
Vérifiez qu'il joint \srvNom.\end{cdc}
{\small\renewcommand{\arraystretch}{1.1}
\begin{tabularx}{\linewidth}{|l|C|C|C|c|}
  \hline
  \textbf{Équipement} & \textbf{Adresse IP} & \textbf{Masque} & \textbf{Passerelle} & \textbf{Visa prof.} \\ \hline
  \hl Laptop-\pUn & \sol{\netUn.50} & \sol{255.255.255.0} & \sol{\netUn.254} & \hspace{1.4cm} \\ \hline
\end{tabularx}}
\bareme{Le point d'accès fait office de pont : le portable est dans \netUn.0/24. Bonus si associé en WPA2-PSK et ping
vers \netSrv.\srvH{} réussi. Note plafonnée à 20.}

\end{document}
"
%   pdflatex -jobname=Eval_Reseau_CorrigeA "\def\variante{A}\def\modecorrige{1}% Évaluation pratique — Réseaux locaux (BTS CIEL 1re année) — sujets A et B
% Compilation (4 PDF) :
%   pdflatex -jobname=Eval_Reseau_SujetA   "\def\variante{A}\def\modecorrige{0}% Évaluation pratique — Réseaux locaux (BTS CIEL 1re année) — sujets A et B
% Compilation (4 PDF) :
%   pdflatex -jobname=Eval_Reseau_SujetA   "\def\variante{A}\def\modecorrige{0}\input{Eval_TP_Reseau_LAN.tex}"
%   pdflatex -jobname=Eval_Reseau_SujetB   "\def\variante{B}\def\modecorrige{0}\input{Eval_TP_Reseau_LAN.tex}"
%   pdflatex -jobname=Eval_Reseau_CorrigeA "\def\variante{A}\def\modecorrige{1}\input{Eval_TP_Reseau_LAN.tex}"
%   pdflatex -jobname=Eval_Reseau_CorrigeB "\def\variante{B}\def\modecorrige{1}\input{Eval_TP_Reseau_LAN.tex}"
\documentclass[10pt,a4paper]{article}

\usepackage[T1]{fontenc}
\usepackage[utf8]{inputenc}
\usepackage{lmodern}
\usepackage[french]{babel}
\usepackage[a4paper,hmargin=1.5cm,top=1.3cm,bottom=1.6cm,footskip=0.7cm]{geometry}
\usepackage[table]{xcolor}
\usepackage{booktabs,tabularx,array}
\usepackage{enumitem}
\usepackage{fancyhdr}
\usepackage[most]{tcolorbox}
\usepackage{tikz}
\usetikzlibrary{positioning}
\usepackage{titlesec}

\providecommand{\modecorrige}{0}
\providecommand{\variante}{A}
\newif\ifcorrige
\ifnum\modecorrige=1 \corrigetrue \else \corrigefalse \fi

\definecolor{ipred}{HTML}{C0392B}
\definecolor{cdcbar}{HTML}{A67C3D}
\definecolor{cdcbg}{HTML}{F4F1EA}
\definecolor{qbar}{HTML}{4A6B5D}
\definecolor{qbg}{HTML}{EDF1EF}
\definecolor{corr}{HTML}{1F5FA8}
\definecolor{corrbg}{HTML}{E8F0FA}
\definecolor{grisfoot}{HTML}{828282}

\newcommand{\ip}[1]{\textcolor{ipred}{\texttt{#1}}}
\newcommand{\pts}[1]{\hfill\mbox{\textbf{(#1)}}}
\newcommand{\sol}[1]{\ifcorrige\textcolor{corr}{#1}\fi}
\newcommand{\bareme}[1]{\ifcorrige\par{\footnotesize\color{corr}#1\par}\fi}
\newcommand{\hl}{\ifcorrige\rule[-1.6mm]{0pt}{5.4mm}\else\rule[-2.2mm]{0pt}{6.6mm}\fi}
\newcolumntype{C}{>{\centering\arraybackslash}X}
\newcolumntype{L}{>{\raggedright\arraybackslash}X}

\titleformat{\section}{\large\bfseries}{}{0pt}{}
\titlespacing*{\section}{0pt}{6pt}{3pt}

\tcbset{qstyle/.style={enhanced,boxrule=0pt,leftrule=2.5pt,arc=0pt,outer arc=0pt,
  left=5pt,right=5pt,top=2pt,bottom=2pt,before skip=4pt,after skip=3pt}}
\newtcolorbox{cdc}{qstyle,colback=cdcbg,colframe=cdcbar}
\newtcolorbox{quest}{qstyle,colback=qbg,colframe=qbar}

% Zone de réponse : vierge dans le sujet, remplie dans le corrigé
\newcommand{\zone}[2]{%
  \ifcorrige
    \begin{tcolorbox}[colback=corrbg,colframe=corr,boxrule=0.5pt,arc=1pt,left=4pt,right=4pt,top=1pt,bottom=1pt,
      before skip=2pt,after skip=2pt,fontupper=\footnotesize\color{corr}]#2\end{tcolorbox}%
  \else
    \begin{tcolorbox}[colback=white,colframe=black!45,boxrule=0.4pt,arc=1pt,height=#1,
      before skip=2pt,after skip=2pt]\end{tcolorbox}%
  \fi}

\newcommand{\question}[3]{\begin{quest}\textbf{#1.} #3\pts{#2}\end{quest}}

%==========================================================================
% Données propres à chaque variante
%==========================================================================
\if\variante A
  \def\entreprise{Le cabinet d'architectes \textbf{ArchiNova}}
  \def\Rnom{R-ArchiNova}
  \def\resUn{Bureau d'études}\def\pUn{BE}\def\netUn{192.168.40}
  \def\resDeux{Accueil}\def\pDeux{ACC}\def\netDeux{172.16}
  \def\srvNom{SRV-Intranet}\def\netSrv{192.168.200}\def\srvH{10}
  \def\ssid{ArchiNova-WiFi}
  \def\AIrows{%
    \hl\ip{192.168.40.25} & \sol{C} & \sol{255.255.255.0} & \sol{192.168.40.0} & \sol{192.168.40.255} & \sol{254} \\ \hline
    \hl\ip{172.16.5.10}   & \sol{B} & \sol{255.255.0.0}   & \sol{172.16.0.0}   & \sol{172.16.255.255} & \sol{65\,534} \\ \hline
    \hl\ip{10.3.2.1}      & \sol{A} & \sol{255.0.0.0}     & \sol{10.0.0.0}     & \sol{10.255.255.255} & \sol{16\,777\,214} \\ \hline}
  \def\apipa{169.254.12.7}
  \def\Cnet{192.168.60}
  \def\Crows{%
    PC-A & \ip{192.168.60.10} & \ip{255.255.255.0} & \ip{192.168.60.254} & — \\
    PC-B & \ip{192.168.60.20} & \ip{255.255.255.0} & \ip{192.168.60.1}   & — \\
    PC-C & \ip{192.168.61.30} & \ip{255.255.255.0} & \ip{192.168.60.254} & — \\
    PC-D & \ip{192.168.60.10} & \ip{255.255.255.0} & \ip{192.168.60.254} & — \\
    SRV-C & \ip{10.0.0.5}     & \ip{255.0.0.0}     & \ip{10.0.0.254}     & — \\
    R1 — G0/0 & \ip{192.168.60.254} & \ip{255.255.255.0} & — & On \\
    R1 — G0/1 & \ip{10.0.0.254}     & \ip{255.0.0.0}     & — & Off \\}
  \def\Csol{%
    \hl 1 & \sol{PC-B} & \sol{Passerelle .60.1 : pas l'adresse du routeur.} & \sol{Joint son réseau, mais aucun autre (serveur).} & \sol{Passerelle 192.168.60.254.} \\ \hline
    \hl 2 & \sol{PC-C} & \sol{IP hors du réseau 192.168.60.0/24.} & \sol{Ne joint ni les postes ni le routeur.} & \sol{Ex. 192.168.60.30.} \\ \hline
    \hl 3 & \sol{PC-D} & \sol{Même IP que PC-A (doublon).} & \sol{Conflit d'adresse : PC-A et PC-D perturbés.} & \sol{Adresse unique, ex. .60.40.} \\ \hline
    \hl 4 & \sol{R1 G0/1} & \sol{Interface désactivée (Off).} & \sol{Réseau du serveur injoignable pour tous.} & \sol{Port Status : On.} \\ \hline}
  \def\Cbareme{PC-A et SRV-C sont correctement configurés ; accepter PC-A cité à la place de PC-D pour le doublon.}
\else
  \def\entreprise{La clinique vétérinaire \textbf{VetoPlus}}
  \def\Rnom{R-VetoPlus}
  \def\resUn{Soins}\def\pUn{SOI}\def\netUn{192.168.70}
  \def\resDeux{Secrétariat}\def\pDeux{SEC}\def\netDeux{172.20}
  \def\srvNom{SRV-Dossiers}\def\netSrv{192.168.250}\def\srvH{20}
  \def\ssid{VetoPlus-WiFi}
  \def\AIrows{%
    \hl\ip{192.168.15.130} & \sol{C} & \sol{255.255.255.0} & \sol{192.168.15.0} & \sol{192.168.15.255} & \sol{254} \\ \hline
    \hl\ip{172.31.200.7}   & \sol{B} & \sol{255.255.0.0}   & \sol{172.31.0.0}   & \sol{172.31.255.255} & \sol{65\,534} \\ \hline
    \hl\ip{10.45.6.9}      & \sol{A} & \sol{255.0.0.0}     & \sol{10.0.0.0}     & \sol{10.255.255.255} & \sol{16\,777\,214} \\ \hline}
  \def\apipa{169.254.88.3}
  \def\Cnet{192.168.80}
  \def\Crows{%
    PC-A & \ip{192.168.8.11}   & \ip{255.255.255.0} & \ip{192.168.80.254} & — \\
    PC-B & \ip{192.168.80.12}  & \ip{255.255.255.0} & \ip{192.168.80.254} & — \\
    PC-C & \ip{192.168.80.13}  & \ip{255.255.255.0} & \ip{192.168.80.100} & — \\
    PC-D & \ip{192.168.80.12}  & \ip{255.255.255.0} & \ip{192.168.80.254} & — \\
    SRV-C & \ip{10.0.0.8}      & \ip{255.0.0.0}     & \ip{10.0.0.254}     & — \\
    R1 — G0/0 & \ip{192.168.80.254} & \ip{255.255.255.0} & — & Off \\
    R1 — G0/1 & \ip{10.0.0.254}     & \ip{255.0.0.0}     & — & On \\}
  \def\Csol{%
    \hl 1 & \sol{PC-A} & \sol{IP hors du réseau 192.168.80.0/24.} & \sol{Ne joint ni les postes ni le routeur.} & \sol{Ex. 192.168.80.11.} \\ \hline
    \hl 2 & \sol{PC-C} & \sol{Passerelle .80.100 : pas l'adresse du routeur.} & \sol{Joint son réseau, mais aucun autre (serveur).} & \sol{Passerelle 192.168.80.254.} \\ \hline
    \hl 3 & \sol{PC-D} & \sol{Même IP que PC-B (doublon).} & \sol{Conflit d'adresse : PC-B et PC-D perturbés.} & \sol{Adresse unique, ex. .80.14.} \\ \hline
    \hl 4 & \sol{R1 G0/0} & \sol{Interface côté LAN désactivée (Off).} & \sol{Aucun poste ne joint le routeur ni le serveur.} & \sol{Port Status : On.} \\ \hline}
  \def\Cbareme{PC-B et SRV-C sont correctement configurés ; accepter PC-B cité à la place de PC-D pour le doublon.}
\fi

\pagestyle{fancy}
\fancyhf{}
\renewcommand{\headrulewidth}{0pt}
\fancyfoot[L]{\footnotesize\itshape\color{grisfoot}Évaluation — Réseaux locaux — Sujet \variante\ifcorrige\ — CORRIGÉ\fi}
\fancyfoot[R]{\footnotesize\color{grisfoot}\thepage/3}

\setlist{itemsep=0pt,topsep=2pt,parsep=0pt}
\setlength{\parindent}{0pt}
\setlength{\parskip}{2pt}

\begin{document}

\begin{center}
  {\Large\bfseries Évaluation pratique — Réseaux locaux \quad
   \fcolorbox{black}{white}{\ SUJET \variante\ }\par}
  \vspace{1pt}
  {\bfseries Adressage IP, commutation et routage — BTS CIEL 1re année — durée : 2 heures\par}
  \ifcorrige\vspace{2pt}{\bfseries\color{corr}CORRIGÉ ET BARÈME — document enseignant}\par\fi
\end{center}
\vspace{-2pt}
{\renewcommand{\arraystretch}{1.35}
\begin{tabularx}{\linewidth}{|X|X|p{2.2cm}|p{2.6cm}|}
  \hline
  \textbf{Nom :} & \textbf{Prénom :} & \textbf{Poste n°} & \textbf{Note :} \hfill /20 \\ \hline
\end{tabularx}}

\textbf{Consignes :} travail \textbf{individuel}, documents et Internet \textbf{non autorisés}. Ce sujet est le
\textbf{document-réponse}. Fichier Packet Tracer : \ip{Nom\_Prenom\_EvalLAN\_\variante.pkt}, à enregistrer
régulièrement dans le dossier de rendu. Les « visas » sont apposés par le professeur, que vous appelez à chaque test réussi.
Durées conseillées : A 30 min — B 1 h 10 — C 20 min.

%==========================================================================
\section{Partie A — Questions de cours et calculs \hfill /6}

\question{A1}{2 pts}{Complétez le tableau en considérant le \textbf{masque par défaut} de la classe de chaque adresse.}
{\small\renewcommand{\arraystretch}{1.15}
\begin{tabularx}{\linewidth}{|l|c|C|C|C|c|}
  \hline
  \textbf{Adresse IP} & \textbf{Classe} & \textbf{Masque par défaut} & \textbf{Adresse réseau} &
  \textbf{Adresse de diffusion} & \textbf{Nb d'hôtes} \\ \hline
  \AIrows
\end{tabularx}}
\bareme{0,5 pt par ligne juste + 0,5 pt si les trois nombres d'hôtes sont justes ($2^{n}-2$, $n$ = nombre de bits d'hôte).}

\question{A2}{1 pt}{Indiquez le type de câble (\textbf{droit} ou \textbf{croisé}) pour chaque liaison.}
{\small\renewcommand{\arraystretch}{1.15}
\begin{tabularx}{\linewidth}{|X|c||X|c|}
  \hline
  \hl PC $\leftrightarrow$ commutateur & \makebox[1.6cm]{\sol{droit}} & Commutateur $\leftrightarrow$ routeur & \makebox[1.6cm]{\sol{droit}} \\ \hline
  \hl PC $\leftrightarrow$ PC (réseau P2P) & \sol{croisé} & Concentrateur $\leftrightarrow$ concentrateur & \sol{croisé} \\ \hline
\end{tabularx}}
\bareme{0,25 pt par case. Équipements de types différents : câble droit ; de même type : câble croisé.}

\question{A3}{1 pt}{Comment un \textbf{concentrateur (hub)} et un \textbf{commutateur (switch)} traitent-ils une trame
reçue ? Lequel limite les collisions, et pourquoi ?}
\zone{2.9cm}{Le \textbf{hub} (couche 1) répète la trame sur \textbf{tous ses autres ports} : un seul domaine de collision
partagé. Le \textbf{switch} (couche 2) lit l'adresse MAC de destination et, grâce à sa \textbf{table MAC}, ne l'envoie
\textbf{que sur le port du destinataire} : un domaine de collision par port, c'est lui qui limite les collisions.
\emph{0,5 pt hub, 0,5 pt switch.}}

\question{A4}{1 pt}{Qu'est-ce que la \textbf{passerelle par défaut} d'un poste ? Quelle condition son adresse doit-elle respecter ?}
\zone{2.4cm}{Adresse IP de l'interface du \textbf{routeur} à laquelle le poste envoie les paquets destinés à un
\textbf{autre réseau}. Elle doit \textbf{appartenir au même réseau} que le poste. \emph{0,5 + 0,5 pt.}}

\question{A5}{0,5 pt}{Un portable configuré en DHCP affiche \ip{\apipa} / \ip{255.255.0.0}. Que s'est-il passé ?}
\zone{1.5cm}{Aucun \textbf{serveur DHCP} n'a répondu : le poste s'est attribué une adresse \textbf{APIPA} (169.254.0.0/16).}

\question{A6}{0,5 pt}{Quel protocole utilise \texttt{ping} ? Nommez les deux messages échangés et la couche OSI des adresses IP.}
\zone{1.5cm}{\textbf{ICMP} ; \textbf{Echo Request} / \textbf{Echo Reply} ; adresses IP en \textbf{couche 3 (Réseau)}.}

%==========================================================================
\newpage
\section{Partie B — Réalisation sous Packet Tracer \hfill /10}

\begin{cdc}
\entreprise{} vous confie son réseau : \textbf{trois réseaux} reliés par un \textbf{routeur unique Cisco 2911}
nommé \ip{\Rnom} (interfaces G0/0, G0/1, G0/2).
\begin{itemize}[leftmargin=12pt]
  \item \textbf{Réseau \resUn} : \textbf{4 postes} PC-\pUn1 à PC-\pUn4, commutateur \ip{SW-\pUn}, adressage statique,
        réseau \ip{\netUn.0/24}.
  \item \textbf{Réseau \resDeux} : \textbf{2 postes} PC-\pDeux1, PC-\pDeux2 et \textbf{1 imprimante} IMP-\pDeux,
        commutateur \ip{SW-\pDeux}, adressage statique, réseau de classe B \ip{\netDeux.0.0/16}.
  \item \textbf{Réseau Serveur} : serveur \ip{\srvNom} d'adresse \ip{\netSrv.\srvH/24}, relié \textbf{directement} au routeur.
  \item La passerelle de chaque réseau est sa \textbf{dernière adresse utilisable}. Tous les équipements communiquent
        entre eux et accèdent à la page web du serveur. Nommage imposé (« Display Name »).
\end{itemize}
\end{cdc}

\question{B1}{2 pts}{Complétez le plan d'adressage \textbf{avant} la réalisation.}
{\footnotesize\setlength{\tabcolsep}{3pt}\renewcommand{\arraystretch}{1.15}
\begin{tabularx}{\linewidth}{|l|C|C|C|C|C|C|}
  \hline
  \textbf{Réseau} & \textbf{Adresse réseau} & \textbf{Masque} & \textbf{1re @ utilisable} &
  \textbf{Dernière @ utilisable} & \textbf{Diffusion} & \textbf{Passerelle} \\ \hline
  \hl \resUn & \sol{\netUn.0} & \sol{255.255.255.0} & \sol{\netUn.1} & \sol{\netUn.254} & \sol{\netUn.255} & \sol{\netUn.254} \\ \hline
  \hl \resDeux & \sol{\netDeux.0.0} & \sol{255.255.0.0} & \sol{\netDeux.0.1} & \sol{\netDeux.255.254} & \sol{\netDeux.255.255} & \sol{\netDeux.255.254} \\ \hline
  \hl Serveur & \sol{\netSrv.0} & \sol{255.255.255.0} & \sol{\netSrv.1} & \sol{\netSrv.254} & \sol{\netSrv.255} & \sol{\netSrv.254} \\ \hline
\end{tabularx}}
\bareme{0,5 pt par ligne juste + 0,5 pt si la convention de passerelle est respectée partout.}

\question{B2}{2 pts}{Construisez la topologie (équipements, câblage, nommage) et relevez l'adressage configuré.}
{\footnotesize\setlength{\tabcolsep}{3pt}\renewcommand{\arraystretch}{1.1}
\begin{tabularx}{\linewidth}{|l|C|C|C||l|C|C|}
  \hline
  \textbf{Hôte} & \textbf{Adresse IP} & \textbf{Masque} & \textbf{Passerelle} & \textbf{Interface} & \textbf{Adresse IP} & \textbf{Masque} \\ \hline
  \hl PC-\pUn1 & \sol{\netUn.1} & \sol{255.255.255.0} & \sol{\netUn.254} & \Rnom{} G0/0 & \sol{\netUn.254} & \sol{255.255.255.0} \\ \hline
  \hl PC-\pDeux1 & \sol{\netDeux.0.1} & \sol{255.255.0.0} & \sol{\netDeux.255.254} & \Rnom{} G0/1 & \sol{\netDeux.255.254} & \sol{255.255.0.0} \\ \hline
  \hl IMP-\pDeux & \sol{\netDeux.0.10} & \sol{255.255.0.0} & \sol{\netDeux.255.254} & \Rnom{} G0/2 & \sol{\netSrv.254} & \sol{255.255.255.0} \\ \hline
  \hl \srvNom & \ip{\netSrv.\srvH} & \sol{255.255.255.0} & \sol{\netSrv.254} & \multicolumn{3}{c|}{\cellcolor{black!8}} \\ \hline
\end{tabularx}}
\bareme{Adresses d'hôtes libres (uniques, dans le bon réseau) ; affectation G0/x libre. Sur le .pkt : équipements 0,5 ;
câblage 1 (câbles droits vers les switchs et switch–routeur, \textbf{croisé} serveur–routeur, liens verts) ; nommage 0,5.}

\question{B3}{3 pts}{Configurez tous les équipements : IP, masques, passerelles des hôtes ; interfaces du routeur
(« Port Status : On »).}
\bareme{Sur le .pkt : interfaces routeur 1 pt ; IP/masques des hôtes 1 pt ; passerelles (imprimante et serveur compris)
1 pt. $-0,25$ par erreur, sans descendre sous 0 par item.}

\question{B4}{3 pts}{Réalisez les tests, notez la commande et le résultat, puis faites viser chaque test réussi.}
{\small\renewcommand{\arraystretch}{1.1}
\begin{tabularx}{\linewidth}{|c|l|l|C|c|c|}
  \hline
  \textbf{N°} & \textbf{Source} & \textbf{Destination} & \textbf{Commande / action} & \textbf{Résultat} & \textbf{Visa prof.} \\ \hline
  \hl 1 & PC-\pUn1 & PC-\pUn4 & \sol{ping \netUn.4} & \sol{réussi} & \hspace{1.4cm} \\ \hline
  \hl 2 & PC-\pDeux1 & IMP-\pDeux & \sol{ping \netDeux.0.10} & \sol{réussi} & \\ \hline
  \hl 3 & PC-\pUn2 & PC-\pDeux2 & \sol{ping \netDeux.0.2} & \sol{réussi} & \\ \hline
  \hl 4 & PC-\pDeux2 & \srvNom & \sol{ping \netSrv.\srvH} & \sol{réussi} & \\ \hline
  \hl 5 & PC-\pUn3 & page web du serveur & \sol{navigateur : {\NoAutoSpacing http://\netSrv.\srvH}} & \sol{page affichée} & \\ \hline
\end{tabularx}}
\bareme{0,5 pt par test visé + 0,5 pt pour B5. Le premier ping inter-réseaux peut échouer (résolution ARP) ; les suivants doivent réussir.}

\question{B5}{incl. B4}{Test n°3 : par quels équipements passe le paquet ? Pourquoi PC-\pUn2 l'envoie-t-il au routeur ?}
\zone{2.3cm}{PC-\pUn2 $\to$ SW-\pUn{} $\to$ \Rnom{} $\to$ SW-\pDeux{} $\to$ PC-\pDeux2. Grâce à son masque, PC-\pUn2 constate que
la destination (\netDeux.0.0/16) n'est \textbf{pas dans son réseau} (\netUn.0/24) : il envoie le paquet à sa
\textbf{passerelle par défaut}, qui le route.}

%==========================================================================
\newpage
\section{Partie C — Diagnostic d'un réseau en panne \hfill /4}

Un technicien constate qu'\textbf{aucun poste ne parvient à joindre le serveur SRV-C} et que certains postes
communiquent mal entre eux. Les câbles sont du bon type.

\begin{center}
\begin{minipage}{0.37\linewidth}\centering
\begin{tikzpicture}[font=\footnotesize,
  eq/.style={draw,fill=qbg,rounded corners=2pt,minimum width=9mm,minimum height=5mm,inner sep=2pt},
  every path/.style={thick}]
  \node[eq] (SW) {SW-1};
  \node[eq,right=9mm of SW] (R) {R1};
  \node[eq,right=7mm of R] (SRV) {SRV-C};
  \node[eq,above left=5mm and -3mm of SW] (A) {PC-A};
  \node[eq,above right=5mm and -3mm of SW] (B) {PC-B};
  \node[eq,below left=5mm and -3mm of SW] (C) {PC-C};
  \node[eq,below right=5mm and -3mm of SW] (D) {PC-D};
  \foreach \n in {A,B,C,D} \draw (\n) -- (SW);
  \draw (SW) -- node[above,font=\tiny]{G0/0} (R);
  \draw (R) -- node[above,font=\tiny]{G0/1} (SRV);
\end{tikzpicture}
\end{minipage}\hfill
\begin{minipage}{0.61\linewidth}\centering
{\small\renewcommand{\arraystretch}{1.05}\setlength{\tabcolsep}{4pt}
\begin{tabular}{lllll}
  \toprule
  \textbf{Équipement} & \textbf{Adresse IP} & \textbf{Masque} & \textbf{Passerelle} & \textbf{Port} \\
  \midrule
  \Crows
  \bottomrule
\end{tabular}}
\end{minipage}
\end{center}

\question{C1}{4 pts}{Relevez les \textbf{quatre anomalies}. Pour chacune : conséquence et correction.}
{\small\renewcommand{\arraystretch}{1.1}
\ifcorrige\def\hc{}\else\def\hc{\rule[-10mm]{0pt}{14mm}}\fi
\begin{tabularx}{\linewidth}{|c|>{\raggedright\arraybackslash}p{2cm}|L|L|L|}
  \hline
  \textbf{N°} & \textbf{Équipement} & \textbf{Anomalie} & \textbf{Conséquence} & \textbf{Correction} \\ \hline
  \ifcorrige\Csol\else
  \hc 1 & & & & \\ \hline
  \hc 2 & & & & \\ \hline
  \hc 3 & & & & \\ \hline
  \hc 4 & & & & \\ \hline
  \fi
\end{tabularx}}
\bareme{1 pt par anomalie (0,25 équipement, anomalie, conséquence, correction ; ordre indifférent). \Cbareme}

%==========================================================================
\section{Bonus — Extension sans fil \hfill +1}

\begin{cdc}Ajoutez au réseau \resUn{} un \textbf{AccessPoint-PT} relié à SW-\pUn{} (SSID \ip{\ssid}, \textbf{WPA2-PSK})
et un portable \ip{Laptop-\pUn} équipé d'un module \texttt{WPC300N}, en adressage \textbf{statique}.
Vérifiez qu'il joint \srvNom.\end{cdc}
{\small\renewcommand{\arraystretch}{1.1}
\begin{tabularx}{\linewidth}{|l|C|C|C|c|}
  \hline
  \textbf{Équipement} & \textbf{Adresse IP} & \textbf{Masque} & \textbf{Passerelle} & \textbf{Visa prof.} \\ \hline
  \hl Laptop-\pUn & \sol{\netUn.50} & \sol{255.255.255.0} & \sol{\netUn.254} & \hspace{1.4cm} \\ \hline
\end{tabularx}}
\bareme{Le point d'accès fait office de pont : le portable est dans \netUn.0/24. Bonus si associé en WPA2-PSK et ping
vers \netSrv.\srvH{} réussi. Note plafonnée à 20.}

\end{document}
"
%   pdflatex -jobname=Eval_Reseau_SujetB   "\def\variante{B}\def\modecorrige{0}% Évaluation pratique — Réseaux locaux (BTS CIEL 1re année) — sujets A et B
% Compilation (4 PDF) :
%   pdflatex -jobname=Eval_Reseau_SujetA   "\def\variante{A}\def\modecorrige{0}\input{Eval_TP_Reseau_LAN.tex}"
%   pdflatex -jobname=Eval_Reseau_SujetB   "\def\variante{B}\def\modecorrige{0}\input{Eval_TP_Reseau_LAN.tex}"
%   pdflatex -jobname=Eval_Reseau_CorrigeA "\def\variante{A}\def\modecorrige{1}\input{Eval_TP_Reseau_LAN.tex}"
%   pdflatex -jobname=Eval_Reseau_CorrigeB "\def\variante{B}\def\modecorrige{1}\input{Eval_TP_Reseau_LAN.tex}"
\documentclass[10pt,a4paper]{article}

\usepackage[T1]{fontenc}
\usepackage[utf8]{inputenc}
\usepackage{lmodern}
\usepackage[french]{babel}
\usepackage[a4paper,hmargin=1.5cm,top=1.3cm,bottom=1.6cm,footskip=0.7cm]{geometry}
\usepackage[table]{xcolor}
\usepackage{booktabs,tabularx,array}
\usepackage{enumitem}
\usepackage{fancyhdr}
\usepackage[most]{tcolorbox}
\usepackage{tikz}
\usetikzlibrary{positioning}
\usepackage{titlesec}

\providecommand{\modecorrige}{0}
\providecommand{\variante}{A}
\newif\ifcorrige
\ifnum\modecorrige=1 \corrigetrue \else \corrigefalse \fi

\definecolor{ipred}{HTML}{C0392B}
\definecolor{cdcbar}{HTML}{A67C3D}
\definecolor{cdcbg}{HTML}{F4F1EA}
\definecolor{qbar}{HTML}{4A6B5D}
\definecolor{qbg}{HTML}{EDF1EF}
\definecolor{corr}{HTML}{1F5FA8}
\definecolor{corrbg}{HTML}{E8F0FA}
\definecolor{grisfoot}{HTML}{828282}

\newcommand{\ip}[1]{\textcolor{ipred}{\texttt{#1}}}
\newcommand{\pts}[1]{\hfill\mbox{\textbf{(#1)}}}
\newcommand{\sol}[1]{\ifcorrige\textcolor{corr}{#1}\fi}
\newcommand{\bareme}[1]{\ifcorrige\par{\footnotesize\color{corr}#1\par}\fi}
\newcommand{\hl}{\ifcorrige\rule[-1.6mm]{0pt}{5.4mm}\else\rule[-2.2mm]{0pt}{6.6mm}\fi}
\newcolumntype{C}{>{\centering\arraybackslash}X}
\newcolumntype{L}{>{\raggedright\arraybackslash}X}

\titleformat{\section}{\large\bfseries}{}{0pt}{}
\titlespacing*{\section}{0pt}{6pt}{3pt}

\tcbset{qstyle/.style={enhanced,boxrule=0pt,leftrule=2.5pt,arc=0pt,outer arc=0pt,
  left=5pt,right=5pt,top=2pt,bottom=2pt,before skip=4pt,after skip=3pt}}
\newtcolorbox{cdc}{qstyle,colback=cdcbg,colframe=cdcbar}
\newtcolorbox{quest}{qstyle,colback=qbg,colframe=qbar}

% Zone de réponse : vierge dans le sujet, remplie dans le corrigé
\newcommand{\zone}[2]{%
  \ifcorrige
    \begin{tcolorbox}[colback=corrbg,colframe=corr,boxrule=0.5pt,arc=1pt,left=4pt,right=4pt,top=1pt,bottom=1pt,
      before skip=2pt,after skip=2pt,fontupper=\footnotesize\color{corr}]#2\end{tcolorbox}%
  \else
    \begin{tcolorbox}[colback=white,colframe=black!45,boxrule=0.4pt,arc=1pt,height=#1,
      before skip=2pt,after skip=2pt]\end{tcolorbox}%
  \fi}

\newcommand{\question}[3]{\begin{quest}\textbf{#1.} #3\pts{#2}\end{quest}}

%==========================================================================
% Données propres à chaque variante
%==========================================================================
\if\variante A
  \def\entreprise{Le cabinet d'architectes \textbf{ArchiNova}}
  \def\Rnom{R-ArchiNova}
  \def\resUn{Bureau d'études}\def\pUn{BE}\def\netUn{192.168.40}
  \def\resDeux{Accueil}\def\pDeux{ACC}\def\netDeux{172.16}
  \def\srvNom{SRV-Intranet}\def\netSrv{192.168.200}\def\srvH{10}
  \def\ssid{ArchiNova-WiFi}
  \def\AIrows{%
    \hl\ip{192.168.40.25} & \sol{C} & \sol{255.255.255.0} & \sol{192.168.40.0} & \sol{192.168.40.255} & \sol{254} \\ \hline
    \hl\ip{172.16.5.10}   & \sol{B} & \sol{255.255.0.0}   & \sol{172.16.0.0}   & \sol{172.16.255.255} & \sol{65\,534} \\ \hline
    \hl\ip{10.3.2.1}      & \sol{A} & \sol{255.0.0.0}     & \sol{10.0.0.0}     & \sol{10.255.255.255} & \sol{16\,777\,214} \\ \hline}
  \def\apipa{169.254.12.7}
  \def\Cnet{192.168.60}
  \def\Crows{%
    PC-A & \ip{192.168.60.10} & \ip{255.255.255.0} & \ip{192.168.60.254} & — \\
    PC-B & \ip{192.168.60.20} & \ip{255.255.255.0} & \ip{192.168.60.1}   & — \\
    PC-C & \ip{192.168.61.30} & \ip{255.255.255.0} & \ip{192.168.60.254} & — \\
    PC-D & \ip{192.168.60.10} & \ip{255.255.255.0} & \ip{192.168.60.254} & — \\
    SRV-C & \ip{10.0.0.5}     & \ip{255.0.0.0}     & \ip{10.0.0.254}     & — \\
    R1 — G0/0 & \ip{192.168.60.254} & \ip{255.255.255.0} & — & On \\
    R1 — G0/1 & \ip{10.0.0.254}     & \ip{255.0.0.0}     & — & Off \\}
  \def\Csol{%
    \hl 1 & \sol{PC-B} & \sol{Passerelle .60.1 : pas l'adresse du routeur.} & \sol{Joint son réseau, mais aucun autre (serveur).} & \sol{Passerelle 192.168.60.254.} \\ \hline
    \hl 2 & \sol{PC-C} & \sol{IP hors du réseau 192.168.60.0/24.} & \sol{Ne joint ni les postes ni le routeur.} & \sol{Ex. 192.168.60.30.} \\ \hline
    \hl 3 & \sol{PC-D} & \sol{Même IP que PC-A (doublon).} & \sol{Conflit d'adresse : PC-A et PC-D perturbés.} & \sol{Adresse unique, ex. .60.40.} \\ \hline
    \hl 4 & \sol{R1 G0/1} & \sol{Interface désactivée (Off).} & \sol{Réseau du serveur injoignable pour tous.} & \sol{Port Status : On.} \\ \hline}
  \def\Cbareme{PC-A et SRV-C sont correctement configurés ; accepter PC-A cité à la place de PC-D pour le doublon.}
\else
  \def\entreprise{La clinique vétérinaire \textbf{VetoPlus}}
  \def\Rnom{R-VetoPlus}
  \def\resUn{Soins}\def\pUn{SOI}\def\netUn{192.168.70}
  \def\resDeux{Secrétariat}\def\pDeux{SEC}\def\netDeux{172.20}
  \def\srvNom{SRV-Dossiers}\def\netSrv{192.168.250}\def\srvH{20}
  \def\ssid{VetoPlus-WiFi}
  \def\AIrows{%
    \hl\ip{192.168.15.130} & \sol{C} & \sol{255.255.255.0} & \sol{192.168.15.0} & \sol{192.168.15.255} & \sol{254} \\ \hline
    \hl\ip{172.31.200.7}   & \sol{B} & \sol{255.255.0.0}   & \sol{172.31.0.0}   & \sol{172.31.255.255} & \sol{65\,534} \\ \hline
    \hl\ip{10.45.6.9}      & \sol{A} & \sol{255.0.0.0}     & \sol{10.0.0.0}     & \sol{10.255.255.255} & \sol{16\,777\,214} \\ \hline}
  \def\apipa{169.254.88.3}
  \def\Cnet{192.168.80}
  \def\Crows{%
    PC-A & \ip{192.168.8.11}   & \ip{255.255.255.0} & \ip{192.168.80.254} & — \\
    PC-B & \ip{192.168.80.12}  & \ip{255.255.255.0} & \ip{192.168.80.254} & — \\
    PC-C & \ip{192.168.80.13}  & \ip{255.255.255.0} & \ip{192.168.80.100} & — \\
    PC-D & \ip{192.168.80.12}  & \ip{255.255.255.0} & \ip{192.168.80.254} & — \\
    SRV-C & \ip{10.0.0.8}      & \ip{255.0.0.0}     & \ip{10.0.0.254}     & — \\
    R1 — G0/0 & \ip{192.168.80.254} & \ip{255.255.255.0} & — & Off \\
    R1 — G0/1 & \ip{10.0.0.254}     & \ip{255.0.0.0}     & — & On \\}
  \def\Csol{%
    \hl 1 & \sol{PC-A} & \sol{IP hors du réseau 192.168.80.0/24.} & \sol{Ne joint ni les postes ni le routeur.} & \sol{Ex. 192.168.80.11.} \\ \hline
    \hl 2 & \sol{PC-C} & \sol{Passerelle .80.100 : pas l'adresse du routeur.} & \sol{Joint son réseau, mais aucun autre (serveur).} & \sol{Passerelle 192.168.80.254.} \\ \hline
    \hl 3 & \sol{PC-D} & \sol{Même IP que PC-B (doublon).} & \sol{Conflit d'adresse : PC-B et PC-D perturbés.} & \sol{Adresse unique, ex. .80.14.} \\ \hline
    \hl 4 & \sol{R1 G0/0} & \sol{Interface côté LAN désactivée (Off).} & \sol{Aucun poste ne joint le routeur ni le serveur.} & \sol{Port Status : On.} \\ \hline}
  \def\Cbareme{PC-B et SRV-C sont correctement configurés ; accepter PC-B cité à la place de PC-D pour le doublon.}
\fi

\pagestyle{fancy}
\fancyhf{}
\renewcommand{\headrulewidth}{0pt}
\fancyfoot[L]{\footnotesize\itshape\color{grisfoot}Évaluation — Réseaux locaux — Sujet \variante\ifcorrige\ — CORRIGÉ\fi}
\fancyfoot[R]{\footnotesize\color{grisfoot}\thepage/3}

\setlist{itemsep=0pt,topsep=2pt,parsep=0pt}
\setlength{\parindent}{0pt}
\setlength{\parskip}{2pt}

\begin{document}

\begin{center}
  {\Large\bfseries Évaluation pratique — Réseaux locaux \quad
   \fcolorbox{black}{white}{\ SUJET \variante\ }\par}
  \vspace{1pt}
  {\bfseries Adressage IP, commutation et routage — BTS CIEL 1re année — durée : 2 heures\par}
  \ifcorrige\vspace{2pt}{\bfseries\color{corr}CORRIGÉ ET BARÈME — document enseignant}\par\fi
\end{center}
\vspace{-2pt}
{\renewcommand{\arraystretch}{1.35}
\begin{tabularx}{\linewidth}{|X|X|p{2.2cm}|p{2.6cm}|}
  \hline
  \textbf{Nom :} & \textbf{Prénom :} & \textbf{Poste n°} & \textbf{Note :} \hfill /20 \\ \hline
\end{tabularx}}

\textbf{Consignes :} travail \textbf{individuel}, documents et Internet \textbf{non autorisés}. Ce sujet est le
\textbf{document-réponse}. Fichier Packet Tracer : \ip{Nom\_Prenom\_EvalLAN\_\variante.pkt}, à enregistrer
régulièrement dans le dossier de rendu. Les « visas » sont apposés par le professeur, que vous appelez à chaque test réussi.
Durées conseillées : A 30 min — B 1 h 10 — C 20 min.

%==========================================================================
\section{Partie A — Questions de cours et calculs \hfill /6}

\question{A1}{2 pts}{Complétez le tableau en considérant le \textbf{masque par défaut} de la classe de chaque adresse.}
{\small\renewcommand{\arraystretch}{1.15}
\begin{tabularx}{\linewidth}{|l|c|C|C|C|c|}
  \hline
  \textbf{Adresse IP} & \textbf{Classe} & \textbf{Masque par défaut} & \textbf{Adresse réseau} &
  \textbf{Adresse de diffusion} & \textbf{Nb d'hôtes} \\ \hline
  \AIrows
\end{tabularx}}
\bareme{0,5 pt par ligne juste + 0,5 pt si les trois nombres d'hôtes sont justes ($2^{n}-2$, $n$ = nombre de bits d'hôte).}

\question{A2}{1 pt}{Indiquez le type de câble (\textbf{droit} ou \textbf{croisé}) pour chaque liaison.}
{\small\renewcommand{\arraystretch}{1.15}
\begin{tabularx}{\linewidth}{|X|c||X|c|}
  \hline
  \hl PC $\leftrightarrow$ commutateur & \makebox[1.6cm]{\sol{droit}} & Commutateur $\leftrightarrow$ routeur & \makebox[1.6cm]{\sol{droit}} \\ \hline
  \hl PC $\leftrightarrow$ PC (réseau P2P) & \sol{croisé} & Concentrateur $\leftrightarrow$ concentrateur & \sol{croisé} \\ \hline
\end{tabularx}}
\bareme{0,25 pt par case. Équipements de types différents : câble droit ; de même type : câble croisé.}

\question{A3}{1 pt}{Comment un \textbf{concentrateur (hub)} et un \textbf{commutateur (switch)} traitent-ils une trame
reçue ? Lequel limite les collisions, et pourquoi ?}
\zone{2.9cm}{Le \textbf{hub} (couche 1) répète la trame sur \textbf{tous ses autres ports} : un seul domaine de collision
partagé. Le \textbf{switch} (couche 2) lit l'adresse MAC de destination et, grâce à sa \textbf{table MAC}, ne l'envoie
\textbf{que sur le port du destinataire} : un domaine de collision par port, c'est lui qui limite les collisions.
\emph{0,5 pt hub, 0,5 pt switch.}}

\question{A4}{1 pt}{Qu'est-ce que la \textbf{passerelle par défaut} d'un poste ? Quelle condition son adresse doit-elle respecter ?}
\zone{2.4cm}{Adresse IP de l'interface du \textbf{routeur} à laquelle le poste envoie les paquets destinés à un
\textbf{autre réseau}. Elle doit \textbf{appartenir au même réseau} que le poste. \emph{0,5 + 0,5 pt.}}

\question{A5}{0,5 pt}{Un portable configuré en DHCP affiche \ip{\apipa} / \ip{255.255.0.0}. Que s'est-il passé ?}
\zone{1.5cm}{Aucun \textbf{serveur DHCP} n'a répondu : le poste s'est attribué une adresse \textbf{APIPA} (169.254.0.0/16).}

\question{A6}{0,5 pt}{Quel protocole utilise \texttt{ping} ? Nommez les deux messages échangés et la couche OSI des adresses IP.}
\zone{1.5cm}{\textbf{ICMP} ; \textbf{Echo Request} / \textbf{Echo Reply} ; adresses IP en \textbf{couche 3 (Réseau)}.}

%==========================================================================
\newpage
\section{Partie B — Réalisation sous Packet Tracer \hfill /10}

\begin{cdc}
\entreprise{} vous confie son réseau : \textbf{trois réseaux} reliés par un \textbf{routeur unique Cisco 2911}
nommé \ip{\Rnom} (interfaces G0/0, G0/1, G0/2).
\begin{itemize}[leftmargin=12pt]
  \item \textbf{Réseau \resUn} : \textbf{4 postes} PC-\pUn1 à PC-\pUn4, commutateur \ip{SW-\pUn}, adressage statique,
        réseau \ip{\netUn.0/24}.
  \item \textbf{Réseau \resDeux} : \textbf{2 postes} PC-\pDeux1, PC-\pDeux2 et \textbf{1 imprimante} IMP-\pDeux,
        commutateur \ip{SW-\pDeux}, adressage statique, réseau de classe B \ip{\netDeux.0.0/16}.
  \item \textbf{Réseau Serveur} : serveur \ip{\srvNom} d'adresse \ip{\netSrv.\srvH/24}, relié \textbf{directement} au routeur.
  \item La passerelle de chaque réseau est sa \textbf{dernière adresse utilisable}. Tous les équipements communiquent
        entre eux et accèdent à la page web du serveur. Nommage imposé (« Display Name »).
\end{itemize}
\end{cdc}

\question{B1}{2 pts}{Complétez le plan d'adressage \textbf{avant} la réalisation.}
{\footnotesize\setlength{\tabcolsep}{3pt}\renewcommand{\arraystretch}{1.15}
\begin{tabularx}{\linewidth}{|l|C|C|C|C|C|C|}
  \hline
  \textbf{Réseau} & \textbf{Adresse réseau} & \textbf{Masque} & \textbf{1re @ utilisable} &
  \textbf{Dernière @ utilisable} & \textbf{Diffusion} & \textbf{Passerelle} \\ \hline
  \hl \resUn & \sol{\netUn.0} & \sol{255.255.255.0} & \sol{\netUn.1} & \sol{\netUn.254} & \sol{\netUn.255} & \sol{\netUn.254} \\ \hline
  \hl \resDeux & \sol{\netDeux.0.0} & \sol{255.255.0.0} & \sol{\netDeux.0.1} & \sol{\netDeux.255.254} & \sol{\netDeux.255.255} & \sol{\netDeux.255.254} \\ \hline
  \hl Serveur & \sol{\netSrv.0} & \sol{255.255.255.0} & \sol{\netSrv.1} & \sol{\netSrv.254} & \sol{\netSrv.255} & \sol{\netSrv.254} \\ \hline
\end{tabularx}}
\bareme{0,5 pt par ligne juste + 0,5 pt si la convention de passerelle est respectée partout.}

\question{B2}{2 pts}{Construisez la topologie (équipements, câblage, nommage) et relevez l'adressage configuré.}
{\footnotesize\setlength{\tabcolsep}{3pt}\renewcommand{\arraystretch}{1.1}
\begin{tabularx}{\linewidth}{|l|C|C|C||l|C|C|}
  \hline
  \textbf{Hôte} & \textbf{Adresse IP} & \textbf{Masque} & \textbf{Passerelle} & \textbf{Interface} & \textbf{Adresse IP} & \textbf{Masque} \\ \hline
  \hl PC-\pUn1 & \sol{\netUn.1} & \sol{255.255.255.0} & \sol{\netUn.254} & \Rnom{} G0/0 & \sol{\netUn.254} & \sol{255.255.255.0} \\ \hline
  \hl PC-\pDeux1 & \sol{\netDeux.0.1} & \sol{255.255.0.0} & \sol{\netDeux.255.254} & \Rnom{} G0/1 & \sol{\netDeux.255.254} & \sol{255.255.0.0} \\ \hline
  \hl IMP-\pDeux & \sol{\netDeux.0.10} & \sol{255.255.0.0} & \sol{\netDeux.255.254} & \Rnom{} G0/2 & \sol{\netSrv.254} & \sol{255.255.255.0} \\ \hline
  \hl \srvNom & \ip{\netSrv.\srvH} & \sol{255.255.255.0} & \sol{\netSrv.254} & \multicolumn{3}{c|}{\cellcolor{black!8}} \\ \hline
\end{tabularx}}
\bareme{Adresses d'hôtes libres (uniques, dans le bon réseau) ; affectation G0/x libre. Sur le .pkt : équipements 0,5 ;
câblage 1 (câbles droits vers les switchs et switch–routeur, \textbf{croisé} serveur–routeur, liens verts) ; nommage 0,5.}

\question{B3}{3 pts}{Configurez tous les équipements : IP, masques, passerelles des hôtes ; interfaces du routeur
(« Port Status : On »).}
\bareme{Sur le .pkt : interfaces routeur 1 pt ; IP/masques des hôtes 1 pt ; passerelles (imprimante et serveur compris)
1 pt. $-0,25$ par erreur, sans descendre sous 0 par item.}

\question{B4}{3 pts}{Réalisez les tests, notez la commande et le résultat, puis faites viser chaque test réussi.}
{\small\renewcommand{\arraystretch}{1.1}
\begin{tabularx}{\linewidth}{|c|l|l|C|c|c|}
  \hline
  \textbf{N°} & \textbf{Source} & \textbf{Destination} & \textbf{Commande / action} & \textbf{Résultat} & \textbf{Visa prof.} \\ \hline
  \hl 1 & PC-\pUn1 & PC-\pUn4 & \sol{ping \netUn.4} & \sol{réussi} & \hspace{1.4cm} \\ \hline
  \hl 2 & PC-\pDeux1 & IMP-\pDeux & \sol{ping \netDeux.0.10} & \sol{réussi} & \\ \hline
  \hl 3 & PC-\pUn2 & PC-\pDeux2 & \sol{ping \netDeux.0.2} & \sol{réussi} & \\ \hline
  \hl 4 & PC-\pDeux2 & \srvNom & \sol{ping \netSrv.\srvH} & \sol{réussi} & \\ \hline
  \hl 5 & PC-\pUn3 & page web du serveur & \sol{navigateur : {\NoAutoSpacing http://\netSrv.\srvH}} & \sol{page affichée} & \\ \hline
\end{tabularx}}
\bareme{0,5 pt par test visé + 0,5 pt pour B5. Le premier ping inter-réseaux peut échouer (résolution ARP) ; les suivants doivent réussir.}

\question{B5}{incl. B4}{Test n°3 : par quels équipements passe le paquet ? Pourquoi PC-\pUn2 l'envoie-t-il au routeur ?}
\zone{2.3cm}{PC-\pUn2 $\to$ SW-\pUn{} $\to$ \Rnom{} $\to$ SW-\pDeux{} $\to$ PC-\pDeux2. Grâce à son masque, PC-\pUn2 constate que
la destination (\netDeux.0.0/16) n'est \textbf{pas dans son réseau} (\netUn.0/24) : il envoie le paquet à sa
\textbf{passerelle par défaut}, qui le route.}

%==========================================================================
\newpage
\section{Partie C — Diagnostic d'un réseau en panne \hfill /4}

Un technicien constate qu'\textbf{aucun poste ne parvient à joindre le serveur SRV-C} et que certains postes
communiquent mal entre eux. Les câbles sont du bon type.

\begin{center}
\begin{minipage}{0.37\linewidth}\centering
\begin{tikzpicture}[font=\footnotesize,
  eq/.style={draw,fill=qbg,rounded corners=2pt,minimum width=9mm,minimum height=5mm,inner sep=2pt},
  every path/.style={thick}]
  \node[eq] (SW) {SW-1};
  \node[eq,right=9mm of SW] (R) {R1};
  \node[eq,right=7mm of R] (SRV) {SRV-C};
  \node[eq,above left=5mm and -3mm of SW] (A) {PC-A};
  \node[eq,above right=5mm and -3mm of SW] (B) {PC-B};
  \node[eq,below left=5mm and -3mm of SW] (C) {PC-C};
  \node[eq,below right=5mm and -3mm of SW] (D) {PC-D};
  \foreach \n in {A,B,C,D} \draw (\n) -- (SW);
  \draw (SW) -- node[above,font=\tiny]{G0/0} (R);
  \draw (R) -- node[above,font=\tiny]{G0/1} (SRV);
\end{tikzpicture}
\end{minipage}\hfill
\begin{minipage}{0.61\linewidth}\centering
{\small\renewcommand{\arraystretch}{1.05}\setlength{\tabcolsep}{4pt}
\begin{tabular}{lllll}
  \toprule
  \textbf{Équipement} & \textbf{Adresse IP} & \textbf{Masque} & \textbf{Passerelle} & \textbf{Port} \\
  \midrule
  \Crows
  \bottomrule
\end{tabular}}
\end{minipage}
\end{center}

\question{C1}{4 pts}{Relevez les \textbf{quatre anomalies}. Pour chacune : conséquence et correction.}
{\small\renewcommand{\arraystretch}{1.1}
\ifcorrige\def\hc{}\else\def\hc{\rule[-10mm]{0pt}{14mm}}\fi
\begin{tabularx}{\linewidth}{|c|>{\raggedright\arraybackslash}p{2cm}|L|L|L|}
  \hline
  \textbf{N°} & \textbf{Équipement} & \textbf{Anomalie} & \textbf{Conséquence} & \textbf{Correction} \\ \hline
  \ifcorrige\Csol\else
  \hc 1 & & & & \\ \hline
  \hc 2 & & & & \\ \hline
  \hc 3 & & & & \\ \hline
  \hc 4 & & & & \\ \hline
  \fi
\end{tabularx}}
\bareme{1 pt par anomalie (0,25 équipement, anomalie, conséquence, correction ; ordre indifférent). \Cbareme}

%==========================================================================
\section{Bonus — Extension sans fil \hfill +1}

\begin{cdc}Ajoutez au réseau \resUn{} un \textbf{AccessPoint-PT} relié à SW-\pUn{} (SSID \ip{\ssid}, \textbf{WPA2-PSK})
et un portable \ip{Laptop-\pUn} équipé d'un module \texttt{WPC300N}, en adressage \textbf{statique}.
Vérifiez qu'il joint \srvNom.\end{cdc}
{\small\renewcommand{\arraystretch}{1.1}
\begin{tabularx}{\linewidth}{|l|C|C|C|c|}
  \hline
  \textbf{Équipement} & \textbf{Adresse IP} & \textbf{Masque} & \textbf{Passerelle} & \textbf{Visa prof.} \\ \hline
  \hl Laptop-\pUn & \sol{\netUn.50} & \sol{255.255.255.0} & \sol{\netUn.254} & \hspace{1.4cm} \\ \hline
\end{tabularx}}
\bareme{Le point d'accès fait office de pont : le portable est dans \netUn.0/24. Bonus si associé en WPA2-PSK et ping
vers \netSrv.\srvH{} réussi. Note plafonnée à 20.}

\end{document}
"
%   pdflatex -jobname=Eval_Reseau_CorrigeA "\def\variante{A}\def\modecorrige{1}% Évaluation pratique — Réseaux locaux (BTS CIEL 1re année) — sujets A et B
% Compilation (4 PDF) :
%   pdflatex -jobname=Eval_Reseau_SujetA   "\def\variante{A}\def\modecorrige{0}\input{Eval_TP_Reseau_LAN.tex}"
%   pdflatex -jobname=Eval_Reseau_SujetB   "\def\variante{B}\def\modecorrige{0}\input{Eval_TP_Reseau_LAN.tex}"
%   pdflatex -jobname=Eval_Reseau_CorrigeA "\def\variante{A}\def\modecorrige{1}\input{Eval_TP_Reseau_LAN.tex}"
%   pdflatex -jobname=Eval_Reseau_CorrigeB "\def\variante{B}\def\modecorrige{1}\input{Eval_TP_Reseau_LAN.tex}"
\documentclass[10pt,a4paper]{article}

\usepackage[T1]{fontenc}
\usepackage[utf8]{inputenc}
\usepackage{lmodern}
\usepackage[french]{babel}
\usepackage[a4paper,hmargin=1.5cm,top=1.3cm,bottom=1.6cm,footskip=0.7cm]{geometry}
\usepackage[table]{xcolor}
\usepackage{booktabs,tabularx,array}
\usepackage{enumitem}
\usepackage{fancyhdr}
\usepackage[most]{tcolorbox}
\usepackage{tikz}
\usetikzlibrary{positioning}
\usepackage{titlesec}

\providecommand{\modecorrige}{0}
\providecommand{\variante}{A}
\newif\ifcorrige
\ifnum\modecorrige=1 \corrigetrue \else \corrigefalse \fi

\definecolor{ipred}{HTML}{C0392B}
\definecolor{cdcbar}{HTML}{A67C3D}
\definecolor{cdcbg}{HTML}{F4F1EA}
\definecolor{qbar}{HTML}{4A6B5D}
\definecolor{qbg}{HTML}{EDF1EF}
\definecolor{corr}{HTML}{1F5FA8}
\definecolor{corrbg}{HTML}{E8F0FA}
\definecolor{grisfoot}{HTML}{828282}

\newcommand{\ip}[1]{\textcolor{ipred}{\texttt{#1}}}
\newcommand{\pts}[1]{\hfill\mbox{\textbf{(#1)}}}
\newcommand{\sol}[1]{\ifcorrige\textcolor{corr}{#1}\fi}
\newcommand{\bareme}[1]{\ifcorrige\par{\footnotesize\color{corr}#1\par}\fi}
\newcommand{\hl}{\ifcorrige\rule[-1.6mm]{0pt}{5.4mm}\else\rule[-2.2mm]{0pt}{6.6mm}\fi}
\newcolumntype{C}{>{\centering\arraybackslash}X}
\newcolumntype{L}{>{\raggedright\arraybackslash}X}

\titleformat{\section}{\large\bfseries}{}{0pt}{}
\titlespacing*{\section}{0pt}{6pt}{3pt}

\tcbset{qstyle/.style={enhanced,boxrule=0pt,leftrule=2.5pt,arc=0pt,outer arc=0pt,
  left=5pt,right=5pt,top=2pt,bottom=2pt,before skip=4pt,after skip=3pt}}
\newtcolorbox{cdc}{qstyle,colback=cdcbg,colframe=cdcbar}
\newtcolorbox{quest}{qstyle,colback=qbg,colframe=qbar}

% Zone de réponse : vierge dans le sujet, remplie dans le corrigé
\newcommand{\zone}[2]{%
  \ifcorrige
    \begin{tcolorbox}[colback=corrbg,colframe=corr,boxrule=0.5pt,arc=1pt,left=4pt,right=4pt,top=1pt,bottom=1pt,
      before skip=2pt,after skip=2pt,fontupper=\footnotesize\color{corr}]#2\end{tcolorbox}%
  \else
    \begin{tcolorbox}[colback=white,colframe=black!45,boxrule=0.4pt,arc=1pt,height=#1,
      before skip=2pt,after skip=2pt]\end{tcolorbox}%
  \fi}

\newcommand{\question}[3]{\begin{quest}\textbf{#1.} #3\pts{#2}\end{quest}}

%==========================================================================
% Données propres à chaque variante
%==========================================================================
\if\variante A
  \def\entreprise{Le cabinet d'architectes \textbf{ArchiNova}}
  \def\Rnom{R-ArchiNova}
  \def\resUn{Bureau d'études}\def\pUn{BE}\def\netUn{192.168.40}
  \def\resDeux{Accueil}\def\pDeux{ACC}\def\netDeux{172.16}
  \def\srvNom{SRV-Intranet}\def\netSrv{192.168.200}\def\srvH{10}
  \def\ssid{ArchiNova-WiFi}
  \def\AIrows{%
    \hl\ip{192.168.40.25} & \sol{C} & \sol{255.255.255.0} & \sol{192.168.40.0} & \sol{192.168.40.255} & \sol{254} \\ \hline
    \hl\ip{172.16.5.10}   & \sol{B} & \sol{255.255.0.0}   & \sol{172.16.0.0}   & \sol{172.16.255.255} & \sol{65\,534} \\ \hline
    \hl\ip{10.3.2.1}      & \sol{A} & \sol{255.0.0.0}     & \sol{10.0.0.0}     & \sol{10.255.255.255} & \sol{16\,777\,214} \\ \hline}
  \def\apipa{169.254.12.7}
  \def\Cnet{192.168.60}
  \def\Crows{%
    PC-A & \ip{192.168.60.10} & \ip{255.255.255.0} & \ip{192.168.60.254} & — \\
    PC-B & \ip{192.168.60.20} & \ip{255.255.255.0} & \ip{192.168.60.1}   & — \\
    PC-C & \ip{192.168.61.30} & \ip{255.255.255.0} & \ip{192.168.60.254} & — \\
    PC-D & \ip{192.168.60.10} & \ip{255.255.255.0} & \ip{192.168.60.254} & — \\
    SRV-C & \ip{10.0.0.5}     & \ip{255.0.0.0}     & \ip{10.0.0.254}     & — \\
    R1 — G0/0 & \ip{192.168.60.254} & \ip{255.255.255.0} & — & On \\
    R1 — G0/1 & \ip{10.0.0.254}     & \ip{255.0.0.0}     & — & Off \\}
  \def\Csol{%
    \hl 1 & \sol{PC-B} & \sol{Passerelle .60.1 : pas l'adresse du routeur.} & \sol{Joint son réseau, mais aucun autre (serveur).} & \sol{Passerelle 192.168.60.254.} \\ \hline
    \hl 2 & \sol{PC-C} & \sol{IP hors du réseau 192.168.60.0/24.} & \sol{Ne joint ni les postes ni le routeur.} & \sol{Ex. 192.168.60.30.} \\ \hline
    \hl 3 & \sol{PC-D} & \sol{Même IP que PC-A (doublon).} & \sol{Conflit d'adresse : PC-A et PC-D perturbés.} & \sol{Adresse unique, ex. .60.40.} \\ \hline
    \hl 4 & \sol{R1 G0/1} & \sol{Interface désactivée (Off).} & \sol{Réseau du serveur injoignable pour tous.} & \sol{Port Status : On.} \\ \hline}
  \def\Cbareme{PC-A et SRV-C sont correctement configurés ; accepter PC-A cité à la place de PC-D pour le doublon.}
\else
  \def\entreprise{La clinique vétérinaire \textbf{VetoPlus}}
  \def\Rnom{R-VetoPlus}
  \def\resUn{Soins}\def\pUn{SOI}\def\netUn{192.168.70}
  \def\resDeux{Secrétariat}\def\pDeux{SEC}\def\netDeux{172.20}
  \def\srvNom{SRV-Dossiers}\def\netSrv{192.168.250}\def\srvH{20}
  \def\ssid{VetoPlus-WiFi}
  \def\AIrows{%
    \hl\ip{192.168.15.130} & \sol{C} & \sol{255.255.255.0} & \sol{192.168.15.0} & \sol{192.168.15.255} & \sol{254} \\ \hline
    \hl\ip{172.31.200.7}   & \sol{B} & \sol{255.255.0.0}   & \sol{172.31.0.0}   & \sol{172.31.255.255} & \sol{65\,534} \\ \hline
    \hl\ip{10.45.6.9}      & \sol{A} & \sol{255.0.0.0}     & \sol{10.0.0.0}     & \sol{10.255.255.255} & \sol{16\,777\,214} \\ \hline}
  \def\apipa{169.254.88.3}
  \def\Cnet{192.168.80}
  \def\Crows{%
    PC-A & \ip{192.168.8.11}   & \ip{255.255.255.0} & \ip{192.168.80.254} & — \\
    PC-B & \ip{192.168.80.12}  & \ip{255.255.255.0} & \ip{192.168.80.254} & — \\
    PC-C & \ip{192.168.80.13}  & \ip{255.255.255.0} & \ip{192.168.80.100} & — \\
    PC-D & \ip{192.168.80.12}  & \ip{255.255.255.0} & \ip{192.168.80.254} & — \\
    SRV-C & \ip{10.0.0.8}      & \ip{255.0.0.0}     & \ip{10.0.0.254}     & — \\
    R1 — G0/0 & \ip{192.168.80.254} & \ip{255.255.255.0} & — & Off \\
    R1 — G0/1 & \ip{10.0.0.254}     & \ip{255.0.0.0}     & — & On \\}
  \def\Csol{%
    \hl 1 & \sol{PC-A} & \sol{IP hors du réseau 192.168.80.0/24.} & \sol{Ne joint ni les postes ni le routeur.} & \sol{Ex. 192.168.80.11.} \\ \hline
    \hl 2 & \sol{PC-C} & \sol{Passerelle .80.100 : pas l'adresse du routeur.} & \sol{Joint son réseau, mais aucun autre (serveur).} & \sol{Passerelle 192.168.80.254.} \\ \hline
    \hl 3 & \sol{PC-D} & \sol{Même IP que PC-B (doublon).} & \sol{Conflit d'adresse : PC-B et PC-D perturbés.} & \sol{Adresse unique, ex. .80.14.} \\ \hline
    \hl 4 & \sol{R1 G0/0} & \sol{Interface côté LAN désactivée (Off).} & \sol{Aucun poste ne joint le routeur ni le serveur.} & \sol{Port Status : On.} \\ \hline}
  \def\Cbareme{PC-B et SRV-C sont correctement configurés ; accepter PC-B cité à la place de PC-D pour le doublon.}
\fi

\pagestyle{fancy}
\fancyhf{}
\renewcommand{\headrulewidth}{0pt}
\fancyfoot[L]{\footnotesize\itshape\color{grisfoot}Évaluation — Réseaux locaux — Sujet \variante\ifcorrige\ — CORRIGÉ\fi}
\fancyfoot[R]{\footnotesize\color{grisfoot}\thepage/3}

\setlist{itemsep=0pt,topsep=2pt,parsep=0pt}
\setlength{\parindent}{0pt}
\setlength{\parskip}{2pt}

\begin{document}

\begin{center}
  {\Large\bfseries Évaluation pratique — Réseaux locaux \quad
   \fcolorbox{black}{white}{\ SUJET \variante\ }\par}
  \vspace{1pt}
  {\bfseries Adressage IP, commutation et routage — BTS CIEL 1re année — durée : 2 heures\par}
  \ifcorrige\vspace{2pt}{\bfseries\color{corr}CORRIGÉ ET BARÈME — document enseignant}\par\fi
\end{center}
\vspace{-2pt}
{\renewcommand{\arraystretch}{1.35}
\begin{tabularx}{\linewidth}{|X|X|p{2.2cm}|p{2.6cm}|}
  \hline
  \textbf{Nom :} & \textbf{Prénom :} & \textbf{Poste n°} & \textbf{Note :} \hfill /20 \\ \hline
\end{tabularx}}

\textbf{Consignes :} travail \textbf{individuel}, documents et Internet \textbf{non autorisés}. Ce sujet est le
\textbf{document-réponse}. Fichier Packet Tracer : \ip{Nom\_Prenom\_EvalLAN\_\variante.pkt}, à enregistrer
régulièrement dans le dossier de rendu. Les « visas » sont apposés par le professeur, que vous appelez à chaque test réussi.
Durées conseillées : A 30 min — B 1 h 10 — C 20 min.

%==========================================================================
\section{Partie A — Questions de cours et calculs \hfill /6}

\question{A1}{2 pts}{Complétez le tableau en considérant le \textbf{masque par défaut} de la classe de chaque adresse.}
{\small\renewcommand{\arraystretch}{1.15}
\begin{tabularx}{\linewidth}{|l|c|C|C|C|c|}
  \hline
  \textbf{Adresse IP} & \textbf{Classe} & \textbf{Masque par défaut} & \textbf{Adresse réseau} &
  \textbf{Adresse de diffusion} & \textbf{Nb d'hôtes} \\ \hline
  \AIrows
\end{tabularx}}
\bareme{0,5 pt par ligne juste + 0,5 pt si les trois nombres d'hôtes sont justes ($2^{n}-2$, $n$ = nombre de bits d'hôte).}

\question{A2}{1 pt}{Indiquez le type de câble (\textbf{droit} ou \textbf{croisé}) pour chaque liaison.}
{\small\renewcommand{\arraystretch}{1.15}
\begin{tabularx}{\linewidth}{|X|c||X|c|}
  \hline
  \hl PC $\leftrightarrow$ commutateur & \makebox[1.6cm]{\sol{droit}} & Commutateur $\leftrightarrow$ routeur & \makebox[1.6cm]{\sol{droit}} \\ \hline
  \hl PC $\leftrightarrow$ PC (réseau P2P) & \sol{croisé} & Concentrateur $\leftrightarrow$ concentrateur & \sol{croisé} \\ \hline
\end{tabularx}}
\bareme{0,25 pt par case. Équipements de types différents : câble droit ; de même type : câble croisé.}

\question{A3}{1 pt}{Comment un \textbf{concentrateur (hub)} et un \textbf{commutateur (switch)} traitent-ils une trame
reçue ? Lequel limite les collisions, et pourquoi ?}
\zone{2.9cm}{Le \textbf{hub} (couche 1) répète la trame sur \textbf{tous ses autres ports} : un seul domaine de collision
partagé. Le \textbf{switch} (couche 2) lit l'adresse MAC de destination et, grâce à sa \textbf{table MAC}, ne l'envoie
\textbf{que sur le port du destinataire} : un domaine de collision par port, c'est lui qui limite les collisions.
\emph{0,5 pt hub, 0,5 pt switch.}}

\question{A4}{1 pt}{Qu'est-ce que la \textbf{passerelle par défaut} d'un poste ? Quelle condition son adresse doit-elle respecter ?}
\zone{2.4cm}{Adresse IP de l'interface du \textbf{routeur} à laquelle le poste envoie les paquets destinés à un
\textbf{autre réseau}. Elle doit \textbf{appartenir au même réseau} que le poste. \emph{0,5 + 0,5 pt.}}

\question{A5}{0,5 pt}{Un portable configuré en DHCP affiche \ip{\apipa} / \ip{255.255.0.0}. Que s'est-il passé ?}
\zone{1.5cm}{Aucun \textbf{serveur DHCP} n'a répondu : le poste s'est attribué une adresse \textbf{APIPA} (169.254.0.0/16).}

\question{A6}{0,5 pt}{Quel protocole utilise \texttt{ping} ? Nommez les deux messages échangés et la couche OSI des adresses IP.}
\zone{1.5cm}{\textbf{ICMP} ; \textbf{Echo Request} / \textbf{Echo Reply} ; adresses IP en \textbf{couche 3 (Réseau)}.}

%==========================================================================
\newpage
\section{Partie B — Réalisation sous Packet Tracer \hfill /10}

\begin{cdc}
\entreprise{} vous confie son réseau : \textbf{trois réseaux} reliés par un \textbf{routeur unique Cisco 2911}
nommé \ip{\Rnom} (interfaces G0/0, G0/1, G0/2).
\begin{itemize}[leftmargin=12pt]
  \item \textbf{Réseau \resUn} : \textbf{4 postes} PC-\pUn1 à PC-\pUn4, commutateur \ip{SW-\pUn}, adressage statique,
        réseau \ip{\netUn.0/24}.
  \item \textbf{Réseau \resDeux} : \textbf{2 postes} PC-\pDeux1, PC-\pDeux2 et \textbf{1 imprimante} IMP-\pDeux,
        commutateur \ip{SW-\pDeux}, adressage statique, réseau de classe B \ip{\netDeux.0.0/16}.
  \item \textbf{Réseau Serveur} : serveur \ip{\srvNom} d'adresse \ip{\netSrv.\srvH/24}, relié \textbf{directement} au routeur.
  \item La passerelle de chaque réseau est sa \textbf{dernière adresse utilisable}. Tous les équipements communiquent
        entre eux et accèdent à la page web du serveur. Nommage imposé (« Display Name »).
\end{itemize}
\end{cdc}

\question{B1}{2 pts}{Complétez le plan d'adressage \textbf{avant} la réalisation.}
{\footnotesize\setlength{\tabcolsep}{3pt}\renewcommand{\arraystretch}{1.15}
\begin{tabularx}{\linewidth}{|l|C|C|C|C|C|C|}
  \hline
  \textbf{Réseau} & \textbf{Adresse réseau} & \textbf{Masque} & \textbf{1re @ utilisable} &
  \textbf{Dernière @ utilisable} & \textbf{Diffusion} & \textbf{Passerelle} \\ \hline
  \hl \resUn & \sol{\netUn.0} & \sol{255.255.255.0} & \sol{\netUn.1} & \sol{\netUn.254} & \sol{\netUn.255} & \sol{\netUn.254} \\ \hline
  \hl \resDeux & \sol{\netDeux.0.0} & \sol{255.255.0.0} & \sol{\netDeux.0.1} & \sol{\netDeux.255.254} & \sol{\netDeux.255.255} & \sol{\netDeux.255.254} \\ \hline
  \hl Serveur & \sol{\netSrv.0} & \sol{255.255.255.0} & \sol{\netSrv.1} & \sol{\netSrv.254} & \sol{\netSrv.255} & \sol{\netSrv.254} \\ \hline
\end{tabularx}}
\bareme{0,5 pt par ligne juste + 0,5 pt si la convention de passerelle est respectée partout.}

\question{B2}{2 pts}{Construisez la topologie (équipements, câblage, nommage) et relevez l'adressage configuré.}
{\footnotesize\setlength{\tabcolsep}{3pt}\renewcommand{\arraystretch}{1.1}
\begin{tabularx}{\linewidth}{|l|C|C|C||l|C|C|}
  \hline
  \textbf{Hôte} & \textbf{Adresse IP} & \textbf{Masque} & \textbf{Passerelle} & \textbf{Interface} & \textbf{Adresse IP} & \textbf{Masque} \\ \hline
  \hl PC-\pUn1 & \sol{\netUn.1} & \sol{255.255.255.0} & \sol{\netUn.254} & \Rnom{} G0/0 & \sol{\netUn.254} & \sol{255.255.255.0} \\ \hline
  \hl PC-\pDeux1 & \sol{\netDeux.0.1} & \sol{255.255.0.0} & \sol{\netDeux.255.254} & \Rnom{} G0/1 & \sol{\netDeux.255.254} & \sol{255.255.0.0} \\ \hline
  \hl IMP-\pDeux & \sol{\netDeux.0.10} & \sol{255.255.0.0} & \sol{\netDeux.255.254} & \Rnom{} G0/2 & \sol{\netSrv.254} & \sol{255.255.255.0} \\ \hline
  \hl \srvNom & \ip{\netSrv.\srvH} & \sol{255.255.255.0} & \sol{\netSrv.254} & \multicolumn{3}{c|}{\cellcolor{black!8}} \\ \hline
\end{tabularx}}
\bareme{Adresses d'hôtes libres (uniques, dans le bon réseau) ; affectation G0/x libre. Sur le .pkt : équipements 0,5 ;
câblage 1 (câbles droits vers les switchs et switch–routeur, \textbf{croisé} serveur–routeur, liens verts) ; nommage 0,5.}

\question{B3}{3 pts}{Configurez tous les équipements : IP, masques, passerelles des hôtes ; interfaces du routeur
(« Port Status : On »).}
\bareme{Sur le .pkt : interfaces routeur 1 pt ; IP/masques des hôtes 1 pt ; passerelles (imprimante et serveur compris)
1 pt. $-0,25$ par erreur, sans descendre sous 0 par item.}

\question{B4}{3 pts}{Réalisez les tests, notez la commande et le résultat, puis faites viser chaque test réussi.}
{\small\renewcommand{\arraystretch}{1.1}
\begin{tabularx}{\linewidth}{|c|l|l|C|c|c|}
  \hline
  \textbf{N°} & \textbf{Source} & \textbf{Destination} & \textbf{Commande / action} & \textbf{Résultat} & \textbf{Visa prof.} \\ \hline
  \hl 1 & PC-\pUn1 & PC-\pUn4 & \sol{ping \netUn.4} & \sol{réussi} & \hspace{1.4cm} \\ \hline
  \hl 2 & PC-\pDeux1 & IMP-\pDeux & \sol{ping \netDeux.0.10} & \sol{réussi} & \\ \hline
  \hl 3 & PC-\pUn2 & PC-\pDeux2 & \sol{ping \netDeux.0.2} & \sol{réussi} & \\ \hline
  \hl 4 & PC-\pDeux2 & \srvNom & \sol{ping \netSrv.\srvH} & \sol{réussi} & \\ \hline
  \hl 5 & PC-\pUn3 & page web du serveur & \sol{navigateur : {\NoAutoSpacing http://\netSrv.\srvH}} & \sol{page affichée} & \\ \hline
\end{tabularx}}
\bareme{0,5 pt par test visé + 0,5 pt pour B5. Le premier ping inter-réseaux peut échouer (résolution ARP) ; les suivants doivent réussir.}

\question{B5}{incl. B4}{Test n°3 : par quels équipements passe le paquet ? Pourquoi PC-\pUn2 l'envoie-t-il au routeur ?}
\zone{2.3cm}{PC-\pUn2 $\to$ SW-\pUn{} $\to$ \Rnom{} $\to$ SW-\pDeux{} $\to$ PC-\pDeux2. Grâce à son masque, PC-\pUn2 constate que
la destination (\netDeux.0.0/16) n'est \textbf{pas dans son réseau} (\netUn.0/24) : il envoie le paquet à sa
\textbf{passerelle par défaut}, qui le route.}

%==========================================================================
\newpage
\section{Partie C — Diagnostic d'un réseau en panne \hfill /4}

Un technicien constate qu'\textbf{aucun poste ne parvient à joindre le serveur SRV-C} et que certains postes
communiquent mal entre eux. Les câbles sont du bon type.

\begin{center}
\begin{minipage}{0.37\linewidth}\centering
\begin{tikzpicture}[font=\footnotesize,
  eq/.style={draw,fill=qbg,rounded corners=2pt,minimum width=9mm,minimum height=5mm,inner sep=2pt},
  every path/.style={thick}]
  \node[eq] (SW) {SW-1};
  \node[eq,right=9mm of SW] (R) {R1};
  \node[eq,right=7mm of R] (SRV) {SRV-C};
  \node[eq,above left=5mm and -3mm of SW] (A) {PC-A};
  \node[eq,above right=5mm and -3mm of SW] (B) {PC-B};
  \node[eq,below left=5mm and -3mm of SW] (C) {PC-C};
  \node[eq,below right=5mm and -3mm of SW] (D) {PC-D};
  \foreach \n in {A,B,C,D} \draw (\n) -- (SW);
  \draw (SW) -- node[above,font=\tiny]{G0/0} (R);
  \draw (R) -- node[above,font=\tiny]{G0/1} (SRV);
\end{tikzpicture}
\end{minipage}\hfill
\begin{minipage}{0.61\linewidth}\centering
{\small\renewcommand{\arraystretch}{1.05}\setlength{\tabcolsep}{4pt}
\begin{tabular}{lllll}
  \toprule
  \textbf{Équipement} & \textbf{Adresse IP} & \textbf{Masque} & \textbf{Passerelle} & \textbf{Port} \\
  \midrule
  \Crows
  \bottomrule
\end{tabular}}
\end{minipage}
\end{center}

\question{C1}{4 pts}{Relevez les \textbf{quatre anomalies}. Pour chacune : conséquence et correction.}
{\small\renewcommand{\arraystretch}{1.1}
\ifcorrige\def\hc{}\else\def\hc{\rule[-10mm]{0pt}{14mm}}\fi
\begin{tabularx}{\linewidth}{|c|>{\raggedright\arraybackslash}p{2cm}|L|L|L|}
  \hline
  \textbf{N°} & \textbf{Équipement} & \textbf{Anomalie} & \textbf{Conséquence} & \textbf{Correction} \\ \hline
  \ifcorrige\Csol\else
  \hc 1 & & & & \\ \hline
  \hc 2 & & & & \\ \hline
  \hc 3 & & & & \\ \hline
  \hc 4 & & & & \\ \hline
  \fi
\end{tabularx}}
\bareme{1 pt par anomalie (0,25 équipement, anomalie, conséquence, correction ; ordre indifférent). \Cbareme}

%==========================================================================
\section{Bonus — Extension sans fil \hfill +1}

\begin{cdc}Ajoutez au réseau \resUn{} un \textbf{AccessPoint-PT} relié à SW-\pUn{} (SSID \ip{\ssid}, \textbf{WPA2-PSK})
et un portable \ip{Laptop-\pUn} équipé d'un module \texttt{WPC300N}, en adressage \textbf{statique}.
Vérifiez qu'il joint \srvNom.\end{cdc}
{\small\renewcommand{\arraystretch}{1.1}
\begin{tabularx}{\linewidth}{|l|C|C|C|c|}
  \hline
  \textbf{Équipement} & \textbf{Adresse IP} & \textbf{Masque} & \textbf{Passerelle} & \textbf{Visa prof.} \\ \hline
  \hl Laptop-\pUn & \sol{\netUn.50} & \sol{255.255.255.0} & \sol{\netUn.254} & \hspace{1.4cm} \\ \hline
\end{tabularx}}
\bareme{Le point d'accès fait office de pont : le portable est dans \netUn.0/24. Bonus si associé en WPA2-PSK et ping
vers \netSrv.\srvH{} réussi. Note plafonnée à 20.}

\end{document}
"
%   pdflatex -jobname=Eval_Reseau_CorrigeB "\def\variante{B}\def\modecorrige{1}% Évaluation pratique — Réseaux locaux (BTS CIEL 1re année) — sujets A et B
% Compilation (4 PDF) :
%   pdflatex -jobname=Eval_Reseau_SujetA   "\def\variante{A}\def\modecorrige{0}\input{Eval_TP_Reseau_LAN.tex}"
%   pdflatex -jobname=Eval_Reseau_SujetB   "\def\variante{B}\def\modecorrige{0}\input{Eval_TP_Reseau_LAN.tex}"
%   pdflatex -jobname=Eval_Reseau_CorrigeA "\def\variante{A}\def\modecorrige{1}\input{Eval_TP_Reseau_LAN.tex}"
%   pdflatex -jobname=Eval_Reseau_CorrigeB "\def\variante{B}\def\modecorrige{1}\input{Eval_TP_Reseau_LAN.tex}"
\documentclass[10pt,a4paper]{article}

\usepackage[T1]{fontenc}
\usepackage[utf8]{inputenc}
\usepackage{lmodern}
\usepackage[french]{babel}
\usepackage[a4paper,hmargin=1.5cm,top=1.3cm,bottom=1.6cm,footskip=0.7cm]{geometry}
\usepackage[table]{xcolor}
\usepackage{booktabs,tabularx,array}
\usepackage{enumitem}
\usepackage{fancyhdr}
\usepackage[most]{tcolorbox}
\usepackage{tikz}
\usetikzlibrary{positioning}
\usepackage{titlesec}

\providecommand{\modecorrige}{0}
\providecommand{\variante}{A}
\newif\ifcorrige
\ifnum\modecorrige=1 \corrigetrue \else \corrigefalse \fi

\definecolor{ipred}{HTML}{C0392B}
\definecolor{cdcbar}{HTML}{A67C3D}
\definecolor{cdcbg}{HTML}{F4F1EA}
\definecolor{qbar}{HTML}{4A6B5D}
\definecolor{qbg}{HTML}{EDF1EF}
\definecolor{corr}{HTML}{1F5FA8}
\definecolor{corrbg}{HTML}{E8F0FA}
\definecolor{grisfoot}{HTML}{828282}

\newcommand{\ip}[1]{\textcolor{ipred}{\texttt{#1}}}
\newcommand{\pts}[1]{\hfill\mbox{\textbf{(#1)}}}
\newcommand{\sol}[1]{\ifcorrige\textcolor{corr}{#1}\fi}
\newcommand{\bareme}[1]{\ifcorrige\par{\footnotesize\color{corr}#1\par}\fi}
\newcommand{\hl}{\ifcorrige\rule[-1.6mm]{0pt}{5.4mm}\else\rule[-2.2mm]{0pt}{6.6mm}\fi}
\newcolumntype{C}{>{\centering\arraybackslash}X}
\newcolumntype{L}{>{\raggedright\arraybackslash}X}

\titleformat{\section}{\large\bfseries}{}{0pt}{}
\titlespacing*{\section}{0pt}{6pt}{3pt}

\tcbset{qstyle/.style={enhanced,boxrule=0pt,leftrule=2.5pt,arc=0pt,outer arc=0pt,
  left=5pt,right=5pt,top=2pt,bottom=2pt,before skip=4pt,after skip=3pt}}
\newtcolorbox{cdc}{qstyle,colback=cdcbg,colframe=cdcbar}
\newtcolorbox{quest}{qstyle,colback=qbg,colframe=qbar}

% Zone de réponse : vierge dans le sujet, remplie dans le corrigé
\newcommand{\zone}[2]{%
  \ifcorrige
    \begin{tcolorbox}[colback=corrbg,colframe=corr,boxrule=0.5pt,arc=1pt,left=4pt,right=4pt,top=1pt,bottom=1pt,
      before skip=2pt,after skip=2pt,fontupper=\footnotesize\color{corr}]#2\end{tcolorbox}%
  \else
    \begin{tcolorbox}[colback=white,colframe=black!45,boxrule=0.4pt,arc=1pt,height=#1,
      before skip=2pt,after skip=2pt]\end{tcolorbox}%
  \fi}

\newcommand{\question}[3]{\begin{quest}\textbf{#1.} #3\pts{#2}\end{quest}}

%==========================================================================
% Données propres à chaque variante
%==========================================================================
\if\variante A
  \def\entreprise{Le cabinet d'architectes \textbf{ArchiNova}}
  \def\Rnom{R-ArchiNova}
  \def\resUn{Bureau d'études}\def\pUn{BE}\def\netUn{192.168.40}
  \def\resDeux{Accueil}\def\pDeux{ACC}\def\netDeux{172.16}
  \def\srvNom{SRV-Intranet}\def\netSrv{192.168.200}\def\srvH{10}
  \def\ssid{ArchiNova-WiFi}
  \def\AIrows{%
    \hl\ip{192.168.40.25} & \sol{C} & \sol{255.255.255.0} & \sol{192.168.40.0} & \sol{192.168.40.255} & \sol{254} \\ \hline
    \hl\ip{172.16.5.10}   & \sol{B} & \sol{255.255.0.0}   & \sol{172.16.0.0}   & \sol{172.16.255.255} & \sol{65\,534} \\ \hline
    \hl\ip{10.3.2.1}      & \sol{A} & \sol{255.0.0.0}     & \sol{10.0.0.0}     & \sol{10.255.255.255} & \sol{16\,777\,214} \\ \hline}
  \def\apipa{169.254.12.7}
  \def\Cnet{192.168.60}
  \def\Crows{%
    PC-A & \ip{192.168.60.10} & \ip{255.255.255.0} & \ip{192.168.60.254} & — \\
    PC-B & \ip{192.168.60.20} & \ip{255.255.255.0} & \ip{192.168.60.1}   & — \\
    PC-C & \ip{192.168.61.30} & \ip{255.255.255.0} & \ip{192.168.60.254} & — \\
    PC-D & \ip{192.168.60.10} & \ip{255.255.255.0} & \ip{192.168.60.254} & — \\
    SRV-C & \ip{10.0.0.5}     & \ip{255.0.0.0}     & \ip{10.0.0.254}     & — \\
    R1 — G0/0 & \ip{192.168.60.254} & \ip{255.255.255.0} & — & On \\
    R1 — G0/1 & \ip{10.0.0.254}     & \ip{255.0.0.0}     & — & Off \\}
  \def\Csol{%
    \hl 1 & \sol{PC-B} & \sol{Passerelle .60.1 : pas l'adresse du routeur.} & \sol{Joint son réseau, mais aucun autre (serveur).} & \sol{Passerelle 192.168.60.254.} \\ \hline
    \hl 2 & \sol{PC-C} & \sol{IP hors du réseau 192.168.60.0/24.} & \sol{Ne joint ni les postes ni le routeur.} & \sol{Ex. 192.168.60.30.} \\ \hline
    \hl 3 & \sol{PC-D} & \sol{Même IP que PC-A (doublon).} & \sol{Conflit d'adresse : PC-A et PC-D perturbés.} & \sol{Adresse unique, ex. .60.40.} \\ \hline
    \hl 4 & \sol{R1 G0/1} & \sol{Interface désactivée (Off).} & \sol{Réseau du serveur injoignable pour tous.} & \sol{Port Status : On.} \\ \hline}
  \def\Cbareme{PC-A et SRV-C sont correctement configurés ; accepter PC-A cité à la place de PC-D pour le doublon.}
\else
  \def\entreprise{La clinique vétérinaire \textbf{VetoPlus}}
  \def\Rnom{R-VetoPlus}
  \def\resUn{Soins}\def\pUn{SOI}\def\netUn{192.168.70}
  \def\resDeux{Secrétariat}\def\pDeux{SEC}\def\netDeux{172.20}
  \def\srvNom{SRV-Dossiers}\def\netSrv{192.168.250}\def\srvH{20}
  \def\ssid{VetoPlus-WiFi}
  \def\AIrows{%
    \hl\ip{192.168.15.130} & \sol{C} & \sol{255.255.255.0} & \sol{192.168.15.0} & \sol{192.168.15.255} & \sol{254} \\ \hline
    \hl\ip{172.31.200.7}   & \sol{B} & \sol{255.255.0.0}   & \sol{172.31.0.0}   & \sol{172.31.255.255} & \sol{65\,534} \\ \hline
    \hl\ip{10.45.6.9}      & \sol{A} & \sol{255.0.0.0}     & \sol{10.0.0.0}     & \sol{10.255.255.255} & \sol{16\,777\,214} \\ \hline}
  \def\apipa{169.254.88.3}
  \def\Cnet{192.168.80}
  \def\Crows{%
    PC-A & \ip{192.168.8.11}   & \ip{255.255.255.0} & \ip{192.168.80.254} & — \\
    PC-B & \ip{192.168.80.12}  & \ip{255.255.255.0} & \ip{192.168.80.254} & — \\
    PC-C & \ip{192.168.80.13}  & \ip{255.255.255.0} & \ip{192.168.80.100} & — \\
    PC-D & \ip{192.168.80.12}  & \ip{255.255.255.0} & \ip{192.168.80.254} & — \\
    SRV-C & \ip{10.0.0.8}      & \ip{255.0.0.0}     & \ip{10.0.0.254}     & — \\
    R1 — G0/0 & \ip{192.168.80.254} & \ip{255.255.255.0} & — & Off \\
    R1 — G0/1 & \ip{10.0.0.254}     & \ip{255.0.0.0}     & — & On \\}
  \def\Csol{%
    \hl 1 & \sol{PC-A} & \sol{IP hors du réseau 192.168.80.0/24.} & \sol{Ne joint ni les postes ni le routeur.} & \sol{Ex. 192.168.80.11.} \\ \hline
    \hl 2 & \sol{PC-C} & \sol{Passerelle .80.100 : pas l'adresse du routeur.} & \sol{Joint son réseau, mais aucun autre (serveur).} & \sol{Passerelle 192.168.80.254.} \\ \hline
    \hl 3 & \sol{PC-D} & \sol{Même IP que PC-B (doublon).} & \sol{Conflit d'adresse : PC-B et PC-D perturbés.} & \sol{Adresse unique, ex. .80.14.} \\ \hline
    \hl 4 & \sol{R1 G0/0} & \sol{Interface côté LAN désactivée (Off).} & \sol{Aucun poste ne joint le routeur ni le serveur.} & \sol{Port Status : On.} \\ \hline}
  \def\Cbareme{PC-B et SRV-C sont correctement configurés ; accepter PC-B cité à la place de PC-D pour le doublon.}
\fi

\pagestyle{fancy}
\fancyhf{}
\renewcommand{\headrulewidth}{0pt}
\fancyfoot[L]{\footnotesize\itshape\color{grisfoot}Évaluation — Réseaux locaux — Sujet \variante\ifcorrige\ — CORRIGÉ\fi}
\fancyfoot[R]{\footnotesize\color{grisfoot}\thepage/3}

\setlist{itemsep=0pt,topsep=2pt,parsep=0pt}
\setlength{\parindent}{0pt}
\setlength{\parskip}{2pt}

\begin{document}

\begin{center}
  {\Large\bfseries Évaluation pratique — Réseaux locaux \quad
   \fcolorbox{black}{white}{\ SUJET \variante\ }\par}
  \vspace{1pt}
  {\bfseries Adressage IP, commutation et routage — BTS CIEL 1re année — durée : 2 heures\par}
  \ifcorrige\vspace{2pt}{\bfseries\color{corr}CORRIGÉ ET BARÈME — document enseignant}\par\fi
\end{center}
\vspace{-2pt}
{\renewcommand{\arraystretch}{1.35}
\begin{tabularx}{\linewidth}{|X|X|p{2.2cm}|p{2.6cm}|}
  \hline
  \textbf{Nom :} & \textbf{Prénom :} & \textbf{Poste n°} & \textbf{Note :} \hfill /20 \\ \hline
\end{tabularx}}

\textbf{Consignes :} travail \textbf{individuel}, documents et Internet \textbf{non autorisés}. Ce sujet est le
\textbf{document-réponse}. Fichier Packet Tracer : \ip{Nom\_Prenom\_EvalLAN\_\variante.pkt}, à enregistrer
régulièrement dans le dossier de rendu. Les « visas » sont apposés par le professeur, que vous appelez à chaque test réussi.
Durées conseillées : A 30 min — B 1 h 10 — C 20 min.

%==========================================================================
\section{Partie A — Questions de cours et calculs \hfill /6}

\question{A1}{2 pts}{Complétez le tableau en considérant le \textbf{masque par défaut} de la classe de chaque adresse.}
{\small\renewcommand{\arraystretch}{1.15}
\begin{tabularx}{\linewidth}{|l|c|C|C|C|c|}
  \hline
  \textbf{Adresse IP} & \textbf{Classe} & \textbf{Masque par défaut} & \textbf{Adresse réseau} &
  \textbf{Adresse de diffusion} & \textbf{Nb d'hôtes} \\ \hline
  \AIrows
\end{tabularx}}
\bareme{0,5 pt par ligne juste + 0,5 pt si les trois nombres d'hôtes sont justes ($2^{n}-2$, $n$ = nombre de bits d'hôte).}

\question{A2}{1 pt}{Indiquez le type de câble (\textbf{droit} ou \textbf{croisé}) pour chaque liaison.}
{\small\renewcommand{\arraystretch}{1.15}
\begin{tabularx}{\linewidth}{|X|c||X|c|}
  \hline
  \hl PC $\leftrightarrow$ commutateur & \makebox[1.6cm]{\sol{droit}} & Commutateur $\leftrightarrow$ routeur & \makebox[1.6cm]{\sol{droit}} \\ \hline
  \hl PC $\leftrightarrow$ PC (réseau P2P) & \sol{croisé} & Concentrateur $\leftrightarrow$ concentrateur & \sol{croisé} \\ \hline
\end{tabularx}}
\bareme{0,25 pt par case. Équipements de types différents : câble droit ; de même type : câble croisé.}

\question{A3}{1 pt}{Comment un \textbf{concentrateur (hub)} et un \textbf{commutateur (switch)} traitent-ils une trame
reçue ? Lequel limite les collisions, et pourquoi ?}
\zone{2.9cm}{Le \textbf{hub} (couche 1) répète la trame sur \textbf{tous ses autres ports} : un seul domaine de collision
partagé. Le \textbf{switch} (couche 2) lit l'adresse MAC de destination et, grâce à sa \textbf{table MAC}, ne l'envoie
\textbf{que sur le port du destinataire} : un domaine de collision par port, c'est lui qui limite les collisions.
\emph{0,5 pt hub, 0,5 pt switch.}}

\question{A4}{1 pt}{Qu'est-ce que la \textbf{passerelle par défaut} d'un poste ? Quelle condition son adresse doit-elle respecter ?}
\zone{2.4cm}{Adresse IP de l'interface du \textbf{routeur} à laquelle le poste envoie les paquets destinés à un
\textbf{autre réseau}. Elle doit \textbf{appartenir au même réseau} que le poste. \emph{0,5 + 0,5 pt.}}

\question{A5}{0,5 pt}{Un portable configuré en DHCP affiche \ip{\apipa} / \ip{255.255.0.0}. Que s'est-il passé ?}
\zone{1.5cm}{Aucun \textbf{serveur DHCP} n'a répondu : le poste s'est attribué une adresse \textbf{APIPA} (169.254.0.0/16).}

\question{A6}{0,5 pt}{Quel protocole utilise \texttt{ping} ? Nommez les deux messages échangés et la couche OSI des adresses IP.}
\zone{1.5cm}{\textbf{ICMP} ; \textbf{Echo Request} / \textbf{Echo Reply} ; adresses IP en \textbf{couche 3 (Réseau)}.}

%==========================================================================
\newpage
\section{Partie B — Réalisation sous Packet Tracer \hfill /10}

\begin{cdc}
\entreprise{} vous confie son réseau : \textbf{trois réseaux} reliés par un \textbf{routeur unique Cisco 2911}
nommé \ip{\Rnom} (interfaces G0/0, G0/1, G0/2).
\begin{itemize}[leftmargin=12pt]
  \item \textbf{Réseau \resUn} : \textbf{4 postes} PC-\pUn1 à PC-\pUn4, commutateur \ip{SW-\pUn}, adressage statique,
        réseau \ip{\netUn.0/24}.
  \item \textbf{Réseau \resDeux} : \textbf{2 postes} PC-\pDeux1, PC-\pDeux2 et \textbf{1 imprimante} IMP-\pDeux,
        commutateur \ip{SW-\pDeux}, adressage statique, réseau de classe B \ip{\netDeux.0.0/16}.
  \item \textbf{Réseau Serveur} : serveur \ip{\srvNom} d'adresse \ip{\netSrv.\srvH/24}, relié \textbf{directement} au routeur.
  \item La passerelle de chaque réseau est sa \textbf{dernière adresse utilisable}. Tous les équipements communiquent
        entre eux et accèdent à la page web du serveur. Nommage imposé (« Display Name »).
\end{itemize}
\end{cdc}

\question{B1}{2 pts}{Complétez le plan d'adressage \textbf{avant} la réalisation.}
{\footnotesize\setlength{\tabcolsep}{3pt}\renewcommand{\arraystretch}{1.15}
\begin{tabularx}{\linewidth}{|l|C|C|C|C|C|C|}
  \hline
  \textbf{Réseau} & \textbf{Adresse réseau} & \textbf{Masque} & \textbf{1re @ utilisable} &
  \textbf{Dernière @ utilisable} & \textbf{Diffusion} & \textbf{Passerelle} \\ \hline
  \hl \resUn & \sol{\netUn.0} & \sol{255.255.255.0} & \sol{\netUn.1} & \sol{\netUn.254} & \sol{\netUn.255} & \sol{\netUn.254} \\ \hline
  \hl \resDeux & \sol{\netDeux.0.0} & \sol{255.255.0.0} & \sol{\netDeux.0.1} & \sol{\netDeux.255.254} & \sol{\netDeux.255.255} & \sol{\netDeux.255.254} \\ \hline
  \hl Serveur & \sol{\netSrv.0} & \sol{255.255.255.0} & \sol{\netSrv.1} & \sol{\netSrv.254} & \sol{\netSrv.255} & \sol{\netSrv.254} \\ \hline
\end{tabularx}}
\bareme{0,5 pt par ligne juste + 0,5 pt si la convention de passerelle est respectée partout.}

\question{B2}{2 pts}{Construisez la topologie (équipements, câblage, nommage) et relevez l'adressage configuré.}
{\footnotesize\setlength{\tabcolsep}{3pt}\renewcommand{\arraystretch}{1.1}
\begin{tabularx}{\linewidth}{|l|C|C|C||l|C|C|}
  \hline
  \textbf{Hôte} & \textbf{Adresse IP} & \textbf{Masque} & \textbf{Passerelle} & \textbf{Interface} & \textbf{Adresse IP} & \textbf{Masque} \\ \hline
  \hl PC-\pUn1 & \sol{\netUn.1} & \sol{255.255.255.0} & \sol{\netUn.254} & \Rnom{} G0/0 & \sol{\netUn.254} & \sol{255.255.255.0} \\ \hline
  \hl PC-\pDeux1 & \sol{\netDeux.0.1} & \sol{255.255.0.0} & \sol{\netDeux.255.254} & \Rnom{} G0/1 & \sol{\netDeux.255.254} & \sol{255.255.0.0} \\ \hline
  \hl IMP-\pDeux & \sol{\netDeux.0.10} & \sol{255.255.0.0} & \sol{\netDeux.255.254} & \Rnom{} G0/2 & \sol{\netSrv.254} & \sol{255.255.255.0} \\ \hline
  \hl \srvNom & \ip{\netSrv.\srvH} & \sol{255.255.255.0} & \sol{\netSrv.254} & \multicolumn{3}{c|}{\cellcolor{black!8}} \\ \hline
\end{tabularx}}
\bareme{Adresses d'hôtes libres (uniques, dans le bon réseau) ; affectation G0/x libre. Sur le .pkt : équipements 0,5 ;
câblage 1 (câbles droits vers les switchs et switch–routeur, \textbf{croisé} serveur–routeur, liens verts) ; nommage 0,5.}

\question{B3}{3 pts}{Configurez tous les équipements : IP, masques, passerelles des hôtes ; interfaces du routeur
(« Port Status : On »).}
\bareme{Sur le .pkt : interfaces routeur 1 pt ; IP/masques des hôtes 1 pt ; passerelles (imprimante et serveur compris)
1 pt. $-0,25$ par erreur, sans descendre sous 0 par item.}

\question{B4}{3 pts}{Réalisez les tests, notez la commande et le résultat, puis faites viser chaque test réussi.}
{\small\renewcommand{\arraystretch}{1.1}
\begin{tabularx}{\linewidth}{|c|l|l|C|c|c|}
  \hline
  \textbf{N°} & \textbf{Source} & \textbf{Destination} & \textbf{Commande / action} & \textbf{Résultat} & \textbf{Visa prof.} \\ \hline
  \hl 1 & PC-\pUn1 & PC-\pUn4 & \sol{ping \netUn.4} & \sol{réussi} & \hspace{1.4cm} \\ \hline
  \hl 2 & PC-\pDeux1 & IMP-\pDeux & \sol{ping \netDeux.0.10} & \sol{réussi} & \\ \hline
  \hl 3 & PC-\pUn2 & PC-\pDeux2 & \sol{ping \netDeux.0.2} & \sol{réussi} & \\ \hline
  \hl 4 & PC-\pDeux2 & \srvNom & \sol{ping \netSrv.\srvH} & \sol{réussi} & \\ \hline
  \hl 5 & PC-\pUn3 & page web du serveur & \sol{navigateur : {\NoAutoSpacing http://\netSrv.\srvH}} & \sol{page affichée} & \\ \hline
\end{tabularx}}
\bareme{0,5 pt par test visé + 0,5 pt pour B5. Le premier ping inter-réseaux peut échouer (résolution ARP) ; les suivants doivent réussir.}

\question{B5}{incl. B4}{Test n°3 : par quels équipements passe le paquet ? Pourquoi PC-\pUn2 l'envoie-t-il au routeur ?}
\zone{2.3cm}{PC-\pUn2 $\to$ SW-\pUn{} $\to$ \Rnom{} $\to$ SW-\pDeux{} $\to$ PC-\pDeux2. Grâce à son masque, PC-\pUn2 constate que
la destination (\netDeux.0.0/16) n'est \textbf{pas dans son réseau} (\netUn.0/24) : il envoie le paquet à sa
\textbf{passerelle par défaut}, qui le route.}

%==========================================================================
\newpage
\section{Partie C — Diagnostic d'un réseau en panne \hfill /4}

Un technicien constate qu'\textbf{aucun poste ne parvient à joindre le serveur SRV-C} et que certains postes
communiquent mal entre eux. Les câbles sont du bon type.

\begin{center}
\begin{minipage}{0.37\linewidth}\centering
\begin{tikzpicture}[font=\footnotesize,
  eq/.style={draw,fill=qbg,rounded corners=2pt,minimum width=9mm,minimum height=5mm,inner sep=2pt},
  every path/.style={thick}]
  \node[eq] (SW) {SW-1};
  \node[eq,right=9mm of SW] (R) {R1};
  \node[eq,right=7mm of R] (SRV) {SRV-C};
  \node[eq,above left=5mm and -3mm of SW] (A) {PC-A};
  \node[eq,above right=5mm and -3mm of SW] (B) {PC-B};
  \node[eq,below left=5mm and -3mm of SW] (C) {PC-C};
  \node[eq,below right=5mm and -3mm of SW] (D) {PC-D};
  \foreach \n in {A,B,C,D} \draw (\n) -- (SW);
  \draw (SW) -- node[above,font=\tiny]{G0/0} (R);
  \draw (R) -- node[above,font=\tiny]{G0/1} (SRV);
\end{tikzpicture}
\end{minipage}\hfill
\begin{minipage}{0.61\linewidth}\centering
{\small\renewcommand{\arraystretch}{1.05}\setlength{\tabcolsep}{4pt}
\begin{tabular}{lllll}
  \toprule
  \textbf{Équipement} & \textbf{Adresse IP} & \textbf{Masque} & \textbf{Passerelle} & \textbf{Port} \\
  \midrule
  \Crows
  \bottomrule
\end{tabular}}
\end{minipage}
\end{center}

\question{C1}{4 pts}{Relevez les \textbf{quatre anomalies}. Pour chacune : conséquence et correction.}
{\small\renewcommand{\arraystretch}{1.1}
\ifcorrige\def\hc{}\else\def\hc{\rule[-10mm]{0pt}{14mm}}\fi
\begin{tabularx}{\linewidth}{|c|>{\raggedright\arraybackslash}p{2cm}|L|L|L|}
  \hline
  \textbf{N°} & \textbf{Équipement} & \textbf{Anomalie} & \textbf{Conséquence} & \textbf{Correction} \\ \hline
  \ifcorrige\Csol\else
  \hc 1 & & & & \\ \hline
  \hc 2 & & & & \\ \hline
  \hc 3 & & & & \\ \hline
  \hc 4 & & & & \\ \hline
  \fi
\end{tabularx}}
\bareme{1 pt par anomalie (0,25 équipement, anomalie, conséquence, correction ; ordre indifférent). \Cbareme}

%==========================================================================
\section{Bonus — Extension sans fil \hfill +1}

\begin{cdc}Ajoutez au réseau \resUn{} un \textbf{AccessPoint-PT} relié à SW-\pUn{} (SSID \ip{\ssid}, \textbf{WPA2-PSK})
et un portable \ip{Laptop-\pUn} équipé d'un module \texttt{WPC300N}, en adressage \textbf{statique}.
Vérifiez qu'il joint \srvNom.\end{cdc}
{\small\renewcommand{\arraystretch}{1.1}
\begin{tabularx}{\linewidth}{|l|C|C|C|c|}
  \hline
  \textbf{Équipement} & \textbf{Adresse IP} & \textbf{Masque} & \textbf{Passerelle} & \textbf{Visa prof.} \\ \hline
  \hl Laptop-\pUn & \sol{\netUn.50} & \sol{255.255.255.0} & \sol{\netUn.254} & \hspace{1.4cm} \\ \hline
\end{tabularx}}
\bareme{Le point d'accès fait office de pont : le portable est dans \netUn.0/24. Bonus si associé en WPA2-PSK et ping
vers \netSrv.\srvH{} réussi. Note plafonnée à 20.}

\end{document}
"
\documentclass[10pt,a4paper]{article}

\usepackage[T1]{fontenc}
\usepackage[utf8]{inputenc}
\usepackage{lmodern}
\usepackage[french]{babel}
\usepackage[a4paper,hmargin=1.5cm,top=1.3cm,bottom=1.6cm,footskip=0.7cm]{geometry}
\usepackage[table]{xcolor}
\usepackage{booktabs,tabularx,array}
\usepackage{enumitem}
\usepackage{fancyhdr}
\usepackage[most]{tcolorbox}
\usepackage{tikz}
\usetikzlibrary{positioning}
\usepackage{titlesec}

\providecommand{\modecorrige}{0}
\providecommand{\variante}{A}
\newif\ifcorrige
\ifnum\modecorrige=1 \corrigetrue \else \corrigefalse \fi

\definecolor{ipred}{HTML}{C0392B}
\definecolor{cdcbar}{HTML}{A67C3D}
\definecolor{cdcbg}{HTML}{F4F1EA}
\definecolor{qbar}{HTML}{4A6B5D}
\definecolor{qbg}{HTML}{EDF1EF}
\definecolor{corr}{HTML}{1F5FA8}
\definecolor{corrbg}{HTML}{E8F0FA}
\definecolor{grisfoot}{HTML}{828282}

\newcommand{\ip}[1]{\textcolor{ipred}{\texttt{#1}}}
\newcommand{\pts}[1]{\hfill\mbox{\textbf{(#1)}}}
\newcommand{\sol}[1]{\ifcorrige\textcolor{corr}{#1}\fi}
\newcommand{\bareme}[1]{\ifcorrige\par{\footnotesize\color{corr}#1\par}\fi}
\newcommand{\hl}{\ifcorrige\rule[-1.6mm]{0pt}{5.4mm}\else\rule[-2.2mm]{0pt}{6.6mm}\fi}
\newcolumntype{C}{>{\centering\arraybackslash}X}
\newcolumntype{L}{>{\raggedright\arraybackslash}X}

\titleformat{\section}{\large\bfseries}{}{0pt}{}
\titlespacing*{\section}{0pt}{6pt}{3pt}

\tcbset{qstyle/.style={enhanced,boxrule=0pt,leftrule=2.5pt,arc=0pt,outer arc=0pt,
  left=5pt,right=5pt,top=2pt,bottom=2pt,before skip=4pt,after skip=3pt}}
\newtcolorbox{cdc}{qstyle,colback=cdcbg,colframe=cdcbar}
\newtcolorbox{quest}{qstyle,colback=qbg,colframe=qbar}

% Zone de réponse : vierge dans le sujet, remplie dans le corrigé
\newcommand{\zone}[2]{%
  \ifcorrige
    \begin{tcolorbox}[colback=corrbg,colframe=corr,boxrule=0.5pt,arc=1pt,left=4pt,right=4pt,top=1pt,bottom=1pt,
      before skip=2pt,after skip=2pt,fontupper=\footnotesize\color{corr}]#2\end{tcolorbox}%
  \else
    \begin{tcolorbox}[colback=white,colframe=black!45,boxrule=0.4pt,arc=1pt,height=#1,
      before skip=2pt,after skip=2pt]\end{tcolorbox}%
  \fi}

\newcommand{\question}[3]{\begin{quest}\textbf{#1.} #3\pts{#2}\end{quest}}

%==========================================================================
% Données propres à chaque variante
%==========================================================================
\if\variante A
  \def\entreprise{Le cabinet d'architectes \textbf{ArchiNova}}
  \def\Rnom{R-ArchiNova}
  \def\resUn{Bureau d'études}\def\pUn{BE}\def\netUn{192.168.40}
  \def\resDeux{Accueil}\def\pDeux{ACC}\def\netDeux{172.16}
  \def\srvNom{SRV-Intranet}\def\netSrv{192.168.200}\def\srvH{10}
  \def\ssid{ArchiNova-WiFi}
  \def\AIrows{%
    \hl\ip{192.168.40.25} & \sol{C} & \sol{255.255.255.0} & \sol{192.168.40.0} & \sol{192.168.40.255} & \sol{254} \\ \hline
    \hl\ip{172.16.5.10}   & \sol{B} & \sol{255.255.0.0}   & \sol{172.16.0.0}   & \sol{172.16.255.255} & \sol{65\,534} \\ \hline
    \hl\ip{10.3.2.1}      & \sol{A} & \sol{255.0.0.0}     & \sol{10.0.0.0}     & \sol{10.255.255.255} & \sol{16\,777\,214} \\ \hline}
  \def\apipa{169.254.12.7}
  \def\Cnet{192.168.60}
  \def\Crows{%
    PC-A & \ip{192.168.60.10} & \ip{255.255.255.0} & \ip{192.168.60.254} & — \\
    PC-B & \ip{192.168.60.20} & \ip{255.255.255.0} & \ip{192.168.60.1}   & — \\
    PC-C & \ip{192.168.61.30} & \ip{255.255.255.0} & \ip{192.168.60.254} & — \\
    PC-D & \ip{192.168.60.10} & \ip{255.255.255.0} & \ip{192.168.60.254} & — \\
    SRV-C & \ip{10.0.0.5}     & \ip{255.0.0.0}     & \ip{10.0.0.254}     & — \\
    R1 — G0/0 & \ip{192.168.60.254} & \ip{255.255.255.0} & — & On \\
    R1 — G0/1 & \ip{10.0.0.254}     & \ip{255.0.0.0}     & — & Off \\}
  \def\Csol{%
    \hl 1 & \sol{PC-B} & \sol{Passerelle .60.1 : pas l'adresse du routeur.} & \sol{Joint son réseau, mais aucun autre (serveur).} & \sol{Passerelle 192.168.60.254.} \\ \hline
    \hl 2 & \sol{PC-C} & \sol{IP hors du réseau 192.168.60.0/24.} & \sol{Ne joint ni les postes ni le routeur.} & \sol{Ex. 192.168.60.30.} \\ \hline
    \hl 3 & \sol{PC-D} & \sol{Même IP que PC-A (doublon).} & \sol{Conflit d'adresse : PC-A et PC-D perturbés.} & \sol{Adresse unique, ex. .60.40.} \\ \hline
    \hl 4 & \sol{R1 G0/1} & \sol{Interface désactivée (Off).} & \sol{Réseau du serveur injoignable pour tous.} & \sol{Port Status : On.} \\ \hline}
  \def\Cbareme{PC-A et SRV-C sont correctement configurés ; accepter PC-A cité à la place de PC-D pour le doublon.}
\else
  \def\entreprise{La clinique vétérinaire \textbf{VetoPlus}}
  \def\Rnom{R-VetoPlus}
  \def\resUn{Soins}\def\pUn{SOI}\def\netUn{192.168.70}
  \def\resDeux{Secrétariat}\def\pDeux{SEC}\def\netDeux{172.20}
  \def\srvNom{SRV-Dossiers}\def\netSrv{192.168.250}\def\srvH{20}
  \def\ssid{VetoPlus-WiFi}
  \def\AIrows{%
    \hl\ip{192.168.15.130} & \sol{C} & \sol{255.255.255.0} & \sol{192.168.15.0} & \sol{192.168.15.255} & \sol{254} \\ \hline
    \hl\ip{172.31.200.7}   & \sol{B} & \sol{255.255.0.0}   & \sol{172.31.0.0}   & \sol{172.31.255.255} & \sol{65\,534} \\ \hline
    \hl\ip{10.45.6.9}      & \sol{A} & \sol{255.0.0.0}     & \sol{10.0.0.0}     & \sol{10.255.255.255} & \sol{16\,777\,214} \\ \hline}
  \def\apipa{169.254.88.3}
  \def\Cnet{192.168.80}
  \def\Crows{%
    PC-A & \ip{192.168.8.11}   & \ip{255.255.255.0} & \ip{192.168.80.254} & — \\
    PC-B & \ip{192.168.80.12}  & \ip{255.255.255.0} & \ip{192.168.80.254} & — \\
    PC-C & \ip{192.168.80.13}  & \ip{255.255.255.0} & \ip{192.168.80.100} & — \\
    PC-D & \ip{192.168.80.12}  & \ip{255.255.255.0} & \ip{192.168.80.254} & — \\
    SRV-C & \ip{10.0.0.8}      & \ip{255.0.0.0}     & \ip{10.0.0.254}     & — \\
    R1 — G0/0 & \ip{192.168.80.254} & \ip{255.255.255.0} & — & Off \\
    R1 — G0/1 & \ip{10.0.0.254}     & \ip{255.0.0.0}     & — & On \\}
  \def\Csol{%
    \hl 1 & \sol{PC-A} & \sol{IP hors du réseau 192.168.80.0/24.} & \sol{Ne joint ni les postes ni le routeur.} & \sol{Ex. 192.168.80.11.} \\ \hline
    \hl 2 & \sol{PC-C} & \sol{Passerelle .80.100 : pas l'adresse du routeur.} & \sol{Joint son réseau, mais aucun autre (serveur).} & \sol{Passerelle 192.168.80.254.} \\ \hline
    \hl 3 & \sol{PC-D} & \sol{Même IP que PC-B (doublon).} & \sol{Conflit d'adresse : PC-B et PC-D perturbés.} & \sol{Adresse unique, ex. .80.14.} \\ \hline
    \hl 4 & \sol{R1 G0/0} & \sol{Interface côté LAN désactivée (Off).} & \sol{Aucun poste ne joint le routeur ni le serveur.} & \sol{Port Status : On.} \\ \hline}
  \def\Cbareme{PC-B et SRV-C sont correctement configurés ; accepter PC-B cité à la place de PC-D pour le doublon.}
\fi

\pagestyle{fancy}
\fancyhf{}
\renewcommand{\headrulewidth}{0pt}
\fancyfoot[L]{\footnotesize\itshape\color{grisfoot}Évaluation — Réseaux locaux — Sujet \variante\ifcorrige\ — CORRIGÉ\fi}
\fancyfoot[R]{\footnotesize\color{grisfoot}\thepage/3}

\setlist{itemsep=0pt,topsep=2pt,parsep=0pt}
\setlength{\parindent}{0pt}
\setlength{\parskip}{2pt}

\begin{document}

\begin{center}
  {\Large\bfseries Évaluation pratique — Réseaux locaux \quad
   \fcolorbox{black}{white}{\ SUJET \variante\ }\par}
  \vspace{1pt}
  {\bfseries Adressage IP, commutation et routage — BTS CIEL 1re année — durée : 2 heures\par}
  \ifcorrige\vspace{2pt}{\bfseries\color{corr}CORRIGÉ ET BARÈME — document enseignant}\par\fi
\end{center}
\vspace{-2pt}
{\renewcommand{\arraystretch}{1.35}
\begin{tabularx}{\linewidth}{|X|X|p{2.2cm}|p{2.6cm}|}
  \hline
  \textbf{Nom :} & \textbf{Prénom :} & \textbf{Poste n°} & \textbf{Note :} \hfill /20 \\ \hline
\end{tabularx}}

\textbf{Consignes :} travail \textbf{individuel}, documents et Internet \textbf{non autorisés}. Ce sujet est le
\textbf{document-réponse}. Fichier Packet Tracer : \ip{Nom\_Prenom\_EvalLAN\_\variante.pkt}, à enregistrer
régulièrement dans le dossier de rendu. Les « visas » sont apposés par le professeur, que vous appelez à chaque test réussi.
Durées conseillées : A 30 min — B 1 h 10 — C 20 min.

%==========================================================================
\section{Partie A — Questions de cours et calculs \hfill /6}

\question{A1}{2 pts}{Complétez le tableau en considérant le \textbf{masque par défaut} de la classe de chaque adresse.}
{\small\renewcommand{\arraystretch}{1.15}
\begin{tabularx}{\linewidth}{|l|c|C|C|C|c|}
  \hline
  \textbf{Adresse IP} & \textbf{Classe} & \textbf{Masque par défaut} & \textbf{Adresse réseau} &
  \textbf{Adresse de diffusion} & \textbf{Nb d'hôtes} \\ \hline
  \AIrows
\end{tabularx}}
\bareme{0,5 pt par ligne juste + 0,5 pt si les trois nombres d'hôtes sont justes ($2^{n}-2$, $n$ = nombre de bits d'hôte).}

\question{A2}{1 pt}{Indiquez le type de câble (\textbf{droit} ou \textbf{croisé}) pour chaque liaison.}
{\small\renewcommand{\arraystretch}{1.15}
\begin{tabularx}{\linewidth}{|X|c||X|c|}
  \hline
  \hl PC $\leftrightarrow$ commutateur & \makebox[1.6cm]{\sol{droit}} & Commutateur $\leftrightarrow$ routeur & \makebox[1.6cm]{\sol{droit}} \\ \hline
  \hl PC $\leftrightarrow$ PC (réseau P2P) & \sol{croisé} & Concentrateur $\leftrightarrow$ concentrateur & \sol{croisé} \\ \hline
\end{tabularx}}
\bareme{0,25 pt par case. Équipements de types différents : câble droit ; de même type : câble croisé.}

\question{A3}{1 pt}{Comment un \textbf{concentrateur (hub)} et un \textbf{commutateur (switch)} traitent-ils une trame
reçue ? Lequel limite les collisions, et pourquoi ?}
\zone{2.9cm}{Le \textbf{hub} (couche 1) répète la trame sur \textbf{tous ses autres ports} : un seul domaine de collision
partagé. Le \textbf{switch} (couche 2) lit l'adresse MAC de destination et, grâce à sa \textbf{table MAC}, ne l'envoie
\textbf{que sur le port du destinataire} : un domaine de collision par port, c'est lui qui limite les collisions.
\emph{0,5 pt hub, 0,5 pt switch.}}

\question{A4}{1 pt}{Qu'est-ce que la \textbf{passerelle par défaut} d'un poste ? Quelle condition son adresse doit-elle respecter ?}
\zone{2.4cm}{Adresse IP de l'interface du \textbf{routeur} à laquelle le poste envoie les paquets destinés à un
\textbf{autre réseau}. Elle doit \textbf{appartenir au même réseau} que le poste. \emph{0,5 + 0,5 pt.}}

\question{A5}{0,5 pt}{Un portable configuré en DHCP affiche \ip{\apipa} / \ip{255.255.0.0}. Que s'est-il passé ?}
\zone{1.5cm}{Aucun \textbf{serveur DHCP} n'a répondu : le poste s'est attribué une adresse \textbf{APIPA} (169.254.0.0/16).}

\question{A6}{0,5 pt}{Quel protocole utilise \texttt{ping} ? Nommez les deux messages échangés et la couche OSI des adresses IP.}
\zone{1.5cm}{\textbf{ICMP} ; \textbf{Echo Request} / \textbf{Echo Reply} ; adresses IP en \textbf{couche 3 (Réseau)}.}

%==========================================================================
\newpage
\section{Partie B — Réalisation sous Packet Tracer \hfill /10}

\begin{cdc}
\entreprise{} vous confie son réseau : \textbf{trois réseaux} reliés par un \textbf{routeur unique Cisco 2911}
nommé \ip{\Rnom} (interfaces G0/0, G0/1, G0/2).
\begin{itemize}[leftmargin=12pt]
  \item \textbf{Réseau \resUn} : \textbf{4 postes} PC-\pUn1 à PC-\pUn4, commutateur \ip{SW-\pUn}, adressage statique,
        réseau \ip{\netUn.0/24}.
  \item \textbf{Réseau \resDeux} : \textbf{2 postes} PC-\pDeux1, PC-\pDeux2 et \textbf{1 imprimante} IMP-\pDeux,
        commutateur \ip{SW-\pDeux}, adressage statique, réseau de classe B \ip{\netDeux.0.0/16}.
  \item \textbf{Réseau Serveur} : serveur \ip{\srvNom} d'adresse \ip{\netSrv.\srvH/24}, relié \textbf{directement} au routeur.
  \item La passerelle de chaque réseau est sa \textbf{dernière adresse utilisable}. Tous les équipements communiquent
        entre eux et accèdent à la page web du serveur. Nommage imposé (« Display Name »).
\end{itemize}
\end{cdc}

\question{B1}{2 pts}{Complétez le plan d'adressage \textbf{avant} la réalisation.}
{\footnotesize\setlength{\tabcolsep}{3pt}\renewcommand{\arraystretch}{1.15}
\begin{tabularx}{\linewidth}{|l|C|C|C|C|C|C|}
  \hline
  \textbf{Réseau} & \textbf{Adresse réseau} & \textbf{Masque} & \textbf{1re @ utilisable} &
  \textbf{Dernière @ utilisable} & \textbf{Diffusion} & \textbf{Passerelle} \\ \hline
  \hl \resUn & \sol{\netUn.0} & \sol{255.255.255.0} & \sol{\netUn.1} & \sol{\netUn.254} & \sol{\netUn.255} & \sol{\netUn.254} \\ \hline
  \hl \resDeux & \sol{\netDeux.0.0} & \sol{255.255.0.0} & \sol{\netDeux.0.1} & \sol{\netDeux.255.254} & \sol{\netDeux.255.255} & \sol{\netDeux.255.254} \\ \hline
  \hl Serveur & \sol{\netSrv.0} & \sol{255.255.255.0} & \sol{\netSrv.1} & \sol{\netSrv.254} & \sol{\netSrv.255} & \sol{\netSrv.254} \\ \hline
\end{tabularx}}
\bareme{0,5 pt par ligne juste + 0,5 pt si la convention de passerelle est respectée partout.}

\question{B2}{2 pts}{Construisez la topologie (équipements, câblage, nommage) et relevez l'adressage configuré.}
{\footnotesize\setlength{\tabcolsep}{3pt}\renewcommand{\arraystretch}{1.1}
\begin{tabularx}{\linewidth}{|l|C|C|C||l|C|C|}
  \hline
  \textbf{Hôte} & \textbf{Adresse IP} & \textbf{Masque} & \textbf{Passerelle} & \textbf{Interface} & \textbf{Adresse IP} & \textbf{Masque} \\ \hline
  \hl PC-\pUn1 & \sol{\netUn.1} & \sol{255.255.255.0} & \sol{\netUn.254} & \Rnom{} G0/0 & \sol{\netUn.254} & \sol{255.255.255.0} \\ \hline
  \hl PC-\pDeux1 & \sol{\netDeux.0.1} & \sol{255.255.0.0} & \sol{\netDeux.255.254} & \Rnom{} G0/1 & \sol{\netDeux.255.254} & \sol{255.255.0.0} \\ \hline
  \hl IMP-\pDeux & \sol{\netDeux.0.10} & \sol{255.255.0.0} & \sol{\netDeux.255.254} & \Rnom{} G0/2 & \sol{\netSrv.254} & \sol{255.255.255.0} \\ \hline
  \hl \srvNom & \ip{\netSrv.\srvH} & \sol{255.255.255.0} & \sol{\netSrv.254} & \multicolumn{3}{c|}{\cellcolor{black!8}} \\ \hline
\end{tabularx}}
\bareme{Adresses d'hôtes libres (uniques, dans le bon réseau) ; affectation G0/x libre. Sur le .pkt : équipements 0,5 ;
câblage 1 (câbles droits vers les switchs et switch–routeur, \textbf{croisé} serveur–routeur, liens verts) ; nommage 0,5.}

\question{B3}{3 pts}{Configurez tous les équipements : IP, masques, passerelles des hôtes ; interfaces du routeur
(« Port Status : On »).}
\bareme{Sur le .pkt : interfaces routeur 1 pt ; IP/masques des hôtes 1 pt ; passerelles (imprimante et serveur compris)
1 pt. $-0,25$ par erreur, sans descendre sous 0 par item.}

\question{B4}{3 pts}{Réalisez les tests, notez la commande et le résultat, puis faites viser chaque test réussi.}
{\small\renewcommand{\arraystretch}{1.1}
\begin{tabularx}{\linewidth}{|c|l|l|C|c|c|}
  \hline
  \textbf{N°} & \textbf{Source} & \textbf{Destination} & \textbf{Commande / action} & \textbf{Résultat} & \textbf{Visa prof.} \\ \hline
  \hl 1 & PC-\pUn1 & PC-\pUn4 & \sol{ping \netUn.4} & \sol{réussi} & \hspace{1.4cm} \\ \hline
  \hl 2 & PC-\pDeux1 & IMP-\pDeux & \sol{ping \netDeux.0.10} & \sol{réussi} & \\ \hline
  \hl 3 & PC-\pUn2 & PC-\pDeux2 & \sol{ping \netDeux.0.2} & \sol{réussi} & \\ \hline
  \hl 4 & PC-\pDeux2 & \srvNom & \sol{ping \netSrv.\srvH} & \sol{réussi} & \\ \hline
  \hl 5 & PC-\pUn3 & page web du serveur & \sol{navigateur : {\NoAutoSpacing http://\netSrv.\srvH}} & \sol{page affichée} & \\ \hline
\end{tabularx}}
\bareme{0,5 pt par test visé + 0,5 pt pour B5. Le premier ping inter-réseaux peut échouer (résolution ARP) ; les suivants doivent réussir.}

\question{B5}{incl. B4}{Test n°3 : par quels équipements passe le paquet ? Pourquoi PC-\pUn2 l'envoie-t-il au routeur ?}
\zone{2.3cm}{PC-\pUn2 $\to$ SW-\pUn{} $\to$ \Rnom{} $\to$ SW-\pDeux{} $\to$ PC-\pDeux2. Grâce à son masque, PC-\pUn2 constate que
la destination (\netDeux.0.0/16) n'est \textbf{pas dans son réseau} (\netUn.0/24) : il envoie le paquet à sa
\textbf{passerelle par défaut}, qui le route.}

%==========================================================================
\newpage
\section{Partie C — Diagnostic d'un réseau en panne \hfill /4}

Un technicien constate qu'\textbf{aucun poste ne parvient à joindre le serveur SRV-C} et que certains postes
communiquent mal entre eux. Les câbles sont du bon type.

\begin{center}
\begin{minipage}{0.37\linewidth}\centering
\begin{tikzpicture}[font=\footnotesize,
  eq/.style={draw,fill=qbg,rounded corners=2pt,minimum width=9mm,minimum height=5mm,inner sep=2pt},
  every path/.style={thick}]
  \node[eq] (SW) {SW-1};
  \node[eq,right=9mm of SW] (R) {R1};
  \node[eq,right=7mm of R] (SRV) {SRV-C};
  \node[eq,above left=5mm and -3mm of SW] (A) {PC-A};
  \node[eq,above right=5mm and -3mm of SW] (B) {PC-B};
  \node[eq,below left=5mm and -3mm of SW] (C) {PC-C};
  \node[eq,below right=5mm and -3mm of SW] (D) {PC-D};
  \foreach \n in {A,B,C,D} \draw (\n) -- (SW);
  \draw (SW) -- node[above,font=\tiny]{G0/0} (R);
  \draw (R) -- node[above,font=\tiny]{G0/1} (SRV);
\end{tikzpicture}
\end{minipage}\hfill
\begin{minipage}{0.61\linewidth}\centering
{\small\renewcommand{\arraystretch}{1.05}\setlength{\tabcolsep}{4pt}
\begin{tabular}{lllll}
  \toprule
  \textbf{Équipement} & \textbf{Adresse IP} & \textbf{Masque} & \textbf{Passerelle} & \textbf{Port} \\
  \midrule
  \Crows
  \bottomrule
\end{tabular}}
\end{minipage}
\end{center}

\question{C1}{4 pts}{Relevez les \textbf{quatre anomalies}. Pour chacune : conséquence et correction.}
{\small\renewcommand{\arraystretch}{1.1}
\ifcorrige\def\hc{}\else\def\hc{\rule[-10mm]{0pt}{14mm}}\fi
\begin{tabularx}{\linewidth}{|c|>{\raggedright\arraybackslash}p{2cm}|L|L|L|}
  \hline
  \textbf{N°} & \textbf{Équipement} & \textbf{Anomalie} & \textbf{Conséquence} & \textbf{Correction} \\ \hline
  \ifcorrige\Csol\else
  \hc 1 & & & & \\ \hline
  \hc 2 & & & & \\ \hline
  \hc 3 & & & & \\ \hline
  \hc 4 & & & & \\ \hline
  \fi
\end{tabularx}}
\bareme{1 pt par anomalie (0,25 équipement, anomalie, conséquence, correction ; ordre indifférent). \Cbareme}

%==========================================================================
\section{Bonus — Extension sans fil \hfill +1}

\begin{cdc}Ajoutez au réseau \resUn{} un \textbf{AccessPoint-PT} relié à SW-\pUn{} (SSID \ip{\ssid}, \textbf{WPA2-PSK})
et un portable \ip{Laptop-\pUn} équipé d'un module \texttt{WPC300N}, en adressage \textbf{statique}.
Vérifiez qu'il joint \srvNom.\end{cdc}
{\small\renewcommand{\arraystretch}{1.1}
\begin{tabularx}{\linewidth}{|l|C|C|C|c|}
  \hline
  \textbf{Équipement} & \textbf{Adresse IP} & \textbf{Masque} & \textbf{Passerelle} & \textbf{Visa prof.} \\ \hline
  \hl Laptop-\pUn & \sol{\netUn.50} & \sol{255.255.255.0} & \sol{\netUn.254} & \hspace{1.4cm} \\ \hline
\end{tabularx}}
\bareme{Le point d'accès fait office de pont : le portable est dans \netUn.0/24. Bonus si associé en WPA2-PSK et ping
vers \netSrv.\srvH{} réussi. Note plafonnée à 20.}

\end{document}
"
%   pdflatex -jobname=Eval_Reseau_CorrigeB "\def\variante{B}\def\modecorrige{1}% Évaluation pratique — Réseaux locaux (BTS CIEL 1re année) — sujets A et B
% Compilation (4 PDF) :
%   pdflatex -jobname=Eval_Reseau_SujetA   "\def\variante{A}\def\modecorrige{0}% Évaluation pratique — Réseaux locaux (BTS CIEL 1re année) — sujets A et B
% Compilation (4 PDF) :
%   pdflatex -jobname=Eval_Reseau_SujetA   "\def\variante{A}\def\modecorrige{0}\input{Eval_TP_Reseau_LAN.tex}"
%   pdflatex -jobname=Eval_Reseau_SujetB   "\def\variante{B}\def\modecorrige{0}\input{Eval_TP_Reseau_LAN.tex}"
%   pdflatex -jobname=Eval_Reseau_CorrigeA "\def\variante{A}\def\modecorrige{1}\input{Eval_TP_Reseau_LAN.tex}"
%   pdflatex -jobname=Eval_Reseau_CorrigeB "\def\variante{B}\def\modecorrige{1}\input{Eval_TP_Reseau_LAN.tex}"
\documentclass[10pt,a4paper]{article}

\usepackage[T1]{fontenc}
\usepackage[utf8]{inputenc}
\usepackage{lmodern}
\usepackage[french]{babel}
\usepackage[a4paper,hmargin=1.5cm,top=1.3cm,bottom=1.6cm,footskip=0.7cm]{geometry}
\usepackage[table]{xcolor}
\usepackage{booktabs,tabularx,array}
\usepackage{enumitem}
\usepackage{fancyhdr}
\usepackage[most]{tcolorbox}
\usepackage{tikz}
\usetikzlibrary{positioning}
\usepackage{titlesec}

\providecommand{\modecorrige}{0}
\providecommand{\variante}{A}
\newif\ifcorrige
\ifnum\modecorrige=1 \corrigetrue \else \corrigefalse \fi

\definecolor{ipred}{HTML}{C0392B}
\definecolor{cdcbar}{HTML}{A67C3D}
\definecolor{cdcbg}{HTML}{F4F1EA}
\definecolor{qbar}{HTML}{4A6B5D}
\definecolor{qbg}{HTML}{EDF1EF}
\definecolor{corr}{HTML}{1F5FA8}
\definecolor{corrbg}{HTML}{E8F0FA}
\definecolor{grisfoot}{HTML}{828282}

\newcommand{\ip}[1]{\textcolor{ipred}{\texttt{#1}}}
\newcommand{\pts}[1]{\hfill\mbox{\textbf{(#1)}}}
\newcommand{\sol}[1]{\ifcorrige\textcolor{corr}{#1}\fi}
\newcommand{\bareme}[1]{\ifcorrige\par{\footnotesize\color{corr}#1\par}\fi}
\newcommand{\hl}{\ifcorrige\rule[-1.6mm]{0pt}{5.4mm}\else\rule[-2.2mm]{0pt}{6.6mm}\fi}
\newcolumntype{C}{>{\centering\arraybackslash}X}
\newcolumntype{L}{>{\raggedright\arraybackslash}X}

\titleformat{\section}{\large\bfseries}{}{0pt}{}
\titlespacing*{\section}{0pt}{6pt}{3pt}

\tcbset{qstyle/.style={enhanced,boxrule=0pt,leftrule=2.5pt,arc=0pt,outer arc=0pt,
  left=5pt,right=5pt,top=2pt,bottom=2pt,before skip=4pt,after skip=3pt}}
\newtcolorbox{cdc}{qstyle,colback=cdcbg,colframe=cdcbar}
\newtcolorbox{quest}{qstyle,colback=qbg,colframe=qbar}

% Zone de réponse : vierge dans le sujet, remplie dans le corrigé
\newcommand{\zone}[2]{%
  \ifcorrige
    \begin{tcolorbox}[colback=corrbg,colframe=corr,boxrule=0.5pt,arc=1pt,left=4pt,right=4pt,top=1pt,bottom=1pt,
      before skip=2pt,after skip=2pt,fontupper=\footnotesize\color{corr}]#2\end{tcolorbox}%
  \else
    \begin{tcolorbox}[colback=white,colframe=black!45,boxrule=0.4pt,arc=1pt,height=#1,
      before skip=2pt,after skip=2pt]\end{tcolorbox}%
  \fi}

\newcommand{\question}[3]{\begin{quest}\textbf{#1.} #3\pts{#2}\end{quest}}

%==========================================================================
% Données propres à chaque variante
%==========================================================================
\if\variante A
  \def\entreprise{Le cabinet d'architectes \textbf{ArchiNova}}
  \def\Rnom{R-ArchiNova}
  \def\resUn{Bureau d'études}\def\pUn{BE}\def\netUn{192.168.40}
  \def\resDeux{Accueil}\def\pDeux{ACC}\def\netDeux{172.16}
  \def\srvNom{SRV-Intranet}\def\netSrv{192.168.200}\def\srvH{10}
  \def\ssid{ArchiNova-WiFi}
  \def\AIrows{%
    \hl\ip{192.168.40.25} & \sol{C} & \sol{255.255.255.0} & \sol{192.168.40.0} & \sol{192.168.40.255} & \sol{254} \\ \hline
    \hl\ip{172.16.5.10}   & \sol{B} & \sol{255.255.0.0}   & \sol{172.16.0.0}   & \sol{172.16.255.255} & \sol{65\,534} \\ \hline
    \hl\ip{10.3.2.1}      & \sol{A} & \sol{255.0.0.0}     & \sol{10.0.0.0}     & \sol{10.255.255.255} & \sol{16\,777\,214} \\ \hline}
  \def\apipa{169.254.12.7}
  \def\Cnet{192.168.60}
  \def\Crows{%
    PC-A & \ip{192.168.60.10} & \ip{255.255.255.0} & \ip{192.168.60.254} & — \\
    PC-B & \ip{192.168.60.20} & \ip{255.255.255.0} & \ip{192.168.60.1}   & — \\
    PC-C & \ip{192.168.61.30} & \ip{255.255.255.0} & \ip{192.168.60.254} & — \\
    PC-D & \ip{192.168.60.10} & \ip{255.255.255.0} & \ip{192.168.60.254} & — \\
    SRV-C & \ip{10.0.0.5}     & \ip{255.0.0.0}     & \ip{10.0.0.254}     & — \\
    R1 — G0/0 & \ip{192.168.60.254} & \ip{255.255.255.0} & — & On \\
    R1 — G0/1 & \ip{10.0.0.254}     & \ip{255.0.0.0}     & — & Off \\}
  \def\Csol{%
    \hl 1 & \sol{PC-B} & \sol{Passerelle .60.1 : pas l'adresse du routeur.} & \sol{Joint son réseau, mais aucun autre (serveur).} & \sol{Passerelle 192.168.60.254.} \\ \hline
    \hl 2 & \sol{PC-C} & \sol{IP hors du réseau 192.168.60.0/24.} & \sol{Ne joint ni les postes ni le routeur.} & \sol{Ex. 192.168.60.30.} \\ \hline
    \hl 3 & \sol{PC-D} & \sol{Même IP que PC-A (doublon).} & \sol{Conflit d'adresse : PC-A et PC-D perturbés.} & \sol{Adresse unique, ex. .60.40.} \\ \hline
    \hl 4 & \sol{R1 G0/1} & \sol{Interface désactivée (Off).} & \sol{Réseau du serveur injoignable pour tous.} & \sol{Port Status : On.} \\ \hline}
  \def\Cbareme{PC-A et SRV-C sont correctement configurés ; accepter PC-A cité à la place de PC-D pour le doublon.}
\else
  \def\entreprise{La clinique vétérinaire \textbf{VetoPlus}}
  \def\Rnom{R-VetoPlus}
  \def\resUn{Soins}\def\pUn{SOI}\def\netUn{192.168.70}
  \def\resDeux{Secrétariat}\def\pDeux{SEC}\def\netDeux{172.20}
  \def\srvNom{SRV-Dossiers}\def\netSrv{192.168.250}\def\srvH{20}
  \def\ssid{VetoPlus-WiFi}
  \def\AIrows{%
    \hl\ip{192.168.15.130} & \sol{C} & \sol{255.255.255.0} & \sol{192.168.15.0} & \sol{192.168.15.255} & \sol{254} \\ \hline
    \hl\ip{172.31.200.7}   & \sol{B} & \sol{255.255.0.0}   & \sol{172.31.0.0}   & \sol{172.31.255.255} & \sol{65\,534} \\ \hline
    \hl\ip{10.45.6.9}      & \sol{A} & \sol{255.0.0.0}     & \sol{10.0.0.0}     & \sol{10.255.255.255} & \sol{16\,777\,214} \\ \hline}
  \def\apipa{169.254.88.3}
  \def\Cnet{192.168.80}
  \def\Crows{%
    PC-A & \ip{192.168.8.11}   & \ip{255.255.255.0} & \ip{192.168.80.254} & — \\
    PC-B & \ip{192.168.80.12}  & \ip{255.255.255.0} & \ip{192.168.80.254} & — \\
    PC-C & \ip{192.168.80.13}  & \ip{255.255.255.0} & \ip{192.168.80.100} & — \\
    PC-D & \ip{192.168.80.12}  & \ip{255.255.255.0} & \ip{192.168.80.254} & — \\
    SRV-C & \ip{10.0.0.8}      & \ip{255.0.0.0}     & \ip{10.0.0.254}     & — \\
    R1 — G0/0 & \ip{192.168.80.254} & \ip{255.255.255.0} & — & Off \\
    R1 — G0/1 & \ip{10.0.0.254}     & \ip{255.0.0.0}     & — & On \\}
  \def\Csol{%
    \hl 1 & \sol{PC-A} & \sol{IP hors du réseau 192.168.80.0/24.} & \sol{Ne joint ni les postes ni le routeur.} & \sol{Ex. 192.168.80.11.} \\ \hline
    \hl 2 & \sol{PC-C} & \sol{Passerelle .80.100 : pas l'adresse du routeur.} & \sol{Joint son réseau, mais aucun autre (serveur).} & \sol{Passerelle 192.168.80.254.} \\ \hline
    \hl 3 & \sol{PC-D} & \sol{Même IP que PC-B (doublon).} & \sol{Conflit d'adresse : PC-B et PC-D perturbés.} & \sol{Adresse unique, ex. .80.14.} \\ \hline
    \hl 4 & \sol{R1 G0/0} & \sol{Interface côté LAN désactivée (Off).} & \sol{Aucun poste ne joint le routeur ni le serveur.} & \sol{Port Status : On.} \\ \hline}
  \def\Cbareme{PC-B et SRV-C sont correctement configurés ; accepter PC-B cité à la place de PC-D pour le doublon.}
\fi

\pagestyle{fancy}
\fancyhf{}
\renewcommand{\headrulewidth}{0pt}
\fancyfoot[L]{\footnotesize\itshape\color{grisfoot}Évaluation — Réseaux locaux — Sujet \variante\ifcorrige\ — CORRIGÉ\fi}
\fancyfoot[R]{\footnotesize\color{grisfoot}\thepage/3}

\setlist{itemsep=0pt,topsep=2pt,parsep=0pt}
\setlength{\parindent}{0pt}
\setlength{\parskip}{2pt}

\begin{document}

\begin{center}
  {\Large\bfseries Évaluation pratique — Réseaux locaux \quad
   \fcolorbox{black}{white}{\ SUJET \variante\ }\par}
  \vspace{1pt}
  {\bfseries Adressage IP, commutation et routage — BTS CIEL 1re année — durée : 2 heures\par}
  \ifcorrige\vspace{2pt}{\bfseries\color{corr}CORRIGÉ ET BARÈME — document enseignant}\par\fi
\end{center}
\vspace{-2pt}
{\renewcommand{\arraystretch}{1.35}
\begin{tabularx}{\linewidth}{|X|X|p{2.2cm}|p{2.6cm}|}
  \hline
  \textbf{Nom :} & \textbf{Prénom :} & \textbf{Poste n°} & \textbf{Note :} \hfill /20 \\ \hline
\end{tabularx}}

\textbf{Consignes :} travail \textbf{individuel}, documents et Internet \textbf{non autorisés}. Ce sujet est le
\textbf{document-réponse}. Fichier Packet Tracer : \ip{Nom\_Prenom\_EvalLAN\_\variante.pkt}, à enregistrer
régulièrement dans le dossier de rendu. Les « visas » sont apposés par le professeur, que vous appelez à chaque test réussi.
Durées conseillées : A 30 min — B 1 h 10 — C 20 min.

%==========================================================================
\section{Partie A — Questions de cours et calculs \hfill /6}

\question{A1}{2 pts}{Complétez le tableau en considérant le \textbf{masque par défaut} de la classe de chaque adresse.}
{\small\renewcommand{\arraystretch}{1.15}
\begin{tabularx}{\linewidth}{|l|c|C|C|C|c|}
  \hline
  \textbf{Adresse IP} & \textbf{Classe} & \textbf{Masque par défaut} & \textbf{Adresse réseau} &
  \textbf{Adresse de diffusion} & \textbf{Nb d'hôtes} \\ \hline
  \AIrows
\end{tabularx}}
\bareme{0,5 pt par ligne juste + 0,5 pt si les trois nombres d'hôtes sont justes ($2^{n}-2$, $n$ = nombre de bits d'hôte).}

\question{A2}{1 pt}{Indiquez le type de câble (\textbf{droit} ou \textbf{croisé}) pour chaque liaison.}
{\small\renewcommand{\arraystretch}{1.15}
\begin{tabularx}{\linewidth}{|X|c||X|c|}
  \hline
  \hl PC $\leftrightarrow$ commutateur & \makebox[1.6cm]{\sol{droit}} & Commutateur $\leftrightarrow$ routeur & \makebox[1.6cm]{\sol{droit}} \\ \hline
  \hl PC $\leftrightarrow$ PC (réseau P2P) & \sol{croisé} & Concentrateur $\leftrightarrow$ concentrateur & \sol{croisé} \\ \hline
\end{tabularx}}
\bareme{0,25 pt par case. Équipements de types différents : câble droit ; de même type : câble croisé.}

\question{A3}{1 pt}{Comment un \textbf{concentrateur (hub)} et un \textbf{commutateur (switch)} traitent-ils une trame
reçue ? Lequel limite les collisions, et pourquoi ?}
\zone{2.9cm}{Le \textbf{hub} (couche 1) répète la trame sur \textbf{tous ses autres ports} : un seul domaine de collision
partagé. Le \textbf{switch} (couche 2) lit l'adresse MAC de destination et, grâce à sa \textbf{table MAC}, ne l'envoie
\textbf{que sur le port du destinataire} : un domaine de collision par port, c'est lui qui limite les collisions.
\emph{0,5 pt hub, 0,5 pt switch.}}

\question{A4}{1 pt}{Qu'est-ce que la \textbf{passerelle par défaut} d'un poste ? Quelle condition son adresse doit-elle respecter ?}
\zone{2.4cm}{Adresse IP de l'interface du \textbf{routeur} à laquelle le poste envoie les paquets destinés à un
\textbf{autre réseau}. Elle doit \textbf{appartenir au même réseau} que le poste. \emph{0,5 + 0,5 pt.}}

\question{A5}{0,5 pt}{Un portable configuré en DHCP affiche \ip{\apipa} / \ip{255.255.0.0}. Que s'est-il passé ?}
\zone{1.5cm}{Aucun \textbf{serveur DHCP} n'a répondu : le poste s'est attribué une adresse \textbf{APIPA} (169.254.0.0/16).}

\question{A6}{0,5 pt}{Quel protocole utilise \texttt{ping} ? Nommez les deux messages échangés et la couche OSI des adresses IP.}
\zone{1.5cm}{\textbf{ICMP} ; \textbf{Echo Request} / \textbf{Echo Reply} ; adresses IP en \textbf{couche 3 (Réseau)}.}

%==========================================================================
\newpage
\section{Partie B — Réalisation sous Packet Tracer \hfill /10}

\begin{cdc}
\entreprise{} vous confie son réseau : \textbf{trois réseaux} reliés par un \textbf{routeur unique Cisco 2911}
nommé \ip{\Rnom} (interfaces G0/0, G0/1, G0/2).
\begin{itemize}[leftmargin=12pt]
  \item \textbf{Réseau \resUn} : \textbf{4 postes} PC-\pUn1 à PC-\pUn4, commutateur \ip{SW-\pUn}, adressage statique,
        réseau \ip{\netUn.0/24}.
  \item \textbf{Réseau \resDeux} : \textbf{2 postes} PC-\pDeux1, PC-\pDeux2 et \textbf{1 imprimante} IMP-\pDeux,
        commutateur \ip{SW-\pDeux}, adressage statique, réseau de classe B \ip{\netDeux.0.0/16}.
  \item \textbf{Réseau Serveur} : serveur \ip{\srvNom} d'adresse \ip{\netSrv.\srvH/24}, relié \textbf{directement} au routeur.
  \item La passerelle de chaque réseau est sa \textbf{dernière adresse utilisable}. Tous les équipements communiquent
        entre eux et accèdent à la page web du serveur. Nommage imposé (« Display Name »).
\end{itemize}
\end{cdc}

\question{B1}{2 pts}{Complétez le plan d'adressage \textbf{avant} la réalisation.}
{\footnotesize\setlength{\tabcolsep}{3pt}\renewcommand{\arraystretch}{1.15}
\begin{tabularx}{\linewidth}{|l|C|C|C|C|C|C|}
  \hline
  \textbf{Réseau} & \textbf{Adresse réseau} & \textbf{Masque} & \textbf{1re @ utilisable} &
  \textbf{Dernière @ utilisable} & \textbf{Diffusion} & \textbf{Passerelle} \\ \hline
  \hl \resUn & \sol{\netUn.0} & \sol{255.255.255.0} & \sol{\netUn.1} & \sol{\netUn.254} & \sol{\netUn.255} & \sol{\netUn.254} \\ \hline
  \hl \resDeux & \sol{\netDeux.0.0} & \sol{255.255.0.0} & \sol{\netDeux.0.1} & \sol{\netDeux.255.254} & \sol{\netDeux.255.255} & \sol{\netDeux.255.254} \\ \hline
  \hl Serveur & \sol{\netSrv.0} & \sol{255.255.255.0} & \sol{\netSrv.1} & \sol{\netSrv.254} & \sol{\netSrv.255} & \sol{\netSrv.254} \\ \hline
\end{tabularx}}
\bareme{0,5 pt par ligne juste + 0,5 pt si la convention de passerelle est respectée partout.}

\question{B2}{2 pts}{Construisez la topologie (équipements, câblage, nommage) et relevez l'adressage configuré.}
{\footnotesize\setlength{\tabcolsep}{3pt}\renewcommand{\arraystretch}{1.1}
\begin{tabularx}{\linewidth}{|l|C|C|C||l|C|C|}
  \hline
  \textbf{Hôte} & \textbf{Adresse IP} & \textbf{Masque} & \textbf{Passerelle} & \textbf{Interface} & \textbf{Adresse IP} & \textbf{Masque} \\ \hline
  \hl PC-\pUn1 & \sol{\netUn.1} & \sol{255.255.255.0} & \sol{\netUn.254} & \Rnom{} G0/0 & \sol{\netUn.254} & \sol{255.255.255.0} \\ \hline
  \hl PC-\pDeux1 & \sol{\netDeux.0.1} & \sol{255.255.0.0} & \sol{\netDeux.255.254} & \Rnom{} G0/1 & \sol{\netDeux.255.254} & \sol{255.255.0.0} \\ \hline
  \hl IMP-\pDeux & \sol{\netDeux.0.10} & \sol{255.255.0.0} & \sol{\netDeux.255.254} & \Rnom{} G0/2 & \sol{\netSrv.254} & \sol{255.255.255.0} \\ \hline
  \hl \srvNom & \ip{\netSrv.\srvH} & \sol{255.255.255.0} & \sol{\netSrv.254} & \multicolumn{3}{c|}{\cellcolor{black!8}} \\ \hline
\end{tabularx}}
\bareme{Adresses d'hôtes libres (uniques, dans le bon réseau) ; affectation G0/x libre. Sur le .pkt : équipements 0,5 ;
câblage 1 (câbles droits vers les switchs et switch–routeur, \textbf{croisé} serveur–routeur, liens verts) ; nommage 0,5.}

\question{B3}{3 pts}{Configurez tous les équipements : IP, masques, passerelles des hôtes ; interfaces du routeur
(« Port Status : On »).}
\bareme{Sur le .pkt : interfaces routeur 1 pt ; IP/masques des hôtes 1 pt ; passerelles (imprimante et serveur compris)
1 pt. $-0,25$ par erreur, sans descendre sous 0 par item.}

\question{B4}{3 pts}{Réalisez les tests, notez la commande et le résultat, puis faites viser chaque test réussi.}
{\small\renewcommand{\arraystretch}{1.1}
\begin{tabularx}{\linewidth}{|c|l|l|C|c|c|}
  \hline
  \textbf{N°} & \textbf{Source} & \textbf{Destination} & \textbf{Commande / action} & \textbf{Résultat} & \textbf{Visa prof.} \\ \hline
  \hl 1 & PC-\pUn1 & PC-\pUn4 & \sol{ping \netUn.4} & \sol{réussi} & \hspace{1.4cm} \\ \hline
  \hl 2 & PC-\pDeux1 & IMP-\pDeux & \sol{ping \netDeux.0.10} & \sol{réussi} & \\ \hline
  \hl 3 & PC-\pUn2 & PC-\pDeux2 & \sol{ping \netDeux.0.2} & \sol{réussi} & \\ \hline
  \hl 4 & PC-\pDeux2 & \srvNom & \sol{ping \netSrv.\srvH} & \sol{réussi} & \\ \hline
  \hl 5 & PC-\pUn3 & page web du serveur & \sol{navigateur : {\NoAutoSpacing http://\netSrv.\srvH}} & \sol{page affichée} & \\ \hline
\end{tabularx}}
\bareme{0,5 pt par test visé + 0,5 pt pour B5. Le premier ping inter-réseaux peut échouer (résolution ARP) ; les suivants doivent réussir.}

\question{B5}{incl. B4}{Test n°3 : par quels équipements passe le paquet ? Pourquoi PC-\pUn2 l'envoie-t-il au routeur ?}
\zone{2.3cm}{PC-\pUn2 $\to$ SW-\pUn{} $\to$ \Rnom{} $\to$ SW-\pDeux{} $\to$ PC-\pDeux2. Grâce à son masque, PC-\pUn2 constate que
la destination (\netDeux.0.0/16) n'est \textbf{pas dans son réseau} (\netUn.0/24) : il envoie le paquet à sa
\textbf{passerelle par défaut}, qui le route.}

%==========================================================================
\newpage
\section{Partie C — Diagnostic d'un réseau en panne \hfill /4}

Un technicien constate qu'\textbf{aucun poste ne parvient à joindre le serveur SRV-C} et que certains postes
communiquent mal entre eux. Les câbles sont du bon type.

\begin{center}
\begin{minipage}{0.37\linewidth}\centering
\begin{tikzpicture}[font=\footnotesize,
  eq/.style={draw,fill=qbg,rounded corners=2pt,minimum width=9mm,minimum height=5mm,inner sep=2pt},
  every path/.style={thick}]
  \node[eq] (SW) {SW-1};
  \node[eq,right=9mm of SW] (R) {R1};
  \node[eq,right=7mm of R] (SRV) {SRV-C};
  \node[eq,above left=5mm and -3mm of SW] (A) {PC-A};
  \node[eq,above right=5mm and -3mm of SW] (B) {PC-B};
  \node[eq,below left=5mm and -3mm of SW] (C) {PC-C};
  \node[eq,below right=5mm and -3mm of SW] (D) {PC-D};
  \foreach \n in {A,B,C,D} \draw (\n) -- (SW);
  \draw (SW) -- node[above,font=\tiny]{G0/0} (R);
  \draw (R) -- node[above,font=\tiny]{G0/1} (SRV);
\end{tikzpicture}
\end{minipage}\hfill
\begin{minipage}{0.61\linewidth}\centering
{\small\renewcommand{\arraystretch}{1.05}\setlength{\tabcolsep}{4pt}
\begin{tabular}{lllll}
  \toprule
  \textbf{Équipement} & \textbf{Adresse IP} & \textbf{Masque} & \textbf{Passerelle} & \textbf{Port} \\
  \midrule
  \Crows
  \bottomrule
\end{tabular}}
\end{minipage}
\end{center}

\question{C1}{4 pts}{Relevez les \textbf{quatre anomalies}. Pour chacune : conséquence et correction.}
{\small\renewcommand{\arraystretch}{1.1}
\ifcorrige\def\hc{}\else\def\hc{\rule[-10mm]{0pt}{14mm}}\fi
\begin{tabularx}{\linewidth}{|c|>{\raggedright\arraybackslash}p{2cm}|L|L|L|}
  \hline
  \textbf{N°} & \textbf{Équipement} & \textbf{Anomalie} & \textbf{Conséquence} & \textbf{Correction} \\ \hline
  \ifcorrige\Csol\else
  \hc 1 & & & & \\ \hline
  \hc 2 & & & & \\ \hline
  \hc 3 & & & & \\ \hline
  \hc 4 & & & & \\ \hline
  \fi
\end{tabularx}}
\bareme{1 pt par anomalie (0,25 équipement, anomalie, conséquence, correction ; ordre indifférent). \Cbareme}

%==========================================================================
\section{Bonus — Extension sans fil \hfill +1}

\begin{cdc}Ajoutez au réseau \resUn{} un \textbf{AccessPoint-PT} relié à SW-\pUn{} (SSID \ip{\ssid}, \textbf{WPA2-PSK})
et un portable \ip{Laptop-\pUn} équipé d'un module \texttt{WPC300N}, en adressage \textbf{statique}.
Vérifiez qu'il joint \srvNom.\end{cdc}
{\small\renewcommand{\arraystretch}{1.1}
\begin{tabularx}{\linewidth}{|l|C|C|C|c|}
  \hline
  \textbf{Équipement} & \textbf{Adresse IP} & \textbf{Masque} & \textbf{Passerelle} & \textbf{Visa prof.} \\ \hline
  \hl Laptop-\pUn & \sol{\netUn.50} & \sol{255.255.255.0} & \sol{\netUn.254} & \hspace{1.4cm} \\ \hline
\end{tabularx}}
\bareme{Le point d'accès fait office de pont : le portable est dans \netUn.0/24. Bonus si associé en WPA2-PSK et ping
vers \netSrv.\srvH{} réussi. Note plafonnée à 20.}

\end{document}
"
%   pdflatex -jobname=Eval_Reseau_SujetB   "\def\variante{B}\def\modecorrige{0}% Évaluation pratique — Réseaux locaux (BTS CIEL 1re année) — sujets A et B
% Compilation (4 PDF) :
%   pdflatex -jobname=Eval_Reseau_SujetA   "\def\variante{A}\def\modecorrige{0}\input{Eval_TP_Reseau_LAN.tex}"
%   pdflatex -jobname=Eval_Reseau_SujetB   "\def\variante{B}\def\modecorrige{0}\input{Eval_TP_Reseau_LAN.tex}"
%   pdflatex -jobname=Eval_Reseau_CorrigeA "\def\variante{A}\def\modecorrige{1}\input{Eval_TP_Reseau_LAN.tex}"
%   pdflatex -jobname=Eval_Reseau_CorrigeB "\def\variante{B}\def\modecorrige{1}\input{Eval_TP_Reseau_LAN.tex}"
\documentclass[10pt,a4paper]{article}

\usepackage[T1]{fontenc}
\usepackage[utf8]{inputenc}
\usepackage{lmodern}
\usepackage[french]{babel}
\usepackage[a4paper,hmargin=1.5cm,top=1.3cm,bottom=1.6cm,footskip=0.7cm]{geometry}
\usepackage[table]{xcolor}
\usepackage{booktabs,tabularx,array}
\usepackage{enumitem}
\usepackage{fancyhdr}
\usepackage[most]{tcolorbox}
\usepackage{tikz}
\usetikzlibrary{positioning}
\usepackage{titlesec}

\providecommand{\modecorrige}{0}
\providecommand{\variante}{A}
\newif\ifcorrige
\ifnum\modecorrige=1 \corrigetrue \else \corrigefalse \fi

\definecolor{ipred}{HTML}{C0392B}
\definecolor{cdcbar}{HTML}{A67C3D}
\definecolor{cdcbg}{HTML}{F4F1EA}
\definecolor{qbar}{HTML}{4A6B5D}
\definecolor{qbg}{HTML}{EDF1EF}
\definecolor{corr}{HTML}{1F5FA8}
\definecolor{corrbg}{HTML}{E8F0FA}
\definecolor{grisfoot}{HTML}{828282}

\newcommand{\ip}[1]{\textcolor{ipred}{\texttt{#1}}}
\newcommand{\pts}[1]{\hfill\mbox{\textbf{(#1)}}}
\newcommand{\sol}[1]{\ifcorrige\textcolor{corr}{#1}\fi}
\newcommand{\bareme}[1]{\ifcorrige\par{\footnotesize\color{corr}#1\par}\fi}
\newcommand{\hl}{\ifcorrige\rule[-1.6mm]{0pt}{5.4mm}\else\rule[-2.2mm]{0pt}{6.6mm}\fi}
\newcolumntype{C}{>{\centering\arraybackslash}X}
\newcolumntype{L}{>{\raggedright\arraybackslash}X}

\titleformat{\section}{\large\bfseries}{}{0pt}{}
\titlespacing*{\section}{0pt}{6pt}{3pt}

\tcbset{qstyle/.style={enhanced,boxrule=0pt,leftrule=2.5pt,arc=0pt,outer arc=0pt,
  left=5pt,right=5pt,top=2pt,bottom=2pt,before skip=4pt,after skip=3pt}}
\newtcolorbox{cdc}{qstyle,colback=cdcbg,colframe=cdcbar}
\newtcolorbox{quest}{qstyle,colback=qbg,colframe=qbar}

% Zone de réponse : vierge dans le sujet, remplie dans le corrigé
\newcommand{\zone}[2]{%
  \ifcorrige
    \begin{tcolorbox}[colback=corrbg,colframe=corr,boxrule=0.5pt,arc=1pt,left=4pt,right=4pt,top=1pt,bottom=1pt,
      before skip=2pt,after skip=2pt,fontupper=\footnotesize\color{corr}]#2\end{tcolorbox}%
  \else
    \begin{tcolorbox}[colback=white,colframe=black!45,boxrule=0.4pt,arc=1pt,height=#1,
      before skip=2pt,after skip=2pt]\end{tcolorbox}%
  \fi}

\newcommand{\question}[3]{\begin{quest}\textbf{#1.} #3\pts{#2}\end{quest}}

%==========================================================================
% Données propres à chaque variante
%==========================================================================
\if\variante A
  \def\entreprise{Le cabinet d'architectes \textbf{ArchiNova}}
  \def\Rnom{R-ArchiNova}
  \def\resUn{Bureau d'études}\def\pUn{BE}\def\netUn{192.168.40}
  \def\resDeux{Accueil}\def\pDeux{ACC}\def\netDeux{172.16}
  \def\srvNom{SRV-Intranet}\def\netSrv{192.168.200}\def\srvH{10}
  \def\ssid{ArchiNova-WiFi}
  \def\AIrows{%
    \hl\ip{192.168.40.25} & \sol{C} & \sol{255.255.255.0} & \sol{192.168.40.0} & \sol{192.168.40.255} & \sol{254} \\ \hline
    \hl\ip{172.16.5.10}   & \sol{B} & \sol{255.255.0.0}   & \sol{172.16.0.0}   & \sol{172.16.255.255} & \sol{65\,534} \\ \hline
    \hl\ip{10.3.2.1}      & \sol{A} & \sol{255.0.0.0}     & \sol{10.0.0.0}     & \sol{10.255.255.255} & \sol{16\,777\,214} \\ \hline}
  \def\apipa{169.254.12.7}
  \def\Cnet{192.168.60}
  \def\Crows{%
    PC-A & \ip{192.168.60.10} & \ip{255.255.255.0} & \ip{192.168.60.254} & — \\
    PC-B & \ip{192.168.60.20} & \ip{255.255.255.0} & \ip{192.168.60.1}   & — \\
    PC-C & \ip{192.168.61.30} & \ip{255.255.255.0} & \ip{192.168.60.254} & — \\
    PC-D & \ip{192.168.60.10} & \ip{255.255.255.0} & \ip{192.168.60.254} & — \\
    SRV-C & \ip{10.0.0.5}     & \ip{255.0.0.0}     & \ip{10.0.0.254}     & — \\
    R1 — G0/0 & \ip{192.168.60.254} & \ip{255.255.255.0} & — & On \\
    R1 — G0/1 & \ip{10.0.0.254}     & \ip{255.0.0.0}     & — & Off \\}
  \def\Csol{%
    \hl 1 & \sol{PC-B} & \sol{Passerelle .60.1 : pas l'adresse du routeur.} & \sol{Joint son réseau, mais aucun autre (serveur).} & \sol{Passerelle 192.168.60.254.} \\ \hline
    \hl 2 & \sol{PC-C} & \sol{IP hors du réseau 192.168.60.0/24.} & \sol{Ne joint ni les postes ni le routeur.} & \sol{Ex. 192.168.60.30.} \\ \hline
    \hl 3 & \sol{PC-D} & \sol{Même IP que PC-A (doublon).} & \sol{Conflit d'adresse : PC-A et PC-D perturbés.} & \sol{Adresse unique, ex. .60.40.} \\ \hline
    \hl 4 & \sol{R1 G0/1} & \sol{Interface désactivée (Off).} & \sol{Réseau du serveur injoignable pour tous.} & \sol{Port Status : On.} \\ \hline}
  \def\Cbareme{PC-A et SRV-C sont correctement configurés ; accepter PC-A cité à la place de PC-D pour le doublon.}
\else
  \def\entreprise{La clinique vétérinaire \textbf{VetoPlus}}
  \def\Rnom{R-VetoPlus}
  \def\resUn{Soins}\def\pUn{SOI}\def\netUn{192.168.70}
  \def\resDeux{Secrétariat}\def\pDeux{SEC}\def\netDeux{172.20}
  \def\srvNom{SRV-Dossiers}\def\netSrv{192.168.250}\def\srvH{20}
  \def\ssid{VetoPlus-WiFi}
  \def\AIrows{%
    \hl\ip{192.168.15.130} & \sol{C} & \sol{255.255.255.0} & \sol{192.168.15.0} & \sol{192.168.15.255} & \sol{254} \\ \hline
    \hl\ip{172.31.200.7}   & \sol{B} & \sol{255.255.0.0}   & \sol{172.31.0.0}   & \sol{172.31.255.255} & \sol{65\,534} \\ \hline
    \hl\ip{10.45.6.9}      & \sol{A} & \sol{255.0.0.0}     & \sol{10.0.0.0}     & \sol{10.255.255.255} & \sol{16\,777\,214} \\ \hline}
  \def\apipa{169.254.88.3}
  \def\Cnet{192.168.80}
  \def\Crows{%
    PC-A & \ip{192.168.8.11}   & \ip{255.255.255.0} & \ip{192.168.80.254} & — \\
    PC-B & \ip{192.168.80.12}  & \ip{255.255.255.0} & \ip{192.168.80.254} & — \\
    PC-C & \ip{192.168.80.13}  & \ip{255.255.255.0} & \ip{192.168.80.100} & — \\
    PC-D & \ip{192.168.80.12}  & \ip{255.255.255.0} & \ip{192.168.80.254} & — \\
    SRV-C & \ip{10.0.0.8}      & \ip{255.0.0.0}     & \ip{10.0.0.254}     & — \\
    R1 — G0/0 & \ip{192.168.80.254} & \ip{255.255.255.0} & — & Off \\
    R1 — G0/1 & \ip{10.0.0.254}     & \ip{255.0.0.0}     & — & On \\}
  \def\Csol{%
    \hl 1 & \sol{PC-A} & \sol{IP hors du réseau 192.168.80.0/24.} & \sol{Ne joint ni les postes ni le routeur.} & \sol{Ex. 192.168.80.11.} \\ \hline
    \hl 2 & \sol{PC-C} & \sol{Passerelle .80.100 : pas l'adresse du routeur.} & \sol{Joint son réseau, mais aucun autre (serveur).} & \sol{Passerelle 192.168.80.254.} \\ \hline
    \hl 3 & \sol{PC-D} & \sol{Même IP que PC-B (doublon).} & \sol{Conflit d'adresse : PC-B et PC-D perturbés.} & \sol{Adresse unique, ex. .80.14.} \\ \hline
    \hl 4 & \sol{R1 G0/0} & \sol{Interface côté LAN désactivée (Off).} & \sol{Aucun poste ne joint le routeur ni le serveur.} & \sol{Port Status : On.} \\ \hline}
  \def\Cbareme{PC-B et SRV-C sont correctement configurés ; accepter PC-B cité à la place de PC-D pour le doublon.}
\fi

\pagestyle{fancy}
\fancyhf{}
\renewcommand{\headrulewidth}{0pt}
\fancyfoot[L]{\footnotesize\itshape\color{grisfoot}Évaluation — Réseaux locaux — Sujet \variante\ifcorrige\ — CORRIGÉ\fi}
\fancyfoot[R]{\footnotesize\color{grisfoot}\thepage/3}

\setlist{itemsep=0pt,topsep=2pt,parsep=0pt}
\setlength{\parindent}{0pt}
\setlength{\parskip}{2pt}

\begin{document}

\begin{center}
  {\Large\bfseries Évaluation pratique — Réseaux locaux \quad
   \fcolorbox{black}{white}{\ SUJET \variante\ }\par}
  \vspace{1pt}
  {\bfseries Adressage IP, commutation et routage — BTS CIEL 1re année — durée : 2 heures\par}
  \ifcorrige\vspace{2pt}{\bfseries\color{corr}CORRIGÉ ET BARÈME — document enseignant}\par\fi
\end{center}
\vspace{-2pt}
{\renewcommand{\arraystretch}{1.35}
\begin{tabularx}{\linewidth}{|X|X|p{2.2cm}|p{2.6cm}|}
  \hline
  \textbf{Nom :} & \textbf{Prénom :} & \textbf{Poste n°} & \textbf{Note :} \hfill /20 \\ \hline
\end{tabularx}}

\textbf{Consignes :} travail \textbf{individuel}, documents et Internet \textbf{non autorisés}. Ce sujet est le
\textbf{document-réponse}. Fichier Packet Tracer : \ip{Nom\_Prenom\_EvalLAN\_\variante.pkt}, à enregistrer
régulièrement dans le dossier de rendu. Les « visas » sont apposés par le professeur, que vous appelez à chaque test réussi.
Durées conseillées : A 30 min — B 1 h 10 — C 20 min.

%==========================================================================
\section{Partie A — Questions de cours et calculs \hfill /6}

\question{A1}{2 pts}{Complétez le tableau en considérant le \textbf{masque par défaut} de la classe de chaque adresse.}
{\small\renewcommand{\arraystretch}{1.15}
\begin{tabularx}{\linewidth}{|l|c|C|C|C|c|}
  \hline
  \textbf{Adresse IP} & \textbf{Classe} & \textbf{Masque par défaut} & \textbf{Adresse réseau} &
  \textbf{Adresse de diffusion} & \textbf{Nb d'hôtes} \\ \hline
  \AIrows
\end{tabularx}}
\bareme{0,5 pt par ligne juste + 0,5 pt si les trois nombres d'hôtes sont justes ($2^{n}-2$, $n$ = nombre de bits d'hôte).}

\question{A2}{1 pt}{Indiquez le type de câble (\textbf{droit} ou \textbf{croisé}) pour chaque liaison.}
{\small\renewcommand{\arraystretch}{1.15}
\begin{tabularx}{\linewidth}{|X|c||X|c|}
  \hline
  \hl PC $\leftrightarrow$ commutateur & \makebox[1.6cm]{\sol{droit}} & Commutateur $\leftrightarrow$ routeur & \makebox[1.6cm]{\sol{droit}} \\ \hline
  \hl PC $\leftrightarrow$ PC (réseau P2P) & \sol{croisé} & Concentrateur $\leftrightarrow$ concentrateur & \sol{croisé} \\ \hline
\end{tabularx}}
\bareme{0,25 pt par case. Équipements de types différents : câble droit ; de même type : câble croisé.}

\question{A3}{1 pt}{Comment un \textbf{concentrateur (hub)} et un \textbf{commutateur (switch)} traitent-ils une trame
reçue ? Lequel limite les collisions, et pourquoi ?}
\zone{2.9cm}{Le \textbf{hub} (couche 1) répète la trame sur \textbf{tous ses autres ports} : un seul domaine de collision
partagé. Le \textbf{switch} (couche 2) lit l'adresse MAC de destination et, grâce à sa \textbf{table MAC}, ne l'envoie
\textbf{que sur le port du destinataire} : un domaine de collision par port, c'est lui qui limite les collisions.
\emph{0,5 pt hub, 0,5 pt switch.}}

\question{A4}{1 pt}{Qu'est-ce que la \textbf{passerelle par défaut} d'un poste ? Quelle condition son adresse doit-elle respecter ?}
\zone{2.4cm}{Adresse IP de l'interface du \textbf{routeur} à laquelle le poste envoie les paquets destinés à un
\textbf{autre réseau}. Elle doit \textbf{appartenir au même réseau} que le poste. \emph{0,5 + 0,5 pt.}}

\question{A5}{0,5 pt}{Un portable configuré en DHCP affiche \ip{\apipa} / \ip{255.255.0.0}. Que s'est-il passé ?}
\zone{1.5cm}{Aucun \textbf{serveur DHCP} n'a répondu : le poste s'est attribué une adresse \textbf{APIPA} (169.254.0.0/16).}

\question{A6}{0,5 pt}{Quel protocole utilise \texttt{ping} ? Nommez les deux messages échangés et la couche OSI des adresses IP.}
\zone{1.5cm}{\textbf{ICMP} ; \textbf{Echo Request} / \textbf{Echo Reply} ; adresses IP en \textbf{couche 3 (Réseau)}.}

%==========================================================================
\newpage
\section{Partie B — Réalisation sous Packet Tracer \hfill /10}

\begin{cdc}
\entreprise{} vous confie son réseau : \textbf{trois réseaux} reliés par un \textbf{routeur unique Cisco 2911}
nommé \ip{\Rnom} (interfaces G0/0, G0/1, G0/2).
\begin{itemize}[leftmargin=12pt]
  \item \textbf{Réseau \resUn} : \textbf{4 postes} PC-\pUn1 à PC-\pUn4, commutateur \ip{SW-\pUn}, adressage statique,
        réseau \ip{\netUn.0/24}.
  \item \textbf{Réseau \resDeux} : \textbf{2 postes} PC-\pDeux1, PC-\pDeux2 et \textbf{1 imprimante} IMP-\pDeux,
        commutateur \ip{SW-\pDeux}, adressage statique, réseau de classe B \ip{\netDeux.0.0/16}.
  \item \textbf{Réseau Serveur} : serveur \ip{\srvNom} d'adresse \ip{\netSrv.\srvH/24}, relié \textbf{directement} au routeur.
  \item La passerelle de chaque réseau est sa \textbf{dernière adresse utilisable}. Tous les équipements communiquent
        entre eux et accèdent à la page web du serveur. Nommage imposé (« Display Name »).
\end{itemize}
\end{cdc}

\question{B1}{2 pts}{Complétez le plan d'adressage \textbf{avant} la réalisation.}
{\footnotesize\setlength{\tabcolsep}{3pt}\renewcommand{\arraystretch}{1.15}
\begin{tabularx}{\linewidth}{|l|C|C|C|C|C|C|}
  \hline
  \textbf{Réseau} & \textbf{Adresse réseau} & \textbf{Masque} & \textbf{1re @ utilisable} &
  \textbf{Dernière @ utilisable} & \textbf{Diffusion} & \textbf{Passerelle} \\ \hline
  \hl \resUn & \sol{\netUn.0} & \sol{255.255.255.0} & \sol{\netUn.1} & \sol{\netUn.254} & \sol{\netUn.255} & \sol{\netUn.254} \\ \hline
  \hl \resDeux & \sol{\netDeux.0.0} & \sol{255.255.0.0} & \sol{\netDeux.0.1} & \sol{\netDeux.255.254} & \sol{\netDeux.255.255} & \sol{\netDeux.255.254} \\ \hline
  \hl Serveur & \sol{\netSrv.0} & \sol{255.255.255.0} & \sol{\netSrv.1} & \sol{\netSrv.254} & \sol{\netSrv.255} & \sol{\netSrv.254} \\ \hline
\end{tabularx}}
\bareme{0,5 pt par ligne juste + 0,5 pt si la convention de passerelle est respectée partout.}

\question{B2}{2 pts}{Construisez la topologie (équipements, câblage, nommage) et relevez l'adressage configuré.}
{\footnotesize\setlength{\tabcolsep}{3pt}\renewcommand{\arraystretch}{1.1}
\begin{tabularx}{\linewidth}{|l|C|C|C||l|C|C|}
  \hline
  \textbf{Hôte} & \textbf{Adresse IP} & \textbf{Masque} & \textbf{Passerelle} & \textbf{Interface} & \textbf{Adresse IP} & \textbf{Masque} \\ \hline
  \hl PC-\pUn1 & \sol{\netUn.1} & \sol{255.255.255.0} & \sol{\netUn.254} & \Rnom{} G0/0 & \sol{\netUn.254} & \sol{255.255.255.0} \\ \hline
  \hl PC-\pDeux1 & \sol{\netDeux.0.1} & \sol{255.255.0.0} & \sol{\netDeux.255.254} & \Rnom{} G0/1 & \sol{\netDeux.255.254} & \sol{255.255.0.0} \\ \hline
  \hl IMP-\pDeux & \sol{\netDeux.0.10} & \sol{255.255.0.0} & \sol{\netDeux.255.254} & \Rnom{} G0/2 & \sol{\netSrv.254} & \sol{255.255.255.0} \\ \hline
  \hl \srvNom & \ip{\netSrv.\srvH} & \sol{255.255.255.0} & \sol{\netSrv.254} & \multicolumn{3}{c|}{\cellcolor{black!8}} \\ \hline
\end{tabularx}}
\bareme{Adresses d'hôtes libres (uniques, dans le bon réseau) ; affectation G0/x libre. Sur le .pkt : équipements 0,5 ;
câblage 1 (câbles droits vers les switchs et switch–routeur, \textbf{croisé} serveur–routeur, liens verts) ; nommage 0,5.}

\question{B3}{3 pts}{Configurez tous les équipements : IP, masques, passerelles des hôtes ; interfaces du routeur
(« Port Status : On »).}
\bareme{Sur le .pkt : interfaces routeur 1 pt ; IP/masques des hôtes 1 pt ; passerelles (imprimante et serveur compris)
1 pt. $-0,25$ par erreur, sans descendre sous 0 par item.}

\question{B4}{3 pts}{Réalisez les tests, notez la commande et le résultat, puis faites viser chaque test réussi.}
{\small\renewcommand{\arraystretch}{1.1}
\begin{tabularx}{\linewidth}{|c|l|l|C|c|c|}
  \hline
  \textbf{N°} & \textbf{Source} & \textbf{Destination} & \textbf{Commande / action} & \textbf{Résultat} & \textbf{Visa prof.} \\ \hline
  \hl 1 & PC-\pUn1 & PC-\pUn4 & \sol{ping \netUn.4} & \sol{réussi} & \hspace{1.4cm} \\ \hline
  \hl 2 & PC-\pDeux1 & IMP-\pDeux & \sol{ping \netDeux.0.10} & \sol{réussi} & \\ \hline
  \hl 3 & PC-\pUn2 & PC-\pDeux2 & \sol{ping \netDeux.0.2} & \sol{réussi} & \\ \hline
  \hl 4 & PC-\pDeux2 & \srvNom & \sol{ping \netSrv.\srvH} & \sol{réussi} & \\ \hline
  \hl 5 & PC-\pUn3 & page web du serveur & \sol{navigateur : {\NoAutoSpacing http://\netSrv.\srvH}} & \sol{page affichée} & \\ \hline
\end{tabularx}}
\bareme{0,5 pt par test visé + 0,5 pt pour B5. Le premier ping inter-réseaux peut échouer (résolution ARP) ; les suivants doivent réussir.}

\question{B5}{incl. B4}{Test n°3 : par quels équipements passe le paquet ? Pourquoi PC-\pUn2 l'envoie-t-il au routeur ?}
\zone{2.3cm}{PC-\pUn2 $\to$ SW-\pUn{} $\to$ \Rnom{} $\to$ SW-\pDeux{} $\to$ PC-\pDeux2. Grâce à son masque, PC-\pUn2 constate que
la destination (\netDeux.0.0/16) n'est \textbf{pas dans son réseau} (\netUn.0/24) : il envoie le paquet à sa
\textbf{passerelle par défaut}, qui le route.}

%==========================================================================
\newpage
\section{Partie C — Diagnostic d'un réseau en panne \hfill /4}

Un technicien constate qu'\textbf{aucun poste ne parvient à joindre le serveur SRV-C} et que certains postes
communiquent mal entre eux. Les câbles sont du bon type.

\begin{center}
\begin{minipage}{0.37\linewidth}\centering
\begin{tikzpicture}[font=\footnotesize,
  eq/.style={draw,fill=qbg,rounded corners=2pt,minimum width=9mm,minimum height=5mm,inner sep=2pt},
  every path/.style={thick}]
  \node[eq] (SW) {SW-1};
  \node[eq,right=9mm of SW] (R) {R1};
  \node[eq,right=7mm of R] (SRV) {SRV-C};
  \node[eq,above left=5mm and -3mm of SW] (A) {PC-A};
  \node[eq,above right=5mm and -3mm of SW] (B) {PC-B};
  \node[eq,below left=5mm and -3mm of SW] (C) {PC-C};
  \node[eq,below right=5mm and -3mm of SW] (D) {PC-D};
  \foreach \n in {A,B,C,D} \draw (\n) -- (SW);
  \draw (SW) -- node[above,font=\tiny]{G0/0} (R);
  \draw (R) -- node[above,font=\tiny]{G0/1} (SRV);
\end{tikzpicture}
\end{minipage}\hfill
\begin{minipage}{0.61\linewidth}\centering
{\small\renewcommand{\arraystretch}{1.05}\setlength{\tabcolsep}{4pt}
\begin{tabular}{lllll}
  \toprule
  \textbf{Équipement} & \textbf{Adresse IP} & \textbf{Masque} & \textbf{Passerelle} & \textbf{Port} \\
  \midrule
  \Crows
  \bottomrule
\end{tabular}}
\end{minipage}
\end{center}

\question{C1}{4 pts}{Relevez les \textbf{quatre anomalies}. Pour chacune : conséquence et correction.}
{\small\renewcommand{\arraystretch}{1.1}
\ifcorrige\def\hc{}\else\def\hc{\rule[-10mm]{0pt}{14mm}}\fi
\begin{tabularx}{\linewidth}{|c|>{\raggedright\arraybackslash}p{2cm}|L|L|L|}
  \hline
  \textbf{N°} & \textbf{Équipement} & \textbf{Anomalie} & \textbf{Conséquence} & \textbf{Correction} \\ \hline
  \ifcorrige\Csol\else
  \hc 1 & & & & \\ \hline
  \hc 2 & & & & \\ \hline
  \hc 3 & & & & \\ \hline
  \hc 4 & & & & \\ \hline
  \fi
\end{tabularx}}
\bareme{1 pt par anomalie (0,25 équipement, anomalie, conséquence, correction ; ordre indifférent). \Cbareme}

%==========================================================================
\section{Bonus — Extension sans fil \hfill +1}

\begin{cdc}Ajoutez au réseau \resUn{} un \textbf{AccessPoint-PT} relié à SW-\pUn{} (SSID \ip{\ssid}, \textbf{WPA2-PSK})
et un portable \ip{Laptop-\pUn} équipé d'un module \texttt{WPC300N}, en adressage \textbf{statique}.
Vérifiez qu'il joint \srvNom.\end{cdc}
{\small\renewcommand{\arraystretch}{1.1}
\begin{tabularx}{\linewidth}{|l|C|C|C|c|}
  \hline
  \textbf{Équipement} & \textbf{Adresse IP} & \textbf{Masque} & \textbf{Passerelle} & \textbf{Visa prof.} \\ \hline
  \hl Laptop-\pUn & \sol{\netUn.50} & \sol{255.255.255.0} & \sol{\netUn.254} & \hspace{1.4cm} \\ \hline
\end{tabularx}}
\bareme{Le point d'accès fait office de pont : le portable est dans \netUn.0/24. Bonus si associé en WPA2-PSK et ping
vers \netSrv.\srvH{} réussi. Note plafonnée à 20.}

\end{document}
"
%   pdflatex -jobname=Eval_Reseau_CorrigeA "\def\variante{A}\def\modecorrige{1}% Évaluation pratique — Réseaux locaux (BTS CIEL 1re année) — sujets A et B
% Compilation (4 PDF) :
%   pdflatex -jobname=Eval_Reseau_SujetA   "\def\variante{A}\def\modecorrige{0}\input{Eval_TP_Reseau_LAN.tex}"
%   pdflatex -jobname=Eval_Reseau_SujetB   "\def\variante{B}\def\modecorrige{0}\input{Eval_TP_Reseau_LAN.tex}"
%   pdflatex -jobname=Eval_Reseau_CorrigeA "\def\variante{A}\def\modecorrige{1}\input{Eval_TP_Reseau_LAN.tex}"
%   pdflatex -jobname=Eval_Reseau_CorrigeB "\def\variante{B}\def\modecorrige{1}\input{Eval_TP_Reseau_LAN.tex}"
\documentclass[10pt,a4paper]{article}

\usepackage[T1]{fontenc}
\usepackage[utf8]{inputenc}
\usepackage{lmodern}
\usepackage[french]{babel}
\usepackage[a4paper,hmargin=1.5cm,top=1.3cm,bottom=1.6cm,footskip=0.7cm]{geometry}
\usepackage[table]{xcolor}
\usepackage{booktabs,tabularx,array}
\usepackage{enumitem}
\usepackage{fancyhdr}
\usepackage[most]{tcolorbox}
\usepackage{tikz}
\usetikzlibrary{positioning}
\usepackage{titlesec}

\providecommand{\modecorrige}{0}
\providecommand{\variante}{A}
\newif\ifcorrige
\ifnum\modecorrige=1 \corrigetrue \else \corrigefalse \fi

\definecolor{ipred}{HTML}{C0392B}
\definecolor{cdcbar}{HTML}{A67C3D}
\definecolor{cdcbg}{HTML}{F4F1EA}
\definecolor{qbar}{HTML}{4A6B5D}
\definecolor{qbg}{HTML}{EDF1EF}
\definecolor{corr}{HTML}{1F5FA8}
\definecolor{corrbg}{HTML}{E8F0FA}
\definecolor{grisfoot}{HTML}{828282}

\newcommand{\ip}[1]{\textcolor{ipred}{\texttt{#1}}}
\newcommand{\pts}[1]{\hfill\mbox{\textbf{(#1)}}}
\newcommand{\sol}[1]{\ifcorrige\textcolor{corr}{#1}\fi}
\newcommand{\bareme}[1]{\ifcorrige\par{\footnotesize\color{corr}#1\par}\fi}
\newcommand{\hl}{\ifcorrige\rule[-1.6mm]{0pt}{5.4mm}\else\rule[-2.2mm]{0pt}{6.6mm}\fi}
\newcolumntype{C}{>{\centering\arraybackslash}X}
\newcolumntype{L}{>{\raggedright\arraybackslash}X}

\titleformat{\section}{\large\bfseries}{}{0pt}{}
\titlespacing*{\section}{0pt}{6pt}{3pt}

\tcbset{qstyle/.style={enhanced,boxrule=0pt,leftrule=2.5pt,arc=0pt,outer arc=0pt,
  left=5pt,right=5pt,top=2pt,bottom=2pt,before skip=4pt,after skip=3pt}}
\newtcolorbox{cdc}{qstyle,colback=cdcbg,colframe=cdcbar}
\newtcolorbox{quest}{qstyle,colback=qbg,colframe=qbar}

% Zone de réponse : vierge dans le sujet, remplie dans le corrigé
\newcommand{\zone}[2]{%
  \ifcorrige
    \begin{tcolorbox}[colback=corrbg,colframe=corr,boxrule=0.5pt,arc=1pt,left=4pt,right=4pt,top=1pt,bottom=1pt,
      before skip=2pt,after skip=2pt,fontupper=\footnotesize\color{corr}]#2\end{tcolorbox}%
  \else
    \begin{tcolorbox}[colback=white,colframe=black!45,boxrule=0.4pt,arc=1pt,height=#1,
      before skip=2pt,after skip=2pt]\end{tcolorbox}%
  \fi}

\newcommand{\question}[3]{\begin{quest}\textbf{#1.} #3\pts{#2}\end{quest}}

%==========================================================================
% Données propres à chaque variante
%==========================================================================
\if\variante A
  \def\entreprise{Le cabinet d'architectes \textbf{ArchiNova}}
  \def\Rnom{R-ArchiNova}
  \def\resUn{Bureau d'études}\def\pUn{BE}\def\netUn{192.168.40}
  \def\resDeux{Accueil}\def\pDeux{ACC}\def\netDeux{172.16}
  \def\srvNom{SRV-Intranet}\def\netSrv{192.168.200}\def\srvH{10}
  \def\ssid{ArchiNova-WiFi}
  \def\AIrows{%
    \hl\ip{192.168.40.25} & \sol{C} & \sol{255.255.255.0} & \sol{192.168.40.0} & \sol{192.168.40.255} & \sol{254} \\ \hline
    \hl\ip{172.16.5.10}   & \sol{B} & \sol{255.255.0.0}   & \sol{172.16.0.0}   & \sol{172.16.255.255} & \sol{65\,534} \\ \hline
    \hl\ip{10.3.2.1}      & \sol{A} & \sol{255.0.0.0}     & \sol{10.0.0.0}     & \sol{10.255.255.255} & \sol{16\,777\,214} \\ \hline}
  \def\apipa{169.254.12.7}
  \def\Cnet{192.168.60}
  \def\Crows{%
    PC-A & \ip{192.168.60.10} & \ip{255.255.255.0} & \ip{192.168.60.254} & — \\
    PC-B & \ip{192.168.60.20} & \ip{255.255.255.0} & \ip{192.168.60.1}   & — \\
    PC-C & \ip{192.168.61.30} & \ip{255.255.255.0} & \ip{192.168.60.254} & — \\
    PC-D & \ip{192.168.60.10} & \ip{255.255.255.0} & \ip{192.168.60.254} & — \\
    SRV-C & \ip{10.0.0.5}     & \ip{255.0.0.0}     & \ip{10.0.0.254}     & — \\
    R1 — G0/0 & \ip{192.168.60.254} & \ip{255.255.255.0} & — & On \\
    R1 — G0/1 & \ip{10.0.0.254}     & \ip{255.0.0.0}     & — & Off \\}
  \def\Csol{%
    \hl 1 & \sol{PC-B} & \sol{Passerelle .60.1 : pas l'adresse du routeur.} & \sol{Joint son réseau, mais aucun autre (serveur).} & \sol{Passerelle 192.168.60.254.} \\ \hline
    \hl 2 & \sol{PC-C} & \sol{IP hors du réseau 192.168.60.0/24.} & \sol{Ne joint ni les postes ni le routeur.} & \sol{Ex. 192.168.60.30.} \\ \hline
    \hl 3 & \sol{PC-D} & \sol{Même IP que PC-A (doublon).} & \sol{Conflit d'adresse : PC-A et PC-D perturbés.} & \sol{Adresse unique, ex. .60.40.} \\ \hline
    \hl 4 & \sol{R1 G0/1} & \sol{Interface désactivée (Off).} & \sol{Réseau du serveur injoignable pour tous.} & \sol{Port Status : On.} \\ \hline}
  \def\Cbareme{PC-A et SRV-C sont correctement configurés ; accepter PC-A cité à la place de PC-D pour le doublon.}
\else
  \def\entreprise{La clinique vétérinaire \textbf{VetoPlus}}
  \def\Rnom{R-VetoPlus}
  \def\resUn{Soins}\def\pUn{SOI}\def\netUn{192.168.70}
  \def\resDeux{Secrétariat}\def\pDeux{SEC}\def\netDeux{172.20}
  \def\srvNom{SRV-Dossiers}\def\netSrv{192.168.250}\def\srvH{20}
  \def\ssid{VetoPlus-WiFi}
  \def\AIrows{%
    \hl\ip{192.168.15.130} & \sol{C} & \sol{255.255.255.0} & \sol{192.168.15.0} & \sol{192.168.15.255} & \sol{254} \\ \hline
    \hl\ip{172.31.200.7}   & \sol{B} & \sol{255.255.0.0}   & \sol{172.31.0.0}   & \sol{172.31.255.255} & \sol{65\,534} \\ \hline
    \hl\ip{10.45.6.9}      & \sol{A} & \sol{255.0.0.0}     & \sol{10.0.0.0}     & \sol{10.255.255.255} & \sol{16\,777\,214} \\ \hline}
  \def\apipa{169.254.88.3}
  \def\Cnet{192.168.80}
  \def\Crows{%
    PC-A & \ip{192.168.8.11}   & \ip{255.255.255.0} & \ip{192.168.80.254} & — \\
    PC-B & \ip{192.168.80.12}  & \ip{255.255.255.0} & \ip{192.168.80.254} & — \\
    PC-C & \ip{192.168.80.13}  & \ip{255.255.255.0} & \ip{192.168.80.100} & — \\
    PC-D & \ip{192.168.80.12}  & \ip{255.255.255.0} & \ip{192.168.80.254} & — \\
    SRV-C & \ip{10.0.0.8}      & \ip{255.0.0.0}     & \ip{10.0.0.254}     & — \\
    R1 — G0/0 & \ip{192.168.80.254} & \ip{255.255.255.0} & — & Off \\
    R1 — G0/1 & \ip{10.0.0.254}     & \ip{255.0.0.0}     & — & On \\}
  \def\Csol{%
    \hl 1 & \sol{PC-A} & \sol{IP hors du réseau 192.168.80.0/24.} & \sol{Ne joint ni les postes ni le routeur.} & \sol{Ex. 192.168.80.11.} \\ \hline
    \hl 2 & \sol{PC-C} & \sol{Passerelle .80.100 : pas l'adresse du routeur.} & \sol{Joint son réseau, mais aucun autre (serveur).} & \sol{Passerelle 192.168.80.254.} \\ \hline
    \hl 3 & \sol{PC-D} & \sol{Même IP que PC-B (doublon).} & \sol{Conflit d'adresse : PC-B et PC-D perturbés.} & \sol{Adresse unique, ex. .80.14.} \\ \hline
    \hl 4 & \sol{R1 G0/0} & \sol{Interface côté LAN désactivée (Off).} & \sol{Aucun poste ne joint le routeur ni le serveur.} & \sol{Port Status : On.} \\ \hline}
  \def\Cbareme{PC-B et SRV-C sont correctement configurés ; accepter PC-B cité à la place de PC-D pour le doublon.}
\fi

\pagestyle{fancy}
\fancyhf{}
\renewcommand{\headrulewidth}{0pt}
\fancyfoot[L]{\footnotesize\itshape\color{grisfoot}Évaluation — Réseaux locaux — Sujet \variante\ifcorrige\ — CORRIGÉ\fi}
\fancyfoot[R]{\footnotesize\color{grisfoot}\thepage/3}

\setlist{itemsep=0pt,topsep=2pt,parsep=0pt}
\setlength{\parindent}{0pt}
\setlength{\parskip}{2pt}

\begin{document}

\begin{center}
  {\Large\bfseries Évaluation pratique — Réseaux locaux \quad
   \fcolorbox{black}{white}{\ SUJET \variante\ }\par}
  \vspace{1pt}
  {\bfseries Adressage IP, commutation et routage — BTS CIEL 1re année — durée : 2 heures\par}
  \ifcorrige\vspace{2pt}{\bfseries\color{corr}CORRIGÉ ET BARÈME — document enseignant}\par\fi
\end{center}
\vspace{-2pt}
{\renewcommand{\arraystretch}{1.35}
\begin{tabularx}{\linewidth}{|X|X|p{2.2cm}|p{2.6cm}|}
  \hline
  \textbf{Nom :} & \textbf{Prénom :} & \textbf{Poste n°} & \textbf{Note :} \hfill /20 \\ \hline
\end{tabularx}}

\textbf{Consignes :} travail \textbf{individuel}, documents et Internet \textbf{non autorisés}. Ce sujet est le
\textbf{document-réponse}. Fichier Packet Tracer : \ip{Nom\_Prenom\_EvalLAN\_\variante.pkt}, à enregistrer
régulièrement dans le dossier de rendu. Les « visas » sont apposés par le professeur, que vous appelez à chaque test réussi.
Durées conseillées : A 30 min — B 1 h 10 — C 20 min.

%==========================================================================
\section{Partie A — Questions de cours et calculs \hfill /6}

\question{A1}{2 pts}{Complétez le tableau en considérant le \textbf{masque par défaut} de la classe de chaque adresse.}
{\small\renewcommand{\arraystretch}{1.15}
\begin{tabularx}{\linewidth}{|l|c|C|C|C|c|}
  \hline
  \textbf{Adresse IP} & \textbf{Classe} & \textbf{Masque par défaut} & \textbf{Adresse réseau} &
  \textbf{Adresse de diffusion} & \textbf{Nb d'hôtes} \\ \hline
  \AIrows
\end{tabularx}}
\bareme{0,5 pt par ligne juste + 0,5 pt si les trois nombres d'hôtes sont justes ($2^{n}-2$, $n$ = nombre de bits d'hôte).}

\question{A2}{1 pt}{Indiquez le type de câble (\textbf{droit} ou \textbf{croisé}) pour chaque liaison.}
{\small\renewcommand{\arraystretch}{1.15}
\begin{tabularx}{\linewidth}{|X|c||X|c|}
  \hline
  \hl PC $\leftrightarrow$ commutateur & \makebox[1.6cm]{\sol{droit}} & Commutateur $\leftrightarrow$ routeur & \makebox[1.6cm]{\sol{droit}} \\ \hline
  \hl PC $\leftrightarrow$ PC (réseau P2P) & \sol{croisé} & Concentrateur $\leftrightarrow$ concentrateur & \sol{croisé} \\ \hline
\end{tabularx}}
\bareme{0,25 pt par case. Équipements de types différents : câble droit ; de même type : câble croisé.}

\question{A3}{1 pt}{Comment un \textbf{concentrateur (hub)} et un \textbf{commutateur (switch)} traitent-ils une trame
reçue ? Lequel limite les collisions, et pourquoi ?}
\zone{2.9cm}{Le \textbf{hub} (couche 1) répète la trame sur \textbf{tous ses autres ports} : un seul domaine de collision
partagé. Le \textbf{switch} (couche 2) lit l'adresse MAC de destination et, grâce à sa \textbf{table MAC}, ne l'envoie
\textbf{que sur le port du destinataire} : un domaine de collision par port, c'est lui qui limite les collisions.
\emph{0,5 pt hub, 0,5 pt switch.}}

\question{A4}{1 pt}{Qu'est-ce que la \textbf{passerelle par défaut} d'un poste ? Quelle condition son adresse doit-elle respecter ?}
\zone{2.4cm}{Adresse IP de l'interface du \textbf{routeur} à laquelle le poste envoie les paquets destinés à un
\textbf{autre réseau}. Elle doit \textbf{appartenir au même réseau} que le poste. \emph{0,5 + 0,5 pt.}}

\question{A5}{0,5 pt}{Un portable configuré en DHCP affiche \ip{\apipa} / \ip{255.255.0.0}. Que s'est-il passé ?}
\zone{1.5cm}{Aucun \textbf{serveur DHCP} n'a répondu : le poste s'est attribué une adresse \textbf{APIPA} (169.254.0.0/16).}

\question{A6}{0,5 pt}{Quel protocole utilise \texttt{ping} ? Nommez les deux messages échangés et la couche OSI des adresses IP.}
\zone{1.5cm}{\textbf{ICMP} ; \textbf{Echo Request} / \textbf{Echo Reply} ; adresses IP en \textbf{couche 3 (Réseau)}.}

%==========================================================================
\newpage
\section{Partie B — Réalisation sous Packet Tracer \hfill /10}

\begin{cdc}
\entreprise{} vous confie son réseau : \textbf{trois réseaux} reliés par un \textbf{routeur unique Cisco 2911}
nommé \ip{\Rnom} (interfaces G0/0, G0/1, G0/2).
\begin{itemize}[leftmargin=12pt]
  \item \textbf{Réseau \resUn} : \textbf{4 postes} PC-\pUn1 à PC-\pUn4, commutateur \ip{SW-\pUn}, adressage statique,
        réseau \ip{\netUn.0/24}.
  \item \textbf{Réseau \resDeux} : \textbf{2 postes} PC-\pDeux1, PC-\pDeux2 et \textbf{1 imprimante} IMP-\pDeux,
        commutateur \ip{SW-\pDeux}, adressage statique, réseau de classe B \ip{\netDeux.0.0/16}.
  \item \textbf{Réseau Serveur} : serveur \ip{\srvNom} d'adresse \ip{\netSrv.\srvH/24}, relié \textbf{directement} au routeur.
  \item La passerelle de chaque réseau est sa \textbf{dernière adresse utilisable}. Tous les équipements communiquent
        entre eux et accèdent à la page web du serveur. Nommage imposé (« Display Name »).
\end{itemize}
\end{cdc}

\question{B1}{2 pts}{Complétez le plan d'adressage \textbf{avant} la réalisation.}
{\footnotesize\setlength{\tabcolsep}{3pt}\renewcommand{\arraystretch}{1.15}
\begin{tabularx}{\linewidth}{|l|C|C|C|C|C|C|}
  \hline
  \textbf{Réseau} & \textbf{Adresse réseau} & \textbf{Masque} & \textbf{1re @ utilisable} &
  \textbf{Dernière @ utilisable} & \textbf{Diffusion} & \textbf{Passerelle} \\ \hline
  \hl \resUn & \sol{\netUn.0} & \sol{255.255.255.0} & \sol{\netUn.1} & \sol{\netUn.254} & \sol{\netUn.255} & \sol{\netUn.254} \\ \hline
  \hl \resDeux & \sol{\netDeux.0.0} & \sol{255.255.0.0} & \sol{\netDeux.0.1} & \sol{\netDeux.255.254} & \sol{\netDeux.255.255} & \sol{\netDeux.255.254} \\ \hline
  \hl Serveur & \sol{\netSrv.0} & \sol{255.255.255.0} & \sol{\netSrv.1} & \sol{\netSrv.254} & \sol{\netSrv.255} & \sol{\netSrv.254} \\ \hline
\end{tabularx}}
\bareme{0,5 pt par ligne juste + 0,5 pt si la convention de passerelle est respectée partout.}

\question{B2}{2 pts}{Construisez la topologie (équipements, câblage, nommage) et relevez l'adressage configuré.}
{\footnotesize\setlength{\tabcolsep}{3pt}\renewcommand{\arraystretch}{1.1}
\begin{tabularx}{\linewidth}{|l|C|C|C||l|C|C|}
  \hline
  \textbf{Hôte} & \textbf{Adresse IP} & \textbf{Masque} & \textbf{Passerelle} & \textbf{Interface} & \textbf{Adresse IP} & \textbf{Masque} \\ \hline
  \hl PC-\pUn1 & \sol{\netUn.1} & \sol{255.255.255.0} & \sol{\netUn.254} & \Rnom{} G0/0 & \sol{\netUn.254} & \sol{255.255.255.0} \\ \hline
  \hl PC-\pDeux1 & \sol{\netDeux.0.1} & \sol{255.255.0.0} & \sol{\netDeux.255.254} & \Rnom{} G0/1 & \sol{\netDeux.255.254} & \sol{255.255.0.0} \\ \hline
  \hl IMP-\pDeux & \sol{\netDeux.0.10} & \sol{255.255.0.0} & \sol{\netDeux.255.254} & \Rnom{} G0/2 & \sol{\netSrv.254} & \sol{255.255.255.0} \\ \hline
  \hl \srvNom & \ip{\netSrv.\srvH} & \sol{255.255.255.0} & \sol{\netSrv.254} & \multicolumn{3}{c|}{\cellcolor{black!8}} \\ \hline
\end{tabularx}}
\bareme{Adresses d'hôtes libres (uniques, dans le bon réseau) ; affectation G0/x libre. Sur le .pkt : équipements 0,5 ;
câblage 1 (câbles droits vers les switchs et switch–routeur, \textbf{croisé} serveur–routeur, liens verts) ; nommage 0,5.}

\question{B3}{3 pts}{Configurez tous les équipements : IP, masques, passerelles des hôtes ; interfaces du routeur
(« Port Status : On »).}
\bareme{Sur le .pkt : interfaces routeur 1 pt ; IP/masques des hôtes 1 pt ; passerelles (imprimante et serveur compris)
1 pt. $-0,25$ par erreur, sans descendre sous 0 par item.}

\question{B4}{3 pts}{Réalisez les tests, notez la commande et le résultat, puis faites viser chaque test réussi.}
{\small\renewcommand{\arraystretch}{1.1}
\begin{tabularx}{\linewidth}{|c|l|l|C|c|c|}
  \hline
  \textbf{N°} & \textbf{Source} & \textbf{Destination} & \textbf{Commande / action} & \textbf{Résultat} & \textbf{Visa prof.} \\ \hline
  \hl 1 & PC-\pUn1 & PC-\pUn4 & \sol{ping \netUn.4} & \sol{réussi} & \hspace{1.4cm} \\ \hline
  \hl 2 & PC-\pDeux1 & IMP-\pDeux & \sol{ping \netDeux.0.10} & \sol{réussi} & \\ \hline
  \hl 3 & PC-\pUn2 & PC-\pDeux2 & \sol{ping \netDeux.0.2} & \sol{réussi} & \\ \hline
  \hl 4 & PC-\pDeux2 & \srvNom & \sol{ping \netSrv.\srvH} & \sol{réussi} & \\ \hline
  \hl 5 & PC-\pUn3 & page web du serveur & \sol{navigateur : {\NoAutoSpacing http://\netSrv.\srvH}} & \sol{page affichée} & \\ \hline
\end{tabularx}}
\bareme{0,5 pt par test visé + 0,5 pt pour B5. Le premier ping inter-réseaux peut échouer (résolution ARP) ; les suivants doivent réussir.}

\question{B5}{incl. B4}{Test n°3 : par quels équipements passe le paquet ? Pourquoi PC-\pUn2 l'envoie-t-il au routeur ?}
\zone{2.3cm}{PC-\pUn2 $\to$ SW-\pUn{} $\to$ \Rnom{} $\to$ SW-\pDeux{} $\to$ PC-\pDeux2. Grâce à son masque, PC-\pUn2 constate que
la destination (\netDeux.0.0/16) n'est \textbf{pas dans son réseau} (\netUn.0/24) : il envoie le paquet à sa
\textbf{passerelle par défaut}, qui le route.}

%==========================================================================
\newpage
\section{Partie C — Diagnostic d'un réseau en panne \hfill /4}

Un technicien constate qu'\textbf{aucun poste ne parvient à joindre le serveur SRV-C} et que certains postes
communiquent mal entre eux. Les câbles sont du bon type.

\begin{center}
\begin{minipage}{0.37\linewidth}\centering
\begin{tikzpicture}[font=\footnotesize,
  eq/.style={draw,fill=qbg,rounded corners=2pt,minimum width=9mm,minimum height=5mm,inner sep=2pt},
  every path/.style={thick}]
  \node[eq] (SW) {SW-1};
  \node[eq,right=9mm of SW] (R) {R1};
  \node[eq,right=7mm of R] (SRV) {SRV-C};
  \node[eq,above left=5mm and -3mm of SW] (A) {PC-A};
  \node[eq,above right=5mm and -3mm of SW] (B) {PC-B};
  \node[eq,below left=5mm and -3mm of SW] (C) {PC-C};
  \node[eq,below right=5mm and -3mm of SW] (D) {PC-D};
  \foreach \n in {A,B,C,D} \draw (\n) -- (SW);
  \draw (SW) -- node[above,font=\tiny]{G0/0} (R);
  \draw (R) -- node[above,font=\tiny]{G0/1} (SRV);
\end{tikzpicture}
\end{minipage}\hfill
\begin{minipage}{0.61\linewidth}\centering
{\small\renewcommand{\arraystretch}{1.05}\setlength{\tabcolsep}{4pt}
\begin{tabular}{lllll}
  \toprule
  \textbf{Équipement} & \textbf{Adresse IP} & \textbf{Masque} & \textbf{Passerelle} & \textbf{Port} \\
  \midrule
  \Crows
  \bottomrule
\end{tabular}}
\end{minipage}
\end{center}

\question{C1}{4 pts}{Relevez les \textbf{quatre anomalies}. Pour chacune : conséquence et correction.}
{\small\renewcommand{\arraystretch}{1.1}
\ifcorrige\def\hc{}\else\def\hc{\rule[-10mm]{0pt}{14mm}}\fi
\begin{tabularx}{\linewidth}{|c|>{\raggedright\arraybackslash}p{2cm}|L|L|L|}
  \hline
  \textbf{N°} & \textbf{Équipement} & \textbf{Anomalie} & \textbf{Conséquence} & \textbf{Correction} \\ \hline
  \ifcorrige\Csol\else
  \hc 1 & & & & \\ \hline
  \hc 2 & & & & \\ \hline
  \hc 3 & & & & \\ \hline
  \hc 4 & & & & \\ \hline
  \fi
\end{tabularx}}
\bareme{1 pt par anomalie (0,25 équipement, anomalie, conséquence, correction ; ordre indifférent). \Cbareme}

%==========================================================================
\section{Bonus — Extension sans fil \hfill +1}

\begin{cdc}Ajoutez au réseau \resUn{} un \textbf{AccessPoint-PT} relié à SW-\pUn{} (SSID \ip{\ssid}, \textbf{WPA2-PSK})
et un portable \ip{Laptop-\pUn} équipé d'un module \texttt{WPC300N}, en adressage \textbf{statique}.
Vérifiez qu'il joint \srvNom.\end{cdc}
{\small\renewcommand{\arraystretch}{1.1}
\begin{tabularx}{\linewidth}{|l|C|C|C|c|}
  \hline
  \textbf{Équipement} & \textbf{Adresse IP} & \textbf{Masque} & \textbf{Passerelle} & \textbf{Visa prof.} \\ \hline
  \hl Laptop-\pUn & \sol{\netUn.50} & \sol{255.255.255.0} & \sol{\netUn.254} & \hspace{1.4cm} \\ \hline
\end{tabularx}}
\bareme{Le point d'accès fait office de pont : le portable est dans \netUn.0/24. Bonus si associé en WPA2-PSK et ping
vers \netSrv.\srvH{} réussi. Note plafonnée à 20.}

\end{document}
"
%   pdflatex -jobname=Eval_Reseau_CorrigeB "\def\variante{B}\def\modecorrige{1}% Évaluation pratique — Réseaux locaux (BTS CIEL 1re année) — sujets A et B
% Compilation (4 PDF) :
%   pdflatex -jobname=Eval_Reseau_SujetA   "\def\variante{A}\def\modecorrige{0}\input{Eval_TP_Reseau_LAN.tex}"
%   pdflatex -jobname=Eval_Reseau_SujetB   "\def\variante{B}\def\modecorrige{0}\input{Eval_TP_Reseau_LAN.tex}"
%   pdflatex -jobname=Eval_Reseau_CorrigeA "\def\variante{A}\def\modecorrige{1}\input{Eval_TP_Reseau_LAN.tex}"
%   pdflatex -jobname=Eval_Reseau_CorrigeB "\def\variante{B}\def\modecorrige{1}\input{Eval_TP_Reseau_LAN.tex}"
\documentclass[10pt,a4paper]{article}

\usepackage[T1]{fontenc}
\usepackage[utf8]{inputenc}
\usepackage{lmodern}
\usepackage[french]{babel}
\usepackage[a4paper,hmargin=1.5cm,top=1.3cm,bottom=1.6cm,footskip=0.7cm]{geometry}
\usepackage[table]{xcolor}
\usepackage{booktabs,tabularx,array}
\usepackage{enumitem}
\usepackage{fancyhdr}
\usepackage[most]{tcolorbox}
\usepackage{tikz}
\usetikzlibrary{positioning}
\usepackage{titlesec}

\providecommand{\modecorrige}{0}
\providecommand{\variante}{A}
\newif\ifcorrige
\ifnum\modecorrige=1 \corrigetrue \else \corrigefalse \fi

\definecolor{ipred}{HTML}{C0392B}
\definecolor{cdcbar}{HTML}{A67C3D}
\definecolor{cdcbg}{HTML}{F4F1EA}
\definecolor{qbar}{HTML}{4A6B5D}
\definecolor{qbg}{HTML}{EDF1EF}
\definecolor{corr}{HTML}{1F5FA8}
\definecolor{corrbg}{HTML}{E8F0FA}
\definecolor{grisfoot}{HTML}{828282}

\newcommand{\ip}[1]{\textcolor{ipred}{\texttt{#1}}}
\newcommand{\pts}[1]{\hfill\mbox{\textbf{(#1)}}}
\newcommand{\sol}[1]{\ifcorrige\textcolor{corr}{#1}\fi}
\newcommand{\bareme}[1]{\ifcorrige\par{\footnotesize\color{corr}#1\par}\fi}
\newcommand{\hl}{\ifcorrige\rule[-1.6mm]{0pt}{5.4mm}\else\rule[-2.2mm]{0pt}{6.6mm}\fi}
\newcolumntype{C}{>{\centering\arraybackslash}X}
\newcolumntype{L}{>{\raggedright\arraybackslash}X}

\titleformat{\section}{\large\bfseries}{}{0pt}{}
\titlespacing*{\section}{0pt}{6pt}{3pt}

\tcbset{qstyle/.style={enhanced,boxrule=0pt,leftrule=2.5pt,arc=0pt,outer arc=0pt,
  left=5pt,right=5pt,top=2pt,bottom=2pt,before skip=4pt,after skip=3pt}}
\newtcolorbox{cdc}{qstyle,colback=cdcbg,colframe=cdcbar}
\newtcolorbox{quest}{qstyle,colback=qbg,colframe=qbar}

% Zone de réponse : vierge dans le sujet, remplie dans le corrigé
\newcommand{\zone}[2]{%
  \ifcorrige
    \begin{tcolorbox}[colback=corrbg,colframe=corr,boxrule=0.5pt,arc=1pt,left=4pt,right=4pt,top=1pt,bottom=1pt,
      before skip=2pt,after skip=2pt,fontupper=\footnotesize\color{corr}]#2\end{tcolorbox}%
  \else
    \begin{tcolorbox}[colback=white,colframe=black!45,boxrule=0.4pt,arc=1pt,height=#1,
      before skip=2pt,after skip=2pt]\end{tcolorbox}%
  \fi}

\newcommand{\question}[3]{\begin{quest}\textbf{#1.} #3\pts{#2}\end{quest}}

%==========================================================================
% Données propres à chaque variante
%==========================================================================
\if\variante A
  \def\entreprise{Le cabinet d'architectes \textbf{ArchiNova}}
  \def\Rnom{R-ArchiNova}
  \def\resUn{Bureau d'études}\def\pUn{BE}\def\netUn{192.168.40}
  \def\resDeux{Accueil}\def\pDeux{ACC}\def\netDeux{172.16}
  \def\srvNom{SRV-Intranet}\def\netSrv{192.168.200}\def\srvH{10}
  \def\ssid{ArchiNova-WiFi}
  \def\AIrows{%
    \hl\ip{192.168.40.25} & \sol{C} & \sol{255.255.255.0} & \sol{192.168.40.0} & \sol{192.168.40.255} & \sol{254} \\ \hline
    \hl\ip{172.16.5.10}   & \sol{B} & \sol{255.255.0.0}   & \sol{172.16.0.0}   & \sol{172.16.255.255} & \sol{65\,534} \\ \hline
    \hl\ip{10.3.2.1}      & \sol{A} & \sol{255.0.0.0}     & \sol{10.0.0.0}     & \sol{10.255.255.255} & \sol{16\,777\,214} \\ \hline}
  \def\apipa{169.254.12.7}
  \def\Cnet{192.168.60}
  \def\Crows{%
    PC-A & \ip{192.168.60.10} & \ip{255.255.255.0} & \ip{192.168.60.254} & — \\
    PC-B & \ip{192.168.60.20} & \ip{255.255.255.0} & \ip{192.168.60.1}   & — \\
    PC-C & \ip{192.168.61.30} & \ip{255.255.255.0} & \ip{192.168.60.254} & — \\
    PC-D & \ip{192.168.60.10} & \ip{255.255.255.0} & \ip{192.168.60.254} & — \\
    SRV-C & \ip{10.0.0.5}     & \ip{255.0.0.0}     & \ip{10.0.0.254}     & — \\
    R1 — G0/0 & \ip{192.168.60.254} & \ip{255.255.255.0} & — & On \\
    R1 — G0/1 & \ip{10.0.0.254}     & \ip{255.0.0.0}     & — & Off \\}
  \def\Csol{%
    \hl 1 & \sol{PC-B} & \sol{Passerelle .60.1 : pas l'adresse du routeur.} & \sol{Joint son réseau, mais aucun autre (serveur).} & \sol{Passerelle 192.168.60.254.} \\ \hline
    \hl 2 & \sol{PC-C} & \sol{IP hors du réseau 192.168.60.0/24.} & \sol{Ne joint ni les postes ni le routeur.} & \sol{Ex. 192.168.60.30.} \\ \hline
    \hl 3 & \sol{PC-D} & \sol{Même IP que PC-A (doublon).} & \sol{Conflit d'adresse : PC-A et PC-D perturbés.} & \sol{Adresse unique, ex. .60.40.} \\ \hline
    \hl 4 & \sol{R1 G0/1} & \sol{Interface désactivée (Off).} & \sol{Réseau du serveur injoignable pour tous.} & \sol{Port Status : On.} \\ \hline}
  \def\Cbareme{PC-A et SRV-C sont correctement configurés ; accepter PC-A cité à la place de PC-D pour le doublon.}
\else
  \def\entreprise{La clinique vétérinaire \textbf{VetoPlus}}
  \def\Rnom{R-VetoPlus}
  \def\resUn{Soins}\def\pUn{SOI}\def\netUn{192.168.70}
  \def\resDeux{Secrétariat}\def\pDeux{SEC}\def\netDeux{172.20}
  \def\srvNom{SRV-Dossiers}\def\netSrv{192.168.250}\def\srvH{20}
  \def\ssid{VetoPlus-WiFi}
  \def\AIrows{%
    \hl\ip{192.168.15.130} & \sol{C} & \sol{255.255.255.0} & \sol{192.168.15.0} & \sol{192.168.15.255} & \sol{254} \\ \hline
    \hl\ip{172.31.200.7}   & \sol{B} & \sol{255.255.0.0}   & \sol{172.31.0.0}   & \sol{172.31.255.255} & \sol{65\,534} \\ \hline
    \hl\ip{10.45.6.9}      & \sol{A} & \sol{255.0.0.0}     & \sol{10.0.0.0}     & \sol{10.255.255.255} & \sol{16\,777\,214} \\ \hline}
  \def\apipa{169.254.88.3}
  \def\Cnet{192.168.80}
  \def\Crows{%
    PC-A & \ip{192.168.8.11}   & \ip{255.255.255.0} & \ip{192.168.80.254} & — \\
    PC-B & \ip{192.168.80.12}  & \ip{255.255.255.0} & \ip{192.168.80.254} & — \\
    PC-C & \ip{192.168.80.13}  & \ip{255.255.255.0} & \ip{192.168.80.100} & — \\
    PC-D & \ip{192.168.80.12}  & \ip{255.255.255.0} & \ip{192.168.80.254} & — \\
    SRV-C & \ip{10.0.0.8}      & \ip{255.0.0.0}     & \ip{10.0.0.254}     & — \\
    R1 — G0/0 & \ip{192.168.80.254} & \ip{255.255.255.0} & — & Off \\
    R1 — G0/1 & \ip{10.0.0.254}     & \ip{255.0.0.0}     & — & On \\}
  \def\Csol{%
    \hl 1 & \sol{PC-A} & \sol{IP hors du réseau 192.168.80.0/24.} & \sol{Ne joint ni les postes ni le routeur.} & \sol{Ex. 192.168.80.11.} \\ \hline
    \hl 2 & \sol{PC-C} & \sol{Passerelle .80.100 : pas l'adresse du routeur.} & \sol{Joint son réseau, mais aucun autre (serveur).} & \sol{Passerelle 192.168.80.254.} \\ \hline
    \hl 3 & \sol{PC-D} & \sol{Même IP que PC-B (doublon).} & \sol{Conflit d'adresse : PC-B et PC-D perturbés.} & \sol{Adresse unique, ex. .80.14.} \\ \hline
    \hl 4 & \sol{R1 G0/0} & \sol{Interface côté LAN désactivée (Off).} & \sol{Aucun poste ne joint le routeur ni le serveur.} & \sol{Port Status : On.} \\ \hline}
  \def\Cbareme{PC-B et SRV-C sont correctement configurés ; accepter PC-B cité à la place de PC-D pour le doublon.}
\fi

\pagestyle{fancy}
\fancyhf{}
\renewcommand{\headrulewidth}{0pt}
\fancyfoot[L]{\footnotesize\itshape\color{grisfoot}Évaluation — Réseaux locaux — Sujet \variante\ifcorrige\ — CORRIGÉ\fi}
\fancyfoot[R]{\footnotesize\color{grisfoot}\thepage/3}

\setlist{itemsep=0pt,topsep=2pt,parsep=0pt}
\setlength{\parindent}{0pt}
\setlength{\parskip}{2pt}

\begin{document}

\begin{center}
  {\Large\bfseries Évaluation pratique — Réseaux locaux \quad
   \fcolorbox{black}{white}{\ SUJET \variante\ }\par}
  \vspace{1pt}
  {\bfseries Adressage IP, commutation et routage — BTS CIEL 1re année — durée : 2 heures\par}
  \ifcorrige\vspace{2pt}{\bfseries\color{corr}CORRIGÉ ET BARÈME — document enseignant}\par\fi
\end{center}
\vspace{-2pt}
{\renewcommand{\arraystretch}{1.35}
\begin{tabularx}{\linewidth}{|X|X|p{2.2cm}|p{2.6cm}|}
  \hline
  \textbf{Nom :} & \textbf{Prénom :} & \textbf{Poste n°} & \textbf{Note :} \hfill /20 \\ \hline
\end{tabularx}}

\textbf{Consignes :} travail \textbf{individuel}, documents et Internet \textbf{non autorisés}. Ce sujet est le
\textbf{document-réponse}. Fichier Packet Tracer : \ip{Nom\_Prenom\_EvalLAN\_\variante.pkt}, à enregistrer
régulièrement dans le dossier de rendu. Les « visas » sont apposés par le professeur, que vous appelez à chaque test réussi.
Durées conseillées : A 30 min — B 1 h 10 — C 20 min.

%==========================================================================
\section{Partie A — Questions de cours et calculs \hfill /6}

\question{A1}{2 pts}{Complétez le tableau en considérant le \textbf{masque par défaut} de la classe de chaque adresse.}
{\small\renewcommand{\arraystretch}{1.15}
\begin{tabularx}{\linewidth}{|l|c|C|C|C|c|}
  \hline
  \textbf{Adresse IP} & \textbf{Classe} & \textbf{Masque par défaut} & \textbf{Adresse réseau} &
  \textbf{Adresse de diffusion} & \textbf{Nb d'hôtes} \\ \hline
  \AIrows
\end{tabularx}}
\bareme{0,5 pt par ligne juste + 0,5 pt si les trois nombres d'hôtes sont justes ($2^{n}-2$, $n$ = nombre de bits d'hôte).}

\question{A2}{1 pt}{Indiquez le type de câble (\textbf{droit} ou \textbf{croisé}) pour chaque liaison.}
{\small\renewcommand{\arraystretch}{1.15}
\begin{tabularx}{\linewidth}{|X|c||X|c|}
  \hline
  \hl PC $\leftrightarrow$ commutateur & \makebox[1.6cm]{\sol{droit}} & Commutateur $\leftrightarrow$ routeur & \makebox[1.6cm]{\sol{droit}} \\ \hline
  \hl PC $\leftrightarrow$ PC (réseau P2P) & \sol{croisé} & Concentrateur $\leftrightarrow$ concentrateur & \sol{croisé} \\ \hline
\end{tabularx}}
\bareme{0,25 pt par case. Équipements de types différents : câble droit ; de même type : câble croisé.}

\question{A3}{1 pt}{Comment un \textbf{concentrateur (hub)} et un \textbf{commutateur (switch)} traitent-ils une trame
reçue ? Lequel limite les collisions, et pourquoi ?}
\zone{2.9cm}{Le \textbf{hub} (couche 1) répète la trame sur \textbf{tous ses autres ports} : un seul domaine de collision
partagé. Le \textbf{switch} (couche 2) lit l'adresse MAC de destination et, grâce à sa \textbf{table MAC}, ne l'envoie
\textbf{que sur le port du destinataire} : un domaine de collision par port, c'est lui qui limite les collisions.
\emph{0,5 pt hub, 0,5 pt switch.}}

\question{A4}{1 pt}{Qu'est-ce que la \textbf{passerelle par défaut} d'un poste ? Quelle condition son adresse doit-elle respecter ?}
\zone{2.4cm}{Adresse IP de l'interface du \textbf{routeur} à laquelle le poste envoie les paquets destinés à un
\textbf{autre réseau}. Elle doit \textbf{appartenir au même réseau} que le poste. \emph{0,5 + 0,5 pt.}}

\question{A5}{0,5 pt}{Un portable configuré en DHCP affiche \ip{\apipa} / \ip{255.255.0.0}. Que s'est-il passé ?}
\zone{1.5cm}{Aucun \textbf{serveur DHCP} n'a répondu : le poste s'est attribué une adresse \textbf{APIPA} (169.254.0.0/16).}

\question{A6}{0,5 pt}{Quel protocole utilise \texttt{ping} ? Nommez les deux messages échangés et la couche OSI des adresses IP.}
\zone{1.5cm}{\textbf{ICMP} ; \textbf{Echo Request} / \textbf{Echo Reply} ; adresses IP en \textbf{couche 3 (Réseau)}.}

%==========================================================================
\newpage
\section{Partie B — Réalisation sous Packet Tracer \hfill /10}

\begin{cdc}
\entreprise{} vous confie son réseau : \textbf{trois réseaux} reliés par un \textbf{routeur unique Cisco 2911}
nommé \ip{\Rnom} (interfaces G0/0, G0/1, G0/2).
\begin{itemize}[leftmargin=12pt]
  \item \textbf{Réseau \resUn} : \textbf{4 postes} PC-\pUn1 à PC-\pUn4, commutateur \ip{SW-\pUn}, adressage statique,
        réseau \ip{\netUn.0/24}.
  \item \textbf{Réseau \resDeux} : \textbf{2 postes} PC-\pDeux1, PC-\pDeux2 et \textbf{1 imprimante} IMP-\pDeux,
        commutateur \ip{SW-\pDeux}, adressage statique, réseau de classe B \ip{\netDeux.0.0/16}.
  \item \textbf{Réseau Serveur} : serveur \ip{\srvNom} d'adresse \ip{\netSrv.\srvH/24}, relié \textbf{directement} au routeur.
  \item La passerelle de chaque réseau est sa \textbf{dernière adresse utilisable}. Tous les équipements communiquent
        entre eux et accèdent à la page web du serveur. Nommage imposé (« Display Name »).
\end{itemize}
\end{cdc}

\question{B1}{2 pts}{Complétez le plan d'adressage \textbf{avant} la réalisation.}
{\footnotesize\setlength{\tabcolsep}{3pt}\renewcommand{\arraystretch}{1.15}
\begin{tabularx}{\linewidth}{|l|C|C|C|C|C|C|}
  \hline
  \textbf{Réseau} & \textbf{Adresse réseau} & \textbf{Masque} & \textbf{1re @ utilisable} &
  \textbf{Dernière @ utilisable} & \textbf{Diffusion} & \textbf{Passerelle} \\ \hline
  \hl \resUn & \sol{\netUn.0} & \sol{255.255.255.0} & \sol{\netUn.1} & \sol{\netUn.254} & \sol{\netUn.255} & \sol{\netUn.254} \\ \hline
  \hl \resDeux & \sol{\netDeux.0.0} & \sol{255.255.0.0} & \sol{\netDeux.0.1} & \sol{\netDeux.255.254} & \sol{\netDeux.255.255} & \sol{\netDeux.255.254} \\ \hline
  \hl Serveur & \sol{\netSrv.0} & \sol{255.255.255.0} & \sol{\netSrv.1} & \sol{\netSrv.254} & \sol{\netSrv.255} & \sol{\netSrv.254} \\ \hline
\end{tabularx}}
\bareme{0,5 pt par ligne juste + 0,5 pt si la convention de passerelle est respectée partout.}

\question{B2}{2 pts}{Construisez la topologie (équipements, câblage, nommage) et relevez l'adressage configuré.}
{\footnotesize\setlength{\tabcolsep}{3pt}\renewcommand{\arraystretch}{1.1}
\begin{tabularx}{\linewidth}{|l|C|C|C||l|C|C|}
  \hline
  \textbf{Hôte} & \textbf{Adresse IP} & \textbf{Masque} & \textbf{Passerelle} & \textbf{Interface} & \textbf{Adresse IP} & \textbf{Masque} \\ \hline
  \hl PC-\pUn1 & \sol{\netUn.1} & \sol{255.255.255.0} & \sol{\netUn.254} & \Rnom{} G0/0 & \sol{\netUn.254} & \sol{255.255.255.0} \\ \hline
  \hl PC-\pDeux1 & \sol{\netDeux.0.1} & \sol{255.255.0.0} & \sol{\netDeux.255.254} & \Rnom{} G0/1 & \sol{\netDeux.255.254} & \sol{255.255.0.0} \\ \hline
  \hl IMP-\pDeux & \sol{\netDeux.0.10} & \sol{255.255.0.0} & \sol{\netDeux.255.254} & \Rnom{} G0/2 & \sol{\netSrv.254} & \sol{255.255.255.0} \\ \hline
  \hl \srvNom & \ip{\netSrv.\srvH} & \sol{255.255.255.0} & \sol{\netSrv.254} & \multicolumn{3}{c|}{\cellcolor{black!8}} \\ \hline
\end{tabularx}}
\bareme{Adresses d'hôtes libres (uniques, dans le bon réseau) ; affectation G0/x libre. Sur le .pkt : équipements 0,5 ;
câblage 1 (câbles droits vers les switchs et switch–routeur, \textbf{croisé} serveur–routeur, liens verts) ; nommage 0,5.}

\question{B3}{3 pts}{Configurez tous les équipements : IP, masques, passerelles des hôtes ; interfaces du routeur
(« Port Status : On »).}
\bareme{Sur le .pkt : interfaces routeur 1 pt ; IP/masques des hôtes 1 pt ; passerelles (imprimante et serveur compris)
1 pt. $-0,25$ par erreur, sans descendre sous 0 par item.}

\question{B4}{3 pts}{Réalisez les tests, notez la commande et le résultat, puis faites viser chaque test réussi.}
{\small\renewcommand{\arraystretch}{1.1}
\begin{tabularx}{\linewidth}{|c|l|l|C|c|c|}
  \hline
  \textbf{N°} & \textbf{Source} & \textbf{Destination} & \textbf{Commande / action} & \textbf{Résultat} & \textbf{Visa prof.} \\ \hline
  \hl 1 & PC-\pUn1 & PC-\pUn4 & \sol{ping \netUn.4} & \sol{réussi} & \hspace{1.4cm} \\ \hline
  \hl 2 & PC-\pDeux1 & IMP-\pDeux & \sol{ping \netDeux.0.10} & \sol{réussi} & \\ \hline
  \hl 3 & PC-\pUn2 & PC-\pDeux2 & \sol{ping \netDeux.0.2} & \sol{réussi} & \\ \hline
  \hl 4 & PC-\pDeux2 & \srvNom & \sol{ping \netSrv.\srvH} & \sol{réussi} & \\ \hline
  \hl 5 & PC-\pUn3 & page web du serveur & \sol{navigateur : {\NoAutoSpacing http://\netSrv.\srvH}} & \sol{page affichée} & \\ \hline
\end{tabularx}}
\bareme{0,5 pt par test visé + 0,5 pt pour B5. Le premier ping inter-réseaux peut échouer (résolution ARP) ; les suivants doivent réussir.}

\question{B5}{incl. B4}{Test n°3 : par quels équipements passe le paquet ? Pourquoi PC-\pUn2 l'envoie-t-il au routeur ?}
\zone{2.3cm}{PC-\pUn2 $\to$ SW-\pUn{} $\to$ \Rnom{} $\to$ SW-\pDeux{} $\to$ PC-\pDeux2. Grâce à son masque, PC-\pUn2 constate que
la destination (\netDeux.0.0/16) n'est \textbf{pas dans son réseau} (\netUn.0/24) : il envoie le paquet à sa
\textbf{passerelle par défaut}, qui le route.}

%==========================================================================
\newpage
\section{Partie C — Diagnostic d'un réseau en panne \hfill /4}

Un technicien constate qu'\textbf{aucun poste ne parvient à joindre le serveur SRV-C} et que certains postes
communiquent mal entre eux. Les câbles sont du bon type.

\begin{center}
\begin{minipage}{0.37\linewidth}\centering
\begin{tikzpicture}[font=\footnotesize,
  eq/.style={draw,fill=qbg,rounded corners=2pt,minimum width=9mm,minimum height=5mm,inner sep=2pt},
  every path/.style={thick}]
  \node[eq] (SW) {SW-1};
  \node[eq,right=9mm of SW] (R) {R1};
  \node[eq,right=7mm of R] (SRV) {SRV-C};
  \node[eq,above left=5mm and -3mm of SW] (A) {PC-A};
  \node[eq,above right=5mm and -3mm of SW] (B) {PC-B};
  \node[eq,below left=5mm and -3mm of SW] (C) {PC-C};
  \node[eq,below right=5mm and -3mm of SW] (D) {PC-D};
  \foreach \n in {A,B,C,D} \draw (\n) -- (SW);
  \draw (SW) -- node[above,font=\tiny]{G0/0} (R);
  \draw (R) -- node[above,font=\tiny]{G0/1} (SRV);
\end{tikzpicture}
\end{minipage}\hfill
\begin{minipage}{0.61\linewidth}\centering
{\small\renewcommand{\arraystretch}{1.05}\setlength{\tabcolsep}{4pt}
\begin{tabular}{lllll}
  \toprule
  \textbf{Équipement} & \textbf{Adresse IP} & \textbf{Masque} & \textbf{Passerelle} & \textbf{Port} \\
  \midrule
  \Crows
  \bottomrule
\end{tabular}}
\end{minipage}
\end{center}

\question{C1}{4 pts}{Relevez les \textbf{quatre anomalies}. Pour chacune : conséquence et correction.}
{\small\renewcommand{\arraystretch}{1.1}
\ifcorrige\def\hc{}\else\def\hc{\rule[-10mm]{0pt}{14mm}}\fi
\begin{tabularx}{\linewidth}{|c|>{\raggedright\arraybackslash}p{2cm}|L|L|L|}
  \hline
  \textbf{N°} & \textbf{Équipement} & \textbf{Anomalie} & \textbf{Conséquence} & \textbf{Correction} \\ \hline
  \ifcorrige\Csol\else
  \hc 1 & & & & \\ \hline
  \hc 2 & & & & \\ \hline
  \hc 3 & & & & \\ \hline
  \hc 4 & & & & \\ \hline
  \fi
\end{tabularx}}
\bareme{1 pt par anomalie (0,25 équipement, anomalie, conséquence, correction ; ordre indifférent). \Cbareme}

%==========================================================================
\section{Bonus — Extension sans fil \hfill +1}

\begin{cdc}Ajoutez au réseau \resUn{} un \textbf{AccessPoint-PT} relié à SW-\pUn{} (SSID \ip{\ssid}, \textbf{WPA2-PSK})
et un portable \ip{Laptop-\pUn} équipé d'un module \texttt{WPC300N}, en adressage \textbf{statique}.
Vérifiez qu'il joint \srvNom.\end{cdc}
{\small\renewcommand{\arraystretch}{1.1}
\begin{tabularx}{\linewidth}{|l|C|C|C|c|}
  \hline
  \textbf{Équipement} & \textbf{Adresse IP} & \textbf{Masque} & \textbf{Passerelle} & \textbf{Visa prof.} \\ \hline
  \hl Laptop-\pUn & \sol{\netUn.50} & \sol{255.255.255.0} & \sol{\netUn.254} & \hspace{1.4cm} \\ \hline
\end{tabularx}}
\bareme{Le point d'accès fait office de pont : le portable est dans \netUn.0/24. Bonus si associé en WPA2-PSK et ping
vers \netSrv.\srvH{} réussi. Note plafonnée à 20.}

\end{document}
"
\documentclass[10pt,a4paper]{article}

\usepackage[T1]{fontenc}
\usepackage[utf8]{inputenc}
\usepackage{lmodern}
\usepackage[french]{babel}
\usepackage[a4paper,hmargin=1.5cm,top=1.3cm,bottom=1.6cm,footskip=0.7cm]{geometry}
\usepackage[table]{xcolor}
\usepackage{booktabs,tabularx,array}
\usepackage{enumitem}
\usepackage{fancyhdr}
\usepackage[most]{tcolorbox}
\usepackage{tikz}
\usetikzlibrary{positioning}
\usepackage{titlesec}

\providecommand{\modecorrige}{0}
\providecommand{\variante}{A}
\newif\ifcorrige
\ifnum\modecorrige=1 \corrigetrue \else \corrigefalse \fi

\definecolor{ipred}{HTML}{C0392B}
\definecolor{cdcbar}{HTML}{A67C3D}
\definecolor{cdcbg}{HTML}{F4F1EA}
\definecolor{qbar}{HTML}{4A6B5D}
\definecolor{qbg}{HTML}{EDF1EF}
\definecolor{corr}{HTML}{1F5FA8}
\definecolor{corrbg}{HTML}{E8F0FA}
\definecolor{grisfoot}{HTML}{828282}

\newcommand{\ip}[1]{\textcolor{ipred}{\texttt{#1}}}
\newcommand{\pts}[1]{\hfill\mbox{\textbf{(#1)}}}
\newcommand{\sol}[1]{\ifcorrige\textcolor{corr}{#1}\fi}
\newcommand{\bareme}[1]{\ifcorrige\par{\footnotesize\color{corr}#1\par}\fi}
\newcommand{\hl}{\ifcorrige\rule[-1.6mm]{0pt}{5.4mm}\else\rule[-2.2mm]{0pt}{6.6mm}\fi}
\newcolumntype{C}{>{\centering\arraybackslash}X}
\newcolumntype{L}{>{\raggedright\arraybackslash}X}

\titleformat{\section}{\large\bfseries}{}{0pt}{}
\titlespacing*{\section}{0pt}{6pt}{3pt}

\tcbset{qstyle/.style={enhanced,boxrule=0pt,leftrule=2.5pt,arc=0pt,outer arc=0pt,
  left=5pt,right=5pt,top=2pt,bottom=2pt,before skip=4pt,after skip=3pt}}
\newtcolorbox{cdc}{qstyle,colback=cdcbg,colframe=cdcbar}
\newtcolorbox{quest}{qstyle,colback=qbg,colframe=qbar}

% Zone de réponse : vierge dans le sujet, remplie dans le corrigé
\newcommand{\zone}[2]{%
  \ifcorrige
    \begin{tcolorbox}[colback=corrbg,colframe=corr,boxrule=0.5pt,arc=1pt,left=4pt,right=4pt,top=1pt,bottom=1pt,
      before skip=2pt,after skip=2pt,fontupper=\footnotesize\color{corr}]#2\end{tcolorbox}%
  \else
    \begin{tcolorbox}[colback=white,colframe=black!45,boxrule=0.4pt,arc=1pt,height=#1,
      before skip=2pt,after skip=2pt]\end{tcolorbox}%
  \fi}

\newcommand{\question}[3]{\begin{quest}\textbf{#1.} #3\pts{#2}\end{quest}}

%==========================================================================
% Données propres à chaque variante
%==========================================================================
\if\variante A
  \def\entreprise{Le cabinet d'architectes \textbf{ArchiNova}}
  \def\Rnom{R-ArchiNova}
  \def\resUn{Bureau d'études}\def\pUn{BE}\def\netUn{192.168.40}
  \def\resDeux{Accueil}\def\pDeux{ACC}\def\netDeux{172.16}
  \def\srvNom{SRV-Intranet}\def\netSrv{192.168.200}\def\srvH{10}
  \def\ssid{ArchiNova-WiFi}
  \def\AIrows{%
    \hl\ip{192.168.40.25} & \sol{C} & \sol{255.255.255.0} & \sol{192.168.40.0} & \sol{192.168.40.255} & \sol{254} \\ \hline
    \hl\ip{172.16.5.10}   & \sol{B} & \sol{255.255.0.0}   & \sol{172.16.0.0}   & \sol{172.16.255.255} & \sol{65\,534} \\ \hline
    \hl\ip{10.3.2.1}      & \sol{A} & \sol{255.0.0.0}     & \sol{10.0.0.0}     & \sol{10.255.255.255} & \sol{16\,777\,214} \\ \hline}
  \def\apipa{169.254.12.7}
  \def\Cnet{192.168.60}
  \def\Crows{%
    PC-A & \ip{192.168.60.10} & \ip{255.255.255.0} & \ip{192.168.60.254} & — \\
    PC-B & \ip{192.168.60.20} & \ip{255.255.255.0} & \ip{192.168.60.1}   & — \\
    PC-C & \ip{192.168.61.30} & \ip{255.255.255.0} & \ip{192.168.60.254} & — \\
    PC-D & \ip{192.168.60.10} & \ip{255.255.255.0} & \ip{192.168.60.254} & — \\
    SRV-C & \ip{10.0.0.5}     & \ip{255.0.0.0}     & \ip{10.0.0.254}     & — \\
    R1 — G0/0 & \ip{192.168.60.254} & \ip{255.255.255.0} & — & On \\
    R1 — G0/1 & \ip{10.0.0.254}     & \ip{255.0.0.0}     & — & Off \\}
  \def\Csol{%
    \hl 1 & \sol{PC-B} & \sol{Passerelle .60.1 : pas l'adresse du routeur.} & \sol{Joint son réseau, mais aucun autre (serveur).} & \sol{Passerelle 192.168.60.254.} \\ \hline
    \hl 2 & \sol{PC-C} & \sol{IP hors du réseau 192.168.60.0/24.} & \sol{Ne joint ni les postes ni le routeur.} & \sol{Ex. 192.168.60.30.} \\ \hline
    \hl 3 & \sol{PC-D} & \sol{Même IP que PC-A (doublon).} & \sol{Conflit d'adresse : PC-A et PC-D perturbés.} & \sol{Adresse unique, ex. .60.40.} \\ \hline
    \hl 4 & \sol{R1 G0/1} & \sol{Interface désactivée (Off).} & \sol{Réseau du serveur injoignable pour tous.} & \sol{Port Status : On.} \\ \hline}
  \def\Cbareme{PC-A et SRV-C sont correctement configurés ; accepter PC-A cité à la place de PC-D pour le doublon.}
\else
  \def\entreprise{La clinique vétérinaire \textbf{VetoPlus}}
  \def\Rnom{R-VetoPlus}
  \def\resUn{Soins}\def\pUn{SOI}\def\netUn{192.168.70}
  \def\resDeux{Secrétariat}\def\pDeux{SEC}\def\netDeux{172.20}
  \def\srvNom{SRV-Dossiers}\def\netSrv{192.168.250}\def\srvH{20}
  \def\ssid{VetoPlus-WiFi}
  \def\AIrows{%
    \hl\ip{192.168.15.130} & \sol{C} & \sol{255.255.255.0} & \sol{192.168.15.0} & \sol{192.168.15.255} & \sol{254} \\ \hline
    \hl\ip{172.31.200.7}   & \sol{B} & \sol{255.255.0.0}   & \sol{172.31.0.0}   & \sol{172.31.255.255} & \sol{65\,534} \\ \hline
    \hl\ip{10.45.6.9}      & \sol{A} & \sol{255.0.0.0}     & \sol{10.0.0.0}     & \sol{10.255.255.255} & \sol{16\,777\,214} \\ \hline}
  \def\apipa{169.254.88.3}
  \def\Cnet{192.168.80}
  \def\Crows{%
    PC-A & \ip{192.168.8.11}   & \ip{255.255.255.0} & \ip{192.168.80.254} & — \\
    PC-B & \ip{192.168.80.12}  & \ip{255.255.255.0} & \ip{192.168.80.254} & — \\
    PC-C & \ip{192.168.80.13}  & \ip{255.255.255.0} & \ip{192.168.80.100} & — \\
    PC-D & \ip{192.168.80.12}  & \ip{255.255.255.0} & \ip{192.168.80.254} & — \\
    SRV-C & \ip{10.0.0.8}      & \ip{255.0.0.0}     & \ip{10.0.0.254}     & — \\
    R1 — G0/0 & \ip{192.168.80.254} & \ip{255.255.255.0} & — & Off \\
    R1 — G0/1 & \ip{10.0.0.254}     & \ip{255.0.0.0}     & — & On \\}
  \def\Csol{%
    \hl 1 & \sol{PC-A} & \sol{IP hors du réseau 192.168.80.0/24.} & \sol{Ne joint ni les postes ni le routeur.} & \sol{Ex. 192.168.80.11.} \\ \hline
    \hl 2 & \sol{PC-C} & \sol{Passerelle .80.100 : pas l'adresse du routeur.} & \sol{Joint son réseau, mais aucun autre (serveur).} & \sol{Passerelle 192.168.80.254.} \\ \hline
    \hl 3 & \sol{PC-D} & \sol{Même IP que PC-B (doublon).} & \sol{Conflit d'adresse : PC-B et PC-D perturbés.} & \sol{Adresse unique, ex. .80.14.} \\ \hline
    \hl 4 & \sol{R1 G0/0} & \sol{Interface côté LAN désactivée (Off).} & \sol{Aucun poste ne joint le routeur ni le serveur.} & \sol{Port Status : On.} \\ \hline}
  \def\Cbareme{PC-B et SRV-C sont correctement configurés ; accepter PC-B cité à la place de PC-D pour le doublon.}
\fi

\pagestyle{fancy}
\fancyhf{}
\renewcommand{\headrulewidth}{0pt}
\fancyfoot[L]{\footnotesize\itshape\color{grisfoot}Évaluation — Réseaux locaux — Sujet \variante\ifcorrige\ — CORRIGÉ\fi}
\fancyfoot[R]{\footnotesize\color{grisfoot}\thepage/3}

\setlist{itemsep=0pt,topsep=2pt,parsep=0pt}
\setlength{\parindent}{0pt}
\setlength{\parskip}{2pt}

\begin{document}

\begin{center}
  {\Large\bfseries Évaluation pratique — Réseaux locaux \quad
   \fcolorbox{black}{white}{\ SUJET \variante\ }\par}
  \vspace{1pt}
  {\bfseries Adressage IP, commutation et routage — BTS CIEL 1re année — durée : 2 heures\par}
  \ifcorrige\vspace{2pt}{\bfseries\color{corr}CORRIGÉ ET BARÈME — document enseignant}\par\fi
\end{center}
\vspace{-2pt}
{\renewcommand{\arraystretch}{1.35}
\begin{tabularx}{\linewidth}{|X|X|p{2.2cm}|p{2.6cm}|}
  \hline
  \textbf{Nom :} & \textbf{Prénom :} & \textbf{Poste n°} & \textbf{Note :} \hfill /20 \\ \hline
\end{tabularx}}

\textbf{Consignes :} travail \textbf{individuel}, documents et Internet \textbf{non autorisés}. Ce sujet est le
\textbf{document-réponse}. Fichier Packet Tracer : \ip{Nom\_Prenom\_EvalLAN\_\variante.pkt}, à enregistrer
régulièrement dans le dossier de rendu. Les « visas » sont apposés par le professeur, que vous appelez à chaque test réussi.
Durées conseillées : A 30 min — B 1 h 10 — C 20 min.

%==========================================================================
\section{Partie A — Questions de cours et calculs \hfill /6}

\question{A1}{2 pts}{Complétez le tableau en considérant le \textbf{masque par défaut} de la classe de chaque adresse.}
{\small\renewcommand{\arraystretch}{1.15}
\begin{tabularx}{\linewidth}{|l|c|C|C|C|c|}
  \hline
  \textbf{Adresse IP} & \textbf{Classe} & \textbf{Masque par défaut} & \textbf{Adresse réseau} &
  \textbf{Adresse de diffusion} & \textbf{Nb d'hôtes} \\ \hline
  \AIrows
\end{tabularx}}
\bareme{0,5 pt par ligne juste + 0,5 pt si les trois nombres d'hôtes sont justes ($2^{n}-2$, $n$ = nombre de bits d'hôte).}

\question{A2}{1 pt}{Indiquez le type de câble (\textbf{droit} ou \textbf{croisé}) pour chaque liaison.}
{\small\renewcommand{\arraystretch}{1.15}
\begin{tabularx}{\linewidth}{|X|c||X|c|}
  \hline
  \hl PC $\leftrightarrow$ commutateur & \makebox[1.6cm]{\sol{droit}} & Commutateur $\leftrightarrow$ routeur & \makebox[1.6cm]{\sol{droit}} \\ \hline
  \hl PC $\leftrightarrow$ PC (réseau P2P) & \sol{croisé} & Concentrateur $\leftrightarrow$ concentrateur & \sol{croisé} \\ \hline
\end{tabularx}}
\bareme{0,25 pt par case. Équipements de types différents : câble droit ; de même type : câble croisé.}

\question{A3}{1 pt}{Comment un \textbf{concentrateur (hub)} et un \textbf{commutateur (switch)} traitent-ils une trame
reçue ? Lequel limite les collisions, et pourquoi ?}
\zone{2.9cm}{Le \textbf{hub} (couche 1) répète la trame sur \textbf{tous ses autres ports} : un seul domaine de collision
partagé. Le \textbf{switch} (couche 2) lit l'adresse MAC de destination et, grâce à sa \textbf{table MAC}, ne l'envoie
\textbf{que sur le port du destinataire} : un domaine de collision par port, c'est lui qui limite les collisions.
\emph{0,5 pt hub, 0,5 pt switch.}}

\question{A4}{1 pt}{Qu'est-ce que la \textbf{passerelle par défaut} d'un poste ? Quelle condition son adresse doit-elle respecter ?}
\zone{2.4cm}{Adresse IP de l'interface du \textbf{routeur} à laquelle le poste envoie les paquets destinés à un
\textbf{autre réseau}. Elle doit \textbf{appartenir au même réseau} que le poste. \emph{0,5 + 0,5 pt.}}

\question{A5}{0,5 pt}{Un portable configuré en DHCP affiche \ip{\apipa} / \ip{255.255.0.0}. Que s'est-il passé ?}
\zone{1.5cm}{Aucun \textbf{serveur DHCP} n'a répondu : le poste s'est attribué une adresse \textbf{APIPA} (169.254.0.0/16).}

\question{A6}{0,5 pt}{Quel protocole utilise \texttt{ping} ? Nommez les deux messages échangés et la couche OSI des adresses IP.}
\zone{1.5cm}{\textbf{ICMP} ; \textbf{Echo Request} / \textbf{Echo Reply} ; adresses IP en \textbf{couche 3 (Réseau)}.}

%==========================================================================
\newpage
\section{Partie B — Réalisation sous Packet Tracer \hfill /10}

\begin{cdc}
\entreprise{} vous confie son réseau : \textbf{trois réseaux} reliés par un \textbf{routeur unique Cisco 2911}
nommé \ip{\Rnom} (interfaces G0/0, G0/1, G0/2).
\begin{itemize}[leftmargin=12pt]
  \item \textbf{Réseau \resUn} : \textbf{4 postes} PC-\pUn1 à PC-\pUn4, commutateur \ip{SW-\pUn}, adressage statique,
        réseau \ip{\netUn.0/24}.
  \item \textbf{Réseau \resDeux} : \textbf{2 postes} PC-\pDeux1, PC-\pDeux2 et \textbf{1 imprimante} IMP-\pDeux,
        commutateur \ip{SW-\pDeux}, adressage statique, réseau de classe B \ip{\netDeux.0.0/16}.
  \item \textbf{Réseau Serveur} : serveur \ip{\srvNom} d'adresse \ip{\netSrv.\srvH/24}, relié \textbf{directement} au routeur.
  \item La passerelle de chaque réseau est sa \textbf{dernière adresse utilisable}. Tous les équipements communiquent
        entre eux et accèdent à la page web du serveur. Nommage imposé (« Display Name »).
\end{itemize}
\end{cdc}

\question{B1}{2 pts}{Complétez le plan d'adressage \textbf{avant} la réalisation.}
{\footnotesize\setlength{\tabcolsep}{3pt}\renewcommand{\arraystretch}{1.15}
\begin{tabularx}{\linewidth}{|l|C|C|C|C|C|C|}
  \hline
  \textbf{Réseau} & \textbf{Adresse réseau} & \textbf{Masque} & \textbf{1re @ utilisable} &
  \textbf{Dernière @ utilisable} & \textbf{Diffusion} & \textbf{Passerelle} \\ \hline
  \hl \resUn & \sol{\netUn.0} & \sol{255.255.255.0} & \sol{\netUn.1} & \sol{\netUn.254} & \sol{\netUn.255} & \sol{\netUn.254} \\ \hline
  \hl \resDeux & \sol{\netDeux.0.0} & \sol{255.255.0.0} & \sol{\netDeux.0.1} & \sol{\netDeux.255.254} & \sol{\netDeux.255.255} & \sol{\netDeux.255.254} \\ \hline
  \hl Serveur & \sol{\netSrv.0} & \sol{255.255.255.0} & \sol{\netSrv.1} & \sol{\netSrv.254} & \sol{\netSrv.255} & \sol{\netSrv.254} \\ \hline
\end{tabularx}}
\bareme{0,5 pt par ligne juste + 0,5 pt si la convention de passerelle est respectée partout.}

\question{B2}{2 pts}{Construisez la topologie (équipements, câblage, nommage) et relevez l'adressage configuré.}
{\footnotesize\setlength{\tabcolsep}{3pt}\renewcommand{\arraystretch}{1.1}
\begin{tabularx}{\linewidth}{|l|C|C|C||l|C|C|}
  \hline
  \textbf{Hôte} & \textbf{Adresse IP} & \textbf{Masque} & \textbf{Passerelle} & \textbf{Interface} & \textbf{Adresse IP} & \textbf{Masque} \\ \hline
  \hl PC-\pUn1 & \sol{\netUn.1} & \sol{255.255.255.0} & \sol{\netUn.254} & \Rnom{} G0/0 & \sol{\netUn.254} & \sol{255.255.255.0} \\ \hline
  \hl PC-\pDeux1 & \sol{\netDeux.0.1} & \sol{255.255.0.0} & \sol{\netDeux.255.254} & \Rnom{} G0/1 & \sol{\netDeux.255.254} & \sol{255.255.0.0} \\ \hline
  \hl IMP-\pDeux & \sol{\netDeux.0.10} & \sol{255.255.0.0} & \sol{\netDeux.255.254} & \Rnom{} G0/2 & \sol{\netSrv.254} & \sol{255.255.255.0} \\ \hline
  \hl \srvNom & \ip{\netSrv.\srvH} & \sol{255.255.255.0} & \sol{\netSrv.254} & \multicolumn{3}{c|}{\cellcolor{black!8}} \\ \hline
\end{tabularx}}
\bareme{Adresses d'hôtes libres (uniques, dans le bon réseau) ; affectation G0/x libre. Sur le .pkt : équipements 0,5 ;
câblage 1 (câbles droits vers les switchs et switch–routeur, \textbf{croisé} serveur–routeur, liens verts) ; nommage 0,5.}

\question{B3}{3 pts}{Configurez tous les équipements : IP, masques, passerelles des hôtes ; interfaces du routeur
(« Port Status : On »).}
\bareme{Sur le .pkt : interfaces routeur 1 pt ; IP/masques des hôtes 1 pt ; passerelles (imprimante et serveur compris)
1 pt. $-0,25$ par erreur, sans descendre sous 0 par item.}

\question{B4}{3 pts}{Réalisez les tests, notez la commande et le résultat, puis faites viser chaque test réussi.}
{\small\renewcommand{\arraystretch}{1.1}
\begin{tabularx}{\linewidth}{|c|l|l|C|c|c|}
  \hline
  \textbf{N°} & \textbf{Source} & \textbf{Destination} & \textbf{Commande / action} & \textbf{Résultat} & \textbf{Visa prof.} \\ \hline
  \hl 1 & PC-\pUn1 & PC-\pUn4 & \sol{ping \netUn.4} & \sol{réussi} & \hspace{1.4cm} \\ \hline
  \hl 2 & PC-\pDeux1 & IMP-\pDeux & \sol{ping \netDeux.0.10} & \sol{réussi} & \\ \hline
  \hl 3 & PC-\pUn2 & PC-\pDeux2 & \sol{ping \netDeux.0.2} & \sol{réussi} & \\ \hline
  \hl 4 & PC-\pDeux2 & \srvNom & \sol{ping \netSrv.\srvH} & \sol{réussi} & \\ \hline
  \hl 5 & PC-\pUn3 & page web du serveur & \sol{navigateur : {\NoAutoSpacing http://\netSrv.\srvH}} & \sol{page affichée} & \\ \hline
\end{tabularx}}
\bareme{0,5 pt par test visé + 0,5 pt pour B5. Le premier ping inter-réseaux peut échouer (résolution ARP) ; les suivants doivent réussir.}

\question{B5}{incl. B4}{Test n°3 : par quels équipements passe le paquet ? Pourquoi PC-\pUn2 l'envoie-t-il au routeur ?}
\zone{2.3cm}{PC-\pUn2 $\to$ SW-\pUn{} $\to$ \Rnom{} $\to$ SW-\pDeux{} $\to$ PC-\pDeux2. Grâce à son masque, PC-\pUn2 constate que
la destination (\netDeux.0.0/16) n'est \textbf{pas dans son réseau} (\netUn.0/24) : il envoie le paquet à sa
\textbf{passerelle par défaut}, qui le route.}

%==========================================================================
\newpage
\section{Partie C — Diagnostic d'un réseau en panne \hfill /4}

Un technicien constate qu'\textbf{aucun poste ne parvient à joindre le serveur SRV-C} et que certains postes
communiquent mal entre eux. Les câbles sont du bon type.

\begin{center}
\begin{minipage}{0.37\linewidth}\centering
\begin{tikzpicture}[font=\footnotesize,
  eq/.style={draw,fill=qbg,rounded corners=2pt,minimum width=9mm,minimum height=5mm,inner sep=2pt},
  every path/.style={thick}]
  \node[eq] (SW) {SW-1};
  \node[eq,right=9mm of SW] (R) {R1};
  \node[eq,right=7mm of R] (SRV) {SRV-C};
  \node[eq,above left=5mm and -3mm of SW] (A) {PC-A};
  \node[eq,above right=5mm and -3mm of SW] (B) {PC-B};
  \node[eq,below left=5mm and -3mm of SW] (C) {PC-C};
  \node[eq,below right=5mm and -3mm of SW] (D) {PC-D};
  \foreach \n in {A,B,C,D} \draw (\n) -- (SW);
  \draw (SW) -- node[above,font=\tiny]{G0/0} (R);
  \draw (R) -- node[above,font=\tiny]{G0/1} (SRV);
\end{tikzpicture}
\end{minipage}\hfill
\begin{minipage}{0.61\linewidth}\centering
{\small\renewcommand{\arraystretch}{1.05}\setlength{\tabcolsep}{4pt}
\begin{tabular}{lllll}
  \toprule
  \textbf{Équipement} & \textbf{Adresse IP} & \textbf{Masque} & \textbf{Passerelle} & \textbf{Port} \\
  \midrule
  \Crows
  \bottomrule
\end{tabular}}
\end{minipage}
\end{center}

\question{C1}{4 pts}{Relevez les \textbf{quatre anomalies}. Pour chacune : conséquence et correction.}
{\small\renewcommand{\arraystretch}{1.1}
\ifcorrige\def\hc{}\else\def\hc{\rule[-10mm]{0pt}{14mm}}\fi
\begin{tabularx}{\linewidth}{|c|>{\raggedright\arraybackslash}p{2cm}|L|L|L|}
  \hline
  \textbf{N°} & \textbf{Équipement} & \textbf{Anomalie} & \textbf{Conséquence} & \textbf{Correction} \\ \hline
  \ifcorrige\Csol\else
  \hc 1 & & & & \\ \hline
  \hc 2 & & & & \\ \hline
  \hc 3 & & & & \\ \hline
  \hc 4 & & & & \\ \hline
  \fi
\end{tabularx}}
\bareme{1 pt par anomalie (0,25 équipement, anomalie, conséquence, correction ; ordre indifférent). \Cbareme}

%==========================================================================
\section{Bonus — Extension sans fil \hfill +1}

\begin{cdc}Ajoutez au réseau \resUn{} un \textbf{AccessPoint-PT} relié à SW-\pUn{} (SSID \ip{\ssid}, \textbf{WPA2-PSK})
et un portable \ip{Laptop-\pUn} équipé d'un module \texttt{WPC300N}, en adressage \textbf{statique}.
Vérifiez qu'il joint \srvNom.\end{cdc}
{\small\renewcommand{\arraystretch}{1.1}
\begin{tabularx}{\linewidth}{|l|C|C|C|c|}
  \hline
  \textbf{Équipement} & \textbf{Adresse IP} & \textbf{Masque} & \textbf{Passerelle} & \textbf{Visa prof.} \\ \hline
  \hl Laptop-\pUn & \sol{\netUn.50} & \sol{255.255.255.0} & \sol{\netUn.254} & \hspace{1.4cm} \\ \hline
\end{tabularx}}
\bareme{Le point d'accès fait office de pont : le portable est dans \netUn.0/24. Bonus si associé en WPA2-PSK et ping
vers \netSrv.\srvH{} réussi. Note plafonnée à 20.}

\end{document}
"
\documentclass[10pt,a4paper]{article}

\usepackage[T1]{fontenc}
\usepackage[utf8]{inputenc}
\usepackage{lmodern}
\usepackage[french]{babel}
\usepackage[a4paper,hmargin=1.5cm,top=1.3cm,bottom=1.6cm,footskip=0.7cm]{geometry}
\usepackage[table]{xcolor}
\usepackage{booktabs,tabularx,array}
\usepackage{enumitem}
\usepackage{fancyhdr}
\usepackage[most]{tcolorbox}
\usepackage{tikz}
\usetikzlibrary{positioning}
\usepackage{titlesec}

\providecommand{\modecorrige}{0}
\providecommand{\variante}{A}
\newif\ifcorrige
\ifnum\modecorrige=1 \corrigetrue \else \corrigefalse \fi

\definecolor{ipred}{HTML}{C0392B}
\definecolor{cdcbar}{HTML}{A67C3D}
\definecolor{cdcbg}{HTML}{F4F1EA}
\definecolor{qbar}{HTML}{4A6B5D}
\definecolor{qbg}{HTML}{EDF1EF}
\definecolor{corr}{HTML}{1F5FA8}
\definecolor{corrbg}{HTML}{E8F0FA}
\definecolor{grisfoot}{HTML}{828282}

\newcommand{\ip}[1]{\textcolor{ipred}{\texttt{#1}}}
\newcommand{\pts}[1]{\hfill\mbox{\textbf{(#1)}}}
\newcommand{\sol}[1]{\ifcorrige\textcolor{corr}{#1}\fi}
\newcommand{\bareme}[1]{\ifcorrige\par{\footnotesize\color{corr}#1\par}\fi}
\newcommand{\hl}{\ifcorrige\rule[-1.6mm]{0pt}{5.4mm}\else\rule[-2.2mm]{0pt}{6.6mm}\fi}
\newcolumntype{C}{>{\centering\arraybackslash}X}
\newcolumntype{L}{>{\raggedright\arraybackslash}X}

\titleformat{\section}{\large\bfseries}{}{0pt}{}
\titlespacing*{\section}{0pt}{6pt}{3pt}

\tcbset{qstyle/.style={enhanced,boxrule=0pt,leftrule=2.5pt,arc=0pt,outer arc=0pt,
  left=5pt,right=5pt,top=2pt,bottom=2pt,before skip=4pt,after skip=3pt}}
\newtcolorbox{cdc}{qstyle,colback=cdcbg,colframe=cdcbar}
\newtcolorbox{quest}{qstyle,colback=qbg,colframe=qbar}

% Zone de réponse : vierge dans le sujet, remplie dans le corrigé
\newcommand{\zone}[2]{%
  \ifcorrige
    \begin{tcolorbox}[colback=corrbg,colframe=corr,boxrule=0.5pt,arc=1pt,left=4pt,right=4pt,top=1pt,bottom=1pt,
      before skip=2pt,after skip=2pt,fontupper=\footnotesize\color{corr}]#2\end{tcolorbox}%
  \else
    \begin{tcolorbox}[colback=white,colframe=black!45,boxrule=0.4pt,arc=1pt,height=#1,
      before skip=2pt,after skip=2pt]\end{tcolorbox}%
  \fi}

\newcommand{\question}[3]{\begin{quest}\textbf{#1.} #3\pts{#2}\end{quest}}

%==========================================================================
% Données propres à chaque variante
%==========================================================================
\if\variante A
  \def\entreprise{Le cabinet d'architectes \textbf{ArchiNova}}
  \def\Rnom{R-ArchiNova}
  \def\resUn{Bureau d'études}\def\pUn{BE}\def\netUn{192.168.40}
  \def\resDeux{Accueil}\def\pDeux{ACC}\def\netDeux{172.16}
  \def\srvNom{SRV-Intranet}\def\netSrv{192.168.200}\def\srvH{10}
  \def\ssid{ArchiNova-WiFi}
  \def\AIrows{%
    \hl\ip{192.168.40.25} & \sol{C} & \sol{255.255.255.0} & \sol{192.168.40.0} & \sol{192.168.40.255} & \sol{254} \\ \hline
    \hl\ip{172.16.5.10}   & \sol{B} & \sol{255.255.0.0}   & \sol{172.16.0.0}   & \sol{172.16.255.255} & \sol{65\,534} \\ \hline
    \hl\ip{10.3.2.1}      & \sol{A} & \sol{255.0.0.0}     & \sol{10.0.0.0}     & \sol{10.255.255.255} & \sol{16\,777\,214} \\ \hline}
  \def\apipa{169.254.12.7}
  \def\Cnet{192.168.60}
  \def\Crows{%
    PC-A & \ip{192.168.60.10} & \ip{255.255.255.0} & \ip{192.168.60.254} & — \\
    PC-B & \ip{192.168.60.20} & \ip{255.255.255.0} & \ip{192.168.60.1}   & — \\
    PC-C & \ip{192.168.61.30} & \ip{255.255.255.0} & \ip{192.168.60.254} & — \\
    PC-D & \ip{192.168.60.10} & \ip{255.255.255.0} & \ip{192.168.60.254} & — \\
    SRV-C & \ip{10.0.0.5}     & \ip{255.0.0.0}     & \ip{10.0.0.254}     & — \\
    R1 — G0/0 & \ip{192.168.60.254} & \ip{255.255.255.0} & — & On \\
    R1 — G0/1 & \ip{10.0.0.254}     & \ip{255.0.0.0}     & — & Off \\}
  \def\Csol{%
    \hl 1 & \sol{PC-B} & \sol{Passerelle .60.1 : pas l'adresse du routeur.} & \sol{Joint son réseau, mais aucun autre (serveur).} & \sol{Passerelle 192.168.60.254.} \\ \hline
    \hl 2 & \sol{PC-C} & \sol{IP hors du réseau 192.168.60.0/24.} & \sol{Ne joint ni les postes ni le routeur.} & \sol{Ex. 192.168.60.30.} \\ \hline
    \hl 3 & \sol{PC-D} & \sol{Même IP que PC-A (doublon).} & \sol{Conflit d'adresse : PC-A et PC-D perturbés.} & \sol{Adresse unique, ex. .60.40.} \\ \hline
    \hl 4 & \sol{R1 G0/1} & \sol{Interface désactivée (Off).} & \sol{Réseau du serveur injoignable pour tous.} & \sol{Port Status : On.} \\ \hline}
  \def\Cbareme{PC-A et SRV-C sont correctement configurés ; accepter PC-A cité à la place de PC-D pour le doublon.}
\else
  \def\entreprise{La clinique vétérinaire \textbf{VetoPlus}}
  \def\Rnom{R-VetoPlus}
  \def\resUn{Soins}\def\pUn{SOI}\def\netUn{192.168.70}
  \def\resDeux{Secrétariat}\def\pDeux{SEC}\def\netDeux{172.20}
  \def\srvNom{SRV-Dossiers}\def\netSrv{192.168.250}\def\srvH{20}
  \def\ssid{VetoPlus-WiFi}
  \def\AIrows{%
    \hl\ip{192.168.15.130} & \sol{C} & \sol{255.255.255.0} & \sol{192.168.15.0} & \sol{192.168.15.255} & \sol{254} \\ \hline
    \hl\ip{172.31.200.7}   & \sol{B} & \sol{255.255.0.0}   & \sol{172.31.0.0}   & \sol{172.31.255.255} & \sol{65\,534} \\ \hline
    \hl\ip{10.45.6.9}      & \sol{A} & \sol{255.0.0.0}     & \sol{10.0.0.0}     & \sol{10.255.255.255} & \sol{16\,777\,214} \\ \hline}
  \def\apipa{169.254.88.3}
  \def\Cnet{192.168.80}
  \def\Crows{%
    PC-A & \ip{192.168.8.11}   & \ip{255.255.255.0} & \ip{192.168.80.254} & — \\
    PC-B & \ip{192.168.80.12}  & \ip{255.255.255.0} & \ip{192.168.80.254} & — \\
    PC-C & \ip{192.168.80.13}  & \ip{255.255.255.0} & \ip{192.168.80.100} & — \\
    PC-D & \ip{192.168.80.12}  & \ip{255.255.255.0} & \ip{192.168.80.254} & — \\
    SRV-C & \ip{10.0.0.8}      & \ip{255.0.0.0}     & \ip{10.0.0.254}     & — \\
    R1 — G0/0 & \ip{192.168.80.254} & \ip{255.255.255.0} & — & Off \\
    R1 — G0/1 & \ip{10.0.0.254}     & \ip{255.0.0.0}     & — & On \\}
  \def\Csol{%
    \hl 1 & \sol{PC-A} & \sol{IP hors du réseau 192.168.80.0/24.} & \sol{Ne joint ni les postes ni le routeur.} & \sol{Ex. 192.168.80.11.} \\ \hline
    \hl 2 & \sol{PC-C} & \sol{Passerelle .80.100 : pas l'adresse du routeur.} & \sol{Joint son réseau, mais aucun autre (serveur).} & \sol{Passerelle 192.168.80.254.} \\ \hline
    \hl 3 & \sol{PC-D} & \sol{Même IP que PC-B (doublon).} & \sol{Conflit d'adresse : PC-B et PC-D perturbés.} & \sol{Adresse unique, ex. .80.14.} \\ \hline
    \hl 4 & \sol{R1 G0/0} & \sol{Interface côté LAN désactivée (Off).} & \sol{Aucun poste ne joint le routeur ni le serveur.} & \sol{Port Status : On.} \\ \hline}
  \def\Cbareme{PC-B et SRV-C sont correctement configurés ; accepter PC-B cité à la place de PC-D pour le doublon.}
\fi

\pagestyle{fancy}
\fancyhf{}
\renewcommand{\headrulewidth}{0pt}
\fancyfoot[L]{\footnotesize\itshape\color{grisfoot}Évaluation — Réseaux locaux — Sujet \variante\ifcorrige\ — CORRIGÉ\fi}
\fancyfoot[R]{\footnotesize\color{grisfoot}\thepage/3}

\setlist{itemsep=0pt,topsep=2pt,parsep=0pt}
\setlength{\parindent}{0pt}
\setlength{\parskip}{2pt}

\begin{document}

\begin{center}
  {\Large\bfseries Évaluation pratique — Réseaux locaux \quad
   \fcolorbox{black}{white}{\ SUJET \variante\ }\par}
  \vspace{1pt}
  {\bfseries Adressage IP, commutation et routage — BTS CIEL 1re année — durée : 2 heures\par}
  \ifcorrige\vspace{2pt}{\bfseries\color{corr}CORRIGÉ ET BARÈME — document enseignant}\par\fi
\end{center}
\vspace{-2pt}
{\renewcommand{\arraystretch}{1.35}
\begin{tabularx}{\linewidth}{|X|X|p{2.2cm}|p{2.6cm}|}
  \hline
  \textbf{Nom :} & \textbf{Prénom :} & \textbf{Poste n°} & \textbf{Note :} \hfill /20 \\ \hline
\end{tabularx}}

\textbf{Consignes :} travail \textbf{individuel}, documents et Internet \textbf{non autorisés}. Ce sujet est le
\textbf{document-réponse}. Fichier Packet Tracer : \ip{Nom\_Prenom\_EvalLAN\_\variante.pkt}, à enregistrer
régulièrement dans le dossier de rendu. Les « visas » sont apposés par le professeur, que vous appelez à chaque test réussi.
Durées conseillées : A 30 min — B 1 h 10 — C 20 min.

%==========================================================================
\section{Partie A — Questions de cours et calculs \hfill /6}

\question{A1}{2 pts}{Complétez le tableau en considérant le \textbf{masque par défaut} de la classe de chaque adresse.}
{\small\renewcommand{\arraystretch}{1.15}
\begin{tabularx}{\linewidth}{|l|c|C|C|C|c|}
  \hline
  \textbf{Adresse IP} & \textbf{Classe} & \textbf{Masque par défaut} & \textbf{Adresse réseau} &
  \textbf{Adresse de diffusion} & \textbf{Nb d'hôtes} \\ \hline
  \AIrows
\end{tabularx}}
\bareme{0,5 pt par ligne juste + 0,5 pt si les trois nombres d'hôtes sont justes ($2^{n}-2$, $n$ = nombre de bits d'hôte).}

\question{A2}{1 pt}{Indiquez le type de câble (\textbf{droit} ou \textbf{croisé}) pour chaque liaison.}
{\small\renewcommand{\arraystretch}{1.15}
\begin{tabularx}{\linewidth}{|X|c||X|c|}
  \hline
  \hl PC $\leftrightarrow$ commutateur & \makebox[1.6cm]{\sol{droit}} & Commutateur $\leftrightarrow$ routeur & \makebox[1.6cm]{\sol{droit}} \\ \hline
  \hl PC $\leftrightarrow$ PC (réseau P2P) & \sol{croisé} & Concentrateur $\leftrightarrow$ concentrateur & \sol{croisé} \\ \hline
\end{tabularx}}
\bareme{0,25 pt par case. Équipements de types différents : câble droit ; de même type : câble croisé.}

\question{A3}{1 pt}{Comment un \textbf{concentrateur (hub)} et un \textbf{commutateur (switch)} traitent-ils une trame
reçue ? Lequel limite les collisions, et pourquoi ?}
\zone{2.9cm}{Le \textbf{hub} (couche 1) répète la trame sur \textbf{tous ses autres ports} : un seul domaine de collision
partagé. Le \textbf{switch} (couche 2) lit l'adresse MAC de destination et, grâce à sa \textbf{table MAC}, ne l'envoie
\textbf{que sur le port du destinataire} : un domaine de collision par port, c'est lui qui limite les collisions.
\emph{0,5 pt hub, 0,5 pt switch.}}

\question{A4}{1 pt}{Qu'est-ce que la \textbf{passerelle par défaut} d'un poste ? Quelle condition son adresse doit-elle respecter ?}
\zone{2.4cm}{Adresse IP de l'interface du \textbf{routeur} à laquelle le poste envoie les paquets destinés à un
\textbf{autre réseau}. Elle doit \textbf{appartenir au même réseau} que le poste. \emph{0,5 + 0,5 pt.}}

\question{A5}{0,5 pt}{Un portable configuré en DHCP affiche \ip{\apipa} / \ip{255.255.0.0}. Que s'est-il passé ?}
\zone{1.5cm}{Aucun \textbf{serveur DHCP} n'a répondu : le poste s'est attribué une adresse \textbf{APIPA} (169.254.0.0/16).}

\question{A6}{0,5 pt}{Quel protocole utilise \texttt{ping} ? Nommez les deux messages échangés et la couche OSI des adresses IP.}
\zone{1.5cm}{\textbf{ICMP} ; \textbf{Echo Request} / \textbf{Echo Reply} ; adresses IP en \textbf{couche 3 (Réseau)}.}

%==========================================================================
\newpage
\section{Partie B — Réalisation sous Packet Tracer \hfill /10}

\begin{cdc}
\entreprise{} vous confie son réseau : \textbf{trois réseaux} reliés par un \textbf{routeur unique Cisco 2911}
nommé \ip{\Rnom} (interfaces G0/0, G0/1, G0/2).
\begin{itemize}[leftmargin=12pt]
  \item \textbf{Réseau \resUn} : \textbf{4 postes} PC-\pUn1 à PC-\pUn4, commutateur \ip{SW-\pUn}, adressage statique,
        réseau \ip{\netUn.0/24}.
  \item \textbf{Réseau \resDeux} : \textbf{2 postes} PC-\pDeux1, PC-\pDeux2 et \textbf{1 imprimante} IMP-\pDeux,
        commutateur \ip{SW-\pDeux}, adressage statique, réseau de classe B \ip{\netDeux.0.0/16}.
  \item \textbf{Réseau Serveur} : serveur \ip{\srvNom} d'adresse \ip{\netSrv.\srvH/24}, relié \textbf{directement} au routeur.
  \item La passerelle de chaque réseau est sa \textbf{dernière adresse utilisable}. Tous les équipements communiquent
        entre eux et accèdent à la page web du serveur. Nommage imposé (« Display Name »).
\end{itemize}
\end{cdc}

\question{B1}{2 pts}{Complétez le plan d'adressage \textbf{avant} la réalisation.}
{\footnotesize\setlength{\tabcolsep}{3pt}\renewcommand{\arraystretch}{1.15}
\begin{tabularx}{\linewidth}{|l|C|C|C|C|C|C|}
  \hline
  \textbf{Réseau} & \textbf{Adresse réseau} & \textbf{Masque} & \textbf{1re @ utilisable} &
  \textbf{Dernière @ utilisable} & \textbf{Diffusion} & \textbf{Passerelle} \\ \hline
  \hl \resUn & \sol{\netUn.0} & \sol{255.255.255.0} & \sol{\netUn.1} & \sol{\netUn.254} & \sol{\netUn.255} & \sol{\netUn.254} \\ \hline
  \hl \resDeux & \sol{\netDeux.0.0} & \sol{255.255.0.0} & \sol{\netDeux.0.1} & \sol{\netDeux.255.254} & \sol{\netDeux.255.255} & \sol{\netDeux.255.254} \\ \hline
  \hl Serveur & \sol{\netSrv.0} & \sol{255.255.255.0} & \sol{\netSrv.1} & \sol{\netSrv.254} & \sol{\netSrv.255} & \sol{\netSrv.254} \\ \hline
\end{tabularx}}
\bareme{0,5 pt par ligne juste + 0,5 pt si la convention de passerelle est respectée partout.}

\question{B2}{2 pts}{Construisez la topologie (équipements, câblage, nommage) et relevez l'adressage configuré.}
{\footnotesize\setlength{\tabcolsep}{3pt}\renewcommand{\arraystretch}{1.1}
\begin{tabularx}{\linewidth}{|l|C|C|C||l|C|C|}
  \hline
  \textbf{Hôte} & \textbf{Adresse IP} & \textbf{Masque} & \textbf{Passerelle} & \textbf{Interface} & \textbf{Adresse IP} & \textbf{Masque} \\ \hline
  \hl PC-\pUn1 & \sol{\netUn.1} & \sol{255.255.255.0} & \sol{\netUn.254} & \Rnom{} G0/0 & \sol{\netUn.254} & \sol{255.255.255.0} \\ \hline
  \hl PC-\pDeux1 & \sol{\netDeux.0.1} & \sol{255.255.0.0} & \sol{\netDeux.255.254} & \Rnom{} G0/1 & \sol{\netDeux.255.254} & \sol{255.255.0.0} \\ \hline
  \hl IMP-\pDeux & \sol{\netDeux.0.10} & \sol{255.255.0.0} & \sol{\netDeux.255.254} & \Rnom{} G0/2 & \sol{\netSrv.254} & \sol{255.255.255.0} \\ \hline
  \hl \srvNom & \ip{\netSrv.\srvH} & \sol{255.255.255.0} & \sol{\netSrv.254} & \multicolumn{3}{c|}{\cellcolor{black!8}} \\ \hline
\end{tabularx}}
\bareme{Adresses d'hôtes libres (uniques, dans le bon réseau) ; affectation G0/x libre. Sur le .pkt : équipements 0,5 ;
câblage 1 (câbles droits vers les switchs et switch–routeur, \textbf{croisé} serveur–routeur, liens verts) ; nommage 0,5.}

\question{B3}{3 pts}{Configurez tous les équipements : IP, masques, passerelles des hôtes ; interfaces du routeur
(« Port Status : On »).}
\bareme{Sur le .pkt : interfaces routeur 1 pt ; IP/masques des hôtes 1 pt ; passerelles (imprimante et serveur compris)
1 pt. $-0,25$ par erreur, sans descendre sous 0 par item.}

\question{B4}{3 pts}{Réalisez les tests, notez la commande et le résultat, puis faites viser chaque test réussi.}
{\small\renewcommand{\arraystretch}{1.1}
\begin{tabularx}{\linewidth}{|c|l|l|C|c|c|}
  \hline
  \textbf{N°} & \textbf{Source} & \textbf{Destination} & \textbf{Commande / action} & \textbf{Résultat} & \textbf{Visa prof.} \\ \hline
  \hl 1 & PC-\pUn1 & PC-\pUn4 & \sol{ping \netUn.4} & \sol{réussi} & \hspace{1.4cm} \\ \hline
  \hl 2 & PC-\pDeux1 & IMP-\pDeux & \sol{ping \netDeux.0.10} & \sol{réussi} & \\ \hline
  \hl 3 & PC-\pUn2 & PC-\pDeux2 & \sol{ping \netDeux.0.2} & \sol{réussi} & \\ \hline
  \hl 4 & PC-\pDeux2 & \srvNom & \sol{ping \netSrv.\srvH} & \sol{réussi} & \\ \hline
  \hl 5 & PC-\pUn3 & page web du serveur & \sol{navigateur : {\NoAutoSpacing http://\netSrv.\srvH}} & \sol{page affichée} & \\ \hline
\end{tabularx}}
\bareme{0,5 pt par test visé + 0,5 pt pour B5. Le premier ping inter-réseaux peut échouer (résolution ARP) ; les suivants doivent réussir.}

\question{B5}{incl. B4}{Test n°3 : par quels équipements passe le paquet ? Pourquoi PC-\pUn2 l'envoie-t-il au routeur ?}
\zone{2.3cm}{PC-\pUn2 $\to$ SW-\pUn{} $\to$ \Rnom{} $\to$ SW-\pDeux{} $\to$ PC-\pDeux2. Grâce à son masque, PC-\pUn2 constate que
la destination (\netDeux.0.0/16) n'est \textbf{pas dans son réseau} (\netUn.0/24) : il envoie le paquet à sa
\textbf{passerelle par défaut}, qui le route.}

%==========================================================================
\newpage
\section{Partie C — Diagnostic d'un réseau en panne \hfill /4}

Un technicien constate qu'\textbf{aucun poste ne parvient à joindre le serveur SRV-C} et que certains postes
communiquent mal entre eux. Les câbles sont du bon type.

\begin{center}
\begin{minipage}{0.37\linewidth}\centering
\begin{tikzpicture}[font=\footnotesize,
  eq/.style={draw,fill=qbg,rounded corners=2pt,minimum width=9mm,minimum height=5mm,inner sep=2pt},
  every path/.style={thick}]
  \node[eq] (SW) {SW-1};
  \node[eq,right=9mm of SW] (R) {R1};
  \node[eq,right=7mm of R] (SRV) {SRV-C};
  \node[eq,above left=5mm and -3mm of SW] (A) {PC-A};
  \node[eq,above right=5mm and -3mm of SW] (B) {PC-B};
  \node[eq,below left=5mm and -3mm of SW] (C) {PC-C};
  \node[eq,below right=5mm and -3mm of SW] (D) {PC-D};
  \foreach \n in {A,B,C,D} \draw (\n) -- (SW);
  \draw (SW) -- node[above,font=\tiny]{G0/0} (R);
  \draw (R) -- node[above,font=\tiny]{G0/1} (SRV);
\end{tikzpicture}
\end{minipage}\hfill
\begin{minipage}{0.61\linewidth}\centering
{\small\renewcommand{\arraystretch}{1.05}\setlength{\tabcolsep}{4pt}
\begin{tabular}{lllll}
  \toprule
  \textbf{Équipement} & \textbf{Adresse IP} & \textbf{Masque} & \textbf{Passerelle} & \textbf{Port} \\
  \midrule
  \Crows
  \bottomrule
\end{tabular}}
\end{minipage}
\end{center}

\question{C1}{4 pts}{Relevez les \textbf{quatre anomalies}. Pour chacune : conséquence et correction.}
{\small\renewcommand{\arraystretch}{1.1}
\ifcorrige\def\hc{}\else\def\hc{\rule[-10mm]{0pt}{14mm}}\fi
\begin{tabularx}{\linewidth}{|c|>{\raggedright\arraybackslash}p{2cm}|L|L|L|}
  \hline
  \textbf{N°} & \textbf{Équipement} & \textbf{Anomalie} & \textbf{Conséquence} & \textbf{Correction} \\ \hline
  \ifcorrige\Csol\else
  \hc 1 & & & & \\ \hline
  \hc 2 & & & & \\ \hline
  \hc 3 & & & & \\ \hline
  \hc 4 & & & & \\ \hline
  \fi
\end{tabularx}}
\bareme{1 pt par anomalie (0,25 équipement, anomalie, conséquence, correction ; ordre indifférent). \Cbareme}

%==========================================================================
\section{Bonus — Extension sans fil \hfill +1}

\begin{cdc}Ajoutez au réseau \resUn{} un \textbf{AccessPoint-PT} relié à SW-\pUn{} (SSID \ip{\ssid}, \textbf{WPA2-PSK})
et un portable \ip{Laptop-\pUn} équipé d'un module \texttt{WPC300N}, en adressage \textbf{statique}.
Vérifiez qu'il joint \srvNom.\end{cdc}
{\small\renewcommand{\arraystretch}{1.1}
\begin{tabularx}{\linewidth}{|l|C|C|C|c|}
  \hline
  \textbf{Équipement} & \textbf{Adresse IP} & \textbf{Masque} & \textbf{Passerelle} & \textbf{Visa prof.} \\ \hline
  \hl Laptop-\pUn & \sol{\netUn.50} & \sol{255.255.255.0} & \sol{\netUn.254} & \hspace{1.4cm} \\ \hline
\end{tabularx}}
\bareme{Le point d'accès fait office de pont : le portable est dans \netUn.0/24. Bonus si associé en WPA2-PSK et ping
vers \netSrv.\srvH{} réussi. Note plafonnée à 20.}

\end{document}
"
%   pdflatex -jobname=Eval_Reseau_CorrigeB "\def\variante{B}\def\modecorrige{1}% Évaluation pratique — Réseaux locaux (BTS CIEL 1re année) — sujets A et B
% Compilation (4 PDF) :
%   pdflatex -jobname=Eval_Reseau_SujetA   "\def\variante{A}\def\modecorrige{0}% Évaluation pratique — Réseaux locaux (BTS CIEL 1re année) — sujets A et B
% Compilation (4 PDF) :
%   pdflatex -jobname=Eval_Reseau_SujetA   "\def\variante{A}\def\modecorrige{0}% Évaluation pratique — Réseaux locaux (BTS CIEL 1re année) — sujets A et B
% Compilation (4 PDF) :
%   pdflatex -jobname=Eval_Reseau_SujetA   "\def\variante{A}\def\modecorrige{0}\input{Eval_TP_Reseau_LAN.tex}"
%   pdflatex -jobname=Eval_Reseau_SujetB   "\def\variante{B}\def\modecorrige{0}\input{Eval_TP_Reseau_LAN.tex}"
%   pdflatex -jobname=Eval_Reseau_CorrigeA "\def\variante{A}\def\modecorrige{1}\input{Eval_TP_Reseau_LAN.tex}"
%   pdflatex -jobname=Eval_Reseau_CorrigeB "\def\variante{B}\def\modecorrige{1}\input{Eval_TP_Reseau_LAN.tex}"
\documentclass[10pt,a4paper]{article}

\usepackage[T1]{fontenc}
\usepackage[utf8]{inputenc}
\usepackage{lmodern}
\usepackage[french]{babel}
\usepackage[a4paper,hmargin=1.5cm,top=1.3cm,bottom=1.6cm,footskip=0.7cm]{geometry}
\usepackage[table]{xcolor}
\usepackage{booktabs,tabularx,array}
\usepackage{enumitem}
\usepackage{fancyhdr}
\usepackage[most]{tcolorbox}
\usepackage{tikz}
\usetikzlibrary{positioning}
\usepackage{titlesec}

\providecommand{\modecorrige}{0}
\providecommand{\variante}{A}
\newif\ifcorrige
\ifnum\modecorrige=1 \corrigetrue \else \corrigefalse \fi

\definecolor{ipred}{HTML}{C0392B}
\definecolor{cdcbar}{HTML}{A67C3D}
\definecolor{cdcbg}{HTML}{F4F1EA}
\definecolor{qbar}{HTML}{4A6B5D}
\definecolor{qbg}{HTML}{EDF1EF}
\definecolor{corr}{HTML}{1F5FA8}
\definecolor{corrbg}{HTML}{E8F0FA}
\definecolor{grisfoot}{HTML}{828282}

\newcommand{\ip}[1]{\textcolor{ipred}{\texttt{#1}}}
\newcommand{\pts}[1]{\hfill\mbox{\textbf{(#1)}}}
\newcommand{\sol}[1]{\ifcorrige\textcolor{corr}{#1}\fi}
\newcommand{\bareme}[1]{\ifcorrige\par{\footnotesize\color{corr}#1\par}\fi}
\newcommand{\hl}{\ifcorrige\rule[-1.6mm]{0pt}{5.4mm}\else\rule[-2.2mm]{0pt}{6.6mm}\fi}
\newcolumntype{C}{>{\centering\arraybackslash}X}
\newcolumntype{L}{>{\raggedright\arraybackslash}X}

\titleformat{\section}{\large\bfseries}{}{0pt}{}
\titlespacing*{\section}{0pt}{6pt}{3pt}

\tcbset{qstyle/.style={enhanced,boxrule=0pt,leftrule=2.5pt,arc=0pt,outer arc=0pt,
  left=5pt,right=5pt,top=2pt,bottom=2pt,before skip=4pt,after skip=3pt}}
\newtcolorbox{cdc}{qstyle,colback=cdcbg,colframe=cdcbar}
\newtcolorbox{quest}{qstyle,colback=qbg,colframe=qbar}

% Zone de réponse : vierge dans le sujet, remplie dans le corrigé
\newcommand{\zone}[2]{%
  \ifcorrige
    \begin{tcolorbox}[colback=corrbg,colframe=corr,boxrule=0.5pt,arc=1pt,left=4pt,right=4pt,top=1pt,bottom=1pt,
      before skip=2pt,after skip=2pt,fontupper=\footnotesize\color{corr}]#2\end{tcolorbox}%
  \else
    \begin{tcolorbox}[colback=white,colframe=black!45,boxrule=0.4pt,arc=1pt,height=#1,
      before skip=2pt,after skip=2pt]\end{tcolorbox}%
  \fi}

\newcommand{\question}[3]{\begin{quest}\textbf{#1.} #3\pts{#2}\end{quest}}

%==========================================================================
% Données propres à chaque variante
%==========================================================================
\if\variante A
  \def\entreprise{Le cabinet d'architectes \textbf{ArchiNova}}
  \def\Rnom{R-ArchiNova}
  \def\resUn{Bureau d'études}\def\pUn{BE}\def\netUn{192.168.40}
  \def\resDeux{Accueil}\def\pDeux{ACC}\def\netDeux{172.16}
  \def\srvNom{SRV-Intranet}\def\netSrv{192.168.200}\def\srvH{10}
  \def\ssid{ArchiNova-WiFi}
  \def\AIrows{%
    \hl\ip{192.168.40.25} & \sol{C} & \sol{255.255.255.0} & \sol{192.168.40.0} & \sol{192.168.40.255} & \sol{254} \\ \hline
    \hl\ip{172.16.5.10}   & \sol{B} & \sol{255.255.0.0}   & \sol{172.16.0.0}   & \sol{172.16.255.255} & \sol{65\,534} \\ \hline
    \hl\ip{10.3.2.1}      & \sol{A} & \sol{255.0.0.0}     & \sol{10.0.0.0}     & \sol{10.255.255.255} & \sol{16\,777\,214} \\ \hline}
  \def\apipa{169.254.12.7}
  \def\Cnet{192.168.60}
  \def\Crows{%
    PC-A & \ip{192.168.60.10} & \ip{255.255.255.0} & \ip{192.168.60.254} & — \\
    PC-B & \ip{192.168.60.20} & \ip{255.255.255.0} & \ip{192.168.60.1}   & — \\
    PC-C & \ip{192.168.61.30} & \ip{255.255.255.0} & \ip{192.168.60.254} & — \\
    PC-D & \ip{192.168.60.10} & \ip{255.255.255.0} & \ip{192.168.60.254} & — \\
    SRV-C & \ip{10.0.0.5}     & \ip{255.0.0.0}     & \ip{10.0.0.254}     & — \\
    R1 — G0/0 & \ip{192.168.60.254} & \ip{255.255.255.0} & — & On \\
    R1 — G0/1 & \ip{10.0.0.254}     & \ip{255.0.0.0}     & — & Off \\}
  \def\Csol{%
    \hl 1 & \sol{PC-B} & \sol{Passerelle .60.1 : pas l'adresse du routeur.} & \sol{Joint son réseau, mais aucun autre (serveur).} & \sol{Passerelle 192.168.60.254.} \\ \hline
    \hl 2 & \sol{PC-C} & \sol{IP hors du réseau 192.168.60.0/24.} & \sol{Ne joint ni les postes ni le routeur.} & \sol{Ex. 192.168.60.30.} \\ \hline
    \hl 3 & \sol{PC-D} & \sol{Même IP que PC-A (doublon).} & \sol{Conflit d'adresse : PC-A et PC-D perturbés.} & \sol{Adresse unique, ex. .60.40.} \\ \hline
    \hl 4 & \sol{R1 G0/1} & \sol{Interface désactivée (Off).} & \sol{Réseau du serveur injoignable pour tous.} & \sol{Port Status : On.} \\ \hline}
  \def\Cbareme{PC-A et SRV-C sont correctement configurés ; accepter PC-A cité à la place de PC-D pour le doublon.}
\else
  \def\entreprise{La clinique vétérinaire \textbf{VetoPlus}}
  \def\Rnom{R-VetoPlus}
  \def\resUn{Soins}\def\pUn{SOI}\def\netUn{192.168.70}
  \def\resDeux{Secrétariat}\def\pDeux{SEC}\def\netDeux{172.20}
  \def\srvNom{SRV-Dossiers}\def\netSrv{192.168.250}\def\srvH{20}
  \def\ssid{VetoPlus-WiFi}
  \def\AIrows{%
    \hl\ip{192.168.15.130} & \sol{C} & \sol{255.255.255.0} & \sol{192.168.15.0} & \sol{192.168.15.255} & \sol{254} \\ \hline
    \hl\ip{172.31.200.7}   & \sol{B} & \sol{255.255.0.0}   & \sol{172.31.0.0}   & \sol{172.31.255.255} & \sol{65\,534} \\ \hline
    \hl\ip{10.45.6.9}      & \sol{A} & \sol{255.0.0.0}     & \sol{10.0.0.0}     & \sol{10.255.255.255} & \sol{16\,777\,214} \\ \hline}
  \def\apipa{169.254.88.3}
  \def\Cnet{192.168.80}
  \def\Crows{%
    PC-A & \ip{192.168.8.11}   & \ip{255.255.255.0} & \ip{192.168.80.254} & — \\
    PC-B & \ip{192.168.80.12}  & \ip{255.255.255.0} & \ip{192.168.80.254} & — \\
    PC-C & \ip{192.168.80.13}  & \ip{255.255.255.0} & \ip{192.168.80.100} & — \\
    PC-D & \ip{192.168.80.12}  & \ip{255.255.255.0} & \ip{192.168.80.254} & — \\
    SRV-C & \ip{10.0.0.8}      & \ip{255.0.0.0}     & \ip{10.0.0.254}     & — \\
    R1 — G0/0 & \ip{192.168.80.254} & \ip{255.255.255.0} & — & Off \\
    R1 — G0/1 & \ip{10.0.0.254}     & \ip{255.0.0.0}     & — & On \\}
  \def\Csol{%
    \hl 1 & \sol{PC-A} & \sol{IP hors du réseau 192.168.80.0/24.} & \sol{Ne joint ni les postes ni le routeur.} & \sol{Ex. 192.168.80.11.} \\ \hline
    \hl 2 & \sol{PC-C} & \sol{Passerelle .80.100 : pas l'adresse du routeur.} & \sol{Joint son réseau, mais aucun autre (serveur).} & \sol{Passerelle 192.168.80.254.} \\ \hline
    \hl 3 & \sol{PC-D} & \sol{Même IP que PC-B (doublon).} & \sol{Conflit d'adresse : PC-B et PC-D perturbés.} & \sol{Adresse unique, ex. .80.14.} \\ \hline
    \hl 4 & \sol{R1 G0/0} & \sol{Interface côté LAN désactivée (Off).} & \sol{Aucun poste ne joint le routeur ni le serveur.} & \sol{Port Status : On.} \\ \hline}
  \def\Cbareme{PC-B et SRV-C sont correctement configurés ; accepter PC-B cité à la place de PC-D pour le doublon.}
\fi

\pagestyle{fancy}
\fancyhf{}
\renewcommand{\headrulewidth}{0pt}
\fancyfoot[L]{\footnotesize\itshape\color{grisfoot}Évaluation — Réseaux locaux — Sujet \variante\ifcorrige\ — CORRIGÉ\fi}
\fancyfoot[R]{\footnotesize\color{grisfoot}\thepage/3}

\setlist{itemsep=0pt,topsep=2pt,parsep=0pt}
\setlength{\parindent}{0pt}
\setlength{\parskip}{2pt}

\begin{document}

\begin{center}
  {\Large\bfseries Évaluation pratique — Réseaux locaux \quad
   \fcolorbox{black}{white}{\ SUJET \variante\ }\par}
  \vspace{1pt}
  {\bfseries Adressage IP, commutation et routage — BTS CIEL 1re année — durée : 2 heures\par}
  \ifcorrige\vspace{2pt}{\bfseries\color{corr}CORRIGÉ ET BARÈME — document enseignant}\par\fi
\end{center}
\vspace{-2pt}
{\renewcommand{\arraystretch}{1.35}
\begin{tabularx}{\linewidth}{|X|X|p{2.2cm}|p{2.6cm}|}
  \hline
  \textbf{Nom :} & \textbf{Prénom :} & \textbf{Poste n°} & \textbf{Note :} \hfill /20 \\ \hline
\end{tabularx}}

\textbf{Consignes :} travail \textbf{individuel}, documents et Internet \textbf{non autorisés}. Ce sujet est le
\textbf{document-réponse}. Fichier Packet Tracer : \ip{Nom\_Prenom\_EvalLAN\_\variante.pkt}, à enregistrer
régulièrement dans le dossier de rendu. Les « visas » sont apposés par le professeur, que vous appelez à chaque test réussi.
Durées conseillées : A 30 min — B 1 h 10 — C 20 min.

%==========================================================================
\section{Partie A — Questions de cours et calculs \hfill /6}

\question{A1}{2 pts}{Complétez le tableau en considérant le \textbf{masque par défaut} de la classe de chaque adresse.}
{\small\renewcommand{\arraystretch}{1.15}
\begin{tabularx}{\linewidth}{|l|c|C|C|C|c|}
  \hline
  \textbf{Adresse IP} & \textbf{Classe} & \textbf{Masque par défaut} & \textbf{Adresse réseau} &
  \textbf{Adresse de diffusion} & \textbf{Nb d'hôtes} \\ \hline
  \AIrows
\end{tabularx}}
\bareme{0,5 pt par ligne juste + 0,5 pt si les trois nombres d'hôtes sont justes ($2^{n}-2$, $n$ = nombre de bits d'hôte).}

\question{A2}{1 pt}{Indiquez le type de câble (\textbf{droit} ou \textbf{croisé}) pour chaque liaison.}
{\small\renewcommand{\arraystretch}{1.15}
\begin{tabularx}{\linewidth}{|X|c||X|c|}
  \hline
  \hl PC $\leftrightarrow$ commutateur & \makebox[1.6cm]{\sol{droit}} & Commutateur $\leftrightarrow$ routeur & \makebox[1.6cm]{\sol{droit}} \\ \hline
  \hl PC $\leftrightarrow$ PC (réseau P2P) & \sol{croisé} & Concentrateur $\leftrightarrow$ concentrateur & \sol{croisé} \\ \hline
\end{tabularx}}
\bareme{0,25 pt par case. Équipements de types différents : câble droit ; de même type : câble croisé.}

\question{A3}{1 pt}{Comment un \textbf{concentrateur (hub)} et un \textbf{commutateur (switch)} traitent-ils une trame
reçue ? Lequel limite les collisions, et pourquoi ?}
\zone{2.9cm}{Le \textbf{hub} (couche 1) répète la trame sur \textbf{tous ses autres ports} : un seul domaine de collision
partagé. Le \textbf{switch} (couche 2) lit l'adresse MAC de destination et, grâce à sa \textbf{table MAC}, ne l'envoie
\textbf{que sur le port du destinataire} : un domaine de collision par port, c'est lui qui limite les collisions.
\emph{0,5 pt hub, 0,5 pt switch.}}

\question{A4}{1 pt}{Qu'est-ce que la \textbf{passerelle par défaut} d'un poste ? Quelle condition son adresse doit-elle respecter ?}
\zone{2.4cm}{Adresse IP de l'interface du \textbf{routeur} à laquelle le poste envoie les paquets destinés à un
\textbf{autre réseau}. Elle doit \textbf{appartenir au même réseau} que le poste. \emph{0,5 + 0,5 pt.}}

\question{A5}{0,5 pt}{Un portable configuré en DHCP affiche \ip{\apipa} / \ip{255.255.0.0}. Que s'est-il passé ?}
\zone{1.5cm}{Aucun \textbf{serveur DHCP} n'a répondu : le poste s'est attribué une adresse \textbf{APIPA} (169.254.0.0/16).}

\question{A6}{0,5 pt}{Quel protocole utilise \texttt{ping} ? Nommez les deux messages échangés et la couche OSI des adresses IP.}
\zone{1.5cm}{\textbf{ICMP} ; \textbf{Echo Request} / \textbf{Echo Reply} ; adresses IP en \textbf{couche 3 (Réseau)}.}

%==========================================================================
\newpage
\section{Partie B — Réalisation sous Packet Tracer \hfill /10}

\begin{cdc}
\entreprise{} vous confie son réseau : \textbf{trois réseaux} reliés par un \textbf{routeur unique Cisco 2911}
nommé \ip{\Rnom} (interfaces G0/0, G0/1, G0/2).
\begin{itemize}[leftmargin=12pt]
  \item \textbf{Réseau \resUn} : \textbf{4 postes} PC-\pUn1 à PC-\pUn4, commutateur \ip{SW-\pUn}, adressage statique,
        réseau \ip{\netUn.0/24}.
  \item \textbf{Réseau \resDeux} : \textbf{2 postes} PC-\pDeux1, PC-\pDeux2 et \textbf{1 imprimante} IMP-\pDeux,
        commutateur \ip{SW-\pDeux}, adressage statique, réseau de classe B \ip{\netDeux.0.0/16}.
  \item \textbf{Réseau Serveur} : serveur \ip{\srvNom} d'adresse \ip{\netSrv.\srvH/24}, relié \textbf{directement} au routeur.
  \item La passerelle de chaque réseau est sa \textbf{dernière adresse utilisable}. Tous les équipements communiquent
        entre eux et accèdent à la page web du serveur. Nommage imposé (« Display Name »).
\end{itemize}
\end{cdc}

\question{B1}{2 pts}{Complétez le plan d'adressage \textbf{avant} la réalisation.}
{\footnotesize\setlength{\tabcolsep}{3pt}\renewcommand{\arraystretch}{1.15}
\begin{tabularx}{\linewidth}{|l|C|C|C|C|C|C|}
  \hline
  \textbf{Réseau} & \textbf{Adresse réseau} & \textbf{Masque} & \textbf{1re @ utilisable} &
  \textbf{Dernière @ utilisable} & \textbf{Diffusion} & \textbf{Passerelle} \\ \hline
  \hl \resUn & \sol{\netUn.0} & \sol{255.255.255.0} & \sol{\netUn.1} & \sol{\netUn.254} & \sol{\netUn.255} & \sol{\netUn.254} \\ \hline
  \hl \resDeux & \sol{\netDeux.0.0} & \sol{255.255.0.0} & \sol{\netDeux.0.1} & \sol{\netDeux.255.254} & \sol{\netDeux.255.255} & \sol{\netDeux.255.254} \\ \hline
  \hl Serveur & \sol{\netSrv.0} & \sol{255.255.255.0} & \sol{\netSrv.1} & \sol{\netSrv.254} & \sol{\netSrv.255} & \sol{\netSrv.254} \\ \hline
\end{tabularx}}
\bareme{0,5 pt par ligne juste + 0,5 pt si la convention de passerelle est respectée partout.}

\question{B2}{2 pts}{Construisez la topologie (équipements, câblage, nommage) et relevez l'adressage configuré.}
{\footnotesize\setlength{\tabcolsep}{3pt}\renewcommand{\arraystretch}{1.1}
\begin{tabularx}{\linewidth}{|l|C|C|C||l|C|C|}
  \hline
  \textbf{Hôte} & \textbf{Adresse IP} & \textbf{Masque} & \textbf{Passerelle} & \textbf{Interface} & \textbf{Adresse IP} & \textbf{Masque} \\ \hline
  \hl PC-\pUn1 & \sol{\netUn.1} & \sol{255.255.255.0} & \sol{\netUn.254} & \Rnom{} G0/0 & \sol{\netUn.254} & \sol{255.255.255.0} \\ \hline
  \hl PC-\pDeux1 & \sol{\netDeux.0.1} & \sol{255.255.0.0} & \sol{\netDeux.255.254} & \Rnom{} G0/1 & \sol{\netDeux.255.254} & \sol{255.255.0.0} \\ \hline
  \hl IMP-\pDeux & \sol{\netDeux.0.10} & \sol{255.255.0.0} & \sol{\netDeux.255.254} & \Rnom{} G0/2 & \sol{\netSrv.254} & \sol{255.255.255.0} \\ \hline
  \hl \srvNom & \ip{\netSrv.\srvH} & \sol{255.255.255.0} & \sol{\netSrv.254} & \multicolumn{3}{c|}{\cellcolor{black!8}} \\ \hline
\end{tabularx}}
\bareme{Adresses d'hôtes libres (uniques, dans le bon réseau) ; affectation G0/x libre. Sur le .pkt : équipements 0,5 ;
câblage 1 (câbles droits vers les switchs et switch–routeur, \textbf{croisé} serveur–routeur, liens verts) ; nommage 0,5.}

\question{B3}{3 pts}{Configurez tous les équipements : IP, masques, passerelles des hôtes ; interfaces du routeur
(« Port Status : On »).}
\bareme{Sur le .pkt : interfaces routeur 1 pt ; IP/masques des hôtes 1 pt ; passerelles (imprimante et serveur compris)
1 pt. $-0,25$ par erreur, sans descendre sous 0 par item.}

\question{B4}{3 pts}{Réalisez les tests, notez la commande et le résultat, puis faites viser chaque test réussi.}
{\small\renewcommand{\arraystretch}{1.1}
\begin{tabularx}{\linewidth}{|c|l|l|C|c|c|}
  \hline
  \textbf{N°} & \textbf{Source} & \textbf{Destination} & \textbf{Commande / action} & \textbf{Résultat} & \textbf{Visa prof.} \\ \hline
  \hl 1 & PC-\pUn1 & PC-\pUn4 & \sol{ping \netUn.4} & \sol{réussi} & \hspace{1.4cm} \\ \hline
  \hl 2 & PC-\pDeux1 & IMP-\pDeux & \sol{ping \netDeux.0.10} & \sol{réussi} & \\ \hline
  \hl 3 & PC-\pUn2 & PC-\pDeux2 & \sol{ping \netDeux.0.2} & \sol{réussi} & \\ \hline
  \hl 4 & PC-\pDeux2 & \srvNom & \sol{ping \netSrv.\srvH} & \sol{réussi} & \\ \hline
  \hl 5 & PC-\pUn3 & page web du serveur & \sol{navigateur : {\NoAutoSpacing http://\netSrv.\srvH}} & \sol{page affichée} & \\ \hline
\end{tabularx}}
\bareme{0,5 pt par test visé + 0,5 pt pour B5. Le premier ping inter-réseaux peut échouer (résolution ARP) ; les suivants doivent réussir.}

\question{B5}{incl. B4}{Test n°3 : par quels équipements passe le paquet ? Pourquoi PC-\pUn2 l'envoie-t-il au routeur ?}
\zone{2.3cm}{PC-\pUn2 $\to$ SW-\pUn{} $\to$ \Rnom{} $\to$ SW-\pDeux{} $\to$ PC-\pDeux2. Grâce à son masque, PC-\pUn2 constate que
la destination (\netDeux.0.0/16) n'est \textbf{pas dans son réseau} (\netUn.0/24) : il envoie le paquet à sa
\textbf{passerelle par défaut}, qui le route.}

%==========================================================================
\newpage
\section{Partie C — Diagnostic d'un réseau en panne \hfill /4}

Un technicien constate qu'\textbf{aucun poste ne parvient à joindre le serveur SRV-C} et que certains postes
communiquent mal entre eux. Les câbles sont du bon type.

\begin{center}
\begin{minipage}{0.37\linewidth}\centering
\begin{tikzpicture}[font=\footnotesize,
  eq/.style={draw,fill=qbg,rounded corners=2pt,minimum width=9mm,minimum height=5mm,inner sep=2pt},
  every path/.style={thick}]
  \node[eq] (SW) {SW-1};
  \node[eq,right=9mm of SW] (R) {R1};
  \node[eq,right=7mm of R] (SRV) {SRV-C};
  \node[eq,above left=5mm and -3mm of SW] (A) {PC-A};
  \node[eq,above right=5mm and -3mm of SW] (B) {PC-B};
  \node[eq,below left=5mm and -3mm of SW] (C) {PC-C};
  \node[eq,below right=5mm and -3mm of SW] (D) {PC-D};
  \foreach \n in {A,B,C,D} \draw (\n) -- (SW);
  \draw (SW) -- node[above,font=\tiny]{G0/0} (R);
  \draw (R) -- node[above,font=\tiny]{G0/1} (SRV);
\end{tikzpicture}
\end{minipage}\hfill
\begin{minipage}{0.61\linewidth}\centering
{\small\renewcommand{\arraystretch}{1.05}\setlength{\tabcolsep}{4pt}
\begin{tabular}{lllll}
  \toprule
  \textbf{Équipement} & \textbf{Adresse IP} & \textbf{Masque} & \textbf{Passerelle} & \textbf{Port} \\
  \midrule
  \Crows
  \bottomrule
\end{tabular}}
\end{minipage}
\end{center}

\question{C1}{4 pts}{Relevez les \textbf{quatre anomalies}. Pour chacune : conséquence et correction.}
{\small\renewcommand{\arraystretch}{1.1}
\ifcorrige\def\hc{}\else\def\hc{\rule[-10mm]{0pt}{14mm}}\fi
\begin{tabularx}{\linewidth}{|c|>{\raggedright\arraybackslash}p{2cm}|L|L|L|}
  \hline
  \textbf{N°} & \textbf{Équipement} & \textbf{Anomalie} & \textbf{Conséquence} & \textbf{Correction} \\ \hline
  \ifcorrige\Csol\else
  \hc 1 & & & & \\ \hline
  \hc 2 & & & & \\ \hline
  \hc 3 & & & & \\ \hline
  \hc 4 & & & & \\ \hline
  \fi
\end{tabularx}}
\bareme{1 pt par anomalie (0,25 équipement, anomalie, conséquence, correction ; ordre indifférent). \Cbareme}

%==========================================================================
\section{Bonus — Extension sans fil \hfill +1}

\begin{cdc}Ajoutez au réseau \resUn{} un \textbf{AccessPoint-PT} relié à SW-\pUn{} (SSID \ip{\ssid}, \textbf{WPA2-PSK})
et un portable \ip{Laptop-\pUn} équipé d'un module \texttt{WPC300N}, en adressage \textbf{statique}.
Vérifiez qu'il joint \srvNom.\end{cdc}
{\small\renewcommand{\arraystretch}{1.1}
\begin{tabularx}{\linewidth}{|l|C|C|C|c|}
  \hline
  \textbf{Équipement} & \textbf{Adresse IP} & \textbf{Masque} & \textbf{Passerelle} & \textbf{Visa prof.} \\ \hline
  \hl Laptop-\pUn & \sol{\netUn.50} & \sol{255.255.255.0} & \sol{\netUn.254} & \hspace{1.4cm} \\ \hline
\end{tabularx}}
\bareme{Le point d'accès fait office de pont : le portable est dans \netUn.0/24. Bonus si associé en WPA2-PSK et ping
vers \netSrv.\srvH{} réussi. Note plafonnée à 20.}

\end{document}
"
%   pdflatex -jobname=Eval_Reseau_SujetB   "\def\variante{B}\def\modecorrige{0}% Évaluation pratique — Réseaux locaux (BTS CIEL 1re année) — sujets A et B
% Compilation (4 PDF) :
%   pdflatex -jobname=Eval_Reseau_SujetA   "\def\variante{A}\def\modecorrige{0}\input{Eval_TP_Reseau_LAN.tex}"
%   pdflatex -jobname=Eval_Reseau_SujetB   "\def\variante{B}\def\modecorrige{0}\input{Eval_TP_Reseau_LAN.tex}"
%   pdflatex -jobname=Eval_Reseau_CorrigeA "\def\variante{A}\def\modecorrige{1}\input{Eval_TP_Reseau_LAN.tex}"
%   pdflatex -jobname=Eval_Reseau_CorrigeB "\def\variante{B}\def\modecorrige{1}\input{Eval_TP_Reseau_LAN.tex}"
\documentclass[10pt,a4paper]{article}

\usepackage[T1]{fontenc}
\usepackage[utf8]{inputenc}
\usepackage{lmodern}
\usepackage[french]{babel}
\usepackage[a4paper,hmargin=1.5cm,top=1.3cm,bottom=1.6cm,footskip=0.7cm]{geometry}
\usepackage[table]{xcolor}
\usepackage{booktabs,tabularx,array}
\usepackage{enumitem}
\usepackage{fancyhdr}
\usepackage[most]{tcolorbox}
\usepackage{tikz}
\usetikzlibrary{positioning}
\usepackage{titlesec}

\providecommand{\modecorrige}{0}
\providecommand{\variante}{A}
\newif\ifcorrige
\ifnum\modecorrige=1 \corrigetrue \else \corrigefalse \fi

\definecolor{ipred}{HTML}{C0392B}
\definecolor{cdcbar}{HTML}{A67C3D}
\definecolor{cdcbg}{HTML}{F4F1EA}
\definecolor{qbar}{HTML}{4A6B5D}
\definecolor{qbg}{HTML}{EDF1EF}
\definecolor{corr}{HTML}{1F5FA8}
\definecolor{corrbg}{HTML}{E8F0FA}
\definecolor{grisfoot}{HTML}{828282}

\newcommand{\ip}[1]{\textcolor{ipred}{\texttt{#1}}}
\newcommand{\pts}[1]{\hfill\mbox{\textbf{(#1)}}}
\newcommand{\sol}[1]{\ifcorrige\textcolor{corr}{#1}\fi}
\newcommand{\bareme}[1]{\ifcorrige\par{\footnotesize\color{corr}#1\par}\fi}
\newcommand{\hl}{\ifcorrige\rule[-1.6mm]{0pt}{5.4mm}\else\rule[-2.2mm]{0pt}{6.6mm}\fi}
\newcolumntype{C}{>{\centering\arraybackslash}X}
\newcolumntype{L}{>{\raggedright\arraybackslash}X}

\titleformat{\section}{\large\bfseries}{}{0pt}{}
\titlespacing*{\section}{0pt}{6pt}{3pt}

\tcbset{qstyle/.style={enhanced,boxrule=0pt,leftrule=2.5pt,arc=0pt,outer arc=0pt,
  left=5pt,right=5pt,top=2pt,bottom=2pt,before skip=4pt,after skip=3pt}}
\newtcolorbox{cdc}{qstyle,colback=cdcbg,colframe=cdcbar}
\newtcolorbox{quest}{qstyle,colback=qbg,colframe=qbar}

% Zone de réponse : vierge dans le sujet, remplie dans le corrigé
\newcommand{\zone}[2]{%
  \ifcorrige
    \begin{tcolorbox}[colback=corrbg,colframe=corr,boxrule=0.5pt,arc=1pt,left=4pt,right=4pt,top=1pt,bottom=1pt,
      before skip=2pt,after skip=2pt,fontupper=\footnotesize\color{corr}]#2\end{tcolorbox}%
  \else
    \begin{tcolorbox}[colback=white,colframe=black!45,boxrule=0.4pt,arc=1pt,height=#1,
      before skip=2pt,after skip=2pt]\end{tcolorbox}%
  \fi}

\newcommand{\question}[3]{\begin{quest}\textbf{#1.} #3\pts{#2}\end{quest}}

%==========================================================================
% Données propres à chaque variante
%==========================================================================
\if\variante A
  \def\entreprise{Le cabinet d'architectes \textbf{ArchiNova}}
  \def\Rnom{R-ArchiNova}
  \def\resUn{Bureau d'études}\def\pUn{BE}\def\netUn{192.168.40}
  \def\resDeux{Accueil}\def\pDeux{ACC}\def\netDeux{172.16}
  \def\srvNom{SRV-Intranet}\def\netSrv{192.168.200}\def\srvH{10}
  \def\ssid{ArchiNova-WiFi}
  \def\AIrows{%
    \hl\ip{192.168.40.25} & \sol{C} & \sol{255.255.255.0} & \sol{192.168.40.0} & \sol{192.168.40.255} & \sol{254} \\ \hline
    \hl\ip{172.16.5.10}   & \sol{B} & \sol{255.255.0.0}   & \sol{172.16.0.0}   & \sol{172.16.255.255} & \sol{65\,534} \\ \hline
    \hl\ip{10.3.2.1}      & \sol{A} & \sol{255.0.0.0}     & \sol{10.0.0.0}     & \sol{10.255.255.255} & \sol{16\,777\,214} \\ \hline}
  \def\apipa{169.254.12.7}
  \def\Cnet{192.168.60}
  \def\Crows{%
    PC-A & \ip{192.168.60.10} & \ip{255.255.255.0} & \ip{192.168.60.254} & — \\
    PC-B & \ip{192.168.60.20} & \ip{255.255.255.0} & \ip{192.168.60.1}   & — \\
    PC-C & \ip{192.168.61.30} & \ip{255.255.255.0} & \ip{192.168.60.254} & — \\
    PC-D & \ip{192.168.60.10} & \ip{255.255.255.0} & \ip{192.168.60.254} & — \\
    SRV-C & \ip{10.0.0.5}     & \ip{255.0.0.0}     & \ip{10.0.0.254}     & — \\
    R1 — G0/0 & \ip{192.168.60.254} & \ip{255.255.255.0} & — & On \\
    R1 — G0/1 & \ip{10.0.0.254}     & \ip{255.0.0.0}     & — & Off \\}
  \def\Csol{%
    \hl 1 & \sol{PC-B} & \sol{Passerelle .60.1 : pas l'adresse du routeur.} & \sol{Joint son réseau, mais aucun autre (serveur).} & \sol{Passerelle 192.168.60.254.} \\ \hline
    \hl 2 & \sol{PC-C} & \sol{IP hors du réseau 192.168.60.0/24.} & \sol{Ne joint ni les postes ni le routeur.} & \sol{Ex. 192.168.60.30.} \\ \hline
    \hl 3 & \sol{PC-D} & \sol{Même IP que PC-A (doublon).} & \sol{Conflit d'adresse : PC-A et PC-D perturbés.} & \sol{Adresse unique, ex. .60.40.} \\ \hline
    \hl 4 & \sol{R1 G0/1} & \sol{Interface désactivée (Off).} & \sol{Réseau du serveur injoignable pour tous.} & \sol{Port Status : On.} \\ \hline}
  \def\Cbareme{PC-A et SRV-C sont correctement configurés ; accepter PC-A cité à la place de PC-D pour le doublon.}
\else
  \def\entreprise{La clinique vétérinaire \textbf{VetoPlus}}
  \def\Rnom{R-VetoPlus}
  \def\resUn{Soins}\def\pUn{SOI}\def\netUn{192.168.70}
  \def\resDeux{Secrétariat}\def\pDeux{SEC}\def\netDeux{172.20}
  \def\srvNom{SRV-Dossiers}\def\netSrv{192.168.250}\def\srvH{20}
  \def\ssid{VetoPlus-WiFi}
  \def\AIrows{%
    \hl\ip{192.168.15.130} & \sol{C} & \sol{255.255.255.0} & \sol{192.168.15.0} & \sol{192.168.15.255} & \sol{254} \\ \hline
    \hl\ip{172.31.200.7}   & \sol{B} & \sol{255.255.0.0}   & \sol{172.31.0.0}   & \sol{172.31.255.255} & \sol{65\,534} \\ \hline
    \hl\ip{10.45.6.9}      & \sol{A} & \sol{255.0.0.0}     & \sol{10.0.0.0}     & \sol{10.255.255.255} & \sol{16\,777\,214} \\ \hline}
  \def\apipa{169.254.88.3}
  \def\Cnet{192.168.80}
  \def\Crows{%
    PC-A & \ip{192.168.8.11}   & \ip{255.255.255.0} & \ip{192.168.80.254} & — \\
    PC-B & \ip{192.168.80.12}  & \ip{255.255.255.0} & \ip{192.168.80.254} & — \\
    PC-C & \ip{192.168.80.13}  & \ip{255.255.255.0} & \ip{192.168.80.100} & — \\
    PC-D & \ip{192.168.80.12}  & \ip{255.255.255.0} & \ip{192.168.80.254} & — \\
    SRV-C & \ip{10.0.0.8}      & \ip{255.0.0.0}     & \ip{10.0.0.254}     & — \\
    R1 — G0/0 & \ip{192.168.80.254} & \ip{255.255.255.0} & — & Off \\
    R1 — G0/1 & \ip{10.0.0.254}     & \ip{255.0.0.0}     & — & On \\}
  \def\Csol{%
    \hl 1 & \sol{PC-A} & \sol{IP hors du réseau 192.168.80.0/24.} & \sol{Ne joint ni les postes ni le routeur.} & \sol{Ex. 192.168.80.11.} \\ \hline
    \hl 2 & \sol{PC-C} & \sol{Passerelle .80.100 : pas l'adresse du routeur.} & \sol{Joint son réseau, mais aucun autre (serveur).} & \sol{Passerelle 192.168.80.254.} \\ \hline
    \hl 3 & \sol{PC-D} & \sol{Même IP que PC-B (doublon).} & \sol{Conflit d'adresse : PC-B et PC-D perturbés.} & \sol{Adresse unique, ex. .80.14.} \\ \hline
    \hl 4 & \sol{R1 G0/0} & \sol{Interface côté LAN désactivée (Off).} & \sol{Aucun poste ne joint le routeur ni le serveur.} & \sol{Port Status : On.} \\ \hline}
  \def\Cbareme{PC-B et SRV-C sont correctement configurés ; accepter PC-B cité à la place de PC-D pour le doublon.}
\fi

\pagestyle{fancy}
\fancyhf{}
\renewcommand{\headrulewidth}{0pt}
\fancyfoot[L]{\footnotesize\itshape\color{grisfoot}Évaluation — Réseaux locaux — Sujet \variante\ifcorrige\ — CORRIGÉ\fi}
\fancyfoot[R]{\footnotesize\color{grisfoot}\thepage/3}

\setlist{itemsep=0pt,topsep=2pt,parsep=0pt}
\setlength{\parindent}{0pt}
\setlength{\parskip}{2pt}

\begin{document}

\begin{center}
  {\Large\bfseries Évaluation pratique — Réseaux locaux \quad
   \fcolorbox{black}{white}{\ SUJET \variante\ }\par}
  \vspace{1pt}
  {\bfseries Adressage IP, commutation et routage — BTS CIEL 1re année — durée : 2 heures\par}
  \ifcorrige\vspace{2pt}{\bfseries\color{corr}CORRIGÉ ET BARÈME — document enseignant}\par\fi
\end{center}
\vspace{-2pt}
{\renewcommand{\arraystretch}{1.35}
\begin{tabularx}{\linewidth}{|X|X|p{2.2cm}|p{2.6cm}|}
  \hline
  \textbf{Nom :} & \textbf{Prénom :} & \textbf{Poste n°} & \textbf{Note :} \hfill /20 \\ \hline
\end{tabularx}}

\textbf{Consignes :} travail \textbf{individuel}, documents et Internet \textbf{non autorisés}. Ce sujet est le
\textbf{document-réponse}. Fichier Packet Tracer : \ip{Nom\_Prenom\_EvalLAN\_\variante.pkt}, à enregistrer
régulièrement dans le dossier de rendu. Les « visas » sont apposés par le professeur, que vous appelez à chaque test réussi.
Durées conseillées : A 30 min — B 1 h 10 — C 20 min.

%==========================================================================
\section{Partie A — Questions de cours et calculs \hfill /6}

\question{A1}{2 pts}{Complétez le tableau en considérant le \textbf{masque par défaut} de la classe de chaque adresse.}
{\small\renewcommand{\arraystretch}{1.15}
\begin{tabularx}{\linewidth}{|l|c|C|C|C|c|}
  \hline
  \textbf{Adresse IP} & \textbf{Classe} & \textbf{Masque par défaut} & \textbf{Adresse réseau} &
  \textbf{Adresse de diffusion} & \textbf{Nb d'hôtes} \\ \hline
  \AIrows
\end{tabularx}}
\bareme{0,5 pt par ligne juste + 0,5 pt si les trois nombres d'hôtes sont justes ($2^{n}-2$, $n$ = nombre de bits d'hôte).}

\question{A2}{1 pt}{Indiquez le type de câble (\textbf{droit} ou \textbf{croisé}) pour chaque liaison.}
{\small\renewcommand{\arraystretch}{1.15}
\begin{tabularx}{\linewidth}{|X|c||X|c|}
  \hline
  \hl PC $\leftrightarrow$ commutateur & \makebox[1.6cm]{\sol{droit}} & Commutateur $\leftrightarrow$ routeur & \makebox[1.6cm]{\sol{droit}} \\ \hline
  \hl PC $\leftrightarrow$ PC (réseau P2P) & \sol{croisé} & Concentrateur $\leftrightarrow$ concentrateur & \sol{croisé} \\ \hline
\end{tabularx}}
\bareme{0,25 pt par case. Équipements de types différents : câble droit ; de même type : câble croisé.}

\question{A3}{1 pt}{Comment un \textbf{concentrateur (hub)} et un \textbf{commutateur (switch)} traitent-ils une trame
reçue ? Lequel limite les collisions, et pourquoi ?}
\zone{2.9cm}{Le \textbf{hub} (couche 1) répète la trame sur \textbf{tous ses autres ports} : un seul domaine de collision
partagé. Le \textbf{switch} (couche 2) lit l'adresse MAC de destination et, grâce à sa \textbf{table MAC}, ne l'envoie
\textbf{que sur le port du destinataire} : un domaine de collision par port, c'est lui qui limite les collisions.
\emph{0,5 pt hub, 0,5 pt switch.}}

\question{A4}{1 pt}{Qu'est-ce que la \textbf{passerelle par défaut} d'un poste ? Quelle condition son adresse doit-elle respecter ?}
\zone{2.4cm}{Adresse IP de l'interface du \textbf{routeur} à laquelle le poste envoie les paquets destinés à un
\textbf{autre réseau}. Elle doit \textbf{appartenir au même réseau} que le poste. \emph{0,5 + 0,5 pt.}}

\question{A5}{0,5 pt}{Un portable configuré en DHCP affiche \ip{\apipa} / \ip{255.255.0.0}. Que s'est-il passé ?}
\zone{1.5cm}{Aucun \textbf{serveur DHCP} n'a répondu : le poste s'est attribué une adresse \textbf{APIPA} (169.254.0.0/16).}

\question{A6}{0,5 pt}{Quel protocole utilise \texttt{ping} ? Nommez les deux messages échangés et la couche OSI des adresses IP.}
\zone{1.5cm}{\textbf{ICMP} ; \textbf{Echo Request} / \textbf{Echo Reply} ; adresses IP en \textbf{couche 3 (Réseau)}.}

%==========================================================================
\newpage
\section{Partie B — Réalisation sous Packet Tracer \hfill /10}

\begin{cdc}
\entreprise{} vous confie son réseau : \textbf{trois réseaux} reliés par un \textbf{routeur unique Cisco 2911}
nommé \ip{\Rnom} (interfaces G0/0, G0/1, G0/2).
\begin{itemize}[leftmargin=12pt]
  \item \textbf{Réseau \resUn} : \textbf{4 postes} PC-\pUn1 à PC-\pUn4, commutateur \ip{SW-\pUn}, adressage statique,
        réseau \ip{\netUn.0/24}.
  \item \textbf{Réseau \resDeux} : \textbf{2 postes} PC-\pDeux1, PC-\pDeux2 et \textbf{1 imprimante} IMP-\pDeux,
        commutateur \ip{SW-\pDeux}, adressage statique, réseau de classe B \ip{\netDeux.0.0/16}.
  \item \textbf{Réseau Serveur} : serveur \ip{\srvNom} d'adresse \ip{\netSrv.\srvH/24}, relié \textbf{directement} au routeur.
  \item La passerelle de chaque réseau est sa \textbf{dernière adresse utilisable}. Tous les équipements communiquent
        entre eux et accèdent à la page web du serveur. Nommage imposé (« Display Name »).
\end{itemize}
\end{cdc}

\question{B1}{2 pts}{Complétez le plan d'adressage \textbf{avant} la réalisation.}
{\footnotesize\setlength{\tabcolsep}{3pt}\renewcommand{\arraystretch}{1.15}
\begin{tabularx}{\linewidth}{|l|C|C|C|C|C|C|}
  \hline
  \textbf{Réseau} & \textbf{Adresse réseau} & \textbf{Masque} & \textbf{1re @ utilisable} &
  \textbf{Dernière @ utilisable} & \textbf{Diffusion} & \textbf{Passerelle} \\ \hline
  \hl \resUn & \sol{\netUn.0} & \sol{255.255.255.0} & \sol{\netUn.1} & \sol{\netUn.254} & \sol{\netUn.255} & \sol{\netUn.254} \\ \hline
  \hl \resDeux & \sol{\netDeux.0.0} & \sol{255.255.0.0} & \sol{\netDeux.0.1} & \sol{\netDeux.255.254} & \sol{\netDeux.255.255} & \sol{\netDeux.255.254} \\ \hline
  \hl Serveur & \sol{\netSrv.0} & \sol{255.255.255.0} & \sol{\netSrv.1} & \sol{\netSrv.254} & \sol{\netSrv.255} & \sol{\netSrv.254} \\ \hline
\end{tabularx}}
\bareme{0,5 pt par ligne juste + 0,5 pt si la convention de passerelle est respectée partout.}

\question{B2}{2 pts}{Construisez la topologie (équipements, câblage, nommage) et relevez l'adressage configuré.}
{\footnotesize\setlength{\tabcolsep}{3pt}\renewcommand{\arraystretch}{1.1}
\begin{tabularx}{\linewidth}{|l|C|C|C||l|C|C|}
  \hline
  \textbf{Hôte} & \textbf{Adresse IP} & \textbf{Masque} & \textbf{Passerelle} & \textbf{Interface} & \textbf{Adresse IP} & \textbf{Masque} \\ \hline
  \hl PC-\pUn1 & \sol{\netUn.1} & \sol{255.255.255.0} & \sol{\netUn.254} & \Rnom{} G0/0 & \sol{\netUn.254} & \sol{255.255.255.0} \\ \hline
  \hl PC-\pDeux1 & \sol{\netDeux.0.1} & \sol{255.255.0.0} & \sol{\netDeux.255.254} & \Rnom{} G0/1 & \sol{\netDeux.255.254} & \sol{255.255.0.0} \\ \hline
  \hl IMP-\pDeux & \sol{\netDeux.0.10} & \sol{255.255.0.0} & \sol{\netDeux.255.254} & \Rnom{} G0/2 & \sol{\netSrv.254} & \sol{255.255.255.0} \\ \hline
  \hl \srvNom & \ip{\netSrv.\srvH} & \sol{255.255.255.0} & \sol{\netSrv.254} & \multicolumn{3}{c|}{\cellcolor{black!8}} \\ \hline
\end{tabularx}}
\bareme{Adresses d'hôtes libres (uniques, dans le bon réseau) ; affectation G0/x libre. Sur le .pkt : équipements 0,5 ;
câblage 1 (câbles droits vers les switchs et switch–routeur, \textbf{croisé} serveur–routeur, liens verts) ; nommage 0,5.}

\question{B3}{3 pts}{Configurez tous les équipements : IP, masques, passerelles des hôtes ; interfaces du routeur
(« Port Status : On »).}
\bareme{Sur le .pkt : interfaces routeur 1 pt ; IP/masques des hôtes 1 pt ; passerelles (imprimante et serveur compris)
1 pt. $-0,25$ par erreur, sans descendre sous 0 par item.}

\question{B4}{3 pts}{Réalisez les tests, notez la commande et le résultat, puis faites viser chaque test réussi.}
{\small\renewcommand{\arraystretch}{1.1}
\begin{tabularx}{\linewidth}{|c|l|l|C|c|c|}
  \hline
  \textbf{N°} & \textbf{Source} & \textbf{Destination} & \textbf{Commande / action} & \textbf{Résultat} & \textbf{Visa prof.} \\ \hline
  \hl 1 & PC-\pUn1 & PC-\pUn4 & \sol{ping \netUn.4} & \sol{réussi} & \hspace{1.4cm} \\ \hline
  \hl 2 & PC-\pDeux1 & IMP-\pDeux & \sol{ping \netDeux.0.10} & \sol{réussi} & \\ \hline
  \hl 3 & PC-\pUn2 & PC-\pDeux2 & \sol{ping \netDeux.0.2} & \sol{réussi} & \\ \hline
  \hl 4 & PC-\pDeux2 & \srvNom & \sol{ping \netSrv.\srvH} & \sol{réussi} & \\ \hline
  \hl 5 & PC-\pUn3 & page web du serveur & \sol{navigateur : {\NoAutoSpacing http://\netSrv.\srvH}} & \sol{page affichée} & \\ \hline
\end{tabularx}}
\bareme{0,5 pt par test visé + 0,5 pt pour B5. Le premier ping inter-réseaux peut échouer (résolution ARP) ; les suivants doivent réussir.}

\question{B5}{incl. B4}{Test n°3 : par quels équipements passe le paquet ? Pourquoi PC-\pUn2 l'envoie-t-il au routeur ?}
\zone{2.3cm}{PC-\pUn2 $\to$ SW-\pUn{} $\to$ \Rnom{} $\to$ SW-\pDeux{} $\to$ PC-\pDeux2. Grâce à son masque, PC-\pUn2 constate que
la destination (\netDeux.0.0/16) n'est \textbf{pas dans son réseau} (\netUn.0/24) : il envoie le paquet à sa
\textbf{passerelle par défaut}, qui le route.}

%==========================================================================
\newpage
\section{Partie C — Diagnostic d'un réseau en panne \hfill /4}

Un technicien constate qu'\textbf{aucun poste ne parvient à joindre le serveur SRV-C} et que certains postes
communiquent mal entre eux. Les câbles sont du bon type.

\begin{center}
\begin{minipage}{0.37\linewidth}\centering
\begin{tikzpicture}[font=\footnotesize,
  eq/.style={draw,fill=qbg,rounded corners=2pt,minimum width=9mm,minimum height=5mm,inner sep=2pt},
  every path/.style={thick}]
  \node[eq] (SW) {SW-1};
  \node[eq,right=9mm of SW] (R) {R1};
  \node[eq,right=7mm of R] (SRV) {SRV-C};
  \node[eq,above left=5mm and -3mm of SW] (A) {PC-A};
  \node[eq,above right=5mm and -3mm of SW] (B) {PC-B};
  \node[eq,below left=5mm and -3mm of SW] (C) {PC-C};
  \node[eq,below right=5mm and -3mm of SW] (D) {PC-D};
  \foreach \n in {A,B,C,D} \draw (\n) -- (SW);
  \draw (SW) -- node[above,font=\tiny]{G0/0} (R);
  \draw (R) -- node[above,font=\tiny]{G0/1} (SRV);
\end{tikzpicture}
\end{minipage}\hfill
\begin{minipage}{0.61\linewidth}\centering
{\small\renewcommand{\arraystretch}{1.05}\setlength{\tabcolsep}{4pt}
\begin{tabular}{lllll}
  \toprule
  \textbf{Équipement} & \textbf{Adresse IP} & \textbf{Masque} & \textbf{Passerelle} & \textbf{Port} \\
  \midrule
  \Crows
  \bottomrule
\end{tabular}}
\end{minipage}
\end{center}

\question{C1}{4 pts}{Relevez les \textbf{quatre anomalies}. Pour chacune : conséquence et correction.}
{\small\renewcommand{\arraystretch}{1.1}
\ifcorrige\def\hc{}\else\def\hc{\rule[-10mm]{0pt}{14mm}}\fi
\begin{tabularx}{\linewidth}{|c|>{\raggedright\arraybackslash}p{2cm}|L|L|L|}
  \hline
  \textbf{N°} & \textbf{Équipement} & \textbf{Anomalie} & \textbf{Conséquence} & \textbf{Correction} \\ \hline
  \ifcorrige\Csol\else
  \hc 1 & & & & \\ \hline
  \hc 2 & & & & \\ \hline
  \hc 3 & & & & \\ \hline
  \hc 4 & & & & \\ \hline
  \fi
\end{tabularx}}
\bareme{1 pt par anomalie (0,25 équipement, anomalie, conséquence, correction ; ordre indifférent). \Cbareme}

%==========================================================================
\section{Bonus — Extension sans fil \hfill +1}

\begin{cdc}Ajoutez au réseau \resUn{} un \textbf{AccessPoint-PT} relié à SW-\pUn{} (SSID \ip{\ssid}, \textbf{WPA2-PSK})
et un portable \ip{Laptop-\pUn} équipé d'un module \texttt{WPC300N}, en adressage \textbf{statique}.
Vérifiez qu'il joint \srvNom.\end{cdc}
{\small\renewcommand{\arraystretch}{1.1}
\begin{tabularx}{\linewidth}{|l|C|C|C|c|}
  \hline
  \textbf{Équipement} & \textbf{Adresse IP} & \textbf{Masque} & \textbf{Passerelle} & \textbf{Visa prof.} \\ \hline
  \hl Laptop-\pUn & \sol{\netUn.50} & \sol{255.255.255.0} & \sol{\netUn.254} & \hspace{1.4cm} \\ \hline
\end{tabularx}}
\bareme{Le point d'accès fait office de pont : le portable est dans \netUn.0/24. Bonus si associé en WPA2-PSK et ping
vers \netSrv.\srvH{} réussi. Note plafonnée à 20.}

\end{document}
"
%   pdflatex -jobname=Eval_Reseau_CorrigeA "\def\variante{A}\def\modecorrige{1}% Évaluation pratique — Réseaux locaux (BTS CIEL 1re année) — sujets A et B
% Compilation (4 PDF) :
%   pdflatex -jobname=Eval_Reseau_SujetA   "\def\variante{A}\def\modecorrige{0}\input{Eval_TP_Reseau_LAN.tex}"
%   pdflatex -jobname=Eval_Reseau_SujetB   "\def\variante{B}\def\modecorrige{0}\input{Eval_TP_Reseau_LAN.tex}"
%   pdflatex -jobname=Eval_Reseau_CorrigeA "\def\variante{A}\def\modecorrige{1}\input{Eval_TP_Reseau_LAN.tex}"
%   pdflatex -jobname=Eval_Reseau_CorrigeB "\def\variante{B}\def\modecorrige{1}\input{Eval_TP_Reseau_LAN.tex}"
\documentclass[10pt,a4paper]{article}

\usepackage[T1]{fontenc}
\usepackage[utf8]{inputenc}
\usepackage{lmodern}
\usepackage[french]{babel}
\usepackage[a4paper,hmargin=1.5cm,top=1.3cm,bottom=1.6cm,footskip=0.7cm]{geometry}
\usepackage[table]{xcolor}
\usepackage{booktabs,tabularx,array}
\usepackage{enumitem}
\usepackage{fancyhdr}
\usepackage[most]{tcolorbox}
\usepackage{tikz}
\usetikzlibrary{positioning}
\usepackage{titlesec}

\providecommand{\modecorrige}{0}
\providecommand{\variante}{A}
\newif\ifcorrige
\ifnum\modecorrige=1 \corrigetrue \else \corrigefalse \fi

\definecolor{ipred}{HTML}{C0392B}
\definecolor{cdcbar}{HTML}{A67C3D}
\definecolor{cdcbg}{HTML}{F4F1EA}
\definecolor{qbar}{HTML}{4A6B5D}
\definecolor{qbg}{HTML}{EDF1EF}
\definecolor{corr}{HTML}{1F5FA8}
\definecolor{corrbg}{HTML}{E8F0FA}
\definecolor{grisfoot}{HTML}{828282}

\newcommand{\ip}[1]{\textcolor{ipred}{\texttt{#1}}}
\newcommand{\pts}[1]{\hfill\mbox{\textbf{(#1)}}}
\newcommand{\sol}[1]{\ifcorrige\textcolor{corr}{#1}\fi}
\newcommand{\bareme}[1]{\ifcorrige\par{\footnotesize\color{corr}#1\par}\fi}
\newcommand{\hl}{\ifcorrige\rule[-1.6mm]{0pt}{5.4mm}\else\rule[-2.2mm]{0pt}{6.6mm}\fi}
\newcolumntype{C}{>{\centering\arraybackslash}X}
\newcolumntype{L}{>{\raggedright\arraybackslash}X}

\titleformat{\section}{\large\bfseries}{}{0pt}{}
\titlespacing*{\section}{0pt}{6pt}{3pt}

\tcbset{qstyle/.style={enhanced,boxrule=0pt,leftrule=2.5pt,arc=0pt,outer arc=0pt,
  left=5pt,right=5pt,top=2pt,bottom=2pt,before skip=4pt,after skip=3pt}}
\newtcolorbox{cdc}{qstyle,colback=cdcbg,colframe=cdcbar}
\newtcolorbox{quest}{qstyle,colback=qbg,colframe=qbar}

% Zone de réponse : vierge dans le sujet, remplie dans le corrigé
\newcommand{\zone}[2]{%
  \ifcorrige
    \begin{tcolorbox}[colback=corrbg,colframe=corr,boxrule=0.5pt,arc=1pt,left=4pt,right=4pt,top=1pt,bottom=1pt,
      before skip=2pt,after skip=2pt,fontupper=\footnotesize\color{corr}]#2\end{tcolorbox}%
  \else
    \begin{tcolorbox}[colback=white,colframe=black!45,boxrule=0.4pt,arc=1pt,height=#1,
      before skip=2pt,after skip=2pt]\end{tcolorbox}%
  \fi}

\newcommand{\question}[3]{\begin{quest}\textbf{#1.} #3\pts{#2}\end{quest}}

%==========================================================================
% Données propres à chaque variante
%==========================================================================
\if\variante A
  \def\entreprise{Le cabinet d'architectes \textbf{ArchiNova}}
  \def\Rnom{R-ArchiNova}
  \def\resUn{Bureau d'études}\def\pUn{BE}\def\netUn{192.168.40}
  \def\resDeux{Accueil}\def\pDeux{ACC}\def\netDeux{172.16}
  \def\srvNom{SRV-Intranet}\def\netSrv{192.168.200}\def\srvH{10}
  \def\ssid{ArchiNova-WiFi}
  \def\AIrows{%
    \hl\ip{192.168.40.25} & \sol{C} & \sol{255.255.255.0} & \sol{192.168.40.0} & \sol{192.168.40.255} & \sol{254} \\ \hline
    \hl\ip{172.16.5.10}   & \sol{B} & \sol{255.255.0.0}   & \sol{172.16.0.0}   & \sol{172.16.255.255} & \sol{65\,534} \\ \hline
    \hl\ip{10.3.2.1}      & \sol{A} & \sol{255.0.0.0}     & \sol{10.0.0.0}     & \sol{10.255.255.255} & \sol{16\,777\,214} \\ \hline}
  \def\apipa{169.254.12.7}
  \def\Cnet{192.168.60}
  \def\Crows{%
    PC-A & \ip{192.168.60.10} & \ip{255.255.255.0} & \ip{192.168.60.254} & — \\
    PC-B & \ip{192.168.60.20} & \ip{255.255.255.0} & \ip{192.168.60.1}   & — \\
    PC-C & \ip{192.168.61.30} & \ip{255.255.255.0} & \ip{192.168.60.254} & — \\
    PC-D & \ip{192.168.60.10} & \ip{255.255.255.0} & \ip{192.168.60.254} & — \\
    SRV-C & \ip{10.0.0.5}     & \ip{255.0.0.0}     & \ip{10.0.0.254}     & — \\
    R1 — G0/0 & \ip{192.168.60.254} & \ip{255.255.255.0} & — & On \\
    R1 — G0/1 & \ip{10.0.0.254}     & \ip{255.0.0.0}     & — & Off \\}
  \def\Csol{%
    \hl 1 & \sol{PC-B} & \sol{Passerelle .60.1 : pas l'adresse du routeur.} & \sol{Joint son réseau, mais aucun autre (serveur).} & \sol{Passerelle 192.168.60.254.} \\ \hline
    \hl 2 & \sol{PC-C} & \sol{IP hors du réseau 192.168.60.0/24.} & \sol{Ne joint ni les postes ni le routeur.} & \sol{Ex. 192.168.60.30.} \\ \hline
    \hl 3 & \sol{PC-D} & \sol{Même IP que PC-A (doublon).} & \sol{Conflit d'adresse : PC-A et PC-D perturbés.} & \sol{Adresse unique, ex. .60.40.} \\ \hline
    \hl 4 & \sol{R1 G0/1} & \sol{Interface désactivée (Off).} & \sol{Réseau du serveur injoignable pour tous.} & \sol{Port Status : On.} \\ \hline}
  \def\Cbareme{PC-A et SRV-C sont correctement configurés ; accepter PC-A cité à la place de PC-D pour le doublon.}
\else
  \def\entreprise{La clinique vétérinaire \textbf{VetoPlus}}
  \def\Rnom{R-VetoPlus}
  \def\resUn{Soins}\def\pUn{SOI}\def\netUn{192.168.70}
  \def\resDeux{Secrétariat}\def\pDeux{SEC}\def\netDeux{172.20}
  \def\srvNom{SRV-Dossiers}\def\netSrv{192.168.250}\def\srvH{20}
  \def\ssid{VetoPlus-WiFi}
  \def\AIrows{%
    \hl\ip{192.168.15.130} & \sol{C} & \sol{255.255.255.0} & \sol{192.168.15.0} & \sol{192.168.15.255} & \sol{254} \\ \hline
    \hl\ip{172.31.200.7}   & \sol{B} & \sol{255.255.0.0}   & \sol{172.31.0.0}   & \sol{172.31.255.255} & \sol{65\,534} \\ \hline
    \hl\ip{10.45.6.9}      & \sol{A} & \sol{255.0.0.0}     & \sol{10.0.0.0}     & \sol{10.255.255.255} & \sol{16\,777\,214} \\ \hline}
  \def\apipa{169.254.88.3}
  \def\Cnet{192.168.80}
  \def\Crows{%
    PC-A & \ip{192.168.8.11}   & \ip{255.255.255.0} & \ip{192.168.80.254} & — \\
    PC-B & \ip{192.168.80.12}  & \ip{255.255.255.0} & \ip{192.168.80.254} & — \\
    PC-C & \ip{192.168.80.13}  & \ip{255.255.255.0} & \ip{192.168.80.100} & — \\
    PC-D & \ip{192.168.80.12}  & \ip{255.255.255.0} & \ip{192.168.80.254} & — \\
    SRV-C & \ip{10.0.0.8}      & \ip{255.0.0.0}     & \ip{10.0.0.254}     & — \\
    R1 — G0/0 & \ip{192.168.80.254} & \ip{255.255.255.0} & — & Off \\
    R1 — G0/1 & \ip{10.0.0.254}     & \ip{255.0.0.0}     & — & On \\}
  \def\Csol{%
    \hl 1 & \sol{PC-A} & \sol{IP hors du réseau 192.168.80.0/24.} & \sol{Ne joint ni les postes ni le routeur.} & \sol{Ex. 192.168.80.11.} \\ \hline
    \hl 2 & \sol{PC-C} & \sol{Passerelle .80.100 : pas l'adresse du routeur.} & \sol{Joint son réseau, mais aucun autre (serveur).} & \sol{Passerelle 192.168.80.254.} \\ \hline
    \hl 3 & \sol{PC-D} & \sol{Même IP que PC-B (doublon).} & \sol{Conflit d'adresse : PC-B et PC-D perturbés.} & \sol{Adresse unique, ex. .80.14.} \\ \hline
    \hl 4 & \sol{R1 G0/0} & \sol{Interface côté LAN désactivée (Off).} & \sol{Aucun poste ne joint le routeur ni le serveur.} & \sol{Port Status : On.} \\ \hline}
  \def\Cbareme{PC-B et SRV-C sont correctement configurés ; accepter PC-B cité à la place de PC-D pour le doublon.}
\fi

\pagestyle{fancy}
\fancyhf{}
\renewcommand{\headrulewidth}{0pt}
\fancyfoot[L]{\footnotesize\itshape\color{grisfoot}Évaluation — Réseaux locaux — Sujet \variante\ifcorrige\ — CORRIGÉ\fi}
\fancyfoot[R]{\footnotesize\color{grisfoot}\thepage/3}

\setlist{itemsep=0pt,topsep=2pt,parsep=0pt}
\setlength{\parindent}{0pt}
\setlength{\parskip}{2pt}

\begin{document}

\begin{center}
  {\Large\bfseries Évaluation pratique — Réseaux locaux \quad
   \fcolorbox{black}{white}{\ SUJET \variante\ }\par}
  \vspace{1pt}
  {\bfseries Adressage IP, commutation et routage — BTS CIEL 1re année — durée : 2 heures\par}
  \ifcorrige\vspace{2pt}{\bfseries\color{corr}CORRIGÉ ET BARÈME — document enseignant}\par\fi
\end{center}
\vspace{-2pt}
{\renewcommand{\arraystretch}{1.35}
\begin{tabularx}{\linewidth}{|X|X|p{2.2cm}|p{2.6cm}|}
  \hline
  \textbf{Nom :} & \textbf{Prénom :} & \textbf{Poste n°} & \textbf{Note :} \hfill /20 \\ \hline
\end{tabularx}}

\textbf{Consignes :} travail \textbf{individuel}, documents et Internet \textbf{non autorisés}. Ce sujet est le
\textbf{document-réponse}. Fichier Packet Tracer : \ip{Nom\_Prenom\_EvalLAN\_\variante.pkt}, à enregistrer
régulièrement dans le dossier de rendu. Les « visas » sont apposés par le professeur, que vous appelez à chaque test réussi.
Durées conseillées : A 30 min — B 1 h 10 — C 20 min.

%==========================================================================
\section{Partie A — Questions de cours et calculs \hfill /6}

\question{A1}{2 pts}{Complétez le tableau en considérant le \textbf{masque par défaut} de la classe de chaque adresse.}
{\small\renewcommand{\arraystretch}{1.15}
\begin{tabularx}{\linewidth}{|l|c|C|C|C|c|}
  \hline
  \textbf{Adresse IP} & \textbf{Classe} & \textbf{Masque par défaut} & \textbf{Adresse réseau} &
  \textbf{Adresse de diffusion} & \textbf{Nb d'hôtes} \\ \hline
  \AIrows
\end{tabularx}}
\bareme{0,5 pt par ligne juste + 0,5 pt si les trois nombres d'hôtes sont justes ($2^{n}-2$, $n$ = nombre de bits d'hôte).}

\question{A2}{1 pt}{Indiquez le type de câble (\textbf{droit} ou \textbf{croisé}) pour chaque liaison.}
{\small\renewcommand{\arraystretch}{1.15}
\begin{tabularx}{\linewidth}{|X|c||X|c|}
  \hline
  \hl PC $\leftrightarrow$ commutateur & \makebox[1.6cm]{\sol{droit}} & Commutateur $\leftrightarrow$ routeur & \makebox[1.6cm]{\sol{droit}} \\ \hline
  \hl PC $\leftrightarrow$ PC (réseau P2P) & \sol{croisé} & Concentrateur $\leftrightarrow$ concentrateur & \sol{croisé} \\ \hline
\end{tabularx}}
\bareme{0,25 pt par case. Équipements de types différents : câble droit ; de même type : câble croisé.}

\question{A3}{1 pt}{Comment un \textbf{concentrateur (hub)} et un \textbf{commutateur (switch)} traitent-ils une trame
reçue ? Lequel limite les collisions, et pourquoi ?}
\zone{2.9cm}{Le \textbf{hub} (couche 1) répète la trame sur \textbf{tous ses autres ports} : un seul domaine de collision
partagé. Le \textbf{switch} (couche 2) lit l'adresse MAC de destination et, grâce à sa \textbf{table MAC}, ne l'envoie
\textbf{que sur le port du destinataire} : un domaine de collision par port, c'est lui qui limite les collisions.
\emph{0,5 pt hub, 0,5 pt switch.}}

\question{A4}{1 pt}{Qu'est-ce que la \textbf{passerelle par défaut} d'un poste ? Quelle condition son adresse doit-elle respecter ?}
\zone{2.4cm}{Adresse IP de l'interface du \textbf{routeur} à laquelle le poste envoie les paquets destinés à un
\textbf{autre réseau}. Elle doit \textbf{appartenir au même réseau} que le poste. \emph{0,5 + 0,5 pt.}}

\question{A5}{0,5 pt}{Un portable configuré en DHCP affiche \ip{\apipa} / \ip{255.255.0.0}. Que s'est-il passé ?}
\zone{1.5cm}{Aucun \textbf{serveur DHCP} n'a répondu : le poste s'est attribué une adresse \textbf{APIPA} (169.254.0.0/16).}

\question{A6}{0,5 pt}{Quel protocole utilise \texttt{ping} ? Nommez les deux messages échangés et la couche OSI des adresses IP.}
\zone{1.5cm}{\textbf{ICMP} ; \textbf{Echo Request} / \textbf{Echo Reply} ; adresses IP en \textbf{couche 3 (Réseau)}.}

%==========================================================================
\newpage
\section{Partie B — Réalisation sous Packet Tracer \hfill /10}

\begin{cdc}
\entreprise{} vous confie son réseau : \textbf{trois réseaux} reliés par un \textbf{routeur unique Cisco 2911}
nommé \ip{\Rnom} (interfaces G0/0, G0/1, G0/2).
\begin{itemize}[leftmargin=12pt]
  \item \textbf{Réseau \resUn} : \textbf{4 postes} PC-\pUn1 à PC-\pUn4, commutateur \ip{SW-\pUn}, adressage statique,
        réseau \ip{\netUn.0/24}.
  \item \textbf{Réseau \resDeux} : \textbf{2 postes} PC-\pDeux1, PC-\pDeux2 et \textbf{1 imprimante} IMP-\pDeux,
        commutateur \ip{SW-\pDeux}, adressage statique, réseau de classe B \ip{\netDeux.0.0/16}.
  \item \textbf{Réseau Serveur} : serveur \ip{\srvNom} d'adresse \ip{\netSrv.\srvH/24}, relié \textbf{directement} au routeur.
  \item La passerelle de chaque réseau est sa \textbf{dernière adresse utilisable}. Tous les équipements communiquent
        entre eux et accèdent à la page web du serveur. Nommage imposé (« Display Name »).
\end{itemize}
\end{cdc}

\question{B1}{2 pts}{Complétez le plan d'adressage \textbf{avant} la réalisation.}
{\footnotesize\setlength{\tabcolsep}{3pt}\renewcommand{\arraystretch}{1.15}
\begin{tabularx}{\linewidth}{|l|C|C|C|C|C|C|}
  \hline
  \textbf{Réseau} & \textbf{Adresse réseau} & \textbf{Masque} & \textbf{1re @ utilisable} &
  \textbf{Dernière @ utilisable} & \textbf{Diffusion} & \textbf{Passerelle} \\ \hline
  \hl \resUn & \sol{\netUn.0} & \sol{255.255.255.0} & \sol{\netUn.1} & \sol{\netUn.254} & \sol{\netUn.255} & \sol{\netUn.254} \\ \hline
  \hl \resDeux & \sol{\netDeux.0.0} & \sol{255.255.0.0} & \sol{\netDeux.0.1} & \sol{\netDeux.255.254} & \sol{\netDeux.255.255} & \sol{\netDeux.255.254} \\ \hline
  \hl Serveur & \sol{\netSrv.0} & \sol{255.255.255.0} & \sol{\netSrv.1} & \sol{\netSrv.254} & \sol{\netSrv.255} & \sol{\netSrv.254} \\ \hline
\end{tabularx}}
\bareme{0,5 pt par ligne juste + 0,5 pt si la convention de passerelle est respectée partout.}

\question{B2}{2 pts}{Construisez la topologie (équipements, câblage, nommage) et relevez l'adressage configuré.}
{\footnotesize\setlength{\tabcolsep}{3pt}\renewcommand{\arraystretch}{1.1}
\begin{tabularx}{\linewidth}{|l|C|C|C||l|C|C|}
  \hline
  \textbf{Hôte} & \textbf{Adresse IP} & \textbf{Masque} & \textbf{Passerelle} & \textbf{Interface} & \textbf{Adresse IP} & \textbf{Masque} \\ \hline
  \hl PC-\pUn1 & \sol{\netUn.1} & \sol{255.255.255.0} & \sol{\netUn.254} & \Rnom{} G0/0 & \sol{\netUn.254} & \sol{255.255.255.0} \\ \hline
  \hl PC-\pDeux1 & \sol{\netDeux.0.1} & \sol{255.255.0.0} & \sol{\netDeux.255.254} & \Rnom{} G0/1 & \sol{\netDeux.255.254} & \sol{255.255.0.0} \\ \hline
  \hl IMP-\pDeux & \sol{\netDeux.0.10} & \sol{255.255.0.0} & \sol{\netDeux.255.254} & \Rnom{} G0/2 & \sol{\netSrv.254} & \sol{255.255.255.0} \\ \hline
  \hl \srvNom & \ip{\netSrv.\srvH} & \sol{255.255.255.0} & \sol{\netSrv.254} & \multicolumn{3}{c|}{\cellcolor{black!8}} \\ \hline
\end{tabularx}}
\bareme{Adresses d'hôtes libres (uniques, dans le bon réseau) ; affectation G0/x libre. Sur le .pkt : équipements 0,5 ;
câblage 1 (câbles droits vers les switchs et switch–routeur, \textbf{croisé} serveur–routeur, liens verts) ; nommage 0,5.}

\question{B3}{3 pts}{Configurez tous les équipements : IP, masques, passerelles des hôtes ; interfaces du routeur
(« Port Status : On »).}
\bareme{Sur le .pkt : interfaces routeur 1 pt ; IP/masques des hôtes 1 pt ; passerelles (imprimante et serveur compris)
1 pt. $-0,25$ par erreur, sans descendre sous 0 par item.}

\question{B4}{3 pts}{Réalisez les tests, notez la commande et le résultat, puis faites viser chaque test réussi.}
{\small\renewcommand{\arraystretch}{1.1}
\begin{tabularx}{\linewidth}{|c|l|l|C|c|c|}
  \hline
  \textbf{N°} & \textbf{Source} & \textbf{Destination} & \textbf{Commande / action} & \textbf{Résultat} & \textbf{Visa prof.} \\ \hline
  \hl 1 & PC-\pUn1 & PC-\pUn4 & \sol{ping \netUn.4} & \sol{réussi} & \hspace{1.4cm} \\ \hline
  \hl 2 & PC-\pDeux1 & IMP-\pDeux & \sol{ping \netDeux.0.10} & \sol{réussi} & \\ \hline
  \hl 3 & PC-\pUn2 & PC-\pDeux2 & \sol{ping \netDeux.0.2} & \sol{réussi} & \\ \hline
  \hl 4 & PC-\pDeux2 & \srvNom & \sol{ping \netSrv.\srvH} & \sol{réussi} & \\ \hline
  \hl 5 & PC-\pUn3 & page web du serveur & \sol{navigateur : {\NoAutoSpacing http://\netSrv.\srvH}} & \sol{page affichée} & \\ \hline
\end{tabularx}}
\bareme{0,5 pt par test visé + 0,5 pt pour B5. Le premier ping inter-réseaux peut échouer (résolution ARP) ; les suivants doivent réussir.}

\question{B5}{incl. B4}{Test n°3 : par quels équipements passe le paquet ? Pourquoi PC-\pUn2 l'envoie-t-il au routeur ?}
\zone{2.3cm}{PC-\pUn2 $\to$ SW-\pUn{} $\to$ \Rnom{} $\to$ SW-\pDeux{} $\to$ PC-\pDeux2. Grâce à son masque, PC-\pUn2 constate que
la destination (\netDeux.0.0/16) n'est \textbf{pas dans son réseau} (\netUn.0/24) : il envoie le paquet à sa
\textbf{passerelle par défaut}, qui le route.}

%==========================================================================
\newpage
\section{Partie C — Diagnostic d'un réseau en panne \hfill /4}

Un technicien constate qu'\textbf{aucun poste ne parvient à joindre le serveur SRV-C} et que certains postes
communiquent mal entre eux. Les câbles sont du bon type.

\begin{center}
\begin{minipage}{0.37\linewidth}\centering
\begin{tikzpicture}[font=\footnotesize,
  eq/.style={draw,fill=qbg,rounded corners=2pt,minimum width=9mm,minimum height=5mm,inner sep=2pt},
  every path/.style={thick}]
  \node[eq] (SW) {SW-1};
  \node[eq,right=9mm of SW] (R) {R1};
  \node[eq,right=7mm of R] (SRV) {SRV-C};
  \node[eq,above left=5mm and -3mm of SW] (A) {PC-A};
  \node[eq,above right=5mm and -3mm of SW] (B) {PC-B};
  \node[eq,below left=5mm and -3mm of SW] (C) {PC-C};
  \node[eq,below right=5mm and -3mm of SW] (D) {PC-D};
  \foreach \n in {A,B,C,D} \draw (\n) -- (SW);
  \draw (SW) -- node[above,font=\tiny]{G0/0} (R);
  \draw (R) -- node[above,font=\tiny]{G0/1} (SRV);
\end{tikzpicture}
\end{minipage}\hfill
\begin{minipage}{0.61\linewidth}\centering
{\small\renewcommand{\arraystretch}{1.05}\setlength{\tabcolsep}{4pt}
\begin{tabular}{lllll}
  \toprule
  \textbf{Équipement} & \textbf{Adresse IP} & \textbf{Masque} & \textbf{Passerelle} & \textbf{Port} \\
  \midrule
  \Crows
  \bottomrule
\end{tabular}}
\end{minipage}
\end{center}

\question{C1}{4 pts}{Relevez les \textbf{quatre anomalies}. Pour chacune : conséquence et correction.}
{\small\renewcommand{\arraystretch}{1.1}
\ifcorrige\def\hc{}\else\def\hc{\rule[-10mm]{0pt}{14mm}}\fi
\begin{tabularx}{\linewidth}{|c|>{\raggedright\arraybackslash}p{2cm}|L|L|L|}
  \hline
  \textbf{N°} & \textbf{Équipement} & \textbf{Anomalie} & \textbf{Conséquence} & \textbf{Correction} \\ \hline
  \ifcorrige\Csol\else
  \hc 1 & & & & \\ \hline
  \hc 2 & & & & \\ \hline
  \hc 3 & & & & \\ \hline
  \hc 4 & & & & \\ \hline
  \fi
\end{tabularx}}
\bareme{1 pt par anomalie (0,25 équipement, anomalie, conséquence, correction ; ordre indifférent). \Cbareme}

%==========================================================================
\section{Bonus — Extension sans fil \hfill +1}

\begin{cdc}Ajoutez au réseau \resUn{} un \textbf{AccessPoint-PT} relié à SW-\pUn{} (SSID \ip{\ssid}, \textbf{WPA2-PSK})
et un portable \ip{Laptop-\pUn} équipé d'un module \texttt{WPC300N}, en adressage \textbf{statique}.
Vérifiez qu'il joint \srvNom.\end{cdc}
{\small\renewcommand{\arraystretch}{1.1}
\begin{tabularx}{\linewidth}{|l|C|C|C|c|}
  \hline
  \textbf{Équipement} & \textbf{Adresse IP} & \textbf{Masque} & \textbf{Passerelle} & \textbf{Visa prof.} \\ \hline
  \hl Laptop-\pUn & \sol{\netUn.50} & \sol{255.255.255.0} & \sol{\netUn.254} & \hspace{1.4cm} \\ \hline
\end{tabularx}}
\bareme{Le point d'accès fait office de pont : le portable est dans \netUn.0/24. Bonus si associé en WPA2-PSK et ping
vers \netSrv.\srvH{} réussi. Note plafonnée à 20.}

\end{document}
"
%   pdflatex -jobname=Eval_Reseau_CorrigeB "\def\variante{B}\def\modecorrige{1}% Évaluation pratique — Réseaux locaux (BTS CIEL 1re année) — sujets A et B
% Compilation (4 PDF) :
%   pdflatex -jobname=Eval_Reseau_SujetA   "\def\variante{A}\def\modecorrige{0}\input{Eval_TP_Reseau_LAN.tex}"
%   pdflatex -jobname=Eval_Reseau_SujetB   "\def\variante{B}\def\modecorrige{0}\input{Eval_TP_Reseau_LAN.tex}"
%   pdflatex -jobname=Eval_Reseau_CorrigeA "\def\variante{A}\def\modecorrige{1}\input{Eval_TP_Reseau_LAN.tex}"
%   pdflatex -jobname=Eval_Reseau_CorrigeB "\def\variante{B}\def\modecorrige{1}\input{Eval_TP_Reseau_LAN.tex}"
\documentclass[10pt,a4paper]{article}

\usepackage[T1]{fontenc}
\usepackage[utf8]{inputenc}
\usepackage{lmodern}
\usepackage[french]{babel}
\usepackage[a4paper,hmargin=1.5cm,top=1.3cm,bottom=1.6cm,footskip=0.7cm]{geometry}
\usepackage[table]{xcolor}
\usepackage{booktabs,tabularx,array}
\usepackage{enumitem}
\usepackage{fancyhdr}
\usepackage[most]{tcolorbox}
\usepackage{tikz}
\usetikzlibrary{positioning}
\usepackage{titlesec}

\providecommand{\modecorrige}{0}
\providecommand{\variante}{A}
\newif\ifcorrige
\ifnum\modecorrige=1 \corrigetrue \else \corrigefalse \fi

\definecolor{ipred}{HTML}{C0392B}
\definecolor{cdcbar}{HTML}{A67C3D}
\definecolor{cdcbg}{HTML}{F4F1EA}
\definecolor{qbar}{HTML}{4A6B5D}
\definecolor{qbg}{HTML}{EDF1EF}
\definecolor{corr}{HTML}{1F5FA8}
\definecolor{corrbg}{HTML}{E8F0FA}
\definecolor{grisfoot}{HTML}{828282}

\newcommand{\ip}[1]{\textcolor{ipred}{\texttt{#1}}}
\newcommand{\pts}[1]{\hfill\mbox{\textbf{(#1)}}}
\newcommand{\sol}[1]{\ifcorrige\textcolor{corr}{#1}\fi}
\newcommand{\bareme}[1]{\ifcorrige\par{\footnotesize\color{corr}#1\par}\fi}
\newcommand{\hl}{\ifcorrige\rule[-1.6mm]{0pt}{5.4mm}\else\rule[-2.2mm]{0pt}{6.6mm}\fi}
\newcolumntype{C}{>{\centering\arraybackslash}X}
\newcolumntype{L}{>{\raggedright\arraybackslash}X}

\titleformat{\section}{\large\bfseries}{}{0pt}{}
\titlespacing*{\section}{0pt}{6pt}{3pt}

\tcbset{qstyle/.style={enhanced,boxrule=0pt,leftrule=2.5pt,arc=0pt,outer arc=0pt,
  left=5pt,right=5pt,top=2pt,bottom=2pt,before skip=4pt,after skip=3pt}}
\newtcolorbox{cdc}{qstyle,colback=cdcbg,colframe=cdcbar}
\newtcolorbox{quest}{qstyle,colback=qbg,colframe=qbar}

% Zone de réponse : vierge dans le sujet, remplie dans le corrigé
\newcommand{\zone}[2]{%
  \ifcorrige
    \begin{tcolorbox}[colback=corrbg,colframe=corr,boxrule=0.5pt,arc=1pt,left=4pt,right=4pt,top=1pt,bottom=1pt,
      before skip=2pt,after skip=2pt,fontupper=\footnotesize\color{corr}]#2\end{tcolorbox}%
  \else
    \begin{tcolorbox}[colback=white,colframe=black!45,boxrule=0.4pt,arc=1pt,height=#1,
      before skip=2pt,after skip=2pt]\end{tcolorbox}%
  \fi}

\newcommand{\question}[3]{\begin{quest}\textbf{#1.} #3\pts{#2}\end{quest}}

%==========================================================================
% Données propres à chaque variante
%==========================================================================
\if\variante A
  \def\entreprise{Le cabinet d'architectes \textbf{ArchiNova}}
  \def\Rnom{R-ArchiNova}
  \def\resUn{Bureau d'études}\def\pUn{BE}\def\netUn{192.168.40}
  \def\resDeux{Accueil}\def\pDeux{ACC}\def\netDeux{172.16}
  \def\srvNom{SRV-Intranet}\def\netSrv{192.168.200}\def\srvH{10}
  \def\ssid{ArchiNova-WiFi}
  \def\AIrows{%
    \hl\ip{192.168.40.25} & \sol{C} & \sol{255.255.255.0} & \sol{192.168.40.0} & \sol{192.168.40.255} & \sol{254} \\ \hline
    \hl\ip{172.16.5.10}   & \sol{B} & \sol{255.255.0.0}   & \sol{172.16.0.0}   & \sol{172.16.255.255} & \sol{65\,534} \\ \hline
    \hl\ip{10.3.2.1}      & \sol{A} & \sol{255.0.0.0}     & \sol{10.0.0.0}     & \sol{10.255.255.255} & \sol{16\,777\,214} \\ \hline}
  \def\apipa{169.254.12.7}
  \def\Cnet{192.168.60}
  \def\Crows{%
    PC-A & \ip{192.168.60.10} & \ip{255.255.255.0} & \ip{192.168.60.254} & — \\
    PC-B & \ip{192.168.60.20} & \ip{255.255.255.0} & \ip{192.168.60.1}   & — \\
    PC-C & \ip{192.168.61.30} & \ip{255.255.255.0} & \ip{192.168.60.254} & — \\
    PC-D & \ip{192.168.60.10} & \ip{255.255.255.0} & \ip{192.168.60.254} & — \\
    SRV-C & \ip{10.0.0.5}     & \ip{255.0.0.0}     & \ip{10.0.0.254}     & — \\
    R1 — G0/0 & \ip{192.168.60.254} & \ip{255.255.255.0} & — & On \\
    R1 — G0/1 & \ip{10.0.0.254}     & \ip{255.0.0.0}     & — & Off \\}
  \def\Csol{%
    \hl 1 & \sol{PC-B} & \sol{Passerelle .60.1 : pas l'adresse du routeur.} & \sol{Joint son réseau, mais aucun autre (serveur).} & \sol{Passerelle 192.168.60.254.} \\ \hline
    \hl 2 & \sol{PC-C} & \sol{IP hors du réseau 192.168.60.0/24.} & \sol{Ne joint ni les postes ni le routeur.} & \sol{Ex. 192.168.60.30.} \\ \hline
    \hl 3 & \sol{PC-D} & \sol{Même IP que PC-A (doublon).} & \sol{Conflit d'adresse : PC-A et PC-D perturbés.} & \sol{Adresse unique, ex. .60.40.} \\ \hline
    \hl 4 & \sol{R1 G0/1} & \sol{Interface désactivée (Off).} & \sol{Réseau du serveur injoignable pour tous.} & \sol{Port Status : On.} \\ \hline}
  \def\Cbareme{PC-A et SRV-C sont correctement configurés ; accepter PC-A cité à la place de PC-D pour le doublon.}
\else
  \def\entreprise{La clinique vétérinaire \textbf{VetoPlus}}
  \def\Rnom{R-VetoPlus}
  \def\resUn{Soins}\def\pUn{SOI}\def\netUn{192.168.70}
  \def\resDeux{Secrétariat}\def\pDeux{SEC}\def\netDeux{172.20}
  \def\srvNom{SRV-Dossiers}\def\netSrv{192.168.250}\def\srvH{20}
  \def\ssid{VetoPlus-WiFi}
  \def\AIrows{%
    \hl\ip{192.168.15.130} & \sol{C} & \sol{255.255.255.0} & \sol{192.168.15.0} & \sol{192.168.15.255} & \sol{254} \\ \hline
    \hl\ip{172.31.200.7}   & \sol{B} & \sol{255.255.0.0}   & \sol{172.31.0.0}   & \sol{172.31.255.255} & \sol{65\,534} \\ \hline
    \hl\ip{10.45.6.9}      & \sol{A} & \sol{255.0.0.0}     & \sol{10.0.0.0}     & \sol{10.255.255.255} & \sol{16\,777\,214} \\ \hline}
  \def\apipa{169.254.88.3}
  \def\Cnet{192.168.80}
  \def\Crows{%
    PC-A & \ip{192.168.8.11}   & \ip{255.255.255.0} & \ip{192.168.80.254} & — \\
    PC-B & \ip{192.168.80.12}  & \ip{255.255.255.0} & \ip{192.168.80.254} & — \\
    PC-C & \ip{192.168.80.13}  & \ip{255.255.255.0} & \ip{192.168.80.100} & — \\
    PC-D & \ip{192.168.80.12}  & \ip{255.255.255.0} & \ip{192.168.80.254} & — \\
    SRV-C & \ip{10.0.0.8}      & \ip{255.0.0.0}     & \ip{10.0.0.254}     & — \\
    R1 — G0/0 & \ip{192.168.80.254} & \ip{255.255.255.0} & — & Off \\
    R1 — G0/1 & \ip{10.0.0.254}     & \ip{255.0.0.0}     & — & On \\}
  \def\Csol{%
    \hl 1 & \sol{PC-A} & \sol{IP hors du réseau 192.168.80.0/24.} & \sol{Ne joint ni les postes ni le routeur.} & \sol{Ex. 192.168.80.11.} \\ \hline
    \hl 2 & \sol{PC-C} & \sol{Passerelle .80.100 : pas l'adresse du routeur.} & \sol{Joint son réseau, mais aucun autre (serveur).} & \sol{Passerelle 192.168.80.254.} \\ \hline
    \hl 3 & \sol{PC-D} & \sol{Même IP que PC-B (doublon).} & \sol{Conflit d'adresse : PC-B et PC-D perturbés.} & \sol{Adresse unique, ex. .80.14.} \\ \hline
    \hl 4 & \sol{R1 G0/0} & \sol{Interface côté LAN désactivée (Off).} & \sol{Aucun poste ne joint le routeur ni le serveur.} & \sol{Port Status : On.} \\ \hline}
  \def\Cbareme{PC-B et SRV-C sont correctement configurés ; accepter PC-B cité à la place de PC-D pour le doublon.}
\fi

\pagestyle{fancy}
\fancyhf{}
\renewcommand{\headrulewidth}{0pt}
\fancyfoot[L]{\footnotesize\itshape\color{grisfoot}Évaluation — Réseaux locaux — Sujet \variante\ifcorrige\ — CORRIGÉ\fi}
\fancyfoot[R]{\footnotesize\color{grisfoot}\thepage/3}

\setlist{itemsep=0pt,topsep=2pt,parsep=0pt}
\setlength{\parindent}{0pt}
\setlength{\parskip}{2pt}

\begin{document}

\begin{center}
  {\Large\bfseries Évaluation pratique — Réseaux locaux \quad
   \fcolorbox{black}{white}{\ SUJET \variante\ }\par}
  \vspace{1pt}
  {\bfseries Adressage IP, commutation et routage — BTS CIEL 1re année — durée : 2 heures\par}
  \ifcorrige\vspace{2pt}{\bfseries\color{corr}CORRIGÉ ET BARÈME — document enseignant}\par\fi
\end{center}
\vspace{-2pt}
{\renewcommand{\arraystretch}{1.35}
\begin{tabularx}{\linewidth}{|X|X|p{2.2cm}|p{2.6cm}|}
  \hline
  \textbf{Nom :} & \textbf{Prénom :} & \textbf{Poste n°} & \textbf{Note :} \hfill /20 \\ \hline
\end{tabularx}}

\textbf{Consignes :} travail \textbf{individuel}, documents et Internet \textbf{non autorisés}. Ce sujet est le
\textbf{document-réponse}. Fichier Packet Tracer : \ip{Nom\_Prenom\_EvalLAN\_\variante.pkt}, à enregistrer
régulièrement dans le dossier de rendu. Les « visas » sont apposés par le professeur, que vous appelez à chaque test réussi.
Durées conseillées : A 30 min — B 1 h 10 — C 20 min.

%==========================================================================
\section{Partie A — Questions de cours et calculs \hfill /6}

\question{A1}{2 pts}{Complétez le tableau en considérant le \textbf{masque par défaut} de la classe de chaque adresse.}
{\small\renewcommand{\arraystretch}{1.15}
\begin{tabularx}{\linewidth}{|l|c|C|C|C|c|}
  \hline
  \textbf{Adresse IP} & \textbf{Classe} & \textbf{Masque par défaut} & \textbf{Adresse réseau} &
  \textbf{Adresse de diffusion} & \textbf{Nb d'hôtes} \\ \hline
  \AIrows
\end{tabularx}}
\bareme{0,5 pt par ligne juste + 0,5 pt si les trois nombres d'hôtes sont justes ($2^{n}-2$, $n$ = nombre de bits d'hôte).}

\question{A2}{1 pt}{Indiquez le type de câble (\textbf{droit} ou \textbf{croisé}) pour chaque liaison.}
{\small\renewcommand{\arraystretch}{1.15}
\begin{tabularx}{\linewidth}{|X|c||X|c|}
  \hline
  \hl PC $\leftrightarrow$ commutateur & \makebox[1.6cm]{\sol{droit}} & Commutateur $\leftrightarrow$ routeur & \makebox[1.6cm]{\sol{droit}} \\ \hline
  \hl PC $\leftrightarrow$ PC (réseau P2P) & \sol{croisé} & Concentrateur $\leftrightarrow$ concentrateur & \sol{croisé} \\ \hline
\end{tabularx}}
\bareme{0,25 pt par case. Équipements de types différents : câble droit ; de même type : câble croisé.}

\question{A3}{1 pt}{Comment un \textbf{concentrateur (hub)} et un \textbf{commutateur (switch)} traitent-ils une trame
reçue ? Lequel limite les collisions, et pourquoi ?}
\zone{2.9cm}{Le \textbf{hub} (couche 1) répète la trame sur \textbf{tous ses autres ports} : un seul domaine de collision
partagé. Le \textbf{switch} (couche 2) lit l'adresse MAC de destination et, grâce à sa \textbf{table MAC}, ne l'envoie
\textbf{que sur le port du destinataire} : un domaine de collision par port, c'est lui qui limite les collisions.
\emph{0,5 pt hub, 0,5 pt switch.}}

\question{A4}{1 pt}{Qu'est-ce que la \textbf{passerelle par défaut} d'un poste ? Quelle condition son adresse doit-elle respecter ?}
\zone{2.4cm}{Adresse IP de l'interface du \textbf{routeur} à laquelle le poste envoie les paquets destinés à un
\textbf{autre réseau}. Elle doit \textbf{appartenir au même réseau} que le poste. \emph{0,5 + 0,5 pt.}}

\question{A5}{0,5 pt}{Un portable configuré en DHCP affiche \ip{\apipa} / \ip{255.255.0.0}. Que s'est-il passé ?}
\zone{1.5cm}{Aucun \textbf{serveur DHCP} n'a répondu : le poste s'est attribué une adresse \textbf{APIPA} (169.254.0.0/16).}

\question{A6}{0,5 pt}{Quel protocole utilise \texttt{ping} ? Nommez les deux messages échangés et la couche OSI des adresses IP.}
\zone{1.5cm}{\textbf{ICMP} ; \textbf{Echo Request} / \textbf{Echo Reply} ; adresses IP en \textbf{couche 3 (Réseau)}.}

%==========================================================================
\newpage
\section{Partie B — Réalisation sous Packet Tracer \hfill /10}

\begin{cdc}
\entreprise{} vous confie son réseau : \textbf{trois réseaux} reliés par un \textbf{routeur unique Cisco 2911}
nommé \ip{\Rnom} (interfaces G0/0, G0/1, G0/2).
\begin{itemize}[leftmargin=12pt]
  \item \textbf{Réseau \resUn} : \textbf{4 postes} PC-\pUn1 à PC-\pUn4, commutateur \ip{SW-\pUn}, adressage statique,
        réseau \ip{\netUn.0/24}.
  \item \textbf{Réseau \resDeux} : \textbf{2 postes} PC-\pDeux1, PC-\pDeux2 et \textbf{1 imprimante} IMP-\pDeux,
        commutateur \ip{SW-\pDeux}, adressage statique, réseau de classe B \ip{\netDeux.0.0/16}.
  \item \textbf{Réseau Serveur} : serveur \ip{\srvNom} d'adresse \ip{\netSrv.\srvH/24}, relié \textbf{directement} au routeur.
  \item La passerelle de chaque réseau est sa \textbf{dernière adresse utilisable}. Tous les équipements communiquent
        entre eux et accèdent à la page web du serveur. Nommage imposé (« Display Name »).
\end{itemize}
\end{cdc}

\question{B1}{2 pts}{Complétez le plan d'adressage \textbf{avant} la réalisation.}
{\footnotesize\setlength{\tabcolsep}{3pt}\renewcommand{\arraystretch}{1.15}
\begin{tabularx}{\linewidth}{|l|C|C|C|C|C|C|}
  \hline
  \textbf{Réseau} & \textbf{Adresse réseau} & \textbf{Masque} & \textbf{1re @ utilisable} &
  \textbf{Dernière @ utilisable} & \textbf{Diffusion} & \textbf{Passerelle} \\ \hline
  \hl \resUn & \sol{\netUn.0} & \sol{255.255.255.0} & \sol{\netUn.1} & \sol{\netUn.254} & \sol{\netUn.255} & \sol{\netUn.254} \\ \hline
  \hl \resDeux & \sol{\netDeux.0.0} & \sol{255.255.0.0} & \sol{\netDeux.0.1} & \sol{\netDeux.255.254} & \sol{\netDeux.255.255} & \sol{\netDeux.255.254} \\ \hline
  \hl Serveur & \sol{\netSrv.0} & \sol{255.255.255.0} & \sol{\netSrv.1} & \sol{\netSrv.254} & \sol{\netSrv.255} & \sol{\netSrv.254} \\ \hline
\end{tabularx}}
\bareme{0,5 pt par ligne juste + 0,5 pt si la convention de passerelle est respectée partout.}

\question{B2}{2 pts}{Construisez la topologie (équipements, câblage, nommage) et relevez l'adressage configuré.}
{\footnotesize\setlength{\tabcolsep}{3pt}\renewcommand{\arraystretch}{1.1}
\begin{tabularx}{\linewidth}{|l|C|C|C||l|C|C|}
  \hline
  \textbf{Hôte} & \textbf{Adresse IP} & \textbf{Masque} & \textbf{Passerelle} & \textbf{Interface} & \textbf{Adresse IP} & \textbf{Masque} \\ \hline
  \hl PC-\pUn1 & \sol{\netUn.1} & \sol{255.255.255.0} & \sol{\netUn.254} & \Rnom{} G0/0 & \sol{\netUn.254} & \sol{255.255.255.0} \\ \hline
  \hl PC-\pDeux1 & \sol{\netDeux.0.1} & \sol{255.255.0.0} & \sol{\netDeux.255.254} & \Rnom{} G0/1 & \sol{\netDeux.255.254} & \sol{255.255.0.0} \\ \hline
  \hl IMP-\pDeux & \sol{\netDeux.0.10} & \sol{255.255.0.0} & \sol{\netDeux.255.254} & \Rnom{} G0/2 & \sol{\netSrv.254} & \sol{255.255.255.0} \\ \hline
  \hl \srvNom & \ip{\netSrv.\srvH} & \sol{255.255.255.0} & \sol{\netSrv.254} & \multicolumn{3}{c|}{\cellcolor{black!8}} \\ \hline
\end{tabularx}}
\bareme{Adresses d'hôtes libres (uniques, dans le bon réseau) ; affectation G0/x libre. Sur le .pkt : équipements 0,5 ;
câblage 1 (câbles droits vers les switchs et switch–routeur, \textbf{croisé} serveur–routeur, liens verts) ; nommage 0,5.}

\question{B3}{3 pts}{Configurez tous les équipements : IP, masques, passerelles des hôtes ; interfaces du routeur
(« Port Status : On »).}
\bareme{Sur le .pkt : interfaces routeur 1 pt ; IP/masques des hôtes 1 pt ; passerelles (imprimante et serveur compris)
1 pt. $-0,25$ par erreur, sans descendre sous 0 par item.}

\question{B4}{3 pts}{Réalisez les tests, notez la commande et le résultat, puis faites viser chaque test réussi.}
{\small\renewcommand{\arraystretch}{1.1}
\begin{tabularx}{\linewidth}{|c|l|l|C|c|c|}
  \hline
  \textbf{N°} & \textbf{Source} & \textbf{Destination} & \textbf{Commande / action} & \textbf{Résultat} & \textbf{Visa prof.} \\ \hline
  \hl 1 & PC-\pUn1 & PC-\pUn4 & \sol{ping \netUn.4} & \sol{réussi} & \hspace{1.4cm} \\ \hline
  \hl 2 & PC-\pDeux1 & IMP-\pDeux & \sol{ping \netDeux.0.10} & \sol{réussi} & \\ \hline
  \hl 3 & PC-\pUn2 & PC-\pDeux2 & \sol{ping \netDeux.0.2} & \sol{réussi} & \\ \hline
  \hl 4 & PC-\pDeux2 & \srvNom & \sol{ping \netSrv.\srvH} & \sol{réussi} & \\ \hline
  \hl 5 & PC-\pUn3 & page web du serveur & \sol{navigateur : {\NoAutoSpacing http://\netSrv.\srvH}} & \sol{page affichée} & \\ \hline
\end{tabularx}}
\bareme{0,5 pt par test visé + 0,5 pt pour B5. Le premier ping inter-réseaux peut échouer (résolution ARP) ; les suivants doivent réussir.}

\question{B5}{incl. B4}{Test n°3 : par quels équipements passe le paquet ? Pourquoi PC-\pUn2 l'envoie-t-il au routeur ?}
\zone{2.3cm}{PC-\pUn2 $\to$ SW-\pUn{} $\to$ \Rnom{} $\to$ SW-\pDeux{} $\to$ PC-\pDeux2. Grâce à son masque, PC-\pUn2 constate que
la destination (\netDeux.0.0/16) n'est \textbf{pas dans son réseau} (\netUn.0/24) : il envoie le paquet à sa
\textbf{passerelle par défaut}, qui le route.}

%==========================================================================
\newpage
\section{Partie C — Diagnostic d'un réseau en panne \hfill /4}

Un technicien constate qu'\textbf{aucun poste ne parvient à joindre le serveur SRV-C} et que certains postes
communiquent mal entre eux. Les câbles sont du bon type.

\begin{center}
\begin{minipage}{0.37\linewidth}\centering
\begin{tikzpicture}[font=\footnotesize,
  eq/.style={draw,fill=qbg,rounded corners=2pt,minimum width=9mm,minimum height=5mm,inner sep=2pt},
  every path/.style={thick}]
  \node[eq] (SW) {SW-1};
  \node[eq,right=9mm of SW] (R) {R1};
  \node[eq,right=7mm of R] (SRV) {SRV-C};
  \node[eq,above left=5mm and -3mm of SW] (A) {PC-A};
  \node[eq,above right=5mm and -3mm of SW] (B) {PC-B};
  \node[eq,below left=5mm and -3mm of SW] (C) {PC-C};
  \node[eq,below right=5mm and -3mm of SW] (D) {PC-D};
  \foreach \n in {A,B,C,D} \draw (\n) -- (SW);
  \draw (SW) -- node[above,font=\tiny]{G0/0} (R);
  \draw (R) -- node[above,font=\tiny]{G0/1} (SRV);
\end{tikzpicture}
\end{minipage}\hfill
\begin{minipage}{0.61\linewidth}\centering
{\small\renewcommand{\arraystretch}{1.05}\setlength{\tabcolsep}{4pt}
\begin{tabular}{lllll}
  \toprule
  \textbf{Équipement} & \textbf{Adresse IP} & \textbf{Masque} & \textbf{Passerelle} & \textbf{Port} \\
  \midrule
  \Crows
  \bottomrule
\end{tabular}}
\end{minipage}
\end{center}

\question{C1}{4 pts}{Relevez les \textbf{quatre anomalies}. Pour chacune : conséquence et correction.}
{\small\renewcommand{\arraystretch}{1.1}
\ifcorrige\def\hc{}\else\def\hc{\rule[-10mm]{0pt}{14mm}}\fi
\begin{tabularx}{\linewidth}{|c|>{\raggedright\arraybackslash}p{2cm}|L|L|L|}
  \hline
  \textbf{N°} & \textbf{Équipement} & \textbf{Anomalie} & \textbf{Conséquence} & \textbf{Correction} \\ \hline
  \ifcorrige\Csol\else
  \hc 1 & & & & \\ \hline
  \hc 2 & & & & \\ \hline
  \hc 3 & & & & \\ \hline
  \hc 4 & & & & \\ \hline
  \fi
\end{tabularx}}
\bareme{1 pt par anomalie (0,25 équipement, anomalie, conséquence, correction ; ordre indifférent). \Cbareme}

%==========================================================================
\section{Bonus — Extension sans fil \hfill +1}

\begin{cdc}Ajoutez au réseau \resUn{} un \textbf{AccessPoint-PT} relié à SW-\pUn{} (SSID \ip{\ssid}, \textbf{WPA2-PSK})
et un portable \ip{Laptop-\pUn} équipé d'un module \texttt{WPC300N}, en adressage \textbf{statique}.
Vérifiez qu'il joint \srvNom.\end{cdc}
{\small\renewcommand{\arraystretch}{1.1}
\begin{tabularx}{\linewidth}{|l|C|C|C|c|}
  \hline
  \textbf{Équipement} & \textbf{Adresse IP} & \textbf{Masque} & \textbf{Passerelle} & \textbf{Visa prof.} \\ \hline
  \hl Laptop-\pUn & \sol{\netUn.50} & \sol{255.255.255.0} & \sol{\netUn.254} & \hspace{1.4cm} \\ \hline
\end{tabularx}}
\bareme{Le point d'accès fait office de pont : le portable est dans \netUn.0/24. Bonus si associé en WPA2-PSK et ping
vers \netSrv.\srvH{} réussi. Note plafonnée à 20.}

\end{document}
"
\documentclass[10pt,a4paper]{article}

\usepackage[T1]{fontenc}
\usepackage[utf8]{inputenc}
\usepackage{lmodern}
\usepackage[french]{babel}
\usepackage[a4paper,hmargin=1.5cm,top=1.3cm,bottom=1.6cm,footskip=0.7cm]{geometry}
\usepackage[table]{xcolor}
\usepackage{booktabs,tabularx,array}
\usepackage{enumitem}
\usepackage{fancyhdr}
\usepackage[most]{tcolorbox}
\usepackage{tikz}
\usetikzlibrary{positioning}
\usepackage{titlesec}

\providecommand{\modecorrige}{0}
\providecommand{\variante}{A}
\newif\ifcorrige
\ifnum\modecorrige=1 \corrigetrue \else \corrigefalse \fi

\definecolor{ipred}{HTML}{C0392B}
\definecolor{cdcbar}{HTML}{A67C3D}
\definecolor{cdcbg}{HTML}{F4F1EA}
\definecolor{qbar}{HTML}{4A6B5D}
\definecolor{qbg}{HTML}{EDF1EF}
\definecolor{corr}{HTML}{1F5FA8}
\definecolor{corrbg}{HTML}{E8F0FA}
\definecolor{grisfoot}{HTML}{828282}

\newcommand{\ip}[1]{\textcolor{ipred}{\texttt{#1}}}
\newcommand{\pts}[1]{\hfill\mbox{\textbf{(#1)}}}
\newcommand{\sol}[1]{\ifcorrige\textcolor{corr}{#1}\fi}
\newcommand{\bareme}[1]{\ifcorrige\par{\footnotesize\color{corr}#1\par}\fi}
\newcommand{\hl}{\ifcorrige\rule[-1.6mm]{0pt}{5.4mm}\else\rule[-2.2mm]{0pt}{6.6mm}\fi}
\newcolumntype{C}{>{\centering\arraybackslash}X}
\newcolumntype{L}{>{\raggedright\arraybackslash}X}

\titleformat{\section}{\large\bfseries}{}{0pt}{}
\titlespacing*{\section}{0pt}{6pt}{3pt}

\tcbset{qstyle/.style={enhanced,boxrule=0pt,leftrule=2.5pt,arc=0pt,outer arc=0pt,
  left=5pt,right=5pt,top=2pt,bottom=2pt,before skip=4pt,after skip=3pt}}
\newtcolorbox{cdc}{qstyle,colback=cdcbg,colframe=cdcbar}
\newtcolorbox{quest}{qstyle,colback=qbg,colframe=qbar}

% Zone de réponse : vierge dans le sujet, remplie dans le corrigé
\newcommand{\zone}[2]{%
  \ifcorrige
    \begin{tcolorbox}[colback=corrbg,colframe=corr,boxrule=0.5pt,arc=1pt,left=4pt,right=4pt,top=1pt,bottom=1pt,
      before skip=2pt,after skip=2pt,fontupper=\footnotesize\color{corr}]#2\end{tcolorbox}%
  \else
    \begin{tcolorbox}[colback=white,colframe=black!45,boxrule=0.4pt,arc=1pt,height=#1,
      before skip=2pt,after skip=2pt]\end{tcolorbox}%
  \fi}

\newcommand{\question}[3]{\begin{quest}\textbf{#1.} #3\pts{#2}\end{quest}}

%==========================================================================
% Données propres à chaque variante
%==========================================================================
\if\variante A
  \def\entreprise{Le cabinet d'architectes \textbf{ArchiNova}}
  \def\Rnom{R-ArchiNova}
  \def\resUn{Bureau d'études}\def\pUn{BE}\def\netUn{192.168.40}
  \def\resDeux{Accueil}\def\pDeux{ACC}\def\netDeux{172.16}
  \def\srvNom{SRV-Intranet}\def\netSrv{192.168.200}\def\srvH{10}
  \def\ssid{ArchiNova-WiFi}
  \def\AIrows{%
    \hl\ip{192.168.40.25} & \sol{C} & \sol{255.255.255.0} & \sol{192.168.40.0} & \sol{192.168.40.255} & \sol{254} \\ \hline
    \hl\ip{172.16.5.10}   & \sol{B} & \sol{255.255.0.0}   & \sol{172.16.0.0}   & \sol{172.16.255.255} & \sol{65\,534} \\ \hline
    \hl\ip{10.3.2.1}      & \sol{A} & \sol{255.0.0.0}     & \sol{10.0.0.0}     & \sol{10.255.255.255} & \sol{16\,777\,214} \\ \hline}
  \def\apipa{169.254.12.7}
  \def\Cnet{192.168.60}
  \def\Crows{%
    PC-A & \ip{192.168.60.10} & \ip{255.255.255.0} & \ip{192.168.60.254} & — \\
    PC-B & \ip{192.168.60.20} & \ip{255.255.255.0} & \ip{192.168.60.1}   & — \\
    PC-C & \ip{192.168.61.30} & \ip{255.255.255.0} & \ip{192.168.60.254} & — \\
    PC-D & \ip{192.168.60.10} & \ip{255.255.255.0} & \ip{192.168.60.254} & — \\
    SRV-C & \ip{10.0.0.5}     & \ip{255.0.0.0}     & \ip{10.0.0.254}     & — \\
    R1 — G0/0 & \ip{192.168.60.254} & \ip{255.255.255.0} & — & On \\
    R1 — G0/1 & \ip{10.0.0.254}     & \ip{255.0.0.0}     & — & Off \\}
  \def\Csol{%
    \hl 1 & \sol{PC-B} & \sol{Passerelle .60.1 : pas l'adresse du routeur.} & \sol{Joint son réseau, mais aucun autre (serveur).} & \sol{Passerelle 192.168.60.254.} \\ \hline
    \hl 2 & \sol{PC-C} & \sol{IP hors du réseau 192.168.60.0/24.} & \sol{Ne joint ni les postes ni le routeur.} & \sol{Ex. 192.168.60.30.} \\ \hline
    \hl 3 & \sol{PC-D} & \sol{Même IP que PC-A (doublon).} & \sol{Conflit d'adresse : PC-A et PC-D perturbés.} & \sol{Adresse unique, ex. .60.40.} \\ \hline
    \hl 4 & \sol{R1 G0/1} & \sol{Interface désactivée (Off).} & \sol{Réseau du serveur injoignable pour tous.} & \sol{Port Status : On.} \\ \hline}
  \def\Cbareme{PC-A et SRV-C sont correctement configurés ; accepter PC-A cité à la place de PC-D pour le doublon.}
\else
  \def\entreprise{La clinique vétérinaire \textbf{VetoPlus}}
  \def\Rnom{R-VetoPlus}
  \def\resUn{Soins}\def\pUn{SOI}\def\netUn{192.168.70}
  \def\resDeux{Secrétariat}\def\pDeux{SEC}\def\netDeux{172.20}
  \def\srvNom{SRV-Dossiers}\def\netSrv{192.168.250}\def\srvH{20}
  \def\ssid{VetoPlus-WiFi}
  \def\AIrows{%
    \hl\ip{192.168.15.130} & \sol{C} & \sol{255.255.255.0} & \sol{192.168.15.0} & \sol{192.168.15.255} & \sol{254} \\ \hline
    \hl\ip{172.31.200.7}   & \sol{B} & \sol{255.255.0.0}   & \sol{172.31.0.0}   & \sol{172.31.255.255} & \sol{65\,534} \\ \hline
    \hl\ip{10.45.6.9}      & \sol{A} & \sol{255.0.0.0}     & \sol{10.0.0.0}     & \sol{10.255.255.255} & \sol{16\,777\,214} \\ \hline}
  \def\apipa{169.254.88.3}
  \def\Cnet{192.168.80}
  \def\Crows{%
    PC-A & \ip{192.168.8.11}   & \ip{255.255.255.0} & \ip{192.168.80.254} & — \\
    PC-B & \ip{192.168.80.12}  & \ip{255.255.255.0} & \ip{192.168.80.254} & — \\
    PC-C & \ip{192.168.80.13}  & \ip{255.255.255.0} & \ip{192.168.80.100} & — \\
    PC-D & \ip{192.168.80.12}  & \ip{255.255.255.0} & \ip{192.168.80.254} & — \\
    SRV-C & \ip{10.0.0.8}      & \ip{255.0.0.0}     & \ip{10.0.0.254}     & — \\
    R1 — G0/0 & \ip{192.168.80.254} & \ip{255.255.255.0} & — & Off \\
    R1 — G0/1 & \ip{10.0.0.254}     & \ip{255.0.0.0}     & — & On \\}
  \def\Csol{%
    \hl 1 & \sol{PC-A} & \sol{IP hors du réseau 192.168.80.0/24.} & \sol{Ne joint ni les postes ni le routeur.} & \sol{Ex. 192.168.80.11.} \\ \hline
    \hl 2 & \sol{PC-C} & \sol{Passerelle .80.100 : pas l'adresse du routeur.} & \sol{Joint son réseau, mais aucun autre (serveur).} & \sol{Passerelle 192.168.80.254.} \\ \hline
    \hl 3 & \sol{PC-D} & \sol{Même IP que PC-B (doublon).} & \sol{Conflit d'adresse : PC-B et PC-D perturbés.} & \sol{Adresse unique, ex. .80.14.} \\ \hline
    \hl 4 & \sol{R1 G0/0} & \sol{Interface côté LAN désactivée (Off).} & \sol{Aucun poste ne joint le routeur ni le serveur.} & \sol{Port Status : On.} \\ \hline}
  \def\Cbareme{PC-B et SRV-C sont correctement configurés ; accepter PC-B cité à la place de PC-D pour le doublon.}
\fi

\pagestyle{fancy}
\fancyhf{}
\renewcommand{\headrulewidth}{0pt}
\fancyfoot[L]{\footnotesize\itshape\color{grisfoot}Évaluation — Réseaux locaux — Sujet \variante\ifcorrige\ — CORRIGÉ\fi}
\fancyfoot[R]{\footnotesize\color{grisfoot}\thepage/3}

\setlist{itemsep=0pt,topsep=2pt,parsep=0pt}
\setlength{\parindent}{0pt}
\setlength{\parskip}{2pt}

\begin{document}

\begin{center}
  {\Large\bfseries Évaluation pratique — Réseaux locaux \quad
   \fcolorbox{black}{white}{\ SUJET \variante\ }\par}
  \vspace{1pt}
  {\bfseries Adressage IP, commutation et routage — BTS CIEL 1re année — durée : 2 heures\par}
  \ifcorrige\vspace{2pt}{\bfseries\color{corr}CORRIGÉ ET BARÈME — document enseignant}\par\fi
\end{center}
\vspace{-2pt}
{\renewcommand{\arraystretch}{1.35}
\begin{tabularx}{\linewidth}{|X|X|p{2.2cm}|p{2.6cm}|}
  \hline
  \textbf{Nom :} & \textbf{Prénom :} & \textbf{Poste n°} & \textbf{Note :} \hfill /20 \\ \hline
\end{tabularx}}

\textbf{Consignes :} travail \textbf{individuel}, documents et Internet \textbf{non autorisés}. Ce sujet est le
\textbf{document-réponse}. Fichier Packet Tracer : \ip{Nom\_Prenom\_EvalLAN\_\variante.pkt}, à enregistrer
régulièrement dans le dossier de rendu. Les « visas » sont apposés par le professeur, que vous appelez à chaque test réussi.
Durées conseillées : A 30 min — B 1 h 10 — C 20 min.

%==========================================================================
\section{Partie A — Questions de cours et calculs \hfill /6}

\question{A1}{2 pts}{Complétez le tableau en considérant le \textbf{masque par défaut} de la classe de chaque adresse.}
{\small\renewcommand{\arraystretch}{1.15}
\begin{tabularx}{\linewidth}{|l|c|C|C|C|c|}
  \hline
  \textbf{Adresse IP} & \textbf{Classe} & \textbf{Masque par défaut} & \textbf{Adresse réseau} &
  \textbf{Adresse de diffusion} & \textbf{Nb d'hôtes} \\ \hline
  \AIrows
\end{tabularx}}
\bareme{0,5 pt par ligne juste + 0,5 pt si les trois nombres d'hôtes sont justes ($2^{n}-2$, $n$ = nombre de bits d'hôte).}

\question{A2}{1 pt}{Indiquez le type de câble (\textbf{droit} ou \textbf{croisé}) pour chaque liaison.}
{\small\renewcommand{\arraystretch}{1.15}
\begin{tabularx}{\linewidth}{|X|c||X|c|}
  \hline
  \hl PC $\leftrightarrow$ commutateur & \makebox[1.6cm]{\sol{droit}} & Commutateur $\leftrightarrow$ routeur & \makebox[1.6cm]{\sol{droit}} \\ \hline
  \hl PC $\leftrightarrow$ PC (réseau P2P) & \sol{croisé} & Concentrateur $\leftrightarrow$ concentrateur & \sol{croisé} \\ \hline
\end{tabularx}}
\bareme{0,25 pt par case. Équipements de types différents : câble droit ; de même type : câble croisé.}

\question{A3}{1 pt}{Comment un \textbf{concentrateur (hub)} et un \textbf{commutateur (switch)} traitent-ils une trame
reçue ? Lequel limite les collisions, et pourquoi ?}
\zone{2.9cm}{Le \textbf{hub} (couche 1) répète la trame sur \textbf{tous ses autres ports} : un seul domaine de collision
partagé. Le \textbf{switch} (couche 2) lit l'adresse MAC de destination et, grâce à sa \textbf{table MAC}, ne l'envoie
\textbf{que sur le port du destinataire} : un domaine de collision par port, c'est lui qui limite les collisions.
\emph{0,5 pt hub, 0,5 pt switch.}}

\question{A4}{1 pt}{Qu'est-ce que la \textbf{passerelle par défaut} d'un poste ? Quelle condition son adresse doit-elle respecter ?}
\zone{2.4cm}{Adresse IP de l'interface du \textbf{routeur} à laquelle le poste envoie les paquets destinés à un
\textbf{autre réseau}. Elle doit \textbf{appartenir au même réseau} que le poste. \emph{0,5 + 0,5 pt.}}

\question{A5}{0,5 pt}{Un portable configuré en DHCP affiche \ip{\apipa} / \ip{255.255.0.0}. Que s'est-il passé ?}
\zone{1.5cm}{Aucun \textbf{serveur DHCP} n'a répondu : le poste s'est attribué une adresse \textbf{APIPA} (169.254.0.0/16).}

\question{A6}{0,5 pt}{Quel protocole utilise \texttt{ping} ? Nommez les deux messages échangés et la couche OSI des adresses IP.}
\zone{1.5cm}{\textbf{ICMP} ; \textbf{Echo Request} / \textbf{Echo Reply} ; adresses IP en \textbf{couche 3 (Réseau)}.}

%==========================================================================
\newpage
\section{Partie B — Réalisation sous Packet Tracer \hfill /10}

\begin{cdc}
\entreprise{} vous confie son réseau : \textbf{trois réseaux} reliés par un \textbf{routeur unique Cisco 2911}
nommé \ip{\Rnom} (interfaces G0/0, G0/1, G0/2).
\begin{itemize}[leftmargin=12pt]
  \item \textbf{Réseau \resUn} : \textbf{4 postes} PC-\pUn1 à PC-\pUn4, commutateur \ip{SW-\pUn}, adressage statique,
        réseau \ip{\netUn.0/24}.
  \item \textbf{Réseau \resDeux} : \textbf{2 postes} PC-\pDeux1, PC-\pDeux2 et \textbf{1 imprimante} IMP-\pDeux,
        commutateur \ip{SW-\pDeux}, adressage statique, réseau de classe B \ip{\netDeux.0.0/16}.
  \item \textbf{Réseau Serveur} : serveur \ip{\srvNom} d'adresse \ip{\netSrv.\srvH/24}, relié \textbf{directement} au routeur.
  \item La passerelle de chaque réseau est sa \textbf{dernière adresse utilisable}. Tous les équipements communiquent
        entre eux et accèdent à la page web du serveur. Nommage imposé (« Display Name »).
\end{itemize}
\end{cdc}

\question{B1}{2 pts}{Complétez le plan d'adressage \textbf{avant} la réalisation.}
{\footnotesize\setlength{\tabcolsep}{3pt}\renewcommand{\arraystretch}{1.15}
\begin{tabularx}{\linewidth}{|l|C|C|C|C|C|C|}
  \hline
  \textbf{Réseau} & \textbf{Adresse réseau} & \textbf{Masque} & \textbf{1re @ utilisable} &
  \textbf{Dernière @ utilisable} & \textbf{Diffusion} & \textbf{Passerelle} \\ \hline
  \hl \resUn & \sol{\netUn.0} & \sol{255.255.255.0} & \sol{\netUn.1} & \sol{\netUn.254} & \sol{\netUn.255} & \sol{\netUn.254} \\ \hline
  \hl \resDeux & \sol{\netDeux.0.0} & \sol{255.255.0.0} & \sol{\netDeux.0.1} & \sol{\netDeux.255.254} & \sol{\netDeux.255.255} & \sol{\netDeux.255.254} \\ \hline
  \hl Serveur & \sol{\netSrv.0} & \sol{255.255.255.0} & \sol{\netSrv.1} & \sol{\netSrv.254} & \sol{\netSrv.255} & \sol{\netSrv.254} \\ \hline
\end{tabularx}}
\bareme{0,5 pt par ligne juste + 0,5 pt si la convention de passerelle est respectée partout.}

\question{B2}{2 pts}{Construisez la topologie (équipements, câblage, nommage) et relevez l'adressage configuré.}
{\footnotesize\setlength{\tabcolsep}{3pt}\renewcommand{\arraystretch}{1.1}
\begin{tabularx}{\linewidth}{|l|C|C|C||l|C|C|}
  \hline
  \textbf{Hôte} & \textbf{Adresse IP} & \textbf{Masque} & \textbf{Passerelle} & \textbf{Interface} & \textbf{Adresse IP} & \textbf{Masque} \\ \hline
  \hl PC-\pUn1 & \sol{\netUn.1} & \sol{255.255.255.0} & \sol{\netUn.254} & \Rnom{} G0/0 & \sol{\netUn.254} & \sol{255.255.255.0} \\ \hline
  \hl PC-\pDeux1 & \sol{\netDeux.0.1} & \sol{255.255.0.0} & \sol{\netDeux.255.254} & \Rnom{} G0/1 & \sol{\netDeux.255.254} & \sol{255.255.0.0} \\ \hline
  \hl IMP-\pDeux & \sol{\netDeux.0.10} & \sol{255.255.0.0} & \sol{\netDeux.255.254} & \Rnom{} G0/2 & \sol{\netSrv.254} & \sol{255.255.255.0} \\ \hline
  \hl \srvNom & \ip{\netSrv.\srvH} & \sol{255.255.255.0} & \sol{\netSrv.254} & \multicolumn{3}{c|}{\cellcolor{black!8}} \\ \hline
\end{tabularx}}
\bareme{Adresses d'hôtes libres (uniques, dans le bon réseau) ; affectation G0/x libre. Sur le .pkt : équipements 0,5 ;
câblage 1 (câbles droits vers les switchs et switch–routeur, \textbf{croisé} serveur–routeur, liens verts) ; nommage 0,5.}

\question{B3}{3 pts}{Configurez tous les équipements : IP, masques, passerelles des hôtes ; interfaces du routeur
(« Port Status : On »).}
\bareme{Sur le .pkt : interfaces routeur 1 pt ; IP/masques des hôtes 1 pt ; passerelles (imprimante et serveur compris)
1 pt. $-0,25$ par erreur, sans descendre sous 0 par item.}

\question{B4}{3 pts}{Réalisez les tests, notez la commande et le résultat, puis faites viser chaque test réussi.}
{\small\renewcommand{\arraystretch}{1.1}
\begin{tabularx}{\linewidth}{|c|l|l|C|c|c|}
  \hline
  \textbf{N°} & \textbf{Source} & \textbf{Destination} & \textbf{Commande / action} & \textbf{Résultat} & \textbf{Visa prof.} \\ \hline
  \hl 1 & PC-\pUn1 & PC-\pUn4 & \sol{ping \netUn.4} & \sol{réussi} & \hspace{1.4cm} \\ \hline
  \hl 2 & PC-\pDeux1 & IMP-\pDeux & \sol{ping \netDeux.0.10} & \sol{réussi} & \\ \hline
  \hl 3 & PC-\pUn2 & PC-\pDeux2 & \sol{ping \netDeux.0.2} & \sol{réussi} & \\ \hline
  \hl 4 & PC-\pDeux2 & \srvNom & \sol{ping \netSrv.\srvH} & \sol{réussi} & \\ \hline
  \hl 5 & PC-\pUn3 & page web du serveur & \sol{navigateur : {\NoAutoSpacing http://\netSrv.\srvH}} & \sol{page affichée} & \\ \hline
\end{tabularx}}
\bareme{0,5 pt par test visé + 0,5 pt pour B5. Le premier ping inter-réseaux peut échouer (résolution ARP) ; les suivants doivent réussir.}

\question{B5}{incl. B4}{Test n°3 : par quels équipements passe le paquet ? Pourquoi PC-\pUn2 l'envoie-t-il au routeur ?}
\zone{2.3cm}{PC-\pUn2 $\to$ SW-\pUn{} $\to$ \Rnom{} $\to$ SW-\pDeux{} $\to$ PC-\pDeux2. Grâce à son masque, PC-\pUn2 constate que
la destination (\netDeux.0.0/16) n'est \textbf{pas dans son réseau} (\netUn.0/24) : il envoie le paquet à sa
\textbf{passerelle par défaut}, qui le route.}

%==========================================================================
\newpage
\section{Partie C — Diagnostic d'un réseau en panne \hfill /4}

Un technicien constate qu'\textbf{aucun poste ne parvient à joindre le serveur SRV-C} et que certains postes
communiquent mal entre eux. Les câbles sont du bon type.

\begin{center}
\begin{minipage}{0.37\linewidth}\centering
\begin{tikzpicture}[font=\footnotesize,
  eq/.style={draw,fill=qbg,rounded corners=2pt,minimum width=9mm,minimum height=5mm,inner sep=2pt},
  every path/.style={thick}]
  \node[eq] (SW) {SW-1};
  \node[eq,right=9mm of SW] (R) {R1};
  \node[eq,right=7mm of R] (SRV) {SRV-C};
  \node[eq,above left=5mm and -3mm of SW] (A) {PC-A};
  \node[eq,above right=5mm and -3mm of SW] (B) {PC-B};
  \node[eq,below left=5mm and -3mm of SW] (C) {PC-C};
  \node[eq,below right=5mm and -3mm of SW] (D) {PC-D};
  \foreach \n in {A,B,C,D} \draw (\n) -- (SW);
  \draw (SW) -- node[above,font=\tiny]{G0/0} (R);
  \draw (R) -- node[above,font=\tiny]{G0/1} (SRV);
\end{tikzpicture}
\end{minipage}\hfill
\begin{minipage}{0.61\linewidth}\centering
{\small\renewcommand{\arraystretch}{1.05}\setlength{\tabcolsep}{4pt}
\begin{tabular}{lllll}
  \toprule
  \textbf{Équipement} & \textbf{Adresse IP} & \textbf{Masque} & \textbf{Passerelle} & \textbf{Port} \\
  \midrule
  \Crows
  \bottomrule
\end{tabular}}
\end{minipage}
\end{center}

\question{C1}{4 pts}{Relevez les \textbf{quatre anomalies}. Pour chacune : conséquence et correction.}
{\small\renewcommand{\arraystretch}{1.1}
\ifcorrige\def\hc{}\else\def\hc{\rule[-10mm]{0pt}{14mm}}\fi
\begin{tabularx}{\linewidth}{|c|>{\raggedright\arraybackslash}p{2cm}|L|L|L|}
  \hline
  \textbf{N°} & \textbf{Équipement} & \textbf{Anomalie} & \textbf{Conséquence} & \textbf{Correction} \\ \hline
  \ifcorrige\Csol\else
  \hc 1 & & & & \\ \hline
  \hc 2 & & & & \\ \hline
  \hc 3 & & & & \\ \hline
  \hc 4 & & & & \\ \hline
  \fi
\end{tabularx}}
\bareme{1 pt par anomalie (0,25 équipement, anomalie, conséquence, correction ; ordre indifférent). \Cbareme}

%==========================================================================
\section{Bonus — Extension sans fil \hfill +1}

\begin{cdc}Ajoutez au réseau \resUn{} un \textbf{AccessPoint-PT} relié à SW-\pUn{} (SSID \ip{\ssid}, \textbf{WPA2-PSK})
et un portable \ip{Laptop-\pUn} équipé d'un module \texttt{WPC300N}, en adressage \textbf{statique}.
Vérifiez qu'il joint \srvNom.\end{cdc}
{\small\renewcommand{\arraystretch}{1.1}
\begin{tabularx}{\linewidth}{|l|C|C|C|c|}
  \hline
  \textbf{Équipement} & \textbf{Adresse IP} & \textbf{Masque} & \textbf{Passerelle} & \textbf{Visa prof.} \\ \hline
  \hl Laptop-\pUn & \sol{\netUn.50} & \sol{255.255.255.0} & \sol{\netUn.254} & \hspace{1.4cm} \\ \hline
\end{tabularx}}
\bareme{Le point d'accès fait office de pont : le portable est dans \netUn.0/24. Bonus si associé en WPA2-PSK et ping
vers \netSrv.\srvH{} réussi. Note plafonnée à 20.}

\end{document}
"
%   pdflatex -jobname=Eval_Reseau_SujetB   "\def\variante{B}\def\modecorrige{0}% Évaluation pratique — Réseaux locaux (BTS CIEL 1re année) — sujets A et B
% Compilation (4 PDF) :
%   pdflatex -jobname=Eval_Reseau_SujetA   "\def\variante{A}\def\modecorrige{0}% Évaluation pratique — Réseaux locaux (BTS CIEL 1re année) — sujets A et B
% Compilation (4 PDF) :
%   pdflatex -jobname=Eval_Reseau_SujetA   "\def\variante{A}\def\modecorrige{0}\input{Eval_TP_Reseau_LAN.tex}"
%   pdflatex -jobname=Eval_Reseau_SujetB   "\def\variante{B}\def\modecorrige{0}\input{Eval_TP_Reseau_LAN.tex}"
%   pdflatex -jobname=Eval_Reseau_CorrigeA "\def\variante{A}\def\modecorrige{1}\input{Eval_TP_Reseau_LAN.tex}"
%   pdflatex -jobname=Eval_Reseau_CorrigeB "\def\variante{B}\def\modecorrige{1}\input{Eval_TP_Reseau_LAN.tex}"
\documentclass[10pt,a4paper]{article}

\usepackage[T1]{fontenc}
\usepackage[utf8]{inputenc}
\usepackage{lmodern}
\usepackage[french]{babel}
\usepackage[a4paper,hmargin=1.5cm,top=1.3cm,bottom=1.6cm,footskip=0.7cm]{geometry}
\usepackage[table]{xcolor}
\usepackage{booktabs,tabularx,array}
\usepackage{enumitem}
\usepackage{fancyhdr}
\usepackage[most]{tcolorbox}
\usepackage{tikz}
\usetikzlibrary{positioning}
\usepackage{titlesec}

\providecommand{\modecorrige}{0}
\providecommand{\variante}{A}
\newif\ifcorrige
\ifnum\modecorrige=1 \corrigetrue \else \corrigefalse \fi

\definecolor{ipred}{HTML}{C0392B}
\definecolor{cdcbar}{HTML}{A67C3D}
\definecolor{cdcbg}{HTML}{F4F1EA}
\definecolor{qbar}{HTML}{4A6B5D}
\definecolor{qbg}{HTML}{EDF1EF}
\definecolor{corr}{HTML}{1F5FA8}
\definecolor{corrbg}{HTML}{E8F0FA}
\definecolor{grisfoot}{HTML}{828282}

\newcommand{\ip}[1]{\textcolor{ipred}{\texttt{#1}}}
\newcommand{\pts}[1]{\hfill\mbox{\textbf{(#1)}}}
\newcommand{\sol}[1]{\ifcorrige\textcolor{corr}{#1}\fi}
\newcommand{\bareme}[1]{\ifcorrige\par{\footnotesize\color{corr}#1\par}\fi}
\newcommand{\hl}{\ifcorrige\rule[-1.6mm]{0pt}{5.4mm}\else\rule[-2.2mm]{0pt}{6.6mm}\fi}
\newcolumntype{C}{>{\centering\arraybackslash}X}
\newcolumntype{L}{>{\raggedright\arraybackslash}X}

\titleformat{\section}{\large\bfseries}{}{0pt}{}
\titlespacing*{\section}{0pt}{6pt}{3pt}

\tcbset{qstyle/.style={enhanced,boxrule=0pt,leftrule=2.5pt,arc=0pt,outer arc=0pt,
  left=5pt,right=5pt,top=2pt,bottom=2pt,before skip=4pt,after skip=3pt}}
\newtcolorbox{cdc}{qstyle,colback=cdcbg,colframe=cdcbar}
\newtcolorbox{quest}{qstyle,colback=qbg,colframe=qbar}

% Zone de réponse : vierge dans le sujet, remplie dans le corrigé
\newcommand{\zone}[2]{%
  \ifcorrige
    \begin{tcolorbox}[colback=corrbg,colframe=corr,boxrule=0.5pt,arc=1pt,left=4pt,right=4pt,top=1pt,bottom=1pt,
      before skip=2pt,after skip=2pt,fontupper=\footnotesize\color{corr}]#2\end{tcolorbox}%
  \else
    \begin{tcolorbox}[colback=white,colframe=black!45,boxrule=0.4pt,arc=1pt,height=#1,
      before skip=2pt,after skip=2pt]\end{tcolorbox}%
  \fi}

\newcommand{\question}[3]{\begin{quest}\textbf{#1.} #3\pts{#2}\end{quest}}

%==========================================================================
% Données propres à chaque variante
%==========================================================================
\if\variante A
  \def\entreprise{Le cabinet d'architectes \textbf{ArchiNova}}
  \def\Rnom{R-ArchiNova}
  \def\resUn{Bureau d'études}\def\pUn{BE}\def\netUn{192.168.40}
  \def\resDeux{Accueil}\def\pDeux{ACC}\def\netDeux{172.16}
  \def\srvNom{SRV-Intranet}\def\netSrv{192.168.200}\def\srvH{10}
  \def\ssid{ArchiNova-WiFi}
  \def\AIrows{%
    \hl\ip{192.168.40.25} & \sol{C} & \sol{255.255.255.0} & \sol{192.168.40.0} & \sol{192.168.40.255} & \sol{254} \\ \hline
    \hl\ip{172.16.5.10}   & \sol{B} & \sol{255.255.0.0}   & \sol{172.16.0.0}   & \sol{172.16.255.255} & \sol{65\,534} \\ \hline
    \hl\ip{10.3.2.1}      & \sol{A} & \sol{255.0.0.0}     & \sol{10.0.0.0}     & \sol{10.255.255.255} & \sol{16\,777\,214} \\ \hline}
  \def\apipa{169.254.12.7}
  \def\Cnet{192.168.60}
  \def\Crows{%
    PC-A & \ip{192.168.60.10} & \ip{255.255.255.0} & \ip{192.168.60.254} & — \\
    PC-B & \ip{192.168.60.20} & \ip{255.255.255.0} & \ip{192.168.60.1}   & — \\
    PC-C & \ip{192.168.61.30} & \ip{255.255.255.0} & \ip{192.168.60.254} & — \\
    PC-D & \ip{192.168.60.10} & \ip{255.255.255.0} & \ip{192.168.60.254} & — \\
    SRV-C & \ip{10.0.0.5}     & \ip{255.0.0.0}     & \ip{10.0.0.254}     & — \\
    R1 — G0/0 & \ip{192.168.60.254} & \ip{255.255.255.0} & — & On \\
    R1 — G0/1 & \ip{10.0.0.254}     & \ip{255.0.0.0}     & — & Off \\}
  \def\Csol{%
    \hl 1 & \sol{PC-B} & \sol{Passerelle .60.1 : pas l'adresse du routeur.} & \sol{Joint son réseau, mais aucun autre (serveur).} & \sol{Passerelle 192.168.60.254.} \\ \hline
    \hl 2 & \sol{PC-C} & \sol{IP hors du réseau 192.168.60.0/24.} & \sol{Ne joint ni les postes ni le routeur.} & \sol{Ex. 192.168.60.30.} \\ \hline
    \hl 3 & \sol{PC-D} & \sol{Même IP que PC-A (doublon).} & \sol{Conflit d'adresse : PC-A et PC-D perturbés.} & \sol{Adresse unique, ex. .60.40.} \\ \hline
    \hl 4 & \sol{R1 G0/1} & \sol{Interface désactivée (Off).} & \sol{Réseau du serveur injoignable pour tous.} & \sol{Port Status : On.} \\ \hline}
  \def\Cbareme{PC-A et SRV-C sont correctement configurés ; accepter PC-A cité à la place de PC-D pour le doublon.}
\else
  \def\entreprise{La clinique vétérinaire \textbf{VetoPlus}}
  \def\Rnom{R-VetoPlus}
  \def\resUn{Soins}\def\pUn{SOI}\def\netUn{192.168.70}
  \def\resDeux{Secrétariat}\def\pDeux{SEC}\def\netDeux{172.20}
  \def\srvNom{SRV-Dossiers}\def\netSrv{192.168.250}\def\srvH{20}
  \def\ssid{VetoPlus-WiFi}
  \def\AIrows{%
    \hl\ip{192.168.15.130} & \sol{C} & \sol{255.255.255.0} & \sol{192.168.15.0} & \sol{192.168.15.255} & \sol{254} \\ \hline
    \hl\ip{172.31.200.7}   & \sol{B} & \sol{255.255.0.0}   & \sol{172.31.0.0}   & \sol{172.31.255.255} & \sol{65\,534} \\ \hline
    \hl\ip{10.45.6.9}      & \sol{A} & \sol{255.0.0.0}     & \sol{10.0.0.0}     & \sol{10.255.255.255} & \sol{16\,777\,214} \\ \hline}
  \def\apipa{169.254.88.3}
  \def\Cnet{192.168.80}
  \def\Crows{%
    PC-A & \ip{192.168.8.11}   & \ip{255.255.255.0} & \ip{192.168.80.254} & — \\
    PC-B & \ip{192.168.80.12}  & \ip{255.255.255.0} & \ip{192.168.80.254} & — \\
    PC-C & \ip{192.168.80.13}  & \ip{255.255.255.0} & \ip{192.168.80.100} & — \\
    PC-D & \ip{192.168.80.12}  & \ip{255.255.255.0} & \ip{192.168.80.254} & — \\
    SRV-C & \ip{10.0.0.8}      & \ip{255.0.0.0}     & \ip{10.0.0.254}     & — \\
    R1 — G0/0 & \ip{192.168.80.254} & \ip{255.255.255.0} & — & Off \\
    R1 — G0/1 & \ip{10.0.0.254}     & \ip{255.0.0.0}     & — & On \\}
  \def\Csol{%
    \hl 1 & \sol{PC-A} & \sol{IP hors du réseau 192.168.80.0/24.} & \sol{Ne joint ni les postes ni le routeur.} & \sol{Ex. 192.168.80.11.} \\ \hline
    \hl 2 & \sol{PC-C} & \sol{Passerelle .80.100 : pas l'adresse du routeur.} & \sol{Joint son réseau, mais aucun autre (serveur).} & \sol{Passerelle 192.168.80.254.} \\ \hline
    \hl 3 & \sol{PC-D} & \sol{Même IP que PC-B (doublon).} & \sol{Conflit d'adresse : PC-B et PC-D perturbés.} & \sol{Adresse unique, ex. .80.14.} \\ \hline
    \hl 4 & \sol{R1 G0/0} & \sol{Interface côté LAN désactivée (Off).} & \sol{Aucun poste ne joint le routeur ni le serveur.} & \sol{Port Status : On.} \\ \hline}
  \def\Cbareme{PC-B et SRV-C sont correctement configurés ; accepter PC-B cité à la place de PC-D pour le doublon.}
\fi

\pagestyle{fancy}
\fancyhf{}
\renewcommand{\headrulewidth}{0pt}
\fancyfoot[L]{\footnotesize\itshape\color{grisfoot}Évaluation — Réseaux locaux — Sujet \variante\ifcorrige\ — CORRIGÉ\fi}
\fancyfoot[R]{\footnotesize\color{grisfoot}\thepage/3}

\setlist{itemsep=0pt,topsep=2pt,parsep=0pt}
\setlength{\parindent}{0pt}
\setlength{\parskip}{2pt}

\begin{document}

\begin{center}
  {\Large\bfseries Évaluation pratique — Réseaux locaux \quad
   \fcolorbox{black}{white}{\ SUJET \variante\ }\par}
  \vspace{1pt}
  {\bfseries Adressage IP, commutation et routage — BTS CIEL 1re année — durée : 2 heures\par}
  \ifcorrige\vspace{2pt}{\bfseries\color{corr}CORRIGÉ ET BARÈME — document enseignant}\par\fi
\end{center}
\vspace{-2pt}
{\renewcommand{\arraystretch}{1.35}
\begin{tabularx}{\linewidth}{|X|X|p{2.2cm}|p{2.6cm}|}
  \hline
  \textbf{Nom :} & \textbf{Prénom :} & \textbf{Poste n°} & \textbf{Note :} \hfill /20 \\ \hline
\end{tabularx}}

\textbf{Consignes :} travail \textbf{individuel}, documents et Internet \textbf{non autorisés}. Ce sujet est le
\textbf{document-réponse}. Fichier Packet Tracer : \ip{Nom\_Prenom\_EvalLAN\_\variante.pkt}, à enregistrer
régulièrement dans le dossier de rendu. Les « visas » sont apposés par le professeur, que vous appelez à chaque test réussi.
Durées conseillées : A 30 min — B 1 h 10 — C 20 min.

%==========================================================================
\section{Partie A — Questions de cours et calculs \hfill /6}

\question{A1}{2 pts}{Complétez le tableau en considérant le \textbf{masque par défaut} de la classe de chaque adresse.}
{\small\renewcommand{\arraystretch}{1.15}
\begin{tabularx}{\linewidth}{|l|c|C|C|C|c|}
  \hline
  \textbf{Adresse IP} & \textbf{Classe} & \textbf{Masque par défaut} & \textbf{Adresse réseau} &
  \textbf{Adresse de diffusion} & \textbf{Nb d'hôtes} \\ \hline
  \AIrows
\end{tabularx}}
\bareme{0,5 pt par ligne juste + 0,5 pt si les trois nombres d'hôtes sont justes ($2^{n}-2$, $n$ = nombre de bits d'hôte).}

\question{A2}{1 pt}{Indiquez le type de câble (\textbf{droit} ou \textbf{croisé}) pour chaque liaison.}
{\small\renewcommand{\arraystretch}{1.15}
\begin{tabularx}{\linewidth}{|X|c||X|c|}
  \hline
  \hl PC $\leftrightarrow$ commutateur & \makebox[1.6cm]{\sol{droit}} & Commutateur $\leftrightarrow$ routeur & \makebox[1.6cm]{\sol{droit}} \\ \hline
  \hl PC $\leftrightarrow$ PC (réseau P2P) & \sol{croisé} & Concentrateur $\leftrightarrow$ concentrateur & \sol{croisé} \\ \hline
\end{tabularx}}
\bareme{0,25 pt par case. Équipements de types différents : câble droit ; de même type : câble croisé.}

\question{A3}{1 pt}{Comment un \textbf{concentrateur (hub)} et un \textbf{commutateur (switch)} traitent-ils une trame
reçue ? Lequel limite les collisions, et pourquoi ?}
\zone{2.9cm}{Le \textbf{hub} (couche 1) répète la trame sur \textbf{tous ses autres ports} : un seul domaine de collision
partagé. Le \textbf{switch} (couche 2) lit l'adresse MAC de destination et, grâce à sa \textbf{table MAC}, ne l'envoie
\textbf{que sur le port du destinataire} : un domaine de collision par port, c'est lui qui limite les collisions.
\emph{0,5 pt hub, 0,5 pt switch.}}

\question{A4}{1 pt}{Qu'est-ce que la \textbf{passerelle par défaut} d'un poste ? Quelle condition son adresse doit-elle respecter ?}
\zone{2.4cm}{Adresse IP de l'interface du \textbf{routeur} à laquelle le poste envoie les paquets destinés à un
\textbf{autre réseau}. Elle doit \textbf{appartenir au même réseau} que le poste. \emph{0,5 + 0,5 pt.}}

\question{A5}{0,5 pt}{Un portable configuré en DHCP affiche \ip{\apipa} / \ip{255.255.0.0}. Que s'est-il passé ?}
\zone{1.5cm}{Aucun \textbf{serveur DHCP} n'a répondu : le poste s'est attribué une adresse \textbf{APIPA} (169.254.0.0/16).}

\question{A6}{0,5 pt}{Quel protocole utilise \texttt{ping} ? Nommez les deux messages échangés et la couche OSI des adresses IP.}
\zone{1.5cm}{\textbf{ICMP} ; \textbf{Echo Request} / \textbf{Echo Reply} ; adresses IP en \textbf{couche 3 (Réseau)}.}

%==========================================================================
\newpage
\section{Partie B — Réalisation sous Packet Tracer \hfill /10}

\begin{cdc}
\entreprise{} vous confie son réseau : \textbf{trois réseaux} reliés par un \textbf{routeur unique Cisco 2911}
nommé \ip{\Rnom} (interfaces G0/0, G0/1, G0/2).
\begin{itemize}[leftmargin=12pt]
  \item \textbf{Réseau \resUn} : \textbf{4 postes} PC-\pUn1 à PC-\pUn4, commutateur \ip{SW-\pUn}, adressage statique,
        réseau \ip{\netUn.0/24}.
  \item \textbf{Réseau \resDeux} : \textbf{2 postes} PC-\pDeux1, PC-\pDeux2 et \textbf{1 imprimante} IMP-\pDeux,
        commutateur \ip{SW-\pDeux}, adressage statique, réseau de classe B \ip{\netDeux.0.0/16}.
  \item \textbf{Réseau Serveur} : serveur \ip{\srvNom} d'adresse \ip{\netSrv.\srvH/24}, relié \textbf{directement} au routeur.
  \item La passerelle de chaque réseau est sa \textbf{dernière adresse utilisable}. Tous les équipements communiquent
        entre eux et accèdent à la page web du serveur. Nommage imposé (« Display Name »).
\end{itemize}
\end{cdc}

\question{B1}{2 pts}{Complétez le plan d'adressage \textbf{avant} la réalisation.}
{\footnotesize\setlength{\tabcolsep}{3pt}\renewcommand{\arraystretch}{1.15}
\begin{tabularx}{\linewidth}{|l|C|C|C|C|C|C|}
  \hline
  \textbf{Réseau} & \textbf{Adresse réseau} & \textbf{Masque} & \textbf{1re @ utilisable} &
  \textbf{Dernière @ utilisable} & \textbf{Diffusion} & \textbf{Passerelle} \\ \hline
  \hl \resUn & \sol{\netUn.0} & \sol{255.255.255.0} & \sol{\netUn.1} & \sol{\netUn.254} & \sol{\netUn.255} & \sol{\netUn.254} \\ \hline
  \hl \resDeux & \sol{\netDeux.0.0} & \sol{255.255.0.0} & \sol{\netDeux.0.1} & \sol{\netDeux.255.254} & \sol{\netDeux.255.255} & \sol{\netDeux.255.254} \\ \hline
  \hl Serveur & \sol{\netSrv.0} & \sol{255.255.255.0} & \sol{\netSrv.1} & \sol{\netSrv.254} & \sol{\netSrv.255} & \sol{\netSrv.254} \\ \hline
\end{tabularx}}
\bareme{0,5 pt par ligne juste + 0,5 pt si la convention de passerelle est respectée partout.}

\question{B2}{2 pts}{Construisez la topologie (équipements, câblage, nommage) et relevez l'adressage configuré.}
{\footnotesize\setlength{\tabcolsep}{3pt}\renewcommand{\arraystretch}{1.1}
\begin{tabularx}{\linewidth}{|l|C|C|C||l|C|C|}
  \hline
  \textbf{Hôte} & \textbf{Adresse IP} & \textbf{Masque} & \textbf{Passerelle} & \textbf{Interface} & \textbf{Adresse IP} & \textbf{Masque} \\ \hline
  \hl PC-\pUn1 & \sol{\netUn.1} & \sol{255.255.255.0} & \sol{\netUn.254} & \Rnom{} G0/0 & \sol{\netUn.254} & \sol{255.255.255.0} \\ \hline
  \hl PC-\pDeux1 & \sol{\netDeux.0.1} & \sol{255.255.0.0} & \sol{\netDeux.255.254} & \Rnom{} G0/1 & \sol{\netDeux.255.254} & \sol{255.255.0.0} \\ \hline
  \hl IMP-\pDeux & \sol{\netDeux.0.10} & \sol{255.255.0.0} & \sol{\netDeux.255.254} & \Rnom{} G0/2 & \sol{\netSrv.254} & \sol{255.255.255.0} \\ \hline
  \hl \srvNom & \ip{\netSrv.\srvH} & \sol{255.255.255.0} & \sol{\netSrv.254} & \multicolumn{3}{c|}{\cellcolor{black!8}} \\ \hline
\end{tabularx}}
\bareme{Adresses d'hôtes libres (uniques, dans le bon réseau) ; affectation G0/x libre. Sur le .pkt : équipements 0,5 ;
câblage 1 (câbles droits vers les switchs et switch–routeur, \textbf{croisé} serveur–routeur, liens verts) ; nommage 0,5.}

\question{B3}{3 pts}{Configurez tous les équipements : IP, masques, passerelles des hôtes ; interfaces du routeur
(« Port Status : On »).}
\bareme{Sur le .pkt : interfaces routeur 1 pt ; IP/masques des hôtes 1 pt ; passerelles (imprimante et serveur compris)
1 pt. $-0,25$ par erreur, sans descendre sous 0 par item.}

\question{B4}{3 pts}{Réalisez les tests, notez la commande et le résultat, puis faites viser chaque test réussi.}
{\small\renewcommand{\arraystretch}{1.1}
\begin{tabularx}{\linewidth}{|c|l|l|C|c|c|}
  \hline
  \textbf{N°} & \textbf{Source} & \textbf{Destination} & \textbf{Commande / action} & \textbf{Résultat} & \textbf{Visa prof.} \\ \hline
  \hl 1 & PC-\pUn1 & PC-\pUn4 & \sol{ping \netUn.4} & \sol{réussi} & \hspace{1.4cm} \\ \hline
  \hl 2 & PC-\pDeux1 & IMP-\pDeux & \sol{ping \netDeux.0.10} & \sol{réussi} & \\ \hline
  \hl 3 & PC-\pUn2 & PC-\pDeux2 & \sol{ping \netDeux.0.2} & \sol{réussi} & \\ \hline
  \hl 4 & PC-\pDeux2 & \srvNom & \sol{ping \netSrv.\srvH} & \sol{réussi} & \\ \hline
  \hl 5 & PC-\pUn3 & page web du serveur & \sol{navigateur : {\NoAutoSpacing http://\netSrv.\srvH}} & \sol{page affichée} & \\ \hline
\end{tabularx}}
\bareme{0,5 pt par test visé + 0,5 pt pour B5. Le premier ping inter-réseaux peut échouer (résolution ARP) ; les suivants doivent réussir.}

\question{B5}{incl. B4}{Test n°3 : par quels équipements passe le paquet ? Pourquoi PC-\pUn2 l'envoie-t-il au routeur ?}
\zone{2.3cm}{PC-\pUn2 $\to$ SW-\pUn{} $\to$ \Rnom{} $\to$ SW-\pDeux{} $\to$ PC-\pDeux2. Grâce à son masque, PC-\pUn2 constate que
la destination (\netDeux.0.0/16) n'est \textbf{pas dans son réseau} (\netUn.0/24) : il envoie le paquet à sa
\textbf{passerelle par défaut}, qui le route.}

%==========================================================================
\newpage
\section{Partie C — Diagnostic d'un réseau en panne \hfill /4}

Un technicien constate qu'\textbf{aucun poste ne parvient à joindre le serveur SRV-C} et que certains postes
communiquent mal entre eux. Les câbles sont du bon type.

\begin{center}
\begin{minipage}{0.37\linewidth}\centering
\begin{tikzpicture}[font=\footnotesize,
  eq/.style={draw,fill=qbg,rounded corners=2pt,minimum width=9mm,minimum height=5mm,inner sep=2pt},
  every path/.style={thick}]
  \node[eq] (SW) {SW-1};
  \node[eq,right=9mm of SW] (R) {R1};
  \node[eq,right=7mm of R] (SRV) {SRV-C};
  \node[eq,above left=5mm and -3mm of SW] (A) {PC-A};
  \node[eq,above right=5mm and -3mm of SW] (B) {PC-B};
  \node[eq,below left=5mm and -3mm of SW] (C) {PC-C};
  \node[eq,below right=5mm and -3mm of SW] (D) {PC-D};
  \foreach \n in {A,B,C,D} \draw (\n) -- (SW);
  \draw (SW) -- node[above,font=\tiny]{G0/0} (R);
  \draw (R) -- node[above,font=\tiny]{G0/1} (SRV);
\end{tikzpicture}
\end{minipage}\hfill
\begin{minipage}{0.61\linewidth}\centering
{\small\renewcommand{\arraystretch}{1.05}\setlength{\tabcolsep}{4pt}
\begin{tabular}{lllll}
  \toprule
  \textbf{Équipement} & \textbf{Adresse IP} & \textbf{Masque} & \textbf{Passerelle} & \textbf{Port} \\
  \midrule
  \Crows
  \bottomrule
\end{tabular}}
\end{minipage}
\end{center}

\question{C1}{4 pts}{Relevez les \textbf{quatre anomalies}. Pour chacune : conséquence et correction.}
{\small\renewcommand{\arraystretch}{1.1}
\ifcorrige\def\hc{}\else\def\hc{\rule[-10mm]{0pt}{14mm}}\fi
\begin{tabularx}{\linewidth}{|c|>{\raggedright\arraybackslash}p{2cm}|L|L|L|}
  \hline
  \textbf{N°} & \textbf{Équipement} & \textbf{Anomalie} & \textbf{Conséquence} & \textbf{Correction} \\ \hline
  \ifcorrige\Csol\else
  \hc 1 & & & & \\ \hline
  \hc 2 & & & & \\ \hline
  \hc 3 & & & & \\ \hline
  \hc 4 & & & & \\ \hline
  \fi
\end{tabularx}}
\bareme{1 pt par anomalie (0,25 équipement, anomalie, conséquence, correction ; ordre indifférent). \Cbareme}

%==========================================================================
\section{Bonus — Extension sans fil \hfill +1}

\begin{cdc}Ajoutez au réseau \resUn{} un \textbf{AccessPoint-PT} relié à SW-\pUn{} (SSID \ip{\ssid}, \textbf{WPA2-PSK})
et un portable \ip{Laptop-\pUn} équipé d'un module \texttt{WPC300N}, en adressage \textbf{statique}.
Vérifiez qu'il joint \srvNom.\end{cdc}
{\small\renewcommand{\arraystretch}{1.1}
\begin{tabularx}{\linewidth}{|l|C|C|C|c|}
  \hline
  \textbf{Équipement} & \textbf{Adresse IP} & \textbf{Masque} & \textbf{Passerelle} & \textbf{Visa prof.} \\ \hline
  \hl Laptop-\pUn & \sol{\netUn.50} & \sol{255.255.255.0} & \sol{\netUn.254} & \hspace{1.4cm} \\ \hline
\end{tabularx}}
\bareme{Le point d'accès fait office de pont : le portable est dans \netUn.0/24. Bonus si associé en WPA2-PSK et ping
vers \netSrv.\srvH{} réussi. Note plafonnée à 20.}

\end{document}
"
%   pdflatex -jobname=Eval_Reseau_SujetB   "\def\variante{B}\def\modecorrige{0}% Évaluation pratique — Réseaux locaux (BTS CIEL 1re année) — sujets A et B
% Compilation (4 PDF) :
%   pdflatex -jobname=Eval_Reseau_SujetA   "\def\variante{A}\def\modecorrige{0}\input{Eval_TP_Reseau_LAN.tex}"
%   pdflatex -jobname=Eval_Reseau_SujetB   "\def\variante{B}\def\modecorrige{0}\input{Eval_TP_Reseau_LAN.tex}"
%   pdflatex -jobname=Eval_Reseau_CorrigeA "\def\variante{A}\def\modecorrige{1}\input{Eval_TP_Reseau_LAN.tex}"
%   pdflatex -jobname=Eval_Reseau_CorrigeB "\def\variante{B}\def\modecorrige{1}\input{Eval_TP_Reseau_LAN.tex}"
\documentclass[10pt,a4paper]{article}

\usepackage[T1]{fontenc}
\usepackage[utf8]{inputenc}
\usepackage{lmodern}
\usepackage[french]{babel}
\usepackage[a4paper,hmargin=1.5cm,top=1.3cm,bottom=1.6cm,footskip=0.7cm]{geometry}
\usepackage[table]{xcolor}
\usepackage{booktabs,tabularx,array}
\usepackage{enumitem}
\usepackage{fancyhdr}
\usepackage[most]{tcolorbox}
\usepackage{tikz}
\usetikzlibrary{positioning}
\usepackage{titlesec}

\providecommand{\modecorrige}{0}
\providecommand{\variante}{A}
\newif\ifcorrige
\ifnum\modecorrige=1 \corrigetrue \else \corrigefalse \fi

\definecolor{ipred}{HTML}{C0392B}
\definecolor{cdcbar}{HTML}{A67C3D}
\definecolor{cdcbg}{HTML}{F4F1EA}
\definecolor{qbar}{HTML}{4A6B5D}
\definecolor{qbg}{HTML}{EDF1EF}
\definecolor{corr}{HTML}{1F5FA8}
\definecolor{corrbg}{HTML}{E8F0FA}
\definecolor{grisfoot}{HTML}{828282}

\newcommand{\ip}[1]{\textcolor{ipred}{\texttt{#1}}}
\newcommand{\pts}[1]{\hfill\mbox{\textbf{(#1)}}}
\newcommand{\sol}[1]{\ifcorrige\textcolor{corr}{#1}\fi}
\newcommand{\bareme}[1]{\ifcorrige\par{\footnotesize\color{corr}#1\par}\fi}
\newcommand{\hl}{\ifcorrige\rule[-1.6mm]{0pt}{5.4mm}\else\rule[-2.2mm]{0pt}{6.6mm}\fi}
\newcolumntype{C}{>{\centering\arraybackslash}X}
\newcolumntype{L}{>{\raggedright\arraybackslash}X}

\titleformat{\section}{\large\bfseries}{}{0pt}{}
\titlespacing*{\section}{0pt}{6pt}{3pt}

\tcbset{qstyle/.style={enhanced,boxrule=0pt,leftrule=2.5pt,arc=0pt,outer arc=0pt,
  left=5pt,right=5pt,top=2pt,bottom=2pt,before skip=4pt,after skip=3pt}}
\newtcolorbox{cdc}{qstyle,colback=cdcbg,colframe=cdcbar}
\newtcolorbox{quest}{qstyle,colback=qbg,colframe=qbar}

% Zone de réponse : vierge dans le sujet, remplie dans le corrigé
\newcommand{\zone}[2]{%
  \ifcorrige
    \begin{tcolorbox}[colback=corrbg,colframe=corr,boxrule=0.5pt,arc=1pt,left=4pt,right=4pt,top=1pt,bottom=1pt,
      before skip=2pt,after skip=2pt,fontupper=\footnotesize\color{corr}]#2\end{tcolorbox}%
  \else
    \begin{tcolorbox}[colback=white,colframe=black!45,boxrule=0.4pt,arc=1pt,height=#1,
      before skip=2pt,after skip=2pt]\end{tcolorbox}%
  \fi}

\newcommand{\question}[3]{\begin{quest}\textbf{#1.} #3\pts{#2}\end{quest}}

%==========================================================================
% Données propres à chaque variante
%==========================================================================
\if\variante A
  \def\entreprise{Le cabinet d'architectes \textbf{ArchiNova}}
  \def\Rnom{R-ArchiNova}
  \def\resUn{Bureau d'études}\def\pUn{BE}\def\netUn{192.168.40}
  \def\resDeux{Accueil}\def\pDeux{ACC}\def\netDeux{172.16}
  \def\srvNom{SRV-Intranet}\def\netSrv{192.168.200}\def\srvH{10}
  \def\ssid{ArchiNova-WiFi}
  \def\AIrows{%
    \hl\ip{192.168.40.25} & \sol{C} & \sol{255.255.255.0} & \sol{192.168.40.0} & \sol{192.168.40.255} & \sol{254} \\ \hline
    \hl\ip{172.16.5.10}   & \sol{B} & \sol{255.255.0.0}   & \sol{172.16.0.0}   & \sol{172.16.255.255} & \sol{65\,534} \\ \hline
    \hl\ip{10.3.2.1}      & \sol{A} & \sol{255.0.0.0}     & \sol{10.0.0.0}     & \sol{10.255.255.255} & \sol{16\,777\,214} \\ \hline}
  \def\apipa{169.254.12.7}
  \def\Cnet{192.168.60}
  \def\Crows{%
    PC-A & \ip{192.168.60.10} & \ip{255.255.255.0} & \ip{192.168.60.254} & — \\
    PC-B & \ip{192.168.60.20} & \ip{255.255.255.0} & \ip{192.168.60.1}   & — \\
    PC-C & \ip{192.168.61.30} & \ip{255.255.255.0} & \ip{192.168.60.254} & — \\
    PC-D & \ip{192.168.60.10} & \ip{255.255.255.0} & \ip{192.168.60.254} & — \\
    SRV-C & \ip{10.0.0.5}     & \ip{255.0.0.0}     & \ip{10.0.0.254}     & — \\
    R1 — G0/0 & \ip{192.168.60.254} & \ip{255.255.255.0} & — & On \\
    R1 — G0/1 & \ip{10.0.0.254}     & \ip{255.0.0.0}     & — & Off \\}
  \def\Csol{%
    \hl 1 & \sol{PC-B} & \sol{Passerelle .60.1 : pas l'adresse du routeur.} & \sol{Joint son réseau, mais aucun autre (serveur).} & \sol{Passerelle 192.168.60.254.} \\ \hline
    \hl 2 & \sol{PC-C} & \sol{IP hors du réseau 192.168.60.0/24.} & \sol{Ne joint ni les postes ni le routeur.} & \sol{Ex. 192.168.60.30.} \\ \hline
    \hl 3 & \sol{PC-D} & \sol{Même IP que PC-A (doublon).} & \sol{Conflit d'adresse : PC-A et PC-D perturbés.} & \sol{Adresse unique, ex. .60.40.} \\ \hline
    \hl 4 & \sol{R1 G0/1} & \sol{Interface désactivée (Off).} & \sol{Réseau du serveur injoignable pour tous.} & \sol{Port Status : On.} \\ \hline}
  \def\Cbareme{PC-A et SRV-C sont correctement configurés ; accepter PC-A cité à la place de PC-D pour le doublon.}
\else
  \def\entreprise{La clinique vétérinaire \textbf{VetoPlus}}
  \def\Rnom{R-VetoPlus}
  \def\resUn{Soins}\def\pUn{SOI}\def\netUn{192.168.70}
  \def\resDeux{Secrétariat}\def\pDeux{SEC}\def\netDeux{172.20}
  \def\srvNom{SRV-Dossiers}\def\netSrv{192.168.250}\def\srvH{20}
  \def\ssid{VetoPlus-WiFi}
  \def\AIrows{%
    \hl\ip{192.168.15.130} & \sol{C} & \sol{255.255.255.0} & \sol{192.168.15.0} & \sol{192.168.15.255} & \sol{254} \\ \hline
    \hl\ip{172.31.200.7}   & \sol{B} & \sol{255.255.0.0}   & \sol{172.31.0.0}   & \sol{172.31.255.255} & \sol{65\,534} \\ \hline
    \hl\ip{10.45.6.9}      & \sol{A} & \sol{255.0.0.0}     & \sol{10.0.0.0}     & \sol{10.255.255.255} & \sol{16\,777\,214} \\ \hline}
  \def\apipa{169.254.88.3}
  \def\Cnet{192.168.80}
  \def\Crows{%
    PC-A & \ip{192.168.8.11}   & \ip{255.255.255.0} & \ip{192.168.80.254} & — \\
    PC-B & \ip{192.168.80.12}  & \ip{255.255.255.0} & \ip{192.168.80.254} & — \\
    PC-C & \ip{192.168.80.13}  & \ip{255.255.255.0} & \ip{192.168.80.100} & — \\
    PC-D & \ip{192.168.80.12}  & \ip{255.255.255.0} & \ip{192.168.80.254} & — \\
    SRV-C & \ip{10.0.0.8}      & \ip{255.0.0.0}     & \ip{10.0.0.254}     & — \\
    R1 — G0/0 & \ip{192.168.80.254} & \ip{255.255.255.0} & — & Off \\
    R1 — G0/1 & \ip{10.0.0.254}     & \ip{255.0.0.0}     & — & On \\}
  \def\Csol{%
    \hl 1 & \sol{PC-A} & \sol{IP hors du réseau 192.168.80.0/24.} & \sol{Ne joint ni les postes ni le routeur.} & \sol{Ex. 192.168.80.11.} \\ \hline
    \hl 2 & \sol{PC-C} & \sol{Passerelle .80.100 : pas l'adresse du routeur.} & \sol{Joint son réseau, mais aucun autre (serveur).} & \sol{Passerelle 192.168.80.254.} \\ \hline
    \hl 3 & \sol{PC-D} & \sol{Même IP que PC-B (doublon).} & \sol{Conflit d'adresse : PC-B et PC-D perturbés.} & \sol{Adresse unique, ex. .80.14.} \\ \hline
    \hl 4 & \sol{R1 G0/0} & \sol{Interface côté LAN désactivée (Off).} & \sol{Aucun poste ne joint le routeur ni le serveur.} & \sol{Port Status : On.} \\ \hline}
  \def\Cbareme{PC-B et SRV-C sont correctement configurés ; accepter PC-B cité à la place de PC-D pour le doublon.}
\fi

\pagestyle{fancy}
\fancyhf{}
\renewcommand{\headrulewidth}{0pt}
\fancyfoot[L]{\footnotesize\itshape\color{grisfoot}Évaluation — Réseaux locaux — Sujet \variante\ifcorrige\ — CORRIGÉ\fi}
\fancyfoot[R]{\footnotesize\color{grisfoot}\thepage/3}

\setlist{itemsep=0pt,topsep=2pt,parsep=0pt}
\setlength{\parindent}{0pt}
\setlength{\parskip}{2pt}

\begin{document}

\begin{center}
  {\Large\bfseries Évaluation pratique — Réseaux locaux \quad
   \fcolorbox{black}{white}{\ SUJET \variante\ }\par}
  \vspace{1pt}
  {\bfseries Adressage IP, commutation et routage — BTS CIEL 1re année — durée : 2 heures\par}
  \ifcorrige\vspace{2pt}{\bfseries\color{corr}CORRIGÉ ET BARÈME — document enseignant}\par\fi
\end{center}
\vspace{-2pt}
{\renewcommand{\arraystretch}{1.35}
\begin{tabularx}{\linewidth}{|X|X|p{2.2cm}|p{2.6cm}|}
  \hline
  \textbf{Nom :} & \textbf{Prénom :} & \textbf{Poste n°} & \textbf{Note :} \hfill /20 \\ \hline
\end{tabularx}}

\textbf{Consignes :} travail \textbf{individuel}, documents et Internet \textbf{non autorisés}. Ce sujet est le
\textbf{document-réponse}. Fichier Packet Tracer : \ip{Nom\_Prenom\_EvalLAN\_\variante.pkt}, à enregistrer
régulièrement dans le dossier de rendu. Les « visas » sont apposés par le professeur, que vous appelez à chaque test réussi.
Durées conseillées : A 30 min — B 1 h 10 — C 20 min.

%==========================================================================
\section{Partie A — Questions de cours et calculs \hfill /6}

\question{A1}{2 pts}{Complétez le tableau en considérant le \textbf{masque par défaut} de la classe de chaque adresse.}
{\small\renewcommand{\arraystretch}{1.15}
\begin{tabularx}{\linewidth}{|l|c|C|C|C|c|}
  \hline
  \textbf{Adresse IP} & \textbf{Classe} & \textbf{Masque par défaut} & \textbf{Adresse réseau} &
  \textbf{Adresse de diffusion} & \textbf{Nb d'hôtes} \\ \hline
  \AIrows
\end{tabularx}}
\bareme{0,5 pt par ligne juste + 0,5 pt si les trois nombres d'hôtes sont justes ($2^{n}-2$, $n$ = nombre de bits d'hôte).}

\question{A2}{1 pt}{Indiquez le type de câble (\textbf{droit} ou \textbf{croisé}) pour chaque liaison.}
{\small\renewcommand{\arraystretch}{1.15}
\begin{tabularx}{\linewidth}{|X|c||X|c|}
  \hline
  \hl PC $\leftrightarrow$ commutateur & \makebox[1.6cm]{\sol{droit}} & Commutateur $\leftrightarrow$ routeur & \makebox[1.6cm]{\sol{droit}} \\ \hline
  \hl PC $\leftrightarrow$ PC (réseau P2P) & \sol{croisé} & Concentrateur $\leftrightarrow$ concentrateur & \sol{croisé} \\ \hline
\end{tabularx}}
\bareme{0,25 pt par case. Équipements de types différents : câble droit ; de même type : câble croisé.}

\question{A3}{1 pt}{Comment un \textbf{concentrateur (hub)} et un \textbf{commutateur (switch)} traitent-ils une trame
reçue ? Lequel limite les collisions, et pourquoi ?}
\zone{2.9cm}{Le \textbf{hub} (couche 1) répète la trame sur \textbf{tous ses autres ports} : un seul domaine de collision
partagé. Le \textbf{switch} (couche 2) lit l'adresse MAC de destination et, grâce à sa \textbf{table MAC}, ne l'envoie
\textbf{que sur le port du destinataire} : un domaine de collision par port, c'est lui qui limite les collisions.
\emph{0,5 pt hub, 0,5 pt switch.}}

\question{A4}{1 pt}{Qu'est-ce que la \textbf{passerelle par défaut} d'un poste ? Quelle condition son adresse doit-elle respecter ?}
\zone{2.4cm}{Adresse IP de l'interface du \textbf{routeur} à laquelle le poste envoie les paquets destinés à un
\textbf{autre réseau}. Elle doit \textbf{appartenir au même réseau} que le poste. \emph{0,5 + 0,5 pt.}}

\question{A5}{0,5 pt}{Un portable configuré en DHCP affiche \ip{\apipa} / \ip{255.255.0.0}. Que s'est-il passé ?}
\zone{1.5cm}{Aucun \textbf{serveur DHCP} n'a répondu : le poste s'est attribué une adresse \textbf{APIPA} (169.254.0.0/16).}

\question{A6}{0,5 pt}{Quel protocole utilise \texttt{ping} ? Nommez les deux messages échangés et la couche OSI des adresses IP.}
\zone{1.5cm}{\textbf{ICMP} ; \textbf{Echo Request} / \textbf{Echo Reply} ; adresses IP en \textbf{couche 3 (Réseau)}.}

%==========================================================================
\newpage
\section{Partie B — Réalisation sous Packet Tracer \hfill /10}

\begin{cdc}
\entreprise{} vous confie son réseau : \textbf{trois réseaux} reliés par un \textbf{routeur unique Cisco 2911}
nommé \ip{\Rnom} (interfaces G0/0, G0/1, G0/2).
\begin{itemize}[leftmargin=12pt]
  \item \textbf{Réseau \resUn} : \textbf{4 postes} PC-\pUn1 à PC-\pUn4, commutateur \ip{SW-\pUn}, adressage statique,
        réseau \ip{\netUn.0/24}.
  \item \textbf{Réseau \resDeux} : \textbf{2 postes} PC-\pDeux1, PC-\pDeux2 et \textbf{1 imprimante} IMP-\pDeux,
        commutateur \ip{SW-\pDeux}, adressage statique, réseau de classe B \ip{\netDeux.0.0/16}.
  \item \textbf{Réseau Serveur} : serveur \ip{\srvNom} d'adresse \ip{\netSrv.\srvH/24}, relié \textbf{directement} au routeur.
  \item La passerelle de chaque réseau est sa \textbf{dernière adresse utilisable}. Tous les équipements communiquent
        entre eux et accèdent à la page web du serveur. Nommage imposé (« Display Name »).
\end{itemize}
\end{cdc}

\question{B1}{2 pts}{Complétez le plan d'adressage \textbf{avant} la réalisation.}
{\footnotesize\setlength{\tabcolsep}{3pt}\renewcommand{\arraystretch}{1.15}
\begin{tabularx}{\linewidth}{|l|C|C|C|C|C|C|}
  \hline
  \textbf{Réseau} & \textbf{Adresse réseau} & \textbf{Masque} & \textbf{1re @ utilisable} &
  \textbf{Dernière @ utilisable} & \textbf{Diffusion} & \textbf{Passerelle} \\ \hline
  \hl \resUn & \sol{\netUn.0} & \sol{255.255.255.0} & \sol{\netUn.1} & \sol{\netUn.254} & \sol{\netUn.255} & \sol{\netUn.254} \\ \hline
  \hl \resDeux & \sol{\netDeux.0.0} & \sol{255.255.0.0} & \sol{\netDeux.0.1} & \sol{\netDeux.255.254} & \sol{\netDeux.255.255} & \sol{\netDeux.255.254} \\ \hline
  \hl Serveur & \sol{\netSrv.0} & \sol{255.255.255.0} & \sol{\netSrv.1} & \sol{\netSrv.254} & \sol{\netSrv.255} & \sol{\netSrv.254} \\ \hline
\end{tabularx}}
\bareme{0,5 pt par ligne juste + 0,5 pt si la convention de passerelle est respectée partout.}

\question{B2}{2 pts}{Construisez la topologie (équipements, câblage, nommage) et relevez l'adressage configuré.}
{\footnotesize\setlength{\tabcolsep}{3pt}\renewcommand{\arraystretch}{1.1}
\begin{tabularx}{\linewidth}{|l|C|C|C||l|C|C|}
  \hline
  \textbf{Hôte} & \textbf{Adresse IP} & \textbf{Masque} & \textbf{Passerelle} & \textbf{Interface} & \textbf{Adresse IP} & \textbf{Masque} \\ \hline
  \hl PC-\pUn1 & \sol{\netUn.1} & \sol{255.255.255.0} & \sol{\netUn.254} & \Rnom{} G0/0 & \sol{\netUn.254} & \sol{255.255.255.0} \\ \hline
  \hl PC-\pDeux1 & \sol{\netDeux.0.1} & \sol{255.255.0.0} & \sol{\netDeux.255.254} & \Rnom{} G0/1 & \sol{\netDeux.255.254} & \sol{255.255.0.0} \\ \hline
  \hl IMP-\pDeux & \sol{\netDeux.0.10} & \sol{255.255.0.0} & \sol{\netDeux.255.254} & \Rnom{} G0/2 & \sol{\netSrv.254} & \sol{255.255.255.0} \\ \hline
  \hl \srvNom & \ip{\netSrv.\srvH} & \sol{255.255.255.0} & \sol{\netSrv.254} & \multicolumn{3}{c|}{\cellcolor{black!8}} \\ \hline
\end{tabularx}}
\bareme{Adresses d'hôtes libres (uniques, dans le bon réseau) ; affectation G0/x libre. Sur le .pkt : équipements 0,5 ;
câblage 1 (câbles droits vers les switchs et switch–routeur, \textbf{croisé} serveur–routeur, liens verts) ; nommage 0,5.}

\question{B3}{3 pts}{Configurez tous les équipements : IP, masques, passerelles des hôtes ; interfaces du routeur
(« Port Status : On »).}
\bareme{Sur le .pkt : interfaces routeur 1 pt ; IP/masques des hôtes 1 pt ; passerelles (imprimante et serveur compris)
1 pt. $-0,25$ par erreur, sans descendre sous 0 par item.}

\question{B4}{3 pts}{Réalisez les tests, notez la commande et le résultat, puis faites viser chaque test réussi.}
{\small\renewcommand{\arraystretch}{1.1}
\begin{tabularx}{\linewidth}{|c|l|l|C|c|c|}
  \hline
  \textbf{N°} & \textbf{Source} & \textbf{Destination} & \textbf{Commande / action} & \textbf{Résultat} & \textbf{Visa prof.} \\ \hline
  \hl 1 & PC-\pUn1 & PC-\pUn4 & \sol{ping \netUn.4} & \sol{réussi} & \hspace{1.4cm} \\ \hline
  \hl 2 & PC-\pDeux1 & IMP-\pDeux & \sol{ping \netDeux.0.10} & \sol{réussi} & \\ \hline
  \hl 3 & PC-\pUn2 & PC-\pDeux2 & \sol{ping \netDeux.0.2} & \sol{réussi} & \\ \hline
  \hl 4 & PC-\pDeux2 & \srvNom & \sol{ping \netSrv.\srvH} & \sol{réussi} & \\ \hline
  \hl 5 & PC-\pUn3 & page web du serveur & \sol{navigateur : {\NoAutoSpacing http://\netSrv.\srvH}} & \sol{page affichée} & \\ \hline
\end{tabularx}}
\bareme{0,5 pt par test visé + 0,5 pt pour B5. Le premier ping inter-réseaux peut échouer (résolution ARP) ; les suivants doivent réussir.}

\question{B5}{incl. B4}{Test n°3 : par quels équipements passe le paquet ? Pourquoi PC-\pUn2 l'envoie-t-il au routeur ?}
\zone{2.3cm}{PC-\pUn2 $\to$ SW-\pUn{} $\to$ \Rnom{} $\to$ SW-\pDeux{} $\to$ PC-\pDeux2. Grâce à son masque, PC-\pUn2 constate que
la destination (\netDeux.0.0/16) n'est \textbf{pas dans son réseau} (\netUn.0/24) : il envoie le paquet à sa
\textbf{passerelle par défaut}, qui le route.}

%==========================================================================
\newpage
\section{Partie C — Diagnostic d'un réseau en panne \hfill /4}

Un technicien constate qu'\textbf{aucun poste ne parvient à joindre le serveur SRV-C} et que certains postes
communiquent mal entre eux. Les câbles sont du bon type.

\begin{center}
\begin{minipage}{0.37\linewidth}\centering
\begin{tikzpicture}[font=\footnotesize,
  eq/.style={draw,fill=qbg,rounded corners=2pt,minimum width=9mm,minimum height=5mm,inner sep=2pt},
  every path/.style={thick}]
  \node[eq] (SW) {SW-1};
  \node[eq,right=9mm of SW] (R) {R1};
  \node[eq,right=7mm of R] (SRV) {SRV-C};
  \node[eq,above left=5mm and -3mm of SW] (A) {PC-A};
  \node[eq,above right=5mm and -3mm of SW] (B) {PC-B};
  \node[eq,below left=5mm and -3mm of SW] (C) {PC-C};
  \node[eq,below right=5mm and -3mm of SW] (D) {PC-D};
  \foreach \n in {A,B,C,D} \draw (\n) -- (SW);
  \draw (SW) -- node[above,font=\tiny]{G0/0} (R);
  \draw (R) -- node[above,font=\tiny]{G0/1} (SRV);
\end{tikzpicture}
\end{minipage}\hfill
\begin{minipage}{0.61\linewidth}\centering
{\small\renewcommand{\arraystretch}{1.05}\setlength{\tabcolsep}{4pt}
\begin{tabular}{lllll}
  \toprule
  \textbf{Équipement} & \textbf{Adresse IP} & \textbf{Masque} & \textbf{Passerelle} & \textbf{Port} \\
  \midrule
  \Crows
  \bottomrule
\end{tabular}}
\end{minipage}
\end{center}

\question{C1}{4 pts}{Relevez les \textbf{quatre anomalies}. Pour chacune : conséquence et correction.}
{\small\renewcommand{\arraystretch}{1.1}
\ifcorrige\def\hc{}\else\def\hc{\rule[-10mm]{0pt}{14mm}}\fi
\begin{tabularx}{\linewidth}{|c|>{\raggedright\arraybackslash}p{2cm}|L|L|L|}
  \hline
  \textbf{N°} & \textbf{Équipement} & \textbf{Anomalie} & \textbf{Conséquence} & \textbf{Correction} \\ \hline
  \ifcorrige\Csol\else
  \hc 1 & & & & \\ \hline
  \hc 2 & & & & \\ \hline
  \hc 3 & & & & \\ \hline
  \hc 4 & & & & \\ \hline
  \fi
\end{tabularx}}
\bareme{1 pt par anomalie (0,25 équipement, anomalie, conséquence, correction ; ordre indifférent). \Cbareme}

%==========================================================================
\section{Bonus — Extension sans fil \hfill +1}

\begin{cdc}Ajoutez au réseau \resUn{} un \textbf{AccessPoint-PT} relié à SW-\pUn{} (SSID \ip{\ssid}, \textbf{WPA2-PSK})
et un portable \ip{Laptop-\pUn} équipé d'un module \texttt{WPC300N}, en adressage \textbf{statique}.
Vérifiez qu'il joint \srvNom.\end{cdc}
{\small\renewcommand{\arraystretch}{1.1}
\begin{tabularx}{\linewidth}{|l|C|C|C|c|}
  \hline
  \textbf{Équipement} & \textbf{Adresse IP} & \textbf{Masque} & \textbf{Passerelle} & \textbf{Visa prof.} \\ \hline
  \hl Laptop-\pUn & \sol{\netUn.50} & \sol{255.255.255.0} & \sol{\netUn.254} & \hspace{1.4cm} \\ \hline
\end{tabularx}}
\bareme{Le point d'accès fait office de pont : le portable est dans \netUn.0/24. Bonus si associé en WPA2-PSK et ping
vers \netSrv.\srvH{} réussi. Note plafonnée à 20.}

\end{document}
"
%   pdflatex -jobname=Eval_Reseau_CorrigeA "\def\variante{A}\def\modecorrige{1}% Évaluation pratique — Réseaux locaux (BTS CIEL 1re année) — sujets A et B
% Compilation (4 PDF) :
%   pdflatex -jobname=Eval_Reseau_SujetA   "\def\variante{A}\def\modecorrige{0}\input{Eval_TP_Reseau_LAN.tex}"
%   pdflatex -jobname=Eval_Reseau_SujetB   "\def\variante{B}\def\modecorrige{0}\input{Eval_TP_Reseau_LAN.tex}"
%   pdflatex -jobname=Eval_Reseau_CorrigeA "\def\variante{A}\def\modecorrige{1}\input{Eval_TP_Reseau_LAN.tex}"
%   pdflatex -jobname=Eval_Reseau_CorrigeB "\def\variante{B}\def\modecorrige{1}\input{Eval_TP_Reseau_LAN.tex}"
\documentclass[10pt,a4paper]{article}

\usepackage[T1]{fontenc}
\usepackage[utf8]{inputenc}
\usepackage{lmodern}
\usepackage[french]{babel}
\usepackage[a4paper,hmargin=1.5cm,top=1.3cm,bottom=1.6cm,footskip=0.7cm]{geometry}
\usepackage[table]{xcolor}
\usepackage{booktabs,tabularx,array}
\usepackage{enumitem}
\usepackage{fancyhdr}
\usepackage[most]{tcolorbox}
\usepackage{tikz}
\usetikzlibrary{positioning}
\usepackage{titlesec}

\providecommand{\modecorrige}{0}
\providecommand{\variante}{A}
\newif\ifcorrige
\ifnum\modecorrige=1 \corrigetrue \else \corrigefalse \fi

\definecolor{ipred}{HTML}{C0392B}
\definecolor{cdcbar}{HTML}{A67C3D}
\definecolor{cdcbg}{HTML}{F4F1EA}
\definecolor{qbar}{HTML}{4A6B5D}
\definecolor{qbg}{HTML}{EDF1EF}
\definecolor{corr}{HTML}{1F5FA8}
\definecolor{corrbg}{HTML}{E8F0FA}
\definecolor{grisfoot}{HTML}{828282}

\newcommand{\ip}[1]{\textcolor{ipred}{\texttt{#1}}}
\newcommand{\pts}[1]{\hfill\mbox{\textbf{(#1)}}}
\newcommand{\sol}[1]{\ifcorrige\textcolor{corr}{#1}\fi}
\newcommand{\bareme}[1]{\ifcorrige\par{\footnotesize\color{corr}#1\par}\fi}
\newcommand{\hl}{\ifcorrige\rule[-1.6mm]{0pt}{5.4mm}\else\rule[-2.2mm]{0pt}{6.6mm}\fi}
\newcolumntype{C}{>{\centering\arraybackslash}X}
\newcolumntype{L}{>{\raggedright\arraybackslash}X}

\titleformat{\section}{\large\bfseries}{}{0pt}{}
\titlespacing*{\section}{0pt}{6pt}{3pt}

\tcbset{qstyle/.style={enhanced,boxrule=0pt,leftrule=2.5pt,arc=0pt,outer arc=0pt,
  left=5pt,right=5pt,top=2pt,bottom=2pt,before skip=4pt,after skip=3pt}}
\newtcolorbox{cdc}{qstyle,colback=cdcbg,colframe=cdcbar}
\newtcolorbox{quest}{qstyle,colback=qbg,colframe=qbar}

% Zone de réponse : vierge dans le sujet, remplie dans le corrigé
\newcommand{\zone}[2]{%
  \ifcorrige
    \begin{tcolorbox}[colback=corrbg,colframe=corr,boxrule=0.5pt,arc=1pt,left=4pt,right=4pt,top=1pt,bottom=1pt,
      before skip=2pt,after skip=2pt,fontupper=\footnotesize\color{corr}]#2\end{tcolorbox}%
  \else
    \begin{tcolorbox}[colback=white,colframe=black!45,boxrule=0.4pt,arc=1pt,height=#1,
      before skip=2pt,after skip=2pt]\end{tcolorbox}%
  \fi}

\newcommand{\question}[3]{\begin{quest}\textbf{#1.} #3\pts{#2}\end{quest}}

%==========================================================================
% Données propres à chaque variante
%==========================================================================
\if\variante A
  \def\entreprise{Le cabinet d'architectes \textbf{ArchiNova}}
  \def\Rnom{R-ArchiNova}
  \def\resUn{Bureau d'études}\def\pUn{BE}\def\netUn{192.168.40}
  \def\resDeux{Accueil}\def\pDeux{ACC}\def\netDeux{172.16}
  \def\srvNom{SRV-Intranet}\def\netSrv{192.168.200}\def\srvH{10}
  \def\ssid{ArchiNova-WiFi}
  \def\AIrows{%
    \hl\ip{192.168.40.25} & \sol{C} & \sol{255.255.255.0} & \sol{192.168.40.0} & \sol{192.168.40.255} & \sol{254} \\ \hline
    \hl\ip{172.16.5.10}   & \sol{B} & \sol{255.255.0.0}   & \sol{172.16.0.0}   & \sol{172.16.255.255} & \sol{65\,534} \\ \hline
    \hl\ip{10.3.2.1}      & \sol{A} & \sol{255.0.0.0}     & \sol{10.0.0.0}     & \sol{10.255.255.255} & \sol{16\,777\,214} \\ \hline}
  \def\apipa{169.254.12.7}
  \def\Cnet{192.168.60}
  \def\Crows{%
    PC-A & \ip{192.168.60.10} & \ip{255.255.255.0} & \ip{192.168.60.254} & — \\
    PC-B & \ip{192.168.60.20} & \ip{255.255.255.0} & \ip{192.168.60.1}   & — \\
    PC-C & \ip{192.168.61.30} & \ip{255.255.255.0} & \ip{192.168.60.254} & — \\
    PC-D & \ip{192.168.60.10} & \ip{255.255.255.0} & \ip{192.168.60.254} & — \\
    SRV-C & \ip{10.0.0.5}     & \ip{255.0.0.0}     & \ip{10.0.0.254}     & — \\
    R1 — G0/0 & \ip{192.168.60.254} & \ip{255.255.255.0} & — & On \\
    R1 — G0/1 & \ip{10.0.0.254}     & \ip{255.0.0.0}     & — & Off \\}
  \def\Csol{%
    \hl 1 & \sol{PC-B} & \sol{Passerelle .60.1 : pas l'adresse du routeur.} & \sol{Joint son réseau, mais aucun autre (serveur).} & \sol{Passerelle 192.168.60.254.} \\ \hline
    \hl 2 & \sol{PC-C} & \sol{IP hors du réseau 192.168.60.0/24.} & \sol{Ne joint ni les postes ni le routeur.} & \sol{Ex. 192.168.60.30.} \\ \hline
    \hl 3 & \sol{PC-D} & \sol{Même IP que PC-A (doublon).} & \sol{Conflit d'adresse : PC-A et PC-D perturbés.} & \sol{Adresse unique, ex. .60.40.} \\ \hline
    \hl 4 & \sol{R1 G0/1} & \sol{Interface désactivée (Off).} & \sol{Réseau du serveur injoignable pour tous.} & \sol{Port Status : On.} \\ \hline}
  \def\Cbareme{PC-A et SRV-C sont correctement configurés ; accepter PC-A cité à la place de PC-D pour le doublon.}
\else
  \def\entreprise{La clinique vétérinaire \textbf{VetoPlus}}
  \def\Rnom{R-VetoPlus}
  \def\resUn{Soins}\def\pUn{SOI}\def\netUn{192.168.70}
  \def\resDeux{Secrétariat}\def\pDeux{SEC}\def\netDeux{172.20}
  \def\srvNom{SRV-Dossiers}\def\netSrv{192.168.250}\def\srvH{20}
  \def\ssid{VetoPlus-WiFi}
  \def\AIrows{%
    \hl\ip{192.168.15.130} & \sol{C} & \sol{255.255.255.0} & \sol{192.168.15.0} & \sol{192.168.15.255} & \sol{254} \\ \hline
    \hl\ip{172.31.200.7}   & \sol{B} & \sol{255.255.0.0}   & \sol{172.31.0.0}   & \sol{172.31.255.255} & \sol{65\,534} \\ \hline
    \hl\ip{10.45.6.9}      & \sol{A} & \sol{255.0.0.0}     & \sol{10.0.0.0}     & \sol{10.255.255.255} & \sol{16\,777\,214} \\ \hline}
  \def\apipa{169.254.88.3}
  \def\Cnet{192.168.80}
  \def\Crows{%
    PC-A & \ip{192.168.8.11}   & \ip{255.255.255.0} & \ip{192.168.80.254} & — \\
    PC-B & \ip{192.168.80.12}  & \ip{255.255.255.0} & \ip{192.168.80.254} & — \\
    PC-C & \ip{192.168.80.13}  & \ip{255.255.255.0} & \ip{192.168.80.100} & — \\
    PC-D & \ip{192.168.80.12}  & \ip{255.255.255.0} & \ip{192.168.80.254} & — \\
    SRV-C & \ip{10.0.0.8}      & \ip{255.0.0.0}     & \ip{10.0.0.254}     & — \\
    R1 — G0/0 & \ip{192.168.80.254} & \ip{255.255.255.0} & — & Off \\
    R1 — G0/1 & \ip{10.0.0.254}     & \ip{255.0.0.0}     & — & On \\}
  \def\Csol{%
    \hl 1 & \sol{PC-A} & \sol{IP hors du réseau 192.168.80.0/24.} & \sol{Ne joint ni les postes ni le routeur.} & \sol{Ex. 192.168.80.11.} \\ \hline
    \hl 2 & \sol{PC-C} & \sol{Passerelle .80.100 : pas l'adresse du routeur.} & \sol{Joint son réseau, mais aucun autre (serveur).} & \sol{Passerelle 192.168.80.254.} \\ \hline
    \hl 3 & \sol{PC-D} & \sol{Même IP que PC-B (doublon).} & \sol{Conflit d'adresse : PC-B et PC-D perturbés.} & \sol{Adresse unique, ex. .80.14.} \\ \hline
    \hl 4 & \sol{R1 G0/0} & \sol{Interface côté LAN désactivée (Off).} & \sol{Aucun poste ne joint le routeur ni le serveur.} & \sol{Port Status : On.} \\ \hline}
  \def\Cbareme{PC-B et SRV-C sont correctement configurés ; accepter PC-B cité à la place de PC-D pour le doublon.}
\fi

\pagestyle{fancy}
\fancyhf{}
\renewcommand{\headrulewidth}{0pt}
\fancyfoot[L]{\footnotesize\itshape\color{grisfoot}Évaluation — Réseaux locaux — Sujet \variante\ifcorrige\ — CORRIGÉ\fi}
\fancyfoot[R]{\footnotesize\color{grisfoot}\thepage/3}

\setlist{itemsep=0pt,topsep=2pt,parsep=0pt}
\setlength{\parindent}{0pt}
\setlength{\parskip}{2pt}

\begin{document}

\begin{center}
  {\Large\bfseries Évaluation pratique — Réseaux locaux \quad
   \fcolorbox{black}{white}{\ SUJET \variante\ }\par}
  \vspace{1pt}
  {\bfseries Adressage IP, commutation et routage — BTS CIEL 1re année — durée : 2 heures\par}
  \ifcorrige\vspace{2pt}{\bfseries\color{corr}CORRIGÉ ET BARÈME — document enseignant}\par\fi
\end{center}
\vspace{-2pt}
{\renewcommand{\arraystretch}{1.35}
\begin{tabularx}{\linewidth}{|X|X|p{2.2cm}|p{2.6cm}|}
  \hline
  \textbf{Nom :} & \textbf{Prénom :} & \textbf{Poste n°} & \textbf{Note :} \hfill /20 \\ \hline
\end{tabularx}}

\textbf{Consignes :} travail \textbf{individuel}, documents et Internet \textbf{non autorisés}. Ce sujet est le
\textbf{document-réponse}. Fichier Packet Tracer : \ip{Nom\_Prenom\_EvalLAN\_\variante.pkt}, à enregistrer
régulièrement dans le dossier de rendu. Les « visas » sont apposés par le professeur, que vous appelez à chaque test réussi.
Durées conseillées : A 30 min — B 1 h 10 — C 20 min.

%==========================================================================
\section{Partie A — Questions de cours et calculs \hfill /6}

\question{A1}{2 pts}{Complétez le tableau en considérant le \textbf{masque par défaut} de la classe de chaque adresse.}
{\small\renewcommand{\arraystretch}{1.15}
\begin{tabularx}{\linewidth}{|l|c|C|C|C|c|}
  \hline
  \textbf{Adresse IP} & \textbf{Classe} & \textbf{Masque par défaut} & \textbf{Adresse réseau} &
  \textbf{Adresse de diffusion} & \textbf{Nb d'hôtes} \\ \hline
  \AIrows
\end{tabularx}}
\bareme{0,5 pt par ligne juste + 0,5 pt si les trois nombres d'hôtes sont justes ($2^{n}-2$, $n$ = nombre de bits d'hôte).}

\question{A2}{1 pt}{Indiquez le type de câble (\textbf{droit} ou \textbf{croisé}) pour chaque liaison.}
{\small\renewcommand{\arraystretch}{1.15}
\begin{tabularx}{\linewidth}{|X|c||X|c|}
  \hline
  \hl PC $\leftrightarrow$ commutateur & \makebox[1.6cm]{\sol{droit}} & Commutateur $\leftrightarrow$ routeur & \makebox[1.6cm]{\sol{droit}} \\ \hline
  \hl PC $\leftrightarrow$ PC (réseau P2P) & \sol{croisé} & Concentrateur $\leftrightarrow$ concentrateur & \sol{croisé} \\ \hline
\end{tabularx}}
\bareme{0,25 pt par case. Équipements de types différents : câble droit ; de même type : câble croisé.}

\question{A3}{1 pt}{Comment un \textbf{concentrateur (hub)} et un \textbf{commutateur (switch)} traitent-ils une trame
reçue ? Lequel limite les collisions, et pourquoi ?}
\zone{2.9cm}{Le \textbf{hub} (couche 1) répète la trame sur \textbf{tous ses autres ports} : un seul domaine de collision
partagé. Le \textbf{switch} (couche 2) lit l'adresse MAC de destination et, grâce à sa \textbf{table MAC}, ne l'envoie
\textbf{que sur le port du destinataire} : un domaine de collision par port, c'est lui qui limite les collisions.
\emph{0,5 pt hub, 0,5 pt switch.}}

\question{A4}{1 pt}{Qu'est-ce que la \textbf{passerelle par défaut} d'un poste ? Quelle condition son adresse doit-elle respecter ?}
\zone{2.4cm}{Adresse IP de l'interface du \textbf{routeur} à laquelle le poste envoie les paquets destinés à un
\textbf{autre réseau}. Elle doit \textbf{appartenir au même réseau} que le poste. \emph{0,5 + 0,5 pt.}}

\question{A5}{0,5 pt}{Un portable configuré en DHCP affiche \ip{\apipa} / \ip{255.255.0.0}. Que s'est-il passé ?}
\zone{1.5cm}{Aucun \textbf{serveur DHCP} n'a répondu : le poste s'est attribué une adresse \textbf{APIPA} (169.254.0.0/16).}

\question{A6}{0,5 pt}{Quel protocole utilise \texttt{ping} ? Nommez les deux messages échangés et la couche OSI des adresses IP.}
\zone{1.5cm}{\textbf{ICMP} ; \textbf{Echo Request} / \textbf{Echo Reply} ; adresses IP en \textbf{couche 3 (Réseau)}.}

%==========================================================================
\newpage
\section{Partie B — Réalisation sous Packet Tracer \hfill /10}

\begin{cdc}
\entreprise{} vous confie son réseau : \textbf{trois réseaux} reliés par un \textbf{routeur unique Cisco 2911}
nommé \ip{\Rnom} (interfaces G0/0, G0/1, G0/2).
\begin{itemize}[leftmargin=12pt]
  \item \textbf{Réseau \resUn} : \textbf{4 postes} PC-\pUn1 à PC-\pUn4, commutateur \ip{SW-\pUn}, adressage statique,
        réseau \ip{\netUn.0/24}.
  \item \textbf{Réseau \resDeux} : \textbf{2 postes} PC-\pDeux1, PC-\pDeux2 et \textbf{1 imprimante} IMP-\pDeux,
        commutateur \ip{SW-\pDeux}, adressage statique, réseau de classe B \ip{\netDeux.0.0/16}.
  \item \textbf{Réseau Serveur} : serveur \ip{\srvNom} d'adresse \ip{\netSrv.\srvH/24}, relié \textbf{directement} au routeur.
  \item La passerelle de chaque réseau est sa \textbf{dernière adresse utilisable}. Tous les équipements communiquent
        entre eux et accèdent à la page web du serveur. Nommage imposé (« Display Name »).
\end{itemize}
\end{cdc}

\question{B1}{2 pts}{Complétez le plan d'adressage \textbf{avant} la réalisation.}
{\footnotesize\setlength{\tabcolsep}{3pt}\renewcommand{\arraystretch}{1.15}
\begin{tabularx}{\linewidth}{|l|C|C|C|C|C|C|}
  \hline
  \textbf{Réseau} & \textbf{Adresse réseau} & \textbf{Masque} & \textbf{1re @ utilisable} &
  \textbf{Dernière @ utilisable} & \textbf{Diffusion} & \textbf{Passerelle} \\ \hline
  \hl \resUn & \sol{\netUn.0} & \sol{255.255.255.0} & \sol{\netUn.1} & \sol{\netUn.254} & \sol{\netUn.255} & \sol{\netUn.254} \\ \hline
  \hl \resDeux & \sol{\netDeux.0.0} & \sol{255.255.0.0} & \sol{\netDeux.0.1} & \sol{\netDeux.255.254} & \sol{\netDeux.255.255} & \sol{\netDeux.255.254} \\ \hline
  \hl Serveur & \sol{\netSrv.0} & \sol{255.255.255.0} & \sol{\netSrv.1} & \sol{\netSrv.254} & \sol{\netSrv.255} & \sol{\netSrv.254} \\ \hline
\end{tabularx}}
\bareme{0,5 pt par ligne juste + 0,5 pt si la convention de passerelle est respectée partout.}

\question{B2}{2 pts}{Construisez la topologie (équipements, câblage, nommage) et relevez l'adressage configuré.}
{\footnotesize\setlength{\tabcolsep}{3pt}\renewcommand{\arraystretch}{1.1}
\begin{tabularx}{\linewidth}{|l|C|C|C||l|C|C|}
  \hline
  \textbf{Hôte} & \textbf{Adresse IP} & \textbf{Masque} & \textbf{Passerelle} & \textbf{Interface} & \textbf{Adresse IP} & \textbf{Masque} \\ \hline
  \hl PC-\pUn1 & \sol{\netUn.1} & \sol{255.255.255.0} & \sol{\netUn.254} & \Rnom{} G0/0 & \sol{\netUn.254} & \sol{255.255.255.0} \\ \hline
  \hl PC-\pDeux1 & \sol{\netDeux.0.1} & \sol{255.255.0.0} & \sol{\netDeux.255.254} & \Rnom{} G0/1 & \sol{\netDeux.255.254} & \sol{255.255.0.0} \\ \hline
  \hl IMP-\pDeux & \sol{\netDeux.0.10} & \sol{255.255.0.0} & \sol{\netDeux.255.254} & \Rnom{} G0/2 & \sol{\netSrv.254} & \sol{255.255.255.0} \\ \hline
  \hl \srvNom & \ip{\netSrv.\srvH} & \sol{255.255.255.0} & \sol{\netSrv.254} & \multicolumn{3}{c|}{\cellcolor{black!8}} \\ \hline
\end{tabularx}}
\bareme{Adresses d'hôtes libres (uniques, dans le bon réseau) ; affectation G0/x libre. Sur le .pkt : équipements 0,5 ;
câblage 1 (câbles droits vers les switchs et switch–routeur, \textbf{croisé} serveur–routeur, liens verts) ; nommage 0,5.}

\question{B3}{3 pts}{Configurez tous les équipements : IP, masques, passerelles des hôtes ; interfaces du routeur
(« Port Status : On »).}
\bareme{Sur le .pkt : interfaces routeur 1 pt ; IP/masques des hôtes 1 pt ; passerelles (imprimante et serveur compris)
1 pt. $-0,25$ par erreur, sans descendre sous 0 par item.}

\question{B4}{3 pts}{Réalisez les tests, notez la commande et le résultat, puis faites viser chaque test réussi.}
{\small\renewcommand{\arraystretch}{1.1}
\begin{tabularx}{\linewidth}{|c|l|l|C|c|c|}
  \hline
  \textbf{N°} & \textbf{Source} & \textbf{Destination} & \textbf{Commande / action} & \textbf{Résultat} & \textbf{Visa prof.} \\ \hline
  \hl 1 & PC-\pUn1 & PC-\pUn4 & \sol{ping \netUn.4} & \sol{réussi} & \hspace{1.4cm} \\ \hline
  \hl 2 & PC-\pDeux1 & IMP-\pDeux & \sol{ping \netDeux.0.10} & \sol{réussi} & \\ \hline
  \hl 3 & PC-\pUn2 & PC-\pDeux2 & \sol{ping \netDeux.0.2} & \sol{réussi} & \\ \hline
  \hl 4 & PC-\pDeux2 & \srvNom & \sol{ping \netSrv.\srvH} & \sol{réussi} & \\ \hline
  \hl 5 & PC-\pUn3 & page web du serveur & \sol{navigateur : {\NoAutoSpacing http://\netSrv.\srvH}} & \sol{page affichée} & \\ \hline
\end{tabularx}}
\bareme{0,5 pt par test visé + 0,5 pt pour B5. Le premier ping inter-réseaux peut échouer (résolution ARP) ; les suivants doivent réussir.}

\question{B5}{incl. B4}{Test n°3 : par quels équipements passe le paquet ? Pourquoi PC-\pUn2 l'envoie-t-il au routeur ?}
\zone{2.3cm}{PC-\pUn2 $\to$ SW-\pUn{} $\to$ \Rnom{} $\to$ SW-\pDeux{} $\to$ PC-\pDeux2. Grâce à son masque, PC-\pUn2 constate que
la destination (\netDeux.0.0/16) n'est \textbf{pas dans son réseau} (\netUn.0/24) : il envoie le paquet à sa
\textbf{passerelle par défaut}, qui le route.}

%==========================================================================
\newpage
\section{Partie C — Diagnostic d'un réseau en panne \hfill /4}

Un technicien constate qu'\textbf{aucun poste ne parvient à joindre le serveur SRV-C} et que certains postes
communiquent mal entre eux. Les câbles sont du bon type.

\begin{center}
\begin{minipage}{0.37\linewidth}\centering
\begin{tikzpicture}[font=\footnotesize,
  eq/.style={draw,fill=qbg,rounded corners=2pt,minimum width=9mm,minimum height=5mm,inner sep=2pt},
  every path/.style={thick}]
  \node[eq] (SW) {SW-1};
  \node[eq,right=9mm of SW] (R) {R1};
  \node[eq,right=7mm of R] (SRV) {SRV-C};
  \node[eq,above left=5mm and -3mm of SW] (A) {PC-A};
  \node[eq,above right=5mm and -3mm of SW] (B) {PC-B};
  \node[eq,below left=5mm and -3mm of SW] (C) {PC-C};
  \node[eq,below right=5mm and -3mm of SW] (D) {PC-D};
  \foreach \n in {A,B,C,D} \draw (\n) -- (SW);
  \draw (SW) -- node[above,font=\tiny]{G0/0} (R);
  \draw (R) -- node[above,font=\tiny]{G0/1} (SRV);
\end{tikzpicture}
\end{minipage}\hfill
\begin{minipage}{0.61\linewidth}\centering
{\small\renewcommand{\arraystretch}{1.05}\setlength{\tabcolsep}{4pt}
\begin{tabular}{lllll}
  \toprule
  \textbf{Équipement} & \textbf{Adresse IP} & \textbf{Masque} & \textbf{Passerelle} & \textbf{Port} \\
  \midrule
  \Crows
  \bottomrule
\end{tabular}}
\end{minipage}
\end{center}

\question{C1}{4 pts}{Relevez les \textbf{quatre anomalies}. Pour chacune : conséquence et correction.}
{\small\renewcommand{\arraystretch}{1.1}
\ifcorrige\def\hc{}\else\def\hc{\rule[-10mm]{0pt}{14mm}}\fi
\begin{tabularx}{\linewidth}{|c|>{\raggedright\arraybackslash}p{2cm}|L|L|L|}
  \hline
  \textbf{N°} & \textbf{Équipement} & \textbf{Anomalie} & \textbf{Conséquence} & \textbf{Correction} \\ \hline
  \ifcorrige\Csol\else
  \hc 1 & & & & \\ \hline
  \hc 2 & & & & \\ \hline
  \hc 3 & & & & \\ \hline
  \hc 4 & & & & \\ \hline
  \fi
\end{tabularx}}
\bareme{1 pt par anomalie (0,25 équipement, anomalie, conséquence, correction ; ordre indifférent). \Cbareme}

%==========================================================================
\section{Bonus — Extension sans fil \hfill +1}

\begin{cdc}Ajoutez au réseau \resUn{} un \textbf{AccessPoint-PT} relié à SW-\pUn{} (SSID \ip{\ssid}, \textbf{WPA2-PSK})
et un portable \ip{Laptop-\pUn} équipé d'un module \texttt{WPC300N}, en adressage \textbf{statique}.
Vérifiez qu'il joint \srvNom.\end{cdc}
{\small\renewcommand{\arraystretch}{1.1}
\begin{tabularx}{\linewidth}{|l|C|C|C|c|}
  \hline
  \textbf{Équipement} & \textbf{Adresse IP} & \textbf{Masque} & \textbf{Passerelle} & \textbf{Visa prof.} \\ \hline
  \hl Laptop-\pUn & \sol{\netUn.50} & \sol{255.255.255.0} & \sol{\netUn.254} & \hspace{1.4cm} \\ \hline
\end{tabularx}}
\bareme{Le point d'accès fait office de pont : le portable est dans \netUn.0/24. Bonus si associé en WPA2-PSK et ping
vers \netSrv.\srvH{} réussi. Note plafonnée à 20.}

\end{document}
"
%   pdflatex -jobname=Eval_Reseau_CorrigeB "\def\variante{B}\def\modecorrige{1}% Évaluation pratique — Réseaux locaux (BTS CIEL 1re année) — sujets A et B
% Compilation (4 PDF) :
%   pdflatex -jobname=Eval_Reseau_SujetA   "\def\variante{A}\def\modecorrige{0}\input{Eval_TP_Reseau_LAN.tex}"
%   pdflatex -jobname=Eval_Reseau_SujetB   "\def\variante{B}\def\modecorrige{0}\input{Eval_TP_Reseau_LAN.tex}"
%   pdflatex -jobname=Eval_Reseau_CorrigeA "\def\variante{A}\def\modecorrige{1}\input{Eval_TP_Reseau_LAN.tex}"
%   pdflatex -jobname=Eval_Reseau_CorrigeB "\def\variante{B}\def\modecorrige{1}\input{Eval_TP_Reseau_LAN.tex}"
\documentclass[10pt,a4paper]{article}

\usepackage[T1]{fontenc}
\usepackage[utf8]{inputenc}
\usepackage{lmodern}
\usepackage[french]{babel}
\usepackage[a4paper,hmargin=1.5cm,top=1.3cm,bottom=1.6cm,footskip=0.7cm]{geometry}
\usepackage[table]{xcolor}
\usepackage{booktabs,tabularx,array}
\usepackage{enumitem}
\usepackage{fancyhdr}
\usepackage[most]{tcolorbox}
\usepackage{tikz}
\usetikzlibrary{positioning}
\usepackage{titlesec}

\providecommand{\modecorrige}{0}
\providecommand{\variante}{A}
\newif\ifcorrige
\ifnum\modecorrige=1 \corrigetrue \else \corrigefalse \fi

\definecolor{ipred}{HTML}{C0392B}
\definecolor{cdcbar}{HTML}{A67C3D}
\definecolor{cdcbg}{HTML}{F4F1EA}
\definecolor{qbar}{HTML}{4A6B5D}
\definecolor{qbg}{HTML}{EDF1EF}
\definecolor{corr}{HTML}{1F5FA8}
\definecolor{corrbg}{HTML}{E8F0FA}
\definecolor{grisfoot}{HTML}{828282}

\newcommand{\ip}[1]{\textcolor{ipred}{\texttt{#1}}}
\newcommand{\pts}[1]{\hfill\mbox{\textbf{(#1)}}}
\newcommand{\sol}[1]{\ifcorrige\textcolor{corr}{#1}\fi}
\newcommand{\bareme}[1]{\ifcorrige\par{\footnotesize\color{corr}#1\par}\fi}
\newcommand{\hl}{\ifcorrige\rule[-1.6mm]{0pt}{5.4mm}\else\rule[-2.2mm]{0pt}{6.6mm}\fi}
\newcolumntype{C}{>{\centering\arraybackslash}X}
\newcolumntype{L}{>{\raggedright\arraybackslash}X}

\titleformat{\section}{\large\bfseries}{}{0pt}{}
\titlespacing*{\section}{0pt}{6pt}{3pt}

\tcbset{qstyle/.style={enhanced,boxrule=0pt,leftrule=2.5pt,arc=0pt,outer arc=0pt,
  left=5pt,right=5pt,top=2pt,bottom=2pt,before skip=4pt,after skip=3pt}}
\newtcolorbox{cdc}{qstyle,colback=cdcbg,colframe=cdcbar}
\newtcolorbox{quest}{qstyle,colback=qbg,colframe=qbar}

% Zone de réponse : vierge dans le sujet, remplie dans le corrigé
\newcommand{\zone}[2]{%
  \ifcorrige
    \begin{tcolorbox}[colback=corrbg,colframe=corr,boxrule=0.5pt,arc=1pt,left=4pt,right=4pt,top=1pt,bottom=1pt,
      before skip=2pt,after skip=2pt,fontupper=\footnotesize\color{corr}]#2\end{tcolorbox}%
  \else
    \begin{tcolorbox}[colback=white,colframe=black!45,boxrule=0.4pt,arc=1pt,height=#1,
      before skip=2pt,after skip=2pt]\end{tcolorbox}%
  \fi}

\newcommand{\question}[3]{\begin{quest}\textbf{#1.} #3\pts{#2}\end{quest}}

%==========================================================================
% Données propres à chaque variante
%==========================================================================
\if\variante A
  \def\entreprise{Le cabinet d'architectes \textbf{ArchiNova}}
  \def\Rnom{R-ArchiNova}
  \def\resUn{Bureau d'études}\def\pUn{BE}\def\netUn{192.168.40}
  \def\resDeux{Accueil}\def\pDeux{ACC}\def\netDeux{172.16}
  \def\srvNom{SRV-Intranet}\def\netSrv{192.168.200}\def\srvH{10}
  \def\ssid{ArchiNova-WiFi}
  \def\AIrows{%
    \hl\ip{192.168.40.25} & \sol{C} & \sol{255.255.255.0} & \sol{192.168.40.0} & \sol{192.168.40.255} & \sol{254} \\ \hline
    \hl\ip{172.16.5.10}   & \sol{B} & \sol{255.255.0.0}   & \sol{172.16.0.0}   & \sol{172.16.255.255} & \sol{65\,534} \\ \hline
    \hl\ip{10.3.2.1}      & \sol{A} & \sol{255.0.0.0}     & \sol{10.0.0.0}     & \sol{10.255.255.255} & \sol{16\,777\,214} \\ \hline}
  \def\apipa{169.254.12.7}
  \def\Cnet{192.168.60}
  \def\Crows{%
    PC-A & \ip{192.168.60.10} & \ip{255.255.255.0} & \ip{192.168.60.254} & — \\
    PC-B & \ip{192.168.60.20} & \ip{255.255.255.0} & \ip{192.168.60.1}   & — \\
    PC-C & \ip{192.168.61.30} & \ip{255.255.255.0} & \ip{192.168.60.254} & — \\
    PC-D & \ip{192.168.60.10} & \ip{255.255.255.0} & \ip{192.168.60.254} & — \\
    SRV-C & \ip{10.0.0.5}     & \ip{255.0.0.0}     & \ip{10.0.0.254}     & — \\
    R1 — G0/0 & \ip{192.168.60.254} & \ip{255.255.255.0} & — & On \\
    R1 — G0/1 & \ip{10.0.0.254}     & \ip{255.0.0.0}     & — & Off \\}
  \def\Csol{%
    \hl 1 & \sol{PC-B} & \sol{Passerelle .60.1 : pas l'adresse du routeur.} & \sol{Joint son réseau, mais aucun autre (serveur).} & \sol{Passerelle 192.168.60.254.} \\ \hline
    \hl 2 & \sol{PC-C} & \sol{IP hors du réseau 192.168.60.0/24.} & \sol{Ne joint ni les postes ni le routeur.} & \sol{Ex. 192.168.60.30.} \\ \hline
    \hl 3 & \sol{PC-D} & \sol{Même IP que PC-A (doublon).} & \sol{Conflit d'adresse : PC-A et PC-D perturbés.} & \sol{Adresse unique, ex. .60.40.} \\ \hline
    \hl 4 & \sol{R1 G0/1} & \sol{Interface désactivée (Off).} & \sol{Réseau du serveur injoignable pour tous.} & \sol{Port Status : On.} \\ \hline}
  \def\Cbareme{PC-A et SRV-C sont correctement configurés ; accepter PC-A cité à la place de PC-D pour le doublon.}
\else
  \def\entreprise{La clinique vétérinaire \textbf{VetoPlus}}
  \def\Rnom{R-VetoPlus}
  \def\resUn{Soins}\def\pUn{SOI}\def\netUn{192.168.70}
  \def\resDeux{Secrétariat}\def\pDeux{SEC}\def\netDeux{172.20}
  \def\srvNom{SRV-Dossiers}\def\netSrv{192.168.250}\def\srvH{20}
  \def\ssid{VetoPlus-WiFi}
  \def\AIrows{%
    \hl\ip{192.168.15.130} & \sol{C} & \sol{255.255.255.0} & \sol{192.168.15.0} & \sol{192.168.15.255} & \sol{254} \\ \hline
    \hl\ip{172.31.200.7}   & \sol{B} & \sol{255.255.0.0}   & \sol{172.31.0.0}   & \sol{172.31.255.255} & \sol{65\,534} \\ \hline
    \hl\ip{10.45.6.9}      & \sol{A} & \sol{255.0.0.0}     & \sol{10.0.0.0}     & \sol{10.255.255.255} & \sol{16\,777\,214} \\ \hline}
  \def\apipa{169.254.88.3}
  \def\Cnet{192.168.80}
  \def\Crows{%
    PC-A & \ip{192.168.8.11}   & \ip{255.255.255.0} & \ip{192.168.80.254} & — \\
    PC-B & \ip{192.168.80.12}  & \ip{255.255.255.0} & \ip{192.168.80.254} & — \\
    PC-C & \ip{192.168.80.13}  & \ip{255.255.255.0} & \ip{192.168.80.100} & — \\
    PC-D & \ip{192.168.80.12}  & \ip{255.255.255.0} & \ip{192.168.80.254} & — \\
    SRV-C & \ip{10.0.0.8}      & \ip{255.0.0.0}     & \ip{10.0.0.254}     & — \\
    R1 — G0/0 & \ip{192.168.80.254} & \ip{255.255.255.0} & — & Off \\
    R1 — G0/1 & \ip{10.0.0.254}     & \ip{255.0.0.0}     & — & On \\}
  \def\Csol{%
    \hl 1 & \sol{PC-A} & \sol{IP hors du réseau 192.168.80.0/24.} & \sol{Ne joint ni les postes ni le routeur.} & \sol{Ex. 192.168.80.11.} \\ \hline
    \hl 2 & \sol{PC-C} & \sol{Passerelle .80.100 : pas l'adresse du routeur.} & \sol{Joint son réseau, mais aucun autre (serveur).} & \sol{Passerelle 192.168.80.254.} \\ \hline
    \hl 3 & \sol{PC-D} & \sol{Même IP que PC-B (doublon).} & \sol{Conflit d'adresse : PC-B et PC-D perturbés.} & \sol{Adresse unique, ex. .80.14.} \\ \hline
    \hl 4 & \sol{R1 G0/0} & \sol{Interface côté LAN désactivée (Off).} & \sol{Aucun poste ne joint le routeur ni le serveur.} & \sol{Port Status : On.} \\ \hline}
  \def\Cbareme{PC-B et SRV-C sont correctement configurés ; accepter PC-B cité à la place de PC-D pour le doublon.}
\fi

\pagestyle{fancy}
\fancyhf{}
\renewcommand{\headrulewidth}{0pt}
\fancyfoot[L]{\footnotesize\itshape\color{grisfoot}Évaluation — Réseaux locaux — Sujet \variante\ifcorrige\ — CORRIGÉ\fi}
\fancyfoot[R]{\footnotesize\color{grisfoot}\thepage/3}

\setlist{itemsep=0pt,topsep=2pt,parsep=0pt}
\setlength{\parindent}{0pt}
\setlength{\parskip}{2pt}

\begin{document}

\begin{center}
  {\Large\bfseries Évaluation pratique — Réseaux locaux \quad
   \fcolorbox{black}{white}{\ SUJET \variante\ }\par}
  \vspace{1pt}
  {\bfseries Adressage IP, commutation et routage — BTS CIEL 1re année — durée : 2 heures\par}
  \ifcorrige\vspace{2pt}{\bfseries\color{corr}CORRIGÉ ET BARÈME — document enseignant}\par\fi
\end{center}
\vspace{-2pt}
{\renewcommand{\arraystretch}{1.35}
\begin{tabularx}{\linewidth}{|X|X|p{2.2cm}|p{2.6cm}|}
  \hline
  \textbf{Nom :} & \textbf{Prénom :} & \textbf{Poste n°} & \textbf{Note :} \hfill /20 \\ \hline
\end{tabularx}}

\textbf{Consignes :} travail \textbf{individuel}, documents et Internet \textbf{non autorisés}. Ce sujet est le
\textbf{document-réponse}. Fichier Packet Tracer : \ip{Nom\_Prenom\_EvalLAN\_\variante.pkt}, à enregistrer
régulièrement dans le dossier de rendu. Les « visas » sont apposés par le professeur, que vous appelez à chaque test réussi.
Durées conseillées : A 30 min — B 1 h 10 — C 20 min.

%==========================================================================
\section{Partie A — Questions de cours et calculs \hfill /6}

\question{A1}{2 pts}{Complétez le tableau en considérant le \textbf{masque par défaut} de la classe de chaque adresse.}
{\small\renewcommand{\arraystretch}{1.15}
\begin{tabularx}{\linewidth}{|l|c|C|C|C|c|}
  \hline
  \textbf{Adresse IP} & \textbf{Classe} & \textbf{Masque par défaut} & \textbf{Adresse réseau} &
  \textbf{Adresse de diffusion} & \textbf{Nb d'hôtes} \\ \hline
  \AIrows
\end{tabularx}}
\bareme{0,5 pt par ligne juste + 0,5 pt si les trois nombres d'hôtes sont justes ($2^{n}-2$, $n$ = nombre de bits d'hôte).}

\question{A2}{1 pt}{Indiquez le type de câble (\textbf{droit} ou \textbf{croisé}) pour chaque liaison.}
{\small\renewcommand{\arraystretch}{1.15}
\begin{tabularx}{\linewidth}{|X|c||X|c|}
  \hline
  \hl PC $\leftrightarrow$ commutateur & \makebox[1.6cm]{\sol{droit}} & Commutateur $\leftrightarrow$ routeur & \makebox[1.6cm]{\sol{droit}} \\ \hline
  \hl PC $\leftrightarrow$ PC (réseau P2P) & \sol{croisé} & Concentrateur $\leftrightarrow$ concentrateur & \sol{croisé} \\ \hline
\end{tabularx}}
\bareme{0,25 pt par case. Équipements de types différents : câble droit ; de même type : câble croisé.}

\question{A3}{1 pt}{Comment un \textbf{concentrateur (hub)} et un \textbf{commutateur (switch)} traitent-ils une trame
reçue ? Lequel limite les collisions, et pourquoi ?}
\zone{2.9cm}{Le \textbf{hub} (couche 1) répète la trame sur \textbf{tous ses autres ports} : un seul domaine de collision
partagé. Le \textbf{switch} (couche 2) lit l'adresse MAC de destination et, grâce à sa \textbf{table MAC}, ne l'envoie
\textbf{que sur le port du destinataire} : un domaine de collision par port, c'est lui qui limite les collisions.
\emph{0,5 pt hub, 0,5 pt switch.}}

\question{A4}{1 pt}{Qu'est-ce que la \textbf{passerelle par défaut} d'un poste ? Quelle condition son adresse doit-elle respecter ?}
\zone{2.4cm}{Adresse IP de l'interface du \textbf{routeur} à laquelle le poste envoie les paquets destinés à un
\textbf{autre réseau}. Elle doit \textbf{appartenir au même réseau} que le poste. \emph{0,5 + 0,5 pt.}}

\question{A5}{0,5 pt}{Un portable configuré en DHCP affiche \ip{\apipa} / \ip{255.255.0.0}. Que s'est-il passé ?}
\zone{1.5cm}{Aucun \textbf{serveur DHCP} n'a répondu : le poste s'est attribué une adresse \textbf{APIPA} (169.254.0.0/16).}

\question{A6}{0,5 pt}{Quel protocole utilise \texttt{ping} ? Nommez les deux messages échangés et la couche OSI des adresses IP.}
\zone{1.5cm}{\textbf{ICMP} ; \textbf{Echo Request} / \textbf{Echo Reply} ; adresses IP en \textbf{couche 3 (Réseau)}.}

%==========================================================================
\newpage
\section{Partie B — Réalisation sous Packet Tracer \hfill /10}

\begin{cdc}
\entreprise{} vous confie son réseau : \textbf{trois réseaux} reliés par un \textbf{routeur unique Cisco 2911}
nommé \ip{\Rnom} (interfaces G0/0, G0/1, G0/2).
\begin{itemize}[leftmargin=12pt]
  \item \textbf{Réseau \resUn} : \textbf{4 postes} PC-\pUn1 à PC-\pUn4, commutateur \ip{SW-\pUn}, adressage statique,
        réseau \ip{\netUn.0/24}.
  \item \textbf{Réseau \resDeux} : \textbf{2 postes} PC-\pDeux1, PC-\pDeux2 et \textbf{1 imprimante} IMP-\pDeux,
        commutateur \ip{SW-\pDeux}, adressage statique, réseau de classe B \ip{\netDeux.0.0/16}.
  \item \textbf{Réseau Serveur} : serveur \ip{\srvNom} d'adresse \ip{\netSrv.\srvH/24}, relié \textbf{directement} au routeur.
  \item La passerelle de chaque réseau est sa \textbf{dernière adresse utilisable}. Tous les équipements communiquent
        entre eux et accèdent à la page web du serveur. Nommage imposé (« Display Name »).
\end{itemize}
\end{cdc}

\question{B1}{2 pts}{Complétez le plan d'adressage \textbf{avant} la réalisation.}
{\footnotesize\setlength{\tabcolsep}{3pt}\renewcommand{\arraystretch}{1.15}
\begin{tabularx}{\linewidth}{|l|C|C|C|C|C|C|}
  \hline
  \textbf{Réseau} & \textbf{Adresse réseau} & \textbf{Masque} & \textbf{1re @ utilisable} &
  \textbf{Dernière @ utilisable} & \textbf{Diffusion} & \textbf{Passerelle} \\ \hline
  \hl \resUn & \sol{\netUn.0} & \sol{255.255.255.0} & \sol{\netUn.1} & \sol{\netUn.254} & \sol{\netUn.255} & \sol{\netUn.254} \\ \hline
  \hl \resDeux & \sol{\netDeux.0.0} & \sol{255.255.0.0} & \sol{\netDeux.0.1} & \sol{\netDeux.255.254} & \sol{\netDeux.255.255} & \sol{\netDeux.255.254} \\ \hline
  \hl Serveur & \sol{\netSrv.0} & \sol{255.255.255.0} & \sol{\netSrv.1} & \sol{\netSrv.254} & \sol{\netSrv.255} & \sol{\netSrv.254} \\ \hline
\end{tabularx}}
\bareme{0,5 pt par ligne juste + 0,5 pt si la convention de passerelle est respectée partout.}

\question{B2}{2 pts}{Construisez la topologie (équipements, câblage, nommage) et relevez l'adressage configuré.}
{\footnotesize\setlength{\tabcolsep}{3pt}\renewcommand{\arraystretch}{1.1}
\begin{tabularx}{\linewidth}{|l|C|C|C||l|C|C|}
  \hline
  \textbf{Hôte} & \textbf{Adresse IP} & \textbf{Masque} & \textbf{Passerelle} & \textbf{Interface} & \textbf{Adresse IP} & \textbf{Masque} \\ \hline
  \hl PC-\pUn1 & \sol{\netUn.1} & \sol{255.255.255.0} & \sol{\netUn.254} & \Rnom{} G0/0 & \sol{\netUn.254} & \sol{255.255.255.0} \\ \hline
  \hl PC-\pDeux1 & \sol{\netDeux.0.1} & \sol{255.255.0.0} & \sol{\netDeux.255.254} & \Rnom{} G0/1 & \sol{\netDeux.255.254} & \sol{255.255.0.0} \\ \hline
  \hl IMP-\pDeux & \sol{\netDeux.0.10} & \sol{255.255.0.0} & \sol{\netDeux.255.254} & \Rnom{} G0/2 & \sol{\netSrv.254} & \sol{255.255.255.0} \\ \hline
  \hl \srvNom & \ip{\netSrv.\srvH} & \sol{255.255.255.0} & \sol{\netSrv.254} & \multicolumn{3}{c|}{\cellcolor{black!8}} \\ \hline
\end{tabularx}}
\bareme{Adresses d'hôtes libres (uniques, dans le bon réseau) ; affectation G0/x libre. Sur le .pkt : équipements 0,5 ;
câblage 1 (câbles droits vers les switchs et switch–routeur, \textbf{croisé} serveur–routeur, liens verts) ; nommage 0,5.}

\question{B3}{3 pts}{Configurez tous les équipements : IP, masques, passerelles des hôtes ; interfaces du routeur
(« Port Status : On »).}
\bareme{Sur le .pkt : interfaces routeur 1 pt ; IP/masques des hôtes 1 pt ; passerelles (imprimante et serveur compris)
1 pt. $-0,25$ par erreur, sans descendre sous 0 par item.}

\question{B4}{3 pts}{Réalisez les tests, notez la commande et le résultat, puis faites viser chaque test réussi.}
{\small\renewcommand{\arraystretch}{1.1}
\begin{tabularx}{\linewidth}{|c|l|l|C|c|c|}
  \hline
  \textbf{N°} & \textbf{Source} & \textbf{Destination} & \textbf{Commande / action} & \textbf{Résultat} & \textbf{Visa prof.} \\ \hline
  \hl 1 & PC-\pUn1 & PC-\pUn4 & \sol{ping \netUn.4} & \sol{réussi} & \hspace{1.4cm} \\ \hline
  \hl 2 & PC-\pDeux1 & IMP-\pDeux & \sol{ping \netDeux.0.10} & \sol{réussi} & \\ \hline
  \hl 3 & PC-\pUn2 & PC-\pDeux2 & \sol{ping \netDeux.0.2} & \sol{réussi} & \\ \hline
  \hl 4 & PC-\pDeux2 & \srvNom & \sol{ping \netSrv.\srvH} & \sol{réussi} & \\ \hline
  \hl 5 & PC-\pUn3 & page web du serveur & \sol{navigateur : {\NoAutoSpacing http://\netSrv.\srvH}} & \sol{page affichée} & \\ \hline
\end{tabularx}}
\bareme{0,5 pt par test visé + 0,5 pt pour B5. Le premier ping inter-réseaux peut échouer (résolution ARP) ; les suivants doivent réussir.}

\question{B5}{incl. B4}{Test n°3 : par quels équipements passe le paquet ? Pourquoi PC-\pUn2 l'envoie-t-il au routeur ?}
\zone{2.3cm}{PC-\pUn2 $\to$ SW-\pUn{} $\to$ \Rnom{} $\to$ SW-\pDeux{} $\to$ PC-\pDeux2. Grâce à son masque, PC-\pUn2 constate que
la destination (\netDeux.0.0/16) n'est \textbf{pas dans son réseau} (\netUn.0/24) : il envoie le paquet à sa
\textbf{passerelle par défaut}, qui le route.}

%==========================================================================
\newpage
\section{Partie C — Diagnostic d'un réseau en panne \hfill /4}

Un technicien constate qu'\textbf{aucun poste ne parvient à joindre le serveur SRV-C} et que certains postes
communiquent mal entre eux. Les câbles sont du bon type.

\begin{center}
\begin{minipage}{0.37\linewidth}\centering
\begin{tikzpicture}[font=\footnotesize,
  eq/.style={draw,fill=qbg,rounded corners=2pt,minimum width=9mm,minimum height=5mm,inner sep=2pt},
  every path/.style={thick}]
  \node[eq] (SW) {SW-1};
  \node[eq,right=9mm of SW] (R) {R1};
  \node[eq,right=7mm of R] (SRV) {SRV-C};
  \node[eq,above left=5mm and -3mm of SW] (A) {PC-A};
  \node[eq,above right=5mm and -3mm of SW] (B) {PC-B};
  \node[eq,below left=5mm and -3mm of SW] (C) {PC-C};
  \node[eq,below right=5mm and -3mm of SW] (D) {PC-D};
  \foreach \n in {A,B,C,D} \draw (\n) -- (SW);
  \draw (SW) -- node[above,font=\tiny]{G0/0} (R);
  \draw (R) -- node[above,font=\tiny]{G0/1} (SRV);
\end{tikzpicture}
\end{minipage}\hfill
\begin{minipage}{0.61\linewidth}\centering
{\small\renewcommand{\arraystretch}{1.05}\setlength{\tabcolsep}{4pt}
\begin{tabular}{lllll}
  \toprule
  \textbf{Équipement} & \textbf{Adresse IP} & \textbf{Masque} & \textbf{Passerelle} & \textbf{Port} \\
  \midrule
  \Crows
  \bottomrule
\end{tabular}}
\end{minipage}
\end{center}

\question{C1}{4 pts}{Relevez les \textbf{quatre anomalies}. Pour chacune : conséquence et correction.}
{\small\renewcommand{\arraystretch}{1.1}
\ifcorrige\def\hc{}\else\def\hc{\rule[-10mm]{0pt}{14mm}}\fi
\begin{tabularx}{\linewidth}{|c|>{\raggedright\arraybackslash}p{2cm}|L|L|L|}
  \hline
  \textbf{N°} & \textbf{Équipement} & \textbf{Anomalie} & \textbf{Conséquence} & \textbf{Correction} \\ \hline
  \ifcorrige\Csol\else
  \hc 1 & & & & \\ \hline
  \hc 2 & & & & \\ \hline
  \hc 3 & & & & \\ \hline
  \hc 4 & & & & \\ \hline
  \fi
\end{tabularx}}
\bareme{1 pt par anomalie (0,25 équipement, anomalie, conséquence, correction ; ordre indifférent). \Cbareme}

%==========================================================================
\section{Bonus — Extension sans fil \hfill +1}

\begin{cdc}Ajoutez au réseau \resUn{} un \textbf{AccessPoint-PT} relié à SW-\pUn{} (SSID \ip{\ssid}, \textbf{WPA2-PSK})
et un portable \ip{Laptop-\pUn} équipé d'un module \texttt{WPC300N}, en adressage \textbf{statique}.
Vérifiez qu'il joint \srvNom.\end{cdc}
{\small\renewcommand{\arraystretch}{1.1}
\begin{tabularx}{\linewidth}{|l|C|C|C|c|}
  \hline
  \textbf{Équipement} & \textbf{Adresse IP} & \textbf{Masque} & \textbf{Passerelle} & \textbf{Visa prof.} \\ \hline
  \hl Laptop-\pUn & \sol{\netUn.50} & \sol{255.255.255.0} & \sol{\netUn.254} & \hspace{1.4cm} \\ \hline
\end{tabularx}}
\bareme{Le point d'accès fait office de pont : le portable est dans \netUn.0/24. Bonus si associé en WPA2-PSK et ping
vers \netSrv.\srvH{} réussi. Note plafonnée à 20.}

\end{document}
"
\documentclass[10pt,a4paper]{article}

\usepackage[T1]{fontenc}
\usepackage[utf8]{inputenc}
\usepackage{lmodern}
\usepackage[french]{babel}
\usepackage[a4paper,hmargin=1.5cm,top=1.3cm,bottom=1.6cm,footskip=0.7cm]{geometry}
\usepackage[table]{xcolor}
\usepackage{booktabs,tabularx,array}
\usepackage{enumitem}
\usepackage{fancyhdr}
\usepackage[most]{tcolorbox}
\usepackage{tikz}
\usetikzlibrary{positioning}
\usepackage{titlesec}

\providecommand{\modecorrige}{0}
\providecommand{\variante}{A}
\newif\ifcorrige
\ifnum\modecorrige=1 \corrigetrue \else \corrigefalse \fi

\definecolor{ipred}{HTML}{C0392B}
\definecolor{cdcbar}{HTML}{A67C3D}
\definecolor{cdcbg}{HTML}{F4F1EA}
\definecolor{qbar}{HTML}{4A6B5D}
\definecolor{qbg}{HTML}{EDF1EF}
\definecolor{corr}{HTML}{1F5FA8}
\definecolor{corrbg}{HTML}{E8F0FA}
\definecolor{grisfoot}{HTML}{828282}

\newcommand{\ip}[1]{\textcolor{ipred}{\texttt{#1}}}
\newcommand{\pts}[1]{\hfill\mbox{\textbf{(#1)}}}
\newcommand{\sol}[1]{\ifcorrige\textcolor{corr}{#1}\fi}
\newcommand{\bareme}[1]{\ifcorrige\par{\footnotesize\color{corr}#1\par}\fi}
\newcommand{\hl}{\ifcorrige\rule[-1.6mm]{0pt}{5.4mm}\else\rule[-2.2mm]{0pt}{6.6mm}\fi}
\newcolumntype{C}{>{\centering\arraybackslash}X}
\newcolumntype{L}{>{\raggedright\arraybackslash}X}

\titleformat{\section}{\large\bfseries}{}{0pt}{}
\titlespacing*{\section}{0pt}{6pt}{3pt}

\tcbset{qstyle/.style={enhanced,boxrule=0pt,leftrule=2.5pt,arc=0pt,outer arc=0pt,
  left=5pt,right=5pt,top=2pt,bottom=2pt,before skip=4pt,after skip=3pt}}
\newtcolorbox{cdc}{qstyle,colback=cdcbg,colframe=cdcbar}
\newtcolorbox{quest}{qstyle,colback=qbg,colframe=qbar}

% Zone de réponse : vierge dans le sujet, remplie dans le corrigé
\newcommand{\zone}[2]{%
  \ifcorrige
    \begin{tcolorbox}[colback=corrbg,colframe=corr,boxrule=0.5pt,arc=1pt,left=4pt,right=4pt,top=1pt,bottom=1pt,
      before skip=2pt,after skip=2pt,fontupper=\footnotesize\color{corr}]#2\end{tcolorbox}%
  \else
    \begin{tcolorbox}[colback=white,colframe=black!45,boxrule=0.4pt,arc=1pt,height=#1,
      before skip=2pt,after skip=2pt]\end{tcolorbox}%
  \fi}

\newcommand{\question}[3]{\begin{quest}\textbf{#1.} #3\pts{#2}\end{quest}}

%==========================================================================
% Données propres à chaque variante
%==========================================================================
\if\variante A
  \def\entreprise{Le cabinet d'architectes \textbf{ArchiNova}}
  \def\Rnom{R-ArchiNova}
  \def\resUn{Bureau d'études}\def\pUn{BE}\def\netUn{192.168.40}
  \def\resDeux{Accueil}\def\pDeux{ACC}\def\netDeux{172.16}
  \def\srvNom{SRV-Intranet}\def\netSrv{192.168.200}\def\srvH{10}
  \def\ssid{ArchiNova-WiFi}
  \def\AIrows{%
    \hl\ip{192.168.40.25} & \sol{C} & \sol{255.255.255.0} & \sol{192.168.40.0} & \sol{192.168.40.255} & \sol{254} \\ \hline
    \hl\ip{172.16.5.10}   & \sol{B} & \sol{255.255.0.0}   & \sol{172.16.0.0}   & \sol{172.16.255.255} & \sol{65\,534} \\ \hline
    \hl\ip{10.3.2.1}      & \sol{A} & \sol{255.0.0.0}     & \sol{10.0.0.0}     & \sol{10.255.255.255} & \sol{16\,777\,214} \\ \hline}
  \def\apipa{169.254.12.7}
  \def\Cnet{192.168.60}
  \def\Crows{%
    PC-A & \ip{192.168.60.10} & \ip{255.255.255.0} & \ip{192.168.60.254} & — \\
    PC-B & \ip{192.168.60.20} & \ip{255.255.255.0} & \ip{192.168.60.1}   & — \\
    PC-C & \ip{192.168.61.30} & \ip{255.255.255.0} & \ip{192.168.60.254} & — \\
    PC-D & \ip{192.168.60.10} & \ip{255.255.255.0} & \ip{192.168.60.254} & — \\
    SRV-C & \ip{10.0.0.5}     & \ip{255.0.0.0}     & \ip{10.0.0.254}     & — \\
    R1 — G0/0 & \ip{192.168.60.254} & \ip{255.255.255.0} & — & On \\
    R1 — G0/1 & \ip{10.0.0.254}     & \ip{255.0.0.0}     & — & Off \\}
  \def\Csol{%
    \hl 1 & \sol{PC-B} & \sol{Passerelle .60.1 : pas l'adresse du routeur.} & \sol{Joint son réseau, mais aucun autre (serveur).} & \sol{Passerelle 192.168.60.254.} \\ \hline
    \hl 2 & \sol{PC-C} & \sol{IP hors du réseau 192.168.60.0/24.} & \sol{Ne joint ni les postes ni le routeur.} & \sol{Ex. 192.168.60.30.} \\ \hline
    \hl 3 & \sol{PC-D} & \sol{Même IP que PC-A (doublon).} & \sol{Conflit d'adresse : PC-A et PC-D perturbés.} & \sol{Adresse unique, ex. .60.40.} \\ \hline
    \hl 4 & \sol{R1 G0/1} & \sol{Interface désactivée (Off).} & \sol{Réseau du serveur injoignable pour tous.} & \sol{Port Status : On.} \\ \hline}
  \def\Cbareme{PC-A et SRV-C sont correctement configurés ; accepter PC-A cité à la place de PC-D pour le doublon.}
\else
  \def\entreprise{La clinique vétérinaire \textbf{VetoPlus}}
  \def\Rnom{R-VetoPlus}
  \def\resUn{Soins}\def\pUn{SOI}\def\netUn{192.168.70}
  \def\resDeux{Secrétariat}\def\pDeux{SEC}\def\netDeux{172.20}
  \def\srvNom{SRV-Dossiers}\def\netSrv{192.168.250}\def\srvH{20}
  \def\ssid{VetoPlus-WiFi}
  \def\AIrows{%
    \hl\ip{192.168.15.130} & \sol{C} & \sol{255.255.255.0} & \sol{192.168.15.0} & \sol{192.168.15.255} & \sol{254} \\ \hline
    \hl\ip{172.31.200.7}   & \sol{B} & \sol{255.255.0.0}   & \sol{172.31.0.0}   & \sol{172.31.255.255} & \sol{65\,534} \\ \hline
    \hl\ip{10.45.6.9}      & \sol{A} & \sol{255.0.0.0}     & \sol{10.0.0.0}     & \sol{10.255.255.255} & \sol{16\,777\,214} \\ \hline}
  \def\apipa{169.254.88.3}
  \def\Cnet{192.168.80}
  \def\Crows{%
    PC-A & \ip{192.168.8.11}   & \ip{255.255.255.0} & \ip{192.168.80.254} & — \\
    PC-B & \ip{192.168.80.12}  & \ip{255.255.255.0} & \ip{192.168.80.254} & — \\
    PC-C & \ip{192.168.80.13}  & \ip{255.255.255.0} & \ip{192.168.80.100} & — \\
    PC-D & \ip{192.168.80.12}  & \ip{255.255.255.0} & \ip{192.168.80.254} & — \\
    SRV-C & \ip{10.0.0.8}      & \ip{255.0.0.0}     & \ip{10.0.0.254}     & — \\
    R1 — G0/0 & \ip{192.168.80.254} & \ip{255.255.255.0} & — & Off \\
    R1 — G0/1 & \ip{10.0.0.254}     & \ip{255.0.0.0}     & — & On \\}
  \def\Csol{%
    \hl 1 & \sol{PC-A} & \sol{IP hors du réseau 192.168.80.0/24.} & \sol{Ne joint ni les postes ni le routeur.} & \sol{Ex. 192.168.80.11.} \\ \hline
    \hl 2 & \sol{PC-C} & \sol{Passerelle .80.100 : pas l'adresse du routeur.} & \sol{Joint son réseau, mais aucun autre (serveur).} & \sol{Passerelle 192.168.80.254.} \\ \hline
    \hl 3 & \sol{PC-D} & \sol{Même IP que PC-B (doublon).} & \sol{Conflit d'adresse : PC-B et PC-D perturbés.} & \sol{Adresse unique, ex. .80.14.} \\ \hline
    \hl 4 & \sol{R1 G0/0} & \sol{Interface côté LAN désactivée (Off).} & \sol{Aucun poste ne joint le routeur ni le serveur.} & \sol{Port Status : On.} \\ \hline}
  \def\Cbareme{PC-B et SRV-C sont correctement configurés ; accepter PC-B cité à la place de PC-D pour le doublon.}
\fi

\pagestyle{fancy}
\fancyhf{}
\renewcommand{\headrulewidth}{0pt}
\fancyfoot[L]{\footnotesize\itshape\color{grisfoot}Évaluation — Réseaux locaux — Sujet \variante\ifcorrige\ — CORRIGÉ\fi}
\fancyfoot[R]{\footnotesize\color{grisfoot}\thepage/3}

\setlist{itemsep=0pt,topsep=2pt,parsep=0pt}
\setlength{\parindent}{0pt}
\setlength{\parskip}{2pt}

\begin{document}

\begin{center}
  {\Large\bfseries Évaluation pratique — Réseaux locaux \quad
   \fcolorbox{black}{white}{\ SUJET \variante\ }\par}
  \vspace{1pt}
  {\bfseries Adressage IP, commutation et routage — BTS CIEL 1re année — durée : 2 heures\par}
  \ifcorrige\vspace{2pt}{\bfseries\color{corr}CORRIGÉ ET BARÈME — document enseignant}\par\fi
\end{center}
\vspace{-2pt}
{\renewcommand{\arraystretch}{1.35}
\begin{tabularx}{\linewidth}{|X|X|p{2.2cm}|p{2.6cm}|}
  \hline
  \textbf{Nom :} & \textbf{Prénom :} & \textbf{Poste n°} & \textbf{Note :} \hfill /20 \\ \hline
\end{tabularx}}

\textbf{Consignes :} travail \textbf{individuel}, documents et Internet \textbf{non autorisés}. Ce sujet est le
\textbf{document-réponse}. Fichier Packet Tracer : \ip{Nom\_Prenom\_EvalLAN\_\variante.pkt}, à enregistrer
régulièrement dans le dossier de rendu. Les « visas » sont apposés par le professeur, que vous appelez à chaque test réussi.
Durées conseillées : A 30 min — B 1 h 10 — C 20 min.

%==========================================================================
\section{Partie A — Questions de cours et calculs \hfill /6}

\question{A1}{2 pts}{Complétez le tableau en considérant le \textbf{masque par défaut} de la classe de chaque adresse.}
{\small\renewcommand{\arraystretch}{1.15}
\begin{tabularx}{\linewidth}{|l|c|C|C|C|c|}
  \hline
  \textbf{Adresse IP} & \textbf{Classe} & \textbf{Masque par défaut} & \textbf{Adresse réseau} &
  \textbf{Adresse de diffusion} & \textbf{Nb d'hôtes} \\ \hline
  \AIrows
\end{tabularx}}
\bareme{0,5 pt par ligne juste + 0,5 pt si les trois nombres d'hôtes sont justes ($2^{n}-2$, $n$ = nombre de bits d'hôte).}

\question{A2}{1 pt}{Indiquez le type de câble (\textbf{droit} ou \textbf{croisé}) pour chaque liaison.}
{\small\renewcommand{\arraystretch}{1.15}
\begin{tabularx}{\linewidth}{|X|c||X|c|}
  \hline
  \hl PC $\leftrightarrow$ commutateur & \makebox[1.6cm]{\sol{droit}} & Commutateur $\leftrightarrow$ routeur & \makebox[1.6cm]{\sol{droit}} \\ \hline
  \hl PC $\leftrightarrow$ PC (réseau P2P) & \sol{croisé} & Concentrateur $\leftrightarrow$ concentrateur & \sol{croisé} \\ \hline
\end{tabularx}}
\bareme{0,25 pt par case. Équipements de types différents : câble droit ; de même type : câble croisé.}

\question{A3}{1 pt}{Comment un \textbf{concentrateur (hub)} et un \textbf{commutateur (switch)} traitent-ils une trame
reçue ? Lequel limite les collisions, et pourquoi ?}
\zone{2.9cm}{Le \textbf{hub} (couche 1) répète la trame sur \textbf{tous ses autres ports} : un seul domaine de collision
partagé. Le \textbf{switch} (couche 2) lit l'adresse MAC de destination et, grâce à sa \textbf{table MAC}, ne l'envoie
\textbf{que sur le port du destinataire} : un domaine de collision par port, c'est lui qui limite les collisions.
\emph{0,5 pt hub, 0,5 pt switch.}}

\question{A4}{1 pt}{Qu'est-ce que la \textbf{passerelle par défaut} d'un poste ? Quelle condition son adresse doit-elle respecter ?}
\zone{2.4cm}{Adresse IP de l'interface du \textbf{routeur} à laquelle le poste envoie les paquets destinés à un
\textbf{autre réseau}. Elle doit \textbf{appartenir au même réseau} que le poste. \emph{0,5 + 0,5 pt.}}

\question{A5}{0,5 pt}{Un portable configuré en DHCP affiche \ip{\apipa} / \ip{255.255.0.0}. Que s'est-il passé ?}
\zone{1.5cm}{Aucun \textbf{serveur DHCP} n'a répondu : le poste s'est attribué une adresse \textbf{APIPA} (169.254.0.0/16).}

\question{A6}{0,5 pt}{Quel protocole utilise \texttt{ping} ? Nommez les deux messages échangés et la couche OSI des adresses IP.}
\zone{1.5cm}{\textbf{ICMP} ; \textbf{Echo Request} / \textbf{Echo Reply} ; adresses IP en \textbf{couche 3 (Réseau)}.}

%==========================================================================
\newpage
\section{Partie B — Réalisation sous Packet Tracer \hfill /10}

\begin{cdc}
\entreprise{} vous confie son réseau : \textbf{trois réseaux} reliés par un \textbf{routeur unique Cisco 2911}
nommé \ip{\Rnom} (interfaces G0/0, G0/1, G0/2).
\begin{itemize}[leftmargin=12pt]
  \item \textbf{Réseau \resUn} : \textbf{4 postes} PC-\pUn1 à PC-\pUn4, commutateur \ip{SW-\pUn}, adressage statique,
        réseau \ip{\netUn.0/24}.
  \item \textbf{Réseau \resDeux} : \textbf{2 postes} PC-\pDeux1, PC-\pDeux2 et \textbf{1 imprimante} IMP-\pDeux,
        commutateur \ip{SW-\pDeux}, adressage statique, réseau de classe B \ip{\netDeux.0.0/16}.
  \item \textbf{Réseau Serveur} : serveur \ip{\srvNom} d'adresse \ip{\netSrv.\srvH/24}, relié \textbf{directement} au routeur.
  \item La passerelle de chaque réseau est sa \textbf{dernière adresse utilisable}. Tous les équipements communiquent
        entre eux et accèdent à la page web du serveur. Nommage imposé (« Display Name »).
\end{itemize}
\end{cdc}

\question{B1}{2 pts}{Complétez le plan d'adressage \textbf{avant} la réalisation.}
{\footnotesize\setlength{\tabcolsep}{3pt}\renewcommand{\arraystretch}{1.15}
\begin{tabularx}{\linewidth}{|l|C|C|C|C|C|C|}
  \hline
  \textbf{Réseau} & \textbf{Adresse réseau} & \textbf{Masque} & \textbf{1re @ utilisable} &
  \textbf{Dernière @ utilisable} & \textbf{Diffusion} & \textbf{Passerelle} \\ \hline
  \hl \resUn & \sol{\netUn.0} & \sol{255.255.255.0} & \sol{\netUn.1} & \sol{\netUn.254} & \sol{\netUn.255} & \sol{\netUn.254} \\ \hline
  \hl \resDeux & \sol{\netDeux.0.0} & \sol{255.255.0.0} & \sol{\netDeux.0.1} & \sol{\netDeux.255.254} & \sol{\netDeux.255.255} & \sol{\netDeux.255.254} \\ \hline
  \hl Serveur & \sol{\netSrv.0} & \sol{255.255.255.0} & \sol{\netSrv.1} & \sol{\netSrv.254} & \sol{\netSrv.255} & \sol{\netSrv.254} \\ \hline
\end{tabularx}}
\bareme{0,5 pt par ligne juste + 0,5 pt si la convention de passerelle est respectée partout.}

\question{B2}{2 pts}{Construisez la topologie (équipements, câblage, nommage) et relevez l'adressage configuré.}
{\footnotesize\setlength{\tabcolsep}{3pt}\renewcommand{\arraystretch}{1.1}
\begin{tabularx}{\linewidth}{|l|C|C|C||l|C|C|}
  \hline
  \textbf{Hôte} & \textbf{Adresse IP} & \textbf{Masque} & \textbf{Passerelle} & \textbf{Interface} & \textbf{Adresse IP} & \textbf{Masque} \\ \hline
  \hl PC-\pUn1 & \sol{\netUn.1} & \sol{255.255.255.0} & \sol{\netUn.254} & \Rnom{} G0/0 & \sol{\netUn.254} & \sol{255.255.255.0} \\ \hline
  \hl PC-\pDeux1 & \sol{\netDeux.0.1} & \sol{255.255.0.0} & \sol{\netDeux.255.254} & \Rnom{} G0/1 & \sol{\netDeux.255.254} & \sol{255.255.0.0} \\ \hline
  \hl IMP-\pDeux & \sol{\netDeux.0.10} & \sol{255.255.0.0} & \sol{\netDeux.255.254} & \Rnom{} G0/2 & \sol{\netSrv.254} & \sol{255.255.255.0} \\ \hline
  \hl \srvNom & \ip{\netSrv.\srvH} & \sol{255.255.255.0} & \sol{\netSrv.254} & \multicolumn{3}{c|}{\cellcolor{black!8}} \\ \hline
\end{tabularx}}
\bareme{Adresses d'hôtes libres (uniques, dans le bon réseau) ; affectation G0/x libre. Sur le .pkt : équipements 0,5 ;
câblage 1 (câbles droits vers les switchs et switch–routeur, \textbf{croisé} serveur–routeur, liens verts) ; nommage 0,5.}

\question{B3}{3 pts}{Configurez tous les équipements : IP, masques, passerelles des hôtes ; interfaces du routeur
(« Port Status : On »).}
\bareme{Sur le .pkt : interfaces routeur 1 pt ; IP/masques des hôtes 1 pt ; passerelles (imprimante et serveur compris)
1 pt. $-0,25$ par erreur, sans descendre sous 0 par item.}

\question{B4}{3 pts}{Réalisez les tests, notez la commande et le résultat, puis faites viser chaque test réussi.}
{\small\renewcommand{\arraystretch}{1.1}
\begin{tabularx}{\linewidth}{|c|l|l|C|c|c|}
  \hline
  \textbf{N°} & \textbf{Source} & \textbf{Destination} & \textbf{Commande / action} & \textbf{Résultat} & \textbf{Visa prof.} \\ \hline
  \hl 1 & PC-\pUn1 & PC-\pUn4 & \sol{ping \netUn.4} & \sol{réussi} & \hspace{1.4cm} \\ \hline
  \hl 2 & PC-\pDeux1 & IMP-\pDeux & \sol{ping \netDeux.0.10} & \sol{réussi} & \\ \hline
  \hl 3 & PC-\pUn2 & PC-\pDeux2 & \sol{ping \netDeux.0.2} & \sol{réussi} & \\ \hline
  \hl 4 & PC-\pDeux2 & \srvNom & \sol{ping \netSrv.\srvH} & \sol{réussi} & \\ \hline
  \hl 5 & PC-\pUn3 & page web du serveur & \sol{navigateur : {\NoAutoSpacing http://\netSrv.\srvH}} & \sol{page affichée} & \\ \hline
\end{tabularx}}
\bareme{0,5 pt par test visé + 0,5 pt pour B5. Le premier ping inter-réseaux peut échouer (résolution ARP) ; les suivants doivent réussir.}

\question{B5}{incl. B4}{Test n°3 : par quels équipements passe le paquet ? Pourquoi PC-\pUn2 l'envoie-t-il au routeur ?}
\zone{2.3cm}{PC-\pUn2 $\to$ SW-\pUn{} $\to$ \Rnom{} $\to$ SW-\pDeux{} $\to$ PC-\pDeux2. Grâce à son masque, PC-\pUn2 constate que
la destination (\netDeux.0.0/16) n'est \textbf{pas dans son réseau} (\netUn.0/24) : il envoie le paquet à sa
\textbf{passerelle par défaut}, qui le route.}

%==========================================================================
\newpage
\section{Partie C — Diagnostic d'un réseau en panne \hfill /4}

Un technicien constate qu'\textbf{aucun poste ne parvient à joindre le serveur SRV-C} et que certains postes
communiquent mal entre eux. Les câbles sont du bon type.

\begin{center}
\begin{minipage}{0.37\linewidth}\centering
\begin{tikzpicture}[font=\footnotesize,
  eq/.style={draw,fill=qbg,rounded corners=2pt,minimum width=9mm,minimum height=5mm,inner sep=2pt},
  every path/.style={thick}]
  \node[eq] (SW) {SW-1};
  \node[eq,right=9mm of SW] (R) {R1};
  \node[eq,right=7mm of R] (SRV) {SRV-C};
  \node[eq,above left=5mm and -3mm of SW] (A) {PC-A};
  \node[eq,above right=5mm and -3mm of SW] (B) {PC-B};
  \node[eq,below left=5mm and -3mm of SW] (C) {PC-C};
  \node[eq,below right=5mm and -3mm of SW] (D) {PC-D};
  \foreach \n in {A,B,C,D} \draw (\n) -- (SW);
  \draw (SW) -- node[above,font=\tiny]{G0/0} (R);
  \draw (R) -- node[above,font=\tiny]{G0/1} (SRV);
\end{tikzpicture}
\end{minipage}\hfill
\begin{minipage}{0.61\linewidth}\centering
{\small\renewcommand{\arraystretch}{1.05}\setlength{\tabcolsep}{4pt}
\begin{tabular}{lllll}
  \toprule
  \textbf{Équipement} & \textbf{Adresse IP} & \textbf{Masque} & \textbf{Passerelle} & \textbf{Port} \\
  \midrule
  \Crows
  \bottomrule
\end{tabular}}
\end{minipage}
\end{center}

\question{C1}{4 pts}{Relevez les \textbf{quatre anomalies}. Pour chacune : conséquence et correction.}
{\small\renewcommand{\arraystretch}{1.1}
\ifcorrige\def\hc{}\else\def\hc{\rule[-10mm]{0pt}{14mm}}\fi
\begin{tabularx}{\linewidth}{|c|>{\raggedright\arraybackslash}p{2cm}|L|L|L|}
  \hline
  \textbf{N°} & \textbf{Équipement} & \textbf{Anomalie} & \textbf{Conséquence} & \textbf{Correction} \\ \hline
  \ifcorrige\Csol\else
  \hc 1 & & & & \\ \hline
  \hc 2 & & & & \\ \hline
  \hc 3 & & & & \\ \hline
  \hc 4 & & & & \\ \hline
  \fi
\end{tabularx}}
\bareme{1 pt par anomalie (0,25 équipement, anomalie, conséquence, correction ; ordre indifférent). \Cbareme}

%==========================================================================
\section{Bonus — Extension sans fil \hfill +1}

\begin{cdc}Ajoutez au réseau \resUn{} un \textbf{AccessPoint-PT} relié à SW-\pUn{} (SSID \ip{\ssid}, \textbf{WPA2-PSK})
et un portable \ip{Laptop-\pUn} équipé d'un module \texttt{WPC300N}, en adressage \textbf{statique}.
Vérifiez qu'il joint \srvNom.\end{cdc}
{\small\renewcommand{\arraystretch}{1.1}
\begin{tabularx}{\linewidth}{|l|C|C|C|c|}
  \hline
  \textbf{Équipement} & \textbf{Adresse IP} & \textbf{Masque} & \textbf{Passerelle} & \textbf{Visa prof.} \\ \hline
  \hl Laptop-\pUn & \sol{\netUn.50} & \sol{255.255.255.0} & \sol{\netUn.254} & \hspace{1.4cm} \\ \hline
\end{tabularx}}
\bareme{Le point d'accès fait office de pont : le portable est dans \netUn.0/24. Bonus si associé en WPA2-PSK et ping
vers \netSrv.\srvH{} réussi. Note plafonnée à 20.}

\end{document}
"
%   pdflatex -jobname=Eval_Reseau_CorrigeA "\def\variante{A}\def\modecorrige{1}% Évaluation pratique — Réseaux locaux (BTS CIEL 1re année) — sujets A et B
% Compilation (4 PDF) :
%   pdflatex -jobname=Eval_Reseau_SujetA   "\def\variante{A}\def\modecorrige{0}% Évaluation pratique — Réseaux locaux (BTS CIEL 1re année) — sujets A et B
% Compilation (4 PDF) :
%   pdflatex -jobname=Eval_Reseau_SujetA   "\def\variante{A}\def\modecorrige{0}\input{Eval_TP_Reseau_LAN.tex}"
%   pdflatex -jobname=Eval_Reseau_SujetB   "\def\variante{B}\def\modecorrige{0}\input{Eval_TP_Reseau_LAN.tex}"
%   pdflatex -jobname=Eval_Reseau_CorrigeA "\def\variante{A}\def\modecorrige{1}\input{Eval_TP_Reseau_LAN.tex}"
%   pdflatex -jobname=Eval_Reseau_CorrigeB "\def\variante{B}\def\modecorrige{1}\input{Eval_TP_Reseau_LAN.tex}"
\documentclass[10pt,a4paper]{article}

\usepackage[T1]{fontenc}
\usepackage[utf8]{inputenc}
\usepackage{lmodern}
\usepackage[french]{babel}
\usepackage[a4paper,hmargin=1.5cm,top=1.3cm,bottom=1.6cm,footskip=0.7cm]{geometry}
\usepackage[table]{xcolor}
\usepackage{booktabs,tabularx,array}
\usepackage{enumitem}
\usepackage{fancyhdr}
\usepackage[most]{tcolorbox}
\usepackage{tikz}
\usetikzlibrary{positioning}
\usepackage{titlesec}

\providecommand{\modecorrige}{0}
\providecommand{\variante}{A}
\newif\ifcorrige
\ifnum\modecorrige=1 \corrigetrue \else \corrigefalse \fi

\definecolor{ipred}{HTML}{C0392B}
\definecolor{cdcbar}{HTML}{A67C3D}
\definecolor{cdcbg}{HTML}{F4F1EA}
\definecolor{qbar}{HTML}{4A6B5D}
\definecolor{qbg}{HTML}{EDF1EF}
\definecolor{corr}{HTML}{1F5FA8}
\definecolor{corrbg}{HTML}{E8F0FA}
\definecolor{grisfoot}{HTML}{828282}

\newcommand{\ip}[1]{\textcolor{ipred}{\texttt{#1}}}
\newcommand{\pts}[1]{\hfill\mbox{\textbf{(#1)}}}
\newcommand{\sol}[1]{\ifcorrige\textcolor{corr}{#1}\fi}
\newcommand{\bareme}[1]{\ifcorrige\par{\footnotesize\color{corr}#1\par}\fi}
\newcommand{\hl}{\ifcorrige\rule[-1.6mm]{0pt}{5.4mm}\else\rule[-2.2mm]{0pt}{6.6mm}\fi}
\newcolumntype{C}{>{\centering\arraybackslash}X}
\newcolumntype{L}{>{\raggedright\arraybackslash}X}

\titleformat{\section}{\large\bfseries}{}{0pt}{}
\titlespacing*{\section}{0pt}{6pt}{3pt}

\tcbset{qstyle/.style={enhanced,boxrule=0pt,leftrule=2.5pt,arc=0pt,outer arc=0pt,
  left=5pt,right=5pt,top=2pt,bottom=2pt,before skip=4pt,after skip=3pt}}
\newtcolorbox{cdc}{qstyle,colback=cdcbg,colframe=cdcbar}
\newtcolorbox{quest}{qstyle,colback=qbg,colframe=qbar}

% Zone de réponse : vierge dans le sujet, remplie dans le corrigé
\newcommand{\zone}[2]{%
  \ifcorrige
    \begin{tcolorbox}[colback=corrbg,colframe=corr,boxrule=0.5pt,arc=1pt,left=4pt,right=4pt,top=1pt,bottom=1pt,
      before skip=2pt,after skip=2pt,fontupper=\footnotesize\color{corr}]#2\end{tcolorbox}%
  \else
    \begin{tcolorbox}[colback=white,colframe=black!45,boxrule=0.4pt,arc=1pt,height=#1,
      before skip=2pt,after skip=2pt]\end{tcolorbox}%
  \fi}

\newcommand{\question}[3]{\begin{quest}\textbf{#1.} #3\pts{#2}\end{quest}}

%==========================================================================
% Données propres à chaque variante
%==========================================================================
\if\variante A
  \def\entreprise{Le cabinet d'architectes \textbf{ArchiNova}}
  \def\Rnom{R-ArchiNova}
  \def\resUn{Bureau d'études}\def\pUn{BE}\def\netUn{192.168.40}
  \def\resDeux{Accueil}\def\pDeux{ACC}\def\netDeux{172.16}
  \def\srvNom{SRV-Intranet}\def\netSrv{192.168.200}\def\srvH{10}
  \def\ssid{ArchiNova-WiFi}
  \def\AIrows{%
    \hl\ip{192.168.40.25} & \sol{C} & \sol{255.255.255.0} & \sol{192.168.40.0} & \sol{192.168.40.255} & \sol{254} \\ \hline
    \hl\ip{172.16.5.10}   & \sol{B} & \sol{255.255.0.0}   & \sol{172.16.0.0}   & \sol{172.16.255.255} & \sol{65\,534} \\ \hline
    \hl\ip{10.3.2.1}      & \sol{A} & \sol{255.0.0.0}     & \sol{10.0.0.0}     & \sol{10.255.255.255} & \sol{16\,777\,214} \\ \hline}
  \def\apipa{169.254.12.7}
  \def\Cnet{192.168.60}
  \def\Crows{%
    PC-A & \ip{192.168.60.10} & \ip{255.255.255.0} & \ip{192.168.60.254} & — \\
    PC-B & \ip{192.168.60.20} & \ip{255.255.255.0} & \ip{192.168.60.1}   & — \\
    PC-C & \ip{192.168.61.30} & \ip{255.255.255.0} & \ip{192.168.60.254} & — \\
    PC-D & \ip{192.168.60.10} & \ip{255.255.255.0} & \ip{192.168.60.254} & — \\
    SRV-C & \ip{10.0.0.5}     & \ip{255.0.0.0}     & \ip{10.0.0.254}     & — \\
    R1 — G0/0 & \ip{192.168.60.254} & \ip{255.255.255.0} & — & On \\
    R1 — G0/1 & \ip{10.0.0.254}     & \ip{255.0.0.0}     & — & Off \\}
  \def\Csol{%
    \hl 1 & \sol{PC-B} & \sol{Passerelle .60.1 : pas l'adresse du routeur.} & \sol{Joint son réseau, mais aucun autre (serveur).} & \sol{Passerelle 192.168.60.254.} \\ \hline
    \hl 2 & \sol{PC-C} & \sol{IP hors du réseau 192.168.60.0/24.} & \sol{Ne joint ni les postes ni le routeur.} & \sol{Ex. 192.168.60.30.} \\ \hline
    \hl 3 & \sol{PC-D} & \sol{Même IP que PC-A (doublon).} & \sol{Conflit d'adresse : PC-A et PC-D perturbés.} & \sol{Adresse unique, ex. .60.40.} \\ \hline
    \hl 4 & \sol{R1 G0/1} & \sol{Interface désactivée (Off).} & \sol{Réseau du serveur injoignable pour tous.} & \sol{Port Status : On.} \\ \hline}
  \def\Cbareme{PC-A et SRV-C sont correctement configurés ; accepter PC-A cité à la place de PC-D pour le doublon.}
\else
  \def\entreprise{La clinique vétérinaire \textbf{VetoPlus}}
  \def\Rnom{R-VetoPlus}
  \def\resUn{Soins}\def\pUn{SOI}\def\netUn{192.168.70}
  \def\resDeux{Secrétariat}\def\pDeux{SEC}\def\netDeux{172.20}
  \def\srvNom{SRV-Dossiers}\def\netSrv{192.168.250}\def\srvH{20}
  \def\ssid{VetoPlus-WiFi}
  \def\AIrows{%
    \hl\ip{192.168.15.130} & \sol{C} & \sol{255.255.255.0} & \sol{192.168.15.0} & \sol{192.168.15.255} & \sol{254} \\ \hline
    \hl\ip{172.31.200.7}   & \sol{B} & \sol{255.255.0.0}   & \sol{172.31.0.0}   & \sol{172.31.255.255} & \sol{65\,534} \\ \hline
    \hl\ip{10.45.6.9}      & \sol{A} & \sol{255.0.0.0}     & \sol{10.0.0.0}     & \sol{10.255.255.255} & \sol{16\,777\,214} \\ \hline}
  \def\apipa{169.254.88.3}
  \def\Cnet{192.168.80}
  \def\Crows{%
    PC-A & \ip{192.168.8.11}   & \ip{255.255.255.0} & \ip{192.168.80.254} & — \\
    PC-B & \ip{192.168.80.12}  & \ip{255.255.255.0} & \ip{192.168.80.254} & — \\
    PC-C & \ip{192.168.80.13}  & \ip{255.255.255.0} & \ip{192.168.80.100} & — \\
    PC-D & \ip{192.168.80.12}  & \ip{255.255.255.0} & \ip{192.168.80.254} & — \\
    SRV-C & \ip{10.0.0.8}      & \ip{255.0.0.0}     & \ip{10.0.0.254}     & — \\
    R1 — G0/0 & \ip{192.168.80.254} & \ip{255.255.255.0} & — & Off \\
    R1 — G0/1 & \ip{10.0.0.254}     & \ip{255.0.0.0}     & — & On \\}
  \def\Csol{%
    \hl 1 & \sol{PC-A} & \sol{IP hors du réseau 192.168.80.0/24.} & \sol{Ne joint ni les postes ni le routeur.} & \sol{Ex. 192.168.80.11.} \\ \hline
    \hl 2 & \sol{PC-C} & \sol{Passerelle .80.100 : pas l'adresse du routeur.} & \sol{Joint son réseau, mais aucun autre (serveur).} & \sol{Passerelle 192.168.80.254.} \\ \hline
    \hl 3 & \sol{PC-D} & \sol{Même IP que PC-B (doublon).} & \sol{Conflit d'adresse : PC-B et PC-D perturbés.} & \sol{Adresse unique, ex. .80.14.} \\ \hline
    \hl 4 & \sol{R1 G0/0} & \sol{Interface côté LAN désactivée (Off).} & \sol{Aucun poste ne joint le routeur ni le serveur.} & \sol{Port Status : On.} \\ \hline}
  \def\Cbareme{PC-B et SRV-C sont correctement configurés ; accepter PC-B cité à la place de PC-D pour le doublon.}
\fi

\pagestyle{fancy}
\fancyhf{}
\renewcommand{\headrulewidth}{0pt}
\fancyfoot[L]{\footnotesize\itshape\color{grisfoot}Évaluation — Réseaux locaux — Sujet \variante\ifcorrige\ — CORRIGÉ\fi}
\fancyfoot[R]{\footnotesize\color{grisfoot}\thepage/3}

\setlist{itemsep=0pt,topsep=2pt,parsep=0pt}
\setlength{\parindent}{0pt}
\setlength{\parskip}{2pt}

\begin{document}

\begin{center}
  {\Large\bfseries Évaluation pratique — Réseaux locaux \quad
   \fcolorbox{black}{white}{\ SUJET \variante\ }\par}
  \vspace{1pt}
  {\bfseries Adressage IP, commutation et routage — BTS CIEL 1re année — durée : 2 heures\par}
  \ifcorrige\vspace{2pt}{\bfseries\color{corr}CORRIGÉ ET BARÈME — document enseignant}\par\fi
\end{center}
\vspace{-2pt}
{\renewcommand{\arraystretch}{1.35}
\begin{tabularx}{\linewidth}{|X|X|p{2.2cm}|p{2.6cm}|}
  \hline
  \textbf{Nom :} & \textbf{Prénom :} & \textbf{Poste n°} & \textbf{Note :} \hfill /20 \\ \hline
\end{tabularx}}

\textbf{Consignes :} travail \textbf{individuel}, documents et Internet \textbf{non autorisés}. Ce sujet est le
\textbf{document-réponse}. Fichier Packet Tracer : \ip{Nom\_Prenom\_EvalLAN\_\variante.pkt}, à enregistrer
régulièrement dans le dossier de rendu. Les « visas » sont apposés par le professeur, que vous appelez à chaque test réussi.
Durées conseillées : A 30 min — B 1 h 10 — C 20 min.

%==========================================================================
\section{Partie A — Questions de cours et calculs \hfill /6}

\question{A1}{2 pts}{Complétez le tableau en considérant le \textbf{masque par défaut} de la classe de chaque adresse.}
{\small\renewcommand{\arraystretch}{1.15}
\begin{tabularx}{\linewidth}{|l|c|C|C|C|c|}
  \hline
  \textbf{Adresse IP} & \textbf{Classe} & \textbf{Masque par défaut} & \textbf{Adresse réseau} &
  \textbf{Adresse de diffusion} & \textbf{Nb d'hôtes} \\ \hline
  \AIrows
\end{tabularx}}
\bareme{0,5 pt par ligne juste + 0,5 pt si les trois nombres d'hôtes sont justes ($2^{n}-2$, $n$ = nombre de bits d'hôte).}

\question{A2}{1 pt}{Indiquez le type de câble (\textbf{droit} ou \textbf{croisé}) pour chaque liaison.}
{\small\renewcommand{\arraystretch}{1.15}
\begin{tabularx}{\linewidth}{|X|c||X|c|}
  \hline
  \hl PC $\leftrightarrow$ commutateur & \makebox[1.6cm]{\sol{droit}} & Commutateur $\leftrightarrow$ routeur & \makebox[1.6cm]{\sol{droit}} \\ \hline
  \hl PC $\leftrightarrow$ PC (réseau P2P) & \sol{croisé} & Concentrateur $\leftrightarrow$ concentrateur & \sol{croisé} \\ \hline
\end{tabularx}}
\bareme{0,25 pt par case. Équipements de types différents : câble droit ; de même type : câble croisé.}

\question{A3}{1 pt}{Comment un \textbf{concentrateur (hub)} et un \textbf{commutateur (switch)} traitent-ils une trame
reçue ? Lequel limite les collisions, et pourquoi ?}
\zone{2.9cm}{Le \textbf{hub} (couche 1) répète la trame sur \textbf{tous ses autres ports} : un seul domaine de collision
partagé. Le \textbf{switch} (couche 2) lit l'adresse MAC de destination et, grâce à sa \textbf{table MAC}, ne l'envoie
\textbf{que sur le port du destinataire} : un domaine de collision par port, c'est lui qui limite les collisions.
\emph{0,5 pt hub, 0,5 pt switch.}}

\question{A4}{1 pt}{Qu'est-ce que la \textbf{passerelle par défaut} d'un poste ? Quelle condition son adresse doit-elle respecter ?}
\zone{2.4cm}{Adresse IP de l'interface du \textbf{routeur} à laquelle le poste envoie les paquets destinés à un
\textbf{autre réseau}. Elle doit \textbf{appartenir au même réseau} que le poste. \emph{0,5 + 0,5 pt.}}

\question{A5}{0,5 pt}{Un portable configuré en DHCP affiche \ip{\apipa} / \ip{255.255.0.0}. Que s'est-il passé ?}
\zone{1.5cm}{Aucun \textbf{serveur DHCP} n'a répondu : le poste s'est attribué une adresse \textbf{APIPA} (169.254.0.0/16).}

\question{A6}{0,5 pt}{Quel protocole utilise \texttt{ping} ? Nommez les deux messages échangés et la couche OSI des adresses IP.}
\zone{1.5cm}{\textbf{ICMP} ; \textbf{Echo Request} / \textbf{Echo Reply} ; adresses IP en \textbf{couche 3 (Réseau)}.}

%==========================================================================
\newpage
\section{Partie B — Réalisation sous Packet Tracer \hfill /10}

\begin{cdc}
\entreprise{} vous confie son réseau : \textbf{trois réseaux} reliés par un \textbf{routeur unique Cisco 2911}
nommé \ip{\Rnom} (interfaces G0/0, G0/1, G0/2).
\begin{itemize}[leftmargin=12pt]
  \item \textbf{Réseau \resUn} : \textbf{4 postes} PC-\pUn1 à PC-\pUn4, commutateur \ip{SW-\pUn}, adressage statique,
        réseau \ip{\netUn.0/24}.
  \item \textbf{Réseau \resDeux} : \textbf{2 postes} PC-\pDeux1, PC-\pDeux2 et \textbf{1 imprimante} IMP-\pDeux,
        commutateur \ip{SW-\pDeux}, adressage statique, réseau de classe B \ip{\netDeux.0.0/16}.
  \item \textbf{Réseau Serveur} : serveur \ip{\srvNom} d'adresse \ip{\netSrv.\srvH/24}, relié \textbf{directement} au routeur.
  \item La passerelle de chaque réseau est sa \textbf{dernière adresse utilisable}. Tous les équipements communiquent
        entre eux et accèdent à la page web du serveur. Nommage imposé (« Display Name »).
\end{itemize}
\end{cdc}

\question{B1}{2 pts}{Complétez le plan d'adressage \textbf{avant} la réalisation.}
{\footnotesize\setlength{\tabcolsep}{3pt}\renewcommand{\arraystretch}{1.15}
\begin{tabularx}{\linewidth}{|l|C|C|C|C|C|C|}
  \hline
  \textbf{Réseau} & \textbf{Adresse réseau} & \textbf{Masque} & \textbf{1re @ utilisable} &
  \textbf{Dernière @ utilisable} & \textbf{Diffusion} & \textbf{Passerelle} \\ \hline
  \hl \resUn & \sol{\netUn.0} & \sol{255.255.255.0} & \sol{\netUn.1} & \sol{\netUn.254} & \sol{\netUn.255} & \sol{\netUn.254} \\ \hline
  \hl \resDeux & \sol{\netDeux.0.0} & \sol{255.255.0.0} & \sol{\netDeux.0.1} & \sol{\netDeux.255.254} & \sol{\netDeux.255.255} & \sol{\netDeux.255.254} \\ \hline
  \hl Serveur & \sol{\netSrv.0} & \sol{255.255.255.0} & \sol{\netSrv.1} & \sol{\netSrv.254} & \sol{\netSrv.255} & \sol{\netSrv.254} \\ \hline
\end{tabularx}}
\bareme{0,5 pt par ligne juste + 0,5 pt si la convention de passerelle est respectée partout.}

\question{B2}{2 pts}{Construisez la topologie (équipements, câblage, nommage) et relevez l'adressage configuré.}
{\footnotesize\setlength{\tabcolsep}{3pt}\renewcommand{\arraystretch}{1.1}
\begin{tabularx}{\linewidth}{|l|C|C|C||l|C|C|}
  \hline
  \textbf{Hôte} & \textbf{Adresse IP} & \textbf{Masque} & \textbf{Passerelle} & \textbf{Interface} & \textbf{Adresse IP} & \textbf{Masque} \\ \hline
  \hl PC-\pUn1 & \sol{\netUn.1} & \sol{255.255.255.0} & \sol{\netUn.254} & \Rnom{} G0/0 & \sol{\netUn.254} & \sol{255.255.255.0} \\ \hline
  \hl PC-\pDeux1 & \sol{\netDeux.0.1} & \sol{255.255.0.0} & \sol{\netDeux.255.254} & \Rnom{} G0/1 & \sol{\netDeux.255.254} & \sol{255.255.0.0} \\ \hline
  \hl IMP-\pDeux & \sol{\netDeux.0.10} & \sol{255.255.0.0} & \sol{\netDeux.255.254} & \Rnom{} G0/2 & \sol{\netSrv.254} & \sol{255.255.255.0} \\ \hline
  \hl \srvNom & \ip{\netSrv.\srvH} & \sol{255.255.255.0} & \sol{\netSrv.254} & \multicolumn{3}{c|}{\cellcolor{black!8}} \\ \hline
\end{tabularx}}
\bareme{Adresses d'hôtes libres (uniques, dans le bon réseau) ; affectation G0/x libre. Sur le .pkt : équipements 0,5 ;
câblage 1 (câbles droits vers les switchs et switch–routeur, \textbf{croisé} serveur–routeur, liens verts) ; nommage 0,5.}

\question{B3}{3 pts}{Configurez tous les équipements : IP, masques, passerelles des hôtes ; interfaces du routeur
(« Port Status : On »).}
\bareme{Sur le .pkt : interfaces routeur 1 pt ; IP/masques des hôtes 1 pt ; passerelles (imprimante et serveur compris)
1 pt. $-0,25$ par erreur, sans descendre sous 0 par item.}

\question{B4}{3 pts}{Réalisez les tests, notez la commande et le résultat, puis faites viser chaque test réussi.}
{\small\renewcommand{\arraystretch}{1.1}
\begin{tabularx}{\linewidth}{|c|l|l|C|c|c|}
  \hline
  \textbf{N°} & \textbf{Source} & \textbf{Destination} & \textbf{Commande / action} & \textbf{Résultat} & \textbf{Visa prof.} \\ \hline
  \hl 1 & PC-\pUn1 & PC-\pUn4 & \sol{ping \netUn.4} & \sol{réussi} & \hspace{1.4cm} \\ \hline
  \hl 2 & PC-\pDeux1 & IMP-\pDeux & \sol{ping \netDeux.0.10} & \sol{réussi} & \\ \hline
  \hl 3 & PC-\pUn2 & PC-\pDeux2 & \sol{ping \netDeux.0.2} & \sol{réussi} & \\ \hline
  \hl 4 & PC-\pDeux2 & \srvNom & \sol{ping \netSrv.\srvH} & \sol{réussi} & \\ \hline
  \hl 5 & PC-\pUn3 & page web du serveur & \sol{navigateur : {\NoAutoSpacing http://\netSrv.\srvH}} & \sol{page affichée} & \\ \hline
\end{tabularx}}
\bareme{0,5 pt par test visé + 0,5 pt pour B5. Le premier ping inter-réseaux peut échouer (résolution ARP) ; les suivants doivent réussir.}

\question{B5}{incl. B4}{Test n°3 : par quels équipements passe le paquet ? Pourquoi PC-\pUn2 l'envoie-t-il au routeur ?}
\zone{2.3cm}{PC-\pUn2 $\to$ SW-\pUn{} $\to$ \Rnom{} $\to$ SW-\pDeux{} $\to$ PC-\pDeux2. Grâce à son masque, PC-\pUn2 constate que
la destination (\netDeux.0.0/16) n'est \textbf{pas dans son réseau} (\netUn.0/24) : il envoie le paquet à sa
\textbf{passerelle par défaut}, qui le route.}

%==========================================================================
\newpage
\section{Partie C — Diagnostic d'un réseau en panne \hfill /4}

Un technicien constate qu'\textbf{aucun poste ne parvient à joindre le serveur SRV-C} et que certains postes
communiquent mal entre eux. Les câbles sont du bon type.

\begin{center}
\begin{minipage}{0.37\linewidth}\centering
\begin{tikzpicture}[font=\footnotesize,
  eq/.style={draw,fill=qbg,rounded corners=2pt,minimum width=9mm,minimum height=5mm,inner sep=2pt},
  every path/.style={thick}]
  \node[eq] (SW) {SW-1};
  \node[eq,right=9mm of SW] (R) {R1};
  \node[eq,right=7mm of R] (SRV) {SRV-C};
  \node[eq,above left=5mm and -3mm of SW] (A) {PC-A};
  \node[eq,above right=5mm and -3mm of SW] (B) {PC-B};
  \node[eq,below left=5mm and -3mm of SW] (C) {PC-C};
  \node[eq,below right=5mm and -3mm of SW] (D) {PC-D};
  \foreach \n in {A,B,C,D} \draw (\n) -- (SW);
  \draw (SW) -- node[above,font=\tiny]{G0/0} (R);
  \draw (R) -- node[above,font=\tiny]{G0/1} (SRV);
\end{tikzpicture}
\end{minipage}\hfill
\begin{minipage}{0.61\linewidth}\centering
{\small\renewcommand{\arraystretch}{1.05}\setlength{\tabcolsep}{4pt}
\begin{tabular}{lllll}
  \toprule
  \textbf{Équipement} & \textbf{Adresse IP} & \textbf{Masque} & \textbf{Passerelle} & \textbf{Port} \\
  \midrule
  \Crows
  \bottomrule
\end{tabular}}
\end{minipage}
\end{center}

\question{C1}{4 pts}{Relevez les \textbf{quatre anomalies}. Pour chacune : conséquence et correction.}
{\small\renewcommand{\arraystretch}{1.1}
\ifcorrige\def\hc{}\else\def\hc{\rule[-10mm]{0pt}{14mm}}\fi
\begin{tabularx}{\linewidth}{|c|>{\raggedright\arraybackslash}p{2cm}|L|L|L|}
  \hline
  \textbf{N°} & \textbf{Équipement} & \textbf{Anomalie} & \textbf{Conséquence} & \textbf{Correction} \\ \hline
  \ifcorrige\Csol\else
  \hc 1 & & & & \\ \hline
  \hc 2 & & & & \\ \hline
  \hc 3 & & & & \\ \hline
  \hc 4 & & & & \\ \hline
  \fi
\end{tabularx}}
\bareme{1 pt par anomalie (0,25 équipement, anomalie, conséquence, correction ; ordre indifférent). \Cbareme}

%==========================================================================
\section{Bonus — Extension sans fil \hfill +1}

\begin{cdc}Ajoutez au réseau \resUn{} un \textbf{AccessPoint-PT} relié à SW-\pUn{} (SSID \ip{\ssid}, \textbf{WPA2-PSK})
et un portable \ip{Laptop-\pUn} équipé d'un module \texttt{WPC300N}, en adressage \textbf{statique}.
Vérifiez qu'il joint \srvNom.\end{cdc}
{\small\renewcommand{\arraystretch}{1.1}
\begin{tabularx}{\linewidth}{|l|C|C|C|c|}
  \hline
  \textbf{Équipement} & \textbf{Adresse IP} & \textbf{Masque} & \textbf{Passerelle} & \textbf{Visa prof.} \\ \hline
  \hl Laptop-\pUn & \sol{\netUn.50} & \sol{255.255.255.0} & \sol{\netUn.254} & \hspace{1.4cm} \\ \hline
\end{tabularx}}
\bareme{Le point d'accès fait office de pont : le portable est dans \netUn.0/24. Bonus si associé en WPA2-PSK et ping
vers \netSrv.\srvH{} réussi. Note plafonnée à 20.}

\end{document}
"
%   pdflatex -jobname=Eval_Reseau_SujetB   "\def\variante{B}\def\modecorrige{0}% Évaluation pratique — Réseaux locaux (BTS CIEL 1re année) — sujets A et B
% Compilation (4 PDF) :
%   pdflatex -jobname=Eval_Reseau_SujetA   "\def\variante{A}\def\modecorrige{0}\input{Eval_TP_Reseau_LAN.tex}"
%   pdflatex -jobname=Eval_Reseau_SujetB   "\def\variante{B}\def\modecorrige{0}\input{Eval_TP_Reseau_LAN.tex}"
%   pdflatex -jobname=Eval_Reseau_CorrigeA "\def\variante{A}\def\modecorrige{1}\input{Eval_TP_Reseau_LAN.tex}"
%   pdflatex -jobname=Eval_Reseau_CorrigeB "\def\variante{B}\def\modecorrige{1}\input{Eval_TP_Reseau_LAN.tex}"
\documentclass[10pt,a4paper]{article}

\usepackage[T1]{fontenc}
\usepackage[utf8]{inputenc}
\usepackage{lmodern}
\usepackage[french]{babel}
\usepackage[a4paper,hmargin=1.5cm,top=1.3cm,bottom=1.6cm,footskip=0.7cm]{geometry}
\usepackage[table]{xcolor}
\usepackage{booktabs,tabularx,array}
\usepackage{enumitem}
\usepackage{fancyhdr}
\usepackage[most]{tcolorbox}
\usepackage{tikz}
\usetikzlibrary{positioning}
\usepackage{titlesec}

\providecommand{\modecorrige}{0}
\providecommand{\variante}{A}
\newif\ifcorrige
\ifnum\modecorrige=1 \corrigetrue \else \corrigefalse \fi

\definecolor{ipred}{HTML}{C0392B}
\definecolor{cdcbar}{HTML}{A67C3D}
\definecolor{cdcbg}{HTML}{F4F1EA}
\definecolor{qbar}{HTML}{4A6B5D}
\definecolor{qbg}{HTML}{EDF1EF}
\definecolor{corr}{HTML}{1F5FA8}
\definecolor{corrbg}{HTML}{E8F0FA}
\definecolor{grisfoot}{HTML}{828282}

\newcommand{\ip}[1]{\textcolor{ipred}{\texttt{#1}}}
\newcommand{\pts}[1]{\hfill\mbox{\textbf{(#1)}}}
\newcommand{\sol}[1]{\ifcorrige\textcolor{corr}{#1}\fi}
\newcommand{\bareme}[1]{\ifcorrige\par{\footnotesize\color{corr}#1\par}\fi}
\newcommand{\hl}{\ifcorrige\rule[-1.6mm]{0pt}{5.4mm}\else\rule[-2.2mm]{0pt}{6.6mm}\fi}
\newcolumntype{C}{>{\centering\arraybackslash}X}
\newcolumntype{L}{>{\raggedright\arraybackslash}X}

\titleformat{\section}{\large\bfseries}{}{0pt}{}
\titlespacing*{\section}{0pt}{6pt}{3pt}

\tcbset{qstyle/.style={enhanced,boxrule=0pt,leftrule=2.5pt,arc=0pt,outer arc=0pt,
  left=5pt,right=5pt,top=2pt,bottom=2pt,before skip=4pt,after skip=3pt}}
\newtcolorbox{cdc}{qstyle,colback=cdcbg,colframe=cdcbar}
\newtcolorbox{quest}{qstyle,colback=qbg,colframe=qbar}

% Zone de réponse : vierge dans le sujet, remplie dans le corrigé
\newcommand{\zone}[2]{%
  \ifcorrige
    \begin{tcolorbox}[colback=corrbg,colframe=corr,boxrule=0.5pt,arc=1pt,left=4pt,right=4pt,top=1pt,bottom=1pt,
      before skip=2pt,after skip=2pt,fontupper=\footnotesize\color{corr}]#2\end{tcolorbox}%
  \else
    \begin{tcolorbox}[colback=white,colframe=black!45,boxrule=0.4pt,arc=1pt,height=#1,
      before skip=2pt,after skip=2pt]\end{tcolorbox}%
  \fi}

\newcommand{\question}[3]{\begin{quest}\textbf{#1.} #3\pts{#2}\end{quest}}

%==========================================================================
% Données propres à chaque variante
%==========================================================================
\if\variante A
  \def\entreprise{Le cabinet d'architectes \textbf{ArchiNova}}
  \def\Rnom{R-ArchiNova}
  \def\resUn{Bureau d'études}\def\pUn{BE}\def\netUn{192.168.40}
  \def\resDeux{Accueil}\def\pDeux{ACC}\def\netDeux{172.16}
  \def\srvNom{SRV-Intranet}\def\netSrv{192.168.200}\def\srvH{10}
  \def\ssid{ArchiNova-WiFi}
  \def\AIrows{%
    \hl\ip{192.168.40.25} & \sol{C} & \sol{255.255.255.0} & \sol{192.168.40.0} & \sol{192.168.40.255} & \sol{254} \\ \hline
    \hl\ip{172.16.5.10}   & \sol{B} & \sol{255.255.0.0}   & \sol{172.16.0.0}   & \sol{172.16.255.255} & \sol{65\,534} \\ \hline
    \hl\ip{10.3.2.1}      & \sol{A} & \sol{255.0.0.0}     & \sol{10.0.0.0}     & \sol{10.255.255.255} & \sol{16\,777\,214} \\ \hline}
  \def\apipa{169.254.12.7}
  \def\Cnet{192.168.60}
  \def\Crows{%
    PC-A & \ip{192.168.60.10} & \ip{255.255.255.0} & \ip{192.168.60.254} & — \\
    PC-B & \ip{192.168.60.20} & \ip{255.255.255.0} & \ip{192.168.60.1}   & — \\
    PC-C & \ip{192.168.61.30} & \ip{255.255.255.0} & \ip{192.168.60.254} & — \\
    PC-D & \ip{192.168.60.10} & \ip{255.255.255.0} & \ip{192.168.60.254} & — \\
    SRV-C & \ip{10.0.0.5}     & \ip{255.0.0.0}     & \ip{10.0.0.254}     & — \\
    R1 — G0/0 & \ip{192.168.60.254} & \ip{255.255.255.0} & — & On \\
    R1 — G0/1 & \ip{10.0.0.254}     & \ip{255.0.0.0}     & — & Off \\}
  \def\Csol{%
    \hl 1 & \sol{PC-B} & \sol{Passerelle .60.1 : pas l'adresse du routeur.} & \sol{Joint son réseau, mais aucun autre (serveur).} & \sol{Passerelle 192.168.60.254.} \\ \hline
    \hl 2 & \sol{PC-C} & \sol{IP hors du réseau 192.168.60.0/24.} & \sol{Ne joint ni les postes ni le routeur.} & \sol{Ex. 192.168.60.30.} \\ \hline
    \hl 3 & \sol{PC-D} & \sol{Même IP que PC-A (doublon).} & \sol{Conflit d'adresse : PC-A et PC-D perturbés.} & \sol{Adresse unique, ex. .60.40.} \\ \hline
    \hl 4 & \sol{R1 G0/1} & \sol{Interface désactivée (Off).} & \sol{Réseau du serveur injoignable pour tous.} & \sol{Port Status : On.} \\ \hline}
  \def\Cbareme{PC-A et SRV-C sont correctement configurés ; accepter PC-A cité à la place de PC-D pour le doublon.}
\else
  \def\entreprise{La clinique vétérinaire \textbf{VetoPlus}}
  \def\Rnom{R-VetoPlus}
  \def\resUn{Soins}\def\pUn{SOI}\def\netUn{192.168.70}
  \def\resDeux{Secrétariat}\def\pDeux{SEC}\def\netDeux{172.20}
  \def\srvNom{SRV-Dossiers}\def\netSrv{192.168.250}\def\srvH{20}
  \def\ssid{VetoPlus-WiFi}
  \def\AIrows{%
    \hl\ip{192.168.15.130} & \sol{C} & \sol{255.255.255.0} & \sol{192.168.15.0} & \sol{192.168.15.255} & \sol{254} \\ \hline
    \hl\ip{172.31.200.7}   & \sol{B} & \sol{255.255.0.0}   & \sol{172.31.0.0}   & \sol{172.31.255.255} & \sol{65\,534} \\ \hline
    \hl\ip{10.45.6.9}      & \sol{A} & \sol{255.0.0.0}     & \sol{10.0.0.0}     & \sol{10.255.255.255} & \sol{16\,777\,214} \\ \hline}
  \def\apipa{169.254.88.3}
  \def\Cnet{192.168.80}
  \def\Crows{%
    PC-A & \ip{192.168.8.11}   & \ip{255.255.255.0} & \ip{192.168.80.254} & — \\
    PC-B & \ip{192.168.80.12}  & \ip{255.255.255.0} & \ip{192.168.80.254} & — \\
    PC-C & \ip{192.168.80.13}  & \ip{255.255.255.0} & \ip{192.168.80.100} & — \\
    PC-D & \ip{192.168.80.12}  & \ip{255.255.255.0} & \ip{192.168.80.254} & — \\
    SRV-C & \ip{10.0.0.8}      & \ip{255.0.0.0}     & \ip{10.0.0.254}     & — \\
    R1 — G0/0 & \ip{192.168.80.254} & \ip{255.255.255.0} & — & Off \\
    R1 — G0/1 & \ip{10.0.0.254}     & \ip{255.0.0.0}     & — & On \\}
  \def\Csol{%
    \hl 1 & \sol{PC-A} & \sol{IP hors du réseau 192.168.80.0/24.} & \sol{Ne joint ni les postes ni le routeur.} & \sol{Ex. 192.168.80.11.} \\ \hline
    \hl 2 & \sol{PC-C} & \sol{Passerelle .80.100 : pas l'adresse du routeur.} & \sol{Joint son réseau, mais aucun autre (serveur).} & \sol{Passerelle 192.168.80.254.} \\ \hline
    \hl 3 & \sol{PC-D} & \sol{Même IP que PC-B (doublon).} & \sol{Conflit d'adresse : PC-B et PC-D perturbés.} & \sol{Adresse unique, ex. .80.14.} \\ \hline
    \hl 4 & \sol{R1 G0/0} & \sol{Interface côté LAN désactivée (Off).} & \sol{Aucun poste ne joint le routeur ni le serveur.} & \sol{Port Status : On.} \\ \hline}
  \def\Cbareme{PC-B et SRV-C sont correctement configurés ; accepter PC-B cité à la place de PC-D pour le doublon.}
\fi

\pagestyle{fancy}
\fancyhf{}
\renewcommand{\headrulewidth}{0pt}
\fancyfoot[L]{\footnotesize\itshape\color{grisfoot}Évaluation — Réseaux locaux — Sujet \variante\ifcorrige\ — CORRIGÉ\fi}
\fancyfoot[R]{\footnotesize\color{grisfoot}\thepage/3}

\setlist{itemsep=0pt,topsep=2pt,parsep=0pt}
\setlength{\parindent}{0pt}
\setlength{\parskip}{2pt}

\begin{document}

\begin{center}
  {\Large\bfseries Évaluation pratique — Réseaux locaux \quad
   \fcolorbox{black}{white}{\ SUJET \variante\ }\par}
  \vspace{1pt}
  {\bfseries Adressage IP, commutation et routage — BTS CIEL 1re année — durée : 2 heures\par}
  \ifcorrige\vspace{2pt}{\bfseries\color{corr}CORRIGÉ ET BARÈME — document enseignant}\par\fi
\end{center}
\vspace{-2pt}
{\renewcommand{\arraystretch}{1.35}
\begin{tabularx}{\linewidth}{|X|X|p{2.2cm}|p{2.6cm}|}
  \hline
  \textbf{Nom :} & \textbf{Prénom :} & \textbf{Poste n°} & \textbf{Note :} \hfill /20 \\ \hline
\end{tabularx}}

\textbf{Consignes :} travail \textbf{individuel}, documents et Internet \textbf{non autorisés}. Ce sujet est le
\textbf{document-réponse}. Fichier Packet Tracer : \ip{Nom\_Prenom\_EvalLAN\_\variante.pkt}, à enregistrer
régulièrement dans le dossier de rendu. Les « visas » sont apposés par le professeur, que vous appelez à chaque test réussi.
Durées conseillées : A 30 min — B 1 h 10 — C 20 min.

%==========================================================================
\section{Partie A — Questions de cours et calculs \hfill /6}

\question{A1}{2 pts}{Complétez le tableau en considérant le \textbf{masque par défaut} de la classe de chaque adresse.}
{\small\renewcommand{\arraystretch}{1.15}
\begin{tabularx}{\linewidth}{|l|c|C|C|C|c|}
  \hline
  \textbf{Adresse IP} & \textbf{Classe} & \textbf{Masque par défaut} & \textbf{Adresse réseau} &
  \textbf{Adresse de diffusion} & \textbf{Nb d'hôtes} \\ \hline
  \AIrows
\end{tabularx}}
\bareme{0,5 pt par ligne juste + 0,5 pt si les trois nombres d'hôtes sont justes ($2^{n}-2$, $n$ = nombre de bits d'hôte).}

\question{A2}{1 pt}{Indiquez le type de câble (\textbf{droit} ou \textbf{croisé}) pour chaque liaison.}
{\small\renewcommand{\arraystretch}{1.15}
\begin{tabularx}{\linewidth}{|X|c||X|c|}
  \hline
  \hl PC $\leftrightarrow$ commutateur & \makebox[1.6cm]{\sol{droit}} & Commutateur $\leftrightarrow$ routeur & \makebox[1.6cm]{\sol{droit}} \\ \hline
  \hl PC $\leftrightarrow$ PC (réseau P2P) & \sol{croisé} & Concentrateur $\leftrightarrow$ concentrateur & \sol{croisé} \\ \hline
\end{tabularx}}
\bareme{0,25 pt par case. Équipements de types différents : câble droit ; de même type : câble croisé.}

\question{A3}{1 pt}{Comment un \textbf{concentrateur (hub)} et un \textbf{commutateur (switch)} traitent-ils une trame
reçue ? Lequel limite les collisions, et pourquoi ?}
\zone{2.9cm}{Le \textbf{hub} (couche 1) répète la trame sur \textbf{tous ses autres ports} : un seul domaine de collision
partagé. Le \textbf{switch} (couche 2) lit l'adresse MAC de destination et, grâce à sa \textbf{table MAC}, ne l'envoie
\textbf{que sur le port du destinataire} : un domaine de collision par port, c'est lui qui limite les collisions.
\emph{0,5 pt hub, 0,5 pt switch.}}

\question{A4}{1 pt}{Qu'est-ce que la \textbf{passerelle par défaut} d'un poste ? Quelle condition son adresse doit-elle respecter ?}
\zone{2.4cm}{Adresse IP de l'interface du \textbf{routeur} à laquelle le poste envoie les paquets destinés à un
\textbf{autre réseau}. Elle doit \textbf{appartenir au même réseau} que le poste. \emph{0,5 + 0,5 pt.}}

\question{A5}{0,5 pt}{Un portable configuré en DHCP affiche \ip{\apipa} / \ip{255.255.0.0}. Que s'est-il passé ?}
\zone{1.5cm}{Aucun \textbf{serveur DHCP} n'a répondu : le poste s'est attribué une adresse \textbf{APIPA} (169.254.0.0/16).}

\question{A6}{0,5 pt}{Quel protocole utilise \texttt{ping} ? Nommez les deux messages échangés et la couche OSI des adresses IP.}
\zone{1.5cm}{\textbf{ICMP} ; \textbf{Echo Request} / \textbf{Echo Reply} ; adresses IP en \textbf{couche 3 (Réseau)}.}

%==========================================================================
\newpage
\section{Partie B — Réalisation sous Packet Tracer \hfill /10}

\begin{cdc}
\entreprise{} vous confie son réseau : \textbf{trois réseaux} reliés par un \textbf{routeur unique Cisco 2911}
nommé \ip{\Rnom} (interfaces G0/0, G0/1, G0/2).
\begin{itemize}[leftmargin=12pt]
  \item \textbf{Réseau \resUn} : \textbf{4 postes} PC-\pUn1 à PC-\pUn4, commutateur \ip{SW-\pUn}, adressage statique,
        réseau \ip{\netUn.0/24}.
  \item \textbf{Réseau \resDeux} : \textbf{2 postes} PC-\pDeux1, PC-\pDeux2 et \textbf{1 imprimante} IMP-\pDeux,
        commutateur \ip{SW-\pDeux}, adressage statique, réseau de classe B \ip{\netDeux.0.0/16}.
  \item \textbf{Réseau Serveur} : serveur \ip{\srvNom} d'adresse \ip{\netSrv.\srvH/24}, relié \textbf{directement} au routeur.
  \item La passerelle de chaque réseau est sa \textbf{dernière adresse utilisable}. Tous les équipements communiquent
        entre eux et accèdent à la page web du serveur. Nommage imposé (« Display Name »).
\end{itemize}
\end{cdc}

\question{B1}{2 pts}{Complétez le plan d'adressage \textbf{avant} la réalisation.}
{\footnotesize\setlength{\tabcolsep}{3pt}\renewcommand{\arraystretch}{1.15}
\begin{tabularx}{\linewidth}{|l|C|C|C|C|C|C|}
  \hline
  \textbf{Réseau} & \textbf{Adresse réseau} & \textbf{Masque} & \textbf{1re @ utilisable} &
  \textbf{Dernière @ utilisable} & \textbf{Diffusion} & \textbf{Passerelle} \\ \hline
  \hl \resUn & \sol{\netUn.0} & \sol{255.255.255.0} & \sol{\netUn.1} & \sol{\netUn.254} & \sol{\netUn.255} & \sol{\netUn.254} \\ \hline
  \hl \resDeux & \sol{\netDeux.0.0} & \sol{255.255.0.0} & \sol{\netDeux.0.1} & \sol{\netDeux.255.254} & \sol{\netDeux.255.255} & \sol{\netDeux.255.254} \\ \hline
  \hl Serveur & \sol{\netSrv.0} & \sol{255.255.255.0} & \sol{\netSrv.1} & \sol{\netSrv.254} & \sol{\netSrv.255} & \sol{\netSrv.254} \\ \hline
\end{tabularx}}
\bareme{0,5 pt par ligne juste + 0,5 pt si la convention de passerelle est respectée partout.}

\question{B2}{2 pts}{Construisez la topologie (équipements, câblage, nommage) et relevez l'adressage configuré.}
{\footnotesize\setlength{\tabcolsep}{3pt}\renewcommand{\arraystretch}{1.1}
\begin{tabularx}{\linewidth}{|l|C|C|C||l|C|C|}
  \hline
  \textbf{Hôte} & \textbf{Adresse IP} & \textbf{Masque} & \textbf{Passerelle} & \textbf{Interface} & \textbf{Adresse IP} & \textbf{Masque} \\ \hline
  \hl PC-\pUn1 & \sol{\netUn.1} & \sol{255.255.255.0} & \sol{\netUn.254} & \Rnom{} G0/0 & \sol{\netUn.254} & \sol{255.255.255.0} \\ \hline
  \hl PC-\pDeux1 & \sol{\netDeux.0.1} & \sol{255.255.0.0} & \sol{\netDeux.255.254} & \Rnom{} G0/1 & \sol{\netDeux.255.254} & \sol{255.255.0.0} \\ \hline
  \hl IMP-\pDeux & \sol{\netDeux.0.10} & \sol{255.255.0.0} & \sol{\netDeux.255.254} & \Rnom{} G0/2 & \sol{\netSrv.254} & \sol{255.255.255.0} \\ \hline
  \hl \srvNom & \ip{\netSrv.\srvH} & \sol{255.255.255.0} & \sol{\netSrv.254} & \multicolumn{3}{c|}{\cellcolor{black!8}} \\ \hline
\end{tabularx}}
\bareme{Adresses d'hôtes libres (uniques, dans le bon réseau) ; affectation G0/x libre. Sur le .pkt : équipements 0,5 ;
câblage 1 (câbles droits vers les switchs et switch–routeur, \textbf{croisé} serveur–routeur, liens verts) ; nommage 0,5.}

\question{B3}{3 pts}{Configurez tous les équipements : IP, masques, passerelles des hôtes ; interfaces du routeur
(« Port Status : On »).}
\bareme{Sur le .pkt : interfaces routeur 1 pt ; IP/masques des hôtes 1 pt ; passerelles (imprimante et serveur compris)
1 pt. $-0,25$ par erreur, sans descendre sous 0 par item.}

\question{B4}{3 pts}{Réalisez les tests, notez la commande et le résultat, puis faites viser chaque test réussi.}
{\small\renewcommand{\arraystretch}{1.1}
\begin{tabularx}{\linewidth}{|c|l|l|C|c|c|}
  \hline
  \textbf{N°} & \textbf{Source} & \textbf{Destination} & \textbf{Commande / action} & \textbf{Résultat} & \textbf{Visa prof.} \\ \hline
  \hl 1 & PC-\pUn1 & PC-\pUn4 & \sol{ping \netUn.4} & \sol{réussi} & \hspace{1.4cm} \\ \hline
  \hl 2 & PC-\pDeux1 & IMP-\pDeux & \sol{ping \netDeux.0.10} & \sol{réussi} & \\ \hline
  \hl 3 & PC-\pUn2 & PC-\pDeux2 & \sol{ping \netDeux.0.2} & \sol{réussi} & \\ \hline
  \hl 4 & PC-\pDeux2 & \srvNom & \sol{ping \netSrv.\srvH} & \sol{réussi} & \\ \hline
  \hl 5 & PC-\pUn3 & page web du serveur & \sol{navigateur : {\NoAutoSpacing http://\netSrv.\srvH}} & \sol{page affichée} & \\ \hline
\end{tabularx}}
\bareme{0,5 pt par test visé + 0,5 pt pour B5. Le premier ping inter-réseaux peut échouer (résolution ARP) ; les suivants doivent réussir.}

\question{B5}{incl. B4}{Test n°3 : par quels équipements passe le paquet ? Pourquoi PC-\pUn2 l'envoie-t-il au routeur ?}
\zone{2.3cm}{PC-\pUn2 $\to$ SW-\pUn{} $\to$ \Rnom{} $\to$ SW-\pDeux{} $\to$ PC-\pDeux2. Grâce à son masque, PC-\pUn2 constate que
la destination (\netDeux.0.0/16) n'est \textbf{pas dans son réseau} (\netUn.0/24) : il envoie le paquet à sa
\textbf{passerelle par défaut}, qui le route.}

%==========================================================================
\newpage
\section{Partie C — Diagnostic d'un réseau en panne \hfill /4}

Un technicien constate qu'\textbf{aucun poste ne parvient à joindre le serveur SRV-C} et que certains postes
communiquent mal entre eux. Les câbles sont du bon type.

\begin{center}
\begin{minipage}{0.37\linewidth}\centering
\begin{tikzpicture}[font=\footnotesize,
  eq/.style={draw,fill=qbg,rounded corners=2pt,minimum width=9mm,minimum height=5mm,inner sep=2pt},
  every path/.style={thick}]
  \node[eq] (SW) {SW-1};
  \node[eq,right=9mm of SW] (R) {R1};
  \node[eq,right=7mm of R] (SRV) {SRV-C};
  \node[eq,above left=5mm and -3mm of SW] (A) {PC-A};
  \node[eq,above right=5mm and -3mm of SW] (B) {PC-B};
  \node[eq,below left=5mm and -3mm of SW] (C) {PC-C};
  \node[eq,below right=5mm and -3mm of SW] (D) {PC-D};
  \foreach \n in {A,B,C,D} \draw (\n) -- (SW);
  \draw (SW) -- node[above,font=\tiny]{G0/0} (R);
  \draw (R) -- node[above,font=\tiny]{G0/1} (SRV);
\end{tikzpicture}
\end{minipage}\hfill
\begin{minipage}{0.61\linewidth}\centering
{\small\renewcommand{\arraystretch}{1.05}\setlength{\tabcolsep}{4pt}
\begin{tabular}{lllll}
  \toprule
  \textbf{Équipement} & \textbf{Adresse IP} & \textbf{Masque} & \textbf{Passerelle} & \textbf{Port} \\
  \midrule
  \Crows
  \bottomrule
\end{tabular}}
\end{minipage}
\end{center}

\question{C1}{4 pts}{Relevez les \textbf{quatre anomalies}. Pour chacune : conséquence et correction.}
{\small\renewcommand{\arraystretch}{1.1}
\ifcorrige\def\hc{}\else\def\hc{\rule[-10mm]{0pt}{14mm}}\fi
\begin{tabularx}{\linewidth}{|c|>{\raggedright\arraybackslash}p{2cm}|L|L|L|}
  \hline
  \textbf{N°} & \textbf{Équipement} & \textbf{Anomalie} & \textbf{Conséquence} & \textbf{Correction} \\ \hline
  \ifcorrige\Csol\else
  \hc 1 & & & & \\ \hline
  \hc 2 & & & & \\ \hline
  \hc 3 & & & & \\ \hline
  \hc 4 & & & & \\ \hline
  \fi
\end{tabularx}}
\bareme{1 pt par anomalie (0,25 équipement, anomalie, conséquence, correction ; ordre indifférent). \Cbareme}

%==========================================================================
\section{Bonus — Extension sans fil \hfill +1}

\begin{cdc}Ajoutez au réseau \resUn{} un \textbf{AccessPoint-PT} relié à SW-\pUn{} (SSID \ip{\ssid}, \textbf{WPA2-PSK})
et un portable \ip{Laptop-\pUn} équipé d'un module \texttt{WPC300N}, en adressage \textbf{statique}.
Vérifiez qu'il joint \srvNom.\end{cdc}
{\small\renewcommand{\arraystretch}{1.1}
\begin{tabularx}{\linewidth}{|l|C|C|C|c|}
  \hline
  \textbf{Équipement} & \textbf{Adresse IP} & \textbf{Masque} & \textbf{Passerelle} & \textbf{Visa prof.} \\ \hline
  \hl Laptop-\pUn & \sol{\netUn.50} & \sol{255.255.255.0} & \sol{\netUn.254} & \hspace{1.4cm} \\ \hline
\end{tabularx}}
\bareme{Le point d'accès fait office de pont : le portable est dans \netUn.0/24. Bonus si associé en WPA2-PSK et ping
vers \netSrv.\srvH{} réussi. Note plafonnée à 20.}

\end{document}
"
%   pdflatex -jobname=Eval_Reseau_CorrigeA "\def\variante{A}\def\modecorrige{1}% Évaluation pratique — Réseaux locaux (BTS CIEL 1re année) — sujets A et B
% Compilation (4 PDF) :
%   pdflatex -jobname=Eval_Reseau_SujetA   "\def\variante{A}\def\modecorrige{0}\input{Eval_TP_Reseau_LAN.tex}"
%   pdflatex -jobname=Eval_Reseau_SujetB   "\def\variante{B}\def\modecorrige{0}\input{Eval_TP_Reseau_LAN.tex}"
%   pdflatex -jobname=Eval_Reseau_CorrigeA "\def\variante{A}\def\modecorrige{1}\input{Eval_TP_Reseau_LAN.tex}"
%   pdflatex -jobname=Eval_Reseau_CorrigeB "\def\variante{B}\def\modecorrige{1}\input{Eval_TP_Reseau_LAN.tex}"
\documentclass[10pt,a4paper]{article}

\usepackage[T1]{fontenc}
\usepackage[utf8]{inputenc}
\usepackage{lmodern}
\usepackage[french]{babel}
\usepackage[a4paper,hmargin=1.5cm,top=1.3cm,bottom=1.6cm,footskip=0.7cm]{geometry}
\usepackage[table]{xcolor}
\usepackage{booktabs,tabularx,array}
\usepackage{enumitem}
\usepackage{fancyhdr}
\usepackage[most]{tcolorbox}
\usepackage{tikz}
\usetikzlibrary{positioning}
\usepackage{titlesec}

\providecommand{\modecorrige}{0}
\providecommand{\variante}{A}
\newif\ifcorrige
\ifnum\modecorrige=1 \corrigetrue \else \corrigefalse \fi

\definecolor{ipred}{HTML}{C0392B}
\definecolor{cdcbar}{HTML}{A67C3D}
\definecolor{cdcbg}{HTML}{F4F1EA}
\definecolor{qbar}{HTML}{4A6B5D}
\definecolor{qbg}{HTML}{EDF1EF}
\definecolor{corr}{HTML}{1F5FA8}
\definecolor{corrbg}{HTML}{E8F0FA}
\definecolor{grisfoot}{HTML}{828282}

\newcommand{\ip}[1]{\textcolor{ipred}{\texttt{#1}}}
\newcommand{\pts}[1]{\hfill\mbox{\textbf{(#1)}}}
\newcommand{\sol}[1]{\ifcorrige\textcolor{corr}{#1}\fi}
\newcommand{\bareme}[1]{\ifcorrige\par{\footnotesize\color{corr}#1\par}\fi}
\newcommand{\hl}{\ifcorrige\rule[-1.6mm]{0pt}{5.4mm}\else\rule[-2.2mm]{0pt}{6.6mm}\fi}
\newcolumntype{C}{>{\centering\arraybackslash}X}
\newcolumntype{L}{>{\raggedright\arraybackslash}X}

\titleformat{\section}{\large\bfseries}{}{0pt}{}
\titlespacing*{\section}{0pt}{6pt}{3pt}

\tcbset{qstyle/.style={enhanced,boxrule=0pt,leftrule=2.5pt,arc=0pt,outer arc=0pt,
  left=5pt,right=5pt,top=2pt,bottom=2pt,before skip=4pt,after skip=3pt}}
\newtcolorbox{cdc}{qstyle,colback=cdcbg,colframe=cdcbar}
\newtcolorbox{quest}{qstyle,colback=qbg,colframe=qbar}

% Zone de réponse : vierge dans le sujet, remplie dans le corrigé
\newcommand{\zone}[2]{%
  \ifcorrige
    \begin{tcolorbox}[colback=corrbg,colframe=corr,boxrule=0.5pt,arc=1pt,left=4pt,right=4pt,top=1pt,bottom=1pt,
      before skip=2pt,after skip=2pt,fontupper=\footnotesize\color{corr}]#2\end{tcolorbox}%
  \else
    \begin{tcolorbox}[colback=white,colframe=black!45,boxrule=0.4pt,arc=1pt,height=#1,
      before skip=2pt,after skip=2pt]\end{tcolorbox}%
  \fi}

\newcommand{\question}[3]{\begin{quest}\textbf{#1.} #3\pts{#2}\end{quest}}

%==========================================================================
% Données propres à chaque variante
%==========================================================================
\if\variante A
  \def\entreprise{Le cabinet d'architectes \textbf{ArchiNova}}
  \def\Rnom{R-ArchiNova}
  \def\resUn{Bureau d'études}\def\pUn{BE}\def\netUn{192.168.40}
  \def\resDeux{Accueil}\def\pDeux{ACC}\def\netDeux{172.16}
  \def\srvNom{SRV-Intranet}\def\netSrv{192.168.200}\def\srvH{10}
  \def\ssid{ArchiNova-WiFi}
  \def\AIrows{%
    \hl\ip{192.168.40.25} & \sol{C} & \sol{255.255.255.0} & \sol{192.168.40.0} & \sol{192.168.40.255} & \sol{254} \\ \hline
    \hl\ip{172.16.5.10}   & \sol{B} & \sol{255.255.0.0}   & \sol{172.16.0.0}   & \sol{172.16.255.255} & \sol{65\,534} \\ \hline
    \hl\ip{10.3.2.1}      & \sol{A} & \sol{255.0.0.0}     & \sol{10.0.0.0}     & \sol{10.255.255.255} & \sol{16\,777\,214} \\ \hline}
  \def\apipa{169.254.12.7}
  \def\Cnet{192.168.60}
  \def\Crows{%
    PC-A & \ip{192.168.60.10} & \ip{255.255.255.0} & \ip{192.168.60.254} & — \\
    PC-B & \ip{192.168.60.20} & \ip{255.255.255.0} & \ip{192.168.60.1}   & — \\
    PC-C & \ip{192.168.61.30} & \ip{255.255.255.0} & \ip{192.168.60.254} & — \\
    PC-D & \ip{192.168.60.10} & \ip{255.255.255.0} & \ip{192.168.60.254} & — \\
    SRV-C & \ip{10.0.0.5}     & \ip{255.0.0.0}     & \ip{10.0.0.254}     & — \\
    R1 — G0/0 & \ip{192.168.60.254} & \ip{255.255.255.0} & — & On \\
    R1 — G0/1 & \ip{10.0.0.254}     & \ip{255.0.0.0}     & — & Off \\}
  \def\Csol{%
    \hl 1 & \sol{PC-B} & \sol{Passerelle .60.1 : pas l'adresse du routeur.} & \sol{Joint son réseau, mais aucun autre (serveur).} & \sol{Passerelle 192.168.60.254.} \\ \hline
    \hl 2 & \sol{PC-C} & \sol{IP hors du réseau 192.168.60.0/24.} & \sol{Ne joint ni les postes ni le routeur.} & \sol{Ex. 192.168.60.30.} \\ \hline
    \hl 3 & \sol{PC-D} & \sol{Même IP que PC-A (doublon).} & \sol{Conflit d'adresse : PC-A et PC-D perturbés.} & \sol{Adresse unique, ex. .60.40.} \\ \hline
    \hl 4 & \sol{R1 G0/1} & \sol{Interface désactivée (Off).} & \sol{Réseau du serveur injoignable pour tous.} & \sol{Port Status : On.} \\ \hline}
  \def\Cbareme{PC-A et SRV-C sont correctement configurés ; accepter PC-A cité à la place de PC-D pour le doublon.}
\else
  \def\entreprise{La clinique vétérinaire \textbf{VetoPlus}}
  \def\Rnom{R-VetoPlus}
  \def\resUn{Soins}\def\pUn{SOI}\def\netUn{192.168.70}
  \def\resDeux{Secrétariat}\def\pDeux{SEC}\def\netDeux{172.20}
  \def\srvNom{SRV-Dossiers}\def\netSrv{192.168.250}\def\srvH{20}
  \def\ssid{VetoPlus-WiFi}
  \def\AIrows{%
    \hl\ip{192.168.15.130} & \sol{C} & \sol{255.255.255.0} & \sol{192.168.15.0} & \sol{192.168.15.255} & \sol{254} \\ \hline
    \hl\ip{172.31.200.7}   & \sol{B} & \sol{255.255.0.0}   & \sol{172.31.0.0}   & \sol{172.31.255.255} & \sol{65\,534} \\ \hline
    \hl\ip{10.45.6.9}      & \sol{A} & \sol{255.0.0.0}     & \sol{10.0.0.0}     & \sol{10.255.255.255} & \sol{16\,777\,214} \\ \hline}
  \def\apipa{169.254.88.3}
  \def\Cnet{192.168.80}
  \def\Crows{%
    PC-A & \ip{192.168.8.11}   & \ip{255.255.255.0} & \ip{192.168.80.254} & — \\
    PC-B & \ip{192.168.80.12}  & \ip{255.255.255.0} & \ip{192.168.80.254} & — \\
    PC-C & \ip{192.168.80.13}  & \ip{255.255.255.0} & \ip{192.168.80.100} & — \\
    PC-D & \ip{192.168.80.12}  & \ip{255.255.255.0} & \ip{192.168.80.254} & — \\
    SRV-C & \ip{10.0.0.8}      & \ip{255.0.0.0}     & \ip{10.0.0.254}     & — \\
    R1 — G0/0 & \ip{192.168.80.254} & \ip{255.255.255.0} & — & Off \\
    R1 — G0/1 & \ip{10.0.0.254}     & \ip{255.0.0.0}     & — & On \\}
  \def\Csol{%
    \hl 1 & \sol{PC-A} & \sol{IP hors du réseau 192.168.80.0/24.} & \sol{Ne joint ni les postes ni le routeur.} & \sol{Ex. 192.168.80.11.} \\ \hline
    \hl 2 & \sol{PC-C} & \sol{Passerelle .80.100 : pas l'adresse du routeur.} & \sol{Joint son réseau, mais aucun autre (serveur).} & \sol{Passerelle 192.168.80.254.} \\ \hline
    \hl 3 & \sol{PC-D} & \sol{Même IP que PC-B (doublon).} & \sol{Conflit d'adresse : PC-B et PC-D perturbés.} & \sol{Adresse unique, ex. .80.14.} \\ \hline
    \hl 4 & \sol{R1 G0/0} & \sol{Interface côté LAN désactivée (Off).} & \sol{Aucun poste ne joint le routeur ni le serveur.} & \sol{Port Status : On.} \\ \hline}
  \def\Cbareme{PC-B et SRV-C sont correctement configurés ; accepter PC-B cité à la place de PC-D pour le doublon.}
\fi

\pagestyle{fancy}
\fancyhf{}
\renewcommand{\headrulewidth}{0pt}
\fancyfoot[L]{\footnotesize\itshape\color{grisfoot}Évaluation — Réseaux locaux — Sujet \variante\ifcorrige\ — CORRIGÉ\fi}
\fancyfoot[R]{\footnotesize\color{grisfoot}\thepage/3}

\setlist{itemsep=0pt,topsep=2pt,parsep=0pt}
\setlength{\parindent}{0pt}
\setlength{\parskip}{2pt}

\begin{document}

\begin{center}
  {\Large\bfseries Évaluation pratique — Réseaux locaux \quad
   \fcolorbox{black}{white}{\ SUJET \variante\ }\par}
  \vspace{1pt}
  {\bfseries Adressage IP, commutation et routage — BTS CIEL 1re année — durée : 2 heures\par}
  \ifcorrige\vspace{2pt}{\bfseries\color{corr}CORRIGÉ ET BARÈME — document enseignant}\par\fi
\end{center}
\vspace{-2pt}
{\renewcommand{\arraystretch}{1.35}
\begin{tabularx}{\linewidth}{|X|X|p{2.2cm}|p{2.6cm}|}
  \hline
  \textbf{Nom :} & \textbf{Prénom :} & \textbf{Poste n°} & \textbf{Note :} \hfill /20 \\ \hline
\end{tabularx}}

\textbf{Consignes :} travail \textbf{individuel}, documents et Internet \textbf{non autorisés}. Ce sujet est le
\textbf{document-réponse}. Fichier Packet Tracer : \ip{Nom\_Prenom\_EvalLAN\_\variante.pkt}, à enregistrer
régulièrement dans le dossier de rendu. Les « visas » sont apposés par le professeur, que vous appelez à chaque test réussi.
Durées conseillées : A 30 min — B 1 h 10 — C 20 min.

%==========================================================================
\section{Partie A — Questions de cours et calculs \hfill /6}

\question{A1}{2 pts}{Complétez le tableau en considérant le \textbf{masque par défaut} de la classe de chaque adresse.}
{\small\renewcommand{\arraystretch}{1.15}
\begin{tabularx}{\linewidth}{|l|c|C|C|C|c|}
  \hline
  \textbf{Adresse IP} & \textbf{Classe} & \textbf{Masque par défaut} & \textbf{Adresse réseau} &
  \textbf{Adresse de diffusion} & \textbf{Nb d'hôtes} \\ \hline
  \AIrows
\end{tabularx}}
\bareme{0,5 pt par ligne juste + 0,5 pt si les trois nombres d'hôtes sont justes ($2^{n}-2$, $n$ = nombre de bits d'hôte).}

\question{A2}{1 pt}{Indiquez le type de câble (\textbf{droit} ou \textbf{croisé}) pour chaque liaison.}
{\small\renewcommand{\arraystretch}{1.15}
\begin{tabularx}{\linewidth}{|X|c||X|c|}
  \hline
  \hl PC $\leftrightarrow$ commutateur & \makebox[1.6cm]{\sol{droit}} & Commutateur $\leftrightarrow$ routeur & \makebox[1.6cm]{\sol{droit}} \\ \hline
  \hl PC $\leftrightarrow$ PC (réseau P2P) & \sol{croisé} & Concentrateur $\leftrightarrow$ concentrateur & \sol{croisé} \\ \hline
\end{tabularx}}
\bareme{0,25 pt par case. Équipements de types différents : câble droit ; de même type : câble croisé.}

\question{A3}{1 pt}{Comment un \textbf{concentrateur (hub)} et un \textbf{commutateur (switch)} traitent-ils une trame
reçue ? Lequel limite les collisions, et pourquoi ?}
\zone{2.9cm}{Le \textbf{hub} (couche 1) répète la trame sur \textbf{tous ses autres ports} : un seul domaine de collision
partagé. Le \textbf{switch} (couche 2) lit l'adresse MAC de destination et, grâce à sa \textbf{table MAC}, ne l'envoie
\textbf{que sur le port du destinataire} : un domaine de collision par port, c'est lui qui limite les collisions.
\emph{0,5 pt hub, 0,5 pt switch.}}

\question{A4}{1 pt}{Qu'est-ce que la \textbf{passerelle par défaut} d'un poste ? Quelle condition son adresse doit-elle respecter ?}
\zone{2.4cm}{Adresse IP de l'interface du \textbf{routeur} à laquelle le poste envoie les paquets destinés à un
\textbf{autre réseau}. Elle doit \textbf{appartenir au même réseau} que le poste. \emph{0,5 + 0,5 pt.}}

\question{A5}{0,5 pt}{Un portable configuré en DHCP affiche \ip{\apipa} / \ip{255.255.0.0}. Que s'est-il passé ?}
\zone{1.5cm}{Aucun \textbf{serveur DHCP} n'a répondu : le poste s'est attribué une adresse \textbf{APIPA} (169.254.0.0/16).}

\question{A6}{0,5 pt}{Quel protocole utilise \texttt{ping} ? Nommez les deux messages échangés et la couche OSI des adresses IP.}
\zone{1.5cm}{\textbf{ICMP} ; \textbf{Echo Request} / \textbf{Echo Reply} ; adresses IP en \textbf{couche 3 (Réseau)}.}

%==========================================================================
\newpage
\section{Partie B — Réalisation sous Packet Tracer \hfill /10}

\begin{cdc}
\entreprise{} vous confie son réseau : \textbf{trois réseaux} reliés par un \textbf{routeur unique Cisco 2911}
nommé \ip{\Rnom} (interfaces G0/0, G0/1, G0/2).
\begin{itemize}[leftmargin=12pt]
  \item \textbf{Réseau \resUn} : \textbf{4 postes} PC-\pUn1 à PC-\pUn4, commutateur \ip{SW-\pUn}, adressage statique,
        réseau \ip{\netUn.0/24}.
  \item \textbf{Réseau \resDeux} : \textbf{2 postes} PC-\pDeux1, PC-\pDeux2 et \textbf{1 imprimante} IMP-\pDeux,
        commutateur \ip{SW-\pDeux}, adressage statique, réseau de classe B \ip{\netDeux.0.0/16}.
  \item \textbf{Réseau Serveur} : serveur \ip{\srvNom} d'adresse \ip{\netSrv.\srvH/24}, relié \textbf{directement} au routeur.
  \item La passerelle de chaque réseau est sa \textbf{dernière adresse utilisable}. Tous les équipements communiquent
        entre eux et accèdent à la page web du serveur. Nommage imposé (« Display Name »).
\end{itemize}
\end{cdc}

\question{B1}{2 pts}{Complétez le plan d'adressage \textbf{avant} la réalisation.}
{\footnotesize\setlength{\tabcolsep}{3pt}\renewcommand{\arraystretch}{1.15}
\begin{tabularx}{\linewidth}{|l|C|C|C|C|C|C|}
  \hline
  \textbf{Réseau} & \textbf{Adresse réseau} & \textbf{Masque} & \textbf{1re @ utilisable} &
  \textbf{Dernière @ utilisable} & \textbf{Diffusion} & \textbf{Passerelle} \\ \hline
  \hl \resUn & \sol{\netUn.0} & \sol{255.255.255.0} & \sol{\netUn.1} & \sol{\netUn.254} & \sol{\netUn.255} & \sol{\netUn.254} \\ \hline
  \hl \resDeux & \sol{\netDeux.0.0} & \sol{255.255.0.0} & \sol{\netDeux.0.1} & \sol{\netDeux.255.254} & \sol{\netDeux.255.255} & \sol{\netDeux.255.254} \\ \hline
  \hl Serveur & \sol{\netSrv.0} & \sol{255.255.255.0} & \sol{\netSrv.1} & \sol{\netSrv.254} & \sol{\netSrv.255} & \sol{\netSrv.254} \\ \hline
\end{tabularx}}
\bareme{0,5 pt par ligne juste + 0,5 pt si la convention de passerelle est respectée partout.}

\question{B2}{2 pts}{Construisez la topologie (équipements, câblage, nommage) et relevez l'adressage configuré.}
{\footnotesize\setlength{\tabcolsep}{3pt}\renewcommand{\arraystretch}{1.1}
\begin{tabularx}{\linewidth}{|l|C|C|C||l|C|C|}
  \hline
  \textbf{Hôte} & \textbf{Adresse IP} & \textbf{Masque} & \textbf{Passerelle} & \textbf{Interface} & \textbf{Adresse IP} & \textbf{Masque} \\ \hline
  \hl PC-\pUn1 & \sol{\netUn.1} & \sol{255.255.255.0} & \sol{\netUn.254} & \Rnom{} G0/0 & \sol{\netUn.254} & \sol{255.255.255.0} \\ \hline
  \hl PC-\pDeux1 & \sol{\netDeux.0.1} & \sol{255.255.0.0} & \sol{\netDeux.255.254} & \Rnom{} G0/1 & \sol{\netDeux.255.254} & \sol{255.255.0.0} \\ \hline
  \hl IMP-\pDeux & \sol{\netDeux.0.10} & \sol{255.255.0.0} & \sol{\netDeux.255.254} & \Rnom{} G0/2 & \sol{\netSrv.254} & \sol{255.255.255.0} \\ \hline
  \hl \srvNom & \ip{\netSrv.\srvH} & \sol{255.255.255.0} & \sol{\netSrv.254} & \multicolumn{3}{c|}{\cellcolor{black!8}} \\ \hline
\end{tabularx}}
\bareme{Adresses d'hôtes libres (uniques, dans le bon réseau) ; affectation G0/x libre. Sur le .pkt : équipements 0,5 ;
câblage 1 (câbles droits vers les switchs et switch–routeur, \textbf{croisé} serveur–routeur, liens verts) ; nommage 0,5.}

\question{B3}{3 pts}{Configurez tous les équipements : IP, masques, passerelles des hôtes ; interfaces du routeur
(« Port Status : On »).}
\bareme{Sur le .pkt : interfaces routeur 1 pt ; IP/masques des hôtes 1 pt ; passerelles (imprimante et serveur compris)
1 pt. $-0,25$ par erreur, sans descendre sous 0 par item.}

\question{B4}{3 pts}{Réalisez les tests, notez la commande et le résultat, puis faites viser chaque test réussi.}
{\small\renewcommand{\arraystretch}{1.1}
\begin{tabularx}{\linewidth}{|c|l|l|C|c|c|}
  \hline
  \textbf{N°} & \textbf{Source} & \textbf{Destination} & \textbf{Commande / action} & \textbf{Résultat} & \textbf{Visa prof.} \\ \hline
  \hl 1 & PC-\pUn1 & PC-\pUn4 & \sol{ping \netUn.4} & \sol{réussi} & \hspace{1.4cm} \\ \hline
  \hl 2 & PC-\pDeux1 & IMP-\pDeux & \sol{ping \netDeux.0.10} & \sol{réussi} & \\ \hline
  \hl 3 & PC-\pUn2 & PC-\pDeux2 & \sol{ping \netDeux.0.2} & \sol{réussi} & \\ \hline
  \hl 4 & PC-\pDeux2 & \srvNom & \sol{ping \netSrv.\srvH} & \sol{réussi} & \\ \hline
  \hl 5 & PC-\pUn3 & page web du serveur & \sol{navigateur : {\NoAutoSpacing http://\netSrv.\srvH}} & \sol{page affichée} & \\ \hline
\end{tabularx}}
\bareme{0,5 pt par test visé + 0,5 pt pour B5. Le premier ping inter-réseaux peut échouer (résolution ARP) ; les suivants doivent réussir.}

\question{B5}{incl. B4}{Test n°3 : par quels équipements passe le paquet ? Pourquoi PC-\pUn2 l'envoie-t-il au routeur ?}
\zone{2.3cm}{PC-\pUn2 $\to$ SW-\pUn{} $\to$ \Rnom{} $\to$ SW-\pDeux{} $\to$ PC-\pDeux2. Grâce à son masque, PC-\pUn2 constate que
la destination (\netDeux.0.0/16) n'est \textbf{pas dans son réseau} (\netUn.0/24) : il envoie le paquet à sa
\textbf{passerelle par défaut}, qui le route.}

%==========================================================================
\newpage
\section{Partie C — Diagnostic d'un réseau en panne \hfill /4}

Un technicien constate qu'\textbf{aucun poste ne parvient à joindre le serveur SRV-C} et que certains postes
communiquent mal entre eux. Les câbles sont du bon type.

\begin{center}
\begin{minipage}{0.37\linewidth}\centering
\begin{tikzpicture}[font=\footnotesize,
  eq/.style={draw,fill=qbg,rounded corners=2pt,minimum width=9mm,minimum height=5mm,inner sep=2pt},
  every path/.style={thick}]
  \node[eq] (SW) {SW-1};
  \node[eq,right=9mm of SW] (R) {R1};
  \node[eq,right=7mm of R] (SRV) {SRV-C};
  \node[eq,above left=5mm and -3mm of SW] (A) {PC-A};
  \node[eq,above right=5mm and -3mm of SW] (B) {PC-B};
  \node[eq,below left=5mm and -3mm of SW] (C) {PC-C};
  \node[eq,below right=5mm and -3mm of SW] (D) {PC-D};
  \foreach \n in {A,B,C,D} \draw (\n) -- (SW);
  \draw (SW) -- node[above,font=\tiny]{G0/0} (R);
  \draw (R) -- node[above,font=\tiny]{G0/1} (SRV);
\end{tikzpicture}
\end{minipage}\hfill
\begin{minipage}{0.61\linewidth}\centering
{\small\renewcommand{\arraystretch}{1.05}\setlength{\tabcolsep}{4pt}
\begin{tabular}{lllll}
  \toprule
  \textbf{Équipement} & \textbf{Adresse IP} & \textbf{Masque} & \textbf{Passerelle} & \textbf{Port} \\
  \midrule
  \Crows
  \bottomrule
\end{tabular}}
\end{minipage}
\end{center}

\question{C1}{4 pts}{Relevez les \textbf{quatre anomalies}. Pour chacune : conséquence et correction.}
{\small\renewcommand{\arraystretch}{1.1}
\ifcorrige\def\hc{}\else\def\hc{\rule[-10mm]{0pt}{14mm}}\fi
\begin{tabularx}{\linewidth}{|c|>{\raggedright\arraybackslash}p{2cm}|L|L|L|}
  \hline
  \textbf{N°} & \textbf{Équipement} & \textbf{Anomalie} & \textbf{Conséquence} & \textbf{Correction} \\ \hline
  \ifcorrige\Csol\else
  \hc 1 & & & & \\ \hline
  \hc 2 & & & & \\ \hline
  \hc 3 & & & & \\ \hline
  \hc 4 & & & & \\ \hline
  \fi
\end{tabularx}}
\bareme{1 pt par anomalie (0,25 équipement, anomalie, conséquence, correction ; ordre indifférent). \Cbareme}

%==========================================================================
\section{Bonus — Extension sans fil \hfill +1}

\begin{cdc}Ajoutez au réseau \resUn{} un \textbf{AccessPoint-PT} relié à SW-\pUn{} (SSID \ip{\ssid}, \textbf{WPA2-PSK})
et un portable \ip{Laptop-\pUn} équipé d'un module \texttt{WPC300N}, en adressage \textbf{statique}.
Vérifiez qu'il joint \srvNom.\end{cdc}
{\small\renewcommand{\arraystretch}{1.1}
\begin{tabularx}{\linewidth}{|l|C|C|C|c|}
  \hline
  \textbf{Équipement} & \textbf{Adresse IP} & \textbf{Masque} & \textbf{Passerelle} & \textbf{Visa prof.} \\ \hline
  \hl Laptop-\pUn & \sol{\netUn.50} & \sol{255.255.255.0} & \sol{\netUn.254} & \hspace{1.4cm} \\ \hline
\end{tabularx}}
\bareme{Le point d'accès fait office de pont : le portable est dans \netUn.0/24. Bonus si associé en WPA2-PSK et ping
vers \netSrv.\srvH{} réussi. Note plafonnée à 20.}

\end{document}
"
%   pdflatex -jobname=Eval_Reseau_CorrigeB "\def\variante{B}\def\modecorrige{1}% Évaluation pratique — Réseaux locaux (BTS CIEL 1re année) — sujets A et B
% Compilation (4 PDF) :
%   pdflatex -jobname=Eval_Reseau_SujetA   "\def\variante{A}\def\modecorrige{0}\input{Eval_TP_Reseau_LAN.tex}"
%   pdflatex -jobname=Eval_Reseau_SujetB   "\def\variante{B}\def\modecorrige{0}\input{Eval_TP_Reseau_LAN.tex}"
%   pdflatex -jobname=Eval_Reseau_CorrigeA "\def\variante{A}\def\modecorrige{1}\input{Eval_TP_Reseau_LAN.tex}"
%   pdflatex -jobname=Eval_Reseau_CorrigeB "\def\variante{B}\def\modecorrige{1}\input{Eval_TP_Reseau_LAN.tex}"
\documentclass[10pt,a4paper]{article}

\usepackage[T1]{fontenc}
\usepackage[utf8]{inputenc}
\usepackage{lmodern}
\usepackage[french]{babel}
\usepackage[a4paper,hmargin=1.5cm,top=1.3cm,bottom=1.6cm,footskip=0.7cm]{geometry}
\usepackage[table]{xcolor}
\usepackage{booktabs,tabularx,array}
\usepackage{enumitem}
\usepackage{fancyhdr}
\usepackage[most]{tcolorbox}
\usepackage{tikz}
\usetikzlibrary{positioning}
\usepackage{titlesec}

\providecommand{\modecorrige}{0}
\providecommand{\variante}{A}
\newif\ifcorrige
\ifnum\modecorrige=1 \corrigetrue \else \corrigefalse \fi

\definecolor{ipred}{HTML}{C0392B}
\definecolor{cdcbar}{HTML}{A67C3D}
\definecolor{cdcbg}{HTML}{F4F1EA}
\definecolor{qbar}{HTML}{4A6B5D}
\definecolor{qbg}{HTML}{EDF1EF}
\definecolor{corr}{HTML}{1F5FA8}
\definecolor{corrbg}{HTML}{E8F0FA}
\definecolor{grisfoot}{HTML}{828282}

\newcommand{\ip}[1]{\textcolor{ipred}{\texttt{#1}}}
\newcommand{\pts}[1]{\hfill\mbox{\textbf{(#1)}}}
\newcommand{\sol}[1]{\ifcorrige\textcolor{corr}{#1}\fi}
\newcommand{\bareme}[1]{\ifcorrige\par{\footnotesize\color{corr}#1\par}\fi}
\newcommand{\hl}{\ifcorrige\rule[-1.6mm]{0pt}{5.4mm}\else\rule[-2.2mm]{0pt}{6.6mm}\fi}
\newcolumntype{C}{>{\centering\arraybackslash}X}
\newcolumntype{L}{>{\raggedright\arraybackslash}X}

\titleformat{\section}{\large\bfseries}{}{0pt}{}
\titlespacing*{\section}{0pt}{6pt}{3pt}

\tcbset{qstyle/.style={enhanced,boxrule=0pt,leftrule=2.5pt,arc=0pt,outer arc=0pt,
  left=5pt,right=5pt,top=2pt,bottom=2pt,before skip=4pt,after skip=3pt}}
\newtcolorbox{cdc}{qstyle,colback=cdcbg,colframe=cdcbar}
\newtcolorbox{quest}{qstyle,colback=qbg,colframe=qbar}

% Zone de réponse : vierge dans le sujet, remplie dans le corrigé
\newcommand{\zone}[2]{%
  \ifcorrige
    \begin{tcolorbox}[colback=corrbg,colframe=corr,boxrule=0.5pt,arc=1pt,left=4pt,right=4pt,top=1pt,bottom=1pt,
      before skip=2pt,after skip=2pt,fontupper=\footnotesize\color{corr}]#2\end{tcolorbox}%
  \else
    \begin{tcolorbox}[colback=white,colframe=black!45,boxrule=0.4pt,arc=1pt,height=#1,
      before skip=2pt,after skip=2pt]\end{tcolorbox}%
  \fi}

\newcommand{\question}[3]{\begin{quest}\textbf{#1.} #3\pts{#2}\end{quest}}

%==========================================================================
% Données propres à chaque variante
%==========================================================================
\if\variante A
  \def\entreprise{Le cabinet d'architectes \textbf{ArchiNova}}
  \def\Rnom{R-ArchiNova}
  \def\resUn{Bureau d'études}\def\pUn{BE}\def\netUn{192.168.40}
  \def\resDeux{Accueil}\def\pDeux{ACC}\def\netDeux{172.16}
  \def\srvNom{SRV-Intranet}\def\netSrv{192.168.200}\def\srvH{10}
  \def\ssid{ArchiNova-WiFi}
  \def\AIrows{%
    \hl\ip{192.168.40.25} & \sol{C} & \sol{255.255.255.0} & \sol{192.168.40.0} & \sol{192.168.40.255} & \sol{254} \\ \hline
    \hl\ip{172.16.5.10}   & \sol{B} & \sol{255.255.0.0}   & \sol{172.16.0.0}   & \sol{172.16.255.255} & \sol{65\,534} \\ \hline
    \hl\ip{10.3.2.1}      & \sol{A} & \sol{255.0.0.0}     & \sol{10.0.0.0}     & \sol{10.255.255.255} & \sol{16\,777\,214} \\ \hline}
  \def\apipa{169.254.12.7}
  \def\Cnet{192.168.60}
  \def\Crows{%
    PC-A & \ip{192.168.60.10} & \ip{255.255.255.0} & \ip{192.168.60.254} & — \\
    PC-B & \ip{192.168.60.20} & \ip{255.255.255.0} & \ip{192.168.60.1}   & — \\
    PC-C & \ip{192.168.61.30} & \ip{255.255.255.0} & \ip{192.168.60.254} & — \\
    PC-D & \ip{192.168.60.10} & \ip{255.255.255.0} & \ip{192.168.60.254} & — \\
    SRV-C & \ip{10.0.0.5}     & \ip{255.0.0.0}     & \ip{10.0.0.254}     & — \\
    R1 — G0/0 & \ip{192.168.60.254} & \ip{255.255.255.0} & — & On \\
    R1 — G0/1 & \ip{10.0.0.254}     & \ip{255.0.0.0}     & — & Off \\}
  \def\Csol{%
    \hl 1 & \sol{PC-B} & \sol{Passerelle .60.1 : pas l'adresse du routeur.} & \sol{Joint son réseau, mais aucun autre (serveur).} & \sol{Passerelle 192.168.60.254.} \\ \hline
    \hl 2 & \sol{PC-C} & \sol{IP hors du réseau 192.168.60.0/24.} & \sol{Ne joint ni les postes ni le routeur.} & \sol{Ex. 192.168.60.30.} \\ \hline
    \hl 3 & \sol{PC-D} & \sol{Même IP que PC-A (doublon).} & \sol{Conflit d'adresse : PC-A et PC-D perturbés.} & \sol{Adresse unique, ex. .60.40.} \\ \hline
    \hl 4 & \sol{R1 G0/1} & \sol{Interface désactivée (Off).} & \sol{Réseau du serveur injoignable pour tous.} & \sol{Port Status : On.} \\ \hline}
  \def\Cbareme{PC-A et SRV-C sont correctement configurés ; accepter PC-A cité à la place de PC-D pour le doublon.}
\else
  \def\entreprise{La clinique vétérinaire \textbf{VetoPlus}}
  \def\Rnom{R-VetoPlus}
  \def\resUn{Soins}\def\pUn{SOI}\def\netUn{192.168.70}
  \def\resDeux{Secrétariat}\def\pDeux{SEC}\def\netDeux{172.20}
  \def\srvNom{SRV-Dossiers}\def\netSrv{192.168.250}\def\srvH{20}
  \def\ssid{VetoPlus-WiFi}
  \def\AIrows{%
    \hl\ip{192.168.15.130} & \sol{C} & \sol{255.255.255.0} & \sol{192.168.15.0} & \sol{192.168.15.255} & \sol{254} \\ \hline
    \hl\ip{172.31.200.7}   & \sol{B} & \sol{255.255.0.0}   & \sol{172.31.0.0}   & \sol{172.31.255.255} & \sol{65\,534} \\ \hline
    \hl\ip{10.45.6.9}      & \sol{A} & \sol{255.0.0.0}     & \sol{10.0.0.0}     & \sol{10.255.255.255} & \sol{16\,777\,214} \\ \hline}
  \def\apipa{169.254.88.3}
  \def\Cnet{192.168.80}
  \def\Crows{%
    PC-A & \ip{192.168.8.11}   & \ip{255.255.255.0} & \ip{192.168.80.254} & — \\
    PC-B & \ip{192.168.80.12}  & \ip{255.255.255.0} & \ip{192.168.80.254} & — \\
    PC-C & \ip{192.168.80.13}  & \ip{255.255.255.0} & \ip{192.168.80.100} & — \\
    PC-D & \ip{192.168.80.12}  & \ip{255.255.255.0} & \ip{192.168.80.254} & — \\
    SRV-C & \ip{10.0.0.8}      & \ip{255.0.0.0}     & \ip{10.0.0.254}     & — \\
    R1 — G0/0 & \ip{192.168.80.254} & \ip{255.255.255.0} & — & Off \\
    R1 — G0/1 & \ip{10.0.0.254}     & \ip{255.0.0.0}     & — & On \\}
  \def\Csol{%
    \hl 1 & \sol{PC-A} & \sol{IP hors du réseau 192.168.80.0/24.} & \sol{Ne joint ni les postes ni le routeur.} & \sol{Ex. 192.168.80.11.} \\ \hline
    \hl 2 & \sol{PC-C} & \sol{Passerelle .80.100 : pas l'adresse du routeur.} & \sol{Joint son réseau, mais aucun autre (serveur).} & \sol{Passerelle 192.168.80.254.} \\ \hline
    \hl 3 & \sol{PC-D} & \sol{Même IP que PC-B (doublon).} & \sol{Conflit d'adresse : PC-B et PC-D perturbés.} & \sol{Adresse unique, ex. .80.14.} \\ \hline
    \hl 4 & \sol{R1 G0/0} & \sol{Interface côté LAN désactivée (Off).} & \sol{Aucun poste ne joint le routeur ni le serveur.} & \sol{Port Status : On.} \\ \hline}
  \def\Cbareme{PC-B et SRV-C sont correctement configurés ; accepter PC-B cité à la place de PC-D pour le doublon.}
\fi

\pagestyle{fancy}
\fancyhf{}
\renewcommand{\headrulewidth}{0pt}
\fancyfoot[L]{\footnotesize\itshape\color{grisfoot}Évaluation — Réseaux locaux — Sujet \variante\ifcorrige\ — CORRIGÉ\fi}
\fancyfoot[R]{\footnotesize\color{grisfoot}\thepage/3}

\setlist{itemsep=0pt,topsep=2pt,parsep=0pt}
\setlength{\parindent}{0pt}
\setlength{\parskip}{2pt}

\begin{document}

\begin{center}
  {\Large\bfseries Évaluation pratique — Réseaux locaux \quad
   \fcolorbox{black}{white}{\ SUJET \variante\ }\par}
  \vspace{1pt}
  {\bfseries Adressage IP, commutation et routage — BTS CIEL 1re année — durée : 2 heures\par}
  \ifcorrige\vspace{2pt}{\bfseries\color{corr}CORRIGÉ ET BARÈME — document enseignant}\par\fi
\end{center}
\vspace{-2pt}
{\renewcommand{\arraystretch}{1.35}
\begin{tabularx}{\linewidth}{|X|X|p{2.2cm}|p{2.6cm}|}
  \hline
  \textbf{Nom :} & \textbf{Prénom :} & \textbf{Poste n°} & \textbf{Note :} \hfill /20 \\ \hline
\end{tabularx}}

\textbf{Consignes :} travail \textbf{individuel}, documents et Internet \textbf{non autorisés}. Ce sujet est le
\textbf{document-réponse}. Fichier Packet Tracer : \ip{Nom\_Prenom\_EvalLAN\_\variante.pkt}, à enregistrer
régulièrement dans le dossier de rendu. Les « visas » sont apposés par le professeur, que vous appelez à chaque test réussi.
Durées conseillées : A 30 min — B 1 h 10 — C 20 min.

%==========================================================================
\section{Partie A — Questions de cours et calculs \hfill /6}

\question{A1}{2 pts}{Complétez le tableau en considérant le \textbf{masque par défaut} de la classe de chaque adresse.}
{\small\renewcommand{\arraystretch}{1.15}
\begin{tabularx}{\linewidth}{|l|c|C|C|C|c|}
  \hline
  \textbf{Adresse IP} & \textbf{Classe} & \textbf{Masque par défaut} & \textbf{Adresse réseau} &
  \textbf{Adresse de diffusion} & \textbf{Nb d'hôtes} \\ \hline
  \AIrows
\end{tabularx}}
\bareme{0,5 pt par ligne juste + 0,5 pt si les trois nombres d'hôtes sont justes ($2^{n}-2$, $n$ = nombre de bits d'hôte).}

\question{A2}{1 pt}{Indiquez le type de câble (\textbf{droit} ou \textbf{croisé}) pour chaque liaison.}
{\small\renewcommand{\arraystretch}{1.15}
\begin{tabularx}{\linewidth}{|X|c||X|c|}
  \hline
  \hl PC $\leftrightarrow$ commutateur & \makebox[1.6cm]{\sol{droit}} & Commutateur $\leftrightarrow$ routeur & \makebox[1.6cm]{\sol{droit}} \\ \hline
  \hl PC $\leftrightarrow$ PC (réseau P2P) & \sol{croisé} & Concentrateur $\leftrightarrow$ concentrateur & \sol{croisé} \\ \hline
\end{tabularx}}
\bareme{0,25 pt par case. Équipements de types différents : câble droit ; de même type : câble croisé.}

\question{A3}{1 pt}{Comment un \textbf{concentrateur (hub)} et un \textbf{commutateur (switch)} traitent-ils une trame
reçue ? Lequel limite les collisions, et pourquoi ?}
\zone{2.9cm}{Le \textbf{hub} (couche 1) répète la trame sur \textbf{tous ses autres ports} : un seul domaine de collision
partagé. Le \textbf{switch} (couche 2) lit l'adresse MAC de destination et, grâce à sa \textbf{table MAC}, ne l'envoie
\textbf{que sur le port du destinataire} : un domaine de collision par port, c'est lui qui limite les collisions.
\emph{0,5 pt hub, 0,5 pt switch.}}

\question{A4}{1 pt}{Qu'est-ce que la \textbf{passerelle par défaut} d'un poste ? Quelle condition son adresse doit-elle respecter ?}
\zone{2.4cm}{Adresse IP de l'interface du \textbf{routeur} à laquelle le poste envoie les paquets destinés à un
\textbf{autre réseau}. Elle doit \textbf{appartenir au même réseau} que le poste. \emph{0,5 + 0,5 pt.}}

\question{A5}{0,5 pt}{Un portable configuré en DHCP affiche \ip{\apipa} / \ip{255.255.0.0}. Que s'est-il passé ?}
\zone{1.5cm}{Aucun \textbf{serveur DHCP} n'a répondu : le poste s'est attribué une adresse \textbf{APIPA} (169.254.0.0/16).}

\question{A6}{0,5 pt}{Quel protocole utilise \texttt{ping} ? Nommez les deux messages échangés et la couche OSI des adresses IP.}
\zone{1.5cm}{\textbf{ICMP} ; \textbf{Echo Request} / \textbf{Echo Reply} ; adresses IP en \textbf{couche 3 (Réseau)}.}

%==========================================================================
\newpage
\section{Partie B — Réalisation sous Packet Tracer \hfill /10}

\begin{cdc}
\entreprise{} vous confie son réseau : \textbf{trois réseaux} reliés par un \textbf{routeur unique Cisco 2911}
nommé \ip{\Rnom} (interfaces G0/0, G0/1, G0/2).
\begin{itemize}[leftmargin=12pt]
  \item \textbf{Réseau \resUn} : \textbf{4 postes} PC-\pUn1 à PC-\pUn4, commutateur \ip{SW-\pUn}, adressage statique,
        réseau \ip{\netUn.0/24}.
  \item \textbf{Réseau \resDeux} : \textbf{2 postes} PC-\pDeux1, PC-\pDeux2 et \textbf{1 imprimante} IMP-\pDeux,
        commutateur \ip{SW-\pDeux}, adressage statique, réseau de classe B \ip{\netDeux.0.0/16}.
  \item \textbf{Réseau Serveur} : serveur \ip{\srvNom} d'adresse \ip{\netSrv.\srvH/24}, relié \textbf{directement} au routeur.
  \item La passerelle de chaque réseau est sa \textbf{dernière adresse utilisable}. Tous les équipements communiquent
        entre eux et accèdent à la page web du serveur. Nommage imposé (« Display Name »).
\end{itemize}
\end{cdc}

\question{B1}{2 pts}{Complétez le plan d'adressage \textbf{avant} la réalisation.}
{\footnotesize\setlength{\tabcolsep}{3pt}\renewcommand{\arraystretch}{1.15}
\begin{tabularx}{\linewidth}{|l|C|C|C|C|C|C|}
  \hline
  \textbf{Réseau} & \textbf{Adresse réseau} & \textbf{Masque} & \textbf{1re @ utilisable} &
  \textbf{Dernière @ utilisable} & \textbf{Diffusion} & \textbf{Passerelle} \\ \hline
  \hl \resUn & \sol{\netUn.0} & \sol{255.255.255.0} & \sol{\netUn.1} & \sol{\netUn.254} & \sol{\netUn.255} & \sol{\netUn.254} \\ \hline
  \hl \resDeux & \sol{\netDeux.0.0} & \sol{255.255.0.0} & \sol{\netDeux.0.1} & \sol{\netDeux.255.254} & \sol{\netDeux.255.255} & \sol{\netDeux.255.254} \\ \hline
  \hl Serveur & \sol{\netSrv.0} & \sol{255.255.255.0} & \sol{\netSrv.1} & \sol{\netSrv.254} & \sol{\netSrv.255} & \sol{\netSrv.254} \\ \hline
\end{tabularx}}
\bareme{0,5 pt par ligne juste + 0,5 pt si la convention de passerelle est respectée partout.}

\question{B2}{2 pts}{Construisez la topologie (équipements, câblage, nommage) et relevez l'adressage configuré.}
{\footnotesize\setlength{\tabcolsep}{3pt}\renewcommand{\arraystretch}{1.1}
\begin{tabularx}{\linewidth}{|l|C|C|C||l|C|C|}
  \hline
  \textbf{Hôte} & \textbf{Adresse IP} & \textbf{Masque} & \textbf{Passerelle} & \textbf{Interface} & \textbf{Adresse IP} & \textbf{Masque} \\ \hline
  \hl PC-\pUn1 & \sol{\netUn.1} & \sol{255.255.255.0} & \sol{\netUn.254} & \Rnom{} G0/0 & \sol{\netUn.254} & \sol{255.255.255.0} \\ \hline
  \hl PC-\pDeux1 & \sol{\netDeux.0.1} & \sol{255.255.0.0} & \sol{\netDeux.255.254} & \Rnom{} G0/1 & \sol{\netDeux.255.254} & \sol{255.255.0.0} \\ \hline
  \hl IMP-\pDeux & \sol{\netDeux.0.10} & \sol{255.255.0.0} & \sol{\netDeux.255.254} & \Rnom{} G0/2 & \sol{\netSrv.254} & \sol{255.255.255.0} \\ \hline
  \hl \srvNom & \ip{\netSrv.\srvH} & \sol{255.255.255.0} & \sol{\netSrv.254} & \multicolumn{3}{c|}{\cellcolor{black!8}} \\ \hline
\end{tabularx}}
\bareme{Adresses d'hôtes libres (uniques, dans le bon réseau) ; affectation G0/x libre. Sur le .pkt : équipements 0,5 ;
câblage 1 (câbles droits vers les switchs et switch–routeur, \textbf{croisé} serveur–routeur, liens verts) ; nommage 0,5.}

\question{B3}{3 pts}{Configurez tous les équipements : IP, masques, passerelles des hôtes ; interfaces du routeur
(« Port Status : On »).}
\bareme{Sur le .pkt : interfaces routeur 1 pt ; IP/masques des hôtes 1 pt ; passerelles (imprimante et serveur compris)
1 pt. $-0,25$ par erreur, sans descendre sous 0 par item.}

\question{B4}{3 pts}{Réalisez les tests, notez la commande et le résultat, puis faites viser chaque test réussi.}
{\small\renewcommand{\arraystretch}{1.1}
\begin{tabularx}{\linewidth}{|c|l|l|C|c|c|}
  \hline
  \textbf{N°} & \textbf{Source} & \textbf{Destination} & \textbf{Commande / action} & \textbf{Résultat} & \textbf{Visa prof.} \\ \hline
  \hl 1 & PC-\pUn1 & PC-\pUn4 & \sol{ping \netUn.4} & \sol{réussi} & \hspace{1.4cm} \\ \hline
  \hl 2 & PC-\pDeux1 & IMP-\pDeux & \sol{ping \netDeux.0.10} & \sol{réussi} & \\ \hline
  \hl 3 & PC-\pUn2 & PC-\pDeux2 & \sol{ping \netDeux.0.2} & \sol{réussi} & \\ \hline
  \hl 4 & PC-\pDeux2 & \srvNom & \sol{ping \netSrv.\srvH} & \sol{réussi} & \\ \hline
  \hl 5 & PC-\pUn3 & page web du serveur & \sol{navigateur : {\NoAutoSpacing http://\netSrv.\srvH}} & \sol{page affichée} & \\ \hline
\end{tabularx}}
\bareme{0,5 pt par test visé + 0,5 pt pour B5. Le premier ping inter-réseaux peut échouer (résolution ARP) ; les suivants doivent réussir.}

\question{B5}{incl. B4}{Test n°3 : par quels équipements passe le paquet ? Pourquoi PC-\pUn2 l'envoie-t-il au routeur ?}
\zone{2.3cm}{PC-\pUn2 $\to$ SW-\pUn{} $\to$ \Rnom{} $\to$ SW-\pDeux{} $\to$ PC-\pDeux2. Grâce à son masque, PC-\pUn2 constate que
la destination (\netDeux.0.0/16) n'est \textbf{pas dans son réseau} (\netUn.0/24) : il envoie le paquet à sa
\textbf{passerelle par défaut}, qui le route.}

%==========================================================================
\newpage
\section{Partie C — Diagnostic d'un réseau en panne \hfill /4}

Un technicien constate qu'\textbf{aucun poste ne parvient à joindre le serveur SRV-C} et que certains postes
communiquent mal entre eux. Les câbles sont du bon type.

\begin{center}
\begin{minipage}{0.37\linewidth}\centering
\begin{tikzpicture}[font=\footnotesize,
  eq/.style={draw,fill=qbg,rounded corners=2pt,minimum width=9mm,minimum height=5mm,inner sep=2pt},
  every path/.style={thick}]
  \node[eq] (SW) {SW-1};
  \node[eq,right=9mm of SW] (R) {R1};
  \node[eq,right=7mm of R] (SRV) {SRV-C};
  \node[eq,above left=5mm and -3mm of SW] (A) {PC-A};
  \node[eq,above right=5mm and -3mm of SW] (B) {PC-B};
  \node[eq,below left=5mm and -3mm of SW] (C) {PC-C};
  \node[eq,below right=5mm and -3mm of SW] (D) {PC-D};
  \foreach \n in {A,B,C,D} \draw (\n) -- (SW);
  \draw (SW) -- node[above,font=\tiny]{G0/0} (R);
  \draw (R) -- node[above,font=\tiny]{G0/1} (SRV);
\end{tikzpicture}
\end{minipage}\hfill
\begin{minipage}{0.61\linewidth}\centering
{\small\renewcommand{\arraystretch}{1.05}\setlength{\tabcolsep}{4pt}
\begin{tabular}{lllll}
  \toprule
  \textbf{Équipement} & \textbf{Adresse IP} & \textbf{Masque} & \textbf{Passerelle} & \textbf{Port} \\
  \midrule
  \Crows
  \bottomrule
\end{tabular}}
\end{minipage}
\end{center}

\question{C1}{4 pts}{Relevez les \textbf{quatre anomalies}. Pour chacune : conséquence et correction.}
{\small\renewcommand{\arraystretch}{1.1}
\ifcorrige\def\hc{}\else\def\hc{\rule[-10mm]{0pt}{14mm}}\fi
\begin{tabularx}{\linewidth}{|c|>{\raggedright\arraybackslash}p{2cm}|L|L|L|}
  \hline
  \textbf{N°} & \textbf{Équipement} & \textbf{Anomalie} & \textbf{Conséquence} & \textbf{Correction} \\ \hline
  \ifcorrige\Csol\else
  \hc 1 & & & & \\ \hline
  \hc 2 & & & & \\ \hline
  \hc 3 & & & & \\ \hline
  \hc 4 & & & & \\ \hline
  \fi
\end{tabularx}}
\bareme{1 pt par anomalie (0,25 équipement, anomalie, conséquence, correction ; ordre indifférent). \Cbareme}

%==========================================================================
\section{Bonus — Extension sans fil \hfill +1}

\begin{cdc}Ajoutez au réseau \resUn{} un \textbf{AccessPoint-PT} relié à SW-\pUn{} (SSID \ip{\ssid}, \textbf{WPA2-PSK})
et un portable \ip{Laptop-\pUn} équipé d'un module \texttt{WPC300N}, en adressage \textbf{statique}.
Vérifiez qu'il joint \srvNom.\end{cdc}
{\small\renewcommand{\arraystretch}{1.1}
\begin{tabularx}{\linewidth}{|l|C|C|C|c|}
  \hline
  \textbf{Équipement} & \textbf{Adresse IP} & \textbf{Masque} & \textbf{Passerelle} & \textbf{Visa prof.} \\ \hline
  \hl Laptop-\pUn & \sol{\netUn.50} & \sol{255.255.255.0} & \sol{\netUn.254} & \hspace{1.4cm} \\ \hline
\end{tabularx}}
\bareme{Le point d'accès fait office de pont : le portable est dans \netUn.0/24. Bonus si associé en WPA2-PSK et ping
vers \netSrv.\srvH{} réussi. Note plafonnée à 20.}

\end{document}
"
\documentclass[10pt,a4paper]{article}

\usepackage[T1]{fontenc}
\usepackage[utf8]{inputenc}
\usepackage{lmodern}
\usepackage[french]{babel}
\usepackage[a4paper,hmargin=1.5cm,top=1.3cm,bottom=1.6cm,footskip=0.7cm]{geometry}
\usepackage[table]{xcolor}
\usepackage{booktabs,tabularx,array}
\usepackage{enumitem}
\usepackage{fancyhdr}
\usepackage[most]{tcolorbox}
\usepackage{tikz}
\usetikzlibrary{positioning}
\usepackage{titlesec}

\providecommand{\modecorrige}{0}
\providecommand{\variante}{A}
\newif\ifcorrige
\ifnum\modecorrige=1 \corrigetrue \else \corrigefalse \fi

\definecolor{ipred}{HTML}{C0392B}
\definecolor{cdcbar}{HTML}{A67C3D}
\definecolor{cdcbg}{HTML}{F4F1EA}
\definecolor{qbar}{HTML}{4A6B5D}
\definecolor{qbg}{HTML}{EDF1EF}
\definecolor{corr}{HTML}{1F5FA8}
\definecolor{corrbg}{HTML}{E8F0FA}
\definecolor{grisfoot}{HTML}{828282}

\newcommand{\ip}[1]{\textcolor{ipred}{\texttt{#1}}}
\newcommand{\pts}[1]{\hfill\mbox{\textbf{(#1)}}}
\newcommand{\sol}[1]{\ifcorrige\textcolor{corr}{#1}\fi}
\newcommand{\bareme}[1]{\ifcorrige\par{\footnotesize\color{corr}#1\par}\fi}
\newcommand{\hl}{\ifcorrige\rule[-1.6mm]{0pt}{5.4mm}\else\rule[-2.2mm]{0pt}{6.6mm}\fi}
\newcolumntype{C}{>{\centering\arraybackslash}X}
\newcolumntype{L}{>{\raggedright\arraybackslash}X}

\titleformat{\section}{\large\bfseries}{}{0pt}{}
\titlespacing*{\section}{0pt}{6pt}{3pt}

\tcbset{qstyle/.style={enhanced,boxrule=0pt,leftrule=2.5pt,arc=0pt,outer arc=0pt,
  left=5pt,right=5pt,top=2pt,bottom=2pt,before skip=4pt,after skip=3pt}}
\newtcolorbox{cdc}{qstyle,colback=cdcbg,colframe=cdcbar}
\newtcolorbox{quest}{qstyle,colback=qbg,colframe=qbar}

% Zone de réponse : vierge dans le sujet, remplie dans le corrigé
\newcommand{\zone}[2]{%
  \ifcorrige
    \begin{tcolorbox}[colback=corrbg,colframe=corr,boxrule=0.5pt,arc=1pt,left=4pt,right=4pt,top=1pt,bottom=1pt,
      before skip=2pt,after skip=2pt,fontupper=\footnotesize\color{corr}]#2\end{tcolorbox}%
  \else
    \begin{tcolorbox}[colback=white,colframe=black!45,boxrule=0.4pt,arc=1pt,height=#1,
      before skip=2pt,after skip=2pt]\end{tcolorbox}%
  \fi}

\newcommand{\question}[3]{\begin{quest}\textbf{#1.} #3\pts{#2}\end{quest}}

%==========================================================================
% Données propres à chaque variante
%==========================================================================
\if\variante A
  \def\entreprise{Le cabinet d'architectes \textbf{ArchiNova}}
  \def\Rnom{R-ArchiNova}
  \def\resUn{Bureau d'études}\def\pUn{BE}\def\netUn{192.168.40}
  \def\resDeux{Accueil}\def\pDeux{ACC}\def\netDeux{172.16}
  \def\srvNom{SRV-Intranet}\def\netSrv{192.168.200}\def\srvH{10}
  \def\ssid{ArchiNova-WiFi}
  \def\AIrows{%
    \hl\ip{192.168.40.25} & \sol{C} & \sol{255.255.255.0} & \sol{192.168.40.0} & \sol{192.168.40.255} & \sol{254} \\ \hline
    \hl\ip{172.16.5.10}   & \sol{B} & \sol{255.255.0.0}   & \sol{172.16.0.0}   & \sol{172.16.255.255} & \sol{65\,534} \\ \hline
    \hl\ip{10.3.2.1}      & \sol{A} & \sol{255.0.0.0}     & \sol{10.0.0.0}     & \sol{10.255.255.255} & \sol{16\,777\,214} \\ \hline}
  \def\apipa{169.254.12.7}
  \def\Cnet{192.168.60}
  \def\Crows{%
    PC-A & \ip{192.168.60.10} & \ip{255.255.255.0} & \ip{192.168.60.254} & — \\
    PC-B & \ip{192.168.60.20} & \ip{255.255.255.0} & \ip{192.168.60.1}   & — \\
    PC-C & \ip{192.168.61.30} & \ip{255.255.255.0} & \ip{192.168.60.254} & — \\
    PC-D & \ip{192.168.60.10} & \ip{255.255.255.0} & \ip{192.168.60.254} & — \\
    SRV-C & \ip{10.0.0.5}     & \ip{255.0.0.0}     & \ip{10.0.0.254}     & — \\
    R1 — G0/0 & \ip{192.168.60.254} & \ip{255.255.255.0} & — & On \\
    R1 — G0/1 & \ip{10.0.0.254}     & \ip{255.0.0.0}     & — & Off \\}
  \def\Csol{%
    \hl 1 & \sol{PC-B} & \sol{Passerelle .60.1 : pas l'adresse du routeur.} & \sol{Joint son réseau, mais aucun autre (serveur).} & \sol{Passerelle 192.168.60.254.} \\ \hline
    \hl 2 & \sol{PC-C} & \sol{IP hors du réseau 192.168.60.0/24.} & \sol{Ne joint ni les postes ni le routeur.} & \sol{Ex. 192.168.60.30.} \\ \hline
    \hl 3 & \sol{PC-D} & \sol{Même IP que PC-A (doublon).} & \sol{Conflit d'adresse : PC-A et PC-D perturbés.} & \sol{Adresse unique, ex. .60.40.} \\ \hline
    \hl 4 & \sol{R1 G0/1} & \sol{Interface désactivée (Off).} & \sol{Réseau du serveur injoignable pour tous.} & \sol{Port Status : On.} \\ \hline}
  \def\Cbareme{PC-A et SRV-C sont correctement configurés ; accepter PC-A cité à la place de PC-D pour le doublon.}
\else
  \def\entreprise{La clinique vétérinaire \textbf{VetoPlus}}
  \def\Rnom{R-VetoPlus}
  \def\resUn{Soins}\def\pUn{SOI}\def\netUn{192.168.70}
  \def\resDeux{Secrétariat}\def\pDeux{SEC}\def\netDeux{172.20}
  \def\srvNom{SRV-Dossiers}\def\netSrv{192.168.250}\def\srvH{20}
  \def\ssid{VetoPlus-WiFi}
  \def\AIrows{%
    \hl\ip{192.168.15.130} & \sol{C} & \sol{255.255.255.0} & \sol{192.168.15.0} & \sol{192.168.15.255} & \sol{254} \\ \hline
    \hl\ip{172.31.200.7}   & \sol{B} & \sol{255.255.0.0}   & \sol{172.31.0.0}   & \sol{172.31.255.255} & \sol{65\,534} \\ \hline
    \hl\ip{10.45.6.9}      & \sol{A} & \sol{255.0.0.0}     & \sol{10.0.0.0}     & \sol{10.255.255.255} & \sol{16\,777\,214} \\ \hline}
  \def\apipa{169.254.88.3}
  \def\Cnet{192.168.80}
  \def\Crows{%
    PC-A & \ip{192.168.8.11}   & \ip{255.255.255.0} & \ip{192.168.80.254} & — \\
    PC-B & \ip{192.168.80.12}  & \ip{255.255.255.0} & \ip{192.168.80.254} & — \\
    PC-C & \ip{192.168.80.13}  & \ip{255.255.255.0} & \ip{192.168.80.100} & — \\
    PC-D & \ip{192.168.80.12}  & \ip{255.255.255.0} & \ip{192.168.80.254} & — \\
    SRV-C & \ip{10.0.0.8}      & \ip{255.0.0.0}     & \ip{10.0.0.254}     & — \\
    R1 — G0/0 & \ip{192.168.80.254} & \ip{255.255.255.0} & — & Off \\
    R1 — G0/1 & \ip{10.0.0.254}     & \ip{255.0.0.0}     & — & On \\}
  \def\Csol{%
    \hl 1 & \sol{PC-A} & \sol{IP hors du réseau 192.168.80.0/24.} & \sol{Ne joint ni les postes ni le routeur.} & \sol{Ex. 192.168.80.11.} \\ \hline
    \hl 2 & \sol{PC-C} & \sol{Passerelle .80.100 : pas l'adresse du routeur.} & \sol{Joint son réseau, mais aucun autre (serveur).} & \sol{Passerelle 192.168.80.254.} \\ \hline
    \hl 3 & \sol{PC-D} & \sol{Même IP que PC-B (doublon).} & \sol{Conflit d'adresse : PC-B et PC-D perturbés.} & \sol{Adresse unique, ex. .80.14.} \\ \hline
    \hl 4 & \sol{R1 G0/0} & \sol{Interface côté LAN désactivée (Off).} & \sol{Aucun poste ne joint le routeur ni le serveur.} & \sol{Port Status : On.} \\ \hline}
  \def\Cbareme{PC-B et SRV-C sont correctement configurés ; accepter PC-B cité à la place de PC-D pour le doublon.}
\fi

\pagestyle{fancy}
\fancyhf{}
\renewcommand{\headrulewidth}{0pt}
\fancyfoot[L]{\footnotesize\itshape\color{grisfoot}Évaluation — Réseaux locaux — Sujet \variante\ifcorrige\ — CORRIGÉ\fi}
\fancyfoot[R]{\footnotesize\color{grisfoot}\thepage/3}

\setlist{itemsep=0pt,topsep=2pt,parsep=0pt}
\setlength{\parindent}{0pt}
\setlength{\parskip}{2pt}

\begin{document}

\begin{center}
  {\Large\bfseries Évaluation pratique — Réseaux locaux \quad
   \fcolorbox{black}{white}{\ SUJET \variante\ }\par}
  \vspace{1pt}
  {\bfseries Adressage IP, commutation et routage — BTS CIEL 1re année — durée : 2 heures\par}
  \ifcorrige\vspace{2pt}{\bfseries\color{corr}CORRIGÉ ET BARÈME — document enseignant}\par\fi
\end{center}
\vspace{-2pt}
{\renewcommand{\arraystretch}{1.35}
\begin{tabularx}{\linewidth}{|X|X|p{2.2cm}|p{2.6cm}|}
  \hline
  \textbf{Nom :} & \textbf{Prénom :} & \textbf{Poste n°} & \textbf{Note :} \hfill /20 \\ \hline
\end{tabularx}}

\textbf{Consignes :} travail \textbf{individuel}, documents et Internet \textbf{non autorisés}. Ce sujet est le
\textbf{document-réponse}. Fichier Packet Tracer : \ip{Nom\_Prenom\_EvalLAN\_\variante.pkt}, à enregistrer
régulièrement dans le dossier de rendu. Les « visas » sont apposés par le professeur, que vous appelez à chaque test réussi.
Durées conseillées : A 30 min — B 1 h 10 — C 20 min.

%==========================================================================
\section{Partie A — Questions de cours et calculs \hfill /6}

\question{A1}{2 pts}{Complétez le tableau en considérant le \textbf{masque par défaut} de la classe de chaque adresse.}
{\small\renewcommand{\arraystretch}{1.15}
\begin{tabularx}{\linewidth}{|l|c|C|C|C|c|}
  \hline
  \textbf{Adresse IP} & \textbf{Classe} & \textbf{Masque par défaut} & \textbf{Adresse réseau} &
  \textbf{Adresse de diffusion} & \textbf{Nb d'hôtes} \\ \hline
  \AIrows
\end{tabularx}}
\bareme{0,5 pt par ligne juste + 0,5 pt si les trois nombres d'hôtes sont justes ($2^{n}-2$, $n$ = nombre de bits d'hôte).}

\question{A2}{1 pt}{Indiquez le type de câble (\textbf{droit} ou \textbf{croisé}) pour chaque liaison.}
{\small\renewcommand{\arraystretch}{1.15}
\begin{tabularx}{\linewidth}{|X|c||X|c|}
  \hline
  \hl PC $\leftrightarrow$ commutateur & \makebox[1.6cm]{\sol{droit}} & Commutateur $\leftrightarrow$ routeur & \makebox[1.6cm]{\sol{droit}} \\ \hline
  \hl PC $\leftrightarrow$ PC (réseau P2P) & \sol{croisé} & Concentrateur $\leftrightarrow$ concentrateur & \sol{croisé} \\ \hline
\end{tabularx}}
\bareme{0,25 pt par case. Équipements de types différents : câble droit ; de même type : câble croisé.}

\question{A3}{1 pt}{Comment un \textbf{concentrateur (hub)} et un \textbf{commutateur (switch)} traitent-ils une trame
reçue ? Lequel limite les collisions, et pourquoi ?}
\zone{2.9cm}{Le \textbf{hub} (couche 1) répète la trame sur \textbf{tous ses autres ports} : un seul domaine de collision
partagé. Le \textbf{switch} (couche 2) lit l'adresse MAC de destination et, grâce à sa \textbf{table MAC}, ne l'envoie
\textbf{que sur le port du destinataire} : un domaine de collision par port, c'est lui qui limite les collisions.
\emph{0,5 pt hub, 0,5 pt switch.}}

\question{A4}{1 pt}{Qu'est-ce que la \textbf{passerelle par défaut} d'un poste ? Quelle condition son adresse doit-elle respecter ?}
\zone{2.4cm}{Adresse IP de l'interface du \textbf{routeur} à laquelle le poste envoie les paquets destinés à un
\textbf{autre réseau}. Elle doit \textbf{appartenir au même réseau} que le poste. \emph{0,5 + 0,5 pt.}}

\question{A5}{0,5 pt}{Un portable configuré en DHCP affiche \ip{\apipa} / \ip{255.255.0.0}. Que s'est-il passé ?}
\zone{1.5cm}{Aucun \textbf{serveur DHCP} n'a répondu : le poste s'est attribué une adresse \textbf{APIPA} (169.254.0.0/16).}

\question{A6}{0,5 pt}{Quel protocole utilise \texttt{ping} ? Nommez les deux messages échangés et la couche OSI des adresses IP.}
\zone{1.5cm}{\textbf{ICMP} ; \textbf{Echo Request} / \textbf{Echo Reply} ; adresses IP en \textbf{couche 3 (Réseau)}.}

%==========================================================================
\newpage
\section{Partie B — Réalisation sous Packet Tracer \hfill /10}

\begin{cdc}
\entreprise{} vous confie son réseau : \textbf{trois réseaux} reliés par un \textbf{routeur unique Cisco 2911}
nommé \ip{\Rnom} (interfaces G0/0, G0/1, G0/2).
\begin{itemize}[leftmargin=12pt]
  \item \textbf{Réseau \resUn} : \textbf{4 postes} PC-\pUn1 à PC-\pUn4, commutateur \ip{SW-\pUn}, adressage statique,
        réseau \ip{\netUn.0/24}.
  \item \textbf{Réseau \resDeux} : \textbf{2 postes} PC-\pDeux1, PC-\pDeux2 et \textbf{1 imprimante} IMP-\pDeux,
        commutateur \ip{SW-\pDeux}, adressage statique, réseau de classe B \ip{\netDeux.0.0/16}.
  \item \textbf{Réseau Serveur} : serveur \ip{\srvNom} d'adresse \ip{\netSrv.\srvH/24}, relié \textbf{directement} au routeur.
  \item La passerelle de chaque réseau est sa \textbf{dernière adresse utilisable}. Tous les équipements communiquent
        entre eux et accèdent à la page web du serveur. Nommage imposé (« Display Name »).
\end{itemize}
\end{cdc}

\question{B1}{2 pts}{Complétez le plan d'adressage \textbf{avant} la réalisation.}
{\footnotesize\setlength{\tabcolsep}{3pt}\renewcommand{\arraystretch}{1.15}
\begin{tabularx}{\linewidth}{|l|C|C|C|C|C|C|}
  \hline
  \textbf{Réseau} & \textbf{Adresse réseau} & \textbf{Masque} & \textbf{1re @ utilisable} &
  \textbf{Dernière @ utilisable} & \textbf{Diffusion} & \textbf{Passerelle} \\ \hline
  \hl \resUn & \sol{\netUn.0} & \sol{255.255.255.0} & \sol{\netUn.1} & \sol{\netUn.254} & \sol{\netUn.255} & \sol{\netUn.254} \\ \hline
  \hl \resDeux & \sol{\netDeux.0.0} & \sol{255.255.0.0} & \sol{\netDeux.0.1} & \sol{\netDeux.255.254} & \sol{\netDeux.255.255} & \sol{\netDeux.255.254} \\ \hline
  \hl Serveur & \sol{\netSrv.0} & \sol{255.255.255.0} & \sol{\netSrv.1} & \sol{\netSrv.254} & \sol{\netSrv.255} & \sol{\netSrv.254} \\ \hline
\end{tabularx}}
\bareme{0,5 pt par ligne juste + 0,5 pt si la convention de passerelle est respectée partout.}

\question{B2}{2 pts}{Construisez la topologie (équipements, câblage, nommage) et relevez l'adressage configuré.}
{\footnotesize\setlength{\tabcolsep}{3pt}\renewcommand{\arraystretch}{1.1}
\begin{tabularx}{\linewidth}{|l|C|C|C||l|C|C|}
  \hline
  \textbf{Hôte} & \textbf{Adresse IP} & \textbf{Masque} & \textbf{Passerelle} & \textbf{Interface} & \textbf{Adresse IP} & \textbf{Masque} \\ \hline
  \hl PC-\pUn1 & \sol{\netUn.1} & \sol{255.255.255.0} & \sol{\netUn.254} & \Rnom{} G0/0 & \sol{\netUn.254} & \sol{255.255.255.0} \\ \hline
  \hl PC-\pDeux1 & \sol{\netDeux.0.1} & \sol{255.255.0.0} & \sol{\netDeux.255.254} & \Rnom{} G0/1 & \sol{\netDeux.255.254} & \sol{255.255.0.0} \\ \hline
  \hl IMP-\pDeux & \sol{\netDeux.0.10} & \sol{255.255.0.0} & \sol{\netDeux.255.254} & \Rnom{} G0/2 & \sol{\netSrv.254} & \sol{255.255.255.0} \\ \hline
  \hl \srvNom & \ip{\netSrv.\srvH} & \sol{255.255.255.0} & \sol{\netSrv.254} & \multicolumn{3}{c|}{\cellcolor{black!8}} \\ \hline
\end{tabularx}}
\bareme{Adresses d'hôtes libres (uniques, dans le bon réseau) ; affectation G0/x libre. Sur le .pkt : équipements 0,5 ;
câblage 1 (câbles droits vers les switchs et switch–routeur, \textbf{croisé} serveur–routeur, liens verts) ; nommage 0,5.}

\question{B3}{3 pts}{Configurez tous les équipements : IP, masques, passerelles des hôtes ; interfaces du routeur
(« Port Status : On »).}
\bareme{Sur le .pkt : interfaces routeur 1 pt ; IP/masques des hôtes 1 pt ; passerelles (imprimante et serveur compris)
1 pt. $-0,25$ par erreur, sans descendre sous 0 par item.}

\question{B4}{3 pts}{Réalisez les tests, notez la commande et le résultat, puis faites viser chaque test réussi.}
{\small\renewcommand{\arraystretch}{1.1}
\begin{tabularx}{\linewidth}{|c|l|l|C|c|c|}
  \hline
  \textbf{N°} & \textbf{Source} & \textbf{Destination} & \textbf{Commande / action} & \textbf{Résultat} & \textbf{Visa prof.} \\ \hline
  \hl 1 & PC-\pUn1 & PC-\pUn4 & \sol{ping \netUn.4} & \sol{réussi} & \hspace{1.4cm} \\ \hline
  \hl 2 & PC-\pDeux1 & IMP-\pDeux & \sol{ping \netDeux.0.10} & \sol{réussi} & \\ \hline
  \hl 3 & PC-\pUn2 & PC-\pDeux2 & \sol{ping \netDeux.0.2} & \sol{réussi} & \\ \hline
  \hl 4 & PC-\pDeux2 & \srvNom & \sol{ping \netSrv.\srvH} & \sol{réussi} & \\ \hline
  \hl 5 & PC-\pUn3 & page web du serveur & \sol{navigateur : {\NoAutoSpacing http://\netSrv.\srvH}} & \sol{page affichée} & \\ \hline
\end{tabularx}}
\bareme{0,5 pt par test visé + 0,5 pt pour B5. Le premier ping inter-réseaux peut échouer (résolution ARP) ; les suivants doivent réussir.}

\question{B5}{incl. B4}{Test n°3 : par quels équipements passe le paquet ? Pourquoi PC-\pUn2 l'envoie-t-il au routeur ?}
\zone{2.3cm}{PC-\pUn2 $\to$ SW-\pUn{} $\to$ \Rnom{} $\to$ SW-\pDeux{} $\to$ PC-\pDeux2. Grâce à son masque, PC-\pUn2 constate que
la destination (\netDeux.0.0/16) n'est \textbf{pas dans son réseau} (\netUn.0/24) : il envoie le paquet à sa
\textbf{passerelle par défaut}, qui le route.}

%==========================================================================
\newpage
\section{Partie C — Diagnostic d'un réseau en panne \hfill /4}

Un technicien constate qu'\textbf{aucun poste ne parvient à joindre le serveur SRV-C} et que certains postes
communiquent mal entre eux. Les câbles sont du bon type.

\begin{center}
\begin{minipage}{0.37\linewidth}\centering
\begin{tikzpicture}[font=\footnotesize,
  eq/.style={draw,fill=qbg,rounded corners=2pt,minimum width=9mm,minimum height=5mm,inner sep=2pt},
  every path/.style={thick}]
  \node[eq] (SW) {SW-1};
  \node[eq,right=9mm of SW] (R) {R1};
  \node[eq,right=7mm of R] (SRV) {SRV-C};
  \node[eq,above left=5mm and -3mm of SW] (A) {PC-A};
  \node[eq,above right=5mm and -3mm of SW] (B) {PC-B};
  \node[eq,below left=5mm and -3mm of SW] (C) {PC-C};
  \node[eq,below right=5mm and -3mm of SW] (D) {PC-D};
  \foreach \n in {A,B,C,D} \draw (\n) -- (SW);
  \draw (SW) -- node[above,font=\tiny]{G0/0} (R);
  \draw (R) -- node[above,font=\tiny]{G0/1} (SRV);
\end{tikzpicture}
\end{minipage}\hfill
\begin{minipage}{0.61\linewidth}\centering
{\small\renewcommand{\arraystretch}{1.05}\setlength{\tabcolsep}{4pt}
\begin{tabular}{lllll}
  \toprule
  \textbf{Équipement} & \textbf{Adresse IP} & \textbf{Masque} & \textbf{Passerelle} & \textbf{Port} \\
  \midrule
  \Crows
  \bottomrule
\end{tabular}}
\end{minipage}
\end{center}

\question{C1}{4 pts}{Relevez les \textbf{quatre anomalies}. Pour chacune : conséquence et correction.}
{\small\renewcommand{\arraystretch}{1.1}
\ifcorrige\def\hc{}\else\def\hc{\rule[-10mm]{0pt}{14mm}}\fi
\begin{tabularx}{\linewidth}{|c|>{\raggedright\arraybackslash}p{2cm}|L|L|L|}
  \hline
  \textbf{N°} & \textbf{Équipement} & \textbf{Anomalie} & \textbf{Conséquence} & \textbf{Correction} \\ \hline
  \ifcorrige\Csol\else
  \hc 1 & & & & \\ \hline
  \hc 2 & & & & \\ \hline
  \hc 3 & & & & \\ \hline
  \hc 4 & & & & \\ \hline
  \fi
\end{tabularx}}
\bareme{1 pt par anomalie (0,25 équipement, anomalie, conséquence, correction ; ordre indifférent). \Cbareme}

%==========================================================================
\section{Bonus — Extension sans fil \hfill +1}

\begin{cdc}Ajoutez au réseau \resUn{} un \textbf{AccessPoint-PT} relié à SW-\pUn{} (SSID \ip{\ssid}, \textbf{WPA2-PSK})
et un portable \ip{Laptop-\pUn} équipé d'un module \texttt{WPC300N}, en adressage \textbf{statique}.
Vérifiez qu'il joint \srvNom.\end{cdc}
{\small\renewcommand{\arraystretch}{1.1}
\begin{tabularx}{\linewidth}{|l|C|C|C|c|}
  \hline
  \textbf{Équipement} & \textbf{Adresse IP} & \textbf{Masque} & \textbf{Passerelle} & \textbf{Visa prof.} \\ \hline
  \hl Laptop-\pUn & \sol{\netUn.50} & \sol{255.255.255.0} & \sol{\netUn.254} & \hspace{1.4cm} \\ \hline
\end{tabularx}}
\bareme{Le point d'accès fait office de pont : le portable est dans \netUn.0/24. Bonus si associé en WPA2-PSK et ping
vers \netSrv.\srvH{} réussi. Note plafonnée à 20.}

\end{document}
"
%   pdflatex -jobname=Eval_Reseau_CorrigeB "\def\variante{B}\def\modecorrige{1}% Évaluation pratique — Réseaux locaux (BTS CIEL 1re année) — sujets A et B
% Compilation (4 PDF) :
%   pdflatex -jobname=Eval_Reseau_SujetA   "\def\variante{A}\def\modecorrige{0}% Évaluation pratique — Réseaux locaux (BTS CIEL 1re année) — sujets A et B
% Compilation (4 PDF) :
%   pdflatex -jobname=Eval_Reseau_SujetA   "\def\variante{A}\def\modecorrige{0}\input{Eval_TP_Reseau_LAN.tex}"
%   pdflatex -jobname=Eval_Reseau_SujetB   "\def\variante{B}\def\modecorrige{0}\input{Eval_TP_Reseau_LAN.tex}"
%   pdflatex -jobname=Eval_Reseau_CorrigeA "\def\variante{A}\def\modecorrige{1}\input{Eval_TP_Reseau_LAN.tex}"
%   pdflatex -jobname=Eval_Reseau_CorrigeB "\def\variante{B}\def\modecorrige{1}\input{Eval_TP_Reseau_LAN.tex}"
\documentclass[10pt,a4paper]{article}

\usepackage[T1]{fontenc}
\usepackage[utf8]{inputenc}
\usepackage{lmodern}
\usepackage[french]{babel}
\usepackage[a4paper,hmargin=1.5cm,top=1.3cm,bottom=1.6cm,footskip=0.7cm]{geometry}
\usepackage[table]{xcolor}
\usepackage{booktabs,tabularx,array}
\usepackage{enumitem}
\usepackage{fancyhdr}
\usepackage[most]{tcolorbox}
\usepackage{tikz}
\usetikzlibrary{positioning}
\usepackage{titlesec}

\providecommand{\modecorrige}{0}
\providecommand{\variante}{A}
\newif\ifcorrige
\ifnum\modecorrige=1 \corrigetrue \else \corrigefalse \fi

\definecolor{ipred}{HTML}{C0392B}
\definecolor{cdcbar}{HTML}{A67C3D}
\definecolor{cdcbg}{HTML}{F4F1EA}
\definecolor{qbar}{HTML}{4A6B5D}
\definecolor{qbg}{HTML}{EDF1EF}
\definecolor{corr}{HTML}{1F5FA8}
\definecolor{corrbg}{HTML}{E8F0FA}
\definecolor{grisfoot}{HTML}{828282}

\newcommand{\ip}[1]{\textcolor{ipred}{\texttt{#1}}}
\newcommand{\pts}[1]{\hfill\mbox{\textbf{(#1)}}}
\newcommand{\sol}[1]{\ifcorrige\textcolor{corr}{#1}\fi}
\newcommand{\bareme}[1]{\ifcorrige\par{\footnotesize\color{corr}#1\par}\fi}
\newcommand{\hl}{\ifcorrige\rule[-1.6mm]{0pt}{5.4mm}\else\rule[-2.2mm]{0pt}{6.6mm}\fi}
\newcolumntype{C}{>{\centering\arraybackslash}X}
\newcolumntype{L}{>{\raggedright\arraybackslash}X}

\titleformat{\section}{\large\bfseries}{}{0pt}{}
\titlespacing*{\section}{0pt}{6pt}{3pt}

\tcbset{qstyle/.style={enhanced,boxrule=0pt,leftrule=2.5pt,arc=0pt,outer arc=0pt,
  left=5pt,right=5pt,top=2pt,bottom=2pt,before skip=4pt,after skip=3pt}}
\newtcolorbox{cdc}{qstyle,colback=cdcbg,colframe=cdcbar}
\newtcolorbox{quest}{qstyle,colback=qbg,colframe=qbar}

% Zone de réponse : vierge dans le sujet, remplie dans le corrigé
\newcommand{\zone}[2]{%
  \ifcorrige
    \begin{tcolorbox}[colback=corrbg,colframe=corr,boxrule=0.5pt,arc=1pt,left=4pt,right=4pt,top=1pt,bottom=1pt,
      before skip=2pt,after skip=2pt,fontupper=\footnotesize\color{corr}]#2\end{tcolorbox}%
  \else
    \begin{tcolorbox}[colback=white,colframe=black!45,boxrule=0.4pt,arc=1pt,height=#1,
      before skip=2pt,after skip=2pt]\end{tcolorbox}%
  \fi}

\newcommand{\question}[3]{\begin{quest}\textbf{#1.} #3\pts{#2}\end{quest}}

%==========================================================================
% Données propres à chaque variante
%==========================================================================
\if\variante A
  \def\entreprise{Le cabinet d'architectes \textbf{ArchiNova}}
  \def\Rnom{R-ArchiNova}
  \def\resUn{Bureau d'études}\def\pUn{BE}\def\netUn{192.168.40}
  \def\resDeux{Accueil}\def\pDeux{ACC}\def\netDeux{172.16}
  \def\srvNom{SRV-Intranet}\def\netSrv{192.168.200}\def\srvH{10}
  \def\ssid{ArchiNova-WiFi}
  \def\AIrows{%
    \hl\ip{192.168.40.25} & \sol{C} & \sol{255.255.255.0} & \sol{192.168.40.0} & \sol{192.168.40.255} & \sol{254} \\ \hline
    \hl\ip{172.16.5.10}   & \sol{B} & \sol{255.255.0.0}   & \sol{172.16.0.0}   & \sol{172.16.255.255} & \sol{65\,534} \\ \hline
    \hl\ip{10.3.2.1}      & \sol{A} & \sol{255.0.0.0}     & \sol{10.0.0.0}     & \sol{10.255.255.255} & \sol{16\,777\,214} \\ \hline}
  \def\apipa{169.254.12.7}
  \def\Cnet{192.168.60}
  \def\Crows{%
    PC-A & \ip{192.168.60.10} & \ip{255.255.255.0} & \ip{192.168.60.254} & — \\
    PC-B & \ip{192.168.60.20} & \ip{255.255.255.0} & \ip{192.168.60.1}   & — \\
    PC-C & \ip{192.168.61.30} & \ip{255.255.255.0} & \ip{192.168.60.254} & — \\
    PC-D & \ip{192.168.60.10} & \ip{255.255.255.0} & \ip{192.168.60.254} & — \\
    SRV-C & \ip{10.0.0.5}     & \ip{255.0.0.0}     & \ip{10.0.0.254}     & — \\
    R1 — G0/0 & \ip{192.168.60.254} & \ip{255.255.255.0} & — & On \\
    R1 — G0/1 & \ip{10.0.0.254}     & \ip{255.0.0.0}     & — & Off \\}
  \def\Csol{%
    \hl 1 & \sol{PC-B} & \sol{Passerelle .60.1 : pas l'adresse du routeur.} & \sol{Joint son réseau, mais aucun autre (serveur).} & \sol{Passerelle 192.168.60.254.} \\ \hline
    \hl 2 & \sol{PC-C} & \sol{IP hors du réseau 192.168.60.0/24.} & \sol{Ne joint ni les postes ni le routeur.} & \sol{Ex. 192.168.60.30.} \\ \hline
    \hl 3 & \sol{PC-D} & \sol{Même IP que PC-A (doublon).} & \sol{Conflit d'adresse : PC-A et PC-D perturbés.} & \sol{Adresse unique, ex. .60.40.} \\ \hline
    \hl 4 & \sol{R1 G0/1} & \sol{Interface désactivée (Off).} & \sol{Réseau du serveur injoignable pour tous.} & \sol{Port Status : On.} \\ \hline}
  \def\Cbareme{PC-A et SRV-C sont correctement configurés ; accepter PC-A cité à la place de PC-D pour le doublon.}
\else
  \def\entreprise{La clinique vétérinaire \textbf{VetoPlus}}
  \def\Rnom{R-VetoPlus}
  \def\resUn{Soins}\def\pUn{SOI}\def\netUn{192.168.70}
  \def\resDeux{Secrétariat}\def\pDeux{SEC}\def\netDeux{172.20}
  \def\srvNom{SRV-Dossiers}\def\netSrv{192.168.250}\def\srvH{20}
  \def\ssid{VetoPlus-WiFi}
  \def\AIrows{%
    \hl\ip{192.168.15.130} & \sol{C} & \sol{255.255.255.0} & \sol{192.168.15.0} & \sol{192.168.15.255} & \sol{254} \\ \hline
    \hl\ip{172.31.200.7}   & \sol{B} & \sol{255.255.0.0}   & \sol{172.31.0.0}   & \sol{172.31.255.255} & \sol{65\,534} \\ \hline
    \hl\ip{10.45.6.9}      & \sol{A} & \sol{255.0.0.0}     & \sol{10.0.0.0}     & \sol{10.255.255.255} & \sol{16\,777\,214} \\ \hline}
  \def\apipa{169.254.88.3}
  \def\Cnet{192.168.80}
  \def\Crows{%
    PC-A & \ip{192.168.8.11}   & \ip{255.255.255.0} & \ip{192.168.80.254} & — \\
    PC-B & \ip{192.168.80.12}  & \ip{255.255.255.0} & \ip{192.168.80.254} & — \\
    PC-C & \ip{192.168.80.13}  & \ip{255.255.255.0} & \ip{192.168.80.100} & — \\
    PC-D & \ip{192.168.80.12}  & \ip{255.255.255.0} & \ip{192.168.80.254} & — \\
    SRV-C & \ip{10.0.0.8}      & \ip{255.0.0.0}     & \ip{10.0.0.254}     & — \\
    R1 — G0/0 & \ip{192.168.80.254} & \ip{255.255.255.0} & — & Off \\
    R1 — G0/1 & \ip{10.0.0.254}     & \ip{255.0.0.0}     & — & On \\}
  \def\Csol{%
    \hl 1 & \sol{PC-A} & \sol{IP hors du réseau 192.168.80.0/24.} & \sol{Ne joint ni les postes ni le routeur.} & \sol{Ex. 192.168.80.11.} \\ \hline
    \hl 2 & \sol{PC-C} & \sol{Passerelle .80.100 : pas l'adresse du routeur.} & \sol{Joint son réseau, mais aucun autre (serveur).} & \sol{Passerelle 192.168.80.254.} \\ \hline
    \hl 3 & \sol{PC-D} & \sol{Même IP que PC-B (doublon).} & \sol{Conflit d'adresse : PC-B et PC-D perturbés.} & \sol{Adresse unique, ex. .80.14.} \\ \hline
    \hl 4 & \sol{R1 G0/0} & \sol{Interface côté LAN désactivée (Off).} & \sol{Aucun poste ne joint le routeur ni le serveur.} & \sol{Port Status : On.} \\ \hline}
  \def\Cbareme{PC-B et SRV-C sont correctement configurés ; accepter PC-B cité à la place de PC-D pour le doublon.}
\fi

\pagestyle{fancy}
\fancyhf{}
\renewcommand{\headrulewidth}{0pt}
\fancyfoot[L]{\footnotesize\itshape\color{grisfoot}Évaluation — Réseaux locaux — Sujet \variante\ifcorrige\ — CORRIGÉ\fi}
\fancyfoot[R]{\footnotesize\color{grisfoot}\thepage/3}

\setlist{itemsep=0pt,topsep=2pt,parsep=0pt}
\setlength{\parindent}{0pt}
\setlength{\parskip}{2pt}

\begin{document}

\begin{center}
  {\Large\bfseries Évaluation pratique — Réseaux locaux \quad
   \fcolorbox{black}{white}{\ SUJET \variante\ }\par}
  \vspace{1pt}
  {\bfseries Adressage IP, commutation et routage — BTS CIEL 1re année — durée : 2 heures\par}
  \ifcorrige\vspace{2pt}{\bfseries\color{corr}CORRIGÉ ET BARÈME — document enseignant}\par\fi
\end{center}
\vspace{-2pt}
{\renewcommand{\arraystretch}{1.35}
\begin{tabularx}{\linewidth}{|X|X|p{2.2cm}|p{2.6cm}|}
  \hline
  \textbf{Nom :} & \textbf{Prénom :} & \textbf{Poste n°} & \textbf{Note :} \hfill /20 \\ \hline
\end{tabularx}}

\textbf{Consignes :} travail \textbf{individuel}, documents et Internet \textbf{non autorisés}. Ce sujet est le
\textbf{document-réponse}. Fichier Packet Tracer : \ip{Nom\_Prenom\_EvalLAN\_\variante.pkt}, à enregistrer
régulièrement dans le dossier de rendu. Les « visas » sont apposés par le professeur, que vous appelez à chaque test réussi.
Durées conseillées : A 30 min — B 1 h 10 — C 20 min.

%==========================================================================
\section{Partie A — Questions de cours et calculs \hfill /6}

\question{A1}{2 pts}{Complétez le tableau en considérant le \textbf{masque par défaut} de la classe de chaque adresse.}
{\small\renewcommand{\arraystretch}{1.15}
\begin{tabularx}{\linewidth}{|l|c|C|C|C|c|}
  \hline
  \textbf{Adresse IP} & \textbf{Classe} & \textbf{Masque par défaut} & \textbf{Adresse réseau} &
  \textbf{Adresse de diffusion} & \textbf{Nb d'hôtes} \\ \hline
  \AIrows
\end{tabularx}}
\bareme{0,5 pt par ligne juste + 0,5 pt si les trois nombres d'hôtes sont justes ($2^{n}-2$, $n$ = nombre de bits d'hôte).}

\question{A2}{1 pt}{Indiquez le type de câble (\textbf{droit} ou \textbf{croisé}) pour chaque liaison.}
{\small\renewcommand{\arraystretch}{1.15}
\begin{tabularx}{\linewidth}{|X|c||X|c|}
  \hline
  \hl PC $\leftrightarrow$ commutateur & \makebox[1.6cm]{\sol{droit}} & Commutateur $\leftrightarrow$ routeur & \makebox[1.6cm]{\sol{droit}} \\ \hline
  \hl PC $\leftrightarrow$ PC (réseau P2P) & \sol{croisé} & Concentrateur $\leftrightarrow$ concentrateur & \sol{croisé} \\ \hline
\end{tabularx}}
\bareme{0,25 pt par case. Équipements de types différents : câble droit ; de même type : câble croisé.}

\question{A3}{1 pt}{Comment un \textbf{concentrateur (hub)} et un \textbf{commutateur (switch)} traitent-ils une trame
reçue ? Lequel limite les collisions, et pourquoi ?}
\zone{2.9cm}{Le \textbf{hub} (couche 1) répète la trame sur \textbf{tous ses autres ports} : un seul domaine de collision
partagé. Le \textbf{switch} (couche 2) lit l'adresse MAC de destination et, grâce à sa \textbf{table MAC}, ne l'envoie
\textbf{que sur le port du destinataire} : un domaine de collision par port, c'est lui qui limite les collisions.
\emph{0,5 pt hub, 0,5 pt switch.}}

\question{A4}{1 pt}{Qu'est-ce que la \textbf{passerelle par défaut} d'un poste ? Quelle condition son adresse doit-elle respecter ?}
\zone{2.4cm}{Adresse IP de l'interface du \textbf{routeur} à laquelle le poste envoie les paquets destinés à un
\textbf{autre réseau}. Elle doit \textbf{appartenir au même réseau} que le poste. \emph{0,5 + 0,5 pt.}}

\question{A5}{0,5 pt}{Un portable configuré en DHCP affiche \ip{\apipa} / \ip{255.255.0.0}. Que s'est-il passé ?}
\zone{1.5cm}{Aucun \textbf{serveur DHCP} n'a répondu : le poste s'est attribué une adresse \textbf{APIPA} (169.254.0.0/16).}

\question{A6}{0,5 pt}{Quel protocole utilise \texttt{ping} ? Nommez les deux messages échangés et la couche OSI des adresses IP.}
\zone{1.5cm}{\textbf{ICMP} ; \textbf{Echo Request} / \textbf{Echo Reply} ; adresses IP en \textbf{couche 3 (Réseau)}.}

%==========================================================================
\newpage
\section{Partie B — Réalisation sous Packet Tracer \hfill /10}

\begin{cdc}
\entreprise{} vous confie son réseau : \textbf{trois réseaux} reliés par un \textbf{routeur unique Cisco 2911}
nommé \ip{\Rnom} (interfaces G0/0, G0/1, G0/2).
\begin{itemize}[leftmargin=12pt]
  \item \textbf{Réseau \resUn} : \textbf{4 postes} PC-\pUn1 à PC-\pUn4, commutateur \ip{SW-\pUn}, adressage statique,
        réseau \ip{\netUn.0/24}.
  \item \textbf{Réseau \resDeux} : \textbf{2 postes} PC-\pDeux1, PC-\pDeux2 et \textbf{1 imprimante} IMP-\pDeux,
        commutateur \ip{SW-\pDeux}, adressage statique, réseau de classe B \ip{\netDeux.0.0/16}.
  \item \textbf{Réseau Serveur} : serveur \ip{\srvNom} d'adresse \ip{\netSrv.\srvH/24}, relié \textbf{directement} au routeur.
  \item La passerelle de chaque réseau est sa \textbf{dernière adresse utilisable}. Tous les équipements communiquent
        entre eux et accèdent à la page web du serveur. Nommage imposé (« Display Name »).
\end{itemize}
\end{cdc}

\question{B1}{2 pts}{Complétez le plan d'adressage \textbf{avant} la réalisation.}
{\footnotesize\setlength{\tabcolsep}{3pt}\renewcommand{\arraystretch}{1.15}
\begin{tabularx}{\linewidth}{|l|C|C|C|C|C|C|}
  \hline
  \textbf{Réseau} & \textbf{Adresse réseau} & \textbf{Masque} & \textbf{1re @ utilisable} &
  \textbf{Dernière @ utilisable} & \textbf{Diffusion} & \textbf{Passerelle} \\ \hline
  \hl \resUn & \sol{\netUn.0} & \sol{255.255.255.0} & \sol{\netUn.1} & \sol{\netUn.254} & \sol{\netUn.255} & \sol{\netUn.254} \\ \hline
  \hl \resDeux & \sol{\netDeux.0.0} & \sol{255.255.0.0} & \sol{\netDeux.0.1} & \sol{\netDeux.255.254} & \sol{\netDeux.255.255} & \sol{\netDeux.255.254} \\ \hline
  \hl Serveur & \sol{\netSrv.0} & \sol{255.255.255.0} & \sol{\netSrv.1} & \sol{\netSrv.254} & \sol{\netSrv.255} & \sol{\netSrv.254} \\ \hline
\end{tabularx}}
\bareme{0,5 pt par ligne juste + 0,5 pt si la convention de passerelle est respectée partout.}

\question{B2}{2 pts}{Construisez la topologie (équipements, câblage, nommage) et relevez l'adressage configuré.}
{\footnotesize\setlength{\tabcolsep}{3pt}\renewcommand{\arraystretch}{1.1}
\begin{tabularx}{\linewidth}{|l|C|C|C||l|C|C|}
  \hline
  \textbf{Hôte} & \textbf{Adresse IP} & \textbf{Masque} & \textbf{Passerelle} & \textbf{Interface} & \textbf{Adresse IP} & \textbf{Masque} \\ \hline
  \hl PC-\pUn1 & \sol{\netUn.1} & \sol{255.255.255.0} & \sol{\netUn.254} & \Rnom{} G0/0 & \sol{\netUn.254} & \sol{255.255.255.0} \\ \hline
  \hl PC-\pDeux1 & \sol{\netDeux.0.1} & \sol{255.255.0.0} & \sol{\netDeux.255.254} & \Rnom{} G0/1 & \sol{\netDeux.255.254} & \sol{255.255.0.0} \\ \hline
  \hl IMP-\pDeux & \sol{\netDeux.0.10} & \sol{255.255.0.0} & \sol{\netDeux.255.254} & \Rnom{} G0/2 & \sol{\netSrv.254} & \sol{255.255.255.0} \\ \hline
  \hl \srvNom & \ip{\netSrv.\srvH} & \sol{255.255.255.0} & \sol{\netSrv.254} & \multicolumn{3}{c|}{\cellcolor{black!8}} \\ \hline
\end{tabularx}}
\bareme{Adresses d'hôtes libres (uniques, dans le bon réseau) ; affectation G0/x libre. Sur le .pkt : équipements 0,5 ;
câblage 1 (câbles droits vers les switchs et switch–routeur, \textbf{croisé} serveur–routeur, liens verts) ; nommage 0,5.}

\question{B3}{3 pts}{Configurez tous les équipements : IP, masques, passerelles des hôtes ; interfaces du routeur
(« Port Status : On »).}
\bareme{Sur le .pkt : interfaces routeur 1 pt ; IP/masques des hôtes 1 pt ; passerelles (imprimante et serveur compris)
1 pt. $-0,25$ par erreur, sans descendre sous 0 par item.}

\question{B4}{3 pts}{Réalisez les tests, notez la commande et le résultat, puis faites viser chaque test réussi.}
{\small\renewcommand{\arraystretch}{1.1}
\begin{tabularx}{\linewidth}{|c|l|l|C|c|c|}
  \hline
  \textbf{N°} & \textbf{Source} & \textbf{Destination} & \textbf{Commande / action} & \textbf{Résultat} & \textbf{Visa prof.} \\ \hline
  \hl 1 & PC-\pUn1 & PC-\pUn4 & \sol{ping \netUn.4} & \sol{réussi} & \hspace{1.4cm} \\ \hline
  \hl 2 & PC-\pDeux1 & IMP-\pDeux & \sol{ping \netDeux.0.10} & \sol{réussi} & \\ \hline
  \hl 3 & PC-\pUn2 & PC-\pDeux2 & \sol{ping \netDeux.0.2} & \sol{réussi} & \\ \hline
  \hl 4 & PC-\pDeux2 & \srvNom & \sol{ping \netSrv.\srvH} & \sol{réussi} & \\ \hline
  \hl 5 & PC-\pUn3 & page web du serveur & \sol{navigateur : {\NoAutoSpacing http://\netSrv.\srvH}} & \sol{page affichée} & \\ \hline
\end{tabularx}}
\bareme{0,5 pt par test visé + 0,5 pt pour B5. Le premier ping inter-réseaux peut échouer (résolution ARP) ; les suivants doivent réussir.}

\question{B5}{incl. B4}{Test n°3 : par quels équipements passe le paquet ? Pourquoi PC-\pUn2 l'envoie-t-il au routeur ?}
\zone{2.3cm}{PC-\pUn2 $\to$ SW-\pUn{} $\to$ \Rnom{} $\to$ SW-\pDeux{} $\to$ PC-\pDeux2. Grâce à son masque, PC-\pUn2 constate que
la destination (\netDeux.0.0/16) n'est \textbf{pas dans son réseau} (\netUn.0/24) : il envoie le paquet à sa
\textbf{passerelle par défaut}, qui le route.}

%==========================================================================
\newpage
\section{Partie C — Diagnostic d'un réseau en panne \hfill /4}

Un technicien constate qu'\textbf{aucun poste ne parvient à joindre le serveur SRV-C} et que certains postes
communiquent mal entre eux. Les câbles sont du bon type.

\begin{center}
\begin{minipage}{0.37\linewidth}\centering
\begin{tikzpicture}[font=\footnotesize,
  eq/.style={draw,fill=qbg,rounded corners=2pt,minimum width=9mm,minimum height=5mm,inner sep=2pt},
  every path/.style={thick}]
  \node[eq] (SW) {SW-1};
  \node[eq,right=9mm of SW] (R) {R1};
  \node[eq,right=7mm of R] (SRV) {SRV-C};
  \node[eq,above left=5mm and -3mm of SW] (A) {PC-A};
  \node[eq,above right=5mm and -3mm of SW] (B) {PC-B};
  \node[eq,below left=5mm and -3mm of SW] (C) {PC-C};
  \node[eq,below right=5mm and -3mm of SW] (D) {PC-D};
  \foreach \n in {A,B,C,D} \draw (\n) -- (SW);
  \draw (SW) -- node[above,font=\tiny]{G0/0} (R);
  \draw (R) -- node[above,font=\tiny]{G0/1} (SRV);
\end{tikzpicture}
\end{minipage}\hfill
\begin{minipage}{0.61\linewidth}\centering
{\small\renewcommand{\arraystretch}{1.05}\setlength{\tabcolsep}{4pt}
\begin{tabular}{lllll}
  \toprule
  \textbf{Équipement} & \textbf{Adresse IP} & \textbf{Masque} & \textbf{Passerelle} & \textbf{Port} \\
  \midrule
  \Crows
  \bottomrule
\end{tabular}}
\end{minipage}
\end{center}

\question{C1}{4 pts}{Relevez les \textbf{quatre anomalies}. Pour chacune : conséquence et correction.}
{\small\renewcommand{\arraystretch}{1.1}
\ifcorrige\def\hc{}\else\def\hc{\rule[-10mm]{0pt}{14mm}}\fi
\begin{tabularx}{\linewidth}{|c|>{\raggedright\arraybackslash}p{2cm}|L|L|L|}
  \hline
  \textbf{N°} & \textbf{Équipement} & \textbf{Anomalie} & \textbf{Conséquence} & \textbf{Correction} \\ \hline
  \ifcorrige\Csol\else
  \hc 1 & & & & \\ \hline
  \hc 2 & & & & \\ \hline
  \hc 3 & & & & \\ \hline
  \hc 4 & & & & \\ \hline
  \fi
\end{tabularx}}
\bareme{1 pt par anomalie (0,25 équipement, anomalie, conséquence, correction ; ordre indifférent). \Cbareme}

%==========================================================================
\section{Bonus — Extension sans fil \hfill +1}

\begin{cdc}Ajoutez au réseau \resUn{} un \textbf{AccessPoint-PT} relié à SW-\pUn{} (SSID \ip{\ssid}, \textbf{WPA2-PSK})
et un portable \ip{Laptop-\pUn} équipé d'un module \texttt{WPC300N}, en adressage \textbf{statique}.
Vérifiez qu'il joint \srvNom.\end{cdc}
{\small\renewcommand{\arraystretch}{1.1}
\begin{tabularx}{\linewidth}{|l|C|C|C|c|}
  \hline
  \textbf{Équipement} & \textbf{Adresse IP} & \textbf{Masque} & \textbf{Passerelle} & \textbf{Visa prof.} \\ \hline
  \hl Laptop-\pUn & \sol{\netUn.50} & \sol{255.255.255.0} & \sol{\netUn.254} & \hspace{1.4cm} \\ \hline
\end{tabularx}}
\bareme{Le point d'accès fait office de pont : le portable est dans \netUn.0/24. Bonus si associé en WPA2-PSK et ping
vers \netSrv.\srvH{} réussi. Note plafonnée à 20.}

\end{document}
"
%   pdflatex -jobname=Eval_Reseau_SujetB   "\def\variante{B}\def\modecorrige{0}% Évaluation pratique — Réseaux locaux (BTS CIEL 1re année) — sujets A et B
% Compilation (4 PDF) :
%   pdflatex -jobname=Eval_Reseau_SujetA   "\def\variante{A}\def\modecorrige{0}\input{Eval_TP_Reseau_LAN.tex}"
%   pdflatex -jobname=Eval_Reseau_SujetB   "\def\variante{B}\def\modecorrige{0}\input{Eval_TP_Reseau_LAN.tex}"
%   pdflatex -jobname=Eval_Reseau_CorrigeA "\def\variante{A}\def\modecorrige{1}\input{Eval_TP_Reseau_LAN.tex}"
%   pdflatex -jobname=Eval_Reseau_CorrigeB "\def\variante{B}\def\modecorrige{1}\input{Eval_TP_Reseau_LAN.tex}"
\documentclass[10pt,a4paper]{article}

\usepackage[T1]{fontenc}
\usepackage[utf8]{inputenc}
\usepackage{lmodern}
\usepackage[french]{babel}
\usepackage[a4paper,hmargin=1.5cm,top=1.3cm,bottom=1.6cm,footskip=0.7cm]{geometry}
\usepackage[table]{xcolor}
\usepackage{booktabs,tabularx,array}
\usepackage{enumitem}
\usepackage{fancyhdr}
\usepackage[most]{tcolorbox}
\usepackage{tikz}
\usetikzlibrary{positioning}
\usepackage{titlesec}

\providecommand{\modecorrige}{0}
\providecommand{\variante}{A}
\newif\ifcorrige
\ifnum\modecorrige=1 \corrigetrue \else \corrigefalse \fi

\definecolor{ipred}{HTML}{C0392B}
\definecolor{cdcbar}{HTML}{A67C3D}
\definecolor{cdcbg}{HTML}{F4F1EA}
\definecolor{qbar}{HTML}{4A6B5D}
\definecolor{qbg}{HTML}{EDF1EF}
\definecolor{corr}{HTML}{1F5FA8}
\definecolor{corrbg}{HTML}{E8F0FA}
\definecolor{grisfoot}{HTML}{828282}

\newcommand{\ip}[1]{\textcolor{ipred}{\texttt{#1}}}
\newcommand{\pts}[1]{\hfill\mbox{\textbf{(#1)}}}
\newcommand{\sol}[1]{\ifcorrige\textcolor{corr}{#1}\fi}
\newcommand{\bareme}[1]{\ifcorrige\par{\footnotesize\color{corr}#1\par}\fi}
\newcommand{\hl}{\ifcorrige\rule[-1.6mm]{0pt}{5.4mm}\else\rule[-2.2mm]{0pt}{6.6mm}\fi}
\newcolumntype{C}{>{\centering\arraybackslash}X}
\newcolumntype{L}{>{\raggedright\arraybackslash}X}

\titleformat{\section}{\large\bfseries}{}{0pt}{}
\titlespacing*{\section}{0pt}{6pt}{3pt}

\tcbset{qstyle/.style={enhanced,boxrule=0pt,leftrule=2.5pt,arc=0pt,outer arc=0pt,
  left=5pt,right=5pt,top=2pt,bottom=2pt,before skip=4pt,after skip=3pt}}
\newtcolorbox{cdc}{qstyle,colback=cdcbg,colframe=cdcbar}
\newtcolorbox{quest}{qstyle,colback=qbg,colframe=qbar}

% Zone de réponse : vierge dans le sujet, remplie dans le corrigé
\newcommand{\zone}[2]{%
  \ifcorrige
    \begin{tcolorbox}[colback=corrbg,colframe=corr,boxrule=0.5pt,arc=1pt,left=4pt,right=4pt,top=1pt,bottom=1pt,
      before skip=2pt,after skip=2pt,fontupper=\footnotesize\color{corr}]#2\end{tcolorbox}%
  \else
    \begin{tcolorbox}[colback=white,colframe=black!45,boxrule=0.4pt,arc=1pt,height=#1,
      before skip=2pt,after skip=2pt]\end{tcolorbox}%
  \fi}

\newcommand{\question}[3]{\begin{quest}\textbf{#1.} #3\pts{#2}\end{quest}}

%==========================================================================
% Données propres à chaque variante
%==========================================================================
\if\variante A
  \def\entreprise{Le cabinet d'architectes \textbf{ArchiNova}}
  \def\Rnom{R-ArchiNova}
  \def\resUn{Bureau d'études}\def\pUn{BE}\def\netUn{192.168.40}
  \def\resDeux{Accueil}\def\pDeux{ACC}\def\netDeux{172.16}
  \def\srvNom{SRV-Intranet}\def\netSrv{192.168.200}\def\srvH{10}
  \def\ssid{ArchiNova-WiFi}
  \def\AIrows{%
    \hl\ip{192.168.40.25} & \sol{C} & \sol{255.255.255.0} & \sol{192.168.40.0} & \sol{192.168.40.255} & \sol{254} \\ \hline
    \hl\ip{172.16.5.10}   & \sol{B} & \sol{255.255.0.0}   & \sol{172.16.0.0}   & \sol{172.16.255.255} & \sol{65\,534} \\ \hline
    \hl\ip{10.3.2.1}      & \sol{A} & \sol{255.0.0.0}     & \sol{10.0.0.0}     & \sol{10.255.255.255} & \sol{16\,777\,214} \\ \hline}
  \def\apipa{169.254.12.7}
  \def\Cnet{192.168.60}
  \def\Crows{%
    PC-A & \ip{192.168.60.10} & \ip{255.255.255.0} & \ip{192.168.60.254} & — \\
    PC-B & \ip{192.168.60.20} & \ip{255.255.255.0} & \ip{192.168.60.1}   & — \\
    PC-C & \ip{192.168.61.30} & \ip{255.255.255.0} & \ip{192.168.60.254} & — \\
    PC-D & \ip{192.168.60.10} & \ip{255.255.255.0} & \ip{192.168.60.254} & — \\
    SRV-C & \ip{10.0.0.5}     & \ip{255.0.0.0}     & \ip{10.0.0.254}     & — \\
    R1 — G0/0 & \ip{192.168.60.254} & \ip{255.255.255.0} & — & On \\
    R1 — G0/1 & \ip{10.0.0.254}     & \ip{255.0.0.0}     & — & Off \\}
  \def\Csol{%
    \hl 1 & \sol{PC-B} & \sol{Passerelle .60.1 : pas l'adresse du routeur.} & \sol{Joint son réseau, mais aucun autre (serveur).} & \sol{Passerelle 192.168.60.254.} \\ \hline
    \hl 2 & \sol{PC-C} & \sol{IP hors du réseau 192.168.60.0/24.} & \sol{Ne joint ni les postes ni le routeur.} & \sol{Ex. 192.168.60.30.} \\ \hline
    \hl 3 & \sol{PC-D} & \sol{Même IP que PC-A (doublon).} & \sol{Conflit d'adresse : PC-A et PC-D perturbés.} & \sol{Adresse unique, ex. .60.40.} \\ \hline
    \hl 4 & \sol{R1 G0/1} & \sol{Interface désactivée (Off).} & \sol{Réseau du serveur injoignable pour tous.} & \sol{Port Status : On.} \\ \hline}
  \def\Cbareme{PC-A et SRV-C sont correctement configurés ; accepter PC-A cité à la place de PC-D pour le doublon.}
\else
  \def\entreprise{La clinique vétérinaire \textbf{VetoPlus}}
  \def\Rnom{R-VetoPlus}
  \def\resUn{Soins}\def\pUn{SOI}\def\netUn{192.168.70}
  \def\resDeux{Secrétariat}\def\pDeux{SEC}\def\netDeux{172.20}
  \def\srvNom{SRV-Dossiers}\def\netSrv{192.168.250}\def\srvH{20}
  \def\ssid{VetoPlus-WiFi}
  \def\AIrows{%
    \hl\ip{192.168.15.130} & \sol{C} & \sol{255.255.255.0} & \sol{192.168.15.0} & \sol{192.168.15.255} & \sol{254} \\ \hline
    \hl\ip{172.31.200.7}   & \sol{B} & \sol{255.255.0.0}   & \sol{172.31.0.0}   & \sol{172.31.255.255} & \sol{65\,534} \\ \hline
    \hl\ip{10.45.6.9}      & \sol{A} & \sol{255.0.0.0}     & \sol{10.0.0.0}     & \sol{10.255.255.255} & \sol{16\,777\,214} \\ \hline}
  \def\apipa{169.254.88.3}
  \def\Cnet{192.168.80}
  \def\Crows{%
    PC-A & \ip{192.168.8.11}   & \ip{255.255.255.0} & \ip{192.168.80.254} & — \\
    PC-B & \ip{192.168.80.12}  & \ip{255.255.255.0} & \ip{192.168.80.254} & — \\
    PC-C & \ip{192.168.80.13}  & \ip{255.255.255.0} & \ip{192.168.80.100} & — \\
    PC-D & \ip{192.168.80.12}  & \ip{255.255.255.0} & \ip{192.168.80.254} & — \\
    SRV-C & \ip{10.0.0.8}      & \ip{255.0.0.0}     & \ip{10.0.0.254}     & — \\
    R1 — G0/0 & \ip{192.168.80.254} & \ip{255.255.255.0} & — & Off \\
    R1 — G0/1 & \ip{10.0.0.254}     & \ip{255.0.0.0}     & — & On \\}
  \def\Csol{%
    \hl 1 & \sol{PC-A} & \sol{IP hors du réseau 192.168.80.0/24.} & \sol{Ne joint ni les postes ni le routeur.} & \sol{Ex. 192.168.80.11.} \\ \hline
    \hl 2 & \sol{PC-C} & \sol{Passerelle .80.100 : pas l'adresse du routeur.} & \sol{Joint son réseau, mais aucun autre (serveur).} & \sol{Passerelle 192.168.80.254.} \\ \hline
    \hl 3 & \sol{PC-D} & \sol{Même IP que PC-B (doublon).} & \sol{Conflit d'adresse : PC-B et PC-D perturbés.} & \sol{Adresse unique, ex. .80.14.} \\ \hline
    \hl 4 & \sol{R1 G0/0} & \sol{Interface côté LAN désactivée (Off).} & \sol{Aucun poste ne joint le routeur ni le serveur.} & \sol{Port Status : On.} \\ \hline}
  \def\Cbareme{PC-B et SRV-C sont correctement configurés ; accepter PC-B cité à la place de PC-D pour le doublon.}
\fi

\pagestyle{fancy}
\fancyhf{}
\renewcommand{\headrulewidth}{0pt}
\fancyfoot[L]{\footnotesize\itshape\color{grisfoot}Évaluation — Réseaux locaux — Sujet \variante\ifcorrige\ — CORRIGÉ\fi}
\fancyfoot[R]{\footnotesize\color{grisfoot}\thepage/3}

\setlist{itemsep=0pt,topsep=2pt,parsep=0pt}
\setlength{\parindent}{0pt}
\setlength{\parskip}{2pt}

\begin{document}

\begin{center}
  {\Large\bfseries Évaluation pratique — Réseaux locaux \quad
   \fcolorbox{black}{white}{\ SUJET \variante\ }\par}
  \vspace{1pt}
  {\bfseries Adressage IP, commutation et routage — BTS CIEL 1re année — durée : 2 heures\par}
  \ifcorrige\vspace{2pt}{\bfseries\color{corr}CORRIGÉ ET BARÈME — document enseignant}\par\fi
\end{center}
\vspace{-2pt}
{\renewcommand{\arraystretch}{1.35}
\begin{tabularx}{\linewidth}{|X|X|p{2.2cm}|p{2.6cm}|}
  \hline
  \textbf{Nom :} & \textbf{Prénom :} & \textbf{Poste n°} & \textbf{Note :} \hfill /20 \\ \hline
\end{tabularx}}

\textbf{Consignes :} travail \textbf{individuel}, documents et Internet \textbf{non autorisés}. Ce sujet est le
\textbf{document-réponse}. Fichier Packet Tracer : \ip{Nom\_Prenom\_EvalLAN\_\variante.pkt}, à enregistrer
régulièrement dans le dossier de rendu. Les « visas » sont apposés par le professeur, que vous appelez à chaque test réussi.
Durées conseillées : A 30 min — B 1 h 10 — C 20 min.

%==========================================================================
\section{Partie A — Questions de cours et calculs \hfill /6}

\question{A1}{2 pts}{Complétez le tableau en considérant le \textbf{masque par défaut} de la classe de chaque adresse.}
{\small\renewcommand{\arraystretch}{1.15}
\begin{tabularx}{\linewidth}{|l|c|C|C|C|c|}
  \hline
  \textbf{Adresse IP} & \textbf{Classe} & \textbf{Masque par défaut} & \textbf{Adresse réseau} &
  \textbf{Adresse de diffusion} & \textbf{Nb d'hôtes} \\ \hline
  \AIrows
\end{tabularx}}
\bareme{0,5 pt par ligne juste + 0,5 pt si les trois nombres d'hôtes sont justes ($2^{n}-2$, $n$ = nombre de bits d'hôte).}

\question{A2}{1 pt}{Indiquez le type de câble (\textbf{droit} ou \textbf{croisé}) pour chaque liaison.}
{\small\renewcommand{\arraystretch}{1.15}
\begin{tabularx}{\linewidth}{|X|c||X|c|}
  \hline
  \hl PC $\leftrightarrow$ commutateur & \makebox[1.6cm]{\sol{droit}} & Commutateur $\leftrightarrow$ routeur & \makebox[1.6cm]{\sol{droit}} \\ \hline
  \hl PC $\leftrightarrow$ PC (réseau P2P) & \sol{croisé} & Concentrateur $\leftrightarrow$ concentrateur & \sol{croisé} \\ \hline
\end{tabularx}}
\bareme{0,25 pt par case. Équipements de types différents : câble droit ; de même type : câble croisé.}

\question{A3}{1 pt}{Comment un \textbf{concentrateur (hub)} et un \textbf{commutateur (switch)} traitent-ils une trame
reçue ? Lequel limite les collisions, et pourquoi ?}
\zone{2.9cm}{Le \textbf{hub} (couche 1) répète la trame sur \textbf{tous ses autres ports} : un seul domaine de collision
partagé. Le \textbf{switch} (couche 2) lit l'adresse MAC de destination et, grâce à sa \textbf{table MAC}, ne l'envoie
\textbf{que sur le port du destinataire} : un domaine de collision par port, c'est lui qui limite les collisions.
\emph{0,5 pt hub, 0,5 pt switch.}}

\question{A4}{1 pt}{Qu'est-ce que la \textbf{passerelle par défaut} d'un poste ? Quelle condition son adresse doit-elle respecter ?}
\zone{2.4cm}{Adresse IP de l'interface du \textbf{routeur} à laquelle le poste envoie les paquets destinés à un
\textbf{autre réseau}. Elle doit \textbf{appartenir au même réseau} que le poste. \emph{0,5 + 0,5 pt.}}

\question{A5}{0,5 pt}{Un portable configuré en DHCP affiche \ip{\apipa} / \ip{255.255.0.0}. Que s'est-il passé ?}
\zone{1.5cm}{Aucun \textbf{serveur DHCP} n'a répondu : le poste s'est attribué une adresse \textbf{APIPA} (169.254.0.0/16).}

\question{A6}{0,5 pt}{Quel protocole utilise \texttt{ping} ? Nommez les deux messages échangés et la couche OSI des adresses IP.}
\zone{1.5cm}{\textbf{ICMP} ; \textbf{Echo Request} / \textbf{Echo Reply} ; adresses IP en \textbf{couche 3 (Réseau)}.}

%==========================================================================
\newpage
\section{Partie B — Réalisation sous Packet Tracer \hfill /10}

\begin{cdc}
\entreprise{} vous confie son réseau : \textbf{trois réseaux} reliés par un \textbf{routeur unique Cisco 2911}
nommé \ip{\Rnom} (interfaces G0/0, G0/1, G0/2).
\begin{itemize}[leftmargin=12pt]
  \item \textbf{Réseau \resUn} : \textbf{4 postes} PC-\pUn1 à PC-\pUn4, commutateur \ip{SW-\pUn}, adressage statique,
        réseau \ip{\netUn.0/24}.
  \item \textbf{Réseau \resDeux} : \textbf{2 postes} PC-\pDeux1, PC-\pDeux2 et \textbf{1 imprimante} IMP-\pDeux,
        commutateur \ip{SW-\pDeux}, adressage statique, réseau de classe B \ip{\netDeux.0.0/16}.
  \item \textbf{Réseau Serveur} : serveur \ip{\srvNom} d'adresse \ip{\netSrv.\srvH/24}, relié \textbf{directement} au routeur.
  \item La passerelle de chaque réseau est sa \textbf{dernière adresse utilisable}. Tous les équipements communiquent
        entre eux et accèdent à la page web du serveur. Nommage imposé (« Display Name »).
\end{itemize}
\end{cdc}

\question{B1}{2 pts}{Complétez le plan d'adressage \textbf{avant} la réalisation.}
{\footnotesize\setlength{\tabcolsep}{3pt}\renewcommand{\arraystretch}{1.15}
\begin{tabularx}{\linewidth}{|l|C|C|C|C|C|C|}
  \hline
  \textbf{Réseau} & \textbf{Adresse réseau} & \textbf{Masque} & \textbf{1re @ utilisable} &
  \textbf{Dernière @ utilisable} & \textbf{Diffusion} & \textbf{Passerelle} \\ \hline
  \hl \resUn & \sol{\netUn.0} & \sol{255.255.255.0} & \sol{\netUn.1} & \sol{\netUn.254} & \sol{\netUn.255} & \sol{\netUn.254} \\ \hline
  \hl \resDeux & \sol{\netDeux.0.0} & \sol{255.255.0.0} & \sol{\netDeux.0.1} & \sol{\netDeux.255.254} & \sol{\netDeux.255.255} & \sol{\netDeux.255.254} \\ \hline
  \hl Serveur & \sol{\netSrv.0} & \sol{255.255.255.0} & \sol{\netSrv.1} & \sol{\netSrv.254} & \sol{\netSrv.255} & \sol{\netSrv.254} \\ \hline
\end{tabularx}}
\bareme{0,5 pt par ligne juste + 0,5 pt si la convention de passerelle est respectée partout.}

\question{B2}{2 pts}{Construisez la topologie (équipements, câblage, nommage) et relevez l'adressage configuré.}
{\footnotesize\setlength{\tabcolsep}{3pt}\renewcommand{\arraystretch}{1.1}
\begin{tabularx}{\linewidth}{|l|C|C|C||l|C|C|}
  \hline
  \textbf{Hôte} & \textbf{Adresse IP} & \textbf{Masque} & \textbf{Passerelle} & \textbf{Interface} & \textbf{Adresse IP} & \textbf{Masque} \\ \hline
  \hl PC-\pUn1 & \sol{\netUn.1} & \sol{255.255.255.0} & \sol{\netUn.254} & \Rnom{} G0/0 & \sol{\netUn.254} & \sol{255.255.255.0} \\ \hline
  \hl PC-\pDeux1 & \sol{\netDeux.0.1} & \sol{255.255.0.0} & \sol{\netDeux.255.254} & \Rnom{} G0/1 & \sol{\netDeux.255.254} & \sol{255.255.0.0} \\ \hline
  \hl IMP-\pDeux & \sol{\netDeux.0.10} & \sol{255.255.0.0} & \sol{\netDeux.255.254} & \Rnom{} G0/2 & \sol{\netSrv.254} & \sol{255.255.255.0} \\ \hline
  \hl \srvNom & \ip{\netSrv.\srvH} & \sol{255.255.255.0} & \sol{\netSrv.254} & \multicolumn{3}{c|}{\cellcolor{black!8}} \\ \hline
\end{tabularx}}
\bareme{Adresses d'hôtes libres (uniques, dans le bon réseau) ; affectation G0/x libre. Sur le .pkt : équipements 0,5 ;
câblage 1 (câbles droits vers les switchs et switch–routeur, \textbf{croisé} serveur–routeur, liens verts) ; nommage 0,5.}

\question{B3}{3 pts}{Configurez tous les équipements : IP, masques, passerelles des hôtes ; interfaces du routeur
(« Port Status : On »).}
\bareme{Sur le .pkt : interfaces routeur 1 pt ; IP/masques des hôtes 1 pt ; passerelles (imprimante et serveur compris)
1 pt. $-0,25$ par erreur, sans descendre sous 0 par item.}

\question{B4}{3 pts}{Réalisez les tests, notez la commande et le résultat, puis faites viser chaque test réussi.}
{\small\renewcommand{\arraystretch}{1.1}
\begin{tabularx}{\linewidth}{|c|l|l|C|c|c|}
  \hline
  \textbf{N°} & \textbf{Source} & \textbf{Destination} & \textbf{Commande / action} & \textbf{Résultat} & \textbf{Visa prof.} \\ \hline
  \hl 1 & PC-\pUn1 & PC-\pUn4 & \sol{ping \netUn.4} & \sol{réussi} & \hspace{1.4cm} \\ \hline
  \hl 2 & PC-\pDeux1 & IMP-\pDeux & \sol{ping \netDeux.0.10} & \sol{réussi} & \\ \hline
  \hl 3 & PC-\pUn2 & PC-\pDeux2 & \sol{ping \netDeux.0.2} & \sol{réussi} & \\ \hline
  \hl 4 & PC-\pDeux2 & \srvNom & \sol{ping \netSrv.\srvH} & \sol{réussi} & \\ \hline
  \hl 5 & PC-\pUn3 & page web du serveur & \sol{navigateur : {\NoAutoSpacing http://\netSrv.\srvH}} & \sol{page affichée} & \\ \hline
\end{tabularx}}
\bareme{0,5 pt par test visé + 0,5 pt pour B5. Le premier ping inter-réseaux peut échouer (résolution ARP) ; les suivants doivent réussir.}

\question{B5}{incl. B4}{Test n°3 : par quels équipements passe le paquet ? Pourquoi PC-\pUn2 l'envoie-t-il au routeur ?}
\zone{2.3cm}{PC-\pUn2 $\to$ SW-\pUn{} $\to$ \Rnom{} $\to$ SW-\pDeux{} $\to$ PC-\pDeux2. Grâce à son masque, PC-\pUn2 constate que
la destination (\netDeux.0.0/16) n'est \textbf{pas dans son réseau} (\netUn.0/24) : il envoie le paquet à sa
\textbf{passerelle par défaut}, qui le route.}

%==========================================================================
\newpage
\section{Partie C — Diagnostic d'un réseau en panne \hfill /4}

Un technicien constate qu'\textbf{aucun poste ne parvient à joindre le serveur SRV-C} et que certains postes
communiquent mal entre eux. Les câbles sont du bon type.

\begin{center}
\begin{minipage}{0.37\linewidth}\centering
\begin{tikzpicture}[font=\footnotesize,
  eq/.style={draw,fill=qbg,rounded corners=2pt,minimum width=9mm,minimum height=5mm,inner sep=2pt},
  every path/.style={thick}]
  \node[eq] (SW) {SW-1};
  \node[eq,right=9mm of SW] (R) {R1};
  \node[eq,right=7mm of R] (SRV) {SRV-C};
  \node[eq,above left=5mm and -3mm of SW] (A) {PC-A};
  \node[eq,above right=5mm and -3mm of SW] (B) {PC-B};
  \node[eq,below left=5mm and -3mm of SW] (C) {PC-C};
  \node[eq,below right=5mm and -3mm of SW] (D) {PC-D};
  \foreach \n in {A,B,C,D} \draw (\n) -- (SW);
  \draw (SW) -- node[above,font=\tiny]{G0/0} (R);
  \draw (R) -- node[above,font=\tiny]{G0/1} (SRV);
\end{tikzpicture}
\end{minipage}\hfill
\begin{minipage}{0.61\linewidth}\centering
{\small\renewcommand{\arraystretch}{1.05}\setlength{\tabcolsep}{4pt}
\begin{tabular}{lllll}
  \toprule
  \textbf{Équipement} & \textbf{Adresse IP} & \textbf{Masque} & \textbf{Passerelle} & \textbf{Port} \\
  \midrule
  \Crows
  \bottomrule
\end{tabular}}
\end{minipage}
\end{center}

\question{C1}{4 pts}{Relevez les \textbf{quatre anomalies}. Pour chacune : conséquence et correction.}
{\small\renewcommand{\arraystretch}{1.1}
\ifcorrige\def\hc{}\else\def\hc{\rule[-10mm]{0pt}{14mm}}\fi
\begin{tabularx}{\linewidth}{|c|>{\raggedright\arraybackslash}p{2cm}|L|L|L|}
  \hline
  \textbf{N°} & \textbf{Équipement} & \textbf{Anomalie} & \textbf{Conséquence} & \textbf{Correction} \\ \hline
  \ifcorrige\Csol\else
  \hc 1 & & & & \\ \hline
  \hc 2 & & & & \\ \hline
  \hc 3 & & & & \\ \hline
  \hc 4 & & & & \\ \hline
  \fi
\end{tabularx}}
\bareme{1 pt par anomalie (0,25 équipement, anomalie, conséquence, correction ; ordre indifférent). \Cbareme}

%==========================================================================
\section{Bonus — Extension sans fil \hfill +1}

\begin{cdc}Ajoutez au réseau \resUn{} un \textbf{AccessPoint-PT} relié à SW-\pUn{} (SSID \ip{\ssid}, \textbf{WPA2-PSK})
et un portable \ip{Laptop-\pUn} équipé d'un module \texttt{WPC300N}, en adressage \textbf{statique}.
Vérifiez qu'il joint \srvNom.\end{cdc}
{\small\renewcommand{\arraystretch}{1.1}
\begin{tabularx}{\linewidth}{|l|C|C|C|c|}
  \hline
  \textbf{Équipement} & \textbf{Adresse IP} & \textbf{Masque} & \textbf{Passerelle} & \textbf{Visa prof.} \\ \hline
  \hl Laptop-\pUn & \sol{\netUn.50} & \sol{255.255.255.0} & \sol{\netUn.254} & \hspace{1.4cm} \\ \hline
\end{tabularx}}
\bareme{Le point d'accès fait office de pont : le portable est dans \netUn.0/24. Bonus si associé en WPA2-PSK et ping
vers \netSrv.\srvH{} réussi. Note plafonnée à 20.}

\end{document}
"
%   pdflatex -jobname=Eval_Reseau_CorrigeA "\def\variante{A}\def\modecorrige{1}% Évaluation pratique — Réseaux locaux (BTS CIEL 1re année) — sujets A et B
% Compilation (4 PDF) :
%   pdflatex -jobname=Eval_Reseau_SujetA   "\def\variante{A}\def\modecorrige{0}\input{Eval_TP_Reseau_LAN.tex}"
%   pdflatex -jobname=Eval_Reseau_SujetB   "\def\variante{B}\def\modecorrige{0}\input{Eval_TP_Reseau_LAN.tex}"
%   pdflatex -jobname=Eval_Reseau_CorrigeA "\def\variante{A}\def\modecorrige{1}\input{Eval_TP_Reseau_LAN.tex}"
%   pdflatex -jobname=Eval_Reseau_CorrigeB "\def\variante{B}\def\modecorrige{1}\input{Eval_TP_Reseau_LAN.tex}"
\documentclass[10pt,a4paper]{article}

\usepackage[T1]{fontenc}
\usepackage[utf8]{inputenc}
\usepackage{lmodern}
\usepackage[french]{babel}
\usepackage[a4paper,hmargin=1.5cm,top=1.3cm,bottom=1.6cm,footskip=0.7cm]{geometry}
\usepackage[table]{xcolor}
\usepackage{booktabs,tabularx,array}
\usepackage{enumitem}
\usepackage{fancyhdr}
\usepackage[most]{tcolorbox}
\usepackage{tikz}
\usetikzlibrary{positioning}
\usepackage{titlesec}

\providecommand{\modecorrige}{0}
\providecommand{\variante}{A}
\newif\ifcorrige
\ifnum\modecorrige=1 \corrigetrue \else \corrigefalse \fi

\definecolor{ipred}{HTML}{C0392B}
\definecolor{cdcbar}{HTML}{A67C3D}
\definecolor{cdcbg}{HTML}{F4F1EA}
\definecolor{qbar}{HTML}{4A6B5D}
\definecolor{qbg}{HTML}{EDF1EF}
\definecolor{corr}{HTML}{1F5FA8}
\definecolor{corrbg}{HTML}{E8F0FA}
\definecolor{grisfoot}{HTML}{828282}

\newcommand{\ip}[1]{\textcolor{ipred}{\texttt{#1}}}
\newcommand{\pts}[1]{\hfill\mbox{\textbf{(#1)}}}
\newcommand{\sol}[1]{\ifcorrige\textcolor{corr}{#1}\fi}
\newcommand{\bareme}[1]{\ifcorrige\par{\footnotesize\color{corr}#1\par}\fi}
\newcommand{\hl}{\ifcorrige\rule[-1.6mm]{0pt}{5.4mm}\else\rule[-2.2mm]{0pt}{6.6mm}\fi}
\newcolumntype{C}{>{\centering\arraybackslash}X}
\newcolumntype{L}{>{\raggedright\arraybackslash}X}

\titleformat{\section}{\large\bfseries}{}{0pt}{}
\titlespacing*{\section}{0pt}{6pt}{3pt}

\tcbset{qstyle/.style={enhanced,boxrule=0pt,leftrule=2.5pt,arc=0pt,outer arc=0pt,
  left=5pt,right=5pt,top=2pt,bottom=2pt,before skip=4pt,after skip=3pt}}
\newtcolorbox{cdc}{qstyle,colback=cdcbg,colframe=cdcbar}
\newtcolorbox{quest}{qstyle,colback=qbg,colframe=qbar}

% Zone de réponse : vierge dans le sujet, remplie dans le corrigé
\newcommand{\zone}[2]{%
  \ifcorrige
    \begin{tcolorbox}[colback=corrbg,colframe=corr,boxrule=0.5pt,arc=1pt,left=4pt,right=4pt,top=1pt,bottom=1pt,
      before skip=2pt,after skip=2pt,fontupper=\footnotesize\color{corr}]#2\end{tcolorbox}%
  \else
    \begin{tcolorbox}[colback=white,colframe=black!45,boxrule=0.4pt,arc=1pt,height=#1,
      before skip=2pt,after skip=2pt]\end{tcolorbox}%
  \fi}

\newcommand{\question}[3]{\begin{quest}\textbf{#1.} #3\pts{#2}\end{quest}}

%==========================================================================
% Données propres à chaque variante
%==========================================================================
\if\variante A
  \def\entreprise{Le cabinet d'architectes \textbf{ArchiNova}}
  \def\Rnom{R-ArchiNova}
  \def\resUn{Bureau d'études}\def\pUn{BE}\def\netUn{192.168.40}
  \def\resDeux{Accueil}\def\pDeux{ACC}\def\netDeux{172.16}
  \def\srvNom{SRV-Intranet}\def\netSrv{192.168.200}\def\srvH{10}
  \def\ssid{ArchiNova-WiFi}
  \def\AIrows{%
    \hl\ip{192.168.40.25} & \sol{C} & \sol{255.255.255.0} & \sol{192.168.40.0} & \sol{192.168.40.255} & \sol{254} \\ \hline
    \hl\ip{172.16.5.10}   & \sol{B} & \sol{255.255.0.0}   & \sol{172.16.0.0}   & \sol{172.16.255.255} & \sol{65\,534} \\ \hline
    \hl\ip{10.3.2.1}      & \sol{A} & \sol{255.0.0.0}     & \sol{10.0.0.0}     & \sol{10.255.255.255} & \sol{16\,777\,214} \\ \hline}
  \def\apipa{169.254.12.7}
  \def\Cnet{192.168.60}
  \def\Crows{%
    PC-A & \ip{192.168.60.10} & \ip{255.255.255.0} & \ip{192.168.60.254} & — \\
    PC-B & \ip{192.168.60.20} & \ip{255.255.255.0} & \ip{192.168.60.1}   & — \\
    PC-C & \ip{192.168.61.30} & \ip{255.255.255.0} & \ip{192.168.60.254} & — \\
    PC-D & \ip{192.168.60.10} & \ip{255.255.255.0} & \ip{192.168.60.254} & — \\
    SRV-C & \ip{10.0.0.5}     & \ip{255.0.0.0}     & \ip{10.0.0.254}     & — \\
    R1 — G0/0 & \ip{192.168.60.254} & \ip{255.255.255.0} & — & On \\
    R1 — G0/1 & \ip{10.0.0.254}     & \ip{255.0.0.0}     & — & Off \\}
  \def\Csol{%
    \hl 1 & \sol{PC-B} & \sol{Passerelle .60.1 : pas l'adresse du routeur.} & \sol{Joint son réseau, mais aucun autre (serveur).} & \sol{Passerelle 192.168.60.254.} \\ \hline
    \hl 2 & \sol{PC-C} & \sol{IP hors du réseau 192.168.60.0/24.} & \sol{Ne joint ni les postes ni le routeur.} & \sol{Ex. 192.168.60.30.} \\ \hline
    \hl 3 & \sol{PC-D} & \sol{Même IP que PC-A (doublon).} & \sol{Conflit d'adresse : PC-A et PC-D perturbés.} & \sol{Adresse unique, ex. .60.40.} \\ \hline
    \hl 4 & \sol{R1 G0/1} & \sol{Interface désactivée (Off).} & \sol{Réseau du serveur injoignable pour tous.} & \sol{Port Status : On.} \\ \hline}
  \def\Cbareme{PC-A et SRV-C sont correctement configurés ; accepter PC-A cité à la place de PC-D pour le doublon.}
\else
  \def\entreprise{La clinique vétérinaire \textbf{VetoPlus}}
  \def\Rnom{R-VetoPlus}
  \def\resUn{Soins}\def\pUn{SOI}\def\netUn{192.168.70}
  \def\resDeux{Secrétariat}\def\pDeux{SEC}\def\netDeux{172.20}
  \def\srvNom{SRV-Dossiers}\def\netSrv{192.168.250}\def\srvH{20}
  \def\ssid{VetoPlus-WiFi}
  \def\AIrows{%
    \hl\ip{192.168.15.130} & \sol{C} & \sol{255.255.255.0} & \sol{192.168.15.0} & \sol{192.168.15.255} & \sol{254} \\ \hline
    \hl\ip{172.31.200.7}   & \sol{B} & \sol{255.255.0.0}   & \sol{172.31.0.0}   & \sol{172.31.255.255} & \sol{65\,534} \\ \hline
    \hl\ip{10.45.6.9}      & \sol{A} & \sol{255.0.0.0}     & \sol{10.0.0.0}     & \sol{10.255.255.255} & \sol{16\,777\,214} \\ \hline}
  \def\apipa{169.254.88.3}
  \def\Cnet{192.168.80}
  \def\Crows{%
    PC-A & \ip{192.168.8.11}   & \ip{255.255.255.0} & \ip{192.168.80.254} & — \\
    PC-B & \ip{192.168.80.12}  & \ip{255.255.255.0} & \ip{192.168.80.254} & — \\
    PC-C & \ip{192.168.80.13}  & \ip{255.255.255.0} & \ip{192.168.80.100} & — \\
    PC-D & \ip{192.168.80.12}  & \ip{255.255.255.0} & \ip{192.168.80.254} & — \\
    SRV-C & \ip{10.0.0.8}      & \ip{255.0.0.0}     & \ip{10.0.0.254}     & — \\
    R1 — G0/0 & \ip{192.168.80.254} & \ip{255.255.255.0} & — & Off \\
    R1 — G0/1 & \ip{10.0.0.254}     & \ip{255.0.0.0}     & — & On \\}
  \def\Csol{%
    \hl 1 & \sol{PC-A} & \sol{IP hors du réseau 192.168.80.0/24.} & \sol{Ne joint ni les postes ni le routeur.} & \sol{Ex. 192.168.80.11.} \\ \hline
    \hl 2 & \sol{PC-C} & \sol{Passerelle .80.100 : pas l'adresse du routeur.} & \sol{Joint son réseau, mais aucun autre (serveur).} & \sol{Passerelle 192.168.80.254.} \\ \hline
    \hl 3 & \sol{PC-D} & \sol{Même IP que PC-B (doublon).} & \sol{Conflit d'adresse : PC-B et PC-D perturbés.} & \sol{Adresse unique, ex. .80.14.} \\ \hline
    \hl 4 & \sol{R1 G0/0} & \sol{Interface côté LAN désactivée (Off).} & \sol{Aucun poste ne joint le routeur ni le serveur.} & \sol{Port Status : On.} \\ \hline}
  \def\Cbareme{PC-B et SRV-C sont correctement configurés ; accepter PC-B cité à la place de PC-D pour le doublon.}
\fi

\pagestyle{fancy}
\fancyhf{}
\renewcommand{\headrulewidth}{0pt}
\fancyfoot[L]{\footnotesize\itshape\color{grisfoot}Évaluation — Réseaux locaux — Sujet \variante\ifcorrige\ — CORRIGÉ\fi}
\fancyfoot[R]{\footnotesize\color{grisfoot}\thepage/3}

\setlist{itemsep=0pt,topsep=2pt,parsep=0pt}
\setlength{\parindent}{0pt}
\setlength{\parskip}{2pt}

\begin{document}

\begin{center}
  {\Large\bfseries Évaluation pratique — Réseaux locaux \quad
   \fcolorbox{black}{white}{\ SUJET \variante\ }\par}
  \vspace{1pt}
  {\bfseries Adressage IP, commutation et routage — BTS CIEL 1re année — durée : 2 heures\par}
  \ifcorrige\vspace{2pt}{\bfseries\color{corr}CORRIGÉ ET BARÈME — document enseignant}\par\fi
\end{center}
\vspace{-2pt}
{\renewcommand{\arraystretch}{1.35}
\begin{tabularx}{\linewidth}{|X|X|p{2.2cm}|p{2.6cm}|}
  \hline
  \textbf{Nom :} & \textbf{Prénom :} & \textbf{Poste n°} & \textbf{Note :} \hfill /20 \\ \hline
\end{tabularx}}

\textbf{Consignes :} travail \textbf{individuel}, documents et Internet \textbf{non autorisés}. Ce sujet est le
\textbf{document-réponse}. Fichier Packet Tracer : \ip{Nom\_Prenom\_EvalLAN\_\variante.pkt}, à enregistrer
régulièrement dans le dossier de rendu. Les « visas » sont apposés par le professeur, que vous appelez à chaque test réussi.
Durées conseillées : A 30 min — B 1 h 10 — C 20 min.

%==========================================================================
\section{Partie A — Questions de cours et calculs \hfill /6}

\question{A1}{2 pts}{Complétez le tableau en considérant le \textbf{masque par défaut} de la classe de chaque adresse.}
{\small\renewcommand{\arraystretch}{1.15}
\begin{tabularx}{\linewidth}{|l|c|C|C|C|c|}
  \hline
  \textbf{Adresse IP} & \textbf{Classe} & \textbf{Masque par défaut} & \textbf{Adresse réseau} &
  \textbf{Adresse de diffusion} & \textbf{Nb d'hôtes} \\ \hline
  \AIrows
\end{tabularx}}
\bareme{0,5 pt par ligne juste + 0,5 pt si les trois nombres d'hôtes sont justes ($2^{n}-2$, $n$ = nombre de bits d'hôte).}

\question{A2}{1 pt}{Indiquez le type de câble (\textbf{droit} ou \textbf{croisé}) pour chaque liaison.}
{\small\renewcommand{\arraystretch}{1.15}
\begin{tabularx}{\linewidth}{|X|c||X|c|}
  \hline
  \hl PC $\leftrightarrow$ commutateur & \makebox[1.6cm]{\sol{droit}} & Commutateur $\leftrightarrow$ routeur & \makebox[1.6cm]{\sol{droit}} \\ \hline
  \hl PC $\leftrightarrow$ PC (réseau P2P) & \sol{croisé} & Concentrateur $\leftrightarrow$ concentrateur & \sol{croisé} \\ \hline
\end{tabularx}}
\bareme{0,25 pt par case. Équipements de types différents : câble droit ; de même type : câble croisé.}

\question{A3}{1 pt}{Comment un \textbf{concentrateur (hub)} et un \textbf{commutateur (switch)} traitent-ils une trame
reçue ? Lequel limite les collisions, et pourquoi ?}
\zone{2.9cm}{Le \textbf{hub} (couche 1) répète la trame sur \textbf{tous ses autres ports} : un seul domaine de collision
partagé. Le \textbf{switch} (couche 2) lit l'adresse MAC de destination et, grâce à sa \textbf{table MAC}, ne l'envoie
\textbf{que sur le port du destinataire} : un domaine de collision par port, c'est lui qui limite les collisions.
\emph{0,5 pt hub, 0,5 pt switch.}}

\question{A4}{1 pt}{Qu'est-ce que la \textbf{passerelle par défaut} d'un poste ? Quelle condition son adresse doit-elle respecter ?}
\zone{2.4cm}{Adresse IP de l'interface du \textbf{routeur} à laquelle le poste envoie les paquets destinés à un
\textbf{autre réseau}. Elle doit \textbf{appartenir au même réseau} que le poste. \emph{0,5 + 0,5 pt.}}

\question{A5}{0,5 pt}{Un portable configuré en DHCP affiche \ip{\apipa} / \ip{255.255.0.0}. Que s'est-il passé ?}
\zone{1.5cm}{Aucun \textbf{serveur DHCP} n'a répondu : le poste s'est attribué une adresse \textbf{APIPA} (169.254.0.0/16).}

\question{A6}{0,5 pt}{Quel protocole utilise \texttt{ping} ? Nommez les deux messages échangés et la couche OSI des adresses IP.}
\zone{1.5cm}{\textbf{ICMP} ; \textbf{Echo Request} / \textbf{Echo Reply} ; adresses IP en \textbf{couche 3 (Réseau)}.}

%==========================================================================
\newpage
\section{Partie B — Réalisation sous Packet Tracer \hfill /10}

\begin{cdc}
\entreprise{} vous confie son réseau : \textbf{trois réseaux} reliés par un \textbf{routeur unique Cisco 2911}
nommé \ip{\Rnom} (interfaces G0/0, G0/1, G0/2).
\begin{itemize}[leftmargin=12pt]
  \item \textbf{Réseau \resUn} : \textbf{4 postes} PC-\pUn1 à PC-\pUn4, commutateur \ip{SW-\pUn}, adressage statique,
        réseau \ip{\netUn.0/24}.
  \item \textbf{Réseau \resDeux} : \textbf{2 postes} PC-\pDeux1, PC-\pDeux2 et \textbf{1 imprimante} IMP-\pDeux,
        commutateur \ip{SW-\pDeux}, adressage statique, réseau de classe B \ip{\netDeux.0.0/16}.
  \item \textbf{Réseau Serveur} : serveur \ip{\srvNom} d'adresse \ip{\netSrv.\srvH/24}, relié \textbf{directement} au routeur.
  \item La passerelle de chaque réseau est sa \textbf{dernière adresse utilisable}. Tous les équipements communiquent
        entre eux et accèdent à la page web du serveur. Nommage imposé (« Display Name »).
\end{itemize}
\end{cdc}

\question{B1}{2 pts}{Complétez le plan d'adressage \textbf{avant} la réalisation.}
{\footnotesize\setlength{\tabcolsep}{3pt}\renewcommand{\arraystretch}{1.15}
\begin{tabularx}{\linewidth}{|l|C|C|C|C|C|C|}
  \hline
  \textbf{Réseau} & \textbf{Adresse réseau} & \textbf{Masque} & \textbf{1re @ utilisable} &
  \textbf{Dernière @ utilisable} & \textbf{Diffusion} & \textbf{Passerelle} \\ \hline
  \hl \resUn & \sol{\netUn.0} & \sol{255.255.255.0} & \sol{\netUn.1} & \sol{\netUn.254} & \sol{\netUn.255} & \sol{\netUn.254} \\ \hline
  \hl \resDeux & \sol{\netDeux.0.0} & \sol{255.255.0.0} & \sol{\netDeux.0.1} & \sol{\netDeux.255.254} & \sol{\netDeux.255.255} & \sol{\netDeux.255.254} \\ \hline
  \hl Serveur & \sol{\netSrv.0} & \sol{255.255.255.0} & \sol{\netSrv.1} & \sol{\netSrv.254} & \sol{\netSrv.255} & \sol{\netSrv.254} \\ \hline
\end{tabularx}}
\bareme{0,5 pt par ligne juste + 0,5 pt si la convention de passerelle est respectée partout.}

\question{B2}{2 pts}{Construisez la topologie (équipements, câblage, nommage) et relevez l'adressage configuré.}
{\footnotesize\setlength{\tabcolsep}{3pt}\renewcommand{\arraystretch}{1.1}
\begin{tabularx}{\linewidth}{|l|C|C|C||l|C|C|}
  \hline
  \textbf{Hôte} & \textbf{Adresse IP} & \textbf{Masque} & \textbf{Passerelle} & \textbf{Interface} & \textbf{Adresse IP} & \textbf{Masque} \\ \hline
  \hl PC-\pUn1 & \sol{\netUn.1} & \sol{255.255.255.0} & \sol{\netUn.254} & \Rnom{} G0/0 & \sol{\netUn.254} & \sol{255.255.255.0} \\ \hline
  \hl PC-\pDeux1 & \sol{\netDeux.0.1} & \sol{255.255.0.0} & \sol{\netDeux.255.254} & \Rnom{} G0/1 & \sol{\netDeux.255.254} & \sol{255.255.0.0} \\ \hline
  \hl IMP-\pDeux & \sol{\netDeux.0.10} & \sol{255.255.0.0} & \sol{\netDeux.255.254} & \Rnom{} G0/2 & \sol{\netSrv.254} & \sol{255.255.255.0} \\ \hline
  \hl \srvNom & \ip{\netSrv.\srvH} & \sol{255.255.255.0} & \sol{\netSrv.254} & \multicolumn{3}{c|}{\cellcolor{black!8}} \\ \hline
\end{tabularx}}
\bareme{Adresses d'hôtes libres (uniques, dans le bon réseau) ; affectation G0/x libre. Sur le .pkt : équipements 0,5 ;
câblage 1 (câbles droits vers les switchs et switch–routeur, \textbf{croisé} serveur–routeur, liens verts) ; nommage 0,5.}

\question{B3}{3 pts}{Configurez tous les équipements : IP, masques, passerelles des hôtes ; interfaces du routeur
(« Port Status : On »).}
\bareme{Sur le .pkt : interfaces routeur 1 pt ; IP/masques des hôtes 1 pt ; passerelles (imprimante et serveur compris)
1 pt. $-0,25$ par erreur, sans descendre sous 0 par item.}

\question{B4}{3 pts}{Réalisez les tests, notez la commande et le résultat, puis faites viser chaque test réussi.}
{\small\renewcommand{\arraystretch}{1.1}
\begin{tabularx}{\linewidth}{|c|l|l|C|c|c|}
  \hline
  \textbf{N°} & \textbf{Source} & \textbf{Destination} & \textbf{Commande / action} & \textbf{Résultat} & \textbf{Visa prof.} \\ \hline
  \hl 1 & PC-\pUn1 & PC-\pUn4 & \sol{ping \netUn.4} & \sol{réussi} & \hspace{1.4cm} \\ \hline
  \hl 2 & PC-\pDeux1 & IMP-\pDeux & \sol{ping \netDeux.0.10} & \sol{réussi} & \\ \hline
  \hl 3 & PC-\pUn2 & PC-\pDeux2 & \sol{ping \netDeux.0.2} & \sol{réussi} & \\ \hline
  \hl 4 & PC-\pDeux2 & \srvNom & \sol{ping \netSrv.\srvH} & \sol{réussi} & \\ \hline
  \hl 5 & PC-\pUn3 & page web du serveur & \sol{navigateur : {\NoAutoSpacing http://\netSrv.\srvH}} & \sol{page affichée} & \\ \hline
\end{tabularx}}
\bareme{0,5 pt par test visé + 0,5 pt pour B5. Le premier ping inter-réseaux peut échouer (résolution ARP) ; les suivants doivent réussir.}

\question{B5}{incl. B4}{Test n°3 : par quels équipements passe le paquet ? Pourquoi PC-\pUn2 l'envoie-t-il au routeur ?}
\zone{2.3cm}{PC-\pUn2 $\to$ SW-\pUn{} $\to$ \Rnom{} $\to$ SW-\pDeux{} $\to$ PC-\pDeux2. Grâce à son masque, PC-\pUn2 constate que
la destination (\netDeux.0.0/16) n'est \textbf{pas dans son réseau} (\netUn.0/24) : il envoie le paquet à sa
\textbf{passerelle par défaut}, qui le route.}

%==========================================================================
\newpage
\section{Partie C — Diagnostic d'un réseau en panne \hfill /4}

Un technicien constate qu'\textbf{aucun poste ne parvient à joindre le serveur SRV-C} et que certains postes
communiquent mal entre eux. Les câbles sont du bon type.

\begin{center}
\begin{minipage}{0.37\linewidth}\centering
\begin{tikzpicture}[font=\footnotesize,
  eq/.style={draw,fill=qbg,rounded corners=2pt,minimum width=9mm,minimum height=5mm,inner sep=2pt},
  every path/.style={thick}]
  \node[eq] (SW) {SW-1};
  \node[eq,right=9mm of SW] (R) {R1};
  \node[eq,right=7mm of R] (SRV) {SRV-C};
  \node[eq,above left=5mm and -3mm of SW] (A) {PC-A};
  \node[eq,above right=5mm and -3mm of SW] (B) {PC-B};
  \node[eq,below left=5mm and -3mm of SW] (C) {PC-C};
  \node[eq,below right=5mm and -3mm of SW] (D) {PC-D};
  \foreach \n in {A,B,C,D} \draw (\n) -- (SW);
  \draw (SW) -- node[above,font=\tiny]{G0/0} (R);
  \draw (R) -- node[above,font=\tiny]{G0/1} (SRV);
\end{tikzpicture}
\end{minipage}\hfill
\begin{minipage}{0.61\linewidth}\centering
{\small\renewcommand{\arraystretch}{1.05}\setlength{\tabcolsep}{4pt}
\begin{tabular}{lllll}
  \toprule
  \textbf{Équipement} & \textbf{Adresse IP} & \textbf{Masque} & \textbf{Passerelle} & \textbf{Port} \\
  \midrule
  \Crows
  \bottomrule
\end{tabular}}
\end{minipage}
\end{center}

\question{C1}{4 pts}{Relevez les \textbf{quatre anomalies}. Pour chacune : conséquence et correction.}
{\small\renewcommand{\arraystretch}{1.1}
\ifcorrige\def\hc{}\else\def\hc{\rule[-10mm]{0pt}{14mm}}\fi
\begin{tabularx}{\linewidth}{|c|>{\raggedright\arraybackslash}p{2cm}|L|L|L|}
  \hline
  \textbf{N°} & \textbf{Équipement} & \textbf{Anomalie} & \textbf{Conséquence} & \textbf{Correction} \\ \hline
  \ifcorrige\Csol\else
  \hc 1 & & & & \\ \hline
  \hc 2 & & & & \\ \hline
  \hc 3 & & & & \\ \hline
  \hc 4 & & & & \\ \hline
  \fi
\end{tabularx}}
\bareme{1 pt par anomalie (0,25 équipement, anomalie, conséquence, correction ; ordre indifférent). \Cbareme}

%==========================================================================
\section{Bonus — Extension sans fil \hfill +1}

\begin{cdc}Ajoutez au réseau \resUn{} un \textbf{AccessPoint-PT} relié à SW-\pUn{} (SSID \ip{\ssid}, \textbf{WPA2-PSK})
et un portable \ip{Laptop-\pUn} équipé d'un module \texttt{WPC300N}, en adressage \textbf{statique}.
Vérifiez qu'il joint \srvNom.\end{cdc}
{\small\renewcommand{\arraystretch}{1.1}
\begin{tabularx}{\linewidth}{|l|C|C|C|c|}
  \hline
  \textbf{Équipement} & \textbf{Adresse IP} & \textbf{Masque} & \textbf{Passerelle} & \textbf{Visa prof.} \\ \hline
  \hl Laptop-\pUn & \sol{\netUn.50} & \sol{255.255.255.0} & \sol{\netUn.254} & \hspace{1.4cm} \\ \hline
\end{tabularx}}
\bareme{Le point d'accès fait office de pont : le portable est dans \netUn.0/24. Bonus si associé en WPA2-PSK et ping
vers \netSrv.\srvH{} réussi. Note plafonnée à 20.}

\end{document}
"
%   pdflatex -jobname=Eval_Reseau_CorrigeB "\def\variante{B}\def\modecorrige{1}% Évaluation pratique — Réseaux locaux (BTS CIEL 1re année) — sujets A et B
% Compilation (4 PDF) :
%   pdflatex -jobname=Eval_Reseau_SujetA   "\def\variante{A}\def\modecorrige{0}\input{Eval_TP_Reseau_LAN.tex}"
%   pdflatex -jobname=Eval_Reseau_SujetB   "\def\variante{B}\def\modecorrige{0}\input{Eval_TP_Reseau_LAN.tex}"
%   pdflatex -jobname=Eval_Reseau_CorrigeA "\def\variante{A}\def\modecorrige{1}\input{Eval_TP_Reseau_LAN.tex}"
%   pdflatex -jobname=Eval_Reseau_CorrigeB "\def\variante{B}\def\modecorrige{1}\input{Eval_TP_Reseau_LAN.tex}"
\documentclass[10pt,a4paper]{article}

\usepackage[T1]{fontenc}
\usepackage[utf8]{inputenc}
\usepackage{lmodern}
\usepackage[french]{babel}
\usepackage[a4paper,hmargin=1.5cm,top=1.3cm,bottom=1.6cm,footskip=0.7cm]{geometry}
\usepackage[table]{xcolor}
\usepackage{booktabs,tabularx,array}
\usepackage{enumitem}
\usepackage{fancyhdr}
\usepackage[most]{tcolorbox}
\usepackage{tikz}
\usetikzlibrary{positioning}
\usepackage{titlesec}

\providecommand{\modecorrige}{0}
\providecommand{\variante}{A}
\newif\ifcorrige
\ifnum\modecorrige=1 \corrigetrue \else \corrigefalse \fi

\definecolor{ipred}{HTML}{C0392B}
\definecolor{cdcbar}{HTML}{A67C3D}
\definecolor{cdcbg}{HTML}{F4F1EA}
\definecolor{qbar}{HTML}{4A6B5D}
\definecolor{qbg}{HTML}{EDF1EF}
\definecolor{corr}{HTML}{1F5FA8}
\definecolor{corrbg}{HTML}{E8F0FA}
\definecolor{grisfoot}{HTML}{828282}

\newcommand{\ip}[1]{\textcolor{ipred}{\texttt{#1}}}
\newcommand{\pts}[1]{\hfill\mbox{\textbf{(#1)}}}
\newcommand{\sol}[1]{\ifcorrige\textcolor{corr}{#1}\fi}
\newcommand{\bareme}[1]{\ifcorrige\par{\footnotesize\color{corr}#1\par}\fi}
\newcommand{\hl}{\ifcorrige\rule[-1.6mm]{0pt}{5.4mm}\else\rule[-2.2mm]{0pt}{6.6mm}\fi}
\newcolumntype{C}{>{\centering\arraybackslash}X}
\newcolumntype{L}{>{\raggedright\arraybackslash}X}

\titleformat{\section}{\large\bfseries}{}{0pt}{}
\titlespacing*{\section}{0pt}{6pt}{3pt}

\tcbset{qstyle/.style={enhanced,boxrule=0pt,leftrule=2.5pt,arc=0pt,outer arc=0pt,
  left=5pt,right=5pt,top=2pt,bottom=2pt,before skip=4pt,after skip=3pt}}
\newtcolorbox{cdc}{qstyle,colback=cdcbg,colframe=cdcbar}
\newtcolorbox{quest}{qstyle,colback=qbg,colframe=qbar}

% Zone de réponse : vierge dans le sujet, remplie dans le corrigé
\newcommand{\zone}[2]{%
  \ifcorrige
    \begin{tcolorbox}[colback=corrbg,colframe=corr,boxrule=0.5pt,arc=1pt,left=4pt,right=4pt,top=1pt,bottom=1pt,
      before skip=2pt,after skip=2pt,fontupper=\footnotesize\color{corr}]#2\end{tcolorbox}%
  \else
    \begin{tcolorbox}[colback=white,colframe=black!45,boxrule=0.4pt,arc=1pt,height=#1,
      before skip=2pt,after skip=2pt]\end{tcolorbox}%
  \fi}

\newcommand{\question}[3]{\begin{quest}\textbf{#1.} #3\pts{#2}\end{quest}}

%==========================================================================
% Données propres à chaque variante
%==========================================================================
\if\variante A
  \def\entreprise{Le cabinet d'architectes \textbf{ArchiNova}}
  \def\Rnom{R-ArchiNova}
  \def\resUn{Bureau d'études}\def\pUn{BE}\def\netUn{192.168.40}
  \def\resDeux{Accueil}\def\pDeux{ACC}\def\netDeux{172.16}
  \def\srvNom{SRV-Intranet}\def\netSrv{192.168.200}\def\srvH{10}
  \def\ssid{ArchiNova-WiFi}
  \def\AIrows{%
    \hl\ip{192.168.40.25} & \sol{C} & \sol{255.255.255.0} & \sol{192.168.40.0} & \sol{192.168.40.255} & \sol{254} \\ \hline
    \hl\ip{172.16.5.10}   & \sol{B} & \sol{255.255.0.0}   & \sol{172.16.0.0}   & \sol{172.16.255.255} & \sol{65\,534} \\ \hline
    \hl\ip{10.3.2.1}      & \sol{A} & \sol{255.0.0.0}     & \sol{10.0.0.0}     & \sol{10.255.255.255} & \sol{16\,777\,214} \\ \hline}
  \def\apipa{169.254.12.7}
  \def\Cnet{192.168.60}
  \def\Crows{%
    PC-A & \ip{192.168.60.10} & \ip{255.255.255.0} & \ip{192.168.60.254} & — \\
    PC-B & \ip{192.168.60.20} & \ip{255.255.255.0} & \ip{192.168.60.1}   & — \\
    PC-C & \ip{192.168.61.30} & \ip{255.255.255.0} & \ip{192.168.60.254} & — \\
    PC-D & \ip{192.168.60.10} & \ip{255.255.255.0} & \ip{192.168.60.254} & — \\
    SRV-C & \ip{10.0.0.5}     & \ip{255.0.0.0}     & \ip{10.0.0.254}     & — \\
    R1 — G0/0 & \ip{192.168.60.254} & \ip{255.255.255.0} & — & On \\
    R1 — G0/1 & \ip{10.0.0.254}     & \ip{255.0.0.0}     & — & Off \\}
  \def\Csol{%
    \hl 1 & \sol{PC-B} & \sol{Passerelle .60.1 : pas l'adresse du routeur.} & \sol{Joint son réseau, mais aucun autre (serveur).} & \sol{Passerelle 192.168.60.254.} \\ \hline
    \hl 2 & \sol{PC-C} & \sol{IP hors du réseau 192.168.60.0/24.} & \sol{Ne joint ni les postes ni le routeur.} & \sol{Ex. 192.168.60.30.} \\ \hline
    \hl 3 & \sol{PC-D} & \sol{Même IP que PC-A (doublon).} & \sol{Conflit d'adresse : PC-A et PC-D perturbés.} & \sol{Adresse unique, ex. .60.40.} \\ \hline
    \hl 4 & \sol{R1 G0/1} & \sol{Interface désactivée (Off).} & \sol{Réseau du serveur injoignable pour tous.} & \sol{Port Status : On.} \\ \hline}
  \def\Cbareme{PC-A et SRV-C sont correctement configurés ; accepter PC-A cité à la place de PC-D pour le doublon.}
\else
  \def\entreprise{La clinique vétérinaire \textbf{VetoPlus}}
  \def\Rnom{R-VetoPlus}
  \def\resUn{Soins}\def\pUn{SOI}\def\netUn{192.168.70}
  \def\resDeux{Secrétariat}\def\pDeux{SEC}\def\netDeux{172.20}
  \def\srvNom{SRV-Dossiers}\def\netSrv{192.168.250}\def\srvH{20}
  \def\ssid{VetoPlus-WiFi}
  \def\AIrows{%
    \hl\ip{192.168.15.130} & \sol{C} & \sol{255.255.255.0} & \sol{192.168.15.0} & \sol{192.168.15.255} & \sol{254} \\ \hline
    \hl\ip{172.31.200.7}   & \sol{B} & \sol{255.255.0.0}   & \sol{172.31.0.0}   & \sol{172.31.255.255} & \sol{65\,534} \\ \hline
    \hl\ip{10.45.6.9}      & \sol{A} & \sol{255.0.0.0}     & \sol{10.0.0.0}     & \sol{10.255.255.255} & \sol{16\,777\,214} \\ \hline}
  \def\apipa{169.254.88.3}
  \def\Cnet{192.168.80}
  \def\Crows{%
    PC-A & \ip{192.168.8.11}   & \ip{255.255.255.0} & \ip{192.168.80.254} & — \\
    PC-B & \ip{192.168.80.12}  & \ip{255.255.255.0} & \ip{192.168.80.254} & — \\
    PC-C & \ip{192.168.80.13}  & \ip{255.255.255.0} & \ip{192.168.80.100} & — \\
    PC-D & \ip{192.168.80.12}  & \ip{255.255.255.0} & \ip{192.168.80.254} & — \\
    SRV-C & \ip{10.0.0.8}      & \ip{255.0.0.0}     & \ip{10.0.0.254}     & — \\
    R1 — G0/0 & \ip{192.168.80.254} & \ip{255.255.255.0} & — & Off \\
    R1 — G0/1 & \ip{10.0.0.254}     & \ip{255.0.0.0}     & — & On \\}
  \def\Csol{%
    \hl 1 & \sol{PC-A} & \sol{IP hors du réseau 192.168.80.0/24.} & \sol{Ne joint ni les postes ni le routeur.} & \sol{Ex. 192.168.80.11.} \\ \hline
    \hl 2 & \sol{PC-C} & \sol{Passerelle .80.100 : pas l'adresse du routeur.} & \sol{Joint son réseau, mais aucun autre (serveur).} & \sol{Passerelle 192.168.80.254.} \\ \hline
    \hl 3 & \sol{PC-D} & \sol{Même IP que PC-B (doublon).} & \sol{Conflit d'adresse : PC-B et PC-D perturbés.} & \sol{Adresse unique, ex. .80.14.} \\ \hline
    \hl 4 & \sol{R1 G0/0} & \sol{Interface côté LAN désactivée (Off).} & \sol{Aucun poste ne joint le routeur ni le serveur.} & \sol{Port Status : On.} \\ \hline}
  \def\Cbareme{PC-B et SRV-C sont correctement configurés ; accepter PC-B cité à la place de PC-D pour le doublon.}
\fi

\pagestyle{fancy}
\fancyhf{}
\renewcommand{\headrulewidth}{0pt}
\fancyfoot[L]{\footnotesize\itshape\color{grisfoot}Évaluation — Réseaux locaux — Sujet \variante\ifcorrige\ — CORRIGÉ\fi}
\fancyfoot[R]{\footnotesize\color{grisfoot}\thepage/3}

\setlist{itemsep=0pt,topsep=2pt,parsep=0pt}
\setlength{\parindent}{0pt}
\setlength{\parskip}{2pt}

\begin{document}

\begin{center}
  {\Large\bfseries Évaluation pratique — Réseaux locaux \quad
   \fcolorbox{black}{white}{\ SUJET \variante\ }\par}
  \vspace{1pt}
  {\bfseries Adressage IP, commutation et routage — BTS CIEL 1re année — durée : 2 heures\par}
  \ifcorrige\vspace{2pt}{\bfseries\color{corr}CORRIGÉ ET BARÈME — document enseignant}\par\fi
\end{center}
\vspace{-2pt}
{\renewcommand{\arraystretch}{1.35}
\begin{tabularx}{\linewidth}{|X|X|p{2.2cm}|p{2.6cm}|}
  \hline
  \textbf{Nom :} & \textbf{Prénom :} & \textbf{Poste n°} & \textbf{Note :} \hfill /20 \\ \hline
\end{tabularx}}

\textbf{Consignes :} travail \textbf{individuel}, documents et Internet \textbf{non autorisés}. Ce sujet est le
\textbf{document-réponse}. Fichier Packet Tracer : \ip{Nom\_Prenom\_EvalLAN\_\variante.pkt}, à enregistrer
régulièrement dans le dossier de rendu. Les « visas » sont apposés par le professeur, que vous appelez à chaque test réussi.
Durées conseillées : A 30 min — B 1 h 10 — C 20 min.

%==========================================================================
\section{Partie A — Questions de cours et calculs \hfill /6}

\question{A1}{2 pts}{Complétez le tableau en considérant le \textbf{masque par défaut} de la classe de chaque adresse.}
{\small\renewcommand{\arraystretch}{1.15}
\begin{tabularx}{\linewidth}{|l|c|C|C|C|c|}
  \hline
  \textbf{Adresse IP} & \textbf{Classe} & \textbf{Masque par défaut} & \textbf{Adresse réseau} &
  \textbf{Adresse de diffusion} & \textbf{Nb d'hôtes} \\ \hline
  \AIrows
\end{tabularx}}
\bareme{0,5 pt par ligne juste + 0,5 pt si les trois nombres d'hôtes sont justes ($2^{n}-2$, $n$ = nombre de bits d'hôte).}

\question{A2}{1 pt}{Indiquez le type de câble (\textbf{droit} ou \textbf{croisé}) pour chaque liaison.}
{\small\renewcommand{\arraystretch}{1.15}
\begin{tabularx}{\linewidth}{|X|c||X|c|}
  \hline
  \hl PC $\leftrightarrow$ commutateur & \makebox[1.6cm]{\sol{droit}} & Commutateur $\leftrightarrow$ routeur & \makebox[1.6cm]{\sol{droit}} \\ \hline
  \hl PC $\leftrightarrow$ PC (réseau P2P) & \sol{croisé} & Concentrateur $\leftrightarrow$ concentrateur & \sol{croisé} \\ \hline
\end{tabularx}}
\bareme{0,25 pt par case. Équipements de types différents : câble droit ; de même type : câble croisé.}

\question{A3}{1 pt}{Comment un \textbf{concentrateur (hub)} et un \textbf{commutateur (switch)} traitent-ils une trame
reçue ? Lequel limite les collisions, et pourquoi ?}
\zone{2.9cm}{Le \textbf{hub} (couche 1) répète la trame sur \textbf{tous ses autres ports} : un seul domaine de collision
partagé. Le \textbf{switch} (couche 2) lit l'adresse MAC de destination et, grâce à sa \textbf{table MAC}, ne l'envoie
\textbf{que sur le port du destinataire} : un domaine de collision par port, c'est lui qui limite les collisions.
\emph{0,5 pt hub, 0,5 pt switch.}}

\question{A4}{1 pt}{Qu'est-ce que la \textbf{passerelle par défaut} d'un poste ? Quelle condition son adresse doit-elle respecter ?}
\zone{2.4cm}{Adresse IP de l'interface du \textbf{routeur} à laquelle le poste envoie les paquets destinés à un
\textbf{autre réseau}. Elle doit \textbf{appartenir au même réseau} que le poste. \emph{0,5 + 0,5 pt.}}

\question{A5}{0,5 pt}{Un portable configuré en DHCP affiche \ip{\apipa} / \ip{255.255.0.0}. Que s'est-il passé ?}
\zone{1.5cm}{Aucun \textbf{serveur DHCP} n'a répondu : le poste s'est attribué une adresse \textbf{APIPA} (169.254.0.0/16).}

\question{A6}{0,5 pt}{Quel protocole utilise \texttt{ping} ? Nommez les deux messages échangés et la couche OSI des adresses IP.}
\zone{1.5cm}{\textbf{ICMP} ; \textbf{Echo Request} / \textbf{Echo Reply} ; adresses IP en \textbf{couche 3 (Réseau)}.}

%==========================================================================
\newpage
\section{Partie B — Réalisation sous Packet Tracer \hfill /10}

\begin{cdc}
\entreprise{} vous confie son réseau : \textbf{trois réseaux} reliés par un \textbf{routeur unique Cisco 2911}
nommé \ip{\Rnom} (interfaces G0/0, G0/1, G0/2).
\begin{itemize}[leftmargin=12pt]
  \item \textbf{Réseau \resUn} : \textbf{4 postes} PC-\pUn1 à PC-\pUn4, commutateur \ip{SW-\pUn}, adressage statique,
        réseau \ip{\netUn.0/24}.
  \item \textbf{Réseau \resDeux} : \textbf{2 postes} PC-\pDeux1, PC-\pDeux2 et \textbf{1 imprimante} IMP-\pDeux,
        commutateur \ip{SW-\pDeux}, adressage statique, réseau de classe B \ip{\netDeux.0.0/16}.
  \item \textbf{Réseau Serveur} : serveur \ip{\srvNom} d'adresse \ip{\netSrv.\srvH/24}, relié \textbf{directement} au routeur.
  \item La passerelle de chaque réseau est sa \textbf{dernière adresse utilisable}. Tous les équipements communiquent
        entre eux et accèdent à la page web du serveur. Nommage imposé (« Display Name »).
\end{itemize}
\end{cdc}

\question{B1}{2 pts}{Complétez le plan d'adressage \textbf{avant} la réalisation.}
{\footnotesize\setlength{\tabcolsep}{3pt}\renewcommand{\arraystretch}{1.15}
\begin{tabularx}{\linewidth}{|l|C|C|C|C|C|C|}
  \hline
  \textbf{Réseau} & \textbf{Adresse réseau} & \textbf{Masque} & \textbf{1re @ utilisable} &
  \textbf{Dernière @ utilisable} & \textbf{Diffusion} & \textbf{Passerelle} \\ \hline
  \hl \resUn & \sol{\netUn.0} & \sol{255.255.255.0} & \sol{\netUn.1} & \sol{\netUn.254} & \sol{\netUn.255} & \sol{\netUn.254} \\ \hline
  \hl \resDeux & \sol{\netDeux.0.0} & \sol{255.255.0.0} & \sol{\netDeux.0.1} & \sol{\netDeux.255.254} & \sol{\netDeux.255.255} & \sol{\netDeux.255.254} \\ \hline
  \hl Serveur & \sol{\netSrv.0} & \sol{255.255.255.0} & \sol{\netSrv.1} & \sol{\netSrv.254} & \sol{\netSrv.255} & \sol{\netSrv.254} \\ \hline
\end{tabularx}}
\bareme{0,5 pt par ligne juste + 0,5 pt si la convention de passerelle est respectée partout.}

\question{B2}{2 pts}{Construisez la topologie (équipements, câblage, nommage) et relevez l'adressage configuré.}
{\footnotesize\setlength{\tabcolsep}{3pt}\renewcommand{\arraystretch}{1.1}
\begin{tabularx}{\linewidth}{|l|C|C|C||l|C|C|}
  \hline
  \textbf{Hôte} & \textbf{Adresse IP} & \textbf{Masque} & \textbf{Passerelle} & \textbf{Interface} & \textbf{Adresse IP} & \textbf{Masque} \\ \hline
  \hl PC-\pUn1 & \sol{\netUn.1} & \sol{255.255.255.0} & \sol{\netUn.254} & \Rnom{} G0/0 & \sol{\netUn.254} & \sol{255.255.255.0} \\ \hline
  \hl PC-\pDeux1 & \sol{\netDeux.0.1} & \sol{255.255.0.0} & \sol{\netDeux.255.254} & \Rnom{} G0/1 & \sol{\netDeux.255.254} & \sol{255.255.0.0} \\ \hline
  \hl IMP-\pDeux & \sol{\netDeux.0.10} & \sol{255.255.0.0} & \sol{\netDeux.255.254} & \Rnom{} G0/2 & \sol{\netSrv.254} & \sol{255.255.255.0} \\ \hline
  \hl \srvNom & \ip{\netSrv.\srvH} & \sol{255.255.255.0} & \sol{\netSrv.254} & \multicolumn{3}{c|}{\cellcolor{black!8}} \\ \hline
\end{tabularx}}
\bareme{Adresses d'hôtes libres (uniques, dans le bon réseau) ; affectation G0/x libre. Sur le .pkt : équipements 0,5 ;
câblage 1 (câbles droits vers les switchs et switch–routeur, \textbf{croisé} serveur–routeur, liens verts) ; nommage 0,5.}

\question{B3}{3 pts}{Configurez tous les équipements : IP, masques, passerelles des hôtes ; interfaces du routeur
(« Port Status : On »).}
\bareme{Sur le .pkt : interfaces routeur 1 pt ; IP/masques des hôtes 1 pt ; passerelles (imprimante et serveur compris)
1 pt. $-0,25$ par erreur, sans descendre sous 0 par item.}

\question{B4}{3 pts}{Réalisez les tests, notez la commande et le résultat, puis faites viser chaque test réussi.}
{\small\renewcommand{\arraystretch}{1.1}
\begin{tabularx}{\linewidth}{|c|l|l|C|c|c|}
  \hline
  \textbf{N°} & \textbf{Source} & \textbf{Destination} & \textbf{Commande / action} & \textbf{Résultat} & \textbf{Visa prof.} \\ \hline
  \hl 1 & PC-\pUn1 & PC-\pUn4 & \sol{ping \netUn.4} & \sol{réussi} & \hspace{1.4cm} \\ \hline
  \hl 2 & PC-\pDeux1 & IMP-\pDeux & \sol{ping \netDeux.0.10} & \sol{réussi} & \\ \hline
  \hl 3 & PC-\pUn2 & PC-\pDeux2 & \sol{ping \netDeux.0.2} & \sol{réussi} & \\ \hline
  \hl 4 & PC-\pDeux2 & \srvNom & \sol{ping \netSrv.\srvH} & \sol{réussi} & \\ \hline
  \hl 5 & PC-\pUn3 & page web du serveur & \sol{navigateur : {\NoAutoSpacing http://\netSrv.\srvH}} & \sol{page affichée} & \\ \hline
\end{tabularx}}
\bareme{0,5 pt par test visé + 0,5 pt pour B5. Le premier ping inter-réseaux peut échouer (résolution ARP) ; les suivants doivent réussir.}

\question{B5}{incl. B4}{Test n°3 : par quels équipements passe le paquet ? Pourquoi PC-\pUn2 l'envoie-t-il au routeur ?}
\zone{2.3cm}{PC-\pUn2 $\to$ SW-\pUn{} $\to$ \Rnom{} $\to$ SW-\pDeux{} $\to$ PC-\pDeux2. Grâce à son masque, PC-\pUn2 constate que
la destination (\netDeux.0.0/16) n'est \textbf{pas dans son réseau} (\netUn.0/24) : il envoie le paquet à sa
\textbf{passerelle par défaut}, qui le route.}

%==========================================================================
\newpage
\section{Partie C — Diagnostic d'un réseau en panne \hfill /4}

Un technicien constate qu'\textbf{aucun poste ne parvient à joindre le serveur SRV-C} et que certains postes
communiquent mal entre eux. Les câbles sont du bon type.

\begin{center}
\begin{minipage}{0.37\linewidth}\centering
\begin{tikzpicture}[font=\footnotesize,
  eq/.style={draw,fill=qbg,rounded corners=2pt,minimum width=9mm,minimum height=5mm,inner sep=2pt},
  every path/.style={thick}]
  \node[eq] (SW) {SW-1};
  \node[eq,right=9mm of SW] (R) {R1};
  \node[eq,right=7mm of R] (SRV) {SRV-C};
  \node[eq,above left=5mm and -3mm of SW] (A) {PC-A};
  \node[eq,above right=5mm and -3mm of SW] (B) {PC-B};
  \node[eq,below left=5mm and -3mm of SW] (C) {PC-C};
  \node[eq,below right=5mm and -3mm of SW] (D) {PC-D};
  \foreach \n in {A,B,C,D} \draw (\n) -- (SW);
  \draw (SW) -- node[above,font=\tiny]{G0/0} (R);
  \draw (R) -- node[above,font=\tiny]{G0/1} (SRV);
\end{tikzpicture}
\end{minipage}\hfill
\begin{minipage}{0.61\linewidth}\centering
{\small\renewcommand{\arraystretch}{1.05}\setlength{\tabcolsep}{4pt}
\begin{tabular}{lllll}
  \toprule
  \textbf{Équipement} & \textbf{Adresse IP} & \textbf{Masque} & \textbf{Passerelle} & \textbf{Port} \\
  \midrule
  \Crows
  \bottomrule
\end{tabular}}
\end{minipage}
\end{center}

\question{C1}{4 pts}{Relevez les \textbf{quatre anomalies}. Pour chacune : conséquence et correction.}
{\small\renewcommand{\arraystretch}{1.1}
\ifcorrige\def\hc{}\else\def\hc{\rule[-10mm]{0pt}{14mm}}\fi
\begin{tabularx}{\linewidth}{|c|>{\raggedright\arraybackslash}p{2cm}|L|L|L|}
  \hline
  \textbf{N°} & \textbf{Équipement} & \textbf{Anomalie} & \textbf{Conséquence} & \textbf{Correction} \\ \hline
  \ifcorrige\Csol\else
  \hc 1 & & & & \\ \hline
  \hc 2 & & & & \\ \hline
  \hc 3 & & & & \\ \hline
  \hc 4 & & & & \\ \hline
  \fi
\end{tabularx}}
\bareme{1 pt par anomalie (0,25 équipement, anomalie, conséquence, correction ; ordre indifférent). \Cbareme}

%==========================================================================
\section{Bonus — Extension sans fil \hfill +1}

\begin{cdc}Ajoutez au réseau \resUn{} un \textbf{AccessPoint-PT} relié à SW-\pUn{} (SSID \ip{\ssid}, \textbf{WPA2-PSK})
et un portable \ip{Laptop-\pUn} équipé d'un module \texttt{WPC300N}, en adressage \textbf{statique}.
Vérifiez qu'il joint \srvNom.\end{cdc}
{\small\renewcommand{\arraystretch}{1.1}
\begin{tabularx}{\linewidth}{|l|C|C|C|c|}
  \hline
  \textbf{Équipement} & \textbf{Adresse IP} & \textbf{Masque} & \textbf{Passerelle} & \textbf{Visa prof.} \\ \hline
  \hl Laptop-\pUn & \sol{\netUn.50} & \sol{255.255.255.0} & \sol{\netUn.254} & \hspace{1.4cm} \\ \hline
\end{tabularx}}
\bareme{Le point d'accès fait office de pont : le portable est dans \netUn.0/24. Bonus si associé en WPA2-PSK et ping
vers \netSrv.\srvH{} réussi. Note plafonnée à 20.}

\end{document}
"
\documentclass[10pt,a4paper]{article}

\usepackage[T1]{fontenc}
\usepackage[utf8]{inputenc}
\usepackage{lmodern}
\usepackage[french]{babel}
\usepackage[a4paper,hmargin=1.5cm,top=1.3cm,bottom=1.6cm,footskip=0.7cm]{geometry}
\usepackage[table]{xcolor}
\usepackage{booktabs,tabularx,array}
\usepackage{enumitem}
\usepackage{fancyhdr}
\usepackage[most]{tcolorbox}
\usepackage{tikz}
\usetikzlibrary{positioning}
\usepackage{titlesec}

\providecommand{\modecorrige}{0}
\providecommand{\variante}{A}
\newif\ifcorrige
\ifnum\modecorrige=1 \corrigetrue \else \corrigefalse \fi

\definecolor{ipred}{HTML}{C0392B}
\definecolor{cdcbar}{HTML}{A67C3D}
\definecolor{cdcbg}{HTML}{F4F1EA}
\definecolor{qbar}{HTML}{4A6B5D}
\definecolor{qbg}{HTML}{EDF1EF}
\definecolor{corr}{HTML}{1F5FA8}
\definecolor{corrbg}{HTML}{E8F0FA}
\definecolor{grisfoot}{HTML}{828282}

\newcommand{\ip}[1]{\textcolor{ipred}{\texttt{#1}}}
\newcommand{\pts}[1]{\hfill\mbox{\textbf{(#1)}}}
\newcommand{\sol}[1]{\ifcorrige\textcolor{corr}{#1}\fi}
\newcommand{\bareme}[1]{\ifcorrige\par{\footnotesize\color{corr}#1\par}\fi}
\newcommand{\hl}{\ifcorrige\rule[-1.6mm]{0pt}{5.4mm}\else\rule[-2.2mm]{0pt}{6.6mm}\fi}
\newcolumntype{C}{>{\centering\arraybackslash}X}
\newcolumntype{L}{>{\raggedright\arraybackslash}X}

\titleformat{\section}{\large\bfseries}{}{0pt}{}
\titlespacing*{\section}{0pt}{6pt}{3pt}

\tcbset{qstyle/.style={enhanced,boxrule=0pt,leftrule=2.5pt,arc=0pt,outer arc=0pt,
  left=5pt,right=5pt,top=2pt,bottom=2pt,before skip=4pt,after skip=3pt}}
\newtcolorbox{cdc}{qstyle,colback=cdcbg,colframe=cdcbar}
\newtcolorbox{quest}{qstyle,colback=qbg,colframe=qbar}

% Zone de réponse : vierge dans le sujet, remplie dans le corrigé
\newcommand{\zone}[2]{%
  \ifcorrige
    \begin{tcolorbox}[colback=corrbg,colframe=corr,boxrule=0.5pt,arc=1pt,left=4pt,right=4pt,top=1pt,bottom=1pt,
      before skip=2pt,after skip=2pt,fontupper=\footnotesize\color{corr}]#2\end{tcolorbox}%
  \else
    \begin{tcolorbox}[colback=white,colframe=black!45,boxrule=0.4pt,arc=1pt,height=#1,
      before skip=2pt,after skip=2pt]\end{tcolorbox}%
  \fi}

\newcommand{\question}[3]{\begin{quest}\textbf{#1.} #3\pts{#2}\end{quest}}

%==========================================================================
% Données propres à chaque variante
%==========================================================================
\if\variante A
  \def\entreprise{Le cabinet d'architectes \textbf{ArchiNova}}
  \def\Rnom{R-ArchiNova}
  \def\resUn{Bureau d'études}\def\pUn{BE}\def\netUn{192.168.40}
  \def\resDeux{Accueil}\def\pDeux{ACC}\def\netDeux{172.16}
  \def\srvNom{SRV-Intranet}\def\netSrv{192.168.200}\def\srvH{10}
  \def\ssid{ArchiNova-WiFi}
  \def\AIrows{%
    \hl\ip{192.168.40.25} & \sol{C} & \sol{255.255.255.0} & \sol{192.168.40.0} & \sol{192.168.40.255} & \sol{254} \\ \hline
    \hl\ip{172.16.5.10}   & \sol{B} & \sol{255.255.0.0}   & \sol{172.16.0.0}   & \sol{172.16.255.255} & \sol{65\,534} \\ \hline
    \hl\ip{10.3.2.1}      & \sol{A} & \sol{255.0.0.0}     & \sol{10.0.0.0}     & \sol{10.255.255.255} & \sol{16\,777\,214} \\ \hline}
  \def\apipa{169.254.12.7}
  \def\Cnet{192.168.60}
  \def\Crows{%
    PC-A & \ip{192.168.60.10} & \ip{255.255.255.0} & \ip{192.168.60.254} & — \\
    PC-B & \ip{192.168.60.20} & \ip{255.255.255.0} & \ip{192.168.60.1}   & — \\
    PC-C & \ip{192.168.61.30} & \ip{255.255.255.0} & \ip{192.168.60.254} & — \\
    PC-D & \ip{192.168.60.10} & \ip{255.255.255.0} & \ip{192.168.60.254} & — \\
    SRV-C & \ip{10.0.0.5}     & \ip{255.0.0.0}     & \ip{10.0.0.254}     & — \\
    R1 — G0/0 & \ip{192.168.60.254} & \ip{255.255.255.0} & — & On \\
    R1 — G0/1 & \ip{10.0.0.254}     & \ip{255.0.0.0}     & — & Off \\}
  \def\Csol{%
    \hl 1 & \sol{PC-B} & \sol{Passerelle .60.1 : pas l'adresse du routeur.} & \sol{Joint son réseau, mais aucun autre (serveur).} & \sol{Passerelle 192.168.60.254.} \\ \hline
    \hl 2 & \sol{PC-C} & \sol{IP hors du réseau 192.168.60.0/24.} & \sol{Ne joint ni les postes ni le routeur.} & \sol{Ex. 192.168.60.30.} \\ \hline
    \hl 3 & \sol{PC-D} & \sol{Même IP que PC-A (doublon).} & \sol{Conflit d'adresse : PC-A et PC-D perturbés.} & \sol{Adresse unique, ex. .60.40.} \\ \hline
    \hl 4 & \sol{R1 G0/1} & \sol{Interface désactivée (Off).} & \sol{Réseau du serveur injoignable pour tous.} & \sol{Port Status : On.} \\ \hline}
  \def\Cbareme{PC-A et SRV-C sont correctement configurés ; accepter PC-A cité à la place de PC-D pour le doublon.}
\else
  \def\entreprise{La clinique vétérinaire \textbf{VetoPlus}}
  \def\Rnom{R-VetoPlus}
  \def\resUn{Soins}\def\pUn{SOI}\def\netUn{192.168.70}
  \def\resDeux{Secrétariat}\def\pDeux{SEC}\def\netDeux{172.20}
  \def\srvNom{SRV-Dossiers}\def\netSrv{192.168.250}\def\srvH{20}
  \def\ssid{VetoPlus-WiFi}
  \def\AIrows{%
    \hl\ip{192.168.15.130} & \sol{C} & \sol{255.255.255.0} & \sol{192.168.15.0} & \sol{192.168.15.255} & \sol{254} \\ \hline
    \hl\ip{172.31.200.7}   & \sol{B} & \sol{255.255.0.0}   & \sol{172.31.0.0}   & \sol{172.31.255.255} & \sol{65\,534} \\ \hline
    \hl\ip{10.45.6.9}      & \sol{A} & \sol{255.0.0.0}     & \sol{10.0.0.0}     & \sol{10.255.255.255} & \sol{16\,777\,214} \\ \hline}
  \def\apipa{169.254.88.3}
  \def\Cnet{192.168.80}
  \def\Crows{%
    PC-A & \ip{192.168.8.11}   & \ip{255.255.255.0} & \ip{192.168.80.254} & — \\
    PC-B & \ip{192.168.80.12}  & \ip{255.255.255.0} & \ip{192.168.80.254} & — \\
    PC-C & \ip{192.168.80.13}  & \ip{255.255.255.0} & \ip{192.168.80.100} & — \\
    PC-D & \ip{192.168.80.12}  & \ip{255.255.255.0} & \ip{192.168.80.254} & — \\
    SRV-C & \ip{10.0.0.8}      & \ip{255.0.0.0}     & \ip{10.0.0.254}     & — \\
    R1 — G0/0 & \ip{192.168.80.254} & \ip{255.255.255.0} & — & Off \\
    R1 — G0/1 & \ip{10.0.0.254}     & \ip{255.0.0.0}     & — & On \\}
  \def\Csol{%
    \hl 1 & \sol{PC-A} & \sol{IP hors du réseau 192.168.80.0/24.} & \sol{Ne joint ni les postes ni le routeur.} & \sol{Ex. 192.168.80.11.} \\ \hline
    \hl 2 & \sol{PC-C} & \sol{Passerelle .80.100 : pas l'adresse du routeur.} & \sol{Joint son réseau, mais aucun autre (serveur).} & \sol{Passerelle 192.168.80.254.} \\ \hline
    \hl 3 & \sol{PC-D} & \sol{Même IP que PC-B (doublon).} & \sol{Conflit d'adresse : PC-B et PC-D perturbés.} & \sol{Adresse unique, ex. .80.14.} \\ \hline
    \hl 4 & \sol{R1 G0/0} & \sol{Interface côté LAN désactivée (Off).} & \sol{Aucun poste ne joint le routeur ni le serveur.} & \sol{Port Status : On.} \\ \hline}
  \def\Cbareme{PC-B et SRV-C sont correctement configurés ; accepter PC-B cité à la place de PC-D pour le doublon.}
\fi

\pagestyle{fancy}
\fancyhf{}
\renewcommand{\headrulewidth}{0pt}
\fancyfoot[L]{\footnotesize\itshape\color{grisfoot}Évaluation — Réseaux locaux — Sujet \variante\ifcorrige\ — CORRIGÉ\fi}
\fancyfoot[R]{\footnotesize\color{grisfoot}\thepage/3}

\setlist{itemsep=0pt,topsep=2pt,parsep=0pt}
\setlength{\parindent}{0pt}
\setlength{\parskip}{2pt}

\begin{document}

\begin{center}
  {\Large\bfseries Évaluation pratique — Réseaux locaux \quad
   \fcolorbox{black}{white}{\ SUJET \variante\ }\par}
  \vspace{1pt}
  {\bfseries Adressage IP, commutation et routage — BTS CIEL 1re année — durée : 2 heures\par}
  \ifcorrige\vspace{2pt}{\bfseries\color{corr}CORRIGÉ ET BARÈME — document enseignant}\par\fi
\end{center}
\vspace{-2pt}
{\renewcommand{\arraystretch}{1.35}
\begin{tabularx}{\linewidth}{|X|X|p{2.2cm}|p{2.6cm}|}
  \hline
  \textbf{Nom :} & \textbf{Prénom :} & \textbf{Poste n°} & \textbf{Note :} \hfill /20 \\ \hline
\end{tabularx}}

\textbf{Consignes :} travail \textbf{individuel}, documents et Internet \textbf{non autorisés}. Ce sujet est le
\textbf{document-réponse}. Fichier Packet Tracer : \ip{Nom\_Prenom\_EvalLAN\_\variante.pkt}, à enregistrer
régulièrement dans le dossier de rendu. Les « visas » sont apposés par le professeur, que vous appelez à chaque test réussi.
Durées conseillées : A 30 min — B 1 h 10 — C 20 min.

%==========================================================================
\section{Partie A — Questions de cours et calculs \hfill /6}

\question{A1}{2 pts}{Complétez le tableau en considérant le \textbf{masque par défaut} de la classe de chaque adresse.}
{\small\renewcommand{\arraystretch}{1.15}
\begin{tabularx}{\linewidth}{|l|c|C|C|C|c|}
  \hline
  \textbf{Adresse IP} & \textbf{Classe} & \textbf{Masque par défaut} & \textbf{Adresse réseau} &
  \textbf{Adresse de diffusion} & \textbf{Nb d'hôtes} \\ \hline
  \AIrows
\end{tabularx}}
\bareme{0,5 pt par ligne juste + 0,5 pt si les trois nombres d'hôtes sont justes ($2^{n}-2$, $n$ = nombre de bits d'hôte).}

\question{A2}{1 pt}{Indiquez le type de câble (\textbf{droit} ou \textbf{croisé}) pour chaque liaison.}
{\small\renewcommand{\arraystretch}{1.15}
\begin{tabularx}{\linewidth}{|X|c||X|c|}
  \hline
  \hl PC $\leftrightarrow$ commutateur & \makebox[1.6cm]{\sol{droit}} & Commutateur $\leftrightarrow$ routeur & \makebox[1.6cm]{\sol{droit}} \\ \hline
  \hl PC $\leftrightarrow$ PC (réseau P2P) & \sol{croisé} & Concentrateur $\leftrightarrow$ concentrateur & \sol{croisé} \\ \hline
\end{tabularx}}
\bareme{0,25 pt par case. Équipements de types différents : câble droit ; de même type : câble croisé.}

\question{A3}{1 pt}{Comment un \textbf{concentrateur (hub)} et un \textbf{commutateur (switch)} traitent-ils une trame
reçue ? Lequel limite les collisions, et pourquoi ?}
\zone{2.9cm}{Le \textbf{hub} (couche 1) répète la trame sur \textbf{tous ses autres ports} : un seul domaine de collision
partagé. Le \textbf{switch} (couche 2) lit l'adresse MAC de destination et, grâce à sa \textbf{table MAC}, ne l'envoie
\textbf{que sur le port du destinataire} : un domaine de collision par port, c'est lui qui limite les collisions.
\emph{0,5 pt hub, 0,5 pt switch.}}

\question{A4}{1 pt}{Qu'est-ce que la \textbf{passerelle par défaut} d'un poste ? Quelle condition son adresse doit-elle respecter ?}
\zone{2.4cm}{Adresse IP de l'interface du \textbf{routeur} à laquelle le poste envoie les paquets destinés à un
\textbf{autre réseau}. Elle doit \textbf{appartenir au même réseau} que le poste. \emph{0,5 + 0,5 pt.}}

\question{A5}{0,5 pt}{Un portable configuré en DHCP affiche \ip{\apipa} / \ip{255.255.0.0}. Que s'est-il passé ?}
\zone{1.5cm}{Aucun \textbf{serveur DHCP} n'a répondu : le poste s'est attribué une adresse \textbf{APIPA} (169.254.0.0/16).}

\question{A6}{0,5 pt}{Quel protocole utilise \texttt{ping} ? Nommez les deux messages échangés et la couche OSI des adresses IP.}
\zone{1.5cm}{\textbf{ICMP} ; \textbf{Echo Request} / \textbf{Echo Reply} ; adresses IP en \textbf{couche 3 (Réseau)}.}

%==========================================================================
\newpage
\section{Partie B — Réalisation sous Packet Tracer \hfill /10}

\begin{cdc}
\entreprise{} vous confie son réseau : \textbf{trois réseaux} reliés par un \textbf{routeur unique Cisco 2911}
nommé \ip{\Rnom} (interfaces G0/0, G0/1, G0/2).
\begin{itemize}[leftmargin=12pt]
  \item \textbf{Réseau \resUn} : \textbf{4 postes} PC-\pUn1 à PC-\pUn4, commutateur \ip{SW-\pUn}, adressage statique,
        réseau \ip{\netUn.0/24}.
  \item \textbf{Réseau \resDeux} : \textbf{2 postes} PC-\pDeux1, PC-\pDeux2 et \textbf{1 imprimante} IMP-\pDeux,
        commutateur \ip{SW-\pDeux}, adressage statique, réseau de classe B \ip{\netDeux.0.0/16}.
  \item \textbf{Réseau Serveur} : serveur \ip{\srvNom} d'adresse \ip{\netSrv.\srvH/24}, relié \textbf{directement} au routeur.
  \item La passerelle de chaque réseau est sa \textbf{dernière adresse utilisable}. Tous les équipements communiquent
        entre eux et accèdent à la page web du serveur. Nommage imposé (« Display Name »).
\end{itemize}
\end{cdc}

\question{B1}{2 pts}{Complétez le plan d'adressage \textbf{avant} la réalisation.}
{\footnotesize\setlength{\tabcolsep}{3pt}\renewcommand{\arraystretch}{1.15}
\begin{tabularx}{\linewidth}{|l|C|C|C|C|C|C|}
  \hline
  \textbf{Réseau} & \textbf{Adresse réseau} & \textbf{Masque} & \textbf{1re @ utilisable} &
  \textbf{Dernière @ utilisable} & \textbf{Diffusion} & \textbf{Passerelle} \\ \hline
  \hl \resUn & \sol{\netUn.0} & \sol{255.255.255.0} & \sol{\netUn.1} & \sol{\netUn.254} & \sol{\netUn.255} & \sol{\netUn.254} \\ \hline
  \hl \resDeux & \sol{\netDeux.0.0} & \sol{255.255.0.0} & \sol{\netDeux.0.1} & \sol{\netDeux.255.254} & \sol{\netDeux.255.255} & \sol{\netDeux.255.254} \\ \hline
  \hl Serveur & \sol{\netSrv.0} & \sol{255.255.255.0} & \sol{\netSrv.1} & \sol{\netSrv.254} & \sol{\netSrv.255} & \sol{\netSrv.254} \\ \hline
\end{tabularx}}
\bareme{0,5 pt par ligne juste + 0,5 pt si la convention de passerelle est respectée partout.}

\question{B2}{2 pts}{Construisez la topologie (équipements, câblage, nommage) et relevez l'adressage configuré.}
{\footnotesize\setlength{\tabcolsep}{3pt}\renewcommand{\arraystretch}{1.1}
\begin{tabularx}{\linewidth}{|l|C|C|C||l|C|C|}
  \hline
  \textbf{Hôte} & \textbf{Adresse IP} & \textbf{Masque} & \textbf{Passerelle} & \textbf{Interface} & \textbf{Adresse IP} & \textbf{Masque} \\ \hline
  \hl PC-\pUn1 & \sol{\netUn.1} & \sol{255.255.255.0} & \sol{\netUn.254} & \Rnom{} G0/0 & \sol{\netUn.254} & \sol{255.255.255.0} \\ \hline
  \hl PC-\pDeux1 & \sol{\netDeux.0.1} & \sol{255.255.0.0} & \sol{\netDeux.255.254} & \Rnom{} G0/1 & \sol{\netDeux.255.254} & \sol{255.255.0.0} \\ \hline
  \hl IMP-\pDeux & \sol{\netDeux.0.10} & \sol{255.255.0.0} & \sol{\netDeux.255.254} & \Rnom{} G0/2 & \sol{\netSrv.254} & \sol{255.255.255.0} \\ \hline
  \hl \srvNom & \ip{\netSrv.\srvH} & \sol{255.255.255.0} & \sol{\netSrv.254} & \multicolumn{3}{c|}{\cellcolor{black!8}} \\ \hline
\end{tabularx}}
\bareme{Adresses d'hôtes libres (uniques, dans le bon réseau) ; affectation G0/x libre. Sur le .pkt : équipements 0,5 ;
câblage 1 (câbles droits vers les switchs et switch–routeur, \textbf{croisé} serveur–routeur, liens verts) ; nommage 0,5.}

\question{B3}{3 pts}{Configurez tous les équipements : IP, masques, passerelles des hôtes ; interfaces du routeur
(« Port Status : On »).}
\bareme{Sur le .pkt : interfaces routeur 1 pt ; IP/masques des hôtes 1 pt ; passerelles (imprimante et serveur compris)
1 pt. $-0,25$ par erreur, sans descendre sous 0 par item.}

\question{B4}{3 pts}{Réalisez les tests, notez la commande et le résultat, puis faites viser chaque test réussi.}
{\small\renewcommand{\arraystretch}{1.1}
\begin{tabularx}{\linewidth}{|c|l|l|C|c|c|}
  \hline
  \textbf{N°} & \textbf{Source} & \textbf{Destination} & \textbf{Commande / action} & \textbf{Résultat} & \textbf{Visa prof.} \\ \hline
  \hl 1 & PC-\pUn1 & PC-\pUn4 & \sol{ping \netUn.4} & \sol{réussi} & \hspace{1.4cm} \\ \hline
  \hl 2 & PC-\pDeux1 & IMP-\pDeux & \sol{ping \netDeux.0.10} & \sol{réussi} & \\ \hline
  \hl 3 & PC-\pUn2 & PC-\pDeux2 & \sol{ping \netDeux.0.2} & \sol{réussi} & \\ \hline
  \hl 4 & PC-\pDeux2 & \srvNom & \sol{ping \netSrv.\srvH} & \sol{réussi} & \\ \hline
  \hl 5 & PC-\pUn3 & page web du serveur & \sol{navigateur : {\NoAutoSpacing http://\netSrv.\srvH}} & \sol{page affichée} & \\ \hline
\end{tabularx}}
\bareme{0,5 pt par test visé + 0,5 pt pour B5. Le premier ping inter-réseaux peut échouer (résolution ARP) ; les suivants doivent réussir.}

\question{B5}{incl. B4}{Test n°3 : par quels équipements passe le paquet ? Pourquoi PC-\pUn2 l'envoie-t-il au routeur ?}
\zone{2.3cm}{PC-\pUn2 $\to$ SW-\pUn{} $\to$ \Rnom{} $\to$ SW-\pDeux{} $\to$ PC-\pDeux2. Grâce à son masque, PC-\pUn2 constate que
la destination (\netDeux.0.0/16) n'est \textbf{pas dans son réseau} (\netUn.0/24) : il envoie le paquet à sa
\textbf{passerelle par défaut}, qui le route.}

%==========================================================================
\newpage
\section{Partie C — Diagnostic d'un réseau en panne \hfill /4}

Un technicien constate qu'\textbf{aucun poste ne parvient à joindre le serveur SRV-C} et que certains postes
communiquent mal entre eux. Les câbles sont du bon type.

\begin{center}
\begin{minipage}{0.37\linewidth}\centering
\begin{tikzpicture}[font=\footnotesize,
  eq/.style={draw,fill=qbg,rounded corners=2pt,minimum width=9mm,minimum height=5mm,inner sep=2pt},
  every path/.style={thick}]
  \node[eq] (SW) {SW-1};
  \node[eq,right=9mm of SW] (R) {R1};
  \node[eq,right=7mm of R] (SRV) {SRV-C};
  \node[eq,above left=5mm and -3mm of SW] (A) {PC-A};
  \node[eq,above right=5mm and -3mm of SW] (B) {PC-B};
  \node[eq,below left=5mm and -3mm of SW] (C) {PC-C};
  \node[eq,below right=5mm and -3mm of SW] (D) {PC-D};
  \foreach \n in {A,B,C,D} \draw (\n) -- (SW);
  \draw (SW) -- node[above,font=\tiny]{G0/0} (R);
  \draw (R) -- node[above,font=\tiny]{G0/1} (SRV);
\end{tikzpicture}
\end{minipage}\hfill
\begin{minipage}{0.61\linewidth}\centering
{\small\renewcommand{\arraystretch}{1.05}\setlength{\tabcolsep}{4pt}
\begin{tabular}{lllll}
  \toprule
  \textbf{Équipement} & \textbf{Adresse IP} & \textbf{Masque} & \textbf{Passerelle} & \textbf{Port} \\
  \midrule
  \Crows
  \bottomrule
\end{tabular}}
\end{minipage}
\end{center}

\question{C1}{4 pts}{Relevez les \textbf{quatre anomalies}. Pour chacune : conséquence et correction.}
{\small\renewcommand{\arraystretch}{1.1}
\ifcorrige\def\hc{}\else\def\hc{\rule[-10mm]{0pt}{14mm}}\fi
\begin{tabularx}{\linewidth}{|c|>{\raggedright\arraybackslash}p{2cm}|L|L|L|}
  \hline
  \textbf{N°} & \textbf{Équipement} & \textbf{Anomalie} & \textbf{Conséquence} & \textbf{Correction} \\ \hline
  \ifcorrige\Csol\else
  \hc 1 & & & & \\ \hline
  \hc 2 & & & & \\ \hline
  \hc 3 & & & & \\ \hline
  \hc 4 & & & & \\ \hline
  \fi
\end{tabularx}}
\bareme{1 pt par anomalie (0,25 équipement, anomalie, conséquence, correction ; ordre indifférent). \Cbareme}

%==========================================================================
\section{Bonus — Extension sans fil \hfill +1}

\begin{cdc}Ajoutez au réseau \resUn{} un \textbf{AccessPoint-PT} relié à SW-\pUn{} (SSID \ip{\ssid}, \textbf{WPA2-PSK})
et un portable \ip{Laptop-\pUn} équipé d'un module \texttt{WPC300N}, en adressage \textbf{statique}.
Vérifiez qu'il joint \srvNom.\end{cdc}
{\small\renewcommand{\arraystretch}{1.1}
\begin{tabularx}{\linewidth}{|l|C|C|C|c|}
  \hline
  \textbf{Équipement} & \textbf{Adresse IP} & \textbf{Masque} & \textbf{Passerelle} & \textbf{Visa prof.} \\ \hline
  \hl Laptop-\pUn & \sol{\netUn.50} & \sol{255.255.255.0} & \sol{\netUn.254} & \hspace{1.4cm} \\ \hline
\end{tabularx}}
\bareme{Le point d'accès fait office de pont : le portable est dans \netUn.0/24. Bonus si associé en WPA2-PSK et ping
vers \netSrv.\srvH{} réussi. Note plafonnée à 20.}

\end{document}
"
\documentclass[10pt,a4paper]{article}

\usepackage[T1]{fontenc}
\usepackage[utf8]{inputenc}
\usepackage{lmodern}
\usepackage[french]{babel}
\usepackage[a4paper,hmargin=1.5cm,top=1.3cm,bottom=1.6cm,footskip=0.7cm]{geometry}
\usepackage[table]{xcolor}
\usepackage{booktabs,tabularx,array}
\usepackage{enumitem}
\usepackage{fancyhdr}
\usepackage[most]{tcolorbox}
\usepackage{tikz}
\usetikzlibrary{positioning}
\usepackage{titlesec}

\providecommand{\modecorrige}{0}
\providecommand{\variante}{A}
\newif\ifcorrige
\ifnum\modecorrige=1 \corrigetrue \else \corrigefalse \fi

\definecolor{ipred}{HTML}{C0392B}
\definecolor{cdcbar}{HTML}{A67C3D}
\definecolor{cdcbg}{HTML}{F4F1EA}
\definecolor{qbar}{HTML}{4A6B5D}
\definecolor{qbg}{HTML}{EDF1EF}
\definecolor{corr}{HTML}{1F5FA8}
\definecolor{corrbg}{HTML}{E8F0FA}
\definecolor{grisfoot}{HTML}{828282}

\newcommand{\ip}[1]{\textcolor{ipred}{\texttt{#1}}}
\newcommand{\pts}[1]{\hfill\mbox{\textbf{(#1)}}}
\newcommand{\sol}[1]{\ifcorrige\textcolor{corr}{#1}\fi}
\newcommand{\bareme}[1]{\ifcorrige\par{\footnotesize\color{corr}#1\par}\fi}
\newcommand{\hl}{\ifcorrige\rule[-1.6mm]{0pt}{5.4mm}\else\rule[-2.2mm]{0pt}{6.6mm}\fi}
\newcolumntype{C}{>{\centering\arraybackslash}X}
\newcolumntype{L}{>{\raggedright\arraybackslash}X}

\titleformat{\section}{\large\bfseries}{}{0pt}{}
\titlespacing*{\section}{0pt}{6pt}{3pt}

\tcbset{qstyle/.style={enhanced,boxrule=0pt,leftrule=2.5pt,arc=0pt,outer arc=0pt,
  left=5pt,right=5pt,top=2pt,bottom=2pt,before skip=4pt,after skip=3pt}}
\newtcolorbox{cdc}{qstyle,colback=cdcbg,colframe=cdcbar}
\newtcolorbox{quest}{qstyle,colback=qbg,colframe=qbar}

% Zone de réponse : vierge dans le sujet, remplie dans le corrigé
\newcommand{\zone}[2]{%
  \ifcorrige
    \begin{tcolorbox}[colback=corrbg,colframe=corr,boxrule=0.5pt,arc=1pt,left=4pt,right=4pt,top=1pt,bottom=1pt,
      before skip=2pt,after skip=2pt,fontupper=\footnotesize\color{corr}]#2\end{tcolorbox}%
  \else
    \begin{tcolorbox}[colback=white,colframe=black!45,boxrule=0.4pt,arc=1pt,height=#1,
      before skip=2pt,after skip=2pt]\end{tcolorbox}%
  \fi}

\newcommand{\question}[3]{\begin{quest}\textbf{#1.} #3\pts{#2}\end{quest}}

%==========================================================================
% Données propres à chaque variante
%==========================================================================
\if\variante A
  \def\entreprise{Le cabinet d'architectes \textbf{ArchiNova}}
  \def\Rnom{R-ArchiNova}
  \def\resUn{Bureau d'études}\def\pUn{BE}\def\netUn{192.168.40}
  \def\resDeux{Accueil}\def\pDeux{ACC}\def\netDeux{172.16}
  \def\srvNom{SRV-Intranet}\def\netSrv{192.168.200}\def\srvH{10}
  \def\ssid{ArchiNova-WiFi}
  \def\AIrows{%
    \hl\ip{192.168.40.25} & \sol{C} & \sol{255.255.255.0} & \sol{192.168.40.0} & \sol{192.168.40.255} & \sol{254} \\ \hline
    \hl\ip{172.16.5.10}   & \sol{B} & \sol{255.255.0.0}   & \sol{172.16.0.0}   & \sol{172.16.255.255} & \sol{65\,534} \\ \hline
    \hl\ip{10.3.2.1}      & \sol{A} & \sol{255.0.0.0}     & \sol{10.0.0.0}     & \sol{10.255.255.255} & \sol{16\,777\,214} \\ \hline}
  \def\apipa{169.254.12.7}
  \def\Cnet{192.168.60}
  \def\Crows{%
    PC-A & \ip{192.168.60.10} & \ip{255.255.255.0} & \ip{192.168.60.254} & — \\
    PC-B & \ip{192.168.60.20} & \ip{255.255.255.0} & \ip{192.168.60.1}   & — \\
    PC-C & \ip{192.168.61.30} & \ip{255.255.255.0} & \ip{192.168.60.254} & — \\
    PC-D & \ip{192.168.60.10} & \ip{255.255.255.0} & \ip{192.168.60.254} & — \\
    SRV-C & \ip{10.0.0.5}     & \ip{255.0.0.0}     & \ip{10.0.0.254}     & — \\
    R1 — G0/0 & \ip{192.168.60.254} & \ip{255.255.255.0} & — & On \\
    R1 — G0/1 & \ip{10.0.0.254}     & \ip{255.0.0.0}     & — & Off \\}
  \def\Csol{%
    \hl 1 & \sol{PC-B} & \sol{Passerelle .60.1 : pas l'adresse du routeur.} & \sol{Joint son réseau, mais aucun autre (serveur).} & \sol{Passerelle 192.168.60.254.} \\ \hline
    \hl 2 & \sol{PC-C} & \sol{IP hors du réseau 192.168.60.0/24.} & \sol{Ne joint ni les postes ni le routeur.} & \sol{Ex. 192.168.60.30.} \\ \hline
    \hl 3 & \sol{PC-D} & \sol{Même IP que PC-A (doublon).} & \sol{Conflit d'adresse : PC-A et PC-D perturbés.} & \sol{Adresse unique, ex. .60.40.} \\ \hline
    \hl 4 & \sol{R1 G0/1} & \sol{Interface désactivée (Off).} & \sol{Réseau du serveur injoignable pour tous.} & \sol{Port Status : On.} \\ \hline}
  \def\Cbareme{PC-A et SRV-C sont correctement configurés ; accepter PC-A cité à la place de PC-D pour le doublon.}
\else
  \def\entreprise{La clinique vétérinaire \textbf{VetoPlus}}
  \def\Rnom{R-VetoPlus}
  \def\resUn{Soins}\def\pUn{SOI}\def\netUn{192.168.70}
  \def\resDeux{Secrétariat}\def\pDeux{SEC}\def\netDeux{172.20}
  \def\srvNom{SRV-Dossiers}\def\netSrv{192.168.250}\def\srvH{20}
  \def\ssid{VetoPlus-WiFi}
  \def\AIrows{%
    \hl\ip{192.168.15.130} & \sol{C} & \sol{255.255.255.0} & \sol{192.168.15.0} & \sol{192.168.15.255} & \sol{254} \\ \hline
    \hl\ip{172.31.200.7}   & \sol{B} & \sol{255.255.0.0}   & \sol{172.31.0.0}   & \sol{172.31.255.255} & \sol{65\,534} \\ \hline
    \hl\ip{10.45.6.9}      & \sol{A} & \sol{255.0.0.0}     & \sol{10.0.0.0}     & \sol{10.255.255.255} & \sol{16\,777\,214} \\ \hline}
  \def\apipa{169.254.88.3}
  \def\Cnet{192.168.80}
  \def\Crows{%
    PC-A & \ip{192.168.8.11}   & \ip{255.255.255.0} & \ip{192.168.80.254} & — \\
    PC-B & \ip{192.168.80.12}  & \ip{255.255.255.0} & \ip{192.168.80.254} & — \\
    PC-C & \ip{192.168.80.13}  & \ip{255.255.255.0} & \ip{192.168.80.100} & — \\
    PC-D & \ip{192.168.80.12}  & \ip{255.255.255.0} & \ip{192.168.80.254} & — \\
    SRV-C & \ip{10.0.0.8}      & \ip{255.0.0.0}     & \ip{10.0.0.254}     & — \\
    R1 — G0/0 & \ip{192.168.80.254} & \ip{255.255.255.0} & — & Off \\
    R1 — G0/1 & \ip{10.0.0.254}     & \ip{255.0.0.0}     & — & On \\}
  \def\Csol{%
    \hl 1 & \sol{PC-A} & \sol{IP hors du réseau 192.168.80.0/24.} & \sol{Ne joint ni les postes ni le routeur.} & \sol{Ex. 192.168.80.11.} \\ \hline
    \hl 2 & \sol{PC-C} & \sol{Passerelle .80.100 : pas l'adresse du routeur.} & \sol{Joint son réseau, mais aucun autre (serveur).} & \sol{Passerelle 192.168.80.254.} \\ \hline
    \hl 3 & \sol{PC-D} & \sol{Même IP que PC-B (doublon).} & \sol{Conflit d'adresse : PC-B et PC-D perturbés.} & \sol{Adresse unique, ex. .80.14.} \\ \hline
    \hl 4 & \sol{R1 G0/0} & \sol{Interface côté LAN désactivée (Off).} & \sol{Aucun poste ne joint le routeur ni le serveur.} & \sol{Port Status : On.} \\ \hline}
  \def\Cbareme{PC-B et SRV-C sont correctement configurés ; accepter PC-B cité à la place de PC-D pour le doublon.}
\fi

\pagestyle{fancy}
\fancyhf{}
\renewcommand{\headrulewidth}{0pt}
\fancyfoot[L]{\footnotesize\itshape\color{grisfoot}Évaluation — Réseaux locaux — Sujet \variante\ifcorrige\ — CORRIGÉ\fi}
\fancyfoot[R]{\footnotesize\color{grisfoot}\thepage/3}

\setlist{itemsep=0pt,topsep=2pt,parsep=0pt}
\setlength{\parindent}{0pt}
\setlength{\parskip}{2pt}

\begin{document}

\begin{center}
  {\Large\bfseries Évaluation pratique — Réseaux locaux \quad
   \fcolorbox{black}{white}{\ SUJET \variante\ }\par}
  \vspace{1pt}
  {\bfseries Adressage IP, commutation et routage — BTS CIEL 1re année — durée : 2 heures\par}
  \ifcorrige\vspace{2pt}{\bfseries\color{corr}CORRIGÉ ET BARÈME — document enseignant}\par\fi
\end{center}
\vspace{-2pt}
{\renewcommand{\arraystretch}{1.35}
\begin{tabularx}{\linewidth}{|X|X|p{2.2cm}|p{2.6cm}|}
  \hline
  \textbf{Nom :} & \textbf{Prénom :} & \textbf{Poste n°} & \textbf{Note :} \hfill /20 \\ \hline
\end{tabularx}}

\textbf{Consignes :} travail \textbf{individuel}, documents et Internet \textbf{non autorisés}. Ce sujet est le
\textbf{document-réponse}. Fichier Packet Tracer : \ip{Nom\_Prenom\_EvalLAN\_\variante.pkt}, à enregistrer
régulièrement dans le dossier de rendu. Les « visas » sont apposés par le professeur, que vous appelez à chaque test réussi.
Durées conseillées : A 30 min — B 1 h 10 — C 20 min.

%==========================================================================
\section{Partie A — Questions de cours et calculs \hfill /6}

\question{A1}{2 pts}{Complétez le tableau en considérant le \textbf{masque par défaut} de la classe de chaque adresse.}
{\small\renewcommand{\arraystretch}{1.15}
\begin{tabularx}{\linewidth}{|l|c|C|C|C|c|}
  \hline
  \textbf{Adresse IP} & \textbf{Classe} & \textbf{Masque par défaut} & \textbf{Adresse réseau} &
  \textbf{Adresse de diffusion} & \textbf{Nb d'hôtes} \\ \hline
  \AIrows
\end{tabularx}}
\bareme{0,5 pt par ligne juste + 0,5 pt si les trois nombres d'hôtes sont justes ($2^{n}-2$, $n$ = nombre de bits d'hôte).}

\question{A2}{1 pt}{Indiquez le type de câble (\textbf{droit} ou \textbf{croisé}) pour chaque liaison.}
{\small\renewcommand{\arraystretch}{1.15}
\begin{tabularx}{\linewidth}{|X|c||X|c|}
  \hline
  \hl PC $\leftrightarrow$ commutateur & \makebox[1.6cm]{\sol{droit}} & Commutateur $\leftrightarrow$ routeur & \makebox[1.6cm]{\sol{droit}} \\ \hline
  \hl PC $\leftrightarrow$ PC (réseau P2P) & \sol{croisé} & Concentrateur $\leftrightarrow$ concentrateur & \sol{croisé} \\ \hline
\end{tabularx}}
\bareme{0,25 pt par case. Équipements de types différents : câble droit ; de même type : câble croisé.}

\question{A3}{1 pt}{Comment un \textbf{concentrateur (hub)} et un \textbf{commutateur (switch)} traitent-ils une trame
reçue ? Lequel limite les collisions, et pourquoi ?}
\zone{2.9cm}{Le \textbf{hub} (couche 1) répète la trame sur \textbf{tous ses autres ports} : un seul domaine de collision
partagé. Le \textbf{switch} (couche 2) lit l'adresse MAC de destination et, grâce à sa \textbf{table MAC}, ne l'envoie
\textbf{que sur le port du destinataire} : un domaine de collision par port, c'est lui qui limite les collisions.
\emph{0,5 pt hub, 0,5 pt switch.}}

\question{A4}{1 pt}{Qu'est-ce que la \textbf{passerelle par défaut} d'un poste ? Quelle condition son adresse doit-elle respecter ?}
\zone{2.4cm}{Adresse IP de l'interface du \textbf{routeur} à laquelle le poste envoie les paquets destinés à un
\textbf{autre réseau}. Elle doit \textbf{appartenir au même réseau} que le poste. \emph{0,5 + 0,5 pt.}}

\question{A5}{0,5 pt}{Un portable configuré en DHCP affiche \ip{\apipa} / \ip{255.255.0.0}. Que s'est-il passé ?}
\zone{1.5cm}{Aucun \textbf{serveur DHCP} n'a répondu : le poste s'est attribué une adresse \textbf{APIPA} (169.254.0.0/16).}

\question{A6}{0,5 pt}{Quel protocole utilise \texttt{ping} ? Nommez les deux messages échangés et la couche OSI des adresses IP.}
\zone{1.5cm}{\textbf{ICMP} ; \textbf{Echo Request} / \textbf{Echo Reply} ; adresses IP en \textbf{couche 3 (Réseau)}.}

%==========================================================================
\newpage
\section{Partie B — Réalisation sous Packet Tracer \hfill /10}

\begin{cdc}
\entreprise{} vous confie son réseau : \textbf{trois réseaux} reliés par un \textbf{routeur unique Cisco 2911}
nommé \ip{\Rnom} (interfaces G0/0, G0/1, G0/2).
\begin{itemize}[leftmargin=12pt]
  \item \textbf{Réseau \resUn} : \textbf{4 postes} PC-\pUn1 à PC-\pUn4, commutateur \ip{SW-\pUn}, adressage statique,
        réseau \ip{\netUn.0/24}.
  \item \textbf{Réseau \resDeux} : \textbf{2 postes} PC-\pDeux1, PC-\pDeux2 et \textbf{1 imprimante} IMP-\pDeux,
        commutateur \ip{SW-\pDeux}, adressage statique, réseau de classe B \ip{\netDeux.0.0/16}.
  \item \textbf{Réseau Serveur} : serveur \ip{\srvNom} d'adresse \ip{\netSrv.\srvH/24}, relié \textbf{directement} au routeur.
  \item La passerelle de chaque réseau est sa \textbf{dernière adresse utilisable}. Tous les équipements communiquent
        entre eux et accèdent à la page web du serveur. Nommage imposé (« Display Name »).
\end{itemize}
\end{cdc}

\question{B1}{2 pts}{Complétez le plan d'adressage \textbf{avant} la réalisation.}
{\footnotesize\setlength{\tabcolsep}{3pt}\renewcommand{\arraystretch}{1.15}
\begin{tabularx}{\linewidth}{|l|C|C|C|C|C|C|}
  \hline
  \textbf{Réseau} & \textbf{Adresse réseau} & \textbf{Masque} & \textbf{1re @ utilisable} &
  \textbf{Dernière @ utilisable} & \textbf{Diffusion} & \textbf{Passerelle} \\ \hline
  \hl \resUn & \sol{\netUn.0} & \sol{255.255.255.0} & \sol{\netUn.1} & \sol{\netUn.254} & \sol{\netUn.255} & \sol{\netUn.254} \\ \hline
  \hl \resDeux & \sol{\netDeux.0.0} & \sol{255.255.0.0} & \sol{\netDeux.0.1} & \sol{\netDeux.255.254} & \sol{\netDeux.255.255} & \sol{\netDeux.255.254} \\ \hline
  \hl Serveur & \sol{\netSrv.0} & \sol{255.255.255.0} & \sol{\netSrv.1} & \sol{\netSrv.254} & \sol{\netSrv.255} & \sol{\netSrv.254} \\ \hline
\end{tabularx}}
\bareme{0,5 pt par ligne juste + 0,5 pt si la convention de passerelle est respectée partout.}

\question{B2}{2 pts}{Construisez la topologie (équipements, câblage, nommage) et relevez l'adressage configuré.}
{\footnotesize\setlength{\tabcolsep}{3pt}\renewcommand{\arraystretch}{1.1}
\begin{tabularx}{\linewidth}{|l|C|C|C||l|C|C|}
  \hline
  \textbf{Hôte} & \textbf{Adresse IP} & \textbf{Masque} & \textbf{Passerelle} & \textbf{Interface} & \textbf{Adresse IP} & \textbf{Masque} \\ \hline
  \hl PC-\pUn1 & \sol{\netUn.1} & \sol{255.255.255.0} & \sol{\netUn.254} & \Rnom{} G0/0 & \sol{\netUn.254} & \sol{255.255.255.0} \\ \hline
  \hl PC-\pDeux1 & \sol{\netDeux.0.1} & \sol{255.255.0.0} & \sol{\netDeux.255.254} & \Rnom{} G0/1 & \sol{\netDeux.255.254} & \sol{255.255.0.0} \\ \hline
  \hl IMP-\pDeux & \sol{\netDeux.0.10} & \sol{255.255.0.0} & \sol{\netDeux.255.254} & \Rnom{} G0/2 & \sol{\netSrv.254} & \sol{255.255.255.0} \\ \hline
  \hl \srvNom & \ip{\netSrv.\srvH} & \sol{255.255.255.0} & \sol{\netSrv.254} & \multicolumn{3}{c|}{\cellcolor{black!8}} \\ \hline
\end{tabularx}}
\bareme{Adresses d'hôtes libres (uniques, dans le bon réseau) ; affectation G0/x libre. Sur le .pkt : équipements 0,5 ;
câblage 1 (câbles droits vers les switchs et switch–routeur, \textbf{croisé} serveur–routeur, liens verts) ; nommage 0,5.}

\question{B3}{3 pts}{Configurez tous les équipements : IP, masques, passerelles des hôtes ; interfaces du routeur
(« Port Status : On »).}
\bareme{Sur le .pkt : interfaces routeur 1 pt ; IP/masques des hôtes 1 pt ; passerelles (imprimante et serveur compris)
1 pt. $-0,25$ par erreur, sans descendre sous 0 par item.}

\question{B4}{3 pts}{Réalisez les tests, notez la commande et le résultat, puis faites viser chaque test réussi.}
{\small\renewcommand{\arraystretch}{1.1}
\begin{tabularx}{\linewidth}{|c|l|l|C|c|c|}
  \hline
  \textbf{N°} & \textbf{Source} & \textbf{Destination} & \textbf{Commande / action} & \textbf{Résultat} & \textbf{Visa prof.} \\ \hline
  \hl 1 & PC-\pUn1 & PC-\pUn4 & \sol{ping \netUn.4} & \sol{réussi} & \hspace{1.4cm} \\ \hline
  \hl 2 & PC-\pDeux1 & IMP-\pDeux & \sol{ping \netDeux.0.10} & \sol{réussi} & \\ \hline
  \hl 3 & PC-\pUn2 & PC-\pDeux2 & \sol{ping \netDeux.0.2} & \sol{réussi} & \\ \hline
  \hl 4 & PC-\pDeux2 & \srvNom & \sol{ping \netSrv.\srvH} & \sol{réussi} & \\ \hline
  \hl 5 & PC-\pUn3 & page web du serveur & \sol{navigateur : {\NoAutoSpacing http://\netSrv.\srvH}} & \sol{page affichée} & \\ \hline
\end{tabularx}}
\bareme{0,5 pt par test visé + 0,5 pt pour B5. Le premier ping inter-réseaux peut échouer (résolution ARP) ; les suivants doivent réussir.}

\question{B5}{incl. B4}{Test n°3 : par quels équipements passe le paquet ? Pourquoi PC-\pUn2 l'envoie-t-il au routeur ?}
\zone{2.3cm}{PC-\pUn2 $\to$ SW-\pUn{} $\to$ \Rnom{} $\to$ SW-\pDeux{} $\to$ PC-\pDeux2. Grâce à son masque, PC-\pUn2 constate que
la destination (\netDeux.0.0/16) n'est \textbf{pas dans son réseau} (\netUn.0/24) : il envoie le paquet à sa
\textbf{passerelle par défaut}, qui le route.}

%==========================================================================
\newpage
\section{Partie C — Diagnostic d'un réseau en panne \hfill /4}

Un technicien constate qu'\textbf{aucun poste ne parvient à joindre le serveur SRV-C} et que certains postes
communiquent mal entre eux. Les câbles sont du bon type.

\begin{center}
\begin{minipage}{0.37\linewidth}\centering
\begin{tikzpicture}[font=\footnotesize,
  eq/.style={draw,fill=qbg,rounded corners=2pt,minimum width=9mm,minimum height=5mm,inner sep=2pt},
  every path/.style={thick}]
  \node[eq] (SW) {SW-1};
  \node[eq,right=9mm of SW] (R) {R1};
  \node[eq,right=7mm of R] (SRV) {SRV-C};
  \node[eq,above left=5mm and -3mm of SW] (A) {PC-A};
  \node[eq,above right=5mm and -3mm of SW] (B) {PC-B};
  \node[eq,below left=5mm and -3mm of SW] (C) {PC-C};
  \node[eq,below right=5mm and -3mm of SW] (D) {PC-D};
  \foreach \n in {A,B,C,D} \draw (\n) -- (SW);
  \draw (SW) -- node[above,font=\tiny]{G0/0} (R);
  \draw (R) -- node[above,font=\tiny]{G0/1} (SRV);
\end{tikzpicture}
\end{minipage}\hfill
\begin{minipage}{0.61\linewidth}\centering
{\small\renewcommand{\arraystretch}{1.05}\setlength{\tabcolsep}{4pt}
\begin{tabular}{lllll}
  \toprule
  \textbf{Équipement} & \textbf{Adresse IP} & \textbf{Masque} & \textbf{Passerelle} & \textbf{Port} \\
  \midrule
  \Crows
  \bottomrule
\end{tabular}}
\end{minipage}
\end{center}

\question{C1}{4 pts}{Relevez les \textbf{quatre anomalies}. Pour chacune : conséquence et correction.}
{\small\renewcommand{\arraystretch}{1.1}
\ifcorrige\def\hc{}\else\def\hc{\rule[-10mm]{0pt}{14mm}}\fi
\begin{tabularx}{\linewidth}{|c|>{\raggedright\arraybackslash}p{2cm}|L|L|L|}
  \hline
  \textbf{N°} & \textbf{Équipement} & \textbf{Anomalie} & \textbf{Conséquence} & \textbf{Correction} \\ \hline
  \ifcorrige\Csol\else
  \hc 1 & & & & \\ \hline
  \hc 2 & & & & \\ \hline
  \hc 3 & & & & \\ \hline
  \hc 4 & & & & \\ \hline
  \fi
\end{tabularx}}
\bareme{1 pt par anomalie (0,25 équipement, anomalie, conséquence, correction ; ordre indifférent). \Cbareme}

%==========================================================================
\section{Bonus — Extension sans fil \hfill +1}

\begin{cdc}Ajoutez au réseau \resUn{} un \textbf{AccessPoint-PT} relié à SW-\pUn{} (SSID \ip{\ssid}, \textbf{WPA2-PSK})
et un portable \ip{Laptop-\pUn} équipé d'un module \texttt{WPC300N}, en adressage \textbf{statique}.
Vérifiez qu'il joint \srvNom.\end{cdc}
{\small\renewcommand{\arraystretch}{1.1}
\begin{tabularx}{\linewidth}{|l|C|C|C|c|}
  \hline
  \textbf{Équipement} & \textbf{Adresse IP} & \textbf{Masque} & \textbf{Passerelle} & \textbf{Visa prof.} \\ \hline
  \hl Laptop-\pUn & \sol{\netUn.50} & \sol{255.255.255.0} & \sol{\netUn.254} & \hspace{1.4cm} \\ \hline
\end{tabularx}}
\bareme{Le point d'accès fait office de pont : le portable est dans \netUn.0/24. Bonus si associé en WPA2-PSK et ping
vers \netSrv.\srvH{} réussi. Note plafonnée à 20.}

\end{document}
"
\documentclass[10pt,a4paper]{article}

\usepackage[T1]{fontenc}
\usepackage[utf8]{inputenc}
\usepackage{lmodern}
\usepackage[french]{babel}
\usepackage[a4paper,hmargin=1.5cm,top=1.3cm,bottom=1.6cm,footskip=0.7cm]{geometry}
\usepackage[table]{xcolor}
\usepackage{booktabs,tabularx,array}
\usepackage{enumitem}
\usepackage{fancyhdr}
\usepackage[most]{tcolorbox}
\usepackage{tikz}
\usetikzlibrary{positioning}
\usepackage{titlesec}

\providecommand{\modecorrige}{0}
\providecommand{\variante}{A}
\newif\ifcorrige
\ifnum\modecorrige=1 \corrigetrue \else \corrigefalse \fi

\definecolor{ipred}{HTML}{C0392B}
\definecolor{cdcbar}{HTML}{A67C3D}
\definecolor{cdcbg}{HTML}{F4F1EA}
\definecolor{qbar}{HTML}{4A6B5D}
\definecolor{qbg}{HTML}{EDF1EF}
\definecolor{corr}{HTML}{1F5FA8}
\definecolor{corrbg}{HTML}{E8F0FA}
\definecolor{grisfoot}{HTML}{828282}

\newcommand{\ip}[1]{\textcolor{ipred}{\texttt{#1}}}
\newcommand{\pts}[1]{\hfill\mbox{\textbf{(#1)}}}
\newcommand{\sol}[1]{\ifcorrige\textcolor{corr}{#1}\fi}
\newcommand{\bareme}[1]{\ifcorrige\par{\footnotesize\color{corr}#1\par}\fi}
\newcommand{\hl}{\ifcorrige\rule[-1.6mm]{0pt}{5.4mm}\else\rule[-2.2mm]{0pt}{6.6mm}\fi}
\newcolumntype{C}{>{\centering\arraybackslash}X}
\newcolumntype{L}{>{\raggedright\arraybackslash}X}

\titleformat{\section}{\large\bfseries}{}{0pt}{}
\titlespacing*{\section}{0pt}{6pt}{3pt}

\tcbset{qstyle/.style={enhanced,boxrule=0pt,leftrule=2.5pt,arc=0pt,outer arc=0pt,
  left=5pt,right=5pt,top=2pt,bottom=2pt,before skip=4pt,after skip=3pt}}
\newtcolorbox{cdc}{qstyle,colback=cdcbg,colframe=cdcbar}
\newtcolorbox{quest}{qstyle,colback=qbg,colframe=qbar}

% Zone de réponse : vierge dans le sujet, remplie dans le corrigé
\newcommand{\zone}[2]{%
  \ifcorrige
    \begin{tcolorbox}[colback=corrbg,colframe=corr,boxrule=0.5pt,arc=1pt,left=4pt,right=4pt,top=1pt,bottom=1pt,
      before skip=2pt,after skip=2pt,fontupper=\footnotesize\color{corr}]#2\end{tcolorbox}%
  \else
    \begin{tcolorbox}[colback=white,colframe=black!45,boxrule=0.4pt,arc=1pt,height=#1,
      before skip=2pt,after skip=2pt]\end{tcolorbox}%
  \fi}

\newcommand{\question}[3]{\begin{quest}\textbf{#1.} #3\pts{#2}\end{quest}}

%==========================================================================
% Données propres à chaque variante
%==========================================================================
\if\variante A
  \def\entreprise{Le cabinet d'architectes \textbf{ArchiNova}}
  \def\Rnom{R-ArchiNova}
  \def\resUn{Bureau d'études}\def\pUn{BE}\def\netUn{192.168.40}
  \def\resDeux{Accueil}\def\pDeux{ACC}\def\netDeux{172.16}
  \def\srvNom{SRV-Intranet}\def\netSrv{192.168.200}\def\srvH{10}
  \def\ssid{ArchiNova-WiFi}
  \def\AIrows{%
    \hl\ip{192.168.40.25} & \sol{C} & \sol{255.255.255.0} & \sol{192.168.40.0} & \sol{192.168.40.255} & \sol{254} \\ \hline
    \hl\ip{172.16.5.10}   & \sol{B} & \sol{255.255.0.0}   & \sol{172.16.0.0}   & \sol{172.16.255.255} & \sol{65\,534} \\ \hline
    \hl\ip{10.3.2.1}      & \sol{A} & \sol{255.0.0.0}     & \sol{10.0.0.0}     & \sol{10.255.255.255} & \sol{16\,777\,214} \\ \hline}
  \def\apipa{169.254.12.7}
  \def\Cnet{192.168.60}
  \def\Crows{%
    PC-A & \ip{192.168.60.10} & \ip{255.255.255.0} & \ip{192.168.60.254} & — \\
    PC-B & \ip{192.168.60.20} & \ip{255.255.255.0} & \ip{192.168.60.1}   & — \\
    PC-C & \ip{192.168.61.30} & \ip{255.255.255.0} & \ip{192.168.60.254} & — \\
    PC-D & \ip{192.168.60.10} & \ip{255.255.255.0} & \ip{192.168.60.254} & — \\
    SRV-C & \ip{10.0.0.5}     & \ip{255.0.0.0}     & \ip{10.0.0.254}     & — \\
    R1 — G0/0 & \ip{192.168.60.254} & \ip{255.255.255.0} & — & On \\
    R1 — G0/1 & \ip{10.0.0.254}     & \ip{255.0.0.0}     & — & Off \\}
  \def\Csol{%
    \hl 1 & \sol{PC-B} & \sol{Passerelle .60.1 : pas l'adresse du routeur.} & \sol{Joint son réseau, mais aucun autre (serveur).} & \sol{Passerelle 192.168.60.254.} \\ \hline
    \hl 2 & \sol{PC-C} & \sol{IP hors du réseau 192.168.60.0/24.} & \sol{Ne joint ni les postes ni le routeur.} & \sol{Ex. 192.168.60.30.} \\ \hline
    \hl 3 & \sol{PC-D} & \sol{Même IP que PC-A (doublon).} & \sol{Conflit d'adresse : PC-A et PC-D perturbés.} & \sol{Adresse unique, ex. .60.40.} \\ \hline
    \hl 4 & \sol{R1 G0/1} & \sol{Interface désactivée (Off).} & \sol{Réseau du serveur injoignable pour tous.} & \sol{Port Status : On.} \\ \hline}
  \def\Cbareme{PC-A et SRV-C sont correctement configurés ; accepter PC-A cité à la place de PC-D pour le doublon.}
\else
  \def\entreprise{La clinique vétérinaire \textbf{VetoPlus}}
  \def\Rnom{R-VetoPlus}
  \def\resUn{Soins}\def\pUn{SOI}\def\netUn{192.168.70}
  \def\resDeux{Secrétariat}\def\pDeux{SEC}\def\netDeux{172.20}
  \def\srvNom{SRV-Dossiers}\def\netSrv{192.168.250}\def\srvH{20}
  \def\ssid{VetoPlus-WiFi}
  \def\AIrows{%
    \hl\ip{192.168.15.130} & \sol{C} & \sol{255.255.255.0} & \sol{192.168.15.0} & \sol{192.168.15.255} & \sol{254} \\ \hline
    \hl\ip{172.31.200.7}   & \sol{B} & \sol{255.255.0.0}   & \sol{172.31.0.0}   & \sol{172.31.255.255} & \sol{65\,534} \\ \hline
    \hl\ip{10.45.6.9}      & \sol{A} & \sol{255.0.0.0}     & \sol{10.0.0.0}     & \sol{10.255.255.255} & \sol{16\,777\,214} \\ \hline}
  \def\apipa{169.254.88.3}
  \def\Cnet{192.168.80}
  \def\Crows{%
    PC-A & \ip{192.168.8.11}   & \ip{255.255.255.0} & \ip{192.168.80.254} & — \\
    PC-B & \ip{192.168.80.12}  & \ip{255.255.255.0} & \ip{192.168.80.254} & — \\
    PC-C & \ip{192.168.80.13}  & \ip{255.255.255.0} & \ip{192.168.80.100} & — \\
    PC-D & \ip{192.168.80.12}  & \ip{255.255.255.0} & \ip{192.168.80.254} & — \\
    SRV-C & \ip{10.0.0.8}      & \ip{255.0.0.0}     & \ip{10.0.0.254}     & — \\
    R1 — G0/0 & \ip{192.168.80.254} & \ip{255.255.255.0} & — & Off \\
    R1 — G0/1 & \ip{10.0.0.254}     & \ip{255.0.0.0}     & — & On \\}
  \def\Csol{%
    \hl 1 & \sol{PC-A} & \sol{IP hors du réseau 192.168.80.0/24.} & \sol{Ne joint ni les postes ni le routeur.} & \sol{Ex. 192.168.80.11.} \\ \hline
    \hl 2 & \sol{PC-C} & \sol{Passerelle .80.100 : pas l'adresse du routeur.} & \sol{Joint son réseau, mais aucun autre (serveur).} & \sol{Passerelle 192.168.80.254.} \\ \hline
    \hl 3 & \sol{PC-D} & \sol{Même IP que PC-B (doublon).} & \sol{Conflit d'adresse : PC-B et PC-D perturbés.} & \sol{Adresse unique, ex. .80.14.} \\ \hline
    \hl 4 & \sol{R1 G0/0} & \sol{Interface côté LAN désactivée (Off).} & \sol{Aucun poste ne joint le routeur ni le serveur.} & \sol{Port Status : On.} \\ \hline}
  \def\Cbareme{PC-B et SRV-C sont correctement configurés ; accepter PC-B cité à la place de PC-D pour le doublon.}
\fi

\pagestyle{fancy}
\fancyhf{}
\renewcommand{\headrulewidth}{0pt}
\fancyfoot[L]{\footnotesize\itshape\color{grisfoot}Évaluation — Réseaux locaux — Sujet \variante\ifcorrige\ — CORRIGÉ\fi}
\fancyfoot[R]{\footnotesize\color{grisfoot}\thepage/3}

\setlist{itemsep=0pt,topsep=2pt,parsep=0pt}
\setlength{\parindent}{0pt}
\setlength{\parskip}{2pt}

\begin{document}

\begin{center}
  {\Large\bfseries Évaluation pratique — Réseaux locaux \quad
   \fcolorbox{black}{white}{\ SUJET \variante\ }\par}
  \vspace{1pt}
  {\bfseries Adressage IP, commutation et routage — BTS CIEL 1re année — durée : 2 heures\par}
  \ifcorrige\vspace{2pt}{\bfseries\color{corr}CORRIGÉ ET BARÈME — document enseignant}\par\fi
\end{center}
\vspace{-2pt}
{\renewcommand{\arraystretch}{1.35}
\begin{tabularx}{\linewidth}{|X|X|p{2.2cm}|p{2.6cm}|}
  \hline
  \textbf{Nom :} & \textbf{Prénom :} & \textbf{Poste n°} & \textbf{Note :} \hfill /20 \\ \hline
\end{tabularx}}

\textbf{Consignes :} travail \textbf{individuel}, documents et Internet \textbf{non autorisés}. Ce sujet est le
\textbf{document-réponse}. Fichier Packet Tracer : \ip{Nom\_Prenom\_EvalLAN\_\variante.pkt}, à enregistrer
régulièrement dans le dossier de rendu. Les « visas » sont apposés par le professeur, que vous appelez à chaque test réussi.
Durées conseillées : A 30 min — B 1 h 10 — C 20 min.

%==========================================================================
\section{Partie A — Questions de cours et calculs \hfill /6}

\question{A1}{2 pts}{Complétez le tableau en considérant le \textbf{masque par défaut} de la classe de chaque adresse.}
{\small\renewcommand{\arraystretch}{1.15}
\begin{tabularx}{\linewidth}{|l|c|C|C|C|c|}
  \hline
  \textbf{Adresse IP} & \textbf{Classe} & \textbf{Masque par défaut} & \textbf{Adresse réseau} &
  \textbf{Adresse de diffusion} & \textbf{Nb d'hôtes} \\ \hline
  \AIrows
\end{tabularx}}
\bareme{0,5 pt par ligne juste + 0,5 pt si les trois nombres d'hôtes sont justes ($2^{n}-2$, $n$ = nombre de bits d'hôte).}

\question{A2}{1 pt}{Indiquez le type de câble (\textbf{droit} ou \textbf{croisé}) pour chaque liaison.}
{\small\renewcommand{\arraystretch}{1.15}
\begin{tabularx}{\linewidth}{|X|c||X|c|}
  \hline
  \hl PC $\leftrightarrow$ commutateur & \makebox[1.6cm]{\sol{droit}} & Commutateur $\leftrightarrow$ routeur & \makebox[1.6cm]{\sol{droit}} \\ \hline
  \hl PC $\leftrightarrow$ PC (réseau P2P) & \sol{croisé} & Concentrateur $\leftrightarrow$ concentrateur & \sol{croisé} \\ \hline
\end{tabularx}}
\bareme{0,25 pt par case. Équipements de types différents : câble droit ; de même type : câble croisé.}

\question{A3}{1 pt}{Comment un \textbf{concentrateur (hub)} et un \textbf{commutateur (switch)} traitent-ils une trame
reçue ? Lequel limite les collisions, et pourquoi ?}
\zone{2.9cm}{Le \textbf{hub} (couche 1) répète la trame sur \textbf{tous ses autres ports} : un seul domaine de collision
partagé. Le \textbf{switch} (couche 2) lit l'adresse MAC de destination et, grâce à sa \textbf{table MAC}, ne l'envoie
\textbf{que sur le port du destinataire} : un domaine de collision par port, c'est lui qui limite les collisions.
\emph{0,5 pt hub, 0,5 pt switch.}}

\question{A4}{1 pt}{Qu'est-ce que la \textbf{passerelle par défaut} d'un poste ? Quelle condition son adresse doit-elle respecter ?}
\zone{2.4cm}{Adresse IP de l'interface du \textbf{routeur} à laquelle le poste envoie les paquets destinés à un
\textbf{autre réseau}. Elle doit \textbf{appartenir au même réseau} que le poste. \emph{0,5 + 0,5 pt.}}

\question{A5}{0,5 pt}{Un portable configuré en DHCP affiche \ip{\apipa} / \ip{255.255.0.0}. Que s'est-il passé ?}
\zone{1.5cm}{Aucun \textbf{serveur DHCP} n'a répondu : le poste s'est attribué une adresse \textbf{APIPA} (169.254.0.0/16).}

\question{A6}{0,5 pt}{Quel protocole utilise \texttt{ping} ? Nommez les deux messages échangés et la couche OSI des adresses IP.}
\zone{1.5cm}{\textbf{ICMP} ; \textbf{Echo Request} / \textbf{Echo Reply} ; adresses IP en \textbf{couche 3 (Réseau)}.}

%==========================================================================
\newpage
\section{Partie B — Réalisation sous Packet Tracer \hfill /10}

\begin{cdc}
\entreprise{} vous confie son réseau : \textbf{trois réseaux} reliés par un \textbf{routeur unique Cisco 2911}
nommé \ip{\Rnom} (interfaces G0/0, G0/1, G0/2).
\begin{itemize}[leftmargin=12pt]
  \item \textbf{Réseau \resUn} : \textbf{4 postes} PC-\pUn1 à PC-\pUn4, commutateur \ip{SW-\pUn}, adressage statique,
        réseau \ip{\netUn.0/24}.
  \item \textbf{Réseau \resDeux} : \textbf{2 postes} PC-\pDeux1, PC-\pDeux2 et \textbf{1 imprimante} IMP-\pDeux,
        commutateur \ip{SW-\pDeux}, adressage statique, réseau de classe B \ip{\netDeux.0.0/16}.
  \item \textbf{Réseau Serveur} : serveur \ip{\srvNom} d'adresse \ip{\netSrv.\srvH/24}, relié \textbf{directement} au routeur.
  \item La passerelle de chaque réseau est sa \textbf{dernière adresse utilisable}. Tous les équipements communiquent
        entre eux et accèdent à la page web du serveur. Nommage imposé (« Display Name »).
\end{itemize}
\end{cdc}

\question{B1}{2 pts}{Complétez le plan d'adressage \textbf{avant} la réalisation.}
{\footnotesize\setlength{\tabcolsep}{3pt}\renewcommand{\arraystretch}{1.15}
\begin{tabularx}{\linewidth}{|l|C|C|C|C|C|C|}
  \hline
  \textbf{Réseau} & \textbf{Adresse réseau} & \textbf{Masque} & \textbf{1re @ utilisable} &
  \textbf{Dernière @ utilisable} & \textbf{Diffusion} & \textbf{Passerelle} \\ \hline
  \hl \resUn & \sol{\netUn.0} & \sol{255.255.255.0} & \sol{\netUn.1} & \sol{\netUn.254} & \sol{\netUn.255} & \sol{\netUn.254} \\ \hline
  \hl \resDeux & \sol{\netDeux.0.0} & \sol{255.255.0.0} & \sol{\netDeux.0.1} & \sol{\netDeux.255.254} & \sol{\netDeux.255.255} & \sol{\netDeux.255.254} \\ \hline
  \hl Serveur & \sol{\netSrv.0} & \sol{255.255.255.0} & \sol{\netSrv.1} & \sol{\netSrv.254} & \sol{\netSrv.255} & \sol{\netSrv.254} \\ \hline
\end{tabularx}}
\bareme{0,5 pt par ligne juste + 0,5 pt si la convention de passerelle est respectée partout.}

\question{B2}{2 pts}{Construisez la topologie (équipements, câblage, nommage) et relevez l'adressage configuré.}
{\footnotesize\setlength{\tabcolsep}{3pt}\renewcommand{\arraystretch}{1.1}
\begin{tabularx}{\linewidth}{|l|C|C|C||l|C|C|}
  \hline
  \textbf{Hôte} & \textbf{Adresse IP} & \textbf{Masque} & \textbf{Passerelle} & \textbf{Interface} & \textbf{Adresse IP} & \textbf{Masque} \\ \hline
  \hl PC-\pUn1 & \sol{\netUn.1} & \sol{255.255.255.0} & \sol{\netUn.254} & \Rnom{} G0/0 & \sol{\netUn.254} & \sol{255.255.255.0} \\ \hline
  \hl PC-\pDeux1 & \sol{\netDeux.0.1} & \sol{255.255.0.0} & \sol{\netDeux.255.254} & \Rnom{} G0/1 & \sol{\netDeux.255.254} & \sol{255.255.0.0} \\ \hline
  \hl IMP-\pDeux & \sol{\netDeux.0.10} & \sol{255.255.0.0} & \sol{\netDeux.255.254} & \Rnom{} G0/2 & \sol{\netSrv.254} & \sol{255.255.255.0} \\ \hline
  \hl \srvNom & \ip{\netSrv.\srvH} & \sol{255.255.255.0} & \sol{\netSrv.254} & \multicolumn{3}{c|}{\cellcolor{black!8}} \\ \hline
\end{tabularx}}
\bareme{Adresses d'hôtes libres (uniques, dans le bon réseau) ; affectation G0/x libre. Sur le .pkt : équipements 0,5 ;
câblage 1 (câbles droits vers les switchs et switch–routeur, \textbf{croisé} serveur–routeur, liens verts) ; nommage 0,5.}

\question{B3}{3 pts}{Configurez tous les équipements : IP, masques, passerelles des hôtes ; interfaces du routeur
(« Port Status : On »).}
\bareme{Sur le .pkt : interfaces routeur 1 pt ; IP/masques des hôtes 1 pt ; passerelles (imprimante et serveur compris)
1 pt. $-0,25$ par erreur, sans descendre sous 0 par item.}

\question{B4}{3 pts}{Réalisez les tests, notez la commande et le résultat, puis faites viser chaque test réussi.}
{\small\renewcommand{\arraystretch}{1.1}
\begin{tabularx}{\linewidth}{|c|l|l|C|c|c|}
  \hline
  \textbf{N°} & \textbf{Source} & \textbf{Destination} & \textbf{Commande / action} & \textbf{Résultat} & \textbf{Visa prof.} \\ \hline
  \hl 1 & PC-\pUn1 & PC-\pUn4 & \sol{ping \netUn.4} & \sol{réussi} & \hspace{1.4cm} \\ \hline
  \hl 2 & PC-\pDeux1 & IMP-\pDeux & \sol{ping \netDeux.0.10} & \sol{réussi} & \\ \hline
  \hl 3 & PC-\pUn2 & PC-\pDeux2 & \sol{ping \netDeux.0.2} & \sol{réussi} & \\ \hline
  \hl 4 & PC-\pDeux2 & \srvNom & \sol{ping \netSrv.\srvH} & \sol{réussi} & \\ \hline
  \hl 5 & PC-\pUn3 & page web du serveur & \sol{navigateur : {\NoAutoSpacing http://\netSrv.\srvH}} & \sol{page affichée} & \\ \hline
\end{tabularx}}
\bareme{0,5 pt par test visé + 0,5 pt pour B5. Le premier ping inter-réseaux peut échouer (résolution ARP) ; les suivants doivent réussir.}

\question{B5}{incl. B4}{Test n°3 : par quels équipements passe le paquet ? Pourquoi PC-\pUn2 l'envoie-t-il au routeur ?}
\zone{2.3cm}{PC-\pUn2 $\to$ SW-\pUn{} $\to$ \Rnom{} $\to$ SW-\pDeux{} $\to$ PC-\pDeux2. Grâce à son masque, PC-\pUn2 constate que
la destination (\netDeux.0.0/16) n'est \textbf{pas dans son réseau} (\netUn.0/24) : il envoie le paquet à sa
\textbf{passerelle par défaut}, qui le route.}

%==========================================================================
\newpage
\section{Partie C — Diagnostic d'un réseau en panne \hfill /4}

Un technicien constate qu'\textbf{aucun poste ne parvient à joindre le serveur SRV-C} et que certains postes
communiquent mal entre eux. Les câbles sont du bon type.

\begin{center}
\begin{minipage}{0.37\linewidth}\centering
\begin{tikzpicture}[font=\footnotesize,
  eq/.style={draw,fill=qbg,rounded corners=2pt,minimum width=9mm,minimum height=5mm,inner sep=2pt},
  every path/.style={thick}]
  \node[eq] (SW) {SW-1};
  \node[eq,right=9mm of SW] (R) {R1};
  \node[eq,right=7mm of R] (SRV) {SRV-C};
  \node[eq,above left=5mm and -3mm of SW] (A) {PC-A};
  \node[eq,above right=5mm and -3mm of SW] (B) {PC-B};
  \node[eq,below left=5mm and -3mm of SW] (C) {PC-C};
  \node[eq,below right=5mm and -3mm of SW] (D) {PC-D};
  \foreach \n in {A,B,C,D} \draw (\n) -- (SW);
  \draw (SW) -- node[above,font=\tiny]{G0/0} (R);
  \draw (R) -- node[above,font=\tiny]{G0/1} (SRV);
\end{tikzpicture}
\end{minipage}\hfill
\begin{minipage}{0.61\linewidth}\centering
{\small\renewcommand{\arraystretch}{1.05}\setlength{\tabcolsep}{4pt}
\begin{tabular}{lllll}
  \toprule
  \textbf{Équipement} & \textbf{Adresse IP} & \textbf{Masque} & \textbf{Passerelle} & \textbf{Port} \\
  \midrule
  \Crows
  \bottomrule
\end{tabular}}
\end{minipage}
\end{center}

\question{C1}{4 pts}{Relevez les \textbf{quatre anomalies}. Pour chacune : conséquence et correction.}
{\small\renewcommand{\arraystretch}{1.1}
\ifcorrige\def\hc{}\else\def\hc{\rule[-10mm]{0pt}{14mm}}\fi
\begin{tabularx}{\linewidth}{|c|>{\raggedright\arraybackslash}p{2cm}|L|L|L|}
  \hline
  \textbf{N°} & \textbf{Équipement} & \textbf{Anomalie} & \textbf{Conséquence} & \textbf{Correction} \\ \hline
  \ifcorrige\Csol\else
  \hc 1 & & & & \\ \hline
  \hc 2 & & & & \\ \hline
  \hc 3 & & & & \\ \hline
  \hc 4 & & & & \\ \hline
  \fi
\end{tabularx}}
\bareme{1 pt par anomalie (0,25 équipement, anomalie, conséquence, correction ; ordre indifférent). \Cbareme}

%==========================================================================
\section{Bonus — Extension sans fil \hfill +1}

\begin{cdc}Ajoutez au réseau \resUn{} un \textbf{AccessPoint-PT} relié à SW-\pUn{} (SSID \ip{\ssid}, \textbf{WPA2-PSK})
et un portable \ip{Laptop-\pUn} équipé d'un module \texttt{WPC300N}, en adressage \textbf{statique}.
Vérifiez qu'il joint \srvNom.\end{cdc}
{\small\renewcommand{\arraystretch}{1.1}
\begin{tabularx}{\linewidth}{|l|C|C|C|c|}
  \hline
  \textbf{Équipement} & \textbf{Adresse IP} & \textbf{Masque} & \textbf{Passerelle} & \textbf{Visa prof.} \\ \hline
  \hl Laptop-\pUn & \sol{\netUn.50} & \sol{255.255.255.0} & \sol{\netUn.254} & \hspace{1.4cm} \\ \hline
\end{tabularx}}
\bareme{Le point d'accès fait office de pont : le portable est dans \netUn.0/24. Bonus si associé en WPA2-PSK et ping
vers \netSrv.\srvH{} réussi. Note plafonnée à 20.}

\end{document}
